% Auto-generated from data/segments.json + layout.json (+ transforms) — do not edit by hand.
% volume: 13Sam02
% mode: reading
% segments: 3440
% transform_rules: 0
% toc_marks: 594

% Volume layout defaults (layout.json ``layout``).
\csromanlayoutapply{1}{1.5}{21.6pt}{5pt}{6.3pt}{65pt}{2.5em}%

% Reading physical-page layout (layout.json ``page_layout_reading_mode``).
\csromanreadingpagelayoutvolume{1}{1.5}{21.6pt}{5pt}{6.3pt}{65pt}{2.5em}%
\csromanreadingpagelayoutenable

\csromanheader{2}{\textbf{สํยุตฺตนิกาย}}
\csromanmatikaanchor{mtk.0}
\csromantocmarknopage{chapter}{ขนฺธวคฺคสํยุตฺตปาฬิ}{mtk.0}
\bookmark[dest={mtk.0},level=0]{ขนฺธวคฺคสํยุตฺตปาฬิ}
\gambhira{\textbf{ขนฺธวคฺคสํยุตฺตปาฬิ}}
\namakkaram{นโม ตสฺส ภควโต อรหโต สมฺมาสมฺพุทฺธสฺส.}
\chapterhead{1. นกุลปิตุวคฺค}
\csromanheader{2}{1. นกุลปิตุสุตฺต}
\proseitem{1}{\csromanfolio{1}เอวํ เม สุตํ\csromandash{}เอกํ สมยํ ภควา ภคฺเคสุ วิหรติ สุสุมารคิเร\footnote{สุํสุมารคิเร (สี, สฺยา, กํ, อิ)} เภสกฬาวเน มิคทาเย.\csromanspacer{} อถ โข นกุลปิตา คหปติ เยน ภควา เตนุปสงฺกมิ, อุปสงฺกมิตฺวา ภควนฺตํ อภิวาเทตฺวา เอกมนฺตํ นิสีทิ, เอกมนฺตํ นิสินฺโน โข นกุลปิตา คหปติ ภควนฺตํ เอตทโวจ\csromandash{}}

\prose{อหมสฺมิ ภนฺเต ชิณฺโณ วุฑฺโฒ มหลฺลโก อทฺธคโต วโย-อนุปฺปตฺโต อาตุรกาโย อภิกฺขณาตงฺโก, อนิจฺจทสฺสาวี โข ปนาหํ ภนฺเต ภควโต มโนภาวนียานญฺจ ภิกฺขูนํ, โอวทตุ มํ ภนฺเต ภควา, อนุสาสตุ มํ ภนฺเต ภควา, ยํ มมสฺส ทีฆรตฺตํ หิตาย สุขายาติ.}

\prose{เอวเมตํ คหปติ, เอวเมตํ คหปติ, อาตุโร หายํ คหปติ กาโย อณฺฑภูโต ปริโยนทฺโธ, โย หิ คหปติ อิมํ กายํ ปริหรนฺโต มุหุตฺตมฺปิ อาโรคฺยํ ปฏิชาเนยฺย, กิมญฺญตฺร พาลฺยา.\csromanspacer{} ตสฺมาติห \csromanfolio{2}เต คหปติ เอวํ สิกฺขิตพฺพํ “อาตุรกายสฺส เม สโต จิตฺตํ อนาตุรํ ภวิสฺสตี”ติ.\csromanspacer{} เอวํ หิ เต คหปติ สิกฺขิตพฺพนฺติ.}

\prose{อถ โข นกุลปิตา คหปติ ภควโต ภาสิตํ อภินนฺทิตฺวา อนุโมทิตฺวา อุฏฺฐายาสนา ภควนฺตํ อภิวาเทตฺวา ปทกฺขิณํ กตฺวา เยนายสฺมา สาริปุตฺโต เตนุปสงฺกมิ, อุปสงฺกมิตฺวา อายสฺมนฺตํ สาริปุตฺตํ อภิวาเทตฺวา เอกมนฺตํ นิสีทิ, เอกมนฺตํ นิสินฺนํ โข นกุลปิตรํ คหปติํ อายสฺมา สาริปุตฺโต เอตทโวจ “วิปฺปสนฺนานิ โข เต คหปติ อินฺทฺริยานิ, ปริสุทฺโธ มุขวณฺโณ ปริโยทาโต, อลตฺถ โน อชฺช ภควโต สมฺมุขา ธมฺมิํ กถํ สวนายา”ติ.}

\prose{กถํ หิ โน สิยา ภนฺเต, อิทานาหํ ภนฺเต ภควตา ธมฺมิยา กถาย อมเตน อภิสิตฺโตติ.\csromanspacer{} ยถา กถํ ปน ตฺวํ คหปติ ภควตา ธมฺมิยา กถาย อมเตน อภิสิตฺโตติ.\csromanspacer{} อิธาหํ ภนฺเต เยน ภควา เตนุปสงฺกมิํ, อุปสงฺกมิตฺวา ภควนฺตํ อภิวาเทตฺวา เอกมนฺตํ นิสีทิํ, เอกมนฺตํ นิสินฺโน ขฺวาหํ ภนฺเต ภควนฺตํ เอตทโวจํ “อหมสฺมิ ภนฺเต ชิณฺโณ วุฑฺโฒ มหลฺลโก อทฺธคโต วโย-อนุปฺปตฺโต อาตุรกาโย อภิกฺขณาตงฺโก, อนิจฺจทสฺสาวี โข ปนาหํ ภนฺเต ภควโต มโนภาวนียานญฺจ ภิกฺขูนํ, โอวทตุ มํ ภนฺเต ภควา, อนุสาสตุ มํ ภนฺเต ภควา, ยํ มมสฺส ทีฆรตฺตํ หิตาย สุขายา”ติ.}

\prose{เอวํ วุตฺเต มํ ภนฺเต ภควา เอตทโวจ “เอวเมตํ คหปติ, เอวเมตํ คหปติ, อาตุโร หายํ คหปติ กาโย อณฺฑภูโต ปริโยนทฺโธ, โย หิ คหปติ อิมํ กายํ ปริหรนฺโต มุหุตฺตมฺปิ อาโรคฺยํ ปฏิชาเนยฺย, กิมญฺญตฺร พาลฺยา.\csromanspacer{} ตสฺมาติห เต คหปติ เอวํ สิกฺขิตพฺพํ “อาตุรกายสฺส เม สโต จิตฺตํ อนาตุรํ ภวิสฺสตี”ติ.\csromanspacer{} เอวํ หิ เต คหปติ สิกฺขิตพฺพนฺติ.\csromanspacer{} เอวํ ขฺวาหํ ภนฺเต ภควตา ธมฺมิยา กถาย อมเตน อภิสิตฺโตติ.}

\prose{น หิ ปน ตํ คหปติ ปฏิภาสิ ภควนฺตํ\footnote{ตํ ภควนฺตํ (สี)} อุตฺตริํ ปฏิปุจฺฉิตุํ “กิตฺตาวตา นุ โข ภนฺเต อาตุรกาโย เจว โหติ อาตุรจิตฺโต จ, กิตฺตาวตา จ ปน อาตุรกาโย หิ โข โหติ โน จ \csromanfolio{3}อาตุรจิตฺโต”ติ.\csromanspacer{} ทูรโตปิ โข มยํ ภนฺเต อาคจฺเฉยฺยาม อายสฺมโต สาริปุตฺตสฺส สนฺติเก เอตสฺส ภาสิตสฺส อตฺถมญฺญาตุํ, สาธุ วตายสฺมนฺตํเยว สาริปุตฺตํ ปฏิภาตุ เอตสฺส ภาสิตสฺส อตฺโถติ.}

\prose{เตน หิ คหปติ สุณาหิ สาธุกํ มนสิ กโรหิ, ภาสิสฺสามีติ. “เอวํ ภนฺเต”ติ โข นกุลปิตา คหปติ อายสฺมโต สาริปุตฺตสฺส ปจฺจสฺโสสิ.\csromanspacer{} อายสฺมา สาริปุตฺโต เอตทโวจ\csromandash{}}

\prose{กถญฺจ คหปติ อาตุรกาโย เจว โหติ อาตุรจิตฺโต จ.\csromanspacer{} อิธ คหปติ อสฺสุตวา ปุถุชฺชโน อริยานํ อทสฺสาวี อริยธมฺมสฺส อโกวิโท อริยธมฺเม อวินีโต, สปฺปุริสานํ อทสฺสาวี สปฺปุริสธมฺมสฺส อโกวิโท สปฺปุริสธมฺเม อวินีโต รูปํ อตฺตโต สมนุปสฺสติ, รูปวนฺตํ วา อตฺตานํ, อตฺตนิ วา รูปํ, รูปสฺมิํ วา อตฺตานํ, “อหํ รูปํ มม รูปนฺ”ติ ปริยุฏฺฐฏฺฐายี โหติ.\csromanspacer{} ตสฺส “อหํ รูปํ มม รูปนฺ”ติ ปริยุฏฺฐฏฺฐายิโน ตํ รูปํ วิปริณมติ อญฺญถา โหติ.\csromanspacer{} ตสฺส รูปวิปริณามญฺญถาภาวา อุปฺปชฺชนฺติ โสกปริเทว ทุกฺข}

\prose{เวทนํ อตฺตโต สมนุปสฺสติ, เวทนาวนฺตํ วา อตฺตานํ, อตฺตนิ วา เวทนํ, เวทนาย วา อตฺตานํ, “อหํ เวทนา มม เวทนา”ติ ปริยุฏฺฐฏฺฐายี โหติ, ตสฺส “อหํ เวทนา มม เวทนา”ติ ปริยุฏฺฐฏฺฐายิโน สา เวทนา วิปริณมติ อญฺญถา โหติ.\csromanspacer{} ตสฺส เวทนาวิปริณามญฺญถาภาวา อุปฺปชฺชนฺติ โสกปริเทว ทุกฺข}

\prose{สญฺญํ อตฺตโต สมนุปสฺสติ, สญฺญาวนฺตํ วา อตฺตานํ, อตฺตนิ วา สญฺญํ, สญฺญาย วา อตฺตานํ, “อหํ สญฺญา มม สญฺญา”ติ ปริยุฏฺฐฏฺฐายี โหติ, ตสฺส “อหํ สญฺญา มม สญฺญา”ติ ปริยุฏฺฐฏฺฐายิโน สา สญฺญา วิปริณมติ อญฺญถา โหติ.\csromanspacer{} ตสฺส สญฺญาวิปริณามญฺญถาภาวา อุปฺปชฺชนฺติ โสก ปริเทว ทุกฺข โทมนสฺสุปายาสา.}

\prose{สงฺขาเร อตฺตโต สมนุปสฺสติ, สงฺขารวนฺตํ วา อตฺตานํ, อตฺตนิ วา สงฺขาเร, สงฺขาเรสุ วา อตฺตานํ, “อหํ สงฺขารา มม สงฺขารา”ติ ปริยุฏฺฐฏฺฐายี โหติ, ตสฺส “อหํ สงฺขารา มม สงฺขารา”ติ ปริยุฏฺฐฏฺฐายิโน เต สงฺขารา วิปริณมนฺติ อญฺญถา โหนฺติ.\csromanspacer{} ตสฺส สงฺขารวิปริณามญฺญถาภาวา อุปฺปชฺชนฺติ โสก ปริเทว ทุกฺข}

\prose{\csromanfolio{4}วิญฺญาณํ อตฺตโต สมนุปสฺสติ, วิญฺญาณวนฺตํ วา อตฺตานํ, อตฺตนิ วา วิญฺญาณํ, วิญฺญาณสฺมิํ วา อตฺตานํ, “อหํ วิญฺญาณํ มม วิญฺญาณนฺ”ติ ปริยุฏฺฐฏฺฐายี โหติ, ตสฺส “อหํ วิญฺญาณํ มม วิญฺญาณนฺ”ติ ปริยุฏฺฐฏฺฐายิโน ตํ วิญฺญาณํ วิปริณมติ อญฺญถา โหติ.\csromanspacer{} ตสฺส วิญฺญาณวิปริณามญฺญถาภาวา อุปฺปชฺชนฺติ โสก ปริเทว ทุกฺข โทมนสฺสุปายาสา.\csromanspacer{} เอวํ โข คหปติ อาตุรกาโย เจว โหติ อาตุรจิตฺโต จ.}

\prose{กถญฺจ คหปติ อาตุรกาโย หิ โข โหติ โน จ อาตุรจิตฺโต.\csromanspacer{} อิธ คหปติ สุตวา อริยสาวโก อริยานํ ทสฺสาวี อริยธมฺมสฺส โกวิโท อริยธมฺเม สุวินีโต, สปฺปุริสานํ ทสฺสาวี สปฺปุริสธมฺมสฺส โกวิโท สปฺปุริสธมฺเม สุวินีโต น รูปํ อตฺตโต สมนุปสฺสติ, น รูปวนฺตํ วา อตฺตานํ, น อตฺตนิ วา รูปํ, น รูปสฺมิํ วา อตฺตานํ, “อหํ รูปํ มม รูปนฺ”ติ น ปริยุฏฺฐฏฺฐายี โหติ, ตสฺส “อหํ รูปํ มม รูปนฺ”ติ อปริยุฏฺฐฏฺฐายิโน ตํ รูปํ วิปริณมติ อญฺญถา โหติ.\csromanspacer{} ตสฺส รูปวิปริณามญฺญถาภาวา นุปฺปชฺชนฺติ โสก ปริเทว ทุกฺข}

\prose{น เวทนํ อตฺตโต สมนุปสฺสติ, น เวทนาวนฺตํ วา อตฺตานํ, น อตฺตนิ วา เวทนํ, น เวทนาย วา อตฺตานํ, “อหํ เวทนา มม เวทนา”ติ น ปริยุฏฺฐฏฺฐายี โหติ, ตสฺส “อหํ เวทนา มม เวทนา”ติ อปริยุฏฺฐฏฺฐายิโน สา เวทนา วิปริณมติ อญฺญถา โหติ.\csromanspacer{} ตสฺส เวทนาวิปริณามญฺญถาภาวา นุปฺปชฺชนฺติ โสกปริเทว ทุกฺข}

\prose{น สญฺญํ อตฺตโต สมนุปสฺสติ, น สญฺญาวนฺตํ วา อตฺตานํ, น อตฺตนิ วา สญฺญํ, น สญฺญาย วา อตฺตานํ, “อหํ สญฺญา มม สญฺญา”ติ นปริยุฏฺฐฏฺฐายี โหติ, ตสฺส “อหํ สญฺญา มม สญฺญา”ติ อปริยุฏฺฐฏฺฐายิโน สา สญฺญาวิปริณมติ อญฺญถา โหติ.\csromanspacer{} ตสฺส สญฺญาวิปริณามญฺญถภาวา นุปฺปชฺชนฺติ โสก ปริเทว ทุกฺขโทมนสฺสุปายาสา.}

\prose{น สงฺขาเร อตฺตโต สมนุปสฺสติ, น สงฺขารวนฺตํ วา อตฺตานํ, น อตฺตนิ วา สงฺขาเร, น สงฺขาเรสุ วา อตฺตานํ, “อหํ สงฺขารา มม สงฺขารา”ติ น ปริยุฏฺฐฏฺฐายี โหติ, ตสฺส “อหํ สงฺขารา มม สงฺขารา”ติ อปริยุฏฺฐฏฺฐายิโน เต สงฺขารา วิปริณมนฺติ อญฺญถา โหนฺติ.\csromanspacer{} ตสฺส สงฺขารวิปริณามญฺญถาภาวา นุปฺปชฺชนฺติ โสกปริเทวทุกฺขโทมนสฺสุปายาสา.}

\prose{\csromanfolio{5}น วิญฺญาณํ อตฺตโต สมนุปสฺสติ, น วิญฺญาณวนฺตํ วา อตฺตานํ, น อตฺตนิ วา วิญฺญาณํ, น วิญฺญาณสฺมิํ วา อตฺตานํ, “อหํ วิญฺญาณํ มม วิญฺญาณนฺ”ติ น ปริยุฏฺฐฏฺฐายี โหติ.\csromanspacer{} ตสฺส “อหํ วิญฺญาณํ มม วิญฺญาณนฺ”ติ อปริยุฏฺฐฏฺฐายิโน ตํ วิญฺญาณํ วิปริณมติ อญฺญถา โหติ.\csromanspacer{} ตสฺส วิญฺญาณวิปริณามญฺญถาภาวา นุปฺปชฺชนฺติ โสกปริเทว ทุกฺขโทมนสฺสุปายาสา.\csromanspacer{} เอวํ โข คหปติ อาตุรกาโย โหติ โน จ อาตุรจิตฺโตติ.}

\prose{อิทมโวจ อายสฺมา สาริปุตฺโต.\csromanspacer{} อตฺตมโน นกุลปิตา คหปติ อายสฺมโต สาริปุตฺตสฺส ภาสิตํ อภินนฺทีติ.\csromanspacer{}\csromanspacer{}ปฐมํ.}

\csromanheader{2}{2. เทวทหสุตฺต}
\proseitem{2}{เอวํ เม สุตํ\csromandash{}เอกํ สมยํ ภควา สกฺเกสุ\footnote{สเกฺยสุ (ก)} วิหรติ เทวทหํ นาม สกฺยานํ นิคโม.\csromanspacer{} อถ โข สมฺพหุลา ปจฺฉาภูมคมิกา ภิกฺขู เยน ภควา เตนุปสงฺกมิํสุ, อุปสงฺกมิตฺวา ภควนฺตํ อภิวาเทตฺวา เอกมนฺตํ นิสีทิํสุ, เอกมนฺตํ นิสินฺนา โข เต ภิกฺขู ภควนฺตํ เอตทโวจุํ “อิจฺฉาม มยํ ภนฺเต ปจฺฉาภูมํ ชนปทํ คนฺตุํ, ปจฺฉาภูเม ชนปเท นิวาสํ กปฺเปตุนฺ”ติ.}

\prose{อปโลกิโต ปน โว ภิกฺขเว สาริปุตฺโตติ.\csromanspacer{} น โข โน ภนฺเต อปโลกิโต อายสฺมา สาริปุตฺโตติ.\csromanspacer{} อปโลเกถ ภิกฺขเว สาริปุตฺตํ, สาริปุตฺโต ภิกฺขเว ปณฺฑิโต ภิกฺขูนํ อนุคฺคาหโก สพฺรหฺมจารีนนฺติ. “เอวํ ภนฺเต”ติ โข เต ภิกฺขู ภควโต ปจฺจสฺโสสุํ.}

\prose{เตน โข ปน สมเยน อายสฺมา สาริปุตฺโต ภควโต อวิทูเร อญฺญตรสฺมิํ เอฬคลาคุมฺเพ นิสินฺโน โหติ.\csromanspacer{} อถ โข เต ภิกฺขู ภควโต ภาสิตํ อภินนฺทิตฺวา อนุโมทิตฺวา อุฏฺฐายาสนา ภควนฺตํ อภิวาเทตฺวา ปทกฺขิณํ กตฺวา เยนายสฺมา สาริปุตฺโต เตนุปสงฺกมิํสุ, อุปสงฺกมิตฺวา อายสฺมตา สาริปุตฺเตน สทฺธิํ สมฺโมทิํสุ, สมฺโมทนียํ กถํ สารณียํ\footnote{สาราณียํ (สี, สฺยา, กํ, อิ)} วีติสาเรตฺวา เอกมนฺตํ นิสีทิํสุ, เอกมนฺตํ นิสินฺนา โข เต ภิกฺขู อายสฺมนฺตํ สาริปุตฺตํ เอตทโวจุํ “อิจฺฉาม มยํ อาวุโส สาริปุตฺต ปจฺฉาภูมํ ชนปทํ คนฺตุํ, ปจฺฉาภูเม ชนปเท นิวาสํ กปฺเปตุํ, อปโลกิโต โน สตฺถา”ติ.}

\prose{\csromanfolio{6}สนฺติ หาวุโส นานาเวรชฺชคตํ ภิกฺขุํ ปญฺหํ ปุจฺฉิตาโร ขตฺติยปณฺฑิตาปิ พฺราหฺมณปณฺฑิตาปิ คหปติปณฺฑิตาปิ สมณปณฺฑิตาปิ.\csromanspacer{} ปณฺฑิตา หาวุโส มนุสฺสา วีมํสกา “กิํ วาที ปนายสฺมนฺตานํ\footnote{กิํ วาทายสฺมนฺตานํ (อิ, ก)} สตฺถา กิมกฺขายี”ติ, กจฺจิ โว อายสฺมนฺตานํ ธมฺมา สุสฺสุตา สุคฺคหิตา สุมนสิกตา สูปธาริตา สุปฺปฏิวิทฺธา ปญฺญาย, ยถา พฺยากรมานา อายสฺมนฺโต วุตฺตวาทิโน เจว ภควโต อสฺสถ, น จ ภควนฺตํ อภูเตน อพฺภาจิกฺเขยฺยาถ, ธมฺมสฺส จานุธมฺมํ พฺยากเรยฺยาถ, น จ โกจิ สหธมฺมิโก วาทานุวาโท\footnote{วาทานุปาโต (อฏฺฐกถายํ ปาฐนฺตรํ)} คารยฺหํ ฐานํ อาคจฺเฉยฺยาติ.}

\prose{ทูรโตปิ โข มยํ อาวุโส อาคจฺเฉยฺยาม อายสฺมโต สาริปุตฺตสฺส สนฺติเก เอตสฺส ภาสิตสฺส อตฺถมญฺญาตุํ, สาธุ วตายสฺมนฺตํเยว สาริปุตฺตํ ปฏิภาตุ เอตสฺส ภาสิตสฺส อตฺโถติ.\csromanspacer{} เตนหาวุโส สุณาถ สาธุกํ มนสิ กโรถ, ภาสิสฺสามีติ. “เอวมาวุโส”ติ โข เต ภิกฺขู อายสฺมโต สาริปุตฺตสฺส ปจฺจสฺโสสุํ.\csromanspacer{} อายสฺมา สาริปุตฺโต เอตทโวจ\csromandash{}}

\prose{สนฺติ หาวุโส นานาเวรชฺชคตํ ภิกฺขุํ ปญฺหํ ปุจฺฉิตาโร ขตฺติยปณฺฑิตาปิ ฯเปฯ สมณปณฺฑิตาปิ.\csromanspacer{} ปณฺฑิตา หาวุโส มนุสฺสา วีมํสกา “กิํ วาที ปนายสฺมนฺตานํ สตฺถา, กิมกฺขายี”ติ, เอวํ ปุฏฺฐา ตุเมฺห อาวุโส เอวํ พฺยากเรยฺยาถ “ฉนฺทราควินยกฺขายี โข โน อาวุโส สตฺถา”ติ.}

\prose{เอวํ พฺยากเตปิ โข อาวุโส อสฺสุเยว อุตฺตริํ ปญฺหํ ปุจฺฉิตาโร ขตฺติยปณฺฑิตาปิ ฯเปฯ สมณปณฺฑิตาปิ.\csromanspacer{} ปณฺฑิตา หาวุโส มนุสฺสา วีมํสกา “กิสฺมิํ ปนายสฺมนฺตานํ ฉนฺทราควินยกฺขายี สตฺถา”ติ, เอวํ ปุฏฺฐา ตุเมฺห อาวุโส เอวํ พฺยากเรยฺยาถ “รูเป โข อาวุโส ฉนฺทราควินยกฺขายี สตฺถา.\csromanspacer{} เวทนาย.\csromanspacer{} สญฺญาย.\csromanspacer{} สงฺขาเรสุ.\csromanspacer{} วิญฺญาเณ ฉนฺทราควินยกฺขายี สตฺถา”ติ.}

\prose{เอวํ พฺยากเตปิ โข อาวุโส อสฺสุเยว อุตฺตริํ ปญฺหํ ปุจฺฉิตาโร ขตฺติยปณฺฑิตาปิ ฯเปฯ สมณปณฺฑิตาปิ.\csromanspacer{} ปณฺฑิตา หาวุโส มนุสฺสา วีมํสกา “กิํ ปนายสฺมนฺตานํ อาทีนวํ ทิสฺวา รูเป ฉนฺทราควินยกฺขายี สตฺถา.\csromanspacer{} เวทนาย.\csromanspacer{} สญฺญาย.\csromanspacer{} สงฺขาเรสุ.\csromanspacer{} วิญฺญาเณ ฉนฺทราควินยกฺขายี สตฺถา”ติ, เอวํ ปุฏฺฐา ตุเมฺห อาวุโส เอวํ พฺยากเรยฺยาถ “รูเป โข \csromanfolio{7}อาวุโส อวิคตราคสฺส\footnote{อวีตราคสฺส (สฺยา, กํ)} อวิคตฉนฺทสฺส อวิคตเปมสฺส อวิคตปิปาสสฺส อวิคตปริฬาหสฺส อวิคตตณฺหสฺส ตสฺส รูปสฺส วิปริณามญฺญถาภาวา อุปฺปชฺชนฺติ โสก ปริเทว ทุกฺข โทมนสฺสุปายาสา.\csromanspacer{} เวทนาย.\csromanspacer{} สญฺญาย.\csromanspacer{} สงฺขาเรสุ อวิคตราคสฺส ฯเปฯ อวิคตตณฺหสฺส เตสํ สงฺขารานํ วิปริณามญฺญถาภาวา อุปฺปชฺชนฺติ โสก ปริเทว ทุกฺข โทมนสฺสุปายาสา.\csromanspacer{} วิญฺญาเณ อวิคตราคสฺส อวิคตฉนฺทสฺส อวิคตเปมสฺส อวิคตปิปาสสฺส อวิคตปริฬาหสฺส อวิคตตณฺหสฺส ตสฺส วิญฺญาณสฺส วิปริณามญฺญถาภาวา อุปฺปชฺชนฺติ โสก ปริเทว ทุกฺข โทมนสฺสุปายาสา.\csromanspacer{} อิทํ โข โน อาวุโส อาทีนวํ ทิสฺวา รูเป ฉนฺทราควินยกฺขายี สตฺถา.\csromanspacer{} เวทนาย.\csromanspacer{} สญฺญาย.\csromanspacer{} สงฺขาเรสุ.\csromanspacer{} วิญฺญาเณ ฉนฺทราควินยกฺขายี สตฺถา”ติ.}

\prose{เอวํ พฺยากเตปิ โข อาวุโส อสฺสุเยว อุตฺตริํ ปญฺหํ ปุจฺฉิตาโร ขตฺติยปณฺฑิตาปิ พฺราหฺมณปณฺฑิตาปิ คหปติปณฺฑิตาปิ สมณปณฺฑิตาปิ.\csromanspacer{} ปณฺฑิตา หาวุโส มนุสฺสา วีมํสกา “กิํ ปนายสฺมนฺตานํ อานิสํสํ ทิสฺวา รูเป ฉนฺทราควินยกฺขายี สตฺถา.\csromanspacer{} เวทนาย.\csromanspacer{} สญฺญาย.\csromanspacer{} สงฺขาเรสุ.\csromanspacer{} วิญฺญาเณ ฉนฺทราควินยกฺขายี สตฺถา”ติ, เอวํ ปุฏฺฐา ตุเมฺห อาวุโส เอวํ พฺยากเรยฺยาถ “รูเป โข อาวุโส วิคตราคสฺส วิคตฉนฺทสฺส วิคตเปมสฺส วิคตปิปาสสฺส วิคตปริฬาหสฺส วิคตตณฺหสฺส ตสฺส รูปสฺส วิปริณามญฺญถาภาวา นุปฺปชฺชนฺติ โสก ปริเทว ทุกฺข โทมนสฺสุปายาสา.\csromanspacer{} เวทนาย.\csromanspacer{} สญฺญาย.\csromanspacer{} สงฺขาเรสุ วิคตราคสฺส วิคตฉนฺทสฺส วิคตเปมสฺส วิคตปิปาสสฺส วิคตปริฬาหสฺส วิคตตณฺหสฺส เตสํ สงฺขารานํ วิปริณามญฺญถาภาวา นุปฺปชฺชนฺติ โสก ปริเทว ทุกฺข โทมนสฺสุปายาสา.\csromanspacer{} วิญฺญาเณ วิคตราคสฺส วิคตฉนฺทสฺส วิคตเปมสฺส วิคตปิปาสสฺส วิคตปริฬาหสฺส วิคตตณฺหสฺส ตสฺส วิญฺญาณสฺส วิปริณามญฺญถาภาวา นุปฺปชฺชนฺติ โสกปริเทวทุกฺขโทมนสฺสุปายาสา.\csromanspacer{} อิทํ โข โน อาวุโส อานิสํสํ ทิสฺวา รูเป ฉนฺทราควินยกฺขายี สตฺถา.\csromanspacer{} เวทนาย.\csromanspacer{} สญฺญาย.\csromanspacer{} สงฺขาเรสุ.\csromanspacer{} วิญฺญาเณ ฉนฺทราควินยกฺขายี สตฺถา”ติ.}

\prose{อกุสเล จาวุโส ธมฺเม อุปสมฺปชฺช วิหรโต ทิฏฺเฐ เจว ธมฺเม สุโข วิหาโร อภวิสฺส อวิฆาโต อนุปายาโส อปริฬาโห, กายสฺส จ เภทา ปรํ มรณา สุคติ ปาฏิกงฺขา, นยิทํ ภควา \csromanfolio{8}อกุสลานํ ธมฺมานํ ปหานํ วณฺเณยฺย, ยสฺมา จ โข อาวุโส อกุสเล ธมฺเม อุปสมฺปชฺช วิหรโต ทิฏฺเฐ เจว ธมฺเม ทุกฺโข วิหาโร สวิฆาโต ส-อุปายาโส สปริฬาโห, กายสฺส จ เภทา ปรํ มรณา ทุคฺคติ ปาฏิกงฺขา, ตสฺมา ภควา อกุสลานํ ธมฺมานํ ปหานํ วณฺเณติ.}

\prose{กุสเล จาวุโส ธมฺเม อุปสมฺปชฺช วิหรโต ทิฏฺเฐ เจว ธมฺเม ทุกฺโข วิหาโร อภวิสฺส สวิฆาโต ส-อุปายาโส สปริฬาโห, กายสฺส จ เภทา ปรํ มรณา ทุคฺคติ ปาฏิกงฺขา, นยิทํ ภควา กุสลานํ ธมฺมานํ อุปสมฺปทํ วณฺเณยฺย, ยสฺมา จ โข อาวุโส กุสเล ธมฺเม อุปสมฺปชฺช วิหรโต ทิฏฺเฐ เจว ธมฺเม สุโข วิหาโร อวิฆาโต อนุปายาโส อปริฬาโห, กายสฺส จ เภทา ปรํ มรณา สุคติ ปาฏิกงฺขา, ตสฺมา ภควา กุสลานํ ธมฺมานํ อุปสมฺปทํ วณฺเณตีติ.}

\prose{อิทมโวจายสฺมา สาริปุตฺโต.\csromanspacer{} อตฺตมนา เต ภิกฺขู อายสฺมโต สาริปุตฺตสฺส ภาสิตํ อภินนฺทุนฺติ.\csromanspacer{}\csromanspacer{}ทุติยํ.}

\csromanheader{2}{3. หาลิทฺทิกานิสุตฺต}
\proseitem{3}{เอวํ เม สุตํ\csromandash{}เอกํ สมยํ อายสฺมา มหากจฺจาโน อวนฺตีสุ วิหรติ กุรรฆเร\footnote{กุลฆเร (ก)} ปปาเต ปพฺพเต.\csromanspacer{} อถ โข หาลิทฺทิกานิ\footnote{หาลิทฺทกานิ (สี), หลิทฺทิกานิ (สฺยา)} คหปติ เยนายสฺมา มหากจฺจาโน เตนุปสงฺกมิ, อุปสงฺกมิตฺวา อายสฺมนฺตํ มหากจฺจานํ อภิวาเทตฺวา เอกมนฺตํ นิสีทิ, เอกมนฺตํ นิสินฺโน โข หาลิทฺทิกานิ คหปติ อายสฺมนฺตํ มหากจฺจานํ เอตทโวจ “วุตฺตมิทํ ภนฺเต ภควตา อฏฺฐกวคฺคิเย มาคณฺฑิยปเญฺห\csromandash{}}

\csromangathameasure{‘โอกํ ปหาย อนิเกตสารี, \\ คาเม อกุพฺพํ มุนิ สนฺถวานิ. \\ กาเมหิ ริตฺโต อปุรกฺขราโน, \\ กถํ น วิคฺคยฺห ชเนน กยิรา’ติ.}
\csromangathagroup{\csromangathastanza{‘โอกํ ปหาย อนิเกตสารี, \\ คาเม อกุพฺพํ\footnote{อกฺรุพฺพํ (ก)} มุนิ สนฺถวานิ\footnote{สนฺธวานิ (ก)}. \\ กาเมหิ ริตฺโต อปุรกฺขราโน\footnote{อกฺรุพฺพํ (ก)}, \\ กถํ น วิคฺคยฺห ชเนน กยิรา’ติ.}}

\prose{อิมสฺส นุ โข ภนฺเต ภควตา สํขิตฺเตน ภาสิตสฺส กถํ วิตฺถาเรน อตฺโถ ทฏฺฐพฺโพ”ติ.}

\prose{\csromanfolio{9}รูปธาตุ โข คหปติ วิญฺญาณสฺส โอโก, รูปธาตุราควินิพนฺธญฺจ\footnote{...วินิพทฺธญฺจ (อิ, สี-ฏฺฐ)} ปน วิญฺญาณํ “โอกสารี”ติ วุจฺจติ.\csromanspacer{} เวทนาธาตุ โข คหปติ วิญฺญาณสฺส โอโก, เวทนาธาตุราควินิพนฺธญฺจ ปน วิญฺญาณํ “โอกสารี”ติ วุจฺจติ.\csromanspacer{} สญฺญาธาตุ โข คหปติ วิญฺญาณสฺส โอโก, สญฺญาธาตุราควินิพนฺธญฺจ ปน วิญฺญาณํ “โอกสารี”ติ วุจฺจติ.\csromanspacer{} สงฺขารธาตุ โข คหปติ วิญฺญาณสฺส โอโก, สงฺขารธาตุราควินิพนฺธญฺจ ปน วิญฺญาณํ “โอกสารี”ติ วุจฺจติ.\csromanspacer{} เอวํ โข คหปติ โอกสารี โหติ.}

\prose{กถญฺจ คหปติ อโนกสารี โหติ.\csromanspacer{} รูปธาตุยา โข คหปติ โย ฉนฺโท โย ราโค ยา นนฺที\footnote{นนฺทิ (สี, สฺยา, กํ, อิ)} ยา ตณฺหา เย อุปยุปาทานา\footnote{อุปายุปาทานา (สี, สฺยา, กํ, อิ)} เจตโส อธิฏฺฐานาภินิเวสานุสยา, เต ตถาคตสฺส ปหีนา อุจฺฉินฺนมูลา ตาลาวตฺถุกตา อนภาวํกตา\footnote{อนภาวกตา (สี, อิ), อนภาวํคตา (สฺยา, กํ)} อายติํ อนุปฺปาทธมฺมา, ตสฺมา ตถาคโต “อโนกสารี”ติ วุจฺจติ.\csromanspacer{} เวทนาธาตุยา โข คหปติ.\csromanspacer{} สญฺญาธาตุยา โข คหปติ.\csromanspacer{} สงฺขารธาตุยา โข คหปติ.\csromanspacer{} วิญฺญาณธาตุยา โข คหปติ โย ฉนฺโท โย ราโค ยา นนฺที ยา ตณฺหา เย อุปยุปาทานา เจตโส อธิฏฺฐานาภินิเวสานุสยา, เต ตถาคตสฺส ปหีนา อุจฺฉินฺนมูลา ตาลาวตฺถุกตา อนภาวํกตา อายติํ อนุปฺปาทธมฺมา, ตสฺมา ตถาคโต “อโนกสารี”ติ วุจฺจติ.\csromanspacer{} เอวํ โข คหปติ อโนกสารี โหติ.}

\prose{กถญฺจ คหปติ นิเกตสารี โหติ.\csromanspacer{} รูปนิมิตฺตนิเกตวิสารวินิพนฺธา โข คหปติ “นิเกตสารี”ติ วุจฺจติ.\csromanspacer{} สทฺทนิมิตฺต ฯเปฯ.\csromanspacer{} คนฺธนิมิตฺต.\csromanspacer{} รสนิมิตฺต.\csromanspacer{} โผฏฺฐพฺพนิมิตฺต.\csromanspacer{} ธมฺมนิมิตฺตนิเกตวิสารวินิพนฺธา โข คหปติ “นิเกตสารี”ติ วุจฺจติ.\csromanspacer{} เอวํ โข คหปติ นิเกตสารีโหติ.}

\prose{กถญฺจ คหปติ อนิเกตสารี โหติ.\csromanspacer{} รูปนิมิตฺตนิเกตวิสารวินิพนฺธา โข คหปติ ตถาคตสฺส ปหีนา อุจฺฉินฺนมูลา ตาลาวตฺถุกตา อนภาวํกตา อายติํ อนุปฺปาทธมฺมา, ตสฺมา ตถาคโต “อนิเกตสารี”ติ วุจฺจติ.\csromanspacer{} สทฺทนิมิตฺต.\csromanspacer{} คนฺธนิมิตฺต.\csromanspacer{} รสนิมิตฺต.\csromanspacer{} โผฏฺฐพฺพนิมิตฺต.\csromanspacer{} ธมฺมนิมิตฺตนิเกตวิสารวินิพนฺธา โข คหปติ ตถาคตสฺส ปหีนา อุจฺฉินฺนมูลา ตาลาวตฺถุกตา อนภาวํกตา อายติํ อนุปฺปาทธมฺมา, \csromanfolio{10}ตสฺมา ตถาคโต “อนิเกตสารี”ติ วุจฺจติ.\csromanspacer{} เอวํ โข คหปติ อนิเกตสารี โหติ.}

\prose{กถญฺจ คหปติ คาเม สนฺถวชาโต\footnote{สนฺธวชาโต (ก)} โหติ.\csromanspacer{} อิธ คหปติ เอกจฺโจ คิหีหิ\footnote{คิหิ (ก)} สํสฏฺโฐ วิหรติ สหนนฺที สหโสกี สุขิเตสุ สุขิโต ทุกฺขิเตสุ ทุกฺขิโต อุปฺปนฺเนสุ กิจฺจกรณีเยสุ อตฺตนา เตสุ โยคํ อาปชฺชติ.\csromanspacer{} เอวํ โข คหปติ คาเม สนฺถวชาโต โหติ.}

\prose{กถญฺจ คหปติ คาเม น สนฺถวชาโต โหติ.\csromanspacer{} อิธ คหปติ ภิกฺขุ คิหีหิ2 อสํสฏฺโฐ วิหรติ น สหนนฺที น สหโสกี น สุขิเตสุ สุขิโต น ทุกฺขิเตสุ ทุกฺขิโต อุปฺปนฺเนสุ กิจฺจกรณีเยสุ น อตฺตนา เตสุ โยคํ อาปชฺชติ.\csromanspacer{} เอวํ โข คหปติ คาเม น สนฺถวชาโต โหติ.}

\prose{กถญฺจ คหปติ กาเมหิ อริตฺโต โหติ.\csromanspacer{} อิธ คหปติ เอกจฺโจ กาเมสุ อวิคตราโค โหติ อวิคตฉนฺโท อวิคตเปโม อวิคตปิปาโส อวิคตปริฬาโห อวิคตตโณฺห.\csromanspacer{} เอวํ โข คหปติ กาเมหิ อริตฺโต โหติ.}

\prose{กถญฺจ คหปติ กาเมหิ ริตฺโต โหติ.\csromanspacer{} อิธ คหปติ เอกจฺโจ กาเมสุ วิคตราโค โหติ วิคตฉนฺโท วิคตเปโม วิคตปิปาโส วิคตปริฬาโห วิคตตโณฺห.\csromanspacer{} เอวํ โข คหปติ กาเมหิ ริตฺโต โหติ.}

\prose{กถญฺจ คหปติ ปุรกฺขราโน โหติ.\csromanspacer{} อิธ คหปติ เอกจฺจสฺส เอวํ โหติ “เอวํรูโป สิยํ อนาคตมทฺธานํ.\csromanspacer{} เอวํเวทโน สิยํ อนาคตมทฺธานํ.\csromanspacer{} เอวํสญฺโญ สิยํ อนาคตมทฺธานํ.\csromanspacer{} เอวํสงฺขาโร สิยํ อนาคตมทฺธานํ.\csromanspacer{} เอวํวิญฺญาโณ สิยํ อนาคตมทฺธานนฺ”ติ.\csromanspacer{} เอวํ โข คหปติ ปุรกฺขราโน โหติ.}

\prose{กถญฺจ คหปติ อปุรกฺขราโน โหติ.\csromanspacer{} อิธ คหปติ เอกจฺจสฺส น เอวํ โหติ “เอวํรูโป สิยํ อนาคตมทฺธานํ.\csromanspacer{} เอวํเวทโน สิยํ อนาคตมทฺธานํ.\csromanspacer{} เอวํสญฺโญ สิยํ อนาคตมทฺธานํ.\csromanspacer{} เอวํสงฺขาโร สิยํ อนาคตมทฺธานํ.\csromanspacer{} เอวํวิญฺญาโณ สิยํ อนาคตมทฺธานนฺ”ติ.\csromanspacer{} เอวํ โข คหปติ อปุรกฺขราโน โหติ.}

\prose{\csromanfolio{11}กถญฺจ คหปติ กถํ วิคฺคยฺห ชเนน กตฺตา โหติ.\csromanspacer{} อิธ คหปติ เอกจฺโจ เอวรูปิํ กถํ กตฺตา โหติ “น ตฺวํ อิมํ ธมฺมวินยํ อาชานาสิ, อหํ อิมํ ธมฺมวินยํ อาชานามิ, กิํ ตฺวํ อิมํ ธมฺมวินยํ อาชานิสฺสสิ, มิจฺฉาปฏิปนฺโน ตฺวมสิ, อหมสฺมิ สมฺมาปฏิปนฺโน, ปุเร วจนียํ ปจฺฉา อวจ, ปจฺฉา วจนียํ ปุเร อวจ, สหิตํ เม, อสหิตํ เต, อธิจิณฺณํ เต วิปราวตฺตํ, อาโรปิโต เต วาโท, จร วาทปฺปโมกฺขาย, นิคฺคหิโตสิ, นิพฺเพเฐหิ วา สเจ ปโหสี”ติ.\csromanspacer{} เอวํ โข คหปติ กถํ วิคฺคยฺห ชเนน กตฺตา โหติ.}

\prose{กถญฺจ คหปติ กถํ น วิคฺคยฺห ชเนน กตฺตา โหติ.\csromanspacer{} อิธ คหปติ ภิกฺขุ น เอวรูปิํ กถํ กตฺตา โหติ “น ตฺวํ อิมํ ธมฺมวินยํ อาชานาสิ ฯเปฯ นิพฺเพเฐหิ วา สเจ ปโหสี”ติ.\csromanspacer{} เอวํ โข คหปติ กถํ น วิคฺคยฺห ชเนน กตฺตา โหติ.}

\prose{อิติ โข คหปติ ยํ ตํ วุตฺตํ ภควตา อฏฺฐกวคฺคิเย มาคณฺฑิยปเญฺห\csromandash{}}

\csromangathameasure{“โอกํ ปหาย อนิเกตสารี, \\ คาเม อกุพฺพํ มุนิ สนฺถวานิ. \\ กาเมหิ ริตฺโต อปุรกฺขราโน, \\ กถํ น วิคฺคยฺห ชเนน กยิรา”ติ.}
\csromangathagroup{\csromangathastanza{“โอกํ ปหาย อนิเกตสารี, \\ คาเม อกุพฺพํ มุนิ สนฺถวานิ. \\ กาเมหิ ริตฺโต อปุรกฺขราโน, \\ กถํ น วิคฺคยฺห ชเนน กยิรา”ติ.}}

\prose{อิมสฺส โข คหปติ ภควตา สํขิตฺเตน ภาสิตสฺส เอวํ วิตฺถาเรน อตฺโถ ทฏฺฐพฺโพติ.\csromanspacer{}\csromanspacer{}ตติยํ.}

\csromanheader{2}{4. ทุติยหาลิทฺทิกานิสุตฺต}
\proseitem{4}{เอวํ เม สุตํ\csromandash{}เอกํ สมยํ อายสฺมา มหากจฺจาโน อวนฺตีสุ วิหรติ กุรรฆเร ปปาเต ปพฺพเต.\csromanspacer{} อถ โข หาลิทฺทิกานิ คหปติ เยนายสฺมา มหากจฺจาโน ฯเปฯ เอกมนฺตํ นิสินฺโน โข หาลิทฺทิกานิ คหปติ อายสฺมนฺตํ มหากจฺจานํ เอตทโวจ “วุตฺตมิทํ ภนฺเต ภควตา สกฺกปเญฺห\csromandash{}’เย เต สมณพฺราหฺมณา ตณฺหาสงฺขยวิมุตฺตา, เต อจฺจนฺตนิฏฺฐา อจฺจนฺตโยคกฺเขมิโน อจฺจนฺตพฺรหฺมจาริโน อจฺจนฺตปริโยสานา เสฏฺฐา เทวมนุสฺสานนฺ’ติ.}

\prose{อิมสฺส นุ โข ภนฺเต ภควตา สํขิตฺเตน ภาสิตสฺส กถํ วิตฺถาเรน อตฺโถ ทฏฺฐพฺโพ”ติ.}

\prose{\csromanfolio{12}รูปธาตุยา โข คหปติ โย ฉนฺโท โย ราโค ยา นนฺที ยา ตณฺหา เย อุปยุปาทานา เจตโส อธิฏฺฐานาภินิเวสานุสยา, เตสํ ขยา วิราคา นิโรธา จาคา ปฏินิสฺสคฺคา “จิตฺตํ สุวิมุตฺตนฺ”ติ วุจฺจติ.}

\prose{เวทนาธาตุยา โข คหปติ ฯเปฯ.\csromanspacer{} สญฺญาธาตุยา โข คหปติ.\csromanspacer{} สงฺขารธาตุยา โข คหปติ.\csromanspacer{} วิญฺญาณธาตุยา โข คหปติ โย ฉนฺโท โย ราโค ยา นนฺที ยา ตณฺหา เย อุปยุปาทานา เจตโส อธิฏฺฐานาภินิเวสานุสยา, เตสํ ขยา วิราคา นิโรธา จาคา ปฏินิสฺสคฺคา “จิตฺตํ สุวิมุตฺตนฺ”ติ วุจฺจติ.}

\prose{อิติ โข คหปติ ยํ ตํ วุตฺตํ ภควตา สกฺกปเญฺห\csromandash{}“เย เต สมณพฺราหฺมณา ตณฺหาสงฺขยวิมุตฺตา, เต อจฺจนฺตนิฏฺฐา อจฺจนฺตโยคกฺเขมิโน อจฺจนฺตพฺรหฺมจาริโน อจฺจนฺตปริโยสานา เสฏฺฐา เทวมนุสฺสานนฺ”ติ.}

\prose{อิมสฺส โข คหปติ ภควตา สํขิตฺเตน ภาสิตสฺส เอวํ วิตฺถาเรน อตฺโถ ทฏฺฐพฺโพติ.\csromanspacer{}\csromanspacer{}จตุตฺถํ.}

\csromanheader{2}{5. สมาธิสุตฺต}
\proseitem{5}{เอวํ เม สุตํ ฯเปฯ.\csromanspacer{} สาวตฺถิยํ ตตฺร โข ฯเปฯ เอตทโวจ\csromandash{} สมาธิํ ภิกฺขเว ภาเวถ, สมาหิโต ภิกฺขเว ภิกฺขุ ยถาภูตํ ปชานาติ.\csromanspacer{} กิญฺจ ยถาภูตํ ปชานาติ, รูปสฺส สมุทยญฺจ อตฺถงฺคมญฺจ.\csromanspacer{} เวทนาย สมุทยญฺจ อตฺถงฺคมญฺจ.\csromanspacer{} สญฺญาย สมุทยญฺจ อตฺถงฺคมญฺจ.\csromanspacer{} สงฺขารานํ สมุทยญฺจ อตฺถงฺคมญฺจ.\csromanspacer{} วิญฺญาณสฺส สมุทยญฺจ อตฺถงฺคมญฺจ.}

\prose{โก จ ภิกฺขเว รูปสฺส สมุทโย, โก เวทนาย สมุทโย, โก สญฺญาย สมุทโย, โก สงฺขารานํ สมุทโย, โก วิญฺญาณสฺส สมุทโย.\csromanspacer{} อิธ ภิกฺขเว ภิกฺขุ อภินนฺทติ อภิวทติ อชฺโฌสาย ติฏฺฐติ.}

\prose{กิญฺจ อภินนฺทติ อภิวทติ อชฺโฌสาย ติฏฺฐติ.\csromanspacer{} รูปํ อภินนฺทติ อภิวทติ อชฺโฌสาย ติฏฺฐติ, ตสฺส รูปํ อภินนฺทโต อภิวทโต อชฺโฌสาย ติฏฺฐโต อุปฺปชฺชติ นนฺที, ยา รูเป นนฺที, ตทุปาทานํ, ตสฺสุปาทานปจฺจยา ภโว, ภวปจฺจยา ชาติ, ชาติปจฺจยา ชรามรณํ โสกปริเทวทุกฺขโทมนสฺสุปายาสา สมฺภวนฺติ.\csromanspacer{} เอวเมตสฺส เกวลสฺส ทุกฺขกฺขนฺธสฺส สมุทโย โหติ.}

\prose{\csromanfolio{13}เวทนํ อภินนฺทติ ฯเปฯ.\csromanspacer{} สญฺญํ อภินนฺทติ.\csromanspacer{} สงฺขาเร อภินนฺทติ.\csromanspacer{} วิญฺญาณํ อภินนฺทติ อภิวทติ อชฺโฌสาย ติฏฺฐติ, ตสฺส วิญฺญาณํ อภินนฺทโต อภิวทโต อชฺโฌสาย ติฏฺฐโต อุปฺปชฺชติ นนฺที, ยา วิญฺญาเณ นนฺที, ตทุปาทานํ, ตสฺสุปาทานปจฺจยา ภโว, ภวปจฺจยา ชาติ, ชาติปจฺจยา ฯเปฯ.\csromanspacer{} เอวเมตสฺส เกวลสฺส ทุกฺขกฺขนฺธสฺส สมุทโย โหติ.}

\prose{อยํ ภิกฺขเว รูปสฺส สมุทโย.\csromanspacer{} อยํ เวทนาย สมุทโย.\csromanspacer{} อยํ สญฺญาย สมุทโย.\csromanspacer{} อยํ สงฺขารานํ สมุทโย.\csromanspacer{} อยํ วิญฺญาณสฺส สมุทโย.}

\prose{โก จ ภิกฺขเว รูปสฺส อตฺถงฺคโม.\csromanspacer{} โก เวทนาย.\csromanspacer{} โก สญฺญาย.\csromanspacer{} โก สงฺขารานํ.\csromanspacer{} โก วิญฺญาณสฺส อตฺถงฺคโม.\csromanspacer{} อิธ ภิกฺขเว นาภินนฺทติ นาภิวทติ นาชฺโฌสาย ติฏฺฐติ.}

\prose{กิญฺจ นาภินนฺทติ นาภิวทติ นาชฺโฌสาย ติฏฺฐติ.\csromanspacer{} รูปํ นาภินนฺทติ นาภิวทติ นาชฺโฌสาย ติฏฺฐติ, ตสฺส รูปํ อนภินนฺทโต อนภิวทโต อนชฺโฌสาย ติฏฺฐโต ยา รูเป นนฺที, สา นิรุชฺฌติ, ตสฺส นนฺทีนิโรธา อุปาทานนิโรโธ, อุปทานนิโรธา ภวนิโรโธ ฯเปฯ.\csromanspacer{} เอวเมตสฺส เกวลสฺส ทุกฺขกฺขนฺธสฺส นิโรโธ โหติ.}

\prose{เวทนํ นาภินนฺทติ นาภิวทติ นาชฺโฌสาย ติฏฺฐติ, ตสฺส เวทนํ อนภินนฺทโต อนภิวทโต อนชฺโฌสาย ติฏฺฐโต ยา เวทนาย นนฺที, สา นิรุชฺฌติ, ตสฺส นนฺทีนิโรธา อุปาทานนิโรโธ, อุปาทานนิโรธา ภวนิโรโธ ฯเปฯ.\csromanspacer{} เอวเมตสฺส เกวลสฺส ทุกฺขกฺขนฺธสฺส นิโรโธ โหติ.}

\prose{สญฺญํ นาภินนฺทติ ฯเปฯ.\csromanspacer{} สงฺขาเร นาภินนฺทติ นาภิวทติ นาชฺโฌสาย ติฏฺฐติ, ตสฺส สงฺขาเร อนภินนฺทโต อนภิวทโต อนชฺโฌสาย ติฏฺฐโต ยา สงฺขาเรสุ นนฺที, สา นิรุชฺฌติ, ตสฺส นนฺทีนิโรธา อุปาทานนิโรโธ, อุปาทานนิโรธา ภวนิโรโธ ฯเปฯ.\csromanspacer{} เอวเมตสฺส เกวลสฺส ทุกฺขกฺขนฺธสฺส นิโรโธ โหติ.}

\prose{วิญฺญาณํ นาภินนฺทติ นาภิวทติ นาชฺโฌสาย ติฏฺฐติ, ตสฺส วิญฺญาณํ อนภินนฺทโต อนภิวทโต อนชฺโฌสาย ติฏฺฐโต ยา วิญฺญาเณ นนฺที, สา นิรุชฺฌติ, ตสฺส นนฺทีนิโรธา อุปาทานนิโรโธ ฯเปฯ.\csromanspacer{} เอวเมตสฺส เกวลสฺส ทุกฺขกฺขนฺทสฺส นิโรโธ โหติ.}

\prose{\csromanfolio{14}อยํ ภิกฺขเว รูปสฺส อตฺถงฺคโม.\csromanspacer{} อยํ เวทนาย อตฺถงฺคโม.\csromanspacer{} อยํ สญฺญาย อตฺถงฺคโม.\csromanspacer{} อยํ สงฺขารานํ อตฺถงฺคโม.\csromanspacer{} อยํ วิญฺญาณสฺส อตฺถงฺคโมติ.\csromanspacer{}\csromanspacer{}ปญฺจมํ.}

\csromanheader{2}{6. ปฏิสลฺลาณสุตฺต}
\proseitem{6}{สาวตฺถินิทานํ.\csromanspacer{} ปฏิสลฺลาเณ ภิกฺขเว โยคมาปชฺชถ, ปฏิสลฺลีโณ ภิกฺขเว ภิกฺขุ ยถาภูตํ ปชานาติ.\csromanspacer{} กิญฺจ ยถาภูตํ ปชานาติ, รูปสฺส สมุทยญฺจ อตฺถงฺคมญฺจ.\csromanspacer{} เวทนาย สมุทยญฺจ อตฺถงฺคมญฺจ.\csromanspacer{} สญฺญาย สมุทยญฺจ อตฺถงฺคมญฺจ.\csromanspacer{} สงฺขารานํ สมุทยญฺจ อตฺถงฺคมญฺจ.\csromanspacer{} วิญฺญาณสฺส สมุทยญฺจ อตฺถงฺคมญฺจ. (ยถา ปฐมสุตฺเต, ตถา วิตฺถาเรตพฺโพ.).\csromanspacer{} ฉฏฺฐํ.}

\csromanheader{2}{7. อุปาทาปริตสฺสนาสุตฺต}
\proseitem{7}{สาวตฺถินิทานํ.\csromanspacer{} อุปาทาปริตสฺสนญฺจ โว ภิกฺขเว เทเสสฺสามิ อนุปาทา-อปริตสฺสนญฺจ, ตํ สุณาถ สาธุกํ มนสิ กโรถ, ภาสิสฺสามีติ. “เอวํ ภนฺเต”ติ โข เต ภิกฺขู ภควโต ปจฺจสฺโสสุํ.\csromanspacer{} ภควา เอตทโวจ\csromandash{}}

\prose{กถญฺจ ภิกฺขเว อุปาทาปริตสฺสนา โหติ.\csromanspacer{} อิธ ภิกฺขเว อสฺสุตวา ปุถุชฺชโน อริยานํ อทสฺสาวี อริยธมฺมสฺส อโกวิโท อริยธมฺเม อวินีโต, สปฺปุริสานํ อทสฺสาวี สปฺปุริสธมฺมสฺส อโกวิโท สปฺปุริสธมฺเม อวินีโต รูปํ อตฺตโต สมนุปสฺสติ, รูปวนฺตํ วา อตฺตานํ, อตฺตนิ วา รูปํ, รูปสฺมิํ วา อตฺตานํ, ตสฺส ตํ รูปํ วิปริณมติ อญฺญถา โหติ, ตสฺส รูปวิปริณามญฺญถาภาวา รูปวิปริณามานุปริวตฺติ วิญฺญาณํ โหติ.\csromanspacer{} ตสฺส รูปวิปริณามานุปริวตฺติชา ปริตสฺสนา ธมฺมสมุปฺปาทา จิตฺตํ ปริยาทาย ติฏฺฐนฺติ.\csromanspacer{} เจตโส ปริยาทานา อุตฺตาสวา จ โหติ วิฆาตวา จ อเปกฺขวา จ, อุปาทาย จ ปริตสฺสติ.}

\prose{เวทนํ อตฺตโต สมนุปสฺสติ, เวทนาวนฺตํ วา อตฺตานํ, อตฺตนิ วา เวทนํ, เวทนาย วา อตฺตานํ, ตสฺส สา เวทนา วิปริณมติ อญฺญถา โหติ, ตสฺส เวทนาวิปริณามญฺญถาภาวา เวทนาวิปริณามานุปริวตฺติ วิญฺญาณํ โหติ, ตสฺส เวทนาวิปริณามานุปริวตฺติชา ปริตสฺสนา ธมฺมสมุปฺปาทา จิตฺตํ ปริยาทาย ติฏฺฐนฺติ, เจตโส ปริยาทานา อุตฺตาสวา จ โหติ วิฆาตวา จ อเปกฺขวา จ, อุปาทาย จ ปริตสฺสติ.}

\prose{\csromanfolio{15}สญฺญํ อตฺตโต สมนุปสฺสติ ฯเปฯ.\csromanspacer{} สงฺขาเร อตฺตโต สมนุปสฺสติ, สงฺขารวนฺตํ วา อตฺตานํ, อตฺตนิ วา สงฺขาเร, สงฺขาเรสุ วา อตฺตานํ, ตสฺส เต สงฺขารา วิปริณมนฺติ อญฺญถา โหนฺติ, ตสฺส สงฺขารวิปริณามญฺญถาภาวา สงฺขารวิปริณามานุปริวตฺติ วิญฺญาณํ โหติ, ตสฺส สงฺขารวิปริณามานุปริวตฺติชา ปริตสฺสนา ธมฺมสมุปฺปาทา จิตฺตํ ปริยาทาย ติฏฺฐนฺติ, เจตโส ปริยาทานา อุตฺตาสวา จ โหติ วิฆาตวา จ อเปกฺขวา จ, อุปาทาย จ ปริตสฺสติ.}

\prose{วิญฺญาณํ อตฺตโต สมนุปสฺสติ, วิญฺญาณวนฺตํ วา อตฺตานํ, อตฺตนิ วา วิญฺญาณํ, วิญฺญาณสฺมิํ วา อตฺตานํ, ตสฺส ตํ วิญฺญาณํ วิปริณมติ อญฺญถา โหติ, ตสฺส วิญฺญาณวิปริณามญฺญถาภาวา วิญฺญาณวิปริณามานุปริวตฺติ วิญฺญาณํ โหติ, ตสฺส วิญฺญาณวิปริณามานุปริวตฺติชา ปริตสฺสนา ธมฺมสมุปฺปาทา จิตฺตํ ปริยาทาย ติฏฺฐนฺติ, เจตโส ปริยาทานา อุตฺตาสวา จ โหติ วิฆาตวา จ อเปกฺขวา จ, อุปาทาย จ ปริตสฺสติ.\csromanspacer{} เอวํ โข ภิกฺขเว อุปาทาปริตสฺสนา โหติ.}

\prose{กถญฺจ ภิกฺขเว อนุปาทา-อปริตสฺสนา โหติ.\csromanspacer{} อิธ ภิกฺขเว สุตวา อริยสาวโก อริยานํ ทสฺสาวี อริยธมฺมสฺส โกวิโท อริยธมฺเม สุวินีโต, สปฺปุริสานํ ทสฺสาวี สปฺปุริสธมฺมสฺส โกวิโท สปฺปุริสธมฺเม สุวินีโต น รูปํ อตฺตโต สมนุปสฺสติ, น รูปวนฺตํ วา อตฺตานํ, น อตฺตนิ วา รูปํ, น รูปสฺมิํ วา อตฺตานํ, ตสฺส ตํ รูปํ วิปริณมติ อญฺญถา โหติ, ตสฺส รูปวิปริณามญฺญถาภาวา น รูปวิปริณามานุปริวตฺติ วิญฺญาณํ โหติ, ตสฺส น รูปวิปริณามานุปริวตฺติชา ปริตสฺสนา ธมฺมสมุปฺปาทา จิตฺตํ ปริยาทาย ติฏฺฐนฺติ, เจตโส อปริยาทานา น เจวุตฺตาสวา\footnote{น เจว อุตฺตาสวา (อิ, ก)} โหติ น จ วิฆาตวา น จ อเปกฺขวา, อนุปาทาย จ น ปริตสฺสติ.}

\prose{น เวทนํ อตฺตโต สมนุปสฺสติ, น เวทนาวนฺตํ วา อตฺตานํ, น อตฺตนิ วา เวทนํ, น เวทนาย วา อตฺตานํ, ตสฺส สา เวทนา วิปริณมติ อญฺญถา โหติ, ตสฺส เวทนาวิปริณามญฺญถาภาวา น เวทนาวิปริณามานุปริวตฺติ วิญฺญาณํ โหติ, ตสฺส น เวทนาวิปริณามานุปริวตฺติชา ปริตสฺสนา ธมฺมสมุปฺปาทา จิตฺตํ ปริยาทาย \csromanfolio{16}ติฏฺฐนฺติ, เจตโส อปริยาทานา น เจวุตฺตาสวา โหติ น จ วิฆาตวา น จ อเปกฺขวา, อนุปาทาย จ น ปริตสฺสติ.}

\prose{น สญฺญํ ฯเปฯ.\csromanspacer{} น สงฺขาเร อตฺตโต สมนุปสฺสติ, น สงฺขารวนฺตํ วา อตฺตานํ, น อตฺตนิ วา สงฺขาเร, น สงฺขาเรสุ วา อตฺตานํ, ตสฺส เต สงฺขารา วิปริณมนฺติ อญฺญถาโหนฺติ, ตสฺส สงฺขารวิปริณมญฺญถาภาวา นสงฺขารวิปริณามานุปริวตฺติ วิญฺญาณํ โหติ, ตสฺส น สงฺขารวิปริณามานุปริวตฺติชา ปริตสฺสนา ธมฺมสมุปฺปาทา จิตฺตํ ปริยาทาย ติฏฺฐนฺติ, เจตโส อปริยาทานา น เจวุตฺตาสวา โหติ น จ วิฆาตวา น จ อเปกฺขวา, อนุปาทาย จ น ปริตสฺสติ.}

\prose{น วิญฺญาณํ อตฺตโต สมนุปสฺสติ, น วิญฺญาณวนฺตํ วา อตฺตานํตฺ ฯเปฯ ตสฺส ตํ วิญฺญาณํ วิปริณมติ อญฺญถา โหติ, ตสฺส วิญฺญาณวิปริณามญฺญถาภาวา น วิญฺญาณวิปริณามานุปริวตฺติ วิญฺญาณํ โหติ, ตสฺส น วิญฺญาณวิปริณามานุปริวตฺติชา ปริตสฺสนา ธมฺมสมุปฺปาทา จิตฺตํ ปริยาทาย ติฏฺฐนฺติ, เจตโส อปริยาทานา น เจวุตฺตาสวา โหติ น จ วิฆาตวา น จ อเปกฺขวา, อนุปาทาย จ น ปริตสฺสติ.\csromanspacer{} เอวํ โข ภิกฺขเว อนุปาทา-อปริตสฺสนํ โหตีติ.\csromanspacer{}\csromanspacer{}สตฺตมํ.}

\csromanheader{2}{8. ทุติย-อุปาทาปริตสฺสนาสุตฺต}
\proseitem{8}{สาวตฺถินิทานํ.\csromanspacer{} อุปาทาปริตสฺสนญฺจ โว ภิกฺขเว เทเสสฺสามิ อนุปาทา-อปริตสฺสนญฺจ, ตํ สุณาถ ฯเปฯ.\csromanspacer{} กถญฺจ ภิกฺขเว อุปาทาปริตสฺสนา โหติ.\csromanspacer{} อิธ ภิกฺขเว อสฺสุตวา ปุถุชฺชโน รูปํ “เอตํ มม, เอโสหมสฺมิ, เอโส เม อตฺตา”ติ สมนุปสฺสติ, ตสฺส ตํ รูปํ วิปริณมติ อญฺญถา โหติ, ตสฺส รูปวิปริณามญฺญถาภาวา อุปฺปชฺชนฺติ โสกปริเทวทุกฺขโทมนสฺสุปายาสา.\csromanspacer{} เวทนํ “เอตํ มม ฯเปฯ สญฺญํ “เอตํ มม.\csromanspacer{} สงฺขาเร “เอตํ มม.\csromanspacer{} วิญฺญาณํ “เอตํ มม, เอโสหมสฺมิ, เอโส เม อตฺตา”ติ สมนุปสฺสติ, ตสฺส ตํ วิญฺญาณํ วิปริณมติ อญฺญถาโหติ, ตสฺส วิญฺญาณวิปริณามญฺญถาภาวา อุปฺปชฺชนฺติ โสกปริเทวทุกฺขโทมนสฺสุปายาสา.\csromanspacer{} เอวํ โข ภิกฺขเว อุปาทาปริตสฺสนา โหติ,}

\prose{กถญฺจ ภิกฺขเว อนุปาทา-อปริตสฺสนา โหติ.\csromanspacer{} อิธ ภิกฺขเว สุตวา อริยสาวโก รูปํ “เนตํ มม, เนโสหมสฺมิ, น เมโส อตฺตา”ติ สมนุปสฺสติ, ตสฺส ตํ รูปํ วิปริณมติ อญฺญถา โหติ, \csromanfolio{17}ตสฺส รูปวิปริณามญฺญถาภาวา นุปฺปชฺชนฺติ โสกปริเทวทุกฺขโทมนสฺสุปายาสา.\csromanspacer{} เวทนํ “เนตํ มม ฯเปฯ สญฺญํ “เนตํ มม.\csromanspacer{} สงฺขาเร “เนตํ มม.\csromanspacer{} วิญฺญาณํ “เนตํ มม, เนโสหมสฺมิ, น เมโส อตฺตา”ติ สมนุปสฺสติ, ตสฺส ตํ วิญฺญาณํ วิปริณมติ อญฺญถา โหติ, ตสฺส วิญฺญาณวิปริณามญฺญถาภาวา นุปฺปชฺชนฺติ โสกปริเทวทุกฺขโทมนสฺสุปายาสา.\csromanspacer{} เอวํ โข ภิกฺขเว อนุปาทา-อปริตสฺสนา โหตีติ.\csromanspacer{}\csromanspacer{}อฏฺฐมํ.}

\csromanheader{2}{9. กาลตฺตย-อนิจฺจสุตฺต}
\proseitem{9}{สาวตฺถินิทานํ.\csromanspacer{} รูปํ ภิกฺขเว อนิจฺจํ อตีตานาคตํ, โก ปน วาโท ปจฺจุปฺปนฺนสฺส.\csromanspacer{} เอวํ ปสฺสํ ภิกฺขเว สุตวา อริยสาวโก อตีตสฺมิํ รูปสฺมิํ อนเปกฺโข โหติ, อนาคตํ รูปํ นาภินนฺทติ, ปจฺจุปฺปนฺนสฺส รูปสฺส นิพฺพิทาย วิราคาย นิโรธาย ปฏิปนฺโน โหติ.\csromanspacer{} เวทนา อนิจฺจา ฯเปฯ.\csromanspacer{} สญฺญา อนิจฺจา.\csromanspacer{} สงฺขารา อนิจฺจา อตีตานาคตา, โก ปน วาโท ปจฺจุปฺปนฺนานํ.\csromanspacer{} เอวํ ปสฺสํ ภิกฺขเว สุตวา อริยสาวโก อตีเตสุ สงฺขาเรสุ อนเปกฺโข โหติ, อนาคเต สงฺขาเร นาภินนฺทติ, ปจฺจุปฺปนฺนานํ สงฺขารานํ นิพฺพิทาย วิราคาย นิโรธาย ปฏิปนฺโน โหติ.\csromanspacer{} วิญฺญาณํ อนิจฺจํ อตีตานาคตํ, โก ปน วาโท ปจฺจุปฺปนฺนสฺส.\csromanspacer{} เอวํ ปสฺสํ ภิกฺขเว สุตวา อริยสาวโก อตีตสฺมิํ วิญฺญาณสฺมิํ อนเปกฺโข โหติ, อนาคตํ วิญฺญาณํ นาภินนฺทติ, ปจฺจุปฺปนฺนสฺส วิญฺญาณสฺส นิพฺพิทาย วิราคาย นิโรธาย ปฏิปนฺโน โหตีติ.\csromanspacer{}\csromanspacer{}นวมํ.}

\csromanheader{2}{10. กาลตฺตยทุกฺขสุตฺต}
\proseitem{10}{สาวตฺถินิทานํ.\csromanspacer{} รูปํ ภิกฺขเว ทุกฺขํ อตีตานาคตํ, โก ปน วาโท ปจฺจุปฺปนฺนสฺส.\csromanspacer{} เอวํ ปสฺสํ ภิกฺขเว สุตวา อริยสาวโก อตีตสฺมิํ รูปสฺมิํ อนเปกฺโข โหติ, อนาคตํ รูปํ นาภินนฺทติ, ปจฺจุปฺปนฺนสฺส รูปสฺส นิพฺพิทาย วิราคาย นิโรธาย ปฏิปนฺโน โหติ.\csromanspacer{} เวทนา ทุกฺขา.\csromanspacer{} สญฺญา ทุกฺขา.\csromanspacer{} สงฺขารา ทุกฺขา.\csromanspacer{} วิญฺญาณํ ทุกฺขํ อตีตานาคตํ, โก ปน วาโท ปจฺจุปฺปนฺนสฺส.\csromanspacer{} เอวํ ปสฺสํ ภิกฺขเว สุตวา อริยสาวโก อตีตสฺมิํ วิญฺญาณสฺมิํ อนเปกฺโข โหติ, อนาคตํ วิญฺญาณํ นาภินนฺทติ, ปจฺจุปฺปนฺนสฺส วิญฺญาณสฺส นิพฺพิทาย วิราคาย นิโรธาย ปฏิปนฺโน โหตีติ.\csromanspacer{}\csromanspacer{}ทสมํ.}

\csromanheader{2}{11. กาลตฺตย-อนตฺตสุตฺต}
\proseitem{11}{\csromanfolio{18}สาวตฺถินิทานํ.\csromanspacer{} รูปํ ภิกฺขเว อนตฺตา อตีตานาคตํ, โก ปน วาโท ปจฺจุปฺปนฺนสฺส.\csromanspacer{} เอวํ ปสฺสํ ภิกฺขเว สุตวา อริยสาวโก อตีตสฺมิํ รูปสฺมิํ อนเปกฺโข โหติ, อนาคตํ รูปํ นาภินนฺทติ, ปจฺจุปฺปนฺนสฺส รูปสฺส นิพฺพิทาย วิราคาย นิโรธาย ปฏิปนฺโน โหติ.\csromanspacer{} เวทนา อนตฺตา.\csromanspacer{} สญฺญา อนตฺตา.\csromanspacer{} สงฺขารา อนตฺตา.\csromanspacer{} วิญฺญาณํ อนตฺตา อตีตานาคตํ, โก ปน วาโท ปจฺจุปฺปนฺนสฺส.\csromanspacer{} เอวํ ปสฺสํ ภิกฺขเว สุตวา อริยสาวโก อตีตสฺมิํ วิญฺญาณสฺมิํ อนเปกฺโข โหติ, อนาคตํ วิญฺญาณํ นาภินนฺทติ, ปจฺจุปฺปนฺนสฺส วิญฺญาณสฺส นิพฺพิทาย วิราคาย นิโรธาย ปฏิปนฺโน โหตีติ.\csromanspacer{}\csromanspacer{}เอกาทสมํ.\csromanspacer{} นกุลปิตุวคฺโค ปฐโม.}

\titlehead{\textbf{ตสฺสุทฺทานํ}}
\csromangathameasure{นกุลปิตา เทวทหา, เทฺวปิ หาลิทฺทิกานิ จ. \\ สมาธิปฏิสลฺลาณา, อุปาทาปริตสฺสนา ทุเว. \\ อตีตานาคตปจฺจุปฺปนฺนา, วคฺโค เตน ปวุจฺจติ.}
\csromangathagroup{\csromangathastanza{นกุลปิตา เทวทหา, เทฺวปิ หาลิทฺทิกานิ จ. \\ สมาธิปฏิสลฺลาณา, อุปาทาปริตสฺสนา ทุเว. \\ อตีตานาคตปจฺจุปฺปนฺนา, วคฺโค เตน ปวุจฺจติ.}}

\csromanmatikaanchor{mtk.1}
\csromanmatikaanchor{mtk.81}
\csromantocmarkref{section}{6. สุทฺธิกสุตฺต}{mtk.1}
\bookmark[dest={mtk.1},level=1]{6. สุทฺธิกสุตฺต}
\csromanheader{1}{2. อนิจฺจวคฺค}
\csromanheader{2}{1. อนิจฺจสุตฺต}
\proseitem{12}{เอวํ เม สุตํ\csromandash{}สาวตฺถิยํ.\csromanspacer{} ตตฺร โข ฯเปฯ.\csromanspacer{} รูปํ ภิกฺขเว อนิจฺจํ, เวทนา อนิจฺจา, สญฺญา อนิจฺจา, สงฺขารา อนิจฺจา, วิญฺญาณํ อนิจฺจํ.\csromanspacer{} เอวํ ปสฺสํ ภิกฺขเว สุตวา อริยสาวโก รูปสฺมิมฺปิ นิพฺพินฺทติ, เวทนายปิ นิพฺพินฺทติ, สญฺญายปิ นิพฺพินฺทติ, สงฺขาเรสุปิ นิพฺพินฺทติ, วิญฺญาณสฺมิมฺปิ นิพฺพินฺทติ, นิพฺพินฺทํ วิรชฺชติ, วิราคา วิมุจฺจติ, วิมุตฺตสฺมิํ “วิมุตฺตมฺ”อิติ ญาณํ โหติ, “ขีณา ชาติ, วุสิตํ พฺรหฺมจริยํ, กตํ กรณียํ นาปรํ อิตฺถตฺตายา”ติ ปชานาตีติ.\csromanspacer{}\csromanspacer{}ปฐมํ.}

\csromanheader{2}{2. ทุกฺขสุตฺต}
\proseitem{13}{สาวตฺถินิทานํ.\csromanspacer{} รูปํ ภิกฺขเว ทุกฺขํ, เวทนา ทุกฺขา, สญฺญา ทุกฺขา, สงฺขารา ทุกฺขา, วิญฺญาณํ ทุกฺขํ.\csromanspacer{} เอวํ ปสฺสํ ฯเปฯ นาปรํ อิตฺถตฺตายาติ ปชานาตีติ.\csromanspacer{}\csromanspacer{}ทุติยํ.}

\csromanheader{2}{3. อนตฺตสุตฺต}
\proseitem{14}{\csromanfolio{19}สาวตฺถินิทานํ.\csromanspacer{} รูปํ ภิกฺขเว อนตฺตา, เวทนา อนตฺตา, สญฺญา อนตฺตา, สงฺขารา อนตฺตา, วิญฺญาณํ อนตฺตา.\csromanspacer{} เอวํ ปสฺสํ ภิกฺขเว สุตวา อริยสาวโก รูปสฺมิมฺปิ นิพฺพินฺทติ, เวทนายปิ นิพฺพินฺทติ, สญฺญายปิ นิพฺพินฺทติ, สงฺขาเรสุปิ นิพฺพินฺทติ, วิญฺญาณสฺมิมฺปิ นิพฺพินฺทติ, นิพฺพินฺทิ วิรชฺชติ, วิราคา วิมุจฺจติ, วิมุตฺตสฺมิํ “วิมุตฺตมฺ”อิติ ญาณํ โหติ, “ขิณา ชาติ, วุสิตํ พฺรหฺมจริยํ, กตํ กรณียํ, นาปรํ อิตฺถตฺตายา”ติ ปชานาตีติ.\csromanspacer{}\csromanspacer{}ตติยํ.}

\csromanheader{2}{4. ยทนิจฺจสุตฺต}
\proseitem{15}{สาวตฺถินิทานํ.\csromanspacer{} รูปํ ภิกฺขเว อนิจฺจํ, ยทนิจฺจํ ตํ ทุกฺขํ, ยํ ทุกฺขํ ตทนตฺตา, ยทนตฺตา ตํ “เนตํ มม, เนโสหมสฺมิ, น เมโส อตฺตา”ติ เอวเมตํ ยถาภูตํ สมฺมปฺปญฺญาย ทฏฺฐพฺพํ.\csromanspacer{} เวทนา อนิจฺจา, ยทนิจฺจํ ตํ ทุกฺขํ, ยํ ทุกฺขํ ตทนตฺตา, ยทนตฺตา ตํ “เนตํ มม, เนโสหมสฺมิ, น เมโส อตฺตา”ติ เอวเมตํ ยถาภูตํ สมฺมปฺปญฺญาย ทฏฺฐพฺพํ.\csromanspacer{} สญฺญา อนิจฺจา ฯเปฯ.\csromanspacer{} สงฺขารา อนิจฺจา.\csromanspacer{} วิญฺญาณํ อนิจฺจํ, ยทนิจฺจํ ตํ ทุกฺขํ, ยํ ทุกฺขํ ตทนตฺตา, ยทนตฺตา ตํ “เนตํ มม, เนโสหมสฺมิ, น เมโส อตฺตา”ติ เอวเมตํ ยถาภูตํ สมฺมปฺปญฺญาย ทฏฺฐพฺพํ.\csromanspacer{} เอวํ ปสฺสํ ฯเปฯ นาปรํ อิตฺถตฺตายาติ ปชานาตีติ.\csromanspacer{}\csromanspacer{}จตุตฺถํ.}

\csromanheader{2}{5. ยํทุกฺขสุตฺต}
\proseitem{16}{สาวตฺถินิทานํ.\csromanspacer{} รูปํ ภิกฺขเว ทุกฺขํ ตทนตฺตา, ยทนตฺตา ตํ “เนตํ มม, เนโสหมสฺมิ, น เมโส อตฺตา”ติ, เอวเมตํ ยถาภูตํ สมฺมปฺปญฺญาย ทฏฺฐพฺพํ.\csromanspacer{} เวทนา ทุกฺขา.\csromanspacer{} สญฺญา ทุกฺขา.\csromanspacer{} สงฺขารา ทุกฺขา.\csromanspacer{} วิญฺญาณํ ทุกฺขํ, ยํ ทุกฺขํ ตทนตฺตา, ยทนตฺตา ตํ “เนตํ มม, เนโสหมสฺมิ, น เมโส อตฺตา”ติ เอวเมตํ ยถาภูตํ สมฺมปฺปญฺญาย ทฏฺฐพฺพํ.\csromanspacer{} เอวํ ปสฺสํ ฯเปฯ นาปรํ อิตฺถตฺตายาติ ปชานาตีติ.\csromanspacer{}\csromanspacer{}ปญฺจมํ.}

\csromanheader{2}{6. ยทนตฺตาสุตฺต}
\proseitem{17}{สาวตฺถินิทานํ.\csromanspacer{} รูปํ ภิกฺขเว อนตฺตา, ยทนตฺตา ตํ “เนตํ มม, เนโสหมสฺมิ, น เมโส อตฺตา”ติ เอวเมตํ ยถาภูตํ สมฺมปฺปญฺญาย ทฏฺฐพฺพํ.\csromanspacer{} เวทนา อนตฺตา.\csromanspacer{} สญฺญา อนตฺตา.\csromanspacer{} สงฺขารา อนตฺตา.\csromanspacer{} วิญฺญาณํ \csromanfolio{20}อนตฺตา, ยทนตฺตา ตํ “เนตํ มม, เนโสหมสฺมิ, น เมโส อตฺตา”ติ เอวเมตํ ยถาภูตํ สมฺมปฺปญฺญาย ทฏฺฐพฺพํ.\csromanspacer{} เอวํ ปสฺสํ ภิกฺขเว ฯเปฯ นาปรํ อิตฺถตฺตายาติ ปชานาตีติ.\csromanspacer{}\csromanspacer{}ฉฏฺฐํ.}

\csromanheader{2}{7. สเหตุ-อนิจฺจสุตฺต}
\proseitem{18}{สาวตฺถินิทานํ.\csromanspacer{} รูปํ ภิกฺขเว อนิจฺจํ, โยปิ เหตุ โยปิ ปจฺจโย รูปสฺส อุปฺปาทาย, โสปิ อนิจฺโจ, อนิจฺจสมฺภูตํ ภิกฺขเว รูปํ, กุโต นิจฺจํ ภวิสฺสติ.\csromanspacer{} เวทนา อนิจฺจา, โยปิ เหตุ โยปิ ปจฺจโย เวทนาย อุปฺปาทาย, โสปิ อนิจฺโจ, อนิจฺจสมฺภูตา ภิกฺขเว เวทนา, กุโต นิจฺจาภวิสฺสติ.\csromanspacer{} สญฺญา อนิจฺจา.\csromanspacer{} สงฺขารา อนิจฺจา, โยปิ เหตุ โยปิ ปจฺจโย สงฺขารานํ อุปฺปาทาย, โสปิ อนิจฺโจ, อนิจฺจสมฺภูตา ภิกฺขเวสงฺขารา, กุโต นิจฺจา ภวิสฺสนฺติ.\csromanspacer{} วิญฺญาณํ อนิจฺจํ, โยปิ เหตุ โยปิ ปจฺจโย วิญฺญาณสฺส อุปฺปาทาย, โสปิ อนิจฺโจ, อนิจฺจสมฺภูตํ ภิกฺขเว วิญฺญาณํ, กุโต นิจฺจํ ภวิสฺสติ.\csromanspacer{} เอวํ ปสฺสํ ฯเปฯ นาปรํ อิตฺถตฺตายาติ ปชานาตีติ.\csromanspacer{}\csromanspacer{}สตฺตมํ.}

\csromanheader{2}{8. สเหตุทุกฺขสุตฺต}
\proseitem{19}{สาวตฺถินิทานํ.\csromanspacer{} รูปํ ภิกฺขเว ทุกฺขํ, โยปิ เหตุ โยปิ ปจฺจโย รูปสฺส อุปฺปาทาย, โสปิ ทุกฺโข, ทุกฺขสมฺภูตํ ภิกฺขเว รูปํ, กุโต สุขํ ภวิสฺสติ.\csromanspacer{} เวทนา ทุกฺขา.\csromanspacer{} สญฺญา ทุกฺขา.\csromanspacer{} สงฺขารา ทุกฺขา.\csromanspacer{} วิญฺญาณํ ทุกฺขํ, โยปิ เหตุ โยปิ ปจฺจโย วิญฺญาณสฺส อุปฺปาทาย, โสปิ ทุกฺโข, ทุกฺขสมฺภูตํ ภิกฺขเว วิญฺญาณํ, กุโต สุขํ ภวิสฺสติ.\csromanspacer{} เอวํ ปสฺสํ ฯเปฯ นาปรํ อิตฺถตฺตายาติ ปชานาตีติ.\csromanspacer{}\csromanspacer{}อฏฺฐมํ.}

\csromanheader{2}{9. สเหตุ-อนตฺตสุตฺต}
\proseitem{20}{สาวตฺถินิทานํ.\csromanspacer{} รูปํ ภิกฺขเว อนตฺตา, โยปิ เหตุ โยปิ ปจฺจโย รูปสฺส อุปฺปาทาย, โสปิ อนตฺตา, อนตฺตสมฺภูตํ ภิกฺขเว รูปํ, กุโต อตฺตา ภวิสฺสติ.\csromanspacer{} เวทนา อนตฺตา.\csromanspacer{} สญฺญา อนตฺตา.\csromanspacer{} สงฺขารา อนตฺตา.\csromanspacer{} วิญฺญาณํ อนตฺตา, โยปิ เหตุ โยปิ ปจฺจโย วิญฺญาณสฺส อุปฺปาทาย, โสปิ อนตฺตา, อนตฺตสมฺภูตํ ภิกฺขเว วิญฺญาณํ, กุโต อตฺตา ภวิสฺสติ.\csromanspacer{} เอวํ ปสฺสํ ฯเปฯ นาปรํ อิตฺถตฺตายาติ ปชานาตีติ.\csromanspacer{}\csromanspacer{}นวมํ.}

\csromanheader{2}{10. อานนฺทสุตฺต}
\proseitem{21}{\csromanfolio{21}สาวตฺถิยํ อาราเม.\csromanspacer{} อถ โข อายสฺมา อานนฺโท เยน ภควา เตนุปสงฺกมิ, อุปสงฺกมิตฺวา ภควนฺตํ อภิวาเทตฺวา เอกมนฺตํ นิสีทิ, เอกมนฺตํ นิสินฺโน โข อายสฺมา อานนฺโท ภควนฺตํ เอตทโวจ “นิโรโธ นิโรโธติ ภนฺเต วุจฺจติ, กตเมสานํ โข ภนฺเต ธมฺมานํ นิโรโธ\footnote{นิโรธา (สี, อิ)} ‘นิโรโธ’ติ วุจฺจตี”ติ.\csromanspacer{} รูปํ โข อานนฺท อนิจฺจํ สงฺขตํ ปฏิจฺจสมุปฺปนฺนํ ขยธมฺมํ วยธมฺมํ วิราคธมฺมํ นิโรธธมฺมํ, ตสฺส นิโรโธ1 “นิโรโธ”ติ วุจฺจติ.\csromanspacer{} เวทนา อนิจฺจา สงฺขตา ปฏิจฺจสมุปฺปนฺนา ขยธมฺมา วยธมฺมา วิราคธมฺมา นิโรธธมฺมา, ตสฺสา นิโรโธ “นิโรโธ”ติ วุจฺจติ.\csromanspacer{} สญฺญา.\csromanspacer{} สงฺขารา อนิจฺจา สงฺขตา ปฏิจฺจสมุปฺปนฺนา ขยธมฺมา วยธมฺมา วิราคธมฺมา นิโรธธมฺมา, เตสํ นิโรโธ “นิโรโธ”ติ วุจฺจติ.\csromanspacer{} วิญฺญาณํ อนิจฺจํ สงฺขตํ ปฏิจฺจสมุปฺปนฺนํ ขยธมฺมํ วยธมฺมํ วิราคธมฺมํ นิโรธธมฺมํ, ตสฺส นิโรโธ “นิโรโธ”ติ วุจฺจติ.\csromanspacer{} อิเมสํ โข อานนฺท ธมฺมานํ นิโรโธ “นิโรโธ”ติ วุจฺจตีติ.\csromanspacer{}\csromanspacer{}ทสมํ.\csromanspacer{} อนิจฺจวคฺโค ทุติโย.}

\titlehead{\textbf{ตสฺสุทฺทานํ}}
\csromangathameasure{อนิจฺจํ ทุกฺขํ อนตฺตา, ยทนิจฺจาปเร ตโย. \\ เหตุนาปิ ตโย วุตฺตา, อานนฺเทน จ เต ทสาติ.}
\csromangathagroup{\csromangathastanza{อนิจฺจํ ทุกฺขํ อนตฺตา, ยทนิจฺจาปเร ตโย. \\ เหตุนาปิ ตโย วุตฺตา, อานนฺเทน จ เต ทสาติ.}}

\chapterhead{3. ภารวคฺค}
\csromanheader{2}{1. ภารสุตฺต}
\proseitem{22}{สาวตฺถิยํ.\csromanspacer{} ตตฺร โข.\csromanspacer{} ภารญฺจ โว ภิกฺขเว เทเสสฺสามิ ภารหารญฺจ ภาราทานญฺจ ภารนิกฺเขปนญฺจ, ตํ สุณาถ.\csromanspacer{} กตโม จ ภิกฺขเว ภาโร, “ปญฺจุปาทานกฺขนฺธา”ติสฺส วจนียํ.\csromanspacer{} กตเม ปญฺจ, รูปุปาทานกฺขนฺโธ เวทนุปาทานกฺขนฺโธ สญฺญุปาทานกฺขนฺโธ สงฺขารุปาทานกฺขนฺโธ วิญฺญาณุปาทานกฺขนฺโธ.\csromanspacer{} อยํ วุจฺจติ ภิกฺขเว ภาโร.\csromanspacer{} กตโม จ ภิกฺขเว \csromanfolio{22}ภารหาโร, “ปุคฺคโล”ติสฺส วจนียํ, ยฺวายํ อายสฺมา เอวํนาโม เอวํโคตฺโต.\csromanspacer{} อยํ วุจฺจติ ภิกฺขเว ภารหาโร.\csromanspacer{} กตมญฺจ ภิกฺขเว ภาราทานํ, ยายํ ตณฺหา โปโนภวิกา\footnote{โปโน พฺภวิกา (สฺยา, กํ, ก)} นนฺทีราคสหคตา\footnote{นนฺทิราคสหคตา (สพฺพตฺถ)} ตตฺรตตฺราภินนฺทินี.\csromanspacer{} เสยฺยถิทํ, กามตณฺหา ภวตณฺหา วิภวตณฺหา.\csromanspacer{} อิทํ วุจฺจติ ภิกฺขเว ภาราทานํ.\csromanspacer{} กตมญฺจ ภิกฺขเว ภารนิกฺเขปนํ, โย ตสฺสาเยว ตณฺหาย อเสสวิราคนิโรโธ จาโค ปฏินิสฺสคฺโค มุตฺติ อนาลโย.\csromanspacer{} อิทํ วุจฺจติ ภิกฺขเว ภารนิกฺเขปนนฺติ.}

\prose{อิทมโวจ ภควา, อิทํ วตฺวาน\footnote{วตฺวา (สี) เอวมีทิเสสุ ฐาเนสุ.} สุคโต อถาปรํ เอตทโวจ สตฺถา\csromandash{}}

\csromangathameasure{“ภารา หเว ปญฺจกฺขนฺธา, ภารหาโร จ ปุคฺคโล. \\ ภาราทานํ ทุขํ โลเก, ภารนิกฺเขปนํ สุขํ. \\ นิกฺขิปิตฺวา ครุํ ภารํ, อญฺญํ ภารํ อนาทิย. \\ สมูลํ ตณฺหมพฺพุยฺห, นิจฺฉาโต ปรินิพฺพุโต”ติ. ปฐมํ.}
\csromangathagroup{\csromangathastanza{“ภารา หเว ปญฺจกฺขนฺธา, ภารหาโร จ ปุคฺคโล. \\ ภาราทานํ ทุขํ โลเก, ภารนิกฺเขปนํ สุขํ.}\csromangathastanza{นิกฺขิปิตฺวา ครุํ ภารํ, อญฺญํ ภารํ อนาทิย. \\ สมูลํ ตณฺหมพฺพุยฺห\footnote{ตณฺหมพฺภุยฺห (อิ, ก)}, นิจฺฉาโต ปรินิพฺพุโต”ติ.\csromanspacer{} ปฐมํ.}}

\csromanheader{2}{2. ปริญฺญสุตฺต}
\proseitem{23}{สาวตฺถินิทานํ.\csromanspacer{} ปริญฺเญยฺเย จ ภิกฺขเว ภิกฺขเว ธมฺเม เทเสสฺสามิ ปริญฺญญฺจ, ตํ สุณาถ.\csromanspacer{} กตเม จ ภิกฺขเว ปริญฺเญยฺยา ธมฺมา, รูปํ ภิกฺขเว ปริญฺเญยฺโย ธมฺโม, เวทนา ปริญฺเญยฺโย ธมฺโม, สญฺญา ปริญฺเญยฺโย ธมฺโม, สงฺขารา ปริญฺเญยฺโย ธมฺโม, วิญฺญาณํ ปริญฺเญยฺโย ธมฺโม.\csromanspacer{} อิเม วุจฺจนฺติ ภิกฺขเว ปริญฺเญยฺยา ธมฺมา.\csromanspacer{} กตมา จ ภิกฺขเว ปริญฺญา, โย ภิกฺขเว ราคกฺขโย โทสกฺขโย โมหกฺขโย.\csromanspacer{} อยํ วุจฺจติ ภิกฺขเว ปริญฺญาติ.\csromanspacer{}\csromanspacer{}ทุติยํ.}

\csromanheader{2}{3. อภิชานสุตฺต}
\proseitem{24}{สาวตฺถินิทานํ.\csromanspacer{} รูปํ ภิกฺขเว อนภิชานํ อปริชานํ อวิราชยํ อปฺปชหํ อภพฺโพ ทุกฺขกฺขยาย.\csromanspacer{} เวทนํ อนภิชานํ อปริชานํ อวิราชยํ อปฺปชหํ อภพฺโพ ทุกฺขกฺขยาย.\csromanspacer{} สญฺญํ อนภิชานํ.\csromanspacer{} สงฺขาเร \csromanfolio{23}อนภิชานํ อปริชานํ อวิราชยํ อปฺปชหํ อภพฺโพ ทุกฺขกฺขยาย.\csromanspacer{} วิญฺญาณํ อนภิชานํ อปริชานํ อวิราชยํ อปฺปชหํ อภพฺโพ ทุกฺขกฺขยาย.\csromanspacer{} รูปญฺจ โข ภิกฺขเว อภิชานํ ปริชานํ วิราชยํ ปชหํ ภพฺโพ ทุกฺขกฺขยาย.\csromanspacer{} เวทนํ อภิชานํ.\csromanspacer{} สญฺญํ.\csromanspacer{} สงฺขาเร.\csromanspacer{} วิญฺญาณํ อภิชานํ ปริชานํ วิราชยํ ปชหํ ภพฺโพ ทุกฺขกฺขยายาติ.\csromanspacer{}\csromanspacer{}ตติยํ.}

\csromanheader{2}{4. ฉนฺทราคสุตฺต}
\proseitem{25}{สาวตฺถินิทานํ.\csromanspacer{} โย ภิกฺขเว รูปสฺมิํ ฉนฺทราโค, ตํ ปชหถ.\csromanspacer{} เอวํ ตํ รูปํ ปหีนํ ภวิสฺสติ อุจฺฉินฺนมูลํ ตาลาวตฺถุกตํ อนภาวํกตํ อายติํ อนุปฺปาทธมฺมํ.\csromanspacer{} โย เวทนาย ฉนฺทราโค, ตํ ปชหถ.\csromanspacer{} เอวํ สา เวทนา ปหีนา ภวิสฺสติ อุจฺฉินฺนมูลา ตาลาวตฺถุกตา อนภาวํกตา อายติํ อนุปฺปาทธมฺมา.\csromanspacer{} โย สญฺญาย ฉนฺทราโค, ตํ ปชหถ.\csromanspacer{} เอวํ สา สญฺญา ปหีนา ภวิสฺสติ อุจฺฉินฺนมูลา ตาลาวตฺถุกตา อนภาวํกตา อายติํ อนุปฺปาทธมฺมา.\csromanspacer{} โย สงฺขาเรสุ ฉนฺทราโค, ตํ ปชหถ.\csromanspacer{} เอวํ เต สงฺขารา ปหีนา ภวิสฺสนฺติ อุจฺฉินฺนมูลา ตาลาวตฺถุกตา อนภาวํกตา อายติํ อนุปฺปาทธมฺมา.\csromanspacer{} โย วิญฺญาณสฺมิํ ฉนฺทราโค, ตํ ปชหถ.\csromanspacer{} เอวํ ตํ วิญฺญาณํ ปหีนํ ภวิสฺสติ อุจฺฉินฺนมูลํ ตาลาวตฺถุกตํ อนภาวํกตํ อายติํ อนุปฺปาทธมฺมนฺติ.\csromanspacer{}\csromanspacer{}จตุตฺถํ.}

\csromanheader{2}{5. อสฺสาทสุตฺต}
\proseitem{26}{สาวตฺถินิทานํ.\csromanspacer{} ปุพฺเพว\footnote{ปุพฺเพ (อิ, ก)} เม ภิกฺขเว สมฺโพธา อนภิสมฺพุทฺธสฺส โพธิสตฺตสฺเสว\footnote{โพธิสตฺตสฺส (อิ, ก)} สโต เอตทโหสิ “โก นุ โข รูปสฺส อสฺสาโท โก อาทีนโว กิํ นิสฺสรณํ, โก เวทนาย อสฺสาโท โก อาทีนโว กิํ นิสฺสรณํ, โก สญฺญาย อสฺสาโท โก อาทีนโว กิํ นิสฺสรณํ, โก สงฺขารานํ อสฺสาโท โก อาทีนโว กิํ นิสฺสรณํ, โก วิญฺญาณสฺส อสฺสาโท โก อาทีนโว กิํ นิสฺสรณนฺ”ติ.\csromanspacer{} ตสฺส มยฺหํ ภิกฺขเว เอตทโหสิ “ยํ โข รูปํ ปฏิจฺจ อุปฺปชฺชติ สุขํ โสมนสฺสํ, อยํ รูปสฺส อสฺสาโท.\csromanspacer{} ยํ รูปํ อนิจฺจํ ทุกฺขํ วิปริณามธมฺมํ, อยํ รูปสฺส อาทีนโว.\csromanspacer{} โย รูปสฺมิํ ฉนฺทราควินโย ฉนฺทราคปฺปหานํ, อิทํ รูปสฺส นิสฺสรณํ.\csromanspacer{} ยํ เวทนํ ปฏิจฺจ อุปฺปชฺชติ สุขํ \csromanfolio{24}โสมนสฺสํ, อยํ เวทนาย อสฺสาโท.\csromanspacer{} ยํ\footnote{ยา (ก)} เวทนา อนิจฺจา ทุกฺขา วิปริณามธมฺมา, อยํ เวทนาย อาทีนโว.\csromanspacer{} โย เวทนาย ฉนฺทราควินโย ฉนฺทราคปฺปหานํ, อิทํ เวทนาย นิสฺสรณํ.\csromanspacer{} ยํ สญฺญํ ปฏิจฺจ อุปฺปชฺชติ ฯเปฯ.\csromanspacer{} ยํ สงฺขาเร ปฏิจฺจ อุปฺปชฺชติ สุขํ โสมนสฺสํ, อยํ สงฺขารานํ อสฺสาโท.\csromanspacer{} ยํ\footnote{เย (สี, ก)} สงฺขารา อนิจฺจา ทุกฺขา วิปริณามธมฺมา, อยํ สงฺขารานํอาทีนโว.\csromanspacer{} โย สงฺขาเรสุ ฉนฺทราควินโย ฉนฺทราคปฺปหานํ, อิทํ สงฺขารานํ นิสฺสรณํ.\csromanspacer{} ยํ วิญฺญาณํ ปฏิจฺจ อุปฺปชฺชติ สุขํ โสมนสฺสํ, อยํ วิญฺญาณสฺส อสฺสาโท.\csromanspacer{} ยํ วิญฺญาณํ อนิจฺจํ ทุกฺขํ วิปริณามธมฺมํ, อยํ วิญฺญาณสฺส อาทีนโว.\csromanspacer{} โย วิญฺญาณสฺมิํ ฉนฺทราควินโย ฉนฺทราคปฺปหานํ, อิทํ วิญฺญาณสฺส นิสฺสรณํ.}

\prose{ยาวกีวญฺจาหํ ภิกฺขเว อิเมสํ ปญฺจนฺนํ อุปาทานกฺขนฺธานํ เอวํ อสฺสาทญฺจ อสฺสาทโต อาทีนวญฺจ อาทีนวโต นิสฺสรณญฺจ นิสฺสรณโต ยถาภูตํ นาพฺภญฺญาสิํ, เนว ตาวาหํ ภิกฺขเว สเทวเก โลเก สมารเก สพฺรหฺมเก สสฺสมณพฺราหฺมณิยา ปชาย สเทวมนุสฺสาย “อนุตฺตรํ สมฺมาสมฺโพธิํ อภิสมฺพุทฺโธ”ติ\footnote{อภิสมฺพุทฺโธ (สี)} ปจฺจญฺญาสิํ.\csromanspacer{} ยโต จ ขฺวาหํ ภิกฺขเว อิเมสํ ปญฺจนฺนํ อุปาทานกฺขนฺธานํ เอวํ อสฺสาทญฺจ อสฺสาทโต อาทีนวญฺจ อาทีนวโต นิสฺสรณญฺจ นิสฺสรณโต ยถาภูตํ อพฺภญฺญาสิํ, อถาหํ ภิกฺขเว สเทวเก โลเก สมารเก สพฺรหฺมเก สสฺสมณพฺราหฺมณิยา ปชาย สเทวมนุสฺสาย “อนุตฺตรํ สมฺมาสมฺโพธิํ อภิสมฺพุทฺโธ”ติ ปจฺจญฺญาสิํ.\csromanspacer{} ญาณญฺจ ปน เม ทสฺสนํ อุทปาทิ, อกุปฺปา เม วิมุตฺติ\footnote{เจโตวิมุตฺติ (สี, อิ, ก)}, อยมนฺติมา ชาติ, นตฺถิ ทานิ ปุนพฺภโวติ.\csromanspacer{}\csromanspacer{}ปญฺจมํ.}

\csromanheader{2}{6. ทุติย-อสฺสาทสุตฺต}
\proseitem{27}{สาวตฺถินิทานํ.\csromanspacer{} รูปสฺสาหํ ภิกฺขเว อสฺสาทปริเยสนํ อจริํ, โย รูปสฺส อสฺสาโท ตทชฺฌคมํ, ยาวตา รูปสฺส อสฺสาโท ปญฺญาย เม โส สุทิฏฺโฐ.\csromanspacer{} รูปสฺสาหํ ภิกฺขเว อาทีนวปริเยสนํ อจริํ, โย รูปสฺส อาทีนโว ตทชฺฌคมํ, ยาวตา รูปสฺส อาทีนโว ปญฺญาย เม โส สุทิฏฺโฐ.\csromanspacer{} รูปสฺสาหํ ภิกฺขเว นิสฺสรณปริเยสนํ อจริํ, ยํ รูปสฺส นิสฺสรณํ ตทชฺฌคมํ, ยาวตา รูปสฺส นิสฺสรณํ ปญฺญาย เม ตํ สุทิฏฺฐํ.\csromanspacer{} เวทนายาหํ ภิกฺขเว.\csromanspacer{} สญฺญายาหํ ภิกฺขเว, \csromanfolio{25}สงฺขารานาหํ ภิกฺขเว.\csromanspacer{} วิญฺญาณสฺสาหํ ภิกฺขเว อสฺสาทปริเยสนํ อจริํ, โย วิญฺญาณสฺส อสฺสาโท ตทชฺฌคมํ, ยาวตา วิญฺญาณสฺส อสฺสาโท ปญฺญาย เม โส สุทิฏฺโฐ.\csromanspacer{} วิญฺญาณสฺสาหํ ภิกฺขเว อาทีนวปริเยสนํ อจริํ, โย วิญฺญาณสฺส อาทีนโว ตทชฺฌคมํ, ยาวตา วิญฺญาณสฺส อาทีนโว ปญฺญาย เม โส สุทิฏฺโฐ.\csromanspacer{} วิญฺญาณสฺสาหํ ภิกฺขเว นิสฺสรณปริเยสนํ อจริํ, ยํ วิญฺญาณสฺส นิสฺสรณํ ตทชฺฌคมํ, ยาวตา วิญฺญาณสฺส นิสฺสรณํ ปญฺญาย เม ตํ สุทิฏฺฐํ.\csromanspacer{} ยาวกีวญฺจาหํ ภิกฺขเว อิเมสํ ปญฺจนฺนํ อุปาทานกฺขนฺธานํ อสฺสาทญฺจ อสฺสาทโต อาทีนวญฺจ อาทีนวโต นิสฺสรณญฺจ นิสฺสรณโต ยถาภูตํ นาพฺภญฺญาสิํ ฯเปฯ อพฺภญฺญาสิํ.\csromanspacer{} ญาณญฺจ ปน เม ทสฺสนํ อุทปาทิ, อกุปฺปา เม วิมุตฺติ\footnote{เจโตวิมุตฺติ (สี, อิ, ก)}, อยมนฺติมา ชาติ, นตฺถิ ทานิ ปุนพฺภโวติ.\csromanspacer{}\csromanspacer{}ฉฏฺฐํ.}

\csromanheader{2}{7. ตติย-อสฺสาทสุตฺต}
\proseitem{28}{สาวตฺถินิทานํ.\csromanspacer{} โน เจทํ ภิกฺขเว รูปสฺส อสฺสาโท อภวิสฺส, นยิทํ สตฺตา รูปสฺมิํ สารชฺเชยฺยุํ.\csromanspacer{} ยสฺมา จ โข ภิกฺขเว อตฺถิ รูปสฺส อสฺสาโท, ตสฺมา สตฺตา รูปสฺมิํ สารชฺชนฺติ.\csromanspacer{} โน เจทํ ภิกฺขเว รูปสฺส อาทีนโว อภวิสฺส, นยิทํ สตฺตา รูปสฺมิํ นิพฺพินฺเทยฺยุํ.\csromanspacer{} ยสฺมา จ โข ภิกฺขเว อตฺถิ รูปสฺส อาทีนโว, ตสฺมา สตฺตา รูปสฺมิํ นิพฺพินฺทนฺติ.\csromanspacer{} โน เจทํ ภิกฺขเว รูปสฺส นิสฺสรณํ อภวิสฺส, นยิทํ สตฺตา รูปสฺมา นิสฺสเรยฺยุํ.\csromanspacer{} ยสฺมา จ โข ภิกฺขเว อตฺถิ รูปสฺส นิสฺสรณํ, กสฺมา สตฺตา รูปสฺมา นิสฺสรนฺติ.\csromanspacer{} โน เจทํ ภิกฺขเว เวทนาย ฯเปฯ.\csromanspacer{} โน เจทํ ภิกฺขเว สญฺญาย.\csromanspacer{} โน เจทํ ภิกฺขเว สงฺขารานํ นิสฺสรณํ อภวิสฺส, นยิทํ สตฺตา สงฺขาเรหิ นิสฺสเรยฺยุํ.\csromanspacer{} ยสฺมา จ โข ภิกฺขเว อตฺถิ สงฺขารานํ นิสฺสรณํ, ตสฺมา สตฺตา สงฺขาเรหิ นิสฺสรนฺติ.\csromanspacer{} โน เจทํ ภิกฺขเว วิญฺญาณสฺส อสฺสาโท อภวิสฺส, นยิทํ สตฺตา วิญฺญาณสฺมิํ สารชฺเชยฺยุํ.\csromanspacer{} ยสฺมา จ โข ภิกฺขเว อตฺถิ วิญฺญาณสฺส อสฺสาโท, ตสฺมา สตฺตา วิญฺญาณสฺมิํ สารชฺชนฺติ.\csromanspacer{} โน เจทํ ภิกฺขเว วิญฺญาณสฺส อาทีนโว อภวิสฺส, นยิทํ สตฺตา วิญฺญาณสฺมิํ นิพฺพินฺเทยฺยุํ.\csromanspacer{} ยสฺมา จ โข ภิกฺขเว อตฺถิ วิญฺญาณสฺส อาทีนโว, ตสฺมา สตฺตา วิญฺญาณสฺมิํ นิพฺพินฺทนฺติ.\csromanspacer{} โน เจทํ ภิกฺขเว วิญฺญาณสฺส นิสฺสรณํ อภวิสฺส, นยิทํ สตฺตา วิญฺญาณสฺมา \csromanfolio{26}นิสฺสเรยฺยุํ.\csromanspacer{} ยสฺมา จ โข ภิกฺขเว อตฺถิ วิญฺญาณสฺส นิสฺสรณํ, ตสฺมา สตฺตา วิญฺญาณสฺมา นิสฺสรนฺติ.}

\prose{ยาวกีวญฺจ ภิกฺขเว สตฺตา อิเมสํ ปญฺจนฺนํ อุปาทานกฺขนฺธานํ อสฺสาทญฺจ อสฺสาทโต อาทีนวญฺจ อาทีนวโต นิสฺสรณญฺจ นิสฺสรณโต ยถาภูตํ นาพฺภญฺญํสุ\footnote{นาพฺภญฺญิํสุ (สี)}.\csromanspacer{} เนว ตาว ภิกฺขเว สตฺตา สเทวกา โลกา สมารกา สพฺรหฺมกา สสฺสมณพฺราหฺมณิยา ปชาย สเทวมนุสฺสาย นิสฺสฏา วิสํยุตฺตา วิปฺปมุตฺตา วิมริยาทีกเตน เจตสา วิหริํสุ.\csromanspacer{} ยโต จ โข ภิกฺขเว สตฺตา อิเมสํ ปญฺจนฺนํ อุปาทานกฺขนฺธานํ อสฺสาทญฺจ อสฺสาทโต อาทีนวญฺจ อาทีนวโต นิสฺสรณญฺจ นิสฺสรณโต ยถาภูตํ อพฺภญฺญํสุ.\csromanspacer{} อถ ภิกฺขเว สตฺตา สเทวกา โลกา สมารกา สพฺรหฺมกา สสฺสมณพฺราหฺมณิยา ปชาย สเทวมนุสฺสาย นิสฺสฏา วิสํยุตฺตา วิปฺปมุตฺตา วิมริยาทีกเตน เจตสา วิหรนฺติ.\csromanspacer{}\csromanspacer{}สตฺตมํ.}

\csromanheader{2}{8. อภินนฺทนสุตฺต}
\proseitem{29}{สาวตฺถินิทานํ.\csromanspacer{} โย ภิกฺขเว รูปํ อภินนฺทติ, ทุกฺขํ โส อภินนฺทติ.\csromanspacer{} โย ทุกฺขํ อภินนฺทติ “อปริมุตฺโต โส ทุกฺขสฺมา”ติ วทามิ.\csromanspacer{} โย เวทนํ อภินนฺทติ.\csromanspacer{} โย สญฺญํ อภินนฺทติ.\csromanspacer{} โย สงฺขาเร อภินนฺทติ.\csromanspacer{} โย วิญฺญาณํ อภินนฺทติ, ทุกฺขํ โส อภินนฺทติ.\csromanspacer{} โย ทุกฺขํ อภินนฺทติ, “อปริมุตฺโต โส ทุกฺขสฺมา”ติ วทามิ.\csromanspacer{} โย จ โข ภิกฺขเว รูปํ นาภินนฺทติ, ทุกฺขํ โส นาภินนฺทติ.\csromanspacer{} โย ทุกฺขํ นาภินนฺทติ, “ปริมุตฺโต โส ทุกฺขสฺมา”ติ วทามิ.\csromanspacer{} โย เวทนํ นาภินนฺทติ.\csromanspacer{} โย สญฺญํ นาภินนฺทติ.\csromanspacer{} โย สงฺขาเร นาภินนฺทติ.\csromanspacer{} โย วิญฺญาณํ นาภินนฺทติ, ทุกฺขํ โส นาภินนฺทติ.\csromanspacer{} โย ทุกฺขํ นาภินนฺทติ, “ปริมุตฺโต โส ทุกฺขสฺมา”ติ วทามีติ.\csromanspacer{}\csromanspacer{}อฏฺฐมํ.}

\csromanheader{2}{9. อุปฺปาทสุตฺต}
\proseitem{30}{สาวตฺถินิทานํ.\csromanspacer{} โย ภิกฺขเว รูปสฺส อุปฺปาโท ฐิติ อภินิพฺพตฺติ ปาตุภาโว, ทุกฺขสฺเสโส อุปฺปาโท, โรคานํ ฐิติ, ชรามรณสฺส ปาตุภาโว.\csromanspacer{} โย เวทนาย.\csromanspacer{} โย สญฺญาย ฯเปฯ.\csromanspacer{} โย สงฺขารานํ.\csromanspacer{} โย \csromanfolio{27}วิญฺญาณสฺส อุปฺปาโท ฐิติ อภินิพฺพตฺติ ปาตุภาโว, ทุกฺขสฺเสโส อุปฺปาโท, โรคานํ ฐิติ, ชารามรณสฺส ปาตุภาโว.\csromanspacer{} โย จ โข ภิกฺขเว รูปสฺส นิโรโธ วูปสโม อตฺถงฺคโม, ทุกฺขสฺเสโส นิโรโธ, โรคานํ วูปสโม, ชรามรณสฺส อตฺถงฺคโม.\csromanspacer{} โย เวทนาย ฯเปฯ.\csromanspacer{} โย สญฺญาย.\csromanspacer{} โย สงฺขารานํ.\csromanspacer{} โย วิญฺญาณสฺส นิโรโธ วูปสโม อตฺถงฺคโม, ทุกฺขสฺเสโส นิโรโธ, โรคานํ วูปสโม, ชรามรณสฺส อตฺถงฺคโมติ.\csromanspacer{}\csromanspacer{}นวมํ.}

\csromanheader{2}{10. อฆมูลสุตฺต}
\proseitem{31}{สาวตฺถินิทานํ.\csromanspacer{} อฆญฺจ ภิกฺขเว เทเสสฺสามิ อฆมูลญฺจ, ตํ สุณาถ.\csromanspacer{} กตมญฺจ ภิกฺขเว อฆํ, รูปํ ภิกฺขเว อฆํ, เวทนา อฆํ, สญฺญา อฆํ, สงฺขารา อฆํ, วิญฺญาณํ อฆํ.\csromanspacer{} อิทํ วุจฺจติ ภิกฺขเว อฆํ.\csromanspacer{} กตมญฺจ ภิกฺขเว อฆมูลํ, ยายํ ตณฺหา โปโนภวิกา นนฺทีราคสหคตา\footnote{นนฺทิราคสหคตา (สพฺพตฺถ)} ตตฺรตตฺราภินนฺทินี.\csromanspacer{} เสยฺยถิทํ, กามตณฺหา ภวตณฺหา วิภวตณฺหา.\csromanspacer{} อิทํ วุจฺจติ ภิกฺขเว อฆมูลนฺติ.\csromanspacer{}\csromanspacer{}ทสมํ.}

\csromanheader{2}{11. ปภงฺคุสุตฺต}
\proseitem{32}{สาวตฺถินิทานํ.\csromanspacer{} ปภงฺคุญฺจ ภิกฺขเว เทเสสฺสามิ อปฺปภงฺคุญฺจ, ตํ สุณาถ.\csromanspacer{} กิญฺจ ภิกฺขเว ปภงฺคุ, กิํ อปฺปภงฺคุ.\csromanspacer{} รูปํ ภิกฺขเว ปภงฺคุ, โย ตสฺส นิโรโธ วูปสโม อตฺถงฺคโม, อิทํ อปฺปภงฺคุ.\csromanspacer{} เวทนา ปภงฺคุ, โย ตสฺสา นิโรโธ วูปสโม อตฺถงฺคโม, อิทํ อปฺปภงฺคุ.\csromanspacer{} สญฺญา ปภงฺคุ.\csromanspacer{} สงฺขารา ปภงฺคุ, โย เตสํ นิโรโธ วูปสโม อตฺถงฺคโม, อิทํ อปฺปภงฺคุ.\csromanspacer{} วิญฺญาณํ ปภงฺคุ, โย ตสฺส นิโรโธ วูปสโม อตฺถงฺคโม, อิทํ อปฺปภงฺคูติ.\csromanspacer{}\csromanspacer{}เอกาทสมํ.\csromanspacer{} ภารวคฺโค ตติโย.}

\titlehead{\textbf{ตสฺสุทฺทานํ}}
\csromangathameasure{ภารํ ปริญฺญํ อภิชานํ, ฉนฺทราคํ จตุตฺถกํ. \\ อสฺสาทา จ ตโย วุตฺตา, อภินนฺทนมฏฺฐมํ. \\ อุปฺปาทํ อฆมูลญฺจ, เอกาทสโม ปภงฺคูติ.}
\csromangathagroup{\csromangathastanza{ภารํ ปริญฺญํ อภิชานํ, ฉนฺทราคํ จตุตฺถกํ. \\ อสฺสาทา จ ตโย วุตฺตา, อภินนฺทนมฏฺฐมํ. \\ อุปฺปาทํ อฆมูลญฺจ, เอกาทสโม ปภงฺคูติ.}}

\chapterheadpageread{4. นตุมฺหากวคฺค}
\csromanheader{2}{1. นตุมฺหากสุตฺต}
\proseitem{33}{\csromanfolio{28}สาวตฺถินิทานํ.\csromanspacer{} ยํ ภิกฺขเว น ตุมฺหากํ ตํ ปชหถ, ตํ โว ปหีนํ หิตาย สุขาย ภวิสฺสติ.\csromanspacer{} กิญฺจ ภิกฺขเว น ตุมฺหากํ, รูปํ ภิกฺขเว น ตุมฺหากํ ตํ ปชหถ, ตํ โว ปหีนํ หิตาย สุขาย ภวิสฺสติ.\csromanspacer{} เวทนา น ตุมฺหากํ ตํ ปชหถ, สา โว ปหีนา หิตาย สุขาย ภวิสฺสติ.\csromanspacer{} สญฺญา น ตุมฺหากํ.\csromanspacer{} สงฺขารา น ตุมฺหากํ เต ปชหถ, เต โว ปหีนา หิตาย สุขาย ภวิสฺสนฺติ.\csromanspacer{} วิญฺญาณํ น ตุมฺหากํ ตํ ปชหถ, ตํ โว ปหีนํ หิตาย สุขาย ภวิสฺสติ.}

\prose{เสยฺยถาปิ ภิกฺขเว ยํ อิมสฺมิํ เชตวเน ติณกฏฺฐสาขาปลาสํ, ตํ ชโน หเรยฺย วา ฑเหยฺย วา ยถาปจฺจยํ วา กเรยฺย, อปิ นุ ตุมฺหากํ เอวมสฺส “อเมฺห ชโน หรติ วา ฑหติ วา ยถาปจฺจยํ วา กโรตี”ติ.\csromanspacer{} โน เหตํ ภนฺเต.\csromanspacer{} ตํ กิสฺส เหตุ, น หิ โน เอตํ ภนฺเต อตฺตา วา อตฺตนิยํ วาติ.\csromanspacer{} เอวเมว โข ภิกฺขเว รูปํ น ตุมฺหากํ ตํ ปชหถ, ตํ โว ปหีนํ หิตาย สุขาย ภวิสฺสติ.\csromanspacer{} เวทนา น ตุมฺหากํ ตํ ปชหถ, สา โว ปหีนา หิตาย สุขาย ภวิสฺสติ.\csromanspacer{} สญฺญา น ตุมฺหากํ.\csromanspacer{} สงฺขารา น ตุมฺหากํ.\csromanspacer{} วิญฺญาณํ น ตุมฺหากํ ตํ ปชหถ, ตํ โว ปหีนํ หิตาย สุขาย ภวิสฺสตีติ.\csromanspacer{}\csromanspacer{}ปฐมํ.}

\csromanheader{2}{2. ทุติยนตุมฺหากสุตฺต}
\proseitem{34}{สาวตฺถินิทานํ.\csromanspacer{} ยํ ภิกฺขเว น ตุมฺหากํ ตํ ปชหถ, ตํ โว ปหีนํ หิตาย สุขาย ภวิสฺสติ.\csromanspacer{} กิญฺจ ภิกฺขเว น ตุมฺหากํ, รูปํ ภิกฺขเว น ตุมฺหากํ ตํ ปชหถ, ตํ โว ปหีนํ หิตาย สุขาย ภวิสฺสติ.\csromanspacer{} เวทนา น ตุมฺหากํ.\csromanspacer{} สญฺญา น ตุมฺหากํ.\csromanspacer{} สงฺขารา น ตุมฺหากํ.\csromanspacer{} วิญฺญาณํ น ตุมฺหากํ ตํ ปชหถ, ตํ โว ปหีนํ หิตาย สุขาย ภวิสฺสติ.\csromanspacer{} ยํ ภิกฺขเว น ตุมฺหากํ ตํ ปชหถ, ตํ โว ปหีนํ หิตาย สุขาย ภวิสฺสตีติ.\csromanspacer{}\csromanspacer{}ทุติยํ.}

\csromanheader{2}{3. อญฺญตรภิกฺขุสุตฺต}
\proseitem{35}{สาวตฺถินิทานํ.\csromanspacer{} อถ โข อญฺญตโร ภิกฺขุ เยน ภควา เตนุปสงฺกมิ, อุปสงฺกมิตฺวา ภควนฺตํ อภิวาเทตฺวา เอกมนฺตํ นิสีทิ, \csromanfolio{29}เอกมนฺตํ นิสินฺโน โข โส ภิกฺขุ ภควนฺตํ เอตทโวจ “สาธุ เม ภนฺเต ภควา สํขิตฺเตน ธมฺมํ เทเสตุ, ยมหํ ภควโต ธมฺมํ สุตฺวา เอโก วูปกฏฺโฐ อปฺปมตฺโต อาตาปี ปหิตตฺโต วิหเรยฺยนฺ”ติ.\csromanspacer{} ยํ โข ภิกฺขุ อนุเสติ เตน สงฺขํ คจฺฉติ, ยํ นานุเสติ น เตน สงฺขํ คจฺฉตีติ.\csromanspacer{} อญฺญาตํ ภควา, อญฺญาตํ สุคตาติ.}

\prose{ยถา กถํ ปน ตฺวํ ภิกฺขุ มยา สํขิตฺเตน ภาสิตสฺส วิตฺถาเรน อตฺถํ อาชานาสีติ.\csromanspacer{} รูปํ เจ ภนฺเต อนุเสติ เตน สงฺขํ คจฺฉติ, เวทนํ เจ อนุเสติ เตน สงฺขํ คจฺฉติ, สญฺญํ เจ อนุเสติ เตน สงฺขํ คจฺฉติ, สงฺขาเร เจ อนุเสติ เตน สงฺขํ คจฺฉติ, วิญฺญาณํ เจ อนุเสติ เตน สงฺขํ คจฺฉติ, รูปํ เจ ภนฺเต นานุเสติ น เตน สงฺขํ คจฺฉติ.\csromanspacer{} เวทนํ เจ.\csromanspacer{} สญฺญํ เจ.\csromanspacer{} สงฺขาเร เจ.\csromanspacer{} วิญฺญาณํ เจ นานุเสติ น เตน สงฺขํ คจฺฉติ.\csromanspacer{} อิมสฺส ขฺวาหํ ภนฺเต ภควตา สํขิตฺเตน ภาสิตสฺส เอวํ วิตฺถาเรน อตฺถํ อาชานามีติ.}

\prose{สาธุ สาธุ ภิกฺขุ, สาธุ โข ตฺวํ ภิกฺขุ มยา สํขิตฺเตน ภาสิตสฺส วิตฺถาเรน อตฺถํ อาชานาสิ.\csromanspacer{} รูปํ เจ ภิกฺขุ อนุเสติ เตน สงฺขํ คจฺฉติ.\csromanspacer{} เวทนํ เจ.\csromanspacer{} สญฺญํ เจ.\csromanspacer{} สงฺขาเร เจ.\csromanspacer{} วิญฺญาณํ เจ อนุเสติ เตน สงฺขํ คจฺฉติ.\csromanspacer{} รูปํ เจ ภิกฺขุ นานุเสติ.\csromanspacer{} น เตน สงฺขํ คจฺฉติ, เวทนํ เจ.\csromanspacer{} สญฺญํ เจ.\csromanspacer{} สงฺขาเร เจ.\csromanspacer{} วิญฺญาณํ เจ นานุเสติ น เตน สงฺขํ คจฺฉติ.\csromanspacer{} อิมสฺส โข ภิกฺขุ มยา สํขิตฺเตน ภาสิตสฺส เอวํ วิตฺถาเรน อตฺโถ ทฏฺฐพฺโพติ.}

\prose{อถ โข โส ภิกฺขุ ภควโต ภาสิตํ อภินนฺทิตฺวา อนุโมทิตฺวา อุฏฺฐายาสนา ภควนฺตํ อภิวาเทตฺวา ปทกฺขิณํ กตฺวา ปกฺกามิ.}

\prose{อถ โข โส ภิกฺขุ เอโก วูปกฏฺโฐ อปฺปมตฺโต อาตาปี ปหิตตฺโต วิหรนฺโต นจิรสฺเสว ยสฺสตฺถาย กุลปุตฺตา สมฺมเทว อคารสฺมา อนคาริยํ ปพฺพชนฺติ, ตทนุตฺตรํ พฺรหฺมจริยปริโยสานํ ทิฏฺเฐว ธมฺเม สยํ อภิญฺญา สจฺฉิกตฺวา อุปสมฺปชฺช วิหาสิ, “ขีณา ชาติ, วุสิตํ พฺราหฺมจริยํ, กตํ กรณียํ, นาปรํ อิตฺถตฺตายา”ติ อพฺภญฺญาสิ.\csromanspacer{} อญฺญตโร จ ปน โส ภิกฺขุ อรหตํ อโหสีติ.\csromanspacer{}\csromanspacer{}ตติยํ.}

\csromanheader{2}{4. ทุติย-อญฺญตรภิกฺขุสุตฺต}
\proseitem{36}{สาวตฺถินิทานํ.\csromanspacer{} อถ โข อญฺญตโร ภิกฺขุ เยน ภควา ฯเปฯ เอกมนฺตํ นิสินฺโน โข โส ภิกฺขุ ภควนฺตํ เอตทโวจ “สาธุ เม ภนฺเต \csromanfolio{30}ภควา สํขิตฺเตน ธมฺมํ เทเสตุ, ยมหํ ภควโต ธมฺมํ สุตฺวา เอโก วูปกฏฺโฐ อปฺปมตฺโต อาตาปี ปหิตตฺโต วิหเรยฺยนฺ”ติ.\csromanspacer{} ยํ โข ภิกฺขุ อนุเสติ ตํ อนุมียติ.\csromanspacer{} ยํ อนุมียติ เตน สงฺขํ คจฺฉติ.\csromanspacer{} ยํ นานุเสติ น ตํ อนุมียติ.\csromanspacer{} ยํ นานุมียติ น เตน สงฺขํ คจฺฉตีติ.\csromanspacer{} อญฺญาตํ ภควา, อญฺญาตํ สุคตาติ.}

\prose{ยถา กถํ ปน ตฺวํ ภิกฺขุ มยา สํขิตฺเตน ภาสิตสฺส วิตฺถาเรน อตฺถํ อาชานาสีติ.\csromanspacer{} รูปํ เจ ภนฺเต อนุเสติ ตํ อนุมียติ.\csromanspacer{} ยํ อนุมียติ เตน สงฺขํ คจฺฉติ.\csromanspacer{} เวทนํ เจ อนุเสติ.\csromanspacer{} สญฺญํ เจ อนุเสติ.\csromanspacer{} สงฺขาเร เจ อนุเสติ.\csromanspacer{} วิญฺญาณํ เจ อนุเสติ ตํ อนุมียติ.\csromanspacer{} ยํ อนุมียติ เตน สงฺขํ คจฺฉติ.\csromanspacer{} รูปํ เจ ภนฺเต นานุเสติ น ตํ อนุมียติ.\csromanspacer{} ยํ นานุมียติ น เตน สงฺขํ คจฺฉติ.\csromanspacer{} เวทนํ เจ นานุเสติ.\csromanspacer{} สญฺญํ เจ นานุเสติ.\csromanspacer{} สงฺขาเร เจ นานุเสติ.\csromanspacer{} วิญฺญาณํ เจ นานุเสติ น ตํ อนุมียติ.\csromanspacer{} ยํ นานุมียติ น เตน สงฺขํ คจฺฉติ.\csromanspacer{} อิมสฺส ขฺวาหํ ภนฺเต ภควตา สํขิตฺเตน ภาสิตสฺส เอวํ วิตฺถาเรน อตฺถํ อาชานามีติ.}

\prose{สาธุ สาธุ ภิกฺขุ, สาธุ โข ตฺวํ ภิกฺขุ มยา สํขิตฺเตน ภาสิตสฺส วิตฺถาเรน อตฺถํ อาชานาสิ.\csromanspacer{} รูปํ เจ ภิกฺขุ อนุเสติ ตํ อนุมียติ.\csromanspacer{} ยํ อนุมียติ เตน สงฺขํ คจฺฉติ.\csromanspacer{} เวทนํ เจ ภิกฺขุ.\csromanspacer{} สญฺญํ เจ ภิกฺขุ.\csromanspacer{} สงฺขาเร เจ ภิกฺขุ.\csromanspacer{} วิญฺญาณํ เจ ภิกฺขุ อนุเสติ ตํ อนุมียติ.\csromanspacer{} ยํ อนุมียติ เตน สงฺขํ คจฺฉติ.\csromanspacer{} รูปํ เจ ภิกฺขุ นานุเสติ น ตํ อนุมียติ.\csromanspacer{} ยํ นานุมียติ น เตน สงฺขํ คจฺฉติ.\csromanspacer{} เวทนํ เจ นานุเสติ.\csromanspacer{} สญฺญํ เจ นานุเสติ.\csromanspacer{} สงฺขาเร เจ นานุเสติ.\csromanspacer{} วิญฺญาณํ เจ นานุเสติ น ตํ อนุมียติ.\csromanspacer{} ยํ นานุมียติ น เตน สงฺขํ คจฺฉติ.\csromanspacer{} อิมสฺส โข ภิกฺขุ มยา สํขิตฺเตน ภาสิตสฺส เอวํ วิตฺถาเรน อตฺโถ ทฏฺฐพฺโพติ ฯเปฯ.\csromanspacer{} อญฺญตโร จ ปน โส ภิกฺขุ อรหตํ อโหสีติ.\csromanspacer{}\csromanspacer{}จตุตฺถํ.}

\csromanheader{2}{5. อานนฺทสุตฺต}
\proseitem{37}{สาวตฺถินิทานํ.\csromanspacer{} อถ โข อายสฺมา อานนฺโท ฯเปฯ เอกมนฺตํ นิสินฺนํ โข อายสฺมนฺตํ อานนฺทํ ภควา เอตทโวจ\csromandash{}}

\prose{\csromanfolio{31}สเจ ตํ อานนฺท เอวํ ปุจฺเฉยฺยุํ “กตเมสํ อาวุโส อานนฺท ธมฺมานํ อุปฺปาโท ปญฺญายติ, วโย ปญฺญายติ, ฐิตสฺส\footnote{ฐิตานํ (สฺยา, กํ, อิ, ก)} อญฺญถตฺตํ ปญฺญายตี”ติ, เอวํ ปุฏฺโฐ ตฺวํ อานนฺท กินฺติ พฺยากเรยฺยาสีติ.\csromanspacer{} สเจ มํ ภนฺเต เอวํ ปุจฺเฉยฺยุํ “กตเมสํ อาวุโส อานนฺท ธมฺมานํ อุปฺปาโท ปญฺญายติ, วโย ปญฺญายติ, ฐิตสฺส อญฺญถตฺตํ ปญฺญายตี”ติ, เอวํ ปุฏฺโฐหํ ภนฺเต เอวํ พฺยากเรยฺยํ “รูปสฺส โข อาวุโส อุปฺปาโท ปญฺญายติ, วโย ปญฺญายติ, ฐิตสฺส อญฺญถตฺตํ ปญฺญายติ.\csromanspacer{} เวทนาย.\csromanspacer{} สญฺญาย.\csromanspacer{} สงฺขารานํ.\csromanspacer{} วิญฺญาณสฺส อุปฺปาโท ปญฺญายติ, วโย ปญฺญายติ, ฐิตสฺส อญฺญถตฺตํ ปญฺญายติ.\csromanspacer{} อิเมสํ โข อาวุโส ธมฺมานํ อุปฺปาโท ปญฺญายติ, วโย ปญฺญายติ, ฐิตสฺส อญฺญถตฺตํ ปญฺญายตี”ติ.\csromanspacer{} เอวํ ปุฏฺโฐหํ ภนฺเต เอวํ พฺยากเรยฺยนฺติ.}

\prose{สาธุ สาธุ อานนฺท, รูปสฺส โข อานนฺท อุปฺปาโท ปญฺญายติ, วโย ปญฺญายติ, ฐิตสฺส อญฺญถตฺตํ ปญฺญายติ.\csromanspacer{} เวทนาย.\csromanspacer{} สญฺญาย.\csromanspacer{} สงฺขารานํ.\csromanspacer{} วิญฺญาณสฺส อุปฺปาโท ปญฺญายติ, วโย ปญฺญายติ, ฐิตสฺส อญฺญถตฺตํ ปญฺญายติ.\csromanspacer{} อิเมสํ โข อานนฺท ธมฺมานํ อุปฺปาโท ปญฺญายติ, วโย ปญฺญายติ, ฐิตสฺส อญฺญถตฺตํ ปญฺญายตีติ.\csromanspacer{} เอวํ ปุฏฺโฐ ตฺวํ อานนฺท เอวํ พฺยากเรยฺยาสีติ.\csromanspacer{}\csromanspacer{}ปญฺจมํ.}

\csromanheader{2}{6. ทุติย-อานนฺทสุตฺต}
\proseitem{38}{สาวตฺถินิทานํ.\csromanspacer{} เอกมนฺตํ นิสินฺนํ โข อายสฺมนฺตํ อานนฺทํ ภควา เอตทโวจ\csromandash{}}

\prose{สเจ ตํ อานนฺท เอวํ ปุจฺเฉยฺยุํ “กตเมสํ อาวุโส อานนฺท ธมฺมานํ อุปฺปาโท ปญฺญายิตฺถ, วโย ปญฺญายิตฺถ, ฐิตสฺส อญฺญถตฺตํ ปญฺญายิตฺถ.\csromanspacer{} กตเมสํ ธมฺมานํ อุปฺปาโท ปญฺญายิสฺสติ, วโย ปญฺญายิสฺสติ, ฐิตสฺส อญฺญถตฺตํ ปญฺญายิสฺสติ.\csromanspacer{} กตเมสํ ธมฺมานํ อุปฺปาโท ปญฺญายติ, วโย ปญฺญายติ, ฐิตสฺส อญฺญถตฺตํ ปญฺญายตี”ติ, เอวํ ปุฏฺโฐ ตฺวํ อานนฺท กินฺติ พฺยากเรยฺยาสีติ.\csromanspacer{} สเจ มํ ภนฺเต เอวํ ปุจฺเฉยฺยุํ “กตเมสํ อาวุโส อานนฺท ธมฺมานํ อุปฺปาโท ปญฺญายิตฺถ, วโย ปญฺญายิตฺถ, ฐิตสฺส อญฺญถตฺตํ ปญฺญายิตฺถ.\csromanspacer{} กตเมสํ ธมฺมานํ อุปฺปาโท \csromanfolio{32}ปญฺญายิสฺสติ, วโย ปญฺญายิสฺสติ, ฐิตสฺส อญฺญถตฺตํ ปญฺญายิสฺสติ.\csromanspacer{} กตเมสํ ธมฺมานํ อุปฺปาโท ปญฺญายติ, วโย ปญฺญายติ, ฐิตสฺส อญฺญถตฺตํ ปญฺญายตี”ติ, เอวํ ปุฏฺโฐหํ ภนฺเต เอวํ พฺยากเรยฺยํ “ยํ โข อาวุโส รูปํ อตีตํ นิรุทฺธํ วิปริณตํ, ตสฺส อุปฺปาโท ปญฺญายิตฺถ, วโย ปญฺญายิตฺถ, ฐิตสฺส อญฺญถตฺตํ ปญฺญายิตฺถ.\csromanspacer{} ยา เวทนา อตีตา นิรุทฺธา วิปริณตา, ตสฺสา อุปฺปาโท ปญฺญายิตฺถ, วโย ปญฺญายิตฺถ, ฐิตาย อญฺญถตฺตํ ปญฺญายิตฺถ.\csromanspacer{} ยา สญฺญา.\csromanspacer{} เย สงฺขารา อตีตา นิรุทฺธา วิปริณตา, เตสํ อุปฺปาโท ปญฺญายิตฺถ, วโย ปญฺญายิตฺถ, ฐิตสฺส อญฺญถตฺตํ ปญฺญายิตฺถ.\csromanspacer{} ยํ วิญฺญาณํ อตีตํ นิรุทฺธํ วิปริณตํ, ตสฺส อุปฺปาโท ปญฺญายิตฺถ, วโย ปญฺญายิตฺถ, ฐิตสฺส อญฺญถตฺตํ ปญฺญายิตฺถ.\csromanspacer{} อิเมสํ โข อาวุโส ธมฺมานํ อุปฺปาโท ปญฺญายิตฺถ, วโย ปญฺญายิตฺถ, ฐิตสฺส อญฺญถตฺตํ ปญฺญายิตฺถ.}

\prose{ยํ โข อาวุโส รูปํ อชาตํ อปาตุภูตํ, ตสฺส อุปฺปาโท ปญฺญายิสฺสติ, วโย ปญฺญายิสฺสติ, ฐิตสฺส อญฺญถตฺตํ ปญฺญายิสฺสติ.\csromanspacer{} ยา เวทนา อชาตา อปาตุภูตา, ตสฺสา อุปฺปาโท ปญฺญายิสฺสติ, วโย ปญฺญายิสฺสติ, ฐิตาย อญฺญถตฺตํ ปญฺญายิสฺสติ, ยา สญฺญา ฯเปฯ.\csromanspacer{} เย สงฺขารา อชาตา อปาตุภูตา, เตสํ อุปฺปาโท ปญฺญายิสฺสติ, วโย ปญฺญายิสฺสติ, ฐิตสฺส อญฺญถตฺตํ ปญฺญายิสฺสติ.\csromanspacer{} ยํ วิญฺญาณํ อชาตํ อปาตุภูตํ, ตสฺส อุปฺปาโท ปญฺญายิสฺสติ, วโย ปญฺญายิสฺสติ, ฐิตสฺส อญฺญถตฺตํ ปญฺญายิสฺสติ.\csromanspacer{} อิเมสํ โข อาวุโส ธมฺมานํ อุปฺปาโท ปญฺญายิสฺสติ, วโย ปญฺญายิสฺสติ, ฐิตสฺส อญฺญถตฺตํ ปญฺญายิสฺสติ.}

\prose{ยํ โข อาวุโส รูปํ ชาตํ ปาตุภูตํ, ตสฺส อุปฺปาโท ปญฺญายติ, วโย ปญฺญายติ, ฐิตสฺส อญฺญถตฺตํ ปญฺญายติ.\csromanspacer{} ยา เวทนา ชาตา ปาตุภูตา ฯเปฯ.\csromanspacer{} ยา สญฺญา.\csromanspacer{} เย สงฺขารา ชาตา ปาตุภูตา, เตสํ อุปฺปาโท ปญฺญายติ, วโย ปญฺญายติ, ฐิตสฺส อญฺญถตฺตํ ปญฺญายติ.\csromanspacer{} ยํ วิญฺญาณํ ชาตํ ปาตุภูตํ, ตสฺส อุปฺปาโท ปญฺญายติ, วโย ปญฺญายติ, ฐิตสฺส อญฺญถตฺตํ ปญฺญายติ.\csromanspacer{} อิเมสํ โข อาวุโส ธมฺมานํ อุปฺปาโท ปญฺญายติ, วโย ปญฺญายติ, ฐิตสฺส อญฺญถตฺตํ ปญฺญายตีติ.\csromanspacer{} เอวํ ปุฏฺโฐหํ ภนฺเต เอวํ พฺยากเรยฺยนฺติ.}

\prose{\csromanfolio{33}สาธุ สาธุ อานนฺท, ยํ โข อานนฺท รูปํ อตีตํ นิรุทฺธํ วิปริณตํ, ตสฺส อุปฺปาโท ปญฺญายิตฺถ, วโย ปญฺญายิตฺถ, ฐิตสฺส อญฺญถตฺตํ ปญฺญายิตฺถ.\csromanspacer{} ยา เวทนา.\csromanspacer{} ยา สญฺญา.\csromanspacer{} เย สงฺขารา.\csromanspacer{} ยํ วิญฺญาณํ อตีตํ นิรุทฺธํ วิปริณตํ, ตสฺส อุปฺปาโท ปญฺญายิตฺถ, วโย ปญฺญายิตฺถ, ฐิตสฺส อญฺญถตฺตํ ปญฺญายิตฺถ.\csromanspacer{} อิเมสํ โข อานนฺท ธมฺมานํ อุปฺปาโท ปญฺญายิตฺถ, วโย ปญฺญายิตฺถ, ฐิตสฺส อญฺญถตฺตํ ปญฺญายิตฺถ.}

\prose{ยํ โข อานนฺท รูปํ อชาตํ อปาตุภูตํ, ตสฺส อุปฺปาโท ปญฺญายิสฺสติ, วโย ปญฺญายิสฺสติ, ฐิตสฺส อญฺญถตฺตํ ปญฺญายิสฺสติ.\csromanspacer{} ยา เวทนา.\csromanspacer{} ยา สญฺญา.\csromanspacer{} เย สงฺขารา.\csromanspacer{} ยํ วิญฺญาณํ อชาตํ อปาตุภูตํ, ตสฺส อุปฺปาโท ปญฺญายิสฺสติ, วโย ปญฺญายิสฺสติ, ฐิตสฺส อญฺญถตฺตํ ปญฺญายิสฺสติ.\csromanspacer{} อิเมสํ โข อานนฺท ธมฺมานํ อุปฺปาโท ปญฺญายิสฺสติ, วโย ปญฺญายิสฺสติ, ฐิตสฺส อญฺญถตฺตํ ปญฺญายิสฺสติ.}

\prose{ยํ โข อานนฺท รูปํ ชาตํ ปาตุภูตํ, ตสฺส อุปฺปาโท ปญฺญายติ, วโย ปญฺญายติ, ฐิตสฺส อญฺญถตฺตํ ปญฺญายติ.\csromanspacer{} ยา เวทนา ชาตา ปาตุภูตา.\csromanspacer{} ยา สญฺญา.\csromanspacer{} เย สงฺขารา.\csromanspacer{} ยํ วิญฺญาณํ ชาตํ ปาตุภูตํ, ตสฺส อุปฺปาโท ปญฺญายติ, วโย ปญฺญายติ, ฐิตสฺส อญฺญถตฺตํ ปญฺญายติ.\csromanspacer{} อิเมสํ โข อานนฺท ธมฺมานํ อุปฺปาโท ปญฺญายติ, วโย ปญฺญายติ, ฐิตสฺส อญฺญถตฺตํ ปญฺญายตีติ.\csromanspacer{} เอวํ ปุฏฺโฐ ตฺวํ อานนฺท เอวํ พฺยากเรยฺยาสีติ.\csromanspacer{}\csromanspacer{}ฉฏฺฐํ.}

\csromanheader{2}{7. อนุธมฺมสุตฺต}
\proseitem{39}{สาวตฺถินิทานํ.\csromanspacer{} ธมฺมานุธมฺมปฺปฏิปนฺนสฺส ภิกฺขเว ภิกฺขุโน อยมนุธมฺโม โหติ “ยํ รูเป นิพฺพิทาพหุโล\footnote{นิพฺพิทาพหุลํ (อิ, ก)} วิหเรยฺย, เวทนาย นิพฺพิทาพหุโล วิหเรยฺย, สญฺญาย นิพฺพิทาพหุโล วิหเรยฺย, สงฺขาเรสุ นิพฺพิทาพหุโล วิหเรยฺย, วิญฺญาเณ นิพฺพิทาพหุโล วิหเรยฺย.\csromanspacer{} โย รูเป นิพฺพิทาพหุโล วิหรนฺโต, เวทนาย.\csromanspacer{} สญฺญาย.\csromanspacer{} สงฺขาเรสุ นิพฺพิทาพหุโล วิหรนฺโต, วิญฺญาเณ นิพฺพิทาพหุโล วิหรนฺโต รูปํ ปริชานาติ.\csromanspacer{} เวทนํ.\csromanspacer{} สญฺญํ.\csromanspacer{} สงฺขาเร.\csromanspacer{} วิญฺญาณํ ปริชานาติ, โส รูปํ ปริชานํ, เวทนํ.\csromanspacer{} สญฺญํ.\csromanspacer{} สงฺขาเร.\csromanspacer{} วิญฺญาณํ ปริชานํ ปริมุจฺจติ รูปมฺหา, ปริมุจฺจติ เวทนาย, ปริมุจฺจติ \csromanfolio{34}สญฺญาย, ปริมุจฺจติ สงฺขาเรหิ, ปริมุจฺจติ วิญฺญาณมฺหา, ปริมุจฺจติ ชาติยา ชารามรเณน โสเกหิ ปริเทเวหิ ทุกฺเขหิ โทมนสฺเสหิ อุปายาเสหิ, ปริมุจฺจติ ทุกฺขสฺมาติ วทามี”ติ.\csromanspacer{}\csromanspacer{}สตฺตมํ.}

\csromanheader{2}{8. ทุติย-อนุธมฺมสุตฺต}
\proseitem{40}{สาวตฺถินิทานํ.\csromanspacer{} ธมฺมานุธมฺมปฺปฏิปนฺนสฺส ภิกฺขเว ภิกฺขุโน อยมนุธมฺโม โหติ “ยํ รูเป อนิจฺจานุปสฺสี วิหเรยฺย ฯเปฯ ปริมุจฺจติ ทุกฺขสฺมาติ วทามี”ติ.\csromanspacer{}\csromanspacer{}อฏฺฐมํ.}

\csromanheader{2}{9. ตติย-อนุธมฺมสุตฺต}
\proseitem{41}{สาวตฺถินิทานํ.\csromanspacer{} ธมฺมานุธมฺมปฺปฏิปนฺนสฺส ภิกฺขเว ภิกฺขุโน อยมนุธมฺโม โหติ “ยํ รูเป ทุกฺขานุปสฺสี วิหเรยฺย ฯเปฯ ปริมุจฺจติ ทุกฺขสฺมาติ วทามี”ติ.\csromanspacer{}\csromanspacer{}นวมํ.}

\csromanheader{2}{10. จตุตฺถ-อนุธมฺมสุตฺต}
\proseitem{42}{สาวตฺถินิทานํ.\csromanspacer{} ธมฺมานุธมฺมปฺปฏิปนฺนสฺส ภิกฺขเว ภิกฺขุโน อยมนุธมฺโม โหติ “ยํ รูเป อนตฺตานุปสฺสี วิหเรยฺย.\csromanspacer{} เวทนาย.\csromanspacer{} สญฺญาย.\csromanspacer{} สงฺขาเรสุ.\csromanspacer{} วิญฺญาเณ อนตฺตานุปสฺสี วิหเรยฺย.\csromanspacer{} โย รูเป อนตฺตานุปสฺสี วิหรนฺโต ฯเปฯ รูปํ ปริชานาติ, เวทนํ.\csromanspacer{} สญฺญํ.\csromanspacer{} สงฺขาเร.\csromanspacer{} วิญฺญาณํ ปริชานาติ.\csromanspacer{} โส รูปํ ปริชานํ, เวทนํ.\csromanspacer{} สญฺญํ.\csromanspacer{} สงฺขาเร.\csromanspacer{} วิญฺญาณํ ปริชานํ ปริมุจฺจติ รูปมฺหา, ปริมุจฺจติ เวทนาย, ปริมุจฺจติ สญฺญาย, ปริมุจฺจติ สงฺขาเรหิ, ปริมุจฺจติ วิญฺญาณมฺหา, ปริมุจฺจติ ชาติยา ชรามรเณน โสเกหิ ปริเทเวหิ ทุกฺเขหิ โทมนสฺเสหิ อุปายาเสหิ, ปริมุจฺจติ ทุกฺขสฺมาติ วทามี”ติ.\csromanspacer{}\csromanspacer{}ทสมํ.\csromanspacer{} นตุมฺหากวคฺโค จตุตฺโถ.}

\titlehead{\textbf{ตสฺสุทฺทานํ}}
\csromangathameasure{นตุมฺหาเกน เทฺว วุตฺตา, ภิกฺขูหิ อปเร ทุเว. \\ อานนฺเทน จ เทฺว วุตฺตา, อนุธมฺเมหิ เทฺว ทุกาติ.}
\csromangathagroup{\csromangathastanza{นตุมฺหาเกน เทฺว วุตฺตา, ภิกฺขูหิ อปเร ทุเว. \\ อานนฺเทน จ เทฺว วุตฺตา, อนุธมฺเมหิ เทฺว ทุกาติ.}}

\chapterheadpageread{5. อตฺตทีปวคฺค}
\csromanheader{2}{1. อตฺตทีปสุตฺต}
\proseitem{43}{\csromanfolio{35}สาวตฺถินิทานํ.\csromanspacer{} อตฺตทีปา ภิกฺขเว วิหรถ อตฺตสรณา อนญฺญสรณา, ธมฺมทีปา ธมฺมสรณา อนญฺญสรณา.\csromanspacer{} อตฺตทีปานํ ภิกฺขเว วิหรตํ อตฺตสรณานํ อนญฺญสรณานํ, ธมฺมทีปานํ ธมฺมสรณานํ อนญฺญสรณานํ โยนิ อุปปริกฺขิตพฺพา “กิํชาติกา โสกปริเทวทุกฺขโทมนสฺสุปายาสา กิํ ปโหติกา”ติ.}

\prose{กิํชาติกา จ ภิกฺขเว โสกปริเทวทุกฺขโทมนสฺสุปายาสา กิํ ปโหติกา.\csromanspacer{} อิธ ภิกฺขเว อสฺสุตวา ปุถุชฺชโน อริยานํ อทสฺสาวี อริยธมฺมสฺส อโกวิโท อริยธมฺเม อวินีโต, สปฺปุริสานํ อทสฺสาวี สปฺปุริสธมฺมสฺส อโกวิโท สปฺปุริสธมฺเม อวินีโต รูปํ อตฺตโต สมนุปสฺสติ, รูปวนฺตํ วา อตฺตานํ, อตฺตนิ วา รูปํ, รูปสฺมิํ วา อตฺตานํ.\csromanspacer{} ตสฺส ตํ รูปํ วิปริณมติ อญฺญถา จ โหติ, ตสฺส รูปวิปริณามญฺญถาภาวา อุปฺปชฺชนฺติ โสกปริเทวทุกฺขโทมนสฺสุปายาสา.\csromanspacer{} เวทนํ อตฺตโต สมนุปสฺสติ, เวทนาวนฺตํ วา อตฺตานํ, อตฺตนิ วา เวทนํ, เวทนาย วา อตฺตานํ, ตสฺส สา เวทนา วิปริณมติ อญฺญถา จ โหติ, ตสฺส เวทนาวิปริณามญฺญถาภาวา อุปฺปชฺชนฺติ โสกปริเทว ฯเปฯ ปายาสา.\csromanspacer{} สญฺญํ อตฺตโต สมนุปสฺสติ.\csromanspacer{} สงฺขาเร อตฺตโต สมนุปสฺสติ.\csromanspacer{} วิญฺญาณํ อตฺตโต สมนุปสฺสติ, วิญฺญาณวนฺตํ วา อตฺตานํ, อตฺตนิ วา วิญฺญาณํ, วิญฺญาณสฺมิํ วา อตฺตานํ.\csromanspacer{} ตสฺส ตํ วิญฺญาณํ วิปริณมติ อญฺญถา จ โหติ, ตสฺส วิญฺญาณวิปริณามญฺญถาภาวา อุปฺปชฺชนฺติ โสกปริเทวทุกฺขโทมนสฺสุปายาสา.}

\prose{รูปสฺส เตฺวว ภิกฺขเว อนิจฺจตํ วิทิตฺวา วิปริณามํ วิราคํ นิโรธํ\footnote{วิปริณาม วิราค นิโรธํ (สี)} “ปุพฺเพ เจว รูปํ เอตรหิ จ สพฺพํ รูปํ อนิจฺจํ ทุกฺขํ วิปริณามธมฺมนฺ”ติ เอวเมตํ ยถาภูตํ สมฺมปฺปญฺญาย ปสฺสโต เย โสกปริเทวทุกฺขโทมนสฺสุปายาสา, เต ปหียนฺติ.\csromanspacer{} เตสํ ปหานา น ปริตสฺสติ, อปริตสฺสํ สุขํ วิหรติ, สุขวิหารี ภิกฺขุ “ตทงฺคนิพฺพุโต”ติ วุจฺจติ.\csromanspacer{} เวทนาย เตฺวว ภิกฺขเว อนิจฺจตํ วิทิตฺวา วิปริณามํ วิราคํ นิโรธํ “ปุพฺเพ เจว เวทนา เอตรหิ จ สพฺพา เวทนา อนิจฺจา ทุกฺขา วิปริณามธมฺมา”ติ เอวเมตํ \csromanfolio{36}ยถาภูตํ สมฺมปฺปญฺญาย ปสฺสโต เย โสกปริเทวทุกฺขโทมนสฺสุปายาสา, เต ปหียนฺติ.\csromanspacer{} เตสํ ปหานา น ปริตสฺสติ, อปริตสฺสํ สุขํ วิหรติ, สุขวิหารี ภิกฺขุ “ตทงฺคนิพฺพุโต”ติ วุจฺจติ.\csromanspacer{} สญฺญาย.\csromanspacer{} สงฺขารานํ เตฺวว ภิกฺขเว อนิจฺจตํ วิทิตฺวา วิปริณามํ วิราคํ นิโรธํ “ปุพฺเพ เจว สงฺขารา เอตรหิ จ สพฺเพ สงฺขารา อนิจฺจา ทุกฺขา วิปริณามธมฺมา”ติ เอวเมตํ ยถาภูตํ สมฺมปฺปญฺญาย ปสฺสโต เย โสกปริเทวทุกฺขโทมนสฺสุปายาสา, เต ปหียนฺติ.\csromanspacer{} เตสํ ปหานา น ปริตสฺสติ, อปริตสฺสํ สุขํ วิหรติ, สุขวิหารี ภิกฺขุ “ตทงฺคนิพฺพุโต”ติ วุจฺจติ.\csromanspacer{} วิญฺญาณสฺส เตฺวว ภิกฺขเว อนิจฺจตํ วิทิตฺวา วิปริณามํ วิราคํ นิโรธํ “ปุพฺเพ เจว วิญฺญาณํ เอตรหิ จ สพฺพํ วิญฺญาณํ อนิจฺจํ ทุกฺขํ วิปริณามธมฺมนฺ”ติ เอวเมตํ ยถาภูตํ สมฺมปฺปญฺญาย ปสฺสโต เย โสกปริเทวทุกฺขโทมนสฺสุปายาสา, เต ปหียนฺติ.\csromanspacer{} เตสํ ปหานา น ปริตสฺสติ, อปริตสฺสํ สุขํ วิหรติ, สุขวิหารี ภิกฺขุ “ตทงฺคนิพฺพุโต”ติ วุจฺจตีติ.\csromanspacer{}\csromanspacer{}ปฐมํ.}

\csromanheader{2}{2. ปฏิปทาสุตฺต}
\proseitem{44}{สาวตฺถินิทานํ.\csromanspacer{} สกฺกายสมุทยคามินิญฺจ โว ภิกฺขเว ปฏิปทํ เทเสสฺสามิ สกฺกายนิโรธคามินิญฺจ ปฏิปทํ, ตํ สุณาถ.\csromanspacer{} กตมา จ ภิกฺขเว สกฺกายสมุทยคามินี ปฏิปทา, อิธ ภิกฺขเว อสฺสุตวา ปุถุชฺชโน อริยานํ อทสฺสาวี อริยธมฺมสฺส อโกวิโท อริยธมฺเม อวินีโต, สปฺปุริสานํ อทสฺสาวี สปฺปุริสธมฺมสฺส อโกวิโท สปฺปุริสธมฺเม อวินีโต รูปํ อตฺตโต สมนุปสฺสติ, รูปวนฺตํ วา อตฺตานํ, อตฺตนิ วา รูปํ, รูปสฺมิํ วา อตฺตานํ.\csromanspacer{} เวทนํ อตฺตโต.\csromanspacer{} สญฺญํ.\csromanspacer{} สงฺขาเร.\csromanspacer{} วิญฺญาณํ อตฺตโต สมนุปสฺสติ, วิญฺญาณวนฺตํ วา อตฺตานํ, อตฺตนิ วา วิญฺญาณํ, วิญฺญาณสฺมิํ วา อตฺตานํ.\csromanspacer{} อยํ วุจฺจติ ภิกฺขเว สกฺกายสมุทยคามินีปฏิปทา. “สกฺกายสมุทยคามินี ปฏิปทา”ติ อิติ หิทํ ภิกฺขเว วุจฺจติ “ทุกฺขสมุทยคามินี สมนุปสฺสนา”ติ อยเมเวตฺถ อตฺโถ.}

\prose{กตมา จ ภิกฺขเว สกฺกายนิโรธคามินี ปฏิปทา, อิธ ภิกฺขเว สุตวา อริยสาวโก อริยานํ ทสฺสาวี อริยธมฺมสฺส โกวิโท อริยธมฺเม สุวินีโต, สปฺปุริสานํ ทสฺสาวี สปฺปุริสธมฺมสฺส โกวิโท สปฺปุริสธมฺเม สุวินีโต น รูปํ อตฺตโต สมนุปสฺสติ, น รูปวนฺตํ วา \csromanfolio{37}อตฺตานํ, น อตฺตนิ วา รูปํ, น รูปสฺมิํ วา อตฺตานํ.\csromanspacer{} น เวทนํ อตฺตโต.\csromanspacer{} น สญฺญํ.\csromanspacer{} น สงฺขาเร.\csromanspacer{} น วิญฺญาณํ อตฺตโต สมนุปสฺสติ, น วิญฺญาณวนฺตํ วา อตฺตานํ, น อตฺตนิ วา วิญฺญาณํ, น วิญฺญาณสฺมิํ วา อตฺตานํ.\csromanspacer{} อยํ วุจฺจติ ภิกฺขเว สกฺกายนิโรธคามินี ปฏิปทา. “สกฺกายนิโรธคามินี ปฏิปทา”ติ อิติ หิทํ ภิกฺขเว วุจฺจติ “ทุกฺขนิโรธคามินี สมนุปสฺสนา”ติ อยเมเวตฺถ อตฺโถติ.\csromanspacer{}\csromanspacer{}ทุติยํ.}

\csromanheader{2}{3. อนิจฺจสุตฺต}
\proseitem{45}{สาวตฺถินิทานํ.\csromanspacer{} รูปํ ภิกฺขเว อนิจฺจํ, ยทนิจฺจํ ตํ ทุกฺขํ, ยํ ทุกฺขํ ตทนตฺตา, ยทนตฺตา ตํ “เนตํ มม, เนโสหมสฺมิ, น เมโส อตฺตา”ติ เอวเมตํ ยถาภูตํ สมฺมปฺปญฺญาย ทฏฺฐพฺพํ, เอวเมตํ ยถาภูตํ สมฺมปฺปญฺญาย ปสฺสโต จิตฺตํ วิรชฺชติ, วิมุจฺจติ อนุปาทาย อาสเวหิ.\csromanspacer{} เวทนา อนิจฺจา.\csromanspacer{} สญฺญา.\csromanspacer{} สงฺขารา.\csromanspacer{} วิญฺญาณํ อนิจฺจํ, ยทนิจฺจํ ตํ ทุกฺขํ, ยํ ทุกฺขํ ตทนตฺตา, ยทนตฺตา ตํ “เนตํ มม, เนโสหมสฺมิ, น เมโส อตฺตา”ติ เอวเมตํ ยถาภูตํ สมฺมปฺปญฺญาย ทฏฺฐพฺพํ, เอวเมตํ ยถาภูตํ สมฺมปฺปญฺญาย ปสฺสโต จิตฺตํ วิรชฺชติ, วิมุจฺจติ อนุปาทาย อาสเวหิ.\csromanspacer{} รูปธาตุยา เจ ภิกฺขเว ภิกฺขุโน จิตฺตํ วิรตฺตํ, วิมุตฺตํ โหติ อนุปาทาย อาสเวหิ.\csromanspacer{} เวทนาธาตุยา ฯเปฯ.\csromanspacer{} สญฺญาธาตุยา.\csromanspacer{} สงฺขารธาตุยา.\csromanspacer{} วิญฺญาณธาตุยา เจ ภิกฺขเว ภิกฺขุโน จิตฺตํ วิรตฺตํ, วิมุตฺตํ โหติ อนุปาทาย อาสเวหิ.\csromanspacer{} วิมุตฺตตฺตา ฐิตํ, ฐิตตฺตา สนฺตุสิตํ 1.\csromanspacer{} สนฺตุสิตตฺตา น ปริตสฺสติ, อปริตสฺสํ ปจฺจตฺตญฺเญว ปรินิพฺพายติ, “ขีณา ชาติ, วุสิตํ พฺรหฺมจริยํ, กตํ กรณียํ, นาปรํ อิตฺถตฺตายา”ติ ปชานาตีติ.\csromanspacer{}\csromanspacer{}ตติยํ.}

\csromanheader{2}{4. ทุติย-อนิจฺจสุตฺต}
\proseitem{46}{สาวตฺถินิทานํ.\csromanspacer{} รูปํ ภิกฺขเว อนิจฺจํ, ยทนิจฺจํ ตํ ทุกฺขํ, ยํ ทุกฺขํ ตทนตฺตา.\csromanspacer{} ยทนตฺตา ตํ “เนตํ มม, เนโสหมสฺมิ, น เมโส อตฺตา”ติ เอวเมตํ ยถาภูตํ สมฺมปฺปญฺญาย ทฏฺฐพฺพํ.\csromanspacer{} เวทนา อนิจฺจา.\csromanspacer{} สญฺญา อนิจฺจา.\csromanspacer{} สงฺขารา อนิจฺจา.\csromanspacer{} วิญฺญาณํ อนิจฺจํ, ยทนิจฺจํ ตํ ทุกฺขํ, ยํ ทุกฺขํ ตทนตฺตา, ยทนตฺตา ตํ “เนตํ มม, เนโสหมสฺมิ, น เมโส อตฺตา”ติ เอวเมตํ ยถาภูตํ สมฺมปฺปญฺญาย ทฏฺฐพฺพํ.}

\prose{\csromanfolio{38}เอวเมตํ ยถาภูตํ สมฺมปฺปญฺญาย ปสฺสโต ปุพฺพนฺตานุทิฏฺฐิโย น โหนฺติ, ปุพฺพนฺตานุทิฏฺฐีนํ อสติ อปรนฺตานุทิฏฺฐิโย น โหนฺติ, อปรนฺตานุทิฏฺฐีนํ อสติ ถามโส\footnote{ถามสา (สี, สฺยา, กํ)} ปรามาโส น โหติ, ถามเส\footnote{ถามสา (สี, สฺยา, กํ), ถามโส (ก)} ปรามาเส อสติ รูปสฺมิํ.\csromanspacer{} เวทนาย.\csromanspacer{} สญฺญาย.\csromanspacer{} สงฺขาเรสุ.\csromanspacer{} วิญฺญาณสฺมิํ จิตฺตํ วิรชฺชติ, วิมุจฺจติ อนุปาทาย อาสเวหิ.\csromanspacer{} วิมุตฺตตฺตา ฐิตํ, ฐิตตฺตา สนฺตุสิตํ, สนฺตุสิตตฺตา น ปริตสฺสติ, อปริตสฺสํ ปจฺจตฺตญฺเญว ปรินิพฺพายติ, “ขีณา ชาติ, วุสิตํ พฺรหฺมจริยํ, กตํ กรณียํ, นาปรํ อิตฺถตฺตายา”ติ ปชานาตีติ.\csromanspacer{}\csromanspacer{}จตุตฺถํ.}

\csromanheader{2}{5. สมนุปสฺสนาสุตฺต}
\proseitem{47}{สาวตฺถินิทานํ.\csromanspacer{} เย หิ เกจิ ภิกฺขเว สมณา วา พฺราหฺมณา วา อเนกวิหิตํ อตฺตานํ สมนุปสฺสมานา สมนุปสฺสนฺติ, สพฺเพเต ปญฺจุปาทานกฺขนฺเธ สมนุปสฺสนฺติ เอเตสํ วา อญฺญตรํ.\csromanspacer{} กตเม ปญฺจ, อิธ ภิกฺขเว อสฺสุตวา ปุถุชฺชโน อริยานํ อทสฺสาวี อริยธมฺมสฺส อโกวิโท อริยธมฺเม อวินีโต, สปฺปุริสานํ อทสฺสาวี สปฺปุริสธมฺมสฺส อโกวิโท สปฺปุริสธมฺเม อวินีโต รูปํ อตฺตโต สมนุปสฺสติ, รูปวนฺตํ วา อตฺตานํ, อตฺตนิ วา รูปํ, รูปสฺมิํ วา อตฺตานํ.\csromanspacer{} เวทนํ.\csromanspacer{} สญฺญํ.\csromanspacer{} สงฺขาเร.\csromanspacer{} วิญฺญาณํ อตฺตโต สมนุปสฺสติ, วิญฺญาณวนฺตํ วา อตฺตานํ, อตฺตนิ วา วิญฺญาณํ, วิญฺญาณสฺมิํ วา อตฺตานํ.}

\prose{อิติ อยญฺเจว สมนุปสฺสนา “อสฺมี”ติ จสฺส อวิคตํ\footnote{อธิคตํ (พหูสุ)} โหติ, “อสฺมี”ติ โข ปน ภิกฺขเว อวิคเต ปญฺจนฺนํ อินฺทฺริยานํ อวกฺกนฺติ โหติ จกฺขุนฺทฺริยสฺส โสตินฺทฺริยสฺส ฆานินฺทฺริยสฺส ชิวฺหินฺทฺริยสฺส กายินฺทฺริยสฺส.\csromanspacer{} อตฺถิ ภิกฺขเว มโน, อตฺถิ ธมฺมา, อตฺถิ อวิชฺชาธาตุ.\csromanspacer{} อวิชฺชาสมฺผสฺสเชน ภิกฺขเว เวทยิเตน ผุฏฺฐสฺส อสฺสุตวโต ปุถุชฺชนสฺส “อสฺมี”ติปิสฺส โหติ, “อยมหมสฺมี”ติปิสฺส โหติ, “ภวิสฺสนฺ”ติปิสฺส โหติ, “น ภวิสฺสนฺ”ติปิสฺส โหติ, “รูปี ภวิสฺสนฺ”ติปิสฺส โหติ, “อรูปี ภวิสฺสนฺ”ติปิสฺส โหติ, “สญฺญี ภวิสฺสนฺ”ติปิสฺส โหติ, “อสญฺญี ภวิสฺสนฺ”ติปิสฺส โหติ, “เนวสญฺญีนาสญฺญี ภวิสฺสนฺ”ติปิสฺส โหติ.}

\prose{\csromanfolio{39}ติฏฺฐนฺเตว โข\footnote{ติฏฺฐนฺติ โข ปน (สี, สฺยา, กํ, อิ)} ภิกฺขเว ตตฺเถว\footnote{ตเถว (กตฺถจิ)} ปญฺจินฺทฺริยานิ, อเถตฺถ สุตวโต อริยสาวกสฺส อวิชฺชา ปหียติ, วิชฺชา อุปฺปชฺชติ, ตสฺส อวิชฺชาวิราคา วิชฺชุปฺปาทา “อสฺมี”ติปิสฺส น โหติ, “อยมหมสฺมี”ติปิสฺส น โหติ, “ภวิสฺสนฺ”ติ. “น ภวิสฺสนฺ”ติ. “รูปี.\csromanspacer{} อรูปี.\csromanspacer{} สญฺญี.\csromanspacer{} อสญฺญี.\csromanspacer{} เนวสญฺญีนาสญฺญี ภวิสฺสนฺ”ติปิสฺส น โหตีติ.\csromanspacer{}\csromanspacer{}ปญฺจมํ.}

\csromanheader{2}{6. ขนฺธสุตฺต}
\proseitem{48}{สาวตฺถินิทานํ.\csromanspacer{} ปญฺจ ภิกฺขเว ขนฺเธ เทเสสฺสามิ ปญฺจุปาทานกฺขนฺเธ จ, ตํ สุณาถ.\csromanspacer{} กตเม จ ภิกฺขเว ปญฺจกฺขนฺธา, ยํ กิญฺจิ ภิกฺขเว รูปํ อตีตานาคตปจฺจุปฺปนฺนํ อชฺฌตฺตํ วา พหิทฺธา วา โอฬาริกํ วา สุขุมํ วา หีนํ วา ปณีตํ วา ยํ ทูเร สนฺติเก วา, อยํ วุจฺจติ รูปกฺขนฺโธ.\csromanspacer{} ยา กาจิ เวทนา ฯเปฯ.\csromanspacer{} ยา กาจิ สญฺญา.\csromanspacer{} เย เกจิ สงฺขารา อตีตานาคตปจฺจุปฺปนฺนา อชฺฌตฺตํ วา พหิทฺธา วา โอฬาริกา วา สุขุมา วา ฯเปฯ.\csromanspacer{} อยํ วุจฺจติ สงฺขารกฺขนฺโธ.\csromanspacer{} ยํ กิญฺจิ วิญฺญาณํ อตีตานาคตปจฺจุปฺปนฺนํ อชฺฌตฺตํ วา พหิทฺธา วา โอฬาริกํ วา สุขุมํ วา หีนํ วา ปณีตํ วา ยํ ทูเร สนฺติเก วา, อยํ วุจฺจติ วิญฺญาณกฺขนฺโธ.\csromanspacer{} อิเม วุจฺจนฺติ ภิกฺขเว ปญฺจกฺขนฺธา.}

\prose{กตเม จ ภิกฺขเว ปญฺจุปาทานกฺขนฺธา, ยํ กิญฺจิ ภิกฺขเว รูปํ อตีตานาคตปจฺจุปฺปนฺนํ ฯเปฯ ยํ ทูเร สนฺติเก วา สาสวํ อุปาทานิยํ, อยํ วุจฺจติ รูปุปาทานกฺขนฺโธ.\csromanspacer{} ยา กาจิ เวทนา ฯเปฯ ยา ทูเร สนฺติเก วา สาสวา อุปาทานิยา, อยํ วุจฺจติ เวทนุปาทานกฺขนฺโธ.\csromanspacer{} ยา กาจิ สญฺญา ฯเปฯ ยา ทูเร สนฺติเก วา สาสวา อุปาทานิยา, อยํ วุจฺจติ สญฺญุปาทานกฺขนฺโธ.\csromanspacer{} เย เกจิ สงฺขารา ฯเปฯ สาสวา อุปาทานิยา, อยํ วุจฺจติ สงฺขารุปาทานกฺขนฺโธ.\csromanspacer{} ยํ กิญฺจิ วิญฺญาณํ อตีตานาคตปจฺจุปฺปนฺนํ ฯเปฯ ยํ ทูเร สนฺติเก วา สาสวํ อุปาทานิยํ, อยํ วุจฺจติ วิญฺญาณุปาทานกฺขนฺโธ.\csromanspacer{} อิเม วุจฺจนฺติ ภิกฺขเว ปญฺจุปาทานกฺขนฺธาติ.\csromanspacer{}\csromanspacer{}ฉฏฺฐํ.}

\csromanheader{2}{7. โสณสุตฺต}
\proseitem{49}{เอวํ เม สุตํ\csromandash{}เอกํ สมยํ ภควา ราชคเห วิหรติ เวฬุวเน กลนฺทกนิวาเป.\csromanspacer{} อถ โข โสโณ คหปติปุตฺโต เยน ภควา \csromanfolio{40}เตนุปสงฺกมิ ฯเปฯ เอกมนฺตํ นิสินฺนํ โข โสณํ คหปติปุตฺตํ ภควา เอตทโวจ\csromandash{}}

\prose{เย หิ เกจิ โสณ สมณา วา พฺราหฺมณา วา อนิจฺเจน รูเปน ทุกฺเขน วิปริณามธมฺเมน “เสยฺโยหมสฺมี”ติ วา สมนุปสฺสนฺติ, “สทิโสหมสฺมี”ติ วา สมนุปสฺสนฺติ, “หีโนหมสฺมี”ติ วา สมนุปสฺสนฺติ, กิมญฺญตฺร ยถาภูตสฺส อทสฺสนา.\csromanspacer{} อนิจฺจาย เวทนาย ทุกฺขาย วิปริณามธมฺมาย “เสยฺโยหมสฺมี”ติ วา สมนุปสฺสนฺติ, “สทิโสหมสฺมี”ติ วา สมนุปสฺสนฺติ, “หีโนหมสฺมี”ติ วา สมนุปสฺสนฺติ, กิมญฺญตฺร ยถาภูตสฺส อทสฺสนา.\csromanspacer{} อนิจฺจาย สญฺญาย.\csromanspacer{} อนิจฺเจหิ สงฺขาเรหิ ทุกฺเขหิ วิปริณามธมฺเมหิ “เสยฺโยหมสฺมี”ติ วา สมนุปสฺสนฺติ, “สทิโสหมสฺมี”ติ วา สมนุปสฺสนฺติ, “หีโนหมสฺมี”ติ วา สมนุปสฺสนฺติ, กิมญฺญตฺร ยถาภูตสฺส อทสฺสนา.\csromanspacer{} อนิจฺเจน วิญฺญาเณน ทุกฺเขน วิปริณามธมฺเมน “เสยฺโยหมสฺมี”ติ วา สมนุปสฺสนฺติ, “สทิโสหมสฺมี”ติ วา สมนุปสฺสนฺติ, “หีโนหมสฺมี”ติ วา สมนุปสฺสนฺติ, กิมญฺญตฺร ยถาภูตสฺส อทสฺสนา.}

\prose{เย จ โข เกจิ โสณ สมณา วา พฺราหฺมณา วา อนิจฺเจน รูเปน ทุกฺเขน วิปริณามธมฺเมน “เสยฺโยหมสฺมี”ติปิ น สมนุปสฺสนฺติ, “สทิโสหมสฺมี”ติปิ น สมนุปสฺสนฺติ, “หีโนหมสฺมี”ติปิ น สมนุปสฺสนฺติ, กิมญฺญตฺร ยถาภูตสฺส ทสฺสนา.\csromanspacer{} อนิจฺจาย เวทนาย.\csromanspacer{} อนิจฺจาย สญฺญาย.\csromanspacer{} อนิจฺเจหิ สงฺขาเรหิ.\csromanspacer{} อนิจฺเจน วิญฺญาเณน ทุกฺเขน วิปริณามธมฺเมน “เสยฺโยหมสฺมี”ติปิ น สมนุปสฺสนฺติ, “สทิโสหมสฺมี”ติปิ น สมนุปสฺสนฺติ, “หีโนหมสฺมี”ติปิ น สมนุปสฺสนฺติ, กิมญฺญตฺร ยถาภูตสฺส ทสฺสนา.}

\prose{ตํ กิํ มญฺญสิ โสณ, รูปํ นิจฺจํ วา อนิจฺจํ วาติ.\csromanspacer{} อนิจฺจํ ภนฺเต.\csromanspacer{} ยํ ปนานิจฺจํ, ทุกฺขํ วา ตํ สุขํ วาติ.\csromanspacer{} ทุกฺขํ ภนฺเต.\csromanspacer{} ยํ ปนานิจฺจํ ทุกฺขํ วิปริณามธมฺมํ, กลฺลํ นุ ตํ สมนุปสฺสิตุํ “เอตํ มม, เอโสหมสฺมิ, เอโส เม อตฺตา”ติ.\csromanspacer{} โน เหตํ ภนฺเต.\csromanspacer{} เวทนา นิจฺจา วา อนิจฺจา วาติ.\csromanspacer{} อนิจฺจา ภนฺเต.\csromanspacer{} สญฺญา.\csromanspacer{} สงฺขารา.\csromanspacer{} วิญฺญาณํ นิจฺจํ วา อนิจฺจํ วาติ.\csromanspacer{} อนิจฺจํ ภนฺเต.\csromanspacer{} ยํ ปนานิจฺจํ, ทุกฺขํ วา ตํ สุขํ วาติ.\csromanspacer{} ทุกฺขํ ภนฺเต.\csromanspacer{} ยํ ปนานิจฺจํ ทุกฺขํ วิปริณมธมฺมํ, กลฺลํ นุ ตํ สมนุปสฺสิตุํ “เอตํ มม, เอโสหมสฺมิ, เอโส เม อตฺตา”ติ.\csromanspacer{} โน เหตํ ภนฺเต.}

\prose{\csromanfolio{41}ตสฺมาติห โสณ ยํ กิญฺจิ รูปํ อตีตานาคตปจฺจุปฺปนฺนํ อชฺฌตฺตํ วา พหิทฺธา วา โอฬาริกํ วา สุขุมํ วา หีนํ วา ปณีตํ วา ยํ ทูเร สนฺติเก วา, สพฺพํ รูปํ “เนตํ มม, เนโสหมสฺมิ, น เมโส อตฺตา”ติ เอวเมตํ ยถาภูตํ สมฺมปฺปญฺญาย ทฏฺฐพฺพํ.}

\prose{ยา กาจิ เวทนา.\csromanspacer{} ยา กาจิ สญฺญา.\csromanspacer{} เย เกจิ สงฺขารา.\csromanspacer{} ยํ กิญฺจิ วิญฺญาณํ อตีตานาคตปจฺจุปฺปนฺนํ อชฺฌตฺตํ วา พหิทฺธา วา โอฬาริกํ วา สุขุมํ วา หีนํ วา ปณีตํ วา ยํ ทูเร สนฺติเก วา, สพฺพํ วิญฺญาณํ “เนตํ มม, เนโสหมสฺมิ, น เมโส อตฺตา”ติ เอวเมตํ ยถาภูตํ สมฺมปฺปญฺญาย ทฏฺฐพฺพํ.}

\prose{เอวํ ปสฺสํ โสณ สุตวา อริยสาวโก รูปสฺมิมฺปิ นิพฺพินฺทติ.\csromanspacer{} เวทนายปิ นิพฺพินฺทติ.\csromanspacer{} สญฺญายปิ นพฺพินฺทติ.\csromanspacer{} สงฺขาเรสุปิ นิพฺพินฺทติ.\csromanspacer{} วิญฺญาณสฺมิมฺปิ นิพฺพินฺทติ, นิพฺพินฺทํ วิรชฺชติ, วิราคา วิมุจฺจติ, วิมุตฺตสฺมิํ “วิมุตฺตมฺ”อิติ ญาณํ โหติ, “ขีณา ชาติ, วุสิตํ พฺราหฺมจริยํ, กตํ กรณียํ, นาปรํ อิตฺถตฺตายา”ติ ปชานาตีติ.\csromanspacer{}\csromanspacer{}สตฺตมํ.}

\csromanheader{2}{8. ทุติยโสณสุตฺต}
\proseitem{50}{เอวํ เม สุตํ\csromandash{}เอกํ สมยํ ภควา ราชคเห วิหรติ เวฬุวเน กลนฺทกนิวาเป.\csromanspacer{} อถ โข โสโณ คหปติปุตฺโต เยน ภควา เตนุปสงฺกมิ, อุปสงฺกมิตฺวา ภควนฺตํ อภิวนฺตํ อภิวาเทตฺวา เอกมนฺตํ นิสีทิ, เอกมนฺตํ นิสินฺนํ โข โสณํ คหปติปุตฺตํ ภควา เอตทโวจ\csromandash{}}

\prose{เย หิ เกจิ โสณ สมณา วา พฺราหฺมณา วา รูปํ นปฺปชานนฺติ, รูปสมุทยํ นปฺปชานนฺติ, รูปนิโรธํ นปฺปชานนฺติ, รูปนิโรธคามินิํ ปฏิปทํ นปฺปชานนฺติ.\csromanspacer{} เวทนํ นปฺปชานนฺติ, เวทนาสมุทยํ นปฺปชานนฺติ, เวทนานิโรธํ นปฺปชานนฺติ, เวทนานิโรธคามินิํ ปฏิปทํ นปฺปชานนฺติ.\csromanspacer{} สญฺญํ นปฺปชานนฺติ ฯเปฯ.\csromanspacer{} สงฺขาเร นปฺปชานนฺติ, สงฺขารสมุทยํ นปฺปชานนฺติ, สงฺขารนิโรธํ นปฺปชานนฺติ, สงฺขารนิโรธคามินิํ ปฏิปทํ นปฺปชานนฺติ.\csromanspacer{} วิญฺญาณํ นปฺปชานนฺติ, วิญฺญาณสมุทยํ นปฺปชานนฺติ, วิญฺญาณนิโรธํ นปฺปชานนฺติ, วิญฺญาณนิโรธคามินิํ ปฏิปทํ นปฺปชานนฺติ.\csromanspacer{} น เม เต โสณ สมณา วา พฺราหฺมณา วา สมเณสุ วา สมณสมฺมตา, พฺราหฺมเณสุ วา พฺราหฺมณสมฺมตา, น จ ปน เต อายสฺมนฺโต สามญฺญตฺถํ วา พฺรหฺมญฺญตฺถํ วา ทิฏฺเฐว ธมฺเม สยํ อภิญฺญา สจฺฉิกตฺวา อุปสมฺปชฺช วิหรนฺติ.}

\prose{\csromanfolio{42}เย จ โข เกจิ โสณ สมณา วา พฺราหฺมณา วา รูปํ ปชานนฺติ, รูปสมุทยํ ปชานนฺติ, รูปนิโรธํ ปชานนฺติ, รูปนิโรธคามินิํ ปฏิปทํ ปชานนฺติ.\csromanspacer{} เวทนํ ปชานนฺติ ฯเปฯ.\csromanspacer{} สญฺญํ ปชานนฺติ.\csromanspacer{} สงฺขาเร ปชานนฺติ.\csromanspacer{} วิญฺญาณํ ปชานนฺติ, วิญฺญาณสมุทยํ ปชานนฺติ, วิญฺญาณนิโรธํ ปชานนฺติ, วิญฺญาณนิโรธคามินิํ ปฏิปทํ ปชานนฺติ.\csromanspacer{} เต จ โข เม โสณ สมณา วา พฺราหฺมณา วาสมเณสุ เจว สมณสมฺมตา, พฺราหฺมเณสุ จ พฺราหฺมณสมฺมตา, เต จ ปนายสฺมนฺโต สามญฺญตฺถญฺจ พฺรหฺมญฺญตฺถญฺจ ทิฏฺเฐว ธมฺเม สยํ อภิญฺญาสจฺฉิกตฺวา อุปสมฺปชฺช วิหรนฺตีติ.\csromanspacer{}\csromanspacer{}อฏฺฐมํ.}

\csromanheader{2}{9. นนฺทิกฺขยสุตฺต}
\proseitem{51}{สาวตฺถินิทานํ.\csromanspacer{} อนิจฺจญฺเญว ภิกฺขเว ภิกฺขุ รูปํ “อนิจฺจนฺ”ติ ปสฺสติ.\csromanspacer{} สาสฺส โหติ สมฺมาทิฏฺฐิ, สมฺมา ปสฺสํ นิพฺพินฺทติ, นนฺทิกฺขยา ราคกฺขโย, ราคกฺขยา นนฺทิกฺขโย, นนฺทิราคกฺขยา จิตฺตํ วิมุตฺตํ “สุวิมุตฺตนฺ”ติ วุจฺจติ.\csromanspacer{} อนิจฺจญฺเญว ภิกฺขเว ภิกฺขุ เวทนํ “อนิจฺจนฺ”ติ ปสฺสติ.\csromanspacer{} สาสฺส โหติ สมฺมาทิฏฺฐิ, สมฺมา ปสฺสํ นิพฺพินฺทติ, นนฺทิกฺขยา ราคกฺขโย, ราคกฺขยา นินฺทิกฺขโย, นนฺทิราคกฺขยา จิตฺตํ วิมุตฺตํ “สุวิมุตฺตนฺ”ติ วุจฺจติ.\csromanspacer{} อนิจฺจญฺเญว ภิกฺขเว ภิกฺขุ สญฺญํ “อนิจฺจนฺ”ติ ปสฺสติ ฯเปฯ.\csromanspacer{} อนิจฺเจเยว ภิกฺขเว ภิกฺขุ สงฺขาเร “อนิจฺจา”ติ ปสฺสติ.\csromanspacer{} สาสฺส โหติ สมฺมาทิฏฺฐิ, สมฺมา ปสฺสํ นิพฺพินฺทติ, นนฺทิกฺขยา ราคกฺขโย, ราคกฺขยา นนฺทิกฺขโย, นนฺทิราคกฺขยา จิตฺตํ วิมุตฺตํ “สุวิมุตฺตนฺ”ติ วุจฺจติ.\csromanspacer{} อนิจฺจญฺเญว ภิกฺขเว ภิกฺขุ วิญฺญาณํ “อนิจฺจนฺ”ติ ปสฺสติ.\csromanspacer{} สาสฺส โหติ สมฺมาทิฏฺฐิ, สมฺมา ปสฺสํ นิพฺพินฺทติ, นนฺทิกฺขยา ราคกฺขโย, ราคกฺขยา นนฺทิกฺขโย, นนฺทิราคกฺขยา จิตฺตํ วิมุตฺตํ “สุวิมุตฺตนฺ”ติ วุจฺจตีติ.\csromanspacer{}\csromanspacer{}นวมํ.}

\csromanheader{2}{10. ทุติยนนฺทิกฺขยสุตฺต}
\proseitem{52}{สาวตฺถินิทานํ.\csromanspacer{} รูปํ ภิกฺขเว โยนิโส มนสิ กโรถ, รูปานิจฺจตญฺจ ยถาภูตํ สมนุปสฺสถ.\csromanspacer{} รูปํ ภิกฺขเว ภิกฺขุ โยนิโส มนสิ กโรนฺโต รูปนิจฺจตญฺจ ยถาภูตํ สมนุปสฺสนฺโต รูปสฺมิํ นิพฺพินฺทติ, นนฺทิกฺขยา ราคกฺขโย, ราคกฺขยา นนฺทิกฺขโย, นนฺทิราคกฺขยา จิตฺตํ วิมุตฺตํ “สุวิมุตฺตนฺ”ติ วุจฺจติ.\csromanspacer{} เวทนํ ภิกฺขเว โยนิโส มนสิ กโรถ, เวทนานิจฺจตญฺจ ยถาภูตํ สมนุปสฺสถ.\csromanspacer{} เวทนํ ภิกฺขเว ภิกฺขุ โยนิโส \csromanfolio{43}มนสิ กโรนฺโต เวทนานิจฺจตญฺจ ยถาภูตํ สมนุปสฺสนฺโต เวทนาย นิพฺพินฺทติ, นนฺทิกฺขยา ราคกฺขโย, ราคกฺขยา นนฺทิกฺขโย, นนฺทิราคกฺขยา จิตฺตํ วิมุตฺตํ “สุวิมุตฺตนฺ”ติ วุจฺจติ.\csromanspacer{} สญฺญํ ภิกฺขเว.\csromanspacer{} สงฺขาเร ภิกฺขเว โยนิโส มนสิ กโรถ, สงฺขารานิจฺจตญฺจ ยถาภูตํ สมนุปสฺสถ.\csromanspacer{} สงฺขาเร ภิกฺขเว ภิกฺขุ โยนิโส มนสิ กโรนฺโต สงฺขารานิจฺจตํ ยถาภูตํ สมนุปสฺสนฺโต สงฺขาเรสุ นิพฺพินฺทติ, นนฺทิกฺขยา ราคกฺขโย, ราคกฺขยา นนฺทิกฺขโย, นนฺทิราคกฺขยา จิตฺตํ วิมุตฺตํ “สุวิมุตฺตนฺ”ติ วุจฺจติ.\csromanspacer{} วิญฺญาณํ ภิกฺขเว โยนิโส มนสิ กโรถ, วิญฺญาณานิจฺจตญฺจ ยถาภูตํ สมนุปสฺสถ.\csromanspacer{} วิญฺญาณํ ภิกฺขเว ภิกฺขุ โยนิโส มนสิ กโรนฺโต วิญฺญาณานิจฺจตญฺจ ยถาภูตํ สมนุปสฺสนฺโต วิญฺญาณสฺมิํ นิพฺพินฺทติ, นนฺทิกฺขยา ราคกฺขโย, ราคกฺขยา นนฺทิกฺขโย, นนฺทิราคกฺขยา จิตฺตํ วิมุตฺตํ “สุวิมุตฺตนฺ”ติ วุจฺจตีติ.\csromanspacer{}\csromanspacer{}ทสมํ.\csromanspacer{} อตฺตทีปวคฺโค ปญฺจโม.}

\titlehead{\textbf{ตสฺสุทฺทานํ}}
\csromangathameasure{อตฺตทีปา ปฏิปทา, เทฺว จ โหนฺติ อนิจฺจตา. \\ สมนุปสฺสนา ขนฺธา, เทฺว โสณา เทฺว นนฺทิกฺขเยน จาติ. มูลปณฺณาสโก สมตฺโต.}
\csromangathagroup{\csromangathastanza{อตฺตทีปา ปฏิปทา, เทฺว จ โหนฺติ อนิจฺจตา. \\ สมนุปสฺสนา ขนฺธา, เทฺว โสณา เทฺว นนฺทิกฺขเยน จาติ.\csromanspacer{} มูลปณฺณาสโก สมตฺโต.}}

\csromanheader{2}{\textbf{ตสฺส มูลปณฺณาสกสฺส วคฺคุทฺทานํ}}
\csromangathameasure{นกุลปิตา อนิจฺโจ จ, ภาโร นตุมฺหาเกน จ. \\ อตฺตทีเปน ปญฺญาโส, ปฐโม เตน ปวุจฺจตีติ.}
\csromangathagroup{\csromangathastanza{นกุลปิตา อนิจฺโจ จ, ภาโร นตุมฺหาเกน จ. \\ อตฺตทีเปน ปญฺญาโส, ปฐโม เตน ปวุจฺจตีติ.}}

\chapterhead{\textbf{(6) 1. อุปยวคฺค}}
\csromanheader{2}{1. อุปยสุตฺต}
\proseitem{53}{สาวตฺถินิทานํ.\csromanspacer{} อุปโย\footnote{อุปาโย (พหูสุ)} ภิกฺขเว อวิมุตฺโต, อนุปโย วิมุตฺโต.\csromanspacer{} รูปุปยํ\footnote{รูปูปายํ (สี, สฺยา, กํ), รูปุปายํ (อิ, ก)} วา ภิกฺขเว วิญฺญาณํ ติฏฺฐมานํ ติฏฺเฐยฺย, รูปารมฺมณํ \csromanfolio{44}รูปปฺปติฏฺฐํ นนฺทูปเสจนํ วุทฺธิํ เวรูฬฺหิํ เวปุลฺลํ อาปชฺเชยฺย.\csromanspacer{} เวทนุปยํ วา ฯเปฯ.\csromanspacer{} สญฺญุปยํ วา ฯเปฯ.\csromanspacer{} สงฺขารุปยํ วา ภิกฺขเว วิญฺญาณํ ติฏฺฐมานํ ติฏฺเฐยฺย, สงฺขารารมฺมณํ สงฺขารปฺปติฏฺฐํ นนฺทูปเสจนํ วุทฺธิํ วิรูฬฺหิํ เวปุลฺลํ อาปชฺเชยฺย.}

\prose{โย ภิกฺขเว เอวํ วเทยฺย “อหมญฺญตฺร รูปา อญฺญตฺร เวทนาย อญฺญตฺร สญฺญาย อญฺญตฺร สงฺขาเรหิ วิญฺญาณสฺส อาคติํ วา คติํ วา จุติํ วา อุปปตฺติํ วา วุทฺธิํ วา วิรูฬฺหิํ วา เวปุลฺลํ วา ปญฺญาเปสฺสามี”ติ เนตํ ฐานํ วิชฺชติ.}

\prose{รูปธาตุยา เจ ภิกฺขเว ภิกฺขุโน ราโค ปหีโน โหติ, ราคสฺส ปหานา โวจฺฉิชฺชตารมฺมณํ, ปติฏฺฐา วิญฺญาณสฺส น โหติ.\csromanspacer{} เวทนาธาตุยา เจ ภิกฺขเว.\csromanspacer{} สญฺญาธาตุยา เจ ภิกฺขเว.\csromanspacer{} สงฺขารธาตุยา เจ ภิกฺขเว.\csromanspacer{} วิญฺญาณธาตุยา เจ ภิกฺขเว ภิกฺขุโน ราโค ปหีโน โหติ, ราคสฺส ปหานา โวจฺฉิชฺชตารมฺมณํ, ปติฏฺฐา วิญฺญาณสฺส น โหติ, ตทปฺปติฏฺฐิตํ วิญฺญาณํ อวิรูฬฺหํ อนภิสงฺขจฺจ วิมุตฺตํ, วิมุตฺตตฺตา ฐิตํ, ฐิตตฺตา สนฺตุสิตํ, สนฺตุสิตตฺตา น ปริตสฺสติ, อปริตสฺสํ ปจฺจตฺตญฺเญว ปรินิพฺพายติ, “ขีณา ชาติ, วุสิตํ พฺรหฺมจริยํ, กตํ กรณียํ นาปรํ อิตฺถตฺตายา”ติ ปชานาตีติ.\csromanspacer{}\csromanspacer{}ปฐมํ.}

\csromanheader{2}{2. พีชสุตฺต}
\proseitem{54}{สาวตฺถินิทานํ.\csromanspacer{} ปญฺจิมานิ ภิกฺขเว พีชชาตานิ.\csromanspacer{} กตมานิ ปญฺจ, มูลพีชํ ขนฺธพีชํ อคฺคพีชํ ผลุพีชํ พีชพีชญฺเญว ปญฺจมํ.\csromanspacer{} อิมานิ จสฺสุ ภิกฺขเว ปญฺจ พีชชาตานิ อขณฺฑานิ อปูติกานิ อวาตาตปหตานิ สาราทานิ\footnote{สาราทายีนิ (กตฺถจิ)} สุขสยิตานิ, ปถวี\footnote{ปฐวี (สี, สฺยา, กํ, อิ)} จ นาสฺส, อาโป จ นาสฺส, อปิ นุมานิ\footnote{อปิ นุ อิมานิ (สี, อิ)} ภิกฺขเว ปญฺจ พีชชาตานิ วุทฺธิํ วิรูฬฺหิํ เวปุลฺลํ อาปชฺเชยฺยุนฺติ.\csromanspacer{} โน เหตํ ภนฺเต.\csromanspacer{} อิมานิ จสฺสุ ภิกฺขเว ปญฺจ พีชชาตานิ อขณฺฑานิ ฯเปฯ สุขสยิตานิ, ปถวี จ อสฺส, อาโป จ อสฺส, อปิ นุมานิ ภิกฺขเว ปญฺจ พีชชาตานิ วุทฺธิํ วีรูฬฺหิํ เวปุลฺลํ อาปชฺเชยฺยุนฺติ.\csromanspacer{} เอวํ ภนฺเต.\csromanspacer{} เสยฺยถาปิ ภิกฺขเว ปถวีธาตุ, เอวํ จตสฺโส วิญฺญาณฏฺฐิติโย ทฏฺฐพฺพา.\csromanspacer{} เสยฺยถาปิ ภิกฺขเว อาโปธาตุ, เอวํ นนฺทิราโค ทฏฺฐพฺโพ.\csromanspacer{} เสยฺยถาปิ ภิกฺขเว ปญฺจ พีชชาตานิ, เอวํ วิญฺญาณํ สาหารํ ทฏฺฐพฺพํ.}

\prose{\csromanfolio{45}รูปุปยํ ภิกฺขเว วิญฺญาณํ ติฏฺฐมานํ ติฏฺเฐยฺย, รูปารมฺมณํ รูปปฺปติฏฺฐํ นนฺทูปเสจนํ วุทฺธิํ วิรูฬฺหิํ เวปุลฺลํ อาปชฺเชยฺย.\csromanspacer{} เวทนุปยํ วา ภิกฺขเว วิญฺญาณํ ติฏฺฐมานํ ติฏฺเฐยฺย ฯเปฯ.\csromanspacer{} สญฺญุปยํ วา ภิกฺขเว วิญฺญาณํ ติฏฺฐมานํ ติฏฺเฐยฺย ฯเปฯ.\csromanspacer{} สงฺขารุปยํ วา ภิกฺขเว วิญฺญาณํ ติฏฺฐมานํ ติฏฺเฐยฺย.\csromanspacer{} สงฺขารารมฺมณํ สงฺขารปฺปติฏฺฐํ นนฺทูปเสจนํ วุทฺธิํ วิรูฬฺหิํ เวปุลฺลํ อาปชฺเชยฺย.}

\prose{โย ภิกฺขเว เอวํ วเทยฺย “อหมญฺญตฺร รูปา อญฺญตฺร เวทนาย อญฺญตฺร สญฺญาย อญฺญตฺร สงฺขาเรหิ วิญฺญาณสฺส อาคติํ วา คติํ วา จุติํ วา อุปปตฺติํ วา วุทฺธิํ วา วิรูฬฺหิํ วา เวปุลฺลํ วา ปญฺญาเปสฺสามี”ติ เนตํ ฐานํ วิชฺชติ.}

\prose{รูปธาตุยา เจ ภิกฺขเว ภิกฺขุโน ราโค ปหีโน โหติ, ราคสฺส ปหานา โวจฺฉิชฺชตารมฺมณํ, ปติฏฺฐา วิญฺญาณสฺส น โหติ.\csromanspacer{} เวทนาธาตุยา เจ.\csromanspacer{} สญฺญาธาตุยา เจ.\csromanspacer{} สงฺขารธาตุยา เจ.\csromanspacer{} วิญฺญาณธาตุยา เจ ภิกฺขเว ภิกฺขุโน ราโค ปหีโน โหติ, ราคสฺส ปหานา โวจฺฉิชฺชตารมฺมณํ, ปติฏฺฐา วิญฺญาณสฺส น โหติ, ตทปฺปติฏฺฐิตํ วิญฺญาณํ อวิรูฬฺหํ อนภิสงฺขจฺจ วิมุตฺตํ, วิมุตฺตตฺตา ฐิตํ, ฐิตตฺตา สนฺตุสิตํ, สนฺตุสิตตฺตา น ปริตสฺสติ, อปริตสฺสํ ปจฺจตฺตญฺเญว ปรินิพฺพายติ, “ขีณา ชาติ ฯเปฯ นาปรํ อิตฺถตฺตายา”ติ ปชานาตีติ.\csromanspacer{}\csromanspacer{}ทุติยํ.}

\csromanheader{2}{3. อุทานสุตฺต}
\proseitem{55}{สาวตฺถินิทานํ.\csromanspacer{} ตตฺร โข ภควา อุทานํ อุทาเนสิ “โน จสฺสํ, โน จ เม สิยา, นาภวิสฺส, น เม ภวิสฺสตีติ เอวํ อธิมุจฺจมาโน ภิกฺขุ ฉินฺเทยฺย โอรมฺภาคิยานิ สํโยชนานี”ติ.\csromanspacer{} เอวํ วุตฺเต อญฺญตโร ภิกฺขุ ภควนฺตํ เอตทโวจ “ยถา กถํ ปน ภนฺเต โน จสฺสํ, โน จ เม สิยา, นาภวิสฺส, น เม ภวิสฺสตีติ เอวํ อธิมุจฺจมาโน ภิกฺขุ ฉินฺเทยฺย โอรมฺภาคิยานิ สํโยชนานี”ติ.}

\prose{อิธ ภิกฺขุ อสฺสุตวา ปุถุชฺชโน อริยานํ อทสฺสาวี ฯเปฯ สปฺปุริสธมฺเม อวินีโต รูปํ อตฺตโต สมนุปสฺสติ, รูปวนฺตํ วา อตฺตานํ, อตฺตนิ วา รูปํ, รูปสฺมิํ วา อตฺตานํ.\csromanspacer{} เวทนํ.\csromanspacer{} สญฺญํ.\csromanspacer{} สงฺขาเร.\csromanspacer{} วิญฺญาณํ อตฺตโต สมนุปสฺสติ, วิญฺญาณวนฺตํ วา อตฺตานํ, อตฺตนิ วา วิญฺญาณํ, วิญฺญาณสฺมิํ วา อตฺตานํ.}

\prose{\csromanfolio{46}โส อนิจฺจํ รูปํ “อนิจฺจํ รูปนฺ”ติ ยถาภูตํ นปฺปชานาติ.\csromanspacer{} อนิจฺจํ เวทนํ “อนิจฺจา เวทนา”ติ ยถาภูตํ นปฺปชานาติ.\csromanspacer{} อนิจฺจํ สญฺญํ “อนิจฺจา สญฺญา”ติ ยถาภูตํ นปฺปชานาติ.\csromanspacer{} อนิจฺเจ สงฺขาเร “อนิจฺจา สงฺขารา”ติ ยถาภูตํ นปฺปชานาติ.\csromanspacer{} อนิจฺจํ วิญฺญาณํ “อนิจฺจํ วิญฺญาณนฺ”ติ ยถาภูตํ นปฺปชานาติ.}

\prose{ทุกฺขํ รูปํ “ทุกฺขํ รูปนฺ”ติ ยถาภูตํ นปฺปชานาติ, ทุกฺขํ เวทนํ.\csromanspacer{} ทุกฺขํ สญฺญํ.\csromanspacer{} ทุกฺเข สงฺขาเร.\csromanspacer{} ทุกฺขํ วิญฺญาณํ “ทุกฺขํ วิญฺญาณนฺ”ติ ยถาภูตํ นปฺปชานาติ.}

\prose{อนตฺตํ รูปํ “อนตฺตา รูปนฺ”ติ ยถาภูตํ นปฺปชานาติ, อนตฺตํ เวทนํ “อนตฺตา เวทนา”ติ ยถาภูตํ นปฺปชานาติ, อนตฺตํ สญฺญํ “อนตฺตา สญฺญา”ติ ยถาภูตํ นปฺปชานาติ, อนตฺเต สงฺขาเร “อนตฺตา สงฺขารา”ติ ยถาภูตํ นปฺปชานาติ, อนตฺตํ วิญฺญาณํ “อนตฺตา วิญฺญาณนฺ”ติ ยถาภูตํ นปฺปชานาติ.}

\prose{สงฺขตํ รูปํ “สงฺขตํ รูปนฺ”ติ ยถาภูตํ นปฺปชานาติ.\csromanspacer{} สงฺขตํ เวทนํ.\csromanspacer{} สงฺขตํ สญฺญํ.\csromanspacer{} สงฺขเต สงฺขาเร.\csromanspacer{} สงฺขตํ วิญฺญาณํ “สงฺขตํ วิญฺญาณนฺ”ติ ยถาภูตํ นปฺปชานาติ. “รูปํ วิภวิสฺสตี”ติ ยถาภูตํ นปฺปชานาติ.\csromanspacer{} เวทนา วิภวิสฺสติ.\csromanspacer{} สญฺญา วิภวิสฺสติ.\csromanspacer{} สงฺขารา วิภวิสฺสนฺติ. “วิญฺญาณํ วิภวิสฺสตี”ติ ยถาภูตํ นปฺปชานาติ.}

\prose{สุตวา จ โข ภิกฺขุ อริยสาวโก อริยานํ ทสฺสาวี อริยธมฺมสฺส โกวิโท อริยธมฺเม สุวินีโต, สปฺปุริสานํ ทสฺสาวี สปฺปุริสธมฺมสฺส โกวิโท สปฺปุริสธมฺเม สุวินีโต น รูปํ อตฺตโต สมนุปสฺสติ ฯเปฯ น เวทนํ.\csromanspacer{} น สญฺญํ.\csromanspacer{} น สงฺขาเร.\csromanspacer{} น วิญฺญาณํ อตฺตโต สมนุปสฺสติ.}

\prose{โส อนิจฺจํ รูปํ “อนิจฺจํ รูปนฺ”ติ ยถาภูตํ ปชานาติ.\csromanspacer{} อนิจฺจํ เวทนํ.\csromanspacer{} อนิจฺจํ สญฺญํ.\csromanspacer{} อนิจฺเจ สงฺขาเร.\csromanspacer{} อนิจฺจํ วิญฺญาณํ “อนิจฺจํ วิญฺญาณนฺ”ติ ยถาภูตํ ปชานาติ.\csromanspacer{} ทุกฺขํ รูปํ ฯเปฯ.\csromanspacer{} ทุกฺขํ วิญฺญาณํ.\csromanspacer{} อนตฺตํ รูปํ ฯเปฯ.\csromanspacer{} อนตฺตํ วิญฺญาณํ.\csromanspacer{} สงฺขตํ รูปํ ฯเปฯ.\csromanspacer{} สงฺขตํ วิญฺญาณํ “สงฺขตํ วิญฺญาณนฺ”ติ ยถาภูตํ ปชานาติ. “รูปํ วิภวิสฺสตี”ติ ยถาภูตํ ปชานาติ.\csromanspacer{} เวทนา.\csromanspacer{} สญฺญา.\csromanspacer{} สงฺขารา. “วิญฺญาณํ วิภวิสฺสตี”ติ ยถาภูตํ ปชานาติ.}

\prose{\csromanfolio{47}โส รูปสฺส วิภวา, เวทนาย วิภวา, สญฺญาย วิภวา, สงฺขารานํ วิภวา, วิญฺญาณสฺส วิภวา เอวํ โข ภิกฺขุ “โน จสฺสํ, โน จ เม สิยา, นาภวิสฺส, น เม ภวิสฺสตี”ติ เอวํ อธิมุจฺจมาโน ภิกฺขุ ฉินฺเทยฺย โอรมฺภาคิยานิ สํโยชนานีติ.\csromanspacer{} เอวํ อธิมุจฺจมาโน ภนฺเต ภิกฺขุ ฉินฺเทยฺย โอรมฺภาคิยานิ สํโยชนานีติ.}

\prose{กถํ ปน ภนฺเต ชานโต กถํ ปสฺสโต อนนฺตรา อาสวานํ ขโย โหตีติ.\csromanspacer{} อิธ ภิกฺขุ อสฺสุตวา ปุถุชฺชโน อตสิตาเย ฐาเน ตาสํ อาปชฺชติ, ตาโส เหโส\footnote{เหสา (ก)} ภิกฺขุ อสฺสุตวโต ปุถุชฺชนสฺส โน จสฺสํ, โน จเม สิยา, นาภวิสฺส, น เม ภวิสฺสตีติ.}

\prose{สุตวา จ โข ภิกฺขุ อริยสาวโก อตสิตาเย ฐาเน น ตาสํ อาปชฺชติ.\csromanspacer{} น เหโส\footnote{น เหสา (ก)} ภิกฺขุ ตาโส สุตวโต อริยสาวกสฺส โน จสฺสํ, โน จ เม สิยา, นาภวิสฺส, น เม ภวิสฺสตีติ.\csromanspacer{} รูปุปยํ วา ภิกฺขุ วิญฺญาณํ ติฏฺฐมานํ ติฏฺเฐยฺย, รูปารมฺมณํ รูปปฺปติฏฺฐํ นนฺทูปเสจนํ วุทฺธิํ วิรูฬฺหิํ เวปุลฺลํ อาปชฺเชยฺย.\csromanspacer{} เวทนุปยํ วา ภิกฺขุ.\csromanspacer{} สญฺญุปยํ วา ภิกฺขุ.\csromanspacer{} สงฺขารุปยํ วา ภิกฺขุ วิญฺญาณํ ติฏฺฐมานํ ติฏฺเฐยฺย, สงฺขารารมฺมณํ สงฺขารปฺปติฏฺฐํ นนฺทูปเสจนํ วุทฺธิํ วิรูฬฺหิํ เวปุลฺลํ อาปชฺเชยฺย.}

\prose{โย\footnote{โส (สพฺพตฺถ)} ภิกฺขุ เอวํ วเทยฺย “อหมญฺญตฺร รูปา อญฺญตฺร เวทนาย อญฺญตฺร สญฺญาย อญฺญตฺร สงฺขาเรหิ วิญฺญาณสฺส อาคติํ วา คติํ วา จุติํ วา อุปปตฺติํ วา วุทฺธิํ วา วิรูฬฺหิํ วา เวปุลฺลํ วา ปญฺญาเปสฺสามี”ติ เนตํ ฐานํ วิชฺชติ.}

\prose{รูปธาตุยา เจ ภิกฺขุ ภิกฺขุโน ราโค ปหีโน โหติ, ราคสฺส ปหานา โวจฺฉิชฺชตารมฺมณํ, ปติฏฺฐา วิญฺญาณสฺส น โหติ.\csromanspacer{} เวทนาธาตุยา เจ ภิกฺขุ ภิกฺขุโน.\csromanspacer{} สญฺญาธาตุยา เจ ภิกฺขุ ภิกฺขุโน.\csromanspacer{} สงฺขารธาตุยา เจ ภิกฺขุ ภิกฺขุโน.\csromanspacer{} วิญฺญาณธาตุยา เจ ภิกฺขุ ภิกฺขุโน ราโค ปหีโน โหติ, ราคสฺส ปหานา โวจฺฉิชฺชตารมฺมณํ, ปติฏฺฐา วิญฺญาณสฺส น โหติ, ตทปฺปติฏฺฐิตํ วิญฺญาณํ อวิรูฬฺหํ อนภิสงฺขจฺจ วิมุตฺตํ, วิมุตฺตตฺตา ฐิตํ, ฐิตตฺตา สนฺตุสิตํ, สนฺตุสิตตฺตา น ปริตสฺสติ, อปริตสฺสํ ปจฺจตฺตญฺเญว ปรินิพฺพายติ, “ขีณา ชาติ ฯเปฯ นาปรํ \csromanfolio{48}อิตฺถตฺตายา”ติ ปชานาติ.\csromanspacer{} เอวํ โข ภิกฺขุ ชานโต เอวํ ปสฺสโต อนนฺตรา อาสวานํ ขโย โหตีติ.\csromanspacer{}\csromanspacer{}ตติยํ.}

\csromanheader{2}{4. อุปาทานปริปวตฺตสุตฺต}
\proseitem{56}{สาวตฺถินิทานํ.\csromanspacer{} ปญฺจิเม ภิกฺขเว อุปาทานกฺขนฺธา.\csromanspacer{} กตเม ปญฺจ, รูปุปาทานกฺขนฺโธ เวทนุปาทานกฺขนฺโธ สญฺญุปาทานกฺขนฺโธ สงฺขารุปาทานกฺขนฺโธ วิญฺญาณุปาทานกฺขนฺโธ.\csromanspacer{} ยาวกีวญฺจาหํ ภิกฺขเว อิเม ปญฺจุปาทานกฺขนฺเธ จตุปริวฏฺฏํ ยถาภูตํ นาพฺภญฺญาสิํ, เนว ตาวาหํ ภิกฺขเว สเทวเก โลเก สมารเก สพฺรหฺมเก สสฺสมณพฺราหฺมณิยา ปชาย สเทวมนุสฺสาย “อนุตฺตรํ สมฺมาสมฺโพธิํ อภิสมฺพุทฺโธ”ติ ปจฺจญฺญาสิํ.\csromanspacer{} ยโต จ ขฺวาหํ ภิกฺขเว อิเม ปญฺจุปาทานกฺขนฺเธ จตุปริวฏฺฏํ ยถาภูตํ อพฺภญฺญาสิํ อถาหํ ภิกฺขเว สเทวเก โลเก ฯเปฯ สเทวมนุสฺสาย “อนุตฺตรํ สมฺมาสมฺโพธิํ อภิสมฺภุทฺโธ”ติ ปจฺจญฺญาสิํ.}

\prose{กถญฺจ จตุปริวฏฺฏํ รูปํ อพฺภญฺญาสิํ, รูปสมุทยํ อพฺภญฺญาสิํ, รูปนิโรธํ อพฺภญฺญาสิํ, รูปนิโรธคามินิํ ปฏิปทํ อพฺภญฺญาสิํ.\csromanspacer{} เวทนํ.\csromanspacer{} สญฺญํ.\csromanspacer{} สงฺขาเร.\csromanspacer{} วิญฺญาณํ อพฺภญฺญาสิํ, วิญฺญาณสมุทยํ อพฺภญฺญาสิํ, วิญฺญาณนิโรธํ อพฺภญฺญาสิํ, วิญฺญาณนิโรธคามินิํ ปฏิปทํ อพฺภญฺญาสิํ.}

\prose{กตมญฺจ ภิกฺขเว รูปํ, จตฺตาโร จ มหาภูตา จตุนฺนญฺจ มหาภูตานํ อุปาทายรูปํ.\csromanspacer{} อิทํ วุจฺจติ ภิกฺขเว รูปํ.\csromanspacer{} อาหารสมุทยา รูปสมุทโย, อาหารนิโรธา รูปนิโรโธ, อยเมว อริโย อฏฺฐงฺคิโก มคฺโค รูปนิโรธคามินี ปฏิปทา.\csromanspacer{} เสยฺยถิทํ, สมฺมาทิฏฺฐิ ฯเปฯ สมฺมาสมาธิ.}

\prose{เย หิ เกจิ ภิกฺขเว สมณา วา พฺราหฺมณา วา เอวํ รูปํ อภิญฺญาย, เอวํ รูปสมุทยํ อภิญฺญาย, เอวํ รูปนิโรธํ อภิญฺญาย, เอวํ รูปนิโรธคามินิํ ปฏิปทํ อภิญฺญาย รูปสฺส นิพฺพิทาย วิราคาย นิโรธาย ปฏิปนฺนา, เต สุปฺปฏิปนฺนา.\csromanspacer{} เย สุปฺปฏิปนฺนา, เต อิมสฺมิํ ธมฺมวินเย คาธนฺติ.}

\prose{เย จ โข เกจิ ภิกฺขเว สมณา วา พฺราหฺมณา วา เอวํ รูปํ อภิญฺญาย ฯเปฯ เอวํ รูปนิโรธคามินิํ ปฏิปทํ อภิญฺญาย รูปสฺส นิพฺพิทา วิราคา นิโรธา อนุปาทา วิมุตฺตา, เต สุวิมุตฺตา.\csromanspacer{} เย สุวิมุตฺตา, เต เกวลิโน.\csromanspacer{} เย เกวลิโน, วฏฺฏํ เตสํ นตฺถิ ปญฺญาปนาย.}

\prose{\csromanfolio{49}กตมา จ ภิกฺขเว เวทนา.\csromanspacer{} ฉยิเม ภิกฺขเว เวทนากายา, จกฺขุสมฺผสฺสชา เวทนา โสตสมฺผสฺสชา เวทนา ฆานสมฺผสฺสชา เวทนา ชิวฺหาสมฺผสฺสชา เวทนา กายสมฺผสฺสชา เวทนา มโนสมฺผสฺสชา เวทนา.\csromanspacer{} อยํ วุจฺจติ ภิกฺขเว เวทนา.\csromanspacer{} ผสฺสสมุทยา เวทนาสมุทโย, ผสฺสนิโรธา เวทนานิโรโธ.\csromanspacer{} อยเมว อริโย อฏฺฐงฺคิโก มคฺโค เวทนานิโรธคามินี ปฏิปทา.\csromanspacer{} เสยฺยถิทํ, สมฺมาทิฏฺฐิ ฯเปฯ สมฺมาสมาธิ.}

\prose{เย หิ เกจิ ภิกฺขเว สมณา วา พฺราหฺมณา วา เอวํ เวทนํ อภิญฺญาย, เอวํ เวทนาสมุทยํ อภิญฺญาย, เอวํ เวทนานิโรธํ อภิญฺญาย, เอวํ เวทนานิโรธคามินิํ ปฏิปทํ อภิญฺญาย เวทนาย นิพฺพิทาย วิราคาย นิโรธาย ปฏิปนฺนา, เต สุปฺปฏิปนฺนา.\csromanspacer{} เย สุปฺปฏิปนฺนา, เต อิมสฺมิํ ธมฺมวินเย คาธนฺติ.}

\prose{เย จ โข เกจิ ภิกฺขเว สมณา วา พฺราหฺมณา วา เอวํ เวทนํ อภิญฺญาย ฯเปฯ เอวํ เวทนานิโรธคามินิํ ปฏิปทํ อภิญฺญาย ฯเปฯ วฏฺฏํ เตสํ นตฺถิ ปญฺญาปนาย.}

\prose{กตมา จ ภิกฺขเว สญฺญา.\csromanspacer{} ฉยิเม ภิกฺขเว สญฺญากายา, รูปสญฺญา สทฺทสญฺญา คนฺธสญฺญา รสสญฺญา โผฏฺฐพฺพสญฺญา ธมฺมสญฺญา.\csromanspacer{} อยํ วุจฺจติ ภิกฺขเว สญฺญา.\csromanspacer{} ผสฺสสมุทยา สญฺญาสมุทโย, ผสฺสนิโรธา สญฺญานิโรโธ.\csromanspacer{} อยเมว อริโย อฏฺฐงฺคิโก มคฺโค สญฺญานิโรธคามินี ปฏิปทา.\csromanspacer{} เสยฺยถิทํ, สมฺมาทิฏฺฐิ ฯเปฯ สมฺมาสมาธิ ฯเปฯ วฏฺฏํ เตสํ นตฺถิ ปญฺญาปนาย.}

\prose{กตเม จ ภิกฺขเว สงฺขารา.\csromanspacer{} ฉยิเม ภิกฺขเว เจตนากายา, รูปสญฺเจตนา สทฺทสญฺเจตนา คนฺธสญฺเจตนา รสสญฺเจตนา โผฏฺฐพฺพสญฺเจตนา ธมฺมสญฺเจตนา.\csromanspacer{} อิเม วุจฺจนฺติ ภิกฺขเว สงฺขารา.\csromanspacer{} ผสฺสสมุทยา สงฺขารสมุทโย, ผสฺสนิโรธา สงฺขารนิโรโธ.\csromanspacer{} อยเมว อริโย อฏฺฐงฺคิโก มคฺโค สงฺขารนิโรธคามินี ปฏิปทา.\csromanspacer{} เสยฺยถิทํ, สมฺมาทิฏฺฐิ ฯเปฯ สมฺมาสมาธิ.}

\prose{เย หิ เกจิ ภิกฺขเว สมณา วา พฺราหฺมณา วา เอวํ สงฺขาเร อภิญฺญาย, เอวํ สงฺขารสมุทยํ อภิญฺญาย, เอวํ สงฺขารนิโรธํ อภิญฺญาย, เอวํ สงฺขารนิโรธคามินิํ ปฏิปทํ อภิญฺญาย สงฺขารานํ นิพฺพิทาย วิราคาย \csromanfolio{50}นิโรธาย ปฏิปนฺนา, เต สุปฺปฏิปนฺนา.\csromanspacer{} เย สุปฺปฏิปนฺนา, เต อิมสฺมิํ ธมฺมวินเย คาธนฺติ.}

\prose{เย จ โข เกจิ ภิกฺขเว สมณา วา พฺราหฺมณา วา เอวํ สงฺขาเร อภิญฺญาย, เอวํ สงฺขารสมุทยํ อภิญฺญาย, เอวํ สงฺขารนิโรธํ อภิญฺญาย, เอวํ สงฺขารนิโรธคามินิํ ปฏิปทํ อภิญฺญาย สงฺขารานํ นิพฺพิทา วิราคา นิโรธา อนุปาทา วิมุตฺตา, เต สุวิมุตฺตา.\csromanspacer{} เย สุวิมุตฺตา, เต เกวลิโน.\csromanspacer{} เย เกวลิโน, วฏฺฏํ เตสํ นตฺถิ ปญฺญาปนาย.}

\prose{กตมญฺจ ภิกฺขเว วิญฺญาณํ.\csromanspacer{} ฉยิเม ภิกฺขเว วิญฺญาณกายา, จกฺขุวิญฺญาณํ โสตวิญฺญาณํ ฆานวิญฺญาณํ ชิวฺหาวิญฺญาณํ กายวิญฺญาณํ มโนวิญฺญาณํ.\csromanspacer{} อิทํ วุจฺจติ ภิกฺขเว วิญฺญาณํ.\csromanspacer{} นามรูปสมุทยา วิญฺญาณสมุทโย, นามรูปนิโรธา วิญฺญาณนิโรโธ.\csromanspacer{} อยเมว อริโย อฏฺฐงฺคิโก มคฺโค วิญฺญาณนิโรธคามินี ปฏิปทา.\csromanspacer{} เสยฺยถิทํ, สมฺมาทิฏฺฐิ ฯเปฯ สมฺมาสมาธิ.}

\prose{เย หิ เกจิ ภิกฺขเว สมณา วา พฺราหฺมณา วา เอวํ วิญฺญาณํ อภิญฺญาย, เอวํ วิญฺญาณสมุทยํ อภิญฺญาย, เอวํ วิญฺญาณนิโรธํ อภิญฺญาย, เอวํ วิญฺญาณนิโรธคามินิํ ปฏิปทํ อภิญฺญาย วิญฺญาณสฺส นิพฺพิทาย วิราคาย นิโรธาย ปฏิปนฺนา, เต สุปฺปฏิปนฺนา.\csromanspacer{} เย สุปฺปฏิปนฺนา, เต อิมสฺมิํ ธมฺมวินเย คาธนฺติ.}

\prose{เย จ โข เกจิ ภิกฺขเว สมณา วา พฺราหฺมณา วา เอวํ วิญฺญาณํ อภิญฺญาย, เอวํ วิญฺญาณสมุทยํ อภิญฺญาย, เอวํ วิญฺญาณนิโรธํ อภิญฺญาย, เอวํ วิญฺญาณนิโรธคามินิํ ปฏิปทํ อภิญฺญาย วิญฺญาณสฺส นิพฺพิทา วิราคา นิโรธา อนุปาทา วิมุตฺตา, เต สุวิมุตฺตา.\csromanspacer{} เย สุวิมุตฺตา, เต เกวลิโน.\csromanspacer{} เย เกวลิโน, วฏฺฏํ เตสํ นตฺถิ ปญฺญาปนายาติ.\csromanspacer{}\csromanspacer{}จตุตฺถํ.}

\csromanheader{2}{5. สตฺตฏฺฐานสุตฺต}
\proseitem{57}{สาวตฺถินิทานํ.\csromanspacer{} สตฺตฏฺฐานกุสโล ภิกฺขเว ภิกฺขุ ติวิธูปปริกฺขี อิมสฺมิํ ธมฺมวินเย เกวลี วุสิตวา อุตฺตมปุริโสติ วุจฺจติ.\csromanspacer{} กถญฺจ ภิกฺขเว ภิกฺขุ สตฺตฏฺฐานกุสโล โหติ, อิธ ภิกฺขเว ภิกฺขุ รูปํ ปชานาติ, รูปสมุทยํ ปชานาติ, รูปนิโรธํ ปชานาติ, \csromanfolio{51}รูปนิโรธคามินิํ ปฏิปทํ ปชานาติ.\csromanspacer{} รูปสฺส อสฺสาทํ ปชานาติ, รูปสฺส อาทีนวํ ปชานาติ, รูปสฺส นิสฺสรณํ ปชานาติ.\csromanspacer{} เวทนํ ปชานาติ.\csromanspacer{} สญฺญํ.\csromanspacer{} สงฺขาเร.\csromanspacer{} วิญฺญาณํ ปชานาติ, วิญฺญาณสมุทยํ ปชานาติ, วิญฺญาณนิโรธํ ปชานาติ, วิญฺญาณนิโรธคามินิํ ปฏิปทํ ปชานาติ.\csromanspacer{} วิญฺญาณสฺส อสฺสาทํ ปชานาติ, วิญฺญาณสฺส อาทีนวํ ปชานาติ, วิญฺญาณสฺส นิสฺสรณํ ปชานาติ.}

\prose{กตมญฺจ ภิกฺขเว รูปํ.\csromanspacer{} จตฺตาโร จ มหาภูตา จตุนฺนญฺจ มหาภูตานํ อุปาทายรูปํ.\csromanspacer{} อิทํ วุจฺจติ ภิกฺขเว รูปํ.\csromanspacer{} อาหารสมุทยา รูปสมุทโย, อาหารนิโรธา รูปนิโรโธ, อยเมว อริโย อฏฺฐงฺคิโก มคฺโค รูปนิโรธคามินี ปฏิปทา.\csromanspacer{} เสยฺยถิทํ, สมฺมาทิฏฺฐิ ฯเปฯ สมฺมาสมาธิ.}

\prose{ยํ รูปํ ปฏิจฺจ อุปฺปชฺชติ สุขํ โสมนสฺสํ, อยํ รูปสฺส อสฺสาโท.\csromanspacer{} ยํ รูปํ อนิจฺจํ ทุกฺขํ วิปริณามธมฺมํ, อยํ รูปสฺส อาทีนโว.\csromanspacer{} โย รูปสฺมิํ ฉนฺทราควินโย ฉนฺทราคปฺปหานํ, อิทํ รูปสฺส นิสฺสรณํ.}

\prose{เย หิ เกจิ ภิกฺขเว สมณา วา พฺราหฺมณา วา เอวํ รูปํ อภิญฺญาย, เอวํ รูปสมุทยํ อภิญฺญาย, เอวํ รูปนิโรธํ อภิญฺญาย, เอวํ รูปนิโรธคามินิํ ปฏิปทํ อภิญฺญาย, เอวํ รูปสฺส อสฺสาทํ อภิญฺญาย เอวํ รูปสฺส อาทีนวํ อภิญฺญาย, เอวํ รูปสฺส นิสฺสรณํ อภิญฺญาย รูปสฺส นิพฺพิทาย วิราคาย นิโรธาย ปฏิปนฺนา, เต สุปฺปฏิปนฺนา.\csromanspacer{} เย สุปฺปฏิปนฺนา, เต อิมสฺมิํ ธมฺมวินเย คาธนฺติ.}

\prose{เย จ โข เกจิ ภิกฺขเว สมณา วา พฺราหฺมณา วา เอวํ รูปํ อภิญฺญาย, เอวํ รูปสมุทยํ อภิญฺญาย, เอวํ รูปนิโรธํ อภิญฺญาย, เอวํ รูปนิโรธคามินิํ ปฏิปทํ อภิญฺญาย.\csromanspacer{} เอวํ รูปสฺส อสฺสาทํ อภิญฺญาย, เอวํ รูปสฺส อาทีนวํ อภิญฺญาย, เอวํ รูปสฺส นิสฺสรณํ อภิญฺญาย รูปสฺส นิพฺพิทา วิราคา นิโรธา อนุปาทา วิมุตฺตา, เต สุวิมุตฺตา.\csromanspacer{} เย สุวิมุตฺตา, เต เกวลิโน.\csromanspacer{} เย เกวลิโน, วฏฺฏํ เตสํ นตฺถิ ปญฺญาปนาย.}

\prose{กตมา จ ภิกฺขเว เวทนา.\csromanspacer{} ฉยิเม ภิกฺขเว เวทนากายา, จกฺขุสมฺผสฺสชา เวทนา ฯเปฯ มโนสมฺผสฺสชา เวทนา.\csromanspacer{} อยํ วุจฺจติ ภิกฺขเว เวทนา.\csromanspacer{} ผสฺสสมุทยา เวทนาสมุทโย, ผสฺสนิโรธา เวทนานิโรโธ.\csromanspacer{} อยเมว อริโย อฏฺฐงฺคิโก มคฺโค เวทนานิโรธคามินี ปฏิปทา.\csromanspacer{} เสยฺยถิทํ, สมฺมาทิฏฺฐิ ฯเปฯ สมฺมาสมาธิ.}

\prose{\csromanfolio{52}ยํ เวทนํ ปฏิจฺจ อุปฺปชฺชติ สุขํ โสมนสฺสํ, อยํ เวทนาย อสฺสาโท.\csromanspacer{} ยา เวทนา อนิจฺจา ทุกฺขา วิปริณามธมฺมา, อยํ เวทนาย อาทีนโว.\csromanspacer{} โย เวทนาย ฉนฺทราควินโย ฉนฺทราคปฺปหานํ, อิทํ เวทนาย นิสฺสรณํ.}

\prose{เย หิ เกจิ ภิกฺขเว สมณา วา พฺราหฺมณา วา เอวํ เวทนํ อภิญฺญาย, เอวํ เวทนาสมุทยํ อภิญฺญาย, เอวํ เวทนานิโรธํ อภิญฺญาย, เอวํ เวทนานิโรธคามินิํ ปฏิปทํ อภิญฺญาย, เอวํ เวทนาย อสฺสาทํ อภิญฺญาย, เอวํ เวทนาย อาทีนวํ อภิญฺญาย, เอวํ เวทนาย นิสฺสรณํ อภิญฺญาย เวทนาย นิพฺพิทาย วิราคาย นิโรธาย ปฏิปนฺนา, เต สุปฺปฏิปนฺนา.\csromanspacer{} เย สุปฺปฏิปนฺนา, เต อิมสฺมิํ ธมฺมวินเย คาธนฺติ.}

\prose{เย จ โข เกจิ ภิกฺขเว สมณา วา พฺราหฺมณา วา เอวํ เวทนํ อภิญฺญาย ฯเปฯ วฏฺฏํ เตสํ นตฺถิ ปญฺญาปนาย.}

\prose{กตมา จ ภิกฺขเว สญฺญา.\csromanspacer{} ฉยิเม ภิกฺขเว สญฺญากายา, รูปสญฺญาสทฺทสญฺญา คนฺธสญฺญา รสสญฺญา โผฏฺฐพฺพสญฺญา ธมฺมสญฺญา.\csromanspacer{} อยํ วุจฺจติ ภิกฺขเว สญฺญา.\csromanspacer{} ผสฺสสมุทยา สญฺญาสมุทโย, ผสฺสนิโรธา สญฺญานิโรโธ.\csromanspacer{} อยเมว อริโย อฏฺฐงฺคิโก มคฺโค สญฺญานิโรธคามินี ปฏิปทา.\csromanspacer{} เสยฺยถิทํ, สมฺมาทิฏฺฐิ ฯเปฯ สมฺมาสมาธิ ฯเปฯ วฏฺฏํ เตสํ นตฺถิ ปญฺญาปนาย.}

\prose{กตเม จ ภิกฺขเว สงฺขารา.\csromanspacer{} ฉยิเม ภิกฺขเว เจตนากายา, รูปสญฺเจตนา ฯเปฯ ธมฺมสญฺเจตนา.\csromanspacer{} อิเม วุจฺจนฺติ ภิกฺขเว สงฺขารา.\csromanspacer{} ผสฺสสมุทยา สงฺขารสมุทโย, ผสฺสนิโรธา สงฺขารนิโรโธ.\csromanspacer{} อยเมว อริโย อฏฺฐงฺคิโก มคฺโค สงฺขารนิโรธคามินี ปฏิปทา.\csromanspacer{} เสยฺยถิทํ, สมฺมาทิฏฺฐิ ฯเปฯ สมฺมาสมาธิ.}

\prose{ยํ สงฺขาเร ปฏิจฺจ อุปฺปชฺชติ สุขํ โสมนสฺสํ, อยํ สงฺขารานํ อสฺสาโท.\csromanspacer{} เย สงฺขารา อนิจฺจา ทุกฺขา วิปริณามธมฺมา, อยํ สงฺขารานํ อาทีนโว.\csromanspacer{} โย สงฺขาเรสุ ฉนฺทราควินโย ฉนฺทราคปฺปหานํ, อิทํ สงฺขารานํ นิสฺสรณํ.}

\prose{เย หิ เกจิ ภิกฺขเว สมณา วา พฺราหฺมณา วา เอวํ สงฺขาเร อภิญฺญาย, เอวํ สงฺขารสมุทยํ อภิญฺญาย, เอวํ สงฺขารนิโรธํ อภิญฺญาย, เอวํ สงฺขารนิโรธคามินิํ ปฏิปทํ อภิญฺญาย ฯเปฯ สงฺขารานํ นิพฺพิทาย วิราคาย \csromanfolio{53}นิโรธาย ปฏิปนฺนา, เต สุปฺปฏิปนฺนา.\csromanspacer{} เย สุปฺปฏิปนฺนา, เต อิมสฺมิํ ธมฺมวินเย คาธนฺติ ฯเปฯ วฏฺฏํ เตสํ นตฺถิ ปญฺญาปนาย.}

\prose{กตมญฺจ ภิกฺขเว วิญฺญาณํ.\csromanspacer{} ฉยิเม ภิกฺขเว วิญฺญาณกายา, จกฺขุวิญฺญาณํ โสตวิญฺญาณํ ฆานวิญฺญาณํ ชิวฺหาวิญฺญาณํ กายวิญฺญาณํ มโนวิญฺญาณํ.\csromanspacer{} อิทํ วุจฺจติ ภิกฺขเว วิญฺญาณํ.\csromanspacer{} นามรูปสมุทยา วิญฺญาณสมุทโย, นามรูปนิโรธา วิญฺญาณนิโรโธ.\csromanspacer{} อยเมว อริโย อฏฺฐงฺคิโก มคฺโค วิญฺญาณนิโรธคามินี ปฏิปทา.\csromanspacer{} เสยฺยถิทํ, สมฺมาทิฏฺฐิ ฯเปฯ สมฺมาสมาธิ.}

\prose{ยํ วิญฺญาณํ ปฏิจฺจ อุปฺปชฺชติ สุขํ โสมนสฺสํ, อยํ วิญฺญาณสฺส อสฺสาโท.\csromanspacer{} ยํ วิญฺญาณํ อนิจฺจํ ทุกฺขํ วิปริณามธมฺมํ, อยํ วิญฺญาณสฺส อาทีนโว.\csromanspacer{} โย วิญฺญาณสฺมิํ ฉนฺทราควินโย ฉนฺทราคปฺปหานํ, อิทํ วิญฺญาณสฺส นิสฺสรณํ.}

\prose{เย หิ เกจิ ภิกฺขเว สมณา วา พฺราหฺมณา วา เอวํ วิญฺญาณํ อภิญฺญาย, เอวํ วิญฺญาณสมุทยํ อภิญฺญาย, เอวํ วิญฺญาณนิโรธํ อภิญฺญาย, เอวํ วิญฺญาณนิโรธคามินิํ ปฏิปทํ อภิญฺญาย, เอวํ วิญฺญาณสฺส อสฺสาทํ อภิญฺญาย, เอวํ วิญฺญาณสฺส อาทีนวํ อภิญฺญาย, เอวํ วิญฺญาณสฺส นิสฺสรณํ อภิญฺญาย วิญฺญาณสฺส นิพฺพิทาย วิราคาย นิโรธาย ปฏิปนฺนา, เต สุปฺปฏิปนฺนา.\csromanspacer{} เย สุปฺปฏิปนฺนา, เต อิมสฺมิํ ธมฺมวินเย คาธนฺติ.}

\prose{เย จ โข เกจิ ภิกฺขเว สมณา วา พฺราหฺมณา วา เอวํ วิญฺญาณํ อภิญฺญาย, เอวํ วิญฺญาณสมุทยํ อภิญฺญาย, เอวํ วิญฺญาณนิโรธํ อภิญฺญาย, เอวํ วิญฺญาณนิโรธคามินิํ ปฏิปทํ อภิญฺญาย, เอวํ วิญฺญาณสฺส อสฺสาทํ อภิญฺญาย, เอวํ วิญฺญาณสฺส อาทีนวํ อภิญฺญาย, เอวํ วิญฺญาณสฺส นิสฺสรณํ อภิญฺญาย วิญฺญาณสฺส นิพฺพิทา วิราคา นิโรธา อนุปาทา วิมุตฺตา, เต สุวิมุตฺตา.\csromanspacer{} เย สุวิมุตฺตา, เต เกวลิโน.\csromanspacer{} เย เกวลิโน, วฏฺฏํ เตสํ นตฺถิ ปญฺญาปนาย.\csromanspacer{} เอวํ โข ภิกฺขเว ภิกฺขุ สตฺตฏฺฐานกุสโล โหติ.}

\prose{กถญฺจ ภิกฺขเว ภิกฺขุ ติวิธูปปริกฺขี โหติ อิธ ภิกฺขเว ภิกฺขุ ธาตุโส อุปปริกฺขติ, อายตนโส อุปปริกฺขติ, ปฏิจฺจสมุปฺปาทโส \csromanfolio{54}อุปปริกฺขติ.\csromanspacer{} เอวํ โข ภิกฺขเว ภิกฺขุ ติวิธูปปริกฺขี โหติ.\csromanspacer{} สตฺตฏฺฐานกุสโล ภิกฺขเว ภิกฺขุ ติวิธูปปริกฺขี อิมสฺมิํ ธมฺมวินเย เกวลี วุสิตวา อุตฺตมปุริโสติ วุจฺจตีติ.\csromanspacer{}\csromanspacer{}ปญฺจมํ.}

\csromanheader{2}{6. สมฺมาสมฺพุทฺธสุตฺต}
\proseitem{58}{สาวตฺถินิทานํ.\csromanspacer{} ตถาคโต ภิกฺขเว อรหํ สมฺมาสมฺพุทฺโธ รูปสฺส นิพฺพิทา วิราคา นิโรธา อนุปาทา วิมุตฺโต สมฺมาสมฺพุทฺโธติ วุจฺจติ.\csromanspacer{} ภิกฺขุปิ ภิกฺขเว ปญฺญาวิมุตฺโต รูปสฺส นิพฺพิทา วิราคา นิโรธา อนุปาทา วิมุตฺโต ปญฺญาวิมุตฺโตติ วุจฺจติ.}

\prose{ตถาคโต ภิกฺขเว อรหํ สมฺมาสมฺพุทฺโธ เวทนาย นิพฺพิทา วิราคา นิโรธา อนุปาทา วิมุตฺโต สมฺมาสมฺพุทฺโธติ วุจฺจติ.\csromanspacer{} ภิกฺขุปิ ภิกฺขเว ปญฺญาวิมุตฺโต เวทนาย นิพฺพิทา ฯเปฯ ปญฺญาวิมุตฺโตติ วุจฺจติ.}

\prose{ตถาคโต ภิกฺขเว อรหํ สมฺมาสมฺพุทฺโธ สญฺญาย.\csromanspacer{} สงฺขารานํ.\csromanspacer{} วิญฺญาณสฺส นิพฺพิทา วิราคา นิโรธา อนุปาทา วิมุตฺโต สมฺมาสมฺพุทฺโธติ วุจฺจติ.\csromanspacer{} ภิกฺขุปิ ภิกฺขเว ปญฺญาวิมุตฺโต วิญฺญณสฺส นิพฺพิทา วิราคา นิโรธา อนุปาทา วิมุตฺโต ปญฺญาวิมุตฺโตติ วุจฺจติ.}

\prose{ตตฺร โข ภิกฺขเว โก วิเสโส, โก อธิปฺปยาโส\footnote{อธิปฺปาโย (สี), อธิปฺปายโส (สฺยา, กํ, อิ, ก)}, กิํ นานากรณํ ตถาคตสฺส อรหโต สมฺมาสมฺพุทฺธสฺส ปญฺญาวิมุตฺเตน ภิกฺขุนาติ.\csromanspacer{} ภควํมูลกา โน ภนฺเต ธมฺมา ภควํเนตฺติกา ภควํปฏิสรณา, สาธุ วต ภนฺเต ภควนฺตญฺเญว ปฏิภาตุ เอตสฺส ภาสิตสฺส อตฺโถ, ภควโต สุตฺวา ภิกฺขู ธาเรสฺสนฺตีติ.\csromanspacer{} เตน หิ ภิกฺขเว สุณาถ สาธุกํ มนสิ กโรถ, ภาสิสฺสามีติ. “เอวํ ภนฺเต”ติ โข เต ภิกฺขู ภควโต ปจฺจสฺโสสุํ.\csromanspacer{} ภควา เอตทโวจ\csromandash{}}

\prose{ตถาคโต ภิกฺขเว อรหํ สมฺมาสมฺพุทฺโธ อนุปฺปนฺนสฺส มคฺคสฺส อุปฺปาเทตา, อสญฺชาตสฺส มคฺคสฺส สญฺชเนตา\footnote{สญฺชาเนตา (สฺยา, กํ)}, อนกฺขาตสฺส มคฺคสฺส อกฺขาตา, มคฺคญฺญู มคฺควิทู มคฺคโกวิโท, มคฺคานุคา จ ภิกฺขเว เอตรหิ สาวกา วิหรนฺติ, ปจฺฉาสมนฺนาคตา.\csromanspacer{} อยํ โข ภิกฺขเว วิเสโส, อยํ อธิปฺปยาโส, อิทํ นานากรณํ ตถาคตสฺส อรหโต สมฺมาสมฺพุทฺธสฺส ปญฺญาวิมุตฺเตน ภิกฺขุนาติ.\csromanspacer{}\csromanspacer{}ฉฏฺฐํ.}

\csromanheader{2}{7. อนตฺตลกฺขณสุตฺต}
\proseitem{59}{\csromanfolio{55}เอกํ สมยํ ภควา พาราณสิยํ วิหรติ อิสิปตเน มิคทาเย.\csromanspacer{} ตตฺร โข ภควา ปญฺจวคฺคิเย ภิกฺขู อามนฺเตสิ ภิกฺขโวติ.\csromanspacer{} ภทนฺเตติ เต ภิกฺขู ภควโต ปจฺจสฺโสสุํ.\csromanspacer{} ภควา เอตทโวจ\csromandash{}}

\prose{รูปํ ภิกฺขเว อนตฺตา, รูปญฺจ หิทํ ภิกฺขเว อตฺตา อภวิสฺส, นยิทํ รูปํ อาพาธาย สํวตฺเตยฺย, ลพฺเภถ จ รูเป “เอวํ เม รูปํ โหตุ, เอวํ เม รูปํ มา อโหสี”ติ.\csromanspacer{} ยสฺมา จ โข ภิกฺขเว รูปํ อนตฺตา, ตสฺมา รูปํ อาพาธาย สํวตฺตติ, น จ ลพฺภติ รูเป “เอวํ เม รูปํ โหตุ, เอวํ เม รูปํ มา อโหสี”ติ.}

\prose{เวทนา อนตฺตา, เวทนา จ หิทํ ภิกฺขเว อตฺตา อภวิสฺส, นยิทํ เวทนา อาพาธาย สํวตฺเตยฺย, ลพฺเภถ จ เวทนาย “เอวํ เม เวทนา โหตุ, เอวํ เม เวทนา มา อโหสี”ติ.\csromanspacer{} ยสฺมา จ โข ภิกฺขเว เวทนา อนตฺตา, ตสฺมา เวทนา อาพาธาย สํวตฺตติ, น จ ลพฺภติ เวทนาย “เอวํ เม เวทนา โหตุ, เอวํ เม เวทนา มา อโหสี”ติ.}

\prose{สญฺญา อนตฺตา ฯเปฯ.\csromanspacer{} สงฺขารา อนตฺตา, สงฺขารา จ หิทํ ภิกฺขเว อตฺตา อภวิสฺสํสุ, นยิทํ สงฺขารา อาพาธาย สํวตฺเตยฺยุํ, ลพฺเภถ จ สงฺขาเรสุ “เอวํ เม สงฺขารา โหนฺตุ, เอวํ เม สงฺขารา มา อเหสุนฺ”ติ.\csromanspacer{} ยสฺมา จ โข ภิกฺขเว สงฺขารา อนตฺตา, ตสฺมา สงฺขารา อาพาธาย สํวตฺตนฺติ, น จ ลพฺภติ สงฺขาเรสุ “เอวํ เม สงฺขารา โหนฺตุ, เอวํ เม สงฺขารา มา อเหสุนฺ”ติ.}

\prose{วิญฺญาณํ อนตฺตา, วิญฺญาณญฺจ หิทํ ภิกฺขเว อตฺตา อภวิสฺส, นยิทํ วิญฺญาณํ อาพาธาย สํวตฺเตยฺย, ลพฺเภถ จ วิญฺญาเณ “เอวํ เม วิญฺญาณํ โหตุ, เอวํ เม วิญฺญาณํ มา อโหสี”ติ.\csromanspacer{} ยสฺมา จ โข ภิกฺขเว วิญฺญาณํ อนตฺตา, ตสฺมา วิญฺญาณํ อาพาธาย สํวตฺตติ, น จ ลพฺภติ วิญฺญาเณ “เอวํ เม วิญฺญาณํ โหตุ, เอวํ เม วิญฺญาณํ มา อโหสี”ติ.}

\prose{ตํ กิํ มญฺญถ ภิกฺขเว, รูปํ นิจฺจํ วา อนิจฺจํ วาติ.\csromanspacer{} อนิจฺจํ ภนฺเต.\csromanspacer{} ยํ ปนานิจฺจํ, ทุกฺขํ วา ตํ สุขํ วาติ.\csromanspacer{} ทุกฺขํ ภนฺเต.\csromanspacer{} ยํ ปนานิจฺจํ ทุกฺขํ วิปริณามธมฺมํ, กลฺลํ นุ ตํ สมนุปสฺสิตุํ “เอตํ มม, เอโสหมสฺมิ, เอโส \csromanfolio{56}เม อตฺตา”ติ.\csromanspacer{} โน เหตํ ภนฺเต.\csromanspacer{} เวทนา.\csromanspacer{} สญฺญา.\csromanspacer{} สงฺขารา.\csromanspacer{} วิญฺญาณํ นิจฺจํ วา อนิจฺจํ วาติ.\csromanspacer{} อนิจฺจํ ภนฺเต.\csromanspacer{} ยํ ปนานิจฺจํ, ทุกฺขํ วา ตํ สุขํ วาติ.\csromanspacer{} ทุกฺขํ ภนฺเต.\csromanspacer{} ยํ ปนานิจฺจํ ทุกฺขํ วิปริณามธมฺมํ, กลฺลํ นุ ตํ สมนุปสฺสิตุํ “เอตํ มม, เอโสหมสฺมิ, เอโส เม อตฺตา”ติ.\csromanspacer{} โน เหตํ ภนฺเต.}

\prose{ตสฺมา ติห ภิกฺขเว ยํ กิญฺจิ รูปํ อตีตานาคตปจฺจุปฺปนฺนํ อชฺฌตฺตํ วา พหิทฺธา วา โอฬาริกํ วา สุขุมํ วา หีนํ วา ปณีตํ วา ยํ ทูเร สนฺติเก วา, สพฺพํ รูปํ “เนตํ มม, เนโสหมสฺมิ, น เมโส อตฺตา”ติ เอวเมตํ ยถาภูตํ สมฺมปฺปญฺญาย ทฏฺฐพฺพํ.\csromanspacer{} ยา กาจิ เวทนา อตีตานาคตปจฺจุปฺปนฺนา อชฺฌตฺตํ วา พหิทฺธา วา ฯเปฯ ยา ทูเร สนฺติเก วา, สพฺพา เวทนา “เนตํ มม, เนโสหมสฺมิ, น เมโส อตฺตา”ติ เอวเมตํ ยถาภูตํ สมฺมปฺปญฺญาย ทฏฺฐพฺพํ.}

\prose{ยา กาจิ สญฺญา ฯเปฯ.\csromanspacer{} เย เกจิ สงฺขารา อตีตานาคตปจฺจุปฺปนฺนา อชฺฌตฺตํ วา พหิทฺธา วา ฯเปฯ เย ทูเร สนฺติเก วา, สพฺเพ สงฺขารา “เนตํ มม, เนโสหมสฺมิ, น เมโส อตฺตา”ติ เอวเมตํ ยถาภูตํ สมฺมปฺปญฺญาย ทฏฺฐพฺพํ.}

\prose{ยํ กิญฺจิ วิญฺญาณํ อตีตานาคตปจฺจุปฺปนฺนํ อชฺฌตฺตํ วา พหิทฺธา วา โอฬาริกํ วา สุขุมํ วา หีนํ วา ปณีตํ วา ยํ ทูเร สนฺติเก วา, สพฺพํ วิญฺญาณํ “เนตํ มม, เนโสหมสฺมิ, น เมโส อตฺตา”ติ เอวเมตํ ยถาภูตํ สมฺมปฺปญฺญาย ทฏฺฐพฺพํ.}

\prose{เอวํ ปสฺสํ ภิกฺขเว สุตวา อริยสาวโก รูปสฺมิมฺปิ นิพฺพินฺทติ, เวทนายปิ นิพฺพินฺทติ, สญฺญายปิ นิพฺพินฺทติ, สงฺขาเรสุปิ นิพฺพินฺทติ, วิญฺญาณสฺมิมฺปิ นิพฺพินฺทติ, นิพฺพินฺทํ วิรชฺชติ, วิราคา วิมุจฺจติ, วิมุตฺตสฺมิํ “วิมุตฺตมฺ”อิติ ญาณํ โหติ, “ขีณา ชาติ, วุสิตํ พฺราหฺมจริยํ, กตํ กรณียํ, นาปรํ อิตฺถตฺตายา”ติ ปชานาตีติ.}

\prose{อิทมโวจ ภควา.\csromanspacer{} อตฺตมนา ปญฺจวคฺคิยา ภิกฺขู ภควโต ภาสิตํ อภินนฺทุํ\footnote{อภินนฺทุนฺติ (ก)}.}

\prose{อิมสฺมิํ จ ปน เวยฺยากรณสฺมิํ ภญฺญมาเน ปญฺจวคฺคิยานํ ภิกฺขูนํ อนุปาทาย อาสเวหิ จิตฺตานิ วิมุจฺจิํสูติ.\csromanspacer{}\csromanspacer{}สตฺตมํ.}

\csromanheader{2}{8. มหาลิสุตฺต}
\proseitem{60}{\csromanfolio{57}เอวํ เม สุตํ\csromandash{}เอกํ สมยํ ภควา เวสาลิยํ วิหรติ มหาวเน กูฏาคารสาลายํ.\csromanspacer{} อถ โข มหาลิ ลิจฺฉวิ เยน ภควา เตนุปสงฺกมิ ฯเปฯ เอกมนฺตํ นิสินฺโน โข มหาลิ ลิจฺฉวิ ภควนฺตํ เอตทโวจ\csromandash{}}

\prose{ปูรโณ ภนฺเต กสฺสโป เอวมาห “นตฺถิ เหตุ นตฺถิ ปจฺจโย สตฺตานํ สํกิเลสาย, อเหตู อปฺปจฺจยา สตฺตา สํกิลิสฺสนฺติ.\csromanspacer{} นตฺถิ เหตุ นตฺถิ ปจฺจโย สตฺตานํ วิสุทฺธิยา, อเหตู อปฺปจฺจยา สตฺตา วิสุชฺฌนฺตี”ติ, อิธ ภควา กิมาหาติ.}

\prose{อตฺถิ มหาลิ เหตุ อตฺถิ ปจฺจโย สตฺตานํ สํกิเลสาย, สเหตู สปฺปจฺจยา สตฺตา สํกิลิสฺสนฺติ.\csromanspacer{} อตฺถิ มหาลิ เหตุ อตฺถิ ปจฺจโย สตฺตานํ วิสุทฺธิยา, สเหตู สปฺปจฺจยา สตฺตา วิสุชฺฌนฺตีติ.}

\prose{กตโม ปน ภนฺเต เหตุ กตโม ปจฺจโย สตฺตานํ สํกิเลสาย, กถํ สเหตู สปฺปจฺจยา สตฺตา สํกิลิสฺสนฺตีติ.}

\prose{รูปญฺจ หิทํ มหาลิ เอกนฺตทุกฺขํ อภวิสฺส ทุกฺขานุปติตํ ทุกฺขาวกฺกนฺตํ อนวกฺกนฺตํ สุเขน, นยิทํ สตฺตา รูปสฺมิํ สารชฺเชยฺยุํ.\csromanspacer{} ยสฺมา จ โข มหาลิ รูปํ สุขํ สุขานุปติตํ สุขาวกฺกนฺตํ อนวกฺกนฺตํ ทุกฺเขน, ตสฺมา สตฺตา รูปสฺมิํ สารชฺชนฺติ, สาราคา สํยุชฺชนฺติ, สํโยคา สํกิลิสฺสนฺติ.\csromanspacer{} อยํ โข มหาลิ เหตุ อยํ ปจฺจโย สตฺตานํ สํกิเลสาย, เอวํ สเหตู สปฺปจฺจยา สตฺตา สํกิลิสฺสนฺติ.}

\prose{เวทนา จ หิทํ มหาลิ เอกนฺตทุกฺขา อภวิสฺส ทุกฺขานุปติตา ทุกฺขาวกฺกานฺตา อนวกฺกนฺตา สุเขน, นยิทํ สตฺตา เวทนาย สารชฺเชยฺยุํ.\csromanspacer{} ยสฺมา จ โข มหาลิ เวทนา สุขา สุขานุปติตา สุขาวกฺกนฺตา อนวกฺกนฺตา ทุกฺเขน, ตสฺมา สตฺตา เวทนาย สารชฺชนฺติ, สาราคา สํยุชฺชนฺติ, สํโยคา สํกิลิสฺสนฺติ.\csromanspacer{} อยมฺปิ โข มหาลิ เหตุ อยํ ปจฺจโย สตฺตานํ สํกิเลสาย, เอวมฺปิ สเหตู สปฺปจฺจยา สตฺตา สํกิลิสฺสนฺติ.}

\prose{สญฺญา จ หิทํ มหาลิ ฯเปฯ.\csromanspacer{} สงฺขารา จ หิทํ มหาลิ เอกนฺตทุกฺขา อภวิสฺสํสุ ทุกฺขานุปติตา ทุกฺขาวกฺกนฺตา อนวกฺกนฺตา สุเขน, นยิทํ \csromanfolio{58}สตฺตา สงฺขาเรสุ สารชฺเชยฺยุํ.\csromanspacer{} ยสฺมา จ โข มหาลิ สงฺขารา สุขา สุขานุปติตา สุขาวกฺกนฺตา อนวกฺกนฺตา ทุกฺเขน, ตสฺมา สตฺตา สงฺขาเรสุ สารชฺชนฺติ, สาราคา สํยุชฺชนฺติ, สํโยคา สํกิลิสฺสนฺติ.\csromanspacer{} อยมฺปิ โข มหาลิ เหตุ อยํ ปจฺจโย สตฺตานํ สํกิเลสาย, เอวมฺปิ สเหตู สปฺปจฺจยา สตฺตา สํกิลิสฺสนฺติ.}

\prose{วิญฺญาณญฺจ หิทํ มหาลิ เอกนฺตทุกฺขํ อภวิสฺส ทุกฺขานุปติตํ ทุกฺขาวกฺขนฺตํ อนวกฺกนฺตํ สุเขน, นยิทํ สตฺตา วิญฺญาณสฺมิํ สารชฺเชยฺยุํ.\csromanspacer{} ยสฺมา จ โข มหาลิ วิญฺญาณํ สุขํ สุขานุปติตํ สุขาวกฺกนฺตํ อนวกฺกนฺตํ ทุกฺเขน, ตสฺมา สตฺตา วิญฺญาณสฺมิํ สารชฺชนฺติ, สาราคาสํยุชฺชนฺติ, สํโยคา สํกิลิสฺสนฺติ.\csromanspacer{} อยมฺปิ โข มหาลิ เหตุ อยํ ปจฺจโย สตฺตานํ สํกิเลสาย, เอวมฺปิ สเหตู สปฺปจฺจยา สตฺตา สํกิลิสฺสนฺตีติ.}

\prose{กตโม ปน ภนฺเต เหตุ กตโม ปจฺจโย สตฺตานํ วิสุทฺธิยา, กถํ สเหตู สปฺปจฺจยา สตฺตา วิสุชฺฌนฺตีติ.\csromanspacer{} รูปญฺจ หิทํ มหาลิ เอกนฺตสุขํ อภวิสฺส สุขานุปติตํ สุขาวกฺกนฺตํ อนวกฺกนฺตํ ทุกฺเขน, นยิทํ สตฺตา รูปสฺมิํ นิพฺพินฺเทยฺยุํ.\csromanspacer{} ยสฺมา จ โข มหาลิ รูปํ ทุกฺขํ ทุกฺขานุปติตํ ทุกฺขาวกฺกนฺตํ อนวกฺกนฺตํ สุเขน, ตสฺมา สตฺตา รูปสฺมิํ นิพฺพินฺทนฺติ, นพฺพินฺทํ วิรชฺชนฺติ, วิราคา วิสุชฺฌนฺติ.\csromanspacer{} อยํ โข มหาลิ เหตุ อยํ ปจฺจโย สตฺตานํ วิสุทฺธิยา.\csromanspacer{} เอวํ สเหตู สปฺปจฺจยา สตฺตา วิสุชฺฌนฺติ.}

\prose{เวทนา จ หิทํ มหาลิ เอกนฺตสุขา อภวิสฺส ฯเปฯ.\csromanspacer{} สญฺญา จ หิทํ มหาลิ ฯเปฯ.\csromanspacer{} สงฺขารา จ หิทํ มหาลิ เอกนฺตสุขา อภวิสฺสํสุ ฯเปฯ.\csromanspacer{} วิญฺญาณญฺจ หิทํ มหาลิ เอกนฺตสุขํ อภวิสฺส สุขานุปติตํ สุขาวกฺกนฺตํ อนวกฺกนฺตํ ทุกฺเขน, นยิทํ สตฺตา วิญฺญาณสฺมิํ นิพฺพินฺเทยฺยุํ.\csromanspacer{} ยสฺมา จ โข มหาลิ วิญฺญาณํ ทุกฺขํ ทุกฺขานุปติตํ ทุกฺขาวกฺกนฺตํ อนวกฺกนฺตํ สุเขน, ตสฺมา สตฺตา วิญฺญาณสฺมิํ นิพฺพินฺทนฺติ, นิพฺพินฺทํ วิรชฺชนฺติ, วิราคา วิสุชฺฌนฺติ.\csromanspacer{} อยํ โข มหาลิ เหตุ อยํ ปจฺจโย สตฺตานํ วิสุทฺธิยา.\csromanspacer{} เอวมฺปิ สเหตู สปฺปจฺจยา สตฺตา วิสุชฺฌนฺตีติ.\csromanspacer{}\csromanspacer{}อฏฺฐมํ.}

\csromanheader{2}{9. อาทิตฺตสุตฺต}
\proseitem{61}{สาวตฺถินิทานํ.\csromanspacer{} รูปํ ภิกฺขเว อาทิตฺตํ, เวทนา อาทิตฺตา, สญฺญา อาทิตฺตา, สงฺขารา อาทิตฺตา, วิญฺญาณํ อาทิตฺตํ.\csromanspacer{} เอวํ ปสฺสํ ภิกฺขเว สุตวา อริยสาวโก รูปสฺมิมฺปิ นิพฺพินฺทติ, เวทนายปิ.\csromanspacer{} สญฺญายปิ.\csromanspacer{} สงฺขาเรสุปิ. \csromanfolio{59}วิญฺญาณสฺมิมฺปิ นิพฺพินฺทติ, นิพฺพินฺทํ วิรชฺชติ, วิราคา วิมุจฺจติ, วิมุตฺตสฺมิํ “วิมุตฺตมฺ”อิติ ญาณํ โหติ, “ขีณา ชาติ, วุสิตํ พฺราหฺมจริยํ, กตํ กรณียํ, นาปรํ อิตฺถตฺตายา”ติ ปชานาตีติ.\csromanspacer{}\csromanspacer{}นวมํ.}

\csromanheader{2}{10. นิรุตฺติปถสุตฺต}
\proseitem{62}{สาวตฺถินิทานํ.\csromanspacer{} ตโยเม ภิกฺขเว นิรุตฺติปถา อธิวจนปถา ปญฺญตฺติปถา อสํกิณฺณา อสํกิณฺณปุพฺพา น สํกียนฺติ น สํกียิสฺสนฺติ อปฺปฏิกุฏฺฐา สมเณหิ พฺราหฺมเณหิ วิญฺญูหิ.\csromanspacer{} กตเม ตโย, ยํ ภิกฺขเว รูปํ อตีตํ นิรุทฺธํ วิปริณตํ, “อโหสี”ติ ตสฺส สงฺขา, “อโหสี”ติ ตสฺส สมญฺญา, “อโหสี”ติ ตสฺส ปญฺญตฺติ, น ตสฺส สงฺขา “อตฺถี”ติ, น ตสฺส สงฺขา “ภวิสฺสตี”ติ.}

\prose{ยา เวทนา อตีตา นิรุทฺธา วิปริณตา, “อโหสี”ติ ตสฺสา สงฺขา, “อโหสี”ติ ตสฺสา สมญฺญา, “อโหสี”ติ ตสฺสา ปญฺญตฺติ, น ตสฺสา สงฺขา “อตฺถี”ติ, น ตสฺสา สงฺขา “ภวิสฺสตี”ติ.}

\prose{ยา สญฺญา.\csromanspacer{} เย สงฺขารา อตีตา นิรุทฺธา วิปริณตา, “อเหสุนฺ”ติ เตสํ สงฺขา, “อเหสุนฺ”ติ เตสํ สมญฺญา, “อเหสุนฺ”ติ เตสํ ปญฺญตฺติ, น เตสํ สงฺขา “อตฺถี”ติ, น เตสํ สงฺขา “ภวิสฺสนฺตี”ติ.}

\prose{ยํ วิญฺญาณํ อตีตํ นิรุทฺธํ วิปริณตํ, “อโหสี”ติ ตสฺส สงฺขา, “อโหสี”ติ ตสฺส สมญฺญา, “อโหสี”ติ ตสฺส ปญฺญตฺติ, น ตสฺส สงฺขา “อตฺถี”ติ, น ตสฺส สงฺขา “ภวิสฺสตี”ติ.}

\prose{ยํ ภิกฺขเว รูปํ อชาตํ อปาตุภูตํ, “ภวิสฺสตี”ติ ตสฺส สงฺขา, “ภวิสฺสตี”ติ ตสฺส สมญฺญา, “ภวิสฺสตี”ติ ตสฺส ปญฺญตฺติ, น ตสฺส สงฺขา “อตฺถี”ติ, น ตสฺส สงฺขา “อโหสี”ติ.}

\prose{ยา เวทนา อชาตา อปาตุภูตา, “ภวิสฺสตี”ติ ตสฺสา สงฺขา, “ภวิสฺสตี”ติ ตสฺสา สมญฺญา, ภวิสฺสตีติ ตสฺสา ปญฺญตฺติ, น ตสฺสา สงฺขา “อตฺถี”ติ, น ตสฺสา สงฺขา “อโหสี”ติ.}

\prose{ยา สญฺญา.\csromanspacer{} เย สงฺขารา อชาตา อปาตุภูตา, “ภวิสฺสนฺตี”ติ เตสํ สงฺขา, “ภวิสฺสนฺตี”ติ เตสํ สมญฺญา, “ภวิสฺสนฺตี”ติ เตสํ ปญฺญตฺติ, น เตสํ สงฺขา “อตฺถี”ติ, นเตสํ สงฺขา “อเหสุนฺ”ติ.}

\prose{\csromanfolio{60}ยํ วิญฺญาณํ อชาตํ อปาตุภูตํ, “ภวิสฺสตี”ติ ตสฺส สงฺขา, “ภวิสฺสตี”ติ ตสฺส สมญฺญา, “ภวิสฺสตี”ติ ตสฺส ปญฺญตฺติ, น ตสฺส สงฺขา “อตฺถี”ติ, น ตสฺส สงฺขา “อโหสี”ติ.}

\prose{ยํ ภิกฺขเว รูปํ ชาตํ ปาตุภูตํ, “อตฺถี”ติ ตสฺส สงฺขา, “อตฺถี”ติ ตสฺส สมญฺญา, “อตฺถี”ติ ตสฺส ปญฺญตฺติ, น ตสฺส สงฺขา “อโหสี”ติ, น ตสฺส สงฺขา “ภวิสฺสตี”ติ.}

\prose{ยา เวทนา ชาตา ปาตุภูตา, “อตฺถี”ติ ตสฺสา สงฺขา, “อตฺถี”ติ ตสฺสา สมญฺญา, “อตฺถี”ติ ตสฺสา ปญฺญตฺติ, น ตสฺสา สงฺขา “อโหสี”ติ, น ตสฺสา สงฺขา “ภวิสฺสตี”ติ.}

\prose{ยา สญฺญา.\csromanspacer{} เย สงฺขารา ชาตา ปาตุภูตา, “อตฺถี”ติ เตสํ สงฺขา, “อตฺถี”ติ เตสํ สมญฺญา, “อตฺถี”ติ เตสํ ปญฺญตฺติ, น เตสํ สงฺขา “อเหสุนฺ”ติ, น เตสํ สงฺขา “ภวิสฺสนฺตี”ติ.}

\prose{ยํ วิญฺญาณํ ชาตํ ปาตุภูตํ, “อตฺถี”ติ ตสฺส สงฺขา, “อตฺถี”ติ ตสฺส สมญฺญา, “อตฺถี”ติ ตสฺส ปญฺญตฺติ, น ตสฺส สงฺขา “อโหสี”ติ, น ตสฺส สงฺขา “ภวิสฺสตี”ติ.}

\prose{อิเม โข ภิกฺขเว ตโย นิรุตฺติปถา อธิวจนปถา ปญฺญตฺติปถา อสํกิณฺณา อสํกิณฺณปุพฺพา น สํกียนฺติ น สํกียิสฺสนฺติ อปฺปฏิกุฏฺฐา สมเณหิ พฺราหฺมเณหิ วิญฺญูหิ.\csromanspacer{} เยปิ เต ภิกฺขเว อเหสุํ อุกฺกลา วสฺสภญฺญา\footnote{โอกฺกลา วยภิญฺญา (ม 3. 121 ปิฏฺเฐ)} อเหตุกวาทา อกิริยวาทา นตฺถิกวาทา, เตปิเม ตโย นิรุตฺติปเถ อธิวจนปเถ ปญฺญตฺติปเถ น ครหิตพฺพํ นปฺปฏิกฺโกสิตพฺพํ อมญฺญิํสุ.\csromanspacer{} ตํ กิสฺส เหตุ, นินฺทาฆฏฺฏนพฺยาโรส-อุปารมฺภภยาติ\footnote{นินฺทาพฺยาโรส-อุปารมฺภภยาติ (สี, สฺยา, กํ, อิ) ม 3. 122 ปิฏฺเฐปิ.}.\csromanspacer{} มชฺฌิมปณฺณาสกสฺส อุปยวคฺโค ปฐโม.}

\titlehead{\textbf{ตสฺสุทฺทานํ}}
\csromangathameasure{อุปโย พีชํ อุทานํ, อุปาทานปริวตฺตํ. \\ สตฺตฏฺฐานํ จ สมฺพุทฺโธ, มญฺจมหาลิ อาทิตฺตา. \\ วคฺโค นิรุตฺติปเถน จาติ.}
\csromangathagroup{\csromangathastanza{อุปโย พีชํ อุทานํ, อุปาทานปริวตฺตํ. \\ สตฺตฏฺฐานํ จ สมฺพุทฺโธ, มญฺจมหาลิ อาทิตฺตา. \\ วคฺโค นิรุตฺติปเถน จาติ.}}

\chapterheadpageread{\textbf{(7) 2. อรหนฺตวคฺค}}
\csromanheader{2}{1. อุปาทิยมานสุตฺต}
\proseitem{63}{\csromanfolio{61}เอวํ เม สุตํ\csromandash{}เอกํ สมยํ ภควา สาวตฺถิยํ วิหรติ เชตวเน อนาถปิณฺฑิกสฺส อาราเม.\csromanspacer{} อถ โข อญฺญตโร ภิกฺขุ เยน ภควา เตนุปสงฺกมิ, อุปสงฺกมิตฺวา ภควนฺตํ อภิวาเทตฺวา เอกมนฺตํ นิสีทิ, เอกมนฺตํ นิสินฺโน โข โส ภิกฺขุ ภควนฺตํ เอตทโวจ “สาธุ เม ภนฺเต ภควา สํขิตฺเตน ธมฺมํ เทเสตุ, ยมหํ ภควโต ธมฺมํ สุตฺวา เอโก วูปกฏฺโฐ อปฺปมตฺโต อาตาปี ปหิตตฺโต วิหเรยฺยนฺ”ติ.\csromanspacer{} อุปาทิยมาโน โข ภิกฺขุ พทฺโธ มารสฺส, อนุปาทิยมาโน มุตฺโต ปาปิมโตติ.\csromanspacer{} อญฺญาตํ ภควา, อญฺญาตํ สุคตาติ.}

\prose{ยถา กถํ ปน ตฺวํ ภิกฺขุ มยา สํขิตฺเตน ภาสิตสฺส วิตฺถาเรน อตฺถํ อาชานาสีติ.\csromanspacer{} รูปํ โข ภนฺเต อุปาทิยมาโน พทฺโธ มารสฺส, อนุปาทิยมาโน มุตฺโต ปาปิมโต.\csromanspacer{} เวทนํ อุปาทิยมาโน พทฺโธ มารสฺส, อนุปาทิยมาโน มุตฺโต ปาปิมโต.\csromanspacer{} สญฺญํ.\csromanspacer{} สงฺขาเร.\csromanspacer{} วิญฺญาณํ อุปาทิยมาโน พทฺโธ มารสฺส, อนุปาทิยมาโน มุตฺโต ปาปิมโต.\csromanspacer{} อิมสฺส ขฺวาหํ ภนฺเต ภควตา สํขิตฺเตน ภาสิตสฺส เอวํ วิตฺถาเรน อตฺถํ อาชานามีติ.}

\prose{สาธุ สาธุ ภิกฺขุ, สาธุ โข ตฺวํ ภิกฺขุ มยา สํขิตฺเตน ภาสิตสฺส วิตฺถาเรน อตฺถํ อาชานาสิ.\csromanspacer{} รูปํ โข ภิกฺขุ อุปาทิยมาโน พทฺโธ มารสฺส, อนุปาทิยมาโน มุตฺโต ปาปิมโต.\csromanspacer{} เวทนํ.\csromanspacer{} สญฺญํ.\csromanspacer{} สงฺขาเร.\csromanspacer{} วิญฺญาณํ อุปาทิยมาโน พทฺโธ มารสฺส, อนุปาทิยมาโน มุตฺโต ปาปิมโต.\csromanspacer{} อิมสฺส โข ภิกฺขุ มยา สํขิตฺเตน ภาสิตสฺส เอวํ วิตฺถาเรน อตฺโถ ทฏฺฐพฺโพติ.}

\prose{อถ โข โส ภิกฺขุ ภควโต ภาสิตํ อภินนฺทิตฺวา อนุโมทิตฺวา อุฏฺฐายาสนา ภควนฺตํ อภิวาเทตฺวา ปทกฺขิณํ กตฺวา ปกฺกามิ.\csromanspacer{} อถ โข โส ภิกฺขุ เอโก วูปกฏฺโฐ อปฺปมตฺโต อาตาปี ปหิตตฺโต วิหรนฺโต นจิรสฺเสว ยสฺสตฺถาย กุลปุตฺตา สมฺมเทว อคารสฺมา อนคาริยํ ปพฺพชนฺติ, ตทนุตฺตรํ พฺรหฺมจริยปริโยสานํ ทิฏฺเฐว ธมฺเม สยํ อภิญฺญา สจฺฉิกตฺวา อุปสมฺปชฺช วิหรติ, “ขีณา ชาติ, วุสิตํ พฺราหฺมจริยํ, \csromanfolio{62}กตํ กรณียํ, นาปรํ อิตฺถตฺตายา”ติ อพฺภญฺญาสิ.\csromanspacer{} อญฺญตโร จ ปน โส ภิกฺขุ อรหตํ อโหสีติ.\csromanspacer{}\csromanspacer{}ปฐมํ.}

\csromanheader{2}{2. มญฺญมานสุตฺต}
\proseitem{64}{สาวตฺถินิทานํ.\csromanspacer{} อถ โข อญฺญตโร ภิกฺขุ ฯเปฯ เอกมนฺตํ นิสินฺโน โข โส ภิกฺขุ ภควนฺตํ เอตทโวจ “สาธุ เม ภนฺเต ภควา สํขิตฺเตน ธมฺมํ เทเสตุ ฯเปฯ อาตาปี ปหิตตฺโต วิหเรยฺยนฺ”ติ.\csromanspacer{} มญฺญมาโน โข ภิกฺขุ พทฺโธ มารสฺส, อมญฺญมาโน มุตฺโต ปาปิมโตติ.\csromanspacer{} อญฺญาตํ ภควา, อญฺญาตํ สุคตาติ.}

\prose{ยถา กถํ ปน ตฺวํ ภิกฺขุ มยา สํขิตฺเตน ภาสิตสฺส วิตฺถาเรน อตฺถํ อาชานาสีติ.\csromanspacer{} รูปํ โข ภนฺเต มญฺญมาโน พทฺโธ มารสฺส, อมญฺญมาโน มุตฺโต ปาปิมโต.\csromanspacer{} เวทนํ.\csromanspacer{} สญฺญํ.\csromanspacer{} สงฺขาเร.\csromanspacer{} วิญฺญาณํ มญฺญมาโน พทฺโธ มารสฺส, อมญฺญมาโน มุตฺโต ปาปิมโต.\csromanspacer{} อิมสฺส ขฺวาหํ ภนฺเต ภควตา สํขิตฺเตน ภาสิตสฺส เอวํ วิตฺถาเรน อตฺถํ อาชานามีติ.}

\prose{สาธุ สาธุ ภิกฺขุ, สาธุ โข ตฺวํ ภิกฺขุ มยา สํขิตฺเตน ภาสิตสฺส วิตฺถาเรน อตฺถํ อาชานาสิ, รูปํ โข ภิกฺขุ มญฺญมาโน พทฺโธ มารสฺส, อมญฺญมาโน มุตฺโต ปาปิมโต.\csromanspacer{} เวทนํ.\csromanspacer{} สญฺญํ.\csromanspacer{} สงฺขาเร.\csromanspacer{} วิญฺญาณํ มญฺญมาโน พทฺโธ มารสฺส, อมญฺญมาโน มุตฺโต ปาปิมตฺโต.\csromanspacer{} อิมสฺส โข ภิกฺขุ มยา สํขิตฺเตน ภาสิตสฺส เอวํ วิตฺถาเรน อตฺโถ ทฏฺฐพฺโพติ ฯเปฯ.\csromanspacer{} อญฺญตโร จ ปน โส ภิกฺขุ อรหตํ อโหสีติ.\csromanspacer{}\csromanspacer{}ทุติยํ.}

\csromanheader{2}{3. อภินนฺทมานสุตฺต}
\proseitem{65}{สาวตฺถินิทานํ.\csromanspacer{} อถ โข อญฺญตโร ภิกฺขุ ฯเปฯ เอกมนฺตํ นิสินฺโน โข โส ภิกฺขุ ภควนฺตํ เอตทโวจ “สาธุ เม ภนฺเต ภควา สํขิตฺเตน ฯเปฯ ปหิตตฺโต วิหเรยฺยนฺ”ติ.\csromanspacer{} อภินนฺทมาโน โข ภิกฺขุ พทฺโธ มารสฺส, อนภินนฺทมาโน มุตฺโต ปาปิมโตติ.\csromanspacer{} อญฺญาตํ ภควา, อญฺญาตํ สุคตาติ.}

\prose{\csromanfolio{63}ยถา กถํ ปน ตฺวํ ภิกฺขุ มยา สํขิตฺเตน ภาสิตสฺส วิตฺถาเรน อตฺถํ อาชานาสีติ.\csromanspacer{} รูปํ โข ภนฺเต อภินนฺทมาโน พทฺโธ มารสฺส, อนภินนฺทมาโน มุตฺโต ปาปิมโต.\csromanspacer{} เวทนํ.\csromanspacer{} สญฺญํ.\csromanspacer{} สงฺขาเร.\csromanspacer{} วิญฺญาณํ อภินนฺทมาโน พทฺโธ มารสฺส, อนภินนฺทมาโน มุตฺโต ปาปิมโต.\csromanspacer{} อิมสฺส ขฺวาหํ ภนฺเต ภควตา สํขิตฺเตน ภาสิตสฺส เอวํ วิตฺถาเรน อตฺถํ อาชานามีติ.}

\prose{สาธุ สาธุ ภิกฺขุ, สาธุ โข ตฺวํ ภิกฺขุ มยา สํขิตฺเตน ภาสิตสฺส วิตฺถาเรน อตฺถํ อาชานาสิ, รูปํ โข ภิกฺขุ อภินนฺทมาโน พทฺโธ มารสฺส, อนภินนฺทมาโน มุตฺโต ปาปิมโต.\csromanspacer{} เวทนํ.\csromanspacer{} สญฺญํ.\csromanspacer{} สงฺขาเร.\csromanspacer{} วิญฺญาณํ อภินนฺทมาโน พทฺโธ มารสฺส, อนภินนฺทมาโน มุตฺโต ปาปิมโต.\csromanspacer{} อิมสฺส โข ภิกฺขุ มยา สํขิตฺเตน ภาสิตสฺส เอวํ วิตฺถาเรน อตฺโถ ทฏฺฐพฺโพติ ฯเปฯ.\csromanspacer{} อญฺญตโร จ ปน โส ภิกฺขุ อรหตํ อโหสีติ.\csromanspacer{}\csromanspacer{}ตติยํ.}

\csromanheader{2}{4. อนิจฺจสุตฺต}
\proseitem{66}{สาวตฺถินิทานํ.\csromanspacer{} อถ โข อญฺญตโร ภิกฺขุ ฯเปฯ เอกมนฺตํ นิสินฺโน โข โส ภิกฺขุ ภควนฺตํ เอตทโวจ “สาธุ เม ภนฺเต ภควา สํขิตฺเตน ธมฺมํ เทเสตุ ฯเปฯ อาตาปี ปหิตตฺโต วิหเรยฺยนฺ”ติ.\csromanspacer{} ยํ โข ภิกฺขุ อนิจฺจํ, ตตฺร เต ฉนฺโท ปหาตพฺโพติ.\csromanspacer{} อญฺญาตํ ภควา, อญฺญาตํ สุคตาติ.}

\prose{ยถา กถํ ปน ตฺวํ ภิกฺขุ มยา สํขิตฺเตน ภาสิตสฺส วิตฺถาเรน อตฺถํ อาชานาสีติ.\csromanspacer{} รูปํ โข ภนฺเต อนิจฺจํ, ตตฺร เม ฉนฺโท ปหาตพฺโพ.\csromanspacer{} เวทนา.\csromanspacer{} สญฺญา.\csromanspacer{} สงฺขารา.\csromanspacer{} วิญฺญาณํ อนิจฺจํ, ตตฺร เม ฉนฺโท ปหาตพฺโพ.\csromanspacer{} อิมสฺส ขฺวาหํ ภนฺเต ภควตา สํขิตฺเตน ภาสิตสฺส เอวํ วิตฺถาเรน อตฺถํ อาชานามีติ.}

\prose{สาธุ สาธุ ภิกฺขุ, สาธุ โข ตฺวํ ภิกฺขุ มยา สํขิตฺเตน ภาสิตสฺส วิตฺถาเรน อตฺถํ อาชานาสิ.\csromanspacer{} รูปํ โข ภิกฺขุ อนิจฺจํ, ตตฺร เต ฉนฺโท ปหาตพฺโพ.\csromanspacer{} เวทนา อนิจฺจา.\csromanspacer{} สญฺญา.\csromanspacer{} สงฺขารา.\csromanspacer{} วิญฺญาณํ อนิจฺจํ, ตตฺร โข เต ฉนฺโท ปหาตพฺโพ.\csromanspacer{} อิมสฺส โข ภิกฺขุ มยา สํขิตฺเตน ภาสิตสฺส เอวํ วิตฺถาเรน อตฺโถ ทฏฺฐพฺโพติ ฯเปฯ.\csromanspacer{} อญฺญตโร จ ปน โส ภิกฺขุ อรหตํ อโหสีติ.\csromanspacer{}\csromanspacer{}จตุตฺถํ.}

\csromanheader{2}{5. ทุกฺขสุตฺต}
\proseitem{67}{\csromanfolio{64}สาวตฺถินิทานํ.\csromanspacer{} อถ โข อญฺญตโร ภิกฺขุ ฯเปฯ เอกมนฺตํ นิสินฺโน โข โส ภิกฺขุ ภควนฺตํ เอตทโวจ “สาธุ เม ภนฺเต ภควา สํขิตฺเตน ธมฺมํ เทเสตุ ฯเปฯ อาตาปี ปหิตตฺโต วิหเรยฺยนฺ”ติ.\csromanspacer{} ยํ โข ภิกฺขุ ทุกฺขํ, ตตฺร เต ฉนฺโท ปหาตพฺโพติ.\csromanspacer{} อญฺญาตํ ภควา, อญฺญาตํ สุคตาติ.}

\prose{ยถา กถํ ปน ตฺวํ ภิกฺขุ มยา สํขิตฺเตน ภาสิตสฺส วิตฺถาเรน อตฺถํ อาชานาสีติ.\csromanspacer{} รูปํ โข ภนฺเต ทุกฺขํ, ตตฺร เม ฉนฺโท ปหาตพฺโพ.\csromanspacer{} เวทนา.\csromanspacer{} สญฺญา.\csromanspacer{} สงฺขารา.\csromanspacer{} วิญฺญาณํ ทุกฺขํ, ตตฺร เม ฉนฺโท ปหาตพฺโพ.\csromanspacer{} อิมสฺส ขฺวาหํ ภนฺเต ภควตา สํขิตฺเตน ภาสิตสฺส เอวํ วิตฺถาเรน อตฺถํ อาชานามีติ.}

\prose{สาธุ สาธุ ภิกฺขุ, สาธุ โข ตฺวํ ภิกฺขุ มยา สํขิตฺเตน ภาสิตสฺส วิตฺถาเรน อตฺถํ อาชานาสิ.\csromanspacer{} รูปํ โข ภิกฺขุ ทุกฺขํ, ตตฺร เต ฉนฺโท ปหาตพฺโพ.\csromanspacer{} เวทนา.\csromanspacer{} สญฺญา.\csromanspacer{} สงฺขารา.\csromanspacer{} วิญฺญาณํ ทุกฺขํ, ตตฺร เต ฉนฺโท ปหาตพฺโพ.\csromanspacer{} อิมสฺส โข ภิกฺขุ มยา สํขิตฺเตน ภาสิตสฺส เอวํ วิตฺถาเรน อตฺโถ ทฏฺฐพฺโพติ ฯเปฯ.\csromanspacer{} อญฺญตโร จ ปน โส ภิกฺขุ อรหตํ อโหสีติ.\csromanspacer{}\csromanspacer{}ปญฺจมํ.}

\csromanheader{2}{6. อนตฺตสุตฺต}
\proseitem{68}{สาวตฺถินิทานํ.\csromanspacer{} อถ โข อญฺญตโร ภิกฺขุ ฯเปฯ เอกมนฺตํ นิสินฺโน โข โส ภิกฺขุ ภควนฺตํ เอตทโวจ “สาธุ เม ภนฺเต ภควา สํขิตฺเตน ธมฺมํ เทเสตุ ฯเปฯ ตาปี ปหิตตฺโต วิหเรยฺยนฺ”ติ.\csromanspacer{} โย โข ภิกฺขุ อนตฺตา, ตตฺร เต ฉนฺโท ปหาตพฺโพติ.\csromanspacer{} อญฺญาตํ ภควา, อญฺญาตํ สุคตาติ.\csromanspacer{} ยถา กถํ ปน ตฺวํ ภิกฺขุ มยา สํขิตฺเตน ภาสิตสฺส วิตฺถาเรน อตฺถํ อาชานาสีติ.\csromanspacer{} รูปํ โข ภนฺเต อนตฺตา, ตตฺร เม ฉนฺโท ปหาตพฺโพ.\csromanspacer{} เวทนา.\csromanspacer{} สญฺญา.\csromanspacer{} สงฺขารา.\csromanspacer{} วิญฺญาณํ อนตฺตา, ตตฺร เม ฉนฺโท ปหาตพฺโพ.\csromanspacer{} อิมสฺส ขฺวาหํ ภนฺเต ภควตา สํขิตฺเตน ภาสิตสฺส เอวํ วิตฺถาเรน อตฺถํ อาชานามีติ.}

\prose{สาธุ สาธุ ภิกฺขุ, สาธุ โข ตฺวํ ภิกฺขุ มยา สํขิตฺเตน ภาสิตสฺส วิตฺถาเรน อตฺถํ อาชานาสิ.\csromanspacer{} รูปํ โข ภิกฺขุ อนตฺตา, ตตฺร เต ฉนฺโท \csromanfolio{65}ปหาตพฺโพ.\csromanspacer{} เวทนา.\csromanspacer{} สญฺญา.\csromanspacer{} สงฺขารา.\csromanspacer{} วิญฺญาณํ อนตฺตา, ตตฺร เต ฉนฺโท ปหาตพฺโพ.\csromanspacer{} อิมสฺส โข ภิกฺขุ มยา สํขิตฺเตน ภาสิตสฺส เอวํ วิตฺถาเรน อตฺโถ ทฏฺฐพฺโพติ ฯเปฯ.\csromanspacer{} อญฺญตโร จ ปน โส ภิกฺขุ อรหตํ อโหสีติ.\csromanspacer{}\csromanspacer{}ฉฏฺฐํ.}

\csromanheader{2}{7. อนตฺตนิยสุตฺต}
\proseitem{69}{สาวตฺถินิทานํ.\csromanspacer{} อถ โข อญฺญตโร ภิกฺขุ ฯเปฯ เอกมนฺตํ นิสินฺโน โข โส ภิกฺขุ ภควนฺตํ เอตทโวจ “สาธุ เม ภนฺเต ภควา สํขิตฺเตน ธมฺมํ เทเสตุ ฯเปฯ วิหเรยฺยนฺ”ติ.\csromanspacer{} ยํ โข ภิกฺขุ อนตฺตนิยํ, ตตฺร เต ฉนฺโท ปหาตพฺโพติ.\csromanspacer{} อญฺญาตํ ภควา, อญฺญาตํ สุคตาติ.\csromanspacer{} ยถา กถํ ปน ตฺวํ ภิกฺขุ มยา สํขิตฺเตน ภาสิตสฺส วิตฺถาเรน อตฺถํ อาชานาสีติ.\csromanspacer{} รูปํ โข ภนฺเต อนตฺตนิยํ, ตตฺร เม ฉนฺโท ปหาตพฺโพ.\csromanspacer{} เวทนา.\csromanspacer{} สญฺญา.\csromanspacer{} สงฺขารา.\csromanspacer{} วิญฺญาณํ อนตฺตนิยํ, ตตฺร เม ฉนฺโท ปหาตพฺโพ.\csromanspacer{} อิมสฺส ขฺวาหํ ภนฺเต ภควตา สํขิตฺเตน ภาสิตสฺส เอวํ วิตฺถาเรน อตฺถํ อาชานามีติ.}

\prose{สาธุ สาธุ ภิกฺขุ, สาธุ โข ตฺวํ ภิกฺขุ มยา สํขิตฺเตน ภาสิตสฺส วิตฺถาเรน อตฺถํ อาชานาสิ.\csromanspacer{} รูปํ โข ภิกฺขุ อนตฺตนิยํ, ตตฺร เต ฉนฺโท ปหาตพฺโพ.\csromanspacer{} เวทนา.\csromanspacer{} สญฺญา.\csromanspacer{} สงฺขารา.\csromanspacer{} วิญฺญาณํ อนตฺตนิยํ, ตตฺร เต ฉนฺโท ปหาตพฺโพ.\csromanspacer{} อิมสฺส โข ภิกฺขุ มยา สํขิตฺเตน ภาสิตสฺส เอวํ วิตฺถาเรน อตฺโถ ทฏฺฐพฺโพติ ฯเปฯ.\csromanspacer{} อญฺญตโร จ ปน โส ภิกฺขุ อรหตํ อโหสีติ.\csromanspacer{}\csromanspacer{}สตฺตมํ.}

\csromanheader{2}{8. รชนียสณฺฐิตสุตฺต}
\proseitem{70}{สาวตฺถินิทานํ.\csromanspacer{} อถ โข อญฺญตโร ภิกฺขุ ฯเปฯ เอกมนฺตํ นิสินฺโน โข โส ภิกฺขุ ภควนฺตํ เอตทโวจ “สาธุ เม ภนฺเต ภควา สํขิตฺเตน ธมฺมํ เทเสตุ, ยมหํ ภควโต ธมฺมํ สุตฺวา ฯเปฯ วิหเรยฺยนฺ”ติ.\csromanspacer{} ยํ โข ภิกฺขุ รชนียสณฺฐิตํ, ตตฺร เต ฉนฺโท ปหาตพฺโพติ.\csromanspacer{} อญฺญาตํ ภควา, อญฺญาตํ สุคตาติ.\csromanspacer{} ยถา กถํ ปน ตฺวํ ภิกฺขุ มยา สํขิตฺเตน ภาสิตสฺส วิตฺถาเรน อตฺถํ อาชานาสีติ.\csromanspacer{} รูปํ โข ภนฺเต รชนียสณฺฐิตํ, ตตฺร เม ฉนฺโท ปหาตพฺโพ.\csromanspacer{} เวทนา.\csromanspacer{} สญฺญา.\csromanspacer{} สงฺขารา.\csromanspacer{} วิญฺญาณํ รชนียสณฺฐิตํ, ตตฺร เม ฉนฺโท ปหาตพฺโพ.\csromanspacer{} อิมสฺส ขฺวาหํ ภนฺเต ภควตา สํขิตฺเตน ภาสิตสฺส เอวํ วิตฺถาเรน อตฺถํ อาชานามีติ.}

\prose{\csromanfolio{66}สาธุ สาธุ ภิกฺขุ, สาธุ โข ตฺวํ ภิกฺขุ มยา สํขิตฺเตน ภาสิตสฺส วิตฺถาเรน อตฺถํ อาชานาสิ.\csromanspacer{} รูปํ โข ภิกฺขุ รชนียสณฺฐิตํ, ตตฺร เต ฉนฺโท ปหาตพฺโพ.\csromanspacer{} เวทนา.\csromanspacer{} สญฺญา.\csromanspacer{} สงฺขารา.\csromanspacer{} วิญฺญาณํ รชนียสณฺฐิตํ, ตตฺร เต ฉนฺโท ปหาตพฺโพ.\csromanspacer{} อิมสฺส โข ภิกฺขุ มยา สํขิตฺเตน ภาสิตสฺส เอวํ วิตฺถาเรน อตฺโถ ทฏฺฐพฺโพติ ฯเปฯ.\csromanspacer{} อญฺญตโร จ ปน โส ภิกฺขุ อรหตํ อโหสีติ.\csromanspacer{}\csromanspacer{}อฏฺฐมํ.}

\csromanheader{2}{9. ราธสุตฺต}
\proseitem{71}{สาวตฺถินิทานํ.\csromanspacer{} อถ โข อายสฺมา ราโธ เยน ภควา เตนุปสงฺกมิ, อุปสงฺกมิตฺวา ภควนฺตํ เอตทโวจ “กถํ นุ โข ภนฺเต ชานโต กถํ ปสฺสโต อิมสฺมิํ จ สวิญฺญาณเก กาเย พหิทฺธา จ สพฺพนิมิตฺเตสุ อหงฺการมมงฺการมานานุสยา น โหนฺตี”ติ.\csromanspacer{} ยํ กิญฺจิ ราธ รูปํ อตีตานาคตปจฺจุปฺปนฺนํ อชฺฌตฺตํ วา พหิทฺธา วา โอฬาริกํ วา สุขุมํ วา หีนํ วา ปณีตํ วา ยํ ทูเร สนฺติเก วา, สพฺพํ รูปํ “เนตํ มม, เนโสหมสฺมิ, น เมโส อตฺตา”ติ เอวเมตํ ยถาภูตํ สมฺมปฺปญฺญาย ปสฺสติ, ยา กาจิ เวทนา.\csromanspacer{} ยา กาจิ สญฺญา.\csromanspacer{} เย เกจิ สงฺขารา.\csromanspacer{} ยํ กิญฺจิ วิญฺญาณํ อตีตานาคตปจฺจุปฺปนฺนํ ฯเปฯ ยํ ทูเร สนฺติเก วา, สพฺพํ วิญฺญาณํ “เนตํ มม, เนโสหมสฺมิ, น เมโส อตฺตา”ติ เอวเมตํ ยถาภูตํ สมฺมปฺปญฺญาย ปสฺสติ.\csromanspacer{} เอวํ โข ราธ ชานโต เอวํ ปสฺสโต อิมสฺมิํ จ สวิญฺญาณเก กาเย พหิทฺธา จ สพฺพนิมิตฺเตสุ อหงฺการมมงฺการมานานุสยา น โหนฺตีติ ฯเปฯ.\csromanspacer{} อญฺญตโร จ ปนายสฺมา ราโธ อรหตํ อโหสีติ.\csromanspacer{}\csromanspacer{}นวมํ.}

\csromanheader{2}{10. สุราธสุตฺต}
\proseitem{72}{อถ โข อายสฺมา สุราโธ ภควนฺตํ เอตทโวจ “กถํ นุ โข ภนฺเต ชานโต กถํ ปสฺสโต อิมสฺมิํ จ สวิญฺญาณเก กาเย พหิทฺธา จ สพฺพนิมิตฺเตสุ อหงฺการมมงฺการมานาปคตํ มานสํ โหติ วิธา สมติกฺกนฺตํ สนฺตํ สุวิมุตฺตนฺ”ติ.\csromanspacer{} ยํ กิญฺจิ สุราธ รูปํ อตีตานาคตปจฺจุปฺปนฺนํ ฯเปฯ ยํ ทูเร สนฺติเก วา, สพฺพํ รูปํ “เนตํ มม, เนโสหมสฺมิ, น เมโส อตฺตา”ติ เอวเมตํ ยถาภูตํ สมฺมปฺปญฺญาย ทิสฺวา อนุปาทาวิมุตฺโต โหติ.\csromanspacer{} ยา กาจิ เวทนา.\csromanspacer{} ยา กาจิ สญฺญา.\csromanspacer{} เย เกจิ สงฺขารา.\csromanspacer{} ยํ กิญฺจิ วิญฺญาณํ อตีตานาคตปจฺจุปฺปนฺนํ อชฺฌตฺตํ วา พหิทฺธา วา โอฬาริกํ \csromanfolio{67}วา สุขุมํ วา หีนํ วา ปณีตํ วา ยํ ทูเร สนฺติเก วา, สพฺพา เวทนา ฯเปฯ สพฺพา สญฺญา.\csromanspacer{} สพฺเพ สงฺขารา.\csromanspacer{} สพฺพํ วิญฺญาณํ “เนตํ มม, เนโสหมสฺมิ, น เมโส อตฺตา”ติ เอวเมตํ ยถาภูตํ สมฺมปฺปญฺญาย ทิสฺวา อนุปาทาวิมุตฺโต โหติ.\csromanspacer{} เอวํ โข สุราธ ชานโต เอวํ ปสฺสโต อิมสฺมิํ จ สวิญฺญาณเก กาเย พหิทฺธา จ สพฺพนิมิตฺเตสุ อหงฺการมมงฺการมานาปคตํ มานสํ โหติ วิธา สมติกฺกนฺตํ สนฺตํ สุวิมุตฺตนฺติ ฯเปฯ.\csromanspacer{} อญฺญตโร จ ปนายสฺมา สุราโธ อรหตํ อโหสีติ.\csromanspacer{}\csromanspacer{}ทสมํ.\csromanspacer{} อรหนฺตวคฺโค ทุติโย.}

\titlehead{\textbf{ตสฺสุทฺทานํ}}
\csromangathameasure{อุปาทิยมญฺญมานา, อถาภินนฺทมาโน จ. \\ อนิจฺจํ ทุกฺขํ อนตฺตา จ, อนตฺตนียํ รชนียสณฺฐิตํ.}
\csromangathagroup{\csromangathastanza{อุปาทิยมญฺญมานา, อถาภินนฺทมาโน จ. \\ อนิจฺจํ ทุกฺขํ อนตฺตา จ, อนตฺตนียํ รชนียสณฺฐิตํ.}}

\csromanheader{2}{ราธสุราเธน เต ทสาติ.}
\csromansectionrule
\chapterhead{\textbf{(8) 3. ขชฺชนียวคฺค}}
\csromanheader{2}{1. อสฺสาทสุตฺต}
\proseitem{73}{สาวตฺถินิทานํ.\csromanspacer{} อสฺสุตวา ภิกฺขเว ปุถุชฺชโน รูปสฺส อสฺสาทญฺจ อาทีนวญฺจ นิสฺสรณญฺจ ยถาภูตํ นปฺปชานาติ.\csromanspacer{} เวทนาย.\csromanspacer{} สญฺญาย.\csromanspacer{} สงฺขารานํ.\csromanspacer{} วิญฺญาณสฺส อสฺสาทญฺจ อาทีนวญฺจ นิสฺสรณญฺจ ยถาภูตํ นปฺปชานาติ.\csromanspacer{} สุตวา จ โข ภิกฺขเว อริยสาวโก รูปสฺส อสฺสาทญฺจ อาทีนวญฺจ นิสฺสรณญฺจ ยถาภูตํ ปชานาติ.\csromanspacer{} เวทนาย.\csromanspacer{} สญฺญาย.\csromanspacer{} สงฺขารานํ.\csromanspacer{} วิญฺญาณสฺส อสฺสาทญฺจ อาทีนวญฺจ นิสฺสรณญฺจ ยถาภูตํ ปชานาตีติ.\csromanspacer{}\csromanspacer{}ปฐมํ.}

\csromanheader{2}{2. สมุทยสุตฺต}
\proseitem{74}{สาวตฺถินิทานํ.\csromanspacer{} อสฺสุตวา ภิกฺขเว ปุถุชฺชโน รูปสฺส สมุทยญฺจ อตฺถงฺคมญฺจ อสฺสาทญฺจ อาทีนวญฺจ นิสฺสรณญฺจ ยถาภูตํ นปฺปชานาติ.\csromanspacer{} เวทนาย.\csromanspacer{} สญฺญาย.\csromanspacer{} สงฺขารานํ.\csromanspacer{} วิญฺญาณสฺส สมุทยญฺจ อตฺถงฺคมญฺจ อสฺสาทญฺจ \csromanfolio{68}อาทีนวญฺจ นิสฺสรณญฺจ ยถาภูตํ นปฺปชานาติ.\csromanspacer{} สุตวา จ โข ภิกฺขเว อริยสาวโก รูปสฺส สมุทยญฺจ อตฺถงฺคมญฺจ อสฺสาทญฺจ อาทีนวญฺจ นิสฺสรณญฺจ ยถาภูตํ ปชานาติ.\csromanspacer{} เวทนาย.\csromanspacer{} สญฺญาย.\csromanspacer{} สงฺขารานํ.\csromanspacer{} วิญฺญาณสฺส สมุทยญฺจ อตฺถงฺคมญฺจ อสฺสาทญฺจ อาทีนวญฺจ นิสฺสรณญฺจ ยถาภูตํ ปชานาตีติ.\csromanspacer{} ทุติยํ.}

\csromanheader{2}{3. ทุติยสมุทยสุตฺต}
\proseitem{75}{สาวตฺถินิทานํ.\csromanspacer{} สุตวา ภิกฺขเว อริยสาวโก รูปสฺส สมุทยญฺจ อตฺถงฺคมญฺจ อสฺสาทญฺจ อาทีนวญฺจ นิสฺสรณญฺจ ยถาภูตํ ปชานาติ.\csromanspacer{} เวทนาย.\csromanspacer{} สญฺญาย.\csromanspacer{} สงฺขารานํ.\csromanspacer{} วิญฺญาณสฺส สมุทยญฺจ อตฺถงฺคมญฺจ อสฺสาทญฺจ อาทีนวญฺจ นิสฺสรณญฺจ ยถาภูตํ ปชานาตีติ.\csromanspacer{}\csromanspacer{}ตติยํ.}

\csromanheader{2}{4. อรหนฺตสุตฺต}
\proseitem{76}{สาวตฺถินิทานํ.\csromanspacer{} รูปํ ภิกฺขเว อนิจฺจํ, ยทนิจฺจํ ตํ ทุกฺขํ, ยํ ทุกฺขํ ตทนตฺตา, ยทนตฺตา, ตํ “เนตํ มม, เนโสหมสฺมิ, น เมโส อตฺตา”ติ เอวเมตํ ยถาภูตํ สมฺมปฺปญฺญาย ทฏฺฐพฺพํ.\csromanspacer{} เวทนา.\csromanspacer{} สญฺญา.\csromanspacer{} สงฺขารา.\csromanspacer{} วิญฺญาณํ อนิจฺจํ, ยทนิจฺจํ ตํ ทุกฺขํ, ยํ ทุกฺขํ ตทนตฺตา, ยทนตฺตา ตํ “เนตํ มม, เนโสหมสฺมิ, น เมโส อตฺตา”ติ เอวเมตํ ยถาภูตํ สมฺมปฺปญฺญาย ทฏฺฐพฺพํ.}

\prose{เอวํ ปสฺสํ ภิกฺขเว สุตวา อริยสาวโก รูปสฺมิมฺปิ นิพฺพินฺทติ.\csromanspacer{} เวทนายปิ.\csromanspacer{} สญฺญายปิ.\csromanspacer{} สงฺขาเรสุปิ.\csromanspacer{} วิญฺญาณสฺมิมฺปิ นิพฺพินฺทติ, นิพฺพินฺทํ วิรชฺชติ, วิราคา วิมุจฺจติ, วิมุตฺตสฺมิํ “วิมุตฺตมฺ”อิติ ญาณํ โหติ, “ขีณา ชาติ, วุสิตํ พฺราหฺมจริยํ, กตํ กรณียํ, นาปรํ อิตฺถตฺตายา”ติ ปชานาติ.\csromanspacer{} ยาวตาภิกฺขเว สตฺตาวาสา ยาวตา ภวคฺคํ, เอเต อคฺคา เอเต เสฏฺฐา โลกสฺมิํ ยทิทํ อรหนฺโตติ.}

\prose{อิทมโวจ ภควา, อิทํ วตฺวาน สุคโต อถาปรํ เอตทโวจ สตฺถา\csromandash{}}

\csromangathameasure{“สุขิโน วต อรหนฺโต, ตณฺหา เตสํ น วิชฺชติ. \\ อสฺมิมาโน สมุจฺฉินฺโน, โมหชาลํ ปทาลิตํ. \\ อเนชํ เต อนุปฺปตฺตา, จิตฺตํ เตสํ อนาวิลํ. \\ โลเก อนุปลิตฺตา เต, พฺรหฺมภูตา อนาสวา. \\ ปญฺจกฺขนฺเธ ปริญฺญาย, สตฺต สทฺธมฺมโคจรา. \\ ปสํสิยา สปฺปุริสา, ปุตฺตา พุทฺธสฺส โอรสา. \\ สตฺตรตนสมฺปนฺนา, ตีสุ สิกฺขาสุ สิกฺขิตา. \\ อนุวิจรนฺติ มหาวีรา, ปหีนภยเภรวา. \\ ทสหงฺเคหิ สมฺปนฺนา, มหานาคา สมาหิตา. \\ เอเต โข เสฏฺฐา โลกสฺมิํ, ตณฺหา เตสํ น วิชฺชติ. \\ อเสขญาณมุปฺปนฺนํ, อนฺติโมยํ สมุสฺสโย. \\ โย สาโร พฺรหฺมจริยสฺส, ตสฺมิํ อปรปจฺจยา. \\ วิธาสุ น วิกมฺปนฺติ, วิปฺปมุตฺตา ปุนพฺภวา. \\ ทนฺตภูมิมนุปฺปตฺตา, เต โลเก วิชตาวิโน. \\ อุทฺธํ ติริยํ อปาจีนํ, นนฺที เตสํ น วิชฺชติ. \\ นทนฺติ เต สีหนาทํ, พุทฺธา โลเก อนุตฺตรา”ติ. จตุตฺถํ.}
\csromangathagroup{\csromangathastanza{“สุขิโน วต อรหนฺโต, ตณฺหา เตสํ น วิชฺชติ. \\ อสฺมิมาโน สมุจฺฉินฺโน, โมหชาลํ ปทาลิตํ.}\csromangathastanza{\csromanfolio{69}อเนชํ เต อนุปฺปตฺตา, จิตฺตํ เตสํ อนาวิลํ. \\ โลเก อนุปลิตฺตา เต, พฺรหฺมภูตา อนาสวา.}\csromangathastanza{ปญฺจกฺขนฺเธ ปริญฺญาย, สตฺต สทฺธมฺมโคจรา. \\ ปสํสิยา สปฺปุริสา, ปุตฺตา พุทฺธสฺส โอรสา.}\csromangathastanza{สตฺตรตนสมฺปนฺนา, ตีสุ สิกฺขาสุ สิกฺขิตา. \\ อนุวิจรนฺติ มหาวีรา, ปหีนภยเภรวา.}\csromangathastanza{ทสหงฺเคหิ สมฺปนฺนา, มหานาคา สมาหิตา. \\ เอเต โข เสฏฺฐา โลกสฺมิํ, ตณฺหา เตสํ น วิชฺชติ.}\csromangathastanza{อเสขญาณมุปฺปนฺนํ, อนฺติโมยํ\footnote{อนฺติมสฺส (ก)} สมุสฺสโย. \\ โย สาโร พฺรหฺมจริยสฺส, ตสฺมิํ อปรปจฺจยา.}\csromangathastanza{วิธาสุ น วิกมฺปนฺติ, วิปฺปมุตฺตา ปุนพฺภวา. \\ ทนฺตภูมิมนุปฺปตฺตา, เต โลเก วิชตาวิโน.}\csromangathastanza{อุทฺธํ ติริยํ อปาจีนํ, นนฺที เตสํ น วิชฺชติ. \\ นทนฺติ เต สีหนาทํ, พุทฺธา โลเก อนุตฺตรา”ติ.\csromanspacer{} จตุตฺถํ.}}

\csromanheader{2}{5. ทุติย-อรหนฺตสุตฺต}
\proseitem{77}{สาวตฺถินิทานํ.\csromanspacer{} รูปํ ภิกฺขเว อนิจฺจํ.\csromanspacer{} ยทนิจฺจํ ตํ ทุกฺขํ, ยํ ทุกฺขํ ตทนตฺตา, ยทนตฺตา ตํ “เนตํ มม, เนโสหมสฺมิ, น เมโส อตฺตา”ติ ฯเปฯ เอวเมตํ ยถาภูตํ สมฺมปฺปญฺญาย ทฏฺฐพฺพํ.}

\prose{เอวํ ปสฺสํ ภิกฺขเว สุตวา อริยสาวโก รูปสฺมิมฺปิ นิพฺพินฺทติ, เวทนายปิ.\csromanspacer{} สญฺญายปิ.\csromanspacer{} สงฺขาเรสุปิ.\csromanspacer{} วิญฺญาณสฺมิมฺปิ นิพฺพินฺทติ, นิพฺพินฺทํ วิรชฺชติ, วิราคา วิมุจฺจติ, วิมุตฺตสฺมิํ “วิมุตฺตมฺ”อิติ ญาณํ โหติ, “ขีณา ชาติ, วุสิตํ พฺรหฺมจริยํ, กตํ กรณียํ, นาปรํ อิตฺถตฺตายา”ติ ปชานาติ.\csromanspacer{} ยาวตา ภิกฺขเว สตฺตาวาสา ยาวตา ภวคฺคํ, เอเต อคฺคา เอเต เสฏฺฐา โลกสฺมิํ ยทิทํ อรหนฺโตติ.\csromanspacer{}\csromanspacer{}ปญฺจมํ.}

\csromanheader{2}{6. สีหสุตฺต}
\proseitem{78}{\csromanfolio{70}สาวตฺถินิทานํ.\csromanspacer{} สีโห ภิกฺขเว มิคราชา สายนฺหสมยํ อาสยา นิกฺขมติ, อาสยา นิกฺขมิตฺวา วิชมฺภติ, วิชมฺภิตฺวา สมนฺตา จตุทฺทิสา อนุวิโลเกติ, สมนฺตา จตุทฺทิสา อนุวิโลเกตฺวา ติกฺขตฺตุํ สีหนาทํ นทติ, ติกฺขตฺตุํ สีหนาทํ นทิตฺวา โคจราย ปกฺกมติ.\csromanspacer{} เย หิ เกจิ ภิกฺขเว ติรจฺฉานคตา ปาณา สีหสฺส มิครญฺโญ นทโต สทฺทํ สุณนฺติ, เยภุยฺเยน ภยํ สํเวคํ สนฺตาสํ อาปชฺชนฺติ.\csromanspacer{} พิลํ พิลาสยา ปวิสนฺติ, ทกํ ทกาสยา ปวิสนฺติ, วนํ วนาสยา ปวิสนฺติ, อากาสํ ปกฺขิโน ภชนฺติ.\csromanspacer{} เยปิ เต ภิกฺขเว รญฺโญ นาคา คามนิคมราชธานีสุ ทเฬฺหหิ วรตฺเตหิ พทฺธา, เตปิ ตานิ พนฺธนานิ สญฺฉินฺทิตฺวา สมฺปทาเลตฺวา ภีตา มุตฺตกรีสํ จชมานา\footnote{โมจนฺตา (อิ, ก)} เยน วา เตน วา ปลายนฺติ.\csromanspacer{} เอวํ มหิทฺธิโก โข ภิกฺขเว สีโห มิคราชา ติรจฺฉานคตานํ ปาณานํ เอวํ มเหสกฺโข เอวํ มหานุภาโว.}

\prose{เอวเมว โข ภิกฺขเว ยทา ตถาคโต โลเก อุปฺปชฺชติ อรหํ สมฺมาสมฺพุทฺโธ วิชฺชาจรณสมฺปนฺโน สุคโต โลกวิทู อนุตฺตโร ปุริสทมฺมสารถิ สตฺถา เทวมนุสฺสานํ พุทฺโธ ภควา, โส ธมฺมํ เทเสติ, “อิติ รูปํ, อิติ รูปสฺส สมุทโย, อิติ รูปสฺส อตฺถงฺคโม.\csromanspacer{} อิติ เวทนา.\csromanspacer{} อิติ สญฺญา.\csromanspacer{} อิติ สงฺขารา.\csromanspacer{} อิติ วิญฺญาณํ, อิติ วิญฺญาณสฺส สมุทโย, อิติ วิญฺญาณสฺส อตฺถงฺคโม”ติ.\csromanspacer{} เยปิ เต ภิกฺขเว เทวา ทีฆายุกา วณฺณวนฺโต สุขพหุลา อุจฺเจสุ วิมาเนสุ จิรฏฺฐิติกา, เตปิ ตถาคตสฺส ธมฺมเทสนํ สุตฺวา เยภุยฺเยน ภยํ สํเวคํ สนฺตาสํ อาปชฺชนฺติ อนิจฺจาว กิร โภ มยํ สมานา “นิจฺจมฺหา”ติ อมญฺญิมฺห, อทฺธุวาว กิร โภ มยํ สมานา “ธุวมฺหา”ติ อมญฺญิมฺห, อสสฺสตาว กิร โภ มยํ สมานา “สสฺสตมฺหา”ติ อมญฺญิมฺห, มยมฺปิ กิร โภ อนิจฺจา อทฺธุวา อสสฺสตา สกฺกายปริยาปนฺนาติ.\csromanspacer{} เอวํ มหิทฺธิโก โข ภิกฺขเว ตถาคโต สเทวกสฺส โลกสฺส เอวํ มเหสกฺโข เอวํ มหานุภาโวติ.\csromanspacer{} อิทมโวจ ภควา ฯเปฯ เอตทโวจ สตฺถา\csromandash{}}

\csromangathameasure{“ยทา พุทฺโธ อภิญฺญาย, ธมฺมจกฺกํ ปวตฺตยิ. \\ สเทวกสฺส โลกสฺส, สตฺถา อปฺปฏิปุคฺคโล. \\ สกฺกายญฺจ นิโรธญฺจ, สกฺกายสฺส จ สมฺภวํ. \\ อริยญฺจฏฺฐงฺคิกํ มคฺคํ, ทุกฺขูปสมคามินํ. \\ เยปิ ทีฆายุกา เทวา, วณฺณวนฺโต ยสสฺสิโน. \\ ภีตา สนฺตาสมาปาทุํ, สีหสฺเสวิตเร มิคา. \\ อวีติวตฺตา สกฺกายํ, อนิจฺจา กิร โภ มยํ. \\ สุตฺวา อรหโต วากฺยํ, วิปฺปมุตฺตสฺส ตาทิโน”ติ. ฉฏฺฐํ.}
\csromangathagroup{\csromangathastanza{“ยทา พุทฺโธ อภิญฺญาย, ธมฺมจกฺกํ ปวตฺตยิ. \\ สเทวกสฺส โลกสฺส, สตฺถา อปฺปฏิปุคฺคโล.}\csromangathastanza{\csromanfolio{71}สกฺกายญฺจ นิโรธญฺจ, สกฺกายสฺส จ สมฺภวํ. \\ อริยญฺจฏฺฐงฺคิกํ มคฺคํ, ทุกฺขูปสมคามินํ.}\csromangathastanza{เยปิ ทีฆายุกา เทวา, วณฺณวนฺโต ยสสฺสิโน. \\ ภีตา สนฺตาสมาปาทุํ, สีหสฺเสวิตเร มิคา.}\csromangathastanza{อวีติวตฺตา สกฺกายํ, อนิจฺจา กิร โภ มยํ. \\ สุตฺวา อรหโต วากฺยํ, วิปฺปมุตฺตสฺส ตาทิโน”ติ.\csromanspacer{} ฉฏฺฐํ.}}

\csromanheader{2}{7. ขชฺชนียสุตฺต}
\proseitem{79}{สาวตฺถินิทานํ.\csromanspacer{} เย หิ เกจิ ภิกฺขเว สมณา วา พฺราหฺมณา วา อเนกวิหิตํ ปุพฺเพนิวาสํ อนุสฺสรมานา อนุสฺสรนฺติ, สพฺเพเต ปญฺจุปาทานกฺขนฺเธ อนุสฺสรนฺติ เอเตสํ วา อญฺญตรํ.\csromanspacer{} กตเม ปญฺจ, “เอวํรูโป อโหสิํ อตีตมทฺธานนฺ”ติ, อิติ วา หิ ภิกฺขเว อนุสฺสรมาโน รูปํเยว อนุสฺสรติ. “เอวํเวทโน อโหสิํ อตีตมทฺธานนฺ”ติ, อิติ วา หิ ภิกฺขเว อนุสฺสรมาโน เวทนํเยว อนุสฺสรติ. “เอวํสญฺโญ อโหสิํ อตีตมทฺธานนฺ”ติ. “เอวํสงฺขาโร อโหสิํ อตีตมทฺธานนฺ”ติ. “เอวํวิญฺญาโณ อโหสิํ อตีตมทฺธานนฺ”ติ, อิติ วา หิ ภิกฺขเว อนุสฺสรมาโน วิญฺญาณเมว อนุสฺสรติ.}

\prose{กิญฺจ ภิกฺขเว รูปํ วเทถ, รุปฺปตีติ โข ภิกฺขเว ตสฺมา “รูปนฺ”ติ วุจฺจติ.\csromanspacer{} เกน รุปฺปติ.\csromanspacer{} สีเตนปิ รุปฺปติ, อุเณฺหนปิ รุปฺปติ, ชิฆจฺฉายปิ รุปฺปติ, ปิปาสายปิ รุปฺปติ, ฑํสมกสวาตาตปสรีสปสมฺผสฺเสนปิ\footnote{...สิริํสปสมฺผสฺเสนปิ (สี, อิ)} รุปฺปติ.\csromanspacer{} รุปฺปตีติ โข ภิกฺขเว ตสฺมา “รูปนฺ”ติ วุจฺจติ.\csromanspacer{} กิญฺจ ภิกฺขเว เวทนํ วเทถ, เวทยตีติ โข ภิกฺขเว ตสฺมา “เวทนา”ติ วุจฺจติ.\csromanspacer{} กิญฺจ เวทยติ.\csromanspacer{} สุขมฺปิ เวทยติ, ทุกฺขมฺปิ เวทยติ, อทุกฺขมสุขมฺปิ เวทยติ, เวทยตีติ โข ภิกฺขเว ตสฺมา “เวทนา”ติ วุจฺจติ.}

\prose{กิญฺจ ภิกฺขเว สญฺญํ วเทถ, สญฺชานาตีติ โข ภิกฺขเว ตสฺมา “สญฺญา”ติ วุจฺจติ.\csromanspacer{} กิญฺจ สญฺชานาติ.\csromanspacer{} นีลมฺปิ สญฺชานาติ, ปีตกมฺปิ สญฺชานาติ, โลหิตกมฺปิ สญฺชานาติ, โอทาตมฺปิ สญฺชานาติ.\csromanspacer{} สญฺชานาตีติ โข ภิกฺขเว ตสฺมา “สญฺญา”ติ วุจฺจติ.}

\prose{\csromanfolio{72}กิญฺจ ภิกฺขเว สงฺขาเร วเทถ, สงฺขตมภิสงฺขโรนฺตีติ โข ภิกฺขเว ตสฺมา “สงฺขารา”ติ วุจฺจติ.\csromanspacer{} กิญฺจ สงฺขตมภิสงฺขโรนฺติ.\csromanspacer{} รูปํ รูปตฺตาย\footnote{รูปตฺถาย (ก)} สงฺขตมภิสงฺขโรนฺติ, เวทนํ เวทนตฺตาย สงฺขตมภิสงฺขโรนฺติ, สญฺญํ สญฺญตฺตาย สงฺขตมภิสงฺขโรนฺติ, สงฺขาเร สงฺขารตฺตาย สงฺขตมภิสงฺขโรนฺติ, วิญฺญาณํ วิญฺญาณตฺตาย สงฺขตมภิสงฺขโรนฺติ.\csromanspacer{} สงฺขตมภิสงฺขโรนฺตีติ โข ภิกฺขเว ตสฺมา “สงฺขารา”ติ วุจฺจติ.}

\prose{กิญฺจ ภิกฺขเว วิญฺญาณํ วเทถ, วิชานาตีติ โข ภิกฺขเว ตสฺมา “วิญฺญาณนฺ”ติ วุจฺจติ.\csromanspacer{} กิญฺจ วิชานาติ.\csromanspacer{} อมฺพิลมฺปิ วิชานาติ, ติตฺตกมฺปิ วิชานาติ, กฏุกมฺปิ วิชานาติ, มธุรมฺปิ วิชานาติ, ขาริกมฺปิ วิชานาติ, อขาริกมฺปิ วิชานาติ, โลณิกมฺปิ วิชานาติ, อโลณิกมฺปิ วิชานาติ, วิชานาตีติ โข ภิกฺขเว ตสฺมา “วิญฺญาณนฺ”ติ วุจฺจติ.}

\prose{ตตฺร ภิกฺขเว สุตวา อริยสาวโก อิติ ปฏิสญฺจิกฺขติ “อหํ โข เอตรหิ รูเปน ขชฺชามิ, อตีตมฺปาหํ อทฺธานํ เอวเมว รูเปน ขชฺชิํ, เสยฺยถาปิ เอตรหิ ปจฺจุปฺปนฺเนน รูเปน ขชฺชามิ.\csromanspacer{} อหญฺเจว โข ปน อนาคตํ รูปํ อภินนฺเทยฺยํ, อนาคตมฺปาหํ อทฺธานํ เอวเมว รูเปน ขชฺเชยฺยํ, เสยฺยถาปิ เอตรหิ ปจฺจุปฺปนฺเนน รูเปน ขชฺชามี”ติ.\csromanspacer{} โส อิติ ปฏิสงฺขาย อตีตสฺมิํ รูปสฺมิํ อนเปกฺโข โหติ, อนาคตํ รูปํ นาภินนฺทติ, ปจฺจุปฺปนฺนสฺส รูปสฺส นิพฺพิทาย วิราคาย นิโรธาย ปฏิปนฺโน โหติ.}

\prose{“อหํ โข เอตรหิ เวทนาย ขชฺชามิ, อตีตมฺปาหํ อทฺธานํ เอวเมว เวทนาย ขชฺชิํ, เสยฺยถาปิ เอตรหิ ปจฺจุปฺปนฺนาย เวทนาย ขชฺชามิ.\csromanspacer{} อหญฺเจว โข ปน อนาคตํ เวทนํ อภินนฺเทยฺยํ, อนาคตมฺปาหํ อทฺธานํ เอวเมว เวทนาย ขชฺเชยฺยํ, เสยฺยถาปิ เอตรหิ ปจฺจุปฺปนฺนาย เวทนาย ขชฺชามี”ติ.\csromanspacer{} โส อิติ ปฏิสงฺขาย อตีตาย เวทนาย อนเปกฺโข โหติ, อนาคตํ เวทนํ นาภินนฺทติ, ปจฺจุปฺปนฺนาย เวทนาย นิพฺพิทาย วิราคาย นิโรธาย ปฏิปนฺโน โหติ.}

\prose{“อหํ โข เอตรหิ สญฺญาย ขชฺชามิ ฯเปฯ. “อหํ โข เอตรหิ สงฺขาเรหิ ขชฺชามิ, อตีตมฺปาหํ อทฺธานํ เอวเมว สงฺขาเรหิ ขชฺชิํ, เสยฺยถาปิ เอตรหิ ปจฺจุปฺปนฺเนหิ สงฺขาเรหิ ขชฺชามีติ.\csromanspacer{} อหญฺเจว โข ปน อนาคเต สงฺขาเร \csromanfolio{73}อภินนฺเทยฺยํ, อนาคตมฺปาหํ อทฺธานํ เอวเมว สงฺขาเรหิ ขชฺเชยฺยํ, เสยฺยถาปิ เอตรหิ ปจฺจุปฺปนฺเนหิ สงฺขาเรหิ ขชฺชามี”ติ.\csromanspacer{} โส อิติ ปฏิสงฺขาย อตีเตสุ สงฺขาเรสุ อนเปกฺโข โหติ, อนาคเต สงฺขาเร นาภินนฺทติ, ปจฺจุปฺปนฺนานํ สงฺขารานํ นิพฺพิทาย วิราคาย นิโรธาย ปฏิปนฺโน โหติ.}

\prose{“อหํ โข เอตรหิ วิญฺญาเณน ขชฺชามิ, อตีตมฺปิ อทฺธานํ เอวเมว วิญฺญาเณน ขิชฺชิํ, เสยฺยถาปิ เอตรหิ ปจฺจุปฺปนฺเนน วิญฺญาเณน ขชฺชามิ.\csromanspacer{} อหญฺเจว โข ปน อนาคตํ วิญฺญาณํ อภินนฺเทยฺยํ, อนาคตมฺปาหํ อทฺธานํ เอวเมว วิญฺญาเณน ขชฺเชยฺยํ, เสยฺยถาปิ เอตรหิ ปจฺจุปฺปนฺเนน วิญฺญาเณน ขชฺชามี”ติ.\csromanspacer{} โส อิติ ปฏิสงฺขาย อตีตสฺมิํ วิญฺญาณสฺมิํ อนเปกฺโข โหติ, อนาคตํ วิญฺญาณํ นาภินนฺทติ, ปจฺจุปฺปนฺนสฺส วิญฺญาณสฺส นิพฺพิทาย วิราคาย นิโรธาย ปฏิปนฺโน โหติ.}

\prose{ตํ กิํ มญฺญถ ภิกฺขเว, รูปํ นิจฺจํ วา อนิจฺจํ วาติ.\csromanspacer{} อนิจฺจํ ภนฺเต.\csromanspacer{} ยํ ปนานิจฺจํ, ทุกฺขํ วา ตํ สุขํ วาติ.\csromanspacer{} ทุกฺขํ ภนฺเต.\csromanspacer{} ยํ ปนานิจฺจํ ทุกฺขํ วิปริณามธมฺมํ, กลฺลํ นุ ตํ สมนุปสฺสิตุํ “เอตํ มม, เอโสหมสฺมิ, เอโส เม อตฺตา”ติ.\csromanspacer{} โน เหตํ ภนฺเต.\csromanspacer{} เวทนา.\csromanspacer{} สญฺญา.\csromanspacer{} สงฺขารา.\csromanspacer{} วิญฺญาณํ นิจฺจํ วา อนิจฺจํ วาติ.\csromanspacer{} อนิจฺจํ ภนฺเต.\csromanspacer{} ยํ ปนานิจฺจํ, ทุกฺขํ วา ตํ สุขํ วาติ.\csromanspacer{} ทุกฺขํ ภนฺเต.\csromanspacer{} ยํ ปนานิจฺจํ ทุกฺขํ วิปริณามธมฺมํ, กลฺลํ นุ ตํ สมนุปสฺสิตุํ “เอตํ มม, เอโสหมสฺมิ, เอโส เม อตฺตา”ติ.\csromanspacer{} โน เหตํ ภนฺเต.\csromanspacer{} ตสฺมาติห ภิกฺขเว ยํ กิญฺจิ รูปํ อตีตานาคตปจฺจุปฺปนฺนํ อชฺฌตฺตํ วา พหิทฺธา วา โอฬาริกํ วา สุขุมํ วา หีนํ วา ปณีตํ วา ยํ ทูเร สนฺติเก วา, สพฺพํ รูปํ “เนตํ มม, เนโสหมสฺมิ, น เมโส อตฺตา”ติ เอวเมตํ ยถาภูตํ สมฺมปฺปญฺญาย ทฏฺฐพฺพํ.\csromanspacer{} ยา กาจิ เวทนา.\csromanspacer{} ยากาจิ สญฺญา.\csromanspacer{} เย เกจิ สงฺขารา.\csromanspacer{} ยํ กิญฺจิ วิญฺญาณํ อตีตานาคตปจฺจุปฺปนฺนํ ฯเปฯ ยํ ทูเร สนฺติเก วา สพฺพํ วิญฺญาณํ “เนตํ มม, เนโสหมสฺมิ, น เมโส อตฺตา”ติ เอวเมตํ ยถาภูตํ สมฺมปฺปญฺญาย ทฏฺฐพฺพํ.}

\prose{อยํ วุจฺจติ ภิกฺขเว อริยสาวโก อปจินาติ โน อาจินาติ, ปชหติ น\footnote{โน (สี)} อุปาทิยติ, วิสิเนติ น1 อุสฺสิเนติ, วิธูเปติ น1 สนฺธูเปติ.\csromanspacer{} กิญฺจ อปจินาติ โน อาจินาติ.\csromanspacer{} รูปํ อปจินาติ โน อาจินาติ. \csromanfolio{74}เวทนํ.\csromanspacer{} สญฺญํ.\csromanspacer{} สงฺขาเร.\csromanspacer{} วิญฺญาณํ อปจินาติ โน อาจินาติ.\csromanspacer{} กิญฺจ ปชหติ น อุปาทิยติ.\csromanspacer{} รูปํ ปชหติ น อุปาทิยติ.\csromanspacer{} เวทนํ.\csromanspacer{} สญฺญํ.\csromanspacer{} สงฺขาเร.\csromanspacer{} วิญฺญาณํ ปชหติ น อุปาทิยติ.\csromanspacer{} กิญฺจ วิสิเนติ น อุสฺสิเนติ.\csromanspacer{} รูปํ วิสิเนติ น อุสฺสิเนติ.\csromanspacer{} เวทนํ.\csromanspacer{} สญฺญํ.\csromanspacer{} สงฺขาเร.\csromanspacer{} วิญฺญาณํ วิสิเนติ น อุสฺสิเนติ.\csromanspacer{} กิญฺจ วิธูเปติ น สนฺธูเปติ.\csromanspacer{} รูปํ วิธูเปติ น สนฺธูเปติ.\csromanspacer{} เวทนํ.\csromanspacer{} สญฺญํ.\csromanspacer{} สงฺขาเร.\csromanspacer{} วิญฺญาณํ วิธูเปติ น สนฺธูเปติ.}

\prose{เอวํ ปสฺสํ ภิกฺขเว สุตวา อริยสาวโก รูปสฺมิมฺปิ นิพฺพินฺทติ.\csromanspacer{} เวทนายปิ.\csromanspacer{} สญฺญายปิ.\csromanspacer{} สงฺขาเรสุปิ.\csromanspacer{} วิญฺญาณสฺมิมฺปิ นิพฺพินฺทติ, นิพฺพินฺทํ วิรชฺชติ, วิราคา วิมุจฺจติ, วิมุตฺตสฺมิํ “วิมุตฺตมฺ”อิติ ญาณํ โหติ, “ขีณา ชาติ, วุสิตํ พฺรหฺมจริยํ, กตํ กรณียํ, นาปรํ อิตฺถตฺตายา”ติ ปชานาติ.}

\prose{อยํ วุจฺจติ ภิกฺขเว ภิกฺขุ เนวาจินาติ น อปจินาติ, อปจินิตฺวา ฐิโต เนว ปชหติ น อุปาทิยติ, ปชหิตฺวา ฐิโต เนว วิสิเนติ น อุสฺสิเนติ, วิสิเนตฺวา ฐิโต เนว วิธูเปติ น สนฺธูเปติ.\csromanspacer{} วิธูเปตฺวา ฐิโต กิญฺจ เนวาจินาติ น อปจินาติ, อปจินิตฺวา ฐิโต รูปํ เนวาจินาติ น อปจินาติ.\csromanspacer{} อปจินิตฺวา ฐิโต เวทนํ.\csromanspacer{} สญฺญํ.\csromanspacer{} สงฺขาเร.\csromanspacer{} วิญฺญาณํ เนวาจินาติ น อปจินาติ, อปจินิตฺวา ฐิโต กิญฺจ เนว ปชหติ น อุปาทิยติ, ปชหิตฺวา ฐิโต รูปํ เนว ปชหติ น อุปาทิยติ.\csromanspacer{} ปชหิตฺวา ฐิโต เวทนํ.\csromanspacer{} สญฺญํ.\csromanspacer{} สงฺขาเร.\csromanspacer{} วิญฺญาณํ เนว ปชหติ น อุปาทิยติ, ปชหิตฺวา ฐิโต กิญฺจ เนว วิสิเนติ น อุสฺสิเนติ, วิสิเนตฺวา ฐิโต รูปํ เนว วิสิเนติ น อุสฺสิเนติ.\csromanspacer{} วิสิเนตฺวา ฐิโต เวทนํ.\csromanspacer{} สญฺญํ.\csromanspacer{} สงฺขาเร.\csromanspacer{} วิญฺญาณํ เนว วิสิเนติ น อุสฺสิเนติ, วิสิเนตฺวา ฐิโต กิญฺจ เนว วิธูเปติ น สนฺธูเปติ.\csromanspacer{} วิธูเปตฺวา ฐิโต รูปํ เนว วิธูเปติ น สนฺธูเปติ.\csromanspacer{} วิธูเปตฺวา ฐิโต เวทนํ.\csromanspacer{} สญฺญํ.\csromanspacer{} สงฺขาเร.\csromanspacer{} วิญฺญาณํ เนว วิธูเปติ น สนฺธูเปติ, วิธูเปตฺวา ฐิโต.\csromanspacer{} เอวํวิมุตฺตจิตฺตํ โข ภิกฺขเว ภิกฺขุํ ส-อินฺทา เทวา สพฺรหฺมกา สปชาปติกา อารกาว นมสฺสนฺติ\csromandash{}}

\csromangathameasure{“นโม เต ปุริสาชญฺญ, นโม เต ปุริสุตฺตม. \\ ยสฺส เต นาภิชานาม, ยมฺปิ นิสฺสาย ฌายสี”ติ. สตฺตมํ.}
\csromangathagroup{\csromangathastanza{“นโม เต ปุริสาชญฺญ, นโม เต ปุริสุตฺตม. \\ ยสฺส เต นาภิชานาม, ยมฺปิ นิสฺสาย ฌายสี”ติ.\csromanspacer{} สตฺตมํ.}}

\csromanheader{2}{8. ปิณฺโฑลฺยสุตฺต}
\proseitem{80}{\csromanfolio{75}เอกํ สมยํ ภควา สกฺเกสุ วิหรติ กปิลวตฺถุสฺมิํ นิโคฺรธาราเม.\csromanspacer{} อถ โข ภควา กิสฺมิญฺจิเทว ปกรเณ ภิกฺขุสํฆํ ปณาเมตฺวา ปุพฺพณฺหสมยํ นิวาเสตฺวา ปตฺตจีวรมาทาย กปิลวตฺถุํ ปิณฺฑาย ปาวิสิ, กปิลวตฺถุสฺมิํ ปิณฺฑาย จริตฺวา ปจฺฉาภตฺตํ ปิณฺฑปาตปฏิกฺกนฺโต เยน มหาวนํ เตนุปสงฺกมิ ทิวาวิหาราย, มหาวนํ อชฺโฌคาเหตฺวา เพลุวลฏฺฐิกาย มูเล ทิวาวิหารํ นิสีทิ.}

\prose{อถ โข ภควโต รโหคตสฺส ปฏิสลฺลีนสฺส เอวํ เจตโส ปริวิตกฺโก อุทปาทิ “มยา โข ภิกฺขุสํโฆ ปพาโฬฺห, สนฺเตตฺถ ภิกฺขู นวา อจิรปพฺพชิตา อธุนาคตา อิมํ ธมฺมวินยํ, เตสํ มมํ อปสฺสนฺตานํ สิยา อญฺญถตฺตํ สิยา วิปริณาโม.\csromanspacer{} เสยฺยถาปิ นาม วจฺฉสฺส ตรุณสฺส มาตรํ อปสฺสนฺตสฺส สิยา อญฺญถตฺตํ สิยา วิปริณาโม.\csromanspacer{} เอวเมว สนฺเตตฺถ ภิกฺขู นวา อจิรปพฺพชิตา อธุนาคตา อิมํ ธมฺมวินยํ, เตสํ มมํ อปสฺสนฺตานํ สิยา อญฺญถตฺตํ สิยา วิปริณาโม.\csromanspacer{} เสยฺยถาปิ นาม พีชานํ ตรุณานํ อุทกํ อลภนฺตานํ สิยา อญฺญถตฺตํ สิยา วิปริณาโม.\csromanspacer{} เอวเมว สนฺเตตฺถ ฯเปฯ เตสํ มมํ อลภนฺตานํ ทสฺสนาย สิยา อญฺญถตฺตํ สิยา วิปริณาโม.\csromanspacer{} ยํนูนาหํ ยเถว มยา ปุพฺเพ ภิกฺขุสํโฆ อนุคฺคหิโต, เอวเมว เอตรหิ อนุคฺคเณฺหยฺยํ ภิกฺขุสํฆนฺ”ติ.}

\prose{อถ โข พฺรหฺมา สหมฺปติ ภควโต เจตสา เจโตปริวิตกฺกมญฺญาย เสยฺยถาปิ นาม พลวา ปุริโส สมิญฺชิตํ\footnote{สมฺมิญฺชิตํ (สี, สฺยา, กํ, อิ)} วา พาลํ พาหํ ปสาเรยฺย, ปสาริตํ วา พาหํ สมิญฺเชยฺย.\csromanspacer{} เอวเมว พฺรหฺมโลเก อนฺตรหิโต ภควโต ปุรโต ปาตุรโหสิ.\csromanspacer{} อถ โข พฺราหฺมา สหมฺปติ เอกํสํ อุตฺตราสงฺคํ กริตฺวา เยน ภควา เตนญฺชลิํ ปณาเมตฺวา ภควนฺตํ เอตทโวจ “เอวเมตํ ภควา, เอวเมตํ สุคต, ภควตา ภนฺเต ภิกฺขุสํโฆ ปพาโฬฺห, สนฺเตตฺถ ภิกฺขู นวา อจิรปพฺพชิตา อธุนาคตา อิมํ ธมฺมวินยํ, เตสํ ภควนฺตํ อปสฺสนฺตานํ สิยา อญฺญถตฺตํ สิยา วิปริณาโม.\csromanspacer{} เสยฺยถาปิ นาม วจฺฉสฺส ตรุณสฺส มาตรํ อปสฺสนฺตสฺส สิยา อญฺญถตฺตํ สิยา วิปริณาโม.\csromanspacer{} เอวเมว สนฺเตตฺถ ภิกฺขู นวา \csromanfolio{76}อจิรปพฺพชิตา อธุนาคตา อิมํ ธมฺมวินยํ, เตสํ ภควนฺตํ อปสฺสนฺตานํ สิยา อญฺญถตฺตํ สิยา วิปริณาโม.\csromanspacer{} เสยฺยถาปิ นาม พีชานํ ตรุณานํ อุทกํ อลภนฺตานํ สิยา อญฺญถตฺตํ สิยา วิปริณาโม.\csromanspacer{} เอวเมว สนฺเตตฺถ ภิกฺขู นวา อจิรปพฺพชิตา อธุนาคตา อิมํ ธมฺมวินยํ, เตสํ ภควนฺตํ อลภนฺตานํ ทสฺสนาย สิยา อญฺญถตฺตํ สิยา วิปริณาโม.\csromanspacer{} อภินนฺทตุ ภนฺเต ภควา ภิกฺขุสํฆํ, อภิวทตุ ภนฺเต ภควา ภิกฺขุสํฆํ.\csromanspacer{} ยเถว ภควตา ปุพฺเพ ภิกฺขุสํโฆ อนุคฺคหิโต, เอวเมว เอตรหิ อนุคฺคณฺหาตุ ภิกฺขุสํฆนฺติ.}

\prose{อธิวาเสสิ ภควา ตุณฺหีภาเวน.\csromanspacer{} อถ โข พฺรหฺมา สหมฺปติ ภควโต อธิวาสนํ วิทิตฺวา ภควนฺตํ อภิวาเทตฺวา ปทกฺขิณํ กตฺวา ตตฺเถวนฺตรธายิ.}

\prose{อถ โข ภควา สายนฺหสมยํ ปฏิสลฺลานา วุฏฺฐิโต เยน นิโคฺรธาราโม เตนุปสงฺกมิ, อุปสงฺกมิตฺวา ปญฺญตฺเต อาสเน นิสีทิ, นิสชฺช โข ภควา ตถารูปํ อิทฺธาภิสงฺขารํ อภิสงฺขาสิ\footnote{อภิสงฺขาเรสิ (สฺยา, กํ), อภิสงฺขายิ (อิ), อภิสงฺขโรติ (ก)}, ยถา เต ภิกฺขู (เอกทฺวีหิกาย สารชฺชมานรูปา เยนาหํ\footnote{เยน ภควา (?)} เตนุปสงฺกเมยฺยุํ.\csromanspacer{} เตปิ ภิกฺขู)\footnote{( ) สี-สฺยา-กํ-โปตฺถเกสุ นตฺถิ.} เอกทฺวีหิกาย สารชฺชมานรูปา เยน ภควา เตนุปสงฺกมิํสุ, อุปสงฺกมิตฺวา ภควนฺตํ อภิวาเทตฺวา เอกมนฺตํ นิสีทิํสุ, เอกมนฺตํ นิสินฺเน โข เต ภิกฺขู ภควา เอตทโวจ\csromandash{}}

\prose{อนฺตมิทํ ภิกฺขเว ชีวิกานํ ยทิทํ ปิณฺโฑลฺยํ, อภิสาโปยํ ภิกฺขเว โลกสฺมิํ “ปิณฺโฑโล วิจรสิ ปตฺตปาณี”ติ, ตญฺจ โข เอตํ ภิกฺขเว กุลปุตฺตา อุเปนฺติ อตฺถวสิกา อตฺถวสํ ปฏิจฺจ, เนว ราชาภินีตา น โจราภินีตา น อิณฏฺฏา น ภยฏฺฏา น อาชีวิกาปกตา, อปิ จ โข “โอติณฺณามฺห ชาติยา ชราย มรเณน โสเกหิ ปริเทเวหิ ทุกฺเขหิ โทมนสฺเสหิ อุปายาเสหิ ทุกฺโขติณฺณา ทุกฺขปเรตา, อปฺเปว นาม อิมสฺส เกวลสฺส ทุกฺขกฺขนฺธสฺส อนฺตกิริยา ปญฺญาเยถา”ติ.}

\prose{เอวํ ปพฺพชิโต จายํ ภิกฺขเว กุลปุตฺโต, โส จ โหติ อภิชฺฌาลุ กาเมสุ ติพฺพสาราโค พฺยาปนฺนจิตฺโต ปทุฏฺฐมนสงฺกปฺโป มุฏฺฐสฺสติ อสมฺปชาโน อสมาหิโต วิพฺภนฺตจิตฺโต ปากตินฺทฺริโย.\csromanspacer{} เสยฺยถาปิ \csromanfolio{77}ภิกฺขเว ฉวาลาตํ อุภโตปทิตฺตํ มชฺเฌ คูถคตํ เนว คาเม กฏฺฐตฺถํ ผรติ, นารญฺเญ กฏฺฐตฺถํ ผรติ.\csromanspacer{} ตถูปมาหํ ภิกฺขเว อิมํ ปุคฺคลํ วทามิ คิหิโภคา จ ปริหีโน, สามญฺญตฺถญฺจ น ปริปูเรติ.}

\prose{ตโยเม ภิกฺขเว อกุสลวิตกฺกา.\csromanspacer{} กามวิตกฺโก พฺยาปาทวิตกฺโก วิหิํสาวิตกฺโก, อิเม จ ภิกฺขเว ตโย อกุสลวิตกฺกา กฺว อปริเสสา นิรุชฺฌนฺติ, จตูสุ วา สติปฏฺฐาเนสุ สุปฺปติฏฺฐิตจิตฺตสฺส วิหรโต อนิมิตฺตํ วา สมาธิํ ภาวยโต.\csromanspacer{} ยาวญฺจิทํ ภิกฺขเว อลเมว อนิมิตฺโต สมาธิ ภาเวตุํ.\csromanspacer{} อนิมิตฺโต ภิกฺขเว สมาธิ ภาวิโต พหุลีกโต มหปฺผโล โหติ มหานิสํโส.}

\prose{เทฺวมา ภิกฺขเว ทิฏฺฐิโย.\csromanspacer{} ภวทิฏฺฐิ จ วิภวทิฏฺฐิ จ.\csromanspacer{} ตตฺร โข ภิกฺขเว สุตวา อริยสาวโก อิติ ปฏิสญฺจิกฺขติ “อตฺถิ นุ โข ตํ กิญฺจิ โลกสฺมิํ, ยมหํ อุปาทิยมาโน น วชฺชวา อสฺสนฺ”ติ.\csromanspacer{} โส เอวํ ปชานาติ “นตฺถิ นุ โข ตํ กิญฺจิ โลกสฺมิํ, ยมหํ อุปาทิยมาโน น วชฺชวา อสฺสํ, อหํ หิ รูปญฺเญว อุปาทิยมาโน อุปาทิเยยฺยํ.\csromanspacer{} เวทนญฺเญว.\csromanspacer{} สญฺญญฺเญว.\csromanspacer{} สงฺขาเรเยว.\csromanspacer{} วิญฺญาณญฺเญว อุปาทิยมาโน อุปาทิเยยฺยํ.\csromanspacer{} ตสฺส เม อสฺส\footnote{อยํ (ก)} อุปาทานปจฺจยา ภโว, ภวปจฺจยา ชาติ, ชาติปจฺจยา ชารามรณํ โสกปริเทวทุกฺขโทมนสฺสุปายาสา สมฺภเวยฺยุํ, เอวเมตสฺส เกวลสฺส ทุกฺขกฺขนฺธสฺส สมุทโย อสฺสา”ติ.}

\prose{ตํ กิํ มญฺญถ ภิกฺขเว, รูปํ นิจฺจํ วา อนิจฺจํ วาติ.\csromanspacer{} อนิจฺจํ ภนฺเต.\csromanspacer{} ยํ ปนานิจฺจํ, ทุกฺขํ วา ตํ สุขํ วาติ.\csromanspacer{} ทุกฺขํ ภนฺเต.\csromanspacer{} ยํ ปนานิจฺจํ ทุกฺขํ วิปริณามธมฺมํ, กลฺลํ นุ ตํ สมนุปสฺสิตุํ “เอตํ มม, เอโสหมสฺมิ, เอโส เม อตฺตา”ติ.\csromanspacer{} โน เหตํ ภนฺเต.\csromanspacer{} เวทนา.\csromanspacer{} สญฺญา.\csromanspacer{} สงฺขารา.\csromanspacer{} วิญฺญาณํ ฯเปฯ.\csromanspacer{} ตสฺมาติห ภิกฺขเว เอวํ ปสฺสํ ฯเปฯ นาปรํ อิตฺถตฺตายาติ ปชานาตีติ.\csromanspacer{}\csromanspacer{}อฏฺฐมํ.}

\csromanheader{2}{9. ปาลิเลยฺยสุตฺต}
\proseitem{81}{เอกํ สมยํ ภควา โกสมฺพิยํ วิหรติ โฆสิตาราเม.\csromanspacer{} อถ โข ภควา ปุพฺพณฺหสมยํ นิวาเสตฺวา ปตฺตจีวรมาทาย โกสมฺพิํ ปิณฺฑาย ปาวิสิ, โกสมฺพิยํ ปิณฺฑาย จริตฺวา ปจฺฉาภตฺตํ ปิณฺฑปาตปฏิกฺกนฺโต สามํ เสนาสนํ สํสาเมตฺวา ปตฺตจีวรมาทาย \csromanfolio{78}อนามนฺเตตฺวา อุปฏฺฐาเก อนปโลเกตฺวา ภิกฺขุสํฆํ เอโก อทุติโย จาริกํ ปกฺกามิ.}

\prose{อถ โข อญฺญตโร ภิกฺขุ อจิรปกฺกนฺตสฺส ภควโต เยนายสฺมา อานนฺโท เตนุปสงฺกมิ, อุปสงฺกมิตฺวา อายสฺมนฺตํ อานนฺทํ เอตทโวจ “เอสาวุโส อานนฺท ภควา สามํ เสนาสนํ สํสาเมตฺวา ปตฺตจีวรมาทาย อนามนฺเตตฺวา อุปฏฺฐาเก อนปโลเกตฺวา ภิกฺขุสํฆํ เอโก อทุติโย จาริกํ ปกฺกนฺโต”ติ.\csromanspacer{} ยสฺมิํ อาวุโส สมเย ภควา สามํ เสนาสนํ สํสาเมตฺวา ปตฺตจีวรมาทาย อนามนฺเตตฺวา อุปฏฺฐาเก อนปโลเกตฺวา ภิกฺขุสํฆํ เอโก อทุติโย จาริกํ ปกฺกมติ, เอโกว ภควา ตสฺมิํ สมเย วิหริตุกาโม โหติ, น ภควา ตสฺมิํ สมเย เกนจิ อนุพนฺธิตพฺโพ โหตีติ.}

\prose{อถ โข ภควา อนุปุพฺเพน จาริกํ จรมาโน เยน ปาลิเลยฺยกํ\footnote{ปาริเลยฺยกํ (สี, อิ)} ตทวสริ, ตตฺร สุทํ ภควา ปาลิเลยฺยเก วิหรติ ภทฺทสาลมูเล.\csromanspacer{} อถ โข สมฺพหุลา ภิกฺขู เยนายสฺมา อานนฺโท เตนุปสงฺกมิํสุ, อุปสงฺกมิตฺวา อายสฺมตา อานนฺเทน สทฺธิํ สมฺโมทิํสุ, สมฺโมทนียํ กถํ สารณียํ วีติสาเรตฺวา เอกมนฺตํ นิสีทิํสุ, เอกมนฺตํ นิสินฺนา โข เต ภิกฺขู อายสฺมนฺตํ อานนฺทํ เอตทโวจุํ “จิรสฺสุตา โข โน อาวุโส อานนฺท ภควโต สมฺมุขา ธมฺมี กถา, อิจฺฉาม มยํ อาวุโส อานนฺท ภควโต สมฺมุขา ธมฺมิํ กถํ โสตุนฺ”ติ.}

\prose{อถ โข อายสฺมา อานนฺโท เตหิ ภิกฺขูหิ สทฺธิํ เยน ปาลิเลยฺยกํ ภทฺทสาลมูลํ เยน ภควา เตนุปสงฺกมิ, อุปสงฺกมิตฺวา ภควนฺตํ อภิวาเทตฺวา เอกมนฺตํ นิสีทิ, เอกมนฺตํ นิสินฺเน โข เต ภิกฺขู ภควา ธมฺมิยา กถาย สนฺทสฺเสสิ สมาทเปสิ สมุตฺเตเชสิ สมฺปหํเสสิ.\csromanspacer{} เตน โข ปน สมเยน อญฺญตรสฺส ภิกฺขุโน เอวํ เจตโส ปริวิตกฺโก อุทปาทิ “กถํ นุ โข ชานโต กถํ ปสฺสโต อนนฺตรา อาสวานํ ขโย โหตี”ติ.\csromanspacer{} อถ โข ภควา ตสฺส ภิกฺขุโน เจตสา เจโตปริวิตกฺกมญฺญาย ภิกฺขู อามนฺเตสิ “วิจยโส เทสิโต ภิกฺขเว มยา ธมฺโม, วิจยโส เทสิตา จตฺตาโร สติปฏฺฐานา, วิจยโส เทสิตา จตฺตาโร สมฺมปฺปธานา, \csromanfolio{79}วิจยโส เทสิตา จตฺตาโร อิทฺธิปาทา, วิจยโส เทสิตานิ ปญฺจินฺทฺริยานิ, วิจยโส เทสิตานิ ปญฺจ พลานิ, วิจยโส เทสิตา สตฺต โพชฺฌงฺคา, วิจยโส เทสิโต อริโย อฏฺฐงฺคิโก มคฺโค, เอวํ วิจยโส เทสิโต ภิกฺขเว มยา ธมฺโม.\csromanspacer{} เอวํ วิจยโส เทสิเต โข ภิกฺขเว มยา ธมฺเม อถ จ ปนิเธกจฺจสฺส ภิกฺขุโน เอวํ เจตโส ปริวิตกฺโก อุทปาทิ “กถํ นุ โข ชานโต กถํ ปสฺสโต อนนฺตรา อาสวานํ ขโย โหตี”ติ.}

\prose{กถญฺจ ภิกฺขเว ชานโต กถํ ปสฺสโต อนนฺตรา อาสวานํ ขโย โหติ, อิธ ภิกฺขเว อสฺสุตวา ปุถุชฺชโน อริยานํ อทสฺสาวี อริยธมฺมสฺส อโกวิโท อริยธมฺเม อวินีโต, สปฺปุริสานํ อทสฺสาวี สปฺปุริสธมฺมสฺส อโกวิโท สปฺปุริสธมฺเม อวินีโต รูปํ อตฺตโต สมนุปสฺสติ, ยา โข ปน สา ภิกฺขเว สมนุปสฺสนา, สงฺขาโร โส.\csromanspacer{} โส ปน สงฺขาโร กิํนิทาโน กิํสมุทโย กิํชาติโก กิํปภโว.\csromanspacer{} อวิชฺชาสมฺผสฺสเชน ภิกฺขเว เวทยิเตน ผุฏฺฐสฺส อสฺสุตวโต ปุถุชฺชนสฺส อุปฺปนฺนา ตณฺหา, ตโตโช โส สงฺขาโร.\csromanspacer{} อิติ โข ภิกฺขเว โสปิ สงฺขาโร อนิจฺโจ สงฺขโต ปฏิจฺจสมุปฺปนฺโน, สาปิ ตณฺหา อนิจฺจา สงฺขตา ปฏิจฺจสมุปฺปนฺนา.\csromanspacer{} สาปิ เวทนา, โสปิ ผสฺโส อนิจฺโจ สงฺขโต ปฏิจฺจสมุปฺปนฺโน, สาปิ อวิชฺชา อนิจฺจา สงฺขตา ปฏิจฺจสมุปฺปนฺนา.\csromanspacer{} เอวมฺปิ โข ภิกฺขเว ชานโต เอวํ ปสฺสโต อนนฺตรา อาสวานํ ขโย โหติ.}

\prose{น เหว โข รูปํ อตฺตโต สมนุปสฺสติ, อปิ จ โข รูปวนฺตํ อตฺตานํ สมนุปสฺสติ.\csromanspacer{} ยา โข ปน สา ภิกฺขเว สมนุปสฺสนา, สงฺขาโร โส.\csromanspacer{} โส ปน สงฺขาโร กิํนิทาโน กิํสมุทโย กิํชาติโก กิํปภโว.\csromanspacer{} อวิชฺชาสมฺผสฺสเชน ภิกฺขเว เวทยิเตน ผุฏฺฐสฺส อสฺสุตวโต ปุถุชฺชนสฺส อุปฺปนฺนา ตณฺหา, ตโตโช โส สงฺขาโร.\csromanspacer{} อิติ โข ภิกฺขเว โสปิ สงฺขาโร อนิจฺโจ สงฺขโต ปฏิจฺจสมุปฺปนฺโน, สาปิ ตณฺหา, สาปิ เวทนา, โสปิ ผสฺโส, สาปิ อวิชฺชา อนิจฺจา สงฺขตา ปฏิจฺจสมุปฺปนฺนา.\csromanspacer{} เอวมฺปิ โข ภิกฺขเว ชานโต เอวํ ปสฺสโต อนนฺตรา อาสวานํ ขโย โหติ.}

\prose{น เหว โข รูปํ อตฺตโต สมนุปสฺสติ, น รูปวนฺตํ อตฺตานํ สมนุปสฺสติ, อปิ จ โข อตฺตนิ รูปํ สมนุปสฺสติ.\csromanspacer{} ยา โข ปน สา \csromanfolio{80}ภิกฺขเว สมนุปสฺสนา, สงฺขาโร โส.\csromanspacer{} โส ปน สงฺขาโร กิํนิทาโน กิํสมุทโย กิํชาติโก กิํปภโว.\csromanspacer{} อวิชฺชาสมฺผสฺสเชน ภิกฺขเว เวทยิเตน ผุฏฺฐสฺส อสฺสุตวโต ปุถุชฺชนสฺส อุปฺปนฺนา ตณฺหา, ตโตโช โส สงฺขาโร.\csromanspacer{} อิติ โข ภิกฺขเว โสปิ สงฺขาโร อนิจฺโจ สงฺขโต ปฏิจฺจสมุปฺปนฺโน.\csromanspacer{} สาปิ ตณฺหา, สาปิ เวทนา, โสปิ ผสฺโส, สาปิ อวิชฺชา อนิจฺจา สงฺขตา ปฏิจฺจสมุปฺปนฺนา.\csromanspacer{} เอวมฺปิ โข ภิกฺขเว ชานโต เอวํ ปสฺสโต อนนฺตรา อาสวานํ ขโย โหติ.}

\prose{น เหว โข รูปํ อตฺตโต สมนุปสฺสติ, น รูปวนฺตํ อตฺตานํ สมนุปสฺสติ, น อตฺตนิ รูปํ สมนุปสฺสติ, อปิ จ โข รูปสฺมิํ อตฺตานํ สมนุปสฺสติ.\csromanspacer{} ยา โข ปน สา ภิกฺขเว สมนุปสฺสนา, สงฺขาโร โส.\csromanspacer{} โส ปน สงฺขาโร กิํนิทาโน กิํสมุทโย กิํชาติโก กิํปภโว, อวิชฺชาสมฺผสฺสเชน ภิกฺขเว เวทยิเตน ผุฏฺฐสฺส อสฺสุตวโต ปุถุชฺชนสฺส อุปฺปนฺนา ตณฺหา, ตโตโช โส สงฺขาโร.\csromanspacer{} อิติ โข ภิกฺขเว โสปิ สงฺขาโร อนิจฺโจ สงฺขโต ปฏิจฺจสมุปฺปนฺโน.\csromanspacer{} สาปิ ตณฺหา, สาปิ เวทนา, โสปิ ผสฺโส, สาปิ อวิชฺชา อนิจฺจา สงฺขตา ปฏิจฺจสมุปฺปนฺนา.\csromanspacer{} เอวมฺปิ โข ภิกฺขเว ชานโต ฯเปฯ อาสวานํ ขโย โหติ.}

\prose{น เหว โข รูปํ อตฺตโต สมนุปสฺสติ, น รูปวนฺตํ อตฺตานํ, น อตฺตนิ รูปํ, น รูปสฺมิํ อตฺตานํ สมนุปสฺสติ, อปิ จ โข เวทนํ อตฺตโต สมนุปสฺสติ, อปิ จ โข เวทนาวนฺตํ อตฺตานํ สมนุปสฺสติ, อปิ จ โข อตฺตนิ เวทนํ สมนุปสฺสติ, อปิ จ โข เวทนาย อตฺตานํ สมนุปสฺสติ.\csromanspacer{} อปิ จ โข สญฺญํ.\csromanspacer{} อปิ จ โข สงฺขาเร อตฺตโต สมนุปสฺสติ, อปิ จ โข สงฺขารวนฺตํ อตฺตานํ สมนุปสฺสติ, อปิ จ โข อตฺตนิ สงฺขาเร สมนุปสฺสติ, อปิ จ โข สงฺขาเรสุ อตฺตานํ สมนุปสฺสติ.\csromanspacer{} อปิ จ โข วิญฺญาณํ อตฺตโต สมนุปสฺสติ, อปิ จ โข วิญฺญาณวนฺตํ อตฺตานํ, อปิ จ โข อตฺตนิ วิญฺญาณํ, อปิ จ โข วิญฺญาณสฺมิํ อตฺตานํ สมนุปสฺสติ.\csromanspacer{} ยา โข ปน สา ภิกฺขเว สมนุปสฺสนา, สงฺขาโร โส.\csromanspacer{} โส ปน สงฺขาโร กิํนิทาโน ฯเปฯ กิํปภโว.\csromanspacer{} อวิชฺชาสมฺผสฺสเชน ภิกฺขเว เวทยิเตน ผุฏฺฐสฺส อสฺสุตวโต ปุถุชฺชนสฺส อุปฺปนฺนา ตณฺหา, ตโตโช โส สงฺขาโร.\csromanspacer{} อิติ โข ภิกฺขเว โสปิ สงฺขาโร อนิจฺโจ สงฺขโต ปฏิจฺจสมุปฺปนฺโน.\csromanspacer{} สาปิ ตณฺหา, สาปิ เวทนา, โสปิ \csromanfolio{81}ผสฺโส, สาปิ อวิชฺชา อนิจฺจา สงฺขตา ปฏิจฺจสมุปฺปนฺนา.\csromanspacer{} เอวํ โข ภิกฺขเว ชานโต เอวํ ปสฺสโต อนนฺตรา อาสวานํ ขโย โหติ.}

\prose{น เหว โข รูปํ อตฺตโต สมนุปสฺสติ, น เวทนํ อตฺตโต สมนุปสฺสติ, น สญฺญํ.\csromanspacer{} น สงฺขาเร.\csromanspacer{} น วิญฺญาณํ อตฺตโต สมนุปสฺสติ, อปิ จ โข เอวํทิฏฺฐิ โหติ “โส อตฺตา, โส โลโก, โส เปจฺจ ภวิสฺสามิ นิจฺโจ ธุโว สสฺสโต อวิปริณามธมฺโม”ติ.\csromanspacer{} ยา โข ปน สา ภิกฺขเว สสฺสตทิฏฺฐิ, สงฺขาโร โส.\csromanspacer{} โส ปน สงฺขาโร กิํนิทาโน ฯเปฯ.\csromanspacer{} เอวมฺปิ โข ภิกฺขเว ชานโต เอวํ ปสฺสโต อนนฺตรา อาสวานํ ขโย โหติ.}

\prose{น เหว โข รูปํ อตฺตโต สมนุปสฺสติ, น เวทนํ.\csromanspacer{} น สญฺญํ.\csromanspacer{} น สงฺขาเร.\csromanspacer{} น วิญฺญาณํ อตฺตโต สมนุปสฺสติ, นาปิ เอวํทิฏฺฐิ โหติ “โส อตฺตา, โส โลโก, โส เปจฺจ ภวิสฺสามิ นิจฺโจ ธุโว สสฺสโต อวิปริณามธมฺโม”ติ, อปิ จ โข เอวํทิฏฺฐิ โหติ “โน จสฺสํ, โน จ เม สิยา, นาภวิสฺสํ, น เม ภวิสฺสตี”ติ.\csromanspacer{} ยา โข ปน สา ภิกฺขเว อุจฺเฉททิฏฺฐิ, สงฺขาโร โส.\csromanspacer{} โส ปน สงฺขาโร กิํนิทาโน กิํสมุทโย กิํชาติโก กิํปภโว.\csromanspacer{} อวิชฺชาสมฺผสฺสเชน ภิกฺขเว เวทยิเตน ผุฏฺฐสฺส อสฺสุตวโต ปุถุชฺชนสฺส อุปฺปนฺนา ตณฺหา, ตโตโช โส สงฺขาโร.\csromanspacer{} อิติ โข ภิกฺขเว โสปิ สงฺขาโร อนิจฺโจ ฯเปฯ.\csromanspacer{} เอวมฺปิ โข ภิกฺขเว ชานโต เอวํ ปสฺสโต อนนฺตรา อาสวานํ ขโย โหติ.}

\prose{น เหว โข รูปํ อตฺตโต สมนุปสฺสติ, น เวทนํ.\csromanspacer{} น สญฺญํ.\csromanspacer{} น สงฺขาเร.\csromanspacer{} น วิญฺญาณํ อตฺตโต สมนุปสฺสติ ฯเปฯ.\csromanspacer{} น วิญฺญาณสฺมิํ อตฺตโต สมนุปสฺสติ, นาปิ เอวํทิฏฺฐิ โหติ “โส อตฺตา, โส โลโก, โส เปจฺจ ภวิสฺสามิ นิจฺโจ ธุโว สสฺสโต อวิปริณามธมฺโม”ติ.\csromanspacer{} นาปิ เอวํทิฏฺฐิ โหติ “โน จสฺสํ, โน จ เม สิยา, นาภวิสฺสํ, น เม ภวิสฺสตี”ติ, อปิ จ โข กงฺขี โหติ วิจิกิจฺฉี อนิฏฺฐงฺคโต สทฺธมฺเม.\csromanspacer{} ยา โข ปน สา ภิกฺขเว กงฺขิตา วิจิกิจฺฉิตา อนิฏฺฐงฺคตตา สทฺธมฺเม, สงฺขาโร โส.\csromanspacer{} โส ปน สงฺขาโร กิํนิทาโน กิํสมุทโย กิํชาติโก กิํปภโว.\csromanspacer{} อวิชฺชาสมฺผสฺสเชน ภิกฺขเว เวทยิเตน ผุฏฺฐสฺส อสฺสุตวโต ปุถุชฺชนสฺส อุปฺปนฺนา ตณฺหา, ตโตโช โส สงฺขาโร.\csromanspacer{} อิติ โข ภิกฺขเว โสปิ สงฺขาโร อนิจฺโจ \csromanfolio{82}สงฺขโต ปฏิจฺจสมุปฺปนฺโน.\csromanspacer{} สาปิ ตณฺหา อนิจฺจา สงฺขตา ปฏิจฺจสมุปฺปนฺนา.\csromanspacer{} สาปิ เวทนา อนิจฺจา สงฺขตา ปฏิจฺจสมุปฺปนฺนา.\csromanspacer{} โสปิ ผสฺโส อนิจฺโจ สงฺขโต ปฏิจฺจสมุปฺปนฺโน.\csromanspacer{} สาปิ อวิชฺชา อนิจฺจา สงฺขตา ปฏิจฺจสมุปฺปนฺนา.\csromanspacer{} เอวํ โข ภิกฺขเว ชานโต เอวํ ปสฺสโต อนนฺตรา อาสวานํ ขโย โหตีติ.\csromanspacer{}\csromanspacer{}นวมํ.}

\csromanheader{2}{10. ปุณฺณมสุตฺต}
\proseitem{82}{เอกํ สมยํ ภควา สาวตฺถิยํ วิหรติ ปุพฺพาราเม มิคารมาตุปาสาเท มหตา ภิกฺขุสํเฆน สทฺธิํ.\csromanspacer{} เตน โข ปน สมเยน ภควา ตทหุโปสเถ ปนฺนรเส ปุณฺณาย ปุณฺณมาย รตฺติยา ภิกฺขุสํฆปริวุโต อชฺโฌกาเส นิสินฺโน โหติ.}

\prose{อถ โข อญฺญตโร ภิกฺขุ อุฏฺฐายาสนา เอกํสํ อุตฺตราสงฺคํ กริตฺวา เยน ภควา เตนญฺชลิํ ปณาเมตฺวา ภควนฺตํ เอตทโวจ “ปุจฺเฉยฺยาหํ ภนฺเต ภควนฺตํ กิญฺจิเทว\footnote{กญฺจิเทว (?)} เทสํ, สเจ เม ภควา โอกาสํ กโรติ ปญฺหสฺส เวยฺยากรณายา”ติ.\csromanspacer{} เตน หิ ตฺวํ ภิกฺขุ สเก อาสเน นิสีทิตฺวา ปุจฺฉ ยทากงฺขสีติ. “เอวํ ภนฺเต”ติ โข โส ภิกฺขุ ภควโต ปฏิสฺสุตฺวา สเก อาสเน นิสีทิตฺวา ภควนฺตํ เอตทโวจ “อิเม นุ โข ภนฺเต ปญฺจุปาทานกฺขนฺธา.\csromanspacer{} เสยฺยถิทํ, รูปุปาทานกฺขนฺโธ เวทนุปาทานกฺขนฺโธ สญฺญุปาทานกฺขนฺโธ สงฺขารุปาทานกฺขนฺโธ วิญฺญาณุปาทานกฺขนฺโธติ.\csromanspacer{} อิเม โข ภิกฺขุ ปญฺจุปาทานกฺขนฺธา.\csromanspacer{} เสยฺยถิทํ, รูปุปาทานกฺขนฺโธ ฯเปฯ วิญฺญาณุปาทานกฺขนฺโธติ. “สาธุ ภนฺเต”ติ โข ภิกฺขุ ภควโต ภาสิตํ อภินนฺทิตฺวา อนุโมทิตฺวา ภควนฺตํ อุตฺตริํ ปญฺหํ อปุจฺฉิ\csromandash{}}

\prose{อิเม โข ปน ภนฺเต ปญฺจุปาทานกฺขนฺธา กิํมูลกาติ.\csromanspacer{} อิเม โข ภิกฺขุ ปญฺจุปาทานกฺขนฺธา ฉนฺทมูลกาติ ฯเปฯ.\csromanspacer{} ตญฺเญว นุ โข ภนฺเต อุปาทานํ เต ปญฺจุปาทานกฺขนฺธา, อุทาหุ อญฺญตฺร ปญฺจหิ อุปาทานกฺขนฺเธหิ อุปาทานนฺติ.\csromanspacer{} น โข ภิกฺขุ ตญฺเญว อุปาทานํ เต ปญฺจุปาทานกฺขนฺธา, นาปิ อญฺญตฺร ปญฺจหิ อุปาทานกฺขนฺเธหิ อุปาทานํ, อปิ จ โย ตตฺถ ฉนฺทราโค, ตํ ตตฺถ อุปาทานนฺติ. “สาธุ ภนฺเต”ติ โข โส ภิกฺขุ ฯเปฯ อุตฺตริํ ปญฺหํ อปุจฺฉิ\csromandash{}}

\prose{\csromanfolio{83}สิยา ปน ภนฺเต ปญฺจุปาทานกฺขนฺเธสุ ฉนฺทราคเวมตฺตตาติ. “สิยา ภิกฺขู”ติ ภควา อโวจ, อิธ ภิกฺขุ เอกจฺจสฺส เอวํ โหติ “เอวํรูโป สิยํ อนาคตมทฺธานํ, เอวํเวทโน สิยํ อนาคตมทฺธานํ, เอวํสญฺโญ สิยํ อนาคตมทฺธานํ, เอวํสงฺขาโร สิยํ อนาคตมทฺธานํ, เอวํวิญฺญาโณ สิยํ อนาคตมทฺธานนฺ”ติ.\csromanspacer{} เอวํ โข ภิกฺขุ สิยา ปญฺจุปาทานกฺขนฺเธสุ ฉนฺทราคเวมตฺตตาติ. “สาธุ ภนฺเต”ติ โข โส ภิกฺขุ ฯเปฯ อุตฺตริํ ปญฺหํ อปุจฺฉิ\csromandash{}}

\prose{กิตฺตาวตา นุ โข ภนฺเต ขนฺธานํ ขนฺธาธิวจนนฺติ.\csromanspacer{} ยํ กิญฺจิ ภิกฺขุ รูปํ อตีตานาคตปจฺจุปฺปนฺนํ อชฺฌตฺตํ วา พหิทฺธา วา โอฬาริกํ วา สุขุมํ วา หีนํ วา ปณีตํ วา ยํ ทูเร สนฺติเก วา, อยํ วุจฺจติ รูปกฺขนฺโธ.\csromanspacer{} ยา กาจิ เวทนา.\csromanspacer{} ยา กาจิ สญฺญา.\csromanspacer{} เย เกจิ สงฺขารา.\csromanspacer{} ยํ กิญฺจิ วิญฺญาณํ อตีตานาคตปจฺจุปฺปนฺนํ อชฺฌตฺตํ วา พหิทฺธา วา โอฬาริกํ วา สุขุมํ วา หีนํ วา ปณีตํ วา ยํ ทูเร สนฺติเก วา, อยํ วุจฺจติ วิญฺญาณกฺขนฺโธ.\csromanspacer{} เอตฺตาวตา โข ภิกฺขุ ขนฺธานํ ขนฺธาธิวจนนฺติ. “สาธุ ภนฺเต”ติ โข โส ภิกฺขุ ฯเปฯ อปุจฺฉิ\csromandash{}}

\prose{โก นุ โข ภนฺเต เหตุ โก ปจฺจโย รูปกฺขนฺธสฺส ปญฺญาปนาย, โก เหตุ โก ปจฺจโย เวทนากฺขนฺธสฺส ปญฺญาปนาย, โก เหตุ โก ปจฺจโย สญฺญากฺขนฺธสฺส ปญฺญาปนาย, โก เหตุ โก ปจฺจโย สงฺขารกฺขนฺธสฺส ปญฺญาปนาย, โก เหตุ โก ปจฺจโย วิญฺญาณกฺขนฺธสฺส ปญฺญาปนายาติ.\csromanspacer{} จตฺตาโร โข ภิกฺขุ มหาภูตา เหตุ จตฺตาโร มหาภูตา ปจฺจโย รูปกฺขนฺธสฺส ปญฺญาปนาย, ผสฺโส เหตุ ผสฺโส ปจฺจโย เวทนากฺขนฺธสฺส ปญฺญาปนาย, ผสฺโส เหตุ ผสฺโส ปจฺจโย สญฺญากฺขนฺธสฺส ปญฺญาปนาย, ผสฺโส เหตุ ผสฺโส ปจฺจโย สงฺขารกฺขนฺธสฺส ปญฺญาปนาย, นามรูปํ เหตุ นามรูปํ ปจฺจโย วิญฺญาณกฺขนฺธสฺส ปญฺญาปนายาติ. “สาธุ ภนฺเต”ติ โข โส ภิกฺขุ ฯเปฯ อปุจฺฉิ\csromandash{}}

\prose{กถํ นุ โข ภนฺเต สกฺกายทิฏฺฐิ โหตีติ.\csromanspacer{} อิธ ภิกฺขุ อสฺสุตวา ปุถุชฺชโน อริยานํ อทสฺสาวี อริยธมฺมสฺส อโกวิโท อริยธมฺเม อวินีโต สปฺปุริสานํ อทสฺสาวี สปฺปุริสธมฺมสฺส อโกวิโท สปฺปุริสธมฺเม อวินีโต รูปํ อตฺตโต สมนุปสฺสติ, รูปวนฺตํ วา อตฺตานํ, อตฺตนิ วา รูปํ, รูปสฺมิํ วา อตฺตานํ.\csromanspacer{} เวทนํ.\csromanspacer{} สญฺญํ.\csromanspacer{} สงฺขาเร.\csromanspacer{} วิญฺญาณํ อตฺตโต สมนุปสฺสติ, วิญฺญาณวนฺตํ วา อตฺตานํ, อตฺตนิ วา วิญฺญาณํ, \csromanfolio{84}วิญฺญาณสฺมิํ วา อตฺตานํ.\csromanspacer{} เอวํ โข ภิกฺขุ สกฺกายทิฏฺฐิ โหตีติ. “สาธุ ภนฺเต”ติ โข โส ภิกฺขุ ฯเปฯ อปุจฺฉิ\csromandash{}}

\prose{กถํ ปน ภนฺเต สกฺกายทิฏฺฐิ น โหตีติ.\csromanspacer{} อิธ ภิกฺขุ สุตวา อริยสาวโก อริยานํ ทสฺสาวี อริยธมฺมสฺส โกวิโท อริยธมฺเม สุวินีโต สปฺปุริสานํ ทสฺสาวี สปฺปุริสธมฺมสฺส โกวิโท สปฺปุริสธมฺเม สุวินีโต น รูปํ อตฺตโต สมนุปสฺสติ, น รูปวนฺตํ วา อตฺตานํ, น อตฺตนิ วา รูปํ, น รูปสฺมิํ วา อตฺตานํ.\csromanspacer{} น เวทนํ.\csromanspacer{} น สญฺญํ.\csromanspacer{} น สงฺขาเร.\csromanspacer{} น วิญฺญาณํ อตฺตโต สมนุปสฺสติ, น วิญฺญาณวนฺตํ วา อตฺตานํ, น อตฺตนิ วา วิญฺญาณํ, น วิญฺญาณสฺมิํ วา อตฺตานํ.\csromanspacer{} เอวํ โข ภิกฺขุ สกฺกายทิฏฺฐิ น โหตีติ. “สาธุ ภนฺเต”ติ โข โส ภิกฺขุ ฯเปฯ อปุจฺฉิ\csromandash{}}

\prose{โก นุ โข ภนฺเต รูปสฺส อสฺสาโท, โก อาทีนโว, กิํ นิสฺสรณํ, โก เวทนาย.\csromanspacer{} โก สญฺญาย.\csromanspacer{} โก สงฺขารานํ.\csromanspacer{} โก วิญฺญาณสฺส อสฺสาโท, โก อาทีนโว, กิํ นิสฺสรณนฺติ.\csromanspacer{} ยํ โข ภิกฺขุ รูปํ ปฏิจฺจ อุปฺปชฺชติ สุขํ โสมนสฺสํ, อยํ รูปสฺส อสฺสาโท.\csromanspacer{} ยํ รูปํ อนิจฺจํ ทุกฺขํ วิปริณามธมฺมํ, อยํ รูปสฺส อาทีนโว.\csromanspacer{} โย รูปสฺมิํ ฉนฺทราควินโย ฉนฺทราคปฺปหานํ, อิทํ รูปสฺส นิสฺสรณํ.\csromanspacer{} ยํ เวทนํ ปฏิจฺจ.\csromanspacer{} ยํ สญฺญํ ปฏิจฺจ.\csromanspacer{} เย สงฺขาเร ปฏิจฺจ.\csromanspacer{} ยํ วิญฺญาณํ ปฏิจฺจ อุปฺปชฺชติ สุขํ โสมนสฺสํ, อยํ วิญฺญาณสฺส อสฺสาโท.\csromanspacer{} ยํ วิญฺญาณํ อนิจฺจํ ทุกฺขํ วิปริณามธมฺมํ, อยํ วิญฺญาณสฺส อาทีนโว.\csromanspacer{} โย วิญฺญาณสฺมิํ ฉนฺทราควินโย ฉนฺทราคปฺปหานํ, อิทํ วิญฺญาณสฺส นิสฺสรณนฺติ. “สาธุ ภนฺเต”ติ โข โส ภิกฺขุ ภควโต ภาสิตํ อภินนฺทิตฺวา อนุโมทิตฺวา ภควนฺตํ อุตฺตริํ ปญฺหํ อปุจฺฉิ\csromandash{}}

\prose{กถํ นุ โข ภนฺเต ชานโต กถํ ปสฺสโต อิมสฺมิํ จ สวิญฺญาณเก กาเย พหิทฺธา จ สพฺพนิมิตฺเตสุ อหงฺการ มมงฺการ มานานุสยา น โหนฺตีติ.\csromanspacer{} ยํ กิญฺจิ ภิกฺขุ รูปํ อตีตานาคตปจฺจุปฺปนฺนํ อชฺฌตฺตํ วา พหิทฺธา วา โอฬาริกํ วา สุขมํ วา หีนํ วา ปณีตํ วา ยํ ทูเร สนฺติเก วา, สพฺพํ รูปํ “เนตํ มม, เนโสหมสฺมิ, น เมโส อตฺตา”ติ เอวเมตํ ยถาภูตํ สมฺมปฺปญฺญาย ปสฺสติ.}

\prose{\csromanfolio{85}ยา กาจิ เวทนา.\csromanspacer{} ยา กาจิ สญฺญา.\csromanspacer{} เย เกจิ สงฺขารา.\csromanspacer{} ยํ กิญฺจิ วิญฺญาณํ อตีตานาคตปจฺจุปฺปนฺนํ อชฺฌตฺตํ วา พหิทฺธา วา โอฬาริกํ วา สุขุมํ วา หีนํ วา ปณีตํ วา ยํ ทูเร สนฺติเก วา, สพฺพํ วิญฺญาณํ “เนตํ มม, เนโสหมสฺมิ, น เมโส อตฺตา”ติ เอวเมตํ ยถาภูตํ สมฺมปฺปญฺญาย ปสฺสติ.\csromanspacer{} เอวํ โข ภิกฺขุ ชานโต เอวํ ปสฺสโต อิมสฺมิํ จ สวิญฺญาณเก กาเย พหิทฺธา จ สพฺพนิมิตฺเตสุ อหงฺการ มมงฺการ มานานุสยา น โหนฺตีติ.}

\prose{เตน โข ปน สมเยน อญฺญตรสฺส ภิกฺขุโน เอวํ เจตโส ปริวิตกฺโก อุทปาทิ “อิติ กิร โภ รูปํ อนตฺตา.\csromanspacer{} เวทนา.\csromanspacer{} สญฺญา.\csromanspacer{} สงฺขารา.\csromanspacer{} วิญฺญาณํ อนตฺตา, อนตฺตกตานิ กมฺมานิ กถมตฺตานํ\footnote{กตมตฺตานํ (อิ), กมฺมตฺตานํ (สฺยา, กํ, ก)} ผุสิสฺสนฺตี”ติ.\csromanspacer{} อถ โข ภควา ตสฺส ภิกฺขุโน เจตสา เจโตปริวิตกฺกมญฺญาย ภิกฺขู อามนฺเตสิ\csromandash{}}

\prose{ฐานํ โข ปเนตํ ภิกฺขเว วิชฺชติ, ยํ อิเธกจฺโจ โมฆปุริโส อวิทฺวา อวิชฺชาคโต ตณฺหาธิปเตยฺเยน เจตสา สตฺถุสาสนํ อติธาวิตพฺพํ มญฺเญยฺย “อิติ กิร โภ รูปํ อนตฺตา, เวทนา.\csromanspacer{} สญฺญา.\csromanspacer{} สงฺขารา.\csromanspacer{} วิญฺญาณํ อนตฺตา, อนตฺตกตานิ กมฺมานิ กถมตฺตานํ ผุสิสฺสนฺตี”ติ, ปฏิปุจฺฉาวินีตา โข เม ตุเมฺห ภิกฺขเว ตตฺร ตตฺร เตสุ เตสุ ธมฺเมสุ.}

\prose{ตํ กิํ มญฺญถ ภิกฺขเว, รูปํ นิจฺจํ วา อนิจฺจํ วาติ.\csromanspacer{} อนิจฺจํ ภนฺเต.\csromanspacer{} เวทนา.\csromanspacer{} สญฺญา.\csromanspacer{} สงฺขารา.\csromanspacer{} วิญฺญาณํ นิจฺจํ วา อนิจฺจํ วาติ.\csromanspacer{} อนิจฺจํ ภนฺเต.\csromanspacer{} ยํ ปนานิจฺจํ, ทุกฺขํ วา ตํ สุขํ วาติ.\csromanspacer{} ทุกฺขํ ภนฺเต.\csromanspacer{} ยํ ปนานิจฺจํ ทุกฺขํ วิปริณามธมฺมํ, กลฺลํ นุ ตํ สมนุปสฺสิตุํ “เอตํ มม, เอโสหสฺมิ, เอโส เม อตฺตา”ติ.\csromanspacer{} โน เหตํ ภนฺเต.\csromanspacer{} ตสฺมา ติห ฯเปฯ.\csromanspacer{} เอวํ ปสฺสํ ฯเปฯ นาปรํ อิตฺถตฺตายาติ ปชานาตีติ.}

\csromangathameasure{เทฺว ขนฺธา ตญฺเญว สิยํ, อธิวจนญฺจ เหตุนา. \\ สกฺกาเยน ทุเว วุตฺตา, อสฺสาทวิญฺญาณเกน จ. \\ เอเต ทสวิธา วุตฺตา, โหติ ภิกฺขุ ปุจฺฉายาติ. ทสมํ. ขชฺชนียวคฺโค ตติโย.}
\csromangathagroup{\csromangathastanza{เทฺว ขนฺธา ตญฺเญว สิยํ, อธิวจนญฺจ เหตุนา. \\ สกฺกาเยน ทุเว วุตฺตา, อสฺสาทวิญฺญาณเกน จ. \\ เอเต ทสวิธา วุตฺตา, โหติ ภิกฺขุ ปุจฺฉายาติ.\csromanspacer{} ทสมํ.\csromanspacer{} ขชฺชนียวคฺโค ตติโย.}}

\titlehead{\textbf{ตสฺสุทฺทานํ}}
\csromangathameasure{อสฺสาโท เทฺว สมุทยา, อรหนฺเตหิ อปเร เทฺว. \\ สีโห ขชฺชนี ปิณฺโฑลฺยํ, ปาลิเลยฺเยน ปุณฺณมาติ.}
\csromangathagroup{\csromangathastanza{\csromanfolio{86}อสฺสาโท เทฺว สมุทยา, อรหนฺเตหิ อปเร เทฺว. \\ สีโห ขชฺชนี ปิณฺโฑลฺยํ, ปาลิเลยฺเยน ปุณฺณมาติ.}}

\csromanmatikaanchor{mtk.2}
\csromanmatikaanchor{mtk.285}
\csromantocmarkref{section}{7-8. คทฺทุลพทฺธสุตฺต}{mtk.2}
\bookmark[dest={mtk.2},level=1]{7-8. คทฺทุลพทฺธสุตฺต}
\csromanheader{1}{\textbf{(9) 4. เถรวคฺค}}
\csromanheader{2}{1. อานนฺทสุตฺต}
\proseitem{83}{สาวตฺถินิทานํ.\csromanspacer{} ตตฺร โข อายสฺมา อานนฺโท ภิกฺขู อามนฺเตสิ “อาวุโส ภิกฺขเว”ติ. “อาวุโส”ติ โข เต ภิกฺขู อายสฺมโต อานนฺทสฺส ปจฺจสฺโสสุํ.\csromanspacer{} อายสฺมา อานนฺโท เอตทโวจ\csromandash{}}

\prose{ปุณฺโณ นาม อาวุโส อายสฺมา มนฺตาณิปุตฺโต\footnote{มนฺตานิปุตฺโต (ก-สี, สฺยา, กํ, อิ, ก)} อมฺหากํ นวกานํ สตํ พหูปกาโร โหติ, โส อเมฺห อิมินา โอวาเทน โอวทติ “อุปาทาย อาวุโส อานนฺท ‘อสฺมี’ติ โหติ โน อนุปาทาย.\csromanspacer{} กิญฺจ อุปาทาย ‘อสฺมี’ติ โหติ โน อนุปาทาย.\csromanspacer{} รูปํ อุปาทาย ‘อสฺมี’ติ โหติ โน อนุปาทาย.\csromanspacer{} เวทนํ.\csromanspacer{} สญฺญํ.\csromanspacer{} สงฺขาเร.\csromanspacer{} วิญฺญาณํ อุปาทาย ‘อสฺมี’ติ โหติ โน อนุปาทาย.}

\prose{เสยฺยถาปิ อาวุโส อานนฺท อิตฺถี วา ปุริโส วา ทหโร ยุวา มณฺฑนกชาติโก อาทาเส วา ปริสุทฺเธ ปริโยทาเต อจฺเฉ วา อุทกปตฺเต สกํ มุขนิมิตฺตํ ปจฺจเวกฺขมาโน อุปาทาย ปสฺเสยฺย โน อนุปาทาย.\csromanspacer{} เอวเมว โข อาวุโส อานนฺท รูปํ อุปาทาย ‘อสฺมี’ติ โหติ โน อนุปาทาย.\csromanspacer{} เวทนํ.\csromanspacer{} สญฺญํ.\csromanspacer{} สงฺขาเร.\csromanspacer{} วิญฺญาณํ อุปาทาย ‘อสฺมี’ติ โหติ โน อนุปาทาย.}

\prose{ตํ กิํ มญฺญสิ อาวุโส อานนฺท, รูปํ นิจฺจํ วา อนิจฺจํ วาติ.\csromanspacer{} อนิจฺจํ อาวุโส.\csromanspacer{} เวทนา.\csromanspacer{} สญฺญา.\csromanspacer{} สงฺขารา.\csromanspacer{} วิญฺญาณํ นิจฺจํ วา อนิจฺจํ วาติ.\csromanspacer{} อนิจฺจํ อาวุโส.\csromanspacer{} ตสฺมาติห ฯเปฯ.\csromanspacer{} เอวํ ปสฺสํ ฯเปฯ นาปรํ อิตฺถตฺตายาติ ปชานาตี”ติ.\csromanspacer{} ปุณฺโณ นาม อาวุโส อายสฺมา มนฺตาณิปุตฺโต อมฺหากํ นวกานํ สตํ พหูปกาโร โหติ, โส อเมฺห อิมินา \csromanfolio{87}โอวาเทน โอวทติ.\csromanspacer{} อิทญฺจ ปน เม อายสฺมโต ปุณฺณสฺส มนฺตาณิปุตฺตสฺส ธมฺมเทสนํ สุตฺวา ธมฺโม อภิสมิโตติ.\csromanspacer{}\csromanspacer{}ปฐมํ.}

\csromanheader{2}{2. ติสฺสสุตฺต}
\proseitem{84}{สาวตฺถินิทานํ.\csromanspacer{} เตน โข ปน สมเยน อายสฺมา ติสฺโส ภควโต ปิตุจฺฉาปุตฺโต สมฺพหุลานํ ภิกฺขูนํ เอวมาโรเจติ “อปิ เม อาวุโส มธุรกชาโต วิย กาโย, ทิสาปิ เม น ปกฺขายนฺติ, ธมฺมาปิ มํ น ปฏิภนฺติ, ถินมิทฺธญฺจ\footnote{ถีนมิทฺธญฺจ (สี, สฺยา, กํ, อิ)} เม จิตฺตํ ปริยาทาย ติฏฺฐติ, อนภิรโต จ พฺรหฺมจริยํ จรามิ, โหติ จ เม ธมฺเมสุ วิจิกิจฺฉา”ติ.}

\prose{อถ โข สมฺพหุลา ภิกฺขู เยน ภควา เตนุปสงฺกมิํสุ, อุปสงฺกมิตฺวา ภควนฺตํ อภิวาเทตฺวา เอกมนฺตํ นิสีทิํสุ, เอกมนฺตํ นิสินฺนา โข เต ภิกฺขู ภควนฺตํ เอตทโวจุํ “อายสฺมา ภนฺเต ติสฺโส ภควโต ปิตุจฺฉาปุตฺโต สมฺพหุลานํ ภิกฺขูนํ เอวมาโรเจติ ‘อปิ เม อาวุโส มธุรกชาโต วิย กาโย, ทิสาปิ เม น ปกฺขายนฺติ, ธมฺมาปิ มํ น ปฏิภนฺติ, ถินมิทฺธญฺจ เม จิตฺตํ ปริยาทาย ติฏฺฐติ, อนภิรโต จ พฺรหฺมจริยํ จรามิ, โหติ จ เม ธมฺเมสุ วิจิกิจฺฉา’ติ”.}

\prose{อถ โข ภควา อญฺญตรํ ภิกฺขุํ อามนฺเตสิ “เอหิ ตฺวํ ภิกฺขุ มม วจเนน ติสฺสํ ภิกฺขุํ อามนฺเตหี”ติ. “เอวํ ภนฺเต”ติ โข โส ภิกฺขุ ภควโต ปฏิสฺสุตฺวา เยนายสฺมา ติสฺโส เตนุปสงฺกมิ, อุปสงฺกมิตฺวา อายสฺมนฺตํ ติสฺสํ เอตทโวจ “สตฺถา ตํ อาวุโส ติสฺส อามนฺเตตี”ติ. “เอวมาวุโส”ติ โข อายสฺมา ติสฺโส ตสฺส ภิกฺขุโน ปฏิสฺสุตฺวา เยน ภควา เตนุปสงฺกมิ, อุปสงฺกมิตฺวา ภควนฺตํ อภิวาเทตฺวา เอกมนฺตํ นิสีทิ, เอกมนฺตํ นิสินฺนํ โข อายสฺมนฺตํ ติสฺสํ ภควา เอตทโวจ “สจฺจํ กิร ตฺวํ ติสฺส สมฺพหุลานํ ภิกฺขูนํ เอวมาโรเจสิ ‘อปิ เม อาวุโส มธุรกชาโต วิย กาโย ฯเปฯ โหติ จ เม ธมฺเมสุ วิจิกิจฺฉา’ติ”.\csromanspacer{} เอวํ ภนฺเต.\csromanspacer{} ตํ กิํ มญฺญสิ ติสฺส, รูเป อวิคตราคสฺส อวิคตจฺฉนฺทสฺส อวิคตเปมสฺส อวิคตปิปาสสฺส อวิคตปริฬาหสฺส อวิคตตณฺหสฺส ตสฺส รูปสฺส วิปริณามญฺญถาภาวา อุปฺปชฺชนฺติ โสกปริเทวทุกฺขโทมนสฺสุปายาสาติ.\csromanspacer{} เอวํ ภนฺเต.}

\prose{\csromanfolio{88}สาธุ สาธุ ติสฺส, เอวํ เหตํ ติสฺส โหติ, ยถา ตํ รูเป อวิคตราคสฺส.\csromanspacer{} เวทนาย.\csromanspacer{} สญฺญาย.\csromanspacer{} สงฺขาเรสุ อวิคตราคสฺส ฯเปฯ เตสํ สงฺขารานํ วิปริณามญฺญถาภาวา อุปฺปชฺชนฺติ โสกปริเทวทุกฺขโทมนสฺสุปายาสาติ.\csromanspacer{} เอวํ ภนฺเต.}

\prose{สาธุ สาธุ ติสฺส, เอวํ เหตํ ติสฺส โหติ, ยถา ตํ สงฺขาเรสุ อวิคตราคสฺส.\csromanspacer{} วิญฺญาเณ อวิคตราคสฺส อวิคตจฺฉนฺทสฺส อวิคตเปมสฺส อวิคตปิปาสสฺส อวิคตปริฬาหสฺส อวิคตณฺหสฺส, ตสฺส วิญฺญาณสฺส วิปริณามญฺญถาภาวา อุปฺปชฺชนฺติ โสกปริเทวทุกฺขโทมนสฺสุปายาสาติ.\csromanspacer{} เอวํ ภนฺเต.}

\prose{สาธุ สาธุ ติสฺส, เอวํ เหตํ ติสฺส โหติ, ยถา ตํ วิญฺญาเณ อวิคตราคสฺส.\csromanspacer{} ตํ กิํ มญฺญสิ ติสฺส, รูเป วิคตราคสฺส วิคตจฺฉนฺทสฺส วิคตเปมสฺส วิคตปิปาสสฺส วิคตปริฬาหสฺส วิคตตณฺหสฺส ตสฺส รูปสฺส วิปริณามญฺญถาวา อุปฺปชฺชนฺติ โสกปริเทวทุกฺขโทมนสฺสุปายาสาติ.\csromanspacer{} โน เหตํ ภนฺเต.}

\prose{สาธุ สาธุ ติสฺส, เอวํ เหตํ ติสฺส โหติ, ยถา ตํ รูเป วิคตราคสฺส.\csromanspacer{} เวทนาย.\csromanspacer{} สญฺญาย.\csromanspacer{} สงฺขาเรสุ วิคตราคสฺส.\csromanspacer{} วิญฺญาเณ วิคตราคสฺส วิคตจฺฉนฺทสฺส วิคตเปมสฺส วิคตปิปาสสฺส วิคตปริฬาหสฺส วิคตตณฺหสฺส, ตสฺส วิญฺญาณสฺส วิปริณามญฺญถาภาวา อุปฺปชฺชนฺติ โสกปริเทวทุกฺขโทมนสฺสุปายาสาติ.\csromanspacer{} โน เหตํ ภนฺเต.}

\prose{สาธุ สาธุ ติสฺส, เอวํ เหตํ ติสฺส โหติ, ยถา ตํ วิญฺญาเณ วิคตราคสฺส.\csromanspacer{} ตํ กิํ มญฺญสิ ติสฺส, รูปํ นิจฺจํ วา อนิจฺจํ วาติ.\csromanspacer{} อนิจฺจํ ภนฺเต.\csromanspacer{} เวทนา.\csromanspacer{} สญฺญา.\csromanspacer{} สงฺขารา.\csromanspacer{} วิญฺญาณํ นิจฺจํ วา อนิจฺจํ วาติ.\csromanspacer{} อนิจฺจํ ภนฺเต.\csromanspacer{} ตสฺมาติห ฯเปฯ.\csromanspacer{} เอวํ ปสฺสํ ฯเปฯ นาปรํ อิตฺถตฺตายาติ ปชานาติ.}

\prose{เสยฺยถาปิ ติสฺส เทฺว ปุริสา เอโก ปุริโส อมคฺคกุสโล เอโก ปุริโส มคฺคกุสโล, ตเมนํ โส อมคฺคกุสโล ปุริโส อมุํ มคฺคกุสลํ ปุริสํ มคฺคํ ปุจฺเฉยฺย.\csromanspacer{} โส เอวํ วเทยฺย “เอหิ โภ ปุริส อยํ มคฺโค, เตน มุหุตฺตํ คจฺฉ, เตน มุหุตฺตํ คนฺตฺวา ทกฺขิสฺสสิ เทฺวธาปถํ.\csromanspacer{} ตตฺถ วามํ มุญฺจิตฺวา ทกฺขิณํ คณฺหาหิ, เตน มุหุตฺตํ คจฺฉ เตน มุหุตฺตํ คนฺตฺวา ทกฺขิสฺสสิ ติพฺพํ วนสณฺฑํ, เตน มุหุตฺตํ คจฺฉ, เตน มุหุตฺตํ คนฺตฺวา ทกฺขิสฺสสิมหนฺตํ นินฺนํ ปลฺลลํ, เตน มุหุตฺตํ คจฺฉ, เตน \csromanfolio{89}มุหุตฺตํ คนฺตฺวา ทกฺขิสฺสสิ โสพฺภํ ปปาตํ, เตน มุหุตฺตํ คจฺฉ, เตน มุหุตฺตํ คนฺตฺวา ทกฺขิสฺสสิ สมํ ภูมิภาคํ รมณียนฺ”ติ.}

\prose{อุปมา โข มฺยายํ ติสฺส กตา อตฺถสฺส วิญฺญาปนาย.\csromanspacer{} อยํ เจเวตฺถ อตฺโถ\csromandash{}“ปุริโส อมคฺคกุสโล”ติ โข ติสฺส ปุถุชฺชนสฺเสตํ อธิวจนํ. “ปุริโส มคฺคกุสโล”ติ โข ติสฺส ตถาคตสฺเสตํ อธิวจนํ อรหโต สมฺมาสมฺพุทฺธสฺส. “เทฺวธาปโถ”ติ โข ติสฺส วิจิกิจฺฉาเยตํ อธิวจนํ. “วาโม มคฺโค”ติ โข ติสฺส อฏฺฐงฺคิกสฺเสตํ มิจฺฉามคฺคสฺส อธิวจนํ.\csromanspacer{} เสยฺยถิทํ, มิจฺฉาทิฏฺฐิยา ฯเปฯ มิจฺฉาสมาธิสฺส. “ทกฺขิโณ มคฺโค”ติ โข ติสฺส อริยสฺเสตํ อฏฺฐงฺคิกสฺส มคฺคสฺส อธิวจนํ.\csromanspacer{} เสยฺยถิทํ, สมฺมาทิฏฺฐิยา ฯเปฯ สมฺมาสมาธิสฺส. “ติพฺโพ วนสณฺโฑ”ติ โข ติสฺส อวิชฺชาเยตํ อธิวจนํ. “มหนฺตํ นินฺนํ ปลฺลลนฺ”ติ โข ติสฺส กามานเมตํ อธิวจนํ. “โสพฺโภ ปปาโต”ติ โข ติสฺส โกธูปายาสสฺเสตํ อธิวจนํ. “สโม ภูมิภาโค รมณีโย”ติ โข ติสฺส นิพฺพานสฺเสตํ อธิวจนํ.\csromanspacer{} อภิรม ติสฺส, อภิรม ติสฺส, อหโมวาเทน อหมนุคฺคเหน อหมนุสาสนิยาติ\footnote{อหมามิสธมฺมานุคฺคเหน มโมวาเทน มมานุสาสนิยาติ (ก)}.}

\prose{อิทมโวจ ภควา.\csromanspacer{} อตฺถมโน อายสฺมา ติสฺโส ภควโต ภาสิตํ อภินนฺทีติ.\csromanspacer{}\csromanspacer{}ทุติยํ.}

\csromanheader{2}{3. ยมกสุตฺต}
\proseitem{85}{เอกํ สมยํ อายสฺมา สาริปุตฺโต สาวตฺถิยํ วิหรติ เชตวเน อนาถปิณฺฑิกสฺส อาราเม.\csromanspacer{} เตน โข ปน สมเยน ยมกสฺส นาม ภิกฺขุโน เอวรูปํ ปาปกํ ทิฏฺฐิคตํ อุปฺปนฺนํ โหติ “ตถาหํ ภควตา ธมฺมํ เทสิตํ อาชานามิ, ยถา ขีณาสโว ภิกฺขุ กายสฺส เภทา อุจฺฉิชฺชติ วินสฺสติ, น โหติ ปรํ มรณา”ติ.}

\prose{อสฺโสสุํ โข สมฺพหุลา ภิกฺขู “ยมกสฺส กิร นาม ภิกฺขุโน เอวรูปํ ปาปกํ ทิฏฺฐิคตํ อุปฺปนฺนํ โหติ ‘ตถาหํ ภควตา ธมฺมํ เทสิตํ อาชานามิ, ยถา ขีณาสโว ภิกฺขุ กายสฺส เภทา อุจฺฉิชฺชติ วินสฺสติ, น โหติ ปรํ มรณา’ติ”.\csromanspacer{} อถ โข เต ภิกฺขู เยนายสฺมา \csromanfolio{90}ยมโก เตนุปสงฺกมิํสุ, อุปสงฺกมิตฺวา อายสฺมตา ยมเกน สทฺธิํ สมฺโมทิํสุ, สมฺโมทนียํ กถํ สารณียํ วีติสาเรตฺวา เอกมนฺตํ นิสีทิํสุ, เอกมนฺตํ นิสินฺนา โข เต ภิกฺขู อายสฺมนฺตํ ยมกํ เอตทโวจุํ\csromandash{}}

\prose{“สจฺจํ กิร เต อาวุโส ยมก เอวรูปํ ปาปกํ ทิฏฺฐิคตํ อุปฺปนฺนํ ‘ตถาหํ ภควตา ธมฺมํ เทสิตํ อาชานามิ, ยถา ขีณาสโว ภิกฺขุ กายสฺส เภทา อุจฺฉิชฺชติ วินสฺสติ, น โหติ ปรํ มรณา’ติ”.\csromanspacer{} เอวํ ขฺวาหํ อาวุโส ภควตา ธมฺมํ เทสิตํ อาชานามิ, ขีณาสโว ภิกฺขุ กายสฺส เภทา อุจฺฉิชฺชติ วินสฺสติ, น โหติ ปรํ มรณาติ.}

\prose{มา อาวุโส ยมก เอวํ อวจ, มา ภควนฺตํ อพฺภาจิกฺขิ, น หิ สาธุ ภควโต อพฺภาจิกฺขนํ, น หิ ภควา เอวํ วเทยฺย “ขีณาสโว ภิกฺขุ กายสฺส เภทา อุจฺฉิชฺชติ วินสฺสติ, น โหติ ปรํ มรณา”ติ.\csromanspacer{} เอวมฺปิ โข อายสฺมา ยมโก เตหิ ภิกฺขูหิ วุจฺจมาโน ตเถว ตํ ปาปกํ ทิฏฺฐิคตํ ถามสา ปรามาสา อภินิวิสฺส โวหรติ “ตถาหํ ภควตา ธมฺมํ เทสิตํ อาชานามิ, ยถา ขีณาสโว ภิกฺขุ กายสฺส เภทา อุจฺฉิชฺชติ วินสฺสติ, น โหติ ปรํ มรณา”ติ.}

\prose{ยถา โข เต ภิกฺขู นาสกฺขิํสุ อายสฺมนฺตํ ยมกํ เอตสฺมา ปาปกา ทิฏฺฐิคตา วิเวเจตุํ.\csromanspacer{} อถ โข เต ภิกฺขู อุฏฺฐายาสนา เยนายสฺมา สาริปุตฺโต เตนุปสงฺกมิํสุ, อุปสงฺกมิตฺวา อายสฺมนฺตํ สาริปุตฺตํ เอตทโวจุํ “ยมกสฺส นาม อาวุโส สาริปุตฺต ภิกฺขุโน เอวรูปํ ปาปกํ ทิฏฺฐิคตํ อุปฺปนฺนํ ‘ตถาหํ ภควตา ธมฺมํ เทสิตํ อาชานามิ, ยถา ขีณาสโว ภิกฺขุ กายสฺส เภทา อุจฺฉิชฺชติ วินสฺสติ, น โหติ ปรํ มรณา’ติ, สาธายสฺมา สาริปุตฺโต เยน ยมโก ภิกฺขุ เตนุปสงฺกมตุ อนุกมฺปํ อุปาทายา”ติ.\csromanspacer{} อธิวาเสสิ โข อายสฺมา สาริปุตฺโต ตุณฺหีภาเวน.\csromanspacer{} อถ โข อายสฺมา สาริปุตฺโต สายนฺหสมยํ ปฏิสลฺลานา วุฏฺฐิโต เยนายสฺมา ยมโก เตนุปสงฺกมิ, อุปสงฺกมิตฺวา อายสฺมตา ยมเกน สทฺธิํ สมฺโมทิ ฯเปฯ เอกมนฺตํ นิสินฺโน โข อายสฺมา สาริปุตฺโต อายสฺมนฺตํ ยมกํ เอตทโวจ\csromandash{}}

\prose{“สจฺจํ กิร เต อาวุโส ยมก เอวรูปํ ปาปกํ ทิฏฺฐิคตํ อุปฺปนฺนํ ‘ตถาหํ ภควตา ธมฺมํ เทสิตํ อาชานามิ, ยถา ขีณาสโว ภิกฺขุ กายสฺส เภทา อุจฺฉิชฺชติ วินสฺสติ, น โหติ ปรํ มรณา’ติ”. \csromanfolio{91}เอวํ ขฺวาหํ อาวุโส ภควตา ธมฺมํ เทสิตํ อาชานามิ, ยถา ขีณาสโว ภิกฺขุ กายสฺส เภทา อุจฺฉิชฺชติ วินสฺสติ, น โหติ ปรํ มรณาติ.}

\prose{ตํ กิํ มญฺญสิ อาวุโส ยมก, รูปํ นิจฺจํ วา อนิจฺจํ วาติ.\csromanspacer{} อนิจฺจํ อาวุโส.\csromanspacer{} เวทนา.\csromanspacer{} สญฺญา.\csromanspacer{} สงฺขารา.\csromanspacer{} วิญฺญาณํ นิจฺจํ วา อนิจฺจํ วาติ.\csromanspacer{} อนิจฺจํ อาวุโส.\csromanspacer{} ตสฺมาติห ฯเปฯ.\csromanspacer{} เอวํ ปสฺสํ ฯเปฯ นาปรํ อิตฺถตฺตายาติ ปชานาติ.}

\prose{ตํ กิํ มญฺญสิ อาวุโส ยมก, รูปํ “ตถาคโต”ติ สมนุปสฺสสีติ.\csromanspacer{} โน เหตํ อาวุโส.\csromanspacer{} เวทนํ “ตถาคโต”ติ สมนุปสฺสสีติ.\csromanspacer{} โน เหตํ อาวุโส.\csromanspacer{} สญฺญํ.\csromanspacer{} สงฺขาเร.\csromanspacer{} วิญฺญาณํ “ตถาคโต”ติ สมนุปสฺสสีติ.\csromanspacer{} โน เหตํ อาวุโส.}

\prose{ตํ กิํ มญฺญสิ อาวุโส ยมก, รูปสฺมิํ “ตถาคโต”ติ สมนุปสฺสสีติ.\csromanspacer{} โน เหตํ อาวุโส.\csromanspacer{} อญฺญตฺร รูปา “ตถาคโต”ติ สมนุปสฺสสีติ.\csromanspacer{} โน เหตํ อาวุโส.\csromanspacer{} เวทนาย.\csromanspacer{} อญฺญตฺร เวทนาย ฯเปฯ สญฺญาย.\csromanspacer{} อญฺญตฺร สญฺญาย.\csromanspacer{} สงฺขาเรสุ.\csromanspacer{} อญฺญตฺร สงฺขาเรหิ.\csromanspacer{} วิญฺญาณสฺมิํ “ตถาคโต”ติ สมนุปสฺสสีติ.\csromanspacer{} โน เหตํ อาวุโส.\csromanspacer{} อญฺญตฺร วิญฺญาณา “ตถาคโต”ติ สมนุปสฺสสีติ.\csromanspacer{} โน เหตํ อาวุโส.}

\prose{ตํ กิํ มญฺญสิ อาวุโส ยมก, รูปํ.\csromanspacer{} เวทนํ.\csromanspacer{} สญฺญํ.\csromanspacer{} สงฺขาเร.\csromanspacer{} วิญฺญาณํ “ตถาคโต”ติ สมนุปสฺสสีติ.\csromanspacer{} โน เหตํ อาวุโส.}

\prose{ตํ กิํ มญฺญสิ อาวุโส ยมก, อยํ โส อรูปี.\csromanspacer{} อเวทโน.\csromanspacer{} อสญฺญี.\csromanspacer{} อสงฺขาโร.\csromanspacer{} อวิญฺญาโณ “ตถาคโต”ติ สมนุปสฺสสีติ.\csromanspacer{} โน เหตํ อาวุโส.\csromanspacer{} เอตฺถ จ เต อาวุโส ยมก ทิฏฺเฐว ธมฺเม สจฺจโต เถตโต\footnote{ตถโต (สฺยา, กํ)} ตถาคเต อนุปลพฺภิยมาเน\footnote{ตถาคเต อนุปลพฺภมาเน (?)} กลฺลํ นุ เต ตํ เวยฺยากรณํ “ตถาหํ ภควตา ธมฺมํ เทสิตํ อาชานามิ, ยถา ขีณาสโว ภิกฺขุ กายสฺส เภทา อุจฺฉิชฺชติ วินสฺสติ, น โหติ ปรํ มรณา”ติ.}

\prose{อหุ โข เม ตํ อาวุโส สาริปุตฺต ปุพฺเพ อวิทฺทสุโน ปาปกํ ทิฏฺฐิคตํ, อิทญฺจ ปนายสฺมโต สาริปุตฺตสฺส ธมฺมเทสนํ สุตฺวา ตญฺเจว ปาปกํ ทิฏฺฐิคตํ ปหีนํ, ธมฺโม จ เม อภิสมิโตติ.}

\prose{\csromanfolio{92}สเจ ตํ อาวุโส ยมก เอวํ ปุจฺเฉยฺยุํ “โย โส อาวุโส ยมก ภิกฺขุ อรหํ ขีณาสโว, โส กายสฺส เภทา ปรํ มรณา กิํ โหตี”ติ, เอวํ ปุฏฺโฐ ตฺวํ อาวุโส ยมก กินฺติ พฺยากเรยฺยาสีติ.\csromanspacer{} สเจ มํ อาวุโส เอวํ ปุจฺเฉยฺยุํ “โย โส อาวุโส ยมก ภิกฺขุ อรหํ ขีณาสโว, โส กายสฺส เภทา ปรํ มรณา กิํ โหตี”ติ, เอวํ ปุฏฺโฐหํ อาวุโส เอวํ พฺยากเรยฺยํ “รูปํ โข อาวุโส อนิจฺจํ, ยทนิจฺจํ ตํ ทุกฺขํ, ยํ ทุกฺขํ ตํ นิรุทฺธํ, ตทตฺถงฺคตํ, เวทนา.\csromanspacer{} สญฺญา.\csromanspacer{} สงฺขารา.\csromanspacer{} วิญฺญาณํ อนิจฺจํ, ยทนิจฺจํ ตํ ทุกฺขํ, ยํ ทุกฺขํ ตํ นิรุทฺธํ, ตทตฺถงฺคตนฺ”ติ.\csromanspacer{} เอวํ ปุฏฺโฐหํ อาวุโส เอวํ พฺยากเรยฺยนฺติ.}

\prose{สาธุ สาธุ อาวุโส ยมก, เตน หาวุโส ยมก อุปมํ เต กริสฺสามิ เอตสฺเสว อตฺถสฺส ภิยฺโยโสมตฺตาย ญาณาย.\csromanspacer{} เสยฺยถาปิ อาวุโส ยมก คหปติ วา คหปติปุตฺโต วา อฑฺโฒ มหทฺธโน มหาโภโค, โส จ อารกฺขสมฺปนฺโน.\csromanspacer{} ตสฺส โกจิเทว ปุริโส อุปฺปชฺเชยฺย อนตฺถกาโม อหิตกาโม อโยคกฺเขมกาโม ชีวิตา โวโรเปตุกาโม.\csromanspacer{} ตสฺส เอวมสฺส “อยํ โข คหปติ วา คหปติปุตฺโต วา อฑฺโฒ มหทฺธโน มหาโภโค, โส จ อารกฺขสมฺปนฺโน, นายํ\footnote{น หายํ (สฺยา, กํ)} สุกโร ปสยฺห ชีวิตา โวโรเปตุํ.\csromanspacer{} ยํนูนาหํ อนุปขชฺช ชีวิตา โวโรเปยฺยนฺ”ติ.\csromanspacer{} โส ตํ คหปติํ วา คหปติปุตฺตํ วา อุปสงฺกมิตฺวา เอวํ วเทยฺย “อุปฏฺฐเหยฺยํ ตํ ภนฺเต”ติ.\csromanspacer{} ตเมนํ โส คหปติ วา คหปติปุตฺโต วา อุปฏฺฐาเปยฺย.\csromanspacer{} โส อุปฏฺฐเหยฺย ปุพฺพุฏฺฐายี ปจฺฉานิปาตี กิํการปฏิสฺสาวี มนาปจารี ปิยวาที, ตสฺส โส คหปติ วา คหปติปุตฺโต วา มิตฺตโตปิ นํ สทฺทเหยฺย\footnote{ทเหยฺย (สฺยา, กํ, อิ, ก)}, สุหชฺชโตปิ นํ สทฺทเหยฺย, ตสฺมิญฺจ วิสฺสาสํ อาปชฺเชยฺย.\csromanspacer{} ยทา โข อาวุโส ตสฺส ปุริสสฺส เอวมสฺส “สํวิสฺสตฺโถ โข มฺยายํ คหปติ วา คหปติปุตฺโต วา”ติ.\csromanspacer{} อถ นํ รโหคตํ วิทิตฺวา ติเณฺหน สตฺเถน ชีวิตา โวโรเปยฺย.}

\prose{ตํ กิํ มญฺญสิ อาวุโส ยมก, ยทา หิ โส ปุริโส อมุํ คหปติํ วา คหปติปุตฺตํ วา อุปสงฺกมิตฺวา เอวํ อาห “อุปฏฺฐเหยฺยํ ตํ ภนฺเต”ติ.\csromanspacer{} ตทาปิ โส วธโกว, วธกญฺจ ปน สนฺตํ น อญฺญาสิ “วธโก เม”ติ.\csromanspacer{} ยทาปิ โส อุปฏฺฐหติ ปุพฺพุฏฺฐายี ปจฺฉานิปาตี กิํการปฏิสฺสาวี มนาปจารี ปิยวาที, ตทาปิ โส วธโกว, วธกญฺจ ปน \csromanfolio{93}สนฺตํ น อญฺญาสิ “วธโก เม”ติ.\csromanspacer{} ยทาปิ นํ รโหคตํ วิทิตฺวา ติเณฺหน สตฺเถน ชีวิตา โวโรเปติ, ตทาปิ โส วธโกว, วธกญฺจ ปน สนฺตํ น อญฺญาสิ “วธโก เม”ติ.\csromanspacer{} เอวมาวุโสติ.\csromanspacer{} เอวเมว โข อาวุโส อสฺสุตวา ปุถุชฺชโน อริยานํ อทสฺสาวี อริยธมฺมสฺส อโกวิโท อริยธมฺเม อวินีโต, สปฺปุริสานํ อทสฺสาวี สปฺปุริสธมฺมสฺส อโกวิโท สปฺปุริสธมฺเม อวินีโต รูปํ อตฺตโต สมนุปสฺสติ, รูปวนฺตํ วา อตฺตานํ, อตฺตนิ วา รูปํ, รูปสฺมิํ วา อตฺตานิ.\csromanspacer{} เวทนํ.\csromanspacer{} สญฺญํ.\csromanspacer{} สงฺขาเร.\csromanspacer{} วิญฺญาณํ อตฺตโต สมนุปสฺสติ, วิญฺญาณวนฺตํ วา อตฺตานํ, อตฺตนิ วา วิญฺญาณํ, วิญฺญาณสฺมิํ วา อตฺตานํ.}

\prose{โส อนิจฺจํ รูปํ “อนิจฺจํ รูปนฺ”ติ ยถาภูตํ นปฺปชานาติ.\csromanspacer{} อนิจฺจํ เวทนํ “อนิจฺจา เวทนา”ติ ยถาภูตํ นปฺปชานาติ.\csromanspacer{} อนิจฺจํ สญฺญํ “อนิจฺจา สญฺญา”ติ ยถาภูตํ นปฺปชานาติ.\csromanspacer{} อนิจฺเจ สงฺขาเร “อนิจฺจา สงฺขารา”ติ ยถาภูตํ นปฺปชานาติ.\csromanspacer{} อนิจฺจํ วิญฺญาณํ “อนิจฺจํ วิญฺญาณนฺ”ติ ยถาภูตํ นปฺปชานาติ.}

\prose{ทุกฺขํ รูปํ “ทุกฺขํ รูปนฺ”ติ ยถาภูตํ นปฺปชานาติ.\csromanspacer{} ทุกฺขํ เวทนํ.\csromanspacer{} ทุกฺขํ สญฺญํ.\csromanspacer{} ทุกฺเข สงฺขาเร.\csromanspacer{} ทุกฺขํ วิญฺญาณํ “ทุกฺขํ วิญฺญาณนฺ”ติ ยถาภูตํ นปฺปชานาติ.}

\prose{อนตฺตํ รูปํ “อนตฺตา รูปนฺ”ติ ยถาภูตํ นปฺปาชานาติ.\csromanspacer{} อนตฺตํ เวทนํ.\csromanspacer{} อนตฺตํ สญฺญํ.\csromanspacer{} อนตฺเต สงฺขาเร.\csromanspacer{} อนตฺตํ วิญฺญาณํ “อนตฺตา วิญฺญาณนฺ”ติ ยถาภูตํ นปฺปชานาติ.}

\prose{สงฺขตํ รูปํ “สงฺขตํ รูปนฺ”ติ ยถาภูตํ นปฺปชานาติ, สงฺขตํ เวทนํ.\csromanspacer{} สงฺขตํ สญฺญํ.\csromanspacer{} สงฺขเต สงฺขาเร.\csromanspacer{} สงฺขตํ วิญฺญาณํ “สงฺขตํ วิญฺญาณนฺ”ติ ยถาภูตํ นปฺปชานาติ.}

\prose{วธกํ รูปํ “วธกํ รูปนฺ”ติ ยถาภูตํ นปฺปชานาติ.\csromanspacer{} วธกํ เวทนํ “วธกา เวทนา”ติ.\csromanspacer{} วธกํ สญฺญํ “วธกา สญฺญา”ติ.\csromanspacer{} วธเก สงฺขาเร “วธกา สงฺขารา”ติ ยถาภูตํ นปฺปชานาติ.\csromanspacer{} วธกํ วิญฺญาณํ “วธกํ วิญฺญาณนฺ”ติ ยถาภูตํ นปฺปชานาติ.}

\prose{โส รูปํ อุเปติ อุปาทิยติ อธิฏฺฐาติ “อตฺตา เม”ติ.\csromanspacer{} เวทนํ.\csromanspacer{} สญฺญํ.\csromanspacer{} สงฺขาเร.\csromanspacer{} วิญฺญาณํ อุเปติ อุปาทิยติ อธิฏฺฐาติ “อตฺตา เม”ติ.\csromanspacer{} ตสฺสิเม ปญฺจุปาทานกฺขนฺธา อุเปตา อุปาทินฺนา ทีฆรตฺตํ อหิตาย ทุกฺขาย สํวตฺตนฺติ.}

\prose{\csromanfolio{94}สุตวา จ โข อาวุโส อริยสาวโก อริยานํ ทสฺสาวี ฯเปฯ สปฺปุริสธมฺเม สุวินีโต น รูปํ อตฺตโต สมนุปสฺสติ, น รูปวนฺตํ อตฺตานํ น อตฺตนิ รูปํ, น รูปสฺมิํ อตฺตานํ.\csromanspacer{} น เวทนํ.\csromanspacer{} น สญฺญํ.\csromanspacer{} น สงฺขาเร.\csromanspacer{} น วิญฺญาณํ อตฺตโต สมนุปสฺสติ, น วิญฺญาณวนฺตํ อตฺตานํ, น อตฺตนิ วิญฺญาณํ น วิญฺญาณสฺมิํ อตฺตานํ.}

\prose{โส อนิจฺจํ รูปํ “อนิจฺจํ รูปนฺ”ติ ยถาภูตํ ปชานาติ.\csromanspacer{} อนิจฺจํ เวทนํ.\csromanspacer{} อนิจฺจํ สญฺญํ.\csromanspacer{} อนิจฺเจ สงฺขาเร.\csromanspacer{} อนิจฺจํ วิญฺญาณํ “อนิจฺจํ วิญฺญาณนฺ”ติ ยถาภูตํ ปชานาติ.}

\prose{ทุกฺขํ รูปํ “ทุกฺขํ รูปนฺ”ติ ยถาภูตํ ปชานาติ.\csromanspacer{} ทุกฺขํ เวทนํ.\csromanspacer{} ทุกฺขํ สญฺญํ.\csromanspacer{} ทุกฺเข สงฺขาเร.\csromanspacer{} ทุกฺขํ วิญฺญาณํ “ทุกฺขํ วิญฺญาณนฺ”ติ ยถาภูตํ ปชานาติ.}

\prose{อนตฺตํ รูปํ “อนตฺตา รูปนฺ”ติ ยถาภูตํ ปชานาติ.\csromanspacer{} อนตฺตํ เวทนํ.\csromanspacer{} อนตฺตํ สญฺญํ.\csromanspacer{} อนตฺเต สงฺขาเร.\csromanspacer{} อนตฺตํ วิญฺญาณํ “อนตฺตา วิญฺญาณนฺ”ติ ยถาภูตํ ปชานาติ.}

\prose{สงฺขตํ รูปํ “สงฺขตํ รูปนฺ”ติ ยถาภูตํ ปชานาติ.\csromanspacer{} สงฺขตํ เวทนํ.\csromanspacer{} สงฺขตํ สญฺญํ.\csromanspacer{} สงฺขเต สงฺขาเร.\csromanspacer{} สงฺขตํ วิญฺญาณํ “สงฺขตํ วิญฺญาณนฺ”ติ ยถาภูตํ ปชานาติ.}

\prose{วธกํ รูปํ “วธกํ รูปนฺ”ติ ยถาภูตํ ปชานาติ.\csromanspacer{} วธกํ เวทนํ.\csromanspacer{} วธกํ สญฺญํ.\csromanspacer{} วธเก สงฺขาเร “วธกา สงฺขารา”ติ ยถาภูตํ ปชานาติ.\csromanspacer{} วธกํ วิญฺญาณํ “วธกํ วิญฺญาณนฺ”ติ ยถาภูตํ ปชานาติ.}

\prose{โส รูปํ น อุเปติ น อุปาทิยติ นาธิฏฺฐาติ “อตฺตา เม”ติ.\csromanspacer{} เวทนํ.\csromanspacer{} สญฺญํ.\csromanspacer{} สงฺขาเร.\csromanspacer{} วิญฺญาณํ น อุเปติ น อุปาทิยติ นาธิฏฺฐาติ “อตฺตา เม”ติ.\csromanspacer{} ตสฺสิเม ปญฺจุปาทานกฺขนฺธา อนุเปตา อนุปาทินฺนา ทีฆรตฺตํ หิตาย สุขาย สํวตฺตนฺตีติ.\csromanspacer{} เอวเมตํ อาวุโส สาริปุตฺต โหติ, เยสํ อายสฺมนฺตานํ ตาทิสา สพฺรหฺมจาริโน อนุกมฺปกา อตฺถกามา โอวาทกา อนุสาสกา.\csromanspacer{} อิทญฺจ ปน เม อายสฺมโต สาริปุตฺตสฺส ธมฺมเทสนํ สุตฺวา อนุปาทาย อาสเวหิ จิตฺตํ วิมุตฺตนฺติ.\csromanspacer{}\csromanspacer{}ตติยํ.}

\csromanheader{2}{4. อนุราธสุตฺต}
\proseitem{86}{เอกํ สมยํ ภควา เวสาลิยํ วิหรติ มหาวเน กูฏาคารสาลายํ.\csromanspacer{} เตน โข ปน สมเยน อายสฺมา อนุราโธ ภควโต \csromanfolio{95}อวิทูเร อรญฺญกุฏิกายํ วิหรติ.\csromanspacer{} อถ โข สมฺพหุลา อญฺญติตฺถิยา ปริพฺพาชกา เยนายสฺมา อนุราโธ เตนุปสงฺกมิํสุ, อุปสงฺกมิตฺวา อายสฺมตา อนุราเธน สทฺธิํ สมฺโมทิํสุ, สมฺโมทนียํ กถํ สารณียํ วีติสาเรตฺวา เอกมนฺตํ นิสีทิํสุ, เอกมนฺตํ นิสินฺนา โข เต อญฺญติตฺถิยา ปริพฺพาชกา อายสฺมนฺตํ อนุราธํ เอตทโวจุํ “โย โส อาวุโส อนุราธ ตถาคโต อุตฺตมปุริโส ปรมปุริโส ปรมปตฺติปตฺโต ตํ ตถาคโต อิเมสุ จตูสุ ฐาเนสุ ปญฺญาปยมาโน ปญฺญาเปติ ‘โหติ ตถาคโต ปรํ มรณา’ติ วา, ‘น โหติ ตถาคโต ปรํ มรณา’ติ วา, โหติ จ น จ โหติ ตถาคโต ปรํ มรณา’ติ วา เนว โหติ น น โหติ ตถาคโต ปรํ มรณา’ติ วา”ติ.}

\prose{เอวํ วุตฺเต อายสฺมา อนุราโธ เต อญฺญติตฺถิเย ปริพฺพาชเก เอตทโวจ “โย โส อาวุโส ตถาคโต อุตฺตมปุริโส ปรมปุริโส ปรมปตฺติปตฺโต ตํ ตถาคโต อญฺญตฺร อิเมหิ จตูหิ ฐาเนหิ ปญฺญาปยมาโน ปญฺญาเปติ ‘โหติ ตถาคโต ปรํ มรณา’ติ วา, ‘น โหติ ตถาคโต ปรํ มรณา’ติ วา, ‘โหติ จ น จ โหติ ตถาคโต ปรํ มรณา’ติ วา, ‘เนว โหติ น น โหติ ตถาคโต ปรํ มรณา’ติ วา”ติ.\csromanspacer{} เอวํ วุตฺเต อญฺญติตฺถิยา ปริพฺพาชกา อายสฺมนฺตํ อนุราธํ เอตทโวจุํ “โส จายํ ภิกฺขุ นโว ภวิสฺสติ อจิรปพฺพชิโต, เถโร วา ปน พาโล อพฺยตฺโต”ติ.\csromanspacer{} อถ โข อญฺญติตฺถิยา ปริพฺพาชกา อายสฺมนฺตํ อนุราธํ นววาเทน จ พาลวาเทน จ อปสาเทตฺวา อุฏฺฐายาสนา ปกฺกมิํสุ.}

\prose{อถ โข อายสฺมโต อนุราธสฺส อจิรปกฺกนฺเตสุ เตสุ อญฺญติตฺถิเยสุ ปริพฺพาชเกสุ เอตทโหสิ “สเจ โข มํ อญฺญติตฺถิยา ปริพฺพาชกา อุตฺตริํ ปญฺหํ ปุจฺเฉยฺยุํ, กถํ พฺยากรมาโน นุ ขฺวาหํ เตสํ อญฺญติตฺถิยานํ ปริพฺพาชกานํ วุตฺตวาที เจว ภควโต อสฺสํ, น จ ภควนฺตํ อภูเตน อพฺภาจิกฺเขยฺยํ, ธมฺมสฺส จานุธมฺมํ พฺยากเรยฺยํ, น จ โกจิ สหธมฺมิโก วาทานุวาโท คารยฺหํ ฐานํ อาคจฺเฉยฺยา”ติ.}

\prose{อถ โข อายสฺมา อนุราโธ เยน ภควา เตนุปสงฺกมิ, อุปสงฺกมิตฺวา ฯเปฯ เอกมนฺตํ นิสินฺโน โข อายสฺมา อนุราโธ ภควนฺตํ เอตทโวจ “อิธาหํ ภนฺเต ภควโต อวิทูเร อรญฺญกุฏิกายํ วิหรามิ.\csromanspacer{} อถ โข ภนฺเต สมฺพหุลา อญฺญติตฺถิยา ปริพฺพาชกา เยนาหํ \csromanfolio{96}เตนุปสงฺกมิํสุ ฯเปฯ มํ เอตทโวจุํ ‘โย โส อาวุโส อนุราธ ตถาคโต อุตฺตมปุริโส ปรมปุริโส ปรมปตฺติปตฺโต, ตํ ตถาคโต อิเมสุ จตูสุ ฐาเนสุ ปญฺญาปยมาโน ปญฺญาเปติ ‘โหติ ตถาคโต ปรํ มรณา’ติ วา, น โหติ.\csromanspacer{} โหติ จ น จ โหติ. ‘เนว โหติ น น โหติ ตถาคโต ปรํ มรณา’ติ วา’ติ.}

\prose{เอวํ วุตฺตาหํ ภนฺเต เต อญฺญติตฺถิเย ปริพฺพาชเก เอตทโวจํ ‘โย โส อาวุโส ตถาคโต อุตฺตมปุริโส ปรมปุริโส ปรมปตฺติปตฺโต, ตํ ตถาคโต อญฺญตฺร อิเมหิ จตูหิ ฐาเนหิ ปญฺญาปยมาโน ปญฺญาเปติ ‘โหติ ตถาคโต ปรํ มรณา’ติ วา ฯเปฯ ‘เนว โหติ น น โหติ ตถาคโต ปรํ มรณา’ติ วา’ติ.\csromanspacer{} เอวํ วุตฺเต ภนฺเต เต อญฺญติตฺถิยา ปริพฺพาชกา มํ เอตทโวจุํ ‘โส จายํ ภิกฺขุ นโว ภวิสฺสติ อจิรปพฺพชิโต, เถโร วา ปน พาโล อพฺยตฺโต’ติ.\csromanspacer{} อถ โข มํ ภนฺเต เต อญฺญติตฺถิยา ปริพฺพาชกา นววาเทน จ พาลวาเทน จ อปสาเทตฺวา อุฏฺฐายาสนา ปกฺกมิํสุ.}

\prose{ตสฺส มยฺหํ ภนฺเต อจิรปกฺกนฺเตสุ เตสุ อญฺญติตฺถิเยสุ ปริพฺพาชเกสุ เอตทโหสิ ‘สเจ โข มํ เต อญฺญติตฺถิยา ปริพฺพาชกา อุตฺตริํ ปญฺหํ ปุจฺเฉยฺยุํ, กถํ พฺยากรมาโน นุ ขฺวาหํ เตสํ อญฺญติตฺถิยานํ ปริพฺพาชกานํ วุตฺตวาที เจว ภควโต อสฺสํ, น จ ภควนฺตํ อภูเตน อพฺภาจิกฺเขยฺยํ, ธมฺมสฺส จานุธมฺมํ พฺยากเรยฺยํ, น จ โกจิ สหธมฺมิโก วาทานุวาโท คารยฺหํ ฐานํ อาคจฺเฉยฺยา’ติ”.}

\prose{ตํ กิํ มญฺญสิ อนุราธ, รูปํ นิจฺจํ วา อนิจฺจํ วาติ.\csromanspacer{} อนิจฺจํ ภนฺเต.\csromanspacer{} ยํ ปนานิจฺจํ, ทุกฺขํ วา ตํ สุขํ วาติ.\csromanspacer{} ทุกฺขํ ภนฺเต.\csromanspacer{} ยํ ปนานิจฺจํ ทุกฺขํ วิปริณามธมฺมํ, กลฺลํ นุ ตํ สมนุปสฺสิตุํ “เอตํ มม, เอโสหมสฺมิ, เอโส เม อตฺตา”ติ.\csromanspacer{} โน เหตํ ภนฺเต.\csromanspacer{} เวทนา.\csromanspacer{} สญฺญา.\csromanspacer{} สงฺขารา.\csromanspacer{} วิญฺญาณํ นิจฺจํ วา อนิจฺจํ วาติ.\csromanspacer{} อนิจฺจํ ภนฺเต ฯเปฯ.\csromanspacer{} ตสฺมาติห ฯเปฯ.\csromanspacer{} เอวํ ปสฺสํ ฯเปฯ นาปรํ อิตฺถตฺตายาติ ปชานาติ.}

\prose{ตํ กิํ มญฺญสิ อนุราธ, รูปํ “ตถาคโต”ติ สมนุปสฺสสีติ.\csromanspacer{} โน เหตํ ภนฺเต.\csromanspacer{} เวทนํ.\csromanspacer{} สญฺญํ.\csromanspacer{} สงฺขาเร.\csromanspacer{} วิญฺญาณํ “ตถาคโต”ติ สมนุปสฺสสีติ.\csromanspacer{} โน เหตํ ภนฺเต.}

\prose{\csromanfolio{97}ตํ กิํ มญฺญสิ อนุราธ, รูปสฺมิํ “ตถาคโต”ติ สมนุปสฺสสีติ.\csromanspacer{} โน เหตํ ภนฺเต.\csromanspacer{} อญฺญตฺร รูปา “ตถาคโต”ติ สมนุปสฺสสีติ.\csromanspacer{} โน เหตํ ภนฺเต.\csromanspacer{} เวทนาย ฯเปฯ.\csromanspacer{} อญฺญตฺร เวทนาย ฯเปฯ สญฺญาย.\csromanspacer{} อญฺญตฺร สญฺญาย.\csromanspacer{} สงฺขาเรสุ.\csromanspacer{} อญฺญตฺร สงฺขาเรหิ.\csromanspacer{} วิญฺญาณสฺมิํ.\csromanspacer{} อญฺญตฺร วิญฺญาณา “ตถาคโต”ติ สมนุปสฺสสีติ.\csromanspacer{} โน เหตํ ภนฺเต.}

\prose{ตํ กิํ มญฺญสิ อนุราธ, รูปํ.\csromanspacer{} เวทนา.\csromanspacer{} สญฺญา.\csromanspacer{} สงฺขารา.\csromanspacer{} วิญฺญาณํ “ตถาคโต”ติ สมนุปสฺสสีติ.\csromanspacer{} โน เหตํ ภนฺเต.}

\prose{ตํ กิํ มญฺญสิ อนุราธ, อยํ โส อรูปี.\csromanspacer{} อเวทโน.\csromanspacer{} อสญฺญี.\csromanspacer{} อสงฺขาโร.\csromanspacer{} อวิญฺญาโณ “ตถาคโต”ติ สมนุปสฺสสีติ.\csromanspacer{} โน เหตํ ภนฺเต.}

\prose{เอตฺถ จ เต อนุราธ ทิฏฺเฐว ธมฺเม สจฺจโต เถตโต ตถาคเต อนุปลพฺภิยมาเน กลฺลํ นุ เต ตํ เวยฺยากรณํ “โย โส อาวุโส ตถาคโต อุตฺตมปุริโส ปรมปุริโส ปรมปตฺติปตฺโต, ตํ ตถาคโต อญฺญตฺร อิเมหิ จตูหิ ฐาเนหิ ปญฺญาปยมาโน ปญฺญาเปติ ‘โหติ ตถาคโต ปรํ มรณา’ติ วา.\csromanspacer{} น โหติ.\csromanspacer{} โหติ จ น จ โหติ. ‘เนว โหติ น น โหติ ตถาคโต ปรํ มรณาติ วา’ติ”.\csromanspacer{} โน เหตํ ภนฺเต.}

\prose{สาธุ สาธุ อนุราธ, ปุพฺเพ จาหํ อนุราธ เอตรหิ จ ทุกฺขญฺเจว ปญฺญเปมิ, ทุกฺขสฺส จ นิโรธนฺติ.\csromanspacer{}\csromanspacer{}จตุตฺถํ.}

\csromanheader{2}{5. วกฺกลิสุตฺต}
\proseitem{87}{เอกํ สมยํ ภควา ราชคเห วิหรติ เวฬุวเน กลนฺทกนิวาเป.\csromanspacer{} เตน โข ปน สมเยน อายสฺมา วกฺกลิ กุมฺภการนิเวสเน วิหรติ อาพาธิโก ทุกฺขิโต พาฬฺหคิลาโน.\csromanspacer{} อถ โข อายสฺมา วกฺกลิ อุปฏฺฐาเก อามนฺเตสิ “เอถ ตุเมฺห อาวุโส เยน ภควา เตนุปสงฺกมถ, อุปสงฺกมิตฺวา มม วจเนน ภควโต ปาเท สิรสา วนฺทถ ‘วกฺกลิ ภนฺเต ภิกฺขุ อาพาธิโก ทุกฺขิโต พาฬฺหคิลาโน, โส ภควโต ปาเท สิรสา วนฺทตี’ติ, เอวญฺจ วเทถ ‘สาธุ กิร ภนฺเต ภควา เยน วกฺกลิ ภิกฺขุ เตนุปสงฺกมตุ อนุกมฺปํ อุปาทายา’ติ”. “เอวมาวุโส”ติ โข เต ภิกฺขู อายสฺมโต วกฺกลิสฺส ปฏิสฺสุตฺวา \csromanfolio{98}เยน ภควา เตนุปสงฺกมิํสุ, อุปสงฺกมิตฺวา ภควนฺตํ อภิวาเทตฺวา เอกมนฺตํ นิสีทิํสุ, เอกมนฺตํ นิสินฺนา โข เต ภิกฺขู ภควนฺตํ เอตทโวจุํ “วกฺกลิ ภนฺเต ภิกฺขุ อาพาธิโก ทุกฺขิโต พาฬฺหคิลาโน, โส ภควโต ปาเท สิรสา วนฺทติ, เอวญฺจ ปน วเทติ ‘สาธุ กิร ภนฺเต ภควา เยน วกฺกลิ ภิกฺขุ เตนุปสงฺกมตุ อนุกมฺปํ อุปาทายา’ติ”.\csromanspacer{} อธิวาเสสิ ภควา ตุณฺหีภาเวน.}

\prose{อถ โข ภควา นิวาเสตฺวา ปตฺตจีวรมาทาย เยนายสฺมา วกฺกลิ เตนุปสงฺกมิ.\csromanspacer{} อทฺทสา โข อายสฺมา วกฺกลิ ภควนฺตํ ทูรโตว อาคจฺฉนฺตํ, ทิสฺวาน มญฺจเก สมโธสิ\footnote{สมญฺโจสิ (สี), สมญฺโจปิ (สฺยา, กํ) สํ + ธู + อี = สมโธสิ.}.\csromanspacer{} อถ โข ภควา อายสฺมนฺตํ วกฺกลิํ เอตทโวจ “อลํ วกฺกลิ มา ตฺวํ มญฺจเก สมโธสิ, สนฺติมานิ อาสนานิ ปญฺญตฺตานิ, ตตฺถาหํ นิสีทิสฺสามี”ติ.\csromanspacer{} นิสีทิ ภควา ปญฺญตฺเต อาสเน, นิสชฺช โข ภควา อายสฺมนฺตํ วกฺกลิํ เอตทโวจ “กจฺจิ เต วกฺกลิ ขมนียํ, กจฺจิ ยาปนียํ, กจฺจิ ทุกฺขา เวทนา ปฏิกฺกมนฺติ โน อภิกฺกมนฺติ, ปฏิกฺกโมสานํ ปญฺญายติ โน อภิกฺกโม”ติ.\csromanspacer{} น เม ภนฺเตขมนียํ, น ยาปนียํ, พาฬฺหา เม ทุกฺขา เวทนา อภิกฺกมนฺติ โน ปฏิกฺกมนฺติ, อภิกฺกโมสานํ ปญฺญายติ โน ปฏิกฺกโมติ.\csromanspacer{} กจฺจิ เต วกฺกลิ น กิญฺจิ กุกฺกุจฺจํ น โกจิ วิปฺปฏิสาโรติ.\csromanspacer{} ตคฺฆ เม ภนฺเต อนปฺปกํ กุกฺกุจฺจํ อนปฺปโก วิปฺปฏิสาโรติ.\csromanspacer{} กจฺจิ ปน ตํ วกฺกลิ อตฺตา สีลโต น อุปวทตีติ.\csromanspacer{} น โข มํ ภนฺเต อตฺตา สีลโต อุปวทตีติ.\csromanspacer{} โน เจ กิร ตํ วกฺกลิ อตฺตา สีลโต อุปวทติ, อถ กิญฺจ เต กุกฺกุจฺจํ โก จ วิปฺปฏิสาโรติ.\csromanspacer{} จิรปฏิกาหํ ภนฺเต ภควนฺตํ ทสฺสนาย อุปสงฺกมิตุกาโม, นตฺถิ จ เม กายสฺมิํ ตาวติกา พลมตฺตา, ยาวตาหํ\footnote{ยาหํ (สี), ยายาหํ (อิ)} ภควนฺตํ ทสฺสนาย อุปสงฺกเมยฺยนฺติ.}

\prose{อลํ วกฺกลิ, กิํ เต อิมินา ปูติกาเยน ทิฏฺเฐน, โย โข วกฺกลิ ธมฺมํ ปสฺสติ, โส มํ ปสฺสติ.\csromanspacer{} โย มํ ปสฺสติ, โส ธมฺมํ ปสฺสติ.\csromanspacer{} ธมฺมํ หิ วกฺกลิ ปสฺสนฺโต มํ ปสฺสติ, มํ ปสฺสนฺโต ธมฺมํ ปสฺสติ.}

\prose{ตํ กิํ มญฺญสิ วกฺกลิ, รูปํ นิจฺจํ วา อนิจฺจํ วาติ.\csromanspacer{} อนิจฺจํ ภนฺเต.\csromanspacer{} ยํ ปนานิจฺจํ, ทุกฺขํ วา ตํ สุขํ วาติ.\csromanspacer{} ทุกฺขํ ภนฺเต.\csromanspacer{} ยํ ปนานิจฺจํ ทุกฺขํ วิปริณามธมฺมํ, \csromanfolio{99}กลฺลํ นุ ตํ สมนุปสฺสิตุํ “เอตํ มม, เอโสหมสฺมิ, เอโส เม อตฺตา”ติ.\csromanspacer{} โน เหตํ ภนฺเต.\csromanspacer{} เวทนา.\csromanspacer{} สญฺญา.\csromanspacer{} สงฺขารา.\csromanspacer{} วิญฺญาณํ นิจฺจํ วา อนิจฺจํ วาติ.\csromanspacer{} อนิจฺจํ ภนฺเต ฯเปฯ.\csromanspacer{} เอโส เม อตฺตาติ.\csromanspacer{} โน เหตํ ภนฺเต.\csromanspacer{} ตสฺมาติห ฯเปฯ.\csromanspacer{} เอวํ ปสฺสํ ฯเปฯ นาปรํ อิตฺถตฺตายาติ ปชานาตีติ.}

\prose{อถ โข ภควา อายสฺมนฺตํ วกฺกลิํ อิมินา โอวาเทน โอวทิตฺวา อุฏฺฐายาสนา เยน คิชฺฌกูโฏ ปพฺพโต เตน ปกฺกามิ.\csromanspacer{} อถ โข อายสฺมา วกฺกลิ อจิรปกฺกนฺตสฺส ภควโต อุปฏฺฐาเก อามนฺเตสิ “เอถ มํ อาวุโส มญฺจกํ อาโรเปตฺวา เยน อิสิคิลิปสฺสํ กาฬสิลา เตนุปสงฺกมถ ‘กถํ หิ นาม มาทิโส อนฺตรฆเร กาลํ กตฺตพฺพํ มญฺเญยฺยา’ติ”. “เอวมาวุโส”ติ โข เต ภิกฺขู อายสฺมโต วกฺกลิสฺส ปฏิสฺสุตฺวา อายสฺมนฺตํ วกฺกลิํ มญฺจกํ อาโรเปตฺวา เยน อิสิคิลิปสฺสํ กาฬสิลา เตนุปสงฺกมิํสุ.\csromanspacer{} อถ โข ภควา ตญฺจ รตฺติํ ตญฺจ ทิวาวเสสํ คิชฺฌกูเฏ ปพฺพเต วิหาสิ.\csromanspacer{} อถ โข เทฺว เทวตาโย อภิกฺกนฺตาย รตฺติยา อภิกฺกนฺตวณฺณา เกวลกปฺปํ คิชฺฌกูฏํ โอภาเสตฺวา เยน ภควา เตนุปสงฺกมิํสุ ฯเปฯ เอกมนฺตํ อฏฺฐํสุ, เอกมนฺตํ ฐิตา โข เอกา เทวตา ภควนฺตํ เอตทโวจ “วกฺกลิ ภนฺเต ภิกฺขุ วิโมกฺขาย เจเตตี”ติ.\csromanspacer{} อปรา เทวตา ภควนฺตํ เอตทโวจ “โส หิ นูน ภนฺเต สุวิมุตฺโต วิมุจฺจิสฺสตี”ติ.\csromanspacer{} อิทมโวจุํ ตา เทวตาโย, อิทํ วตฺวา ภควนฺตํ อภิวาเทตฺวา ปทกฺขิณํ กตฺวา ตตฺเถวนฺตรธายิํสุ.}

\prose{อถ โข ภควา ตสฺสา รตฺติยา อจฺจเยน ภิกฺขู อามนฺเตสิ “เอถ ตุเมฺห ภิกฺขเว เยน วกฺกลิ ภิกฺขุ เตนุปสงฺกมถ, อุปสงฺกมิตฺวา วกฺกลิํ ภิกฺขุํ เอวํ วเทถ\csromandash{}}

\prose{สุณาวุโส ตฺวํ วกฺกลิ ภควโต วจนํ ทฺวินฺนญฺจ เทวตานํ, อิมํ อาวุโส รตฺติํ เทฺว เทวตาโย อภิกฺกนฺตาย รตฺติยา อภิกฺกนฺตวณฺณา เกวลกปฺปํ คิชฺฌกูฏํ โอภาเสตฺวา เยน ภควา เตนุปสงฺกมิํสุ, อุปสงฺกมิตฺวา ภควนฺตํ อภิวาเทตฺวา เอกมนฺตํ อฏฺฐํสุ, เอกมนฺตํ ฐิตา โข อาวุโส เอกา เทวตา ภควนฺตํ เอตทโวจ ‘วกฺกลิ ภนฺเต ภิกฺขุ วิโมกฺขาย เจเตตี’ติ.\csromanspacer{} อปรา เทวตา ภควนฺตํ เอตทโวจ ‘โส หิ นูน ภนฺเต สุวิมุตฺโต วิมุจฺจิสฺสตี’ติ.\csromanspacer{} ภควา จ ตํ อาวุโส วกฺกลิ เอวมาห ‘มา ภายิ วกฺกลิ, มา ภายิ วกฺกลิ, อปาปกํ เต มรณํ ภวิสฺสติ, อปาปิกา กาลกิริยา’ติ”. “เอวํ ภนฺเต”ติ โข \csromanfolio{100}เต ภิกฺขู ภควโต ปฏิสฺสุตฺวา เยนายสฺมา วกฺกลิ เตนุปสงฺกมิํสุ, อุปสงฺกมิตฺวา อายสฺมนฺตํ วกฺกลิํ เอตทโวจุํ “สุณาวุโส วกฺกลิ ภควโต วจนํ ทฺวินฺนญฺจ เทวตานนฺ”ติ.}

\prose{อถ โข อายสฺมา วกฺกลิ อุปฏฺฐาเก อามนฺเตสิ “เอถ มํ อาวุโส มญฺจกา โอโรเปถ ‘กถํ หิ นาม มาทิโส อุจฺเจ อาสเน นิสีทิตฺวา ตสฺส ภควโต สาสนํ โสตพฺพํ มญฺเญยฺยา’ติ”. “เอวมาวุโส”ติ โข เต ภิกฺขู อายสฺมโต วกฺกลิสฺส ปฏิสฺสุตฺวา อายสฺมนฺตํ วกฺกลิํ มญฺจกา โอโรเปสุํ.\csromanspacer{} อิมํ อาวุโส รตฺติํ เทฺว เทวตาโย อภิกฺกนฺตาย รตฺติยา ฯเปฯ เอกมนฺตํ อฏฺฐํสุ, เอกมนฺตํ ฐิตา โข อาวุโส เอกา เทวตา ภควนฺตํ เอตทโวจ “วกฺกลิ ภนฺเต ภิกฺขุ วิโมกฺขาย เจเตตี”ติ.\csromanspacer{} อปรา เทวตา ภควนฺตํ เอตทโวจ “โส หิ นูน ภนฺเต สุวิมุตฺโต วิมุจฺจิสฺสตี”ติ.\csromanspacer{} ภควา จ ตํ อาวุโส วกฺกลิ เอวมาห “มา ภายิ วกฺกลิ, มา ภายิ วกฺกลิ, อปาปกํ เต มรณํ ภวิสฺสติ, อปาปิกา กาลกิริยา”ติ.\csromanspacer{} เตนหาวุโส มม วจเนน ภควโต ปาเท สิรสา วนฺทถ “วกฺกลิ ภนฺเต ภิกฺขุ อาพาธิโก ทุกฺขิโต พาฬฺหคิลาโน, โส ภควโต ปาเท สิรสา วนฺทตี”ติ, เอวญฺจ วเทถ “รูปํ อนิจฺจํ, ตาหํ ภนฺเต น กงฺขามิ, ยทนิจฺจํ ตํ ทุกฺขนฺติ น วิจิกิจฺฉามิ, ยทนิจฺจํ ทุกฺขํ วิปริณามธมฺมํ, นตฺถิ เม ตตฺถ ฉนฺโท วา ราโค วา เปมํ วาติ น วิจิกิจฺฉามิ.\csromanspacer{} เวทนา อนิจฺจา, ตาหํ ภนฺเต น กงฺขามิ, ยทนิจฺจํ ตํ ทุกฺขนฺติ วิจิกิจฺฉามิ, ยทนิจฺจํ ทุกฺขํ วิปริณามธมฺมํ, นตฺถิ เม ตตฺถ ฉนฺโท วา ราโค วา เปมํ วาติ น วิจิกิจฺฉามิ.\csromanspacer{} สญฺญา.\csromanspacer{} สงฺขารา อนิจฺจา, ตาหํ ภนฺเต น กงฺขามิ, ยทนิจฺจํ ตํ ทุกฺขนฺติ น วิจิกิจฺฉามิ, ยทนิจฺจํ ทุกฺขํ วิปริณามธมฺมํ, นตฺถิ เม ตตฺถ ฉนฺโท วา ราโค วา เปมํ วาติ น วิจิกิจฺฉามิ.\csromanspacer{} วิญฺญาณํ อนิจฺจํ, ตาหํ ภนฺเต น กงฺขามิ, ยทนิจฺจํ ตํ ทุกฺขนฺติ น วิจิกิจฺฉามิ, ยทนิจฺจํ ทุกฺขํ วิปริณามธมฺมํ, นตฺถิ เม ตตฺถ ฉนฺโท วา ราโค วา เปมํ วาติ น วิจิกิจฺฉามี”ติ. “เอวมาวุโส”ติ โข เต ภิกฺขู อายสฺมโต วกฺกลิสฺส ปฏิสฺสุตฺวา ปกฺกมิํสุ.\csromanspacer{} อถ โข อายสฺมา วกฺกลิ อจิรปกฺกนฺเตสุ เตสุ ภิกฺขูสุ สตฺถํ อาหเรสิ.}

\prose{อถ โข เต ภิกฺขู เยน ภควา เตนุปสงฺกมิํสุ, อุปสงฺกมิตฺวา เอกมนฺตํ นิสีทิํสุ, เอกมนฺตํ นิสินฺนา โข เต ภิกฺขู ภควนฺตํ เอตทโวจุํ “วกฺกลิ ภนฺเต ภิกฺขุ อาพาธิโก ทุกฺขิโต พาฬฺหคิลาโน, โส \csromanfolio{101}ภควโต ปาเท สิรสา วนฺทติ, เอวญฺจ วเทติ ‘รูปํ อนิจฺจํ, ตาหํ ภนฺเต น กงฺขามิ, ยทนิจฺจํ ตํ ทุกฺขนฺติ น วิจิกิจฺฉามิ, ยทนิจฺจํ ทุกฺขํ วิปริณามธมฺมํ, นตฺถิ เม ตตฺถ ฉนฺโท วา ราโค วา เปมํ วาติ น วิจิกิจฺฉามิ.\csromanspacer{} เวทนา.\csromanspacer{} สญฺญา.\csromanspacer{} สงฺขารา.\csromanspacer{} วิญฺญาณํ อนิจฺจํ, ตาหํ ภนฺเต น กงฺขามิ, ยทนิจฺจํ ตํ ทุกฺขนฺติ น วิจิกิจฺฉามิ, ยทนิจฺจํ ทุกฺขํ วิปริณามธมฺมํ, นตฺถิ เม ตตฺถ ฉนฺโท วา ราโค วา เปมํ วาติ น วิจิกิจฺฉามี’ติ”.}

\prose{อถ โข ภควา ภิกฺขู อามนฺเตสิ “อายาม ภิกฺขเว เยน อิสิคิลิปสฺสํ กาฬสิลา เตนุปสงฺกมิสฺสาม, ยตฺถ วกฺกลินา กุลปุตฺเตน สตฺถมาหริตนฺ”ติ. “เอวํ ภนฺเต”ติ โข เต ภิกฺขู ภควโต ปจฺจสฺโสสุํ.\csromanspacer{} อถ โข ภควา สมฺพหุเลหิ ภิกฺขูหิ สทฺธิํ เยน อิสิคิลิปสฺสํ กาฬสิลา เตนุปสงฺกมิ.\csromanspacer{} อทฺทสา โข ภควา อายสฺมนฺตํ วกฺกลิํ ทูรโตว มญฺจเก วิวตฺตกฺขนฺธํ เสมานํ.}

\prose{เตน โข ปน สมเยน ธูมายิตตฺตํ ติมิรายิตตฺตํ คจฺฉเตว ปุริมํ ทิสํ, คจฺฉติ ปจฺฉิมํ ทิสํ, คจฺฉติ อุตฺตรํ ทิสํ, คจฺฉติ ทกฺขิณํ ทิสํ, คจฺฉติ อุทฺธํ ทิสํ, คจฺฉติ อโธ ทิสํ, คจฺฉติ อนุทิสํ.\csromanspacer{} อถ โข ภควา ภิกฺขู อามนฺเตสิ “ปสฺสถ โน ตุเมฺห ภิกฺขเว เอตํ ธูมายิตตฺตํ ติมิรายิตตฺตํ คจฺฉเตว ปุริมํ ทิสํ ฯเปฯ คจฺฉติ อนุทิสนฺ”ติ.\csromanspacer{} เอวํ ภนฺเต.\csromanspacer{} เอโส โข ภิกฺขเว มาโร ปาปิมา วกฺกลิสฺส กุลปุตฺตสฺส วิญฺญาณํ สมเนฺวสติ\footnote{สมนฺเนสติ (ก-สี, อิ)} “กตฺถ วกฺกลิสฺส กุลปุตฺตสฺส วิญฺญาณํ ปติฏฺฐิตนฺ”ติ, อปฺปติฏฺฐิเตน จ ภิกฺขเว วิญฺญาเณน วกฺกลิ กุลปุตฺโต ปรินิพฺพุโตสิ.\csromanspacer{}\csromanspacer{}ปญฺจมํ.}

\csromanheader{2}{6. อสฺสชิสุตฺต}
\proseitem{88}{เอกํ สมยํ ภควา ราชคเห วิหรติ เวฬุวเน กลนฺทกนิวาเป.\csromanspacer{} เตน โข ปน สมเยน อายสฺมา อสฺสชิ กสฺสปการาเม วิหรติ อาพาธิโก ทุกฺขิโต พาฬฺหคิลาโน.\csromanspacer{} อถ โข อายสฺมา อสฺสชิ อุปฏฺฐาเก อามนฺเตสิ “เอถ ตุเมฺห อาวุโส เยน ภควา เตนุปสงฺกมถ, อุปสงฺกมิตฺวา มม วจเนน ภควโต ปาเท สิรสา วนฺทถ ‘อสฺสชิ ภนฺเต ภิกฺขุ อาพาธิโก ทุกฺขิโต พาฬฺหคิลาโน, โส ภควโต ปาเท สิรสา วนฺทตี’ติ, เอวญฺจ วเทถ ‘สาธุ กิร ภนฺเต’ \csromanfolio{102}ภควา เยน อสฺสชิ ภิกฺขุ เตนุปสงฺกมตุ อนุกมฺปํ อุปาทายา’ติ”. “เอวมาวุโส”ติ โข เต ภิกฺขู อายสฺมโต อสฺสชิสฺส ปฏิสฺสุตฺวา เยน ภควา เตนุปสงฺกมิํสุ, อุปสงฺกมิตฺวา ภควนฺตํ อภิวาเทตฺวา เอกมนฺตํ นิสีทิํสุ, เอกมนฺตํ นิสินฺนา โข เต ภิกฺขู ภควนฺตํ เอตทโวจุํ “อสฺสชิ ภนฺเต ภิกฺขุ อาพาธิโก ฯเปฯ “สาธุ กิร ภนฺเต ภควา เยน อสฺสชิ ภิกฺขุ เตนุปสงฺกมตุ อนุกมฺปํ อุปาทายา”ติ.\csromanspacer{} อธิวาเสสิ ภควา ตุณฺหีภาเวน.}

\prose{อถ โข ภควา สายนฺหสมยํ ปฏิสลฺลานา วุฏฺฐิโต เยนายสฺมา อสฺสชิ เตนุปสงฺกมิ.\csromanspacer{} อทฺทสา โข อายสฺมา อสฺสชิ ภควนฺตํ ทูรโตว อาคจฺฉนฺตํ, ทิสฺวาน มญฺจเก สมโธสิ.\csromanspacer{} อถ โข ภควา อายสฺมนฺตํ อสฺสชิํ\footnote{อายสฺมโต อสฺสชิสฺส (อิ, ก)} เอตทโวจ “อลํ อสฺสชิ มา ตฺวํ มญฺจเก สมโธสิ, สนฺติมานิ อาสนานิ ปญฺญตฺตานิ, ตตฺถาหํ นิสีทิสฺสามี”ติ.\csromanspacer{} นิสีทิ ภควา ปญฺญตฺเต อาสเน, นิสชฺช โข ภควา อายสฺมนฺตํ อสฺสชิํ เอตทโวจ “กจฺจิ เต อาสฺสชิ ขมนียํ, กจฺจิ ยาปนียํ ฯเปฯ ปฏิกฺกโมสานํ ปญฺญายติ โน อภิกฺกโม”ติ.}

\prose{น เม ภนฺเต ขมนียํ ฯเปฯ อภิกฺกโมสานํ ปญฺญายติ โน ปฏิกฺกโมติ.\csromanspacer{} กจฺจิ เต อสฺสชิ น กิญฺจิ กุกฺกุจฺจํ น โกจิ วิปฺปฏิสาโรติ.\csromanspacer{} ตคฺฆ เม ภนฺเต อนปฺปกํ กุกฺกุจฺจํ อนปฺปโก วิปฺปฏิสาโรติ.\csromanspacer{} กจฺจิ ปน ตํ อสฺสชิ อตฺตา สีลโต น อุปวทตีติ.\csromanspacer{} น โข มํ ภนฺเต อตฺตา สีลโต อุปวทตีติ.\csromanspacer{} โน เจ กิร ตํ อสฺสชิ อตฺตา สีลโต อุปวทติ, อถ กิญฺจ เต กุกฺกุจฺจํ โก จ วิปฺปฏิสาโรติ.\csromanspacer{} ปุพฺเพ ขฺวาหํ ภนฺเต เคลญฺเญ ปสฺสมฺเภตฺวา ปสฺสมฺเภตฺวา กายสงฺขาเร วิหรามิ, โสหํ สมาธิํ นปฺปฏิลภามิ.\csromanspacer{} ตสฺส มยฺหํ ภนฺเต ตํ สมาธิํ อปฺปฏิลภโต เอวํ โหติ “โน จสฺสาหํ ปริหายามี”ติ.\csromanspacer{} เย เต อสฺสชิ สมณพฺราหฺมณา สมาธิสารกา สมาธิสามญฺญา, เตสํ ตํ สมาธิํ อปฺปฏิลภตํ เอวํ โหติ “โน จสฺสุ มยํ ปริหายามา”ติ.}

\prose{ตํ กิํ มญฺญสิ อสฺสชิ, รูปํ นิจฺจํ วา อนิจฺจํ วาติ.\csromanspacer{} อนิจฺจํ ภนฺเต ฯเปฯ.\csromanspacer{} วิญฺญาณํ ฯเปฯ.\csromanspacer{} ตสฺมาติห ฯเปฯ.\csromanspacer{} เอวํ ปสฺสํ ฯเปฯ นาปรํ อิตฺถตฺตายาติ ปชานาติ.\csromanspacer{} โส สุขํ เจ เวทนํ เวทยติ\footnote{เวทิยติ (สี, อิ)}, สา “อนิจฺจา”ติ ปชานาติ, \csromanfolio{103}“อนชฺโฌสิตา”ติ ปชานาติ, “อนิภินนฺทิตา”ติ ปชานาติ.\csromanspacer{} ทุกฺขํ เจ เวทนํ เวทยติ, สา “อนิจฺจา”ติ ปชานาติ, “อนชฺโฌสิตา”ติ ปชานาติ, “อนภินนฺทิตา”ติ ปชานาติ.\csromanspacer{} อทุกฺขมสุขํ เจ เวทนํ เวทยติ, สา “อนิจฺจา”ติ ปชานาติ ฯเปฯ “อนภินนฺทิตา”ติ ปชานาติ.\csromanspacer{} โส สุขํ เจ เวทนํ เวทยติ, วิสํยุตฺโต นํ เวทยติ, ทุกฺขํ เจ เวทนํ เวทยติ, วิสํยุตฺโต นํ เวทยติ.\csromanspacer{} อทุกฺขมสุขํ เจ เวทนํ เวทยติ, วิสํยุตฺโต นํ เวทยติ.\csromanspacer{} โส กายปริยนฺติกํ เจ เวทนํ เวทยมาโน “กายปริยนฺติกํ เวทนํ เวทยามี”ติ ปชานาติ, ชีวิตปริยนฺติกํ เจ เวทนํ เวทยมาโน “ชีวิตปริยนฺติกํ เวทนํ เวทยามี”ติ ปชานาติ, “กายสฺส เภทา อุทฺธํ ชีวิตปริยาทานา อิเธว สพฺพเวทยิตานิ อนภินนฺทิตานิ สีตีภวิสฺสนฺตี”ติ ปชานาติ.}

\prose{เสยฺยถาปิ อสฺสชิ เตลญฺจ ปฏิจฺจ วฏฺฏิญฺจ ปฏิจฺจ เตลปฺปทีโป ฌาเยยฺย, ตสฺเสว เตลสฺส จ วฏฺฏิยา จ ปริยาทานา อนาหาโร นิพฺพาเยยฺย.\csromanspacer{} เอวเมว โข อสฺสชิ ภิกฺขุ กายปริยนฺติกํ เวทนํ เวทยมาโน “กายปริยนฺติกํ เวทนํ เวทยามี”ติ ปชานาติ, ชีวิตปริยนฺติกํ เวทนํ เวทยมาโน “ชีวิตปริยนฺติกํ เวทนํ เวทยามี”ติ ปชานาติ, “กายสฺส เภทา อุทฺธํ ชีวิตปริยาทานา อิเธว สพฺพเวทยิตานิ อนภินนฺทิตานิ สีตีภวิสฺสนฺตี”ติ ปชานาตีติ.\csromanspacer{}\csromanspacer{}ฉฏฺฐํ.}

\csromanheader{2}{7. เขมกสุตฺต}
\proseitem{89}{เอกํ สมยํ สมฺพหุลา เถรา ภิกฺขู โกสมฺพิยํ วิหรนฺติ โฆสิตาราเม.\csromanspacer{} เตน โข ปน สมเยน อายสฺมา เขมโก พทริการาเม วิหรติ อาพาธิโก ทุกฺขิโต พาฬฺหคิลาโน.\csromanspacer{} อถ โข เถรา ภิกฺขู สายนฺหสมยํ ปฏิสลฺลานา วุฏฺฐิตา อายสฺมนฺตํ ทาสกํ อามนฺเตสุํ “เอหิ ตฺวํ อาวุโส ทาสก เยน เขมโก ภิกฺขุ เตนุปสงฺกม, อุปสงฺกมิตฺวา เขมกํ ภิกฺขุํ เอวํ วเทหิ ‘เถรา ตํ อาวุโส เขมก เอวมาหํสุ ‘กจฺจิ เต อาวุโส ขมนียํ, กจฺจิ ยาปนียํ, กจฺจิ ทุกฺขา เวทนา ปฏิกฺกมนฺติ โน อภิกฺกมนฺติ, ปฏิกฺกโมสานํ ปญฺญายติ โน อภิกฺกโม’ติ”. “เอวมาวุโส”ติ โข อายสฺมา ทาสโก เถรานํ ภิกฺขูนํ ปฏิสฺสุตฺวา เยนายสฺมา เขมโก เตนุปสงฺกมิ, อุปสงฺกมิตฺวา อายสฺมนฺตํ เขมกํ เอตทโวจ “เถรา ตํ \csromanfolio{104}อาวุโส เขมก เอวมาหํสุ ‘กจฺจิ เต อาวุโส ขมนียํ ฯเปฯ โน อภิกฺกโม’ติ”.\csromanspacer{} น เม อาวุโส ขมนียํ, น ยาปนียํ ฯเปฯ อภิกฺกโมสานํ ปญฺญายติ โน ปฏิกฺกโมติ.}

\prose{อถ โข อายสฺมา ทาสโก เยน เถรา ภิกฺขู เตนุปสงฺกมิ, อุปสงฺกมิตฺวา เถเร ภิกฺขู เอตทโวจ “เขมโก อาวุโส ภิกฺขุ เอวมาห ‘น เม อาวุโส ขมนียํ ฯเปฯ อภิกฺกโมสานํ ปญฺญายติ โน ปฏิกฺกโม’ติ”.\csromanspacer{} เอหิ ตฺวํ อาวุโส ทาสก เยน เขมโก ภิกฺขุ เตนุปสงฺกม, อุปสงฺกมิตฺวา เขมกํ ภิกฺขุํ เอวํ วเทหิ “เถรา ตํ อาวุโส เขมก เอวมาหํสุ ‘ปญฺจิเม อาวุโส อุปาทานกฺขนฺธา วุตฺตา ภควตา.\csromanspacer{} เสยฺยถิทํ, รูปุปาทานกฺขนฺโธ เวทนุปาทานกฺขนฺโธ สญฺญุปาทานกฺขนฺโธ สงฺขารุปาทานกฺขนฺโธ วิญฺญาณุปาทานกฺขนฺโธ, อิเมสุ อายสฺมา เขมโก ปญฺจสุ อุปาทานกฺขนฺเธสุ กิญฺจิ อตฺตํ วา อตฺตนิยํ วา สมนุปสฺสตี’ติ”.}

\prose{“เอวมาวุโส”ติ โข อายสฺมา ทาสโก เถรานํ ภิกฺขูนํ ปฏิสฺสุตฺวา เยนายสฺมา เขมโก เตนุปสงฺกมิ, อุปสงฺกมิตฺวา ฯเปฯ.\csromanspacer{} เถรา ตํ อาวุโส เขมก เอวมาหํสุ “ปญฺจิเม อาวุโส อุปาทานกฺขนฺธา วุตฺตา ภควตา.\csromanspacer{} เสยฺยถิทํ, รูปุปาทานกฺขนฺโธ ฯเปฯ วิญฺญาณุปาทานกฺขนฺโธ, อิเมสุ อายสฺมา เขมโก ปญฺจสุ อุปาทานกฺขนฺเธสุ กิญฺจิ อตฺตํ วา อตฺตนิยํ วา สมนุปสฺสตี”ติ.\csromanspacer{} ปญฺจิเม อาวุโส อุปาทานกฺขนฺธา วุตฺตา ภควตา.\csromanspacer{} เสยฺยถิทํ, รูปุปาทานกฺขนฺโธ ฯเปฯ วิญฺญาณุปาทานกฺขนฺโธ, อิเมสุ ขฺวาหํ อาวุโส ปญฺจสุ อุปาทานกฺขนฺเธสุ น กิญฺจิ อตฺตํ วา อตฺตนิยํ วา สมนุปสฺสามีติ.}

\prose{อถ โข อายสฺมา ทาสโก เยน เถรา ภิกฺขู เตนุปสงฺกมิ, อุปสงฺกมิตฺวา เถเร ภิกฺขู เอตทโวจ “เขมโก อาวุโส ภิกฺขุ เอวมาห ‘ปญฺจิเม อาวุโส อุปาทานกฺขนฺธา วุตฺตา ภควตา.\csromanspacer{} เสยฺยถิทํ, รูปุปาทานกฺขนฺโธ ฯเปฯ วิญฺญาณุปาทานกฺขนฺโธ, อิเมสุ ขฺวาหํ อาวุโส ปญฺจสุ อุปาทานกฺขนฺเธสุ น กิญฺจิ อตฺตํ วา อตฺตนิยํ วา สมนุปสฺสามี’ติ”.\csromanspacer{} เอหิ ตฺวํ อาวุโส ทาสก เยน เขมโก ภิกฺขุ เตนุปสงฺกม, อุปสงฺกมิตฺวา เขมกํ ภิกฺขุํ เอวํ วเทหิ “เถรา ตํ อาวุโส เขมก เอวมาหํสุ ‘ปญฺจิเม อาวุโส อุปาทานกฺขนฺธา วุตฺตา ภควตา.\csromanspacer{} เสยฺยถิทํ, รูปุปาทานกฺขนฺโธ ฯเปฯ วิญฺญาณุปาทานกฺขนฺโธ, โน เจ กิรายสฺมา เขมโก \csromanfolio{105}อิเมสุ ปญฺจสุ อุปาทานกฺขนฺเธสุ กิญฺจิ อตฺตํ วา อตฺตนิยํ วา สมนุปสฺสติ, เตนหายสฺมา เขมโก อรหํ ขีณาสโว’ติ”.}

\prose{“เอวมาวุโส”ติ โข อายสฺมา ทาสโก เถรานํ ภิกฺขูนํ ปฏิสฺสุตฺวา เยนายสฺมา เขมโก ฯเปฯ.\csromanspacer{} เถรา ตํ อาวุโส เขมก เอวมาหํสุ “ปญฺจิเม อาวุโส อุปาทานกฺขนฺธา วุตฺตา ภควตา.\csromanspacer{} เสยฺยถิทํ, รูปุปาทานกฺขนฺโธ ฯเปฯ วิญฺญาณุปาทานกฺขนฺโธ, โน เจ กิรายสฺมา เขมโก อิเมสุ ปญฺจสุ อุปาทานกฺขนฺเธสุ กิญฺจิ อตฺตํ วา อตฺตนิยํ วา สมนุปสฺสติ, เตนหายสฺมา เขมโก อรหํ ขีณาสโว’ติ”.\csromanspacer{} ปญฺจิเม อาวุโส อุปาทานกฺขนฺธา วุตฺตา ภควตา.\csromanspacer{} เสยฺยถิทํ, รูปุปาทานกฺขนฺโธ ฯเปฯ วิญฺญาณุปาทานกฺขนฺโธ, อิเมสุ ขฺวาหํ อาวุโส ปญฺจสุ อุปาทานกฺขนฺเธสุ น กิญฺจิ อตฺตํ วา อตฺตนิยํ วา สมนุปสฺสามิ, น จมฺหิ อรหํ ขีณาสโว, อปิ จ เม อาวุโส ปญฺจสุ อุปาทานกฺขนฺเธสุ “อสฺมี”ติ อธิคตํ, “อยมหมสฺมี”ติ น จ สมนุปสฺสามีติ.}

\prose{อถ โข อายสฺมา ทาสโก เยน เถรา ภิกฺขู ฯเปฯ เถเร ภิกฺขู เอตทโวจ\csromandash{}เขมโก อาวุโส ภิกฺขุ เอวมาห “ปญฺจิเม อาวุโส อุปาทานกฺขนฺธา วุตฺตา ภควตา.\csromanspacer{} เสยฺยถิทํ, รูปุปาทานกฺขนฺโธ ฯเปฯ วิญฺญาณุปาทานกฺขนฺโธ, อิเมสุ ขฺวาหํ อาวุโส ปญฺจสุ อุปาทานกฺขนฺเธสุ น กิญฺจิ อตฺตํ วา อตฺตนิยํ วา สมนุปสฺสามิ, น จมฺหิ อรหํ ขีณาสโว, อปิ จ เม อาวุโส ปญฺจสุ อุปาทานกฺขนฺเธสุ ‘อสฺมี’ติ อธิคตํ, ‘อยมหมสฺมี’ติ น จ สมนุปสฺสามี”ติ.}

\prose{เอหิ ตฺวํ อาวุโส ทาสก เยน เขมโก ภิกฺขุ เตนุปสงฺกม, อุปสงฺกมิตฺวา เขมกํ ภิกฺขุํ เอวํ วเทหิ\csromandash{}เถรา ตํ อาวุโส เขมก เอวมาหํสุ “ยเมตํ อาวุโส เขมก ‘อสฺมี’ติ วเทสิ, กิเมตํ ‘อสฺมี’ติ วเทสิ, รูปํ ‘อสฺมี’ติ วเทสิ, อญฺญตฺร รูปา ‘อสฺมี’ติ วเทสิ.\csromanspacer{} เวทนํ.\csromanspacer{} สญฺญํ.\csromanspacer{} สงฺขาเร.\csromanspacer{} วิญฺญาณํ ‘อสฺมี’ติ วเทสิ, อญฺญตฺร วิญฺญาณา ‘อสฺมี’ติ วเทสิ.\csromanspacer{} ยเมตํ อาวุโส เขมก ‘อสฺมี’ติ วเทสิ, กิเมตํ ‘อสฺมี’ติ วเทสี”ติ.}

\prose{“เอวมาวุโส”ติ โข อายสฺมา ทาสโก เถรานํ ภิกฺขูนํ ปฏิสฺสุตฺวา เยนายสฺมา เขมโก เตนุปสงฺกมิ, อุปสงฺกมิตฺวา อายสฺมนฺตํ เขมกํ เอตทโวจ\csromandash{}เถรา ตํ อาวุโส เขมก เอวมาหํสุ “ยเมตํ \csromanfolio{106}อาวุโส เขมก ‘อสฺมี’ติ วเทสิ, กิเมตํ ‘อสฺมี’ติ วเทสิ, รูปํ ‘อสฺมี’ติ วเทสิ, อญฺญตฺร รูปา ‘อสฺมี’ติ วเทสิ.\csromanspacer{} เวทนํ.\csromanspacer{} สญฺญํ.\csromanspacer{} สงฺขาเร.\csromanspacer{} วิญฺญาณํ ‘อสฺมี’ติ วเทสิ, อญฺญตฺร วิญฺญาณา ‘อสฺมี’ติ วเทสิ.\csromanspacer{} ยเมตํ อาวุโส เขมก ‘อสฺมี’ติ วเทสิ, กิเมตํ ‘อสฺมี’ติ วเทสี”ติ.\csromanspacer{} อลํ อาวุโส ทาสก กิํ อิมาย สนฺธาวนิกาย, อาหราวุโส ทณฺฑํ, อหเมว เยน เถรา ภิกฺขู เตนุปสงฺกมิสฺสามีติ.}

\prose{อถ โข อายสฺมา เขมโก ทณฺฑโมลุพฺภ เยน เถรา ภิกฺขุ เตนุปสงฺกมิ, อุปสงฺกมิตฺวา เถเรหิ ภิกฺขูหิ สทฺธิํ สมฺโมทิ, สมฺโมทนียํ กถํ สารณียํ วีติสาเรตฺวา เอกมนฺตํ นิสีทิ, เอกมนฺตํ นิสินฺนํ โข อายสฺมนฺตํ เขมกํ เถรา ภิกฺขู เอตทโวจุํ “ยเมตํ อาวุโส เขมก ‘อสฺมี’ติ วเทสิ, กิเมตํ ‘อสฺมี’ติ วเทสิ, รูปํ ‘อสฺมี’ติ วเทสิ, อญฺญตฺร รูปา ‘อสฺมี’ติ วเทสิ.\csromanspacer{} เวทนํ.\csromanspacer{} สญฺญํ.\csromanspacer{} สงฺขาเร.\csromanspacer{} วิญฺญาณํ ‘อสฺมี’ติ วเทสิ, อญฺญตฺร วิญฺญาณา ‘อสฺมี’ติ วเทสิ.\csromanspacer{} ยเมตํ อาวุโส เขมก ‘อสฺมี’ติ วเทสิ, กิเมตํ ‘อสฺมี’ติ วเทสี”ติ.\csromanspacer{} น ขฺวาหํ อาวุโส รูปํ “อสฺมี”ติ วทามิ, นปิ อญฺญตฺร รูปา “อสฺมี”ติ วทามิ.\csromanspacer{} น เวทนํ.\csromanspacer{} น สญฺญํ.\csromanspacer{} น สงฺขาเร.\csromanspacer{} น วิญฺญาณํ “อสฺมี”ติ วทามิ, นปิ อญฺญตฺร วิญฺญาณา “อสฺมี”ติ วทามิ, อปิ จ เม อาวุโส ปญฺจสุ อุปาทานกฺขนฺเธสุ “อสฺมี”ติ อธิคตํ, “อยมหมสฺมี”ติ น จ สมนุปสฺสามิ.}

\prose{เสยฺยถาปิ อาวุโส อุปฺปลสฺส วา ปทุมสฺส วา ปุณฺฑรีกสฺส วา คนฺโธ, โย นุ โข เอวํ วเทยฺย “ปตฺตสฺส คนฺโธ”ติ วา “วณฺณสฺส\footnote{วณฺฏสฺส (กตฺถจิ)} คนฺโธ”ติ วา “กิญฺชกฺขสฺส คนฺโธ”ติ วา, สมฺมา นุ โข โส วทมาโน วเทยฺยาติ.\csromanspacer{} โน เหตํ อาวุโส.\csromanspacer{} ยถา กถํ ปนาวุโส สมฺมา พฺยากรมาโน พฺยากเรยฺยาติ. “ปุปฺผสฺส คนฺโธ”ติ โข อาวุโส สมฺมา พฺยากรมาโน พฺยากเรยฺยาติ.\csromanspacer{} เอวเมว ขฺวาหํ อาวุโส น รูปํ “อสฺมี”ติ วทามิ, นปิ อญฺญตฺร รูปา “อสฺมี”ติ วทามิ.\csromanspacer{} น เวทนํ.\csromanspacer{} น สญฺญํ.\csromanspacer{} น สงฺขาเร.\csromanspacer{} น วิญฺญาณํ “อสฺมี”ติ วทามิ, นปิ อญฺญตฺร วิญฺญาณา “อสฺมี”ติ วทามิ, อปิ จ เม อาวุโส ปญฺจสุ อุปาทานกฺขนฺเธสุ “อสฺมี”ติ อธิคตํ, “อยมหมสฺมี”ติ น จ สมนุปสฺสามิ.}

\prose{กิญฺจาปิ อาวุโส อริยสาวกสฺส ปญฺโจรมฺภาคิยานิ สํโยชนานิ ปหีนานิ ภวนฺติ, อถ ขฺวสฺส โหติ โย จ ปญฺจสุ อุปาทานกฺขนฺเธสุ \csromanfolio{107}อนุสหคโต “อสฺมี”ติ มาโน “อสฺมี”ติ ฉนฺโท “อสฺมี”ติ อนุสโย อสมูหโต, โส อปเรน สมเยน ปญฺจสุ อุปาทานกฺขนฺเธสุ อุทยพฺพยานุปสฺสี วิหรติ “อิติ รูปํ, อิติ รูปสฺส สมุทโย, อิติ รูปสฺส อตฺถงฺคโม.\csromanspacer{} อิติ เวทนา.\csromanspacer{} อิติ สญฺญา.\csromanspacer{} อิติ สงฺขารา.\csromanspacer{} อิติ วิญฺญาณํ, อิติ วิญฺญาณสฺส สมุทโย, อิติ วิญฺญาณสฺส อตฺถงฺคโม”ติ.\csromanspacer{} ตสฺสิเมสุ ปญฺจสุ อุปาทานกฺขนฺเธสุ อุทยพฺพยานุปสฺสิโน วิหรโต โยปิสฺส โหติ ปญฺจสุ อุปาทานกฺขนฺเธสุ อนุสหคโต “อสฺมี”ติ มาโน “อสฺมี”ติ ฉนฺโท “อสฺมี”ติ อนุสโย อสมูหโต, โสปิ สมุคฺฆาตํ คจฺฉติ.}

\prose{เสยฺยถาปิ อาวุโส วตฺถํ สํกิลิฏฺฐํ มลคฺคหิตํ, ตเมนํ สามิกา รชกสฺส อนุปทชฺชุํ, ตเมนํ รชโก อูเส วา ขาเร วา โคมเย วา สมฺมทฺทิตฺวา อจฺเฉ อุทเก วิกฺขาเลติ, กิญฺจาปิ ตํ โหติ วตฺถํ ปริสุทฺธํ ปริโยทาตํ, อถ ขฺวสฺส โหติ เยว อนุสหคโต อูสคนฺโธ วา ขารคนฺโธ วา โคมยคนฺโธ วา อสมูหโต, ตเมนํ รชโก สามิกานํ เทติ, ตเมนํ สามิกา คนฺธปริภาวิเต กรณฺฑเก นิกฺขิปนฺติ, โยปิสฺส โหติ อนุสหคโต อูสคนฺโธ วา ขารคนฺโธ วา โคมยคนฺโธ วา อสมูหโต, โสปิ สมุคฺฆาตํ คจฺฉติ.\csromanspacer{} เอวเมว โข อาวุโส กิญฺจาปิ อริยสาวกสฺส ปญฺโจรมฺภาคิยานิ สํโยชนานิ ปหีนานิ ภวนฺติ, อถ ขฺวสฺส โหติ เยว ปญฺจสุ อุปาทานกฺขนฺเธสุ อนุสหคโต “อสฺมี”ติ มาโน “อสฺมี”ติ ฉนฺโท “อสฺมี”ติ อนุสโย อสมูหโต, โส อปเรน สมเยน ปญฺจสุ อุปาทานกฺขนฺเธสุ อุทยพฺพยานุปสฺสี วิหรติ “อิติ รูปํ, อิติ รูปสฺส สมุทโย, อิติ รูปสฺส อตฺถงฺคโม.\csromanspacer{} อิติ เวทนา.\csromanspacer{} อิติ สญฺญา.\csromanspacer{} อิติ สงฺขารา.\csromanspacer{} อิติ วิญฺญาณํ, อิติ วิญฺญาณสฺส สมุทโย, อิติ วิญฺญาณสฺส อตฺถงฺคโม”ติ, ตสฺส อิเมสุ ปญฺจสุ อุปาทานกฺขนฺเธสุ อุทยพฺพยานุปสฺสิโน วิหรโต โยปิสฺส โหติ ปญฺจสุ อุปาทานกฺขนฺเธสุ อนุสหคโต “อสฺมี”ติ มาโน “อสฺมี”ติ ฉนฺโท “อสฺมี”ติ อนุสโย อสมูหโต, โสปิ สมุคฺฆาตํ คจฺฉตีติ.}

\prose{เอวํ วุตฺเต เถรา ภิกฺขู อายสฺมนฺตํ เขมกํ เอตทโวจุํ “น โข\footnote{น โข ปน (ก)} มยํ อายสฺมนฺตํ เขมกํ วิเหสาเปขา ปุจฺฉิมฺห, อปิ จายสฺมา เขมโก ปโหสิ ตสฺส ภควโต สาสนํ วิตฺถาเรน อาจิกฺขิตุํ เทเสตุํ \csromanfolio{108}ปญฺญาเปตุํ ปฏฺฐเปตุํ วิวริตุํ วิภชิตุํ อุตฺตานีกาตุํ, ตยิทํ อายสฺมตา เขมเกน ตสฺส ภควโต สาสนํ วิตฺถาเรน อาจิกฺขิตํ เทสิตํ ปญฺญาปิตํ ปฏฺฐปิตํ วิวริตํ วิภชิตํ อุตฺตานีกตนฺ”ติ.}

\prose{อิทมโวจ อายสฺมา เขมโก.\csromanspacer{} อตฺตมนา เถรา ภิกฺขู อายสฺมโต เขมกสฺส ภาสิตํ อภินนฺทุํ.\csromanspacer{} อิมสฺมิํ จ ปน เวยฺยากรณสฺมิํ ภญฺญมาเน สฏฺฐิมตฺตานํ เถรานํ ภิกฺขูนํ อนุปาทาย อาสเวหิ จิตฺตานิ วิมุจฺจิํสุ อายสฺมโต เขมกสฺส จาติ.\csromanspacer{}\csromanspacer{}สตฺตมํ.}

\csromanheader{2}{8. ฉนฺนสุตฺต}
\proseitem{90}{เอกํ สมยํ สมฺพหุลา เถรา ภิกฺขู พาราณสิยํ วิหรนฺติ อิสิปตเน มิคทาเย.\csromanspacer{} อถ โข อายสฺมา ฉนฺโน สายนฺหสมยํ ปฏิสลฺลานา วุฏฺฐิโต อวาปุรณํ\footnote{อปาปุรณํ (สี, สฺยา, กํ)} อาทาย วิหาเรน วิหารํ อุปสงฺกมิตฺวา เถเร ภิกฺขู เอตทโวจ “โอวทนฺตุ มํ อายสฺมนฺโต เถรา, อนุสาสนฺตุ มํ อายสฺมนฺโต เถรา, กโรนฺตุ เม อายสฺมนฺโต เถรา ธมฺมิํ กถํ, ยถาหํ ธมฺมํ ปสฺเสยฺยนฺ”ติ.}

\prose{เอวํ วุตฺเต เถรา ภิกฺขู อายสฺมนฺตํ ฉนฺนํ เอตทโวจุํ “รูปํ โข อาวุโส ฉนฺน อนิจฺจํ, เวทนา อนิจฺจา, สญฺญา อนิจฺจา, สงฺขารา อนิจฺจา, วิญฺญาณํ อนิจฺจํ.\csromanspacer{} รูปํ อนตฺตา.\csromanspacer{} เวทนา.\csromanspacer{} สญฺญา.\csromanspacer{} สงฺขารา.\csromanspacer{} วิญฺญาณํ อนตฺตา.\csromanspacer{} สพฺเพ สงฺขารา อนิจฺจา, สพฺเพ ธมฺมา อนตฺตา”ติ.}

\prose{อถ โข อายสฺมโต ฉนฺนสฺส เอตทโหสิ\csromandash{}“มยฺหมฺปิ โข เอตํ เอวํ\footnote{มยฺหมฺปิ โข เอวํ (สฺยา, กํ)} โหติ ‘รูปํ อนิจฺจํ.\csromanspacer{} เวทนา.\csromanspacer{} สญฺญา.\csromanspacer{} สงฺขารา.\csromanspacer{} วิญฺญาณํ อนิจฺจํ.\csromanspacer{} รูปํ อนตฺตา.\csromanspacer{} เวทนา.\csromanspacer{} สญฺญา.\csromanspacer{} สงฺขารา.\csromanspacer{} วิญฺญาณํ อนตฺตา.\csromanspacer{} สพฺเพ สงฺขารา อนิจฺจา, สพฺเพ ธมฺมา อนตฺตา’ติ.\csromanspacer{} อถ จ ปน เม สพฺพสงฺขารสมเถ สพฺพูปธิปฏินิสฺสคฺเค ตณฺหากฺขเย วิราเค นิโรเธ นิพฺพาเน จิตฺตํ น ปกฺขนฺทติ นปฺปสีทติ น สนฺติฏฺฐติ นาธิมุจฺจติ, ปริตสฺสนา อุปาทานํ อุปฺปชฺชติ, ปจฺจุทาวตฺตติ มานสํ, ‘อถ โก จรหิ เม อตฺตา’ติ.\csromanspacer{} น โข ปเนวํ ธมฺมํ ปสฺสโต โหติ.\csromanspacer{} โก นุ โข เม ตถา ธมฺมํ เทเสยฺย, ยถาหํ ธมฺมํ ปสฺเสยฺยนฺ”ติ.}

\prose{อถ โข อายสฺมโต ฉนฺนสฺส เอตทโหสิ “อยํ โข อายสฺมา อานนฺโท โกสมฺพิยํ วิหรติ โฆสิตาราเม, สตฺถุ เจว สํวณฺณิโต, \csromanfolio{109}สมฺภาวิโต จ วิญฺญูนํ สพฺรหฺมจารีนํ, ปโหติ จ เม อายสฺมา อานนฺโท ตถา ธมฺมํ เทเสตุํ, ยถาหํ ธมฺมํ ปสฺเสยฺยํ.\csromanspacer{} อตฺถิ จ เม อายสฺมนฺเต อานนฺเท ตาวติกา วิสฺสฏฺฐิ\footnote{วิสฺสตฺถิ (?)}.\csromanspacer{} ยํนูนาหํ เยนายสฺมา อานนฺโท เตนุปสงฺกเมยฺยนฺ”ติ.\csromanspacer{} อถ โข อายสฺมา ฉนฺโน เสนาสนํ สํสาเมตฺวา ปตฺตจีวรมาทาย เยน โกสมฺพี โฆสิตาราโม, เยนายสฺมา อานนฺโท เตนุปสงฺกมิ, อุปสงฺกมิตฺวา อายสฺมตา อานนฺเทสุ สทฺธิํ สมฺโมทิ ฯเปฯ เอกมนฺตํ นิสินฺโน โข อายสฺมา ฉนฺโน อายสฺมนฺตํ อานนฺทํ เอตทโวจ\csromandash{}}

\prose{เอกมิทาหํ อาวุโส อานนฺท สมยํ พาราณสิยํ วิหรามิ อิสิปตเน มิคทาเย, อถ ขฺวาหํ อาวุโส สายนฺหสมยํ ปฏิสลฺลานา วุฏฺฐิโต อวาปุรณํ อาทาย วิหาเรน วิหารํ อุปสงฺกมิํ, อุปสงฺกมิตฺวา เถเร ภิกฺขู เอตทโวจํ “โอวทนฺตุ มํ อายสฺมนฺโต เถรา, อนุสาสนฺตุ มํ อายสฺมนฺโต เถรา, กโรนฺตุ เม อายสฺมนฺโต เถรา ธมฺมิํ กถํ, ยถาหํ ธมฺมํ ปสฺเสยฺยนฺ”ติ.\csromanspacer{} เอวํ วุตฺเต มํ อาวุโส เถรา ภิกฺขู เอตทโวจุํ “รูปํ โข อาวุโส ฉนฺน อนิจฺจํ.\csromanspacer{} เวทนา.\csromanspacer{} สญฺญา.\csromanspacer{} สงฺขารา.\csromanspacer{} วิญฺญาณํ อนิจฺจํ.\csromanspacer{} รูปํ อนตฺตา ฯเปฯ.\csromanspacer{} วิญฺญาณํ อนตฺตา.\csromanspacer{} สพฺเพ สงฺขารา อนิจฺจา, สพฺเพ ธมฺมา อนตฺตา”ติ.}

\prose{ตสฺส มยฺหํ อาวุโส เอตทโหสิ “มยฺหมฺปิ โข เอตํ เอวํ โหติ ‘รูปํ อนิจฺจํ ฯเปฯ วิญฺญาณํ อนิจฺจํ.\csromanspacer{} รูปํ อนตฺตา.\csromanspacer{} เวทนา.\csromanspacer{} สญฺญา.\csromanspacer{} สงฺขารา.\csromanspacer{} วิญฺญาณํ อนตฺตา.\csromanspacer{} สพฺเพ สงฺขารา อนิจฺจา, สพฺเพ ธมฺมา อนตฺตา’ติ.\csromanspacer{} อถ จ ปน เม สพฺพสงฺขารสมเถ สพฺพูปธิปฏินิสฺสคฺเค ตณฺหากฺขเย วิราเค นิโรเธ นิพฺพาเน จิตฺตํ น ปกฺขนฺทติ นปฺปสีทติ น สนฺติฏฺฐติ นาธิมุจฺจติ, ปริตสฺสนา อุปาทานํ อุปฺปชฺชติ, ปจฺจุทาวตฺตติ มานสํ, ‘อถ โก จรหิ เม อตฺตา’ติ.\csromanspacer{} น โข ปเนวํ ธมฺมํ ปสฺสโต โหติ.\csromanspacer{} โก นุ โข เม ตถา ธมฺมํ เทเสยฺย, ยถาหํ ธมฺมํ ปสฺเสยฺยนฺ”ติ.}

\prose{ตสฺส มยฺหํ อาวุโส เอตทโหสิ “อยํ โข อายสฺมา อานนฺโท โกสมฺพิยํ วิหรติ โฆสิตาราเม, สตฺถุ เจว สํวณฺณิโต, สมฺภาวิโต จ วิญฺญูนํ สพฺรหฺมจารีนํ, ปโหติ จ เม อายสฺมา อานนฺโท ตถา ธมฺมํ เทเสตุํ, ยถาหํ ธมฺมํ ปสฺเสยฺยํ, อตฺถิ จ เม อายสฺมนฺเต \csromanfolio{110}อานนฺเท ตาวติกา วิสฺสฏฺฐิ.\csromanspacer{} ยํนูนาหํ เยนายสฺมา อานนฺโท เตนุปสงฺกเมยฺยนฺ”ติ.\csromanspacer{} โอวทตุ มํ อายสฺมา อานนฺโท, อนุสาสตุ มํ อายสฺมา อานนฺโท, กโรตุ เม อายสฺมา อานนฺโท ธมฺมิํ กถํ, ยถาหํ ธมฺมํ ปสฺเสยฺยนฺติ.}

\prose{เอตฺตเกนปิ มยํ อายสฺมโต ฉนฺนสฺส อตฺตมนา, อปิ นาม ตํ\footnote{อตฺตมนา อภิรทฺธา, ตํ (สี, สฺยา, กํ)} อายสฺมา ฉนฺโน อาวิ อกาสิ, ขีลํ ฉินฺทิ\footnote{ปภินฺทิ (สี, สฺยา, กํ, อิ)}.\csromanspacer{} โอทหาวุโส ฉนฺน โสตํ, ภพฺโพสิ\footnote{ภพฺโพ ตฺวํ (ก) 4. ตาวเทว (สี)} ธมฺมํ วิญฺญาตุนฺติ.\csromanspacer{} อถ โข อายสฺมโต ฉนฺนสฺส ตาวตเกเนว4 อุฬารํ ปีติปาโมชฺชํ อุปฺปชฺชิ “ภพฺโพ กิรสฺมิ ธมฺมํ วิญฺญาตุนฺ”ติ.}

\prose{สมฺมุขา เมตํ อาวุโส ฉนฺน ภควโต สุตํ, สมฺมุขา ปฏิคฺคหิตํ กจฺจานโคตฺตํ ภิกฺขุํ โอวทนฺตสฺส, ทฺวยนิสฺสิโต ขฺวายํ กจฺจาน โลโก เยภุยฺเยน อตฺถิตญฺเจว นตฺถิตญฺจ, โลกสมุทยํ โข กจฺจาน ยถาภูตํ สมฺมปฺปญฺญาย ปสฺสโต ยา โลเก นตฺถิตา, สา น โหติ.\csromanspacer{} โลกนิโรธํ โข กจฺจาน ยถาภูตํ สมฺมปฺปญฺญาย ปสฺสโต ยา โลเก อตฺถิตา, สา น โหติ.\csromanspacer{} อุปยุปาทานาภินิเวสวินิพนฺโธ ขฺวายํ กจฺจาน โลโก เยภุยฺเยน, ตญฺจายํ อุปยุปาทานํ เจตโส อธิฏฺฐานาภินิเวสานุสยํ น อุเปติ น อุปาทิยติ, นาธิฏฺฐาติ “อตฺตา เม”ติ. “ทุกฺขเมว อุปฺปชฺชมานํ อุปฺปชฺชติ, ทุกฺขํ นิรุชฺฌมานํ นิรุชฺฌตี”ติ น กงฺขติ น วิจิกิจฺฉติ อปรปฺปจฺจยา ญาณเมวสฺส เอตฺถ โหติ.\csromanspacer{} เอตฺตาวตา โข กจฺจาน สมฺมาทิฏฺฐิ โหติ. “สพฺพมตฺถี”ติ โข กจฺจาน อยเมโก อนฺโต, “สพฺพํ นตฺถี”ติ อยํ ทุติโย อนฺโต, เอเต เต กจฺจาน อุโภ อนฺเต อนุปคมฺม มชฺเฌน ตถาคโต ธมฺมํ เทเสติ\csromandash{}อวิชฺชาปจฺจยา สงฺขารา, สงฺขารปจฺจยา วิญฺญาณํ ฯเปฯ เอวเมตสฺส เกวลสฺส ทุกฺขกฺขนฺธสฺส สมุทโย โหติ.\csromanspacer{} อวิชฺชาย เตฺวว อเสสวิราคนิโรธา สงฺขารนิโรโธ ฯเปฯ เอวเมตสฺส เกวลสฺส ทุกฺขกฺขนฺธสฺส นิโรโธ โหตีติ.}

\prose{เอวเมตํ อาวุโส อานนฺท โหติ, เยสํ อายสฺมนฺตานํ ตาทิสา สพฺรหฺมจารโย อนุกมฺปกา อตฺถกามา โอวาทกา อนุสาสกา, อิทญฺจ ปน เม อายสฺมโต อานนฺทสฺส ธมฺมเทสนํ สุตฺวา ธมฺโม อภิสมิโตติ.\csromanspacer{}\csromanspacer{}อฏฺฐมํ.}

\csromanheader{2}{9. ราหุลสุตฺต}
\proseitem{91}{\csromanfolio{111}สาวตฺถินิทานํ.\csromanspacer{} อถ โข อายสฺมา ราหุโล เยน ภควา เตนุปสงฺกมิ, อุปสงฺกมิตฺวา ฯเปฯ เอกมนฺตํ นิสินฺโน โข อายสฺมา ราหุโล ภควนฺตํ เอตทโวจ “กถํ นุ โข ภนฺเต ชานโต กถํ ปสฺสโต อิมสฺมิํ จ สวิญฺญาณเก กาเย พหิทฺธา จ สพฺพนิมิตฺเตสุ อหงฺการมมงฺการมานานุสยา น โหนฺตี”ติ.}

\prose{ยํ กิญฺจิ ราหุล รูปํ อตีตานาคตปจฺจุปฺปนฺนํ อชฺฌตฺตํ วา พหิทฺธา วา โอฬาริกํ วา สุขุมํ วา หีนํ วา ปณีตํ วา ยํ ทูเร สนฺติเก วา, สพฺพํ รูปํ “เนตํ มม, เนโสหมสฺมิ, น เมโส อตฺตา”ติ เอวเมตํ ยถาภูตํ สมฺมปฺปญฺญาย ปสฺสติ.\csromanspacer{} ยา กาจิ เวทนา.\csromanspacer{} ยา กาจิ สญฺญา.\csromanspacer{} เย เกจิ สงฺขารา ฯเปฯ.\csromanspacer{} ยํ กิญฺจิ วิญฺญาณํ อตีตานาคตปจฺจุปฺปนฺนํ อชฺฌตฺตํ วา พหิทฺธา วา ฯเปฯ สพฺพํ วิญฺญาณํ “เนตํ มม, เนโสหสฺมิ, น เมโส อตฺตา”ติ เอวเมตํ ยถาภูตํ สมฺมปฺปญฺญาย ปสฺสติ.\csromanspacer{} เอวํ โข ราหุล ชานโต เอวํ ปสฺสโต อิมสฺมิํ จ สวิญฺญาณเก กาเย พหิทฺธา จ สพฺพนิมิตฺเตสุ อหงฺการมมงฺการมานานุสยา น โหนฺตีติ.\csromanspacer{}\csromanspacer{}นวมํ.}

\csromanheader{2}{10. ทุติยราหุลสุตฺต}
\proseitem{92}{สาวตฺถินิทานํ.\csromanspacer{} เอกมนฺตํ นิสินฺโน โข อายสฺมา ราหุโล ภควนฺตํ เอตทโวจ “กถํ นุ โข ภนฺเต ชานโต กถํ ปสฺสโต อิมสฺมิํ จ สวิญฺญาณเก กาเย พหิทฺธา จ สพฺพนิมิตฺเตสุ อหงฺการมมงฺการมานาปคตํ มานสํ โหติ วิธาสมติกฺกนฺตํ สนฺตํ สุวิมุตฺตนฺ”ติ.\csromanspacer{} ยํกิญฺจิ ราหุล รูปํ อตีตานาคตปจฺจุปฺปนฺนํ อชฺฌตฺตํ วา พหิทฺธา วา ฯเปฯ ยํ ทูเร สนฺติเก วา, สพฺพํ รูปํ “เนตํ มม, เนโสหมสฺมิ, น เมโส อตฺตา”ติ เอวเมตํ ยถาภูตํ สมฺมปฺปญฺญาย ทิสฺวา อนุปาทาวิมุตฺโต โหติ.\csromanspacer{} ยา กาจิ เวทนา.\csromanspacer{} ยา กาจิ สญฺญา.\csromanspacer{} เย เกจิ สงฺขารา.\csromanspacer{} ยํ กิญฺจิ วิญฺญาณํ อตีตานาคตปจฺจุปฺปนฺนํ อชฺฌตฺตํ วา พหิทฺธา วา โอฬาริกํ วา สุขุมํ วา หีนํ วา ปณีตํ วา ยํ ทูเร สนฺติเก วา, สพฺพํ วิญฺญาณํ “เนตํ มม, เนโสหมสฺมิ, น เมโส อตฺตา”ติ เอวเมตํ ยถาภูตํ สมฺมปฺปญฺญาย ทิสฺวา อนุปาทาวิมุตฺโต โหติ.\csromanspacer{} เอวํ โข \csromanfolio{112}ราหุล ชานโต เอวํ ปสฺสโต อิมสฺมิํ จ สวิญฺญาณเก กาเย พหิทฺธา จ สพฺพนิมิตฺเตสุ อหงฺการมมงฺการมานาปคตํ มานสํ โหติ วิธา สมติกฺกนฺตํ สนฺตํ สุวิมุตฺตนฺติ.\csromanspacer{}\csromanspacer{}ทสมํ.\csromanspacer{} เถรวคฺโค จตุตฺโถ.}

\titlehead{\textbf{ตสฺสุทฺทานํ}}
\csromangathameasure{อานนฺโท ติสฺโส ยมโก, อนุราโธ จ วกฺกลิ. \\ อสฺสชิ เขมโก ฉนฺโน, ราหุลา อปเร ทุเว.}
\csromangathagroup{\csromangathastanza{อานนฺโท ติสฺโส ยมโก, อนุราโธ จ วกฺกลิ. \\ อสฺสชิ เขมโก ฉนฺโน, ราหุลา อปเร ทุเว.}}

\chapterhead{\textbf{(10) 5. ปุปฺผวคฺค}}
\csromanheader{2}{1. นทีสุตฺต}
\proseitem{93}{สาวตฺถินิทานํ.\csromanspacer{} เสยฺยถาปิ ภิกฺขเว นที ปพฺพเตยฺยา โอหารินี ทูรงฺคมา สีฆโสตา, ตสฺสา อุโภสุ ตีเรสุ\footnote{อุภโต ตีเร (สี), อุภโต ตีเรสุ (สฺยา, กํ)} กาสา เจปิ ชาตา อสฺสุ, เต นํ อชฺโฌลมฺเพยฺยุํ, กุสา เจปิ ชาตา อสฺสุ, เต นํ อชฺโฌลมฺเพยฺยุํ, ปพฺพชา\footnote{พพฺพชา (สี, อิ)} เจปิ ชาตา อสฺสุ, เต นํ อชฺโฌลมฺเพยฺยุํ, พีรณา เจปิ ชาตา อสฺสุ, เต นํ อชฺโฌลมฺเพยฺยุํ, รุกฺขา เจปิ ชาตา อสฺสุ, เต นํ อชฺโฌลมฺเพยฺยุํ.\csromanspacer{} ตสฺสา ปุริโส โสเตน วุยฺหมาโน กาเส เจปิ คเณฺหยฺย, เต ปลุชฺเชยฺยุํ, โส ตโตนิทานํ อนยพฺยสนํ อาปชฺเชยฺย, กุเส เจปิ คเณฺหยฺย, ปพฺพเช เจปิ คเณฺหยฺย, พีรเณ เจปิ คเณฺหยฺย, รุกฺเข เจปิ คเณฺหยฺย, เต ปลุชฺเชยฺยุํ, โส ตโตนิทานํ อนยพฺยสนํ อาปชฺเชยฺย.\csromanspacer{} เอวเมว โข ภิกฺขเว อสฺสุตวา ปุถุชฺชโน อริยานํ อทสฺสาวี อริยธมฺมสฺส อโกวิโท อริยธมฺเม อวินีโต, สปฺปุริสานํ อทสฺสาวี สปฺปุริสธมฺมสฺส อโกวิโท สปฺปุริสธมฺเม อวินีโต รูปํ อตฺตโต สมนุปสฺสติ, รูปวนฺตํ วา อตฺตานํ, อตฺตนิ วา รูปํ, รูปสฺมิํ วา อตฺตานํ, ตสฺส ตํ รูปํ ปลุชฺชติ, โส ตโตนิทานํ อนยพฺยสนํ อาปชฺชติ.\csromanspacer{} เวทนํ.\csromanspacer{} สญฺญํ.\csromanspacer{} สงฺขาเร.\csromanspacer{} วิญฺญาณํ อตฺตโต สมนุปสฺสติ, วิญฺญาณวนฺตํ วา \csromanfolio{113}อตฺตานํ, อตฺตนิ วา วิญฺญาณํ, วิญฺญาณสฺมิํ วา อตฺตานํ, ตสฺส ตํ วิญฺญาณํ ปลุชฺชติ.\csromanspacer{} โส ตโตนิทานํ อนยพฺยสนํ อาปชฺชติ.\csromanspacer{} ตํ กิํ มญฺญถ ภิกฺขเว, รูปํ นิจฺจํ วา อนิจฺจํ วาติ.\csromanspacer{} อนิจฺจํ ภนฺเต.\csromanspacer{} เวทนา.\csromanspacer{} สญฺญา.\csromanspacer{} สงฺขารา.\csromanspacer{} วิญฺญาณํ นิจฺจํ วา อนิจฺจํ วาติ.\csromanspacer{} อนิจฺจํ ภนฺเต.\csromanspacer{} ตสฺมาติห ฯเปฯ.\csromanspacer{} เอวํ ปสฺสํ ฯเปฯ นาปรํ อิตฺถตฺตายาติ ปชานาตีติ.\csromanspacer{}\csromanspacer{}ปฐมํ.}

\csromanheader{2}{2. ปุปฺผสุตฺต}
\proseitem{94}{สาวตฺถินิทานํ.\csromanspacer{} นาหํ ภิกฺขเว โลเกน วิวทามิ, โลโกว มยา วิวทติ.\csromanspacer{} น ภิกฺขเว ธมฺมวาที เกนจิ โลกสฺมิํ วิวทติ.\csromanspacer{} ยํ ภิกฺขเว นตฺถิสมฺมตํ โลเก ปณฺฑิตานํ, อหมฺปิ ตํ “นตฺถี”ติ วทามิ.\csromanspacer{} ยํ ภิกฺขเว อตฺถิสมฺมตํ โลเก ปณฺฑิตานํ, อหมฺปิ ตํ “อตฺถี”ติ วทามิ.}

\prose{กิญฺจ ภิกฺขเว นตฺถิสมฺมตํ โลเก ปณฺฑิตานํ, ยมหํ “นตฺถี”ติ วทามิ.\csromanspacer{} รูปํ ภิกฺขเว นิจฺจํ ธุวํ สสฺสตํ อวิปริณามธมฺมํ นตฺถิสมฺมตํ โลเก ปณฺฑิตานํ, อหมฺปิ ตํ “นตฺถี”ติ วทามิ.\csromanspacer{} เวทนา.\csromanspacer{} สญฺญา.\csromanspacer{} สงฺขารา.\csromanspacer{} วิญฺญาณํ นิจฺจํ ธุวํ สสฺสตํ อวิปริณามธมฺมํ นตฺถิสมฺมตํ โลเก ปณฺฑิตานํ, อหมฺปิ ตํ “นตฺถี”ติ วทามิ.\csromanspacer{} อิทํ โข ภิกฺขเว นตฺถิสมฺมตํ โลเก ปณฺฑิตานํ, อหมฺปิ ตํ “นตฺถี”ติ วทามิ.}

\prose{กิญฺจ ภิกฺขเว อตฺถิสมฺมตํ โลเก ปณฺฑิตานํ, ยมหํ “อตฺถี”ติ วทามิ.\csromanspacer{} รูปํ ภิกฺขเว อนิจฺจํ ทุกฺขํ วิปริณามธมฺมํ อตฺถิสมฺมตํ โลเก ปณฺฑิตานํ, อหมฺปิ ตํ “อตฺถี”ติ วทามิ.\csromanspacer{} เวทนา อนิจฺจา ฯเปฯ.\csromanspacer{} วิญฺญาณํ อนิจฺจํ ทุกฺขํ วิปริณามธมฺมํ อตฺถิสมฺมตํ โลเก ปณฺฑิตานํ, อหมฺปิ ตํ “อตฺถี”ติ วทามิ.\csromanspacer{} อิทํ โข ภิกฺขเว อตฺถิสมฺมตํ โลเก ปณฺฑิตานํ, อหมฺปิ ตํ “อตฺถี”ติ วทามิ.}

\prose{อตฺถิ ภิกฺขเว โลเก โลกธมฺโม, ตํ ตถาคโต อภิสมฺพุชฺฌติ อภิสเมติ, อภิสมฺพุชฺฌิตฺวา อภิสเมตฺวา ตํ อาจิกฺขติ เทเสติ ปญฺญเปติ ปฏฺฐเปติ วิวรติ วิภชติ อุตฺตานีกโรติ.}

\prose{กิญฺจ ภิกฺขเว โลเก โลกธมฺโม, ตํ ตถาคโต อภิสมฺพุชฺฌติ อภิสเมติ, อภิสมฺพุชฺฌิตฺวา อภิสเมตฺวา อาจิกฺขติ เทเสติ ปญฺญเปติ ปฏฺฐเปติ วิวรติ วิภชติ อุตฺตานีกโรติ.\csromanspacer{} รูปํ ภิกฺขเว โลเก โลกธมฺโม, ตํ ตถาคโต อภิสมฺพุชฺฌติ อภิสเมติ, \csromanfolio{114}อภิสมฺพุชฺฌิตฺวา อภิสเมตฺวา อาจิกฺขติ เทเสติ ปญฺญเปติ ปฏฺฐเปติ วิวรติ วิภชติ อุตฺตานีกโรติ.}

\prose{โย ภิกฺขเว ตถาคเตน เอวํ อาจิกฺขิยมาเน เทสิยมาเน ปญฺญปิยมาเน ปฏฺฐปิยมาเน วิวริยมาเน วิภชิยมาเน อุตฺตานีกริยมาเน น ชานาติ น ปสฺสติ, ตมหํ ภิกฺขเว พาลํ ปุถุชฺชนํ อนฺธํ อจกฺขุกํ อชานนฺตํ อปสฺสนฺตํ กินฺติ กโรมิ.\csromanspacer{} เวทนา ภิกฺขเว โลเก โลกธมฺโม ฯเปฯ.\csromanspacer{} สญฺญา ภิกฺขเว.\csromanspacer{} สงฺขารา ภิกฺขเว.\csromanspacer{} วิญฺญาณํ ภิกฺขเว โลเก โลกธมฺโม, ตํ ตถาคโต อภิสมฺพุชฺฌติ อภิสเมติ, อภิสมฺพุชฺฌิตฺวา อภิสเมตฺวา อาจิกฺขติ เทเสติ ปญฺญเปติ ปฏฺฐเปติ วิวรติ วิภชติ อุตฺตานีกโรติ.}

\prose{โย ภิกฺขเว ตถาคเตน เอวํ อาจิกฺขิยมาเน เทสิยมาเน ปญฺญปิยมาเน ปฏฺฐปิยมาเน วิวริยมาเน วิภชิยมาเน อุตฺตานีกริยมาเน น ชานาติ น ปสฺสติ, ตมหํ ภิกฺขเว พาลํ ปุถุชฺชนํ อนฺธํ อจกฺขุกํ อชานนฺตํ อปสฺสนฺตํ กินฺติ กโรมิ.}

\prose{เสยฺยถาปิ ภิกฺขเว อุปฺปลํ วา ปทุมํ วา ปุณฺฑรีกํ วา อุทเก ชาตํ อุทเก สํวฑฺฒํ อุทกา อจฺจุคฺคมฺม ฐาติ\footnote{ติฏฺฐนฺตํ (ก)} อนุปลิตฺตํ อุทเกน.\csromanspacer{} เอวเมว โข ภิกฺขเว ตถาคโต โลเก ชาโต โลเก สํวฑฺโฒ โลกํ อภิภุยฺย วิหรติ อนุปลิตฺโต โลเกนาติ.\csromanspacer{}\csromanspacer{}ทุติยํ.}

\csromanheader{2}{3. เผณปิณฺฑูปมสุตฺต}
\proseitem{95}{เอกํ สมยํ ภควา อยุชฺฌายํ\footnote{อโยชฺฌายํ (สี, อิ)} วิหรติ คงฺคาย นทิยา ตีเร.\csromanspacer{} ตตฺร โข ภควา ภิกฺขู อามนฺเตสิ “เสยฺยถาปิ ภิกฺขเว อยํ คงฺคา นที มหนฺตํ เผณปิณฺฑํ อาวเหยฺย, ตเมนํ จกฺขุมา ปุริโส ปสฺเสยฺย นิชฺฌาเยยฺย โยนิโส อุปปริกฺเขยฺย, ตสฺส ตํ ปสฺสโต นิชฺฌายโต โยนิโส อุปปริกฺขโต ริตฺตกญฺเญว ขาเยยฺย, ตุจฺฉกญฺเญว ขาเยยฺย, อสารกญฺเญว ขาเยยฺย, กิํ หิ สิยา ภิกฺขเว เผณปิณฺเฑ สาโร.\csromanspacer{} เอวเมว โข ภิกฺขเว ยํ กิญฺจิ รูปํ อตีตานาคตปจฺจุปฺปนฺนํ ฯเปฯ ยํ ทูเร สนฺติเก วา, ตํ ภิกฺขุ ปสฺสติ นิชฺฌายติ โยนิโส อุปปริกฺขติ, ตสฺส ตํ ปสฺสโต นิชฺฌายโต โยนิโส อุปปริกฺขโต ริตฺตกญฺเญว \csromanfolio{115}ขายติ, ตุจฺฉกญฺเญว ขายติ, อสารกญฺเญว ขายติ, กิํ หิ สิยา ภิกฺขเว รูเป สาโร.}

\prose{เสยฺยถาปิ ภิกฺขเว สรทสมเย ถุลฺลผุสิตเก เทเว วสฺสนฺเต อุทเก อุทกปุพฺพุฬํ\footnote{อุทกพุพฺพุฬํ (สี, อิ)} อุปฺปชฺชติ เจว นิรุชฺฌติ จ, ตเมนํ จกฺขุมา ปุริโส ปสฺเสยฺย นิชฺฌาเยยฺย โยนิโส อุปปริกฺเขยฺย, ตสฺส ตํ ปสฺสโต นิชฺฌายโต โยนิโส อุปปริกฺขโต ริตฺตกญฺเญว ขาเยยฺย, ตุจฺฉกญฺเญว ขาเยยฺย, อสารกญฺเญว ขาเยยฺย, กิํ หิ สิยา ภิกฺขเว อุทกปุพฺพุเฬ สาโร.\csromanspacer{} เอวเมว โข ภิกฺขเว ยา กาจิ เวทนา อตีตานาคตปจฺจุปฺปนฺนา ฯเปฯ ยา ทูเร สนฺติเก วา, ตํ ภิกฺขุ ปสฺสติ นิชฺฌายติ โยนิโส อุปปริกฺขติ, ตสฺส ตํ ปสฺสโต นิชฺฌายโต โยนิโส อุปปริกฺขโต ริตฺตกญฺเญว ขายติ, ตุจฺฉกญฺเญว ขายติ, อสารกญฺเญว ขายติ, กิํ หิ สิยา ภิกฺขเว เวทนาย สาโร.}

\prose{เสยฺยถาปิ ภิกฺขเว คิมฺหานํ ปจฺฉิเม มาเส ฐิเต มชฺฌนฺหิเก กาเล มรีจิกา ผนฺทติ, ตเมนํ จกฺขุมา ปุริโส ปสฺเสยฺย นิชฺฌาเยยฺย โยนิโส อุปปริกฺเขยฺย, ตสฺส ตํ ปสฺสโต นิชฺฌายโต โยนิโส อุปปริกฺขโต ริตฺตกญฺเญว ขาเยยฺย, ตุจฺฉกญฺเญว ขาเยยฺย ฯเปฯ กิํ หิ สิยา ภิกฺขเว มรีจิกาย สาโร.\csromanspacer{} เอวเมว โข ภิกฺขเว ยา กาจิ สญฺญา ฯเปฯ.}

\prose{เสยฺยถาปิ ภิกฺขเว ปุริโส สารตฺถิโก สารคเวสี สารปริเยสนํ จรมาโน ติณฺหํ กุฐาริํ\footnote{กุธาริํ (สฺยา, กํ, ก)} อาทาย วนํ ปวิเสยฺย.\csromanspacer{} โส ตตฺถ ปสฺเสยฺย มหนฺตํ กทลิกฺขนฺธํ อุชุํ นวํ อกุกฺกุกชาตํ\footnote{อกุกฺกชาตํ (ก-สี, อิ), อกุสชาตํ (ก-สี), อกุกฺกุชกชาตํ (ก)}, ตเมนํ มูเล ฉินฺเทยฺย, มูเล เฉตฺวา อคฺเค ฉินฺเทยฺย, อคฺเค เฉตฺวา ปตฺตวฏฺฏิํ วินิพฺภุเชยฺย, โส ตสฺส ปตฺตวฏฺฏิํ วินิพฺภุชนฺโต เผคฺคุมฺปิ นาธิคจฺเฉยฺย, กุโต สารํ, ตเมนํ จกฺขุมา ปุริโส ปสฺเสยฺย นิชฺฌาเยยฺย โยนิโส อุปปริกฺเขยฺย, ตสฺส ตํ ปสฺสโต นิชฺฌายโต โยนิโส อุปปริกฺขโต ริตฺตกญฺเญว ขาเยยฺย, ตุจฺฉกญฺเญว ขาเยยฺย, อสารกญฺเญว ขาเยยฺย, กิํ หิ สิยา ภิกฺขเว กทลิกฺขนฺเธ สาโร.\csromanspacer{} เอวเมว โข ภิกฺขเว เย เกจิ สงฺขารา อตีตานาคตปจฺจุปฺปนฺนา ฯเปฯ \csromanfolio{116}เย ทูเร สนฺติเก วา, ตํ ภิกฺขุ ปสฺสติ นิชฺฌายติ โยนิโส อุปปริกฺขติ, ตสฺส ตํ ปสฺสโต นิชฺฌายโต โยนิโส อุปปริกฺขโต ริตฺตกญฺเญว ขายติ, ตุจฺฉกญฺเญว ขายติ, อสารกญฺเญว ขายติ, กิํ หิ สิยา ภิกฺขเว สงฺขาเรสุ สาโร.}

\prose{เสยฺยถาปิ ภิกฺขเว มายากาโร วา มายาการนฺเตวาสี วา จตุมหาปเถ\footnote{จาตุมฺมหาปเถ (สี, สฺยา, กํ, อิ)} มายํ วิทํเสยฺย, ตเมนํ จกฺขุมา ปุริโส ปสฺเสยฺย นิชฺฌาเยยฺย โยนิโส อุปปริกฺเขยฺย, ตสฺส ตํ ปสฺสโต นิชฺฌายโต โยนิโส อุปปริกฺขโต ริตฺตกญฺเญว ขาเยยฺย, ตุจฺฉกญฺเญว ขาเยยฺย, อสารกญฺเญว ขาเยยฺย, กิํ หิ สิยา ภิกฺขเว มายาย สาโร.\csromanspacer{} เอวเมว โข ภิกฺขเว ยํ กิญฺจิ วิญฺญาณํ อตีตานาคตปจฺจุปฺปนฺนํ ฯเปฯ ยํ ทูเร สนฺติเก วา, ตํ ภิกฺขุ ปสฺสติ นิชฺฌายติ โยนิโส อุปปริกฺขติ, ตสฺส ตํ ปสฺสโต นิชฺฌายโต โยนิโส อุปปริกฺขโต ริตฺตกญฺเญว ขายติ, ตุจฺฉกญฺเญว ขายติ, อสารกญฺเญว ขายติ, กิํ หิ สิยา ภิกฺขเว วิญฺญาเณ สาโร.}

\prose{เอวํ ปสฺสํ ภิกฺขเว สุตวา อริยสาวโก รูปสฺมิมฺปิ นิพฺพินฺทติ.\csromanspacer{} เวทนายปิ.\csromanspacer{} สญฺญายปิ.\csromanspacer{} สงฺขาเรสุปิ.\csromanspacer{} วิญฺญาณสฺมิมฺปิ นิพฺพินฺทติ, นิพฺพินฺทํ วิรชฺชติ, วิราคา วิมุจฺจติ, วิมุตฺตสฺมิํ “วิมุตฺตมฺ”อิติ ญาณํ โหติ ฯเปฯ นาปรํ อิตฺถตฺตายาติ ปชานาติ.}

\prose{อิทมโวจ ภควา, อิทํ วตฺวาน สุคโต อถาปรํ เอตทโวจ สตฺถา\csromandash{}}

\csromangathameasure{เผณปิณฺฑูปมํ รูปํ, เวทนา ปุพฺพุฬูปมา. \\ มรีจิกูปมา สญฺญา, สงฺขารา กทลูปมา. \\ มายูปมญฺจ วิญฺญาณํ, เทสิตาทิจฺจพนฺธุนา. \\ ยถา ยถา นิชฺฌายติ, โยนิโส อุปปริกฺขติ. \\ ริตฺตกํ ตุจฺฉกํ โหติ, โย นํ ปสฺสติ โยนิโส. \\ อิมญฺจ กายํ อารพฺภ, ภูริปญฺเญน เทสิตํ. \\ ปหานํ ติณฺณํ ธมฺมานํ, รูปํ ปสฺสถ ฉฑฺฑิตํ. \\ อายุ อุสฺมา จ วิญฺญาณํ, ยทา กายํ ชหนฺติมํ. \\ อปวิทฺโธ ตทา เสติ, ปรภตฺตํ อเจตนํ. \\ เอตาทิสายํ สนฺตาโน, มายายํ พาลลาปินี. \\ วธโก เอส อกฺขาโต, สาโร เอตฺถ น วิชฺชติ. \\ เอวํ ขนฺเธ อเวกฺเขยฺย, ภิกฺขุ อารทฺธวีริโย. \\ ทิวา วา ยทิ วา รตฺติํ, สมฺปชาโน ปฏิสฺสโต. \\ ชเหยฺย สพฺพสํโยคํ, กเรยฺย สรณตฺตโน. \\ จเรยฺยาทิตฺตสีโสว, ปตฺถยํ อจฺจุตํ ปทนฺติ. ตติยํ.}
\csromangathagroup{\csromangathastanza{เผณปิณฺฑูปมํ รูปํ, เวทนา ปุพฺพุฬูปมา\footnote{พุพฺพุลูปมา (สี), ปุพฺพุโฬปมา (ก)}. \\ มรีจิกูปมา สญฺญา, สงฺขารา กทลูปมา.}\csromangathastanza{มายูปมญฺจ วิญฺญาณํ, เทสิตาทิจฺจพนฺธุนา. \\ ยถา ยถา นิชฺฌายติ, โยนิโส อุปปริกฺขติ.}\csromangathastanza{ริตฺตกํ ตุจฺฉกํ โหติ, โย นํ ปสฺสติ โยนิโส. \\ อิมญฺจ กายํ อารพฺภ, ภูริปญฺเญน เทสิตํ.}\csromangathastanza{ปหานํ ติณฺณํ ธมฺมานํ, รูปํ ปสฺสถ\footnote{ปสฺเสถ (สี)} ฉฑฺฑิตํ. \\ อายุ อุสฺมา จ วิญฺญาณํ, ยทา กายํ ชหนฺติมํ.}\csromangathastanza{\csromanfolio{117}อปวิทฺโธ\footnote{อปวิฏฺโฐ (สฺยา, กํ)} ตทา เสติ, ปรภตฺตํ อเจตนํ. \\ เอตาทิสายํ สนฺตาโน, มายายํ พาลลาปินี.}\csromangathastanza{วธโก เอส อกฺขาโต, สาโร เอตฺถ น วิชฺชติ. \\ เอวํ ขนฺเธ อเวกฺเขยฺย, ภิกฺขุ อารทฺธวีริโย.}\csromangathastanza{ทิวา วา ยทิ วา รตฺติํ, สมฺปชาโน ปฏิสฺสโต. \\ ชเหยฺย สพฺพสํโยคํ, กเรยฺย สรณตฺตโน. \\ จเรยฺยาทิตฺตสีโสว, ปตฺถยํ อจฺจุตํ ปทนฺติ.\csromanspacer{} ตติยํ.}}

\csromanheader{2}{4. โคมยปิณฺฑสุตฺต}
\proseitem{96}{สาวตฺถินิทานํ.\csromanspacer{} เอกมนฺตํ นิสินฺโน โข โส ภิกฺขุ ภควนฺตํ เอตทโวจ “อตฺถิ นุ โข ภนฺเต กิญฺจิ รูปํ, ยํ รูปํ นิจฺจํ ธุวํ สสฺสตํ อวิปริณามธมฺมํ สสฺสติสมํ ตเถว ฐสฺสติ.\csromanspacer{} อตฺถิ นุ โข ภนฺเต กาจิ เวทนา, ยา เวทนา นิจฺจา ธุวา สสฺสตา อวิปริณามธมฺมา สสฺสติสมํ ตเถว ฐสฺสติ.\csromanspacer{} อตฺถิ นุ โข ภนฺเต กาจิ สญฺญา, ยา สญฺญา ฯเปฯ.\csromanspacer{} อตฺถิ นุ โข ภนฺเต เกจิ สงฺขารา, เย สงฺขารา นิจฺจา ธุวา สสฺสตา อวิปริณามธมฺมา สสฺสติสมํ ตเถว ฐสฺสนฺติ.\csromanspacer{} อตฺถิ นุ โข ภนฺเต กิญฺจิ วิญฺญาณํ, ยํ วิญฺญาณํ นิจฺจํ ธุวํ สสฺสตํ อวิปริณามธมฺมํ สสฺสติสมํ ตเถว ฐสฺสตี”ติ.\csromanspacer{} นตฺถิ โข ภิกฺขุ กิญฺจิ รูปํ, ยํ รูปํ นิจฺจํ ธุวํ สสฺสตํ อวิปริณามธมฺมํ สสฺสติสมํ ตเถว ฐสฺสติ.\csromanspacer{} นตฺถิ โข ภิกฺขุ กาจิ เวทนา.\csromanspacer{} กาจิ สญฺญา.\csromanspacer{} เกจิ สงฺขารา.\csromanspacer{} กิญฺจิ วิญฺญาณํ, ยํ วิญฺญาณํ นิจฺจํ ธุวํ สสฺสตํ อวิปริณามธมฺมํ สสฺสติสมํ ตเถว ฐสฺสตีติ.}

\prose{อถ โข ภควา ปริตฺตํ โคมยปิณฺฑํ ปาณินา คเหตฺวา ตํ ภิกฺขุํ เอตทโวจ\csromandash{}เอตฺตโกปิ โข ภิกฺขุ อตฺตภาวปฏิลาโภ นตฺถิ นิจฺโจ ธุโว สสฺสโต อวิปริณามธมฺโม สสฺสติสมํ ตเถว ฐสฺสติ.\csromanspacer{} เอตฺตโก เจปิ ภิกฺขุ อตฺตภาวปฏิลาโภ อภวิสฺส นิจฺโจ ธุโว \csromanfolio{118}สสฺสโต อวิปริณามธมฺโม, น ยิทํ พฺรหฺมจริยวาโส ปญฺญาเยถ สมฺมา ทุกฺขกฺขยาย.\csromanspacer{} ยสฺมา จ โข ภิกฺขุ เอตฺตโกปิ อตฺตภาวปฏิลาโภ นตฺถิ นิจฺโจ ธุโว สสฺสโต อวิปริณามธมฺโม, ตสฺมา พฺรหฺมจริยวาโส ปญฺญายติ สมฺมา ทุกฺขกฺขยาย.}

\prose{ภูตปุพฺพาหํ ภิกฺขุ ราชา อโหสิํ ขตฺติโย มุทฺธาวสิตฺโต, ตสฺส มยฺหํ ภิกฺขุ รญฺโญ สโต ขตฺติยสฺส มุทฺธาวสิตฺตสฺส จตุราสีตินครสหสฺสานิ อเหสุํ กุสาวตีราชธานิปฺปมุขานิ.\csromanspacer{} ตสฺส มยฺหํ ภิกฺขุ รญฺโญ สโต ขตฺติยสฺส มุทฺธาวสิตฺตสฺส จตุราสีติปาสาทสหสฺสานิ อเหสุํ ธมฺมปาสาทปฺปมุขานิ.\csromanspacer{} ตสฺส มยฺหํ ภิกฺขุ รญฺโญ สโต ขตฺติยสฺส มุทฺธาวสิตฺตสฺส จตุราสีติกูฏาคารสหสฺสานิ อเหสุํ มหาพฺยูหกูฏาคารปฺปมุขานิ\footnote{มหาวิยูหกูฏาคารปฺปมุขานิ (ที 2. 153 ปิฏฺเฐ)}.\csromanspacer{} ตสฺส มยฺหํ ภิกฺขุ รญฺโญ สโต ขตฺติยสฺส มุทฺธาวสิตฺตสฺส จตุราสีติปลฺลงฺกสหสฺสานิ อเหสุํ ทนฺตมยานิ สารมยานิ โสวณฺณมยานิ โคณกตฺถตานิ ปฏิกตฺถตานิ ปฏลิกตฺถตานิ กทลิมิคปวรปจฺจตฺถรณานิ\footnote{กาทลิมิคปวรปจฺจตฺถรณานิ (สี)} ส-อุตฺตรจฺฉทานิ อุภโตโลหิตกูปธานานิ.\csromanspacer{} ตสฺส มยฺหํ ภิกฺขุ รญฺโญ สโต ขตฺติยสฺส มุทฺธาวสิตฺตสฺส จตุราสีตินาคสหสฺสานิ อเหสุํ โสวณฺณาลงฺการานิ โสวณฺณทฺธชานิ เหมชาลปฏิจฺฉนฺนานิ อุโปสถนาคราชปฺปมุขานิ.\csromanspacer{} ตสฺส มยฺหํ ภิกฺขุ รญฺโญ สโต ขตฺติยสฺส มุทฺธาวสิตฺตสฺส จตุราสีตสฺสสหสฺสานิ อเหสุํ โสวณฺณาลงฺการานิ โสวณฺณทฺธชานิ เหมชาลปฏิจฺฉนฺนานิ วลาหก-อสฺสราชปฺปมุขานิ.\csromanspacer{} ตสฺส มยฺหํ ภิกฺขุ รญฺโญ สโต ขตฺติยสฺส มุทฺธาวสิตฺตสฺส จตุราสีติรถสหสฺสานิ อเหสุํ โสวณฺณาลงฺการานิ โสวณฺณทฺธชานิ เหมชาลปฏิจฺฉนฺนานิ เวชยนฺตรถปฺปมุขานิ.\csromanspacer{} ตสฺส มยฺหํ ภิกฺขุ รญฺโญ สโต ขตฺติยสฺส มุทฺธาวสิตฺตสฺส จตุราสีติมณิสหสฺสานิ อเหสุํ มณิรตนปฺปมุขานิ.\csromanspacer{} ตสฺส มยฺหํ ภิกฺขุ ฯเปฯ จตุราสีติ-อิตฺถิสหสฺสานิ อเหสุํ สุภทฺทาเทวิปฺปมุขานิ.\csromanspacer{} ตสฺส มยฺหํ ภิกฺขุ ฯเปฯ จตุราสีติขตฺติยสหสฺสานิ อเหสุํ อนุยนฺตานิ ปริณายกรตนปฺปมุขานิ.\csromanspacer{} ตสฺส มยฺหํ ภิกฺขุ ฯเปฯ จตุราสีติเธนุสหสฺสานิ อเหสุํ ทุกูลสนฺทานานิ กํสูปธารณานิ.\csromanspacer{} ตสฺส มยฺหํ ภิกฺขุ ฯเปฯ จตุราสีติวตฺถโกฏิสหสฺสานิ อเหสุํ โขมสุขุมานิ โกเสยฺยสุขุมานิ \csromanfolio{119}กมฺพลสุขุมานิ กปฺปาสิกสุขุมานิ.\csromanspacer{} ตสฺส มยฺหํ ภิกฺขุ ฯเปฯ จตุราสีติถาลิปากสหสฺสานิ อเหสุํ, สายํ ปาตํ ภตฺตาภิหาโร อภิหริยิตฺถ.}

\prose{เตสํ โข ปน ภิกฺขุ จตุราสีติยา นครสหสฺสานํ เอกญฺเญว ตํ นครํ โหติ, ยมหํ เตน สมเยน อชฺฌาวสามิ กุสาวตี ราชธานี.\csromanspacer{} เตสํ โข ปน ภิกฺขุ จตุราสีติยา ปาสาทสหสฺสานํ เอโกเยว โส ปาสาโท โหติ, ยมหํ เตน สมเยน อชฺฌาวสามิ ธมฺโม ปาสาโท.\csromanspacer{} เตสํ โข ปน ภิกฺขุ จตุราสีติยา กูฏาคารสหสฺสานํ เอกญฺเญว ตํ กูฏาคารํ โหติ, ยมหํ เตน สมเยน อชฺฌาวสามิ มหาพฺยูหํ กูฏาคารํ.\csromanspacer{} เตสํ โข ปน ภิกฺขุ จตุราสีติยา ปลฺลงฺกสหสฺสานํ เอโกเยว โส ปลฺลงฺโก โหติ, ยมหํ เตน สมเยน ปริภุญฺชามิ ทนฺตมโย วา สารมโย วา โสวณฺณมโย วา รูปิยมโย วา.\csromanspacer{} เตสํ โข ปน ภิกฺขุ จตุราสีติยา นาคสหสฺสานํ เอโกเยว โส นาโค โหติ, ยมหํ เตน สมเยน อภิรุหามิ อุโปสโถ นาคราชา.\csromanspacer{} เตสํ โข ปน ภิกฺขุ จตุราสีติยา อสฺสสหสฺสานํ เอโกเยว โส อสฺโส โหติ, ยมหํ เตน สมเยน อภิรุหามิ วลาหโก อสฺสราชา.\csromanspacer{} เตสํ โข ปน ภิกฺขุ จตุราสีติยา รถสหสฺสานํ เอโกเยว โส รโถ โหติ, ยมหํ เตน สมเยน อภิรุหามิ เวชยนฺโต รโถ.\csromanspacer{} เตสํ โข ปน ภิกฺขุ จตุราสีติยา อิตฺถิสหสฺสานํ เอกาเยว สา อิตฺถิ โหติ, ยา มํ เตน สมเยน ปจฺจุปฏฺฐาติ ขตฺติยานี วา เวลามิกา วา.\csromanspacer{} เตสํ โข ปน ภิกฺขุ จตุราสีติยา วตฺถโกฏิสหสฺสานํ เอกญฺเญว ตํ วตฺถยุคํ โหติ, ยมหํ เตน สมเยน ปริทหามิ โขมสุขุมํ วา โกเสยฺยสุขุมํ วา กมฺพลสุขุมํ วา กปฺปาสิกสุขุมํ วา.\csromanspacer{} เตสํ โข ปน ภิกฺขุ จตุราสีติยา ถาลิปากสหสฺสานํ เอโกเยว โส ถาลิปาโก โหติ, ยโต นาฬิโกทนปรมํ ภุญฺชามิ, ตทุปิยญฺจ สูเปยฺยํ\footnote{สูปพฺยญฺชนํ (สฺยา, กํ)}.\csromanspacer{} อิติ โข ภิกฺขุ สพฺเพเต สงฺขารา อตีตา นิรุทฺธา วิปริณตา.\csromanspacer{} เอวํ อนิจฺจา โข ภิกฺขุ สงฺขารา, เอวํ อทฺธุวา โข ภิกฺขุ สงฺขารา, เอวํ อนสฺสาสิกา โข ภิกฺขุ สงฺขารา.\csromanspacer{} ยาวญฺจิทํ ภิกฺขุ อลเมว สพฺพสงฺขาเรสุ นิพฺพินฺทิตุํ, อลํ วิรชฺชิตุํ, อลํ วิมุจฺจิตุนฺติ.\csromanspacer{}\csromanspacer{}จตุตฺถํ.}

\csromanheader{2}{5. นขสิขาสุตฺต}
\proseitem{97}{\csromanfolio{120}สาวตฺถินิทานํ.\csromanspacer{} เอกมนฺตํ นิสินฺโน โข โส ภิกฺขุ ภควนฺตํ เอตทโวจ “อตฺถิ นุ โข ภนฺเต กิญฺจิ รูปํ, ยํ รูปํ นิจฺจํ ธุวํ สสฺสตํ อวิปริณามธมฺมํ สสฺสติสมํ ตเถว ฐสฺสติ.\csromanspacer{} อตฺถิ นุ โข ภนฺเต กาจิ เวทนา, ยา เวทนา นิจฺจา ธุวา สสฺสตา อวิปริณามธมฺมา สสฺสติสมํ ตเถว ฐสฺสติ.\csromanspacer{} อตฺถิ นุ โข ภนฺเต กาจิ สญฺญา ฯเปฯ เกจิ สงฺขารา, เย สงฺขารา นิจฺจา ธุวา สสฺสตา อวิปริณามธมฺมา สสฺสติสมํ ตเถว ฐสฺสนฺติ.\csromanspacer{} อตฺถิ นุ โข ภนฺเต กิญฺจิ วิญฺญาณํ, ยํ วิญฺญาณํ นิจฺจํ ธุวํ สสฺสตํ อวิปริณามธมฺมํ สสฺสติสมํ ตเถว ฐสฺสตี”ติ.\csromanspacer{} นตฺถิ โข ภิกฺขุ กิญฺจิ รูปํ, ยํ รูปํ นิจฺจํ ธุวํ สสฺสตํ อวิปริณามธมฺมํ สสฺสติสมํ ตเถว ฐสฺสติ.\csromanspacer{} นตฺถิ โข ภิกฺขุ กาจิ เวทนา.\csromanspacer{} กาจิ สญฺญา.\csromanspacer{} เกจิ สงฺขารา ฯเปฯ กิญฺจิ วิญฺญาณํ, ยํ วิญฺญาณํ นิจฺจํ ธุวํ สสฺสตํ อวิปริณามธมฺมํ สสฺสติสมํ ตเถว ฐสฺสตีติ.}

\prose{อถ โข ภควา ปริตฺตํ นขสิขายํ ปํสุํ อาโรเปตฺวา ตํ ภิกฺขุํ เอตทโวจ\csromandash{}เอตฺตกมฺปิ โข ภิกฺขุ รูปํ นตฺถิ นิจฺจํ ธุวํ สสฺสตํ อวิปริณามธมฺมํ สสฺสติสมํ ตเถว ฐสฺสติ.\csromanspacer{} เอตฺตกํ เจปิ ภิกฺขุ รูปํ อภวิสฺส นิจฺจํ ธุวํ สสฺสตํ อวิปริณามธมฺมํ, น ยิทํ พฺรหฺมจริยวาโส ปญฺญาเยถ สมฺมา ทุกฺขกฺขยาย.\csromanspacer{} ยสฺมา จ โข ภิกฺขุ เอตฺตกมฺปิ รูปํ นตฺถิ นิจฺจํ ธุวํ สสฺสตํ อวิปริณามธมฺมํ, ตสฺมา พฺรหฺมจริยวาโส ปญฺญายติ สมฺมา ทุกฺขกฺขยาย.}

\prose{เอตฺตกาปิ โข ภิกฺขุ เวทนา นตฺถิ นิจฺจา ธุวา สสฺสตา อวิปริณามธมฺมา สสฺสติสมํ ตเถว ฐสฺสติ.\csromanspacer{} เอตฺตกา เจปิ ภิกฺขุ เวทนา อภวิสฺส นิจฺจา ธุวา สสฺสตา อวิปริณามธมฺมา, น ยิทํ พฺรหฺมจริยวาโส ปญฺญาเยถ สมฺมา ทุกฺขกฺขยาย.\csromanspacer{} ยสฺมา จ โข ภิกฺขุ เอตฺตกาปิ เวทนา นตฺถิ นิจฺจา ธุวา สสฺสตา อวิปริณามธมฺมา, ตสฺมา พฺรหฺมจริยวาโส ปญฺญายติ สมฺมา ทุกฺขกฺขยาย.}

\prose{เอตฺตกาปิ โข ภิกฺขุ สญฺญา นตฺถิ ฯเปฯ.\csromanspacer{} เอตฺตกาปิ โข ภิกฺขุ สงฺขารา นตฺถิ นิจฺจา ธุวา สสฺสตา อวิปริณามธมฺมา สสฺสติสมํ ตเถว ฐสฺสนฺติ.\csromanspacer{} เอตฺตกา เจปิ ภิกฺขุ สงฺขารา อภวิสฺสํสุ นิจฺจา ธุวา สสฺสตา อวิปริณามธมฺมา, น ยิทํ พฺรหฺมจริยวาโส ปญฺญาเยถ สมฺมา ทุกฺขกฺขยาย.\csromanspacer{} ยสฺมา จ โข ภิกฺขุ เอตฺตกาปิ สงฺขารา นตฺถิ นิจฺจา ธุวา \csromanfolio{121}สสฺสตา อวิปริณามธมฺมา, ตสฺมา พฺรหฺมจริยวาโส ปญฺญายติ สมฺมา ทุกฺขกฺขยาย.}

\prose{เอตฺตกมฺปิ โข ภิกฺขุ วิญฺญาณํ นตฺถิ นิจฺจํ ธุวํ สสฺสตํ อวิปริณามธมฺมํ สสฺสติสมํ ตเถว ฐสฺสติ.\csromanspacer{} เอตฺตกมฺปิ โข ภิกฺขุ วิญฺญาณํ อภวิสฺส นิจฺจํ ธุวํ สสฺสตํ อวิปริณามธมฺมํ, น ยิทํ พฺรหฺมจริยวาโส ปญฺญาเยถ สมฺมา ทุกฺขกฺขยาย.\csromanspacer{} ยสฺมา จ โข ภิกฺขุ เอตฺตกมฺปิ วิญฺญาณํ นตฺถิ นิจฺจํ ธุวํ สสฺสตํ อวิปริณามธมฺมํ, ตสฺมา พฺรหฺมจริยวาโส ปญฺญายติ สมฺมา ทุกฺขกฺขยาย.}

\prose{ตํ กิํ มญฺญสิ ภิกฺขุ, รูปํ นิจฺจํ วา อนิจฺจํ วาติ.\csromanspacer{} อนิจฺจํ ภนฺเต.\csromanspacer{} เวทนา.\csromanspacer{} สญฺญา.\csromanspacer{} สงฺขารา.\csromanspacer{} วิญฺญาณํ นิจฺจํ วา อนิจฺจํ วาติ.\csromanspacer{} อนิจฺจํ ภนฺเต ฯเปฯ.\csromanspacer{} ตสฺมาติห.\csromanspacer{} เอวํ ปสฺสํ ฯเปฯ นาปรํ อิตฺถตฺตายาติ ปชานาตีติ.\csromanspacer{}\csromanspacer{}ปญฺจมํ.}

\csromanheader{1}{6. สุทฺธิกสุตฺต}
\proseitem{98}{สาวตฺถินิทานํ.\csromanspacer{} เอกมนฺตํ นิสินฺโน โข โส ภิกฺขุ ภควนฺตํ เอตทโวจ “อตฺถิ นุ โข ภนฺเต กิญฺจิ รูปํ, ยํ รูปํ นิจฺจํ ธุวํ สสฺสตํ อวิปริณามธมฺมํ สสฺสติสมํ ตเถว ฐสฺสติ.\csromanspacer{} อตฺถิ นุ โข ภนฺเต กาจิ เวทนา ฯเปฯ กาจิ สญฺญา.\csromanspacer{} เกจิ สงฺขารา.\csromanspacer{} กิญฺจิ วิญฺญาณํ, ยํ วิญฺญาณํ นิจฺจํ ธุวํ สสฺสตํ อวิปริณามธมฺมํ สสฺสติสมํ ตเถว ฐสฺสตี”ติ.\csromanspacer{} นตฺถิ โข ภิกฺขุ กิญฺจิ รูปํ, ยํ รูปํ นิจฺจํ ธุวํ สสฺสตํ อวิปริณามธมฺมํ สสฺสติสมํ ตเถว ฐสฺสติ.\csromanspacer{} นตฺถิ โข ภิกฺขุ กาจิ เวทนา ฯเปฯ กาจิ สญฺญา.\csromanspacer{} เกจิ สงฺขารา.\csromanspacer{} กิญฺจิ วิญฺญาณํ, ยํ วิญฺญาณํ นิจฺจํ ธุวํ สสฺสตํ อวิปริณามธมฺมํ สสฺสติสมํ ตเถว ฐสฺสตีติ.\csromanspacer{}\csromanspacer{}ฉฏฺฐํ.}

\csromanheader{1}{7. คทฺทุลพทฺธสุตฺต}
\proseitem{99}{สาวตฺถินิทานํ.\csromanspacer{} อนมตคฺโคยํ ภิกฺขเว สํสาโร, ปุพฺพา โกฏิ น ปญฺญายติ อวิชฺชานีวรณานํ สตฺตานํ ตณฺหาสํโยชนานํ สนฺธาวตํ สํสรตํ.\csromanspacer{} โหติ โส ภิกฺขเว สมโย ยํ มหาสมุทฺโท อุสฺสุสฺสติ วิสุสฺสติ น ภวติ, น เตฺววาหํ ภิกฺขเว อวิชฺชานีวรณานํ สตฺตานํ ตณฺหาสํโยชนานํ สนฺธาวตํ สํสรตํ ทุกฺขสฺส อนฺตกิริยํ วทามิ.\csromanspacer{} โหติ โส ภิกฺขเว สมโย ยํ สิเนรุ ปพฺพตราชา ฑยฺหติ วินสฺสติ \csromanfolio{122}น ภวติ, น เตฺววาหํ ภิกฺขเว อวิชฺชานีวรณานํ สตฺตานํ ตณฺหาสํโยชนานํ สนฺธาวตํ สํสรตํ ทุกฺขสฺส อนฺตกิริยํ วทามิ.\csromanspacer{} โหติ โส ภิกฺขเว สมโย ยํ มหาปถวี ฑยฺหติ วินสฺสติ น ภวติ, น เตฺววาหํ ภิกฺขเว อวิชฺชานีวรณานํ สตฺตานํ ตณฺหาสํโยชนานํ สนฺธาวตํ สํสรตํ ทุกฺขสฺส อนฺตกิริยํ วทามิ.}

\prose{เสยฺยถาปิ ภิกฺขเว สา คทฺทุลพทฺโธ\footnote{คทฺทูลพนฺโธ (สฺยา, กํ)} ทเฬฺห ขีเล วา ถมฺเภ วา อุปนิพทฺโธ, ตเมว ขีลํ วา ถมฺภํ วา อนุปริธาวติ อนุปริวตฺตติ.\csromanspacer{} เอวเมว โข ภิกฺขเว อสฺสุตวา ปุถุชฺชโน อริยานํ อทสฺสาวี ฯเปฯ สปฺปุริสธมฺเม อวินีโต รูปํ อตฺตโต สมนุปสฺสติ ฯเปฯ.\csromanspacer{} เวทนํ อตฺตโต สมนุปสฺสติ.\csromanspacer{} สญฺญํ อตฺตโต สมนุปสฺสติ.\csromanspacer{} สงฺขาเร อตฺตโต สมนุปสฺสติ.\csromanspacer{} วิญฺญาณํ อตฺตโต สมนุปสฺสติ, วิญฺญาณวนฺตํ วา อตฺตานํ, อตฺตนิ วา วิญฺญาณํ, วิญฺญาณสฺมิํ วา อตฺตานํ.\csromanspacer{} โส รูปญฺเญว อนุปริธาวติ อนุปริวตฺตติ.\csromanspacer{} เวทนญฺเญว ฯเปฯ.\csromanspacer{} สญฺญญฺเญว.\csromanspacer{} สงฺขาเรเยว.\csromanspacer{} วิญฺญาณญฺเญว อนุปริธาวติ อนุปริวตฺตติ.\csromanspacer{} โส รูปํ อนุปริธาวํ อนุปริวตฺตํ.\csromanspacer{} เวทนํ ฯเปฯ.\csromanspacer{} สญฺญํ.\csromanspacer{} สงฺขาเร.\csromanspacer{} วิญฺญาณํ อนุปริธาวํ อนุปริวตฺตํ น ปริมุจฺจติ รูปมฺหา, น ปริมุจฺจติ เวทนาย, น ปริมุจฺจติ สญฺญาย, น ปริมุจฺจติ สงฺขาเรหิ, น ปริมุจฺจติ วิญฺญาณมฺหา, น ปริมุจฺจติ ชาติยา ชรามรเณน โสเกหิ ปริเทเวหิ ทุกฺเขหิ โทมนสฺเสหิ อุปายาเสหิ, น ปริมุจฺจติ ทุกฺขสฺมาติ วทามิ.}

\prose{สุตวา จ โข ภิกฺขเว อริยสาวโก อริยานํ ทสฺสาวี ฯเปฯ สปฺปุริสธมฺเม สุวินีโต น รูปํ อตฺตโต สมนุปสฺสติ ฯเปฯ น เวทนํ.\csromanspacer{} น สญฺญํ.\csromanspacer{} น สงฺขาเร.\csromanspacer{} น วิญฺญาณํ อตฺตโต สมนุปสฺสติ, น วิญฺญาณวนฺตํ วา อตฺตานํ, น อตฺตนิ วา วิญฺญาณํ, น วิญฺญาณสฺมิํ วา อตฺตานํ.\csromanspacer{} โส รูปํ นานุปริธาวติ นานุปริวตฺตติ.\csromanspacer{} เวทนํ.\csromanspacer{} สญฺญํ.\csromanspacer{} สงฺขาเร.\csromanspacer{} วิญฺญาณํ นานุปริธาวติ นานุปริวตฺตติ.\csromanspacer{} โส รูปํ อนนุปริธาวํ อนนุปริวตฺตํ.\csromanspacer{} เวทนํ.\csromanspacer{} สญฺญํ.\csromanspacer{} สงฺขาเร.\csromanspacer{} วิญฺญาณํ อนนุปริธาวํ อนนุปริวตฺตํ ปริมุจฺจติ รูปมฺหา, ปริมุจฺจติ เวทนาย, ปริมุจฺจติ สญฺญาย, ปริมุจฺจติ สงฺขาเรหิ, ปริมุจฺจติ วิญฺญาณมฺหา, ปริมุจฺจติ ชาติยา ชรามรเณน โสเกหิ ปริเทเวหิ ทุกฺเขหิ โทมนสฺเสหิ อุปายาเสหิ, ปริมุจฺจติ ทุกฺขสฺมาติ วทามีติ.\csromanspacer{}\csromanspacer{}สตฺตมํ.}

\csromanheader{2}{8. ทุติยคทฺทุลพทฺธสุตฺต}
\proseitem{100}{\csromanfolio{123}สาวตฺถินิทานํ.\csromanspacer{} อนมตคฺโคยํ ภิกฺขเว สํสาโร, ปุพฺพา โกฏิ น ปญฺญายติ อวิชฺชานีวรณานํ สตฺตานํ ตณฺหาสํโยชนานํ สนฺธาวตํ สํสรตํ.\csromanspacer{} เสยฺยถาปิ ภิกฺขเว สา คทฺทุลพทฺโธ ทเฬฺห ขีเล วา ถมฺเภ วา อุปนิพทฺโธ.\csromanspacer{} โส คจฺฉติ เจปิ, ตเมว ขีลํ วา ถมฺภํ วา อุปคจฺฉติ.\csromanspacer{} ติฏฺฐติ เจปิ, ตเมว ขีลํ วา ถมฺภํ วา อุปติฏฺฐติ.\csromanspacer{} นิสีทติ เจปิ, ตเมว ขีลํ วา ถมฺภํ วา อุปนิสีทติ.\csromanspacer{} นิปชฺชติ เจปิ, ตเมว ขีลํ วา ถมฺภํ วา อุปนิปชฺชติ.\csromanspacer{} เอวเมว โข ภิกฺขเว อสฺสุตวา ปุถุชฺชโน รูปํ “เอตํ มม, เอโสหมสฺมิ, เอโส เม อตฺตา”ติ สมนุปสฺสติ.\csromanspacer{} เวทนํ.\csromanspacer{} สญฺญํ.\csromanspacer{} สงฺขาเร.\csromanspacer{} วิญฺญาณํ “เอตํ มม, เอโสหมสฺมิ, เอโส เม อตฺตา”ติ สมนุปสฺสติ.\csromanspacer{} โส คจฺฉติ เจปิ, อิเม ปญฺจุปาทานกฺขนฺเธ อุปคจฺฉติ.\csromanspacer{} ติฏฺฐติ เจปิ, อิเม ปญฺจุปาทานกฺขนฺเธ อุปติฏฺฐติ.\csromanspacer{} นิสีทติ เจปิ, อิเม ปญฺจุปาทานกฺขนฺเธ อุปนิสีทติ.\csromanspacer{} นิปชฺชติ เจปิ, อิเม ปญฺจุปาทานกฺขนฺเธ อุปนิปชฺชติ.\csromanspacer{} ตสฺมาติห ภิกฺขเว อภิกฺขณํ สกํ จิตฺตํ ปจฺจเวกฺขิตพฺพํ “ทีฆรตฺตมิทํ จิตฺตํ สํกิลิฏฺฐํ ราเคน โทเสน โมเหนา”ติ.\csromanspacer{} จิตฺตสํกิเลสา ภิกฺขเว สตฺตา สํกิลิสฺสนฺติ, จิตฺตโวทานา สตฺตา วิสุชฺฌนฺติ.}

\prose{ทิฏฺฐํ โว ภิกฺขเว จรณํ นาม จิตฺตนฺติ.\csromanspacer{} เอวํ ภนฺเต.\csromanspacer{} ตมฺปิ โข ภิกฺขเว จรณํ นาม จิตฺตํ จิตฺเตเนว จิตฺติตํ, เตนปิ โข ภิกฺขเว จรเณน จิตฺเตน จิตฺตญฺเญว จิตฺตตรํ.\csromanspacer{} ตสฺมาติห ภิกฺขเว อภิกฺขณํ สกํ จิตฺตํ ปจฺจเวกฺขิตพฺพํ “ทีฆรตฺตมิทํ จิตฺตํ สํกิลิฏฺฐํ ราเคน โทเสน โมเหนา”ติ.\csromanspacer{} จิตฺตสํกิเลสา ภิกฺขเว สตฺตา สํกิลิสฺสนฺติ, จิตฺตโวทานา สตฺตา วิสุชฺฌนฺติ.}

\prose{นาหํ ภิกฺขเว อญฺญํ เอกนิกายมฺปิ สมนุปสฺสามิ เอวํ จิตฺตํ ยถยิทํ ภิกฺขเว ติรจฺฉานคตา ปาณา, เตปิ โข ภิกฺขเว ติรจฺฉานคตา ปาณา จิตฺเตเนว จิตฺติตา.\csromanspacer{} เตหิปิ โข ภิกฺขเว ติรจฺฉานคเตหิ ปาเณหิ จิตฺตญฺเญว จิตฺตตรํ.\csromanspacer{} ตสฺมาติห ภิกฺขเว อภิกฺขณํ สกํ จิตฺตํ ปจฺจเวกฺขิตพฺพํ “ทีฆรตฺตมิทํ จิตฺตํ สํกิลิฏฺฐํ ราเคน โทเสน โมเหนา”ติ.\csromanspacer{} จิตฺตสํกิเลสา ภิกฺขเว สตฺตา สํกิลิสฺสนฺติ, จิตฺตโวทานา สตฺตา วิสุชฺฌนฺติ.}

\prose{\csromanfolio{124}เสยฺยถาปิ ภิกฺขเว รชโก วา จิตฺตการโก วา รชนาย วา ลาขาย วา หลิทฺทิยา วา นีลิยา วา มญฺชิฏฺฐาย\footnote{มญฺเชฏฺฐาย (สี, สฺยา, กํ), มญฺเชฏฺฐิยา (อิ)} วา สุปริมฏฺเฐ ผลเก วา ภิตฺติยา วา ทุสฺสปฏฺเฏ วา อิตฺถิรูปํ วา ปุริสรูปํ วา อภินิมฺมิเนยฺย สพฺพงฺคปจฺจงฺคิํ.\csromanspacer{} เอวเมว โข ภิกฺขเว อสฺสุตวา ปุถุชฺชโน รูปญฺเญว อภินิพฺพตฺเตนฺโต อภินิพฺพตฺเตติ.\csromanspacer{} เวทนญฺเญว ฯเปฯ.\csromanspacer{} สญฺญญฺเญว.\csromanspacer{} สงฺขาเรเยว.\csromanspacer{} วิญฺญาณญฺเญว อภินิพฺพตฺเตนฺโต อภินิพฺพตฺเตติ.\csromanspacer{} ตํ กิํ มญฺญถ ภิกฺขเว, รูปํ นิจฺจํ วา อนิจฺจํ วาติ.\csromanspacer{} อนิจฺจํ ภนฺเต.\csromanspacer{} เวทนา.\csromanspacer{} สญฺญา.\csromanspacer{} สงฺขารา.\csromanspacer{} วิญฺญาณํ ฯเปฯ.\csromanspacer{} ตสฺมาติห ภิกฺขเว ฯเปฯ.\csromanspacer{} เอวํ ปสฺสํ ฯเปฯ นาปรํ อิตฺถตฺตายาติ ปชานาตีติ.\csromanspacer{}\csromanspacer{}อฏฺฐมํ.}

\csromanmatikaanchor{mtk.3}
\csromantocmarkref{section}{9. วาสิชฏสุตฺต}{mtk.3}
\bookmark[dest={mtk.3},level=1]{9. วาสิชฏสุตฺต}
\csromanheader{1}{9. วาสิชฏสุตฺต}
\proseitem{101}{สาวตฺถินิทานํ.\csromanspacer{} ชานโต อหํ ภิกฺขเว ปสฺสโต อาสวานํ ขยํ วทามิ โน อชานโต โน อปสฺสโต.\csromanspacer{} กิญฺจิ ภิกฺขเว ชานโต กิํ ปสฺสโต อาสวานํ ขโย โหติ. “อิติ รูปํ, อิติ รูปสฺส สมุทโย, อิติ รูปสฺส อตฺถงฺคโม.\csromanspacer{} อิติ เวทนา ฯเปฯ อิติ สญฺญา.\csromanspacer{} อิติ สงฺขารา.\csromanspacer{} อิติ วิญฺญาณํ, อิติ วิญฺญาณสฺส สมุทโย, อิติ วิญฺญาณสฺส อตฺถงฺคโม”ติ, เอวํ โข ภิกฺขเว ชานโต เอวํ ปสฺสโต อาสวานํ ขโย โหติ.}

\prose{ภาวนานุโยคํ อนนุยุตฺตสฺส ภิกฺขเว ภิกฺขุโน วิหรโต กิญฺจาปิ เอวํ อิจฺฉา อุปฺปชฺเชยฺย “อโห วต เม อนุปาทาย อาสเวหิ จิตฺตํ วิมุจฺเจยฺยา”ติ, อถ ขฺวสฺส เนว อนุปาทาย อาสเวหิ จิตฺตํ วิมุจฺจติ.\csromanspacer{} ตํ กิสฺส เหตุ, “อภาวิตตฺตา”ติสฺส วจนียํ.\csromanspacer{} กิสฺส อภาวิตตฺตา, อภาวิตตฺตา จตุนฺนํ สติปฏฺฐานานํ, อภาวิตตฺตา จตุนฺนํ สมฺมปฺปธานํ, อภาวิตตฺตา จตุนฺนํ อิทฺธิปาทานํ, อภาวิตตฺตา ปญฺจนฺนํ อินฺทฺริยานํ, อภาวิตตฺตา ปญฺจนฺนํ พลานํ, อภาวิตตฺตา สตฺตนฺนํ โพชฺฌงฺคานํ, อภาวิตตฺตา อริยสฺส อฏฺฐงฺคิกสฺส มคฺคสฺส.}

\prose{เสยฺยถาปิ ภิกฺขเว กุกฺกุฏิยา อณฺฑานิ อฏฺฐ วา ทส วา ทฺวาทส วา, ตานสฺสุ กุกฺกุฏิยา น สมฺมา อธิสยิตานิ, น สมฺมา ปริเสทิตานิ, น สมฺมา ปริภาวิตานิ.\csromanspacer{} กิญฺจาปิ ตสฺสา กุกฺกุฏิยา เอวํ อิจฺฉา อุปฺปชฺเชยฺย “อโห วต \csromanfolio{125}เม กุกฺกุฏโปตกา ปาทนขสิขาย วา มุขตุณฺฑเกน วา อณฺฑโกสํ ปทาเลตฺวา โสตฺถินา อภินิพฺภิชฺเชยฺยุนฺ”ติ.\csromanspacer{} อถ โข อภพฺพาว เต กุกฺกุฏโปตกา ปาทนขสิขาย วา มุขตุณฺฑเกน วา อณฺฑโกสํ ปทาเลตฺวา โสตฺถินา อภินิพฺภิชฺชิตุํ.\csromanspacer{} ตํ กิสฺส เหตุ, ตถา หิ ปน ภิกฺขเว กุกฺกุฏิยา อณฺฑานิ อฏฺฐ วา ทส วา ทฺวาทส วา, ตานิ กุกฺกุฏิยา น สมฺมา อธิสยิตานิ, น สมฺมา ปริเสทิตานิ, น สมฺมา ปริภาวิตานิ.\csromanspacer{} เอวเมว โข ภิกฺขเว ภาวนานุโยคํ อนนุยุตฺตสฺส ภิกฺขุโน วิหรโต กิญฺจาปิ เอวํ อิจฺฉา อุปฺปชฺเชยฺย “อโห วต เม อนุปาทาย อาสเวหิ จิตฺตํ วิมุจฺเจยฺยา”ติ, อถ ขฺวสฺส เนว อนุปาทาย อาสเวหิ จิตฺตํ วิมุจฺจติ.\csromanspacer{} ตํ กิสฺส เหตุ, “อภาวิตตฺตา”ติสฺส วจนียํ.\csromanspacer{} กิสฺส อภาวิตตฺตา, อภาวิตตฺตา จตุนฺนํ สติปฏฺฐานานํ ฯเปฯ อฏฺฐงฺคิกสฺส มคฺคสฺส}

\prose{ภาวนานุโยคํ อนุยุตฺตสฺส ภิกฺขเว ภิกฺขุโน วิหรโต กิญฺจาปิ น เอวํ อิจฺฉา อุปฺปชฺเชยฺย “อโห วต เม อนุปาทาย อาสเวหิ จิตฺตํ วิมุจฺเจยฺยา”ติ, อถ ขฺวสฺส อนุปาทาย อาสเวหิ จิตฺตํ วิมุจฺจติ.\csromanspacer{} ตํ กิสฺส เหตุ, “ภาวิตตฺตา”ติสฺส วจนียํ.\csromanspacer{} กิสฺส ภาวิตตฺตา, ภาวิตตฺตา จตุนฺนํ สติปฏฺฐานานํ, ภาวิตตฺตา จตุนฺนํ สมฺมปฺปธานานํ, ภาวิตตฺตา จตุนฺนํ อิทฺธิปาทานํ, ภาวิตตฺตา ปญฺจนฺนํ อินฺทฺริยานํ, ภาวิตตฺตา ปญฺจนฺนํ พลานํ, ภาวิตตฺตา สตฺตนฺนํ โพชฺฌงฺคานํ, ภาวิตตฺตา อริยสฺส อฏฺฐงฺคิกสฺส มคฺคสฺส.}

\prose{เสยฺยถาปิ ภิกฺขเว กุกฺกุฏิยา อณฺฑานิ อฏฺฐ วา ทส วา ทฺวาทส วา, ตานสฺสุ กุกฺกุฏิยา สมฺมา อธิสยิตานิ, สมฺมา ปริเสทิตานิ, สมฺมา ปริภาวิตานิ, กิญฺจาปิ ตสฺสา กุกฺกุฏิยา น เอวํ อิจฺฉา อุปฺปชฺเชยฺย “อโห วต เม กุกฺกุฏโปตกา ปาทนขสิขาย วา มุขตุณฺฑเกน วา อณฺฑโกสํ ปทาเลตฺวา โสตฺถินา อภินิพฺภิชฺเชยฺยุนฺ”ติ.\csromanspacer{} อถ โข ภพฺพาว เต กุกฺกุฏโปตกา ปาทนขสิขาย วา มุขตุณฺฑเกน วา อณฺฑโกสํ ปทาเลตฺวา โสตฺถินา อภินิพฺภิชฺชิตุํ.\csromanspacer{} ตํ กิสฺส เหตุ, ตถา หิ ปน ภิกฺขเว กุกฺกุฏิยา อณฺฑานิ อฏฺฐ วา ทส วา ทฺวาทส วา, ตานสฺสุ กุกฺกุฏิยา สมฺมา อธิสยิตานิ, สมฺมา ปริเสทิตานิ, สมฺมา ปริภาวิตานิ.\csromanspacer{} เอวเมว โข ภิกฺขเว ภาวนานุโยคํ อนุยุตฺตสฺส ภิกฺขุโน วิหรโต กิญฺจาปิ น เอวํ อิจฺฉา อุปฺปชฺเชยฺย “อโห วต เม อนุปาทาย อาสเวหิ จิตฺตํ วิมุจฺเจยฺยา”ติ, อถ ขฺวสฺส อนุปาทาย อาสเวหิ จิตฺตํ วิมุจฺจติ.\csromanspacer{} ตํ กิสฺส เหตุ, “ภาวิตตฺตา”ติสฺส วจนียํ. \csromanfolio{126}กิสฺส ภาวิตตฺตา, ภาวิตตฺตา จตุนฺนํ สติปฏฺฐานานํ ฯเปฯ ภาวิตตฺตา อริยสฺส อฏฺฐงฺคิกสฺส มคฺคสฺส.}

\prose{เสยฺยถาปิ ภิกฺขเว ปลคณฺฑสฺส วา ปลคณฺฑนฺเตวาสิสฺส วา วาสิชเฏ ทิสฺสนฺเตว องฺคุลิปทานิ, ทิสฺสติ องฺคุฏฺฐปทํ.\csromanspacer{} โน จ ขฺวสฺส เอวํ ญาณํ โหติ “เอตฺตกํ วต เม อชฺช วาสิชฏสฺส ขีณํ, เอตฺตกํ หิยฺโย, เอตฺตกํ ปเร”ติ.\csromanspacer{} อถ ขฺวสฺส ขีเณ ขีณนฺเตฺวว ญาณํ โหติ.\csromanspacer{} เอวเมว โข ภิกฺขเว ภาวนานุโยคํ อนุยฺยุตฺตสฺส ภิกฺขุโน วิหรโต กิญฺจาปิ น เอวํ ญาณํ โหติ “เอตฺตกํ วต เม อชฺช อาสวานํ ขีณํ, เอตฺตกํ หิยฺโย, เอตฺตกํ ปเร”ติ.\csromanspacer{} อถ ขฺวสฺส ขีเณ ขีณนฺเตฺวว ญาณํ โหติ.\csromanspacer{} เสยฺยถาปิ ภิกฺขเว สามุทฺทิกาย นาวาย เวตฺตพนฺธนพทฺธาย วสฺสมาสานิ อุทเก ปริยาทาย เหมนฺติเกน ถลํ อุกฺขิตฺตาย วาตาตปปเรตานิ เวตฺตพนฺธนานิ, ตานิ ปาวุสเกน เมเฆน อภิปฺปวุฏฺฐานิ อปฺปกสิเรเนว ปฏิปฺปสฺสมฺภนฺติ, ปูติกานิ ภวนฺติ.\csromanspacer{} เอวเมว โข ภิกฺขเว ภาวนานุโยคํ อนุยุตฺตสฺส ภิกฺขุโน วิหรโต อปฺปกสิเรเนว สํโยชนานิ ปฏิปฺปสฺสมฺภนฺติ, ปูติกานิ ภวนฺตีติ.\csromanspacer{}\csromanspacer{}นวมํ.}

\csromanmatikaanchor{mtk.4}
\csromantocmarkref{section}{10. อนิจฺจสญฺญาสุตฺต}{mtk.4}
\bookmark[dest={mtk.4},level=1]{10. อนิจฺจสญฺญาสุตฺต}
\csromanheader{1}{10. อนิจฺจสญฺญาสุตฺต}
\proseitem{102}{สาวตฺถินิทานํ.\csromanspacer{} อนิจฺจสญฺญา ภิกฺขเว ภาวิตา พหุลีกตา สพฺพํ กามราคํ ปริยาทิยติ, สพฺพํ รูปราคํ ปริยาทิยติ, สพฺพํ ภวราคํ ปริยาทิยติ, สพฺพํ อวิชฺชํ ปริยาทิยติ, สพฺพํ อสฺมิมานํ สมูหนติ.}

\prose{เสยฺยถาปิ ภิกฺขเว สรทสมเย กสฺสโก มหานงฺคเลน กสนฺโต สพฺพานิ มูลสนฺตานกานิ สมฺปทาเลนฺโต กสติ.\csromanspacer{} เอวเมว โข ภิกฺขเว อนิจฺจสญฺญา ภาวิตา พหุลีกตา สพฺพํ กามราคํ ปริยาทิยติ, สพฺพํ รูปราคํ ปริยาทิยติ, สพฺพํ ภวราคํ ปริยาทิยติ, สพฺพํ อวิชฺชํ ปริยาทิยติ, สพฺพํ อสฺมิมานํ สมูหนติ.}

\prose{เสยฺยถาปิ ภิกฺขเว ปพฺพชลายโก ปพฺพชํ ลายิตฺวา อคฺเค คเหตฺวา โอธุนาติ นิทฺธุนาติ นิจฺโฉเฏติ.\csromanspacer{} เอวเมว โข ภิกฺขเว อนิจฺจสญฺญา ภาวิตา พหุลีกตา สพฺพํ กามราคํ ปริยาทิยติ ฯเปฯ สพฺพํ อสฺมิมานํ สมูหนติ.}

\prose{\csromanfolio{127}เสยฺยถาปิ ภิกฺขเว อมฺพปิณฺฑิยา วณฺฏจฺฉินฺนาย ยานิ ตตฺถ อมฺพานิ วณฺฏปฏิพนฺธานิ, สพฺพานิ ตานิ ตทนฺวยานิ ภวนฺติ.\csromanspacer{} เอวเมว โข ภิกฺขเว อนิจฺจสญฺญา ภาวิตา ฯเปฯ สพฺพํ อสฺมิมานํ สมูหนติ.}

\prose{เสยฺยถาปิ ภิกฺขเว กูฏาคารสฺส ยา กาจิ โคปานสิโย สพฺพา ตา กูฏงฺคมา กูฏนินฺนา กูฏสโมสรณา, กูฏํ ตาสํ อคฺคมกฺขายติ.\csromanspacer{} เอวเมว โข ภิกฺขเว อนิจฺจสญฺญา ภาวิตา ฯเปฯ สพฺพํ อสฺมิมานํ สมูหนติ.}

\prose{เสยฺยถาปิ ภิกฺขเว เย เกจิ มูลคนฺธา, กาฬานุสาริคนฺโธ เตสํ อคฺคมกฺขายติ.\csromanspacer{} เอวเมว โข ภิกฺขเว อนิจฺจสญฺญา ฯเปฯ สพฺพํ อสฺมิมานํ สมูหนติ.}

\prose{เสยฺยถาปิ ภิกฺขเว เย เกจิ สารคนฺธา, โลหิตจนฺทนํ เตสํ อคฺคมกฺขายติ.\csromanspacer{} เอวเมว โข ภิกฺขเว อนิจฺจสญฺญา ฯเปฯ สพฺพํ อสฺมิมานํ สมูหนติ.}

\prose{เสยฺยถาปิ ภิกฺขเว เย เกจิ ปุปฺผคนฺธา, วสฺสิกํ เตสํ อคฺคมกฺขายติ.\csromanspacer{} เอวเมว โข ภิกฺขเว อนิจฺจสญฺญา ฯเปฯ สพฺพํ อสฺมิมานํ สมูหนติ.}

\prose{เสยฺยถาปิ ภิกฺขเว เย เกจิ กุฏฺฏราชาโน\footnote{กุฑฺฑราชาโน (สี)}, สพฺเพเต รญฺโญ จกฺกวตฺติสฺส อนุยนฺตา ภวนฺติ, ราชา เตสํ จกฺกวตฺติ อคฺคมกฺขายติ.\csromanspacer{} เอวเมว โข ภิกฺขเว อนิจฺจสญฺญา ฯเปฯ สพฺพํ อสฺมิมานํ สมูหนติ.}

\prose{เสยฺยถาปิ ภิกฺขเว ยา กาจิ ตารกรูปานํ ปภา, สพฺพา ตา จนฺทิมปฺปภาย กลํ นาคฺฆนฺติ โสฬสิํ, จนฺทปฺปภา ตาสํ อคฺคมกฺขายติ.\csromanspacer{} เอวเมว โข ภิกฺขเว อนิจฺจสญฺญา ฯเปฯ สพฺพํ อสฺมิมานํ สมูหนติ.}

\prose{เสยฺยถาปิ ภิกฺขเว สรทสมเย วิทฺเธ วิคตวลาหเก เทเว อาทิจฺโจ นตํ อพฺภุสฺสกฺกมาโน สพฺพํ อากาสคตํ ตมคตํ อภิวิหจฺจ ภาสเต จ ตปเต จ วิโรจเต จ.\csromanspacer{} เอวเมว โข ภิกฺขเว อนิจฺจสญฺญา ภาวิตา พหุลีกตา สพฺพํ กามราคํ ปริยาทิยติ, สพฺพํ รูปราคํ ปริยาทิยติ, สพฺพํ ภวราคํ ปริยาทิยติ, สพฺพํ อวิชฺชํ ปริยาทิยติ, สพฺพํ อสฺมิมานํ สมูหนติ.}

\prose{กถํ ภาวิตา จ ภิกฺขเว อนิจฺจสญฺญา กถํ พหุลีกตา สพฺพํ กามราคํ ปริยาทิยติ ฯเปฯ สพฺพํ อสฺมิมานํ สมูหนติ.\csromanspacer{} อิติ รูปํ, อิติ \csromanfolio{128}รูปสฺส สมุทโย, อิติ รูปสฺส อตฺถงฺคโม.\csromanspacer{} อิติ เวทนา.\csromanspacer{} อิติ สญฺญา.\csromanspacer{} อิติ สงฺขารา.\csromanspacer{} อิติ วิญฺญาณํ, อิติ วิญฺญาณสฺส สมุทโย, อิติ วิญฺญาณสฺส อตฺถงฺคโมติ.\csromanspacer{} เอวํ ภาวิตา โข ภิกฺขเว อนิจฺจสญฺญา เอวํ พหุลีกตา สพฺพํ กามราคํ ปริยาทิยติ, สพฺพํ รูปราคํ ปริยาทิยติ, สพฺพํ ภวราคํ ปริยาทิยติ, สพฺพํ อวิชฺชํ ปริยาทิยติ, สพฺพํ อสฺมิมานํ สมูหนตีติ.\csromanspacer{}\csromanspacer{}ทสมํ.\csromanspacer{} ปุปฺผวคฺโค ปญฺจโม.}

\titlehead{\textbf{ตสฺสุทฺทานํ}}
\csromangathameasure{นที ปุปฺผญฺจ เผณญฺจ, โคมยญฺจ นขาสิขํ. \\ สุทฺธิกํ เทฺว จ คทฺทุลา, วาสีชฏํ อนิจฺจตาติ. มชฺฌิมปณฺณาสโก สมตฺโต.}
\csromangathagroup{\csromangathastanza{นที ปุปฺผญฺจ เผณญฺจ, โคมยญฺจ นขาสิขํ. \\ สุทฺธิกํ เทฺว จ คทฺทุลา, วาสีชฏํ อนิจฺจตาติ.\csromanspacer{} มชฺฌิมปณฺณาสโก สมตฺโต.}}

\csromanheader{2}{\textbf{ตสฺส มชฺฌิมปณฺณาสกสฺส วคฺคุทฺทานํ}}
\csromangathameasure{อุปโย อรหนฺโต จ, ขชฺชนี เถรสวฺหยํ. \\ ปุปฺผวคฺเคน ปณฺณาส, ทุติโย เตน วุจฺจตีติ.}
\csromangathagroup{\csromangathastanza{อุปโย อรหนฺโต จ, ขชฺชนี เถรสวฺหยํ. \\ ปุปฺผวคฺเคน ปณฺณาส, ทุติโย เตน วุจฺจตีติ.}}

\csromanmatikaanchor{mtk.5}
\csromantocmarknopage{section}{(11) 1. อนฺตวคฺค}{mtk.5}
\bookmark[dest={mtk.5},level=1]{(11) 1. อนฺตวคฺค}
\csromanheader{1}{\textbf{(11) 1. อนฺตวคฺค}}
\csromanmatikaanchor{mtk.6}
\csromantocmarkref{subsection}{1. อนฺตสุตฺต}{mtk.6}
\bookmark[dest={mtk.6},level=2]{1. อนฺตสุตฺต}
\csromanheader{2}{1. อนฺตสุตฺต}
\proseitem{103}{สาวตฺถินิทานํ.\csromanspacer{} จตฺตาโรเม ภิกฺขเว อนฺตา.\csromanspacer{} กตเม จตฺตาโร, สกฺกายนฺโต สกฺกายสมุทยนฺโต สกฺกายนิโรธนฺโต สกฺกายนิโรธคามินิปฺปฏิปทนฺโต.\csromanspacer{} กตโม จ ภิกฺขเว สกฺกายนฺโต, ปญฺจุปาทานกฺขานฺธาติสฺส วจนียํ.\csromanspacer{} กตเม ปญฺจ, เสยฺยถิทํ, รูปุปาทานกฺขนฺโธ เวทนุปาทานกฺขนฺโธ สญฺญุปาทานกฺขนฺโธ สงฺขารุปาทานกฺขนฺโธ วิญฺญาณุปาทานกฺขนฺโธ.\csromanspacer{} อยํ วุจฺจติ ภิกฺขเว สกฺกายนฺโต.\csromanspacer{} กตโม จ ภิกฺขเว สกฺกายสมุทยนฺโต, ยายํ ตณฺหา โปโนภวิกา นนฺทิราคสหคตา ตตฺรตตฺราภินนฺทินี.\csromanspacer{} เสยฺยถิทํ, กามตณฺหา ภวตณฺหา วิภวตณฺหา. \csromanfolio{129}อยํ วุจฺจติ ภิกฺขเว สกฺกายสมุทยนฺโต.\csromanspacer{} กตโม จ ภิกฺขเว สกฺกายนิโรธนฺโต, โย ตสฺสาเยว ตณฺหาย อเสสวิราคนิโรโธ จาโค ปฏินิสฺสคฺโค มุตฺติ อนาลโย.\csromanspacer{} อยํ วุจฺจติ ภิกฺขเว สกฺกายนิโรธนฺโต.\csromanspacer{} กตโม จ ภิกฺขเว สกฺกายนิโรธคามินิปฺปฏิปทนฺโต, อยเมว อริโย อฏฺฐงฺคิโก มคฺโค.\csromanspacer{} เสยฺยถิทํ, สมฺมาทิฏฺฐิ ฯเปฯ สมฺมาสมาธิ.\csromanspacer{} อยํ วุจฺจติ ภิกฺขเว สกฺกายนิโรธคามินิปฺปฏิปทนฺโต.\csromanspacer{} อิเม โข ภิกฺขเว จตฺตาโร อนฺตาติ.\csromanspacer{}\csromanspacer{}ปฐมํ.}

\csromanmatikaanchor{mtk.7}
\csromantocmarkref{subsection}{2. ทุกฺขสุตฺต}{mtk.7}
\bookmark[dest={mtk.7},level=2]{2. ทุกฺขสุตฺต}
\csromanheader{2}{2. ทุกฺขสุตฺต}
\proseitem{104}{สาวตฺถินิทานํ.\csromanspacer{} ทุกฺขญฺจ โว ภิกฺขเว เทเสสฺสามิ ทุกฺขสมุทยญฺจ ทุกฺขนิโรธญฺจ ทุกฺขนิโรธคามินิญฺจ ปฏิปทํ, ตํ สุณาถ.\csromanspacer{} กตมญฺจ ภิกฺขเว ทุกฺขํ, ปญฺจุปาทานกฺขนฺธาติสฺส วจนียํ.\csromanspacer{} กตเม ปญฺจ, เสยฺยถิทํ, รูปุปาทานกฺขนฺโธ ฯเปฯ วิญฺญาณุปาทานกฺขนฺโธ.\csromanspacer{} อิทํ วุจฺจติ ภิกฺขเว ทุกฺขํ.\csromanspacer{} กตโม จ ภิกฺขเว ทุกฺขสมุทโย ยายํ ตณฺหา โปโนภวิกา ฯเปฯ วิภวตณฺหา.\csromanspacer{} อยํ วุจฺจติ ภิกฺขเว ทุกฺขสมุทโย.\csromanspacer{} กตโม จ ภิกฺขเว ทุกฺขนิโรโธ, โย ตสฺสาเยว ตณฺหาย อเสสวิราคนิโรโธ จาโค ปฏินิสฺสคฺโค มุตฺติ อนาลโย.\csromanspacer{} อยํ วุจฺจติ ภิกฺขเว ทุกฺขนิโรโธ.\csromanspacer{} กตมา จ ภิกฺขเว ทุกฺขนิโรธคามินี ปฏิปทา, อยเมว อริโย อฏฺฐงฺคิโก มคฺโค.\csromanspacer{} เสยฺยถิทํ, สมฺมาทิฏฺฐิ ฯเปฯ สมฺมาสมาธิ.\csromanspacer{} อยํ วุจฺจติ ภิกฺขเว ทุกฺขนิโรธคามินี ปฏิปทาติ.\csromanspacer{}\csromanspacer{}ทุติยํ.}

\csromanmatikaanchor{mtk.8}
\csromantocmarkref{subsection}{3. สกฺกายสุตฺต}{mtk.8}
\bookmark[dest={mtk.8},level=2]{3. สกฺกายสุตฺต}
\csromanheader{2}{3. สกฺกายสุตฺต}
\proseitem{105}{สาวตฺถินิทานํ.\csromanspacer{} สกฺกายญฺจ โว ภิกฺขเว เทเสสฺสามิ สกฺกายสมุทยญฺจ สกฺกายนิโรธญฺจ สกฺกายนิโรธคามินิญฺจ ปฏิปทํ, ตํ สุณาถ.\csromanspacer{} กตโม จ ภิกฺขเว สกฺกาโย, ปญฺจุปาทานกฺขนฺธาติสฺส วจนียํ.\csromanspacer{} กตเม ปญฺจ, เสยฺยถิทํ, รูปุปาทานกฺขนฺโธ เวทนุปาทานกฺขนฺโธ สญฺญุปาทานกฺขนฺโธ สงฺขารุปาทานกฺขนฺโธ วิญฺญาณุปาทานกฺขนฺโธ.\csromanspacer{} อยํ วุจฺจติ ภิกฺขเว สกฺกาโย.\csromanspacer{} กตโม จ ภิกฺขเว สกฺกายสมุทโย, ยายํ ตณฺหา โปโนภวิกา ฯเปฯ.\csromanspacer{} อยํ วุจฺจติ ภิกฺขเว สกฺกายสมุทโย.\csromanspacer{} กตโม จ ภิกฺขเว สกฺกายนิโรโธ, โย ตสฺสาเยว ตณฺหาย ฯเปฯ.\csromanspacer{} อยํ \csromanfolio{130}วุจฺจติ ภิกฺขเว สกฺกายนิโรโธ.\csromanspacer{} กถมา จ ภิกฺขเว สกฺกายนิโรธคามินี ปฏิปทา, อยเมว อริโย อฏฺฐงฺคิโก มคฺโค.\csromanspacer{} เสยฺยถิทํ สมฺมาทิฏฺฐิ ฯเปฯ สมฺมาสมาธิ.\csromanspacer{} อยํ วุจฺจติ ภิกฺขเว สกฺกายนิโรธคามินี ปฏิปทาติ.\csromanspacer{}\csromanspacer{}ตติยํ.}

\csromanmatikaanchor{mtk.9}
\csromantocmarkref{subsection}{4. ปริญฺเญยฺยสุตฺต}{mtk.9}
\bookmark[dest={mtk.9},level=2]{4. ปริญฺเญยฺยสุตฺต}
\csromanheader{2}{4. ปริญฺเญยฺยสุตฺต}
\proseitem{106}{สาวตฺถินิทานํ.\csromanspacer{} ปริญฺเญยฺเย จ ภิกฺขเว ธมฺเม เทเสสฺสามิ ปริญฺญญฺจ ปริญฺญาตาวิญฺจ ปุคฺคลํ, ตํ สุณาถ.\csromanspacer{} กตเม จ ภิกฺขเว ปริญฺเญยฺยา ธมฺมา, รูปํ ภิกฺขเว ปริญฺเญยฺโย ธมฺโม, เวทนา ฯเปฯ สญฺญา.\csromanspacer{} สงฺขารา.\csromanspacer{} วิญฺญาณํ ปริญฺเญยฺโย ธมฺโม.\csromanspacer{} อิเม วุจฺจนฺติ ภิกฺขเว ปริญฺเญยฺยา ธมฺมา.\csromanspacer{} กตมา จ ภิกฺขเว ปริญฺญา, ราคกฺขโย โทสกฺขโย โมหกฺขโย.\csromanspacer{} อยํ วุจฺจติ ภิกฺขเว ปริญฺญา.\csromanspacer{} กตโม จ ภิกฺขเว ปริญฺญาตาวี ปุคฺคโล, อรหาติสฺส วจนียํ, ยฺวายํ อายสฺมา เอวํนาโม เอวํโคตฺโต.\csromanspacer{} อยํ วุจฺจติ ภิกฺขเว ปริญฺญาตาวี ปุคฺคโลติ.\csromanspacer{}\csromanspacer{}จตุตฺถํ.}

\csromanmatikaanchor{mtk.10}
\csromantocmarkref{subsection}{5-6. สมณสุตฺต}{mtk.10}
\bookmark[dest={mtk.10},level=2]{5-6. สมณสุตฺต}
\csromanheader{2}{5. สมณสุตฺต}
\proseitem{107}{สาวตฺถินิทานํ.\csromanspacer{} ปญฺจิเม ภิกฺขเว อุปาทานกฺขนฺธา.\csromanspacer{} กตเม ปญฺจ, เสยฺยถิทํ, รูปุปาทานกฺขนฺโธ ฯเปฯ วิญฺญาณุปาทานกฺขนฺโธ.\csromanspacer{} เย หิ เกจิ ภิกฺขเว สมณา วา พฺรหฺมณา วา อิเมสํ ปญฺจนฺนํ อุปาทานกฺขนฺธานํ อสฺสาทญฺจ อาทีนวญฺจ นิสฺสรณญฺจ ยถาภูตํ นปฺปชานนฺติ ฯเปฯ ปชานนฺติ สยํ อภิญฺญา สจฺฉิกตฺวา อุปสมฺปชฺช วิหรนฺตีติ.\csromanspacer{}\csromanspacer{}ปญฺจมํ.}

\csromanheader{2}{6. ทุติยสมณสุตฺต}
\proseitem{108}{สาวตฺถินิทานํ.\csromanspacer{} ปญฺจิเม ภิกฺขเว อุปาทานกฺขนฺธา.\csromanspacer{} กตเม ปญฺจ, เสยฺยถิทํ, รูปุปาทานกฺขนฺโธ เวทนุปาทานกฺขนฺโธ สญฺญุปาทานกฺขนฺโธ สงฺขารุปาทานกฺขนฺโธ วิญฺญาณุปาทานกฺขนฺโธ.\csromanspacer{} เย หิ เกจิ ภิกฺขเว สมณา วา พฺรหฺมณา วา อิเมสํ ปญฺจนฺนํ อุปาทานกฺขนฺธานํ สมุทยญฺจ อตฺถงฺคมญฺจ อสฺสาทญฺจ อาทีนวญฺจ นิสฺสรณญฺจ ยถาภูตํ นปฺปชานนฺติ ฯเปฯ ปชานนฺติ สยํ อภิญฺญา สจฺฉิกตฺวา อุปสมฺปชฺช วิหรนฺตีติ.\csromanspacer{}\csromanspacer{}ฉฏฺฐํ.}

\csromanmatikaanchor{mtk.11}
\csromantocmarkref{subsection}{7. โสตาปนฺนสุตฺต}{mtk.11}
\bookmark[dest={mtk.11},level=2]{7. โสตาปนฺนสุตฺต}
\csromanheader{2}{7. โสตาปนฺนสุตฺต}
\proseitem{109}{\csromanfolio{131}สาวตฺถินิทานํ.\csromanspacer{} ปญฺจิเม ภิกฺขเว อุปาทานกฺขนฺธา.\csromanspacer{} กตเม ปญฺจ, เสยฺยถิทํ, รูปุปาทานกฺขนฺโธ ฯเปฯ วิญฺญาณุปาทานกฺขนฺโธ.\csromanspacer{} ยโต โข ภิกฺขเว อริยสาวโก อิเมสํ ปญฺจนฺนํ อุปาทานกฺขนฺธานํ สมุทยญฺจ อตฺถงฺคมญฺจ อสฺสาทญฺจ อาทีนวญฺจ นิสฺสรณญฺจ ยถาภูตํ ปชานาติ.\csromanspacer{} อยํ วุจฺจติ ภิกฺขเว อริยสาวโก โสตาปนฺโน อวินิปาตธมฺโม นิยโต สมฺโพธิปรายโนติ.\csromanspacer{}\csromanspacer{}สตฺตมํ.}

\csromanmatikaanchor{mtk.12}
\csromantocmarkref{subsection}{8. อรหนฺตสุตฺต}{mtk.12}
\bookmark[dest={mtk.12},level=2]{8. อรหนฺตสุตฺต}
\csromanheader{2}{8. อรหนฺตสุตฺต}
\proseitem{110}{สาวตฺถินิทานํ.\csromanspacer{} ปญฺจิเม ภิกฺขเว อุปาทานกฺขนฺธา.\csromanspacer{} กตเม ปญฺจ, เสยฺยถิทํ, รูปุปาทานกฺขนฺโธ ฯเปฯ วิญฺญาณุปาทานกฺขนฺโธ.\csromanspacer{} ยโต โข ภิกฺขเว ภิกฺขุ อิเมสํ ปญฺจนฺนํ อุปาทานกฺขนฺธานํ สมุทยญฺจ อตฺถงฺคมญฺจ อสฺสาทญฺจ อาทีนวญฺจ นิสฺสรณญฺจ ยถาภูตํ วิทิตฺวา อนุปาทาวิมุตฺโต โหติ.\csromanspacer{} อยํ วุจฺจติ ภิกฺขเว ภิกฺขุ อรหํ ขีณาสโว วุสิตวา กตกรณีโย โอหิตภาโร อนุปฺปตฺตสทตฺโถ ปริกฺขีณภวสํโยชโน สมฺมทญฺญาวิมุตฺโตติ.\csromanspacer{}\csromanspacer{}อฏฺฐมํ.}

\csromanmatikaanchor{mtk.13}
\csromantocmarkref{subsection}{9-10. ฉนฺทปฺปหานสุตฺต}{mtk.13}
\bookmark[dest={mtk.13},level=2]{9-10. ฉนฺทปฺปหานสุตฺต}
\csromanheader{2}{9. ฉนฺทปฺปหานสุตฺต}
\proseitem{111}{สาวตฺถินิทานํ.\csromanspacer{} รูเป ภิกฺขเว โย ฉนฺโท โย ราโค ยา นนฺที ยา ตณฺหา, ตํ ปชหถ.\csromanspacer{} เอวํ ตํ รูปํ ปหีนํ ภวิสฺสติ อุจฺฉินฺนมูลํ ตาลาวตฺถุกตํ อนภาวํกตํ อายติํ อนุปฺปาทธมฺมํ.\csromanspacer{} เวทนาย ฯเปฯ.\csromanspacer{} สญฺญาย.\csromanspacer{} สงฺขาเรสุ.\csromanspacer{} วิญฺญาเณ โย ฉนฺโท โย ราโค ยา นนฺที ยา ตณฺหา, ตํ ปชหถ.\csromanspacer{} เอวํ ตํ วิญฺญาณํ ปหีนํ ภวิสฺสติ อุจฺฉินฺนมูลํ ตาลาวตฺถุกตํ อนภาวํกตํ อายติํ อนุปฺปาทธมฺมนฺติ.\csromanspacer{}\csromanspacer{}นวมํ.}

\csromanheader{2}{10. ทุติยฉนฺทปฺปหานสุตฺต}
\proseitem{112}{สาวตฺถินิทานํ.\csromanspacer{} รูเป ภิกฺขเว โย ฉนฺโท โย ราโค ยา นนฺที ยา ตณฺหา เย อุปยุปาทานา เจตโส อธิฏฺฐานาภินิเวสานุสยา, เต ปชหถ.\csromanspacer{} เอวํ ตํ รูปํ ปหีนํ ภวิสฺสติ อุจฺฉินฺนมูลํ ฯเปฯ.\csromanspacer{} เวทนาย.\csromanspacer{} สญฺญาย.\csromanspacer{} สงฺขาเรสุ.\csromanspacer{} โย ฉนฺโท ฯเปฯ.\csromanspacer{} เอวํ เต สงฺขารา ปหีนา ภวิสฺสนฺติ อุจฺฉินฺนมูลา ตาลาวตฺถุกตา อนภาวํกตา \csromanfolio{132}อายติํ อนุปฺปาทธมฺมา.\csromanspacer{} วิญฺญาเณ โย ฉนฺโท โย ราโค ยา นนฺที ยา ตณฺหา เย อุปยุปาทานา เจตโส อธิฏฺฐานาภินิเวสานุสยา, เต ปชหถ.\csromanspacer{} เอวํ ตํ วิญฺญาณํ ปหีนํ ภวิสฺสติ อุจฺฉินฺนมูลํ ตาลาวตฺถุกตํ อนภาวํกตํ อายติํ อนุปฺปาทธมฺมนฺติ.\csromanspacer{}\csromanspacer{}ทสมํ.\csromanspacer{} อนฺตวคฺโค ปฐโม.}

\titlehead{\textbf{ตสฺสุทฺทานํ}}
\csromangathameasure{อนฺโต ทุกฺขญฺจ สกฺกาโย, ปริญฺเญยฺยา สมณา ทุเว. \\ โสตาปนฺโน อรหา จ, ทุเว จ ฉนฺทปฺปหานาติ.}
\csromangathagroup{\csromangathastanza{อนฺโต ทุกฺขญฺจ สกฺกาโย, ปริญฺเญยฺยา สมณา ทุเว. \\ โสตาปนฺโน อรหา จ, ทุเว จ ฉนฺทปฺปหานาติ.}}

\csromanmatikaanchor{mtk.14}
\csromantocmarknopage{section}{(12) 2. ธมฺมกถิกวคฺค}{mtk.14}
\bookmark[dest={mtk.14},level=1]{(12) 2. ธมฺมกถิกวคฺค}
\csromanheader{1}{\textbf{(12) 2. ธมฺมกถิกวคฺค}}
\csromanmatikaanchor{mtk.15}
\csromantocmarkref{subsection}{1. อวิชฺชาสุตฺต}{mtk.15}
\bookmark[dest={mtk.15},level=2]{1. อวิชฺชาสุตฺต}
\csromanheader{2}{1. อวิชฺชาสุตฺต}
\proseitem{113}{สาวตฺถินิทานํ.\csromanspacer{} อถ โข อญฺญตโร ภิกฺขุ เยน ภควา เตนุปสงฺกมิ ฯเปฯ เอกมนฺตํ นิสินฺโน โข โส ภิกฺขุ ภควนฺตํ เอตทโวจ “อวิชฺชา อวิชฺชาติ ภนฺเต วุจฺจติ, กตมา นุ โข ภนฺเต อวิชฺชา, กิตฺตาวตา จ อวิชฺชาคโต โหตี”ติ.\csromanspacer{} อิธ ภิกฺขุ อสฺสุตวา ปุถุชฺชโน รูปํ นปฺปชานาติ, รูปสมุทยํ นปฺปชานาติ, รูปนิโรธํ นปฺปชานาติ, รูปนิโรธคามินิํ ปฏิปทํ นปฺปชานาติ.\csromanspacer{} เวทนํ นปฺปชานาติ.\csromanspacer{} สญฺญํ.\csromanspacer{} สงฺขาเร นปฺปชานาติ ฯเปฯ วิญฺญาณนิโรธคามินิํ ปฏิปทํ นปฺปชานาติ.\csromanspacer{} อยํ วุจฺจติ ภิกฺขุ อวิชฺชา, เอตฺตาวตา จ อวิชฺชาคโต โหตีติ.\csromanspacer{}\csromanspacer{}ปฐมํ.}

\csromanmatikaanchor{mtk.16}
\csromantocmarkref{subsection}{2. วิชฺชาสุตฺต}{mtk.16}
\bookmark[dest={mtk.16},level=2]{2. วิชฺชาสุตฺต}
\csromanheader{2}{2. วิชฺชาสุตฺต}
\proseitem{114}{สาวตฺถินิทานํ.\csromanspacer{} เอกมนฺตํ นิสินฺโน โข โส ภิกฺขุ ภควนฺตํ เอตทโวจ “วิชฺชา วิชฺชาติ ภนฺเต วุจฺจติ, กตมา นุ โข ภนฺเต วิชฺชา, กิตฺตาวตา จ วิชฺชาคโต โหตี”ติ.\csromanspacer{} อิธ ภิกฺขุ สุตวา อริยสาวโก รูปํ ปชานาติ, รูปสมุทยํ.\csromanspacer{} รูปนิโรธํ.\csromanspacer{} รูปนิโรธคามินิํ ปฏิปทํ ปชานาติ.\csromanspacer{} เวทนํ.\csromanspacer{} สญฺญํ.\csromanspacer{} สงฺขาเร ปชานาติ ฯเปฯ วิญฺญาณนิโรธคามินิํ ปฏิปทํ ปชานาติ.\csromanspacer{} อยํ วุจฺจติ ภิกฺขุ วิชฺชา, เอตฺตาวตา จ วิชฺชาคโต โหตีติ.\csromanspacer{}\csromanspacer{}ทุติยํ.}

\csromanmatikaanchor{mtk.17}
\csromantocmarkref{subsection}{3-4. ธมฺมกถิกสุตฺต}{mtk.17}
\bookmark[dest={mtk.17},level=2]{3-4. ธมฺมกถิกสุตฺต}
\csromanheader{2}{3. ธมฺมกถิกสุตฺต}
\proseitem{115}{\csromanfolio{133}สาวตฺถินิทานํ.\csromanspacer{} เอกมนฺตํ นิสินฺโน โข โส ภิกฺขุ ภควนฺตํ เอตทโวจ “ธมฺมกถิโก ธมฺมกถิโกติ ภนฺเต วุจฺจติ, กิตฺตาวตา นุ โข ภนฺเต ธมฺมกถิโก โหตี”ติ.\csromanspacer{} รูปสฺส เจ ภิกฺขุ นิพฺพิทาย วิราคาย นิโรธาย ธมฺมํ เทเสติ, “ธมฺมกถิโก ภิกฺขู”ติ อลํ วจนาย.\csromanspacer{} รูปสฺส เจ ภิกฺขุ นิพฺพิทาย วิราคาย นิโรธาย ปฏิปนฺโน โหติ, “ธมฺมานุธมฺมปฺปฏิปนฺโน ภิกฺขู”ติ อลํ วจนาย.\csromanspacer{} รูปสฺส เจ ภิกฺขุ นิพฺพิทา วิราคา นิโรธา อนุปาทาวิมุตฺโต โหติ, “ทิฏฺฐธมฺมนิพฺพานปฺปตฺโต ภิกฺขู”ติ อลํ วจนาย.\csromanspacer{} เวทนาย เจ ภิกฺขุ ฯเปฯ.\csromanspacer{} สญฺญาย เจ ภิกฺขุ.\csromanspacer{} สงฺขารานํ เจ ภิกฺขุ.\csromanspacer{} วิญฺญาณสฺส เจ ภิกฺขุ นิพฺพิทาย วิราคาย นิโรธาย ธมฺมํ เทเสติ, “ธมฺมกถิโก ภิกฺขู”ติ อลํ วจนาย.\csromanspacer{} วิญฺญาณสฺส เจ ภิกฺขุ นิพฺพิทาย วิราคาย นิโรธาย ปฏิปนฺโน โหติ, “ธมฺมานุธมฺมปฺปฏิปนฺโน ภิกฺขู”ติ อลํ วจนาย.\csromanspacer{} วิญฺญาณสฺส เจ ภิกฺขุ นิพฺพิทา วิราคา นิโรธา อนุปาทาวิมุตฺโต โหติ, “ทิฏฺฐธมฺมนิพฺพานปฺปตฺโต ภิกฺขู”ติ อลํ วจนายาติ.\csromanspacer{}\csromanspacer{}ตติยํ.}

\csromanheader{2}{4. ทุติยธมฺมกถิกสุตฺต}
\csromanmatikaanchor{mtk.18}
\csromantocmarkref{subsection}{5. พนฺทนสุตฺต}{mtk.18}
\bookmark[dest={mtk.18},level=2]{5. พนฺทนสุตฺต}
\proseitem{116}{สาวตฺถินิทานํ.\csromanspacer{} เอกมนฺตํ นิสินฺโน โข โส ภิกฺขุ ภควนฺตํ เอตทโวจ “ธมฺมกถิโก ธมฺมกถิโกติ ภนฺเต วุจฺจติ, กิตฺตาวตา นุ โข ภนฺเต ธมฺมกถิโก โหติ, กิตฺตาวตา ธมฺมานุธมฺมปฺปฏิปนฺโน โหติ, กิตฺตาวตา ทิฏฺฐธมฺมนิพฺพานปฺปตฺโต โหตี”ติ.\csromanspacer{} รูปสฺส เจ ภิกฺขุ นิพฺพิทาย วิราคาย นิโรธาย ธมฺมํ เทเสติ, “ธมฺมกถิโก ภิกฺขู”ติ อลํ วจนาย.\csromanspacer{} รูปสฺส เจ ภิกฺขุ นิพฺพิทาย วิราคาย นิโรธาย ปฏิปนฺโน โหติ, “ธมฺมานุธมฺมปฺปฏิปนฺโน ภิกฺขู”ติ อลํ วจนาย.\csromanspacer{} รูปสฺส เจ ภิกฺขุ นิพฺพิทา วิราคา นิโรธา อนุปาทาวิมุตฺโต โหติ, “ทิฏฺฐธมฺมนิพฺพานปฺปตฺโต ภิกฺขู”ติ อลํ วจนาย.\csromanspacer{} เวทนาย เจ ภิกฺขุ ฯเปฯ.\csromanspacer{} สญฺญาย เจ ภิกฺขุ.\csromanspacer{} สงฺขารานํ เจ ภิกฺขุ.\csromanspacer{} วิญฺญาณสฺส เจ ภิกฺขุ นิพฺพิทาย วิราคาย นิโรธาย ธมฺมํ เทเสติ, “ธมฺมกถิโก ภิกฺขู”ติ อลํ วจนาย.\csromanspacer{} วิญฺญาณสฺส เจ ภิกฺขุ นิพฺพิทาย วิราคาย นิโรธาย ปฏิปนฺโน โหติ, “ธมฺมานุธมฺมปฺปฏิปนฺโน ภิกฺขู”ติ อลํ วจนาย.\csromanspacer{} วิญฺญาณสฺส เจ ภิกฺขุ นิพฺพิทา \csromanfolio{134}วิราคา นิโรธา อนุปาทาวิมุตฺโต โหติ, “ทิฏฺฐธมฺมนิพฺพานปฺปตฺโต ภิกฺขู”ติ อลํ วจนายาติ.\csromanspacer{}\csromanspacer{}จตุตฺถํ.}

\csromanheader{2}{5. พนฺธนสุตฺต}
\proseitem{117}{สาวตฺถินิทานํ.\csromanspacer{} อิธ ภิกฺขเว อสฺสุตวา ปุถุชฺชโน อริยานํ อทสฺสาวี ฯเปฯ สปฺปุริสธมฺเม อวินีโต รูปํ อตฺตโต สมนุปสฺสติ, รูปวนฺตํ วา อตฺตานํ, อตฺตนิ วา รูปํ, รูปสฺมิํ วา อตฺตานํ.\csromanspacer{} อยํ วุจฺจติ ภิกฺขเว อสฺสุตวา ปุถุชฺชโน รูปพนฺธนพทฺโธ สนฺตรพาหิรพนฺธนพทฺโธ อตีรทสฺสี อปารทสฺสี พทฺโธ ชียติ\footnote{พทฺโธ ชายติ (สี, อิ) พทฺโธ ชายติ พทฺโธ ชียติ (สี-ฐ, สฺยา-ฐ)} พทฺโธ มียติ พทฺโธ อสฺมา โลกา ปรํ โลกํ คจฺฉติ.\csromanspacer{} เวทนํ อตฺตโต สมนุปสฺสติ ฯเปฯ.\csromanspacer{} เวทนาย วา อตฺตานํ.\csromanspacer{} อยํ วุจฺจติ ภิกฺขเว อสฺสุตวา ปุถุชฺชโน เวทนาพนฺธนพทฺโธ สนฺตรพาหิรพนฺธนพทฺโธ อตีรทสฺสี อปารทสฺสี พทฺโธ ชียติ พทฺโธ มียติ พทฺโธ อสฺมา โลกา ปรํ โลกํ คจฺฉติ.\csromanspacer{} สญฺญํ.\csromanspacer{} สงฺขาเร.\csromanspacer{} วิญฺญาณํ อตฺตโต สมนุปสฺสติ ฯเปฯ.\csromanspacer{} อยํ วุจฺจติ ภิกฺขเว อสฺสุตวา ปุถุชฺชโน วิญฺญาณพนฺธนพทฺโธ สนฺตรพาหิรพนฺธนพทฺโธ อตีรทสฺสี อปารทสฺสี พทฺโธ ชียติ พทฺโธ มียติ พทฺโธ อสฺมา โลกา ปรํ โลกํ คจฺฉติ.}

\prose{สุตวา จ โข ภิกฺขเว อริยสาวโก อริยานํ ทสฺสาวี ฯเปฯ สปฺปุริสธมฺเม สุวินีโต น รูปํ อตฺตโต สมนุปสฺสติ, น รูปวนฺตํ วา อตฺตานํ, น อตฺตนิ วา รูปํ, น รูปสฺมิํ วา อตฺตานํ.\csromanspacer{} อยํ วุจฺจติ ภิกฺขเว สุตวา อริยสาวโก น รูปพนฺธนพทฺโธ น สนฺตรพาหิรพนฺธนพทฺโธ ตีรทสฺสี ปารทสฺสี, ปริมุตฺโต โส ทุกฺขสฺมาติ วทามิ.\csromanspacer{} น เวทนํ อตฺตโต ฯเปฯ.\csromanspacer{} น สญฺญํ อตฺตโต ฯเปฯ.\csromanspacer{} น สงฺขาเร อตฺตโต ฯเปฯ.\csromanspacer{} น วิญฺญาณํ อตฺตโต สมนุปสฺสติ ฯเปฯ.\csromanspacer{} อยํ วุจฺจติ ภิกฺขเว สุตวา อริยสาวโก น วิญฺญาณพนฺธนพทฺโธ น สนฺตรพาหิรพนฺธนพทฺโธ ตีรทสฺสี ปารทสฺสี, ปริมุตฺโต โส ทุกฺขสฺมาติ วทามีติ.\csromanspacer{}\csromanspacer{}ปญฺจมํ.}

\csromanmatikaanchor{mtk.19}
\csromantocmarkref{subsection}{6-7. ปริปุจฺฉิตสุตฺต}{mtk.19}
\bookmark[dest={mtk.19},level=2]{6-7. ปริปุจฺฉิตสุตฺต}
\csromanheader{2}{6. ปริปุจฺฉิตสุตฺต}
\proseitem{118}{สาวตฺถินิทานํ.\csromanspacer{} ตํ กิํ มญฺญถ ภิกฺขเว, รูปํ “เอตํ มม, เอโสหมสฺมิ, เอโส เม อตฺตา”ติ สมนุปสฺสถาติ.\csromanspacer{} โน เหตํ ภนฺเต.\csromanspacer{} สาธุ ภิกฺขเว, รูปํ ภิกฺขเว “เนตํ มม, เนโสหมสฺมิ, น เมโส อตฺตา”ติ \csromanfolio{135}เอวเมตํ ยถาภูตํ สมฺมปฺปญฺญาย ทฏฺฐพฺพํ.\csromanspacer{} เวทนํ.\csromanspacer{} สญฺญํ.\csromanspacer{} สงฺขาเร.\csromanspacer{} วิญฺญาณํ “เอตํ มม, เอโสหมสฺมิ, เอโส เม อตฺตา”ติ สมนุปสฺสถาติ.\csromanspacer{} โน เหตํ ภนฺเต.\csromanspacer{} สาธุ ภิกฺขเว, วิญฺญาณํ ภิกฺขเว “เนตํ มม, เนโสหมสฺมิ, น เมโส อตฺตา”ติ เอวเมตํ ยถาภูตํ สมฺมปฺปญฺญาย ทฏฺฐพฺพํ ฯเปฯ.\csromanspacer{} เอวํ ปสฺสํ ฯเปฯ นาปรํ อิตฺถตฺตายาติ ปชานาตีติ.\csromanspacer{}\csromanspacer{}ฉฏฺฐํ.}

\csromanheader{2}{7. ทุติยปริปุจฺฉิตสุตฺต}
\proseitem{119}{สาวตฺถินิทานํ.\csromanspacer{} ตํ กิํ มญฺญถ ภิกฺขเว, รูปํ “เนตํ มม, เนโสหมสฺมิ, น เมโส อตฺตา”ติ สมนุปสฺสถาติ.\csromanspacer{} เอวํ ภนฺเต.\csromanspacer{} สาธุ ภิกฺขเว, รูปํ ภิกฺขเว “เนตํ มม, เนโสหมสฺมิ, น เมโส อตฺตา”ติ เอวเมตํ ยถาภูตํ สมฺมปฺปญฺญาย ทฏฺฐพฺพํ.\csromanspacer{} เวทนํ.\csromanspacer{} สญฺญํ.\csromanspacer{} สงฺขาเร.\csromanspacer{} วิญฺญาณํ “เนตํ มม, เนโสหมสฺมิ, น เมโส อตฺตา”ติ สมนุปสฺสถาติ.\csromanspacer{} เอวํ ภนฺเต.\csromanspacer{} สาธุ ภิกฺขเว, วิญฺญาณํ ภิกฺขเว “เนตํ มม, เนโสหมสฺมิ, น เมโส อตฺตา”ติ เอวเมตํ ยถาภูตํ สมฺมปฺปญฺญาย ทฏฺฐพฺพํ.\csromanspacer{} เอวํ ฯเปฯ นาปรํ อิตฺถตฺตายาติ ปชานาตีติ.\csromanspacer{}\csromanspacer{}สตฺตมํ.}

\csromanmatikaanchor{mtk.20}
\csromantocmarkref{subsection}{8. สํโยชนิยสุตฺต}{mtk.20}
\bookmark[dest={mtk.20},level=2]{8. สํโยชนิยสุตฺต}
\csromanheader{2}{8. สํโยชนิยสุตฺต}
\proseitem{120}{สาวตฺถินิทานํ.\csromanspacer{} สํโยชนิเย จ ภิกฺขเว ธมฺเม เทเสสฺสามิ สํโยชนญฺจ, ตํ สุณาถ.\csromanspacer{} กตเม จ ภิกฺขเว สํโยชนิยา ธมฺมา, กตมํ สํโยชนํ.\csromanspacer{} รูปํ ภิกฺขเว สํโยชนิโย ธมฺโม.\csromanspacer{} โย ตตฺถ ฉนฺทราโค, ตํ ตตฺถ สํโยชนํ.\csromanspacer{} เวทนา ฯเปฯ.\csromanspacer{} สญฺญา.\csromanspacer{} สงฺขารา.\csromanspacer{} วิญฺญาณํ สํโยชนิโย ธมฺโม.\csromanspacer{} โย ตตฺถ ฉนฺทราโค, ตํ ตตฺถ สํโยชนํ.\csromanspacer{} อิเม วุจฺจนฺติ ภิกฺขเว สํโยชนิยา ธมฺมา, อิทํ สํโยชนนฺติ.\csromanspacer{}\csromanspacer{}อฏฺฐมํ.}

\csromanmatikaanchor{mtk.21}
\csromantocmarkref{subsection}{9. อุปาทานิยสุตฺต}{mtk.21}
\bookmark[dest={mtk.21},level=2]{9. อุปาทานิยสุตฺต}
\csromanheader{2}{9. อุปาทานิยสุตฺต}
\proseitem{121}{สาวตฺถินิทานํ.\csromanspacer{} อุปาทานิเย จ ภิกฺขเว ธมฺเม เทเสสฺสามิ อุปาทานญฺจ, ตํ สุณาถ.\csromanspacer{} กตเม จ ภิกฺขเว อุปาทานิยา ธมฺมา, กตมํ อุปาทานํ.\csromanspacer{} รูปํ ภิกฺขเว อุปาทานิโย ธมฺโม.\csromanspacer{} โย ตตฺถ ฉนฺทราโค, ตํ ตตฺถ อุปาทานํ.\csromanspacer{} เวทนา ฯเปฯ.\csromanspacer{} สญฺญา.\csromanspacer{} สงฺขารา.\csromanspacer{} วิญฺญาณํ อุปาทานิโย ธมฺโม.\csromanspacer{} โย ตตฺถ ฉนฺทราโค, ตํ ตตฺถ อุปาทานํ.\csromanspacer{} อิเม วุจฺจนฺติ ภิกฺขเว อุปาทานิยา ธมฺมา, อิทํ อุปาทานนฺติ.\csromanspacer{}\csromanspacer{}นวมํ.}

\csromanmatikaanchor{mtk.22}
\csromantocmarkref{subsection}{10. สีลวนฺตสุตฺต}{mtk.22}
\bookmark[dest={mtk.22},level=2]{10. สีลวนฺตสุตฺต}
\csromanheader{2}{10. สีลวนฺตสุตฺต}
\proseitem{122}{\csromanfolio{136}เอกํ สมยํ อายสฺมา จ สาริปุตฺโต อายสฺมา จ มหาโกฏฺฐิโก\footnote{มหาโกฏฺฐิโต (สี, สฺยา, กํ, อิ)} พาราณสิยํ วิหรนฺติ อิสิปตเน มิคทาเย.\csromanspacer{} อถ โข อายสฺมา มหาโกฏฺฐิโก สายนฺหสมยํ ปฏิสลฺลานา วุฏฺฐิโต เยนายสฺมา สาริปุตฺโต เตนุปสงฺกมิ ฯเปฯ เอตทโวจ “สีลวตาวุโส สาริปุตฺต ภิกฺขุนา กตเม ธมฺมา โยนิโส มนสิ กาตพฺพา”ติ.\csromanspacer{} สีลวตาวุโส โกฏฺฐิก ภิกฺขุนา ปญฺจุปาทานกฺขนฺธา อนิจฺจโต ทุกฺขโต โรคโต คณฺฑโต สลฺลโต อฆโต อาพาธโต ปรโต ปโลกโต สุญฺญโต อนตฺตโต โยนิโส มนสิ กาตพฺพา.\csromanspacer{} กตเม ปญฺจ, เสยฺยถิทํ, รูปุปาทานกฺขนฺโธ เวทนุปาทานกฺขนฺโธ สญฺญุปาทานกฺขนฺโธ สงฺขารุปาทานกฺขนฺโธ วิญฺญาณุปาทานกฺขนฺโธ.\csromanspacer{} สีลวตาวุโส โกฏฺฐิก ภิกฺขุนา อิเม ปญฺจุปาทานกฺขนฺธา อนิจฺจโต ทุกฺขโต โรคโต คณฺฑโต สลฺลโต อฆโต อาพาธโต ปรโต ปโลกโต สุญฺญโต อนตฺตโต โยนิโส มนสิ กาตพฺพา.\csromanspacer{} ฐานํ โข ปเนตํ อาวุโส วิชฺชติ, ยํ สีลวา ภิกฺขุ อิเม ปญฺจุปาทานกฺขนฺเธ อนิจฺจโต ฯเปฯ อนตฺตโต โยนิโส มนสิ กโรนฺโต โสตาปตฺติผลํ สจฺฉิกเรยฺยาติ.}

\prose{โสตาปนฺเนน ปนาวุโส สาริปุตฺต ภิกฺขุนา กตเม ธมฺมา โยนิโส มนสิ กาตพฺพาติ.\csromanspacer{} โสตาปนฺเนนปิ โข อาวุโส โกฏฺฐิก ภิกฺขุนา อิเม ปญฺจุปาทานกฺขนฺธา อนิจฺจโต ฯเปฯ อนตฺตโต โยนิโส มนสิ กาตพฺพา.\csromanspacer{} ฐานํ โข ปเนตํ อาวุโส วิชฺชติ, ยํ โสตาปนฺโน ภิกฺขุ อิเม ปญฺจุปาทานกฺขนฺเธ อนิจฺจโต ฯเปฯ อนตฺตโต โยนิโส มนสิ กาโรนฺโต สกทาคามิผลํ สจฺฉิกเรยฺยาติ.}

\prose{สกทาคามินา ปนาวุโส สาริปุตฺต ภิกฺขุนา กตเม ธมฺมา โยนิโส มนสิ กาตพฺพาติ.\csromanspacer{} สกทาคามินาปิ โข อาวุโส โกฏฺฐิก ภิกฺขุนา อิเม ปญฺจุปาทานกฺขนฺธา อนิจฺจโต ฯเปฯ อนตฺตโต โยนิโส มนสิ กาตพฺพา.\csromanspacer{} ฐานํ โข ปเนตํ อาวุโส วิชฺชติ, ยํ สกทาคามี ภิกฺขุ อิเม ปญฺจุปาทานกฺขนฺเธ อนิจฺจโต ฯเปฯ อนตฺตโต โยนิโส มนสิ กโรนฺโต อนาคามิผลํ สจฺฉิกเรยฺยาติ.}

\prose{\csromanfolio{137}อนาคามินา ปนาวุโส สาริปุตฺต ภิกฺขุนา กตเม ธมฺมา โยนิโส มนสิ กาตพฺพาติ.\csromanspacer{} อนาคามินาปิ โข อาวุโส โกฏฺฐิก ภิกฺขุนา อิเม ปญฺจุปาทานกฺขนฺธา อนิจฺจโต ฯเปฯ อนตฺตโต โยนิโส มนสิ กาตพฺพา.\csromanspacer{} ฐานํ โข ปเนตํ อาวุโส วิชฺชติ, ยํ อนาคามี ภิกฺขุ อิเม ปญฺจุปาทานกฺขนฺเธ อนิจฺจโต ฯเปฯ อนตฺตโต โยนิโส มนสิ กโรนฺโต อรหตฺตํ สจฺฉิกเรยฺยาติ.}

\prose{อรหตา ปนาวุโส สาริปุตฺต กตเม ธมฺมา โยนิโส มนสิ กาตพฺพาติ.\csromanspacer{} อรหตาปิ โข อาวุโส โกฏฺฐิก อิเม ปญฺจุปาทานกฺขนฺเธ อนิจฺจโต ทุกฺขโต โรคโต คณฺฑโต สลฺลโต อฆโต อาพาธโต ปรโต ปโลกโต สุญฺญโต อนตฺตโต โยนิโส มนสิ กาตพฺพา.\csromanspacer{} นตฺถิ ขฺวาวุโส อรหโต อุตฺตริ กรณียํ, กตสฺส วา ปติจโย, อปิ จ อิเม ธมฺมา ภาวิตา พหุลีกตา ทิฏฺฐธมฺมสุขวิหาราย เจว สํวตฺตนฺติ สติสมฺปชญฺญาย จาติ.\csromanspacer{}\csromanspacer{}ทสมํ.}

\csromanmatikaanchor{mtk.23}
\csromantocmarkref{subsection}{11. สุตวนฺตสุตฺต}{mtk.23}
\bookmark[dest={mtk.23},level=2]{11. สุตวนฺตสุตฺต}
\csromanheader{2}{11. สุตวนฺตสุตฺต}
\proseitem{123}{เอกํ สมยํ อายสฺมา จ สาริปุตฺโต อายสฺมา จ มหาโกฏฺฐิโก พาราณสิยํ วิหรนฺติ อิสิปตเน มิคทาเย.\csromanspacer{} อถ โข อายสฺมา มหาโกฏฺฐิโก สายนฺหสมยํ ปฏิสลฺลานา วุฏฺฐิโต เยนายสฺมา สาริปุตฺโต เตนุปสงฺกมิ, อุปสงฺกมิตฺวา ฯเปฯ เอตทโวจ\csromandash{}}

\prose{สุตวตาวุโส สาริปุตฺต ภิกฺขุนา กตเม ธมฺมา โยนิโส มนสิ กาตพฺพาติ.\csromanspacer{} สุตวตาวุโส โกฏฺฐิก ภิกฺขุนา ปญฺจุปาทานกฺขนฺธา อนิจฺจโต ฯเปฯ อนตฺตโต โยนิโส มนสิ กาตพฺพา.\csromanspacer{} กตเม ปญฺจ, เสยฺยถิทํ, รูปุปาทานกฺขนฺโธ ฯเปฯ วิญฺญาณุปาทานกฺขนฺโธ.\csromanspacer{} สุตวตาวุโส โกฏฺฐิก ภิกฺขุนา อิเม ปญฺจุปาทานกฺขนฺธา อนิจฺจโต ฯเปฯ อนตฺตโต โยนิโส มนสิ กาตพฺพา.\csromanspacer{} ฐานํ โข ปเนตํ อาวุโส วิชฺชติ, ยํ สุตวา ภิกฺขุ อิเม ปญฺจุปาทานกฺขนฺเธ อนิจฺจโต ฯเปฯ อนตฺตโต โยนิโส มนสิ กโรนฺโต โสตาปตฺติผลํ สจฺฉิกเรยฺยาติ.}

\prose{โสตาปนฺเนน ปนาวุโส สาริปุตฺต ภิกฺขุนา กตเม ธมฺมา โยนิโส มนสิ กาตพฺพาติ.\csromanspacer{} โสตาปนฺเนนปิ โข อาวุโส \csromanfolio{138}โกฏฺฐิก ภิกฺขุนา อิเม ปญฺจุปาทานกฺขนฺธา อนิจฺจโต ฯเปฯ อนตฺตโต โยนิโส มนสิ กาตพฺพา.\csromanspacer{} ฐานํ โข ปเนตํ อาวุโส วิชฺชติ, ยํ โสตาปนฺโน ภิกฺขุ อิเม ปญฺจุปาทานกฺขนฺเธ อนิจฺจโต ฯเปฯ อนตฺตโต โยนิโส มนสิ กโรนฺโต สกทาคามิผลํ ฯเปฯ อนาคามิผลํ ฯเปฯ อรหตฺตผลํ สจฺฉิกเรยฺยาติ.}

\prose{อรหตา ปนาวุโส สาริปุตฺต กตเม ธมฺมา โยนิโส มนสิ กาตพฺพาติ.\csromanspacer{} อรหตาปิ ขฺวาวุโส โกฏฺฐิก อิเม ปญฺจุปาทานกฺขนฺธา อนิจฺจโต ทุกฺขโต โรคโต คณฺฑโต สลฺลโต อฆโต อาพาธโต ปรโต ปโลกโต สุญฺญโต อนตฺตโต โยนิโส มนสิ กาตพฺพา.\csromanspacer{} นตฺถิ ขฺวาวุโส อรหโต อุตฺตริ กรณียํ, กตสฺส วา ปติจโย, อปิ จ โข อิเม ธมฺมา ภาวิตา พหุลีกตา ทิฏฺฐธมฺมสุขวิหาราย เจว สํวตฺตนฺติ สติสมฺปชญฺญาย จาติ.\csromanspacer{}\csromanspacer{}เอกาทสมํ.}

\csromanmatikaanchor{mtk.24}
\csromantocmarkref{subsection}{12. กปฺปสุตฺต}{mtk.24}
\bookmark[dest={mtk.24},level=2]{12. กปฺปสุตฺต}
\csromanheader{2}{12. กปฺปสุตฺต}
\proseitem{124}{สาวตฺถินิทานํ.\csromanspacer{} อถ โข อายสฺมา กปฺโป เยน ภควา เตนุปสงฺกมิ.\csromanspacer{} เอกมนฺตํ นิสินฺโน โข อายสฺมา กปฺโป ภควนฺตํ เอตทโวจ “กถํ นุ โข ภนฺเต ชานโต กถํ ปสฺสโต อิมสฺมิํ จ สวิญฺญาณเก กาเย พหิทฺธา จ สพฺพนิมิตฺเตสุ อหงฺการมมงฺการมานานุสยา น โหนฺตี”ติ.}

\prose{ยํ กิญฺจิ กปฺป รูปํ อตีตานาคตปจฺจุปฺปนฺนํ อชฺฌตฺตํ วา พหิทฺธา วา โอฬาริกํ วา สุขุมํ วา หีนํ วา ปณีตํ วา ยํ ทูเร สนฺติเก วา, สพฺพํ รูปํ “เนตํ มม, เนโสหมสฺมิ, น เมโส อตฺตา”ติ เอวเมตํ ยถาภูตํ สมฺมปฺปญฺญาย ปสฺสติ.\csromanspacer{} ยา กาจิ เวทนา ฯเปฯ.\csromanspacer{} ยา กาจิ สญฺญา, เย เกจิ สงฺขารา, ยํ กิญฺจิ วิญฺญาณํ อตีตานาคตปจฺจุปฺปนฺนํ อชฺฌตฺตํ วา พหิทฺธา วา โอฬาริกํ วา สุขุมํ วา หีนํ วา ปณีตํ วา ยํ ทูเร สนฺติเก วา, สพฺพํ วิญฺญาณํ “เนตํ มม, เนโสหมสฺมิ, น เมโส อตฺตา”ติ เอวเมตํ ยถาภูตํ สมฺมปฺปญฺญาย ปสฺสติ.\csromanspacer{} เอวํ โข กปฺป ชานโต เอวํ ปสฺสโต อิมสฺมิํ จ สวิญฺญาณเก กาเย พหิทฺธา จ สพฺพนิมิตฺเตสุ อหงฺการมมงฺการมานานุสยา น โหนฺตีติ.\csromanspacer{}\csromanspacer{}ทฺวาทสมํ.}

\csromanmatikaanchor{mtk.25}
\csromantocmarkref{subsection}{13. ทุติยกปฺปสุตฺต}{mtk.25}
\bookmark[dest={mtk.25},level=2]{13. ทุติยกปฺปสุตฺต}
\csromanheader{2}{13. ทุติยกปฺปสุตฺต}
\proseitem{125}{\csromanfolio{139}สาวตฺถินิทานํ.\csromanspacer{} เอกมนฺตํ นิสินฺโน โข อายสฺมา กปฺโป ภควนฺตํ เอตทโวจ “กถํ นุ โข ภนฺเต ชานโต กถํ ปสฺสโต อิมสฺมิํ จ สวิญฺญาณเก กาเย พหิทฺธา จ สพฺพนิมิตฺเตสุ อหงฺกรมมงฺการมานาปคตํ มานสํ โหติ วิธา สมติกฺกนฺตํ สนฺตํ สุวิมุตฺตนฺ”ติ.}

\prose{ยํ กิญฺจิ กปฺป รูปํ อตีตานาคตปจฺจุปฺปนฺนํ ฯเปฯ สพฺพํ รูปํ “เนตํ มม, เนโสหมสฺมิ, น เมโส อตฺตา”ติ เอวเมตํ ยถาภูตํ สมฺมปฺปญฺญาย ทิสฺวา อนุปาทาวิมุตฺโต โหติ.\csromanspacer{} ยา กาจิ เวทนา.\csromanspacer{} ยา กาจิ สญฺญา.\csromanspacer{} เย เกจิ สงฺขารา.\csromanspacer{} ยํ กิญฺจิ วิญฺญาณํ อตีตานาคตปจฺจุปฺปนฺนํ อชฺฌตฺตํ วา พหิทฺธา วา โอฬาริกํ วา สุขุมํ วา หีนํ วา ปณีตํ วา ยํ ทูเร สนฺติเก วา, สพฺพํ วิญฺญาณํ “เนตํ มม, เนโสหมสฺมิ, น เมโส อตฺตา”ติ เอวเมตํ ยถาภูตํ สมฺมปฺปญฺญาย ทิสฺวา อนุปาทาวิมุตฺโต โหติ.\csromanspacer{} เอวํ โข กปฺป ชานโต เอวํ ปสฺสโต อิมสฺมิํ จ สวิญฺญาณเก กาเย พหิทฺธา จ สพฺพนิมิตฺเตสุ อหงฺการมมงฺการมานาปคตํ มานสํ โหติ วิธา สมติกฺกนฺตํ สนฺตํ สุวิมุตฺตนฺติ.\csromanspacer{}\csromanspacer{}เตรสมํ.\csromanspacer{} ธมฺมกถิกวคฺโค ทุติโย.}

\titlehead{\textbf{ตสฺสุทฺทานํ}}
\csromangathameasure{อวิชฺชา วิชฺชา เทฺว กถิกา, พนฺธนา ปริปุจฺฉิตา ทุเว. \\ สํโยชนํ อุปาทานํ, สีลํ สุตวา เทฺว จ กปฺเปนาติ.}
\csromangathagroup{\csromangathastanza{อวิชฺชา วิชฺชา เทฺว กถิกา, พนฺธนา ปริปุจฺฉิตา ทุเว. \\ สํโยชนํ อุปาทานํ, สีลํ สุตวา เทฺว จ กปฺเปนาติ.}}

\csromanmatikaanchor{mtk.26}
\csromantocmarknopage{section}{(13) 3. อวิชฺชาวคฺค}{mtk.26}
\bookmark[dest={mtk.26},level=1]{(13) 3. อวิชฺชาวคฺค}
\csromanheader{1}{\textbf{(13) 3. อวิชฺชาวคฺค}}
\csromanmatikaanchor{mtk.27}
\csromantocmarkref{subsection}{1-3. สมุทยธมฺมสุตฺต}{mtk.27}
\bookmark[dest={mtk.27},level=2]{1-3. สมุทยธมฺมสุตฺต}
\csromanheader{2}{1. สมุทยธมฺมสุตฺต}
\proseitem{126}{สาวตฺถินิทานํ.\csromanspacer{} อถ โข อญฺญตโร ภิกฺขุ เยน ภควา เตนุปสงฺกมิ, อุปสงฺกมิตฺวา ฯเปฯ เอกมนฺตํ นิสินฺโน โข โส ภิกฺขุ ภควนฺตํ เอตทโวจ “อวิชฺชา อวิชฺชาติ ภนฺเต วุจฺจติ, กตมา นุ โข ภนฺเต อวิชฺชา, กิตฺตาวตา จ อวิชฺชาคโต โหตี”ติ.}

\prose{\csromanfolio{140}อิธ ภิกฺขุ อสฺสุตวา ปุถุชฺชโน สมุทยธมฺมํ รูปํ “สมุทยธมฺมํ รูปนฺ”ติ ยถาภูตํ นปฺปชานาติ, วยธมฺมํ รูปํ “วยธมฺมํ รูปนฺ”ติ ยถาภูตํ นปฺปชานาติ, สมุทยวยธมฺมํ รูปํ “สมุทยวยธมฺมํ รูปนฺ”ติ ยถาภูตํ นปฺปชานาติ, สมุทยธมฺมํ เวทนํ “สมุทยธมฺมา เวทนา”ติ ยถาภูตํ นปฺปชานาติ, วยธมฺมํ เวทนํ “วยธมฺมา เวทนา”ติ ยถาภูตํ นปฺปชานาติ, สมุทยวยธมฺมํ เวทนํ “สมุทยวยธมฺมา เวทนา”ติ ยถาภูตํ นปฺปชานาติ, สมุทยธมฺมํ สญฺญํ ฯเปฯ สมุทยธมฺเม สงฺขาเร “สมุทยธมฺมา สงฺขารา”ติ ยถาภูตํ นปฺปชานาติ, วยธมฺเม สงฺขาเร “วยธมฺมา สงฺขารา”ติ ยถาภูตํ นปฺปชานาติ, สมุทยวยธมฺเม สงฺขาเร “สมุทยวยธมฺมา สงฺขารา”ติ ยถาภูตํ นปฺปชานาติ, สมุทยธมฺมํ วิญฺญาณํ “สมุทยธมฺมํ วิญฺญาณนฺ”ติ ยถาภูตํ นปฺปชานาติ, วยธมฺมํ วิญฺญาณํ “วยธมฺมํ วิญฺญาณนฺ”ติ ยถาภูตํ นปฺปชานาติ, สมุทยวยธมฺมํ วิญฺญาณํ “สมุทยวยธมฺมํ วิญฺญาณนฺ”ติ ยถาภูตํ นปฺปชานาติ.\csromanspacer{} อยํ วุจฺจติ ภิกฺขุ อวิชฺชา, เอตฺตาวตา จ อวิชฺชาคโต โหตีติ.}

\prose{เอวํ วุตฺเต โส ภิกฺขุ ภควนฺตํ เอตทโวจ “วิชฺชา วิชฺชาติ ภนฺเต วุจฺจติ, กตมา นุ โข ภนฺเต วิชฺชา, กิตฺตาวตา จ วิชฺชาคโต โหตี”ติ.}

\prose{อิธ ภิกฺขุ สุตวา อริยสาวโก สมุทยธมฺมํ รูปํ “สมุทยธมฺมํ รูปนฺ”ติ ยถาภูตํ ปชานาติ, วยธมฺมํ รูปํ “วยธมฺมํ รูปนฺ”ติ ยถาภูตํ ปชานาติ, สมุทยวยธมฺมํ รูปํ “สมุทยวยธมฺมํ รูปนฺ”ติ ยถาภูตํ ปชานาติ, สมุทยธมฺมํ เวทนํ “สมุทยธมฺมา เวทนา”ติ ยถาภูตํ ปชานาติ, วยธมฺมํ เวทนํ “วยธมฺมา เวทนา”ติ ยถาภูตํ ปชานาติ, สมุทยวยธมฺมํ เวทนํ “สมุทยวยธมฺมา เวทนา”ติ ยถาภูตํ ปชานาติ, สมุทยธมฺมํ สญฺญํ.\csromanspacer{} สมุทยธมฺเม สงฺขาเร “สมุทยธมฺมา สงฺขารา”ติ ยถาภูตํ ปชานาติ, วยธมฺเม สงฺขาเร “วยธมฺมา สงฺขารา”ติ ยถาภูตํ ปชานาติ, สมุทยวยธมฺเม สงฺขาเร “สมุทยวยธมฺมา สงฺขารา”ติ ยถาภูตํ ปชานาติ, สมุทยธมฺมํ วิญฺญาณํ “สมุทยธมฺมํ วิญฺญาณนฺ”ติ ยถาภูตํ ปชานาติ, วยธมฺมํ วิญฺญาณํ “วยธมฺมํ วิญฺญาณนฺ”ติ ยถาภูตํ ปชานาติ, สมุทยวยธมฺมํ วิญฺญาณํ “สมุทยวยธมฺมํ วิญฺญาณนฺ”ติ ยถาภูตํ ปชานาติ.\csromanspacer{} อยํ วุจฺจติ ภิกฺขุ วิชฺชา, เอตฺตาวตา จ วิชฺชาคโต โหตีติ.\csromanspacer{}\csromanspacer{}ปฐมํ.}

\csromanheader{2}{2. ทุติยสมุทยธมฺมสุตฺต}
\proseitem{127}{\csromanfolio{141}เอกํ สมยํ อายสฺมา จ สาริปุตฺโต อายสฺมา จ มหาโกฏฺฐิโก พาราณสิยํ วิหรนฺติ อิสิปตเน มิคทาเย.\csromanspacer{} อถ โข อายสฺมา มหาโกฏฺฐิโก สายนฺหสมยํ ปฏิสลฺลานา วุฏฺฐิโต ฯเปฯ เอกมนฺตํ นิสินฺโน โข อายสฺมา มหาโกฏฺฐิโก อายสฺมนฺตํ สาริปุตฺตํ เอตทโวจ “อวิชฺชา อวิชฺชาติ อาวุโส สาริปุตฺต วุจฺจติ, กตมา นุ โข อาวุโส อวิชฺชา, กิตฺตาวตา จ อวิชฺชาคโต โหตี”ติ.}

\prose{อิธาวุโส อสฺสุตวา ปุถุชฺชโน สมุทยธมฺมํ รูปํ “สมุทยธมฺมํ รูปนฺ”ติ ยถาภูตํ นปฺปชานาติ, วยธมฺมํ รูปํ ฯเปฯ “สมุทยวธมฺมํ รูปนฺ”ติ ยถาภูตํ นปฺปชานาติ, สมุทยธมฺมํ เวทนํ ฯเปฯ วยธมฺมํ เวทนํ ฯเปฯ “สมุทยวยธมฺมา เวทนา”ติ ยถาภูตํ นปฺปชานาติ, สมุทยธมฺมํ สญฺญํ ฯเปฯ สมุทยธมฺเม สงฺขาเร ฯเปฯ วยธมฺเม สงฺขาเร.\csromanspacer{} สมุทยวยธมฺเม สงฺขาเร “สมุทยวยธมฺมา สงฺขารา”ติ ยถาภูตํ นปฺปชานาติ, สมุทยธมฺมํ วิญฺญาณํ ฯเปฯ สมุทยวยธมฺมํ วิญฺญาณํ “สมุทยวยธมฺมํ วิญฺญาณนฺ”ติ ยถาภูตํ นปฺปชานาติ.\csromanspacer{} อยํ วุจฺจติ อาวุโส อวิชฺชา, เอตฺตาวตา จ อวิชฺชาคโต โหตีติ.\csromanspacer{}\csromanspacer{}ทุติยํ.}

\csromanheader{2}{3. ตติยสมุทยธมฺมสุตฺต}
\proseitem{128}{เอกํ สมยํ อายสฺมา จ สาริปุตฺโต อายสฺมา จ มหาโกฏฺฐิโก พาราณสิยํ วิหรนฺติ อิสิปตเน มิคทาเย.\csromanspacer{} เอกมนฺตํ นิสินฺโน โข อายสฺมา มหาโกฏฺฐิโก อายสฺมนฺตํ สาริปุตฺตํ เอตทโวจ “วิชฺชา วิชฺชาติ อาวุโส สาริปุตฺต วุจฺจติ, กตมา นุ โข อาวุโส วิชฺชา, กิตฺตาวตา จ วิชฺชาคโต โหตี”ติ.}

\prose{อิธาวุโส สุตวา อริยสาวโก สมุทยธมฺมํ รูปํ “สมุทยธมฺมํ รูปนฺ”ติ ยถาภูตํ ปชานาติ, วยธมฺมํ รูปํ ฯเปฯ สมุทยวยธมฺมํ รูปํ “สมุทยวยธมฺมํ รูปนฺ”ติ ยถาภูตํ ปชานาติ, สมุทยธมฺมํ เวทนํ ฯเปฯ สมุทยวยธมฺมา เวทนา.\csromanspacer{} สมุทยธมฺมํ สญฺญํ ฯเปฯ สมุทยธมฺเม สงฺขาเร.\csromanspacer{} วยธมฺเม สงฺขาเร.\csromanspacer{} สมุทยวยธมฺเม สงฺขาเร “สมุทยวยธมฺมา สงฺขารา”ติ ยถาภูตํ ปชานาติ, สมุทยธมฺมํ วิญฺญาณํ.\csromanspacer{} วยธมฺมํ วิญฺญาณํ.\csromanspacer{} สมุทยวยธมฺมํ วิญฺญาณํ “สมุทยวยธมฺมํ วิญฺญาณนฺ”ติ ยถาภูตํ ปชานาติ.\csromanspacer{} อยํ วุจฺจตาวุโส วิชฺชา, เอตฺตาวตา จ วิชฺชาคโต โหตีติ.\csromanspacer{}\csromanspacer{}ตติยํ.}

\csromanmatikaanchor{mtk.28}
\csromantocmarkref{subsection}{4-5. อสฺสาทสุตฺต}{mtk.28}
\bookmark[dest={mtk.28},level=2]{4-5. อสฺสาทสุตฺต}
\csromanheader{2}{4. อสฺสาทสุตฺต}
\proseitem{129}{\csromanfolio{142}พาราณสิยํ วิหรนฺติ อิสิปตเน มิคทาเย ฯเปฯ เอกมนฺตํ นิสินฺโน โข อายสฺมา มหาโกฏฺฐิโก อายสฺมนฺตํ สาริปุตฺตํ เอตทโวจ “อวิชฺชา อวิชฺชาติ อาวุโส สาริปุตฺต วุจฺจติ, กตมา นุ โข อาวุโส อวิชฺชา, กิตฺตาวตา จ อวิชฺชาคโต โหตี”ติ.}

\prose{อิธาวุโส อสฺสุตวา ปุถุชฺชโน รูปสฺส อสฺสาทญฺจ อาทีนวญฺจ นิสฺสรณญฺจ ยถาภูตํ นปฺปชานาติ, เวทนาย ฯเปฯ สญฺญาย.\csromanspacer{} สงฺขารานํ.\csromanspacer{} วิญฺญาณสฺส อสฺสาทญฺจ อาทีนวญฺจ นิสฺสรณญฺจ ยถาภูตํ นปฺปชานาติ.\csromanspacer{} อยํ วุจฺจตาวุโส อวิชฺชา, เอตฺตาวตา จ อวิชฺชาคโต โหตีติ.\csromanspacer{}\csromanspacer{}จตุตฺถํ.}

\csromanheader{2}{5. ทุติย-อสฺสาทสุตฺต}
\proseitem{130}{พาราณสิยํ วิหรนฺติ อิสิปตเน มิคทาเย.\csromanspacer{} วิชฺชา วิชฺชาติ อาวุโส สาริปุตฺต วุจฺจติ, กตมา นุ โข อาวุโส วิชฺชา, กิตฺตาวตา จ วิชฺชาคโต โหตีติ.}

\prose{อิธาวุโส สุตวา อริยสาวโก รูปสฺส อสฺสาทญฺจ อาทีนวญฺจ นิสฺสรณญฺจ ยถาภูตํ ปชานาติ, เวทนาย ฯเปฯ สญฺญาย.\csromanspacer{} สงฺขารานํ.\csromanspacer{} วิญฺญาณสฺส อสฺสาทญฺจ อาทีนวญฺจ นิสฺสรณญฺจ ยถาภูตํ ปชานาติ.\csromanspacer{} อยํ วุจฺจตาวุโส วิชฺชา, เอตฺตาวตา จ วิชฺชาคโต โหตีติ.\csromanspacer{}\csromanspacer{}ปญฺจมํ.}

\csromanmatikaanchor{mtk.29}
\csromantocmarkref{subsection}{6-7. สมุทยสุตฺต}{mtk.29}
\bookmark[dest={mtk.29},level=2]{6-7. สมุทยสุตฺต}
\csromanheader{2}{6. สมุทยสุตฺต}
\proseitem{131}{พาราณสิยํ วิหรนฺติ อิสิปตเน มิคทาเย ฯเปฯ “อวิชฺชา อวิชฺชาติ อาวุโส สาริปุตฺต วุจฺจติ, กตมา นุ โข อาวุโส อวิชฺชา, กิตฺตาวตา จ อวิชฺชาคโต โหตีติ.}

\prose{อิธาวุโส อสฺสุตวา ปุถุชฺชโน รูปสฺส สมุทยญฺจ อตฺถงฺคมญฺจ อสฺสาทญฺจ อาที นวญฺจ นิสฺสรณญฺจ ยถาภูตํ นปฺปชานาติ, เวทนาย ฯเปฯ สญฺญาย.\csromanspacer{} สงฺขารานํ.\csromanspacer{} วิญฺญาณสฺส สมุทยญฺจ อตฺถงฺคมญฺจ อสฺสาทญฺจ อาทีนวญฺจ นิสฺสรณญฺจ ยถาภูตํ นปฺปชานาติ.\csromanspacer{} อยํ วุจฺจตาวุโส อวิชฺชา เอตฺตาวตา จ อวิชฺชาคโต โหตีติ.\csromanspacer{}\csromanspacer{}ฉฏฺฐํ.}

\csromanheader{2}{7. ทุติยสมุทยสุตฺต}
\proseitem{132}{\csromanfolio{143}พาราณสิยํ วิหรนฺติ อิสิปตเน มิคทาเย ฯเปฯ เอกมนฺตํ นิสินฺโน โข อายสฺมา มหาโกฏฺฐิโก อายสฺมนฺตํ สาริปุตฺตํ เอตทโวจ “วิชฺชา วิชฺชาติ อาวุโส สาริปุตฺต วุจฺจติ, กตมา นุ โข อาวุโส วิชฺชา, กิตฺตาวตา จ วิชฺชาคโต โหตี”ติ.}

\prose{อิธาวุโส สุตวา อริยสาวโก รูปสฺส สมุทยญฺจ อตฺถงฺคมญฺจ อสฺสาทญฺจ อาทีนวญฺจ นิสฺสรณญฺจ ยถาภูตํ ปชานาติ, เวทนาย ฯเปฯ สญฺญาย.\csromanspacer{} สงฺขารานํ.\csromanspacer{} วิญฺญาณสฺส สมุทยญฺจ อตฺถงฺคมญฺจ อสฺสาทญฺจ อาทีนวญฺจ นิสฺสรณญฺจ ยถาภูตํ ปชานาติ.\csromanspacer{} อยํ วุจฺจตาวุโส วิชฺชา, เอตฺตาวตา จ วิชฺชาคโต โหตีติ.\csromanspacer{}\csromanspacer{}สตฺตมํ.}

\csromanmatikaanchor{mtk.30}
\csromantocmarkref{subsection}{8-10. โกฏฺฐิกสุตฺต}{mtk.30}
\bookmark[dest={mtk.30},level=2]{8-10. โกฏฺฐิกสุตฺต}
\csromanheader{2}{8. โกฏฺฐิกสุตฺต}
\proseitem{133}{พาราณสิยํ วิหรนฺติ อิสิปตเน มิคทาเย.\csromanspacer{} อถ โข อายสฺมา สาริปุตฺโต สายนฺหสมยํ ฯเปฯ เอกมนฺตํ นิสินฺโน โข อายสฺมา สาริปุตฺโต อายสฺมนฺตํ มหาโกฏฺฐิกํ เอตทโวจ “อวิชฺชา อวิชฺชาติ อาวุโส โกฏฺฐิก วุจฺจติ, กตมา นุ โข อาวุโส อวิชฺชา, กิตฺตาวตา จ อวิชฺชาคโต โหตี”ติ.}

\prose{อิธาวุโส อสฺสุตวา ปุถุชฺชโน รูปสฺส อสฺสาทญฺจ อาทีนวญฺจ นิสฺสรณญฺจ ยถาภูตํ นปฺปชานาติ, เวทนาย ฯเปฯ สญฺญาย.\csromanspacer{} สงฺขารานํ.\csromanspacer{} วิญฺญาณสฺส อสฺสาทญฺจ อาทีนวญฺจ นิสฺสรณญฺจ ยถาภูตํ นปฺปชานาติ.\csromanspacer{} อยํ วุจฺจตาวุโส อวิชฺชา, เอตฺตาวตา จ อวิชฺชาคโต โหตีติ.}

\prose{เอวํ วุตฺเต อายสฺมา สาริปุตฺโต อายสฺมนฺตํ มหาโกฏฺฐิกํ เอตทโวจ “วิชฺชา วิชฺชาติ อาวุโส โกฏฺฐิก วุจฺจติ, กตมา นุ โข อาวุโส วิชฺชา, กิตฺตาวตา จ วิชฺชาคโต โหตี”ติ.}

\prose{อิธาวุโส สุตวา อริยสาวโก รูปสฺส อสฺสาทญฺจ อาทีนวญฺจ นิสฺสรณญฺจ ยถาภูตํ ปชานาติ, เวทนาย ฯเปฯ สญฺญาย.\csromanspacer{} สงฺขารานํ.\csromanspacer{} วิญฺญาณสฺส อสฺสาทญฺจ อาทีนวญฺจ นิสฺสรณญฺจ ยถาภูตํ ปชานาติ.\csromanspacer{} อยํ วุจฺจตาวุโส วิชฺชา, เอตฺตาวตา จ วิชฺชาคโต โหตีติ.\csromanspacer{}\csromanspacer{}อฏฺฐมํ.}

\csromanheader{2}{9. ทุติยโกฏฺฐิกสุตฺต}
\proseitem{134}{\csromanfolio{144}พาราณสิยํ วิหรนฺติ อิสิปตเน มิคทาเย.\csromanspacer{} อวิชฺชา อวิชฺชาติ อาวุโส โกฏฺฐิก วุจฺจติ, กตมา นุ โข อาวุโส อวิชฺชา, กิตฺตาวตา จ อวิชฺชาคโต โหตีติ.}

\prose{อิธาวุโส อสฺสุตวา ปุถุชฺชโน รูปสฺส สมุทยญฺจ อตฺถงฺคมญฺจ อสฺสาทญฺจ อาทีนวญฺจ นิสฺสรณญฺจ ยถาภูตํ นปฺปชานาติ, เวทนาย ฯเปฯ สญฺญาย.\csromanspacer{} สงฺขารานํ.\csromanspacer{} วิญฺญาณสฺส สมุทยญฺจ อตฺถงฺคมญฺจ อสฺสาทญฺจ อาทีนวญฺจ นิสฺสรณญฺจ ยถาภูตํ นปฺปชานาติ.\csromanspacer{} อยํ วุจฺจตาวุโส อวิชฺชา, เอตฺตาวตา จ อวิชฺชาคโต โหตีติ.}

\prose{เอวํ วุตฺเต อายสฺมา สาริปุตฺโต อายสฺมนฺตํ มหาโกฏฺฐิกํ เอตทโวจ “วิชฺชา วิชฺชาติ อาวุโส โกฏฺฐิก วุจฺจติ, กตมา นุ โข อาวุโส วิชฺชา, กิตฺตาวตา จ วิชฺชาคโต โหตี”ติ.}

\prose{อิธาวุโส สุตวา อริยสาวโก รูปสฺส สมุทยญฺจ อตฺถงฺคมญฺจ อสฺสาทญฺจ อาทีนวญฺจ นิสฺสรณญฺจ ยถาภูตํ ปชานาติ, เวทนาย ฯเปฯ สญฺญาย.\csromanspacer{} สงฺขารานํ.\csromanspacer{} วิญฺญาณสฺส สมุทยญฺจ อตฺถงฺคมญฺจ อสฺสาทญฺจ อาทีนวญฺจ นิสฺสรณญฺจ ยถาภูตํ ปชานาติ.\csromanspacer{} อยํ วุจฺจตาวุโส วิชฺชา.\csromanspacer{} เอตฺตาวตา จ วิชฺชาคโต โหตีติ.\csromanspacer{}\csromanspacer{}นวมํ.}

\csromanheader{2}{10. ตติยโกฏฺฐิกสุตฺต}
\proseitem{135}{ตญฺเญว นิทานํ.\csromanspacer{} เอกมนฺตํ นิสินฺโน โข อายสฺมา สาริปุตฺโต อายสฺมนฺตํ มหาโกฏฺฐิกํ เอตทโวจ “อวิชฺชา อวิชฺชาติ อาวุโส โกฏฺฐิก วุจฺจติ, กตมา นุ โข อาวุโส อวิชฺชา, กิตฺตาวตา จ อวิชฺชาคโต โหตี”ติ.}

\prose{อิธาวุโส อสฺสุตวา ปุถุชฺชโน รูปํ นปฺปชานาติ, รูปสมุทยํ นปฺปชานาติ, รูปนิโรธํ นปฺปชานาติ, รูปนิโรธคามินิํ ปฏิปทํ นปฺปชานาติ.\csromanspacer{} เวทนํ นปฺปชานาติ ฯเปฯ สญฺญํ.\csromanspacer{} สงฺขาเร.\csromanspacer{} วิญฺญาณํ นปฺปชานาติ, วิญฺญาณสมุทยํ นปฺปชานาติ, วิญฺญาณนิโรธํ นปฺปชานาติ, วิญฺญาณนิโรธคามินิํ ปฏิปทํ นปฺปชานาติ.\csromanspacer{} อยํ วุจฺจตาวุโส อวิชฺชา, เอตฺตาวตา จ อวิชฺชาคโต โหตีติ.}

\prose{\csromanfolio{145}เอวํ วุตฺเต อายสฺมา สาริปุตฺโต อายสฺมนฺตํ มหาโกฏฺฐิกํ เอตทโวจ\csromandash{} “วิชฺชา วิชฺชา”ติ อาวุโส โกฏฺฐิก วุจฺจติ, กตมา นุ โข อาวุโส วิชฺชา, กิตฺตาวตา จ วิชฺชาคโต โหตีติ.\csromanspacer{} อิธาวุโส สุตวา อริยสาวโก รูปํ ปชานาติ, รูปสมุทยํ ปชานาติ, รูปนิโรธํ ปชานาติ, รูปนิโรธคามินิํ ปฏิปทํ ปชานาติ.\csromanspacer{} เวทนํ.\csromanspacer{} สญฺญํ.\csromanspacer{} สงฺขาเร.\csromanspacer{} วิญฺญาณํ ปชานาติ, วิญฺญาณสมุทยํ ปชานาติ, วิญฺญาณนิโรธํ ปชานาติ, วิญฺญาณนิโรธคามินิํ ปฏิปทํ ปชานาติ.\csromanspacer{} อยํ วุจฺจตาวุโส วิชฺชา, เอตฺตาวตา จ วิชฺชาคโต โหตีติ.\csromanspacer{}\csromanspacer{}ทสมํ.\csromanspacer{} อวิชฺชาวคฺโค ตติโย.}

\titlehead{\textbf{ตสฺสุทฺทานํ}}
\csromangathameasure{สมุทยธมฺเม ตีณิ, อสฺสาโท อปเร ทุเว. \\ สมุทเย จ เทฺว วุตฺตา, โกฏฺฐิเก อปเร ตโยติ.}
\csromangathagroup{\csromangathastanza{สมุทยธมฺเม ตีณิ, อสฺสาโท อปเร ทุเว. \\ สมุทเย จ เทฺว วุตฺตา, โกฏฺฐิเก อปเร ตโยติ.}}

\csromanmatikaanchor{mtk.31}
\csromantocmarknopage{section}{(14) 4. กุกฺกุฬวคฺค}{mtk.31}
\bookmark[dest={mtk.31},level=1]{(14) 4. กุกฺกุฬวคฺค}
\csromanheader{1}{\textbf{(14) 4. กุกฺกุฬวคฺค}}
\csromanmatikaanchor{mtk.32}
\csromantocmarkref{subsection}{1. กุกฺกุฬสุตฺต}{mtk.32}
\bookmark[dest={mtk.32},level=2]{1. กุกฺกุฬสุตฺต}
\csromanheader{2}{1. กุกฺกุฬสุตฺต}
\proseitem{136}{สาวตฺถิ นิทานํ.\csromanspacer{} รูปํ ภิกฺขเว กุกฺกุฬํ, เวทนา กุกฺกุฬา, สญฺญา กุกฺกุฬา, สงฺขารา กุกฺกุฬา, วิญฺญาณํ กุกฺกุฬํ, เอวํ ปสฺสํ ภิกฺขเว สุตวา อริยสาวโก รูปสฺมิมฺปิ นิพฺพินฺทติ, เวทนายปิ นิพฺพินฺทติ, สญฺญายปิ นิพฺพินฺทติ, สงฺขาเรสุปิ นิพฺพินฺทติ, วิญฺญาณสฺมิมฺปิ นิพฺพินฺทติ, นิพฺพินฺทํ วิรชฺชติ, วิราคา วิมุจฺจติ วิมุตฺตสฺมิํ “วิมุตฺตมฺ”อิติ ญาณํ โหติ “ขีณา ชาติ, วุสิตํ พฺรหฺมจริยํ, กตํ กรณียํ, นาปรํ อิตฺถตฺตายาติ ปชานาตี”ติ.\csromanspacer{}\csromanspacer{}ปฐมํ.}

\csromanmatikaanchor{mtk.33}
\csromantocmarkref{subsection}{2-4. อนิจฺจสุตฺต}{mtk.33}
\bookmark[dest={mtk.33},level=2]{2-4. อนิจฺจสุตฺต}
\csromanheader{2}{2. อนิจฺจสุตฺต}
\proseitem{137}{สาวตฺถินิทานํ.\csromanspacer{} ยํ ภิกฺขเว อนิจฺจํ, ตตฺร โว ฉนฺโท ปหาตพฺโพ.\csromanspacer{} กิญฺจ ภิกฺขเว อนิจฺจํ.\csromanspacer{} รูปํ ภิกฺขเว อนิจฺจํ, ตตฺร โว ฉนฺโท ปหาตพฺโพ.\csromanspacer{} เวทนา อนิจฺจา ฯเปฯ.\csromanspacer{} สญฺญา.\csromanspacer{} สงฺขารา.\csromanspacer{} วิญฺญาณํ อนิจฺจํ, ตตฺร โว ฉนฺโท ปหาตพฺโพ.\csromanspacer{} ยํ ภิกฺขเว อนิจฺจํ, ตตฺร โว ฉนฺโท ปหาตพฺโพติ.\csromanspacer{}\csromanspacer{}ทุติยํ.}

\csromanheader{2}{3. ทุติย-อนิจฺจสุตฺต}
\proseitem{138}{\csromanfolio{146}สาวตฺถินิทานํ.\csromanspacer{} ยํ ภิกฺขเว อนิจฺจํ, ตตฺร โว ราโค ปหาตพฺโพ.\csromanspacer{} กิญฺจ ภิกฺขเว อนิจฺจํ.\csromanspacer{} รูปํ ภิกฺขเว อนิจฺจํ, ตตฺร โว ราโค ปหาตพฺโพ.\csromanspacer{} เวทนา อนิจฺจา.\csromanspacer{} สญฺญา.\csromanspacer{} สงฺขารา.\csromanspacer{} วิญฺญาณํ อนิจฺจํ, ตตฺร โว ราโค ปหาตพฺโพ.\csromanspacer{} ยํ ภิกฺขเว อนิจฺจํ, ตตฺร โว ราโค ปหาตพฺโพติ.\csromanspacer{}\csromanspacer{}ตติยํ.}

\csromanheader{2}{4. ตติย-อนิจฺจสุตฺต}
\proseitem{139}{สาวตฺถินิทานํ.\csromanspacer{} ยํ ภิกฺขเว อนิจฺจํ, ตตฺร โว ฉนฺทราโค ปหาตพฺโพ.\csromanspacer{} กิญฺจ ภิกฺขเว อนิจฺจํ.\csromanspacer{} รูปํ ภิกฺขเว อนิจฺจํ, ตตฺร โว ฉนฺทราโค ปหาตพฺโพ.\csromanspacer{} เวทนา อนิจฺจา.\csromanspacer{} สญฺญา.\csromanspacer{} สงฺขารา.\csromanspacer{} วิญฺญาณํ อนิจฺจํ, ตตฺร โว ฉนฺทราโค ปหาตพฺโพ.\csromanspacer{} ยํ ภิกฺขเว อนิจฺจํ, ตตฺร โว ฉนฺทราโค ปหาตพฺโพติ.\csromanspacer{}\csromanspacer{}จตุตฺถํ.}

\csromanmatikaanchor{mtk.34}
\csromantocmarkref{subsection}{5-7. ทุกฺขสุตฺต}{mtk.34}
\bookmark[dest={mtk.34},level=2]{5-7. ทุกฺขสุตฺต}
\csromanheader{2}{5. ทุกฺขสุตฺต}
\proseitem{140}{สาวตฺถินิทานํ.\csromanspacer{} ยํ ภิกฺขเว ทุกฺขํ, ตตฺร โว ฉนฺโท ปหาตพฺโพ ฯเปฯ.\csromanspacer{} ยํ ภิกฺขเว ทุกฺขํ, ตตฺร โว ฉนฺโท ปหาตพฺโพติ.\csromanspacer{}\csromanspacer{}ปญฺจมํ.}

\csromanheader{2}{6. ทุติยทุกฺขสุตฺต}
\proseitem{141}{สาวตฺถินิทานํ.\csromanspacer{} ยํ ภิกฺขเว ทุกฺขํ, ตตฺร โว ราโค ปหาตพฺโพ ฯเปฯ.\csromanspacer{} ยํ ภิกฺขเว ทุกฺขํ, ตตฺร โว ราโค ปหาตพฺโพติ.\csromanspacer{}\csromanspacer{}ฉฏฺฐํ.}

\csromanheader{2}{7. ตติยทุกฺขสุตฺต}
\proseitem{142}{สาวตฺถินิทานํ.\csromanspacer{} ยํ ภิกฺขเว ทุกฺขํ, ตตฺร โว ฉนฺทราโคปหาตพฺโพ ฯเปฯ.\csromanspacer{} ยํ ภิกฺขเว ทุกฺขํ, ตตฺร โว ฉนฺทราโค ปหาตพฺโพติ.\csromanspacer{}\csromanspacer{}สตฺตมํ.}

\csromanmatikaanchor{mtk.35}
\csromantocmarkref{subsection}{8-10. อนตฺตสุตฺต}{mtk.35}
\bookmark[dest={mtk.35},level=2]{8-10. อนตฺตสุตฺต}
\csromanheader{2}{8. อนตฺตสุตฺต}
\proseitem{143}{สาวตฺถินิทานํ.\csromanspacer{} โย ภิกฺขเว อนตฺตา, ตตฺร โว ฉนฺโท ปหาตพฺโพ.\csromanspacer{} โก จ ภิกฺขเว อนตฺตา.\csromanspacer{} รูปํ ภิกฺขเว อนตฺตา, ตตฺร โว ฉนฺโท ปหาตพฺโพ.\csromanspacer{} เวทนา อนตฺตา.\csromanspacer{} สญฺญา.\csromanspacer{} สงฺขารา.\csromanspacer{} วิญฺญาณํ อนตฺตา, ตตฺร โว ฉนฺโท ปหาตพฺโพ.\csromanspacer{} โย ภิกฺขเว อนตฺตา, ตตฺร โว ฉนฺโท ปหาตพฺโพติ.\csromanspacer{}\csromanspacer{}อฏฺฐมํ.}

\csromanheader{2}{9. ทุติย-อนตฺตสุตฺต}
\csromanmatikaanchor{mtk.36}
\csromanmatikaanchor{mtk.37}
\csromantocmarkref{subsection}{11. นิพฺพิทาพหุลสุตฺต}{mtk.36}
\bookmark[dest={mtk.36},level=2]{11. นิพฺพิทาพหุลสุตฺต}
\csromantocmarkref{subsection}{12-14. อนิจฺจานุปสฺสีสุตฺตาทิ}{mtk.37}
\bookmark[dest={mtk.37},level=2]{12-14. อนิจฺจานุปสฺสีสุตฺตาทิ}
\proseitem{144}{\csromanfolio{147}สาวตฺถินิทานํ.\csromanspacer{} โย ภิกฺขเว อนตฺตา, ตตฺร โว ราโค ปหาตพฺโพ.\csromanspacer{} โก จ ภิกฺขเว อนตฺตา.\csromanspacer{} รูปํ ภิกฺขเว อนตฺตา, ตตฺร โว ราโค ปหาตพฺโพ.\csromanspacer{} เวทนา อนตฺตา.\csromanspacer{} สญฺญา.\csromanspacer{} สงฺขารา.\csromanspacer{} วิญฺญาณํ อนตฺตา, ตตฺร โว ราโค ปหาตพฺโพ.\csromanspacer{} โย ภิกฺขเว อนตฺตา, ตตฺร โว ราโค ปหาตพฺโพติ.\csromanspacer{}\csromanspacer{}นวมํ.}

\csromanheader{2}{10. ตติย-อนตฺตสุตฺต}
\proseitem{145}{สาวตฺถินิทานํ.\csromanspacer{} โย ภิกฺขเว อนตฺตา, ตตฺร โว ฉนฺทราโค ปหาตพฺโพ.\csromanspacer{} โก จ ภิกฺขเว อนตฺตา.\csromanspacer{} รูปํ ภิกฺขเว อนตฺตา, ตตฺร โว ฉนฺทราโค ปหาตพฺโพ.\csromanspacer{} เวทนา อนตฺตา.\csromanspacer{} สญฺญา.\csromanspacer{} สงฺขารา.\csromanspacer{} วิญฺญาณํ อนตฺตา, ตตฺร โว ฉนฺทราโค ปหาตพฺโพ.\csromanspacer{} โย ภิกฺขเว อนตฺตา, ตตฺร โว ฉนฺทราโค ปหาตพฺโพติ.\csromanspacer{}\csromanspacer{}ทสมํ.}

\csromanheader{2}{11. นิพฺพิทาพหุลสุตฺต}
\proseitem{146}{สาวตฺถินิทานํ.\csromanspacer{} สทฺธาปพฺพชิตสฺส ภิกฺขเว กุลปุตฺตสฺส อยมนุธมฺโม โหติ.\csromanspacer{} ยํ รูเป นิพฺพิทาพหุโล\footnote{นิพฺพิทาพหุลํ (สฺยา, กํ, อิ, ก)} วิหเรยฺย.\csromanspacer{} เวทนาย ฯเปฯ.\csromanspacer{} สญฺญาย.\csromanspacer{} สงฺขาเรสุ.\csromanspacer{} วิญฺญาเณ นิพฺพิทาพหุโล วิหเรยฺย.\csromanspacer{} โย รูเป นิพฺพิทาพหุโล วิหรนฺโต.\csromanspacer{} เวทนาย.\csromanspacer{} สญฺญาย.\csromanspacer{} สงฺขาเรสุ.\csromanspacer{} วิญฺญาเณ นิพฺพิทาพหุโล วิหรนฺโต รูปํ ปริชานาติ.\csromanspacer{} เวทนํ.\csromanspacer{} สญฺญํ.\csromanspacer{} สงฺขาเร.\csromanspacer{} วิญฺญาณํ ปริชานาติ.\csromanspacer{} โส รูปํ ปริชานํ, เวทนํ ปริชานํ, สญฺญํ ปริชานํ, สงฺขาเร ปริชานํ, วิญฺญาณํ ปริชานํ ปริมุจฺจติ รูปมฺหา, ปริมุจฺจติ เวทนาย, ปริมุจฺจติ สญฺญาย, ปริมุจฺจติ สงฺขาเรหิ, ปริมุจฺจติ วิญฺญาณมฺหา, ปริมุจฺจติ ชาติยา ชราย มรเณน โสเกหิ ปริเทเวหิ ทุกฺเขหิ โทมนสฺเสหิ อุปายาเสหิ, “ปริมุจฺจติ ทุกฺขสฺมา”ติ วทามีติ.\csromanspacer{}\csromanspacer{}เอกาทสมํ.}

\csromanheader{2}{12. อนิจฺจานุปสฺสีสุตฺต}
\proseitem{147}{สาวตฺถินิทานํ.\csromanspacer{} สทฺธาปพฺพชิตสฺส ภิกฺขเว กุลปุตฺตสฺส อยมนุ ธมฺโม โหติ.\csromanspacer{} ยํ รูเป อนิจฺจานุปสฺสี วิหเรยฺย.\csromanspacer{} เวทนาย.\csromanspacer{} สญฺญาย.\csromanspacer{} สงฺขาเรสุ.\csromanspacer{} วิญฺญาเณ อนิจฺจานุปสฺสี วิหเรยฺย ฯเปฯ “ปริมุจฺจติ ทุกฺขสฺมา”ติ วทามีติ.\csromanspacer{}\csromanspacer{}ทฺวาทสมํ.}

\csromanheader{2}{13. ทุกฺขานุปสฺสีสุตฺต}
\proseitem{148}{\csromanfolio{148}สาวตฺถินิทานํ.\csromanspacer{} สทฺธาปพฺพชิตสฺส ภิกฺขเว กุลปุตฺตสฺส อยมนุธมฺโม โหติ.\csromanspacer{} ยํ รูเป ทุกฺขานุปสฺสี วิหเรยฺย.\csromanspacer{} เวทนาย.\csromanspacer{} สญฺญาย.\csromanspacer{} สงฺขาเรสุ.\csromanspacer{} วิญฺญาเณ ทุกฺขานุปสฺสี วิหเรยฺย ฯเปฯ “ปริมุจฺจติ ทุกฺขสฺมา”ติ วทามีติ.\csromanspacer{}\csromanspacer{}เตรสมํ.}

\csromanheader{2}{14. อนตฺตานุปสฺสีสุตฺต}
\proseitem{149}{สาวตฺถินิทานํ.\csromanspacer{} สทฺธาปพฺพชิตสฺส ภิกฺขเว กุลปุตฺตสฺส อยมนุธมฺโม โหติ.\csromanspacer{} ยํ รูเป อนตฺตานุปสฺสี วิหเรยฺย.\csromanspacer{} เวทนาย.\csromanspacer{} สญฺญาย.\csromanspacer{} สงฺขาเรสุ.\csromanspacer{} วิญฺญาเณ อนตฺตานุปสฺสี วิหเรยฺย.\csromanspacer{} อนตฺตานุปสฺสี วิหรนฺโต.\csromanspacer{} เวทนาย.\csromanspacer{} สญฺญาย.\csromanspacer{} สงฺขาเรสุ.\csromanspacer{} วิญฺญาเณ อนตฺตานุปสฺสี วิหรนฺโต รูปํ ปริชานาติ.\csromanspacer{} เวทนํ ฯเปฯ.\csromanspacer{} สญฺญํ.\csromanspacer{} สงฺขาเร.\csromanspacer{} วิญฺญาณํ ปริชานาติ.\csromanspacer{} โส รูปํ ปริชานํ, เวทนํ ปริชานํ, สญฺญํ ปริชานํ, สงฺขาเร ปริชานํ, วิญฺญาณํ ปริชานํ ปริมุจฺจติ รูปมฺหา, ปริมุจฺจติ เวทนาย, ปริมุจฺจติ สญฺญาย, ปริมุจฺจติ สงฺขาเรหิ, ปริมุจฺจติ วิญฺญาณมฺหา, ปริมุจฺจติ ชาติยา ชาราย มรเณน โสเกหิ ปริเทเวหิ ทุกฺเขหิ โทมนสฺเสหิ อุปายาเสหิ, “ปริมุจฺจติ ทุกฺขสฺมา”ติ วทามีติ.\csromanspacer{}\csromanspacer{}จุทฺทสมํ.\csromanspacer{} กุกฺกุฬวคฺโค จตุตฺโถ.}

\titlehead{\textbf{ตสฺสุทฺทานํ}}
\csromangathameasure{กุกฺกุฬา ตโย อนิจฺเจน, ทุกฺเขน อปเร ตโย. \\ อนตฺเตน ตโย วุตฺตา, กุลปุตฺเตน เทฺว ทุกาติ.}
\csromangathagroup{\csromangathastanza{กุกฺกุฬา ตโย อนิจฺเจน, ทุกฺเขน อปเร ตโย. \\ อนตฺเตน ตโย วุตฺตา, กุลปุตฺเตน เทฺว ทุกาติ.}}

\csromanmatikaanchor{mtk.38}
\csromantocmarknopage{section}{(15) 5. ทิฏฺฐิวคฺค}{mtk.38}
\bookmark[dest={mtk.38},level=1]{(15) 5. ทิฏฺฐิวคฺค}
\csromanheader{1}{\textbf{(15) 5. ทิฏฺฐิวคฺค}}
\csromanmatikaanchor{mtk.39}
\csromantocmarkref{subsection}{1. อชฺฌตฺตสุตฺต}{mtk.39}
\bookmark[dest={mtk.39},level=2]{1. อชฺฌตฺตสุตฺต}
\csromanheader{2}{1. อชฺฌตฺตสุตฺต}
\proseitem{150}{สาวตฺถินิทานํ.\csromanspacer{} กิสฺมิํ นุ โข ภิกฺขเว สติ กิํ อุปาทาย อุปฺปชฺชติ อชฺฌตฺตํ สุขทุกฺขนฺติ.\csromanspacer{} ภควํมูลกา โน ภนฺเต ธมฺมา ฯเปฯ.\csromanspacer{} รูเป โข ภิกฺขเว สติ รูปํ อุปาทาย อุปฺปชฺชติ อชฺฌตฺตํ สุขทุกฺขํ.\csromanspacer{} เวทนาย สติ ฯเปฯ. \csromanfolio{149}สญฺญาย สติ.\csromanspacer{} สงฺขาเรสุ สติ.\csromanspacer{} วิญฺญาเณ สติ วิญฺญาณํ อุปาทาย อุปฺปชฺชติ อชฺฌตฺตํ สุขทุกฺขํ.\csromanspacer{} ตํ กิํ มญฺญถ ภิกฺขเว, รูปํ นิจฺจํ วา อนิจฺจํ วาติ.\csromanspacer{} อนิจฺจํ ภนฺเต.\csromanspacer{} ยํ ปนานิจฺจํ, ทุกฺขํ วา ตํ สุขํ วาติ.\csromanspacer{} ทุกฺขํ ภนฺเต.\csromanspacer{} ยํ ปนานิจฺจํ ทุกฺขํ วิปริณามธมฺมํ, อปิ นุ ตํ อนุปาทาย อุปฺปชฺเชยฺย อชฺฌตฺตํ สุขทุกฺขนฺติ.\csromanspacer{} โน เหตํ ภนฺเต.\csromanspacer{} เวทนา ฯเปฯ.\csromanspacer{} สญฺญา.\csromanspacer{} สงฺขารา.\csromanspacer{} วิญฺญาณํ นิจฺจํ วา อนิจฺจํ วาติ.\csromanspacer{} อนิจฺจํ ภนฺเต.\csromanspacer{} ยํ ปนานิจฺจํ, ทุกฺขํ วา ตํ สุขํ วาติ.\csromanspacer{} ทุกฺขํ ภนฺเต.\csromanspacer{} ยํ ปนานิจฺจํ ทุกฺขํ วิปริณามธมฺมํ, อปิ นุ ตํ อนุปาทาย อุปฺปชฺเชยฺย อชฺฌตฺตํ สุขทุกฺขนฺติ.\csromanspacer{} โน เหตํ ภนฺเต.\csromanspacer{} เอวํ ปสฺสํ ฯเปฯ นาปรํ อิตฺถตฺตายาติ ปชานาตีติ.\csromanspacer{}\csromanspacer{}ปฐมํ.}

\csromanmatikaanchor{mtk.40}
\csromantocmarkref{subsection}{2. เอตํมมสุตฺต}{mtk.40}
\bookmark[dest={mtk.40},level=2]{2. เอตํมมสุตฺต}
\csromanheader{2}{2. เอตํมมสุตฺต}
\proseitem{151}{สาวตฺถินิทานํ.\csromanspacer{} กิสฺมิํ นุ โข ภิกฺขเว สติ กิํ อุปาทาย กิํ อภินิวิสฺส “เอตํ มม, เอโสหมสฺมิ, เอโส เม อตฺตา”ติ สมนุปสฺสตีติ.\csromanspacer{} ภควํมูลกา โน ภนฺเต ธมฺมา ฯเปฯ.\csromanspacer{} รูเป โข ภิกฺขเว สติ รูปํ อุปาทาย รูปํ อภินิวิสฺส ฯเปฯ.\csromanspacer{} วิญฺญาเณ สติ วิญฺญาณํ อุปาทาย วิญฺญาณํ อภินิวิสฺส “เอตํ มม, เอโสหมสฺมิ, เอโส เม อตฺตา”ติ สมนุปสฺสติ.\csromanspacer{} ตํ กิํ มญฺญถ ภิกฺขเว, รูปํ นิจฺจํ วา อนิจฺจํ วาติ.\csromanspacer{} อนิจฺจํ ภนฺเต ฯเปฯ วิปริณามธมฺมํ, อปิ นุ ตํ อนุปาทาย “เอตํ มม, เอโสหมสฺมิ, เอโส เม อตฺตา”ติ สมนุปสฺเสยฺยาติ.\csromanspacer{} โน เหตํ ภนฺเต.\csromanspacer{} เวทนา.\csromanspacer{} สญฺญา.\csromanspacer{} สงฺขารา.\csromanspacer{} วิญฺญาณํ นิจฺจํ วา อนิจฺจํ วาติ.\csromanspacer{} อนิจฺจํ ภนฺเต ฯเปฯ วิปริณามธมฺมํ, อปิ นุ ตํ อนุปาทาย “เอตํ มม, เอโสหมสฺมิ, เอโส เม อตฺตา”ติ สมนุปสฺเสยฺยาติ.\csromanspacer{} โน เหตํ ภนฺเต.\csromanspacer{} เอวํ ปสฺสํ ฯเปฯ นาปรํ อิตฺถตฺตายาติ ปชานาตีติ.\csromanspacer{}\csromanspacer{}ทุติยํ.}

\csromanmatikaanchor{mtk.41}
\csromantocmarkref{subsection}{3. โส-อตฺตาสุตฺต}{mtk.41}
\bookmark[dest={mtk.41},level=2]{3. โส-อตฺตาสุตฺต}
\csromanheader{2}{3. โส-อตฺตาสุตฺต}
\proseitem{152}{สาวตฺถินิทานํ.\csromanspacer{} กิสฺมิํ นุ โข ภิกฺขเว สติ กิํ อุปาทาย กิํ อภินิวิสฺส เอวํ ทิฏฺฐิ อุปฺปชฺชติ “โส อตฺตา โส โลโก โส เปจฺจ ภวิสฺสามิ นิจฺโจ ธุโว สสฺสโต อวิปริณามธมฺโม”ติ.\csromanspacer{} ภควํมูลกา โน ภนฺเต ธมฺมา ฯเปฯ.\csromanspacer{} รูเป โข ภิกฺขเว สติ รูปํ อุปาทาย รูปํ อภินิวิสฺส เอวํ ทิฏฺฐิ อุปฺปชฺชติ “โส อตฺตา โส โลโก โส เปจฺจ ภวิสฺสามิ นิจฺโจ ธุโว สสฺสโต อวิปริณามธมฺโม”ติ.\csromanspacer{} เวทนาย ฯเปฯ.\csromanspacer{} สญฺญาย. \csromanfolio{150}สงฺขาเรสุ.\csromanspacer{} วิญฺญาเณ สติ วิญฺญาณํ อุปาทาย วิญฺญาณํ อภินิวิสฺส เอวํ ทิฏฺฐิ อุปฺปชฺชติ “โส อตฺตา โส โลโก โส เปจฺจ ภวิสฺสามิ นิจฺโจ ธุโว สสฺสโต อวิปริณามธมฺโม”ติ.}

\prose{ตํ กิํ มญฺญถ ภิกฺขเว, รูปํ นิจฺจํ วา อนิจฺจํ วาติ.\csromanspacer{} อนิจฺจํ ภนฺเต.\csromanspacer{} ยํ ปนานิจฺจํ, ทุกฺขํ วา ตํ สุขํ วาติ.\csromanspacer{} ทุกฺขํ ภนฺเต.\csromanspacer{} ยํ ปนานิจฺจํ ทุกฺขํ วิปริณามธมฺมํ, อปิ นุ ตํ อนุปาทาย เอวํ ทิฏฺฐิ อุปฺปชฺเชยฺย “โส อตฺตา โส โลโก โส เปจฺจ ภวิสฺสามิ นิจฺโจ ธุโว สสฺสโต อวิปริณามธมฺโม”ติ.\csromanspacer{} โน เหตํ ภนฺเต.\csromanspacer{} เวทนา.\csromanspacer{} สญฺญา.\csromanspacer{} สงฺขารา.\csromanspacer{} วิญฺญาณํ นิจฺจํ วา อนิจฺจํ วาติ.\csromanspacer{} อนิจฺจํ ภนฺเต.\csromanspacer{} ยํ ปนานิจฺจํ ทุกฺขํ วา ตํ สุขํ วาติ.\csromanspacer{} ทุกฺขํ ภนฺเต.\csromanspacer{} ยํ ปนานิจฺจํ ทุกฺขํ วิปริณามธมฺมํ, อปิ นุ ตํ อนุปาทาย เอวํ ทิฏฺฐิ อุปฺปชฺเชยฺย “โส อตฺตา โส โลโก โส เปจฺจ ภวิสฺสามิ นิจฺโจ ธุโว สสฺสโต อวิปริณามธมฺโม”ติ.\csromanspacer{} โน เหตํ ภนฺเต.\csromanspacer{} เอวํ ปสฺสํ ฯเปฯ นาปรํ อิตฺถตฺตายาติ ปชานาตีติ.\csromanspacer{}\csromanspacer{}ตติยํ.}

\csromanmatikaanchor{mtk.42}
\csromantocmarkref{subsection}{4. โนจเมสิยาสุตฺต}{mtk.42}
\bookmark[dest={mtk.42},level=2]{4. โนจเมสิยาสุตฺต}
\csromanheader{2}{4. โนจเมสิยาสุตฺต}
\proseitem{153}{สาวตฺถินิทานํ.\csromanspacer{} กิสฺมิํ นุ โข ภิกฺขเว สติ กิํ อุปาทาย กิํ อภินิวิสฺส เอวํ ทิฏฺฐิ อุปฺปชฺชติ “โน จสฺสํ, โน จ เม สิยา, นาภวิสฺส, น เม ภวิสฺสตี”ติ.\csromanspacer{} ภควํมูลกา โน ภนฺเต ธมฺมา ฯเปฯ.\csromanspacer{} รูเป โข ภิกฺขเว สติ รูปํ อุปาทาย รูปํ อภินิวิสฺส เอวํ ทิฏฺฐิ อุปฺปชฺชติ “โน จสฺสํ, โน จ เม สิยา, นาภวิสฺส, น เม ภวิสฺสตี”ติ.\csromanspacer{} เวทนาย สติ.\csromanspacer{} สญฺญาย สติ.\csromanspacer{} สงฺขาเรสุ สติ.\csromanspacer{} วิญฺญาเณ สติ วิญฺญาณํ อุปาทาย วิญฺญาณํ อภินิวิสฺส เอวํ ทิฏฺฐิ อุปฺปชฺชติ “โน จสฺสํ, โน จ เม สิยา, นาภวิสฺส, น เม ภวิสฺสตี”ติ.\csromanspacer{} ตํ กิํ มญฺญถ ภิกฺขเว, รูปํ นิจฺจํ วา อนิจฺจํ วาติ.\csromanspacer{} อนิจฺจํ ภนฺเต.\csromanspacer{} ยํ ปนานิจฺจํ, ทุกฺขํ วาตํ สุขํ วาติ.\csromanspacer{} ทุกฺขํ ภนฺเต.\csromanspacer{} ยํ ปนานิจฺจํ ทุกฺขํ วิปริณามธมฺมํ, อปิ นุ ตํ อนุปาทาย เอวํ ทิฏฺฐิ อุปฺปชฺเชยฺย “โน จสฺสํ, โน จ เม สิยา, นาภวิสฺส น เม ภวิสฺสตี”ติ.\csromanspacer{} โน เหตํ ภนฺเต.\csromanspacer{} เวทนา.\csromanspacer{} สญฺญา.\csromanspacer{} สงฺขารา.\csromanspacer{} วิญฺญาณํ นิจฺจํ วา อนิจฺจํ วาติ.\csromanspacer{} อนิจฺจํ ภนฺเต.\csromanspacer{} ยํ ปนานิจฺจํ, ทุกฺขํ วา ตํ สุขํ วาติ.\csromanspacer{} ทุกฺขํ ภนฺเต.\csromanspacer{} ยํ ปนานิจฺจํ ทุกฺขํ วิปริณามธมฺมํ, อปิ นุ ตํ อนุปาทาย เอวํ ทิฏฺฐิ อุปฺปชฺเชยฺย “โน จสฺสํ, โน จ เม สิยา, นาภวิสฺส, น เม ภวิสฺสตี”ติ.\csromanspacer{} โน เหตํ ภนฺเต.\csromanspacer{} เอวํ ปสฺสํ ฯเปฯ นาปรํ อิตฺถตฺตายาติ ปชานาตีติ.\csromanspacer{}\csromanspacer{}จตุตฺถํ.}

\csromanmatikaanchor{mtk.43}
\csromantocmarkref{subsection}{5. มิจฺฉาทิฏฺฐิสุตฺต}{mtk.43}
\bookmark[dest={mtk.43},level=2]{5. มิจฺฉาทิฏฺฐิสุตฺต}
\csromanheader{2}{5. มิจฺฉาทิฏฺฐิสุตฺต}
\proseitem{154}{\csromanfolio{151}สาวตฺถินิทานํ.\csromanspacer{} กิสฺมิํ นุ โข ภิกฺขเว สติ กิํ อุปาทาย กิํ อภินิวิสฺส มิจฺฉาทิฏฺฐิ อุปฺปชฺชตีติ.\csromanspacer{} ภควํมูลกา โน ภนฺเต ธมฺมา ฯเปฯ.\csromanspacer{} รูเป โข ภิกฺขเว สติ รูปํ อุปาทาย รูปํ อภินิวิสฺส มิจฺฉาทิฏฺฐิ อุปฺปชฺชติ.\csromanspacer{} เวทนาย สติ.\csromanspacer{} สญฺญาย สติ.\csromanspacer{} สงฺขาเรสุ สติ.\csromanspacer{} วิญฺญาเณ สติ วิญฺญาณํ อุปาทาย วิญฺญาณํ อภินิวิสฺส มิจฺฉาทิฏฺฐิ อุปฺปชฺชติ.\csromanspacer{} ตํ กิํ มญฺญถ ภิกฺขเว, รูปํ นิจฺจํ วา อนิจฺจํ วาติ.\csromanspacer{} อนิจฺจํ ภนฺเต.\csromanspacer{} ยํ ปนานิจฺจํ ฯเปฯ อปิ นุตํ อนุปาทาย มิจฺฉาทิฏฺฐิ อุปฺปชฺเชยฺยาติ.\csromanspacer{} โน เหตํ ภนฺเต.\csromanspacer{} เวทนา.\csromanspacer{} สญฺญา.\csromanspacer{} สงฺขารา.\csromanspacer{} วิญฺญาณํ นิจฺจํ วา อนิจฺจํ วาติ.\csromanspacer{} อนิจฺจํ ภนฺเต.\csromanspacer{} ยํ ปนานิจฺจํ, ทุกฺขํ วา ตํ สุขํ วาติ.\csromanspacer{} ทุกฺขํ ภนฺเต.\csromanspacer{} ยํ ปนานิจฺจํ ทุกฺขํ วิปริณามธมฺมํ, อปิ นุ ตํ อนุปาทาย มิจฺฉาทิฏฺฐิ อุปฺปชฺเชยฺยาติ.\csromanspacer{} โน เหตํ ภนฺเต.\csromanspacer{} เอวํ ปสฺสํ ฯเปฯ นาปรํ อิตฺถตฺตายาติ ปชานาตีติ.\csromanspacer{}\csromanspacer{}ปญฺจมํ.}

\csromanmatikaanchor{mtk.44}
\csromantocmarkref{subsection}{6. สกฺกายทิฏฺฐิสุตฺต}{mtk.44}
\bookmark[dest={mtk.44},level=2]{6. สกฺกายทิฏฺฐิสุตฺต}
\csromanheader{2}{6. สกฺกายทิฏฺฐิสุตฺต}
\proseitem{155}{สาวตฺถินิทานํ.\csromanspacer{} กิสฺมิํ นุ โข ภิกฺขเว สติ กิํ อุปาทาย กิํ อภินิวิสฺส สกฺกายทิฏฺฐิ อุปฺปชฺชตีติ.\csromanspacer{} ภควํมูลกา โน ภนฺเต ธมฺมา ฯเปฯ.\csromanspacer{} รูเป โข ภิกฺขเว สติ รูปํ อุปาทาย รูปํ อภินิวิสฺส สกฺกายทิฏฺฐิ อุปฺปชฺชติ.\csromanspacer{} เวทนาย สติ.\csromanspacer{} สญฺญาย สติ.\csromanspacer{} สงฺขาเรสุ สติ.\csromanspacer{} วิญฺญาเณ สติ.\csromanspacer{} วิญฺญาณํ อุปาทาย วิญฺญาณํ อภินิวิสฺส สกฺกายทิฏฺฐิ อุปฺปชฺชติ.\csromanspacer{} ตํ กิํ มญฺญถ ภิกฺขเว, รูปํ นิจฺจํ วา อนิจฺจํ วาติ.\csromanspacer{} อนิจฺจํ ภนฺเต.\csromanspacer{} ยํ ปนานิจฺจํ ฯเปฯ อปิ นุ ตํ อนุปาทาย สกฺกายทิฏฺฐิ อุปฺปชฺเชยฺยาติ.\csromanspacer{} โน เหตํ ภนฺเต.\csromanspacer{} เวทนา.\csromanspacer{} สญฺญา.\csromanspacer{} สงฺขารา.\csromanspacer{} วิญฺญาณํ นิจฺจํ วา อนิจฺจํ วาติ.\csromanspacer{} อนิจฺจํ ภนฺเต.\csromanspacer{} ยํ ปนานิจฺจํ ฯเปฯ อปิ นุ ตํ อนุปาทาย สกฺกายทิฏฺฐิ อุปฺปชฺเชยฺยาติ.\csromanspacer{} โน เหตํ ภนฺเต.\csromanspacer{} เอวํ ปสฺสํ ฯเปฯ นาปรํ อิตฺถตฺตายาติ ปชานาตีติ.\csromanspacer{}\csromanspacer{}ฉฏฺฐํ.}

\csromanmatikaanchor{mtk.45}
\csromantocmarkref{subsection}{7. อตฺตานุทิฏฺฐิสุตฺต}{mtk.45}
\bookmark[dest={mtk.45},level=2]{7. อตฺตานุทิฏฺฐิสุตฺต}
\csromanheader{2}{7. อตฺตานุทิฏฺฐิสุตฺต}
\proseitem{156}{สาวตฺถินิทานํ.\csromanspacer{} กิสฺมิํ นุ โข ภิกฺขเว สติ กิํ อุปาทาย กิํ อภินิวิสฺส อตฺตานุทิฏฺฐิ อุปฺปชฺชตีติ.\csromanspacer{} ภควํมูลกา โน ภนฺเต ธมฺมา ฯเปฯ.\csromanspacer{} รูเป โข ภิกฺขเว สติ.\csromanspacer{} รูปํ อุปาทาย รูปํ อภินิวิสฺส อตฺตานุทิฏฺฐิ อุปฺปชฺชติ.\csromanspacer{} เวทนาย สติ.\csromanspacer{} สญฺญาย สติ.\csromanspacer{} สงฺขาเรสุ สติ.\csromanspacer{} วิญฺญาเณ สติ วิญฺญาณํ อุปาทาย วิญฺญาณํ อภินิวิสฺส อตฺตานุทิฏฺฐิ \csromanfolio{152}อุปฺปชฺชติ.\csromanspacer{} ตํ กิํ มญฺญถ ภิกฺขเว, รูปํ นิจฺจํ วา อนิจฺจํ วาติ.\csromanspacer{} อนิจฺจํ ภนฺเต.\csromanspacer{} ยํ ปนานิจฺจํ ฯเปฯ อปิ นุ ตํ อนุปาทาย อตฺตานุทิฏฺฐิ อุปฺปชฺเชยฺยาติ.\csromanspacer{} โน เหตํ ภนฺเต.\csromanspacer{} เวทนา.\csromanspacer{} สญฺญา.\csromanspacer{} สงฺขารา.\csromanspacer{} วิญฺญาณํ นิจฺจํ วา อนิจฺจํ วาติ.\csromanspacer{} อนิจฺจํ ภนฺเต.\csromanspacer{} ยํ ปนานิจฺจํ ฯเปฯ อปิ นุ ตํ อนุปาทาย อตฺตานุทิฏฺฐิ อุปฺปชฺเชยฺยาติ.\csromanspacer{} โน เหตํ ภนฺเต.\csromanspacer{} เอวํ ปสฺสํ ฯเปฯ นาปรํ อิตฺถตฺตายาติ ปชานาตีติ.\csromanspacer{}\csromanspacer{}สตฺตมํ.}

\csromanmatikaanchor{mtk.46}
\csromantocmarkref{subsection}{8-9. อภินิเวสสุตฺต}{mtk.46}
\bookmark[dest={mtk.46},level=2]{8-9. อภินิเวสสุตฺต}
\csromanheader{2}{8. อภินิเวสสุตฺต}
\proseitem{157}{สาวตฺถินิทานํ.\csromanspacer{} กิสฺมิํ นุ โข ภิกฺขเว สติ กิํ อุปาทาย กิํ อภินิวิสฺส อุปฺปชฺชนฺติ สํโยชนาภินิเวสวินิพนฺธาติ.\csromanspacer{} ภควํมูลกา โน ภนฺเต ธมฺมา ฯเปฯ.\csromanspacer{} รูเป โข ภิกฺขเว สติ รูปํ อุปาทาย รูปํ อภินิวิสฺส อุปฺปชฺชนฺติ สํโยชนาภินิเวสวินิพนฺธา.\csromanspacer{} เวทนาย สติ.\csromanspacer{} สญฺญาย สติ.\csromanspacer{} สงฺขาเรสุ สติ.\csromanspacer{} วิญฺญาเณ สติ วิญฺญาณํ อุปาทาย วิญฺญาณํ อภินิวิสฺส อุปฺปชฺชนฺติ สํโยชนาภินิเวสวินิพนฺธา.\csromanspacer{} ตํ กิํ มญฺญถ ภิกฺขเว, รูปํ นิจฺจํ วา อนิจฺจํ วาติ.\csromanspacer{} อนิจฺจํ ภนฺเต.\csromanspacer{} ยํ ปนานิจฺจํ ฯเปฯ อปิ นุ ตํ อนุปาทาย อุปฺปชฺเชยฺยุํ สํโยชนาภินิเวสวินิพนฺธาติ.\csromanspacer{} โน เหตํ ภนฺเต.\csromanspacer{} เอวํ ปสฺสํ ฯเปฯ นาปรํ อิตฺถตฺตายาติ ปชานาตีติ.\csromanspacer{}\csromanspacer{}อฏฺฐมํ.}

\csromanheader{2}{9. ทุติย-อภินิเวสสุตฺต}
\proseitem{158}{สาวตฺถินิทานํ.\csromanspacer{} กิสฺมิํ นุ โข ภิกฺขเว สติ กิํ อุปาทาย กิํ อภินิวิสฺส อุปฺปชฺชนฺติ สํโยชนาภินิเวสวินิพนฺธาชฺโฌสานาติ\footnote{วินิพนฺธา อชฺโฌสานาติ (สี, ก)}.\csromanspacer{} ภควํมูลกา โน ภนฺเต ธมฺมา ฯเปฯ.\csromanspacer{} รูเป โข ภิกฺขเว สติ รูปํ อุปาทาย รูปํ อภินิวิสฺส อุปฺปชฺชนฺติ สํโยชนาภินิเวสวินิพนฺธาชฺโฌสานา.\csromanspacer{} เวทนาย สติ.\csromanspacer{} สญฺญาย สติ.\csromanspacer{} สงฺขาเรสุ สติ.\csromanspacer{} วิญฺญาเณ สติ วิญฺญาณํ อุปาทาย วิญฺญาณํ อภินิวิสฺส อุปฺปชฺชนฺติ สํโยชนาภินิเวสวินิพนฺธาชฺโฌสานา.\csromanspacer{} ตํ กิํ มญฺญถ ภิกฺขเว, รูปํ นิจฺจํ วา อนิจฺจํ วาติ.\csromanspacer{} อนิจฺจํ ภนฺเต.\csromanspacer{} ยํ ปนานิจฺจํ ฯเปฯ อปิ นุ ตํ อนุปาทาย อุปฺปชฺเชยฺยุํ สํโยชนาภินิเวสวินิพนฺธาชฺโฌสานาติ.\csromanspacer{} โน เหตํ ภนฺเต.\csromanspacer{} เอวํ ปสฺสํ ฯเปฯ นาปรํ อิตฺถตฺตายาติ ปชานาตีติ.\csromanspacer{}\csromanspacer{}นวมํ.}

\csromanmatikaanchor{mtk.47}
\csromantocmarkref{subsection}{10. อานนฺทสุตฺต}{mtk.47}
\bookmark[dest={mtk.47},level=2]{10. อานนฺทสุตฺต}
\csromanheader{2}{10. อานนฺทสุตฺต}
\proseitem{159}{\csromanfolio{153}สาวตฺถินิทานํ.\csromanspacer{} อถ โข อายสฺมา อานนฺโท เยน ภควา เตนุปสงฺกมิ, อุปสงฺกมิตฺวา ฯเปฯ ภควนฺตํ เอตทโวจ “สาธุ เม ภนฺเต ภควา สํขิตฺเตน ธมฺมํ เทเสตุ, ยมหํ ภควโต ธมฺมํ สุตฺวา เอโก วูปกฏฺโฐ อปฺปมตฺโต อาตาปี ปหิตตฺโต วิหเรยฺยนฺ”ติ.}

\prose{ตํ กิํ มญฺญสิ อานนฺท, รูปํ นิจฺจํ วา อนิจฺจํ วาติ.\csromanspacer{} อนิจฺจํ ภนฺเต.\csromanspacer{} ยํ ปนานิจฺจํ, ทุกฺขํ วา ตํ สุขํ วาติ.\csromanspacer{} ทุกฺขํ ภนฺเต.\csromanspacer{} ยํ ปนานิจฺจํ ทุกฺขํ วิปริณามธมฺมํ, กลฺลํ นุ ตํ สมนุปสฺสิตุํ “เอตํ มม, เอโสหมสฺมิ, เอโส เม อตฺตา”ติ.\csromanspacer{} โน เหตํ ภนฺเต.\csromanspacer{} เวทนา.\csromanspacer{} สญฺญา.\csromanspacer{} สงฺขารา.\csromanspacer{} วิญฺญาณํ นิจฺจํ วา อนิจฺจํ วาติ.\csromanspacer{} อนิจฺจํ ภนฺเต.\csromanspacer{} ยํ ปนานิจฺจํ, ทุกฺขํ วา ตํ สุขํ วาติ.\csromanspacer{} ทุกฺขํ ภนฺเต.\csromanspacer{} ยํ ปนานิจฺจํ ทุกฺขํ วิปริณามธมฺมํ, กลฺลํ นุ ตํ สมนุปสฺสิตุํ “เอตํ มม, เอโสหมสฺมิ, เอโส เม อตฺตา”ติ.\csromanspacer{} โน เหตํ ภนฺเต\footnote{โน เหตํ ภนฺเต. ตสฺมาติหานนฺท ยํ กิญฺจิ รูปํ อตีตานาคตปจฺจุปฺปนฺนํ ฯเปฯ ทฏฺฐพฺพํ. (สี, สฺยา, กํ, อิ)}.\csromanspacer{} เอวํ ปสฺสํ ฯเปฯ นาปรํ อิตฺถตฺตายาติ ปชานาตีติ.\csromanspacer{}\csromanspacer{}ทสมํ.\csromanspacer{} ทิฏฺฐิวคฺโค ปญฺจโม.}

\titlehead{\textbf{ตสฺสุทฺทานํ}}
\csromangathameasure{อชฺฌตฺติกํ เอตํมม, โส-อตฺตา โนจเมสิยา. \\ มิจฺฉาสกฺกายตฺตานุ เทฺว, อภินิเวสา อานนฺเทนาติ. อุปริปณฺณาสโก สมตฺโต.}
\csromangathagroup{\csromangathastanza{อชฺฌตฺติกํ เอตํมม, โส-อตฺตา โนจเมสิยา. \\ มิจฺฉาสกฺกายตฺตานุ เทฺว, อภินิเวสา อานนฺเทนาติ.\csromanspacer{} อุปริปณฺณาสโก สมตฺโต.}}

\csromanheader{2}{\textbf{ตสฺส อุปริปณฺณาสกสฺส วคฺคุทฺทานํ}}
\csromangathameasure{อนฺโต ธมฺมกถิกาวิชฺชา, กุกฺกุฬํ ทิฏฺฐิปญฺจมํ. \\ ตติโย ปณฺณาสโก วุตฺโต, นิปาโตติ ปวุจฺจตีติ. \\ \textbf{ขนฺธสํยุตฺตํ สมตฺตํ.}}
\csromangathagroup{\csromangathastanza{อนฺโต ธมฺมกถิกาวิชฺชา, กุกฺกุฬํ ทิฏฺฐิปญฺจมํ. \\ ตติโย ปณฺณาสโก วุตฺโต, นิปาโตติ ปวุจฺจตีติ\footnote{นิปาโต เตน วุจฺจตีติ (สี, สฺยา, กํ)}. \\ \textbf{ขนฺธสํยุตฺตํ สมตฺตํ.}}}

\csromanmatikaanchor{mtk.48}
\csromantocmarknopage{chapter}{2. ราธสํยุตฺต}{mtk.48}
\bookmark[dest={mtk.48},level=0]{2. ราธสํยุตฺต}
\chapterheadpageread{2. ราธสํยุตฺต}
\chapterhead{1. ปฏฺฐมวคฺค}
\csromanheader{2}{1. มารสุตฺต}
\csromanmatikaanchor{mtk.49}
\csromantocmarkref{section}{1. ปฐมวคฺค}{mtk.49}
\bookmark[dest={mtk.49},level=1]{1. ปฐมวคฺค}
\proseitem{160}{\csromanfolio{154}สาวตฺถินิทานํ.\csromanspacer{} อถ โข อายสฺมา ราโธ เยน ภควา เตนุปสงฺกมิ, อุปสงฺกมิตฺวา ภควนฺตํ อภิวาเทตฺวา เอกมนฺตํ นิสีทิ, เอกมนฺตํ นิสินฺโน โข อายสฺมา ราโธ ภควนฺตํ เอตทโวจ.}

\prose{“มาโร มาโร”ติ ภนฺเต วุจฺจติ, กิตฺตาวตา นุ โข ภนฺเต มาโรติ.\csromanspacer{} รูเป โข ราธ สติ มาโร วา อสฺส มาเรตา วา, โย วา ปน มียติ, ตสฺมาติห ตฺวํ ราธ รูปํ “มาโร”ติ ปสฺส, “มาเรตา”ติ ปสฺส, “มียตี”ติ ปสฺส, “โรโค”ติ ปสฺส, “คณฺโฑ”ติ ปสฺส, “สลฺลนฺ”ติ ปสฺส, “อฆนฺ”ติ ปสฺส, “อฆภูตนฺ”ติ ปสฺส.\csromanspacer{} เย นํ เอวํ ปสฺสนฺติ, เต สมฺมา ปสฺสนฺติ.\csromanspacer{} เวทนาย สติ ฯเปฯ.\csromanspacer{} สญฺญาย สติ.\csromanspacer{} สงฺขาเรสุ สติ.\csromanspacer{} วิญฺญาเณ สติ มาโร วา อสฺส มาเรตา วา, โย วา ปน มียติ, ตสฺมาติห ตฺวํ ราธ วิญฺญาณํ “มาโร”ติ ปสฺส, “มาเรตา”ติ ปสฺส, “มียตี”ติ ปสฺส, “โรโค”ติ ปสฺส, “คณฺโฑ”ติ ปสฺส, “สลฺลนฺ”ติ ปสฺส, “อฆนฺ”ติ ปสฺส, “อฆภูตนฺ”ติ ปสฺส.\csromanspacer{} เย นํ เอวํ ปสฺสนฺติ, เต สมฺมา ปสฺสนฺตีติ.}

\prose{สมฺมาทสฺสนํ ปน ภนฺเต กิมตฺถิยนฺติ.\csromanspacer{} สมฺมา ทสฺสนํ โข ราธ นิพฺพิทตฺถํ.\csromanspacer{} นิพฺพิทา ปน ภนฺเต กิมตฺถิยาติ.\csromanspacer{} นิพฺพิทา โข ราธ วิราคตฺถา.\csromanspacer{} วิราโค ปน ภนฺเต กิมตฺถิโยติ.\csromanspacer{} วิราโค โข ราธ วิมุตฺตตฺโถ.\csromanspacer{} วิมุตฺติ ปน ภนฺเต กิมตฺถิยาติ.\csromanspacer{} วิมุตฺติ โข ราธ นิพฺพานตฺถา.\csromanspacer{} นิพฺพานํ ปน ภนฺเต กิมตฺถิยนฺติ.\csromanspacer{} อจฺจยาสิ\footnote{อจฺจสรา (สี, สฺยา, กํ), อสฺส (อิ), อจฺจยา (ก)} ราธ ปญฺหํ, นาสกฺขิ ปญฺหสฺส ปริยนฺตํ คเหตุํ, นิพฺพาโนคธํ หิ ราธ พฺรหฺมจริยํ วุสฺสติ นิพฺพานปรายนํ นิพฺพานปริโยสานนฺติ.\csromanspacer{}\csromanspacer{}ปฐมํ.}

\csromanheader{2}{2. สตฺตสุตฺต}
\proseitem{161}{สาวตฺถินิทานํ.\csromanspacer{} เอกมนฺตํ นิสินฺโน โข อายสฺมา ราโธ ภควนฺตํ เอตทโวจ “สตฺโต สตฺโต”ติ ภนฺเต วุจฺจติ, กิตฺตาวตา นุ โข}

\csromanheader{2}{2. ราธสํยุตฺต}
\prose{\csromanfolio{155}ภนฺเต “สตฺโต”ติ วุจฺจตีติ.\csromanspacer{} รูเป โข ราธ โย ฉนฺโท โย ราโค ยา นนฺที ยา ตณฺหา, ตตฺร สตฺโต ตตฺร วิสตฺโต, ตสฺมา “สตฺโต”ติ วุจฺจติ.\csromanspacer{} เวทนาย.\csromanspacer{} สญฺญาย.\csromanspacer{} สงฺขาเรสุ.\csromanspacer{} วิญฺญาเณ โย ฉนฺโท โย ราโค ยา นนฺที ยา ตณฺหา, ตตฺร สตฺโต ตตฺร วิสตฺโต, ตสฺมา “สตฺโต”ติ วุจฺจติ.}

\prose{เสยฺยถาปิ ราธ กุมารกา วา กุมาริกาโย วา ปํสฺวาคารเกหิ กีฬนฺติ.\csromanspacer{} ยาวกีวญฺจ เตสุ ปํสฺวาคารเกสุ อวิคตราคา โหนฺติ อวิคตจฺฉนฺทา อวิคตเปมา อวิคตปิปาสา อวิคตปริฬาหา อวิคตตณฺหา.\csromanspacer{} ตาว ตานิ ปํสฺวาคารกานิ อลฺลียนฺติ เกฬายนฺติ ธนายนฺติ\footnote{มนายนฺติ (สี, อิ, ก)} มมายนฺติ.\csromanspacer{} ยโต จ โข ราธ กุมารกา วา กุมาริกาโย วา เตสุ ปํสฺวาคารเกสุ วิคตราคา โหนฺติ วิคตจฺฉนฺทา วิคตเปมา วิคตปิปาสา วิคตปริฬาหา วิคตตณฺหา.\csromanspacer{} อถ โข ตานิ ปํสฺวาคารกานิ หตฺเถหิ จ ปาเทหิ จ วิกิรนฺติ วิธมนฺติ วิทฺธํเสนฺติ วิกีฬนิยํ\footnote{วิกีฬนิกํ (สี, สฺยา, กํ, อิ)} กโรนฺติ.\csromanspacer{} เอวเมว โข ราธ ตุเมฺหปิ รูปํ วิกิรถ วิธมถ วิทฺธํเสถ วิกีฬนิยํ กโรถ, ตณฺหากฺขยาย ปฏิปชฺชถ.\csromanspacer{} เวทนํ วิกิรถ วิธมถ วิทฺธํเสถ วิกีฬนิยํ กโรถ, ตณฺหากฺขยาย ปฏิปชฺชถ.\csromanspacer{} สญฺญํ.\csromanspacer{} สงฺขาเร วิกิรถ วิธมถ วิทฺธํเสถ วิกีฬนิยํ กโรถ, ตณฺหากฺขยาย ปฏิปชฺชถ.\csromanspacer{} วิญฺญาณํ วิกิรถ วิธมถ วิทฺธํเสถ วิกีฬนิยํ กโรถ, ตณฺหากฺขยาย ปฏิปชฺชถ, ตณฺหากฺขโย หิ ราธ นิพฺพานนฺติ.\csromanspacer{}\csromanspacer{}ทุติยํ.}

\csromanheader{2}{3. ภวเนตฺติสุตฺต}
\proseitem{162}{สาวตฺถินิทานํ.\csromanspacer{} เอกมนฺตํ นิสินฺโน โข อายสฺมา ราโธ ภควนฺตํ เอตทโวจ “ภวเนตฺตินิโรโธ\footnote{ภวเนตฺติ (สี, สฺยา, กํ, อิ)} ภวเนตฺตินิโรโธ”ติ\footnote{ภวเนตฺตีติ (สี, สฺยา, กํ)} ภนฺเต วุจฺจติ, กตมา นุ โข ภนฺเต ภวเนตฺติ, กตโม ภวเนตฺตินิโรโธติ.\csromanspacer{} รูเป โข ราธ โย ฉนฺโท โย ราโค ยา นนฺที ยา ตณฺหา เย อุปยุปาทานา เจตโส อธิฏฺฐานาภินิเวสานุสยา, อยํ วุจฺจติ ภวเนตฺติ, เตสํ นิโรโธ\footnote{นิโรธา (สี, สฺยา, กํ, อิ)} ภวเนตฺตินิโรโธ.\csromanspacer{} เวทนาย.\csromanspacer{} สญฺญาย. \csromanfolio{156}สงฺขาเรสุ.\csromanspacer{} วิญฺญาเณ โย ฉนฺโท ฯเปฯ อธิฏฺฐานาภินิเวสานุสยา.\csromanspacer{} อยํ วุจฺจติ ภวเนตฺติ, เตสํ นิโรโธ ภวเนตฺตินิโรโธติ.\csromanspacer{}\csromanspacer{}ตติยํ.}

\csromanheader{2}{4. ปริญฺเญยฺยสุตฺต}
\proseitem{163}{สาวตฺถินิทานํ.\csromanspacer{} อายสฺมา ราโธ เยน ภควา เตนุปสงฺกมิ, อุปสงฺกมิตฺวา ภควนฺตํ อภิวาเทตฺวา เอกมนฺตํ นิสีทิ, เอกมนฺตํ นิสินฺนํ โข อายสฺมนฺตํ ราธํ ภควา เอตทโวจ\csromandash{}}

\prose{ปริญฺเญยฺเย จ ราธ ธมฺเม เทเสสฺสามิ ปริญฺญญฺจ ปริญฺญาตาวิํ ปุคฺคลญฺจ.\csromanspacer{} ตํ สุณาหิ สาธุกํ มนสิ กโรหิ, ภาสิสฺสามีติ. “เอวํ ภนฺเต”ติ โข อายสฺมา ราโธ ภควโต ปจฺจสฺโสสิ.\csromanspacer{} ภควา เอตทโวจ\csromandash{} กตเม จ ราธ ปริญฺเญยฺยา ธมฺมา.\csromanspacer{} รูปํ โข ราธ ปริญฺเญยฺโย ธมฺโม, เวทนา ปริญฺเญยฺโย ธมฺโม, สญฺญา ปริญฺเญยฺโย ธมฺโม, สงฺขารา ปริญฺเญยฺโย ธมฺโม, วิญฺญาณํ ปริญฺเญยฺโย ธมฺโม.\csromanspacer{} อิเม วุจฺจนฺติ ราธ ปริญฺเญยฺยา ธมฺมา.\csromanspacer{} กตมา จ ราธ ปริญฺญา.\csromanspacer{} โย โข ราธ ราคกฺขโย โทสกฺขโย โมหกฺขโย, อยํ วุจฺจติ ราธ ปริญฺญา.\csromanspacer{} กตโม จ ราธ ปริญฺญาตาวี ปุคฺคโล. “อรหา”ติสฺส วจนียํ.\csromanspacer{} ยฺวายํ อายสฺมา เอวํนาโม เอวํโคตฺโต, อยํ วุจฺจติ ราธ ปริญฺญาตาวี ปุคฺคโลติ.\csromanspacer{}\csromanspacer{}จตุตฺถํ.}

\csromanheader{2}{5. สมณสุตฺต}
\proseitem{164}{สาวตฺถินิทานํ.\csromanspacer{} เอกมนฺตํ นิสินฺนํ โข อายสฺมนฺตํ ราธํ ภควา เอตทโวจ\csromandash{}ปญฺจิเม ราธ อุปาทานกฺขนฺธา.\csromanspacer{} กตเม ปญฺจ, รูปุปาทานกฺขนฺโธ เวทนุปาทานกฺขนฺโธ สญฺญุปาทานกฺขนฺโธ สงฺขารุปาทานกฺขนฺโธ วิญฺญาณุปาทานกฺขนฺโธ.\csromanspacer{} เย หิ เกจิ ราธ สมณา วา พฺราหฺมณา วา อิเมสํ ปญฺจนฺนํ อุปาทานกฺขนฺธานํ อสฺสาทญฺจ อาทีนวญฺจ นิสฺสรณญฺจ ยถาภูตํ นปฺปชานนฺติ.\csromanspacer{} น เม เต ราธ สมณา วา พฺราหฺมณ วา สมเณสุ วา สมณสมฺมตา พฺราหฺมเณสุ วา พฺราหฺมณสมฺมตา, น จ ปน เต อายสฺมนฺโต สามญฺญตฺถํ วา พฺราหฺมญฺญตฺถํ วา ทิฏฺเฐว ธมฺเม สยํ อภิญฺญา สจฺฉิกตฺวา อุปสมฺปชฺช วิหรนฺติ.\csromanspacer{} เย จ โข เกจิ ราธ สมณา วา พฺราหฺมณา วา อิเมสํ ปญฺจนฺนํ อุปาทานกฺขนฺธานํ อสฺสาทญฺจ อาทีนวญฺจ นิสฺสรณญฺจ ยถาภูตํ ปชานนฺติ, เต โข เม ราธ สมณา วา พฺราหฺมณา วา สมเณสุ เจว สมณสมฺมตา พฺราหฺมเณสุ จ}

\csromanheader{2}{2. ราธสํยุตฺต}
\prose{\csromanfolio{157}พฺราหฺมณสมฺมตา, เต จ ปนายสฺมนฺโต สามญฺญตฺถญฺจ พฺรหฺมญฺญตฺถญฺจ ทิฏฺเฐว ธมฺเม สยํ อภิญฺญา สจฺฉิกตฺวา อุปสมฺปชฺช วิหรนฺตีติ.\csromanspacer{}\csromanspacer{}ปญฺจมํ.}

\csromanheader{2}{6. ทุติยสมณสุตฺต}
\proseitem{165}{สาวตฺถินิทานํ.\csromanspacer{} เอกมนฺตํ นิสินฺนํ โข อายสฺมนฺตํ ราธํ ภควา เอตทโวจ\csromandash{}ปญฺจิเม ราธ อุปาทานกฺขนฺธา.\csromanspacer{} กตเม ปญฺจ, รูปุปาทานกฺขนฺโธ ฯเปฯ วิญฺญาณุปาทานกฺขนฺโธ.\csromanspacer{} เย หิ เกจิ ราธ สมณา วา พฺราหฺมณา วา อิเมสํ ปญฺจนฺนํ อุปาทานกฺขนฺธานํ สมุทยญฺจ อตฺถงฺคมญฺจ อสฺสาทญฺจ อาทีนวญฺจ นิสฺสรณญฺจ ยถาภูตํ นปฺปชานนฺติ ฯเปฯ สยํ อภิญฺญา สจฺฉิกตฺวา อุปสมฺปชฺช วิหรนฺตีติ.\csromanspacer{}\csromanspacer{}ฉฏฺฐํ.}

\csromanheader{2}{7. โสตาปนฺนสุตฺต}
\proseitem{166}{สาวตฺถินิทานํ.\csromanspacer{} เอกมนฺตํ นิสินฺนํ โข อายสฺมนฺตํ ราธํ ภควา เอตทโวจ\csromandash{}ปญฺจิเม ราธ อุปาทานกฺขนฺธา.\csromanspacer{} กตเม ปญฺจ, รูปุปาทานกฺขนฺโธ ฯเปฯ วิญฺญาณุปาทานกฺขนฺโธ.\csromanspacer{} ยโต โข ราธ อริยสาวโก อิเมสํ ปญฺจนฺนํ อุปาทานกฺขนฺธานํ สมุทยญฺจ อตฺถงฺคมญฺจ อสฺสาทญฺจ อาทีนวญฺจ นิสฺสรณญฺจ ยถาภูตํ ปชานาติ, อยํ วุจฺจติ ราธ อริยสาวโก โสตาปนฺโน อวินิปาตธมฺโม นิยโต สมฺโพธิปรายโนติ.\csromanspacer{}\csromanspacer{}สตฺตมํ.}

\csromanheader{2}{8. อรหนฺตสุตฺต}
\proseitem{167}{สาวตฺถินิทานํ.\csromanspacer{} เอกมนฺตํ นิสินฺนํ โข อายสฺมนฺตํ ราธํ ภควา เอตทโวจ\csromandash{}ปญฺจิเม ราธ อุปาทานกฺขนฺธา.\csromanspacer{} กตเม ปญฺจ, รูปุปาทานกฺขนฺโธ ฯเปฯ วิญฺญาณุปาทานกฺขนฺโธ.\csromanspacer{} ยโต โข ราธ ภิกฺขุ อิเมสํ ปญฺจนฺนํ อุปาทานกฺขนฺธานํ สมุทยญฺจ อตฺถงฺคมญฺจ อสฺสาทญฺจ อาทีนวญฺจ นิสฺสรณญฺจ ยถาภูตํ วิทิตฺวา อนุปาทาวิมุตฺโต โหติ, อยํ วุจฺจติ ภิกฺขุ อรหํ ขีณาสโว วุสิตวา กตกรณีโย โอหิตภาโร อนุปฺปตฺตสทตฺโถ ปริกฺขีณภวสํโยชโน สมฺมทญฺญาวิมุตฺโตติ.\csromanspacer{}\csromanspacer{}อฏฺฐมํ.}

\csromanheader{2}{9. ฉนฺทราคสุตฺต}
\proseitem{168}{สาวตฺถิทานํ.\csromanspacer{} เอกมนฺตํ นิสินฺนํ โข อายสฺมนฺตํ ราธํ ภควา เอตทโวจ\csromandash{}รูเป โข ราธ โย ฉนฺโท โย ราโค ยา นนฺที ยา ตณฺหา, ตํ ปชหถ, เอวํ ตํ รูปํ ปหีนํ ภวิสฺสติ อุจฺฉินฺนมูลํ ตาลาวตฺถุกตํ อนภาวํกตํ อายติํ อนุปฺปาทธมฺมํ.\csromanspacer{} เวทนาย \csromanfolio{158}โย ฉนฺโท โย ราโค ยา นนฺที ยา ตณฺหา, ตํ ปชหถ, เอวํ สา เวทนา ปหีนา ภวิสฺสติ อุจฺฉินฺนมูลา ตาลาวตฺถุกตา อนภาวํกตา อายติํ อนุปฺปาทธมฺมา.\csromanspacer{} สญฺญาย.\csromanspacer{} สงฺขาเรสุ โย ฉนฺโท โย ราโค ยา นนฺที ยา ตณฺหา, ตํ ปชหถ, เอวํ เต สงฺขารา ปหีนา ภวิสฺสนฺติ อุจฺฉินฺนมูลา ตาลาวตฺถุกตา อนภาวํกตา อายติํ อนุปฺปาทธมฺมา.\csromanspacer{} วิญฺญาเณ โย ฉนฺโท โย ราโค ยา นนฺที ยา ตณฺหา, ตํ ปชหถ, เอวํ ตํ วิญฺญาณํ ปหีนํ ภวิสฺสติ ฯเปฯ อนุปฺปาทธมฺมนฺติ.\csromanspacer{}\csromanspacer{}นวมํ.}

\csromanheader{2}{10. ทุติยฉนฺทราคสุตฺต}
\proseitem{169}{สาวตฺถินิทานํ.\csromanspacer{} เอกมนฺตํ นิสินฺนํ โข อายสฺมนฺตํ ราธํ ภควา เอตทโวจ\csromandash{}รูเป โข ราธ โย ฉนฺโท โย ราโค ยา นนฺที ยา ตณฺหา เย อุปยุปาทานา เจตโส อธิฏฺฐานาภินิเวสานุสยา, เต ปชหถ, เอวํ ตํ รูปํ ปหีนํ ภวิสฺสติ อุจฺฉินฺนมูลํ ตาลาวตฺถุกตํ อนภาวํกตํ อายติํ อนุปฺปาทธมฺมํ.\csromanspacer{} เวทนาย โย ฉนฺโท โย ราโค ยา นนฺที ยา ตณฺหา เย อุปยุปาทานา เจตโส อธิฏฺฐานาภินิเวสานุสยา, เต ปชหถ.\csromanspacer{} เอวํ สา เวทนา ปหีนา ภวิสฺสติ, อุจฺฉินฺนมูลา ตาลาวตฺถุกตา อนภาวํกตา อายติํ อนุปฺปาทธมฺมา.\csromanspacer{} สญฺญาย.\csromanspacer{} สงฺขาเรสุ โย ฉนฺโท โย ราโค ยา นนฺที ยา ตณฺหา เย อุปยุปาทานา เจตโส อธิฏฺฐานาภินิเวสานุสยา, เต ปชหถ, เอวํ เต สงฺขารา ปหีนา ภวิสฺสนฺติ อุจฺฉินฺนมูลา ตาลาวตฺถุกตา อนภาวํกตา อายติํ อนุปฺปาทธมฺมา.\csromanspacer{} วิญฺญาเณ โย ฉนฺโท โย ราโค ยา นนฺที ยา ตณฺหา เย อุปยุปาทานา เจตโส อธิฏฺฐานาภินิเวสานุสยา, เต ปชหถ, เอวํ ตํ วิญฺญาณํ ปหีนํ ภวิสฺสติ อุจฺฉินฺนมูลํ ตาลาวตฺถุกตํ อนภาวํกตํ อายติํ อนุปฺปาทธมฺมนฺติ.\csromanspacer{}\csromanspacer{}ทสมํ.\csromanspacer{} ราธสํยุตฺตสฺส ปฐโม วคฺโค.}

\titlehead{\textbf{ตสฺสุทฺทานํ}}
\csromangathameasure{มาโร สตฺโต ภวเนตฺติ, ปริญฺเญยฺยา สมณา ทุเว. \\ โสตาปนฺโน อรหา จ, ฉนฺทราคาปเร ทุเวติ.}
\csromangathagroup{\csromangathastanza{มาโร สตฺโต ภวเนตฺติ, ปริญฺเญยฺยา สมณา ทุเว. \\ โสตาปนฺโน อรหา จ, ฉนฺทราคาปเร ทุเวติ.}}

\csromanheader{2}{2. ราธสํยุตฺต}
\csromanmatikaanchor{mtk.50}
\csromantocmarkref{section}{2. ทุติยวคฺค}{mtk.50}
\bookmark[dest={mtk.50},level=1]{2. ทุติยวคฺค}
\csromanheader{1}{2. ทุติยวคฺค}
\csromanheader{2}{1. มารสุตฺต}
\proseitem{170}{\csromanfolio{159}สาวตฺถินิทานํ.\csromanspacer{} เอกมนฺตํ นิสินฺโน โข อายสฺมา ราโธ ภควนฺตํ เอตทโวจ\csromandash{}“มาโร มาโร”ติ ภนฺเต วุจฺจติ, กตโม นุ โข ภนฺเต มาโรติ.\csromanspacer{} รูปํ โข ราธ มาโร เวทนา มาโร สญฺญา มาโร สงฺขารา มาโร วิญฺญาณํ มาโร, เอวํ ปสฺสํ ราธ สุตวา อริยสาวโก รูปสฺมิมฺปิ นิพฺพินฺทติ, เวทนายปิ นิพฺพินฺทติ, สญฺญายปิ นิพฺพินฺทติ, สงฺขาเรสุปิ นิพฺพินฺทติ, วิญฺญาณสฺมิมฺปิ นิพฺพินฺทติ, นิพฺพินฺทํ วิรชฺชติ, วิราคา วิมุจฺจติ, วิมุตฺตสฺมิํ “วิมุตฺตมฺ”อิติ ญาณํ โหติ, “ขีณา ชาติ, วุสิตํ พฺราหฺมจริยํ, กตํ กรณียํ, นาปรํ อิตฺถตฺตายา”ติ ปชานาตีติ.\csromanspacer{}\csromanspacer{}ปฐมํ.}

\csromanheader{2}{2. มารธมฺมสุตฺต}
\proseitem{171}{สาวตฺถินิทานํ.\csromanspacer{} เอกมนฺตํ นิสินฺโน โข อายสฺมา ราโธ ภควนฺตํ เอตทโวจ\csromandash{}“มารธมฺโม มารธมฺโม”ติ ภนฺเต วุจฺจติ, กตโม นุ โข ภนฺเต มารธมฺโมติ.\csromanspacer{} รูปํ โข ราธ มารธมฺโม, เวทนา มารธมฺโม, สญฺญา มารธมฺโม, สงฺขารา มารธมฺโม, วิญฺญาณํ มารธมฺโม.\csromanspacer{} เอวํ ปสฺสํ ฯเปฯ นาปรํ อิตฺถตฺตายาติ ปชานาตีติ.\csromanspacer{}\csromanspacer{}ทุติยํ.}

\csromanheader{2}{3. อนิจฺจสุตฺต}
\proseitem{172}{สาวตฺถินิทานํ.\csromanspacer{} เอกมนฺตํ นิสินฺโน โข อายสฺมา ราโธ ภควนฺตํ เอตทโวจ\csromandash{}“อนิจฺจํ อนิจฺจนฺ”ติ ภนฺเต วุจฺจติ, กตมํ นุ โข ภนฺเต อนิจฺจนฺติ.\csromanspacer{} รูปํ โข ราธ อนิจฺจํ, เวทนา อนิจฺจา, สญฺญา อนิจฺจา, สงฺขารา อนิจฺจา, วิญฺญาณํ อนิจฺจํ.\csromanspacer{} เอวํ ปสฺสํ ฯเปฯ นาปรํ อิตฺถตฺตายาติ ปชานาตีติ.\csromanspacer{}\csromanspacer{}ตติยํ.}

\csromanheader{2}{4. อนิจฺจธมฺมสุตฺต}
\proseitem{173}{สาวตฺถินิทานํ.\csromanspacer{} เอกมนฺตํ นิสินฺโน โข อายสฺมา ราโธ ภควนฺตํ เอตทโวจ\csromandash{}“อนิจฺจธมฺโม อนิจฺจธมฺโม”ติ ภนฺเต วุจฺจติ, กตโม นุ โข ภนฺเต อนิจฺจธมฺโมติ.\csromanspacer{} รูปํ โข ราธ อนิจฺจธมฺโม, เวทนา อนิจฺจธมฺโม, สญฺญา อนิจฺจธมฺโม, สงฺขารา อนิจฺจธมฺโม, วิญฺญาณํ อนิจฺจธมฺโม.\csromanspacer{} เอวํ ปสฺสํ ฯเปฯ นาปรํ อิตฺถตฺตายาติ ปชานาตีติ.\csromanspacer{}\csromanspacer{}จตุตฺถํ.}

\csromanheader{2}{5. ทุกฺขสุตฺต}
\proseitem{174}{\csromanfolio{160}สาวตฺถินิทานํ.\csromanspacer{} เอกมนฺตํ นิสินฺโน โข อายสฺมา ราโธ ภควนฺตํ เอตทโวจ\csromandash{}“ทุกฺขํ ทุกฺขนฺ”ติ ภนฺเต วุจฺจติ, กตมํ นุ โข ภนฺเต ทุกฺขนฺติ.\csromanspacer{} รูปํ โข ราธ ทุกฺขํ, เวทนา ทุกฺขา, สญฺญา ทุกฺขา, สงฺขารา ทุกฺขา, วิญฺญาณํ ทุกฺขํ.\csromanspacer{} เอวํ ปสฺสํ ฯเปฯ นาปรํ อิตฺถตฺตายาติ ปชานาตีติ.\csromanspacer{}\csromanspacer{}ปญฺจมํ.}

\csromanheader{2}{6. ทุกฺขธมฺมสุตฺต}
\proseitem{175}{สาวตฺถินิทานํ.\csromanspacer{} เอกมนฺตํ นิสินฺโน โข อายสฺมา ราโธ ภควนฺตํ เอตทโวจ\csromandash{}“ทุกฺขธมฺโม ทุกฺขธมฺโม”ติ ภนฺเต วุจฺจติ, กตโม นุ โข ภนฺเต ทุกฺขธมฺโมติ.\csromanspacer{} รูปํ โข ราธ ทุกฺขธมฺโม, เวทนา ทุกฺขธมฺโม, สญฺญา ทุกฺขธมฺโม, สงฺขารา ทุกฺขธมฺโม, วิญฺญาณํ ทุกฺขธมฺโม.\csromanspacer{} เอวํ ปสฺสํ ฯเปฯ นาปรํ อิตฺถตฺตายาติ ปชานาตีติ.\csromanspacer{}\csromanspacer{}ฉฏฺฐํ.}

\csromanheader{2}{7. อนตฺตสุตฺต}
\proseitem{176}{สาวตฺถินิทานํ.\csromanspacer{} เอกมนฺตํ นิสินฺโน โข อายสฺมา ราโธ ภควนฺตํ เอตทโวจ\csromandash{}“อนตฺตา อนตฺตา”ติ ภนฺเต วุจฺจติ, กตโม นุ โข ภนฺเต อนตฺตาติ.\csromanspacer{} รูปํ โข ราธ อนตฺตา, เวทนา อนตฺตา, สญฺญา อนตฺตา, สงฺขารา อนตฺตา, วิญฺญาณํ อนตฺตา.\csromanspacer{} เอวํ ปสฺสํ ฯเปฯ นาปรํ อิตฺถตฺตายาติ ปชานาตีติ.\csromanspacer{}\csromanspacer{}สตฺตมํ.}

\csromanheader{2}{8. อนตฺตธมฺมสุตฺต}
\proseitem{177}{สาวตฺถินิทานํ.\csromanspacer{} เอกมนฺตํ นิสินฺโน โข อายสฺมา ราโธ ภควนฺตํ เอตทโวจ\csromandash{}“อนตฺตธมฺโม อนตฺตธมฺโม”ติ ภนฺเต วุจฺจติ, กตโม นุ โข ภนฺเต อนตฺตธมฺโมติ.\csromanspacer{} รูปํ โข ราธ อนตฺตธมฺโม, เวทนา อนตฺตธมฺโม, สญฺญา อนตฺตธมฺโม, สงฺขารา อนตฺตธมฺโม, วิญฺญาณํ อนตฺตธมฺโม.\csromanspacer{} เอวํ ปสฺสํ ฯเปฯ นาปรํ อิตฺถตฺตายาติ ปชานาตีติ.\csromanspacer{}\csromanspacer{}อฏฺฐมํ.}

\csromanheader{2}{9. ขยธมฺมสุตฺต}
\proseitem{178}{สาวตฺถินิทานํ.\csromanspacer{} เอกมนฺตํ นิสินฺโน โข อายสฺมา ราโธ ภควนฺตํ เอตทโวจ\csromandash{}“ขยธมฺโม ขยธมฺโม”ติ ภนฺเต วุจฺจติ, กตโม นุ โข ภนฺเต ขยธมฺโมติ.\csromanspacer{} รูปํ โข ราธ ขยธมฺโม, เวทนา ขยธมฺโม,}

\csromanheader{2}{2. ราธสํยุตฺต}
\prose{\csromanfolio{161}สญฺญา ขยธมฺโม, สงฺขารา ขยธมฺโม, วิญฺญาณํ ขยธมฺโม.\csromanspacer{} เอวํ ปสฺสํ ฯเปฯ นาปรํ อิตฺถตฺตายาติ ปชานาตีติ.\csromanspacer{}\csromanspacer{}นวมํ.}

\csromanheader{2}{10. วยธมฺมสุตฺต}
\proseitem{179}{สาวตฺถินิทานํ.\csromanspacer{} เอกมนฺตํ นิสินฺโน โข อายสฺมา ราโธ ภควนฺตํ เอตทโวจ\csromandash{}“วยธมฺโม วยธมฺโม”ติ ภนฺเต วุจฺจติ, กตโม นุ โข ภนฺเต วยธมฺโมติ.\csromanspacer{} รูปํ โข ราธ วยธมฺโม, เวทนา วยธมฺโม, สญฺญา วยธมฺโม, สงฺขารา วยธมฺโม, วิญฺญาณํ วยธมฺโม.\csromanspacer{} เอวํ ปสฺสํ ฯเปฯ นาปรํ อิตฺถตฺตายาติ ปชานาตีติ.\csromanspacer{}\csromanspacer{}ทสมํ.}

\csromanheader{2}{11. สมุทยธมฺมสุตฺต}
\proseitem{180}{สาวตฺถินิทานํ.\csromanspacer{} เอกมนฺตํ นิสินฺโน โข อายสฺมา ราโธ ภควนฺตํ เอตทโวจ\csromandash{}“สมุทยธมฺโม สมุทยธมฺโม”ติ ภนฺเต วุจฺจติ, กตโม นุ โข ภนฺเต สมุทยธมฺโมติ.\csromanspacer{} รูปํ โข ราธ สมุทยธมฺโม, เวทนา สมุทยธมฺโม, สญฺญา สมุทยธมฺโม, สงฺขารา สมุทยธมฺโม, วิญฺญาณํ สมุทยธมฺโม.\csromanspacer{} เอวํ ปสฺสํ ฯเปฯ นาปรํ อิตฺถตฺตายาติ ปชานาตีติ.\csromanspacer{}\csromanspacer{}เอกาทสมํ.}

\csromanheader{2}{12. นิโรธธมฺมสุตฺต}
\proseitem{181}{สาวตฺถินิทานํ.\csromanspacer{} เอกมนฺตํ นิสินฺโน โข อายสฺมา ราโธ ภควนฺตํ เอตทโวจ\csromandash{}“นิโรธธมฺโม นิโรธธมฺโม”ติ ภนฺเต วุจฺจติ, กตโม นุ โข ภนฺเต นิโรธธมฺโมติ.\csromanspacer{} รูปํ โข ราธ นิโรธธมฺโม, เวทนา นิโรธธมฺโม, สญฺญา นิโรธธมฺโม, สงฺขารา นิโรธธมฺโม, วิญฺญาณํ นิโรธธมฺโม.\csromanspacer{} เอวํ ปสฺสํ ฯเปฯ นาปรํ อิตฺถตฺตายาติ ปชานาตีติ.\csromanspacer{}\csromanspacer{}ทฺวาทสมํ.\csromanspacer{} ทุติโย วคฺโค.}

\titlehead{\textbf{ตสฺสุทฺทานํ}}
\csromangathameasure{มาโร จ มารธมฺโม จ, อนิจฺเจน อปเร ทุเว. \\ ทุกฺเขน จ ทุเว วุตฺตา, อนตฺเตน ตเถว จ. \\ ขยวยสมุทยํ, นิโรธธมฺเมน ทฺวาทสาติ.}
\csromangathagroup{\csromangathastanza{มาโร จ มารธมฺโม จ, อนิจฺเจน อปเร ทุเว. \\ ทุกฺเขน จ ทุเว วุตฺตา, อนตฺเตน\footnote{อนตฺเตหิ (สี, สฺยา, กํ)} ตเถว จ. \\ ขยวยสมุทยํ, นิโรธธมฺเมน ทฺวาทสาติ.}}

\csromanmatikaanchor{mtk.51}
\csromantocmarkref{section}{3. อายาจนวคฺค}{mtk.51}
\bookmark[dest={mtk.51},level=1]{3. อายาจนวคฺค}
\csromanheader{1}{3. อายาจนวคฺค}
\csromanheader{2}{1-11. มาราทิสุตฺต-เอกาทสก}
\proseitem{182}{\csromanfolio{162}สาวตฺถินิทานํ.\csromanspacer{} เอกมนฺตํ นิสินฺโน โข อายสฺมา ราโธ ภควนฺตํ เอตทโวจ\csromandash{}“สาธุ เม ภนฺเต ภควา สํขิตฺเตน ธมฺมํ เทเสตุ, ยมหํ ภควโต ธมฺมํ สุตฺวา เอโก วูปกฏฺโฐ อปฺปมตฺโต อาตาปี ปหิตตฺโต วิหเรยฺยนฺ”ติ.}

\prose{โย โข ราธ มาโร, ตตฺร เต ฉนฺโท ปหาตพฺโพ, ราโค ปหาตพฺโพ, ฉนฺทราโค ปหาตพฺโพ\footnote{สีหฬ-โปตฺถเก ปน “ตตฺร เต ฉนฺโท ปหาตพฺโพติ เอกํ สุตฺตํ, ตตฺร เต ราโค ปหาตพฺโพติ เอกํ สุตฺตํ, ตตฺร เต ฉนฺทราโค ปหาตพฺโพติ เอกํ สุตฺตนฺ”ติ เอวํ วิสุํ วิสุํ ตีณิ สุตฺตานิ วิภชิตฺวา ทสฺสิตานิ. เอวมุปริสุตฺเตสุปิ.}.\csromanspacer{} โก จ ราธ มาโร.\csromanspacer{} รูปํ โข ราธ มาโร, ตตฺร เต ฉนฺโท ปหาตพฺโพ, ราโค ปหาตพฺโพ, ฉนฺทราโค ปหาตพฺโพ1.\csromanspacer{} เวทนา มาโร, ตตฺร เต ฉนฺโท ปหาตพฺโพ ฯเปฯ.\csromanspacer{} สญฺญา มาโร, ตตฺร เต ฉนฺโท ปหาตพฺโพ ฯเปฯ.\csromanspacer{} สงฺขารา มาโร, ตตฺร เต ฉนฺโท ปหาตพฺโพ ฯเปฯ.\csromanspacer{} วิญฺญาณํ มาโร, ตตฺร เต ฉนฺโท ปหาตพฺโพ ฯเปฯ.\csromanspacer{} โย โข ราธ มาโร, ตตฺร เต ฉนฺโท ปหาตพฺโพ, ราโค ปหาตพฺโพ, ฉนฺทราโค ปหาตพฺโพติ.}

\proseitem{183}{โย โข ราธ มารธมฺโม, ตตฺร เต ฉนฺโท ปหาตพฺโพ, ราโค ปหาตพฺโพ, ฉนฺทราโค ปหาตพฺโพ ฯเปฯ.}

\proseitem{184}{ยํ โข ราธ อนิจฺจํ ฯเปฯ.}

\proseitem{185}{โย โข ราธ อนิจฺจธมฺโม ฯเปฯ.}

\proseitem{186}{ยํ โข ราธ ทุกฺขํ ฯเปฯ.}

\proseitem{187}{โย โข ราธ ทุกฺขธมฺโม ฯเปฯ.}

\proseitem{188}{โย โข ราธ อนตฺตา ฯเปฯ.}

\csromanheader{2}{2. ราธสํยุตฺต}
\proseitem{189}{\csromanfolio{163}โย โข ราธ อนตฺตธมฺโม ฯเปฯ.}

\proseitem{190}{โย โข ราธ ขยธมฺโม ฯเปฯ.}

\proseitem{191}{โย โข ราธ วยธมฺโม ฯเปฯ.}

\proseitem{192}{โย โข ราธ สมุทยธมฺโม, ตตฺร เต ฉนฺโท ปหาตพฺโพ, ราโค ปหาตพฺโพ, ฉนฺทราโค ปหาตพฺโพ ฯเปฯ.}

\csromanheader{2}{12. นิโรธธมฺมสุตฺต}
\proseitem{193}{สาวตฺถินิทานํ.\csromanspacer{} อายสฺมา ราโธ ภควนฺตํ เอตทโวจ\csromandash{}“สาธุ เม ภนฺเต ภควา สํขิตฺเตน ธมฺมํ เทเสตุ, ยมหํ ภควโต ธมฺมํ สุตฺวา เอโก วูปกฏฺโฐ อปฺปมตฺโต อาตาปี ปหิตตฺโต วิหเรยฺยนฺ”ติ.}

\prose{โย โข ราธ นิโรธธมฺโม, ตตฺร เต ฉนฺโท ปหาตพฺโพ, ราโค ปหาตพฺโพ, ฉนฺทราโค ปหาตพฺโพ.\csromanspacer{} โก จ ราธ นิโรธธมฺโม.\csromanspacer{} รูปํ โข ราธ นิโรธธมฺโม, ตตฺร เต ฉนฺโท ปหาตพฺโพ ฯเปฯ.\csromanspacer{} เวทนา นิโรธธมฺโม, ตตฺร เต ฉนฺโท ปหาตพฺโพ ฯเปฯ.\csromanspacer{} สญฺญา นิโรธธมฺโม, ตตฺร เต ฉนฺโท ปหาตพฺโพ ฯเปฯ.\csromanspacer{} สงฺขารา นิโรธธมฺโม, ตตฺร เต ฉนฺโท ปหาตพฺโพ ฯเปฯ.\csromanspacer{} วิญฺญาณํ นิโรธธมฺโม, ตตฺร เต ฉนฺโท ปหาตพฺโพ ฯเปฯ.\csromanspacer{} โย โข ราธ นิโรธธมฺโม, ตตฺร เต ฉนฺโท ปหาตพฺโพ, ราโค ปหาตพฺโพ, ฉนฺทราโค ปหาตพฺโพติ\footnote{สีหฬ-โปตฺถเก ปน อิมสฺมิํ วคฺเค ฉตฺติํส สุตฺตานิ วิภตฺตานิ เอเกกํ สุตฺตํ ตีณิ ตีณิ กตฺวา, เอวํ จตุตฺถวคฺเคปิ.}.\csromanspacer{} อายาจนวคฺโค ตติโย.}

\titlehead{\textbf{ตสฺสุทฺทานํ}}
\csromangathameasure{มาโรจ มารธมฺโม จ, อนิจฺเจน อปเร ทุเว. \\ ทุกฺเขน จ ทุเว วุตฺตา, อนตฺเตน ตเถว จ. \\ ขยวยสมุทยํ, นิโรธธมฺเมน ทฺวาทสาติ.}
\csromangathagroup{\csromangathastanza{มาโรจ มารธมฺโม จ, อนิจฺเจน อปเร ทุเว. \\ ทุกฺเขน จ ทุเว วุตฺตา, อนตฺเตน ตเถว จ. \\ ขยวยสมุทยํ, นิโรธธมฺเมน ทฺวาทสาติ.}}

\csromanmatikaanchor{mtk.52}
\csromantocmarkref{section}{4. อุปนิสินฺนวคฺค}{mtk.52}
\bookmark[dest={mtk.52},level=1]{4. อุปนิสินฺนวคฺค}
\csromanheader{1}{4. อุปนิสินฺนวคฺค}
\csromanheader{2}{1-11. มาราทิสุตฺต-เอกาทสก}
\proseitem{194}{\csromanfolio{164}สาวตฺถินิทานํ.\csromanspacer{} เอกมนฺตํ นิสินฺนํ โข อายสฺมนฺตํ ราธํ ภควา เอตทโวจ\csromandash{}โย โข ราธ มาโร, ตตฺร เต ฉนฺโท ปหาตพฺโพ, ราโค ปหาตพฺโพ, ฉนฺทราโค ปหาตพฺโพ.\csromanspacer{} โก จ ราธ มาโร.\csromanspacer{} รูปํ โข ราธ มาโร, ตตฺร เต ฉนฺโท ปหาตพฺโพ ฯเปฯ.\csromanspacer{} วิญฺญาณํ มาโร, ตตฺร เต ฉนฺโท ปหาตพฺโพ ฯเปฯ.\csromanspacer{} โย โข ราธ มาโร, ตตฺร เต ฉนฺโท ปหาตพฺโพ, ราโค ปหาตพฺโพ, ฉนฺทราโค ปหาตพฺโพติ.}

\proseitem{195}{โย โข ราธ มารธมฺโม, ตตฺร เต ฉนฺโท ปหาตพฺโพ, ราโค ปหาตพฺโพ, ฉนฺทราโค ปหาตพฺโพ ฯเปฯ.}

\proseitem{196}{ยํ โข ราธ อนิจฺจํ ฯเปฯ.}

\proseitem{197}{โย โข ราธ อนิจฺจธมฺโม ฯเปฯ.}

\proseitem{198}{ยํ โข ราธ ทุกฺขํ ฯเปฯ.}

\proseitem{199}{โย โข ราธ ทุกฺขธมฺโม ฯเปฯ.}

\proseitem{200}{โย โข ราธ อนตฺตา ฯเปฯ.}

\proseitem{201}{โย โข ราธ อนตฺตธมฺโม ฯเปฯ.}

\proseitem{202}{โย โข ราธ ขยธมฺโม ฯเปฯ.}

\proseitem{203}{โย โข ราธ วยธมฺโม ฯเปฯ.}

\proseitem{204}{โย โข ราธ สมุทยธมฺโม, ตตฺร เต ฉนฺโท ปหาตพฺโพ, ราโค ปหาตพฺโพ, ฉนฺทราโค ปหาตพฺโพ ฯเปฯ.}

\csromanheader{2}{12. นิโรธธมฺมสุตฺต}
\proseitem{205}{สาวตฺถินิทานํ.\csromanspacer{} เอกมนฺตํ นิสินฺนํ โข อายสฺมนฺตํ ราธํ ภควา เอตทโวจ\csromandash{}โย โข ราธ นิโรธธมฺโม, ตตฺร เต ฉนฺโท ปหาตพฺโพ, ราโค ปหาตพฺโพ, ฉนฺทราโค ปหาตพฺโพ.\csromanspacer{} โก จ ราธ นิโรธธมฺโม.}

\csromanheader{2}{2. ราธสํยุตฺต}
\prose{\csromanfolio{165}รูปํ โข ราธ นิโรธธมฺโม, ตตฺร เต ฉนฺโท ปหาตพฺโพ, ราโค ปหาตพฺโพ, ฉนฺทราโค ปหาตพฺโพ.\csromanspacer{} เวทนา ฯเปฯ.\csromanspacer{} สญฺญา ฯเปฯ.\csromanspacer{} สงฺขารา ฯเปฯ.\csromanspacer{} วิญฺญาณํ นิโรธธมฺโม, ตตฺร เต ฉนฺโท ปหาตพฺโพ, ราโค ปหาตพฺโพ, ฉนฺทราโค ปหาตพฺโพ.\csromanspacer{} โย โข ราธ นิโรธธมฺโม, ตตฺร เต ฉนฺโท ปหาตพฺโพ, ราโค ปหาตพฺโพ, ฉนฺทราโค ปหาตพฺโพติ.\csromanspacer{} อุปนิสินฺนวคฺโค จตุตฺโถ.}

\titlehead{\textbf{ตสฺสุทฺทานํ}}
\csromangathameasure{มาโร จ มารธมฺโม จ, อนิจฺเจน อปเร ทุเว. \\ ทุกฺเขน จ ทุเว วุตฺตา, อนตฺเตน ตเถว จ. \\ ขยวยสมุทยํ, นิโรธธมฺเมน ทฺวาทสาติ. \\ \textbf{ราธสํยุตฺตํ สมตฺตํ.}}
\csromangathagroup{\csromangathastanza{มาโร จ มารธมฺโม จ, อนิจฺเจน อปเร ทุเว. \\ ทุกฺเขน จ ทุเว วุตฺตา, อนตฺเตน ตเถว จ.}\csromangathastanza{ขยวยสมุทยํ, นิโรธธมฺเมน ทฺวาทสาติ. \\ \textbf{ราธสํยุตฺตํ สมตฺตํ.}}}

\csromanmatikaanchor{mtk.53}
\csromantocmarknopage{chapter}{3. ทิฏฺฐิสํยุตฺต}{mtk.53}
\bookmark[dest={mtk.53},level=0]{3. ทิฏฺฐิสํยุตฺต}
\chapterheadpageread{3. ทิฏฺฐิสํยุตฺต}
\csromanmatikaanchor{mtk.54}
\csromantocmarkref{section}{1. โสตาปตฺติวคฺค}{mtk.54}
\bookmark[dest={mtk.54},level=1]{1. โสตาปตฺติวคฺค}
\csromanheader{1}{1. โสตาปตฺติวคฺค}
\csromanheader{2}{1. วาตสุตฺต}
\proseitem{206}{\csromanfolio{166}เอกํ สมยํ ภควา สาวตฺถิยํ วิหรติ เชตวเน.\csromanspacer{} ภควา เอตทโวจ\csromandash{}กิสฺมิํ นุ โข ภิกฺขเว สติ กิํ อุปาทาย กิํ อภินิวิสฺส เอวํ ทิฏฺฐิ อุปฺปชฺชติ “น วาตา วายนฺติ, น นชฺโช สนฺทนฺติ, น คพฺภินิโย วิชายนฺติ, น จนฺทิมสูริยา อุเทนฺติ วา อเปนฺติ วา, เอสิกฏฺฐายิฏฺฐิตา”ติ.}

\prose{ภควํมูลกา โน ภนฺเต ธมฺมา ภควํเนตฺติกา ภควํปฏิสรณา, สาธุ วต ภนฺเต ภควนฺตญฺเญว ปฏิภาตุ เอตสฺส ภาสิตสฺส อตฺโถ, ภควโต สุตฺวา ภิกฺขู ธาเรสฺสนฺตีติ.\csromanspacer{} เตน หิ ภิกฺขเว สุณาถ สาธุกํ มนสิ กโรถ, ภาสิสฺสามีติ. “เอวํ ภนฺเต”ติ โข เต ภิกฺขู ภควโต ปจฺจสฺโสสุํ.\csromanspacer{} ภควา เอตทโวจ\csromandash{}}

\prose{รูเป โข ภิกฺขเว สติ รูปํ อุปาทาย รูปํ อภินิวิสฺส เอวํ ทิฏฺฐิ อุปฺปชฺชติ “น วาตา วายนฺติ, น นชฺโช สนฺทนฺติ, น คพฺภินิโย วิชายนฺติ, น จนฺทิมสูริยา อุเทนฺติ วา อเปนฺติ วา, เอสิกฏฺฐายิฏฺฐิตา”ติ.\csromanspacer{} เวทนาย สติ ฯเปฯ.\csromanspacer{} สญฺญาย สติ.\csromanspacer{} สงฺขาเรสุ สติ.\csromanspacer{} วิญฺญาเณ สติ วิญฺญาณํ อุปาทาย วิญฺญาณํ อภินิวิสฺส เอวํ ทิฏฺฐิ อุปฺปชฺชติ “น วาตา วายนฺติ, น นชฺโช สนฺทนฺติ, น คพฺภินิโย วิชายนฺติ, น จนฺทิมสูริยา อุเทนฺติ วา อเปนฺติ วา, เอสิกฏฺฐายิฏฺฐิตา”ติ.\csromanspacer{} ตํ กิํ มญฺญถ ภิกฺขเว, รูปํ นิจฺจํ วา อนิจฺจํ วาติ.\csromanspacer{} อนิจฺจํ ภนฺเต.\csromanspacer{} ยํ ปนานิจฺจํ, ทุกฺขํ วา ตํ สุขํ วาติ.\csromanspacer{} ทุกฺขํ ภนฺเต.\csromanspacer{} ยํ ปนานิจฺจํ ทุกฺขํ วิปริณามธมฺมํ, อปิ นุ ตํ อนุปาทาย เอวํ ทิฏฺฐิ อุปฺปชฺเชยฺย “น วาตา วายนฺติ, น นชฺโช สนฺทนฺติ, น คพฺภินิโย วิชายนฺติ, น จนฺทิมสูริยา อุเทนฺติ วา อเปนฺติ วา เอสิกฏฺฐายิฏฺฐิตา”ติ.\csromanspacer{} โน เหตํ ภนฺเต.}

\prose{เวทนา นิจฺจา วา อนิจฺจา วาติ.\csromanspacer{} สญฺญา.\csromanspacer{} สงฺขารา.\csromanspacer{} วิญฺญาณํ นิจฺจํ วา อนิจฺจํ วาติ.\csromanspacer{} อนิจฺจํ ภนฺเต.\csromanspacer{} ยํ ปนานิจฺจํ, ทุกฺขํ วา ตํ สุขํ วาติ.\csromanspacer{} ทุกฺขํ ภนฺเต.\csromanspacer{} ยํ ปนานิจฺจํ, ทุกฺขํ วิปริณามธมฺมํ, อปิ นุ ตํ อนุปาทาย เอวํ ทิฏฺฐิ อุปฺปชฺเชยฺย “น วาตา วายนฺติ, น นชฺโช สนฺทนฺติ, น คพฺภินิโย วิชายนฺติ, น จนฺทิมสูริยา}

\csromanheader{2}{3. ทิฏฺฐิสํยุตฺต}
\prose{\csromanfolio{167}อุเทนฺติ วา อเปนฺติวา, เอสิกฏฺฐายิฏฺฐาตา”ติ.\csromanspacer{} โน เหตํ ภนฺเต.\csromanspacer{} ยมฺปิทํ\footnote{ยมิทํ (อญฺญตฺถ)} ทิฏฺฐิ สุตํ มุตํ วิญฺญาตํ ปตฺตํ ปริเยสิตํ อนุวิจริตํ มนสา, ตมฺปิ นิจฺจํ วา อนิจฺจํ วาติ.\csromanspacer{} อนิจฺจํ ภนฺเต.\csromanspacer{} ยํ ปนานิจฺจํ, ทุกฺขํ วา ตํ สุขํ วาติ.\csromanspacer{} ทุกฺขํ ภนฺเต.\csromanspacer{} ยํ ปนานิจฺจํ ทุกฺขํ วิปริณามธมฺมํ, อปิ นุ ตํ อนุปาทาย เอวํ ทิฏฺฐิ อุปฺปชฺเชยฺย “น วาตา วายนฺติ, น นชฺโช สนฺทนฺติ, น คพฺภินิโย วิชายนฺติ, น จนฺทิมสูริยา อุเทนฺติ วา อเปนฺติ วา, เอสิกฏฺฐายิฏฺฐิตา”ติ.\csromanspacer{} โน เหตํ ภนฺเต.}

\prose{ยโต โข ภิกฺขเว อริยสาวกสฺส อิเมสุ จ\footnote{อิเมสุ ฉสุ (สี, สฺยา, กํ, อิ) เอวมุปริปิ.} ฐาเนสุ กงฺขา ปหีนา โหติ, ทุกฺเขปิสฺส กงฺขา ปหีนา โหติ, ทุกฺขสมุทเยปิสฺส กงฺขา ปหีนา โหติ, ทุกฺขนิโรเธปิสฺส กงฺขา ปหีนา โหติ, ทุกฺขนิโรธคามินิยา ปฏิปทายปิสฺส กงฺขา ปหีนา โหติ.\csromanspacer{} อยํ วุจฺจติ ภิกฺขเว อริยสาวโก โสตาปนฺโน อวินิปาตธมฺโม นิยโต สมฺโภธิปรายโนติ.\csromanspacer{}\csromanspacer{}ปฐมํ.}

\csromanheader{2}{2. เอตํมมสุตฺต}
\proseitem{207}{สาวตฺถินิทานํ.\csromanspacer{} กิสฺมิํ นุ โข ภิกฺขเว สติ กิํ อุปาทาย กิํ อภินิวิสฺส เอวํ ทิฏฺฐิ อุปฺปชฺชติ “เอตํ มม, เอโสหมสฺมิ, เอโส เม อตฺตา”ติ.\csromanspacer{} ภควํมูลกา โน ภนฺเต ธมฺมา ฯเปฯ.\csromanspacer{} รูเป โข ภิกฺขเว สติ รูปํ อุปาทาย รูปํ อภินิวิสฺส เอวํ ทิฏฺฐิ อุปฺปชฺชติ “เอตํ มม, เอโสหมสฺมิ, เอโส เม อตฺตา”ติ.\csromanspacer{} เวทนาย สติ ฯเปฯ.\csromanspacer{} สญฺญาย สติ.\csromanspacer{} สงฺขาเรสุ สติ.\csromanspacer{} วิญฺญาเณ สติ วิญฺญาณํ อุปาทาย วิญฺญาณํ อภินิวิสฺส เอวํ ทิฏฺฐิ อุปฺปชฺชติ “เอตํ มม, เอโสหมสฺมิ, เอโส เม อตฺตา”ติ.}

\prose{ตํ กิํ มญฺญถ ภิกฺขเว, รูปํ นิจฺจํ วา อนิจฺจํ วาติ.\csromanspacer{} อนิจฺจํ ภนฺเต ฯเปฯ.\csromanspacer{} เวทนา.\csromanspacer{} สญฺญา.\csromanspacer{} สงฺขารา.\csromanspacer{} วิญฺญาณํ นิจฺจํ วา อนิจฺจํ วาติ.\csromanspacer{} อนิจฺจํ ภนฺเต ฯเปฯ.\csromanspacer{} อปิ นุ ตํ อนุปาทาย เอวํ ทิฏฺฐิ อุปฺปชฺเชยฺย “เอตํ มม, เอโสหมสฺมิ, เอโส เม อตฺตา”ติ.\csromanspacer{} โน เหตํ ภนฺเต.\csromanspacer{} ยมฺปิทํ ทิฏฺฐํ สุตํ มุตํ วิญฺญาตํ ปตฺตํ ปริเยสิตํ อนุวิจริตํ มนสา.\csromanspacer{} ตมฺปิ นิจฺจํ วา อนิจฺจํ วาติ.\csromanspacer{} อนิจฺจํ ภนฺเต.\csromanspacer{} ยํ ปนานิจฺจํ, ทุกฺขํ วา ตํ สุขํ วาติ.\csromanspacer{} ทุกฺขํ ภนฺเต.\csromanspacer{} ยํ ปนานิจฺจํ ทุกฺขํ วิปริณามธมฺมํ, อปิ นุ ตํ อนุปาทาย เอวํ ทิฏฺฐิ อุปฺปชฺเชยฺย “เอตํ มม, เอโสหมสฺมิ, เอโส เม อตฺตา”ติ.\csromanspacer{} โน เหตํ ภนฺเต.}

\prose{\csromanfolio{168}ยโต โข ภิกฺขเว อริยสาวกสฺส อิเมสุ จ ฐาเนสุ กงฺขา ปหีนา โหติ, ทุกฺเขปิสฺส กงฺขา ปหีนา โหติ ฯเปฯ ทุกฺขนิโรธคามินิยา ปฏิปทายปิสฺส กงฺขา ปหีนา โหติ.\csromanspacer{} อยํ วุจฺจติ ภิกฺขเว อริยสาวโก โสตาปนฺโน อวินิปาตธมฺโม นิยโต สมฺโพธิปรายโนติ.\csromanspacer{}\csromanspacer{}ทุติยํ.}

\csromanheader{2}{3. โส-อตฺตาสุตฺต}
\proseitem{208}{สาวตฺถินิทานํ.\csromanspacer{} กิสฺมิํ นุ โข ภิกฺขเว สติ กิํ อุปาทาย กิํ อภินิวิสฺส เอวํ ทิฏฺฐิ อุปฺปชฺชติ “โส อตฺตา โส โลโก โส เปจฺจ ภวิสฺสามิ นิจฺโจ ธุโว สสฺสโต อวิปริณามธมฺโม”ติ.\csromanspacer{} ภควํมูลกา โน ภนฺเต ธมฺมา ฯเปฯ.}

\prose{รูเป โข ภิกฺขเว สติ รูปํ อุปาทาย รูปํ อภินิวิสฺส เอวํ ทิฏฺฐิ อุปฺปชฺชติ “โส อตฺตา โส โลโก โส เปจฺจ ภวิสฺสามิ นิจฺโจ ธุโว สสฺสโต อวิปริณามธมฺโม”ติ.\csromanspacer{} เวทนาย สติ ฯเปฯ.\csromanspacer{} สญฺญาย สติ.\csromanspacer{} สงฺขาเรสุ สติ.\csromanspacer{} วิญฺญาเณ สติ วิญฺญาณํ อุปาทาย วิญฺญาณํ อภินิวิสฺส เอวํ ทิฏฺฐิ อุปฺปชฺชติ “โส อตฺตา โส โลโก โส เปจฺจ ภวิสฺสามิ นิจฺโจ ธุโว สสฺสโต อวิปริณามธมฺโม”ติ.}

\prose{ตํ กิํ มญฺญถ ภิกฺขเว, รูปํ นิจฺจํ วา อนิจฺจํ วาติ.\csromanspacer{} อนิจฺจํ ภนฺเต ฯเปฯ อปิ นุ ตํ อนุปาทาย เอวํ ทิฏฺฐิ อุปฺปชฺเชยฺย “โส อตฺตา ฯเปฯ อวิปริณามธมฺโม”ติ.\csromanspacer{} โน เหตํ ภนฺเต.\csromanspacer{} เวทนา.\csromanspacer{} สญฺญา.\csromanspacer{} สงฺขารา.\csromanspacer{} วิญฺญณํ นิจฺจํ วา อนิจฺจํ วาติ.\csromanspacer{} อนิจฺจํ ภนฺเต ฯเปฯ อปิ นุ ตํ อนุปาทาย เอวํ ทิฏฺฐิ อุปฺปชฺเชยฺย “โส อตฺตา ฯเปฯ อวิปริณามธมฺโม”ติ.\csromanspacer{} โน เหตํ ภนฺเต.\csromanspacer{} ยมฺปิทํ ทิฏฺฐํ สุตํ มุตํ วิญฺญาตํ ปตฺตํ ปริเยสิตํ อนุวิจริตํ มนสา, ตมฺปิ นิจฺจํ วา อนิจฺจํ วาติ.\csromanspacer{} อนิจฺจํ ภนฺเต ฯเปฯ อปิ นุ ตํ อนุปาทาย เอวํ ทิฏฺฐิ อุปฺปชฺเชยฺย “โส อตฺตา โส โลโก โส เปจฺจ ภวิสฺสามิ นิจฺโจ ธุโว สสฺสโต อวิปริณามธมฺโม”ติ.\csromanspacer{} โน เหตํ ภนฺเต.}

\prose{ยโต โข ภิกฺขเว อริยสาวกสฺส อิเมสุ จ ฐาเนสุ กงฺขา ปหีนา โหติ, ทุกฺเขปิสฺส กงฺขา ปหีนา โหติ ฯเปฯ ทุกฺขนิโรธคามินิยา ปฏิปทายปิสฺส กงฺขา ปหีนา โหติ.\csromanspacer{} อยํ วุจฺจติ ภิกฺขเว อริยสาวโก โสตาปนฺโน อวินิปาตธมฺโม นิยโต สมฺโพธิปรายโนติ.\csromanspacer{}\csromanspacer{}ตติยํ.}

\csromanheader{2}{3. ทิฏฺฐิสํยุตฺต}
\csromanheader{2}{4. โนจเมสิยาสุตฺต}
\proseitem{209}{\csromanfolio{169}สาวตฺถินิทานํ.\csromanspacer{} กิสฺมิํ นุ โข ภิกฺขเว สติ กิํ อุปาทาย กิํ อภินิวิสฺส เอวํ ทิฏฺฐิ อุปฺปชฺชติ “โน จสฺสํ, โน จ เม สิยา, นาภวิสฺส, น เม ภวิสฺสตี”ติ.\csromanspacer{} ภควํมูลกา โน ภนฺเต ธมฺมา ฯเปฯ.}

\prose{รูเป โข ภิกฺขเว สติ รูปํ อุปาทาย รูปํ อภินิวิสฺส เอวํ ทิฏฺฐิ อุปฺปชฺชติ “โน จสฺสํ, โน จ เม สิยา, นาภวิสฺส, น เม ภวิสฺสตี”ติ.\csromanspacer{} เวทนาย สติ.\csromanspacer{} สญฺญาย สติ.\csromanspacer{} สงฺขาเรสุ สติ.\csromanspacer{} วิญฺญาเณ สติ วิญฺญาณํ อุปาทาย วิญฺญาณํ อภินิวิสฺส เอวํ ทิฏฺฐิ อุปฺปชฺชติ “โน จสฺสํ, โน จ เม สิยา, นาภวิสฺส, น เม ภวิสฺสตี”ติ.}

\prose{ตํ กิํ มญฺญถ ภิกฺขเว, รูปํ นิจฺจํ วา อนิจฺจํ วาติ.\csromanspacer{} อนิจฺจํ ภนฺเต ฯเปฯ อปิ นุ ตํ อนุปาทาย เอวํ ทิฏฺฐิ อุปฺปชฺเชยฺย “โน จสฺสํ, โน จ เม สิยา, นาภวิสฺส, น เม ภวิสฺสตี”ติ.\csromanspacer{} โน เหตํ ภนฺเต.\csromanspacer{} เวทนา.\csromanspacer{} สญฺญา.\csromanspacer{} สงฺขารา.\csromanspacer{} วิญฺญาณํ นิจฺจํ วา อนิจฺจํ วาติ.\csromanspacer{} อนิจฺจํ ภนฺเต ฯเปฯ อปิ นุ ตํ อนุปาทาย เอวํ ทิฏฺฐิ อุปฺปชฺเชยฺย “โน จสฺสํ, โน จ เม สิยา, นาภวิสฺส, น เม ภวิสฺสตี”ติ.\csromanspacer{} โน เหตํ ภนฺเต.\csromanspacer{} ยมฺปิทํ ทิฏฺฐํ สุตํ มุตํ วิญฺญาตํ ปตฺตํ ปริเยสิตํ อนุวิจริตํ มนสา, ตมฺปิ นิจฺจํ วา อนิจฺจํ วาติ, อนิจฺจํ ภนฺเต ฯเปฯ อปิ นุ ตํ อนุปาทาย เอวํ ทิฏฺฐิ อุปฺปชฺเชยฺย “โน จสฺสํ, โน จ เม สิยา, นาภวิสฺส, น เม ภวิสฺสตี”ติ.\csromanspacer{} โน เหตํ ภนฺเต.}

\prose{ยโต โข ภิกฺขเว อริยสาวกสฺส อิเมสุ จ ฐาเนสุ กงฺขา ปหีนา โหติ, ทุกฺเขปิสฺส กงฺขา ปหีนา โหตินฺ ฯเปฯ ทุกฺขนิโรธคามินิยา ปฏิปทายปิสฺส กงฺขา ปหีนา โหติ.\csromanspacer{} อยํ วุจฺจติ ภิกฺขเว อริยสาวโก โสตาปนฺโน อวินิปาตธมฺโม นิยโต สมฺโพธิปรายโนติ.\csromanspacer{}\csromanspacer{}จตุตฺถํ.}

\csromanheader{2}{5. นตฺถิทินฺนสุตฺต}
\proseitem{210}{สาวตฺถินิทานํ.\csromanspacer{} กิสฺมิํ นุ โข ภิกฺขเว สติ กิํ อุปาทาย กิํ อภินิวิสฺส เอวํ ทิฏฺฐิ อุปฺปชฺชติ “นตฺถิ ทินฺนํ, นตฺถิ ยิฏฺฐํ, นตฺถิ หุตํ, นตฺถิ สุกตทุกฺกฏานํ\footnote{สุกฺกฏทุกฺกฏานํ (สี, อิ)} กมฺมานํ ผลํ วิปาโก, นตฺถิ อยํ โลโก, นตฺถิ ปโร โลโก, นตฺถิ มาตา, นตฺถิ ปิตา, นตฺถิ สตฺตา โอปปาติกา, นตฺถิ \csromanfolio{170}โลเก สมณพฺราหฺมณา สมฺมคฺคตา\footnote{สมคฺคตา (ก)} สมฺมาปฏิปนฺนา, เย อิมญฺจ โลกํ ปรญฺจ โลกํ สยํ อภิญฺญา สจฺฉิกตฺวา ปเวเทนฺติ.\csromanspacer{} จาตุมหาภูติโก\footnote{จาตุมฺมหาภูติโก (สี, สฺยา, กํ, อิ)} อยํ ปุริโส ยทา กาลํ กโรติ, ปถวี ปถวีกายํ อนุเปติ อนุปคจฺฉติ, อาโป อาโปกายํ อนุเปติ อนุปคจฺฉติ, เตโช เตโชกายํ อนุเปติ อนุปคจฺฉติ, วาโย วาโยกายํ อนุเปติ อนุปคจฺฉติ, อากาสํ อินฺทฺริยานิ สงฺกมนฺติ, อาสนฺทิปญฺจมา ปุริสา มตํ อาทาย คจฺฉนฺติ ยาว อาฬาหนา ปทานิ ปญฺญายนฺติ, กาโปตกานิ อฏฺฐีนิ ภวนฺติ, ภสฺสนฺตา อาหุติโย, ทตฺตุปญฺญตฺตํ ยทิทํ ทานํ\footnote{ทตฺตุปญฺญตฺตมิทํ ทานํ นาม (สพฺพตฺถ)}, เตสํ ตุจฺฉํ มุสา วิลาโป, เย เกจิ อตฺถิกวาทํ วทนฺติ, พาเล จ ปณฺฑิเต จ กายสฺส เภทา อุจฺฉิชฺชนฺติ วินสฺสนฺติ, น โหนฺติ ปรํ มรณา”ติ.\csromanspacer{} ภควํมูลกา โน ภนฺเต ธมฺมา ฯเปฯ.\csromanspacer{} รูเป โข ภิกฺขเว สติ รูปํ อุปาทาย รูปํ อภินิวิสฺส เอวํ ทิฏฺฐิ อุปฺปชฺชติ “นตฺถิ ทินฺนํ, นตฺถิ ยิฏฺฐํ ฯเปฯ กายสฺส เภทา อุจฺฉิชฺชนฺติ วินสฺสนฺติ, น โหนฺติ ปรํ มรณา”ติ.\csromanspacer{} เวทนาย สติ ฯเปฯ.\csromanspacer{} สญฺญาย สติ.\csromanspacer{} สงฺขาเรสุ สติ.\csromanspacer{} วิญฺญาเณ สติ วิญฺญาณํ อุปาทาย วิญฺญาณํ อภินิวิสฺส เอวํ ทิฏฺฐิ อุปฺปชฺชติ “นตฺถิ ทินฺนํ, นตฺถิ ยิฏฺฐํ ฯเปฯ กายสฺส เภทา อุจฺฉิชฺชนฺติ วินสฺสนฺติ, น โหนฺติ ปรํ มรณา”ติ.}

\prose{ตํ กิํ มญฺญถ ภิกฺขเว, รูปํ นิจฺจํ วา อนิจฺจํ วาติ.\csromanspacer{} อนิจฺจํ ภนฺเต ฯเปฯ อปิ นุ ตํ อนุปาทาย เอวํ ทิฏฺฐิ อุปฺปชฺเชยฺย “นตฺถิ ทินฺนํ, นตฺถิ ยิฏฺฐํ ฯเปฯ กายสฺส เภทา อุจฺฉิชฺชนฺติ วินสฺสนฺติ, น โหนฺติ ปรํ มรณา”ติ.\csromanspacer{} โน เหตํ ภนฺเต.\csromanspacer{} เวทนา.\csromanspacer{} สญฺญา.\csromanspacer{} สงฺขารา.\csromanspacer{} วิญฺญาณํ นิจฺจํ วา อนิจฺจํ วาติ.\csromanspacer{} อนิจฺจํ ภนฺเต ฯเปฯ อปิ นุ ตํ อนุปาทาย เอวํ ทิฏฺฐิ อุปฺปชฺเชยฺย “นตฺถิ ทินฺนํ, นตฺถิ ยิฏฺฐํ ฯเปฯ กายสฺส เภทา อุจฺฉิชฺชนฺติ วินสฺสนฺติ, น โหนฺติ ปรํ มรณา”ติ.\csromanspacer{} โน เหตํ ภนฺเต.\csromanspacer{} ยมฺปิทํ ทิฏฺฐํ สุตํ มุตํ วิญฺญาตํ ปตฺตํ ปริเยสิตํ อนุวิจริตํ มนสา, ตมฺปิ นิจฺจํ วา อนิจฺจํ วาติ.\csromanspacer{} อนิจฺจํ ภนฺเต ฯเปฯ อปิ นุ ตํ อนุปาทาย เอวํ ทิฏฺฐิ อุปฺปชฺเชยฺย “นตฺถิ ทินฺนํ, นตฺถิ ยิฏฺฐํ ฯเปฯ เย เกจิ อตฺถิกวาทํ วทนฺติ, พาเล จ ปณฺฑิเต จ กายสฺส เภทา อุจฺฉิชฺชนฺติ วินสฺสนฺติ, น โหนฺติ ปรํ มรณา”ติ.\csromanspacer{} โน เหตํ ภนฺเต.}

\prose{ยโต โข ภิกฺขเว อริยสาวกสฺส อิเมสุ จ ฐาเนสุ กงฺขา ปหีนา โหติ, ทุกฺเขปิสฺส กงฺขา ปหีนา โหติ ฯเปฯ ทุกฺขนิโรธคามินิยา}

\csromanheader{2}{3. ทิฏฺฐิสํยุตฺต}
\prose{\csromanfolio{171}ปฏิปทายปิสฺส กงฺขา ปหีนา โหติ.\csromanspacer{} อยํ วุจฺจติ ภิกฺขเว อริยสาวโก โสตาปนฺโน อวินิปาตธมฺโม นิยโต สมฺโพธิปรายโนติ.\csromanspacer{}\csromanspacer{}ปญฺจมํ.}

\csromanheader{2}{6. กโรโตสุตฺต}
\proseitem{211}{สาวตฺถินิทานํ.\csromanspacer{} กิสฺมิํ นุ โข ภิกฺขเว สติ กิํ อุปาทาย กิํ อภินิวิสฺส เอวํ ทิฏฺฐิ อุปฺปชฺชติ “กโรโต การยโต, ฉินฺทโต เฉทาปยโต, ปจโต ปาจาปยโต, โสจโต โสจาปยโต, กิลมโต กิลมาปยโต, ผนฺทโต ผนฺทาปยโต, ปาณมติปาตยโต, อทินฺนํ อาทิยโต, สนฺธิํ ฉนฺทโต, นิลฺโลปํ หรโต, เอกาคาริกํ กโรโต, ปริปนฺเถ ติฏฺฐโต, ปรทารํ คจฺฉโต, มุสา ภณโต, กโรโต น กรียติ ปาปํ, ขุรปริยนฺเตน เจปิ จกฺเกน โย อิมิสฺสา ปถวิยา ปาเณ เอกมํสขลํ เอกมํสปุญฺชํ กเรยฺย, นตฺถิ ตโตนิทานํ ปาปํ, นตฺถิ ปาปสฺส อาคโม.\csromanspacer{} ทกฺขิณญฺเจปิ คงฺคาย ตีรํ คจฺเฉยฺย, หนนฺโต ฆาเตนฺโต ฉินฺทนฺโต เฉทาเปนฺโต ปจนฺโต ปาเจนฺโต, นตฺถิ ตโตนิทานํ ปาปํ, นตฺถิ ปาปสฺส อาคโม.\csromanspacer{} อุตฺตรญฺเจปิ คงฺคาย ตีรํ คจฺเฉยฺย, ททนฺโต ทาเปนฺโต, ยชนฺโต ยชาเปนฺโต, นตฺถิ ตโตนิทานํ ปุญฺญํ, นตฺถิ ปุญฺญสฺส อาคโม.\csromanspacer{} ทาเนน ทเมน สํยเมน สจฺจวชฺเชน นตฺถิ ปุญฺญํ, นตฺถิ ปุญฺญสฺส อาคโมติ.\csromanspacer{} ภควํมูลกา โน ภนฺเต ธมฺมา ฯเปฯ.\csromanspacer{} รูเป โข ภิกฺขเว สติ รูปํ อุปาทาย รูปํ อภินิวิสฺส เอวํ ทิฏฺฐิ อุปฺปชฺชติ “กโรโต การยโต ฯเปฯ นตฺถิ ปุญฺญํ นตฺถิ ปุญฺญสฺส อาคโม”ติ.\csromanspacer{} เวทนาย สติ ฯเปฯ.\csromanspacer{} สญฺญาย สติ.\csromanspacer{} สงฺขาเรสุ สติ.\csromanspacer{} วิญฺญาเณ สติ วิญฺญาณํ อุปาทาย วิญฺญาณํ อภินิวิสฺส เอวํ ทิฏฺฐิ อุปฺปชฺชติ “กโรโต การยโต ฯเปฯ นตฺถิ ปุญฺญํ, นตฺถิ ปุญฺญสฺส อาคโม”ติ.}

\prose{ตํ กิํ มญฺญถ ภิกฺขเว, รูปํ นิจฺจํ วา อนิจฺจํ วาติ.\csromanspacer{} อนิจฺจํ ภนฺเต ฯเปฯ อปิ นุ ตํ อนุปาทาย เอวํ ทิฏฺฐิ อุปฺปชฺเชยฺย “กโรโต ฯเปฯ นตฺถิ ปุญฺญํ, นตฺถิ ปุญฺญสฺส อาคโม”ติ.\csromanspacer{} โน เหตํ ภนฺเต.\csromanspacer{} เวทนา.\csromanspacer{} สญฺญา.\csromanspacer{} สงฺขารา.\csromanspacer{} วิญฺญาณํ นิจฺจํ วา อนิจฺจํ วาติ.\csromanspacer{} อนิจฺจํ ภนฺเต ฯเปฯ อปิ นุ ตํ อนุปาทาย เอวํ ทิฏฺฐิ อุปฺปชฺเชยฺย “กโรโต การยโต ฯเปฯ นตฺถิ ปุญฺญํ, นตฺถิ ปุญฺญสฺส อาคโม”ติ.\csromanspacer{} โน เหตํ ภนฺเต.\csromanspacer{} ยมฺปิทํ ทิฏฺฐํ สุตํ มุตํ วิญฺญาตํ ปตฺตํ ปริเยสิตํ อนุวิจริตํ มนสา, ตมฺปิ นิจฺจํ วา อนิจฺจํ วาติ.\csromanspacer{} อนิจฺจํ ภนฺเต ฯเปฯ อปิ นุ ตํ อนุปาทาย เอวํ ทิฏฺฐิ อุปฺปชฺเชยฺย “กโรโต \csromanfolio{172}การยโต ฯเปฯ นตฺถิ ปุญฺญํ, นตฺถิ ปุญฺญสฺส อาคโม”ติ.\csromanspacer{} โน เหตํ ภนฺเต.}

\prose{ยโต โข ภิกฺขเว อริยสาวกสฺส อิเมสุ จ ฐาเนสุ กงฺขา ปหีนา โหติ, ทุกฺเขปิสฺส กงฺขา ปหีนา โหติ ฯเปฯ ทุกฺขนิโรธคามินิยา ปฏิปทายปิสฺส กงฺขา ปหีนา โหติ.\csromanspacer{} อยํ วุจฺจติ ภิกฺขเว อริยสาวโก โสตาปนฺโน อวินิปาตธมฺโม นิยโต สมฺโพธิปรายโนติ.\csromanspacer{}\csromanspacer{}ฉฏฺฐํ.}

\csromanheader{2}{7. เหตุสุตฺต}
\proseitem{212}{สาวตฺถินิทานํ.\csromanspacer{} กิสฺมิํ นุ โข ภิกฺขเว สติ กิํ อุปาทาย กิํ อภินิวิสฺส เอวํ ทิฏฺฐิ อุปฺปชฺชติ “นตฺถิ เหตุ นตฺถิ ปจฺจโย สตฺตานํ สํกิเลสาย, อเหตู อปฺปจฺจยา สตฺตา สํกิลิสฺสนฺติ.\csromanspacer{} นตฺถิ เหตุ นตฺถิ ปจฺจโย สตฺตานํ วิสุทฺธิยา, อเหตู อปฺปจฺจยา สตฺตา วิสุชฺฌนฺติ.\csromanspacer{} นตฺถิ พลํ, นตฺถิ วีริยํ, นตฺถิ ปุริสถาโม, นตฺถิ ปุริสปรกฺกโม, สพฺเพ สตฺตา สพฺเพ ปาณา สพฺเพ ภูตา สพฺเพ ชีวา อวสา อพลา อวีริยา นิยติ สงฺคติ ภาว ปริณตา ฉเสฺววาภิชาตีสุ สุขทุกฺขํ ปฏิสํเวเทนฺตี”ติ.\csromanspacer{} ภควํมูลกา โน ภนฺเต ธมฺมา ฯเปฯ.\csromanspacer{} รูเป โข ภิกฺขเว สติ รูปํ อุปาทาย รูปํ อภินิวิสฺส เอวํ ทิฏฺฐิ อุปฺปชฺชติ “นตฺถิ เหตุ, นตฺถิ ปจฺจโย ฯเปฯ สุขทุกฺขํ ปฏิสํเวเทนฺตี”ติ.\csromanspacer{} เวทนาย สติ ฯเปฯ.\csromanspacer{} สญฺญาย สติ.\csromanspacer{} สงฺขาเรสุ สติ.\csromanspacer{} วิญฺญาเณ สติ วิญฺญาณํ อุปาทาย วิญฺญาณํ อภินิวิสฺส เอวํ ทิฏฺฐิ อุปฺปชฺชติ “นตฺถิ เหตุ, นตฺถิ ปจฺจโย ฯเปฯ สุขทุกฺขํ ปฏิสํเวเทนฺตี”ติ.}

\prose{ตํ กิํ มญฺญถ ภิกฺขเว, รูปํ นิจฺจํ วา อนิจฺจํ วาติ.\csromanspacer{} อนิจฺจํ ภนฺเต ฯเปฯ วิปริณามธมฺมํ.\csromanspacer{} อปิ นุ ตํ อนุปาทาย เอวํ ทิฏฺฐิ อุปฺปชฺเชยฺย “นตฺถิ เหตุ, นตฺถิ ปจฺจโย ฯเปฯ สุขทุกฺขํ ปฏิสํเวเทนฺตี”ติ.\csromanspacer{} โน เหตํ ภนฺเต.\csromanspacer{} เวทนา.\csromanspacer{} สญฺญา.\csromanspacer{} สงฺขารา.\csromanspacer{} วิญฺญาณํ นิจฺจํ วา อนิจฺจํ วาติ.\csromanspacer{} อนิจฺจํ ภนฺเต ฯเปฯ อปิ นุ ตํ อนุปาทาย เอวํ ทิฏฺฐิ อุปฺปชฺเชยฺย “นตฺถิ เหตุ นตฺถิ ปจฺจโย ฯเปฯ สุขทุกฺขํ ปฏิสํเวเทนฺตี”ติ.\csromanspacer{} โน เหตํ ภนฺเต.\csromanspacer{} ยมฺปิทํ ทิฏฺฐํ สุตํ มุตํ วิญฺญาตํ ปตฺตํ ปริเยสิตํ อนุวิจริตํ มนสา, ตมฺปิ นิจฺจํ วา อนิจฺจํ วาติ.\csromanspacer{} อนิจฺจํ ภนฺเต ฯเปฯ อปิ นุ ตํ อนุปาทาย เอวํ ทิฏฺฐิ อุปฺปชฺเชยฺย “นตฺถิ เหตุ นตฺถิ ปจฺจโย ฯเปฯ สุขทุกฺขํ ปฏิสํเวเทนฺตี”ติ.\csromanspacer{} โน เหตํ ภนฺเต.}

\csromanheader{2}{3. ทิฏฺฐิสํยุตฺต}
\prose{\csromanfolio{173}ยโต โข ภิกฺขเว อริยสาวกสฺส อิเมสุ จ ฐาเนสุ กงฺขา ปหีนา โหติ, ทุกฺเขปิสฺส กงฺขา ปหีนา โหติ ฯเปฯ ทุกฺขนิโรธคามินิยา ปฏิปทายปิสฺส กงฺขา ปหีนา โหติ.\csromanspacer{} อยํ วุจฺจติ ภิกฺขเว อริยสาวโก โสตาปนฺโน อวินิปาตธมฺโม นิยโต สมฺโพธิปรายโนติ.\csromanspacer{}\csromanspacer{}สตฺตมํ.}

\csromanheader{2}{8. มหาทิฏฺฐิสุตฺต}
\proseitem{213}{สาวตฺถินิทานํ.\csromanspacer{} กิสฺมิํ นุ โข ภิกฺขเว สติ กิํ อุปาทาย กิํ อภินิวิสฺส เอวํ ทิฏฺฐิ อุปฺปชฺชติ “สตฺติเม กายา อกฏา อกฏวิธา อนิมฺมิตา อนิมฺมาตา วญฺฌา กูฏฏฺฐา เอสิกฏฺฐายิฏฺฐิตา, เต น อิญฺชนฺติ, น วิปริณมนฺติ\footnote{น วิปริณาเมนฺติ (อิ, ก)}, น อญฺญมญฺญํ พฺยาพาเธนฺติ, นาลํ อญฺญมญฺญสฺส สุขาย วา ทุกฺขาย วา สุขทุกฺขาย วา.\csromanspacer{} กตเม สตฺต.\csromanspacer{} ปถวีกาโย อาโปกาโย เตโชกาโย วาโยกาโย สุเข ทุกฺเข ชีเว สตฺตเม.\csromanspacer{} อิเม สตฺต\footnote{ชีเว. สตฺติเม (พหูสุ)} กายา อกฏา อกฏวิธา อนิมฺมิตา อนิมฺมาตา วญฺฌา กูฏฏฺฐา เอสิกฏฺฐายิฏฺฐิตา, เต น อิญฺชนฺติ, น วิปริณมนฺติ, น อญฺญมญฺญํ พฺยาพาเธนฺติ, นาลํ อญฺญมญฺญสฺส สุขาย วา ทุกฺขาย วา สุขทุกฺขาย วา.\csromanspacer{} โยปิ ติเณฺหน สตฺเถน สีสํ ฉินฺทติ, น โสปิ กิญฺจิ\footnote{น โกจิ กญฺจิ (สี, สฺยา, กํ), น โกจิ ตํ (อิ, ก)} ชีวิตา โวโรเปติ, สตฺตนฺนํเตฺวว กายานมนฺตเรน สตฺถํ วิวรมนุปวิสติ\footnote{วิวรมนุปตติ (กตฺถจิ) ทีฆมชฺฌิเมสุปิ.}.\csromanspacer{} จุทฺทส โข ปนิมานิ โยนิปมุขสตสหสฺสานิ, สฏฺฐิ จ สตานิ ฉ จ สตานิ ปญฺจ จ กมฺมุโน สตานิ ปญฺจ จ กมฺมานิ ตีณิ จ กมฺมานิ กมฺเม จ อฑฺฒกมฺเม จ ทฺวฏฺฐิปฏิปทา ทฺวฏฺฐนฺตรกปฺปา ฉฬาภิชาติโย อฏฺฐ ปุริสภูมิโย เอกูนปญฺญาส อาชีวกสเต เอกูนปญฺญาส ปริพฺพาชกสเต เอกูนปญฺญาส นาควาสสเต วีเส อินฺทฺริยสเต ติํเส นิรยสเต ฉตฺติํสรโชธาตุโย สตฺต สญฺญีคพฺภา สตฺต อสญฺญีคพฺภา สตฺต นิคณฺฐิคพฺภา สตฺต เทวา สตฺต มานุสา สตฺต เปสาจา สตฺต สรา สตฺต ปวุฏา\footnote{สปุฏา (ก), ปวุธา (อิ)} สตฺต ปปาตา สตฺต จ ปปาตสตานิ สตฺต สุปินา สตฺต สุปินสตานิ จุลฺลาสีติ มหากปฺปิโน\footnote{มหากปฺปุโน (สี, อิ)} สตสหสฺสานิ.\csromanspacer{} ยานิ พาเล จ ปณฺฑิเต จ \csromanfolio{174}สนฺธาวิตฺวา สํสริตฺวา ทุกฺขสฺสนฺตํ กริสฺสนฺติ.\csromanspacer{} ตตฺถ นตฺถิ “อิมินาหํ สีเลน วา วเตน วา ตเปน วา พฺรหฺมจริเยน วา อปริปกฺกํ วา กมฺมํ ปริปาเจสฺสามิ, ปริปกฺกํ วา กมฺมํ ผุสฺส ผุสฺส พฺยนฺตีกริสฺสามี”ติ, เหวํ นตฺถิ.\csromanspacer{} โทณมิเต สุขทุกฺเข, ปริยนฺตกเต สํสาเร, นตฺถิ หายนวฑฺฒเน, นตฺถิ อุกฺกํสาวกํเส.\csromanspacer{} เสยฺยถาปิ นาม สุตฺตคุเฬ ขิตฺเต นิพฺเพฐิยมานเมว ปเลติ, เอวเมว พาเล จ ปณฺฑิเต จ นิพฺเพฐิยมานา สุขทุกฺขํ ปเลนฺตี”ติ.}

\prose{ภควํมูลกา โน ภนฺเต ธมฺมา ฯเปฯ.\csromanspacer{} รูเป โข ภิกฺขเว สติ รูปํ อุปาทาย รูปํ อภินิวิสฺส เอวํ ทิฏฺฐิ อุปฺปชฺชติ “สตฺติเม กายา อกฏา อกฏวิธา ฯเปฯ สุขทุกฺขํ ปเลนฺตี”ติ.\csromanspacer{} เวทนาย สติ ฯเปฯ.\csromanspacer{} สญฺญาย สติ.\csromanspacer{} สงฺขาเรสุ สติ.\csromanspacer{} วิญฺญาเณ สติ วิญฺญาณํ อุปาทาย วิญฺญาณํ อภินิวิสฺส เอวํ ทิฏฺฐิ อุปฺปชฺชติ “สตฺติเม กายา อกฏา อกฏวิธา ฯเปฯ สุขทุกฺขํ ปเลนฺตี”ติ.\csromanspacer{} ตํ กิํ มญฺญถ ภิกฺขเว, รูปํ นิจฺจํ วา อนิจฺจํ วาติ.\csromanspacer{} อนิจฺจํ ภนฺเต ฯเปฯ ยํ ปนานิจฺจํ ทุกฺขํ วิปริณามธมฺมํ, อปิ นุ ตํ อนุปาทาย เอวํ ทิฏฺฐิ อุปฺปชฺเชยฺย “สตฺติเม กายา อกฏา อกฏวิธา ฯเปฯ สุขทุกฺขํ ปเลนฺตี”ติ.\csromanspacer{} โน เหตํ ภนฺเต.\csromanspacer{} ยมฺปิทํ ทิฏฺฐํ สุตํ มุตํ วิญฺญาตํ ปตฺตํ ปริเยสิตํ อนุวิจริตํ มนสา, ตมฺปิ นิจฺจํ วา อนิจฺจํ วาติ.\csromanspacer{} อนิจฺจํ ภนฺเต ฯเปฯ อปิ นุ ตํ อนุปาทาย เอวํ ทิฏฺฐิ อุปฺปชฺเชยฺย “สตฺติเม กายา อกฏา อกฏวิธา ฯเปฯ นิพฺเพฐิยมานา สุขทุกฺขํ ปเลนฺตี”ติ.\csromanspacer{} โน เหตํ ภนฺเต.}

\prose{ยโต โข ภิกฺขเว อริยสาวกสฺส อิเมสุ จ ฐาเนสุ กงฺขา ปหีนา โหติ, ทุกฺเขปิสฺส กงฺขา ปหีนา โหติ ฯเปฯ ทฺกฺขนิโรธคามินิยา ปฏิปทายปิสฺส กงฺขา ปหีนา โหติ.\csromanspacer{} อยํ วุจฺจติ ภิกฺขเว อริยสาวโก โสตาปนฺโน อวินิปาตธมฺโม นิยโต สมฺโพธิปรายโนติ.\csromanspacer{}\csromanspacer{}อฏฺฐมํ.}

\csromanheader{2}{9. สสฺสตทิฏฺฐิสุตฺต}
\proseitem{214}{สาวตฺถินิทานํ.\csromanspacer{} กิสฺมิํ นุ โข ภิกฺขเว สติ กิํ อุปาทาย กิํ อภินิวิสฺส เอวํ ทิฏฺฐิ อุปฺปชฺชติ “สสฺสโต โลโก”ติ.\csromanspacer{} ภควํมูลกาโน ภนฺเต ธมฺมา ฯเปฯ.\csromanspacer{} รูเป โข ภิกฺขเว สติ รูปํ อุปาทาย รูปํ อภินิวิสฺส เอวํ ทิฏฺฐิ อุปฺปชฺชติ “สสฺสโต โลโก”ติ.\csromanspacer{} เวทนาย สติ ฯเปฯ.\csromanspacer{} สงฺขาเรสุ สติ.}

\csromanheader{2}{3. ทิฏฺฐิสํยุตฺต}
\prose{\csromanfolio{175}วิญฺญาเณ สติ วิญฺญาณํ อุปาทาย วิญฺญาณํ อภินิวิสฺส เอวํ ทิฏฺฐิ อุปฺปชฺชติ “สสฺสโต โลโก”ติ.}

\prose{ตํ กิํ มญฺญถ ภิกฺขเว, รูปํ นิจฺจํ วา อนิจฺจํ วาติ.\csromanspacer{} อนิจฺจํ ภนฺเต ฯเปฯ วิปริณามธมฺมํ, อปิ นุ ตํ อนุปาทาย เอวํ ทิฏฺฐิ อุปฺปชฺเชยฺย “สสฺสโต โลโก”ติ.\csromanspacer{} โน เหตํ ภนฺเต.\csromanspacer{} เวทนา.\csromanspacer{} สญฺญา.\csromanspacer{} สงฺขารา.\csromanspacer{} วิญฺญาณํ นิจฺจํ วา อนิจฺจํ วาติ.\csromanspacer{} อนิจฺจํ ภนฺเต ฯเปฯ อปิ นุ ตํ อนุปาทาย เอวํ ทิฏฺฐิ อุปฺปชฺเชยฺย “สสฺสโต โลโก”ติ.\csromanspacer{} โน เหตํ ภนฺเต.\csromanspacer{} ยมฺปิทํ ทิฏฺฐํ สุตํ มุตํ วิญฺญาตํ ปตฺตํ ปริเยสิตํ อนุวิจริตํ มนสา, ตมฺปิ นิจฺจํ วา อนิจฺจํ วาติ.\csromanspacer{} อนิจฺจํ ภนฺเต.\csromanspacer{} ยํ ปนานิจฺจํ, ทุกฺขํ วา ตํ สุขํ วาติ.\csromanspacer{} ทุกฺขํ ภนฺเต.\csromanspacer{} ยํ ปนานิจฺจํ ทุกฺขํ วิปริณามธมฺมํ, อปิ นุ ตํ อนุปาทาย เอวํ ทิฏฺฐิ อุปฺปชฺเชยฺย “สสฺสโต โลโก”ติ.\csromanspacer{} โน เหตํ ภนฺเต.}

\prose{ยโต โข ภิกฺขเว อริยสาวกสฺส อิเมสุ จ ฐาเนสุ กงฺขา ปหีนา โหติ, ทุกฺเขปิสฺส กงฺขา ปหีนา โหติ ฯเปฯ ทุกฺขนิโรธคามินิยา ปฏิปทายปิสฺส กงฺขา ปหีนา โหติ.\csromanspacer{} อยํ วุจฺจติ ภิกฺขเว อริยสาวโก โสตาปนฺโน อวินิปาตธมฺโม นิยโต สมฺโพธิปรายโนติ.\csromanspacer{}\csromanspacer{}นวมํ.}

\csromanheader{2}{10. อสสฺสตทิฏฺฐิสุตฺต}
\proseitem{215}{สาวตฺถินิทานํ.\csromanspacer{} กิสฺมิํ นุ โข ภิกฺขเว สติ กิํ อุปาทาย กิํ อภินิวิสฺส เอวํ ทิฏฺฐิ อุปฺปชฺชติ “อสสฺสโต โลโก”ติ.\csromanspacer{} ภควํมูลกา โน ภนฺเต ธมฺมา ฯเปฯ.\csromanspacer{} รูเป โข ภิกฺขเว สติ ฯเปฯ.\csromanspacer{} วิญฺญาณํ นิจฺจํ วา อนิจฺจํ วาติ.\csromanspacer{} อนิจฺจํ ภนฺเต ฯเปฯ อปิ นุ ตํ อนุปาทาย เอวํ ทิฏฺฐิ อุปฺปชฺเชยฺย “อสสฺสโต โลโก”ติ.\csromanspacer{} โน เหตํ ภนฺเต.\csromanspacer{} ยมฺปิทํ ทิฏฺฐํ สุตํ มุตํ วิญฺญาตํ ปตฺตํ ปริเยสิตํ อนุวิจริตํ มนสา, ตมฺปิ นิจฺจํ วา อนิจฺจํ วาติ.\csromanspacer{} อนิจฺจํ ภนฺเต ฯเปฯ อปิ นุ ตํ อนุปาทาย เอวํ ทิฏฺฐิ อุปฺปชฺเชยฺย “อสสฺสโต โลโก”ติ.\csromanspacer{} โน เหตํ ภนฺเต.}

\prose{ยโต โข ภิกฺขเว อริยสาวกสฺส อิเมสุ จ ฐาเนสุ กงฺขา ปหีนา โหติ, ทุกฺเขปิสฺส กงฺขา ปหีนา โหตฺ ฯเปฯ ทุกฺขนิโรธคามินิยา ปฏิปทายปิสฺส กงฺขา ปหีนา โหติ.\csromanspacer{} อยํ วุจฺจติ ภิกฺขเว อริยสาวโก โสตาปนฺโน อวินิปาตธมฺโม นิยโต สมฺโพธิปรายโนติ.\csromanspacer{}\csromanspacer{}ทสมํ.}

\csromanheader{2}{11. อนฺตวาสุตฺต}
\proseitem{216}{\csromanfolio{176}สาวตฺถินิทานํ.\csromanspacer{} กิสฺมิํ นุ โข ภิกฺขเว สติ กิํ อุปาทาย กิํ อภินิวิสฺส เอวํ ทิฏฺฐิ อุปฺปชฺชติ “อนฺตวา โลโก”ติ.\csromanspacer{} ภควํมูลกา โน ภนฺเต ธมฺมา ฯเปฯ นิยโต สมฺโพธิปรายโนติ.\csromanspacer{}\csromanspacer{}เอกาทสมํ.}

\csromanheader{2}{12. อนนฺตวาสุตฺต}
\proseitem{217}{สาวตฺถินิทานํ.\csromanspacer{} กิสฺมิํ นุ โข ภิกฺขเว สติ กิํ อุปาทาย กิํ อภินิวิสฺส เอวํ ทิฏฺฐิ อุปฺปชฺชติ “อนนฺตวา โลโก”ติ.\csromanspacer{} ภควํมูลกา โน ภนฺเต ธมฺมา ฯเปฯ นิยโต สมฺโพธิปรายโนติ.\csromanspacer{}\csromanspacer{}ทฺวาทสมํ.}

\csromanheader{2}{13. ตํชีวํตํสรีรํสุตฺต}
\proseitem{218}{สาวตฺถินิทานํ.\csromanspacer{} กิสฺมิํ นุ โข ภิกฺขเว สติ กิํ อุปาทาย กิํ อภินิวิสฺส เอวํ ทิฏฺฐิ อุปฺปชฺชติ “ตํ ชีวํ ตํ สรีรนฺ”ติ.\csromanspacer{} ภควํมูลกา โน ภนฺเต ธมฺมา ฯเปฯ นิยโต สมฺโพธิปรายโนติ.\csromanspacer{}\csromanspacer{}เตรสมํ.}

\csromanheader{2}{14. อญฺญํชีวํ-อญฺญํสรีรํสุตฺต}
\proseitem{219}{สาวตฺถินิทานํ.\csromanspacer{} กิสฺมิํ นุ โข ภิกฺขเว สติ กิํ อุปาทาย กิํ อภินิวิสฺส เอวํ ทิฏฺฐิ อุปฺปชฺชติ “อญฺญํ ชีวํ อญฺญํ สรีรนฺ”ติ.\csromanspacer{} ภควํมูลกา โน ภนฺเต ธมฺมา ฯเปฯ นิยโต สมฺโพธิปรายโนติ.\csromanspacer{}\csromanspacer{}จุทฺทสมํ.}

\csromanheader{2}{15. โหติตถาคโตสุตฺต}
\proseitem{220}{สาวตฺถินิทานํ.\csromanspacer{} กิสฺมิํ นุ โข ภิกฺขเว สติ กิํ อุปาทาย กิํ อภินิวิสฺส เอวํ ทิฏฺฐิ อุปฺปชฺชติ “โหติ ตถาคโต ปรํ มรณา”ติ.\csromanspacer{} ภควํมูลกา โน ภนฺเต ธมฺมา ฯเปฯ นิยโต สมฺโพธิปรายโนติ.\csromanspacer{}\csromanspacer{}ปนฺนรสมํ.}

\csromanheader{2}{16. นโหติตถาคโตสุตฺต}
\proseitem{221}{สาวตฺถินิทานํ.\csromanspacer{} กิสฺมิํ นุ โข ภิกฺขเว สติ กิํ อุปาทาย กิํ อภินิวิสฺส เอวํ ทิฏฺฐิ อุปฺปชฺชติ “น โหติ ตถาคโต ปรํ มรณา”ติ.\csromanspacer{} ภควํมูลกา โน ภนฺเต ธมฺมา ฯเปฯ นิยโต สมฺโพธิปรายโนติ.\csromanspacer{}\csromanspacer{}โสฬสมํ.}

\csromanheader{2}{3. ทิฏฺฐิสํยุตฺต}
\csromanheader{2}{17. โหติจนจโหติตถาคโตสุตฺต}
\proseitem{222}{\csromanfolio{177}สาวตฺถินิทานํ.\csromanspacer{} กิสฺมิํ นุ โข ภิกฺขเว สติ กิํ อุปาทาย กิํ อภินิวิสฺส เอวํ ทิฏฺฐิ อุปฺปชฺชติ “โหติ จ น จ โหติ ตถาคโต ปรํ มรณา”ติ.\csromanspacer{} ภควํมูลกา โน ภนฺเต ธมฺมา ฯเปฯ นิยโต สมฺโพธิปรายโนติ.\csromanspacer{}\csromanspacer{}สตฺตรสมํ.}

\csromanheader{2}{18. เนวโหตินนโหติตถาคโตสุตฺต}
\proseitem{223}{สาวตฺถินิทานํ.\csromanspacer{} กิสฺมิํ นุ โข ภิกฺขเว สติ กิํ อุปาทาย กิํ อภินิวิสฺส เอวํ ทิฏฺฐิ อุปฺปชฺชติ “เนว โหติ น น โหติ ตถาคโต ปรํ มรณา”ติ.\csromanspacer{} ภควํมูลกา โน ภนฺเต ธมฺมา ฯเปฯ.\csromanspacer{} รูเป โข ภิกฺขเว สติ รูปํ อุปาทาย รูปํ อภินิวิสฺส เอวํ ทิฏฺฐิ อุปฺปชฺชติ “เนว โหติ น น โหติ ตถาคโต ปรํ มรณา”ติ ฯเปฯ.}

\prose{ตํ กิํ มญฺญถ ภิกฺขเว, รูปํ นิจฺจํ วา อนิจฺจํ วาติ.\csromanspacer{} อนิจฺจํ ภนฺเต ฯเปฯ วิปริณามธมฺมํ, อปิ นุ ตํ อนุปาทาย เอวํ ทิฏฺฐิ อุปฺปชฺเชยฺย “เนว โหติ น น โหติ ตถาคโต ปรํ มรณา”ติ.\csromanspacer{} โน เหตํ ภนฺเต.\csromanspacer{} ยมฺปิทํ ทิฏฺฐํ สุตํ มุตํ วิญฺญาตํ ปตฺตํ ปริเยสิตํ อนุวิจริตํ มนสา, ตมฺปิ นิจฺจํ วา อนิจฺจํ วาติ.\csromanspacer{} อนิจฺจํ ภนฺเต.\csromanspacer{} ยํ ปนานิจฺจํ, ทุกฺขํ วา ตํ สุขํ วาติ.\csromanspacer{} ทุกฺขํ ภนฺเต.\csromanspacer{} ยํ ปนานิจฺจํ ทุกฺขํ วิปริณามธมฺมํ, อปิ นุ ตํ อนุปาทาย เอวํ ทิฏฺฐิ อุปฺปชฺเชยฺย “เนว โหติ น น โหติ ถถาคโต ปรํ มรณา”ติ.\csromanspacer{} โน เหตํ ภนฺเต.}

\prose{ยโต โข ภิกฺขเว อริยสาวกสฺส อิเมสุ จ ฐาเนสุ กงฺขา ปหีนา โหติ, ทุกฺเขปิสฺส กงฺขา ปหีนา โหติ, ทุกฺขสมุทเยปิสฺส กงฺขา ปหีนา โหติ, ทุกฺขนิโรเธปิสฺส กงฺขา ปหีนา โหติ, ทุกฺขนิโรธคามินิยา ปฏิปทายปิสฺส กงฺขา ปหีนา โหติ.\csromanspacer{} อยํ วุจฺจติ ภิกฺขเว อริยสาวโก โสตาปนฺโน อวินิปาตธมฺโม นิยโต สมฺโพธิปรายโนติ.\csromanspacer{}\csromanspacer{}อฏฺฐารสมํ.}

\vspace{\tipitakaparskip}
\csromancenter{โสตาปตฺติวคฺโค.}
\nitthitam{อฏฺฐารสเวยฺยากรณํ นิฏฺฐิตํ.}
\titlehead{\textbf{ตสฺสุทฺทานํ}}
\csromangathameasure{วาตํ เอตํ มม, โส อตฺตา โน จ เม สิยา. \\ นตฺถิ กโรโต เหตุ จ, มหาทิฏฺเฐน อฏฺฐมํ. \\ สสฺสโต โลโก จ, \\ อสสฺสโต จ อนฺตวา จ. \\ อนนฺตวา จ ตํ ชีวํ ตํ สรีรนฺติ, \\ อญฺญํ ชีวํ อญฺญํ สรีรนฺติ จ. \\ โหติ ตถาคโต ปรํ มรณาติ, \\ น โหติ ตถาคโต ปรํ มรณาติ,}
\csromangathagroup{\csromangathastanza{\csromanfolio{178}วาตํ เอตํ มม, โส อตฺตา โน จ เม สิยา. \\ นตฺถิ กโรโต เหตุ จ, มหาทิฏฺเฐน อฏฺฐมํ.}\csromangathastanza{สสฺสโต โลโก จ, \\ อสสฺสโต จ อนฺตวา จ. \\ อนนฺตวา จ ตํ ชีวํ ตํ สรีรนฺติ, \\ อญฺญํ ชีวํ อญฺญํ สรีรนฺติ จ. \\ โหติ ตถาคโต ปรํ มรณาติ, \\ น โหติ ตถาคโต ปรํ มรณาติ,}}

\csromanheader{2}{เนว โหติ น น โหติ ตถาคโต ปรํ มรณาติ.}
\csromansectionrule
\csromanmatikaanchor{mtk.55}
\csromantocmarkref{section}{2. ทุติยคมนวคฺค}{mtk.55}
\bookmark[dest={mtk.55},level=1]{2. ทุติยคมนวคฺค}
\csromanheader{1}{2. ทุติยคมนวคฺค}
\csromanheader{2}{1. วาตสุตฺต}
\proseitem{224}{สาวตฺถินิทานํ.\csromanspacer{} กิสฺมิํ นุ โข ภิกฺขเว สติ กิํ อุปาทาย กิํ อภินิวิสฺส เอวํ ทิฏฺฐิ อุปฺปชฺชติ “น วาตา วายนฺติ, น นชฺโช สนฺทนฺติ, น คพฺภินิโย วิชายนฺติ, น จนฺทิมสูริยา อุเทนฺติ วา อเปนฺติ วา, เอสิกฏฺฐายิฏฺฐิตา”ติ.\csromanspacer{} ภควํมูลกา โน ภนฺเต ธมฺมา.\csromanspacer{} รูเป โข ภิกฺขเว สติ รูปํ อุปาทาย รูปํ อภินิวิสฺส เอวํ ทิฏฺฐิ อุปฺปชฺชติ “น วาตา วายนฺติ ฯเปฯ เอสิกฏฺฐายิฏฺฐิตา”ติ.\csromanspacer{} เวทนาย สติ ฯเปฯ.\csromanspacer{} สญฺญาย สติ ฯเปฯ.\csromanspacer{} สงฺขาเรสุ สติ.\csromanspacer{} วิญฺญาเณ สติ วิญฺญาณํ อุปาทาย วิญฺญาณํ อภินิวิสฺส เอวํ ทิฏฺฐิ อุปฺปชฺชติ “น วาตา วายนฺติ ฯเปฯ เอสิกฏฺฐายิฏฺฐิตา”ติ.}

\prose{ตํ กิํ มญฺญถ ภิกฺขเว, รูปํ นิจฺจํ วา อนิจฺจํ วาติ.\csromanspacer{} อนิจฺจํ ภนฺเต ฯเปฯ วิปริณามธมฺมํ, อปิ นุ ตํ อนุปาทาย เอวํ ทิฏฺฐิ อุปฺปชฺเชยฺย “น วาตา วายนฺติ ฯเปฯ เอสิกฏฺฐายิฏฺฐิตา”ติ.\csromanspacer{} โน เหตํ ภนฺเต.\csromanspacer{} อิติ โข ภิกฺขเว ทุกฺเข สติ ทุกฺขํ อุปาทาย ทุกฺขํ อภินิวิสฺส เอวํ ทิฏฺฐิ อุปฺปชฺชติ “น วาตา วายนฺติ ฯเปฯ เอสิกฏฺฐายิฏฺฐิตา”ติ.\csromanspacer{} เวทนา.\csromanspacer{} สญฺญา.\csromanspacer{} สงฺขารา.\csromanspacer{} วิญฺญาณํ นิจฺจํ วา อนิจฺจํ วาติ.\csromanspacer{} อนิจฺจํ ภนฺเต ฯเปฯ วิปริณามธมฺมํ, อปิ นุ ตํ อนุปาทาย เอวํ ทิฏฺฐิ อุปฺปชฺเชยฺย “น วาตา วายนฺติ ฯเปฯ เอสิกฏฺฐายิฏฺฐิตา”ติ.\csromanspacer{} โน เหตํ ภนฺเต.\csromanspacer{} อิติ โข ภิกฺขเว ทุกฺเข สติ ทุกฺขํ อุปาทาย ทุกฺขํ อภินิวิสฺส เอวํ ทิฏฺฐิ อุปฺปชฺเชยฺย “น วาตา วายนฺติ, น นชฺโช สนฺทนฺติ, น คพฺภินิโย วิชายนฺติ, น จนฺทิมสูริยา อุเทนฺติ วา อเปนฺติ วา, เอสิกฏฺฐายิฏฺฐิตา”ติ.\csromanspacer{}\csromanspacer{}ปฐมํ.}

\csromanheader{2}{3. ทิฏฺฐิสํยุตฺต}
\proseitem{225-240}{\csromanfolio{179}(ปุริมวคฺเค วิย อฏฺฐารส เวยฺยากรณานิ วิตฺถาเรตพฺพานีติ.) .\csromanspacer{} สตฺตรสมํ.}

\csromanheader{2}{18. เนวโหตินนโหติสุตฺต}
\proseitem{241}{สาวตฺถินิทานํ.\csromanspacer{} กิสฺมิํ นุ โข ภิกฺขเว สติ กิํ อุปาทาย กิํ อภินิวิสฺส เอวํ ทิฏฺฐิ อุปฺปชฺชติ “เนว โหติ น น โหติ ตถาคโต ปรํ มรณา”ติ.\csromanspacer{} ภควํมูลกา โน ภนฺเต ธมฺมา ฯเปฯ.\csromanspacer{} รูเป โข ภิกฺขเว สติ รูปํ อุปาทาย รูปํ อภินิวิสฺส เอวํ ทิฏฺฐิ อุปฺปชฺชติ “เนว โหติ น น โหติ ตถาคโต ปรํ มรณา”ติ.\csromanspacer{} เวทนาย สติ.\csromanspacer{} สญฺญาย สติ.\csromanspacer{} สงฺขาเรสุ สติ.\csromanspacer{} วิญฺญาเณ สติ วิญฺญาณํ อุปาทาย วิญฺญาณํ อภินิวิสฺส เอวํ ทิฏฺฐิ อุปฺปชฺชติ “เนว โหติ น น โหติ ตถาคโต ปรํ มรณา”ติ.}

\prose{ตํ กิํ มญฺญถ ภิกฺขเว, รูปํ นิจฺจํ วา อนิจฺจํ วาติ.\csromanspacer{} อนิจฺจํ ภนฺเต ฯเปฯ วิปริณามธมฺมํ, อปิ นุ ตํ อนุปาทาย เอวํ ทิฏฺฐิ อุปฺปชฺเชยฺย “เนว โหติ น น โหติ ตถาคโต ปรํ มรณา”ติ.\csromanspacer{} โน เหตํ ภนฺเต.\csromanspacer{} อิติ โข ภิกฺขเว ทุกฺเข สติ ทุกฺขํ อุปาทาย ทุกฺขํ อภินิวิสฺส เอวํ ทิฏฺฐิ อุปฺปชฺชติ “เนว โหติ น น โหติ ตถาคโต ปรํ มรณา”ติ.\csromanspacer{} เวทนา.\csromanspacer{} สญฺญา.\csromanspacer{} สงฺขารา.\csromanspacer{} วิญฺญาณํ นิจฺจํ วา อนิจฺจํ วาติ.\csromanspacer{} อนิจฺจํ ภนฺเต ฯเปฯ วิปริณามธมฺมํ, อปิ นุ ตํ อนุปาทาย เอวํ ทิฏฺฐิ อุปฺปชฺเชยฺย “เนว โหติ น น โหติ ตถาคโต ปรํ มรณา”ติ.\csromanspacer{} โน เหตํ ภนฺเต.\csromanspacer{} อิติ โข ภิกฺขเว ทุกฺเข สติ ทุกฺขํ อุปาทาย ทุกฺขํ อภินิวิสฺส เอวํ ทิฏฺฐิ อุปฺปชฺชติ “เนว โหติ น น โหติ ตถาคโต ปรํ มรณา”ติ.\csromanspacer{}\csromanspacer{}อฏฺฐารสมํ.}

\csromanheader{2}{19. รูปี-อตฺตาสุตฺต}
\proseitem{242}{สาวตฺถินิทานํ.\csromanspacer{} กิสฺมิํ นุ โข ภิกฺขเว สติ กิํ อุปาทาย กิํ อภินิวิสฺส เอวํ ทิฏฺฐิ อุปฺปชฺชติ “รูปี อตฺตา โหติ อโรโค ปรํ มรณา”ติ.\csromanspacer{} ภควํมูลกา โน ภนฺเต ธมฺมา ฯเปฯ.\csromanspacer{} รูเป โข ภิกฺขเว สติ รูปํ อุปาทาย รูปํ อภินิวิสฺส เอวํ ทิฏฺฐิ อุปฺปชฺชติ “รูปี อตฺตา โหติ อโรโค ปรํ มรณา”ติ.\csromanspacer{} เวทนาย สติ.\csromanspacer{} สญฺญาย สติ.\csromanspacer{} สงฺขาเรสุ สติ.\csromanspacer{} วิญฺญาเณ สติ วิญฺญาณํ อุปาทาย วิญฺญาณํ อภินิวิสฺส เอวํ ทิฏฺฐิ อุปฺปชฺชติ “รูปี อตฺตา โหติ อโรโค ปรํ มรณา”ติ.}

\prose{\csromanfolio{180}ตํ กิํ มญฺญถ ภิกฺขเว, รูปํ นิจฺจํ วา อนิจฺจํ วาติ.\csromanspacer{} อนิจฺจํ ภนฺเต ฯเปฯ วิปริณามธมฺมํ, อปิ นุ ตํ อนุปาทาย เอวํ ทิฏฺฐิ อุปฺปชฺเชยฺย “รูปี อตฺตา โหติ อโรโค ปรํ มรณา”ติ.\csromanspacer{} โน เหตํ ภนฺเต.\csromanspacer{} อิติ โข ภิกฺขเว ทุกฺเข สติ ทุกฺขํ อุปาทาย ทุกฺขํ อภินิวิสฺส เอวํ ทิฏฺฐิ อุปฺปชฺชติ “รูปี อตฺตา โหติ อโรโค ปรํ มรณา”ติ.\csromanspacer{} เวทนา ฯเปฯ.\csromanspacer{} โน เหตํ ภนฺเต.\csromanspacer{} อิติ โข ภิกฺขเว ทุกฺเข สติ ทุกฺขํ อุปาทาย ทุกฺขํ อภินิวิสฺส เอวํ ทิฏฺฐิ อุปฺปชฺชติ “รูปี อตฺตา โหติ อโรโค ปรํ มรณา”ติ.\csromanspacer{}\csromanspacer{}เอกูนวีสติมํ.}

\csromanheader{2}{20. อรูปี-อตฺตาสุตฺต}
\proseitem{243}{สาวตฺถินิทานํ.\csromanspacer{} กิสฺมิํ นุ โข ภิกฺขเว สติ กิํ อุปาทาย กิํ อภินิวิสฺส เอวํ ทิฏฺฐิ อุปฺปชฺชติ “อรูปี อตฺตา โหติ อโรโค ปรํ มรณา”ติ. (เปยฺยาโล.).\csromanspacer{} วีสติมํ.}

\csromanheader{2}{21. รูปีจ-อรูปีจ-อตฺตาสุตฺต}
\proseitem{244}{สาวตฺถินิทานํ.\csromanspacer{} รูปี จ อรูปี จ อตฺตา โหติ อโรโค ปรํ มรณาติ.\csromanspacer{}\csromanspacer{}เอกวีสติมํ.}

\csromanheader{2}{22. เนวรูปีนารูปี-อตฺตาสุตฺต}
\vspace{\tipitakaparskip}
\csromancenter{เนว รูปี นารูปี อตฺตา โหติ อโรโค ปรํ มรณาติ.\csromanspacer{}\csromanspacer{}พาวีสติมํ.}
\csromanheader{2}{23. เอกนฺตสุขีสุตฺต}
\vspace{\tipitakaparskip}
\csromancenter{เอกนฺตสุขี อตฺตา โหติ อโรโค ปรํ มรณาติ.\csromanspacer{}\csromanspacer{}เตวีสติมํ.}
\csromanheader{2}{24. เอกนฺตทุกฺขีสุตฺต}
\vspace{\tipitakaparskip}
\csromancenter{เอกนฺตทุกฺขี อตฺตา โหติ อโรโค ปรํ มรณาติ.\csromanspacer{}\csromanspacer{}จตุวีสติมํ.}
\csromanheader{2}{25. สุขทุกฺขีสุตฺต}
\vspace{\tipitakaparskip}
\csromancenter{สุขทุกฺขี อตฺตา โหติ อโรโค ปรํ มรณาติ.\csromanspacer{}\csromanspacer{}ปญฺจวีสติมํ.}
\csromanheader{2}{3. ทิฏฺฐิสํยุตฺต}
\csromanheader{2}{26. อทุกฺขมสุขีสุตฺต}
\proseitem{249}{\csromanfolio{181}อทุกฺขมสุขี อตฺตา โหติ อโรโค ปรํ มรณาติ.\csromanspacer{} ภควํมูลกา โน ภนฺเต ธมฺมา ฯเปฯ.\csromanspacer{} รูเป โข ภิกฺขเว สติ รูปํ อุปาทาย รูปํ อภินิวิสฺส เอวํ ทิฏฺฐิ อุปฺปชฺชติ “อทุกฺขมสุขี อตฺตา โหติ อโรโค ปรํ มรณา”ติ.\csromanspacer{} เวทนาย สติ.\csromanspacer{} สญฺญาย สติ.\csromanspacer{} สงฺขาเรสุ สติ.\csromanspacer{} วิญฺญาเณ สติ วิญฺญาณํ อุปาทาย วิญฺญาณํ อภินิวิสฺส เอวํ ทิฏฺฐิ อุปฺปชฺชติ “อทุกฺขมสุขี อตฺตา โหติ อโรโค ปรํ มรณา”ติ.}

\prose{ตํ กิํ มญฺญถ ภิกฺขเว, รูปํ นิจฺจํ วา อนิจฺจํ วาติ.\csromanspacer{} อนิจฺจํ ภนฺเต ฯเปฯ วิปริณามธมฺมํ, อปิ นุ ตํ อนุปาทาย เอวํ ทิฏฺฐิ อุปฺปชฺเชยฺย “อทุกฺขมสุขี อตฺตา โหติ อโรโค ปรํ มรณา”ติ.\csromanspacer{} โน เหตํ ภนฺเต.\csromanspacer{} อิติ โข ภิกฺขเว ทุกฺเข สติ ทุกฺขํ อุปาทาย ทุกฺขํ อภินิวิสฺส เอวํ ทิฏฺฐิ อุปฺปชฺชติ “อทุกฺขมสุขี อตฺตา โหติ อโรโค ปรํ มรณา”ติ.\csromanspacer{} เวทนา.\csromanspacer{} สญฺญา.\csromanspacer{} สงฺขารา.\csromanspacer{} วิญฺญาณํ นิจฺจํ วา อนิจฺจํ วาติ.\csromanspacer{} อนิจฺจํ ภนฺเต ฯเปฯ วิปริณามธมฺมํ, อปิ นุ ตํ อนุปาทาย เอวํ ทิฏฺฐิ อุปฺปชฺเชยฺย “อทุกฺขมสุขี อตฺตา โหติ อโรโค ปรํ มรณา”ติ.\csromanspacer{} โน เหตํ ภนฺเต.\csromanspacer{} อิติ โข ภิกฺขเว ทุกฺเข สติ ทุกฺขํ อุปาทาย ทุกฺขํ อภินิวิสฺส เอวํ ทิฏฺฐิ อุปฺปชฺชติ “อทุกฺขมสุขี อตฺตา โหติ อโรโค ปรํ มรณา”ติ.\csromanspacer{}\csromanspacer{}ฉพฺพีสติมํ.\csromanspacer{} ทุติยเปยฺยาโล.}

\titlehead{\textbf{ตสฺสุทฺทานํ}}
\csromangathameasure{วาตํ เอตํ มม, โส อตฺตา โน จ เม สิยา. \\ นตฺถิ กโรโต เหตุ จ, มหาทิฏฺเฐน อฏฺฐมํ. \\ สสฺสโต อสสฺสโต เจว, อนฺตานนฺตวา จ วุจฺจติ. \\ ตํ ชีวํ อญฺญํ ชีวญฺจ, ตถาคเตน จตฺตาโร. \\ รูปี อตฺตา โหติ อรูปี จ อตฺตา โหติ, \\ รูปี จ อรูปี จ อตฺตา โหติ. \\ เนว รูปี นารูปี อตฺตา โหติ, \\ เอกนฺตสุขี อตฺตา โหติ. \\ เอกนฺตทุกฺขี อตฺตา โหติ, สุขทุกฺขี อตฺตา โหติ. \\ อทุกฺขมสุขี อตฺตา โหติ, อโรโค ปรํ มรณาติ. \\ อิเม ฉพฺพีสติ สุตฺตา, ทุติยวาเรน เทสิตา.}
\csromangathagroup{\csromangathastanza{วาตํ เอตํ มม, โส อตฺตา โน จ เม สิยา. \\ นตฺถิ กโรโต เหตุ จ, มหาทิฏฺเฐน อฏฺฐมํ.}\csromangathastanza{สสฺสโต อสสฺสโต เจว, อนฺตานนฺตวา จ วุจฺจติ. \\ ตํ ชีวํ อญฺญํ ชีวญฺจ, ตถาคเตน จตฺตาโร.}\csromangathastanza{รูปี อตฺตา โหติ อรูปี จ อตฺตา โหติ, \\ รูปี จ อรูปี จ อตฺตา โหติ. \\ เนว รูปี นารูปี อตฺตา โหติ, \\ เอกนฺตสุขี อตฺตา โหติ.}\csromangathastanza{เอกนฺตทุกฺขี อตฺตา โหติ, สุขทุกฺขี อตฺตา โหติ. \\ อทุกฺขมสุขี อตฺตา โหติ, อโรโค ปรํ มรณาติ. \\ อิเม ฉพฺพีสติ สุตฺตา, ทุติยวาเรน เทสิตา.}}

\csromanmatikaanchor{mtk.56}
\csromantocmarkref{section}{3. ตติยคมนวคฺค}{mtk.56}
\bookmark[dest={mtk.56},level=1]{3. ตติยคมนวคฺค}
\csromanheader{1}{3. ตติยคมนวคฺค}
\csromanheader{2}{1. นวาตสุตฺต}
\proseitem{250}{\csromanfolio{182}สาวตฺถินิทานํ.\csromanspacer{} กิสฺมิํ นุ โข ภิกฺขเว สติ กิํ อุปาทาย กิํ อภินิวิสฺส เอวํ ทิฏฺฐิ อุปฺปชฺชติ “น วาตา วายนฺติ, น นชฺโช สนฺทนฺติ, น คพฺภินิโย วิชายนฺติ, น จนฺทิมสูริยา อุเทนฺติ วา อเปนฺติ วา, เอสิกฏฺฐายิฏฺฐิตา”ติ.\csromanspacer{} ภควํมูลกา โน ภนฺเต ธมฺมา ฯเปฯ.}

\prose{รูเป โข ภิกฺขเว สติ รูปํ อุปาทาย รูปํ อภินิวิสฺส เอวํ ทิฏฺฐิ อุปฺปชฺชติ “น วาตา วายนฺติ ฯเปฯ.\csromanspacer{} เวทนาย สติ.\csromanspacer{} สญฺญาย สติ.\csromanspacer{} สงฺขาเรสุ สติ.\csromanspacer{} วิญฺญาเณ สติ วิญฺญาณํ อุปาทาย วิญฺญาณํ อภินิวิสฺส เอวํ ทิฏฺฐิ อุปฺปชฺชติ “น วาตา วายนฺติ ฯเปฯ เอสิกฏฺฐายิฏฺฐิตา”ติ.}

\prose{ตํ กิํ มญฺญถ ภิกฺขเว, รูปํ นิจฺจํ วา อนิจฺจํ วาติ.\csromanspacer{} อนิจฺจํ ภนฺเต ฯเปฯ วิปริณามธมฺมํ, อปิ นุ ตํ อนุปาทาย เอวํ ทิฏฺฐิ อุปฺปชฺเชยฺย “น วาตา วายนฺติ ฯเปฯ เอสิกฏฺฐายิฏฺฐิตา”ติ.\csromanspacer{} โน เหตํ ภนฺเต.\csromanspacer{} อิติ โข ภิกฺขเว ยทนิจฺจํ, ตํ ทุกฺขํ, ตสฺมิํ สติ ตทุปาทาย เอวํ ทิฏฺฐิ อุปฺปชฺชติ “น วาตา วายนฺติ, น นชฺโช สนฺทนฺติ, น คพฺภินิโย วิชายนฺติ, น จนฺทิมสูริยา อุเทนฺติ วา อเปนฺติ วา, เอสิกฏฺฐายิฏฺฐิตา”ติ.\csromanspacer{} เวทนา.\csromanspacer{} สญฺญา.\csromanspacer{} สงฺขารา.\csromanspacer{} วิญฺญาณํ นิจฺจํ วา อนิจฺจํ วาติ.\csromanspacer{} อนิจฺจํ ภนฺเต ฯเปฯ วิปริณามธมฺมํ, อปิ นุ ตํ อนุปาทาย เอวํ ทิฏฺฐิ อุปฺปชฺเชยฺย “น วาตา วายนฺติ ฯเปฯ เอสิกฏฺฐายิฏฺฐิตา”ติ.\csromanspacer{} โน เหตํ ภนฺเต.\csromanspacer{} อิติ โข ภิกฺขเว ยทนิจฺจํ, ตํ ทุกฺขํ, ตสฺมิํ สติ ตทุปาทาย เอวํ ทิฏฺฐิ อุปฺปชฺชติ “น วาตา วายนฺติ ฯเปฯ เอสิกฏฺฐายิฏฺฐิตา”ติ.\csromanspacer{}\csromanspacer{}ปฐมํ.}

\proseitem{251-274}{(ทุติยวคฺเค วิย จตุวีสติ สุตฺตานิ ปูเรตพฺพานิ.).\csromanspacer{} ปญฺจวีสติมํ.}

\csromanheader{2}{26. อทุกฺขมสุขีสุตฺต}
\proseitem{275}{สาวตฺถินิทานํ.\csromanspacer{} กิสฺมิํ นุ โข ภิกฺขเว สติ กิํ อุปาทาย กิํ อภินิวิสฺส เอวํ ทิฏฺฐิ อุปฺปชฺชติ “อทุกฺขมสุขี อตฺตา โหติ อโรโค ปรํ มรณา”ติ.\csromanspacer{} ภควํมูลกา โน ภนฺเต ธมฺมา ฯเปฯ.}

\prose{รูเป โข ภิกฺขเว สติ รูปํ อุปาทาย รูปํ อภินิวิสฺส เอวํ ทิฏฺฐิ อุปฺปชฺชติ “อทุกฺขมสุขี อตฺตา โหติ อโรโค ปรํ มรณา”ติ.\csromanspacer{} เวทนาย}

\csromanheader{2}{3. ทิฏฺฐิสํยุตฺต}
\prose{\csromanfolio{183}สติ ฯเปฯ.\csromanspacer{} สญฺญาย สติ.\csromanspacer{} สงฺขาเรสุ สติ.\csromanspacer{} วิญฺญาเณ สติ วิญฺญาณํ อุปาทาย วิญฺญาณํ อภินิวิสฺส เอวํ ทิฏฺฐิ อุปฺปชฺชติ “อทุกฺขมสุขี อตฺตา โหติ อโรโค ปรํ มรณา”ติ.}

\prose{ตํ กิํ มญฺญถ ภิกฺขเว, รูปํ นิจฺจํ วา อนิจฺจํ วาติ.\csromanspacer{} อนิจฺจํ ภนฺเต ฯเปฯ วิปริณามธมฺมํ, อปิ นุ ตํ อนุปาทาย เอวํ ทิฏฺฐิ อุปฺปชฺเชยฺย “อทุกฺขมสุขี อตฺตา โหติ อโรโค ปรํ มรณา”ติ.\csromanspacer{} โน เหตํ ภนฺเต.\csromanspacer{} อิติ โข ภิกฺขเว ยทนิจฺจํ, ตํ ทุกฺขํ, ตสฺมิํ สติ ตทุปาทาย เอวํ ทิฏฺฐิ อุปฺปชฺชติ “อทุกฺขมสุขี อตฺตา โหติ อโรโค ปรํ มรณา”ติ.\csromanspacer{} เวทนา.\csromanspacer{} สญฺญา.\csromanspacer{} สงฺขารา.\csromanspacer{} วิญฺญาณํ นิจฺจํ วา อนิจฺจํ วาติ.\csromanspacer{} อนิจฺจํ ภนฺเต ฯเปฯ วิปริณามธมฺมํ, อปิ นุ ตํ อนุปาทาย เอวํ ทิฏฺฐิ อุปฺปชฺเชยฺย “อทุกฺขมสุขี อตฺตา โหติ อโรโค ปรํ มรณา”ติ.\csromanspacer{} โน เหตํ ภนฺเต.\csromanspacer{} อิติ โข ภิกฺขเว ยทนิจฺจํ, ตํ ทุกฺขํ, ตสฺมิํ สติ ตทุปาทาย เอวํ อุปฺปชฺชติ “อทุกฺขมสุขี อตฺตา โหติ อโรโค ปรํ มรณา”ติ.\csromanspacer{}\csromanspacer{}ฉพฺพีสติมํ.\csromanspacer{} ตติยเปยฺยาโล.}
\csromansectionrule

\csromanmatikaanchor{mtk.57}
\csromantocmarkref{section}{4. จตุตฺถคมนวคฺค}{mtk.57}
\bookmark[dest={mtk.57},level=1]{4. จตุตฺถคมนวคฺค}
\csromanheader{1}{4. จตุตฺถคมนวคฺค}
\csromanheader{2}{1. นวาตสุตฺต}
\proseitem{276}{สาวตฺถินิทานํ.\csromanspacer{} กิสฺมิํ นุ โข ภิกฺขเว สติ กิํ อุปาทาย กิํ อภินิวิสฺส เอวํ ทิฏฺฐิ อุปฺปชฺชติ “น วาตา วายนฺติ, น นชฺโช สนฺทนฺติ, น คพฺภินิโย วิชายนฺติ, น จนฺทิมสูริยา อุเทนฺติ วา อเปนฺติ วา, เอสิกฏฺฐายิฏฺฐิตา”ติ.\csromanspacer{} ภควํมูลกา โน ภนฺเต ธมฺมาฯเปฯ.}

\prose{รูเป โข ภิกฺขเว สติ รูปํ อุปาทาย รูปํ อภินิวิสฺส เอวํ ทิฏฺฐิ อุปฺปชฺชติ “น วาตา วายนฺติ ฯเปฯ เอสิกฏฺฐายิฏฺฐิตา”ติ.\csromanspacer{} เวทนาย สติ ฯเปฯ.\csromanspacer{} สญฺญาย สติ.\csromanspacer{} สงฺขาเรสุ สติ.\csromanspacer{} วิญฺญาเณ สติ วิญฺญาณํ อุปาทาย วิญฺญาณํ อภินิวิสฺส เอวํ ทิฏฺฐิ อุปฺปชฺชติ “น วาตา วายนฺติ ฯเปฯ เอสิกฏฺฐายิฏฺฐิตา”ติ.\csromanspacer{} ตํ กิํ มญฺญถ ภิกฺขเว, รูปํ นิจฺจํ วา อนิจฺจํ วาติ.\csromanspacer{} อนิจฺจํ ภนฺเต.\csromanspacer{} ยํ ปนานิจฺจํ, ทุกฺขํ วา ตํ สุขํ วาติ.\csromanspacer{} ทุกฺขํ ภนฺเต.\csromanspacer{} ยํ ปนานิจฺจํ ทุกฺขํ วิปริณามธมฺมํ, กลฺลํ นุ ตํ สมนุปสฺสิตุํ “เอตํ มม, เอโสหมสฺมิ, \csromanfolio{184}เอโส เม อตฺตา”ติ.\csromanspacer{} โน เหตํ ภนฺเต.\csromanspacer{} เวทนา.\csromanspacer{} สญฺญา.\csromanspacer{} สงฺขารา.\csromanspacer{} วิญฺญาณํ นิจฺจํ วา อนิจฺจํ วาติ.\csromanspacer{} อนิจฺจํ ภนฺเต.\csromanspacer{} ยํ ปนานิจฺจํ, ทุกฺขํ วา ตํ สุขํ วาติ.\csromanspacer{} ทุกฺขํ ภนฺเต.\csromanspacer{} ยํ ปนานิจฺจํ ทุกฺขํ วิปริณามธมฺมํ, กลฺลํ นุ ตํ สมนุปสฺสิตุํ “เอตํ มม, เอโสหมสฺมิ, เอโส เม อตฺตา”ติ.\csromanspacer{} โน เหตํ ภนฺเต.}

\prose{ตสฺมาติห ภิกฺขเว ยํ กิญฺจิ รูปํ อตีตานาคตปจฺจุปฺปนฺนํ อชฺฌตฺตํ วา พหิทฺธา วา โอฬาริกํ วา สุขุมํ วา หีนํ วา ปณีตํ วา ยํ ทูเร สนฺติเก วา, สพฺพํ รูปํ “เนตํ มม, เนโสหมสฺมิ, น เมโส อตฺตา”ติ เอวเมตํ ยถาภูตํ สมฺมปฺปญฺญาย ทฏฺฐพฺพํ.\csromanspacer{} ยา กาจิ เวทนา.\csromanspacer{} ยา กาจิ สญฺญา.\csromanspacer{} เย เกจิ สงฺขารา.\csromanspacer{} ยํ กิญฺจิ วิญฺญาณํ อตีตานาคตปจฺจุปฺปนฺนํ อชฺฌตฺตํ วา พหิทฺธา วา โอฬาริกํ วา สุขุมํ วา หีนํ วา ปณีตํ วา ยํ ทูเร สนฺติเก วา, สพฺพํ วิญฺญาณํ “เนตํ มม, เนโสหมสฺมิ, น เมโส อตฺตา”ติ เอวเมตํ ยถาภูตํ สมฺมปฺปญฺญาย ทฏฺฐพฺพํ.\csromanspacer{} เอวํ ปสฺสํ ฯเปฯ นาปรํ อิตฺถตฺตายาติ ปชานาตีติ.\csromanspacer{}\csromanspacer{}ปฐมํ.}

\proseitem{277-300}{(ทุติยวคฺเค วิย จตุวีสติ สุตฺตานิ ปูเรตพฺพานิ.).\csromanspacer{} ปญฺจวีสติมํ.}

\csromanheader{2}{26. อทุกฺขมสุขีสุตฺต}
\proseitem{301}{สาวตฺถินิทานํ.\csromanspacer{} กิสฺมิํ นุ โข ภิกฺขเว สติ กิํ อุปาทาย กิํ อภินิวิสฺส เอวํ ทิฏฺฐิ อุปฺปชฺชติ “อทุกฺขมสุขี อตฺตา โหติ อโรโค ปรํ มรณา”ติ.\csromanspacer{} ภควํมูลกา โน ภนฺเต ธมฺมา ฯเปฯ.}

\prose{รูเป โข ภิกฺขเว สติ รูปํ อุปาทาย รูปํ อภินิวิสฺส เอวํ ทิฏฺฐิ อุปฺปชฺชติ “อทุกฺขมสุขี อตฺตา โหติ อโรโค ปรํ มรณา”ติ.\csromanspacer{} เวทนาย สติ.\csromanspacer{} สญฺญาย สติ.\csromanspacer{} สงฺขาเรสุ สติ.\csromanspacer{} วิญฺญาเณ สติ วิญฺญาณํ อุปาทาย วิญฺญาณํ อภินิวิสฺส เอวํ ทิฏฺฐิ อุปฺปชฺชติ “อทุกฺขมสุขี อตฺตา โหติ อโรโค ปรํ มรณา”ติ.}

\prose{ตํ กิํ มญฺญถ ภิกฺขเว, รูปํ นิจฺจํ วา อนิจฺจํ วาติ.\csromanspacer{} อนิจฺจํ ภนฺเต.\csromanspacer{} ยํ ปนานิจฺจํ, ทุกฺขํ วา ตํ สุขํ วาติ.\csromanspacer{} ทุกฺขํ ภนฺเต.\csromanspacer{} ยํ ปนานิจฺจํ ทุกฺขํ วิปริณามธมฺมํ, กลฺลํ นุ ตํ สมนุปสฺสิตุํ “เอตํ มม, เอโสหมสฺมิ, เอโส เม อตฺตา”ติ.\csromanspacer{} โน เหตํ ภนฺเต.\csromanspacer{} เวทนา.\csromanspacer{} สญฺญา.\csromanspacer{} สงฺขารา.\csromanspacer{} วิญฺญาณํ นิจฺจํ วา อนิจฺจํ วาติ.\csromanspacer{} อนิจฺจํ ภนฺเต.\csromanspacer{} ยํ ปนานิจฺจํ, ทุกฺขํ วา ตํ สุขํ วาติ.\csromanspacer{} ทุกฺขํ ภนฺเต.\csromanspacer{} ยํ ปนานิจฺจํ}

\csromanheader{2}{3. ทิฏฺฐิสํยุตฺต}
\prose{\csromanfolio{185}ทุกฺขํ วิปริณามธมฺมํ, กลฺลํ นุ ตํ สมนุปสฺสิตุํ “เอตํ มม, เอโสหมสฺมิ, เอโส เม อตฺตา”ติ.\csromanspacer{} โน เหตํ ภนฺเต.}

\prose{ตสฺมาติห ภิกฺขเว ยํ กิญฺจิ รูปํ อตีตานาคตปจฺจุปฺปนฺนํ อชฺฌตฺตํ วา พหิทฺธา วา โอฬาริกํ วา สุขุมํ วา หีนํ วา ปณีตํ วา ยํ ทูเร สนฺติเก วา, สพฺพํ รูปํ “เนตํ มม, เนโสหมสฺมิ, น เมโส อตฺตา”ติ เอวเมตํ ยถาภูตํ สมฺมปฺปญฺญาย ทฏฺฐพฺพํ.\csromanspacer{} ยา กาจิ เวทนา.\csromanspacer{} ยา กาจิ สญฺญา.\csromanspacer{} เย เกจิ สงฺขารา.\csromanspacer{} ยํ กิญฺจิ วิญฺญาณํ อตีตานาคตปจฺจุปฺปนฺนํ อชฺฌตฺตํ วา พหิทฺธา วา โอฬาริกํ วา สุขุมํ วา หีนํ วา ปณีตํ วา ยํ ทูเร สนฺติเก วา, สพฺพํ วิญฺญาณํ “เนตํ มม, เนโสหมสฺมิ, น เมโส อตฺตา”ติ เอวเมตํ ยถาภูตํ สมฺมปฺปญฺญาย ทฏฺฐพฺพํ.}

\prose{เอวํ ปสฺสํ ภิกฺขเว สุตวา อริยสาวโก รูปสฺมิมฺปิ นิพฺพินฺทติ, เวทนายปิ นิพฺพินฺทติ, สญฺญายปิ นพฺพินฺทติ, สงฺขาเรสุปิ นิพฺพินฺทติ, วิญฺญาณสฺมิมฺปิ นิพฺพินฺทติ, นิพฺพินฺทํ วิรชฺชติ, วิราคา วิมุจฺจติ, วิมุตฺตสฺมิํ “วิมุตฺตมฺ”อิติ ญาณํ โหติ, “ขีณา ชาติ, วุสิตํ พฺรหฺมจริยํ, กตํ กรณียํ, นาปรํ อิตฺถตฺตายา”ติ ปชานาตีติ.}

\csromangathameasure{ปุริมคมเน อฏฺฐารส เวยฺยากรณา, \\ ทุติยคมเน ฉพฺพีสํ วิตฺถาเรตพฺพานิ. \\ ตติยคมเน ฉพฺพีสํ วิตฺถาเรตพฺพานิ, \\ จตุตฺถคมเน ฉพฺพีสํ วิตฺถาเรตพฺพานิ.}
\csromangathagroup{\csromangathastanza{ปุริมคมเน อฏฺฐารส เวยฺยากรณา, \\ ทุติยคมเน ฉพฺพีสํ วิตฺถาเรตพฺพานิ. \\ ตติยคมเน ฉพฺพีสํ วิตฺถาเรตพฺพานิ, \\ จตุตฺถคมเน ฉพฺพีสํ วิตฺถาเรตพฺพานิ.}}

\vspace{\tipitakaparskip}
\csromancenter{\textbf{ทิฏฺฐิสํยุตฺตํ สมตฺตํ.}}
\csromanmatikaanchor{mtk.58}
\csromantocmarknopage{chapter}{4. โอกฺกนฺตสํยุตฺต}{mtk.58}
\bookmark[dest={mtk.58},level=0]{4. โอกฺกนฺตสํยุตฺต}
\chapterheadpageread{4. โอกฺกนฺตสํยุตฺต}
\csromanheader{2}{1. จกฺขุสุตฺต}
\csromanmatikaanchor{mtk.59}
\csromantocmarkref{section}{1-10. จกฺขุสุตฺตาทิ}{mtk.59}
\bookmark[dest={mtk.59},level=1]{1-10. จกฺขุสุตฺตาทิ}
\proseitem{302}{\csromanfolio{186}สาวตฺถินิทานํ.\csromanspacer{} จกฺขุํ ภิกฺขเว อนิจฺจํ วิปริณามิ อญฺญถาภาวิ, โสตํ อนิจฺจํ วิปริณามิ อญฺญถาภาวิ, ฆานํ อนิจฺจํ วิปริณามิ อญฺญถาภาวิ, ชิวฺหา อนิจฺจา วิปริณามี อญฺญถาภาวี\footnote{วิปริณามินี อญฺญถาภาวินี (?)}, กาโย อนิจฺโจ วิปริณามิ อญฺญถภาวิ, มโน อนิจฺโจ วิปริณามี อญฺญถาภาวี.\csromanspacer{} โย ภิกฺขเว อิเม ธมฺเม เอวํ สทฺทหติ อธิมุจฺจติ.\csromanspacer{} อยํ วุจฺจติ สทฺธานุสารี โอกฺกนฺโต สมฺมตฺตนิยามํ, สปฺปุริสภูมิํ โอกฺกนฺโต, วีติวตฺโต ปุถุชฺชนภูมิํ, อภพฺโพ ตํ กมฺมํ กาตุํ, ยํ กมฺมํ กตฺวา นิรยํ วา ติรจฺฉานโยนิํ วา เปตฺติวิสยํ วา อุปปชฺเชยฺย, อภพฺโพ จ\footnote{อภพฺโพว (สี, สฺยา, กํ)} ตาว กาลํ กาตุํ, ยาว น โสตาปตฺติผลํ สจฺฉิกโรติ.}

\prose{ยสฺส โข ภิกฺขเว อิเม ธมฺมา เอวํ ปญฺญาย มตฺตโส นิชฺฌานํ ขมนฺติ.\csromanspacer{} อยํ วุจฺจติ ธมฺมานุสารี โอกฺกนฺโต สมฺมตฺตนิยามํ, สปฺปุริสภูมิํ โอกฺกนฺโต, วีติวตฺโต ปุถุชฺชนภูมิํ, อภพฺโพ ตํ กมฺมํ กาตุํ, ยํ กมฺมํ กตฺวา นิรยํ วา ติรจฺฉานโยนิํ วา เปตฺติวิสยํ วา อุปปชฺเชยฺย, อภพฺโพ จ ตาว กาลํ กาตุํ, ยาว น โสตาปตฺติผลํ สจฺฉิกโรติ.\csromanspacer{} โย ภิกฺขเว อิเม ธมฺเม เอวํ ปชานาติ เอวํ ปสฺสติ, อยํ วุจฺจติ โสตาปนฺโน อวินิปาตธมฺโม นิยโต สมฺโพธิปรายโนติ.\csromanspacer{}\csromanspacer{}ปฐมํ.}

\csromanheader{2}{2. รูปสุตฺต}
\proseitem{303}{สาวตฺถินิทานํ.\csromanspacer{} รูปา ภิกฺขเว อนิจฺจา วิปริณามิโน อญฺญถา ภาวิโน, สทฺทา อนิจฺจา วิปริณามิโน อญฺญถาภาวิโน, คนฺธา อนิจฺจา วิปริณามิโน อญฺญถาภาวิโน, รสา อนิจฺจา วิปริณามิโน อญฺญถาภาวิโน, โผฏฺฐพฺพา อนิจฺจา วิปริณามิโน อญฺญถาภาวิโน ธมฺมา อนิจฺจา วิปริณามิโน อญฺญถาภาวิโน.\csromanspacer{} โย ภิกฺขเว อิเม ธมฺเม เอวํ สทฺทหติ อธิมุจฺจติ.\csromanspacer{} อยํ วุจฺจติ สทฺธานุสารี โอกฺกนฺโต สมฺมตฺตนิยามํ, สปฺปุริสภูมิํ โอกฺกนฺโต, วีติวตฺโต ปุถุชฺชนภูมิํ, อภพฺโพ ตํ กมฺมํ กาตุํ, ยํ กมฺมํ กตฺวา นิรยํ วา ติรจฺฉานโยนิํ วา}

\csromanheader{2}{4. โอกฺกนฺตสํยุตฺต}
\prose{\csromanfolio{187}เปตฺติวิสยํ วา อุปปชฺเชยฺย, อภพฺโพ จ ตาว กาลํ กาตุํ, ยาว น โสตาปตฺติผลํ สจฺฉิกโรติ.}

\prose{ยสฺส โข ภิกฺขเว อิเม ธมฺมา เอวํ ปญฺญาย มตฺตโส นิชฺฌานํ ขมนฺติ.\csromanspacer{} อยํ วุจฺจติ ธมฺมานุสารี โอกฺกนฺโต สมฺมตฺตนิยามํ, สปฺปุริสภูมิํ โอกฺกนฺโต, วีติวตฺโต ปุถุชฺชนภูมิํ, อภพฺโพ ตํ กมฺมํ กาตุํ, ยํ กมฺมํ กตฺวา นิรยํ วา ติรจฺฉานโยนิํ วา เปตฺติวิสยํ วา อุปฺปชฺเชยฺย, อภพฺโพ จ ตาว กาลํ กาตุํ, ยาว น โสตาปตฺติผลํ สจฺฉิกโรติ.\csromanspacer{} โย ภิกฺขเว อิเม ธมฺเม เอวํ ปชานาติ เอวํ ปสฺสติ, อยํ วุจฺจติ โสตาปนฺโน อวินิปาตธมฺโม นิยโต สมฺโพธิปรายโนติ.\csromanspacer{}\csromanspacer{}ทุติยํ.}

\csromanheader{2}{3. วิญฺญาณสุตฺต}
\proseitem{304}{สาวตฺถินิทานํ.\csromanspacer{} จกฺขุวิญฺญาณํ ภิกฺขเว อนิจฺจํ วิปริณามิ อญฺญถาภาวิ.\csromanspacer{} โสตวิญฺญาณํ.\csromanspacer{} ฆานวิญฺญาณํ.\csromanspacer{} ชิวฺหาวิญฺญาณํ.\csromanspacer{} กายวิญฺญาณํ.\csromanspacer{} มโนวิญฺญาณํ อนิจฺจํ วิปริณามิ อญฺญถาภาวิ.\csromanspacer{} โย ภิกฺขเว ฯเปฯ สมฺโพธิปรายโนติ.\csromanspacer{}\csromanspacer{}ตติยํ.}

\csromanheader{2}{4. สมฺผสฺสสุตฺต}
\proseitem{305}{สาวตฺถินิทานํ.\csromanspacer{} จกฺขุสมฺผสฺโส ภิกฺขเว อนิจฺโจ วิปริณามี อญฺญถาภาวี, โสตสมฺผสฺโส.\csromanspacer{} ฆานสมฺผสฺโส.\csromanspacer{} ชิวฺหาสมฺผสฺโส.\csromanspacer{} กายสมฺผสฺโส.\csromanspacer{} มโนสมฺผสฺโส อนิจฺโจ วิปริณามี อญฺญถาภาวี.\csromanspacer{} โย ภิกฺขเว อิเม ธมฺเม เอวํ สทฺทหติ อธิมุจฺจติ.\csromanspacer{} อยํ วุจฺจติ สทฺธานุสารี ฯเปฯ สมฺพาธิปรายโนติ.\csromanspacer{}\csromanspacer{}จตุตฺถํ.}

\csromanheader{2}{5. สมฺผสฺสชาสุตฺต}
\proseitem{306}{สาวตฺถินิทานํ.\csromanspacer{} จกฺขุสมฺผสฺสชา ภิกฺขเว เวทนา อนิจฺจา วิปริณามี อญฺญถาภาวี, โสตสมฺผสฺสชา เวทนา ฯเปฯ.\csromanspacer{} ฆานสมฺผสฺสชา เวทนา ฯเปฯ.\csromanspacer{} ชิวฺหาสมฺผสฺสชา เวทนา ฯเปฯ.\csromanspacer{} กายสมฺผสฺสชา เวทนา ฯเปฯ.\csromanspacer{} มโนสมฺผสฺสชา เวทนา อนิจฺจา วิปริณามี อญฺญถาภาวี.\csromanspacer{} โย ภิกฺขเว อิเม ธมฺเม เอวํ สทฺทหติ อธิมุจฺจติ.\csromanspacer{} อยํ วุจฺจติ สทฺธานุสารี ฯเปฯ สมฺโพธิปรายโนติ.\csromanspacer{}\csromanspacer{}ปญฺจมํ.}

\csromanheader{2}{6. รูปสญฺญาสุตฺต}
\proseitem{307}{\csromanfolio{188}สาวตฺถินิทานํ.\csromanspacer{} รูปสญฺญา ภิกฺขเว อนิจฺจา วิปริณามี อญฺญถาภาวี, สทฺทสญฺญา.\csromanspacer{} คนฺธสญฺญา.\csromanspacer{} รสสญฺญา.\csromanspacer{} โผฏฺฐพฺพสญฺญา.\csromanspacer{} ธมฺมสญฺญา อนิจฺจา วิปริณามี อญฺญถาภาวี.\csromanspacer{} โย ภิกฺขเว อิเม ธมฺเม เอวํ สทฺทหติ อธิมุจฺจติ.\csromanspacer{} อยํ วุจฺจติ สทฺธานุสารี ฯเปฯ สมฺโพธิปรายโนติ.\csromanspacer{}\csromanspacer{}ฉฏฺฐํ.}

\csromanheader{2}{7. รูปสญฺเจตนาสุตฺต}
\proseitem{308}{สาวตฺถินิทานํ.\csromanspacer{} รูปสญฺเจตนา ภิกฺขเว อนิจฺจา วิปริณามี อญฺญถาภาวี, สทฺทสญฺเจตนา.\csromanspacer{} คนฺธสญฺเจตนา.\csromanspacer{} รสสญฺเจตนา.\csromanspacer{} โผฏฺฐพฺพสญฺเจตนา.\csromanspacer{} ธมฺมสญฺเจตนา อนิจฺจา วิปริณามี อญฺญถาภาวี.\csromanspacer{} โย ภิกฺขเว อิเม ธมฺเม เอวํ สทฺทหติ อธิมุจฺจติ.\csromanspacer{} อยํ วุจฺจติ สทฺธานุสารี ฯเปฯ สมฺโพธิปรายโนติ.\csromanspacer{}\csromanspacer{}สตฺตมํ.}

\csromanheader{2}{8. รูปตณฺหาสุตฺต}
\proseitem{309}{สาวตฺถินิทานํ.\csromanspacer{} รูปตณฺหา ภิกฺขเว อนิจฺจา วิปริณามี อญฺญถาภาวี, สทฺทตณฺหา.\csromanspacer{} คนฺธตณฺหา.\csromanspacer{} รสตณฺหา.\csromanspacer{} โผฏฺฐพฺพตณฺหา.\csromanspacer{} ธมฺมตณฺหา อนิจฺจา วิปริณามี อญฺญถาภาวี.\csromanspacer{} โย ภิกฺขเว อิเม ธมฺเม เอวํ สทฺทหติ อธิมุจฺจติ.\csromanspacer{} อยํ วุจฺจติ สทฺธานุสารี ฯเปฯ สมฺโพธิปรายโนติ.\csromanspacer{}\csromanspacer{}อฏฺฐมํ.}

\csromanheader{2}{9. ปถวีธาตุสุตฺต}
\proseitem{310}{สาวตฺถินิทานํ.\csromanspacer{} ปถวีธาตุ ภิกฺขเว อนิจฺจา วิปริณามี อญฺญถาภาวี, อาโปธาตุ.\csromanspacer{} เตโชธาตุ.\csromanspacer{} วาโยธาตุ.\csromanspacer{} อากาสธาตุ.\csromanspacer{} วิญฺญาณธาตุ อนิจฺจา วิปริณามี อญฺญถาภาวี.\csromanspacer{} โย ภิกฺขเว อิเม ธมฺเม เอวํ สทฺทหติ อธิมุจฺจติ.\csromanspacer{} อยํ วุจฺจติ สทฺธานุสารี ฯเปฯ สมฺโพธิปรายโนติ.\csromanspacer{}\csromanspacer{}นวมํ.}

\csromanheader{2}{10. ขนฺธสุตฺต}
\proseitem{311}{สาวตฺถินิทานํ.\csromanspacer{} รูปํ ภิกฺขเว อนิจฺจํ วิปริณามิ อญฺญถาภาวิ, เวทนา อนิจฺจา วิปริณามี อญฺญถาภาวี, สญฺญา.\csromanspacer{} สงฺขารา อนิจฺจา วิปริณามิโน อญฺญถาภาวิโน, วิญฺญาณํ อนิจฺจํ วิปริณามิ}

\csromanheader{2}{4. โอกฺกนฺตสํยุตฺต}
\prose{\csromanfolio{189}อญฺญถาภาวิ.\csromanspacer{} โย ภิกฺขเว อิเม ธมฺเม เอวํ สทฺทหติ อธิมุจฺจติ.\csromanspacer{} อยํ วุจฺจติ สทฺธานุสารี โอกฺกนฺโต สมฺมตฺตนิยามํ, สปฺปุริสภูมิํ โอกฺกนฺโต, วีติวตฺโต ปุถุชฺชนภูมิํ, อภพฺโพ ตํ กมฺมํ กาตุํ, ยํ กมฺมํ กตฺวา นิรยํ วา ติรจฺฉานโยนิํ วา เปตฺติวิสยํ วา อุปปชฺเชยฺย, อภพฺโพ จ ตาว กาลํ กาตุํ, ยาว น โสตาปตฺติผลํ สจฺฉิกโรติ.}

\prose{ยสฺส โข ภิกฺขเว อิเม ธมฺมา เอวํ ปญฺญาย มตฺตโส นิชฺฌานํ ขมนฺติ.\csromanspacer{} อยํ วุจฺจติ ธมฺมานุสารี โอกฺกนฺโต สมฺมตฺตนิยามํ, สปฺปุริสภูมิํ โอกฺกนฺโต, วีติวตฺโต ปุถุชฺชนภูมิํ, อภพฺโพ ตํ กมฺมํ กาตุํ, ยํ กมฺมํ กตฺวา นิรยํ วา ติรจฺฉานโยนิํ วา เปตฺติวิสยํ วา อุปปชฺเชยฺย, อภพฺโพ จ ตาว กาลํ กาตุํ, ยาว น โสตาปตฺติผลํ สจฺฉิกโรติ.\csromanspacer{} โย ภิกฺขเว อิเม ธมฺเม เอวํ ปชานาติ เอวํ ปสฺสติ.\csromanspacer{} อยํ วุจฺจติ โสตาปนฺโน อวินิปาตธมฺโม นิยโต สมฺโพธิปรายโนติ.\csromanspacer{}\csromanspacer{}ทสมํ.}

\vspace{\tipitakaparskip}
\csromancenter{\textbf{โอกฺกนฺตสํยุตฺตํ}\footnote{โอกฺกนฺติกสํยุตฺตํ (อิ, ก)} \textbf{สมตฺตํ.}}
\titlehead{\textbf{ตสฺสุทฺทานํ}}
\csromangathameasure{จกฺขุ รูปญฺจ วิญฺญาณํ, ผสฺโส จ เวทนาย จ. \\ สญฺญา จ เจตนา ตณฺหา, ธาตุ ขนฺเธน เต ทสาติ.}
\csromangathagroup{\csromangathastanza{จกฺขุ รูปญฺจ วิญฺญาณํ, ผสฺโส จ เวทนาย จ. \\ สญฺญา จ เจตนา ตณฺหา, ธาตุ ขนฺเธน เต ทสาติ.}}

\csromanmatikaanchor{mtk.60}
\csromantocmarknopage{chapter}{5. อุปฺปาทสํยุตฺต}{mtk.60}
\bookmark[dest={mtk.60},level=0]{5. อุปฺปาทสํยุตฺต}
\chapterheadpageread{5. อุปฺปาทสํยุตฺต}
\csromanheader{2}{1. จกฺขุสุตฺต}
\csromanmatikaanchor{mtk.61}
\csromantocmarkref{section}{1-10. จกฺขุสุตฺตาทิ}{mtk.61}
\bookmark[dest={mtk.61},level=1]{1-10. จกฺขุสุตฺตาทิ}
\proseitem{312}{\csromanfolio{190}สาวตฺถินิทานํ.\csromanspacer{} โย โข ภิกฺขเว จกฺขุสฺส อุปฺปาโท ฐิติ อภินิพฺพตฺติ ปาตุภาโว, ทุกฺขสฺเสโส อุปฺปาโท, โรคานํ ฐิติ, ชรามรณสฺส ปาตุภาโว.\csromanspacer{} โย โสตสฺส อุปฺปาโท ฐิติ ฯเปฯ.\csromanspacer{} โย ฆานสฺส อุปฺปาโท ฐิติ.\csromanspacer{} โย ชิวฺหาย อุปฺปาโท ฐิติ.\csromanspacer{} โย กายสฺส อุปฺปาโท ฐิติ.\csromanspacer{} โย มนสฺส อุปฺปาโท ฐิติ อภินิพฺพตฺติ ปาตุภาโว, ทุกฺขสฺเสโส อุปฺปาโท, โรคานํ ฐิติ, ชารามรณสฺส ปาตุภาโว.\csromanspacer{} โย จ ภิกฺขเว จกฺขุสฺส นิโรโธ วูปสโม อตฺถงฺคโม, ทุกฺขสฺเสโส นิโรโธ, โรคานํ วูปสโม, ชรามรณสฺส อตฺถงฺคโม.\csromanspacer{} โย โสตสฺส นิโรโธ ฯเปฯ.\csromanspacer{} โย ฆานสฺส นิโรโธ.\csromanspacer{} โย ชิวฺหาย นิโรโธ.\csromanspacer{} โย กายสฺส นิโรโธ.\csromanspacer{} โย มนสา นิโรโธ วูปสโม อตฺถงฺคโม, ทุกฺขสฺเสโส นิโรโธ, โรคานํ วูปสโม, ชรามรณสฺส อตฺถงฺคโมติ.\csromanspacer{}\csromanspacer{}ปฐมํ.}

\csromanheader{2}{2. รูปสุตฺต}
\proseitem{313}{สาวตฺถินิทานํ.\csromanspacer{} โย โข ภิกฺขเว รูปานํ อุปฺปาโท ฐิติ อภินิพฺพตฺติ ปาตุภาโว, ทุกฺขสฺเสโส อุปฺปาโท, โรคานํ ฐิติ, ชรามรณสฺส ปาตุภาโว.\csromanspacer{} โย สทฺทานํ.\csromanspacer{} โย คนฺธานํ.\csromanspacer{} โย รสานํ.\csromanspacer{} โย โผฏฺฐพฺพานํ.\csromanspacer{} โย ธมฺมานํ อุปฺปาโท ฐิติ อภินิพฺพตฺติ ปาตุภาโว, ทุกฺขสฺเสโส อุปฺปาโท, โรคานํ ฐิติ, ชรามรณสฺส ปาตุภาโว.\csromanspacer{} โย จ โข ภิกฺขเว รูปานํ นิโรโธ วูปสโม อตฺถงฺคโม, ทุกฺขสฺเสโส นิโรโธ, โรคานํ วูปสโม, ชรามรณสฺส อตฺถงฺคโม.\csromanspacer{} โย สทฺทานํ.\csromanspacer{} โย คนฺธานํ.\csromanspacer{} โย รสานํ.\csromanspacer{} โย โผฏฺฐพฺพานํ.\csromanspacer{} โย ธมฺมานํ นิโรโธ วูปสโม อตฺถงฺคโม, ทุกฺขสฺเสโส นิโรโธ, โรคานํ วูปสโม, ชรามรณสฺส อตฺถงฺคโมติ.\csromanspacer{}\csromanspacer{}ทุติยํ.}

\csromanheader{2}{5. อุปฺปาทสํยุตฺต \textbf{3. วิญฺญาณสุตฺต}}
\proseitem{314}{\csromanfolio{191}สาวตฺถินิทานํ.\csromanspacer{} โย โข ภิกฺขเว จกฺขุวิญฺญาณสฺส อุปฺปาโท ฐิติ ฯเปฯ ชรามรณสฺส ปาตุภาโว ฯเปฯ.\csromanspacer{} โย มโนวิญฺญาณสฺส อุปฺปาโท ฐิติ ฯเปฯ ชรามรณสฺส ปาตุภาโว.\csromanspacer{} โย จ โข ภิกฺขเว จกฺขุวิญฺญาณสฺส นิโรโธ ฯเปฯ ชรามรณสฺส อตฺถงฺคโม ฯเปฯ.\csromanspacer{} โย มโนวิญฺญาณสฺส นิโรโธ ฯเปฯ ชรามรณสฺส อตฺถงฺคโมติ.\csromanspacer{}\csromanspacer{}ตติยํ.}

\csromanheader{2}{4. สมฺผสฺสสุตฺต}
\proseitem{315}{สาวตฺถินิทานํ.\csromanspacer{} โย โข ภิกฺขเว จกฺขุสมฺผสฺสสฺส อุปฺปาโท ฐิติ ฯเปฯ ชรามรณสฺส ปาตุภาโว ฯเปฯ.\csromanspacer{} โย มโนสมฺผสฺสสฺส อุปฺปาโท ฐิติ ฯเปฯ ชรามรณสฺส ปาตุภาโว.\csromanspacer{} โย จ โข ภิกฺขเว จกฺขุสมฺผสฺสสฺส นิโรโธ ฯเปฯ ชรามรณสฺส อตฺถงฺคโม ฯเปฯ.\csromanspacer{} โย มโนสมฺผสฺสสฺส นิโรโธ ฯเปฯ ชรามรณสฺส อตฺถงฺคโมติ.\csromanspacer{}\csromanspacer{}จตุตฺถํ.}

\csromanheader{2}{5. สมฺผสฺสชสุตฺต}
\proseitem{316}{สาวตฺถินิทานํ.\csromanspacer{} โย โข ภิกฺขเว จกฺขุสมฺผสฺสชาย เวทนาย อุปฺปาโท ฐิติ ฯเปฯ ชรามรณสฺส ปาตุภาโว ฯเปฯ.\csromanspacer{} โย มโนสมฺผสฺสชาย เวทนาย อุปฺปาโท ฐิติ ฯเปฯ ชรามรณสฺส ปาตุภาโว.\csromanspacer{} โย จ โข ภิกฺขเว จกฺขุสมฺผสฺสชาย เวทนาย นิโรโธ วูปสโม ฯเปฯ ชรามรณสฺส อตฺถงฺคโม ฯเปฯ.\csromanspacer{} โย มโนสมฺผสฺสชาย เวทนาย นิโรโธ วูปสโม อตฺถงฺคโม, ทุกฺขสฺเสโส นิโรโธ, โรคานํ วูปสโม, ชรามรณสฺส อตฺถงฺคโมติ.\csromanspacer{}\csromanspacer{}ปญฺจมํ.}

\csromanheader{2}{6. สญฺญาสุตฺต}
\proseitem{317}{สาวตฺถินิทานํ.\csromanspacer{} โย โข ภิกฺขเว รูปสญฺญาย อุปฺปาโท ฐิติ ฯเปฯ ชรามรณสฺส ปาตุภาโว ฯเปฯ.\csromanspacer{} โย ธมฺมสญฺญาย อุปฺปาโท ฐิติ อภินิพฺพตฺติ ปาตุภาโว, ทุกฺขสฺเสโส อุปฺปาโท, โรคานํ ฐิติ, ชรามรณสฺส ปาตุภาโว.\csromanspacer{} โย จ โข ภิกฺขเว รูปสญฺญาย นิโรโธ ฯเปฯ ชรามรณสฺส อตฺถงฺคโม ฯเปฯ.\csromanspacer{} โย ธมฺมสญฺญาย นิโรโธ วูปสโม อตฺถงฺคโม, ทุกฺขสฺเสโส นิโรโธ, โรคานํ วูปสโม, ชรามรณสฺส อตฺถงฺคโมติ.\csromanspacer{}\csromanspacer{}ฉฏฺฐํ.}

\csromanheader{2}{7. สญฺเจตนาสุตฺต}
\proseitem{318}{\csromanfolio{192}สาวตฺถินิทานํ.\csromanspacer{} โย โข ภิกฺขเว รูปสญฺเจตนาย อุปฺปาโท ฐิติ ฯเปฯ ชรามรณสฺส ปาตุภาโว ฯเปฯ.\csromanspacer{} โย ธมฺมสญฺเจตนาย อุปฺปาโท ฐิติ อภินิพฺพตฺติ ปาตุภาโว, ทุกฺขสฺเสโส อุปฺปาโท, โรคานํ ฐิติ, ชรามรณสฺส ปาตุภาโว.\csromanspacer{} โย จ โข ภิกฺขเว รูปสญฺเจตนาย นิโรโธ ฯเปฯ ชรามรณสฺส อตฺถงฺคโม ฯเปฯ.\csromanspacer{} โย ธมฺมสญฺเจตนาย นิโรโธ วูปสโม อตฺถงฺคโม, ทุกฺขสฺเสโส นิโรโธ, โรคานํ วูปสโม, ชรามรณสฺส อตฺถงฺคโมติ.\csromanspacer{}\csromanspacer{}สตฺตมํ.}

\csromanheader{2}{8. ตณฺหาสุตฺต}
\proseitem{319}{สาวตฺถินิทานํ.\csromanspacer{} โย โข ภิกฺขเว รูปตณฺหาย อุปฺปาโท ฐิติ ฯเปฯ ชรามรณสฺส ปาตุภาโว ฯเปฯ.\csromanspacer{} โย ธมฺมตณฺหาย อุปฺปาโท ฐิติ อภินิพฺพตฺติ ปาตุภาโว, ทุกฺขสฺเสโส อุปฺปาโท, โรคานํ ฐิติ, ชรามรณสฺส ปาตุภาโว.\csromanspacer{} โย จ โข ภิกฺขเว รูปตณฺหาย นิโรโธ ฯเปฯ ชรามรณสฺส อตฺถงฺคโม ฯเปฯ.\csromanspacer{} โย ธมฺมตณฺหาย นิโรโธ วูปสโม อตฺถงฺคโม, ทุกฺขสฺเสโส นิโรโธ, โรคานํ วูปสโม, ชรามรณสฺส อตฺถงฺคโมติ.\csromanspacer{}\csromanspacer{}อฏฺฐมํ.}

\csromanheader{2}{9. ธาตุสุตฺต}
\proseitem{320}{สาวตฺถินิทานํ.\csromanspacer{} โย โข ภิกฺขเว ปถวีธาตุยา อุปฺปาโท ฐิติ อภินิพฺพตฺติ ปาตุภาโว ฯเปฯ ชรามรณสฺส ปาตุภาโว.\csromanspacer{} โย อาโปธาตุยา.\csromanspacer{} โย เตโชธาตุยา.\csromanspacer{} โย วาโยธาตุยา.\csromanspacer{} โย อากาสธาตุยา.\csromanspacer{} โย วิญฺญาณธาตุยา อุปฺปาโท ฐิติ อภินิพฺพตฺติ ปาตุภาโว, ทุกฺขสฺเสโส อุปฺปาโท, โรคานํ ฐิติ, ชรามรณสฺส ปาตุภาโว.\csromanspacer{} โย จ โข ภิกฺขเว ปถวีธาตุยา นิโรโธ ฯเปฯ ชรามรณสฺส อตฺถงฺคโม.\csromanspacer{} โย อาโปธาตุยา นิโรโธ.\csromanspacer{} โย เตโชธาตุยา นิโรโธ.\csromanspacer{} โย วาโยธาตุยา นิโรโธ.\csromanspacer{} โย อากาสธาตุยา นิโรโธ.\csromanspacer{} โย วิญฺญาณธาตุยา นิโรโธ วูปสโม อตฺถงฺคโม, ทุกฺขสฺเสโส นิโรโธ, โรคานํ วูปสโม, ชรามรณสฺส อตฺถงฺคโมติ.\csromanspacer{}\csromanspacer{}นวมํ.}

\csromanheader{2}{5. อุปฺปาทสํยุตฺต \textbf{10. ขนฺธสุตฺต}}
\proseitem{321}{\csromanfolio{193}สาวตฺถินิทานํ.\csromanspacer{} โย โข ภิกฺขเว รูปสฺส อุปฺปาโท ฐิติ อภินิพฺพตฺติ ปาตุภาโว, ทุกฺขสฺเสโส อุปฺปาโท โรคานํ ฐิติ ชรามรณสฺส ปาตุภาโว.\csromanspacer{} โย เวทนาย.\csromanspacer{} โย สญฺญาย.\csromanspacer{} โย สงฺขารานํ.\csromanspacer{} โย วิญฺญาณสฺส อุปฺปาโท ฐิติ อภินิพฺพตฺติ ปาตุภาโว, ทุกฺขสฺเสโส อุปฺปาโท โรคานํ ฐิติ ชรามรณสฺส ปาตุภาโว.\csromanspacer{} โย จ โข ภิกฺขเว รูปสฺส นิโรโธ วูปสโม อตฺถงฺคโม, ทุกฺขสฺเสโส นิโรโธ โรคานํ วูปสโม ชรามรณสฺส อตฺถงฺคโม.\csromanspacer{} โย เวทนาย.\csromanspacer{} โย สญฺญาย.\csromanspacer{} โย สงฺขารานํ.\csromanspacer{} โย วิญฺญาณสฺส นิโรโธ วูปสโม อตฺถงฺคโม, ทุกฺขสฺเสโส นิโรโธ โรคานํ วูปสโม ชรามรณสฺส อตฺถงฺคโมติ.\csromanspacer{}\csromanspacer{}ทสมํ.}

\vspace{\tipitakaparskip}
\csromancenter{\textbf{อุปฺปาทสํยุตฺตํ สมตฺตํ.}}
\titlehead{\textbf{ตสฺสุทฺทานํ}}
\csromangathameasure{จกฺขุ รูปญฺจ วิญฺญาณํ, ผสฺโส จ เวทนาย จ. \\ สญฺญา จ เจตนา ตณฺหา, ธาตุ ขนฺเธน เต ทสาติ.}
\csromangathagroup{\csromangathastanza{จกฺขุ รูปญฺจ วิญฺญาณํ, ผสฺโส จ เวทนาย จ. \\ สญฺญา จ เจตนา ตณฺหา, ธาตุ ขนฺเธน เต ทสาติ.}}

\csromanmatikaanchor{mtk.62}
\csromantocmarknopage{chapter}{6. กิเลสสํยุตฺต}{mtk.62}
\bookmark[dest={mtk.62},level=0]{6. กิเลสสํยุตฺต}
\chapterheadpageread{6. กิเลสสํยุตฺต}
\csromanheader{2}{1. จกฺขุสุตฺต}
\csromanmatikaanchor{mtk.63}
\csromantocmarkref{section}{1-10. จกฺขุสุตฺตาทิ}{mtk.63}
\bookmark[dest={mtk.63},level=1]{1-10. จกฺขุสุตฺตาทิ}
\proseitem{322}{\csromanfolio{194}สาวตฺถินิทานํ.\csromanspacer{} โย ภิกฺขเว จกฺขุสฺมิํ ฉนฺทราโค, จิตฺตสฺเสโส อุปกฺกิเลโส.\csromanspacer{} โย โสตสฺมิํ ฉนฺทราโค, จิตฺตสฺเสโส อุปกฺกิเลโส.\csromanspacer{} โย ฆานสฺมิํ ฉนฺทราโค, จิตฺตสฺเสโส อุปกฺกิเลโส.\csromanspacer{} โย ชิวฺหาย ฉนฺทราโค, จิตฺตสฺเสโส อุปกฺกิเลโส.\csromanspacer{} โย กายสฺมิํ ฉนฺทราโค, จิตฺตสฺเสโส อุปกฺกิเลโส.\csromanspacer{} โย มนสฺมิํ ฉนฺทราโค, จิตฺตสฺเสโส อุปกฺกิเลโส.\csromanspacer{} ยโต โข ภิกฺขเว ภิกฺขุโน อิเมสุ ฉสุ ฐาเนสุ เจตโส อุปกฺกิเลโส ปหีโน โหติ, เนกฺขมฺมนินฺนํ จสฺส จิตฺตํ โหติ, เนกฺขมฺมปริภาวิตํ จิตฺตํ กมฺมนิยํ ขายติ อภิญฺญา สจฺฉิกรณีเยสุ ธมฺเมสูติ.\csromanspacer{}\csromanspacer{}ปฐมํ.}

\csromanheader{2}{2. รูปสุตฺต}
\proseitem{323}{สาวตฺถินิทานํ.\csromanspacer{} โย ภิกฺขเว รูเปสุ ฉนฺทราโค, จิตฺตสฺเสโส อุปกฺกิเลโส.\csromanspacer{} โย สทฺเทสุ.\csromanspacer{} โย คนฺเธสุ.\csromanspacer{} โย รเสสุ.\csromanspacer{} โย โผฏฺฐพฺเพสุ.\csromanspacer{} โย ธมฺเมสุ ฉนฺทราโค, จิตฺตสฺเสโส อุปกฺกิเลโส.\csromanspacer{} ยโต โข ภิกฺขเว ภิกฺขุโน อิเมสุ ฉสุ ฐาเนสุ เจตโส อุปกฺกิเลโส ปหีโน โหติ, เนกฺขมฺมนินฺนํ จสฺส จิตฺตํ โหติ, เนกฺขมฺมปริภาวิตํ จิตฺตํ กมฺมนิยํ ขายติ อภิญฺญา สจฺฉิกรณีเยสุ ธมฺเมสูติ.\csromanspacer{}\csromanspacer{}ทุติยํ.}

\csromanheader{2}{3. วิญฺญาณสุตฺต}
\proseitem{324}{สาวตฺถินิทานํ.\csromanspacer{} โย ภิกฺขเว จกฺขวิญฺญาณสฺมิํ ฉนฺทราโค, จิตฺตสฺเสโส อุปกฺกิเลโส.\csromanspacer{} โย โสตวิญฺญาณสฺมิํ.\csromanspacer{} โย ฆานวิญฺญาณสฺมิํ.\csromanspacer{} โย ชิวฺหาวิญฺญาณสฺมิํ.\csromanspacer{} โย กายวิญฺญาณสฺมิํ.\csromanspacer{} โย มโนวิญฺญาณสฺมิํ ฉนฺทราโค, จิตฺตสฺเสโส อุปกฺกิเลโส.\csromanspacer{} ยโต โข ภิกฺขเว ภิกฺขุโน อิเมสุ ฉสุ ฐาเนสุ เจตโส อุปกฺกิเลโส ปหีโน โหติ, เนกฺขมฺมนินฺนํ จสฺส จิตฺตํ โหติ, เนกฺขมฺมปริภาวิตํ จิตฺตํ กมฺมนิยํ ขายติ อภิญฺญา สจฺฉิกรณีเยสุ ธมฺเมสูติ.\csromanspacer{}\csromanspacer{}ตติยํ.}

\csromanheader{2}{6. กิเลสสํยุตฺต \textbf{4. สมฺผสฺสสุตฺต}}
\proseitem{325}{\csromanfolio{195}สาวตฺถินิทานํ.\csromanspacer{} โย ภิกฺขเว จกฺขุสมฺผสฺสสฺมิํ ฉนฺทราโค, จิตฺตสฺเสโส อุปกฺกิเลโส.\csromanspacer{} โย โสตสมฺผสฺสสฺมิํ.\csromanspacer{} โย ฆานสมฺผสฺสสฺมิํ.\csromanspacer{} โย ชิวฺหาสมฺผสฺสสฺมิํ.\csromanspacer{} โย กายสมฺผสฺสสฺมิํ.\csromanspacer{} โย มโนสมฺผสฺสสฺมิํ ฉนฺทราโค, จิตฺตสฺเสโส อุปกฺกิเลโส.\csromanspacer{} ยโต โข ภิกฺขเว ภิกฺขุโน ฯเปฯ อภิญฺญา สจฺฉิกรณีเยสุ ธมฺเมสูติ.\csromanspacer{}\csromanspacer{}จตุตฺถํ.}

\csromanheader{2}{5. สมฺผสฺสชสุตฺต}
\proseitem{326}{สาวตฺถินิทานํ.\csromanspacer{} โย ภิกฺขเว จกฺขุสมฺผสฺสชาย เวทนาย ฉนฺทราโค, จิตฺตสฺเสโส อุปกฺกิเลโส.\csromanspacer{} โย โสตสมฺผสฺสชาย เวทนาย.\csromanspacer{} โย ฆานสมฺผสฺสชาย เวทนาย.\csromanspacer{} โย ชิวฺหาสมฺผสฺสชาย เวทนาย.\csromanspacer{} โย กายสมฺผสฺสชาย เวทนาย.\csromanspacer{} โย มโนสมฺผสฺสชาย เวทนาย ฉนฺทราโค, จิตฺตสฺเสโส อุปกฺกิเลโส.\csromanspacer{} ยโต โข ภิกฺขเว ภิกฺขุโน ฯเปฯ อภิญฺญา สจฺฉิกรณีเยสุ ธมฺเมสูติ.\csromanspacer{}\csromanspacer{}ปญฺจมํ.}

\csromanheader{2}{6. สญฺญาสุตฺต}
\proseitem{327}{สาวตฺถินิทานํ.\csromanspacer{} โย ภิกฺขเว รูปสญฺญาย ฉนฺทราโค, จิตฺตสฺเสโส อุปกฺกิเลโส.\csromanspacer{} โย สทฺทสญฺญาย.\csromanspacer{} โย คนฺธสญฺญาย.\csromanspacer{} โย รสสญฺญาย.\csromanspacer{} โย โผฏฺฐพฺพสญฺญาย.\csromanspacer{} โย ธมฺมสญฺญาย ฉนฺทราโค, จิตฺตสฺเสโส อุปกฺกิเลโส.\csromanspacer{} ยโต โข ภิกฺขเว ภิกฺขุโน ฯเปฯ อภิญฺญา สจฺฉิกรณีเยสุ ธมฺเมสูติ.\csromanspacer{}\csromanspacer{}ฉฏฺฐํ.}

\csromanheader{2}{7. สญฺเจตนาสุตฺต}
\proseitem{328}{สาวตฺถินิทานํ.\csromanspacer{} โย ภิกฺขเว รูปสญฺเจตนาย ฉนฺทราโค, จิตฺตสฺเสโส อุปกฺกิเลโส.\csromanspacer{} โย สทฺทสญฺเจตนาย.\csromanspacer{} โย คนฺธสญฺเจตนาย.\csromanspacer{} โย รสสญฺเจตนาย.\csromanspacer{} โย โผฏฺฐพฺพสญฺเจตนาย.\csromanspacer{} โย ธมฺมสญฺเจตนาย ฉนฺทราโค, จิตฺตสฺเสโส อุปกฺกิเลโส.\csromanspacer{} ยโต โข ภิกฺขเว ภิกฺขุโน ฯเปฯ อภิญฺญา สจฺฉิกรณีเยสุ ธมฺเมสูติ.\csromanspacer{}\csromanspacer{}สตฺตมํ.}

\csromanheader{2}{8. ตณฺหาสุตฺต}
\proseitem{329}{\csromanfolio{196}สาวตฺถินิทานํ.\csromanspacer{} โย ภิกฺขเว รูปตณฺหาย ฉนฺทราโค, จิตฺตสฺเสโส อุปกฺกิเลโส.\csromanspacer{} โย สทฺทตณฺหาย.\csromanspacer{} โย คนฺธตณฺหาย.\csromanspacer{} โย รสตณฺหาย.\csromanspacer{} โย โผฏฺฐพฺพตณฺหาย.\csromanspacer{} โย ธมฺมตณฺหาย ฉนฺทราโค, จิตฺตสฺเสโส อุปกฺกิเลโส.\csromanspacer{} ยโต โข ภิกฺขเว ภิกฺขุโน ฯเปฯ อภิญฺญา สจฺฉิกรณีเยสุ ธมฺเมสูติ.\csromanspacer{}\csromanspacer{}อฏฺฐมํ.}

\csromanheader{2}{9. ธาตุสุตฺต}
\proseitem{330}{สาวตฺถินิทานํ.\csromanspacer{} โย ภิกฺขเว ปถวีธาตุยา ฉนฺทราโค, จิตฺตสฺเสโส อุปกฺกิเลโส.\csromanspacer{} โย อาโปธาตุยา.\csromanspacer{} โย เตโชธาตุยา.\csromanspacer{} โย วาโยธาตุยา.\csromanspacer{} โย อากาสธาตุยา.\csromanspacer{} โย วิญฺญาณธาตุยา ฉนฺทราโค, จิตฺตสฺเสโส อุปกฺกิเลโส.\csromanspacer{} ยโต โข ภิกฺขเว ภิกฺขุโน อิเมสุ ฉสุ ฐาเนสุ เจตโส อุปกฺกิเลโส ปหีโน โหติ, เนกฺขมฺมนินฺนํ จสฺส จิตฺตํ โหติ, เนกฺขมฺมปริภาวิตํ จิตฺตํ กมฺมนิยํ ขายติ อภิญฺญา สจฺฉิกรณีเยสุ ธมฺเมสูติ.\csromanspacer{}\csromanspacer{}นวมํ.}

\csromanheader{2}{10. ขนฺธสุตฺต}
\proseitem{331}{สาวตฺถินิทานํ.\csromanspacer{} โย ภิกฺขเว รูปสฺมิํ ฉนฺทราโค, จิตฺตสฺเสโส อุปกฺกิเลโส ฯเปฯ.\csromanspacer{} โย วิญฺญาณสฺมิํ ฉนฺทราโค, จิตฺตสฺเสโส อุปกฺกิเลโส.\csromanspacer{} ยโต โข ภิกฺขเว ภิกฺขุโน อิเมสุ ปญฺจสุ ฐาเนสุ เจตโส อุปกฺกิเลโส ปหีโน โหติ, เนกฺขมฺมนินฺนํ จสฺส จิตฺตํ โหติ, เนกฺขมฺมปริภาวิตํ จิตฺตํ กมฺมนิยํ ขายติ อภิญฺญา สจฺฉิกรณีเยสุ ธมฺเมสูติ.\csromanspacer{}\csromanspacer{}ทสมํ.}

\vspace{\tipitakaparskip}
\csromancenter{\textbf{กิเลสสํยุตฺตํ สมตฺตํ.}}
\titlehead{\textbf{ตสฺสุทฺทานํ}}
\csromangathameasure{จกฺขุ รูปญฺจ วิญฺญาณํ, ผสฺโส จ เวทนาย จ. \\ สญฺญา จ เจตนา ตณฺหา, ธาตุ ขนฺเธน เต ทสาติ.}
\csromangathagroup{\csromangathastanza{จกฺขุ รูปญฺจ วิญฺญาณํ, ผสฺโส จ เวทนาย จ. \\ สญฺญา จ เจตนา ตณฺหา, ธาตุ ขนฺเธน เต ทสาติ.}}

\csromanmatikaanchor{mtk.64}
\csromantocmarknopage{chapter}{7. สาริปุตฺตสํยุตฺต}{mtk.64}
\bookmark[dest={mtk.64},level=0]{7. สาริปุตฺตสํยุตฺต}
\chapterheadpageread{7. สาริปุตฺตสํยุตฺต}
\csromanheader{2}{1. วิเวกชสุตฺต}
\csromanmatikaanchor{mtk.65}
\csromantocmarkref{section}{1-9. วิเวกชสุตฺตาทิ}{mtk.65}
\bookmark[dest={mtk.65},level=1]{1-9. วิเวกชสุตฺตาทิ}
\proseitem{332}{\csromanfolio{197}เอกํ สมยํ อายสฺมา สาริปุตฺโต สาวตฺถิยํ วิหรติ เชตวเน อนาถปิณฺฑิกสฺส อาราเม.\csromanspacer{} อถ โข อายสฺมา สาริปุตฺโต ปุพฺพณฺหสมยํ นิวาเสตฺวา ปตฺตจีวรมาทาย สาวตฺถิํ ปิณฺฑาย ปาวิสิ, สาวตฺถิยํ ปิณฺฑาย จริตฺวา ปจฺฉาภตฺตํ ปิณฺฑปาตปฏิกฺกนฺโต เยน อนฺธวนํ เตนุปสงฺกมิ ทิวาวิหาราย, อนฺธวนํ อชฺโฌคาเหตฺวา อญฺญตรสฺมิํ รุกฺขมูเล ทิวาวิหารํ นิสีทิ.}

\prose{อถ โข อายสฺมา สาริปุตฺโต สายนฺหสมยํ ปฏิสลฺลานา วุฏฺฐิโต เยน เชตวนํ อนาถปิณฺฑิกสฺส อาราโม เตนุปสงฺกมิ.\csromanspacer{} อทฺทสา โข อายสฺมา อานนฺโท อายสฺมนฺตํ สาริปุตฺตํ ทูรโตว อาคจฺฉนฺตํ, ทิสฺวาน อายสฺมนฺตํ สาริปุตฺตํ เอตทโวจ “วิปฺปสนฺนานิ โข เต อาวุโส สาริปุตฺต อินฺทฺริยานิ ปริสุทฺโธ มุขวณฺโณ ปริโยทาโต, กตเมนายสฺมา สาริปุตฺโต อชฺช วิหาเรน วิหาสี”ติ.}

\prose{อิธาหํ อาวุโส วิวิจฺเจว กาเมหิ วิวิจฺจ อกุสเลหิ ธมฺเมหิ สวิตกฺกํ สวิจารํ วิเวกชํ ปีติสุขํ ปฐมํ ฌานํ อุปสมฺปชฺช วิหรามิ.\csromanspacer{} ตสฺส มยฺหํ อาวุโส น เอวํ โหติ “อหํ ปฐมํ ฌานํ สมาปชฺชามี”ติ วา, “อหํ ปฐมํ ฌานํ สมาปนฺโน”ติ วา, “อหํ ปฐมา ฌานา วุฏฺฐิโต”ติ วาติ.\csromanspacer{} ตถา หิ ปนายสฺมโต สาริปุตฺตสฺส ทีฆรตฺตํ อหงฺการมมงฺการมานานุสยา สุสมูหตา, ตสฺมา อายสฺมโต สาริปุตฺตสฺส น เอวํ โหติ “อหํ ปฐมํ ฌานํ สมาปชฺชามี”ติ วา, “อหํ ปฐมํ ฌานํ สมาปนฺโน”ติ วา, “อหํ ปฐมา ฌานา วุฏฺฐิโต”ติ วาติ.\csromanspacer{}\csromanspacer{}ปฐมํ.}

\csromanheader{2}{2. อวิตกฺกสุตฺต}
\proseitem{333}{สาวตฺถินิทานํ.\csromanspacer{} อทฺทสา โข อายสฺมา อานนฺโท ฯเปฯ ±ยสฺมนฺตํ สาริปุตฺตํ เอตทโวจ “วิปฺปสนฺนานิ โข เต อาวุโส สาริปุตฺต อินฺทฺริยานิ ปริสุทฺโธ มุขวณฺโณ ปริโยทาโต, กตเมนายสฺมา สาริปุตฺโต อชฺช วิหาเรน วิหาสี”ติ.}

\prose{\csromanfolio{198}อิธาหํ อาวุโส วิตกฺกวิจารานํ วูปสมา อชฺฌตฺตํ สมฺปสาทนํ เจตโส เอโกทิภาวํ อวิตกฺกํ อวิจารํ สมาธิชํ ปีติสุขํ ทุติยํ ฌานํ อุปสมฺปชฺช วิหารามิ.\csromanspacer{} ตสฺส มยฺหํ อาวุโส น เอวํ โหติ “อหํ ทุติยํ ฌานํ สมาปชฺชามี”ติ วา, “อหํ ทุติยํ ฌานํ สมาปนฺโน”ติ วา, “อหํ ทุติยา ฌานา วุฏฺฐิโต”ติ วาติ.\csromanspacer{} ตถา หิ ปนายสฺมโต สาริปุตฺตสฺส ทีฆรตฺตํ อหงฺการมมงฺการมานานุสยา สุสมูหตา, ตสฺมา อายสฺมโต สาริปุตฺตสฺส น เอวํ โหติ “อหํ ทุติยํ ฌานํ สมาปชฺชามี”ติ วา, “อหํ ทุติยํ ฌานํ สมาปนฺโน”ติ วา, “อหํ ทุติยา ฌานา วุฏฺฐิโต”ติ วาติ.\csromanspacer{}\csromanspacer{}ทุติยํ.}

\csromanheader{2}{3. ปีติสุตฺต}
\proseitem{334}{สาวตฺถินิทานํ.\csromanspacer{} อทฺทสา โข อายสฺมา อานนฺโท ฯเปฯ วิปฺปสนฺนานิ โข เต อาวุโส สาริปุตฺต อินฺทฺริยานิ ปริสุทฺโธ มุขวณฺโณ ปริโยทาโต, กตเมนายสฺมา สาริปุตฺโต อชฺช วิหาเรน วิหาสีติ.}

\prose{อิธาหํ อาวุโส ปีติยา จ วิราคา อุเปกฺขโก จ วิหาสิํ สโต จ สมฺปชาโน สุขญฺจ กาเยน ปฏิสํเวเทมิ, ยํ ตํ อริยา อาจิกฺขนฺติ “อุเปกฺขโก สติมา สุขวิหารี”ติ, ตติยํ ฌานํ อุปสมฺปชฺช วิหรามิ.\csromanspacer{} ตสฺส มยฺหํ อาวุโส น เอวํ โหติ “อหํ ตติยํ ฌานํ สมาปชฺชามี”ติ วา, “อหํ ตติยํ ฌานํ สมาปนฺโน”ติ วา, “อหํ ตติยา ฌานา วุฏฺฐิโต”ติ วาติ.\csromanspacer{} ตถา หิ ปนายสฺมโต สาริปุตฺตสฺส ทีฆรตฺตํ อหงฺการมมงฺการมานานุสยา สุสมูหตา, ตสฺมา อายสฺมโต สาริปุตฺตสฺส น เอวํ โหติ “อหํ ตติยํ ฌานํ สมาปชฺชามี”ติ วา, “อหํ ตติยํ ฌานํ สมาปนฺโน”ติ วา, “อหํ ตติยา ฌานา วุฏฺฐิโต”ติ วาติ.\csromanspacer{}\csromanspacer{}ตติยํ.}

\csromanheader{2}{4. อุเปกฺขาสุตฺต}
\proseitem{335}{สาวตฺถินิทานํ.\csromanspacer{} อทฺทสา โข อายสฺมา อานนฺโท ฯเปฯ วิปฺปสนฺนานิ โข เต อาวุโส สาริปุตฺต อินฺทฺริยานิ ปริสุทฺโธ มุขวณฺโณ ปริโยทาโต, กตเมนายสฺมา สาริปุตฺโต อชฺช วิหาเรน วิหาสีติ.}

\prose{อิธาหํ อาวุโส สุขสฺส จ ปหานา ทุกฺขสฺส จ ปหานา ปุพฺเพว โสมนสฺสโทมนสฺสานํ อตฺถงฺคมา อทุกฺขมสุขํ อุเปกฺขาสติปาริสุทฺธิํ}

\vspace{\tipitakaparskip}
\csromancenter{สาริปุตฺตสํยุตฺต จตุตฺถํ ฌานํ อุปสมฺปชฺช วิหรามิ.\csromanspacer{} ตสฺส มยฺหํ อาวุโส น เอวํ โหติ “อหํ จตุตฺถํ ฌานํ สมาปชฺชามี”ติ วา, “อหํ จตุตฺถํ ฌานํ สมาปนฺโน”ติ วา, “อหํ จตุตฺถา ฌานา วุฏฺฐิโต”ติ วาติ.\csromanspacer{} ตถา หิ ปนายสฺมโต สาริปุตฺตสฺส ทีฆรตฺตํ อหงฺการมมงฺการมานานุสยา สุสมูหตา, ตสฺมา อายสฺมโต สาริปุตฺตสฺส น เอวํ โหติ “อหํ จตุตฺถํ ฌานํ สมาปชฺชามี”ติ วา, “อหํ จตุตฺถํ ฌานํ สมาปนฺโน”ติ วา, “อหํ จตุตฺถา ฌานา วุฏฺฐิโต”ติ วาติ.\csromanspacer{}\csromanspacer{}จตุตฺถํ.}
\csromanheader{2}{5. อากาสานญฺจายตนสุตฺต}
\proseitem{336}{\csromanfolio{199}สาวตฺถินิทานํ.\csromanspacer{} อทฺทสา โข อายสฺมา อานนฺโท ฯเปฯ.\csromanspacer{} อิธาหํ อาวุโส สพฺพโส รูปสญฺญานํ สมติกฺกมา ปฏิฆสญฺญานํ อตฺถงฺคมา นานตฺตสญฺญานํ อมนสิการา “อนนฺโต อากาโส”ติ อากาสานญฺจายตนํ อุปสมฺปชฺช วิหรามิ ฯเปฯ วุฏฺฐิโตติ วาติ.\csromanspacer{}\csromanspacer{}ปญฺจมํ.}

\csromanheader{2}{6. วิญฺญาณญฺจายตนสุตฺต}
\proseitem{337}{สาวตฺถินิทานํ.\csromanspacer{} อทฺทสา โข อายสฺมา อานนฺโท ฯเปฯ.\csromanspacer{} อิธาหํ อาวุโส สพฺพโส อากาสานญฺจายตนํ สมติกฺกมฺม “อนนฺตํ วิญฺญาณนฺ”ติ วิญฺญาณญฺจายตนํ อุปสมฺปชฺช วิหรามิ ฯเปฯ.\csromanspacer{} วุฏฺฐิโตติ วาติ.\csromanspacer{}\csromanspacer{}ฉฏฺฐํ.}

\csromanheader{2}{7. อากิญฺจญฺญายตนสุตฺต}
\proseitem{338}{สาวตฺถินิทานํ.\csromanspacer{} อถ โข อายสฺมา สาริปุตฺโต ฯเปฯ.\csromanspacer{} อิธาหํ อาวุโส สพฺพโส วิญฺญาณญฺจายตนํ สมติกฺกมฺม “นตฺถิ กิญฺจี”ติ อากิญฺจญฺญายตนํ อุปสมฺปชฺช วิหรามิ ฯเปฯ วุฏฺฐิโตติ วาติ.\csromanspacer{}\csromanspacer{}สตฺตมํ.}

\csromanheader{2}{8. เนวสญฺญานาสญฺญายตนสุตฺต}
\proseitem{339}{สาวตฺถินิทานํ.\csromanspacer{} อถ โข อายสฺมา สาริปุตฺโต ฯเปฯ.\csromanspacer{} อิธาหํ อาวุโส อากิญฺจญฺญายตนํ สมติกฺกมฺม เนวสญฺญานาสญฺญายตนํ อุปสมฺปชฺช วิหรามิ ฯเปฯ วุฏฺฐิโตติ วาติ.\csromanspacer{}\csromanspacer{}อฏฺฐมํ.}

\csromanheader{2}{9. นิโรธสมาปตฺติสุตฺต}
\proseitem{340}{สาวตฺถินิทานํ.\csromanspacer{} อถ โข อายสฺมา สาริปุตฺโต ฯเปฯ.\csromanspacer{} อิธาหํ อาวุโส สพฺพโส เนวสญฺญานาสญฺญายตนํ สมติกฺกมฺม \csromanfolio{200}สญฺญาเวทยิตนิโรธํ อุปสมฺปชฺช วิหรามิ.\csromanspacer{} ตสฺส มยฺหํ อาวุโส น เอวํ โหติ “อหํ สญฺญาเวทยิตนิโรธํ สมาปชฺชามี”ติ วา, “อหํ สญฺญาเวทยิตนิโรธํ สมาปนฺโน”ติ วา, “อหํ สญฺญาเวทยิตนิโรธา วุฏฺฐิโต”ติ วาติ.\csromanspacer{} ตถา หิ ปนายสฺมโต สาริปุตฺตสฺส ทีฆรตฺตํ อหงฺการมมงฺการมานานุสยา สุสมูหตา, ตสฺมา อายสฺมโต สาริปุตฺตสฺส น เอวํ โหติ “อหํ สญฺญาเวทยิตนิโรธํ สมาปชฺชามี”ติ วา, “อหํ สญฺญาเวทยิตนิโรธํ สมาปนฺโน”ติ วา, “อหํ สญฺญาเวทยิตนิโรธา วุฏฺฐิโต”ติ วาติ.\csromanspacer{}\csromanspacer{}นวมํ.}

\csromanmatikaanchor{mtk.66}
\csromantocmarkref{section}{10. สุจิมุขีสุตฺต}{mtk.66}
\bookmark[dest={mtk.66},level=1]{10. สุจิมุขีสุตฺต}
\csromanheader{1}{10. สุจิมุขีสุตฺต}
\proseitem{341}{เอกํ สมยํ อายสฺมา สาริปุตฺโต ราชคเห วิหรติ เวฬุวเน กลนฺทกนิวาเป.\csromanspacer{} อถ โข อายสฺมา สาริปุตฺโต ปุพฺพณฺหสมยํ นิวาเสตฺวา ปตฺตจีวรมาทาย ราชคเห ปิณฺฑาย ปาวิสิ, ราชคเห สปทานํ ปิณฺฑาย จริตฺวา ตํ ปิณฺฑปาตํ อญฺญตรํ กุฏฺฏมูลํ\footnote{กุฑฺฑมูลํ (สี, สฺยา, กํ), กุฑฺฑํ (อิ)} นิสฺสาย ปริภุญฺชติ.\csromanspacer{} อถ โข สุจิมุขี ปริพฺพาชิกา เยนายสฺมา สาริปุตฺโต เตนุปสงฺกมิ, อุปสงฺกมิตฺวา อายสฺมนฺตํ สาริปุตฺตํ เอตทโวจ\csromandash{}}

\prose{กิํ นุ โข สมณ อโธมุโข ภุญฺชสีติ.\csromanspacer{} น ขฺวาหํ ภคินิ อโธมุโข ภุญฺชามีติ.\csromanspacer{} เตน หิ สมณ อุพฺภมุโข\footnote{อุทฺธํมุโข (สี-ฐ)} ภุญฺชสีติ.\csromanspacer{} น ขฺวาหํ ภคินิ อุพฺภมุโข ภุญฺชามีติ.\csromanspacer{} เตน หิ สมณ ทิสามุโข ภุญฺชสีติ.\csromanspacer{} น ขฺวาหํ ภคินิ ทิสามุโข ภุญฺชามีติ.\csromanspacer{} เตน หิ สมณ วิทิสามุโข ภุญฺชสีติ.\csromanspacer{} น ขฺวาหํ ภคินิ วิทิสามุโข ภุญฺชามีติ.}

\prose{“กิํ นุ สมณ อโธมุโข ภุญฺชสี”ติ อิติ ปุฏฺโฐ สมาโน “น ขฺวาหํ ภคินิ อโธมุโข ภุญฺชามี”ติ วเทสิ, “เตน หิ สมณ อุพฺภมุโข ภุญฺชสี”ติ อิติ ปุฏฺโฐ สมาโน “น ขฺวาหํ ภคินิ อุพฺภมุโข ภุญฺชามี”ติ วเทสิ, “เตน หิ สมณ ทิสามุโข ภุญฺชสี”ติ อิติ ปุฏฺโฐ สมาโน “น ขฺวาหํ ภคินิ ทิสามุโข ภุญฺชามี”ติ วเทสิ, “เตน หิ สมณ วิทิสามุโข ภุญฺชสี”ติ อิติ ปุฏฺโฐ สมาโน “น ขฺวาหํ ภคินิ วิทิสามุโข ภุญฺชามี”ติ วเทสิ.}

\csromanheader{2}{7. สาริปุตฺตสํยุตฺต}
\prose{\csromanfolio{201}กถญฺจรหิ สมณ ภุญฺชสีติ.\csromanspacer{} เย หิ เกจิ ภคินิ สมณพฺราหฺมณา\footnote{สมณา วา พฺราหฺมณา วา (สี) นิคมนวาเกฺย ปน สพฺพตฺถาปิ สมาโสเยว ทิสฺสติ.} วตฺถุวิชฺชาติรจฺฉานวิชฺชาย มิจฺฉาชีเวน ชีวิกํ\footnote{ชีวิตํ (ก)} กปฺเปนฺติ, อิเม วุจฺจนฺติ ภคินิ สมณพฺราหฺมณา อโธมุขา ภุญฺชนฺตีติ.\csromanspacer{} เย หิ เกจิ ภคินิ สมณพฺราหฺมณา นกฺขตฺตวิชฺชาติรจฺฉานวิชฺชาย มิจฺฉาชีเวน ชีวิกํ กปฺเปนฺติ, อิเม วุจฺจนฺติ ภคินิ สมณพฺราหฺมณา อุพฺภมุขา ภุญฺชนฺตีติ.\csromanspacer{} เย หิ เกจิ ภคินิ สมณพฺราหฺมณา ทูเตยฺยปหิณคมนานุโยคาย\footnote{...นุโยคา (สี, สฺยา, กํ, อิ), ...นุโยเคน (?)} มิจฺฉาชีเวน ชีวิกํ กปฺเปนฺติ, อิเม วุจฺจนฺติ ภคินิ สมณพฺราหฺมณา ทิสามุขา ภุญฺชนฺตีติ.\csromanspacer{} เย หิ เกจิ ภคินิ สมณาพฺราหฺมณา องฺควิชฺชาติรจฺฉานวิชฺชาย มิจฺฉาชีเวน ชีวิกํ กปฺเปนฺติ, อิเม วุจฺจนฺติ ภคินิ สมณพฺราหฺมณา วิทิสามุขา ภุญฺชนฺตีติ.}

\prose{โส ขฺวาหํ ภคินิ น วตฺถุวิชฺชาติรจฺฉานวิชฺชาย มิจฺฉาชีเวน ชีวิกํ กปฺเปมิ, น นกฺขตฺตวิชฺชาติรจฺฉานวิชฺชาย มิจฺฉาชีเวน ชีวิกํ กปฺเปมิ, น ทูเตยฺยปหิณคมนานุโยคาย มิจฺฉาชีเวน ชีวิกํ กปฺเปมิ, น องฺควิชฺชาติรจฺฉานวิชฺชาย มิจฺฉาชีเวน ชีวิกํ กปฺเปมิ, ธมฺเมน ภิกฺขํ ปริเยสามิ, ธมฺเมน ภิกฺขํ ปริเยสิตฺวา ภุญฺชามีติ.}

\prose{อถ โข สุจิมุขี ปริพฺพาชิกา ราชคเห รถิยาย รถิยํ สิงฺฆาฏเกน สิงฺฆาฏกํ อุปสงฺกมิตฺวา เอวมาโรเจสิ “ธมฺมิกํ สมณา สกฺยปุตฺติยา อาหารํ อาหาเรนฺติ, อนวชฺชํ\footnote{อนวชฺเชน (ก)} สมณา สกฺยปุตฺติยา อาหารํ อาหาเรนฺติ, เทถ สมณานํ สกฺยปุตฺติยานํ ปิณฺฑนฺ”ติ.\csromanspacer{}\csromanspacer{}ทสมํ.}

\vspace{\tipitakaparskip}
\csromancenter{\textbf{สาริปุตฺตสํยุตฺตํ สมตฺตํ.}}
\titlehead{\textbf{ตสฺสุทฺทานํ}}
\csromangathameasure{วิเวกชํ อวิตกฺกํ, ปีติ อุเปกฺขา จตุตฺถกํ. \\ อากาสญฺเจว วิญฺญาณํ, อากิญฺจํ เนวสญฺญินา. \\ นิโรโธ นวโม วุตฺโต, ทสมํ สุจิมุขี จาติ.}
\csromangathagroup{\csromangathastanza{วิเวกชํ อวิตกฺกํ, ปีติ อุเปกฺขา จตุตฺถกํ. \\ อากาสญฺเจว วิญฺญาณํ, อากิญฺจํ เนวสญฺญินา. \\ นิโรโธ นวโม วุตฺโต, ทสมํ สุจิมุขี จาติ.}}

\csromanmatikaanchor{mtk.67}
\csromantocmarknopage{chapter}{8. นาคสํยุตฺต}{mtk.67}
\bookmark[dest={mtk.67},level=0]{8. นาคสํยุตฺต}
\chapterheadpageread{8. นาคสํยุตฺต}
\csromanheader{2}{1. สุทฺธิกสุตฺต}
\csromanmatikaanchor{mtk.68}
\csromantocmarkref{section}{1-50. สุทฺธิกสุตฺตาทิ}{mtk.68}
\bookmark[dest={mtk.68},level=1]{1-50. สุทฺธิกสุตฺตาทิ}
\proseitem{342}{\csromanfolio{202}สาวตฺถินิทานํ.\csromanspacer{} จตสฺโส อิมา ภิกฺขเว นาคโยนิโย.\csromanspacer{} กตมา จตสฺโส, อณฺฑชา นาคา ชลาพุชา นาคา สํเสทชา นาคา โอปปาติกา นาคา.\csromanspacer{} อิมา โข ภิกฺขเว จตสฺโส นาคนิโยติ.\csromanspacer{}\csromanspacer{}ปฐมํ.}

\csromanheader{2}{2. ปณีตตรสุตฺต}
\proseitem{343}{สาวตฺถินิทานํ.\csromanspacer{} จตสฺโส อิมา ภิกฺขเว นาคโยนิโย.\csromanspacer{} กตมา จตสฺโส, อณฺฑชา นาคา ชลาพุชา นาคา สํเสทชา นาคา โอปปาติกา นาคา.\csromanspacer{} ตตฺร ภิกฺขเว อณฺฑเชหิ นาเคหิ ชลาพุชา จ สํเสทชา จ โอปปาติกา จ นาคา ปณีตตรา.\csromanspacer{} ตตฺร ภิกฺขเว อณฺฑเชหิ จ ชลาพุเชหิ จ นาเคหิ สํเสทชา จ โอปปาติกา จ นาคา ปณีตตรา.\csromanspacer{} ตตฺร ภิกฺขเว อณฺฑเชหิ จ ชลาพุเชหิ จ สํเสทเชหิ จ นาเคหิ โอปปาติกา นาคา ปณีตตรา.\csromanspacer{} อิมา โข ภิกฺขเว จตสฺโส นาคโยนิโยติ.\csromanspacer{}\csromanspacer{}ทุติยํ.}

\csromanheader{2}{3. อุโปสถสุตฺต}
\proseitem{344}{เอกํ สมยํ ภควา สาวตฺถิยํ วิหรติ เชตวเน อนาถปิณฺฑิกสฺส อาราเม.\csromanspacer{} อถ โข อญฺญตโร ภิกฺขุ เยน ภควา เตนุปสงฺกมิ, อุปสงฺกมิตฺวา เอกมนฺตํ นิสีทิ, เอกมนฺตํ นิสินฺโน โข โส ภิกฺขุ ภควนฺตํ เอตทโวจ “โก นุ โข ภนฺเต เหตุ โก ปจฺจโย, เยน มิเธกจฺเจ อณฺฑชา นาคา อุโปสถํ อุปวสนฺติ, โวสฺสฏฺฐกายา จ ภวนฺตี”ติ.}

\prose{อิธ ภิกฺขุ เอกจฺจานํ อณฺฑชานํ นาคานํ เอวํ โหติ “มยํ โข ปุพฺเพ กาเยน ทฺวยการิโน อหุมฺห วาจาย ทฺวยการิโน มนสา ทฺวยการิโน, เต มยํ กาเยน ทฺวยการิโน วาจาย ทฺวยการิโน มนสา ทฺวยการิโน กายสฺส เภทา ปรํ มรณา อณฺฑชานํ นาคานํ สหพฺยตํ อุปปนฺนา, สจชฺช มยํ กาเยน สุจริตํ จเรยฺยาม, วาจาย สุจริตํ จเรยฺยาม, มนสา สุจริตํ จเรยฺยาม, เอวํ}

\vspace{\tipitakaparskip}
\csromancenter{นาคสํยุตฺต มยํ กายสฺส เภทา ปรํ มรณา สุคติํ สคฺคํ โลกํ อุปปชฺเชยฺยาม, หนฺท มยํ เอตรหิ กาเยน สุจริตํ จราม, วาจาย สุจริตํ จราม, มนสา สุจริตํ จรามา”ติ.\csromanspacer{} อยํ โข ภิกฺขุ เหตุ อยํ ปจฺจโย, เยน มิเธกจฺเจ อณฺฑชา นาคา อุโปสถํ อุปวสนฺติ, โวสฺสฏฺฐกายา จ ภวนฺตีติ.\csromanspacer{}\csromanspacer{}ตติยํ.}
\csromanheader{2}{4. ทุติย-อุโปสถสุตฺต}
\proseitem{345}{\csromanfolio{203}สาวตฺถินิทานํ.\csromanspacer{} อถ โข อญฺญตโร ภิกฺขุ เยน ภควา ฯเปฯ เอกมนฺตํ นิสินฺโน โข โส ภิกฺขุ ภควนฺตํ เอตทโวจ “โก นุ โข ภนฺเต เหตุ โก ปจฺจโย, เยน มิเธกจฺเจ ชลาพุชา นาคา อุโปสถํ อุปวสนฺติ, โวสฺสฏฺฐกายา จ ภวนฺตี”ติ. (สพฺพํ วิตฺถาเรตพฺพํ.) อยํ โข ภิกฺขุ เหตุ อยํ ปจฺจโย, เยน มิเธกจฺเจ ชลาพุชา นาคา อุโปสถํ อุปวสนฺติ, โวสฺสฏฺฐกายา จ ภวนฺตีติ.\csromanspacer{}\csromanspacer{}จตุตฺถํ.}

\csromanheader{2}{5. ตติย-อุโปสถสุตฺต}
\proseitem{346}{สาวตฺถินิทานํ.\csromanspacer{} เอกมนฺตํ นิสินฺโน โข โส ภิกฺขุ ภควนฺตํ เอตทโวจ “โก นุ โข ภนฺเต เหตุ โก ปจฺจโย, เยน มิเธกจฺเจ สํเสทชา นาคา อุโปสถํ อุปวสนฺติ, โวสฺสฏฺฐกายา จ ภวนฺตี”ติ. (สพฺพํ วิตฺถาเรตพฺพํ.) อยํ โข ภิกฺขุ เหตุ อยํ ปจฺจโย, เยน มิเธกจฺเจ สํสทชา นาคา อุโปสถํ อุโปวสนฺติ, โวสฺสฏฺฐกายา จ ภวนฺตีติ.\csromanspacer{}\csromanspacer{}ปญฺจมํ.}

\csromanheader{2}{6. จตุตฺถ-อุโปสถสุตฺต}
\proseitem{347}{สาวตฺถินิทานํ.\csromanspacer{} เอกมนฺตํ นิสินฺโน โข โส ภิกฺขุ ภควนฺตํ เอตทโวจ “โก นุ โข ภนฺเต เหตุ โก ปจฺจโย, เยน มิเธกจฺเจ โอปปาติกา นาคา อุโปสถํ อุปวสนฺติ, โวสฺสฏฺฐกายา จ ภวนฺตี”ติ.}

\prose{อิธ ภิกฺขุ เอกจฺจานํ โอปปาติกานํ นาคานํ เอวํ โหติ “มยํ โข ปุพฺเพ กาเยน ทฺวยการิโน อหุมฺห วาจาย ทฺวยการิโน มนสา ทฺวยการิโน, เต มยํ กาเยน ทฺวยการิโน วาจาย ทฺวยการิโน มนสา ทฺวยการิโน กายสฺส เภทา ปรํ มรณา โอปปาติกานํ นาคานํ สหพฺยตํ อุปปนฺนา, สจชฺช มยํ กาเยน สุจริตํ \csromanfolio{204}จเรยฺยาม, วาจาย.\csromanspacer{} มนสา สุจริตํ จเรยฺยาม, เอวํ มยํ กายสฺส เภทา ปรํ มรณา สุคติํ สคฺคํ โลกํ อุปปชฺเชยฺยาม, หนฺท มยํ เอตรหิ กาเยน สุจริตํ จราม, วาจาย.\csromanspacer{} มนสา สุจริตํ จรามา”ติ.\csromanspacer{} อยํ โข ภิกฺขุ เหตุ อยํ ปจฺจโย, เยน มิเธกจฺเจ โอปปาติกา นาคา อุโปสถํ อุปวสนฺติ, โวสฺสฏฺฐกายา จ ภวนฺตีติ.\csromanspacer{}\csromanspacer{}ฉฏฺฐํ.}

\csromanheader{2}{7. สุตสุตฺต}
\proseitem{348}{สาวตฺถินิทานํ.\csromanspacer{} เอกมนฺตํ นิสินฺโน โข โส ภิกฺขุ ภควนฺตํ เอตทโวจ “โก นุ โข ภนฺเต เหตุ โก ปจฺจโย, เยน มิเธกจฺโจ กายสฺส เภทา ปรํ มรณา อณฺฑชานํ นาคานํ สหพฺยตํ อุปปชฺชตี”ติ.}

\prose{อิธ ภิกฺขุ เอกจฺโจ กาเยน ทฺวยการี โหติ, วาจาย ทฺวยการี โหติ, มนสา ทฺวยการี โหติ, ตสฺส สุตํ โหติ “อณฺฑชา นาคา ทีฆายุกา วณฺณวนฺโต สุขพหุลา”ติ.\csromanspacer{} ตสฺส เอวํ โหติ “อโห วตาหํ กายสฺส เภทา ปรํ มรณา อณฺฑชานํ นาคานํ สหพฺยตํ อุปปชฺเชยฺยนฺ”ติ.\csromanspacer{} โส กายสฺส เภทา ปรํ มรณา อณฺฑชานํ นาคานํ สหพฺยตํ อุปปชฺชติ, อยํ โข ภิกฺขุ เหตุ อยํ ปจฺจโย, เยน มิเธกจฺโจ กายสฺส เภทา ปรํ มรณา อณฺฑชานํ นาคานํ สหพฺยตํ อุปปชฺชตีติ.\csromanspacer{}\csromanspacer{}สตฺตมํ.}

\csromanheader{2}{8. ทุติยสุตสุตฺต}
\proseitem{349}{สาวตฺถินิทานํ.\csromanspacer{} เอกมนฺตํ นิสินฺโน โข โส ภิกฺขุ ภควนฺตํ เอตทโวจ “โก นุ โข ภนฺเต เหตุ โก ปจฺจโย, เยน มิเธกจฺโจ กายสฺส เภทา ปรํ มรณา ชลาพุชานํ นาคานํ สหพฺยตํ อุปปชฺชตี”ติ ฯเปฯ อยํ โข ภิกฺขุ เหตุ อยํ ปจฺจโย, เยน มิเธกจฺโจ กายสฺส เภทา ปรํ มรณา ชลาพุชานํ นาคานํ สหพฺยตํ อุปปชฺชตีติ.\csromanspacer{}\csromanspacer{}อฏฺฐมํ.}

\csromanheader{2}{9. ตติยสุตสุตฺต}
\proseitem{350}{สาวตฺถินิทานํ.\csromanspacer{} เอกมนฺตํ นิสินฺโน โข โส ภิกฺขุ ภควนฺตํ เอตทโวจ “โก นุ โข ภนฺเต เหตุ โก ปจฺจโย, เยน มิเธกจฺโจ}

\vspace{\tipitakaparskip}
\csromancenter{นาคสํยุตฺต กายสฺส เภทา ปรํ มรณา สํเสทชานํ นาคานํ สหพฺยตํ อุปปชฺชตี”ติ ฯเปฯ อยํ โข ภิกฺขุ เหตุ อยํ ปจฺจโย, เยน มิเธกจฺโจ กายสฺส เภทา ปรํ มรณา สํเสทชานํ นาคานํ สหพฺยตํ อุปปชฺชตีติ.\csromanspacer{}\csromanspacer{}นวมํ.}
\csromanheader{2}{10. จตุตฺถสุตสุตฺต}
\proseitem{351}{\csromanfolio{205}สาวตฺถินิทานํ.\csromanspacer{} เอกมนฺตํ นิสินฺโน โข โส ภิกฺขุ ภควนฺตํ เอตทโวจ “โก นุ โข ภนฺเต เหตุ โก ปจฺจโย, เยน มิเธกจฺโจ กายสฺส เภทา ปรํ มรณา โอปปาติกานํ นาคานํ สหพฺยตํ อุปปชฺชตี”ติ.}

\prose{อิธ ภิกฺขุ เอกจฺโจ กาเยน ทฺวยการี โหติ วาจาย ทฺวยการี มนสา ทฺวยการี, ตสฺส สุตํ โหติ “โอปปาติกา นาคา ทีฆายุกา วณฺณวนฺโต สุขพหุลา”ติ.\csromanspacer{} ตสฺส เอวํ โหติ “อโห วตาหํ กายสฺส เภทา ปรํ มรณา โอปปาติกานํ นาคานํ สหพฺยตํ อุปปชฺเชยฺยนฺ”ติ.\csromanspacer{} โส กายสฺส เภทา ปรํ มรณา โอปปาติกานํ นาคานํ สหพฺยตํ อุปปชฺชติ.\csromanspacer{} อยํ โข ภิกฺขุ เหตุ อยํ ปจฺจโย, เยน มิเธกจฺโจ กายสฺส เภทา ปรํ มรณา โอปปาติกานํ นาคานํ สหพฺยตํ อุปปชฺชตีติ.\csromanspacer{}\csromanspacer{}ทสมํ.}

\csromanheader{2}{11-20. อณฺฑชทานูปการสุตฺตทสก}
\proseitem{352-361}{เอกมนฺตํ นิสินฺโน โข โส ภิกฺขุ ภควนฺตํ เอตทโวจ “โก นุ โข ภนฺเต เหตุ โก ปจฺจโย, เยน มิเธกจฺโจ กายสฺส เภทา ปรํ มรณา อณฺฑชานํ นาคานํ สหพฺยตํ อุปปชฺชตี”ติ.}

\prose{อิธ ภิกฺขุ เอกจฺโจ กาเยน ทฺวยการี โหติ วาจาย ทฺวยการี มนสา ทฺวยการี, ตสฺส สุตํ โหติ “อณฺฑชา นาคา ทีฆายุกา วณฺณวนฺโต สุขพหุลา”ติ.\csromanspacer{} ตสฺส เอวํ โหติ “อโห วตาหํ กายสฺส เภทา ปรํ มรณา อณฺฑชานํ นาคานํ สหพฺยตํ อุปปชฺเชยฺยนฺ”ติ.\csromanspacer{} โส อนฺนํ เทติ, โส กายสฺส เภทา ปรํ มรณา อณฺฑชานํ นาคานํ สหพฺยตํ อุปปชฺชติ.\csromanspacer{} อยํ โข ภิกฺขุ เหตุ ฯเปฯ อุปปชฺชตีติ ฯเปฯ.\csromanspacer{} โส ปานํ เทติ ฯเปฯ.\csromanspacer{} วตฺถํ เทติ ฯเปฯ.\csromanspacer{} ยานํ เทติ ฯเปฯ.\csromanspacer{} มาลํ เทติ ฯเปฯ.\csromanspacer{} คนฺธํ เทติ ฯเปฯ.\csromanspacer{} วิเลปนํ เทติ ฯเปฯ.\csromanspacer{} เสยฺยํ เทติ ฯเปฯ.\csromanspacer{} อาวสถํ เทติ ฯเปฯ. \csromanfolio{206}ปทีเปยฺยํ เทติ, โส กายสฺส เภทา ปรํ มรณา อณฺฑชานํ นาคานํ สหพฺยตํ อุปปชฺชติ.\csromanspacer{} อยํ โข ภิกฺขุ เหตุ อยํ ปจฺจโย, เยน มิเธกจฺโจ กายสฺส เภทา ปรํ มรณา อณฺฑชานํ นาคานํ สหพฺยตํ อุปปชฺชตีติ.\csromanspacer{}\csromanspacer{}วีสติมํ.}

\csromanheader{2}{21-50. ชลาพุชาทิทานูปการสุตฺตตฺติํสก}
\proseitem{362-391}{สาวตฺถินิทานํ.\csromanspacer{} เอกมนฺตํ นิสินฺโน โข โส ภิกฺขุ ภควนฺตํ เอตทโวจ “โก นุ โข ภนฺเต เหตุ โก ปจฺจโย, เยน มิเธกจฺโจ กายสฺส เภทา ปรํ มรณา ชลาพุชานํ นาคานํ ฯเปฯ สํเสทชานํ นาคานํ ฯเปฯ โอปปาติกานํ นาคานํ สหพฺยตํ อุปปชฺชตี”ติ.}

\prose{อิธ ภิกฺขุ เอกจฺโจ กาเยน ทฺวยการี โหติ วาจาย ทฺวยการี มนสา ทฺวยการี, ตสฺส สุตํ โหติ “โอปปาติกา นาคา ทีฆายุกา วณฺณวนฺโต สุขพหุลา”ติ.\csromanspacer{} ตสฺส เอวํ โหติ “อโห วตาหํ กายสฺส เภทา ปรํ มรณา โอปปาติกานํ นาคานํ สหพฺยตํ อุปปชฺเชยฺยนฺ”ติ.\csromanspacer{} โส อนฺนํ เทติ ฯเปฯ.\csromanspacer{} ปานํ เทติ ฯเปฯ.\csromanspacer{} ปทีเปยฺยํ เทติ, โส กายสฺส เภทา ปรํ มรณา โอปปาติกานํ นาคานํ สหพฺยตํ อุปปชฺชติ.\csromanspacer{} อยํ โข ภิกฺขุ เหตุ อยํ ปจฺจโย, เยน มิเธกจฺโจ กายสฺส เภทา ปรํ มรณา โอปปาติกานํ นาคานํ สหพฺยตํ อุปปชฺชตีติ.}

\prose{(อิมินา เปยฺยาเลน ทส ทส สุตฺตนฺตา กาตพฺพา.\csromanspacer{} เอวํ จตูสุ โยนีสุ จตฺตาลีสํ เวยฺยากรณา โหนฺติ.\csromanspacer{} ปุริเมหิ ปน ทสหิ สุตฺตนฺเตหิ สห โหนฺติ ปณฺณาสสุตฺตนฺตาติ.)}

\vspace{\tipitakaparskip}
\csromancenter{\textbf{นาคสํยุตฺตํ สมตฺตํ.}}
\titlehead{\textbf{ตสฺสุทฺทานํ}}
\csromangathameasure{สุทฺธิกํ ปณีตตรํ, จตุโร จ อุโปสถา. \\ ตสฺส สุตํ จตุโร จ, ทานูปการา จ ตาลีสํ. \\ ปณฺณาส ปิณฺฑโต สุตฺตา, นาคมฺหิ สุปฺปกาสิตาติ.}
\csromangathagroup{\csromangathastanza{สุทฺธิกํ ปณีตตรํ, จตุโร จ อุโปสถา. \\ ตสฺส สุตํ จตุโร จ, ทานูปการา จ ตาลีสํ. \\ ปณฺณาส ปิณฺฑโต สุตฺตา, นาคมฺหิ สุปฺปกาสิตาติ.}}

\csromanmatikaanchor{mtk.69}
\csromantocmarknopage{chapter}{9. สุปณฺณสํยุตฺต}{mtk.69}
\bookmark[dest={mtk.69},level=0]{9. สุปณฺณสํยุตฺต}
\chapterheadpageread{9. สุปณฺณสํยุตฺต}
\csromanheader{2}{1. สุทฺธิกสุตฺต}
\csromanmatikaanchor{mtk.70}
\csromantocmarkref{section}{1-46. สุทฺธิกสุตฺตาทิ}{mtk.70}
\bookmark[dest={mtk.70},level=1]{1-46. สุทฺธิกสุตฺตาทิ}
\proseitem{392}{\csromanfolio{207}สาวตฺถินิทานํ.\csromanspacer{} จตสฺโส อิมา ภิกฺขเว สุปณฺณโยนิโย.\csromanspacer{} กตมา จตสฺโส, อณฺฑชา สุปณฺณา ชลาพุชา สุปณฺณา สํเสทชา สุปณฺณา โอปปาติกา สุปณฺณา.\csromanspacer{} อิมา โข ภิกฺขเว จตสฺโส สุปณฺณโยนิโยติ.\csromanspacer{}\csromanspacer{}ปฐมํ.}

\csromanheader{2}{2. หรนฺติสุตฺต}
\proseitem{393}{สาวตฺถินิทานํ.\csromanspacer{} จตสฺโส อิมา ภิกฺขเว สุปณฺณโยนิโย.\csromanspacer{} กตมา จตสฺโส, อณฺฑชา ฯเปฯ.\csromanspacer{} อิมา โข ภิกฺขเว จตสฺโส สุปณฺณโยนิโย.\csromanspacer{} ตตฺร ภิกฺขเว อณฺฑชา สุปณฺณา อณฺฑเชว นาเค หรนฺติ, น ชลาพุเช น สํเสทเช น โอปปาติเก.\csromanspacer{} ตตฺร ภิกฺขเว ชลาพุชา สุปณฺณา อณฺฑเช จ ชลาพุเช จ นาเค หรนฺติ, น สํเสทเช น โอปปาติเก.\csromanspacer{} ตตฺร ภิกฺขเว สํเสทชา สุปณฺณา อณฺฑเช จ ชลาพุเช จ สํเสทเช จ นาเค หรนฺติ, น โอปปาติเก.\csromanspacer{} ตตฺร ภิกฺขเว โอปปาติกา สุปณฺณา อณฺฑเช จ ชลาพุเช จ สํเสทเช จ โอปปาติเก จ นาเค หรนฺติ.\csromanspacer{} อิมา โข ภิกฺขเว จตสฺโส สุปณฺณโยนิโยติ.\csromanspacer{}\csromanspacer{}ทุติยํ.}

\csromanheader{2}{3. ทฺวยการีสุตฺต}
\proseitem{394}{สาวตฺถินิทานํ.\csromanspacer{} อญฺญตโร ภิกฺขุ เยน ภควา เตนุปสงฺกมิ, อุปสงฺกมิตฺวา ภควนฺตํ อภิวาเทตฺวา เอกมนฺตํ นิสีทิ, เอกมนฺตํ นิสินฺโน โข โส ภิกฺขุ ภควนฺตํ เอตทโวจ “โก นุ โข ภนฺเต เหตุ โก ปจฺจโย, เยน มิเธกจฺโจ กายสฺส เภทา ปรํ มรณา อณฺฑชานํ สุปณฺณานํ สหพฺยตํ อุปปชฺชตี”ติ.\csromanspacer{} อิธ ภิกฺขุ เอกจฺโจ กาเยน ทฺวยการี โหติ วาจาย ทฺวยการี มนสา ทฺวยการี, ตสฺส สุตํ โหติ “อณฺฑชา สุปณฺณา ทีฆายุกา วณฺณวนฺโต สุขพหุลา”ติ.\csromanspacer{} ตสฺส เอวํ โหติ “อโห วตาหํ กายสฺส เภทา ปรํ มรณา อณฺฑชานํ สุปณฺณานํ สหพฺยตํ อุปปชฺเชยฺยนฺ”ติ.\csromanspacer{} โส กายสฺส เภทา ปรํ มรณา อณฺฑชานํ สุปณฺณานํ สหพฺยตํ อุปปชฺชติ.\csromanspacer{} อยํ โข ภิกฺขุ เหตุ อยํ \csromanfolio{208}ปจฺจโย, เยน มิเธกจฺโจ กายสฺส เภทา ปรํ มรณา อณฺฑชานํ สุปณฺณานํ สหพฺยตํ อุปปชฺชตีติ.\csromanspacer{}\csromanspacer{}ตติยํ.}

\csromanheader{2}{4-6. ทุติยาทิทฺวยการีสุตฺตตฺติก}
\proseitem{395-397}{สาวตฺถินิทานํ.\csromanspacer{} เอกมนฺตํ นิสินฺโน โข โส ภิกฺขุ ภควนฺตํ เอตทโวจ “โก นุ โข ภนฺเต เหตุ โก ปจฺจโย, เยน มิเธกจฺโจ กายสฺส เภทา ปรํ มรณา ชลาพุชานํ สุปณฺณานํ ฯเปฯ สํเสทชานํ สุปณฺณานํ ฯเปฯ โอปปาติกานํ สุปณฺณานํ สหพฺยตํ อุปปชฺชตี”ติ.\csromanspacer{} อิธ ภิกฺขุ เอกจฺโจ กาเยน ทฺวยการี โหติ วาจาย ทฺวยการี มนสา ทฺวยการี, ตสฺส สุตํ โหติ “โอปปาติกา สุปณฺณา ทีฆายุกา วณฺณวนฺโต สุขพหุลา”ติ.\csromanspacer{} ตสฺส เอวํ โหติ “อโห วตาหํ กายสฺส เภทา ปรํ มรณา โอปปาติกานํ สุปณฺณานํ สหพฺยตํ อุปปชฺเชยฺยนฺ”ติ.\csromanspacer{} โส กายสฺส เภทา ปรํ มรณา โอปปาติกานํ สุปณฺณานํ สหพฺยตํ อุปปชฺชติ.\csromanspacer{} อยํ โข ภิกฺขุ เหตุ อยํ ปจฺจโย, เยน มิเธกจฺโจ กายสฺส เภทา ปรํ มรณา โอปปาติกานํ สุปณฺณานํ สหพฺยตํ อุปปชฺชตีติ.\csromanspacer{}\csromanspacer{}ฉฏฺฐํ.}

\csromanheader{2}{7-16. อณฺฑชทานูปการสุตฺตทสก}
\proseitem{398-407}{สาวตฺถินิทานํ.\csromanspacer{} เอกมนฺตํ นิสินฺโน โข โส ภิกฺขุ ภควนฺตํ เอตทโวจ “โก นุ โข ภนฺเต เหตุ โก ปจฺจโย, เยน มิเธกจฺโจ กายสฺส เภทา ปรํ มรณา อณฺฑชานํ สุปณฺณานํ สหพฺยตํ อุปปชฺชตี”ติ.\csromanspacer{} อิธ ภิกฺขุ เอกจฺโจ กาเยน ทฺวยการี โหติ วาจาย ทฺวยการี มนสา ทฺวยการี, ตสฺส สุตํ โหติ “อณฺฑชา สุปณฺณา ทีฆายุกา วณฺณวนฺโต สุขพหุลา”ติ.\csromanspacer{} ตสฺส เอวํ โหติ “อโห วตาหํ กายสฺส เภทา ปรํ มรณา อณฺฑชานํ สุปณฺณานํ สหพฺยตํ อุปปชฺเชยฺยนฺ”ติ.\csromanspacer{} โส อนฺนํ เทติ ฯเปฯ ปานํ เทติ.\csromanspacer{} วตฺถํ เทติ.\csromanspacer{} ยานํ เทติ.\csromanspacer{} มาลํ เทติ.\csromanspacer{} คนฺธํ เทติ.\csromanspacer{} วิเลปนํ เทติ.\csromanspacer{} เสยฺยํ เทติ.\csromanspacer{} อาวสถํ เทติ.\csromanspacer{} ปทีเปยฺยํ เทติ, โส กายสฺส เภทา ปรํ มรณา อณฺฑชานํ สุปณฺณานํ สหพฺยตํ อุปปชฺชติ.\csromanspacer{} อยํ โข ภิกฺขุ เหตุ อยํ ปจฺจโย, เยน มิเธกจฺโจ กายสฺส เภทา ปรํ มรณา อณฺฑชานํ สุปณฺณานํ สหพฺยตํ อุปปชฺชตีติ.\csromanspacer{}\csromanspacer{}โสฬสมํ.}

\csromanheader{2}{9. สุปณฺณสํยุตฺต \textbf{17-46. ชลาพุชทานูปการสุตฺต}}
\proseitem{408-437}{\csromanfolio{209}สาวตฺถินิทานํ.\csromanspacer{} เอกมนฺตํ นิสินฺโน โข โส ภิกฺขุ ภควนฺตํ เอตทโวจ “โก นุ โข ภนฺเต เหตุ โก ปจฺจโย, เยน มิเธกจฺโจ กายสฺส เภทา ปรํ มรณา ชลาพุชานํ สุปณฺณานํ ฯเปฯ สํเสทชานํ สุปณฺณานํ ฯเปฯ โอปปาติกานํ สุปณฺณานํ สหพฺยตํ อุปปชฺชตี”ติ.\csromanspacer{} อิธ ภิกฺขุ เอกจฺโจ กาเยน ทฺวยการี โหติ วาจาย ทฺวยการี มนสา ทฺวยการี, ตสฺส สุตํ โหติ “โอปปาติกา สุปณฺณา ทีฆายุกา วณฺณวนฺโต สุขพหุลา”ติ.\csromanspacer{} ตสฺส เอวํ โหติ “อโห วตาหํ กายสฺส เภทา ปรํ มรณา โอปปาติกานํ สุปณฺณานํ สหพฺยตํ อุปปชฺเชยฺยนฺ”ติ, โส อนฺนํ เทติ ฯเปฯ ปานํ เทติ ฯเปฯ.\csromanspacer{} ปทีเปยฺยํ เทติ.\csromanspacer{} โส กายสฺส เภทา ปรํ มรณา โอปปาติกานํ สุปณฺณานํ สหพฺยตํ อุปปชฺชติ.\csromanspacer{} อยํ โข ภิกฺขุ เหตุ อยํ ปจฺจโย, เยน มิเธกจฺโจ กายสฺส เภทา ปรํ มรณา โอปปาติกานํ สุปณฺณานํ สหพฺยตํ อุปปชฺชตีติ.\csromanspacer{}\csromanspacer{}ฉจตฺตาลีสมํ.}

\vspace{\tipitakaparskip}
\csromancenter{(เอวํ ปิณฺฑเกน ฉจตฺตาลีสํ สุตฺตนฺตา โหนฺติ.)}
\csromancenter{\textbf{สุปณฺณสํยุตฺตํ สมตฺตํ.}}
\titlehead{\textbf{ตสฺสุทฺทานํ}}
\csromangathameasure{สุทฺธิกํ หรนฺติ เจว, ทฺวยการี จ จตุโร. \\ ทานูปการา ตาลีสํ, สุปณฺเณ สุปฺปกาสิตาติ.}
\csromangathagroup{\csromangathastanza{สุทฺธิกํ หรนฺติ เจว, ทฺวยการี จ จตุโร. \\ ทานูปการา ตาลีสํ, สุปณฺเณ สุปฺปกาสิตาติ.}}

\csromanmatikaanchor{mtk.71}
\csromantocmarknopage{chapter}{10. คนฺธพฺพกายสํยุตฺต}{mtk.71}
\bookmark[dest={mtk.71},level=0]{10. คนฺธพฺพกายสํยุตฺต}
\chapterheadpageread{10. คนฺธพฺพกายสํยุตฺต}
\csromanheader{2}{1. สุทฺธิกสุตฺต}
\csromanmatikaanchor{mtk.72}
\csromantocmarkref{section}{1-112. สุทฺธิกสุตฺตาทิ}{mtk.72}
\bookmark[dest={mtk.72},level=1]{1-112. สุทฺธิกสุตฺตาทิ}
\proseitem{438}{\csromanfolio{210}เอกํ สมยํ ภควา สาวตฺถิยํ วิหรติ เชตวเน อนาถปิณฺฑิกสฺส อาราเม ฯเปฯ ภควา เอตทโวจ “คนฺธพฺพกายิเก โว ภิกฺขเว เทเว เทเสสฺสามิ, ตํ สุณาถ.\csromanspacer{} กตมา จ ภิกฺขเว คนฺธพฺพกายิกา เทวา, สนฺติ ภิกฺขเว มูลคนฺเธ อธิวตฺถา เทวา, สนฺติ ภิกฺขเว สารคนฺเธ อธิวตฺถา เทวา, สนฺติ ภิกฺขเว เผคฺคุคนฺเธ อธิวตฺถา เทวา, สนฺติ ภิกฺขเว ตจคนฺเธ อธิวตฺถา เทวา, สนฺติ ภิกฺขเว ปปฏิกคนฺเธ อธิวตฺถา เทวา, สนฺติ ภิกฺขเว ปตฺตคนฺเธ อธิวตฺถา เทวา, สนฺติ ภิกฺขเว ปุปฺผคนฺเธ อธิวตฺถา เทวา, สนฺติ ภิกฺขเว ผลคนฺเธ อธิวตฺถา เทวา, สนฺติ ภิกฺขเว รสคนฺเธ อธิวตฺถา เทวา, สนฺติ ภิกฺขเว คนฺธคนฺเธ อธิวตฺถา เทวา.\csromanspacer{} อิเม วุจฺจนฺติ ภิกฺขเว คนฺธพฺพกายิกา เทวา”ติ.\csromanspacer{}\csromanspacer{}ปฐมํ.}

\csromanheader{2}{2. สุจริตสุตฺต}
\proseitem{439}{สาวตฺถินิทานํ.\csromanspacer{} เอกมนฺตํ นิสินฺโน โข โส ภิกฺขุ ภควนฺตํ เอตทโวจ “โก นุ โข ภนฺเต เหตุ โก ปจฺจโย, เยน มิเธกจฺโจ กายสฺส เภทา ปรํ มรณา คนฺธพฺพกายิกานํ เทวานํ สหพฺยตํ อุปปชฺชตี”ติ.\csromanspacer{} อิธ ภิกฺขุ เอกจฺโจ กาเยน สุจริตํ จรติ, วาจาย สุจริตํ จรติ, มนสา สุจริตํ จรติ, ตสฺส สุตํ โหติ “คนฺธพฺพกายิกา เทวา ทีฆายุกา วณฺณวนฺโต สุขพหุลา”ติ.\csromanspacer{} ตสฺส เอวํ โหติ “อโห วตาหํ กายสฺส เภทา ปรํ มรณา คนฺธพฺพกายิกานํ เทวานํ สหพฺยตํ อุปปชฺเชยฺยนฺ”ติ.\csromanspacer{} โส กายสฺส เภทา ปรํ มรณา คนฺธพฺพกายิกานํ เทวานํ สหพฺยตํ อุปปชฺชติ.\csromanspacer{} อยํ โข ภิกฺขุ เหตุ อยํ ปจฺจโย, เยน มิเธกจฺโจ กายสฺส เภทา ปรํ มรณา คนฺธพฺพกายิกานํ เทวานํ สหพฺยตํ อุปปชฺชตีติ.\csromanspacer{}\csromanspacer{}ทุติยํ.}

\csromanheader{2}{3. มูลคนฺธทาตาสุตฺต}
\proseitem{440}{สาวตฺถินิทานํ.\csromanspacer{} เอกมนฺตํ นิสินฺโน โข โส ภิกฺขุ ภควนฺตํ เอตทโวจ “โก นุ โข ภนฺเต เหตุ โก ปจฺจโย, เยน มิเธกจฺโจ กายสฺส เภทา ปรํ มรณา มูลคนฺเธ อธิวตฺถานํ เทวานํ สหพฺยตํ}

\vspace{\tipitakaparskip}
\csromancenter{คนฺธพฺพกายสํยุตฺต อุปปชฺชตี”ติ.\csromanspacer{} อิธ ภิกฺขุ เอกจฺโจ กาเยน สุจริตํ จรติ, วาจาย สุจริตํ จรติ, มนสา สุจริตํ จรติ, ตสฺส สุตํ โหติ “มูลคนฺเธ อธิวตฺถา เทวา ทีฆายุกา วณฺณวนฺโต สุขพหุลา”ติ.\csromanspacer{} ตสฺส เอวํ โหติ “อโห วตาหํ กายสฺส เภทา ปรํ มรณา มูลคนฺเธ อธิวตฺถานํ เทวานํ สหพฺยตํ อุปปชฺเชยฺยนฺ”ติ, โส ทาตา โหติ มูลคนฺธานํ.\csromanspacer{} โส กายสฺส เภทา ปรํ มรณา มูลคนฺเธ อธิวตฺถานํ เทวานํ สหพฺยตํ อุปปชฺชติ.\csromanspacer{} อยํ โข ภิกฺขุ เหตุ ฯเปฯ เยน มิเธกจฺโจ กายสฺส เภทา ปรํ มรณา มูลคนฺเธ อธิวตฺถานํ เทวานํ สหพฺยตํ อุปปชฺชตีติ.\csromanspacer{}\csromanspacer{}ตติยํ.}
\csromanheader{2}{4-12. สารคนฺธาทิทาตาสุตฺตนวก}
\proseitem{441-449}{\csromanfolio{211}สาวตฺถินิทานํ.\csromanspacer{} เอกมนฺตํ นิสินฺโน โข โส ภิกฺขุ ภควนฺตํ เอตทโวจ “โก นุ โข ภนฺเต เหตุ โก ปจฺจโย, เยน มิเธกจฺโจ กายสฺส เภทา ปรํ มรณา สารคนฺเธ อธิวตฺถานํ เทวานํ ฯเปฯ เผคฺคุคนฺเธ อธิวตฺถานํ เทวานํ.\csromanspacer{} ตจคนฺเธ อธิวตฺถานํ เทวานํ.\csromanspacer{} ปปฏิกคนฺเธ อธิวตฺถานํ เทวานํ.\csromanspacer{} ปตฺตคนฺเธ อธิวตฺถานํ เทวานํ.\csromanspacer{} ปุปฺผคนฺเธ อธิวตฺถานํ เทวานํ.\csromanspacer{} ผลคนฺเธ อธิวตฺถานํ เทวานํ.\csromanspacer{} รสคนฺเธ อธิวตฺถานํ เทวานํ.\csromanspacer{} คนฺธคนฺเธ อธิวตฺถานํ เทวานํ สหพฺยตํ อุปปชฺชตี”ติ.\csromanspacer{} อิธ ภิกฺขุ เอกจฺโจ กาเยน สุจริตํ จรติ, วาจาย สุจริตํ จรติ, มนสา สุจริตํ จรติ, ตสฺส สุตํ โหติ “สารคนฺเธ อธิวตฺถา เทวา ทีฆายุกา วณฺณวนฺโต สุขพหุลา”ติ.\csromanspacer{} ตสฺส เอวํ โหติ “อโห วตาหํ กายสฺส เภทา ปรํ มรณา สารคนฺเธ อธิวตฺถานํ เทวานํ ฯเปฯ เผคฺคุคนฺเธ อธิวตฺถานํ เทวานํ.\csromanspacer{} ตจคนฺเธ อธิวตฺถานํ เทวานํ.\csromanspacer{} ปปฏิกคนฺเธ อธิวตฺถานํ เทวานํ.\csromanspacer{} ปตฺตคนฺเธ อธิวตฺถานํ เทวานํ.\csromanspacer{} ปุปฺผคนฺเธ อธิวตฺถานํ เทวานํ.\csromanspacer{} ผลคนฺเธ อธิวตฺถานํ เทวานํ.\csromanspacer{} รสคนฺเธ อธิวตฺถานํ เทวานํ.\csromanspacer{} คนฺธคนฺเธ อธิวตฺถานํ เทวานํ สหพฺยตํ อุปปชฺเชยฺยนฺ”ติ, โส ทาตา โหติ สารคนฺธานํ ฯเปฯ โส ทาตา โหติ เผคฺคุคนฺธานํ.\csromanspacer{} โส ทาตา โหติ ตจคนฺธานํ.\csromanspacer{} โส ทาตา โหติ ปปฏิกคนฺธานํ.\csromanspacer{} โส ทาตา โหติ ปตฺตคนฺธานํ.\csromanspacer{} โส ทาตา โหติ ปุปฺผคนฺธานํ.\csromanspacer{} โส ทาตา โหติ ผลคนฺธานํ.\csromanspacer{} โส ทาตา โหติ รสคนฺธานํ.\csromanspacer{} โส ทาตา โหติ คนฺธคนฺธานํ.\csromanspacer{} โส กายสฺส เภทา ปรํ มรณา คนฺธคนฺเธ อธิวตฺถานํ เทวานํ สหพฺยตํ อุปปชฺชติ.\csromanspacer{} อยํ โข ภิกฺขุ เหตุ \csromanfolio{212}อยํ ปจฺจโย, เยน มิเธกจฺโจ กายสฺส เภทา ปรํ มรณา คนฺธคนฺเธ อธิวตฺถานํ เทวานํ สหพฺยตํ อุปปชฺชตีติ.\csromanspacer{}\csromanspacer{}ทฺวาทสมํ.}

\csromanheader{2}{13-22. มูลคนฺธทานูปการสุตฺตทสก}
\proseitem{450-459}{สาวตฺถินิทานํ.\csromanspacer{} เอกมนฺตํ นิสินฺโน โข โส ภิกฺขุ ภควนฺตํ เอตทโวจ “โก นุ โข ภนฺเต เหตุ โก ปจฺจโย, เยน มิเธกจฺโจ กายสฺส เภทา ปรํ มรณา มูลคนฺเธ อธิวตฺถานํ เทวานํ สหพฺยตํ อุปปชฺชตี”ติ.\csromanspacer{} อิธ ภิกฺขุ เอกจฺโจ กาเยน สุจริตํ จรติ, วาจาย สุจริตํ จรติ, มนสา สุจริตํ จรติ, ตสฺส สุตํ โหติ “มูลคนฺเธ อธิวตฺถา เทวา ทีฆายุกา วณฺณวนฺโต สุขพหุลา”ติ.\csromanspacer{} ตสฺส เอวํ โหติ “อโห วตาหํ กายสฺส เภทา ปรํ มรณา มูลคนฺเธ อธิวตฺถานํ เทวานํ สหพฺยตํ อุปปชฺเชยฺยนฺ”ติ, โส อนฺนํ เทติ ฯเปฯ ปานํ เทติ.\csromanspacer{} วตฺถํ เทติ.\csromanspacer{} ยานํ เทติ.\csromanspacer{} มาลํ เทติ.\csromanspacer{} คนฺธํ เทติ.\csromanspacer{} วิเลปนํ เทติ.\csromanspacer{} เสยฺยํ เทติ.\csromanspacer{} อาวสถํ เทติ.\csromanspacer{} ปทีเปยฺยํ เทติ.\csromanspacer{} โส กายสฺส เภทา ปรํ มรณา มูลคนฺเธ อธิวตฺถานํ เทวานํ สหพฺยตํ อุปปชฺชติ.\csromanspacer{} อยํ โข ภิกฺขุ เหตุ อยํ ปจฺจโย, เยน มิเธกจฺโจ กายสฺส เภทา ปรํ มรณา มูลคนฺเธ อธิวตฺถานํ เทวานํ สหพฺยตํ อุปปชฺชตีติ.\csromanspacer{}\csromanspacer{}พาวีสติมํ.}

\csromanheader{2}{23-112. สารคนฺธาทิทานูปการสุตฺตนวุติก}
\proseitem{460-549}{สาวตฺถินิทานํ.\csromanspacer{} เอกมนฺตํ นิสินฺโน โข โส ภิกฺขุ ภควนฺตํ เอตทโวจ “โก นุ โข ภนฺเต เหตุ โก ปจฺจโย, เยน มิเธกจฺโจ กายสฺส เภทา ปรํ มรณา สารคนฺเธ อธิวตฺถานํ เทวานํ ฯเปฯ เผคฺคุคนฺเธ อธิวตฺถานํ เทวานํ.\csromanspacer{} ตจคนฺเธ อธิวตฺถานํ เทวานํ.\csromanspacer{} ปปฏิกคนฺเธ อธิวตฺถานํ เทวานํ.\csromanspacer{} ปตฺตคนฺเธ อธิวตฺถานํ เทวานํ.\csromanspacer{} ปุปฺผคนฺเธ อธิวตฺถานํ เทวานํ.\csromanspacer{} ผลคนฺเธ อธิวตฺถานํ เทวานํ.\csromanspacer{} รสคนฺเธ อธิวตฺถานํ เทวานํ.\csromanspacer{} คนฺธคนฺเธ อธิวตฺถานํ เทวานํ สหพฺยตํ อุปปชฺชตี”ติ.\csromanspacer{} อิธ ภิกฺขุ เอกจฺโจ กาเยน สุจริตํ จรติ, วาจาย สุจริตํ จรติ, มนสา สุจริตํ จรติ, ตสฺส สุตํ โหติ “คนฺธคนฺเธ อธิวตฺถา เทวา ทีฆายุกา วณฺณวนฺโต สุขพหุลา”ติ.\csromanspacer{} ตสฺส เอวํ โหติ “อโห วตาหํ กายสฺส เภทา ปรํ มรณา คนฺธคนฺเธ อธิวตฺถานํ เทวานํ สหพฺยตํ อุปปชฺเชยฺยนฺ”ติ, โส อนฺนํ เทติ ฯเปฯ ปานํ เทติ.\csromanspacer{} วตฺถํ เทติ.\csromanspacer{} ยานํ เทติ.\csromanspacer{} มาลํ เทติ.\csromanspacer{} คนฺธํ เทติ.}

\vspace{\tipitakaparskip}
\csromancenter{คนฺธพฺพกายสํยุตฺต วิเลปนํ เทติ.\csromanspacer{} เสยฺยํ เทติ.\csromanspacer{} อาวสถํ เทติ.\csromanspacer{} ปทีเปยฺยํ เทติ.\csromanspacer{} โส กายสฺส เภทา ปรํ มรณา คนฺธคนฺเธ อธิวตฺถานํ เทวานํ สหพฺยตํ อุปปชฺชติ.\csromanspacer{} อยํ โข ภิกฺขุ เหตุ อยํ ปจฺจโย, เยน มิเธกจฺโจ กายสฺส เภทา ปรํ มรณา คนฺธคนฺเธ อธิวตฺถานํ เทวานํ สหพฺยตํ อุปปชฺชตีติ.\csromanspacer{}\csromanspacer{}ทฺวาทสสติมํ.}
\csromancenter{(เอวํ ปิณฺฑเกน เอกสตญฺจ ทฺวาทส จ สุตฺตนฺตา โหนฺติ.)}
\csromancenter{\textbf{คนฺธพฺพกายสํยุตฺตํ สมตฺตํ.}}
\titlehead{\textbf{ตสฺสุทฺทานํ}}
\csromangathameasure{สุทฺธิกํ จ สุจริตํ, ทาตา หิ อปเร ทส. \\ ทานูปการา สตธา, คนฺธพฺเพ สุปฺปกาสิตาติ.}
\csromangathagroup{\csromangathastanza{\csromanfolio{213}สุทฺธิกํ จ สุจริตํ, ทาตา หิ อปเร ทส. \\ ทานูปการา สตธา, คนฺธพฺเพ สุปฺปกาสิตาติ.}}

\csromanmatikaanchor{mtk.73}
\csromantocmarknopage{chapter}{11. วลาหกสํยุตฺต}{mtk.73}
\bookmark[dest={mtk.73},level=0]{11. วลาหกสํยุตฺต}
\chapterheadpageread{11. วลาหกสํยุตฺต}
\csromanheader{2}{1. สุทฺธิกสุตฺต}
\csromanmatikaanchor{mtk.74}
\csromantocmarkref{section}{1-57. สุทฺธิกสุตฺตาทิ}{mtk.74}
\bookmark[dest={mtk.74},level=1]{1-57. สุทฺธิกสุตฺตาทิ}
\proseitem{550}{\csromanfolio{214}สาวตฺถินิทานํ.\csromanspacer{} วลาหกกายิเก โว ภิกฺขเว เทเว เทเสสฺสามิ, ตํ สุณาถ.\csromanspacer{} กตเม จ ภิกฺขเว วลาหกกายิกา เทวา, สนฺติ ภิกฺขเว สีตวลาหกา เทวา, สนฺติ อุณฺหวลาหกา เทวา, สนฺติ อพฺภวลาหกา เทวา, สนฺติ วาตวลาหกา เทวา, สนฺติ วสฺสวลาหกา เทวา.\csromanspacer{} อิเม วุจฺจนฺติ ภิกฺขเว วลาหกกายิกา เทวาติ.\csromanspacer{}\csromanspacer{}ปฐมํ.}

\csromanheader{2}{2. สุจริตสุตฺต}
\proseitem{551}{สาวตฺถินิทานํ.\csromanspacer{} เอกมนฺตํ นิสินฺโน โข โส ภิกฺขุ ภควนฺตํ เอตทโวจ “โก นุ โข ภนฺเต เหตุ โก ปจฺจโย, เยน มิเธกจฺโจ กายสฺส เภทา ปรํ มรณา วลาหกกายิกานํ เทวานํ สหพฺยตํ อุปปชฺชตี”ติ.\csromanspacer{} อิธ ภิกฺขุ เอกจฺโจ กาเยน สุจริตํ จรติ, วาจาย สุจริตํ จรติ, มนสา สุจริตํ จรติ, ตสฺส สุตํ โหติ “วลาหกกายิกา เทวา ทีฆายุกา วณฺณวนฺโต สุขพหุลา”ติ.\csromanspacer{} ตสฺส เอวํ โหติ “อโห วตาหํ กายสฺส เภทา ปรํ มรณา วลาหกกายิกานํ เทวานํ สหพฺยตํ อุปปชฺเชยฺยนฺ”ติ.\csromanspacer{} โส กายสฺส เภทา ปรํ มรณา วลาหกกายิกานํ เทวานํ สหพฺยตํ อุปปชฺชติ.\csromanspacer{} อยํ โข ภิกฺขุ เหตุ อยํ ปจฺจโย, เยน มิเธกจฺโจ กายสฺส เภทา ปรํ มรณา วลาหกกายิกานํ เทวานํ สหพฺยตํ อุปปชฺชตีติ.\csromanspacer{}\csromanspacer{}ทุติยํ.}

\csromanheader{2}{3-12. สีตวลาหกทานูปการสุตฺตทสก}
\proseitem{552-561}{สาวตฺถินิทานํ.\csromanspacer{} เอกมนฺตํ นิสินฺโน โข โส ภิกฺขุ ภควนฺตํ เอตทโวจ “โก นุ โข ภนฺเต เหตุ โก ปจฺจโย, เยน มิเธกจฺโจ กายสฺส เภทา ปรํ มรณา สีตวลาหกานํ เทวานํ สหพฺยตํ อุปปชฺชตี”ติ.\csromanspacer{} อิธ ภิกฺขุ เอกจฺโจ กาเยน สุจริตํ จรติ, วาจาย สุจริตํ จรติ, มนสา สุจริตํ จรติ, ตสฺส สุตํ โหติ “สีตวลาหกา เทวา ทีฆายุกา วณฺณวนฺโต สุขพหุลา”ติ.\csromanspacer{} ตสฺส}

\csromanheader{2}{11. วลาหกสํยุตฺต}
\prose{\csromanfolio{215}เอวํ โหติ “อโห วตาหํ กายสฺส เภทา ปรํ มรณา สีตวลาหกานํ เทวานํ สหพฺยตํ อุปปชฺเชยฺยนฺ”ติ, โส อนฺนํ เทติ ฯเปฯ ปทีเปยฺยํ เทติ.\csromanspacer{} โส กายสฺส เภทา ปรํ มรณา สีตวลาหกานํ เทวานํ สหพฺยตํ อุปปชฺชติ.\csromanspacer{} อยํ โข ภิกฺขุ เหตุ อยํ ปจฺจโย, เยน มิเธกจฺโจ กายสฺส เภทา ปรํ มรณา สีตวลาหกานํ เทวานํ สหพฺยตํ อุปปชฺชตีติ.\csromanspacer{}\csromanspacer{}ทฺวาทสมํ.}

\csromanheader{2}{13-52. อุณฺหวลาหกทานูปการสุตฺต}
\proseitem{562-601}{สาวตฺถินิทานํ.\csromanspacer{} เอกมนฺตํ นิสินฺโน โข โส ภิกฺขุ ภควนฺตํ เอตทโวจ “โก นุ โข ภนฺเต เหตุ โก ปจฺจโย, เยน มิเธกจฺโจ กายสฺส เภทา ปรํ มรณา อุณฺหวลาหกานํ เทวานํ ฯเปฯ อพฺภวลาหกานํ เทวานํ วาตวลาหกานํ เทวานํ ฯเปฯ วสฺสวลาหกานํ เทวานํ สหพฺยตํ อุปปชฺชตี”ติ.\csromanspacer{} อิธ ภิกฺขุ เอกจฺโจ กาเยน สุจริตํ จรติ, วาจาย สุจริตํ จรติ, มนสา สุจริตํ จรติ, ตสฺส สุตํ โหติ “วสฺสวลาหกา เทวา ทีฆายุกา วณฺณวนฺโต สุขพหุลา”ติ.\csromanspacer{} ตสฺส เอวํ โหติ “อโห วตาหํ กายสฺส เภทา ปรํ มรณา วสฺสวลาหกานํ เทวานํ สหพฺยตํ อุปปชฺเชยฺยนฺ”ติ, โส อนฺนํ เทติ ฯเปฯ ปทีเปยฺยํ เทติ.\csromanspacer{} โส กายสฺส เภทา ปรํ มรณา วสฺสวลาหกานํ เทวานํ สหพฺยตํ อุปปชฺชติ.\csromanspacer{} อยํ โข ภิกฺขุ เหตุ อยํ ปจฺจโย, เยน มิเธกจฺโจ กายสฺส เภทา ปรํ มรณา วสฺสวลาหกานํ เทวานํ สหพฺยตํ อุปปชฺชตีติ.\csromanspacer{}\csromanspacer{}เทฺวปญฺญาสมํ.}

\csromanheader{2}{53. สีตวลาหกสุตฺต}
\proseitem{602}{สาวตฺถินิทานํ.\csromanspacer{} เอกมนฺตํ นิสินฺโน โข โส ภิกฺขุ ภควนฺตํ เอตทโวจ “โก นุ โข ภนฺเต เหตุ โก ปจฺจโย, เยเนกทา สีตํ โหตี”ติ.\csromanspacer{} สนฺติ ภิกฺขุ สีตวลาหกา นาม เทวา, เตสํ ยทา เอวํ โหติ “ยํนูน มยํ สกาย รติยา วเสยฺยามา”ติ\footnote{รเมยฺยามาติ (สี, สฺยา, กํ, อิ) เอวมุปริปิ.}, เตสํ ตํ เจโตปณิธิมนฺวาย สีตํ โหติ.\csromanspacer{} อยํ โข ภิกฺขุ เหตุ อยํ ปจฺจโย, เยเนกทา สีตํ โหตีติ.\csromanspacer{}\csromanspacer{}เตปญฺญาสมํ.}

\csromanheader{2}{54. อุณฺหวลาหกสุตฺต}
\proseitem{603}{\csromanfolio{216}สาวตฺถินิทานํ.\csromanspacer{} เอกมนฺตํ นิสินฺโน โข โส ภิกฺขุ ภควนฺตํ เอตทโวจ “โก นุ โข ภนฺเต เหตุ โก ปจฺจโย, เยเนกทา อุณฺหํ โหตี”ติ.\csromanspacer{} สนฺติ ภิกฺขุ อุณฺหวลาหกา นาม เทวา, เตสํ ยทา เอวํ โหติ “ยํนูน มยํ สกาย รติยา วเสยฺยามา”ติ, เตสํ ตํ เจโตปณิธิมนฺวาย อุณฺหํ โหติ.\csromanspacer{} อยํ โข ภิกฺขุ เหตุ อยํ ปจฺจโย, เยเนกทา อุณฺหํ โหตีติ.\csromanspacer{}\csromanspacer{}จตุปญฺญาสมํ.}

\csromanheader{2}{55. อพฺภวลาหกสุตฺต}
\proseitem{604}{สาวตฺถินิทานํ.\csromanspacer{} เอกมนฺตํ นิสินฺโน โข โส ภิกฺขุ ภควนฺตํ เอตทโวจ “โก นุ โข ภนฺเต เหตุ โก ปจฺจโย, เยเนกทา อพฺภํ โหตี”ติ.\csromanspacer{} สนฺติ ภิกฺขุ อพฺภวลาหกา นาม เทวา, เตสํ ยทา เอวํ โหติ “ยํนูน มยํ สกาย รติยา วเสยฺยามา”ติ, เตสํ ตํ เจโตปณิธิมนฺวาย อพฺภํ โหติ.\csromanspacer{} อยํ โข ภิกฺขุ เหตุ อยํ ปจฺจโย, เยเนกทา อพฺภํ โหตีติ.\csromanspacer{}\csromanspacer{}ปญฺจปญฺญาสมํ.}

\csromanheader{2}{56. วาตวลาหกสุตฺต}
\proseitem{605}{สาวตฺถินิทานํ.\csromanspacer{} เอกมนฺตํ นิสินฺโน โข โส ภิกฺขุ ภควนฺตํ เอตทโวจ “โก นุ โข ภนฺเต เหตุ โก ปจฺจโย, เยเนกทา วาโต โหตี”ติ.\csromanspacer{} สนฺติ ภิกฺขุ วาตวลาหกา นาม เทวา, เตสํ ยทา เอวํ โหติ “ยํนูน มยํ สกาย รติยา วเสยฺยามา”ติ, เตสํ ตํ เจโตปณิธิมนฺวาย วาโต โหติ.\csromanspacer{} อยํ โข ภิกฺขุ เหตุ อยํ ปจฺจโย, เยเนกทา วาโต โหตีติ.\csromanspacer{}\csromanspacer{}ฉปฺปญฺญาสมํ.}

\csromanheader{2}{57. วสฺสวลาหกสุตฺต}
\proseitem{606}{สาวตฺถินิทานํ.\csromanspacer{} เอกมนฺตํ นิสินฺโน โข โส ภิกฺขุ ภควนฺตํ เอตทโวจ “โก นุ โข ภนฺเต เหตุ โก ปจฺจโย, เยเนกทา เทโว วสฺสตี”ติ.\csromanspacer{} สนฺติ ภิกฺขุ วสฺสวลาหกา นาม เทวา, เตสํ}

\csromanheader{2}{11. วลาหกสํยุตฺต}
\prose{\csromanfolio{217}ยทา เอวํ โหติ “ยํนูน มยํ สกาย รติยา วเสยฺยามา”ติ, เตสํ ตํ เจโตปณิธิมนฺวาย เทโว วสฺสติ.\csromanspacer{} อยํ โข ภิกฺขุ เหตุ อยํ ปจฺจโย, เยเนกทา เทโว วสฺสตีติ.\csromanspacer{}\csromanspacer{}สตฺตปญฺญาสมํ.}

\nitthitam{สตฺตปญฺญาสสุตฺตนฺตํ นิฏฺฐิตํ.}
\vspace{\tipitakaparskip}
\csromancenter{\textbf{วลาหกสํยุตฺตํ สมตฺตํ.}}
\titlehead{\textbf{ตสฺสุทฺทานํ}}
\csromangathameasure{สุทฺธิกํ สุจริตญฺจ, ทานูปการปญฺญาสํ. \\ สีตํ อุณฺหญฺจ อพฺภญฺจ, วาต วสฺสวลาหกาติ.}
\csromangathagroup{\csromangathastanza{สุทฺธิกํ สุจริตญฺจ, ทานูปการปญฺญาสํ. \\ สีตํ อุณฺหญฺจ อพฺภญฺจ, วาต วสฺสวลาหกาติ.}}

\csromanmatikaanchor{mtk.75}
\csromantocmarknopage{chapter}{12. วจฺฉโคตฺตสํยุตฺต}{mtk.75}
\bookmark[dest={mtk.75},level=0]{12. วจฺฉโคตฺตสํยุตฺต}
\chapterheadpageread{12. วจฺฉโคตฺตสํยุตฺต}
\csromanmatikaanchor{mtk.76}
\csromantocmarkref{section}{1-55. รูป-อญฺญาณสุตฺตาทิ}{mtk.76}
\bookmark[dest={mtk.76},level=1]{1-55. รูป-อญฺญาณสุตฺตาทิ}
\csromanheader{1}{1. รูป-อญฺญาณสุตฺต}
\proseitem{607}{\csromanfolio{218}เอกํ สมยํ ภควา สาวตฺถิยํ วิหรติ เชตวเน อนาถปิณฺฑิกสฺส อาราเม.\csromanspacer{} อถ โข วจฺฉโคตฺโต ปริพฺพาชโก เยน ภควา เตนุปสงฺกมิ, อุปสงฺกมิตฺวา ภควตา สทฺธิํ สมฺโมทิ, สมฺโมทนียํ กถํ สารณียํ วีติสาเรตฺวา เอกมนฺตํ นิสีทิ, เอกมนฺตํ นิสินฺโน โข วจฺฉโคตฺโต ปริพฺพาชโก ภควนฺตํ เอตทโวจ “โก นุ โข โภ โคตม เหตุ โก ปจฺจโย, ยานิมานิ\footnote{เยนิมานิ (?)} อเนกวิหิตานิ ทิฏฺฐิคตานิ โลเก อุปฺปชฺชนฺติ ‘สสฺสโต โลโก’ติ วา, ‘อสสฺสโต โลโก’ติ วา, ‘อนฺตวา โลโก’ติ วา, ‘อนนฺตวา โลโก’ติ วา, ‘ตํ ชีวํ ตํ สรีรนฺ’ติ วา, ‘อญฺญํ ชีวํ อญฺญํ สรีรนฺ’ติ วา, ‘โหติ ตถาคโต ปรํ มรณา’ติ วา, ‘น โหติ ตถาคโต ปรํ มรณา’ติ วา, ‘โหติ จ น จ โหติ ตถาคโต ปรํ มรณา’ติ วา, ‘เนว โหติ น น โหติ ตถาคโต ปรํ มรณา’ติ วา”ติ.\csromanspacer{} รูเป โข วจฺฉ อญฺญาณา รูปสมุทเย อญฺญาณา รูปนิโรเธ อญฺญาณา รูปนิโรธคามินิยา ปฏิปทาย อญฺญาณา, เอวมิมานิ อเนกวิหิตานิ ทิฏฺฐิคตานิ โลเก อุปฺปชฺชนฺติ “สสฺสโต โลโก”ติ วา ฯเปฯ “เนว โหติ น น โหติ ตถาคโต ปรํ มรณา”ติ วาติ.\csromanspacer{} อยํ โข วจฺฉ เหตุ อยํ ปจฺจโย, ยานิมานิ\footnote{เยน (สี), เยนิมานิ (?)} อเนกวิหิตานิ ทิฏฺฐิคตานิ โลเก อุปฺปชฺชนฺติ “สสฺสโต โลโก”ติ วา, “อสสฺสโต โลโก”ติ วา ฯเปฯ “เนว โหติ น น โหติ ตถาคโต ปรํ มรณา”ติ วาติ.\csromanspacer{}\csromanspacer{}ปฐมํ.}

\csromanheader{2}{2. เวทนา-อญฺญาณสุตฺต}
\proseitem{608}{สาวตฺถินิทานํ.\csromanspacer{} เอกมนฺตํ นิสินฺโน โข วจฺฉโคตฺโต ปริพฺพาชโก ภควนฺตํ เอตทโวจ “โก นุ โข โภ โคตม เหตุ โก ปจฺจโย, ยานิมานิ อเนกวิหิตานิ ทิฏฺฐิคตานิ โลเก อุปฺปชฺชนฺติ ‘สสฺสโต โลโก’ติ วา, ‘อสสฺสโต โลโก’ติ วา ฯเปฯ ‘เนว โหติ น น โหติ ตถาคโต ปรํ มรณา’ติ วา”ติ.\csromanspacer{} เวทนาย โข วจฺฉ อญฺญาณา เวทนาสมุทเย อญฺญาณา เวทนานิโรเธ อญฺญาณา}

\csromanheader{2}{12. วจฺฉโคตฺตสํยุตฺต}
\prose{\csromanfolio{219}เวทนานิโรธคามินิยา ปฏิปทาย อญฺญาณา, เอวมิมานิ อเนกวิหิตานิ ทิฏฺฐิคตานิ โลเก อุปฺปชฺชนฺติ “สสฺสโต โลโก”ติ วา, “อสสฺสโต โลโก”ติ วา ฯเปฯ “เนว โหติ น น โหติ ตถาคโต ปรํ มรณา”ติ วาติ.\csromanspacer{} อยํ โข วจฺฉ เหตุ อยํ ปจฺจโย, ยานิมานิ อเนกวิหิตานิ ทิฏฺฐิคตานิ โลเก อุปฺปชฺชนฺติ “สสฺสโต โลโก”ติ วา, “อสสฺสโต โลโก”ติ วา ฯเปฯ “เนว โหติ น น โหติ ตถาคโต ปรํ มรณา”ติ วาติ.\csromanspacer{}\csromanspacer{}ทุติยํ.}

\csromanheader{2}{3. สญฺญา-อญฺญาณสุตฺต}
\proseitem{609}{สาวตฺถินิทานํ.\csromanspacer{} เอกมนฺตํ นิสินฺโน โข วจฺฉ โคตฺโต ปริพฺพาชโก ภควนฺตํ เอตทโวจ “โก นุ โข โภ โคตม เหตุ โก ปจฺจโย, ยานิมานิ อเนกวิหิตานิ ทิฏฺฐิคตานิ โลเก อุปฺปชฺชนฺติ ‘สสฺสโต โลโก’ติ วา, ‘อสสฺสโต โลโก’ติ วา ฯเปฯ ‘เนว โหติ น น โหติ ตถาคโต ปรํ มรณา’ติ วา”ติ.\csromanspacer{} สญฺญาย โข วจฺฉ อญฺญาณา สญฺญาสมุทเย อญฺญาณา สญฺญานิโรเธ อญฺญาณา สญฺญานิโรธคามินิยา ปฏิปทาย อญฺญาณา, เอวมิมานิ อเนกวิหิตานิ ทิฏฺฐิคตานิ โลเก อุปฺปชฺชนฺติ “สสฺสโต โลโก”ติ วา, “อสสฺสโต โลโก”ติ วา ฯเปฯ “เนว โหติ น น โหติ ตถาคโต ปรํ มรณา”ติ วาติ.\csromanspacer{} อยํ โข วจฺฉ เหตุ อยํ ปจฺจโย, ยานิมานิ อเนกวิหิตานิ ทิฏฺฐิคตานิ โลเก อุปฺปชฺชนฺติ “สสฺสโต โลโก”ติ วา, “อสสฺสโต โลโก”ติ วา ฯเปฯ “เนว โหติ น น โหติ ตถาคโต ปรํ มรณา”ติ วาติ.\csromanspacer{}\csromanspacer{}ตติยํ.}

\csromanheader{2}{4. สงฺขาร-อญฺญาณสุตฺต}
\proseitem{610}{สาวตฺถินิทานํ.\csromanspacer{} เอกมนฺตํ นิสินฺโน โข วจฺฉโคตฺโต ปริพฺพาชโก ภควนฺตํ เอตทโวจ “โก นุ โข โภ โคตม เหตุ โก ปจฺจโย, ยานิมานิ อเนกวิหิตานิ ทิฏฺฐิคตานิ โลเก อุปฺปชฺชนฺติ ‘สสฺสโต โลโก’ติ วา, ‘อสสฺสโต โลโก’ติ วา ฯเปฯ “เนว โหติ น น โหติ ตถาคโต ปรํ มรณา’ติ วา”ติ.\csromanspacer{} สงฺขาเรสุ โข วจฺฉ อญฺญาณา สงฺขารสมุทเย อญฺญาณา สงฺขารนิโรเธ อญฺญาณา สงฺขารนิโรธคามินิยา ปฏิปทาย อญฺญาณา, เอวมิมานิ อเนกวิหิตานิ ทิฏฺฐิคตานิ \csromanfolio{220}โลเก อุปฺปชฺชนฺติ “สสฺสโต โลโก”ติ วา, “อสสฺสโต โลโก”ติ วา ฯเปฯ “เนว โหติ น น โหติ ตถาคโต ปรํ มรณา”ติ วาติ.\csromanspacer{} อยํ โข วจฺฉ เหตุ อยํ ปจฺจโย, ยานิมานิ อเนกวิหิตานิ ทิฏฺฐิคตานิ โลเก อุปฺปชฺชนฺติ “สสฺสโต โลโก”ติ วา, “อสสฺสโต โลโก”ติ วา ฯเปฯ “เนว โหติ น น โหติ ตถาคโต ปรํ มรณา”ติ วาติ.\csromanspacer{}\csromanspacer{}จตุตฺถํ.}

\csromanheader{2}{5. วิญฺญาณ-อญฺญาณสุตฺต}
\proseitem{611}{สาวตฺถินิทานํ.\csromanspacer{} เอกมนฺตํ นิสินฺโน โข วจฺฉโคตฺโต ปริพฺพาชโก ภควนฺตํ เอตทโวจ “โก นุ โข โภ โคตม เหตุ โก ปจฺจโย, ยานิมานิ อเนกวิหิตานิ ทิฏฺฐิคตานิ โลเก อุปฺปชฺชนฺติ ‘สสฺสโต โลโก’ติ วา, ‘อสสฺสโต โลโก’ติ วา ฯเปฯ ‘เนว โหติ น น โหติ ตถาคโต ปรํ มรณา’ติ วา”ติ.\csromanspacer{} วิญฺญาเณ โข วจฺฉ อญฺญาณา วิญฺญาณสมุทเย อญฺญาณา วิญฺญาณนิโรเธ อญฺญาณา วิญฺญาณนิโรธคามินิยา ปฏิปทาย อญฺญาณา, เอวมิมานิ อเนกวิหิตานิ ทิฏฺฐิคตานิ โลเก อุปฺปชฺชนฺติ “สสฺสโต โลโก”ติ วา, “อสสฺสโต โลโก”ติ วา ฯเปฯ “เนว โหติ น น โหติ ตถาคโต ปรํ มรณา”ติ วาติ.\csromanspacer{} อยํ โข วจฺฉ เหตุ อยํ ปจฺจโย, ยานิมานิ อเนกวิหิตานิ ทิฏฺฐิคตานิ โลเก อุปฺปชฺชนฺติ “สสฺสโต โลโก”ติ วา, “อสสฺสโต โลโก”ติ วา ฯเปฯ “เนว โหติ น น โหติ ตถาคโต ปรํ มรณา”ติ วาติ.\csromanspacer{}\csromanspacer{}ปญฺจมํ.}

\csromanheader{2}{6-10. รูป-อทสฺสนาทิสุตฺตปญฺจก}
\proseitem{612-616}{สาวตฺถินิทานํ.\csromanspacer{} เอกมนฺตํ นิสินฺโน โข วจฺฉโคตฺโต ปริพฺพาชโก ภควนฺตํ เอตทโวจ “โก นุ โข โภ โคตม เหตุ โก ปจฺจโย, ยานิมานิ อเนกวิหิตานิ ทิฏฺฐิคตานิ โลเก อุปฺปชฺชนฺติ ‘สสฺสโต โลโก’ติ วา, ‘อสสฺสโต โลโก’ติ วา ฯเปฯ ‘เนว โหติ น น โหติ ตถาคโต ปรํ มรณา’ติ วา”ติ.\csromanspacer{} รูเป โข วจฺฉ อทสฺสนา ฯเปฯ รูปนิโรธคามินิยา ปฏิปทาย อทสฺสนา ฯเปฯ.\csromanspacer{} เวทนาย.\csromanspacer{} สญฺญาย.\csromanspacer{} สงฺขาเรสุ โข วจฺฉ อทสฺสนา ฯเปฯ.\csromanspacer{} วิญฺญาเณ โข วจฺฉ อทสฺสนา ฯเปฯ วิญฺญาณนิโรธคามินิยา ปฏิปทาย อทสฺสนา ฯเปฯ.\csromanspacer{}\csromanspacer{}ทสมํ.}

\csromanheader{2}{12. วจฺฉโคตฺตสํยุตฺต}
\csromanheader{2}{11-15. รูป-อนภิสมยาทิสุตฺตปญฺจก}
\proseitem{617-621}{\csromanfolio{221}สาวตฺถินิทานํ.\csromanspacer{} รูเป โข วจฺฉ อนภิสมยา ฯเปฯ.\csromanspacer{} รูปนิโรธคามินิยา ปฏิปทาย อนภิสมยา ฯเปฯ.}

\prose{สาวตฺถินิทานํ.\csromanspacer{} เวทนาย โข วจฺฉ อนภิสมยา ฯเปฯ.}

\prose{สาวตฺถินิทานํ.\csromanspacer{} สญฺญาย โข วจฺฉ อนภิสมยา ฯเปฯ.}

\prose{สาวตฺถินิทานํ.\csromanspacer{} สงฺขาเรสุ โข วจฺฉ อนภิสมยา ฯเปฯ.}

\prose{สาวตฺถินิทานํ.\csromanspacer{} วิญฺญาเณ โข วจฺฉ อนภิสมยา ฯเปฯ.\csromanspacer{}\csromanspacer{}ปนฺนรสมํ.}

\csromanheader{2}{16-20. รูป-อนนุโพธาทิสุตฺตปญฺจก}
\proseitem{622}{626.\csromanspacer{} สาวตฺถินิทานํ.\csromanspacer{} เอกมนฺตํ นิสินฺโน โข วจฺฉโคตฺโต ปริพฺพาชโก ภควนฺตํ เอตทโวจ\textbf{\csromandash{}}โก นุ โข โภ โคตม เหตุ โก ปจฺจโย ฯเปฯ.\csromanspacer{} รูเป โข วจฺฉ อนนุโพธา ฯเปฯ รูปนิโรธคามินิยา ปฏิปทาย อนนุโพธา ฯเปฯ.}

\prose{สาวตฺถินิทานํ.\csromanspacer{} เวทนาย โข วจฺฉ ฯเปฯ.}

\prose{สาวตฺถินิทานํ.\csromanspacer{} สญฺญาย โข วจฺฉ ฯเปฯ.}

\prose{สาวตฺถินิทานํ.\csromanspacer{} สงฺขาเรสุ โข วจฺฉ ฯเปฯ.}

\prose{สาวตฺถินิทานํ.\csromanspacer{} วิญฺญาเณ โข วจฺฉ อนนุโพธา ฯเปฯ วิญฺญาณนิโรธคามินิยา ปฏิปทาย อนนุโพธา ฯเปฯ.\csromanspacer{}\csromanspacer{}วีสติมํ.}

\csromanheader{2}{21-25. รูป-อปฺปฏิเวธาทิสุตฺตปญฺจก}
\proseitem{627-631}{สาวตฺถินิทานํ.\csromanspacer{} โก นุ โข โภ โคตม เหตุ โก ปจฺจโย ฯเปฯ.\csromanspacer{} รูเป โข วจฺฉ อปฺปฏิเวธา ฯเปฯ.\csromanspacer{} วิญฺญาเณ โข วจฺฉ อปฺปฏิเวธา ฯเปฯ.\csromanspacer{}\csromanspacer{}ปญฺจวีสติมํ.}

\csromanheader{2}{26-30. รูป-อสลฺลกฺขณาทิสุตฺตปญฺจก}
\proseitem{632-636}{สาวตฺถินิทานํ.\csromanspacer{} รูเป โข วจฺฉ อสลฺลกฺขณา ฯเปฯ.\csromanspacer{} วิญฺญาเณ โข วจฺฉ อสลฺลกฺขณา ฯเปฯ.\csromanspacer{}\csromanspacer{}ติํสติมํ.}

\csromanheader{2}{31-35. รูป-อนุปลกฺขณาทิสุตฺตปญฺจก}
\proseitem{637-641}{\csromanfolio{222}สาวตฺถินิทานํ.\csromanspacer{} รูเป โข วจฺฉ อนุปลกฺขณา ฯเปฯ.\csromanspacer{} วิญฺญาเณ โข วจฺฉ อนุปลกฺขณา ฯเปฯ.\csromanspacer{}\csromanspacer{}ปญฺจติํสติมํ.}

\csromanheader{2}{36-40. รูป-อปฺปจฺจุปลกฺขณาทิสุตฺตปญฺจก}
\proseitem{642-646}{สาวตฺถินิทานํ.\csromanspacer{} รูเป โข วจฺฉ อปฺปจฺจุปลกฺขณา ฯเปฯ.\csromanspacer{} วิญฺญาเณ โข วจฺฉ อปฺปจฺจุปลกฺขณา ฯเปฯ.\csromanspacer{}\csromanspacer{}จตฺตาลีสมํ.}

\csromanheader{2}{41-45. รูป-อสมเปกฺขณาทิสุตฺตปญฺจก}
\proseitem{647-651}{สาวตฺถินิทานํ.\csromanspacer{} รูเป โข วจฺฉ อสมเปกฺขณา ฯเปฯ.\csromanspacer{} วิญฺญาเณ โข วจฺฉ อสมเปกฺขณา ฯเปฯ.\csromanspacer{}\csromanspacer{}ปญฺจจตฺตาลีสมํ.}

\csromanheader{2}{46-50. รูป-อปฺปจฺจุเปกฺขณาทิสุตฺตปญฺจก}
\proseitem{652-656}{สาวตฺถินิทานํ.\csromanspacer{} รูเป โข วจฺฉ อปฺปจฺจุเปกฺขณา ฯเปฯ.\csromanspacer{} วิญฺญาเณ โข วจฺฉ อปฺปจฺจุเปกฺขณา ฯเปฯ.\csromanspacer{} ปญฺญาสมํ.}

\csromanheader{2}{51-54. รูป-อปฺปจฺจกฺขกมฺมาทิสุตฺตจตุกฺก}
\prose{657-66 สาวตฺถินิทานํ.\csromanspacer{} อถ โข วจฺฉโคตฺโต ปริพฺพาชโก เยน ภควา เตนุปสงฺกมิ, อุปสงฺกมิตฺวา ภควตา สทฺธิํ สมฺโมทิ, สมฺโมทนียํ กถํ สารณียํ วีติสาเรตฺวา เอกมนฺตํ นิสีทิ, เอกมนฺตํ นิสินฺโน โข วจฺฉโคตฺโต ปริพฺพาชโก ภควนฺตํ เอตทโวจ “โก นุ โข โภ โคตม เหตุ โก ปจฺจโย, ยานิมานิ อเนกวิหิตานิ ทิฏฺฐิคตานิ โลเก อุปฺปชฺชนฺติ ‘สสฺสโต โลโก’ติ วา ฯเปฯ “เนว โหติ น น โหติ ตถาคโต ปรํ มรณา’ติ วา”ติ.\csromanspacer{} รูเป โข วจฺฉ อปฺปจฺจกฺขกมฺมา รูปสมุทเย อปฺปจฺจกฺขกมฺมา รูปนิโรเธ อปฺปจฺจกฺขกมฺมา รูปนิโรธคามินิยา ปฏิปทาย อปฺปจฺจกฺขกมฺมา ฯเปฯ.}

\prose{สาวตฺถินิทานํ.\csromanspacer{} เวทนาย โข วจฺฉ อปฺปจฺจกฺขกมฺมา ฯเปฯ เวทนานิโรธคามินิยา ปฏิปทาย อปฺปจฺจกฺขกมฺมา ฯเปฯ.}

\prose{สาวตฺถินิทานํ.\csromanspacer{} สญฺญาย โข วจฺฉ อปฺปจฺจกฺขกมฺมา ฯเปฯ สญฺญานิโรธคามินิยา ปฏิปทาย อปฺปจฺจกฺขกมฺมา ฯเปฯ.}

\csromanheader{2}{12. วจฺฉโคตฺตสํยุตฺต}
\prose{\csromanfolio{223}สาวตฺถินิทานํ.\csromanspacer{} สงฺขาเรสุ โข วจฺฉ อปฺปจฺจกฺขกมฺมา ฯเปฯ.\csromanspacer{} สงฺขารนิโรธคามินิยา ปฏิปทาย อปฺปจฺจกฺขกมฺมา ฯเปฯ.\csromanspacer{}\csromanspacer{}จตุปญฺญาสมํ.}

\csromanheader{2}{55. วิญฺญาณ-อปฺปจฺจกฺขกมฺมสุตฺต}
\proseitem{661}{สาวตฺถินิทานํ.\csromanspacer{} วิญฺญาเณ โข วจฺฉ อปฺปจฺจกฺขกมฺมา วิญฺญาณสมุทเย อปฺปจฺจกฺขกมฺมา วิญฺญาณนิโรเธ อปฺปจฺจกฺขกมฺมา วิญฺญาณนิโรธคามินิยา ปฏิปทาย อปฺปจฺจกฺขกมฺมา, เอวมิมานิ อเนกวิหิตานิ ทิฏฺฐิคตานิ โลเก อุปฺปชฺชนฺติ “สสฺสโต โลโก”ติ วา, “อสสฺสโต โลโก”ติ วา ฯเปฯ “เนว โหติ น น โหติ ตถาคโต ปรํ มรณา”ติ วาติ.\csromanspacer{} อยํ โข วจฺฉ เหตุ อยํ ปจฺจโย, ยานิมานิ อเนกวิหิตานิ ทิฏฺฐิคตานิ โลเก อุปฺปชฺชนฺติ “สสฺสโต โลโก”ติ วา, “อสสฺสโต โลโก”ติ วา, “อนฺตวา โลโก”ติ วา, “อนนฺตวา โลโก”ติ วา, “ตํ ชีวํ ตํ สรีรนฺ”ติ วา, “อญฺญํ ชีวํ อญฺญํ สรีรนฺ”ติ วา, “โหติ ตถาคโต ปรํ มรณา”ติ วา, “น โหติ ตถาคโต ปรํ มรณา”ติ วา, “โหติ จ น จ โหติ ตถาคโต ปรํ มรณา”ติ วา, “เนว โหติ น น โหติ ตถาคโต ปรํ มรณา”ติ วาติ.\csromanspacer{}\csromanspacer{}ปญฺจปญฺญาสมํ.}

\vspace{\tipitakaparskip}
\csromancenter{\textbf{วจฺฉโคตฺตสํยุตฺตํ สมตฺตํ.}}
\titlehead{\textbf{ตสฺสุทฺทานํ}}
\csromangathameasure{อญฺญาณา อทสฺสนา เจว, อนภิสมยา อนนุโพธา. \\ อปฺปฏิเวธา อสลฺลกฺขณา, อนุปลกฺขเณน อปฺปจฺจุปลกฺขณา. อสมเปกฺขณา อปฺปจฺจุเปกฺขณา อปฺปจฺจกฺขกมฺมนฺติ.}
\csromangathagroup{\csromangathastanza{อญฺญาณา อทสฺสนา เจว, อนภิสมยา อนนุโพธา. \\ อปฺปฏิเวธา อสลฺลกฺขณา, อนุปลกฺขเณน อปฺปจฺจุปลกฺขณา.\csromanspacer{} อสมเปกฺขณา อปฺปจฺจุเปกฺขณา อปฺปจฺจกฺขกมฺมนฺติ.}}

\csromanmatikaanchor{mtk.77}
\csromantocmarknopage{chapter}{13. ฌานสํยุตฺต}{mtk.77}
\bookmark[dest={mtk.77},level=0]{13. ฌานสํยุตฺต}
\chapterheadpageread{13. ฌานสํยุตฺต}
\csromanheader{2}{1. สมาธิมูลกสมาปตฺติสุตฺต}
\csromanmatikaanchor{mtk.78}
\csromantocmarkref{section}{1-55. สมาธิมูลกสมาปตฺติสุตฺตาทิ}{mtk.78}
\bookmark[dest={mtk.78},level=1]{1-55. สมาธิมูลกสมาปตฺติสุตฺตาทิ}
\proseitem{662}{\csromanfolio{224}สาวตฺถินิทานํ.\csromanspacer{} จตฺตาโรเม ภิกฺขเว ฌายี.\csromanspacer{} กตเม จตฺตาโร, อิธ ภิกฺขเว เอกจฺโจ ฌายี สมาธิสฺมิํ สมาธิกุสโล โหติ, น สมาธิสฺมิํ สมาปตฺติกุสโล.\csromanspacer{} อิธ ปน ภิกฺขเว เอกจฺโจ ฌายี สมาธิสฺมิํ สมาปตฺติกุสโล โหติ, น สมาธิสฺมิํ สมาธิกุสโล.\csromanspacer{} อิธ ปน ภิกฺขเว เอกจฺโจ ฌายี เนว สมาธิสฺมิํ สมาธิกุสโล โหติ, น จ สมาธิสฺมิํ สมาปตฺติกุสโล.\csromanspacer{} อิธ ปน ภิกฺขเว เอกจฺโจ ฌายี สมาธิสฺมิํ สมาธิกุสโล จ โหติ, สมาธิสฺมิํ สมาปตฺติกุสโล จ.\csromanspacer{} ตตฺร ภิกฺขเว ยฺวายํ ฌายี สมาธิสฺมิํ สมาธิกุสโล จ โหติ, สมาธิสฺมิํ สมาปตฺติกุสโล จ, อยํ อิเมสํ จตุนฺนํ ฌายีนํ อคฺโค จ เสฏฺโฐ จ โมกฺโข\footnote{ปาโมกฺโข (สฺยา, กํ) เอวมุปริปิ.} จ อุตฺตโม จ ปวโร จ.\csromanspacer{} เสยฺยถาปิ ภิกฺขเว ควา ขีรํ ขีรมฺหา ทธิ ทธิมฺหา นวนีตํ นวนีตมฺหา สปฺปิ สปฺปิมฺหา สปฺปิมณฺโฑ ตตฺร อคฺคมกฺขายติ.\csromanspacer{} เอวเมว โข ภิกฺขเว ยฺวายํ ฌายี สมาธิสฺมิํ สมาธิกุสโล จ โหติ, สมาธิสฺมิํ สมาปตฺติกุสโล จ.\csromanspacer{} อยํ อิเมสํ จตุนฺนํ ฌายีนํ อคฺโค จ เสฏฺโฐ จ โมกฺโข1 จ อุตฺตโม จ ปวโร จาติ.\csromanspacer{}\csromanspacer{}ปฐมํ.}

\csromanheader{2}{2. สมาธิมูลกฐิติสุตฺต}
\proseitem{663}{สาวตฺถินิทานํ.\csromanspacer{} จตฺตาโรเม ภิกฺขเว ฌายี.\csromanspacer{} กตเม จตฺตาโร, อิธ ภิกฺขเว เอกจฺโจ ฌายี สมาธิสฺมิํ สมาธิกุสโล โหติ, น สมาธิสฺมิํ ฐิติกุสโล.\csromanspacer{} อิธ ปน ภิกฺขเว เอกจฺโจ ฌายี สมาธิสฺมิํ ฐิติกุสโล โหติ, น สมาธิสฺมิํ สมาธิกุสโล.\csromanspacer{} อิธ ปน ภิกฺขเว เอกจฺโจ ฌายี เนว สมาธิสฺมิํ สมาธิกุสโล โหติ, น จ สมาธิสฺมิํ ฐิติกุสโล.\csromanspacer{} อิธ ปน ภิกฺขเว เอกจฺโจ ฌายี สมาธิสฺมิํ สมาธิกุสโล จ โหติ, สมาธิสฺมิํ ฐิติกุสโล จ.\csromanspacer{} ตตฺร ภิกฺขเว ยฺวายํ ฌายี สมาธิสฺมิํ สมาธิกุสโล จ โหติ, สมาธิสฺมิํ ฐิติกุสโล จ, อยํ อิเมสํ จตุนฺนํ ฌายีนํ อคฺโค จ เสฏฺโฐ จ โมกฺโข จ อุตฺตโม จ ปวโร จ.\csromanspacer{} เสยฺยถาปิ ภิกฺขเว ควา ขีรํ ขีรมฺหา ทธิ ทธิมฺหา นวนีตํ}

\csromanheader{2}{13. ฌานสํยุตฺต}
\prose{\csromanfolio{225}นวนีตมฺหา สปฺปิ สปฺปิมฺหา สปฺปิมณฺโฑ ตตฺร อคฺคมกฺขายติ.\csromanspacer{} เอวเมว โข ภิกฺขเว ยฺวายํ ฌายี สมาธิสฺมิํ สมาธิกุสโล จ โหติ, สมาธิสฺมิํ ฐิติกุสโล จ.\csromanspacer{} อยํ อิเมสํ จตุนฺนํ ฌายีนํ อคฺโค จ เสฏฺโฐ จ โมกฺโข จ อุตฺตโม จ ปวโร จาติ.\csromanspacer{}\csromanspacer{}ทุติยํ.}

\csromanheader{2}{3. สมาธิมูลกวุฏฺฐานสุตฺต}
\proseitem{664}{สาวตฺถินิทานํ.\csromanspacer{} จตฺตาโรเม ภิกฺขเว ฌายี.\csromanspacer{} กตเม จตฺตาโร, อิธ ภิกฺขเว เอกจฺโจ ฌายี สมาธิสฺมิํ สมาธิกุสโล โหติ, น สมาธิสฺมิํ วุฏฺฐานกุสโล.\csromanspacer{} อิธ ปน ภิกฺขเว เอกจฺโจ ฌายี สมาธิสฺมิํ วุฏฺฐานกุสโล โหติ, น สมาธิสฺมิํ สมาธิกุสโล.\csromanspacer{} อิธ ปน ภิกฺขเว เอกจฺโจ ฌายี เนว สมาธิสฺมิํ สมาธิกุสโล โหติ, น จ สมาธิสฺมิํ วุฏฺฐานกุสโล.\csromanspacer{} อิธ ปน ภิกฺขเว เอกจฺโจ ฌายี สมาธิสฺมิํ สมาธิกุสโล จ โหติ, สมาธิสฺมิํ วุฏฺฐานกุสโล จ.\csromanspacer{} ตตฺร ภิกฺขเว ยฺวายํ ฌายี สมาธิสฺมิํ สมาธิกุสโล จ โหติ, สมาธิสิํ วุฏฺฐานกุสโล จ, อยํ อิเมสํ จตุนฺนํ ฌายีนํ อคฺโค จ เสฏฺโฐ จ โมกฺโข จ อุตฺตโม จ ปวโร จ.\csromanspacer{} เสยฺยถาปิ ภิกฺขเว ควา ขีรํ ฯเปฯ ปวโร จาติ.\csromanspacer{}\csromanspacer{}ตติยํ.}

\csromanheader{2}{4. สมาธิมูลกกลฺลิตสุตฺต}
\proseitem{665}{สาวตฺถินิทานํ.\csromanspacer{} จตฺตาโรเม ภิกฺขเว ฌายี.\csromanspacer{} กตเม จตฺตาโร, อิธ ภิกฺขเว เอกจฺโจ ฌายี สมาธิสฺมิํ สมาธิกุสโล โหติ, น สมาธิสฺมิํ กลฺลิตกุสโล.\csromanspacer{} อิธ ปน ภิกฺขเว เอกจฺโจ ฌายี สมาธิสฺมิํ กลฺลิตกุสโล โหติ, น สมาธิสฺมิํ สมาธิกุสโล.\csromanspacer{} อิธ ปน ภิกฺขเว เอกจฺโจ ฌายี เนว สมาธิสฺมิํ สมาธิกุสโล โหติ, น จ สมาธิสฺมิํ กลฺลิตกุสโล.\csromanspacer{} อิธ ปน ภิกฺขเว เอกจฺโจ ฌายี สมาธิสฺมิํ สมาธิกุสโล จ โหติ, สมาธิสฺมิํ กลฺลิตกุสโล จ.\csromanspacer{} ตตฺร ภิกฺขเว ยฺวายํ ฌายี สมาธิสฺมิํ สมาธิกุสโล จ โหติ, สมาธิสฺมิํ กลฺลิตกุสโล จ.\csromanspacer{} อยํ อิเมสํ จตุนฺนํ ฌายีนํ อคฺโค จ เสฏฺโฐ จ โมกฺโข จ อุตฺตโม จ ปวโร จ.\csromanspacer{} เสยฺยถาปิ ภิกฺขเว ควา ขีรํ ฯเปฯ ปวโร จาติ.\csromanspacer{}\csromanspacer{}จตุตฺถํ.}

\csromanheader{2}{5. สมาธิมูลก-อารมฺมณสุตฺต}
\proseitem{666}{\csromanfolio{226}สาวตฺถินิทานํ.\csromanspacer{} จตฺตาโรเม ภิกฺขเว ฌายี.\csromanspacer{} กตเม จตฺตาโร, อิธ ภิกฺขเว เอกจฺโจ ฌายี สมาธิสฺมิํ สมาธิกุสโล โหติ, น สมาธิสฺมิํ อารมฺมณกุสโล.\csromanspacer{} อิธ ปน ภิกฺขเว เอกจฺโจ ฌายี สมาธิสฺมิํ อารมฺมณกุสโล โหติ, น สมาธิสฺมิํ สมาธิกุสโล.\csromanspacer{} อิธ ปน ภิกฺขเว เอกจฺโจ ฌายี เนว สมาธิสฺมิํ สมาธิกุสโล โหติ, น จ สมาธิสฺมิํ อารมฺมณกุสโล.\csromanspacer{} อิธ ปน ภิกฺขเว เอกจฺโจ ฌายี สมาธิสฺมิํ สมาธิกุสโล จ โหติ, สมาธิสฺมิํ อารมฺมณกุสโล จ.\csromanspacer{} ตตฺร ภิกฺขเว ยฺวายํ ฌายี สมาธิสฺมิํ สมาธิกุสโล จ โหติ, สมาธิสฺมิํ อารมฺมณกุสโล จ.\csromanspacer{} อยํ อิเมสํ จตุนฺนํ ฌายีนํ อคฺโค จ เสฏฺโฐ จ โมกฺโข จ อุตฺตโม จ ปวโร จ.\csromanspacer{} เสยฺยถาปิ ภิกฺขเว ควา ขีรํ ฯเปฯ ปวโร จาติ.\csromanspacer{}\csromanspacer{}ปญฺจมํ.}

\csromanheader{2}{6. สมาธิมูลกโคจรสุตฺต}
\proseitem{667}{สาวตฺถินิทานํ.\csromanspacer{} จตฺตาโรเม ภิกฺขเว ฌายี.\csromanspacer{} กตเม จตฺตาโร, อิธ ภิกฺขเว เอกจฺโจ ฌายี สมาธิสฺมิํ สมาธิกุสโล โหติ, น สมาธิสฺมิํ โคจรกุสโล.\csromanspacer{} อิธ ปน ภิกฺขเว เอกจฺโจ ฌายี สมาธิสฺมิํ โคจรกุสโล โหติ, น สมาธิสฺมิํ สมาธิกุสโล.\csromanspacer{} อิธ ปน ภิกฺขเว เอกจฺโจ ฌายี เนว สมาธิสฺมิํ สมาธิกุสโล โหติ, น จ สมาธิสฺมิํ โคจรกุสโล.\csromanspacer{} อิธ ปน ภิกฺขเว เอกจฺโจ ฌายี สมาธิสฺมิํ สมาธิกุสโล จ โหติ, สมาธิสฺมิํ โคจรกุสโล จ.\csromanspacer{} ตตฺร ภิกฺขเว ยฺวายํ ฌายี สมาธิสฺมิํ สมาธิกุสโล จ โหติ, สมาธิสฺมิํ โคจรกุสโล จ.\csromanspacer{} อยํ อิเมสํ จตุนฺนํ ฌายีนํ อคฺโค จ เสฏฺโฐ จ โมกฺโข จ อุตฺตโม จ ปวโร จ.\csromanspacer{} เสยฺยถาปิ ภิกฺขเว ควา ขีรํ ฯเปฯ ปวโร จาติ.\csromanspacer{}\csromanspacer{}ฉฏฺฐํ.}

\csromanheader{2}{7. สมาธิมูลก-อภินีหารสุตฺต}
\proseitem{668}{สาวตฺถินิทานํ.\csromanspacer{} จตฺตาโรเม ภิกฺขเว ฌายี.\csromanspacer{} กตเม จตฺตาโร, อิธ ภิกฺขเว เอกจฺโจ ฌายี สมาธิสฺมิํ สมาธิกุสโล โหติ, น สมาธิสฺมิํ อภินีหารกุสโล.\csromanspacer{} อิธ ปน ภิกฺขเว เอกจฺโจ ฌายี สมาธิสฺมิํ อภินีหารกุสโล โหติ, น สมาธิสฺมิํ สมาธิกุสโล.}

\vspace{\tipitakaparskip}
\csromancenter{ฌานสํยุตฺต อิธ ปน ภิกฺขเว เอกจฺโจ ฌายี เนว สมาธิสฺมิํ สมาธิกุสโล โหติ, น จ สมาธิสฺมิํ อภินีหารกุสโล.\csromanspacer{} อิธ ปน ภิกฺขเว เอกจฺโจ ฌายี สมาธิสฺมิํ สมาธิกุสโล จ โหติ, สมาธิสฺมิํ อภินีหารกุสโล จ.\csromanspacer{} ตตฺร ภิกฺขเว ยฺวายํ ฌายี สมาธิสฺมิํ สมาธิกุสโล จ โหติ, สมาธิสฺมิํ อภินีหารกุสโล จ.\csromanspacer{} อยํ อิเมสํ จตุนฺนํ ฌายีนํ อคฺโค จ เสฏฺโฐ จ โมกฺโข จ อุตฺตโม จ ปวโร จ.\csromanspacer{} เสยฺยถาปิ ภิกฺขเว ควา ขีรํ ฯเปฯ ปวโร จาติ.\csromanspacer{}\csromanspacer{}สตฺตมํ.}
\csromanheader{2}{8. สมาธิมูลกสกฺกจฺจการีสุตฺต}
\proseitem{669}{\csromanfolio{227}สาวตฺถินิทานํ.\csromanspacer{} จตฺตาโรเม ภิกฺขเว ฌายี.\csromanspacer{} กตเม จตฺตาโร, อิธ ภิกฺขเว เอกจฺโจ ฌายี สมาธิสฺมิํ สมาธิกุสโล โหติ, น สมาธิสฺมิํ สกฺกจฺจการี.\csromanspacer{} อิธ ปน ภิกฺขเว เอกจฺโจ ฌายี สมาธิสฺมิํ สกฺกจฺจการี โหติ, น สมาธิสฺมิํ สมาธิกุสโล.\csromanspacer{} อิธ ปน ภิกฺขเว เอกจฺโจ ฌายี เนว สมาธิสฺมิํ สมาธิกุสโล โหติ, น จ สมาธิสฺมิํ สกฺกจฺจการี.\csromanspacer{} อิธ ปน ภิกฺขเว เอกจฺโจ ฌายี สมาธิสฺมิํ สมาธิกุสโล จ โหติ, สมาธิสฺมิํ สกฺกจฺจการี จ.\csromanspacer{} ตตฺร ภิกฺขเว ยฺวายํ ฌายี สมาธิสฺมิํ สมาธิกุสโล จ โหติ, สมาธิสฺมิํ สกฺกจฺจการี จ.\csromanspacer{} อยํ อิเมสํ จตุนฺนํ ฌายีนํ อคฺโค จ เสฏฺโฐ จ โมกฺโข จ อุตฺตโม จ ปวโร จ.\csromanspacer{} เสยฺยถาปิ ภิกฺขเว ควา ขีรํ ฯเปฯ ปวโร จาติ.\csromanspacer{}\csromanspacer{}อฏฺฐมํ.}

\csromanheader{2}{9. สมาธิมูลกสาตจฺจการีสุตฺต}
\proseitem{670}{สาวตฺถินิทานํ.\csromanspacer{} จตฺตาโรเม ภิกฺขเว ฌายี.\csromanspacer{} กตเม จตฺตาโร, อิธ ภิกฺขเว เอกจฺโจ ฌายี สมาธิสฺมิํ สมาธิกุสโล โหติ, น สมาธิสฺมิํ สาตจฺจการี.\csromanspacer{} อิธ ปน ภิกฺขเว เอกจฺโจ ฌายี สมาธิสฺมิํ สาตจฺจการี โหติ, น สมาธิสฺมิํ สมาธิกุสโล.\csromanspacer{} อิธ ปน ภิกฺขเว เอกจฺโจ ฌายี เนว สมาธิสฺมิํ สมาธิกุสโล โหติ, น จ สมาธิสฺมิํ สาตจฺจการี.\csromanspacer{} อิธ ปน ภิกฺขเว เอกจฺโจ ฌายี สมาธิสฺมิํ สมาธิกุสโล จ โหติ, สมาธิสฺมิํ สาตจฺจการี จ.\csromanspacer{} ตตฺร ภิกฺขเว ยฺวายํ ฌายี สมาธิสฺมิํ สมาธิกุสโล จ โหติ, สมาธิสฺมิํ สาตจฺจการี จ.\csromanspacer{} อยํ อิเมสํ จตุนฺนํ ฌายีนํ อคฺโค จ เสฏฺโฐ จ โมกฺโข จ อุตฺตโม จ ปวโร จ.\csromanspacer{} เสยฺยถาปิ ภิกฺขเว ควา ขีรํ ฯเปฯ ปวโร จาติ.\csromanspacer{}\csromanspacer{}นวมํ.}

\csromanheader{2}{10. สมาธิมูลกสปฺปายการีสุตฺต}
\proseitem{671}{\csromanfolio{228}สาวตฺถินิทานํ.\csromanspacer{} จตฺตาโรเม ภิกฺขเว ฌายี.\csromanspacer{} กตเม จตฺตาโร, อิธ ภิกฺขเว เอกจฺโจ ฌายี สมาธิสฺมิํ สมาธิกุสโล โหติ, น สมาธิสฺมิํ สปฺปายการี.\csromanspacer{} อิธ ปน ภิกฺขเว เอกจฺโจ ฌายี สมาธิสฺมิํ สปฺปายการี โหติ, น สมาธิสฺมิํ สมาธิกุสโล.\csromanspacer{} อิธ ปน ภิกฺขเว เอกจฺโจ ฌายี เนว สมาธิสฺมิํ สมาธิกุสโล โหติ, น จ สมาธิสฺมิํ สปฺปายการี.\csromanspacer{} อิธ ปน ภิกฺขเว เอกจฺโจ ฌายี สมาธิสฺมิํ สมาธิกุสโล จ โหติ, สมาธิสฺมิํ สปฺปายการี จ.\csromanspacer{} ตตฺร ภิกฺขเว ยฺวายํ ฌายี สมาธิสฺมิํ สมาธิกุสโล จ โหติ, สมาธิสฺมิํ สปฺปายการี จ.\csromanspacer{} อยํ อิเมสํ จตุนฺนํ ฌายีนํ อคฺโค จ เสฏฺโฐ จ โมกฺโข จ อุตฺตโม จ ปวโร จ.\csromanspacer{} เสยฺยถาปิ ภิกฺขเว ควา ขีรํ ฯเปฯ ปวโร จาติ.\csromanspacer{}\csromanspacer{}ทสมํ.\csromanspacer{}\csromanspacer{}(สมาธิมูลกํ.)}

\csromanheader{2}{11. สมาปตฺติมูลกฐิติสุตฺต}
\proseitem{672}{สาวตฺถินิทานํ.\csromanspacer{} จตฺตาโรเม ภิกฺขเว ฌายี.\csromanspacer{} กตเม จตฺตาโร, อิธ ภิกฺขเว เอกจฺโจ ฌายี สมาธิสฺมิํ สมาปตฺติกุสโล โหติ, น สมาธิสฺมิํ ฐิติกุสโล.\csromanspacer{} อิธ ปน ภิกฺขเว เอกจฺโจ ฌายี สมาธิสฺมิํ ฐิติกุสโล โหติ, น สมาธิสฺมิํ สมาปตฺติกุสโล.\csromanspacer{} อิธ ปน ภิกฺขเว เอกจฺโจ ฌายี เนว สมาธิสฺมิํ สมาปตฺติกุสโล โหติ, น จ สมาธิสฺมิํ ฐิติกุสโล.\csromanspacer{} อิธ ปน ภิกฺขเว เอกจฺโจ ฌายี สมาธิสฺมิํ สมาปตฺติกุสโล จ โหติ, สมาธิสฺมิํ ฐิติกุสโล จ.\csromanspacer{} ตตฺร ภิกฺขเว ยฺวายํ ฌายี สมาธิสฺมิํ สมาปตฺติกุสโล จ โหติ, สมาธิสฺมิํ ฐิติกุสโล จ.\csromanspacer{} อยํ อิเมสํ จตุนฺนํ ฌายีนํ อคฺโค จ เสฏฺโฐ จ โมกฺโข จ อุตฺตโม จ ปวโร จ.\csromanspacer{} เสยฺยถาปิ ภิกฺขเว ควา ขีรํ ฯเปฯ ปวโร จาติ.\csromanspacer{}\csromanspacer{}เอกาทสมํ.}

\csromanheader{2}{12. สมาปตฺติมูลกวุฐานสุตฺต}
\proseitem{673}{สาวตฺถินิทานํ.\csromanspacer{} จตฺตาโรเม ภิกฺขเว ฌายี.\csromanspacer{} กตเม จตฺตาโร, อิธ ภิกฺขเว เอกจฺโจ ฌายี สมาธิสฺมิํ สมาปตฺติกุสโล โหติ, น สมาธิสฺมิํ วุฏฺฐานกุสโล.\csromanspacer{} อิธ ปน ภิกฺขเว เอกจฺโจ ฌายี สมาธิสฺมิํ วุฏฺฐานกุสโล โหติ, น สมาธิสฺมิํ สมาปตฺติกุสโล.\csromanspacer{} อิธ ปน ภิกฺขเว เอกจฺโจ ฌายี เนว สมาธิสฺมิํ สมาปตฺติกุสโล โหติ, น}

\vspace{\tipitakaparskip}
\csromancenter{ฌานสํยุตฺต จ สมาธิสฺมิํ วุฏฺฐานกุสโล.\csromanspacer{} อิธ ปน ภิกฺขเว เอกจฺโจ ฌายี สมาธิสฺมิํ สมาปตฺติกุสโล จ โหติ, สมาธิสฺมิํ วุฏฺฐานกุสโล จ.\csromanspacer{} ตตฺร ภิกฺขเว ยฺวายํ ฌายี ฯเปฯ ปวโร จาติ.\csromanspacer{}\csromanspacer{}ทฺวาทสมํ.}
\csromanheader{2}{13. สมาปตฺติมูลกกลฺลิตสุตฺต}
\proseitem{674}{\csromanfolio{229}สาวตฺถินิทานํ.\csromanspacer{} จตฺตาโรเม ภิกฺขเว ฌายี.\csromanspacer{} กตเม จตฺตาโร, อิธ ภิกฺขเว เอกจฺโจ ฌายี สมาธิสฺมิํ สมาปตฺติกุสโล โหติ, น สมาธิสฺมิํ กลฺลิตกุสโล.\csromanspacer{} อิธ ปน ภิกฺขเว เอกจฺโจ ฌายี สมาธิสฺมิํ กลฺลิตกุสโล โหติ, น สมาธิสฺมิํ สมาปตฺติกุสโล.\csromanspacer{} อิธ ปน ภิกฺขเว เอกจฺโจ ฌายี เนว สมาธิสฺมิํ สมาปตฺติกุสโล โหติ, น จ สมาธิสฺมิํ กลฺลิตกุสโล.\csromanspacer{} อิธ ปน ภิกฺขเว เอกจฺโจ ฌายี สมาธิสฺมิํ สมาปตฺติกุสโล จ โหติ, สมาธิสฺมิํ กลฺลิตกุสโล จ.\csromanspacer{} ตตฺร ฯเปฯ ปวโร จาติ.\csromanspacer{}\csromanspacer{}เตรสมํ.}

\csromanheader{2}{14. สมาปตฺติมูลก-อารมฺมณสุตฺต}
\proseitem{675}{สาวตฺถินิทานํ.\csromanspacer{} จตฺตาโรเม ภิกฺขเว ฌายี.\csromanspacer{} กตเม จตฺตาโร, อิธ ภิกฺขเว เอกจฺโจ ฌายี สมาธิสฺมิํ สมาปตฺติกุสโล โหติ, น สมาธิสฺมิํ อารมฺมณกุสโล.\csromanspacer{} อิธ ปน ภิกฺขเว เอกจฺโจ ฌายี สมาธิสฺมิํ อารมฺมณกุสโล โหติ, น สมาธิสฺมิํ สมาปตฺติกุสโล.\csromanspacer{} อิธ ปน ภิกฺขเว เอกจฺโจ ฌายี เนว สมาธิสฺมิํ สมาปตฺติกุสโล โหติ, น จ สมาธิสฺมิํ อารมฺมณกุสโล.\csromanspacer{} อิธ ปน ภิกฺขเว เอกจฺโจ ฌายี สมาธิสฺมิํ สมาปตฺติกุสโล จ โหติ, สมาธิสฺมิํ อารมฺมณกุสโล จ.\csromanspacer{} ตตฺร ฯเปฯ ปวโร จาติ.\csromanspacer{}\csromanspacer{}จุทฺทสมํ.}

\csromanheader{2}{15. สมาปตฺติมูลกโคจรสุตฺต}
\proseitem{676}{สาวตฺถินิทานํ.\csromanspacer{} จตฺตาโรเม ภิกฺขเว ฌายี.\csromanspacer{} กตเม จตฺตาโร, อิธ ภิกฺขเว เอกจฺโจ ฌายี สมาธิสฺมิํ สมาปตฺติกุสโล โหติ, น สมาธิสฺมิํ โคจรกุสโล.\csromanspacer{} อิธ ปน ภิกฺขเว เอกจฺโจ ฌายี สมาธิสฺมิํ โคจรกุสโล โหติ, น สมาธิสฺมิํ สมาปตฺติกุสโล.\csromanspacer{} อิธ ปน ภิกฺขเว เอกจฺโจ ฌายี เนว สมาธิสฺมิํ สมาปตฺติกุสโล โหติ, น จ สมาธิสฺมิํ โคจรกุสโล.\csromanspacer{} อิธ ปน ภิกฺขเว เอกจฺโจ ฌายี สมาธิสฺมิํ \csromanfolio{230}สมาปตฺติกุสโล จ โหติ, สมาธิสฺมิํ โคจรกุสโล จ.\csromanspacer{} ตตฺร ฯเปฯ ปวโร จาติ.\csromanspacer{}\csromanspacer{}ปนฺนรสมํ.}

\csromanheader{2}{16. สมาปตฺติมูลก-อภินีหารสุตฺต}
\proseitem{677}{สาวตฺถินิทานํ.\csromanspacer{} จตฺตาโรเม ภิกฺขเว ฌายี.\csromanspacer{} กตเม จตฺตาโร, อิธ ภิกฺขเว เอกจฺโจ ฌายี สมาธิสฺมิํ สมาปตฺติกุสโล โหติ, น สมาธิสฺมิํ อภินีหารกุสโล.\csromanspacer{} อิธ ปน ภิกฺขเว เอกจฺโจ ฌายี สมาธิสฺมิํ อภินีหารกุสโล โหติ, น สมาธิสฺมิํ สมาปตฺติกุสโล.\csromanspacer{} อิธ ปน ภิกฺขเว เอกจฺโจ ฌายี เนว สมาธิสฺมิํ สมาปตฺติกุสโล โหติ, น จ สมาธิสฺมิํ อภินีหารกุสโล.\csromanspacer{} อิธ ปน ภิกฺขเว เอกจฺโจ ฌายี สมาธิสฺมิํ สมาปตฺติกุสโล จ โหติ, สมาธิสฺมิํ อภินีหารกุสโล จ.\csromanspacer{} ตตฺร ฯเปฯ ปวโร จาติ.\csromanspacer{}\csromanspacer{}โสฬสมํ.}

\csromanheader{2}{17. สมาปตฺติมูลกสกฺกจฺจสุตฺต}
\proseitem{678}{สาวตฺถินิทานํ.\csromanspacer{} จตฺตาโรเม ภิกฺขเว ฌายี.\csromanspacer{} กตเม จตฺตาโร, อิธ ภิกฺขเว เอกจฺโจ ฌายี สมาธิสฺมิํ สมาปตฺติกุสโล โหติ, น สมาธิสฺมิํ สกฺกจฺจการี.\csromanspacer{} อิธ ปน ภิกฺขเว เอกจฺโจ ฌายี สมาธิสฺมิํ สกฺกจฺจการี โหติ, น สมาธิสฺมิํ สมาปตฺติกุสโล.\csromanspacer{} อิธ ปน ภิกฺขเว เอกจฺโจ ฌายี เนว สมาธิสฺมิํ สมาปตฺติกุสโล โหติ, น จ สมาธิสฺมิํ สกฺกจฺจการี.\csromanspacer{} อิธ ปน ภิกฺขเว เอกจฺโจ ฌายี สมาธิสฺมิํ สมาปตฺติกุสโล จ โหติ, สมาธิสฺมิํ สกฺกจฺจการี จ.\csromanspacer{} ตตฺร ฯเปฯ ปวโร จาติ.\csromanspacer{}\csromanspacer{}สตฺตรสมํ.}

\csromanheader{2}{18. สมาปตฺติมูลกสาตจฺจสุตฺต}
\proseitem{679}{สาวตฺถินิทานํ.\csromanspacer{} จตฺตาโรเม ภิกฺขเว ฌายี.\csromanspacer{} กตเม จตฺตาโร, อิธ ภิกฺขเว เอกจฺโจ ฌายี สมาธิสฺมิํ สมาปตฺติกุสโล โหติ, น สมาธิสฺมิํ สาตจฺจการี.\csromanspacer{} อิธ ปน ภิกฺขเว เอกจฺโจ ฌายี สมาธิสฺมิํ สาตจฺจการี โหติ, น สมาธิสฺมิํ สมาปตฺติกุสโล.\csromanspacer{} อิธ ปน ภิกฺขเว เอกจฺโจ ฌายี เนว สมาธิสฺมิํ สมาปตฺติกุสโล โหติ, น จ สมาธิสฺมิํ สาตจฺจการี.\csromanspacer{} อิธ ปน ภิกฺขเว เอกจฺโจ ฌายี สมาธิสฺมิํ สมาปตฺติกุสโล จ โหติ, สมาธิสฺมิํ สาตจฺจการี จ.\csromanspacer{} ตตฺร ฯเปฯ ปวโร จาติ.\csromanspacer{}\csromanspacer{}อฏฺฐารสมํ.}

\csromanheader{2}{13. ฌานสํยุตฺต \textbf{19. สมาปตฺติมูลกสปฺปายการีสุตฺต}}
\proseitem{680}{\csromanfolio{231}สาวตฺถินิทานํ.\csromanspacer{} จตฺตาโรเม ภิกฺขเว ฌายี.\csromanspacer{} กตเม จตฺตาโร, อิธ ภิกฺขเว เอกจฺโจ ฌายี สมาธิสฺมิํ สมาปตฺติกุสโล โหติ, น สมาธิสฺมิํ สปฺปายการี.\csromanspacer{} อิธ ปน ภิกฺขเว เอกจฺโจ ฌายี สมาธิสฺมิํ สปฺปายการี โหติ, น สมาธิสฺมิํ สมาปตฺติกุสโล.\csromanspacer{} อิธ ปน ภิกฺขเว เอกจฺโจ ฌายี เนว สมาธิสฺมิํ สมาปตฺติกุสโล โหติ, น จ สมาธิสฺมิํ สปฺปายการี.\csromanspacer{} อิธ ปน ภิกฺขเว เอกจฺโจ ฌายี สมาธิสฺมิํ สมาปตฺติกุสโล จ โหติ, สมาธิสฺมิํ สปฺปายการี จ.\csromanspacer{} ตตฺร ภิกฺขเว ยฺวายํ ฌายี สมาธิสฺมิํ สมาปตฺติกุสโล จ โหติ, สมาธิสฺมิํ สปฺปายการี จ.\csromanspacer{} อยํ อิเมสํ จตุนฺนํ ฌายีนํ อคฺโค จ เสฏฺโฐ จ โมกฺโข จ อุตฺตโม จ ปวโร จ.\csromanspacer{} เสยฺยถาปิ ภิกฺขเว ควา ขีรํ ขีรมฺหา ทธิ ทธิมฺหา นวนีตํ นวนีตมฺหา สปฺปิ สปฺปิมฺหา สปฺปิมณฺโฑ ตตฺร อคฺคมกฺขายติ.\csromanspacer{} เอวเมว โข ภิกฺขเว ยฺวายํ ฌายี สมาธิสฺมิํ สมาปตฺติกุสโล จ โหติ, สมาธิสฺมิํ สปฺปายการี จ.\csromanspacer{} อยํ อิเมสํ จตุนฺนํ ฌายีนํ อคฺโค จ เสฏฺโฐ จ โมกฺโข จ อุตฺตโม จ ปวโร จาติ.\csromanspacer{}\csromanspacer{}เอกูนวีสติมํ.\csromanspacer{}\csromanspacer{}(สมาปตฺติมูลกํ.)}

\csromanheader{2}{20-27. ฐิติมูลกวุฏฺฐานสุตฺตาทิ-อฏฺฐก}
\proseitem{681}{688.\csromanspacer{} สาวตฺถินิทานํ.\csromanspacer{} จตฺตาโรเม ภิกฺขเว ฌายี.\csromanspacer{} กตเม จตฺตาโร, อิธ ภิกฺขเว เอกจฺโจ ฌายี สมาธิสฺมิํ ฐิติกุสโล โหติ, น สมาธิสฺมิํ วุฏฺฐานกุสโล.\csromanspacer{} อิธ ปน ภิกฺขเว เอกจฺโจ ฌายี สมาธิสฺมิํ วุฏฺฐานกุสโล โหติ, น สมาธิสฺมิํ ฐิติกุสโล.\csromanspacer{} อิธ ปน ภิกฺขเว เอกจฺโจ ฌายี เนว สมาธิสฺมิํ ฐิติกุสโล โหติ, น จ สมาธิสฺมิํ วุฏฺฐานกุสโล.\csromanspacer{} อิธ ปน ภิกฺขเว เอกจฺโจ ฌายี สมาธิสฺมิํ ฐิติกุสโล จ โหติ, สมาธิสฺมิํ วุฏฺฐานกุสโล จ.\csromanspacer{} ตตฺร ภิกฺขเว ยฺวายํ ฌายี ฯเปฯ อุตฺตโม จ ปวโร จาติ.\csromanspacer{}\csromanspacer{}วีสติมํ.\csromanspacer{}\csromanspacer{}(ปุริมมูลกานิ วิย ยาว สตฺตวีสติมา ฐิติมูลกสปฺปายการีสุตฺตา อฏฺฐ สุตฺตานิ ปูเรตพฺพานิ.\csromanspacer{} ฐิติมูลกํ.)}

\csromanheader{2}{28-34. วุฏฺฐานมูลกกลฺลิตสุตฺตาทิสตฺตก}
\proseitem{689-695}{สาวตฺถินิทานํ.\csromanspacer{} จตฺตาโรเม ภิกฺขเว ฌายี.\csromanspacer{} กตเม จตฺตาโร, อิธ ภิกฺขเว เอกจฺโจ ฌายี สมาธิสฺมิํ วุฏฺฐานกุสโล โหติ, \csromanfolio{232}น สมาธิสฺมิํ กลฺลิตกุสโล.\csromanspacer{} สมาธิสฺมิํ กลฺลิตกุสโล โหติ, น สมาธิสฺมิํ วุฏฺฐานกุสโล.\csromanspacer{} เนว สมาธิสฺมิํ วุฏฺฐานกุสโล โหติ, น จ สมาธิสฺมิํ กลฺลิตกุสโล.\csromanspacer{} สมาธิสฺมิํ วุฏฺฐานกุสโล จ โหติ, สมาธิสฺมิํ กลฺลิตกุสโล จ.\csromanspacer{} ตตฺร ภิกฺขเว ยฺวายํ ฌายี ฯเปฯ อุตฺตโม จ ปวโร จาติ.\csromanspacer{}\csromanspacer{}อฏฺฐวีสติมํ.\csromanspacer{}\csromanspacer{}(ปุริมมูลกานิ วิย ยาว จตุตฺติํ สติมา วุฏฺฐานมูลกสปฺปายการีสุตฺตา สตฺต สุตฺตานิ ปูเรตพฺพานิ.\csromanspacer{} วุฏฺฐานมูลกํ.)}

\csromanheader{2}{35-48. กลฺลิตมูลก-อารมฺมณสุตฺตาทิฉกฺก}
\proseitem{696-701}{สาวตฺถินิทานํ.\csromanspacer{} สมาธิสฺมิํ กลฺลิตกุสโล โหติ, น สมาธิสฺมิํ อารมฺมณกุสโล.\csromanspacer{} สมาธิสฺมิํ อารมฺมณกุสโล โหติ, น สมาธิสฺมิํ กลฺลิตกุสโล.\csromanspacer{} เนว สมาธิสฺมิํ กลฺลิตกุสโล โหติ, น จ สมาธิสฺมิํ อารมฺมณกุสโล.\csromanspacer{} สมาธิสฺมิํ กลฺลิตกุสโล จ โหติ, สมาธิสฺมิํ อารมฺมณกุสโล จ.\csromanspacer{} ตตฺร ภิกฺขเว ยฺวายํ ฌายี ฯเปฯ อุตฺตโม จ ปวโร จาติ.\csromanspacer{}\csromanspacer{}ปญฺจติํสติมํ.\csromanspacer{}\csromanspacer{}(ปุริมมูลกานิ วิย ยาว จตฺตาลีสมา กลฺลิตมูลกสปฺปายการีสุตฺตา ฉ สุตฺตานิ ปูเรตพฺพานิ.\csromanspacer{} กลฺลิตมูลกํ.)}

\csromanheader{2}{41-45. อารมฺมณมูลกโคจรสุตฺตาทิปญฺจก}
\proseitem{702-706}{สาวตฺถินิทานํ.\csromanspacer{} สมาธิสฺมิํ อารมฺมณกุสโล โหติ, น สมาธิสฺมิํ โคจรกุสโล.\csromanspacer{} สมาธิสฺมิํ โคจรกุสโล โหติ, น สมาธิสฺมิํ อารมฺมณกุสโล.\csromanspacer{} เนว สมาธิสฺมิํ อารมฺมณกุสโล โหติ, น จ สมาธิสฺมิํ โคจรกุสโล.\csromanspacer{} สมาธิสฺมิํ อารมฺมณกุสโล จ โหติ, สมาธิสฺมิํ โคจรกุสโล จ.\csromanspacer{} ตตฺร ภิกฺขเว ยฺวายํ ฌายี ฯเปฯ อุตฺตโม จ ปวโร จาติ.\csromanspacer{}\csromanspacer{}เอกจตฺตาลีสมํ.\csromanspacer{}\csromanspacer{}(ปุริมมูลกานิ วิย ยาว ปญฺจจตฺตาลีสมา อารมฺมณมูลกสปฺปายการีสุตฺตา ปญฺจ สุตฺตานิ ปูเรตพฺพานิ.\csromanspacer{} อารมฺมณมูลกํ.)}

\csromanheader{2}{46-49. โคจรมูลก-อภินีหารสุตฺตาทิจตุกฺก}
\proseitem{707}{สาวตฺถินิทานํ.\csromanspacer{} สมาธิสฺมิํ โคจรกุสโล โหติ, น สมาธิสฺมิํ อภินีหารกุสโล.\csromanspacer{} สมาธิสฺมิํ อภินีหารกุสโล}

\vspace{\tipitakaparskip}
\csromancenter{ฌานสํยุตฺต โหติ, น สมาธิสฺมิํ โคจรกุสโล.\csromanspacer{} เนว สมาธิสฺมิํ โคจรกุสโล โหติ, น จ สมาธิสฺมิํ อภินีหารกุสโล.\csromanspacer{} สมาธิสฺมิํ โคจรกุสโล จ โหติ, สมาธิสฺมิํ อภินีหารกุสโล จ.\csromanspacer{} เสยฺยถาปิ ภิกฺขเว ควา ขีรํ ขีรมฺหา ทธิ ทธิมฺหา นวนีตํ นวนีตมฺหา สปฺปิ สปฺปิมฺหา สปฺปิมณฺโฑ ตตฺร อคฺคมกฺขายติ.\csromanspacer{} เอวเมว โข ภิกฺขเว ยฺวายํ ฌายี สมาธิสฺมิํ โคจรกุสโล จ โหติ, สมาธิสฺมิํ อภินีหารกุสโล จ.\csromanspacer{} อยํ อิเมสํ จตุนฺนํ ฌายีนํ ฯเปฯ อุตฺตโม จ ปวโร จาติ.\csromanspacer{}\csromanspacer{}ฉจตฺตาลีสมํ.}
\proseitem{708}{\csromanfolio{233}สมาธิสฺมิํ โคจรกุสโล โหติ, น สมาธิสฺมิํ สกฺกจฺจการี ฯเปฯ.\csromanspacer{} วิตฺถาเรตพฺพํ.\csromanspacer{}\csromanspacer{}สตฺตจตฺตาลีสมํ.}

\proseitem{709}{สมาธิสฺมิํ โคจรกุสโล โหติ, น สมาธิสฺมิํ สาตจฺจการี ฯเปฯ.\csromanspacer{}\csromanspacer{}อฏฺฐจตฺตาลีสมํ.}

\proseitem{710}{สมาธิสฺมิํ โคจรกุสโล โหติ, น สมาธิสฺมิํ สปฺปายการี ฯเปฯ.\csromanspacer{}\csromanspacer{}เอกูนปญฺญาสมํ.\csromanspacer{}\csromanspacer{}(โคจรมูลกํ.)}

\csromanheader{2}{50-52. อภินีหารมูลกสกฺกจฺจสุตฺตาทิติก}
\proseitem{711}{สาวตฺถินิทานํ.\csromanspacer{} สมาธิสฺมิํ อภินีหารกุสโล โหติ, น สมาธิสฺมิํ สกฺกจฺจการี.\csromanspacer{} สมาธิสฺมิํ สกฺกจฺจการี โหติ, น สมาธิสฺมิํ อภินีหารกุสโล.\csromanspacer{} เนว สมาธิสฺมิํ อภินีหารกุสโล โหติ, น จ สมาธิสฺมิํ สกฺกจฺจการี.\csromanspacer{} สมาธิสฺมิํ อภินีหารกุสโล จ โหติ, สมาธิสฺมิํ สกฺกจฺจการี จ.\csromanspacer{} ตตฺร ภิกฺขเว ยฺวายํ ฌายี ฯเปฯ อุตฺตโม จ ปวโร จาติ.\csromanspacer{}\csromanspacer{}ปญฺญาสมํ.}

\proseitem{712}{สมาธิสฺมิํ อภินีหารกุสโล โหติ, น สมาธิสฺมิํ สาตจฺจการี ฯเปฯ.\csromanspacer{}\csromanspacer{}เอกปญฺญาสมํ.}

\proseitem{713}{สมาธิสฺมิํ อภินีหารกุสโล โหติ, น สมาธิสฺมิํ สปฺปายการี ฯเปฯ.\csromanspacer{}\csromanspacer{}เทฺวปญฺญาสมํ.\csromanspacer{}\csromanspacer{}(อภินีหารมูลกํ.)}

\csromanheader{2}{53-54. สกฺกจฺจมูลกสาตจฺจการีสุตฺตาทิทุก}
\proseitem{714}{\csromanfolio{234}สาวตฺถินิทานํ.\csromanspacer{} สมาธิสฺมิํ สกฺกจฺจการี โหติ, น สมาธิสฺมิํ สาตจฺจการี.\csromanspacer{} สมาธิสฺมิํ สาตจฺจการี โหติ, น สมาธิสฺมิํ สกฺกจฺจการี.\csromanspacer{} เนว สมาธิสฺมิํ สกฺกจฺจการี โหติ, น จ สมาธิสฺมิํ สาตจฺจการี.\csromanspacer{} สมาธิสฺมิํ สกฺกจฺจการี จ โหติ, สมาธิสฺมิํ สาตจฺจการี จ.\csromanspacer{} ตตฺร ภิกฺขเว ยฺวายํ ฯเปฯ อุตฺตโม จ ปวโร จาติ.\csromanspacer{}\csromanspacer{}เตปญฺญาสมํ.}

\proseitem{715}{สมาธิสฺมิํ สกฺกจฺจการี โหติ, น สมาธิสฺมิํ สปฺปายการี ฯเปฯ.\csromanspacer{}\csromanspacer{}จตุปญฺญาสมํ.}

\csromanheader{2}{55. สาตจฺจมูลกสปฺปายการีสุตฺต}
\proseitem{716}{สาวตฺถินิทานํ.\csromanspacer{} จตฺตาโรเม ภิกฺขเว ฌายี.\csromanspacer{} กตเม จตฺตาโร, อิธ ภิกฺขเว เอกจฺโจ ฌายี สมาธิสฺมิํ สาตจฺจการี โหติ, น สมาธิสฺมิํ สปฺปายการี.\csromanspacer{} อิธ ปน ภิกฺขเว เอกจฺโจ ฌายี สมาธิสฺมิํ สปฺปายการี โหติ, น สมาธิสฺมิํ สาตจฺจการี.\csromanspacer{} อิธ ปน ภิกฺขเว เอกจฺโจ ฌายี เนว สมาธิสฺมิํ สาตจฺจการี โหติ, น จ สมาธิสฺมิํ สปฺปายการี.\csromanspacer{} อิธ ปน ภิกฺขเว เอกจฺโจ ฌายี สมาธิสฺมิํ สาตจฺจการี จ โหติ, สมาธิสฺมิํ สปฺปายการี จ.\csromanspacer{} ตตฺร ภิกฺขเว ยฺวายํ ฌายี สมาธิสฺมิํ สาตจฺจการี จ โหติ, สมาธิสฺมิํ สปฺปายการี จ.\csromanspacer{} อยํ อิเมสํ จตุนฺนํ ฌายีนํ อคฺโค จ เสฏฺโฐ จ โมกฺโข จ อุตฺตโม จ ปวโร จ.\csromanspacer{} เสยฺยถาปิ ภิกฺขเว ควา ขีรํ ขีรมฺหา ทธิ ทธิมฺหา นวนีตํ นวนีตมฺหา สปฺปิ สปฺปิมฺหา สปฺปิมณฺโฑ ตตฺร อคฺคมกฺขายติ.\csromanspacer{} เอวเมว โข ภิกฺขเว ยฺวายํ ฌายี สมาธิสฺมิํ สาตจฺจการี จ โหติ, สมาธิสฺมิํ สปฺปายการี จ.\csromanspacer{} อยํ อิเมสํ จตุนฺนํ ฌายีนํ อคฺโค จ เสฏฺโฐ จ โมกฺโข จ อุตฺตโม จ ปวโร จาติ.\csromanspacer{} อิทมโวจ ภควา, อตฺตมนา เต ภิกฺขู ภควโต ภาสิตํ อภินนฺทุนฺติ.\csromanspacer{}\csromanspacer{}ปญฺจปญฺญาสมํ.}

\prose{(ยถา ปญฺจปญฺญาสํ เวยฺยากรณานิ โหนฺติ, ตถา วิตฺถาเรตพฺพานิ.)}

\vspace{\tipitakaparskip}
\csromancenter{\textbf{ฌานสํยุตฺตํ}\footnote{สมาธิสํยุตฺตํ (สฺยา, กํ)} \textbf{สมตฺตํ.}}
\titlehead{13. ฌานสํยุตฺต \textbf{ตสฺสุทฺทานํ}}
\csromangathameasure{สมาธิ สมาปตฺติ ฐิติ จ, วุฏฺฐานํ กลฺลิตารมฺมเณน จ. \\ โคจรา อภินีหาโร สกฺกจฺจ, สาตจฺจ อโถปิ สปฺปายนฺติ. ขนฺธวคฺโค ตติโย.}
\csromangathagroup{\csromangathastanza{\csromanfolio{235}สมาธิ สมาปตฺติ ฐิติ จ, วุฏฺฐานํ กลฺลิตารมฺมเณน จ. \\ โคจรา อภินีหาโร สกฺกจฺจ, สาตจฺจ อโถปิ สปฺปายนฺติ.\csromanspacer{} ขนฺธวคฺโค ตติโย.}}

\titlehead{\textbf{ตสฺสุทฺทานํ}}
\csromangathameasure{ขนฺธ ราธสํยุตฺตญฺจ, ทิฏฺฐิ โอกฺกนฺต อุปฺปาทา. \\ กิเลส สาริปุตฺตา จ, นาคา สุปณฺณ คนฺธพฺพา. \\ วลาห วจฺฉ ฌานนฺติ, ขนฺธวคฺคมฺหิ เตรสาติ.}
\csromangathagroup{\csromangathastanza{ขนฺธ ราธสํยุตฺตญฺจ, ทิฏฺฐิ โอกฺกนฺต\footnote{โอกฺกนฺติ (สพฺพตฺถ)} อุปฺปาทา. \\ กิเลส สาริปุตฺตา จ, นาคา สุปณฺณ คนฺธพฺพา. \\ วลาห วจฺฉ ฌานนฺติ, ขนฺธวคฺคมฺหิ เตรสาติ.}}

\nitthitam{\textbf{ขนฺธวคฺคสํยุตฺตปาฬิ นิฏฺฐิตา.}}
\csromanheader{2}{\textbf{สํยุตฺตนิกาย}}
\csromanmatikaanchor{mtk.79}
\csromantocmarknopage{chapter}{สฬายตนวคฺคสํยุตฺตปาฬิ}{mtk.79}
\bookmark[dest={mtk.79},level=0]{สฬายตนวคฺคสํยุตฺตปาฬิ}
\gambhira{\textbf{สฬายตนวคฺคสํยุตฺตปาฬิ}}
\namakkaram{นโม ตสฺส ภควโต อรหโต สมฺมาสมฺพุทฺธสฺส.}
\csromanmatikaanchor{mtk.80}
\csromantocmarknopage{chapter}{1. สฬายตนสํยุตฺต}{mtk.80}
\bookmark[dest={mtk.80},level=0]{1. สฬายตนสํยุตฺต}
\csromantocmarknopage{section}{1. อนิจฺจวคฺค}{mtk.81}
\bookmark[dest={mtk.81},level=1]{1. อนิจฺจวคฺค}
\chapterhead{1. สฬายตนสํยุตฺต}
\chapterhead{1. อนิจฺจวคฺค}
\csromanheader{2}{1. อชฺฌตฺตานิจฺจสุตฺต}
\csromanmatikaanchor{mtk.82}
\csromantocmarkref{subsection}{1-3. อชฺฌตฺตานิจฺจสุตฺตาทิ}{mtk.82}
\bookmark[dest={mtk.82},level=2]{1-3. อชฺฌตฺตานิจฺจสุตฺตาทิ}
\proseitem{1}{\csromanfolio{236}เอวํ เม สุตํ\csromandash{}เอกํ สมยํ ภควา สาวตฺถิยํ วิหรติ เชตวเน อนาถปิณฺฑิกสฺส อาราเม.\csromanspacer{} ตตฺร โข ภควา ภิกฺขู อามนฺเตสิ “ภิกฺขโว”ติ. “ภทนฺเต”ติ เต ภิกฺขู ภควโต ปจฺจสฺโสสุํ, ภควา เอตทโวจ\csromandash{}}

\prose{จกฺขุํ ภิกฺขเว อนิจฺจํ, ยทนิจฺจํ ตํ ทุกฺขํ, ยํ ทุกฺขํ ตทนตฺตา, ยทนตฺตา ตํ “เนตํ มม, เนโสหมสฺมิ, น เมโส อตฺตา”ติ เอวเมตํ ยถาภูตํ สมฺมปฺปญฺญาย ทฏฺฐพฺพํ.\csromanspacer{} โสตํ อนิจฺจํ, ยทนิจฺจํ ฯเปฯ.\csromanspacer{} ฆานํ อนิจฺจํ, ยทนิจฺจํ ฯเปฯ.\csromanspacer{} ชิวฺหา อนิจฺจา, ยทนิจฺจํ ตํ ทุกฺขํ, ยํ ทุกฺขํ ตทนตฺตา, ยทนตฺตา ตํ “เนตํ มม, เนโสหมสฺมิ, น เมโส อตฺตา”ติ เอวเมตํ ยถาภูตํ สมฺมปฺปญฺญาย ทฏฺฐพฺพํ.\csromanspacer{} กาโย อนิจฺโจ, ยทนิจฺจํ ฯเปฯ.\csromanspacer{} มโน อนิจฺโจ, ยทนิจฺจํ ตํ ทุกฺขํ, ยํ ทุกฺขํ ตทนตฺตา, ยทนตฺตา ตํ “เนตํ มม, เนโสหมสฺมิ, น เมโส อตฺตา”ติ เอวเมตํ ยถาภูตํ สมฺมปฺปญฺญาย ทฏฺฐพฺพํ.\csromanspacer{} เอวํ ปสฺสํ ภิกฺขเว สุตวา อริยสาวโก จกฺขุสฺมิมฺปิ นิพฺพินฺทติ, โสตสฺมิมฺปิ นิพฺพินฺทติ, ฆานสฺมิมฺปิ นิพฺพินฺทติ, ชิวฺหายปิ นิพฺพินฺทติ, กายสฺมิมฺปิ}

\vspace{\tipitakaparskip}
\csromancenter{สฬายตนสํยุตฺต นิพฺพินฺทติ, มนสฺมิมฺปิ นิพฺพินฺทติ, นิพฺพินฺทํ วิรชฺชติ, วิราคา วิมุจฺจติ, วิมุตฺตสฺมิํ “วิมุตฺตมฺ”อิติ ญาณํ โหติ, “ขีณา ชาติ, วุสิตํ พฺรหฺมจริยํ, กตํ กรณียํ, นาปรํ อิตฺถตฺตายา”ติ ปชานาตีติ.\csromanspacer{}\csromanspacer{}ปฐมํ.}
\csromanheader{2}{2. อชฺฌตฺตทุกฺขสุตฺต}
\csromanmatikaanchor{mtk.83}
\csromantocmarkref{subsection}{4-6. พาหิรานิจฺจสุตฺตาทิ}{mtk.83}
\bookmark[dest={mtk.83},level=2]{4-6. พาหิรานิจฺจสุตฺตาทิ}
\proseitem{2}{\csromanfolio{237}จกฺขุํ ภิกฺขเว ทุกฺขํ, ยํ ทุกฺขํ ตทนตฺตา, ยทนตฺตา ตํ “เนตํ มม, เนโสหมสฺมิ, น เมโส อตฺตา”ติ เอวเมตํ ยถาภูตํ สมฺมปฺปญฺญาย ทฏฺฐพฺพํ.\csromanspacer{} โสตํ ทุกฺขํ ฯเปฯ.\csromanspacer{} ฆานํ ทุกฺขํ.\csromanspacer{} ชิวฺหา ทุกฺขา.\csromanspacer{} กาโย ทุกฺโข.\csromanspacer{} มโน ทุกฺโข, ยํ ทุกฺขํ ตทนตฺตา, ยทนตฺตา ตํ “เนตํ มม, เนโสหมสฺมิ, น เมโส อตฺตา”ติ เอวเมตํ ยถาภูตํ สมฺมปฺปญฺญาย ทฏฺฐพฺพํ.\csromanspacer{} เอวํ ปสฺสํ ฯเปฯ นาปรํ อิตฺถตฺตายาติ ปชานาตีติ.\csromanspacer{}\csromanspacer{}ทุติยํ.}

\csromanheader{2}{3. อชฺฌตฺตานตฺตสุตฺต}
\proseitem{3}{จกฺขุํ ภิกฺขเว อนตฺตา, ยทนตฺตา ตํ “เนตํ มม, เนโสหมสฺมิ, น เมโส อตฺตา”ติ เอวเมตํ ยถาภูตํ สมฺมปฺปญฺญาย ทฏฺฐพฺพํ.\csromanspacer{} โสตํ อนตฺตา ฯเปฯ.\csromanspacer{} ฆานํ อนตฺตา.\csromanspacer{} ชิวฺหา อนตฺตา.\csromanspacer{} กาโย อนตฺตา.\csromanspacer{} มโน อนตฺตา, ยทนตฺตา ตํ “เนตํ มม, เนโสหมสฺมิ, น เมโส อตฺตา”ติ เอวเมตํ ยถาภูตํ สมฺมปฺปญฺญาย ทฏฺฐพฺพํ.\csromanspacer{} เอวํ ปสฺสํ ฯเปฯ นาปรํ อิตฺถตฺตายาติ ปชานาตีติ.\csromanspacer{}\csromanspacer{}ตติยํ.}

\csromanheader{2}{4. พาหิรานิจฺจสุตฺต}
\proseitem{4}{รูปา ภิกฺขเว อนิจฺจา.\csromanspacer{} ยทนิจฺจํ ตํ ทุกฺขํ, ยํ ทุกฺขํ ตทนตฺตา, ยทนตฺตา ตํ “เนตํ มม, เนโสหมสฺมิ, น เมโส อตฺตา”ติ เอวเมตํ ยถาภูตํ สมฺมปฺปญฺญาย ทฏฺฐพฺพํ.\csromanspacer{} สทฺทา.\csromanspacer{} คนฺธา.\csromanspacer{} รสา.\csromanspacer{} โผฏฺฐพฺพา.\csromanspacer{} ธมฺมา อนิจฺจา, ยทนิจฺจํ ตํ ทุกฺขํ, ยํ ทุกฺขํ ตทนตฺตา, ยทนตฺตา ตํ “เนตํ มม, เนโสหมสฺมิ, น เมโส อตฺตา”ติ เอวเมตํ ยถาภูตํ สมฺมปฺปญฺญาย ทฏฺฐพฺพํ.\csromanspacer{} เอวํ ปสฺสํ ภิกฺขเว สุตวา อริยสาวโก รูเปสุปิ นิพฺพินฺทติ, สทฺเทสุปิ นิพฺพินฺทติ, คนฺเธสุปิ นิพฺพินฺทติ, รเสสุปิ นิพฺพินฺทติ, โผฏฺฐพฺเพสุปิ นิพฺพินฺทติ, ธมฺเมสุปิ นิพฺพินฺทติ, นิพฺพินฺทํ วิรชฺชติ, วิราคา วิมุจฺจติ, วิมุตฺตสฺมิํ “วิมุตฺตมฺ”อิติ ญาณํ โหติ, “ขีณา ชาติ, วุสิตํ พฺรหฺมจริยํ, กตํ กรณียํ, นาปรํ อิตฺถตฺตายา”ติ ปชานาตีติ.\csromanspacer{}\csromanspacer{}จตุตฺถํ.}

\gambhira{สฬายตนวคฺคสํยุตฺตปาฬิ}
\csromanheader{2}{5. พาหิรทุกฺขสุตฺต}
\csromanmatikaanchor{mtk.84}
\csromantocmarkref{subsection}{7-9. อชฺฌตฺตานิจฺจาตีตานาคตสุตฺตาทิ}{mtk.84}
\bookmark[dest={mtk.84},level=2]{7-9. อชฺฌตฺตานิจฺจาตีตานาคตสุตฺตาทิ}
\proseitem{5}{\csromanfolio{238}รูปา ภิกฺขเว ทุกฺขา, ยํ ทุกฺขํ ตทนตฺตา, ยทนตฺตา ตํ “เนตํ มม, เนโสหมสฺมิ, น เมโส อตฺตา”ติ เอวเมตํ ยถาภูตํ สมฺมปฺปญฺญาย ทฏฺฐพฺพํ.\csromanspacer{} สทฺทา.\csromanspacer{} คนฺธา.\csromanspacer{} รสา.\csromanspacer{} โผฏฺฐพฺพา.\csromanspacer{} ธมฺมา ทุกฺขา, ยํ ทุกฺขํ ตทนตฺตา, ยทนตฺตา ตํ “เนตํ มม, เนโสหมสฺมิ, น เมโส อตฺตา”ติ เอวเมตํ ยถาภูตํ สมฺมปฺปญฺญาย ทฏฺฐพฺพํ.\csromanspacer{} เอวํ ปสฺสํ ฯเปฯ นาปรํ อิตฺถตฺตายาติ ปชานาตีติ.\csromanspacer{}\csromanspacer{}ปญฺจมํ.}

\csromanheader{2}{6. พาหิรานตฺตสุตฺต}
\proseitem{6}{รูปา ภิกฺขเว อนตฺตา, ยทนตฺตา ตํ “เนตํ มม, เนโสหมสฺมิ, น เมโส อตฺตา”ติ เอวเมตํ ยถาภูตํ สมฺมปฺปญฺญาย ทฏฺฐพฺพํ ฯเปฯ.\csromanspacer{} สทฺทา.\csromanspacer{} คนฺธา.\csromanspacer{} รสา.\csromanspacer{} โผฏฺฐพฺพา.\csromanspacer{} ธมฺมา อนตฺตา, ยทนตฺตา ตํ “เนตํ มม, เนโสหมสฺมิ, น เมโส อตฺตา”ติ เอวเมตํ ยถาภูตํ สมฺมปฺปญฺญาย ทฏฺฐพฺพํ.\csromanspacer{} เอวํ ปสฺสํ ฯเปฯ นาปรํ อิตฺถตฺตายาติ ปชานาตีติ.\csromanspacer{}\csromanspacer{}ฉฏฺฐํ.}

\csromanheader{2}{7. อชฺฌตฺตานิจฺจาตีตานาคตสุตฺต}
\proseitem{7}{จกฺขุํ ภิกฺขเว อนิจฺจํ อตีตานาคตํ, โก ปน วาโท ปจฺจุปฺปนฺนสฺส.\csromanspacer{} เอวํ ปสฺสํ ภิกฺขเว สุตวา อริยสาวโก อตีตสฺมิํ จกฺขุสฺมิํ อนเปกฺโข โหติ, อนาคตํ จกฺขุํ นาภินนฺทติ, ปจฺจุปฺปนฺนสฺส จกฺขุสฺส นิพฺพิทาย วิราคาย นิโรธาย ปฏิปนฺโน โหติ.\csromanspacer{} โสตํ อนิจฺจํ.\csromanspacer{} ฆานํ อนิจฺจํ.\csromanspacer{} ชิวฺหา อนิจฺจา อตีตานาคตา, โก ปน วาโท ปจฺจุปฺปนฺนาย.\csromanspacer{} เอวํ ปสฺสํ ภิกฺขเว สุตวา อริยสาวโก อตีตาย ชิวฺหาย อนเปกฺโข โหติ, อนาคตํ ชิวฺหํ นาภินนฺทติ, ปจฺจุปฺปนฺนาย ชิวฺหาย นิพฺพิทาย วิราคาย นิโรธาย ปฏิปนฺโน โหติ.\csromanspacer{} กาโย อนิจฺโจ ฯเปฯ.\csromanspacer{} มโน อนิจฺโจ อตีตานาคโต, โก ปน วาโท ปจฺจุปฺปนฺนสฺส.\csromanspacer{} เอวํ ปสฺสํ ภิกฺขเว สุตวา อริยสาวโก อตีตสฺมิํ มนสฺมิํ อนเปกฺโข โหติ, อนาคตํ มนํ นาภินนฺทติ, ปจฺจุปฺปนฺนสฺส มนสฺส นิพฺพิทาย วิราคาย นิโรธาย ปฏิปนฺโน โหตีติ.\csromanspacer{}\csromanspacer{}สตฺตมํ.}

\csromanheader{2}{1. สฬายตนสํยุตฺต \textbf{8. อชฺฌตฺตทุกฺขาตีตานาคตสุตฺต}}
\csromanmatikaanchor{mtk.85}
\csromantocmarkref{subsection}{10-12. พาหิรานิจฺจาตีตานาคตสุตฺตาทิ}{mtk.85}
\bookmark[dest={mtk.85},level=2]{10-12. พาหิรานิจฺจาตีตานาคตสุตฺตาทิ}
\proseitem{8}{\csromanfolio{239}จกฺขุํ ภิกฺขเว ทุกฺขํ อตีตานาคตํ, โก ปน วาโท ปจฺจุปฺปนฺนสฺส.\csromanspacer{} เอวํ ปสฺสํ ภิกฺขเว สุตวา อริยสาวโก อตีตสฺมิํ จกฺขุสฺมิํ อนเปกฺโข โหติ, อนาคตํ จกฺขุํ นาภินนฺทติ, ปจฺจุปฺปนฺนสฺส จกฺขุสฺส นิพฺพิทาย วิราคาย นิโรธาย ปฏิปนฺโน โหติ.\csromanspacer{} โสตํ ทุกฺขํ ฯเปฯ.\csromanspacer{} ฆานํ ทุกฺขํ ฯเปฯ.\csromanspacer{} ชิวฺหา ทุกฺขา อตีตานาคตา, โก ปน วาโท ปจฺจุปฺปนฺนาย.\csromanspacer{} เอวํ ปสฺสํ ภิกฺขเว สุตวา อริยสาวโก อตีตาย ชิวฺหาย อนเปกฺโข โหติ, อนาคตํ ชิวฺหํ นาภินนฺทติ, ปจฺจุปฺปนฺนาย ชิวฺหาย นิพฺพิทาย วิราคาย นิโรธาย ปฏิปนฺโน โหติ.\csromanspacer{} กาโย ทุกฺโข ฯเปฯ.\csromanspacer{} มโน ทุกฺโข อตีตานาคโต, โก ปน วาโท ปจฺจุปฺปนฺนสฺส.\csromanspacer{} เอวํ ปสฺสํ ภิกฺขเว สุตวา อริยสาวโก อตีตสฺมิํ มนสฺมิํ อนเปกฺโข โหติ, อนาคตํ มนํ นาภินนฺทติ, ปจฺจุปฺปนฺนสฺส มนสฺส นิพฺพิทาย วิราคาย นิโรธาย ปฏิปนฺโน โหตีติ.\csromanspacer{}\csromanspacer{}อฏฺฐมํ.}

\csromanheader{2}{9. อชฺฌตฺตานตฺตาตีตานาคตสุตฺต}
\proseitem{9}{จกฺขุํ ภิกฺขเว อนตฺตา อตีตานาคตํ, โก ปน วาโท ปจฺจุปฺปนฺนสฺส.\csromanspacer{} เอวํ ปสฺสํ ภิกฺขเว สุตวา อริยสาวโก อตีตสฺมิํ จกฺขุสฺมิํ อนเปกฺโข โหติ, อนาคตํ จกฺขุํ นาภินนฺทติ, ปจฺจุปฺปนฺนสฺส จกฺขุสฺส นิพฺพิทาย วิราคาย นิโรธาย ปฏิปนฺโน โหติ.\csromanspacer{} โสตํ อนตฺตา ฯเปฯ.\csromanspacer{} ฆานํ อนตฺตา ฯเปฯ.\csromanspacer{} ชิวฺหา อนตฺตา อตีตานาคตา, โก ปน วาโท ปจฺจุปฺปนฺนาย.\csromanspacer{} เอวํ ปสฺสํ ภิกฺขเว สุตวา อริยสาวโก อตีตาย ชิวฺหาย อนเปกฺโข โหติ, อนาคตํ ชิวฺหํ นาภินนฺทติ, ปจฺจุปฺปนฺนาย ชิวฺหาย นิพฺพิทาย วิราคาย นิโรธาย ปฏิปนฺโน โหติ.\csromanspacer{} กาโย อนตฺตา ฯเปฯ.\csromanspacer{} มโน อนตฺตา อตีตานาคโต, โก ปน วาโท ปจฺจุปฺปนฺนสฺส.\csromanspacer{} เอวํ ปสฺสํ ภิกฺขเว สุตวา อริยสาวโก อตีตสฺมิํ มนสฺมิํ อนเปกฺโข โหติ, อนาคตํ มนํ นาภินนฺทติ, ปจฺจปฺปุนฺนสฺส มนสฺส นิพฺพิทาย วิราคาย นิโรธาย ปฏิปนฺโน โหตีติ.\csromanspacer{}\csromanspacer{}นวมํ.}

\csromanheader{2}{10. พาหิรานิจฺจาตีตานาคตสุตฺต}
\proseitem{10}{รูปา ภิกฺขเว อนิจฺจา อตีตานาคตา, โก ปน วาโท ปจฺจุปฺปนฺนานํ.\csromanspacer{} เอวํ ปสฺสํ ภิกฺขเว สุตวา อริยสาวโก อตีเตสุ}

\gambhira{สฬายตนวคฺคสํยุตฺตปาฬิ}
\prose{\csromanfolio{240}รูเปสุ อนเปกฺโข โหติ, อนาคเต รูเป นาภินนฺทติ, ปจฺจุปฺปนฺนานํ รูปานํ นิพฺพิทาย วิราคาย นิโรธาย ปฏิปนฺโน โหติ.\csromanspacer{} สทฺทา.\csromanspacer{} คนฺธา.\csromanspacer{} รสา.\csromanspacer{} โผฏฺฐพฺพา.\csromanspacer{} ธมฺมา อนิจฺจา อตีตานาคตา, โก ปน วาโท ปจฺจุปฺปนฺนานํ.\csromanspacer{} เอวํ ปสฺสํ ภิกฺขเว สุตวา อริยสาวโก อตีเตสุ ธมฺเมสุ อนเปกฺโข โหติ, อนาคเต ธมฺเม นาภินนฺทติ, ปจฺจุปฺปนฺนานํ ธมฺมานํ นิพฺพิทาย วิราคาย นิโรธาย ปฏิปนฺโน โหตีติ.\csromanspacer{}\csromanspacer{}ทสมํ.}

\csromanheader{2}{11. พาหิรทุกฺขาตีตานาคตสุตฺต}
\proseitem{11}{รูปา ภิกฺขเว ทุกฺขา อตีตานาคตา, โก ปน วาโท ปจฺจุปฺปนฺนานํ.\csromanspacer{} เอวํ ปสฺสํ ภิกฺขเว สุตวา อริยสาวโก อตีเตสุ รูเปสุ อนเปกฺโข โหติ, อนาคเต รูเป นาภินนฺทติ, ปจฺจุปฺปนฺนานํ รูปานํ นิพฺพิทาย วิราคาย นิโรธาย ปฏิปนฺโน โหติ ฯเปฯ.\csromanspacer{} เอโกทสมํ.}

\csromanheader{2}{12. พาหิรานตฺตาตีตานาคตสุตฺต}
\proseitem{12}{รูปา ภิกฺขเว อนตฺตา อตีตานาคตา, โก ปน วาโท ปจฺจุปฺปนฺนานํ.\csromanspacer{} เอวํ ปสฺสํ ภิกฺขเว สุตวา อริยสาวโก อตีเตสุ รูเปสุ อนเปกฺโข โหติ, อนาคเต รูเป นาภินนฺทติ, ปจฺจุปฺปนฺนานํ รูปานํ นิพฺพิทาย วิราคาย นิโรธาย ปฏิปนฺโน โหติ.\csromanspacer{} สทฺทา.\csromanspacer{} คนฺธา.\csromanspacer{} รสา.\csromanspacer{} โผฏฺฐพฺพา.\csromanspacer{} ธมฺมา อนตฺตา อตีตานาคตา, โก ปน วาโท ปจฺจุปฺปนฺนานํ.\csromanspacer{} เอวํ ปสฺสํ ภิกฺขเว สุตวา อริยสาวโก อตีเตสุ ธมฺเมสุ อนเปกฺโข โหติ, อนาคเต ธมฺเม นาภินนฺทติ, ปจฺจุปฺปนฺนานํ ธมฺมานํ นิพฺพิทาย วิราคาย นิโรธาย ปฏิปนฺโน โหตีติ.\csromanspacer{}\csromanspacer{}ทฺวาทสมํ.\csromanspacer{} อนิจฺจวคฺโค ปฐโม.}

\titlehead{\textbf{ตสฺสุทฺทานํ}}
\csromangathameasure{อนิจฺจํ ทุกฺขํ อนตฺตา จ, ตโย อชฺฌตฺตพาหิรา. \\ ยทนิจฺเจน ตโย วุตฺตา, เต เต อชฺฌตฺตพาหิราติ.}
\csromangathagroup{\csromangathastanza{อนิจฺจํ ทุกฺขํ อนตฺตา จ, ตโย อชฺฌตฺตพาหิรา. \\ ยทนิจฺเจน ตโย วุตฺตา, เต เต อชฺฌตฺตพาหิราติ.}}

\csromanmatikaanchor{mtk.86}
\csromantocmarknopage{section}{2. ยมกวคฺค}{mtk.86}
\bookmark[dest={mtk.86},level=1]{2. ยมกวคฺค}
\csromanheader{1}{1. สฬายตนสํยุตฺต \textbf{2. ยมกวคฺค}}
\csromanheader{2}{1. ปฐมปุพฺเพสมฺโพธสุตฺต}
\csromanmatikaanchor{mtk.87}
\csromantocmarkref{subsection}{1-2. ปุพฺเพสมฺโพธสุตฺต}{mtk.87}
\bookmark[dest={mtk.87},level=2]{1-2. ปุพฺเพสมฺโพธสุตฺต}
\proseitem{13}{\csromanfolio{241}สาวตฺถินิทานํ.\csromanspacer{} ปุพฺเพว เม ภิกฺขเว สมฺโพธา อนภิสมฺพุทฺธสฺส โพธิสตฺตสฺเสว สโต เอตทโหสิ “โก นุ โข จกฺขุสฺส อสฺสาโท, โก อาทีนโว, กิํนิสฺสรณํ.\csromanspacer{} โก โสตสฺส ฯเปฯ.\csromanspacer{} โก ฆานสฺส.\csromanspacer{} โก ชิวฺหาย.\csromanspacer{} โก กายสฺส.\csromanspacer{} โก มนสฺส อสฺสาโท, โก อาทีนโว, กิํ นิสฺสรณนฺ”ติ.\csromanspacer{} ตสฺส มยฺหํ ภิกฺขเว เอตทโหสิ, ยํ โข จกฺขุํ ปฏิจฺจ อุปฺปชฺชติ สุขํ โสมนสฺสํ, อยํ จกฺขุสฺส อสฺสาโท.\csromanspacer{} ยํ จกฺขุํ อนิจฺจํ ทุกฺขํ วิปริณามธมฺมํ, อยํ จกฺขุสฺส อาทีนโว.\csromanspacer{} โย จกฺขุสฺมิํ ฉนฺทราควินโย ฉนฺทราคปฺปหานํ, อิทํ จกฺขุสฺส นิสฺสรณํ.\csromanspacer{} ยํ โสตํ ฯเปฯ.\csromanspacer{} ยํ ฆานํ ฯเปฯ.\csromanspacer{} ยํ ชิวฺหํ ปฏิจฺจ อุปฺปชฺชติ สุขํ โสมนสฺสํ, อยํ ชิวฺหาย อสฺสาโท.\csromanspacer{} ยํ\footnote{ยา (สี, สฺยา, กํ, ก)} ชิวฺหา อนิจฺจา ทุกฺขา วิปริณามธมฺมา, อยํ ชิวฺหาย อาทีนโว.\csromanspacer{} โย ชิวฺหาย ฉนฺทราควินโย ฉนฺทราคปฺปหานํ, อิทํ ชิวฺหาย นิสฺสรณํ.\csromanspacer{} ยํ กายํ ฯเปฯ.\csromanspacer{} ยํ มนํ ปฏิจฺจ อุปฺปชฺชติ สุขํ โสมนสฺสํ, อยํ มนสฺส อสฺสาโท.\csromanspacer{} ยํ\footnote{โย (สี, สฺยา, กํ, ก)} มโน อนิจฺโจ ทุกฺโข วิปริณามธมฺโม, อยํ มนสฺส อาทีนโว.\csromanspacer{} โย มนสฺมิํ ฉนฺทราควินโย ฉนฺทราคปฺปหานํ, อิทํ มนสฺส นิสฺสรณนฺติ.\csromanspacer{} ยาวกีวญฺจาหํ ภิกฺขเว อิเมสํ ฉนฺนํ อชฺฌตฺติกานํ อายตนานํ เอวํ อสฺสาทญฺจ อสฺสาทโต อาทีนวญฺจ อาทีนวโต นิสฺสรณญฺจ นิสฺสรณโต ยถาภูตํ นาพฺภญฺญาสิํ, เนว ตาวาหํ ภิกฺขเว สเทวเก โลเก สมารเก สพฺรหฺมเก สสฺสมณพฺราหฺมณิยา ปชาย สเทวมนุสฺสาย อนุตฺตรํ สมฺมาสมฺโพธิํ อภิสมฺพุทฺโธติ\footnote{สพฺพตฺถาปิ เอวเมว อิติสทฺเทน สห ทิสฺสติ.} ปจฺจญฺญาสิํ.\csromanspacer{} ยโต จ ขฺวาหํ ภิกฺขเว อิเมสํ ฉนฺนํ อชฺฌตฺติกานํ อายตนานํ เอวํ อสฺสาทญฺจ อสฺสาทโต อาทีนวญฺจ อาทีนวโต นิสฺสรณญฺจ นิสฺสรณโต ยถาภูตํ อพฺภญฺญาสิํ, อถาหํ ภิกฺขเว สเทวเก โลเก สมารเก สสฺสมณพฺราหฺมณิยา ปชาย สเทวมนุสฺสาย อนุตฺตรํ สมฺโมสมฺโพธิํ อภิสมฺพุทฺโธติ ปจฺจญฺญาสิํ.\csromanspacer{} ญาณญฺจ ปน เม ทสฺสนํ อุทปาทิ, อกุปฺปา เม วิมุตฺติ\footnote{เจโตวิมุตฺติ (สี, อิ, ก) เอวมุปริปิ.}, อยมนฺติมา ชาติ, นตฺถิ ทานิ ปุนพฺภโวติ.\csromanspacer{}\csromanspacer{}ปฐมํ.}

\gambhira{สฬายตนวคฺคสํยุตฺตปาฬิ}
\csromanheader{2}{2. ทุติยปุพฺเพสมฺโพธสุตฺต}
\proseitem{14}{\csromanfolio{242}ปุพฺเพว เม ภิกฺขเว สมฺโพธา อนภิสมฺพุทฺธสฺส โพธิสตฺตสฺเสว สโต เอตทโหสิ “โก นุ โข รูปานํ อสฺสาโท, โก อาทีนโว, กิํ นิสฺสรณํ, โก สทฺทานํ ฯเปฯ.\csromanspacer{} โก คนฺธานํ.\csromanspacer{} โก รสานํ.\csromanspacer{} โก โผฏฺฐพฺพานํ.\csromanspacer{} โก ธมฺมานํ อสฺสาโท, โก อาทีนโว, กิํ นิสฺสรณนฺ”ติ.\csromanspacer{} ตสฺส มยฺหํ ภิกฺขเว เอตทโหสิ ยํ โข รูเป ปฏิจฺจ อุปฺปชฺชติ สุขํ โสมนสฺสํ, อยํ รูปานํ อสฺสาโท.\csromanspacer{} ยํ รูปา อนิจฺจา ทุกฺขา วิปริณามธมฺมา, อยํ รูปานํ อาทีนโว.\csromanspacer{} โย รูเปสุ ฉนฺทราควินโย ฉนฺทราคปฺปหานํ, อิทํ รูปานํ นิสฺสรณํ.\csromanspacer{} ยํ สทฺเท.\csromanspacer{} คนฺเธ.\csromanspacer{} รเส.\csromanspacer{} โผฏฺฐพฺเพ.\csromanspacer{} ยํ ธมฺเม ปฏิจฺจ อุปฺปชฺชติ สุขํ โสมนสฺสํ, อยํ ธมฺมานํ อสฺสาโท.\csromanspacer{} ยํ ธมฺมา อนิจฺจา ทุกฺขา วิปริณามธมฺมา, อยํ ธมฺมานํ อาทีนโว.\csromanspacer{} โย ธมฺเมสุ ฉนฺทราควินโย ฉนฺทราคปฺปหานํ, อิทํ ธมฺมานํ นิสฺสรณนฺติ.\csromanspacer{} ยาวกีวญฺจาหํ ภิกฺขเว อิเมสํ ฉนฺนํ พาหิรานํ อายตนานํ เอวํ อสฺสาทญฺจ อสฺสาทโต อาทีนวญฺจ อาทีนวโต นิสฺสรณญฺจ นิสฺสรณโต ยถาภูตํ นาพฺภญฺญาสิํ, เนว ตาวาหํ ภิกฺขเว สเทวเก โลเก สมารเก สพฺรหฺมเก สสฺสมณพฺราหฺมณิยา ปชาย สเทวมนุสฺสาย อนุตฺตรํ สมฺมาสมฺโพธิํ อภิสมฺพุทฺโธติ ปจฺจญฺญาสิํ.\csromanspacer{} ยโต จ ขฺวาหํ ภิกฺขเว อิเมสํ ฉนฺนํ พาหิรานํ อายตนานํ เอวํ อสฺสาทญฺจ อสฺสาทโต อาทีนวญฺจ อาทีนวโต นิสฺสรณญฺจ นิสฺสรณโต ยถาภูตํ อพฺภญฺญาสิํ, อถาหํ ภิกฺขเว สเทวเก โลเก สมารเก สพฺรหฺมเก สสฺสมณพฺราหฺมณิยา ปชาย สเทวมนุสฺสาย อนุตฺตรํ สมฺมาสมฺโพธิํ อภิสมฺพุทฺโธติ ปจฺจญฺญาสิํ.\csromanspacer{} ญาณญฺจ ปน เม ทสฺสนํ อุทปาทิ, อกุปฺปา เม วิมุตฺติ, อยมนฺติมา ชาติ, นตฺถิ ทานิ ปุนพฺภโวติ.\csromanspacer{}\csromanspacer{}ทุติยํ.}

\csromanmatikaanchor{mtk.88}
\csromantocmarkref{subsection}{3-4. อสฺสาทปริเยสนสุตฺต}{mtk.88}
\bookmark[dest={mtk.88},level=2]{3-4. อสฺสาทปริเยสนสุตฺต}
\csromanheader{2}{3. ปฐม-อสฺสาทปริเยสนสุตฺต}
\proseitem{15}{จกฺขุสฺสาหํ ภิกฺขเว อสฺสาทปริเยสนํ อจริํ, โย จกฺขุสฺส อสฺสาโท, ตทชฺฌคมํ, ยาวตา จกฺขุสฺส อสฺสาโท, ปญฺญาย เม โส สุทิฏฺโฐ.\csromanspacer{} จกฺขุสฺสาหํ ภิกฺขเว อาทีนวปริเยสนํ อจริํ, โย จกฺขุสฺส อาทีนโว, ตทชฺฌคมํ, ยาวตา จกฺขุสฺส อาทีนโว, ปญฺญาย เม โส สุทิฏฺโฐ.\csromanspacer{} จกฺขุสฺสาหํ ภิกฺขเว นิสฺสรณปริเยสนํ อจริํ, ยํ จกฺขุสฺส นิสฺสรณํ, ตทชฺฌคมํ, ยาวตา จกฺขุสฺส นิสฺสรณํ, ปญฺญาย เม ตํ สุทิฏฺฐํ.}

\vspace{\tipitakaparskip}
\csromancenter{สฬายตนสํยุตฺต โสตสฺสาหํ ภิกฺขเว.\csromanspacer{} ฆานสฺสาหํ ภิกฺขเว.\csromanspacer{} ชิวฺหายาหํ ภิกฺขเว อสฺสาทปริเยสนํ อจริํ, โย ชิวฺหาย อสฺสาโท, ตทชฺฌคมํ, ยาวตา ชิวฺหาย อสฺสาโท, ปญฺญาย เม โส สุทิฏฺโฐ.\csromanspacer{} ชิวฺหายาหํ ภิกฺขเว อาทีนวปริเยสนํ อจริํ, โย ชิวฺหาย อาทีนโว, ตทชฺฌคมํ, ยาวตา ชิวฺหาย อาทีนโว, ปญฺญาย เม โส สุทิฏฺโฐ.\csromanspacer{} ชิวฺหายาหํ ภิกฺขเว นิสฺสรณปริเยสนํ อจริํ, ยํ ชิวฺหาย นิสฺสรณํ, ตทชฺฌคมํ, ยาวตา ชิวฺหาย นิสฺสรณํ, ปญฺญาย เม ตํ สุทิฏฺฐํ ฯเปฯ.\csromanspacer{} มนสฺสาหํ ภิกฺขเว อสฺสาทปริเยสนํ อจริํ, โย มนสฺส อสฺสาโท, ตทชฺฌคมํ, ยาวตา มนสฺส อสฺสาโท, ปญฺญาย เม โส สุทิฏฺโฐ.\csromanspacer{} มนสฺสาหํ ภิกฺขเว อาทีนวปริเยสนํ อจริํ, โย มนสฺส อาทีนโว, ตทชฺฌคมํ, ยาวตา มนสฺส อาทีนโว, ปญฺญาย เม โส สุทิฏฺโฐ.\csromanspacer{} มนสฺสาหํ ภิกฺขเว นิสฺสรณปริเยสนํ อจริํ, ยํ มนสฺส นิสฺสรณํ, ตทชฺฌคมํ, ยาวตา มนสฺส นิสฺสรณํ, ปญฺญาย เม ตํ สุทิฏฺฐํ.\csromanspacer{} ยาวกีวญฺจาหํ ภิกฺขเว อิเมสํ ฉนฺนํ อชฺฌตฺติกานํ อายตนานํ อสฺสาทญฺจ อสฺสาทโต อาทีนวญฺจ อาทีนวโต นิสฺสรณญฺจ นิสฺสรณโต ยถาภูตํ นาพฺภญฺญาสิํ ฯเปฯ ปจฺจญฺญาสิํ.\csromanspacer{} ญาณญฺจ ปน เม ทสฺสนํ อุทปาทิ, อกุปฺปา เม วิมุตฺติ, อยมนฺติมา ชาติ, นตฺถิ ทานิ ปุนพฺภโวติ.\csromanspacer{}\csromanspacer{}ตติยํ.}
\csromanheader{2}{4. ทุติย-อสฺสาทปริเยสนสุตฺต}
\proseitem{16}{\csromanfolio{243}รูปานาหํ ภิกฺขเว อสฺสาทปริเยสนํ อจริํ, โย รูปานํ อสฺสาโท, ตทชฺฌคมํ, ยาวตา รูปานํ อสฺสาโท, ปญฺญาย เม โส สุทิฏฺโฐ.\csromanspacer{} รูปานาหํ ภิกฺขเว อาทีนวปริเยสนํ อจริํ, โย รูปานํ อาทีนโว, ตทชฺฌคมํ, ยาวตา รูปานํ อาทีนโว, ปญฺญาย เม โส สุทิฏฺโฐ.\csromanspacer{} รูปานาหํ ภิกฺขเว นิสฺสรณปริเยสนํ อจริํ, ยํ รูปานํ นิสฺสรณํ, ตทชฺฌคมํ, ยาวตา รูปานํ นิสฺสรณํ, ปญฺญาย เม ตํ สุทิฏฺฐํ.\csromanspacer{} สทฺทานาหํ ภิกฺขเว.\csromanspacer{} คนฺธานาหํ ภิกฺขเว.\csromanspacer{} รสานาหํ ภิกฺขเว.\csromanspacer{} โผฏฺฐพฺพานาหํ ภิกฺขเว.\csromanspacer{} ธมฺมานาหํ ภิกฺขเว อสฺสาทปริเยสนํ อจริํ, โย ธมฺมานํ อสฺสาโท, ตทชฺฌคมํ, ยวตา ธมฺมานํ อสฺสาโท, ปญฺญาย เม โส สุทิฏฺโฐ.\csromanspacer{} ธมฺมานาหํ ภิกฺขเว อาทีนวปริเยสนํ อจริํ, โย ธมฺมานํ อาทีนโว, ตทชฺฌคมํ, ยาวตา ธมฺมานํ อาทีนโว, ปญฺญาย เม โส สุทิฏฺโฐ.\csromanspacer{} ธมฺมานาหํ ภิกฺขเว นิสฺสรณปริเยสนํ อจริํ, ยํ ธมฺมานํ นิสฺสรณํ, ตทชฺฌคมํ, ยาวตา ธมฺมานํ นิสฺสรณํ,}

\gambhira{สฬายตนวคฺคสํยุตฺตปาฬิ}
\prose{\csromanfolio{244}ปญฺญาย เม ตํ สุทิฏฺฐํ.\csromanspacer{} ยาวกีวญฺจาหํ ภิกฺขเว อิเมสํ ฉนฺนํ พาหิรานํ อายตนานํ อสฺสาทญฺจ อสฺสาทโต อาทีนวญฺจ อาทีนวโต นิสฺสรณญฺจ นิสฺสรณโต ยถาภูตํ นาพฺภญฺญาสิํ ฯเปฯ ปจฺจญฺญาสิํ.\csromanspacer{} ญาณญฺจ ปน เม ทสฺสนํ อุทปาทิ, อกุปฺปา เม วิมุตฺติ, อยมนฺติมา ชาติ, นตฺถิ ทานิ ปุนพฺภโวติ.\csromanspacer{}\csromanspacer{}จตุตฺถํ.}

\csromanmatikaanchor{mtk.89}
\csromantocmarkref{subsection}{5-6. โนเจ-อสฺสาทสุตฺต}{mtk.89}
\bookmark[dest={mtk.89},level=2]{5-6. โนเจ-อสฺสาทสุตฺต}
\csromanheader{2}{5. ปฐมโนเจ-อสฺสาทสุตฺต}
\proseitem{17}{โน เจทํ ภิกฺขเว จกฺขุสฺส อสฺสาโท อภวิสฺส, น ยิทํ สตฺตา จกฺขุสฺมิํ สารชฺเชยฺยุํ, ยสฺมา จ โข ภิกฺขเว อตฺถิ จกฺขุสฺส อสฺสาโท, ตสฺมา สตฺตา จกฺขุสฺมิํ สารชฺชนฺติ.\csromanspacer{} โน เจทํ ภิกฺขเว จกฺขุสฺส อาทีนโว อภวิสฺส, น ยิทํ สตฺตา จกฺขุสฺมิํ นิพฺพินฺเทยฺยุํ, ยสฺมา จ โข ภิกฺขเว อตฺถิ จกฺขุสฺส อาทีนโว, ตสฺมา สตฺตา จกฺขุสฺมิํ นิพฺพินฺทนฺติ.\csromanspacer{} โน เจทํ ภิกฺขเว จกฺขุสฺส นิสฺสรณํ อภวิสฺส, น ยิทํ สตฺตา จกฺขุสฺมา นิสฺสเรยฺยุํ, ยสฺมา จ โข ภิกฺขเว อตฺถิ จกฺขุสฺส นิสฺสรณํ, ตสฺมา สตฺตา จกฺขุสฺมา นิสฺสรนฺติ.\csromanspacer{} โน เจทํ ภิกฺขเว โสตสฺส อสฺสาโท อภวิสฺส.\csromanspacer{} โน เจทํ ภิกฺขเว ฆานสฺส อสฺสาโท อภวิสฺส.\csromanspacer{} โน เจทํ ภิกฺขเว ชิวฺหาย อสฺสาโท อภวิสฺส, น ยิทํ สตฺตา ชิวฺหาย สารชฺเชยฺยุํ, ยสฺมา จ โข ภิกฺขเว อตฺถิ ชิวฺหาย อสฺสาโท, ตสฺมา สตฺตา ชิวฺหาย สารชฺชนฺติ.\csromanspacer{} โน เจทํ ภิกฺขเว ชิวฺหาย อาทีนโว อภวิสฺส, น ยิทํ สตฺตา ชิวฺหาย นิพฺพินฺเทยฺยุํ, ยสฺมา จ โข ภิกฺขเว อตฺถิ ชิวฺหาย อาทีนโว, ตสฺมา สตฺตา ชิวฺหาย นิพฺพินฺทนฺติ.\csromanspacer{} โน เจทํ ภิกฺขเว ชิวฺหาย นิสฺสรณํ อภวิสฺส, น ยิทํ สตฺตา ชิวฺหาย นิสฺสเรยฺยุํ, ยสฺมา จ โข ภิกฺขเว อตฺถิ ชิวฺหาย นิสฺสรณํ, ตสฺมา สตฺตา ชิวฺหาย นิสฺสรนฺติ.\csromanspacer{} โน เจทํ ภิกฺขเว กายสฺส อสฺสาโท อภวิสฺส.\csromanspacer{} โน เจทํ ภิกฺขเว มนสฺส อสฺสาโท อภวิสฺส, น ยิทํ สตฺตา มนสฺมิํ สารชฺเชยฺยุํ, ยสฺมา จ โข ภิกฺขเว อตฺถิ มนสฺส อสฺสาโท, ตสฺมา สตฺตา มนสฺมิํ สารชฺชนฺติ.\csromanspacer{} โน เจทํ ภิกฺขเว มนสฺส อาทีนโว อภวิสฺส, น ยิทํ สตฺตา มนสฺมิํ นิพฺพินฺเทยฺยุํ, ยสฺมา จ โข ภิกฺขเว อตฺถิ มนสฺส อาทีนโว, ตสฺมา สตฺตา มนสฺมิํ นิพฺพินฺทนฺติ.\csromanspacer{} โน เจทํ ภิกฺขเว มนสฺส นิสฺสรณํ อภวิสฺส, น ยิทํ สตฺตา มนสฺมา นิสฺสเรยฺยุํ, ยสฺมา จ โข ภิกฺขเว อตฺถิ มนสฺส นิสฺสรณํ, ตสฺมา สตฺตา มนสฺมา นิสฺสรนฺติ.\csromanspacer{} ยาวกีวญฺจ ภิกฺขเว สตฺตา อิเมสํ ฉนฺนํ อชฺฌตฺติกานํ}

\vspace{\tipitakaparskip}
\csromancenter{สฬายตนสํยุตฺต อายตนานํ อสฺสาทญฺจ อสฺสาทโต อาทีนวญฺจ อาทีนวโต นิสฺสรณญฺจ นิสฺสรณโต ยถาภูตํ นาพฺภญฺญํสุ, เนว ตาว ภิกฺขเว สตฺตา สเทวกา โลกา สมารกา สพฺรหฺมกา สสฺสมณพฺราหฺมณิยา ปชาย สเทวมนุสฺสาย นิสฺสฏา วิสญฺญุตฺตา วิปฺปมุตฺตา วิมริยาทีกเตน\footnote{วิมริยาทิกเตน (สี, อิ), วิปริยาทิกเตน (สฺยา, กํ, ก)} เจตสา วิหริํสุ.\csromanspacer{} ยโต จ โข ภิกฺขเว สตฺตา อิเมสํ ฉนฺนํ อชฺฌตฺติกานํ อายตนานํ อสฺสาทญฺจ อสฺสาทโต อาทีนวญฺจ อาทีนวโต นิสฺสรณญฺจ นิสฺสรณโต ยถาภูตํ อพฺภญฺญํสุ, อถ ภิกฺขเว สตฺตา สเทวกา โลกา สมารกา สพฺรหฺมกา สสฺสมณพฺราหฺมณิยา ปชาย สเทวมนุสฺสาย นิสฺสฏา วิสญฺญุตฺตา วิปฺปมุตฺตา วิมริยาทีกเตน เจตสา วิหรนฺตีติ.\csromanspacer{}\csromanspacer{}ปญฺจมํ.}
\csromanheader{2}{6. ทุติยโนเจ-อสฺสาทสุตฺต}
\proseitem{18}{\csromanfolio{245}โน เจทํ ภิกฺขเว รูปานํ อสฺสาโท อภวิสฺส, น ยิทํ สตฺตา รูเปสุ สารชฺเชยฺยุํ, ยสฺมา จ โข ภิกฺขเว อตฺถิ รูปานํ อสฺสาโท, ตสฺมา สตฺตา รูเปสุ สารชฺชนฺติ.\csromanspacer{} โน เจทํ ภิกฺขเว รูปานํ อาทีนโว อภวิสฺส, น ยิทํ สตฺตา รูเปสุ นิพฺพินฺเทยฺยุํ, ยสฺมา จ โข ภิกฺขเว อตฺถิ รูปานํ อาทีนโว, ตสฺมา สตฺตา รูเปสุ นิพฺพินฺทนฺติ.\csromanspacer{} โน เจทํ ภิกฺขเว รูปานํ นิสฺสรณํ อภวิสฺส, น ยิทํ สตฺตา รูเปหิ นิสฺสเรยฺยุํ, ยสฺมา จ โข ภิกฺขเว อตฺถิ รูปานํ นิสฺสรณํ, ตสฺมา สตฺตา รูเปหิ นิสฺสรนฺติ.\csromanspacer{} โน เจทํ ภิกฺขเว สทฺทานํ.\csromanspacer{} คนฺธานํ.\csromanspacer{} รสานํ.\csromanspacer{} โผฏฺฐพฺพานํ.\csromanspacer{} ธมฺมานํ อสฺสาโท อภวิสฺส, น ยิทํ สตฺตา ธมฺเมสุ สารชฺเชยฺยุํ, ยสฺมา จ โข ภิกฺขเว อตฺถิ ธมฺมานํ อสฺสาโท, ตสฺมา สตฺตา ธมฺเมสุ สารชฺชนฺติ.\csromanspacer{} โน เจทํ ภิกฺขเว ธมฺมานํ อาทีนโว อภวิสฺส, น ยิทํ สตฺตา ธมฺเมสุ นิพฺพินฺเทยฺยุํ, ยสฺมา จ โข ภิกฺขเว อตฺถิ ธมฺมานํ อาทีนโว, ตสฺมา สตฺตา ธมฺเมสุ นิพฺพินฺทนฺติ.\csromanspacer{} โน เจทํ ภิกฺขเว ธมฺมานํ นิสฺสรณํ อภวิสฺส, น ยิทํ สตฺตา ธมฺเมหิ นิสฺสเรยฺยุํ, ยสฺมา จ โข ภิกฺขเว อตฺถิ ธมฺมานํ นิสฺสรณํ, ตสฺมา สตฺตา ธมฺเมหิ นิสฺสรนฺติ.\csromanspacer{} ยาวกีวญฺจ ภิกฺขเว สตฺตา อิเมสํ ฉนฺนํ พาหิรานํ อายตนานํ อสฺสาทญฺจ อสฺสาทโต อาทีนวญฺจ อาทีนวโต นิสฺสรณญฺจ นิสฺสรณโต ยถาภูตํ นาพฺภญฺญํสุ, เนว ตาว ภิกฺขเว สตฺตา สเทวกา โลกา สมารกา สพฺรหฺมกา}

\gambhira{สฬายตนวคฺคสํยุตฺตปาฬิ}
\csromanmatikaanchor{mtk.90}
\csromantocmarkref{subsection}{7-8. อภินนฺทสุตฺต}{mtk.90}
\bookmark[dest={mtk.90},level=2]{7-8. อภินนฺทสุตฺต}
\prose{\csromanfolio{246}สสฺสมณพฺราหฺมณิยา ปชาย สเทวมนุสฺสาย นิสฺสฏา วิสญฺญุตฺตา วิปฺปมุตฺตา วิมริยาทีกเตน เจตสา วิหริํสุ.\csromanspacer{} ยโต จ โข ภิกฺขเว สตฺตา อิเมสํ ฉนฺนํ พาหิรานํ อายตนานํ อสฺสาทญฺจ อสฺสาทโต อาทีนวญฺจ อาทีนวโต นิสฺสรณญฺจ นิสฺสรณโต ยถาภูตํ อพฺภญฺญํสุ, อถ ภิกฺขเว สตฺตา สเทวกา โลกา สมารกา สพฺรหฺมกา สสฺสมณพฺราหฺมณิยา ปชาย สเทวมนุสฺสาย นิสฺสฏา วิสญฺญุตฺตา วิปฺปมุตฺตา วิมริยาทีกเตน เจตสา วิหรนฺตีติ.\csromanspacer{}\csromanspacer{}ฉฏฺฐํ.}

\csromanheader{2}{7. ปฐมาภินนฺทสุตฺต}
\proseitem{19}{โย ภิกฺขเว จกฺขุํ อภินนฺทติ, ทุกฺขํ โส อภินนฺทติ.\csromanspacer{} โย ทุกฺขํ อภินนฺทติ, อปริมุตฺโต โส ทุกฺขสฺมาติ วทามิ.\csromanspacer{} โย โสตํ ฯเปฯ.\csromanspacer{} โย ฆานํ ฯเปฯ.\csromanspacer{} โย ชิวฺหํ อภินนฺทติ, ทุกฺขํ โส อภินนฺทติ.\csromanspacer{} โย ทุกฺขํ อภินนฺทติ, อปริมุตฺโต โส ทุกฺขสฺมาติ วทามิ.\csromanspacer{} โย กายํ ฯเปฯ.\csromanspacer{} โย มนํ อภินนฺทติ, ทุกฺขํ โส อภินนฺทติ.\csromanspacer{} โย ทุกฺขํ อภินนฺทติ, อปริมุตฺโต โส ทุกฺขสฺมาติ วทามิ.}

\prose{โย จ โข ภิกฺขเว จกฺขุํ นาภินนฺทติ, ทุกฺขํ โส นาภินนฺทติ.\csromanspacer{} โย ทุกฺขํ นาภินนฺทติ, ปริมุตฺโต โส ทุกฺขสฺมาติ วทามิ.\csromanspacer{} โย โสตํ ฯเปฯ.\csromanspacer{} โย ฆานํ ฯเปฯ.\csromanspacer{} โย ชิวฺหํ นาภินนฺทติ, ทุกฺขํ โส นาภินนฺทติ.\csromanspacer{} โย ทุกฺขํ นาภินนฺทติ, ปริมุตฺโต โส ทุกฺขสฺมาติ วทามิ.\csromanspacer{} โย กายํ ฯเปฯ.\csromanspacer{} โย มนํ นาภินนฺทติ, ทุกฺขํ โส นาภินนฺทติ.\csromanspacer{} โย ทุกฺขํ นาภินนฺทติ, ปริมุตฺโต โส ทุกฺขสฺมาติ วทามิ.\csromanspacer{}\csromanspacer{}สตฺตมํ.}

\csromanheader{2}{8. ทุติยาภินนฺทสุตฺต}
\proseitem{20}{โย ภิกฺขเว รูเป อภินนฺทติ, ทุกฺขํ โส อภินนฺทติ.\csromanspacer{} โย ทุกฺขํ อภินนฺทติ, อปริมุตฺโต โส ทุกฺขสฺมาติ วทามิ.\csromanspacer{} โย สทฺเท ฯเปฯ.\csromanspacer{} คนฺเธ.\csromanspacer{} รเส.\csromanspacer{} โผฏฺฐพฺเพ.\csromanspacer{} ธมฺเม อภินนฺทติ, ทุกฺขํ โส อภินนฺทติ.\csromanspacer{} โย ทุกฺขํ อภินนฺทติ, อปริมุตฺโต โส ทุกฺขสฺมาติ วทามิ.}

\prose{โย จ โข ภิกฺขเว รูเป นาภินนฺทติ, ทุกฺขํ โส นาภินนฺทติ.\csromanspacer{} โย ทุกฺขํ นาภินนฺทติ, ปริมุตฺโต โส ทุกฺขสฺมาติ วทามิ.\csromanspacer{} โย สทฺเท ฯเปฯ.\csromanspacer{} คนฺเธ.\csromanspacer{} รเส.\csromanspacer{} โผฏฺฐพฺเพ.\csromanspacer{} ธมฺเม นาภินนฺทติ, ทุกฺขํ โส นาภินนฺทติ.\csromanspacer{} โย ทุกฺขํ นาภินนฺทติ, ปริมุตฺโต โส ทุกฺขสฺมาติ วทามิ.\csromanspacer{}\csromanspacer{}อฏฺฐมํ.}

\csromanheader{2}{1. สฬายตนสํยุตฺต \textbf{9. ปฐมทุกฺขุปฺปาทสุตฺต}}
\csromanmatikaanchor{mtk.91}
\csromantocmarkref{subsection}{9-10. ทุกฺขุปฺปาทสุตฺต}{mtk.91}
\bookmark[dest={mtk.91},level=2]{9-10. ทุกฺขุปฺปาทสุตฺต}
\proseitem{21}{\csromanfolio{247}โย ภิกฺขเว จกฺขุสฺส อุปฺปาโท ฐิติ อภินิพฺพตฺติ ปาตุภาโว, ทุกฺขสฺเสโส อุปฺปาโท, โรคานํ ฐิติ, ชรามรณสฺส ปาตุภาโว.\csromanspacer{} โย โสตสฺส ฯเปฯ.\csromanspacer{} โย ฆานสฺส.\csromanspacer{} โย ชิวฺหาย.\csromanspacer{} โย กายสฺส.\csromanspacer{} โย มนสฺส อุปฺปาโท ฐิติ อภินิพฺพตฺติ ปาตุภาโว, ทุกฺขสฺเสโส อุปฺปาโท, โรคานํ ฐิติ, ชรามรณสฺส ปาตุภาโว.}

\prose{โย จ โข ภิกฺขเว จกฺขุสฺส นิโรโธ วูปสโม อตฺถงฺคโม, ทุกฺขสฺเสโส นิโรโธ, โรคานํ วูปสโม, ชรามณสฺส อตฺถงฺคโม.\csromanspacer{} โย โสตสฺส.\csromanspacer{} โย ฆานสฺส.\csromanspacer{} โย ชิวฺหาย.\csromanspacer{} โย กายสฺส.\csromanspacer{} โย มนสฺส นิโรโธ วูปสโม อตฺถงฺคโม, ทุกฺขสฺเสโส นิโรโธ, โรคานํ วูปสโม, ชรามรณสฺส อตฺถงฺคโมติ.\csromanspacer{}\csromanspacer{}นวมํ.}

\csromanheader{2}{10. ทุติยทุกฺขุปฺปาทสุตฺต}
\proseitem{22}{โย ภิกฺขเว รูปานํ อุปฺปาโท ฐิติ อภินิพฺพตฺติ ปาตุภาโว, ทุกฺขสฺเสโส อุปฺปาโท, โรคานํ ฐิติ, ชรามรณสฺส ปาตุภาโว.\csromanspacer{} โย สทฺทานํ ฯเปฯ.\csromanspacer{} โย คนฺธานํ.\csromanspacer{} โย รสานํ.\csromanspacer{} โย โผฏฺฐพฺพานํ.\csromanspacer{} โย ธมฺมานํ อุปฺปาโท ฐิติ อภินิพฺพตฺติ ปาตุภาโว, ทุกฺขสฺเสโส อุปฺปาโท, โรคานํ ฐิติ, ชรามรณสฺส ปาตุภาโว.}

\prose{โย จ โข ภิกฺขเว รูปานํ นิโรโธ วูปสโม อตฺถงฺคโม, ทุกฺขสฺเสโส นิโรโธ, โรคานํ วูปสโม, ชรามรณสฺส อตฺถงฺคโม.\csromanspacer{} โย สทฺทานํ ฯเปฯ.\csromanspacer{} โย คนฺธานํ.\csromanspacer{} โย รสานํ.\csromanspacer{} โย โผฏฺฐพฺพานํ.\csromanspacer{} โย ธมฺมานํ นิโรโธ วูปสโม อตฺถงฺคโม, ทุกฺขสฺเสโส นิโรโธ, โรคานํ วุปสโม, ชรามรณสฺส อตฺถงฺคโมติ.\csromanspacer{}\csromanspacer{}ทสมํ.\csromanspacer{} ยมกวคฺโค ทุติโย.}

\titlehead{\textbf{ตสฺสุทฺทานํ}}
\csromangathameasure{สมฺโพเธน ทุเว วุตฺตา, อสฺสาเทน อปเร ทุเว. \\ โน เจเตน ทุเว วุตฺตา, อภินนฺเทน อปเร ทุเว. \\ อุปฺปาเทน ทุเว วุตฺตา, วคฺโค เตน ปวุจฺจตีติ.}
\csromangathagroup{\csromangathastanza{สมฺโพเธน ทุเว วุตฺตา, อสฺสาเทน อปเร ทุเว. \\ โน เจเตน ทุเว วุตฺตา, อภินนฺเทน อปเร ทุเว. \\ อุปฺปาเทน ทุเว วุตฺตา, วคฺโค เตน ปวุจฺจตีติ.}}

\gambhira{สฬายตนวคฺคสํยุตฺตปาฬิ}
\csromanmatikaanchor{mtk.92}
\csromantocmarknopage{section}{3. สพฺพวคฺค}{mtk.92}
\bookmark[dest={mtk.92},level=1]{3. สพฺพวคฺค}
\csromanheader{1}{3. สพฺพวคฺค}
\csromanmatikaanchor{mtk.93}
\csromantocmarkref{subsection}{1. สพฺพสุตฺต}{mtk.93}
\bookmark[dest={mtk.93},level=2]{1. สพฺพสุตฺต}
\csromanheader{2}{1. สพฺพสุตฺต}
\proseitem{23}{\csromanfolio{248}สาวตฺถินิทานํ.\csromanspacer{} สพฺพํ โว ภิกฺขเว เทเสสฺสามิ, ตํ สุณาถ.\csromanspacer{} กิญฺจ ภิกฺขเว สพฺพํ.\csromanspacer{} จกฺขุํ เจว รูปา จ โสตญฺจ\footnote{โสตญฺเจว (?) เอวมิตรยุคเลสุปิ.} สทฺทา จ ฆานญฺจ คนฺธา จ ชิวฺหา จ รสา จ กาโย จ โผฏฺฐพฺพา จ มโน จ ธมฺมา จ, อิทํ วุจฺจติ ภิกฺขเว สพฺพํ.\csromanspacer{} โย ภิกฺขเว เอวํ วเทยฺย “อหเมตํ สพฺพํ ปจฺจกฺขาย อญฺญํ สพฺพํ ปญฺญาเปสฺสามี”ติ, ตสฺส วาจาวตฺถุกเมวสฺส\footnote{วาจาวตฺถุเรวสฺส (สี, อิ), วาจาวตฺถุเทวสฺส (สฺยา, กํ)}, ปุฏฺโฐ จ น สมฺปาเยยฺย, อุตฺตริํ จ วิฆาตํ อาปชฺเชยฺย.\csromanspacer{} ตํ กิสฺส เหตุ, ยถา ตํ ภิกฺขเว อวิสยสฺมินฺติ.\csromanspacer{}\csromanspacer{}ปฐมํ.}

\csromanmatikaanchor{mtk.94}
\csromantocmarkref{subsection}{2. ปหานสุตฺต}{mtk.94}
\bookmark[dest={mtk.94},level=2]{2. ปหานสุตฺต}
\csromanheader{2}{2. ปหานสุตฺต}
\proseitem{24}{สพฺพปฺปหานาย\footnote{สพฺพํ ปหานาย (สฺยา, กํ, ก)} โว ภิกฺขเว ธมฺมํ เทเสสฺสามิ, ตํ สุณาถ.\csromanspacer{} กตโม จ ภิกฺขเว สพฺพปฺปหานาย ธมฺโม.\csromanspacer{} จกฺขุํ ภิกฺขเว ปหาตพฺพํ, รูปา ปหาตพฺพา, จกฺขุวิญฺญาณํ ปหาตพฺพํ, จกฺขุสมฺผสฺโส ปหาตพฺโพ, ยมฺปิทํ จกฺขุสมฺผสฺสปจฺจยา อุปฺปชฺชติ เวทยิตํ สุขํ วา ทุกฺขํ วา อทุกฺขมสุขํ วา, ตมฺปิ ปหาตพฺพํ ฯเปฯ ชิวฺหา ปหาตพฺพา, รสา ปหาตพฺพา, ชิวฺหาวิญฺญาณํ ปหาตพฺพํ, ชิวฺหาสมฺผสฺโส ปหาตพฺโพ, ยมฺปิทํ ชิวฺหาสมฺผสฺสปจฺจยา อุปฺปชฺชติ เวทยิตํ สุขํ วา ทุกฺขํ วา อทุกฺขมสุขํ วา, ตมฺปิ ปหาตพฺพํ.\csromanspacer{} กาโย ปหาตพฺโพ.\csromanspacer{} มโน ปหาตพฺโพ, ธมฺมา ปหาตพฺพา, มโนวิญฺญาณํ ปหาตพฺพํ, มโนสมฺผสฺโส ปหาตพฺโพ, ยมฺปิทํ มโนสมฺผสฺสปจฺจยา อุปฺปชฺชติ เวทยิตํ สุขํ วา ทุกฺขํ วา อทุกฺขมสุขํ วา, ตมฺปิ ปหาตพฺพํ.\csromanspacer{} อยํ โข ภิกฺขเว สพฺพปฺปหานาย ธมฺโมติ.\csromanspacer{}\csromanspacer{}ทุติยํ.}

\csromanmatikaanchor{mtk.95}
\csromantocmarkref{subsection}{3. อภิญฺญาปริญฺญาปหานสุตฺต}{mtk.95}
\bookmark[dest={mtk.95},level=2]{3. อภิญฺญาปริญฺญาปหานสุตฺต}
\csromanheader{2}{3. อภิญฺญาปริญฺญาปหานสุตฺต}
\proseitem{25}{สพฺพํ อภิญฺญา ปริญฺญา ปหานาย โว ภิกฺขเว ธมฺมํ เทเสสฺสามิ, ตํ สุณาถ.\csromanspacer{} กตโม จ ภิกฺขเว สพฺพํ อภิญฺญา ปริญฺญา ปหานาย ธมฺโม.\csromanspacer{} จกฺขุํ ภิกฺขเว อภิญฺญา ปริญฺญา ปหาตพฺพํ, รูปา อภิญฺญา ปริญฺญา}

\vspace{\tipitakaparskip}
\csromancenter{สฬายตนสํยุตฺต ปหาตพฺพา, จกฺขุวิญฺญาณํ อภิญฺญา ปริญฺญา ปหาตพฺพํ, จกฺขุสมฺผสฺโส อภิญฺญา ปริญฺญา ปหาตพฺโพ, ยมฺปิทํ จกฺขุสมฺผสฺสปจฺจยา อุปฺปชฺชติ เวทยิตํ สุขํ วา ทุกฺขํ วา อทุกฺขมสุขํ วา, ตมฺปิ อภิญฺญา ปริญฺญา ปหาตพฺพํ ฯเปฯ.\csromanspacer{} ชิวฺหา อภิญฺญา ปริญฺญา ปหาตพฺพา, รสา อภิญฺญา ปริญฺญา ปหาตพฺพา, ชิวฺหาวิญฺญาณํ อภิญฺญา ปริญฺญา ปหาตพฺพํ, ชิวฺหาสมฺผสฺโส อภิญฺญา ปริญฺญา ปหาตพฺโพ, ยมฺปิทํ ชิวฺหาสมฺผสฺสปจฺจยา อุปฺปชฺชติ เวทยิตํ สุขํ วา ทุกฺขํ วา อทุกฺขมสุขํ วา, ตมฺปิ อภิญฺญา ปริญฺญา ปหาตพฺพํ.\csromanspacer{} กาโย อภิญฺญา ปริญฺญา ปหาตพฺโพ.\csromanspacer{} มโน อภิญฺญา ปริญฺญา ปหาตพฺโพ, ธมฺมา อภิญฺญา ปริญฺญา ปหาตพฺพา, มโนวิญฺญาณํ อภิญฺญา ปริญฺญา ปหาตพฺพํ, มโนสมฺผสฺโส อภิญฺญา ปริญฺญา ปหาตพฺโพ, ยมฺปิทํ มโนสมฺผสฺสปจฺจยา อุปฺปชฺชติ เวทยิตํ สุขํ วา ทุกฺขํ วา อทุกฺขมสุขํ วา, ตมฺปิ อภิญฺญา ปริญฺญา ปหาตพฺพํ.\csromanspacer{} อยํ โข ภิกฺขเว สพฺพํ อภิญฺญา ปริญฺญา ปหานาย ธมฺโมติ.\csromanspacer{}\csromanspacer{}ตติยํ.}
\csromanmatikaanchor{mtk.96}
\csromantocmarkref{subsection}{4-5. อปริชานนสุตฺต}{mtk.96}
\bookmark[dest={mtk.96},level=2]{4-5. อปริชานนสุตฺต}
\csromanheader{2}{4. ปฐม-อปริชานนสุตฺต}
\proseitem{26}{\csromanfolio{249}สพฺพํ ภิกฺขเว อนภิชานํ อปริชานํ อวิราชยํ อปฺปชหํ อภพฺโพ ทุกฺขกฺขยาย.\csromanspacer{} กิญฺจ ภิกฺขเว อนภิชานํ อปริชานํ อวิราชยํ อปฺปชหํ อภพฺโพ ทุกฺขกฺขยาย.\csromanspacer{} จกฺขุํ ภิกฺขเว อนภิชานํ อปริชานํ อวิราชยํ อปฺปชหํ อภพฺโพ ทุกฺขกฺขยาย.\csromanspacer{} รูเป อนภิชานํ อปริชานํ อวิราชยํ อปฺปชหํ อภพฺโพ ทุกฺขกฺขยาย.\csromanspacer{} จกฺขุวิญฺญาณํ อนภิชานํ อปริชานํ อวิราชยํ อปฺปชหํ อภพฺโพ ทุกฺขกฺขยาย.\csromanspacer{} จกฺขุสมฺผสฺสํ อนภิชานํ อปริชานํ อวิราชยํ อปฺปชหํ อภพฺโพ ทุกฺขกฺขยาย.\csromanspacer{} ยมฺปิทํ จกฺขุสมฺผสฺสปจฺจยา อุปฺปชฺชติ เวทยิตํ สุขํ วา ทุกฺขํ วา อทุกฺขมสุขํ วา, ตมฺปิ อนภิชานํ อปริชานํ อวิราชยํ อปฺปชหํ อภพฺโพ ทุกฺขกฺขยาย ฯเปฯ.\csromanspacer{} ชิวฺหํ อนภิชานํ อปริชานํ อวิราชยํ อปฺปชหํ อภพฺโพ ทุกฺขกฺขยาย.\csromanspacer{} รเส ฯเปฯ.\csromanspacer{} ชิวฺหาวิญฺญาณํ ฯเปฯ ชิวฺหาสมฺผสฺสํ ฯเปฯ.\csromanspacer{} ยมฺปิทํ ชิวฺหาสมฺผสฺสปจฺจยา อุปฺปชฺชติ เวทยิตํ สุขํ วา ทุกฺขํ วา อทุกฺขมสุขํ วา, ตมฺปิ อนภิชานํ อปริชานํ อวิราชยํ อปฺปชหํ อภพฺโพ ทุกฺขกฺขยาย.\csromanspacer{} กายํ ฯเปฯ.\csromanspacer{} มนํ อนภิชานํ อปริชานํ อวิราชยํ อปฺปชหํ อภพฺโพ ทุกฺขกฺขยาย.\csromanspacer{} ธมฺเม ฯเปฯ.\csromanspacer{} มโนวิญฺญาณํ ฯเปฯ มโนสมฺผสฺสํ ฯเปฯ ยมฺปิทํ มโนสมฺผสฺสปจฺจยา อุปฺปชฺชติ เวทยิตํ สุขํ วา ทุกฺขํ วา อทุกฺขมสุขํ วา, ตมฺปิ อนภิชานํ อปริชานํ}

\gambhira{สฬายตนวคฺคสํยุตฺตปาฬิ}
\prose{\csromanfolio{250}อวิราชยํ อปฺปชหํ อภพฺโพ ทุกฺขกฺขยาย.\csromanspacer{} อิทํ โข ภิกฺขเว สพฺพํ อนภิชานํ อปริชานํ อวิราชยํ อปฺปชหํ อภพฺโพ ทุกฺขกฺขยาย.}

\prose{สพฺพํ จ โข ภิกฺขเว อภิชานํ ปริชานํ วิราชยํ ปชหํ ภพฺโพ ทุกฺขกฺขยาย.\csromanspacer{} กิญฺจ ภิกฺขเว สพฺพํ อภิชานํ ปริชานํ วิราชยํ ปชหํ ภพฺโพ ทุกฺขกฺขยาย.\csromanspacer{} จกฺขุํ ภิกฺขเว อภิชานํ ปริชานํ วิราชยํ ปชหํ ภพฺโพ ทุกฺขกฺขยาย.\csromanspacer{} รูเป อภิชานํ ปริชานํ วิราชยํ ปชหํ ภพฺโพ ทุกฺขกฺขยาย.\csromanspacer{} จกฺขุวิญฺญาณํ อภิชานํ ปริชานํ วิราชยํ ปชหํ ภพฺโพ ทุกฺขกฺขยาย.\csromanspacer{} จกฺขุสมฺผสฺสํ อภิชานานํ ปริชานํ วิราชยํ ปชหํ ภพฺโพ ทุกฺขกฺขยาย.\csromanspacer{} ยมฺปิทํ จกฺขุสมฺผสฺสปจฺจยา อุปฺปชฺชติ เวทยิตํ สุขํ วา ทุกฺขํ วา อทุกฺขมสุขํ วา, ตมฺปิ อภิชานํ ปริชานํ วิราชยํ ปชหํ ภพฺโพ ทุกฺขกฺขยาย ฯเปฯ.\csromanspacer{} ชิวฺหํ อภิชานํ ปริชานํ วิราชยํ ปชหํ ภพฺโพ ทุกฺขกฺขยาย.\csromanspacer{} รเส ฯเปฯ.\csromanspacer{} ชิวฺหาวิญฺญาณํ ฯเปฯ.\csromanspacer{} ชิวฺหาสมฺผสฺสํ ฯเปฯ.\csromanspacer{} ยมฺปิทํ ชิวฺหาสมฺผสฺสปจฺจยา อุปฺปชฺชติ เวทยิตํ สุขํ วา ทุกฺขํ วา อทุกฺขมสุขํ วา, ตมฺปิ อภิชานํ ปริชานํ วิราชยํ ปชหํ ภพฺโพ ทุกฺขกฺขยาย.\csromanspacer{} กายํ ฯเปฯ.\csromanspacer{} มนํ อภิชานํ ปริชานํ วิราชยํ ปชหํ ภพฺโพ ทุกฺขกฺขยาย.\csromanspacer{} ธมฺเม ฯเปฯ.\csromanspacer{} มโนวิญฺญาณํ ฯเปฯ.\csromanspacer{} มโนสมฺผสฺสํ ฯเปฯ.\csromanspacer{} ยมฺปิทํ มโนสมฺผสฺสปจฺจยา อุปฺปชฺชติ เวทยิตํ สุขํ วา ทุกฺขํ วา อทุกฺขมสุขํ วา, ตมฺปิ อภิชานํ ปริชานํ วิราชยํ ปชหํ ภพฺโพ ทุกฺขกฺขยาย.\csromanspacer{} อิทํ โข ภิกฺขเว สพฺพํ อภิชานํ ปริชานํ วิราชยํ ปชหํ ภพฺโพ ทุกฺขกฺขยายาติ.\csromanspacer{}\csromanspacer{}จตุตฺถํ.}

\csromanheader{2}{5. ทุติย-อปริชานนสุตฺต}
\proseitem{27}{สพฺพํ ภิกฺขเว อนภิชานํ อปริชานํ อวิราชยํ อปฺปชหํ อภพฺโพ ทุกฺขกฺขยาย.\csromanspacer{} กิญฺจ ภิกฺขเว สพฺพํ อนภิชานํ อปริชานํ อวิราชยํ อปฺปชหํ อภพฺโพ ทุกฺขกฺขยาย.\csromanspacer{} ยญฺจ ภิกฺขเว จกฺขุ เย จ รูปา ยญฺจ จกฺขุวิญฺญาณํ เย จ จกฺขุวิญฺญาณวิญฺญาตพฺพา ธมฺมา ฯเปฯ.\csromanspacer{} ยา จ ชิวฺหา เย จ รสา ยญฺจ ชิวฺหาวิญฺญาณํ เย จ ชิวฺหาวิญฺญาณวิญฺญาตพฺพา ธมฺมา.\csromanspacer{} โย จ กาโย เย จ โผฏฺฐพฺพา ยญฺจ กายวิญฺญาณํ เย จ กายวิญฺญาณวิญฺญาตพฺพา ธมฺมา.\csromanspacer{} โย จ มโน เย จ ธมฺมา ยญฺจ มโนวิญฺญาณํ เย จ มโนวิญฺญาณวิญฺญาตพฺพา ธมฺมา.\csromanspacer{} อิทํ โข ภิกฺขเว สพฺพํ อนภิชานํ อปริชานํ อวิราชยํ อปฺปชหํ อภพฺโพ ทุกฺขกฺขยาย.}

\csromanheader{2}{1. สฬายตนสํยุตฺต}
\prose{\csromanfolio{251}สพฺพํ ภิกฺขเว อภิชานํ ปริชานํ วิราชยํ ปชหํ ภพฺโพ ทุกฺขกฺขยาย.\csromanspacer{} กิญฺจ ภิกฺขเว สพฺพํ อภิชานํ ปริชานํ วิราชยํ ปชหํ ภพฺโพ ทุกฺขกฺขยาย.\csromanspacer{} ยญฺจ ภิกฺขเว จกฺขุ เย จ รูปา ยญฺจ จกฺขุวิญฺญาณํ เย จ จกฺขุวิญฺญาณวิญฺญาตพฺพา ธมฺมา ฯเปฯ.\csromanspacer{} ยา จ ชิวฺหา เย จ รสา ยญฺจ ชิวฺหาวิญฺญาณํ เย จ ชิวฺหาวิญฺญาณวิญฺญาตพฺพา ธมฺมา.\csromanspacer{} โย จ กาโย เย จ โผฏฺฐพฺพา ยญฺจ กายวิญฺญาณํ เย จ กายวิญฺญาณวิญฺญาตพฺพา ธมฺมา.\csromanspacer{} โย จ มโน เย จ ธมฺมา ยญฺจ มโนวิญฺญาณํ เย จ มโนวิญฺญาณวิญฺญาตพฺพา ธมฺมา.\csromanspacer{} อิทํ โข ภิกฺขเว สพฺพํ อภิชานํ ปริชานํ วิราชยํ ปชหํ ภพฺโพ ทุกฺขกฺขยายาติ.\csromanspacer{}\csromanspacer{}ปญฺจมํ.}

\csromanmatikaanchor{mtk.97}
\csromantocmarkref{subsection}{6. อาทิตฺตสุตฺต}{mtk.97}
\bookmark[dest={mtk.97},level=2]{6. อาทิตฺตสุตฺต}
\csromanheader{2}{6. อาทิตฺตสุตฺต}
\proseitem{28}{เอกํ สมยํ ภควา คยายํ วิหรติ คยาสีเส สทฺธิํ ภิกฺขุสหสฺเสน.\csromanspacer{} ตตฺร โข ภควา ภิกฺขู อามนฺเตสิ\csromandash{}สพฺพํ ภิกฺขเว อาทิตฺตํ.\csromanspacer{} กิญฺจ ภิกฺขเว สพฺพํ อาทิตฺตํ.\csromanspacer{} จกฺขุ\footnote{จกฺขุํ (สี, สฺยา, กํ, อิ)} ภิกฺขเว อาทิตฺตํ, รูปา อาทิตฺตา, จกฺขุวิญฺญาณํ อาทิตฺตํ, จกฺขุสมฺผสฺโส อาทิตฺโต, ยมฺปิทํ จกฺขุสมฺผสฺสปจฺจยา อุปฺปชฺชติ เวทยิตํ สุขํ วา ทุกฺขํ วา อทุกฺขมสุขํ วา, ตมฺปิ อาทิตฺตํ.\csromanspacer{} เกน อาทิตฺตํ.\csromanspacer{} ราคคฺคินา โทสคฺคินา โมหคฺคินา อาทิตฺตํ, ชาติยา ชราย มรเณน โสเกหิ ปริเทเวหิ ทุกฺเขหิ โทมนสฺเสหิ อุปายาเสหิ อาทิตฺตนฺติ วทามิ ฯเปฯ.\csromanspacer{} ชิวฺหา อาทิตฺตา, รสา อาทิตฺตา, ชิวฺหาวิญฺญาณํ อาทิตฺตํ, ชิวฺหาสมฺผสฺโส อาทิตฺโต, ยมฺปิทํ ชิวฺหาสมฺผสฺสปจฺจยา อุปฺปชฺชติ เวทยิตํ สุขํ วา ทุกฺขํ วา อทุกฺขมสุขํ วา, ตมฺปิ อาทิตฺตํ.\csromanspacer{} เกน อาทิตฺตํ.\csromanspacer{} ราคคฺคินา โทสคฺคินา โมหคฺคินา อาทิตฺตํ, ชาติยา ชราย มรเณน โสเกหิ ปริเทเวหิ ทุกฺเขหิ โทมนสฺเสหิ อุปายาเสหิ อาทิตฺตนฺติ วทามิ ฯเปฯ.\csromanspacer{} มโน อาทิตฺโต, ธมฺมา อาทิตฺตา, มโนวิญฺญาณํ อาทิตฺตํ, มโนสมฺผสฺโส อาทิตฺโต, ยมฺปิทํ มโนสมฺผสฺสปจฺจยา อุปฺปชฺชติ เวทยิตํ สุขํ วา ทุกฺขํ วา อทุกฺขมสุขํ วา, ตมฺปิ อาทิตฺตํ.\csromanspacer{} เกน อาทิตฺตํ.\csromanspacer{} ราคคฺคินา โทสคฺคินา โมหคฺคินา อาทิตฺตํ, ชาติยา ชราย มรเณน โสเกหิ ปริเทเวหิ ทุกฺเขหิ โทมนสฺเสหิ อุปายาเสหิ อาทิตฺตนฺติ วทามิ.\csromanspacer{} เอวํ ปสฺสํ ภิกฺขเว สุตวา อริยสาวโก จกฺขุสฺมิมฺปิ นิพฺพินฺทติ, รูเปสุปิ นิพฺพินฺทติ, จกฺขุวิญฺญาเณปิ นิพฺพินฺทติ, จกฺขุสมฺผสฺเสปิ นิพฺพินฺทติ, ยมฺปิทํ จกฺขุสมฺผสฺสปจฺจยา}

\gambhira{สฬายตนวคฺคสํยุตฺตปาฬิ}
\prose{\csromanfolio{252}อุปฺปชฺชติ เวทยิตํ สุขํ วา ทุกฺขํ วา อทุกฺขมสุขํ วา, ตสฺมิมฺปิ นิพฺพินฺทติ ฯเปฯ ยมฺปิทํ มโนสมฺผสฺสปจฺจยา อุปฺปชฺชติ เวทยิตํ สุขํ วา ทุกฺขํ วา อทุกฺขมสุขํ วา, ตสฺมิมฺปิ นิพฺพินฺทติ, นิพฺพินฺทํ วิรชฺชติ, วิราคา วิมุจฺจติ, วิมุตฺตสฺมิํ “วิมุตฺตมฺ”อิติ ญาณํ โหติ, “ขีณา ชาติ, วุสิตํ พฺรหฺมจริยํ, กตํ กรณียํ, นาปรํ อิตฺถตฺตายา”ติ ปชานาตีติ.\csromanspacer{} อิทมโวจ ภควา.\csromanspacer{} อตฺตมนา เต ภิกฺขู ภควโต ภาสิตํ อภินนฺทุํ.\csromanspacer{} อิมสฺมิํ จ ปน เวยฺยากรณสฺมิํ ภญฺญมาเน ตสฺส ภิกฺขุสหสฺสสฺส อนุปาทาย อาสเวหิ จิตฺตานิ วิมุจฺจิํสูติ.\csromanspacer{}\csromanspacer{}ฉฏฺฐํ.}

\csromanmatikaanchor{mtk.98}
\csromantocmarkref{subsection}{7. อทฺธภูตสุตฺต}{mtk.98}
\bookmark[dest={mtk.98},level=2]{7. อทฺธภูตสุตฺต}
\csromanheader{2}{7. อทฺธภูตสุตฺต}
\proseitem{29}{เอวํ เม สุตํ\csromandash{}เอกํ สมยํ ภควา ราชคเห วิหรติ เวฬุวเน กลนฺทกนิวาเป.\csromanspacer{} ตตฺร โข ภควา ภิกฺขู อามนฺเตสิ\csromandash{}สพฺพํ ภิกฺขเว อทฺธภูตํ\footnote{อนฺธภูตํ (สี, สฺยา, กํ)}.\csromanspacer{} กิญฺจิ ภิกฺขเว สพฺพํ อทฺธภูตํ.\csromanspacer{} จกฺขุ ภิกฺขเว อทฺธภูตํ, รูปา อทฺธภูตา, จกฺขุวิญฺญาณํ อทฺธภูตํ, จกฺขุสมฺผสฺโส อทฺธภูโต, ยมฺปิทํ จกฺขุสมฺผสฺสปจฺจยา อุปฺปชฺชติ เวทยิตํ สุขํ วา ทุกฺขํ วา อทุกฺขมสุขํ วา, ตมฺปิ อทฺธภูตํ.\csromanspacer{} เกน อทฺธภูตํ.\csromanspacer{} ชาติยา ชราย มรเณน โสเกหิ ปริเทเวหิ ทุกฺเขหิ โทมนสฺเสหิ อุปายาเสหิ อทฺธภูตนฺติ วทามิ ฯเปฯ.\csromanspacer{} ชิวฺหา อทฺธภูโต, รสา อทฺธภูตา, ชิวฺหาวิญฺญาณํ อทฺธภูตํ, ชิวฺหาสมฺผสฺโส อทฺธภูโต, ยมฺปิทํ ชิวฺหาสมฺผสฺสปจฺจยา อุปฺปชฺชติ เวทยิตํ สุขํ วา ทุกฺขํ วา อทุกฺขมสุขํ วา, ตมฺปิ อทฺธภูตํ.\csromanspacer{} เกน อทฺธภูตํ.\csromanspacer{} ชาติยา ชราย มรเณน โสเกหิ ปริเทเวหิ ทุกฺเขหิ โทมนสฺเสหิ อุปายาเสหิ อทฺธภูตนฺติ วทามิ.\csromanspacer{} กาโย อทฺธภูโต ฯเปฯ.\csromanspacer{} มโน อทฺธภูโต, ธมฺมา อทฺธภูตา, มโนวิญฺญาณํ อทฺธภูตํ, มโนสมฺผสฺโส อทฺธภูโต, ยมฺปิทํ มโนสมฺผสฺสปจฺจยา อุปฺปชฺชติ เวทยิตํ สุขํ วา ทุกฺขํ วา อทุกฺขมสุขํ วา, ตมฺปิ อทฺธภูตํ.\csromanspacer{} เกน อทฺธภูตํ.\csromanspacer{} ชาติยา ชราย มรเณน โสเกหิ ปริเทเวหิ ทุกฺเขหิ โทมนสฺเสหิ อุปายาเสหิ อทฺธภูตนฺติ, วทามิ.\csromanspacer{} เอวํ ปสฺสํ ภิกฺขเว สุตวา อริยสาวโก จกฺขุสฺมิมฺปิ นิพฺพินฺทติ, รูเปสุปิ นิพฺพินฺทติ, จกฺขุวิญฺญาเณปิ นิพฺพินฺทติ, จกฺขุสมฺผสฺเสปิ นิพฺพินฺทติ ฯเปฯ ยมฺปิทํ มโนสมฺผสฺสปจฺจยา อุปฺปชฺชติ เวทยิตํ สุขํ วา ทุกฺขํ วา อทุกฺขมสุขํ วา, ตสฺมิมฺปิ นิพฺพินฺทติ, นิพฺพินฺทํ วิรชฺชติ, วิราคา วิมุจฺจติ, วิมุตฺตสฺมิํ “วิมุตฺตมฺ”อิติ ญาณํ}

\vspace{\tipitakaparskip}
\csromancenter{สฬายตนสํยุตฺต โหติ, ขีณา ชาติ, วุสิตํ พฺรหฺมจริยํ, กตํ กรณียํ, นาปรํ อิตฺถตฺตายาติ ปชานาตีติ.\csromanspacer{}\csromanspacer{}สตฺตมํ.}
\csromanheader{2}{8. สมุคฺฆาตสารุปฺปสุตฺต}
\csromanmatikaanchor{mtk.99}
\csromantocmarkref{subsection}{8-10. สมุคฺฆาตสารุปฺปสุตฺตาทิ}{mtk.99}
\bookmark[dest={mtk.99},level=2]{8-10. สมุคฺฆาตสารุปฺปสุตฺตาทิ}
\proseitem{30}{\csromanfolio{253}สพฺพมญฺญิตสมุคฺฆาตสารุปฺปํ โว ภิกฺขเว ปฏิปทํ เทเสสฺสามิ, ตํ สุณาถ สาธุกํ มนสิ กโรถ, ภาสิสฺสามีติ.\csromanspacer{} กตมา จ สา ภิกฺขเว สพฺพมญฺญิตสมุคฺฆาตสารุปฺปา ปฏิปทา.\csromanspacer{} อิธ ภิกฺขเว ภิกฺขุ จกฺขุํ น มญฺญติ, จกฺขุสฺมิํ น มญฺญติ, จกฺขุโต น มญฺญติ, จกฺขุํ เมติ น มญฺญติ.\csromanspacer{} รูเป น มญฺญติ, รูเปสุ น มญฺญติ, รูปโต น มญฺญติ, รูปา เมติ น มญฺญติ.\csromanspacer{} จกฺขุวิญฺญาณํ น มญฺญติ, จกฺขุวิญฺญาณสฺมิํ น มญฺญติ, จกฺขุวิญฺญาณโต น มญฺญติ, จกฺขุวิญฺญาณํ เมติ น มญฺญติ.\csromanspacer{} จกฺขุสมฺผสฺสํ น มญฺญติ, จกฺขุสมฺผสฺสสฺมิํ น มญฺญติ, จกฺขุสมฺผสฺสโต น มญฺญติ, จกฺขุสมฺผสฺโส เมติ น มญฺญติ.\csromanspacer{} ยมฺปิทํ จกฺขุสมฺผสฺสปจฺจยา อุปฺปชฺชติ เวทยิตํ สุขํ วา ทุกฺขํ วา อทุกฺขมสุขํ วา, ตมฺปิ น มญฺญติ, ตสฺมิมฺปิ น มญฺญติ, ตโตปิ น มญฺญติ, ตํ เมติ น มญฺญติ ฯเปฯ.\csromanspacer{} ชิวฺหํ น มญฺญติ, ชิวฺหาย น มญฺญติ, ชิวฺหาโต น มญฺญติ, ชิวฺหา เมติ น มญฺญติ.\csromanspacer{} รเส น มญฺญติ, รเสสุ น มญฺญติ, รสโต น มญฺญติ, รสา เมติ น มญฺญติ.\csromanspacer{} ชิวฺหาวิญฺญาณํ น มญฺญติ, ชิวฺหาวิญฺญาณสฺมิํ น มญฺญติ, ชิวฺหาวิญฺญาณโต น มญฺญติ, ชิวฺหาวิญฺญาณํ เมติ น มญฺญติ.\csromanspacer{} ชิวฺหาสมฺผสฺสํ น มญฺญติ, ชิวฺหาสมฺผสฺสสฺมิํ น มญฺญติ, ชิวฺหาสมฺผสฺสโต น มญฺญติ.\csromanspacer{} ชิวฺหาสมฺผสฺโส เมติ น มญฺญติ.\csromanspacer{} ยมฺปิทํ ชิวฺหาสมฺผสฺสปจฺจยา อุปฺปชฺชติ เวทยิตํ สุขํ วา ทุกฺขํ วา อทุกฺขมสุขํ วา, ตมฺปิ น มญฺญติ, ตสฺมิํปิ น มญฺญติ, ตโตปิ น มญฺญติ, ตํ เมติ น มญฺญติ ฯเปฯ.\csromanspacer{} มนํ น มญฺญติ, มนสฺมิํ น มญฺญติ, มนโต น มญฺญติ, มโน เมติ น มญฺญติ.\csromanspacer{} ธมฺเม น มญฺญติ, ธมฺเมสุ น มญฺญติ, ธมฺมโต น มญฺญติ, ธมฺมา เมติ น มญฺญติ.\csromanspacer{} มโนวิญฺญาณํ น มญฺญติ, มโนวิญฺญาณสฺมิํ น มญฺญติ, มโนวิญฺญาณโต น มญฺญติ, มโนวิญฺญาณํ เมติ น มญฺญติ.\csromanspacer{} มโนสมฺผสฺสํ น มญฺญติ, มโนสมฺผสฺสสฺมิํ น มญฺญติ, มโนสมฺผสฺสโต น มญฺญติ, มโนสมฺผสฺโส เมติ น มญฺญติ.\csromanspacer{} ยมฺปิทํ มโนสมฺผสฺสปจฺจยา อุปฺปชฺชติ เวทยิตํ สุขํ วา ทุกฺขํ วา อทุกฺขมสุขํ วา, ตมฺปิ น มญฺญติ, ตสฺมิมฺปิ น มญฺญติ, ตโตปิ น มญฺญติ, ตํ เมติ น มญฺญติ.\csromanspacer{} สพฺพํ น มญฺญติ, สพฺพสฺมิํ น มญฺญติ, สพฺพโต น มญฺญติ, สพฺพํ เมติ น มญฺญติ.\csromanspacer{} โส เอวํ}

\gambhira{สฬายตนวคฺคสํยุตฺตปาฬิ}
\prose{\csromanfolio{254}อมญฺญมาโน น จ กิญฺจิ โลเก อุปาทิยติ, อนุปาทิยํ น ปริตสฺสติ, อปริตสฺสํ ปจฺจตฺตญฺเญว ปรินิพฺพายติ, “ขีณา ชาติ, วุสิตํ พฺรหฺมจริยํ, กตํ กรณียํ, นาปรํ อิตฺถตฺตายา”ติ ปชานาติ.\csromanspacer{} อยํ โข สา ภิกฺขเว สพฺพมญฺญิตสมุคฺฆาตสารุปฺปา ปฏิปทาติ.\csromanspacer{}\csromanspacer{}อฏฺฐมํ.}

\csromanheader{2}{9. ปฐมสมุคฺฆาตสปฺปายสุตฺต}
\proseitem{31}{สพฺพมญฺญิตสมุคฺฆาตสปฺปายํ โว ภิกฺขเว ปฏิปทํ เทเสสฺสามิ, ตํ สุณาถ.\csromanspacer{} กตมา จ สา ภิกฺขเว สพฺพมญฺญิตสมุคฺฆาตสปฺปายา ปฏิปทา.\csromanspacer{} อิธ ภิกฺขเว ภิกฺขุ จกฺขุํ น มญฺญติ, จกฺขุสฺมิํ น มญฺญติ, จกฺขุโต น มญฺญติ, จกฺขุํ เมติ น มญฺญติ.\csromanspacer{} รูเป น มญฺญติ ฯเปฯ.\csromanspacer{} จกฺขุวิญฺญาณํ น มญฺญติ.\csromanspacer{} จกฺขุสมฺผสฺสํ น มญฺญติ.\csromanspacer{} ยมฺปิทํ จกฺขุสมฺผสฺสปจฺจยา อุปฺปชฺชติ เวทยิตํ สุขํ วา ทุกฺขํ วา อทุกฺขมสุขํ วา, ตมฺปิ น มญฺญติ, ตสฺมิมฺปิ น มญฺญติ, ตโตปิ น มญฺญติ, ตํ เมติ น มญฺญติ.\csromanspacer{} ยํ หิ ภิกฺขเว มญฺญติ, ยสฺมิํ มญฺญติ, ยโต มญฺญติ, ยํ เมติ มญฺญติ, ตโต ตํ โหติ อญฺญถา, อญฺญถาภาวี ภวสตฺโต โลโก ภวเมวาภินนฺทติ ฯเปฯ.\csromanspacer{} ชิวฺหํ น มญฺญติ, ชิวฺหาย น มญฺญติ, ชิวฺหาโต น มญฺญติ, ชิวฺหาเมติ น มญฺญติ.\csromanspacer{} รเส น มญฺญติ ฯเปฯ.\csromanspacer{} ชิวฺหาวิญฺญาณํ น มญฺญติ.\csromanspacer{} ชิวฺหาสมฺผสฺสํ น มญฺญติ.\csromanspacer{} ยมฺปิทํ ชิวฺหาสมฺผสฺสปจฺจยา อุปฺปชฺชติ เวทยิตํ สุขํ วา ทุกฺขํ วา อทุกฺขมสุขํ วา, ตมฺปิ น มญฺญติ, ตสฺมิมฺปิ น มญฺญติ, ตโตปิ น มญฺญติ, ตํ เมติ น มญฺญติ.\csromanspacer{} ยํ หิ ภิกฺขเว มญฺญติ, ยสฺมิํ มญฺญติ, ยโต มญฺญติ, ยํ เมติ มญฺญติ, ตโต ตํ โหติ อญฺญถา, อญฺญถาภาวี ภวสตฺโต โลโก ภวเมวาภินนฺทติ ฯเปฯ.\csromanspacer{} มนํ น มญฺญติ, มนสฺมิํ น มญฺญติ, มนโต น มญฺญติ, มโน เมติ น มญฺญติ.\csromanspacer{} ธมฺเม น มญฺญติ ฯเปฯ.\csromanspacer{} มโนวิญฺญาณํ น มญฺญติ.\csromanspacer{} มโนสมฺผสฺสํ น มญฺญติ.\csromanspacer{} ยมฺปิทํ มโนสมฺผสฺสปจฺจยา อุปฺปชฺชติ เวทยิตํ สุขํ วา ทุกฺขํ วา อทุกฺขมสุขํ วา, ตมฺปิ น มญฺญติ, ตสฺมิมฺปิ น มญฺญติ, ตโตปิ น มญฺญติ, ตํ เมติ น มญฺญติ.\csromanspacer{} ยํ หิ ภิกฺขเว มญฺญติ, ยสฺมิํ มญฺญติ, ยโต มญฺญติ, ยํ เมติ มญฺญติ, ตโต ตํ โหติ อญฺญถา, อญฺญถาภาวี ภวสตฺโต โลโก ภวเมวาภินนฺทติ.\csromanspacer{} ยาวตา ภิกฺขเว ขนฺธธาตุ-อายตนํ, ตมฺปิ น มญฺญติ, ตสฺมิมฺปิ น มญฺญติ, ตโตปิ น มญฺญติ, ตํ เมติ น มญฺญติ.\csromanspacer{} โส เอวํ อมญฺญมาโน น จ กิญฺจิ โลเก อุปาทิยติ, อนุปาทิยํ น ปริตสฺสติ, อปริตสฺสํ ปจฺจตฺตญฺเญว ปรินิพฺพายติ, “ขีณา ชาติ, วุสิตํ พฺรหฺมจริยํ, กตํ กรณียํ, นาปรํ อิตฺถตฺตายา”ติ}

\vspace{\tipitakaparskip}
\csromancenter{สฬายตนสํยุตฺต ปชานาติ.\csromanspacer{} อยํ โข สา ภิกฺขเว สพฺพมญฺญิตสมุคฺฆาตสปฺปายา ปฏิปทาติ.\csromanspacer{} นวมํ.}
\csromanheader{2}{10. ทุติยสมุคฺฆาตสปฺปายสุตฺต}
\proseitem{32}{\csromanfolio{255}สพฺพมญฺญิตสมุคฺฆาตสปฺปายํ โว ภิกฺขเว ปฏิปทํ เทเสสฺสามิ, ตํ สุณาถ.\csromanspacer{} กตมา จ สา ภิกฺขเว สพฺพมญฺญิตสมุคฺฆาตสปฺปายา ปฏิปทา.\csromanspacer{} ตํ กิํ มญฺญถ ภิกฺขเว จกฺขุ นิจฺจํ วา อนิจฺจํ วาติ.\csromanspacer{} อนิจฺจํ ภนฺเต.\csromanspacer{} ยํ ปนานิจฺจํ ทุกฺขํ วา ตํ สุขํ วาติ.\csromanspacer{} ทุกฺขํ ภนฺเต.\csromanspacer{} ยํ ปนานิจฺจํ ทุกฺขํ วิปริณามธมฺมํ, กลฺลํ นุ ตํ สมนุปสฺสิตุํ “เอตํ มม, เอโสหมสฺมิ, เอโส เม อตฺตา”ติ.\csromanspacer{} โน เหตํ ภนฺเต.\csromanspacer{} รูปา ฯเปฯ.\csromanspacer{} จกฺขุวิญฺญาณํ.\csromanspacer{} จกฺขุสมฺผสฺโส นิจฺโจ วา อนิจฺโจ วาติ.\csromanspacer{} อนิจฺโจ ภนฺเต ฯเปฯ.\csromanspacer{} ยมฺปิทํ จกฺขุสมฺผสฺสปจฺจยา อุปฺปชฺชติ เวทยิตํ สุขํ วา ทุกฺขํ วา อทุกฺขมสุขํ วา, ตมฺปิ นิจฺจํ วา อนิจฺจํ วาติ.\csromanspacer{} อนิจฺจํ ภนฺเต.\csromanspacer{} ยํ ปนานิจฺจํ, ทุกฺขํ วา ตํ สุขํ วาติ.\csromanspacer{} ทุกฺขํ ภนฺเต.\csromanspacer{} ยํ ปนานิจฺจํ ทุกฺขํ วิปริณามธมฺมํ, กลฺลํ นุ ตํ สมนุปสฺสิตุํ “เอตํ มม, เอโสหมสฺมิ, เอโส เม อตฺตา”ติ.\csromanspacer{} โน เหตํ ภนฺเต ฯเปฯ.\csromanspacer{} ชิวฺหา นิจฺจาวา อนิจฺจา วาติ.\csromanspacer{} อนิจฺจา ภนฺเต ฯเปฯ.\csromanspacer{} รสา.\csromanspacer{} ชิวฺหาวิญฺญานํ.\csromanspacer{} ชิวฺหาสมฺผสฺโส ฯเปฯ.\csromanspacer{} ยมฺปิทํ ชิวฺหาสมฺผสฺสปจฺจยา อุปฺปชฺชติ เวทยิตํ สุขํ วา ทุกฺขํ วา อทุกฺขมสุขํ วา, ตมฺปิ นิจฺจํ วา อนิจฺจํ วาติ.\csromanspacer{} อนิจฺจํ ภนฺเต ฯเปฯ.\csromanspacer{} ธมฺมา.\csromanspacer{} มโนวิญฺญาณํ.\csromanspacer{} มโนสมฺผสฺโส นิจฺโจ วา อนิจฺโจ วาติ.\csromanspacer{} อนิจฺโจ ภนฺเต.\csromanspacer{} ยมฺปิทํ มโนสมฺผสฺสปจฺจยา อุปฺปชฺชติ เวทยิตํ สุขํ วา ทุกฺขํ วา อทุกฺขมสุขํ วา, ตมฺปิ นิจฺจํ วา อนิจฺจํ วาติ.\csromanspacer{} อนิจฺจํ ภนฺเต.\csromanspacer{} ยํ ปนานิจฺจํ, ทุกฺขํ วา ตํ สุขํ วาติ.\csromanspacer{} ทุกฺขํ ภนฺเต.\csromanspacer{} ยํ ปนานิจฺจํ ทุกฺขํ วิปริณามธมฺมํ, กลฺลํ นุ ตํ สมนุปสฺสิตุํ “เอตํ มม, เอโสหมสฺมิ, เอโส เม อตฺตา”ติ.\csromanspacer{} โน เหตํ ภนฺเต.\csromanspacer{} เอวํ ปสฺสํ ภิกฺขเว สุตวา อริยสาวโก จกฺขุสฺมิมฺปิ นิพฺพินฺทติ, รูเปสุปิ นิพฺพินฺทติ, จกฺขุวิญฺญาเณปิ นิพฺพินฺทติ, จกฺขุสมฺผสฺเสปิ นิพฺพินฺทติ, ยมฺปิทํ จกฺขุสมฺผสฺสปจฺจยา อุปฺปชฺชติ เวทยิตํ สุขํ วา ทุกฺขํ วา อทุกฺขมสุขํ วา, ตสฺมิมฺปิ นิพฺพินฺทติ ฯเปฯ.\csromanspacer{} ชิวฺหายปิ นิพฺพินฺทติ, รเสสุปิ ฯเปฯ.\csromanspacer{} ยมฺปิทํ ชิวฺหาสมฺผสฺสปจฺจยา อุปฺปชฺชติ เวทยิตํ สุขํ วา ทุกฺขํ วา อทุกฺขมสุขํ วา, ตสฺมิมฺปิ นิพฺพินฺทติ, มนสฺมิมฺปิ นิพฺพินฺทติ, ธมฺเมสุปิ นิพฺพินฺทติ, มโนวิญฺญาเณปิ นิพฺพินฺทติ, มโนสมฺผสฺเสปิ นิพฺพินฺทติ, ยมฺปิทํ มโนสมฺผสฺสปจฺจยา อุปฺปชฺชติ เวทยิตํ สุขํ วา ทุกฺขํ วา อทุกฺขมสุขํ วา, ตสฺมิมฺปิ นิพฺพินฺทติ, นิพฺพินฺทํ วิรชฺชติ, วิราคา วิมุจฺจติ, วิมุตฺตสฺมิํ “วิมุตฺตมฺ”อิติ ญาณํ โหติ, “ขีณา ชาติ, วุสิตํ พฺรหฺมจริยํ, กตํ กรณียํ, นาปรํ อิตฺถตฺตายา”ติ}

\gambhira{สฬายตนวคฺคสํยุตฺตปาฬิ}
\prose{\csromanfolio{256}ปชานาติ.\csromanspacer{} อยํ โข สา ภิกฺขเว สพฺพมญฺญิตสมุคฺฆาตสปฺปายา ปฏิปทาติ.\csromanspacer{}\csromanspacer{}ทสมํ.\csromanspacer{} สพฺพวคฺโค ตติโย.}

\titlehead{\textbf{ตสฺสุทฺทานํ}}
\csromangathameasure{สพฺพํ จ เทฺวปิ ปหานา, ปริชานาปเร ทุเว. \\ อาทิตฺตํ อทฺธภูตญฺจ, สารุปฺปา เทฺว จ สปฺปายา.}
\csromangathagroup{\csromangathastanza{สพฺพํ จ เทฺวปิ ปหานา, ปริชานาปเร ทุเว. \\ อาทิตฺตํ อทฺธภูตญฺจ, สารุปฺปา เทฺว จ สปฺปายา.}}

\csromanheader{2}{วคฺโค เตน ปวุจฺจตีติ.}
\csromansectionrule
\csromanmatikaanchor{mtk.100}
\csromantocmarknopage{section}{4. ชาติธมฺมวคฺค}{mtk.100}
\bookmark[dest={mtk.100},level=1]{4. ชาติธมฺมวคฺค}
\csromanheader{1}{4. ชาติธมฺมวคฺค}
\csromanmatikaanchor{mtk.101}
\csromantocmarkref{subsection}{1-10. ชาติธมฺมาทิสุตฺตทสก}{mtk.101}
\bookmark[dest={mtk.101},level=2]{1-10. ชาติธมฺมาทิสุตฺตทสก}
\csromanheader{2}{1-10. ชาติธมฺมาทิสุตฺตทสก}
\proseitem{33}{สาวตฺถินิทานํ.\csromanspacer{} ตตฺร โข ฯเปฯ.\csromanspacer{} สพฺพํ ภิกฺขเว ชาติธมฺมํ.\csromanspacer{} กิญฺจ ภิกฺขเว สพฺพํ ชาติธมฺมํ.\csromanspacer{} จกฺขุ ภิกฺขเว ชาติธมฺมํ.\csromanspacer{} รูปา.\csromanspacer{} จกฺขุวิญฺญาณํ.\csromanspacer{} จกฺขุสมฺผสฺโส ชาติธมฺโม.\csromanspacer{} ยมฺปิทํ จกฺขุสมฺผสฺสปจฺจยา อุปฺปชฺชติ เวทยิตํ สุขํ วา ทุกฺขํ วา อทุกฺขมสุขํ วา, ตมฺปิ ชาติธมฺมํ ฯเปฯ.\csromanspacer{} ชิวฺหา.\csromanspacer{} รสา.\csromanspacer{} ชิวฺหาวิญฺญาณํ.\csromanspacer{} ชิวฺหาสมฺผสฺโส.\csromanspacer{} ยมฺปิทํ ชิวฺหาสมฺผสฺสปจฺจยา อุปฺปชฺชติ เวทยิตํ สุขํ วา ทุกฺขํ วา อทุกฺขมสุขํ วา, ตมฺปิ ชาติธมฺมํ.\csromanspacer{} กาโย ฯเปฯ.\csromanspacer{} มโน ชาติธมฺโม.\csromanspacer{} ธมฺมา ชาติธมฺมา.\csromanspacer{} มโนวิญฺญาณํ ชาติธมฺมํ.\csromanspacer{} มโนสมฺผสฺโส ชาติธมฺโม.\csromanspacer{} ยมฺปิทํ มโนสมฺผสฺสปจฺจยา อุปฺปชฺชติ เวทยิตํ สุขํ วา ทุกฺขํ วา อทุกฺขมสุขํ วา, ตมฺปิ ชาติธมฺมํ.\csromanspacer{} เอวํ ปสฺสํ ภิกฺขเว สุตวา อริยสาวโก จกฺขุสฺมิมฺปิ นิพฺพินฺทติ.\csromanspacer{} รูเปสุปิ.\csromanspacer{} จกฺขุวิญฺญาเณปิ.\csromanspacer{} จกฺขุสมฺผสฺเสปิ ฯเปฯ นาปรํ อิตฺถตฺตายาติ ปชานาตีติ.\csromanspacer{}\csromanspacer{}ปฐมํ.}

\csromanheader{2}{34. สพฺพํ ภิกฺขเว ชราธมฺมํ. (สํขิตฺตํ.)}
\csromanheader{2}{35. สพฺพํ ภิกฺขเว พฺยาธิธมฺมํ. (สํขิตฺตํ.)}
\csromanheader{2}{36. สพฺพํ ภิกฺขเว มรณธมฺมํ. (สํขิตฺตํ.)}
\csromanheader{2}{37. สพฺพํ ภิกฺขเว โสกธมฺมํ. (สํขิตฺตํ.)}
\csromanheader{2}{38. สพฺพํ ภิกฺขเว สํกิเลสิกธมฺมํ. (สํขิตฺตํ.)}
\csromanheader{2}{1. สฬายตนสํยุตฺต}
\csromanheader{2}{39. สพฺพํ ภิกฺขเว ขยธมฺมํ. . (สํขิตฺตํ.)}
\csromanheader{2}{40. สพฺพํ ภิกฺขเว วยธมฺมํ. . (สํขิตฺตํ.)}
\csromanheader{2}{41. สพฺพํ ภิกฺขเว สมุทยธมฺมํ. . (สํขิตฺตํ.)}
\proseitem{42}{\csromanfolio{257}สพฺพํ ภิกฺขเว นิโรธธมฺมํ.\csromanspacer{}\csromanspacer{}(สํขิตฺตํ.).\csromanspacer{} ทสมํ.\csromanspacer{} ชาติธมฺมวคฺโค จตุตฺโถ.}

\titlehead{\textbf{ตสฺสุทฺทานํ}}
\csromangathameasure{ชาติชราพฺยาธิมรณํ, โสโก จ สํกิเลสิกํ. \\ ขยวยสมุทยํ, นิโรธธมฺเมน เต ทสาติ.}
\csromangathagroup{\csromangathastanza{ชาติชราพฺยาธิมรณํ, โสโก จ สํกิเลสิกํ. \\ ขยวยสมุทยํ, นิโรธธมฺเมน เต ทสาติ.}}

\csromanmatikaanchor{mtk.102}
\csromantocmarknopage{section}{5. สพฺพ-อนิจฺจวคฺค}{mtk.102}
\bookmark[dest={mtk.102},level=1]{5. สพฺพ-อนิจฺจวคฺค}
\csromanheader{1}{5. สพฺพ-อนิจฺจวคฺค}
\csromanmatikaanchor{mtk.103}
\csromantocmarkref{subsection}{1-9. อนิจฺจาทิสุตฺตนวก}{mtk.103}
\bookmark[dest={mtk.103},level=2]{1-9. อนิจฺจาทิสุตฺตนวก}
\csromanheader{2}{1-9. อนิจฺจาทิสุตฺตนวก}
\proseitem{43}{สาวตฺถินิทานํ.\csromanspacer{} ตตฺร โข.\csromanspacer{} สพฺพํ ภิกฺขเว อนิจฺจํ.\csromanspacer{} กิญฺจ ภิกฺขเว สพฺพํ อนิจฺจํ.\csromanspacer{} จกฺขุ ภิกฺขเว อนิจฺจํ, รูปา อนิจฺจา, จกฺขุวิญฺญาณํ อนิจฺจํ, จกฺขุสมฺผสฺโส อนิจฺโจ, ยมฺปิทํ จกฺขุสมฺผสฺสปจฺจยา อุปฺปชฺชติ เวทยิตํ สุขํ วา ทุกฺขํ วา อทุกฺขมสุขํ วา, ตมฺปิ อนิจฺจํ ฯเปฯ.\csromanspacer{}\csromanspacer{}ชิวฺหา อนิจฺจา, รสา อนิจฺจา, ชิวฺหาวิญฺญาณํ อนิจฺจํ, ชิวฺหาสมฺผสฺโส อนิจฺโจ, ยมฺปิทํ ชิวฺหาสมฺผสฺสปจฺจยา อุปฺปชฺชติ เวทยิตํ สุขํ วา ทุกฺขํ วา อทุกฺขมสุขํ วา, ตมฺปิ อนิจฺจํ.\csromanspacer{} กาโย อนิจฺโจ ฯเปฯ.\csromanspacer{} มโน อนิจฺโจ, ธมฺมา อนิจฺจา, มโนวิญฺญาณํ อนิจฺจํ, มโนสมฺผสฺโส อนิจฺโจ, ยมฺปิทํ มโนสมฺผสฺสปจฺจยา อุปฺปชฺชติ เวทยิตํ สุขํ วา ทุกฺขํ วา อทุกฺขมสุขํ วา, ตมฺปิ อนิจฺจํ.\csromanspacer{} เอวํ ปสฺสํ ภิกฺขเว สุตวา อริยสาวโก จกฺขุสฺมิมฺปิ นิพฺพินฺทติ, รูเปสุปิ นิพฺพินฺทติ, จกฺขุวิญฺญาเณปิ นิพฺพินฺทติ, จกฺขุสมฺผสฺเสปิ นิพฺพินฺทติ, ยมฺปิทํ จกฺขุสมฺผสฺสปจฺจยา อุปฺปชฺชติ เวทยิตํ สุขํ วา ทุกฺขํ วา อทุกฺขมสุขํ วา, ตสฺมิมฺปิ นิพฺพินฺทติ ฯเปฯ.\csromanspacer{} มนสฺมิมฺปิ นิพฺพินฺทติ, ธมฺเมสุปิ นิพฺพินฺทติ, มโนวิญฺญาเณปิ นิพฺพินฺทติ, มโนสมฺผสฺเสปิ นิพฺพินฺทติ, ยมฺปิทํ มโนสมฺผสฺสปจฺจยา อุปฺปชฺชติ เวทยิตํ สุขํ วา ทุกฺขํ วา อทุกฺขมสุขํ วา, ตสฺมิมฺปิ นิพฺพินฺทติ, นิพฺพินฺทํ วิรชฺชติ, วิราคา วิมุจฺจติ, วิมุตฺตสฺมิํ}

\gambhira{สฬายตนวคฺคสํยุตฺตปาฬิ}
\prose{\csromanfolio{258}“วิมุตฺตมฺ”อิติ ญาณํ โหติ, “ขีณา ชาติ, วุสิตํ พฺรหฺมจริยํ, กตํ กรณียํ, นาปรํ อิตฺถตฺตายา”ติ ปชานาตีติ.\csromanspacer{}\csromanspacer{}ปฐมํ.}

\proseitem{44}{สพฺพํ ภิกฺขเว ทุกฺขํ ฯเปฯ.\csromanspacer{}\csromanspacer{}ทุติยํ.}

\proseitem{45}{สพฺพํ ภิกฺขเว อนตฺตา ฯเปฯ.\csromanspacer{}\csromanspacer{}ตติยํ.}

\proseitem{46}{สพฺพํ ภิกฺขเว อภิญฺเญยฺยํ ฯเปฯ.\csromanspacer{}\csromanspacer{}จตุตฺถํ.}

\proseitem{47}{สพฺพํ ภิกฺขเว ปริญฺเญยฺยํ ฯเปฯ.\csromanspacer{}\csromanspacer{}ปญฺจมํ.}

\proseitem{48}{สพฺพํ ภิกฺขเว ปหาตพฺพํ ฯเปฯ.\csromanspacer{}\csromanspacer{}ฉฏฺฐํ.}

\proseitem{49}{สพฺพํ ภิกฺขเว สจฺฉิกาตพฺพํ ฯเปฯ.\csromanspacer{}\csromanspacer{}สตฺตมํ.}

\proseitem{50}{สพฺพํ ภิกฺขเว อภิญฺญาปริญฺเญยฺยํ ฯเปฯ.\csromanspacer{}\csromanspacer{}อฏฺฐมํ.}

\proseitem{51}{สพฺพํ ภิกฺขเว อุปทฺทุตํ ฯเปฯ.\csromanspacer{}\csromanspacer{}นวมํ.}

\csromanmatikaanchor{mtk.104}
\csromantocmarkref{subsection}{10. อุปสฺสฏฺฐสุตฺต}{mtk.104}
\bookmark[dest={mtk.104},level=2]{10. อุปสฺสฏฺฐสุตฺต}
\csromanheader{2}{10. อุปสฺสฏฺฐสุตฺต}
\proseitem{52}{สพฺพํ ภิกฺขเว อุปสฺสฏฺฐํ\footnote{อุปสฏฺฐํ.}.\csromanspacer{} กิญฺจ ภิกฺขเว สพฺพํ อุปสฺสฏฺฐํ.\csromanspacer{} จกฺขุ ภิกฺขเว อุปสฺสฏฺฐํ, รูปา อุปสฺสฏฺฐา, จกฺขุวิญฺญาณํ อุปสฺสฏฺฐํ, จกฺขุสมฺผสฺโส อุปสฺสฏฺโฐ, ยมฺปิทํ จกฺขุสมฺผสฺสปจฺจยา อุปฺปชฺชติ เวทยิตํ สุขํ วา ทุกฺขํ วา อทุกฺขมสุขํ วา, ตมฺปิ อุปสฺสฏฺฐํ ฯเปฯ.\csromanspacer{} ชิวฺหา อุปสฺสฏฺฐา, รสา อุปสฺสฏฺฐา, ชิวฺหาวิญฺญาณํ อุปสฺสฏฺฐํ, ชิวฺหาสมฺผสฺโส อุปสฺสฏฺโฐ, ยมฺปิทํ ชิวฺหาสมฺผสฺสปจฺจยา อุปฺปชฺชติ เวทยิตํ สุขํ วา ทุกฺขํ วา อทุกฺขมสุขํ วา, ตมฺปิ อุปสฺสฏฺฐํ.\csromanspacer{} กาโย อุปสฺสฏฺโฐ.\csromanspacer{} มโน อุปสฺสฏฺโฐ, ธมฺมา อุปสฺสฏฺฐา, มโนวิญฺญาณํ อุปสฺสฏฺฐํ, มโนสมฺผสฺโส อุปสฺสฏฺโฐ, ยมฺปิทํ มโนสมฺผสฺสปจฺจยา อุปฺปชฺชติ เวทยิตํ สุขํ วา ทุกฺขํ วา อทุกฺขมสุขํ วา, ตมฺปิ อุปสฺสฏฺฐํ.\csromanspacer{} เอวํ ปสฺสํ ภิกฺขเว สุตวา อริยสาวโก จกฺขุสฺมิมฺปิ นิพฺพินฺทติ, รูเปสุปิ นิพฺพินฺทติ, จกฺขุวิญฺญาเณปิ นิพฺพินฺทติ, จกฺขุสมฺผสฺเสปิ นิพฺพินฺทติ, ยมฺปิทํ จกฺขุสมฺผสฺสปจฺจยา อุปฺปชฺชติ เวทยิตํ สุขํ วา ทุกฺขํ วา อทุกฺขมสุขํ วา, ตสฺมิมฺปิ นิพฺพินฺทติ ฯเปฯ.\csromanspacer{} มนสฺมิมฺปิ นิพฺพินฺทติ, ธมฺเมสุปิ นิพฺพินฺทติ, มโนวิญฺญาเณปิ นิพฺพินฺทติ, มโนสมฺผสฺเสปิ นิพฺพินฺทติ, ยมฺปิทํ มโนสมฺผสฺสปจฺจยา อุปฺปชฺชติ เวทยิตํ สุขํ วา ทุกฺขํ วา อทุกฺขมสุขํ วา, ตสฺมิมฺปิ นิพฺพินฺทติ, นิพฺพินฺทํ วิรชฺชติ, วิราคา วิมุจฺจติ, วิมุตฺตสฺมิํ “วิมุตฺตมฺ”อิติ ญาณํ}

\vspace{\tipitakaparskip}
\csromancenter{สฬายตนสํยุตฺต โหติ, “ขีณา ชาติ, วุสิตํ พฺรหฺมจริยํ, กตํ กรณียํ, นาปรํ อิตฺถตฺตายา”ติ ปชานาตีติ.\csromanspacer{}\csromanspacer{}ทสมํ.\csromanspacer{} สพฺพ-อนิจฺจวคฺโค ปญฺจโม.}
\titlehead{\textbf{ตสฺสุทฺทานํ}}
\csromanmatikaanchor{mtk.105}
\csromanmatikaanchor{mtk.106}
\csromantocmarknopage{section}{(6) 1. อวิชฺชาวคฺค}{mtk.105}
\bookmark[dest={mtk.105},level=1]{(6) 1. อวิชฺชาวคฺค}
\csromantocmarkref{subsection}{1-7. อวิชฺชาปหานสุตฺตาทิ}{mtk.106}
\bookmark[dest={mtk.106},level=2]{1-7. อวิชฺชาปหานสุตฺตาทิ}
\csromangathameasure{อนิจฺจํ ทุกฺขํ อนตฺตา, อภิญฺเญยฺยํ ปริญฺเญยฺยํ. \\ ปหาตพฺพํ สจฺฉิกาตพฺพํ, อภิญฺเญยฺยปริญฺเญยฺยํ. \\ อุปทฺทุตํ อุปสฺสฏฺฐํ, วคฺโค เตน ปวุจฺจตีติ. สฬายตนวคฺเค ปฐมปณฺณาสโก สมตฺโต.}
\csromangathagroup{\csromangathastanza{\csromanfolio{259}อนิจฺจํ ทุกฺขํ อนตฺตา, อภิญฺเญยฺยํ ปริญฺเญยฺยํ. \\ ปหาตพฺพํ สจฺฉิกาตพฺพํ, อภิญฺเญยฺยปริญฺเญยฺยํ\footnote{อภิญฺเญยฺยํ ปริญฺเญยฺยํ (สี, สฺยา, กํ), อภิญฺญาตํ ปริญฺเญยฺยํ (อิ, ก)}. \\ อุปทฺทุตํ อุปสฺสฏฺฐํ, วคฺโค เตน ปวุจฺจตีติ.\csromanspacer{} สฬายตนวคฺเค ปฐมปณฺณาสโก สมตฺโต.}}

\csromanheader{2}{\textbf{ตสฺส วคฺคุทฺทานํ}}
\csromangathameasure{อนิจฺจวคฺคํ ยมกํ, สพฺพํ วคฺคํ ชาติธมฺมํ. \\ อนิจฺจวคฺเคน ปญฺญาสํ, ปญฺจโม เตน ปวุจฺจตีติ.}
\csromangathagroup{\csromangathastanza{อนิจฺจวคฺคํ ยมกํ, สพฺพํ วคฺคํ ชาติธมฺมํ. \\ อนิจฺจวคฺเคน ปญฺญาสํ, ปญฺจโม เตน ปวุจฺจตีติ.}}

\csromanheader{1}{\textbf{(6) 1. อวิชฺชาวคฺค}}
\csromanheader{2}{1. อวิชฺชาปหานสุตฺต}
\proseitem{53}{สาวตฺถินิทานํ.\csromanspacer{} อถ โข อญฺญตโร ภิกฺขุ เยน ภควา เตนุปสงฺกมิ, อุปสงฺกมิตฺวา ภควนฺตํ อภิวาเทตฺวา เอกมนฺตํ นิสีทิ, เอกมนฺตํ นิสินฺโน โข โส ภิกฺขุ ภควนฺตํ เอตทโวจ “กถํ นุ โข ภนฺเต ชานโต กถํ ปสฺสโต อวิชฺชา ปหียติ, วิชฺชา อุปฺปชฺชตี”ติ.}

\prose{จกฺขุํ โข ภิกฺขุ อนิจฺจโต ชานโต ปสฺสโต อวิชฺชา ปหียติ, วิชฺชา อุปฺปชฺชติ.\csromanspacer{} รูเป อนิจฺจโต ชานโต ปสฺสโต อวิชฺชา ปหียติ, วิชฺชา อุปฺปชฺชติ.\csromanspacer{} จกฺขุวิญฺญาณํ.\csromanspacer{} จกฺขุสมฺผสฺสํ.\csromanspacer{} ยมฺปิทํ จกฺขุสมฺผสฺสปจฺจยา อุปฺปชฺชติ เวทยิตํ สุขํ วา ทุกฺขํ วา อทุกฺขมสุขํ วา, ตมฺปิ อนิจฺจโต ชานโต ปสฺสโต อวิชฺชา ปหียติ, วิชฺชา อุปฺปชฺชติ.\csromanspacer{} โสตํ.\csromanspacer{} ฆานํ.\csromanspacer{} ชิวฺหํ.\csromanspacer{} กายํ.\csromanspacer{} มนํ อนิจฺจโต ชานโต ปสฺสโต อวิชฺชา ปหียติ, วิชฺชา อุปฺปชฺชติ.}

\gambhira{สฬายตนวคฺคสํยุตฺตปาฬิ}
\prose{\csromanfolio{260}ธมฺเม.\csromanspacer{} มโนวิญฺญาณํ.\csromanspacer{} มโนสมฺผสฺสํ.\csromanspacer{} ยมฺปิทํ มโนสมฺผสฺสปจฺจยา อุปฺปชฺชติ เวทยิตํ สุขํ วา ทุกฺขํ วา อทุกฺขมสุขํ วา, ตมฺปิ อนิจฺจโต ชานโต ปสฺสโต อวิชฺชา ปหียติ, วิชฺชา อุปฺปชฺชติ.\csromanspacer{} เอวํ โข ภิกฺขุ ชานโต เอวํ ปสฺสโต อวิชฺชา ปหียติ, วิชฺชา อุปฺปชฺชตีติ.\csromanspacer{}\csromanspacer{}ปฐมํ.}

\csromanheader{2}{2. สํโยชนปหานสุตฺต}
\proseitem{54}{กถํ นุ โข ภนฺเต ชานโต กถํ ปสฺสโต สํโยชนา ปหียนฺตีติ.\csromanspacer{} จกฺขุํ โข ภิกฺขุ อนิจฺจโต ชานโต ปสฺสโต สํโยชนา ปหียนฺติ.\csromanspacer{} รูเป.\csromanspacer{} จกฺขุวิญฺญาณํ.\csromanspacer{} จกฺขุสมฺผสฺสํ.\csromanspacer{} ยมฺปิทํ จกฺขุสมฺผสฺสปจฺจยา อุปฺปชฺชติ เวทยิตํ สุขํ วา ทุกฺขํ วา อทุกฺขมสุขํ วา, ตมฺปิ อนิจฺจโต ชานโต ปสฺสโต สํโยชนา ปหียนฺติ.\csromanspacer{} โสตํ.\csromanspacer{} ฆานํ.\csromanspacer{} ชิวฺหํ.\csromanspacer{} กายํ.\csromanspacer{} มนํ.\csromanspacer{} ธมฺเม.\csromanspacer{} มโนวิญฺญาณํ.\csromanspacer{} มโนสมฺผสฺสํ.\csromanspacer{} ยมฺปิทํ มโนสมฺผสฺสปจฺจยา อุปฺปชฺชติ เวทยิตํ สุขํ วา ทุกฺขํ วา อทุกฺขมสุขํ วา, ตมฺปิ อนิจฺจโต ชานโต ปสฺสโต สํโยชนา ปหียนฺติ.\csromanspacer{} เอวํ โข ภิกฺขุ ชานโต เอวํ ปสฺสโต สํโยชนา ปหียนฺตีติ.\csromanspacer{}\csromanspacer{}ทุติยํ.}

\csromanheader{2}{3. สํโยชนสมุคฺฆาตสุตฺต}
\proseitem{55}{กถํ นุ โข ภนฺเต ชานโต กถํ ปสฺสโต สํโยชนา สมุคฺฆาตํ คจฺฉนฺตีติ.\csromanspacer{} จกฺขุํ โข ภิกฺขุ อนตฺตโต ชานโต ปสฺสโต สํโยชนา สมุคฺฆาตํ คจฺฉนฺติ.\csromanspacer{} รูเป อนตฺตโต.\csromanspacer{} จกฺขุวิญฺญาณํ.\csromanspacer{} จกฺขุสมฺผสฺสํ.\csromanspacer{} ยมฺปิทํ จกฺขุสมฺผสฺสปจฺจยา อุปฺปชฺชติ เวทยิตํ สุขํ วา ทุกฺขํ วา อทุกฺขมสุขํ วา, ตมฺปิ อนตฺตโต ชานโต ปสฺสโต สํโยชนา สมุคฺฆาตํ คจฺฉนฺติ.\csromanspacer{} โสตํ.\csromanspacer{} ฆานํ.\csromanspacer{} ชิวฺหํ.\csromanspacer{} กายํ.\csromanspacer{} มนํ.\csromanspacer{} ธมฺเม.\csromanspacer{} มโนวิญฺญาณํ.\csromanspacer{} มโนสมฺผสฺสํ.\csromanspacer{} ยมฺปิทํ มโนสมฺผสฺสปจฺจยา อุปฺปชฺชติ เวทยิตํ สุขํ วา ทุกฺขํ วา อทุกฺขมสุขํ วา, ตมฺปิ อนตฺตโต ชานโต ปสฺสโต สํโยชนา สมุคฺฆาตํ คจฺฉนฺติ.\csromanspacer{} เอวํ โข ภิกฺขุ ชานโต เอวํ ปสฺสโต สํโยชนา สมุคฺฆาตํ คจฺฉนฺตีติ.\csromanspacer{}\csromanspacer{}ตติยํ.}

\csromanheader{2}{4. อาสวปหานสุตฺต}
\vspace{\tipitakaparskip}
\csromancenter{กถํ นุ โข ภนฺเต ชานโต กถํ ปสฺสโต อาสวา ปหียนฺตีติ ฯเปฯ.\csromanspacer{}\csromanspacer{}จตุตฺถํ.}
\csromanheader{2}{1. สฬายตนสํยุตฺต \textbf{5. อาสวสมุคฺฆาตสุตฺต}}
\proseitem{57}{\csromanfolio{261}กถํ นุ โข ภนฺเต ชานโต กถํ ปสฺสโต อาสวา สมุคฺฆาตํ คจฺฉนฺตีติ ฯเปฯ.\csromanspacer{}\csromanspacer{}ปญฺจมํ.}

\csromanheader{2}{6. อนุสยปหานสุตฺต}
\vspace{\tipitakaparskip}
\csromancenter{กถํ นุ โข ฯเปฯ อนุสยา ปหียนฺตีติ ฯเปฯ.\csromanspacer{}\csromanspacer{}ฉฏฺฐํ.}
\csromanheader{2}{7. อนุสยสมุคฺฆาตสุตฺต}
\proseitem{59}{กถํ นุ โข ฯเปฯ อนุสยา สมุคฺฆาตํ คจฺฉนฺตีติ.\csromanspacer{} จกฺขุํ โข ภิกฺขุ อนตฺตโต ชานโต ปสฺสโต อนุสยา สมุคฺฆาตํ คจฺฉนฺติ ฯเปฯ.\csromanspacer{} โสตํ.\csromanspacer{} ฆานํ.\csromanspacer{} ชิวฺหํ.\csromanspacer{} กายํ.\csromanspacer{} มนํ.\csromanspacer{} ธมฺเม.\csromanspacer{} มโนวิญฺญาณํ.\csromanspacer{} มโนสมฺผสฺสํ.\csromanspacer{} ยมฺปิทํ มโนสมฺผสฺสปจฺจยา อุปฺปชฺชติ เวทยิตํ สุขํ วา ทุกฺขํ วา อทุกฺขมสุขํ วา, ตมฺปิ อนตฺตโต ชานโต ปสฺสโต อนุสยา สมุคฺฆาตํ คจฺฉนฺติ.\csromanspacer{} เอวํ โข ภิกฺขุ ชานโต เอวํ ปสฺสโต อนุสยา สมุคฺฆาตํ คจฺฉนฺตีติ.\csromanspacer{}\csromanspacer{}สตฺตมํ.}

\csromanmatikaanchor{mtk.107}
\csromantocmarkref{subsection}{8. สพฺพุปาทานปริญฺญาสุตฺต}{mtk.107}
\bookmark[dest={mtk.107},level=2]{8. สพฺพุปาทานปริญฺญาสุตฺต}
\csromanheader{2}{8. สพฺพุปาทานปริญฺญาสุตฺต}
\proseitem{60}{สพฺพุปาทานปริญฺญาย โว ภิกฺขเว ธมฺมํ เทเสสฺสามิ, ตํ สุณาถ.\csromanspacer{} กตโม จ ภิกฺขเว สพฺพุปาทานปริญฺญาย ธมฺโม.\csromanspacer{} จกฺขุญฺจ ปฏิจฺจ รูเป จ อุปฺปชฺชติ จกฺขุวิญฺญาณํ, ติณฺณํ สงฺคติ ผสฺโส, ผสฺสปจฺจยา เวทนา, เอวํ ปสฺสํ ภิกฺขเว สุตวา อริยสาวโก จกฺขุสฺมิมฺปิ นิพฺพินฺทติ, รูเปสุปิ นิพฺพินฺทติ, จกฺขุวิญฺญาเณปิ นิพฺพินฺทติ, จกฺขุสมฺผสฺเสปิ นิพฺพินฺทติ, เวทนายปิ นิพฺพินฺทติ, นิพฺพินฺทํ วิรชฺชติ, วิราคา วิมุจฺจติ, วิโมกฺขา\footnote{วิโมกฺขํ (ก), วิโมกฺข (สฺยา, กํ)} “ปริญฺญาตํ เม อุปาทานนฺ”ติ ปชานาติ.\csromanspacer{} โสตญฺจ ปฏิจฺจ สทฺเท จ อุปฺปชฺชติ.\csromanspacer{} ฆานญฺจ ปฏิจฺจ คนฺเธ จ.\csromanspacer{} ชิวฺหญฺจ ปฏิจฺจ รเส จ.\csromanspacer{} กายญฺจ ปฏิจฺจ โผฏฺฐพฺเพ จ.\csromanspacer{} มนญฺจ ปฏิจฺจ ธมฺเม จ อุปฺปชฺชติ มโนวิญฺญาณํ, ติณฺณํ สงฺคติ ผสฺโส, ผสฺสปจฺจยา เวทนา.\csromanspacer{} เอวํ ปสฺสํ ภิกฺขเว สุตวา อริยสาวโก มนสฺมิมฺปิ นิพฺพินฺทติ, ธมฺเมสุปิ นิพฺพินฺทติ, มโนวิญฺญาเณปิ นิพฺพินฺทติ, มโนสมฺผสฺเสปิ นิพฺพินฺทติ, เวทนายปิ นิพฺพินฺทติ, นิพฺพินฺทํ วิรชฺชติ, วิราคา วิมุจฺจติ, วิโมกฺขา1 “ปริญฺญาตํ เม}

\gambhira{สฬายตนวคฺคสํยุตฺตปาฬิ}
\csromanmatikaanchor{mtk.108}
\csromantocmarkref{subsection}{9-10. สพฺพุปาทานปริยาทานสุตฺต}{mtk.108}
\bookmark[dest={mtk.108},level=2]{9-10. สพฺพุปาทานปริยาทานสุตฺต}
\prose{\csromanfolio{262}อุปาทานนฺ”ติ ปชานาติ.\csromanspacer{} อยํ โข ภิกฺขเว สพฺพุปาทานปริญฺญาย ธมฺโมติ.\csromanspacer{}\csromanspacer{}อฏฺฐมํ.}

\csromanheader{2}{9. ปฐมสพฺพุปาทานปริยาทานสุตฺต}
\proseitem{61}{สพฺพุปาทานปริยาทานาย โว ภิกฺขเว ธมฺมํ เทเสสฺสามิ, ตํ สุณาถ.\csromanspacer{} กตโม จ ภิกฺขเว สพฺพุปาทานปริยาทานาย ธมฺโม.\csromanspacer{} จกฺขุญฺจ ปฏิจฺจ รูเป จ อุปฺปชฺชติ จกฺขุวิญฺญาณํ, ติณฺณํ สงฺคติ ผสฺโส, ผสฺสปจฺจยา เวทนา.\csromanspacer{} เอวํ ปสฺสํ ภิกฺขเว สุตวา อริยสาวโก จกฺขุสฺมิมฺปิ นิพฺพินฺทติ, รูเปสุปิ นิพฺพินฺทติ, จกฺขุวิญฺญาเณปิ นิพฺพินฺทติ, จกฺขุสมฺผสฺเสปิ นิพฺพินฺทติ, เวทนายปิ นิพฺพินฺทติ, นิพฺพินฺทํ วิรชฺชติ, วิราคา วิมุจฺจติ, วิโมกฺขา “ปริยาทินฺนํ\footnote{สพฺพตฺถปิ เอวเมว ทิสฺสติ ทนฺตช-นกาเรเนว.} เม อุปาทานนฺ”ติ ปชานาติ ฯเปฯ.\csromanspacer{} ชิวฺหญฺจ ปฏิจฺจ รเส จ อุปฺปชฺชติ ชิวฺหาวิญฺญาณํ ฯเปฯ.\csromanspacer{} มนญฺจ ปฏิจฺจ ธมฺเม จ อุปฺปชฺชติ มโนวิญฺญาณํ, ติณฺณํ สงฺคติ ผสฺโส, ผสฺสปจฺจยา เวทนา.\csromanspacer{} เอวํ ปสฺสํ ภิกฺขเว สุตวา อริยสาวโก มนสฺมิมฺปิ นิพฺพินฺทติ, ธมฺเมสุปิ นิพฺพินฺทติ, มโนวิญฺญาเณปิ นิพฺพินฺทติ, มโนสมฺผสฺเสปิ นิพฺพินฺทติ, เวทนายปิ นิพฺพินฺทติ, นิพฺพินฺทํ วิรชฺชติ, วิราคา วิมุจฺจติ, วิโมกฺขา “ปริยาทินฺนํ เม อุปาทานนฺ”ติ ปชานาติ.\csromanspacer{} อยํ โข ภิกฺขเว สพฺพุปาทานปริยาทานาย ธมฺโมติ.\csromanspacer{}\csromanspacer{}นวมํ.}

\csromanheader{2}{10. ทุติยสพฺพุปาทานปริยาทานสุตฺต}
\proseitem{62}{สพฺพุปาทานปริยาทานาย โว ภิกฺขเว ธมฺมํ เทเสสฺสามิ, ตํ สุณาถ.\csromanspacer{} กตโม จ ภิกฺขเว สพฺพุปาทานปริยาทานาย ธมฺโม.\csromanspacer{} ตํ กิํ มญฺญถ ภิกฺขเว, จกฺขุ นิจฺจํ วา อนิจฺจํ วาติ.\csromanspacer{} อนิจฺจํ ภนฺเต.\csromanspacer{} ยํ ปนานิจฺจํ, ทุกฺขํ วา ตํ สุขํ วาติ.\csromanspacer{} ทุกฺขํ ภนฺเต.\csromanspacer{} ยํ ปนานิจฺจํ ทุกฺขํ วิปริณามธมฺมํ, กลฺลํ นุ ตํ สมนุปสฺสิตุํ “เอตํ มม, เอโสหมสฺมิ, เอโส เม อตฺตา”ติ.\csromanspacer{} โน เหตํ ภนฺเต.\csromanspacer{} รูปา ฯเปฯ.\csromanspacer{} จกฺขุวิญฺญาณํ นิจฺจํ วา อนิจฺจํ วาติ.\csromanspacer{} อนิจฺจํ ภนฺเต.\csromanspacer{} จกฺขุสมฺผสฺโส นิจฺโจ วา อนิจฺโจ วาติ.\csromanspacer{} อนิจฺโจ ภนฺเต ฯเปฯ.\csromanspacer{} ยมฺปิทํ จกฺขุสมฺผสฺสปจฺจยา อุปฺปชฺชติ เวทยิตํ สุขํ วา ทุกฺขํ วา อทุกฺขมสุขํ วา, ตมฺปิ นิจฺจํ วา อนิจฺจํ วาติ.\csromanspacer{} อนิจฺจํ ภนฺเต.\csromanspacer{} โสตํ.\csromanspacer{} ฆานํ.\csromanspacer{} ชิวฺหา.\csromanspacer{} กาโย.\csromanspacer{} มโน.\csromanspacer{} ธมฺมา.\csromanspacer{} มโนวิญฺญาณํ.\csromanspacer{} มโนสมฺผสฺโส.\csromanspacer{} ยมฺปิทํ มโนสมฺผสฺสปจฺจยา อุปฺปชฺชติ เวทยิตํ สุขํ วา ทุกฺขํ วา อทุกฺขมสุขํ วา, ตมฺปิ นิจฺจํ วา อนิจฺจํ วาติ.}

\vspace{\tipitakaparskip}
\csromancenter{สฬายตนสํยุตฺต อนิจฺจํ ภนฺเต.\csromanspacer{} ยํ ปนานิจฺจํ, ทุกฺขํ วา ตํ สุขํ วาติ.\csromanspacer{} ทุกฺขํ ภนฺเต.\csromanspacer{} ยํ ปนานิจฺจํ ทุกฺขํ วิปริณามธมฺมํ, กลฺลํ นุ ตํ สมนุปสฺสิตุํ “เอตํ มม, เอโสหมสฺมิ, เอโส เม อตฺตา”ติ.\csromanspacer{} โน เหตํ ภนฺเต.\csromanspacer{} เอวํ ปสฺสํ ภิกฺขเว สุตวา อริยสาวโก จกฺกุสฺมิมฺปิ นิพฺพินฺทติ, รูเปสุปิ นิพฺพินฺทติ, จกฺขุวิญฺญาเณปิ นิพฺพินฺทติ, จกฺขุสมฺผสฺเสปิ นิพฺพินฺทติ.\csromanspacer{} ยมฺปิทํ จกฺขุสมฺผสฺสปจฺจยา อุปฺปชฺชติ เวทยิตํ สุขํ วา ทุกฺขํ วา อทุกฺขมสุขํ วา, ตสฺมิมฺปิ นิพฺพินฺทติ ฯเปฯ ชิวฺหายปิ นิพฺพินฺทติ, รเสสุปิ นิพฺพินฺทติ, ชิวฺหาวิญฺญาเณปิ นิพฺพินฺทติ, ชิวฺหาสมฺผสฺเสปิ นิพฺพินฺทติ.\csromanspacer{} ยมฺปิทํ ชิวฺหาสมฺผสฺสปจฺจยา อุปฺปชฺชติ ฯเปฯ มนสฺมิมฺปิ นิพฺพินฺทติ, ธมฺเมสุปิ นิพฺพินฺทติ, มโนวิญฺญาเณปิ นิพฺพินฺทติ, มโนสมฺผสฺเสปิ นิพฺพินฺทติ.\csromanspacer{} ยมฺปิทํ มโนสมฺผสฺสปจฺจยา อุปฺปชฺชติ เวทยิตํ สุขํ วา ทุกฺขํ วา อทุกฺขมสุขํ วา, ตสฺมิมฺปิ นิพฺพินฺทติ, นิพฺพินฺทํ วิรชฺชติ, วิราคา วิมุจฺจติ, วิมุตฺตสฺมิํ “วิมุตฺตมฺ”อิติ ญาณํ โหติ, “ขีณา ชาติ, วุสิตํ พฺรหฺมจริยํ, กตํ กรณียํ, นาปรํ อิตฺถตฺตายา”ติ ปชานาติ.\csromanspacer{} อยํ โข ภิกฺขเว สพฺพุปาทาปริยาทานาย ธมฺโมติ.\csromanspacer{}\csromanspacer{}ทสมํ.\csromanspacer{} อวิชฺชาวคฺโค ปฐโม.}
\titlehead{\textbf{ตสฺสุทฺทานํ}}
\csromanmatikaanchor{mtk.109}
\csromanmatikaanchor{mtk.110}
\csromantocmarknopage{section}{(7) 2. มิคชาลวคฺค}{mtk.109}
\bookmark[dest={mtk.109},level=1]{(7) 2. มิคชาลวคฺค}
\csromantocmarkref{subsection}{1-2. มิคชาลสุตฺต}{mtk.110}
\bookmark[dest={mtk.110},level=2]{1-2. มิคชาลสุตฺต}
\csromangathameasure{อวิชฺชา สํโยชนา เทฺว, อาสเวน ทุเว วุตฺตา. \\ อนุสยา อปเร เทฺว, ปริญฺญา เทฺว ปริยาทินฺนํ.}
\csromangathagroup{\csromangathastanza{\csromanfolio{263}อวิชฺชา สํโยชนา เทฺว, อาสเวน ทุเว วุตฺตา. \\ อนุสยา อปเร เทฺว, ปริญฺญา เทฺว ปริยาทินฺนํ.}}

\csromanheader{2}{วคฺโค เตน ปวุจฺจตีติ.}
\csromansectionrule
\csromanheader{1}{\textbf{(7) 2. มิคชาลวคฺค}}
\csromanheader{2}{1. ปฐมมิคชาลสุตฺต}
\proseitem{63}{สาวตฺถินิทานํ.\csromanspacer{} อถ โข อายสฺมา มิคชาโล เยน ภควา ฯเปฯ เอกมนฺตํ นิสินฺโน โข อายสฺมา มิคชาโล ภควนฺตํ เอตทโวจ\csromandash{} “เอกวิหารี เอกวิหารี”ติ ภนฺเต วุจฺจติ, กิตฺตาวตา นุ โข ภนฺเต เอกวิหารี โหติ, กิตฺตาวตา จ ปน สทุติยวิหารี โหตีติ.}

\prose{สนฺติ โข มิคชาล จกฺขุวิญฺเญยฺยา รูปา อิฏฺฐา กนฺตา มนาปา ปิยรูปา กามูปสํหิตา รชนียา, ตญฺเจ ภิกฺขุ อภินนฺทติ อภิวทติ}

\gambhira{สฬายตนวคฺคสํยุตฺตปาฬิ}
\prose{\csromanfolio{264}อชฺโฌสาย ติฏฺฐติ, ตสฺส ตํ อภินนฺทโต อภิวทโต อชฺโฌสาย ติฏฺฐโต อุปฺปชฺชติ นนฺที\footnote{นนฺทิ (สี, สฺยา, กํ, อิ)}, นนฺทิยา สติ สาราโค โหติ, สาราเค สติ สํโยโค โหติ, นนฺทิสํโยชนสํยุตฺโต โข มิคชาล ภิกฺขุ “สทุติยวิหารี”ติ วุจฺจติ ฯเปฯ.\csromanspacer{} สนฺติ โข มิคชาล ชิวฺหาวิญฺเญยฺยารสา อิฏฺฐา กนฺตา มนาปา ปิยรูปา กามูปสํหิตา รชนียา, ตญฺเจ ภิกฺขุ อภินนฺทติ อภิวทติ อชฺโฌสาย ติฏฺฐติ, ตสฺส ตํ อภินนฺทโต อภิวทโต อชฺโฌสาย ติฏฺฐโต อุปฺปชฺชติ นนฺที, นนฺทิยา สติ สาราโค โหติ, สาราเค สติ สํโยโค โหติ, นนฺทิสํโยชนสํยุตฺโต โข มิคชาล ภิกฺขุ “สทุติยวิหารี”ติ วุจฺจติ.\csromanspacer{} เอวํวิหารี จ มิคชาล ภิกฺขุ กิญฺจาปิ อรญฺญวนปตฺถานิ ปนฺตานิ เสนาสนานิ ปฏิเสวติ อปฺปสทฺทานิ อปฺปนิคฺโฆสานิ วิชนวาตานิ มนุสฺสราหสฺเสยฺยกานิ\footnote{มนุสฺสราหเสยฺยกานิ (สี, สฺยา, กํ, อิ)} ปฏิสลฺลานสารุปฺปานิ, อถ โข “สทุติยวิหารี”ติ วุจฺจติ.\csromanspacer{} ตํ กิสฺส เหตุ, ตณฺหา หิสฺส ทุติยา, สาสฺส อปฺปหีนา, ตสฺมา “สทุติยวิหารี”ติ วุจฺจติ.}

\prose{สนฺติ จ โข มิคชาล จกฺขุวิญฺเญยฺยา รูปา อิฏฺฐา กนฺตา มนาปา ปิยรูปา กามูปสํหิตา รชนียา, ตญฺเจ ภิกฺขุ นาภินนฺทติ นาภิวทติ นาชฺโฌสาย ติฏฺฐติ, ตสฺส ตํ อนภินนฺทโต อนภิวทโต อนชฺโฌสาย ติฏฺฐโต นนฺที นิรุชฺฌติ, นนฺทิยา อสติ สาราโค น โหติ, สาราเค อสติ สํโยโค น โหติ, นนฺทิสํโยชนวิสํยุตฺโต โข มิคชาล ภิกฺขุ “เอกวิหารี”ติ วุจฺจติ ฯเปฯ.\csromanspacer{} สนฺติ จ โข มิคชาล ชิวฺหาวิญฺเญยฺยา รสา ฯเปฯ.\csromanspacer{} สนฺติ จ โข มิคชาล มโนวิญฺเญยฺยา ธมฺมา อิฏฺฐา กนฺตา มนาปา ปิยรูปา กามูปสํหิตา รชนียา, ตญฺเจ ภิกฺขุ นาภินนฺทติ นาภิวทติ นาชฺโฌสาย ติฏฺฐติ, ตสฺส ตํ อนภินนฺทโต อนภิวทโต อนชฺโฌสาย ติฏฺฐโต นนฺที นิรุชฺฌติ, นนฺทิยา อสติ สาราโค น โหติ, สาราเค อสติ สํโยโค น โหติ, นนฺทิสํโยชนวิปฺปยุตฺโต โข มิคชาล ภิกฺขุ “เอกวิหารี”ติ วุจฺจติ.\csromanspacer{} เอวํวิหารี จ มิคชาล ภิกฺขุ กิญฺจาปิ คามนฺเต วิหรติ อากิณฺโณ ภิกฺขูหิ ภิกฺขุนีหิ อุปาสเกหิ อุปาสิกาหิ ราชูหิ ราชมหามตฺเตหิ ติตฺถิเยหิ ติตฺถิยสาวเกหิ, อถ โข “เอกวิหารี”ติ วุจฺจติ.\csromanspacer{} ตํ กิสฺส เหตุ, ตณฺหา หิสฺส ทุติยา, สาสฺส ปหีนา, ตสฺมา “เอกวิหารี”ติ วุจฺจตีติ.\csromanspacer{}\csromanspacer{}ปฐมํ.}

\csromanheader{2}{1. สฬายตนสํยุตฺต \textbf{2. ทุติยมิคชาลสุตฺต}}
\proseitem{64}{\csromanfolio{265}อถ โข อายสฺมา มิคชาโล เยน ภควา เตนุปสงฺกมิ ฯเปฯ เอกมนฺตํ นิสินฺโน โข อายสฺมา มิคชาโล ภควนฺตํ เอตทโวจ\csromandash{}“สาธุ เม ภนฺเต ภควา สํขิตฺเตน ธมฺมํ เทเสตุ, ยมหํ ภควโต ธมฺมํ สุตฺวา เอโก วูปกฏฺโฐ อปฺปมตฺโต อาตาปี ปหิตตฺโต วิหเรยฺยนฺ”ติ.}

\prose{สนฺติ โข มิคชาล จกฺขุวิญฺเญยฺยา รูปา อิฏฺฐา กนฺตา มนาปา ปิยรูปา กามูปสํหิตา รชนียา, ตญฺเจ ภิกฺขุ อภินนฺทติ อภิวทติ อชฺโฌสาย ติฏฺฐติ, ตสฺส ตํ อภินนฺทโต อภิวทโต อชฺโฌสาย ติฏฺฐโต อุปฺปชฺชติ นนฺที, นนฺทิสมุทยา ทุกฺขสมุทโย มิคชาลาติ วทามิ ฯเปฯ.\csromanspacer{} สนฺติ จ โข มิคชาล ชิวฺหาวิญฺเญยฺยา รสา ฯเปฯ.\csromanspacer{} สนฺติ จ โข มิคชาล มโนวิญฺเญยฺยา ธมฺมา อิฏฺฐา กนฺตา มนาปา ปิยรูปา กามูปสํหิตา รชนียา, ตญฺเจ ภิกฺขุ อภินนฺทติ อภิวทติ อชฺโฌสาย ติฏฺฐติ, ตสฺส ตํ อภินนฺทโต อภิวทโต อชฺโฌสาย ติฏฺฐโต อุปฺปชฺชติ นนฺที, นนฺทิสมุทยา ทุกฺขสมุทโย มิคชาลาติ วทามิ.}

\prose{สนฺติ จ โข มิคชาล จกฺขุวิญฺเญยฺยา รูปา อิฏฺฐา กนฺตา มนาปา ปิยรูปา กามูปสํหิตา รชนียา, ตญฺเจ ภิกฺขุ นาภินนฺทติ นาภิวทติ นาชฺโฌสาย ติฏฺฐติ, ตสฺส ตํ อนภินนฺทโต อนภิวทโต อนชฺโฌสาย ติฏฺฐโต นนฺที นิรุชฺฌติ, นนฺทินิโรธา ทุกฺขนิโรโธ มิคชาลาติ วทามิ ฯเปฯ.\csromanspacer{} สนฺติ จ โข มิคชาล ชิวฺหาวิญฺเญยฺยา รสา อิฏฺฐา กนฺตา ฯเปฯ.\csromanspacer{} สนฺติ จ โข มิคชาล มโนวิญฺเญยฺยา ธมฺมา อิฏฺฐา กนฺตา มนาปา ปิยรูปา กามูปสํหิตา รชนียา, ตญฺเจ ภิกฺขุ นาภินนฺทติ นาภิวทติ นาชฺโฌสาย ติฏฺฐติ, ตสฺส ตํ อนภินนฺทโต อนภิวทโต อนชฺโฌสาย ติฏฺฐโต นนฺที นิรุชฺฌติ, นนฺทินิโรธา ทุกฺขนิโรโธ มิคชาลาติ วทามีติ.}

\prose{อถ โข อายสฺมา มิคชาโล ภควโต ภาสิตํ อภินนฺทิตฺวา อนุโมทิตฺวา อุฏฺฐายาสนา ภควนฺตํ อภิวาเทตฺวา ปทกฺขิณํ กตฺวา ปกฺกามิ.\csromanspacer{} อถ โข อายสฺมา มิคชาโล เอโก วูปกฏฺโฐ อปฺปมตฺโต อาตาปี ปหิตตฺโต วิหรโต นจิรสฺเสว ยสฺสตฺถาย กุลปุตฺตา สมฺมเทว อคารสฺมา อนคาริยํ ปพฺพชนฺติ, ตทนุตฺตรํ พฺรหฺมจริยปริโยสานํ}

\gambhira{สฬายตนวคฺคสํยุตฺตปาฬิ}
\csromanmatikaanchor{mtk.111}
\csromantocmarkref{subsection}{3-6. ปฐมสมิทฺธิมารปญฺหาสุตฺตาทิ}{mtk.111}
\bookmark[dest={mtk.111},level=2]{3-6. ปฐมสมิทฺธิมารปญฺหาสุตฺตาทิ}
\prose{\csromanfolio{266}ทิฏฺเฐว ธมฺเม สยํ อภิญฺญา สจฺฉิกตฺวา อุปสมฺปชฺช วิหาสิ, “ขีณา ชาติ, วุสิตํ พฺรหฺมจริยํ, กตํ กรณียํ, นาปรํ อิตฺถตฺตายา”ติ อพฺภญฺญาสิ.\csromanspacer{} อญฺญตโร จ ปนายสฺมา มิคชาโล อรหตํ อโหสีติ.\csromanspacer{}\csromanspacer{}ทุติยํ.}

\csromanheader{2}{3. ปฐมสมิทฺธิมารปญฺหาสุตฺต}
\proseitem{65}{เอกํ สมยํ ภควา ราชคเห วิหรติ เวฬุวเน กลนฺทกนิวาเป.\csromanspacer{} อถ โข อายสฺมา สมิทฺธิ เยน ภควา ฯเปฯ ภควนฺตํ เอตทโวจ\csromandash{}“มาโร มาโร”ติ ภนฺเต วุจฺจติ, กิตฺตาวตา นุ โข ภนฺเต มาโร วา อสฺสมารปญฺญตฺติ วาติ.}

\prose{ยตฺถ โข สมิทฺธิ อตฺถิ จกฺขุ, อตฺถิ รูปา, อตฺถิ จกฺขุ วิญฺญาณํ, อตฺถิ จกฺขุวิญฺญาณวิญฺญาตพฺพา ธมฺมา, อตฺถิ ตตฺถ มาโร วา มารปญฺญตฺติ วา.\csromanspacer{} อตฺถิ โสตํ, อตฺถิ สทฺทา, อตฺถิ โสตวิญฺญาณํ, อตฺถิ โสตวิญฺญาณวิญฺญาตพฺพา ธมฺมา, อตฺถิ ตตฺถ มาโร วา มารปญฺญตฺติ วา.\csromanspacer{} อตฺถิ ฆานํ, อตฺถิ คนฺธา, อตฺถิ ฆานวิญฺญาณํ, อตฺถิ ฆานวิญฺญาณวิญฺญาตพฺพา ธมฺมา, อตฺถิ ตตฺถ มาโร วา มารปญฺญตฺติ วา.\csromanspacer{} อตฺถิ ชิวฺหา, อตฺถิ รสา, อตฺถิ ชิวฺหาวิญฺญาณํ, อตฺถิ ชิวฺหาวิญฺญาณวิญฺญาตพฺพา ธมฺมา, อตฺถิ ตตฺถ มาโร วา มารปญฺญตฺติ วา.\csromanspacer{} อตฺถิ กาโย, อตฺถิ โผฏฺฐพฺพา, อตฺถิ กายวิญฺญาณํ, อตฺถิ กายวิญฺญาณวิญฺญาตพฺพา ธมฺมา, อตฺถิ ตตฺถ มาโร วา มารปญฺญตฺติ วา.\csromanspacer{} อตฺถิ มโน, อตฺถิ ธมฺมา, อตฺถิ มโนวิญฺญาณํ, อตฺถิ มโนวิญฺญาณวิญฺญาตพฺพา ธมฺมา, อตฺถิ ตตฺถ มาโร วา มารปญฺญตฺติ วา.}

\prose{ยตฺถ จ โข สมิทฺธิ นตฺถิ จกฺขุ, นตฺถิ รูปา, นตฺถิ จกฺขุวิญฺญาณํ, นตฺถิ จกฺขุวิญฺญาณวิญฺญาตพฺพา ธมฺมา, นตฺถิ ตตฺถ มาโร วา มารปญฺญตฺติ วา.\csromanspacer{} นตฺถิ โสตํ ฯเปฯ.\csromanspacer{} นตฺถิ ฆานํ ฯเปฯ.\csromanspacer{} นตฺถิ ชิวฺหา, นตฺถิ รสา, นตฺถิ ชิวฺหาวิญฺญาณํ, นตฺถิ ชิวฺหาวิญฺญาณวิญฺญาตพฺพา ธมฺมา, นตฺถิ ตตฺถ มาโร วา มารปญฺญตฺติ วา.\csromanspacer{} นตฺถิ กาโย ฯเปฯ.\csromanspacer{} นตฺถิ มโน, นตฺถิ ธมฺมา, นตฺถิ มโนวิญฺญาณํ, นตฺถิ มโนวิญฺญาณวิญฺญาตพฺพา ธมฺมา, นตฺถิ ตตฺถ มาโร วา มารปญฺญตฺติ วาติ.\csromanspacer{}\csromanspacer{}ตติยํ.}

\csromanheader{2}{4. สมิทฺธิสตฺตปญฺหาสุตฺต}
\proseitem{66}{“สตฺโต สตฺโต”ติ ภนฺเต วุจฺจติ, กิตฺตาวตา นุ โข ภนฺเต สตฺโต วา อสฺส สตฺตปญฺญตฺติ วาติ ฯเปฯ.\csromanspacer{}\csromanspacer{}จตุตฺถํ.}

\csromanheader{2}{1. สฬายตนสํยุตฺต \textbf{5. สมิทฺธิทุกฺขปญฺหาสุตฺต}}
\proseitem{67}{\csromanfolio{267}“ทุกฺขํ ทุกฺขนฺ”ติ ภนฺเต วุจฺจติ, กิตฺตาวตา นุ โข ภนฺเต ทุกฺขํ วา อสฺส ทุกฺขปญฺญตฺติ วาติ ฯเปฯ.\csromanspacer{}\csromanspacer{}ปญฺจมํ.}

\csromanheader{2}{6. สมิทฺธิโลกปญฺหาสุตฺต}
\proseitem{68}{“โลโก โลโก”ติ ภนฺเต วุจฺจติ, กิตฺตาวตา นุ โข ภนฺเต โลโก วา อสฺส โลกปญฺญตฺติ วาติ.\csromanspacer{} ยตฺถ โข สมิทฺธิ อตฺถิ จกฺขุ, อตฺถิ รูปา, อตฺถิ จกฺขุวิญฺญาณํ, อตฺถิ จกฺขุวิญฺญาณวิญฺญาตพฺพา ธมฺมา, อตฺถิ ตตฺถ โลโก วา โลกปญฺญตฺติ วาติ ฯเปฯ อตฺถิ ชิวฺหา ฯเปฯ อตฺถิ มโน, อตฺถิ ธมฺมา, อตฺถิ มโนวิญฺญาณํ, อตฺถิ มโนวิญฺญาณวิญฺญาตพฺพา ธมฺมา, อตฺถิ ตตฺถ โลโก วา โลกปญฺญตฺติ วา.}

\prose{ยตฺถ จ โข สมิทฺธิ นตฺถิ จกฺขุ, นตฺถิ รูปา, นตฺถิ จกฺขุวิญฺญาณํ, นตฺถิ จกฺขุวิญฺญาณวิญฺญาตพฺพา ธมฺมา, นตฺถิ ตตฺถ โลโก วา โลกปญฺญตฺติ วา ฯเปฯ นตฺถิ ชิวฺหา ฯเปฯ นตฺถิ มโน, นตฺถิ ธมฺมา, นตฺถิ มโนวิญฺญาณํ, นตฺถิ มโนวิญฺญาณวิญฺญาตพฺพา ธมฺมา, นตฺถิ ตตฺถ โลโก วา โลกปญฺญตฺติ วาติ.\csromanspacer{}\csromanspacer{}ฉฏฺฐํ.}

\csromanmatikaanchor{mtk.112}
\csromantocmarkref{subsection}{7. อุปเสน-อาสีวิสสุตฺต}{mtk.112}
\bookmark[dest={mtk.112},level=2]{7. อุปเสน-อาสีวิสสุตฺต}
\csromanheader{2}{7. อุปเสน-อาสีวิสสุตฺต}
\proseitem{69}{เอกํ สมยํ อายสฺมา จ สาริปุตฺโต อายสฺมา จ อุปเสโน ราชคเห วิหรนฺติ สีตวเน สปฺปโสณฺฑิกปพฺภาเร.\csromanspacer{} เตน โข ปน สมเยน อายสฺมโต อุปเสนสฺส กาเย อาสีวิโส ปติโต โหติ.\csromanspacer{} อถ โข อายสฺมา อุปเสโน ภิกฺขู อามนฺเตสิ “เอถ เม อาวุโส อิมํ กายํ มญฺจกํ อาโรเปตฺวา พหิทฺธา นีหรถ, ปุรายํ กาโย อิเธว วิกิรติ เสยฺยถาปิ ภุสมุฏฺฐี”ติ.}

\prose{เอวํ วุตฺเต อายสฺมา สาริปุตฺโต อายสฺมนฺตํ อุปเสนํ เอตทโวจ “น โข ปน มยํ ปสฺสาม อายสฺมโต อุปเสนสฺส กายสฺส วา อญฺญถตฺตํ อินฺทฺริยานํ วา วิปริณามํ, อถ จ ปนายสฺมา อุปเสโน เอวมาห ‘เอถ เม อาวุโส อิมํ กายํ มญฺจกํ อาโรเปตฺวา พหิทฺธา นีหรถ, ปุรายํ กาโย อิเธว วิกิรติ เสยฺยถาปิ ภุสมุฏฺฐี’ติ”.\csromanspacer{} ยสฺส นูน อาวุโส สาริปุตฺต เอวมสฺส “อหํ จกฺขูติ วา มม จกฺขูติ วา ฯเปฯ อหํ ชิวฺหาติ วา}

\gambhira{สฬายตนวคฺคสํยุตฺตปาฬิ}
\prose{\csromanfolio{268}มม ชิวฺหาติ วา, อหํ มโนติ วา มม มโนติ วา”, ตสฺส อาวุโส สาริปุตฺต สิยา กายสฺส วา อญฺญถตฺตํ อินฺทฺริยานํ วา วิปริณาโม.\csromanspacer{} มยฺหญฺจ โข อาวุโส สาริปุตฺต น เอวํ โหติ “อหํ จกฺขูติ วา มม จกฺขูติ วา ฯเปฯ อหํ ชิวฺหาติ วา มม ชิวฺหาติ วา, อหํ มโนติ วา มม มโนติ วา”, ตสฺส มยฺหญฺจ โข อาวุโส สาริปุตฺต กิํ กายสฺส วา อญฺญถตฺตํ ภวิสฺสติ อินฺทฺริยานํ วา วิปริณาโมติ.}

\prose{ตถา หิ ปนายสฺมโต อุปเสนสฺส ทีฆรตฺตํ อหงฺการมมงฺการมานานุสโย สุสมูหโต, ตสฺมา อายสฺมโต อุปเสนสฺส น เอวํ โหติ “อหํ จกฺขูติ วา มม จกฺขูติ วา ฯเปฯ อหํ ชิวฺหาติ วา มม ชิวฺหาติ วา, อหํ มโนติ วา มม มโนติ วา”ติ.\csromanspacer{} อถ โข เต ภิกฺขู อายสฺมโต อุปเสนสฺส กายํ มญฺจกํ อาโรเปตฺวา พหิทฺธา นีหริํสุ.\csromanspacer{} อถ โข อายสฺมโต อุปเสนสฺส กาโย ตตฺเถว วิกิริ เสยฺยถาปิ ภุสมุฏฺฐีติ.\csromanspacer{}\csromanspacer{}สตฺตมํ.}

\csromanmatikaanchor{mtk.113}
\csromantocmarkref{subsection}{8. อุปวาณสนฺทิฏฺฐิกสุตฺต}{mtk.113}
\bookmark[dest={mtk.113},level=2]{8. อุปวาณสนฺทิฏฺฐิกสุตฺต}
\csromanheader{2}{8. อุปวาณสนฺทิฏฺฐิกสุตฺต}
\proseitem{70}{อถ โข อายสฺมา อุปวาโณ เยน ภควา เตนุปสงฺกมิ ฯเปฯ เอกมนฺตํ นิสินฺโน โข อายสฺมา อุปวาโณ ภควนฺตํ เอตทโวจ\csromandash{} “สนฺทิฏฺฐิโก ธมฺโม สนฺทิฏฺฐิโก ธมฺโม”ติ ภนฺเต วุจฺจติ, กิตฺตาวตา นุ โข ภนฺเต สนฺทิฏฺฐิโก ธมฺโม โหติ อกาลิโก เอหิปสฺสิโก โอปเนยฺยิโก ปจฺจตฺตํ เวทิตพฺโพ วิญฺญูหีติ.}

\prose{อิธ ปน อุปวาณ ภิกฺขุ จกฺขุนา รูปํ ทิสฺวา รูปปฺปฏิสํเวที จ โหติ รูปราคปฺปฏิสํเวที จ, สนฺตญฺจ อชฺฌตฺตํ รูเปสุ ราคํ “อตฺถิ เม อชฺฌตฺตํ รูเปสุ ราโค”ติ ปชานาติ, ยํ ตํ อุปวาณ ภิกฺขุ จกฺขุนา รูปํ ทิสฺวา รูปปฺปฏิสํเวที จ โหติ รูปราคปฺปฏิสํเวที จ, สนฺตญฺจ อชฺฌตฺตํ รูเปสุ ราคํ “อตฺถิ เม อชฺฌตฺตํ รูเปสุ ราโค”ติ ปชานาติ.\csromanspacer{} เอวมฺปิ โข อุปวาณ สนฺทิฏฺฐิโก ธมฺโม โหติ อกาลิโก เอหิปสฺสิโก โอปเนยฺยิโก ปจฺจตฺตํ เวทิตพฺโพ วิญฺญูหิ ฯเปฯ.}

\prose{ปุน จปรํ อุปวาณ ภิกฺขุ ชิวฺหาย รสํ สายิตฺวา รสปฺปฏิสํเวที จ โหติ รสราคปฺปฏิสํเวที จ, สนฺตญฺจ อชฺฌตฺตํ รเสสุ ราคํ “อตฺถิ เม อชฺฌตฺตํ รเสสุ ราโค”ติ ปชานาติ, ยํ ตํ อุปวาณ ภิกฺขุ ชิวฺหาย รสํ สายิตฺวา รสปฺปฏิสํเวที จ โหติ รสราคปฺปฏิสํเวที จ, สนฺตญฺจ อชฺฌตฺตํ}

\vspace{\tipitakaparskip}
\csromancenter{สฬายตนสํยุตฺต รเสสุ ราคํ “อตฺถิ เม อชฺฌตฺตํ รเสสุ ราโค”ติ ปชานาติ.\csromanspacer{} เอวมฺปิ โข อุปวาณ สนฺทิฏฺฐิโก ธมฺโม โหติ อกาลิโก เอหิปสฺสิโก โอปเนยฺยิโก ปจฺจตฺตํ เวทิตพฺโพ วิญฺญูหิ ฯเปฯ.}
\csromanmatikaanchor{mtk.114}
\csromantocmarkref{subsection}{9-11. ฉผสฺสายตนสุตฺต}{mtk.114}
\bookmark[dest={mtk.114},level=2]{9-11. ฉผสฺสายตนสุตฺต}
\prose{\csromanfolio{269}ปุน จปรํ อุปวาณ ภิกฺขุ มนสา ธมฺมํ วิญฺญาย ธมฺมปฺปฏิสํเวที จ โหติ ธมฺมราคปฺปฏิสํเวที จ, สนฺตญฺจ อชฺฌตฺตํ ธมฺเมสุ ราคํ “อตฺถิ เม อชฺฌตฺตํ ธมฺเมสุ ราโค”ติ ปชานาติ, ยํ ตํ อุปวาณ ภิกฺขุ มนสา ธมฺมํ วิญฺญาย ธมฺมปฺปฏิสํเวที จ โหติ ธมฺมราคปฺปฏิสํเวที จ, สนฺตญฺจ อชฺฌตฺตํ ธมฺเมสุ ราคํ “อตฺถิ เม อชฺฌตฺตํ ธมฺเมสุ ราโค”ติ ปชานาติ.\csromanspacer{} เอวมฺปิ โข อุปวาณ สนฺทิฏฺฐิโก ธมฺโม โหติ ฯเปฯ ปจฺจตฺตํ เวทิตพฺโพ วิญฺญูหิ ฯเปฯ.}

\prose{อิธ ปน อุปวาณ ภิกฺขุ จกฺขุนา รูปํ ทิสฺวา รูปปฺปฏิสํเวที จ โหติ โน จ รูปราคปฺปฏิสํเวที, อสนฺตญฺจ อชฺฌตฺตํ รูเปสุ ราคํ “นตฺถิ เม อชฺฌตฺตํ รูเปสุ ราโค”ติ ปชานาติ, ยํ ตํ อุปวาณ ภิกฺขุ จกฺขุนา รูปํ ทิสฺวา รูปปฺปฏิสํเวทีหิ โข โหติ โน จ รูปราคปฺปฏิสํเวที, อสนฺตญฺจ อชฺฌตฺตํ รูเปสุ ราคํ “นตฺถิ เม อชฺฌตฺตํ รูเปสุ ราโค”ติ ปชานาติ.\csromanspacer{} เอวมฺปิ โข อุปวาณ สนฺทิฏฺฐิโก ธมฺโม โหติ อกาลิโก เอหิปสฺสิโก โอปเนยฺยิโก ปจฺจตฺตํ เวทิตพฺโพ วิญฺญูหิ ฯเปฯ.}

\prose{ปุน จปรํ อุปวาณ ภิกฺขุ ชิวฺหาย รสํ สายิตฺวา รสปฺปฏิสํเวทีหิ โข โหติ โน จ รสราคปฺปฏิสํเวที, อสนฺตญฺจ อชฺฌตฺตํ รเสสุ ราคํ “นตฺถิ เม อชฺฌตฺตํ รเสสุ ราโค”ติ ปชานาติ ฯเปฯ.}

\prose{ปุน จปรํ อุปาวาณ ภิกฺขุ มนสา ธมฺมํ วิญฺญาย ธมฺมปฺปฏิสํเวทีหิ โข โหติ โน จ ธมฺมราคปฺปฏิสํเวที, อสนฺตญฺจ อชฺฌตฺตํ ธมฺเมสุ ราคํ “นตฺถิ เม อชฺฌตฺตํ ธมฺเมสุ ราโค”ติ ปชานาติ, ยํ ตํ อุปวาณ ภิกฺขุ มนสา ธมฺมํ วิญฺญาย ธมฺมปฺปฏิสํเวทีหิ โข โหติ โน จ ธมฺมราคปฺปฏิสํเวที, อสนฺตญฺจ อชฺฌตฺตํ ธมฺเมสุ ราคํ “นตฺถิ เม อชฺฌตฺตํ ธมฺเมสุ ราโค”ติ ปชานาติ.\csromanspacer{} เอวมฺปิ โข อุปวาณ สนฺทิฏฺฐิโก ธมฺโม โหติ, อกาลิโก เอหิปสฺสิโก โอปเนยฺยิโก ปจฺจตฺตํ เวทิตพฺโพ วิญฺญูหีติ.\csromanspacer{}\csromanspacer{}อฏฺฐมํ.}

\csromanheader{2}{9. ปฐมฉผสฺสายตนสุตฺต}
\proseitem{71}{โย หิ โกจิ ภิกฺขเว ภิกฺขุ ฉนฺนํ ผสฺสายตนานํ สมุทยญฺจ อตฺถงฺคมญฺจ อสฺสาทญฺจ อาทีนวญฺจ นิสฺสรณญฺจ ยถาภูตํ นปฺปชานาติ,}

\gambhira{สฬายตนวคฺคสํยุตฺตปาฬิ}
\prose{\csromanfolio{270}อวุสิตํ เตน พฺรหฺมจริยํ, อารกา โส อิมสฺมา ธมฺมวินยาติ.\csromanspacer{} เอวํ วุตฺเต อญฺญตโร ภิกฺขุ ภควนฺตํ เอตทโวจ\csromandash{}“เอตฺถาหํ ภนฺเต อนสฺสสํ\footnote{อนสฺสสิํ (สี), อนสฺสาสํ (สฺยา, กํ), อนสฺสาสิํ (อิ)}, อหํ หิ ภนฺเต ฉนฺนํ ผสฺสายตนานํ สมุทยญฺจ อตฺถงฺคมญฺจ อสฺสาทญฺจ อาทีนวญฺจ นิสฺสรณญฺจ ยถาภูตํ นปฺปชานามี”ติ.}

\prose{ตํ กิํ มญฺญสิ ภิกฺขุ, จกฺขุํ “เอตํ มม, เอโสหมสฺมิ, เอโส เม อตฺตา”ติ สมนุปสฺสสีติ.\csromanspacer{} โน เหตํ ภนฺเต.\csromanspacer{} สาธุ ภิกฺขุ เอตฺถ จ เต ภิกฺขุ จกฺขุ “เนตํ มม, เนโสหมสฺมิ, น เมโส อตฺตา”ติ เอวเมตํ ยถาภูตํ สมฺมปฺปญฺญาย สุทิฏฺฐํ ภวิสฺสติ, เอเสวนฺโต ทุกฺขสฺส ฯเปฯ ชิวฺหํ “เอตํ มม, เอโสหมสฺมิ, เอโส เม อตฺตา”ติ สมนุปสฺสสีติ.\csromanspacer{} โน เหตํ ภนฺเต.\csromanspacer{} สาธุ ภิกฺขุ เอตฺถ จ เต ภิกฺขุ ชิวฺหา “เนตํ มม, เนโสหมสฺมิ, น เมโส อตฺตา”ติ เอวเมตํ ยถาภูตํ สมฺมปฺปญฺญาย สุทิฏฺฐํ ภวิสฺสติ, เอเสวนฺโต ทุกฺขสฺส ฯเปฯ มนํ “เอตํ มม, เอโสหมสฺมิ, เอโส เม อตฺตา”ติ สมนุปสฺสสีติ.\csromanspacer{} โน เหตํ ภนฺเต.\csromanspacer{} สาธุ ภิกฺขุ เอตฺถ จ เต ภิกฺขุ มโน “เนตํ มม, เนโสหมสฺมิ, น เมโส อตฺตา”ติ เอวเมตํ ยถาภูตํ สมฺมปฺปญฺญาย สุทิฏฺฐํ ภวิสฺสติ, เอเสวนฺโต ทุกฺขสฺสาติ.\csromanspacer{}\csromanspacer{}นวมํ.}

\csromanheader{2}{10. ทุติยฉผสฺสายตนสุตฺต}
\proseitem{72}{โย หิ โกจิ ภิกฺขเว ภิกฺขุ ฉนฺนํ ผสฺสายตนานํ สมุทยญฺจ อตฺถงฺคมญฺจ อสฺสาทญฺจ อาทีนวญฺจ นิสฺสรณญฺจ ยถาภูตํ นปฺปชานาติ, อวุสิตํ เตน พฺรหฺมจริยํ, อารกา โส อิมสฺมา ธมฺมวินยาติ.\csromanspacer{} เอวํ วุตฺเต อญฺญตโร ภิกฺขุ ภควนฺตํ เอตทโวจ\csromandash{} “เอตฺถาหํ ภนฺเต อนสฺสสํ ปนสฺสสํ, อหํ หิ ภนฺเต ฉนฺนํ ผสฺสายตนานํ สมุทยญฺจ อตฺถงฺคมญฺจ อสฺสาทญฺจ อาทีนวญฺจ นิสฺสรณญฺจ ยถาภูตํ นปฺปชานามี”ติ.}

\prose{ตํ กิํ มญฺญสิ ภิกฺขุ, จกฺขุํ “เนตํ มม, เนโสหมสฺมิ, น เมโส อตฺตา”ติ สมนุปสฺสสีติ.\csromanspacer{} เอวํ ภนฺเต.\csromanspacer{} สาธุ ภิกฺขุ เอตฺถ จ เต ภิกฺขุ จกฺขุ “เนตํ มม, เนโสหมสฺมิ, น เมโส อตฺตา”ติ เอวเมตํ ยถาภูตํ สมฺมปฺปญฺญาย สุทิฏฺฐํ ภวิสฺสติ, เอวํ เต เอตํ ปฐมํ ผสฺสายตนํ ปหีนํ}

\vspace{\tipitakaparskip}
\csromancenter{สฬายตนสํยุตฺต ภวิสฺสติ อายติํ อปุนพฺภวาย ฯเปฯ ชิวฺหํ “เนตํ มม, เนโสหมสฺมิ, น เมโส อตฺตา”ติ สมนุปสฺสสีติ.\csromanspacer{} เอวํ ภนฺเต.\csromanspacer{} สาธุ ภิกฺขุ เอตฺถ จ เต ภิกฺขุ ชิวฺหา “เนตํ มม, เนโสหมสฺมิ, น เมโส อตฺตา”ติ เอวเมตํ ยถาภูตํ สมฺมปฺปญฺญาย สุทิฏฺฐํ ภวิสฺสติ, เอวํ เต เอตํ จตุตฺถํ ผสฺสายตนํ ปหีนํ ภวิสฺสติ อายติํ อปุนพฺภวาย ฯเปฯ มนํ “เนตํ มม, เนโสหมสฺมิ, น เมโส อตฺตา”ติ สมนุปสฺสสีติ.\csromanspacer{} เอวํ ภนฺเต.\csromanspacer{} สาธุ ภิกฺขุ เอตฺถ จ เต ภิกฺขุ มโน “เนตํ มม, เนโสหมสฺมิ, น เมโส อตฺตา”ติ เอวเมตํ ยถาภูตํ สมฺมปฺปญฺญาย สุทิฏฺฐํ ภวิสฺสติ, เอวํ เต เอตํ ฉฏฺฐํ ผสฺสายตนํ ปหีนํ ภวิสฺสติ อายติํ อปุนพฺภวายาติ.\csromanspacer{}\csromanspacer{}ทสมํ.}
\csromanheader{2}{11. ตติยฉผสฺสายตนสุตฺต}
\proseitem{73}{\csromanfolio{271}โย หิ โกจิ ภิกฺขเว ภิกฺขุ ฉนฺนํ ผสฺสายตนานํ สมุทยญฺจ อตฺถงฺคมญฺจ อสฺสาทญฺจ อาทีนวญฺจ นิสฺสรณญฺจ ยถาภูตํ นปฺปชานาติ, อวุสิตํ เตน พฺรหฺมจริยํ, อารกา โส อิมสฺมา ธมฺมวินยาติ.\csromanspacer{} เอวํ วุตฺเต อญฺญตโร ภิกฺขุ ภควนฺตํ เอตทโวจ\csromandash{} “เอตฺถาหํ ภนฺเต อนสฺสสํ ปนสฺสสํ, อหํ หิ ภนฺเต ฉนฺนํ ผสฺสายตนานํ สมุทยญฺจ อตฺถงฺคมญฺจ อสฺสาทญฺจ อาทีนวญฺจ นิสฺสรณญฺจ ยถาภูตํ นปฺปชานามี”ติ.}

\prose{ตํ กิํ มญฺญสิ ภิกฺขุ, จกฺขุ นิจฺจํ วา อนิจฺจํ วาติ.\csromanspacer{} อนิจฺจํ ภนฺเต.\csromanspacer{} ยํ ปนานิจฺจํ, ทุกฺขํ วา ตํ สุขํ วาติ.\csromanspacer{} ทุกฺขํ ภนฺเต.\csromanspacer{} ยํ ปนานิจฺจํ ทุกฺขํ วิปริณามธมฺมํ, กลฺลํ นุ ตํ สมนุปสฺสิตุํ “เอตํ มม, เอโสหมสฺมิ, เอโส เม อตฺตา”ติ.\csromanspacer{} โน เหตํ ภนฺเต.\csromanspacer{} โสตํ.\csromanspacer{} ฆานํ.\csromanspacer{} ชิวฺหา.\csromanspacer{} กาโย.\csromanspacer{} มโน นิจฺโจ วา อนิจฺโจ วาติ.\csromanspacer{} อนิจฺโจ ภนฺเต.\csromanspacer{} ยํ ปนานิจฺจํ, ทุกฺขํ วา ตํ สุขํ วาติ.\csromanspacer{} ทุกฺขํ ภนฺเต.\csromanspacer{} ยํ ปนานิจฺจํ ทุกฺขํ วิปริณามธมฺมํ, กลฺลํ นุ ตํ สมนุปสฺสิตุํ “เอตํ มม, เอโสหมสฺมิ, เอโส เม อตฺตา”ติ.\csromanspacer{} โน เหตํ ภนฺเต.\csromanspacer{} เอวํ ปสฺสํ ภิกฺขุ สุตวา อริยสาวโก จกฺขุสฺมิมฺปิ นิพฺพินฺทติ, โสตสฺมิมฺปิ นิพฺพินฺทติ, ฆานสฺมิมฺปิ นิพฺพินฺทติ, ชิวฺหายปิ นิพฺพินฺทติ, กายสฺมิมฺปิ นิพฺพินฺทติ, มนสฺมิมฺปิ นิพฺพินฺทติ, นิพฺพินฺทํ วิรชฺชติ, วิราคา วิมุจฺจติ, วิมุตฺตสฺมิํ “วิมุตฺตมฺ”อิติ ญาณํ โหติ, “ขีณา ชาติ, วุสิตํ พฺรหฺมจริยํ, กตํ กรณียํ, นาปรํ อิตฺถตฺตายา”ติ ปชานาตีติ.\csromanspacer{}\csromanspacer{}เอกาทสมํ.}

\vspace{\tipitakaparskip}
\csromancenter{มิคชาลวคฺโค ทุติโย.}
\gambhira{สฬายตนวคฺคสํยุตฺตปาฬิ}
\titlehead{\textbf{ตสฺสุทฺทานํ}}
\csromanmatikaanchor{mtk.115}
\csromanmatikaanchor{mtk.116}
\csromantocmarknopage{section}{(8) 3. คิลานวคฺค}{mtk.115}
\bookmark[dest={mtk.115},level=1]{(8) 3. คิลานวคฺค}
\csromantocmarkref{subsection}{1-2. คิลานสุตฺต}{mtk.116}
\bookmark[dest={mtk.116},level=2]{1-2. คิลานสุตฺต}
\csromangathameasure{มิคชาเลน เทฺว วุตฺตา, จตฺตาโร จ สมิทฺธินา. \\ อุปเสโน อุปวาโณ, ฉผสฺสายตนิกา ตโยติ.}
\csromangathagroup{\csromangathastanza{\csromanfolio{272}มิคชาเลน เทฺว วุตฺตา, จตฺตาโร จ สมิทฺธินา. \\ อุปเสโน อุปวาโณ, ฉผสฺสายตนิกา ตโยติ.}}

\csromanheader{1}{\textbf{(8) 3. คิลานวคฺค}}
\csromanheader{2}{1. ปฐมคิลานสุตฺต}
\proseitem{74}{สาวตฺถินิทานํ.\csromanspacer{} อถ โข อญฺญตโร ภิกฺขุ เยน ภควา เตนุปสงฺกมิ ฯเปฯ เอกมนฺตํ นิสินฺโน โข โส ภิกฺขุ ภควนฺตํ เอตทโวจ\csromandash{}“อมุกสฺมิํ ภนฺเต วิหาเร อญฺญตโร ภิกฺขุ นโว อปฺปญฺญาโต อาพาธิโก ทุกฺขิโต พาฬฺหคิลาโน, สาธุ ภนฺเต ภควา, เยน โส ภิกฺขุ เตนุปสงฺกมตุ อนุกมฺปํ อุปาทายา”ติ.}

\prose{อถ โข ภควา นววาทญฺจ สุตฺวา คิลานวาทญฺจ “อปฺปญฺญาโต ภิกฺขู”ติ อิติ วิทิตฺวา เยน โส ภิกฺขุ เตนุปสงฺกมิ.\csromanspacer{} อทฺทสา โข โส ภิกฺขุ ภควนฺตํ ทูรโตว อาคจฺฉนฺตํ, ทิสฺวาน มญฺจเก สมโธสิ\footnote{สมญฺโจสิ (สี), สมเตสิ (สฺยา, กํ), สมญฺโจปิ (อิ)}.\csromanspacer{} อถ โข ภควา ตํ ภิกฺขุํ เอตทโวจ “อลํ ภิกฺขุ มา ตฺวํ มญฺจเก สมโธสิ, สนฺติมานิ อาสนานิ ปญฺญตฺตานิ, ตตฺถาหํ นิสีทิสฺสามี”ติ.\csromanspacer{} นิสีทิ ภควา ปญฺญตฺเต อาสเน, นิสชฺช โข ภควา ตํ ภิกฺขุํ เอตทโวจ “กจฺจิ เต ภิกฺขุ ขมนียํ, กจฺจิ ยาปนียํ, กจฺจิ ทุกฺขา เวทนา ปฏิกฺกมนฺติ, โน อภิกฺกมนฺติ, ปฏิกฺกโมสานํ ปญฺญายติ, โน อภิกฺกโม”ติ.\csromanspacer{} น เม ภนฺเต ขมนียํ, น ยาปนียํ, พาฬฺหา เม ทุกฺขา เวทนา อภิกฺกมนฺติ, โน ปฏิกฺกมนฺติ, อภิกฺกโมสานํ ปญฺญายติ, โน ปฏิกฺกโมติ.\csromanspacer{} กจฺจิ เต ภิกฺขุ น กิญฺจิ กุกฺกุจฺจํ น โกจิ วิปฺปฏิสาโรติ.\csromanspacer{} ตคฺฆ เม ภนฺเต อนปฺปกํ กุกฺกุจฺจํ, อนปฺปโก วิปฺปฏิสาโรติ.\csromanspacer{} กจฺจิ ปน ตํ\footnote{ตฺวํ (สี), เต (สฺยา, กํ, ก)} ภิกฺขุ อตฺตา สีลโต อุปวทตีติ.\csromanspacer{} น โข มํ ภนฺเต อตฺตา สีลโต อุปวทตีติ\footnote{โน เหตํ ภนฺเต (อิ, ก)}.\csromanspacer{} โน เจ กิร เต ภิกฺขุ อตฺตา สีลโต อุปวทติ, อถ กิญฺจ\footnote{อถ กิสฺมิญฺจ (สี), อถ ภิกฺขุ กิสฺมิญฺจ (สฺยา, กํ, อิ, ก)} เต กุกฺกุจฺจํ, โก จ วิปฺปฏิสาโรติ.\csromanspacer{} น ขฺวาหํ ภนฺเต สีลวิสุทฺธตฺถํ ภควตา ธมฺมํ เทสิตํ}

\vspace{\tipitakaparskip}
\csromancenter{สฬายตนสํยุตฺต อาชานามีติ.\csromanspacer{} โน เจ กิร ตฺวํ ภิกฺขุ สีลวิสุทฺธตฺถํ มยา ธมฺมํ เทสิตํ อาชานาสิ, อถ กิมตฺถํ จรหิ ตฺวํ ภิกฺขุ มยา ธมฺมํ เทสิตํ อาชานาสีติ.\csromanspacer{} ราควิราคตฺถํ ขฺวาหํ ภนฺเต ภควตา ธมฺมํ เทสิตํ อาชานามีติ.}
\prose{\csromanfolio{273}สาธุ สาธุ ภิกฺขุ, สาธุ โข ตฺวํ ภิกฺขุ ราควิราคตฺถํ มยา ธมฺมํ เทสิตํ อาชานาสิ, ราควิราคตฺโถ หิ ภิกฺขุ มยา ธมฺโม เทสิโต.\csromanspacer{} ตํ กิํ มญฺญสิ ภิกฺขุ, จกฺขุ นิจฺจํ วา อนิจฺจํ วาติ.\csromanspacer{} อนิจฺจํ ภนฺเต ฯเปฯ.\csromanspacer{} โสตํ.\csromanspacer{} ฆานํ.\csromanspacer{} ชิวฺหา.\csromanspacer{} กาโย.\csromanspacer{} มโน นิจฺโจ วา อนิจฺโจ วาติ.\csromanspacer{} อนิจฺโจ ภนฺเต.\csromanspacer{} ยํ ปนานิจฺจํ, ทุกฺขํ วา ตํ สุขํ วาติ.\csromanspacer{} ทุกฺขํ ภนฺเต.\csromanspacer{} ยํ ปนานิจฺจํ ทุกฺขํ วิปริณามธมฺมํ, กลฺลํ นุ ตํ สมนุปสฺสิตุํ “เอตํ มม, เอโสหมสฺมิ, เอโส เม อตฺตา”ติ.\csromanspacer{} โน เหตํ ภนฺเต.\csromanspacer{} เอวํ ปสฺสํ ภิกฺขุ สุตวา อริยสาวโก จกฺขุสฺมิมฺปิ นิพฺพินฺทติ, โสตสฺมิมฺปิ นิพฺพินฺทติ ฯเปฯ มนสฺมิมฺปิ นิพฺพินฺทติ, นิพฺพินฺทํ วิรชฺชติ, วิราคา วิมุจฺจติ, วิมุตฺตสฺมิํ “วิมุตฺตมฺ”อิติ ญาณํ โหติ, “ขีณา ชาติ ฯเปฯ นาปรํ อิตฺถตฺตายา”ติ ปชานาตีติ.\csromanspacer{} อิทมโวจ ภควา.\csromanspacer{} อตฺตมโน โส ภิกฺขุ ภควโต ภาสิตํ อภินนฺทิ.\csromanspacer{} อิมสฺมิํ จ ปน เวยฺยากรณสฺมิํ ภญฺญมาเน ตสฺส ภิกฺขุโน วิรชํ วีตมลํ ธมฺมจกฺขุํ อุทปาทิ “ยํ กิญฺจิ สมุทยธมฺมํ, สพฺพํ ตํ นิโรธธมฺมนฺ”ติ.\csromanspacer{}\csromanspacer{}ปฐมํ.}

\csromanheader{2}{2. ทุติยคิลานสุตฺต}
\proseitem{75}{อถ โข อญฺญตโร ภิกฺขุ ฯเปฯ ภควนฺตํ เอตทโวจ\csromandash{} อมุกสฺมิํ ภนฺเต วิหาเร อญฺญตโร ภิกฺขุ นโว อปฺปญฺญาโต อาพาธิโก ทุกฺขิโต พาฬฺหคิลาโน, สาธุ ภนฺเต ภควา, เยน โส ภิกฺขุ เตนุปสงฺกมตุ อนุกมฺปํ อุปาทายาติ.}

\prose{อถ โข ภควา นววาทญฺจ สุตฺวา คิลานวาทญฺจ “อปฺปญฺญาโต ภิกฺขู”ติ อิติ วิทิตฺวา เยน โส ภิกฺขุ เตนุปสงฺกมิ, อทฺทสา โข โส ภิกฺขุ ภควนฺตํ ทูรโตว อาคจฺฉนฺตํ, ทิสฺวาน มญฺจเก สมโธสิ.\csromanspacer{} อถ โข ภควา ตํ ภิกฺขุํ เอตทโวจ “อลํ ภิกฺขุ มา ตฺวํ มญฺจเก สมโธสิ, สนฺติมานิ อาสนานิ ปญฺญตฺตานิ, ตตฺถาหํ นิสีทิสฺสามี”ติ.\csromanspacer{} นิสีทิ ภควา ปญฺญตฺเต อาสเน, นิสชฺช โข ภควา ตํ ภิกฺขุํ เอตทโวจ “กจฺจิ เต ภิกฺขุ ขมนียํ, กจฺจิ ยาปนียํ, กจฺจิ ทุกฺขา เวทนา ปฏิกฺกมนฺติ, โน อภิกฺกมนฺติ,}

\gambhira{สฬายตนวคฺคสํยุตฺตปาฬิ}
\prose{\csromanfolio{274}ปฏิกฺกโมสานํ ปญฺญายติ, โน อภิกฺกโม”ติ.\csromanspacer{} น เม ภนฺเต ขมนียํ, น ยาปนียํ ฯเปฯ.\csromanspacer{} น โข มํ\footnote{เม (สพฺพตฺถ)} ภนฺเต อตฺตา สีลโต อุปวทตีติ.\csromanspacer{} โน เจ กิร เต ภิกฺขุ อตฺตา สีลโต อุปวทติ.\csromanspacer{} อถ กิญฺจ เต กุกฺกุจฺจํ, โก จ วิปฺปฏิสาโรติ.\csromanspacer{} น ขฺวาหํ ภนฺเต สีลวิสุทฺธตฺถํ ภควตา ธมฺมํ เทสิตํ อาชานามีติ.\csromanspacer{} โน เจ กิร ตฺวํ ภิกฺขุ สีลวิสุทฺธตฺถํ มยา ธมฺมํ เทสิตํ อาชานาสิ, อถ กิมตฺถํ จรหิ ตฺวํ ภิกฺขุ มยา ธมฺมํ เทสิตํ อาชานาสีติ.\csromanspacer{} อนุปาทาปรินิพฺพานตฺถํ ขฺวาหํ ภนฺเต ภควตา ธมฺมํ เทสิตํ อาชานามีติ.}

\prose{สาธุ สาธุ ภิกฺขุ, สาธุ โข ตฺวํ ภิกฺขุ อนุปาทาปรินิพฺพานตฺถํ มยา ธมฺมํ เทสิตํ อาชานาสิ, อนุปาทาปรินิพฺพานตฺโถ หิ ภิกฺขุ มยา ธมฺโม เทสิโต.\csromanspacer{} ตํ กิํ มญฺญสิ ภิกฺขุ, จกฺขุ นิจฺจํ วา อนิจฺจํ วาติ.\csromanspacer{} อนิจฺจํ ภนฺเต ฯเปฯ.\csromanspacer{} โสตํ.\csromanspacer{} ฆานํ.\csromanspacer{} ชิวฺหา.\csromanspacer{} กาโย.\csromanspacer{} มโน.\csromanspacer{} มโนวิญฺญาณํ.\csromanspacer{} มโนสมฺผสฺโส.\csromanspacer{} ยมฺปิทํ มโนสมฺผสฺสปจฺจยา อุปฺปชฺชติ เวทยิตํ สุขํ วา ทุกฺขํ วา อทุกฺขมสุขํ วา, ตมฺปิ นิจฺจํ วา อนิจฺจํ วาติ.\csromanspacer{} อนิจฺจํ ภนฺเต.\csromanspacer{} ยํ ปนานิจฺจํ, ทุกฺขํ วา ตํ สุขํ วาติ.\csromanspacer{} ทุกฺขํ ภนฺเต.\csromanspacer{} ยํ ปนานิจฺจํ ทุกฺขํ วิปริณามธมฺมํ, กลฺลํ นุ ตํ สมนุปสฺสิตุํ “เอตํ มม.\csromanspacer{} เอโสหมสฺมิ, เอโส เม อตฺตา”ติ.\csromanspacer{} โน เหตํ ภนฺเต.\csromanspacer{} เอวํ ปสฺสํ ภิกฺขุ สุตวา อริยสาวโก จกฺขุสฺมิมฺปิ นิพฺพินฺทติ ฯเปฯ มนสฺมิมฺปิ.\csromanspacer{} มโนวิญฺญาเณปิ.\csromanspacer{} มโนสมฺผสฺเสปิ นิพฺพินฺทติ.\csromanspacer{} ยมฺปิทํ มโนสมฺผสฺสปจฺจยา อุปฺปชฺชติ เวทยิตํ สุขํ วา ทุกฺขํ วา อทุกฺขมสุขํ วา, ตสฺมิมฺปิ นิพฺพินฺทติ, นิพฺพินฺทํ วิรชฺชติ, วิราคา วิมุจฺจติ, วิมุตฺตสฺมิํ “วิมุตฺตมฺ”อิติ ญาณํ โหติ, “ขีณา ชาติ, วุสิตํ พฺรหฺมจริยํ, กตํ กรณียํ, นาปรํ อิตฺถตฺตายา”ติ ปชานาตีติ.\csromanspacer{} อิทมโวจ ภควา.\csromanspacer{} อตฺตมโน โส ภิกฺขุ ภควโต ภาสิตํ อภินนฺทิ.\csromanspacer{} อิมสฺมิํ จ ปน เวยฺยากรณสฺมิํ ภญฺญมาเน ตสฺส ภิกฺขุสฺส อนุปาทาย อาสเวหิ จิตฺตํ วิมุจฺจีติ\footnote{วิมุจฺจตีติ (สพฺพตฺถ)}.\csromanspacer{}\csromanspacer{}ทุติยํ.}

\csromanmatikaanchor{mtk.117}
\csromantocmarkref{subsection}{3-5. ราธ-อนิจฺจสุตฺตาทิ}{mtk.117}
\bookmark[dest={mtk.117},level=2]{3-5. ราธ-อนิจฺจสุตฺตาทิ}
\csromanheader{2}{3. ราธ-อนิจฺจสุตฺต}
\proseitem{76}{อถ โข อายสฺมา ราโธ ฯเปฯ เอกมนฺตํ นิสินฺโน โข อายสฺมา ราโธ ภควนฺตํ เอตทโวจ “สาธุ เม ภนฺเต ภควา สํขิตฺเตน ธมฺมํ}

\vspace{\tipitakaparskip}
\csromancenter{สฬายตนสํยุตฺต เทเสตุ, ยมหํ ภควโต ธมฺมํ สุตฺวา เอโก วูปกฏฺโฐ อปฺปมตฺโต อาตาปี ปหิตตฺโต วิหเรยฺยนฺ”ติ.\csromanspacer{} ยํ โข ราธ อนิจฺจํ, ตตฺร เต ฉนฺโท ปหาตพฺโพ.\csromanspacer{} กิญฺจ ราธ อนิจฺจํ, ตตฺร เต ฉนฺโท ปหาตพฺโพ.\csromanspacer{} จกฺขุ อนิจฺจํ, รูปา อนิจฺจา, จกฺขุวิญฺญาณํ.\csromanspacer{} จกฺขุสมฺผสฺโส.\csromanspacer{} ยมฺปิทํ จกฺขุสมฺผสฺสปจฺจยา อุปฺปชฺชติ เวทยิตํ สุขํ วา ทุกฺขํ วา อทุกฺขมสุขํ วา, ตมฺปิ อนิจฺจํ, ตตฺร เต ฉนฺโท ปหาตพฺโพ ฯเปฯ.\csromanspacer{} ชิวฺหา.\csromanspacer{} กาโย.\csromanspacer{} มโน อนิจฺโจ, ตตฺร เต ฉนฺโท ปหาตพฺโพ.\csromanspacer{} ธมฺมา.\csromanspacer{} มโนวิญฺญาณํ.\csromanspacer{} มโนสมฺผสฺโส.\csromanspacer{} ยมฺปิทํ มโนสมฺผสฺสปจฺจยา อุปฺปชฺชติ เวทยิตํ สุขํ วา ทุกฺขํ วา อทุกฺขมสุขํ วา, ตมฺปิ อนิจฺจํ, ตตฺร เต ฉนฺโท ปหาตพฺโพ.\csromanspacer{} ยํ โข ราธ อนิจฺจํ, ตตฺร เต ฉนฺโท ปหาตพฺโพติ.\csromanspacer{}\csromanspacer{}ตติยํ.}
\csromanheader{2}{4. ราธทุกฺขสุตฺต}
\proseitem{77}{\csromanfolio{275}ยํ โข ราธ ทุกฺขํ, ตตฺร เต ฉนฺโท ปหาตพฺโพ.\csromanspacer{} กิญฺจ ราธ ทุกฺขํ.\csromanspacer{} จกฺขุ โข ราธ ทุกฺขํ, ตตฺร เต ฉนฺโท ปหาตพฺโพ.\csromanspacer{} รูปา.\csromanspacer{} จกฺขุวิญฺญาณํ.\csromanspacer{} จกฺขุสมฺผสฺโส.\csromanspacer{} ยมฺปิทํ จกฺขุสมฺผสฺส ฯเปฯ.\csromanspacer{} อทุกฺขมสุขํ วา, ตมฺปิ ทุกฺขํ, ตตฺร เต ฉนฺโท ปหาตพฺโพ ฯเปฯ.\csromanspacer{} มโน ทุกฺโข.\csromanspacer{} ธมฺมา.\csromanspacer{} มโนวิญฺญาณํ.\csromanspacer{} มโนสมฺผสฺโส.\csromanspacer{} ยมฺปิทํ มโนสมฺผสฺสปจฺจยา อุปฺปชฺชติ เวทยิตํ สุขํ วา ทุกฺขํ วา อทุกฺขมสุขํ วา, ตมฺปิ ทุกฺขํ, ตตฺร เต ฉนฺโท ปหาตพฺโพ.\csromanspacer{} ยํ โข ราธ ทุกฺขํ, ตตฺร เต ฉนฺโท ปหาตพฺโพติ.\csromanspacer{}\csromanspacer{}จตุตฺถํ.}

\csromanheader{2}{5. ราธ-อนตฺตสุตฺต}
\proseitem{78}{โย โข ราธ อนตฺตา, ตตฺร เต ฉนฺโท ปหาตพฺโพ.\csromanspacer{} โก จ ราธ อนตฺตา.\csromanspacer{} จกฺขุ โข ราธ อนตฺตา, ตตฺร เต ฉนฺโท ปหาตพฺโพ.\csromanspacer{} รูปา.\csromanspacer{} จกฺขุวิญฺญาณํ.\csromanspacer{} จกฺขุสมฺผสฺโส.\csromanspacer{} ยมฺปิทํ จกฺขุสมฺผสฺสปจฺจยา ฯเปฯ.\csromanspacer{} มโน อนตฺตา.\csromanspacer{} ธมฺมา.\csromanspacer{} มโนวิญฺญาณํ.\csromanspacer{} มโนสมฺผสฺโส.\csromanspacer{} ยมฺปิทํ มโนสมฺผสฺสปจฺจยา อุปฺปชฺชติ เวทยิตํ สุขํ วา ทุกฺขํ วา อทุกฺขมสุขํ วา, ตมฺปิ อนตฺตา, ตตฺร เต ฉนฺโท ปหาตพฺโพ.\csromanspacer{} โย โข ราธ อนตฺตา, ตตฺร เต ฉนฺโท ปหาตพฺโพติ.\csromanspacer{}\csromanspacer{}ปญฺจมํ.}

\gambhira{สฬายตนวคฺคสํยุตฺตปาฬิ}
\csromanmatikaanchor{mtk.118}
\csromantocmarkref{subsection}{6-7. อวิชฺชาปหานสุตฺต}{mtk.118}
\bookmark[dest={mtk.118},level=2]{6-7. อวิชฺชาปหานสุตฺต}
\csromanheader{2}{6. ปฐม-อวิชฺชาปหานสุตฺต}
\proseitem{79}{\csromanfolio{276}อถ โข อญฺญตโร ภิกฺขุ เยน ภควา เตนุปสงฺกมิ ฯเปฯ เอกมนฺตํ นิสินฺโน โข โส ภิกฺขุ ภควนฺตํ เอตทโวจ “อตฺถิ นุ โข ภนฺเต เอโก ธมฺโม, ยสฺส ปหานา ภิกฺขุโน อวิชฺชา ปหียติ, วิชฺชา อุปฺปชฺชตี”ติ.\csromanspacer{} อตฺถิ โข ภิกฺขุ เอโก ธมฺโม, ยสฺส ปหานา ภิกฺขุโน อวิชฺชา ปหียติ, วิชฺชา อุปฺปชฺชตีติ.\csromanspacer{} กตโม ปน ภนฺเต เอโก ธมฺโม, ยสฺส ปหานา ภิกฺขุโน อวิชฺชา ปหียติ, วิชฺชา อุปฺปชฺชตีติ.\csromanspacer{} อวิชฺชา โข ภิกฺขุ เอโก ธมฺโม, ยสฺส ปหานา ภิกฺขุโน อวิชฺชา ปหียติ, วิชฺชา อุปฺปชฺชตีติ.\csromanspacer{} กถํ ปน ภนฺเต ชานโต กถํ ปสฺสโต ภิกฺขุโน อวิชฺชา ปหียติ, วิชฺชา อุปฺปชฺชตีติ.\csromanspacer{} จกฺขุํ โข ภิกฺขุ อนิจฺจโต ชานโต ปสฺสโต ภิกฺขุโน อวิชฺชา ปหียติ, วิชฺชา อุปฺปชฺชติ.\csromanspacer{} รูเป.\csromanspacer{} จกฺขุวิญฺญาณํ.\csromanspacer{} จกฺขุสมฺผสฺสํ.\csromanspacer{} ยมฺปิทํ จกฺขุสมฺผสฺสปจฺจยา อุปฺปชฺชติ เวทยิตํ สุขํ วา ทุกฺขํ วา อทุกฺขมสุขํ วา, ตมฺปิ อนิจฺจโต ชานโต ปสฺสโต ภิกฺขุโน อวิชฺชา ปหียติ, วิชฺชา อุปฺปชฺชติ ฯเปฯ.\csromanspacer{} มนํ อนิจฺจโต ชานโต ปสฺสโต ภิกฺขุโน อวิชฺชา ปหียติ, วิชฺชา อุปฺปชฺชติ.\csromanspacer{} ธมฺเม.\csromanspacer{} มโนวิญฺญาณํ.\csromanspacer{} มโนสมฺผสฺสํ.\csromanspacer{} ยมฺปิทํ มโนสมฺผสฺสปจฺจยา อุปฺปชฺชติ เวทยิตํ สุขํ วา ทุกฺขํ วา อทุกฺขมสุขํ วา, ตมฺปิ อนิจฺจโต ชานโต ปสฺสโต ภิกฺขุโน อวิชฺชา ปหียติ, วิชฺชา อุปฺปชฺชติ.\csromanspacer{} เอวํ โข ภิกฺขุ ชานโต เอวํ ปสฺสโต ภิกฺขุโน อวิชฺชา ปหียติ, วิชฺชา อุปฺปชฺชตีติ.\csromanspacer{}\csromanspacer{}ฉฏฺฐํ.}

\csromanheader{2}{7. ทุติย-อวิชฺชาปหานสุตฺต}
\proseitem{80}{อถ โข อญฺญตโร ภิกฺขุ ฯเปฯ เอตทโวจ “อตฺถิ นุ โข ภนฺเต เอโก ธมฺโม, ยสฺส ปหานา ภิกฺขุโน อวิชฺชา ปหียติ, วิชฺชา อุปฺปชฺชตี”ติ.\csromanspacer{} อตฺถิ โข ภิกฺขุ เอโก ธมฺโม, ยสฺส ปหานา ภิกฺขุโน อวิชฺชา ปหียติ, วิชฺชา อุปฺปชฺชตีติ.\csromanspacer{} กตโม ปน ภนฺเต เอโก ธมฺโม, ยสฺส ปหานา ภิกฺขุโน อวิชฺชา ปหียติ, วิชฺชา อุปฺปชฺชตีติ.\csromanspacer{} อวิชฺชา โข ภิกฺขุ เอโก ธมฺโม, ยสฺส ปหานา ภิกฺขุโน อวิชฺชา ปหียติ, วิชฺชา อุปฺปชฺชตีติ.\csromanspacer{} กถํ ปน ภนฺเต ชานโต กถํ ปสฺสโต อวิชฺชา ปหียติ, วิชฺชา อุปฺปชฺชตีติ.\csromanspacer{} อิธ ภิกฺขุ ภิกฺขุโน สุตํ โหติ “สพฺเพ ธมฺมา นาลํ อภินิเวสายา”ติ.\csromanspacer{} เอวญฺเจตํ ภิกฺขุ ภิกฺขุโน สุตํ โหติ “สพฺเพ}

\vspace{\tipitakaparskip}
\csromancenter{สฬายตนสํยุตฺต ธมฺมา นาลํ อภินิเวสายา”ติ.\csromanspacer{} โส สพฺพํ ธมฺมํ อภิชานาติ, สพฺพํ ธมฺมํ อภิญฺญาย สพฺพํ ธมฺมํ ปริชานาติ, สพฺพํ ธมฺมํ ปริญฺญาย สพฺพนิมิตฺตานิ อญฺญโต ปสฺสติ.\csromanspacer{} จกฺขุํ อญฺญโต ปสฺสติ.\csromanspacer{} รูเป.\csromanspacer{} จกฺขุวิญฺญาณํ.\csromanspacer{} จกฺขุสมฺผสฺสํ.\csromanspacer{} ยมฺปิทํ จกฺขุสมฺผสฺสปจฺจยา อุปฺปชฺชติ เวทยิตํ สุขํ วา ทุกฺขํ วา อทุกฺขมสุขํ วา, ตมฺปิ อญฺญโต ปสฺสติ ฯเปฯ.\csromanspacer{} มนํ อญฺญโต ปสฺสติ.\csromanspacer{} ธมฺเม.\csromanspacer{} มโนวิญฺญาณํ.\csromanspacer{} มโนสมฺผสฺสํ.\csromanspacer{} ยมฺปิทํ มโนสมฺผสฺสปจฺจยา อุปฺปชฺชติ เวทยิตํ สุขํ วา ทุกฺขํ วา อทุกฺขมสุขํ วา, ตมฺปิ อญฺญโต ปสฺสติ.\csromanspacer{} เอวํ โข ภิกฺขุ ชานโต เอวํ ปสฺสโต ภิกฺขุโน อวิชฺชา ปหียติ, วิชฺชา อุปฺปชฺชตีติ.\csromanspacer{}\csromanspacer{}สตฺตมํ.}
\csromanmatikaanchor{mtk.119}
\csromantocmarkref{subsection}{8. สมฺพหุลภิกฺขุสุตฺต}{mtk.119}
\bookmark[dest={mtk.119},level=2]{8. สมฺพหุลภิกฺขุสุตฺต}
\csromanheader{2}{8. สมฺพหุลภิกฺขุสุตฺต}
\proseitem{81}{\csromanfolio{277}อถ โข สมฺพหุลา ภิกฺขู เยน ภควา เตนุปสงฺกมิํสุ ฯเปฯ เอกมนฺตํ นิสินฺนา โข เต ภิกฺขู ภควนฺตํ เอตทโวจุํ “อิธ โน ภนฺเต อญฺญติตฺถิยา ปริพฺพาชกา อเมฺห เอวํ ปุจฺฉนฺติ ‘กิมตฺถิยํ อาวุโส สมเณ โคตเม พฺรหฺมจริยํ วุสฺสตี”ติ, เอวํ ปุฏฺฐา มยํ ภนฺเต เตสํ อญฺญติตฺถิยานํ ปริพฺพาชกานํ เอวํ พฺยากโรม, ทุกฺขสฺส โข อาวุโส ปริญฺญตฺถํ ภควติ พฺรหฺมจริยํ วุสฺสตี”ติ.\csromanspacer{} กจฺจิ มยํ ภนฺเต เอวํ ปุฏฺฐา เอวํ พฺยากรมานา วุตฺตวาทิโน เจว ภควโต โหม, น จ ภควนฺตํ อภูเตน อพฺภาจิกฺขาม, ธมฺมสฺส จานุธมฺมํ พฺยากโรม, น จ โกจิ สหธมฺมิโก วาทานุวาโท คารยฺหํ ฐานํ อาคจฺฉตีติ.}

\prose{ตคฺฆ ตุเมฺห ภิกฺขเว เอวํ ปุฏฺฐา เอวํ พฺยากรมานา วุตฺตวาทิโน เจว เม โหถ, น จ มํ อภูเตน อพฺภาจิกฺขถ, ธมฺมสฺส จานุธมฺมํ พฺยากโรถ, น จ โกจิ สหธมฺมิโก วาทานุวาโท คารยฺหํ ฐานํ อาคจฺฉติ.\csromanspacer{} ทุกฺขสฺส หิ ภิกฺขเว ปริญฺญตฺถํ มยิ พฺรหฺมจริยํ วุสฺสติ, สเจ ปน โว ภิกฺขเว อญฺญติตฺถิยา ปริพฺพาชกา เอวํ ปุจฺเฉยฺยุํ “กตมํ ปน ตํ อาวุโส ทุกฺขํ, ยสฺส ปริญฺญาย สมเณ โคตเม พฺรหฺมจริยํ วุสฺสตี”ติ, เอวํ ปุฏฺฐา ตุเมฺห ภิกฺขเว เตสํ อญฺญติตฺถิยานํ ปริพฺพาชกานํ เอวํ พฺยากเรยฺยาถ “จกฺขุ โข อาวุโส ทุกฺขํ, ตสฺส ปริญฺญาย ภควติ พฺรหฺมจริยํ วุสฺสติ.\csromanspacer{} รูปา ฯเปฯ.\csromanspacer{} ยมฺปิทํ จกฺขุสมฺผสฺสปจฺจยา อุปฺปชฺชติ เวทยิตํ สุขํ วา ทุกฺขํ วา อทุกฺขมสุขํ วา, ตมฺปิ ทุกฺขํ, ตสฺส ปริญฺญาย ภควติ พฺรหฺมจริยํ วุสฺสติ ฯเปฯ.\csromanspacer{} มโน ทุกฺโข ฯเปฯ.}

\gambhira{สฬายตนวคฺคสํยุตฺตปาฬิ}
\prose{\csromanfolio{278}ยมฺปิทํ มโนสมฺผสฺสปจฺจยา อุปฺปชฺชติ เวทยิตํ สุขํ วา ทุกฺขํ วา อทุกฺขมสุขํ วา, ตมฺปิ ทุกฺขํ, ตสฺส ปริญฺญาย ภควติ พฺรหฺมจริยํ วุสฺสติ.\csromanspacer{} อิทํ โข ตํ อาวุโส ทุกฺขํ, ตสฺส ปริญฺญาย ภควติ พฺรหฺมจริยํ วุสฺสตี”ติ, เอวํ ปุฏฺฐา ตุเมฺห ภิกฺขเว เตสํ อญฺญติตฺถิยานํ ปริพฺพาชกานํ เอวํ พฺยากเรยฺยาถาติ.\csromanspacer{}\csromanspacer{}อฏฺฐมํ.}

\csromanmatikaanchor{mtk.120}
\csromantocmarkref{subsection}{9. โลกปญฺหาสุตฺต}{mtk.120}
\bookmark[dest={mtk.120},level=2]{9. โลกปญฺหาสุตฺต}
\csromanheader{2}{9. โลกปญฺหาสุตฺต}
\proseitem{82}{อถ โข อญฺญตโร ภิกฺขุ เยน ภควา ฯเปฯ เอกมนฺตํ นิสินฺโน โข โส ภิกฺขุ ภควนฺตํ เอตทโวจ\csromandash{}}

\prose{“โลโก โลโก”ติ ภนฺเต วุจฺจติ, กิตฺตาวตา นุ โข ภนฺเต “โลโก”ติ วุจฺจตีติ.\csromanspacer{} ลุชฺชตีติ โข ภิกฺขุ ตสฺมา “โลโก”ติ วุจฺจติ.\csromanspacer{} กิญฺจ ลุชฺชติ.\csromanspacer{} จกฺขุ โข ภิกฺขุ ลุชฺชติ, รูปา ลุชฺชนฺติ, จกฺขุวิญฺญาณํ ลุชฺชติ, จกฺขุสมฺผสฺโส ลุชฺชติ, ยมฺปิทํ จกฺขุสมฺผสฺสปจฺจยา อุปฺปชฺชติ เวทยิตํ สุขํ วา ทุกฺขํ วา อทุกฺขมสุขํ วา, ตมฺปิ ลุชฺชติ ฯเปฯ.\csromanspacer{} ชิวฺหา ลุชฺชติ ฯเปฯ.\csromanspacer{} มโน ลุชฺชติ.\csromanspacer{} ธมฺมา ลุชฺชนฺติ.\csromanspacer{} มโนวิญฺญาณํ ลุชฺชติ.\csromanspacer{} มโนสมฺผสฺโส ลุชฺชติ.\csromanspacer{} ยมฺปิทํ มโนสมฺผสฺสปจฺจยา อุปฺปชฺชติ เวทยิตํ สุขํ วา ทุกฺขํ วา อทุกฺขมสุขํ วา, ตมฺปิ ลุชฺชติ.\csromanspacer{} ลุชฺชตีติ โข ภิกฺขุ ตสฺมา “โลโก”ติ วุจฺจตีติ.\csromanspacer{}\csromanspacer{}นวมํ.}

\csromanmatikaanchor{mtk.121}
\csromantocmarkref{subsection}{10. ผคฺคุนปญฺหาสุตฺต}{mtk.121}
\bookmark[dest={mtk.121},level=2]{10. ผคฺคุนปญฺหาสุตฺต}
\csromanheader{2}{10. ผคฺคุนปญฺหาสุตฺต}
\proseitem{83}{อถ โข อายสฺมา ผคฺคุโน ฯเปฯ เอกมนฺตํ นิสินฺโน โข อายสฺมา ผคฺคุโน ภควนฺตํ เอตทโวจ\csromandash{}}

\prose{อตฺถิ นุ โข ภนฺเต ตํ จกฺขุ, เยน จกฺขุนา อตีเต พุทฺเธ ปรินิพฺพุเต ฉินฺนปปญฺเจ ฉินฺนวฏุเม ปริยาทินฺนวฏฺเฏ สพฺพทุกฺขวีติวฏฺเฏ ปญฺญาปยมาโน ปญฺญาเปยฺย ฯเปฯ.\csromanspacer{} อตฺถิ นุ โข ภนฺเต สา ชิวฺหา, ยาย ชิวฺหาย อตีเต พุทฺเธ ปรินิพฺพุเต ฉินฺนปปญฺเจ ฉินฺนวฏุเม ปริยาทินฺนวฏฺเฏ สพฺพทุกฺขวีติวฏฺเฏ ปญฺญาปยมาโน ปญฺญาเปยฺย ฯเปฯ.\csromanspacer{} อตฺถิ นุ โข โส ภนฺเต มโน, เยน มเนน อตีเต พุทฺเธ ปรินิพฺพุเต ฉินฺนปปญฺเจ ฉินฺนวฏุเม ปริยาทินฺนวฏฺเฏ สพฺพทุกฺขวีติวฏฺเฏ ปญฺญาปยมาโน ปญฺญาเปยฺยาติ.}

\prose{นตฺถิ โข ตํ ผคฺคุน จกฺขุ, เยน จกฺขุนา อตีเต พุทฺเธ ปรินิพฺพุเต ฉินฺนปปญฺเจ ฉินฺนวฏุเม ปริยาทินฺนวฏฺเฏ สพฺพทุกฺขวีติวฏฺเฏ ปญฺญาปยมาโน}

\vspace{\tipitakaparskip}
\csromancenter{สฬายตนสํยุตฺต ปญฺญาเปยฺย ฯเปฯ.\csromanspacer{} นตฺถิ โข สา ผคฺคุน ชิวฺหา, ยาย ชิวฺหาย อตีเต พุทฺเธ ปรินิพฺพุเต ฉินฺนปปญฺเจ ฉินฺนวฏุเม ปริยาทินฺนวฏฺเฏ สพฺพทุกฺขวีติวฏฺเฏ ปญฺญาปยมาโน ปญฺญาเปยฺย ฯเปฯ.\csromanspacer{} นตฺถิ โข โส ผคฺคุน มโน, เยน มเนน อตีเต พุทฺเธ ปรินิพฺพุเต ฉินฺนปปญฺเจ ฉินฺนวฏุเม ปริยาทินฺนวฏฺเฏ สพฺพทุกฺขวีติวฏฺเฏ ปญฺญาปยมาโน ปญฺญาเปยฺยาติ.\csromanspacer{}\csromanspacer{}ทสมํ.\csromanspacer{} คิลานวคฺโค ตติโย.}
\titlehead{\textbf{ตสฺสุทฺทานํ}}
\csromangathameasure{คิลาเนน ทุเว วุตฺตา, ราเธน อปเร ตโย. \\ อวิชฺชาย จ เทฺว วุตฺตา, ภิกฺขุ โลโก จ ผคฺคุโนติ.}
\csromangathagroup{\csromangathastanza{\csromanfolio{279}คิลาเนน ทุเว วุตฺตา, ราเธน อปเร ตโย. \\ อวิชฺชาย จ เทฺว วุตฺตา, ภิกฺขุ โลโก จ ผคฺคุโนติ.}}

\csromanmatikaanchor{mtk.122}
\csromantocmarknopage{section}{(9) 4. ฉนฺนวคฺค}{mtk.122}
\bookmark[dest={mtk.122},level=1]{(9) 4. ฉนฺนวคฺค}
\csromanheader{1}{\textbf{(9) 4. ฉนฺนวคฺค}}
\csromanmatikaanchor{mtk.123}
\csromantocmarkref{subsection}{1. ปโลกธมฺมสุตฺต}{mtk.123}
\bookmark[dest={mtk.123},level=2]{1. ปโลกธมฺมสุตฺต}
\csromanheader{2}{1. ปโลกธมฺมสุตฺต}
\proseitem{84}{สาวตฺถินิทานํ.\csromanspacer{} อถ โข อายสฺมา อานนฺโท เยน ภควา เตนุปสงฺกมิ ฯเปฯ เอกมนฺตํ นิสินฺโน โข อายสฺมา อานนฺโท ภควนฺตํ เอตทโวจ\csromandash{}}

\prose{“โลโก โลโก”ติ ภนฺเต วุจฺจติ, กิตฺตาวตา นุ โข ภนฺเต “โลโก”ติ วุจฺจตีติ.\csromanspacer{} ยํ โข อานนฺท ปโลกธมฺมํ, อยํ วุจฺจติ อริยสฺส วินเย โลโก.\csromanspacer{} กิญฺจานนฺท ปโลกธมฺมํ.\csromanspacer{} จกฺขุ โข อานนฺท ปโลกธมฺมํ, รูปา ปโลกธมฺมา, จกฺขุวิญฺญาณํ ปโลกธมฺมํ, จกฺขุสมฺผสฺโส ปโลกธมฺโม, ยมฺปิทํ จกฺขุสมฺผสฺสปจฺจยา ฯเปฯ.\csromanspacer{} ตมฺปิ ปโลกธมฺมํ ฯเปฯ.\csromanspacer{} ชิวฺหา ปโลกธมฺมา, รสา ปโลกธมฺมา, ชิวฺหาวิญฺญาณํ ปโลกธมฺมํ, ชิวฺหาสมฺผสฺโส ปโลกธมฺโม, ยมฺปิทํ ชิวฺหาสมฺผสฺสปจฺจยา ฯเปฯ.\csromanspacer{} ตมฺปิ ปโลกธมฺมํ ฯเปฯ.\csromanspacer{} มโน ปโลกธมฺโม, ธมฺมา ปโลกธมฺมา, มโนวิญฺญาณํ ปโลกธมฺมํ, มโนสมฺผสฺโส ปโลกธมฺโม, ยมฺปิทํ มโนสมฺผสฺสปจฺจยา อุปฺปชฺชติ เวทยิตํ สุขํ วา ทุกฺขํ วา อทุกฺขมสุขํ วา, ตมฺปิปโลกธมฺมํ.\csromanspacer{} ยํ โข อานนฺท ปโลกธมฺมํ, อยํ วุจฺจติ อริยสฺส วินเย โลโกติ.\csromanspacer{}\csromanspacer{}ปฐมํ.}

\gambhira{สฬายตนวคฺคสํยุตฺตปาฬิ}
\csromanmatikaanchor{mtk.124}
\csromantocmarkref{subsection}{2. สุญฺญตโลกสุตฺต}{mtk.124}
\bookmark[dest={mtk.124},level=2]{2. สุญฺญตโลกสุตฺต}
\csromanheader{2}{2. สุญฺญตโลกสุตฺต}
\proseitem{85}{\csromanfolio{280}อถ โข อายสฺมา อานนฺโท ฯเปฯ ภควนฺตํ เอตทโวจ\csromandash{}“สุญฺโญ โลโก สุญฺโญ โลโก”ติ ภนฺเต วุจฺจติ, กิตฺตาวตา นุ โข ภนฺเต “สุญฺโญ โลโก”ติ วุจฺจตีติ.\csromanspacer{} ยสฺมา จ โข อานนฺท สุญฺญํ อตฺเตน วา อตฺตนิเยน วา, ตสฺมา “สุญฺโญ โลโก”ติ วุจฺจติ.\csromanspacer{} กิญฺจ อานนฺท สุญฺญํ อตฺเตน วา อตฺตนิเยน วา.\csromanspacer{} จกฺขุ โข อานนฺท สุญฺญํ อตฺเตน วา อตฺตนิเยน วา, รูปา สุญฺญา อตฺเตน วา อตฺตนิเยน วา, จกฺขุวิญฺญาณํ สุญฺญํ อตฺเตน วา อตฺตนิเยน วา, จกฺขุสมฺผสฺโส สุญฺโญ อตฺเตน วา อตฺตนิเยน วา ฯเปฯ.\csromanspacer{} ยมฺปิทํ มโนสมฺผสฺสปจฺจยา อุปฺปชฺชติ เวทยิตํ สุขํ วา ทุกฺขํ วา อทุกฺขมสุขํ วา, ตมฺปิ สุญฺญํ อตฺเตน วา อตฺตนิเยน วา.\csromanspacer{} ยสฺมา จ โข อานนฺท สุญฺญํ อตฺเตน วา อตฺตนิเยน วา, ตสฺมา “สุญฺโญ โลโก”ติ วุจฺจตีติ.\csromanspacer{}\csromanspacer{}ทุติยํ.}

\csromanmatikaanchor{mtk.125}
\csromantocmarkref{subsection}{3. สํขิตฺตธมฺมสุตฺต}{mtk.125}
\bookmark[dest={mtk.125},level=2]{3. สํขิตฺตธมฺมสุตฺต}
\csromanheader{2}{3. สํขิตฺตธมฺมสุตฺต}
\proseitem{86}{เอกมนฺตํ นิสินฺโน โข อายสฺมา อานนฺโท ภควนฺตํ เอตทโวจ “สาธุ เม ภนฺเต ภควา สํขิตฺเตน ธมฺมํ เทเสตุ, ยมหํ ภควโต ธมฺมํ สุตฺวา เอโก วูปกฏฺโฐ อปฺปมตฺโต อาตาปี ปหิตตฺโต วิหเรยฺยนฺ”ติ.}

\prose{ตํ กิํ มญฺญสิ อานนฺท, จกฺขุ นิจฺจํ วา อนิจฺจํ วาติ.\csromanspacer{} อนิจฺจํ ภนฺเต.\csromanspacer{} ยํ ปนานิจฺจํ, ทุกฺขํ วา ตํ สุขํ วาติ.\csromanspacer{} ทุกฺขํ ภนฺเต.\csromanspacer{} ยํ ปนานิจฺจํ ทุกฺขํ วิปริณามธมฺมํ, กลฺลํ นุ ตํ สมนุปสฺสิตุํ “เอตํ มม, เอโสหมสฺมิ, เอโส เม อตฺตา”ติ.\csromanspacer{} โน เหตํ ภนฺเต.\csromanspacer{} รูปา นิจฺจา วา อนิจฺจา วาติ.\csromanspacer{} อนิจฺจา ภนฺเต.\csromanspacer{} จกฺขุวิญฺญาณํ ฯเปฯ.\csromanspacer{} ยมฺปิทํ จกฺขุสมฺผสฺสปจฺจยา อุปฺปชฺชติ เวทยิตํ สุขํ วา ทุกฺขํ วา อทุกฺขมสุขํ วา, ตมฺปิ นิจฺจํ วา อนิจฺจํ วาติ.\csromanspacer{} อนิจฺจํ ภนฺเต.\csromanspacer{} ยํ ปนานิจฺจํ, ทุกฺขํ วา ตํ สุขํ วาติ.\csromanspacer{} ทุกฺขํ ภนฺเต.\csromanspacer{} ยํ ปนานิจฺจํ ทุกฺขํ วิปริณามธมฺมํ, กลฺลํ นุ ตํ สมนุปสฺสิตุํ “เอตํ มม, เอโสหมสฺมิ, เอโส เม อตฺตา”ติ.\csromanspacer{} โน เหตํ ภนฺเต ฯเปฯ.\csromanspacer{} ชิวฺหา นิจฺจา วา อนิจฺจา วาติ.\csromanspacer{} อนิจฺจา ภนฺเต ฯเปฯ.\csromanspacer{} ชิวฺหาวิญฺญาณํ.\csromanspacer{} ชิวฺหาสมฺผสฺโส ฯเปฯ.\csromanspacer{} ยมฺปิทํ มโนสมฺผสฺสปจฺจยา อุปฺปชฺชติ เวทยิตํ สุขํ วา ทุกฺขํ วา อทุกฺขมสุขํ วา, ตมฺปิ นิจฺจํ วา อนิจฺจํ วาติ.\csromanspacer{} อนิจฺจํ ภนฺเต.\csromanspacer{} ยํ ปนานิจฺจํ, ทุกฺขํ วา ตํ สุขํ วาติ.\csromanspacer{} ทุกฺขํ ภนฺเต.\csromanspacer{} ยํ ปนานิจฺจํ ทุกฺขํ วิปริณามธมฺมํ, กลฺลํ นุ ตํ สมนุปสฺสิตุํ “เอตํ มม, เอโสหมสฺมิ,}

\vspace{\tipitakaparskip}
\csromancenter{สฬายตนสํยุตฺต เอโส เม อตฺตา”ติ.\csromanspacer{} โน เหตํ ภนฺเต.\csromanspacer{} เอวํ ปสฺสํ อานนฺท สุตวา อริยสาวโก จกฺขุสฺมิมฺปิ นิพฺพินฺทติ ฯเปฯ จกฺขุสมฺผสฺเสปิ นิพฺพินฺทติ ฯเปฯ.\csromanspacer{} ยมฺปิทํ มโนสมฺผสฺสปจฺจยา อุปฺปชฺชติ เวทยิตํ สุขํ วา ทุกฺขํ วา อทุกฺขมสุขํ วา, ตสฺมิมฺปิ นิพฺพินฺทติ, นิพฺพินฺทํ วิรชฺชติ, วิราคา วิมุจฺจติ, วิมุตฺตสฺมิํ “วิมุตฺตมฺ”อิติ ญาณํ โหติ, “ขีณา ชาติ, วุสิตํ พฺรหฺมจริยํ, กตํ กรณียํ, นาปรํ อิตฺถตฺตายา”ติ ปชานาตีติ.\csromanspacer{}\csromanspacer{}ตติยํ.}
\csromanmatikaanchor{mtk.126}
\csromantocmarkref{subsection}{4. ฉนฺนสุตฺต}{mtk.126}
\bookmark[dest={mtk.126},level=2]{4. ฉนฺนสุตฺต}
\csromanheader{2}{4. ฉนฺนสุตฺต}
\proseitem{87}{\csromanfolio{281}เอกํ สมยํ ภควา ราชคเห วิหรติ เวฬุวเน กลนฺทกนิวาเป.\csromanspacer{} เตน โข ปน สมเยน อายสฺมา จ สาริปุตฺโต อายสฺมา จ มหาจุนฺโท อายสฺมา จ ฉนฺโน คิชฺฌกูเฏ ปพฺพเต วิหรนฺติ.\csromanspacer{} เตน โข ปน สมเยน อายสฺมา ฉนฺโน อาพาธิโก โหติ ทุกฺขิโต พาฬฺหคิลาโน.\csromanspacer{} อถ โข อายสฺมา สาริปุตฺโต สายนฺหสมยํ ปฏิสลฺลานา วุฏฺฐิโต เยนายสฺมา มหาจุนฺโท เตนุปสงฺกมิ, อุปสงฺกมิตฺวา อายสฺมนฺตํ มหาจุนฺทํ เอตทโวจ “อายามาวุโส จุนฺท เยนายสฺมา ฉนฺโน เตนุปสงฺกมิสฺสาม คิลานปุจฺฉกา”ติ. “เอวมาวุโส”ติ โข อายสฺมา มหาจุนฺโท อายสฺมโต สาริปุตฺตสฺส ปจฺจสฺโสสิ.}

\prose{อถ โข อายสฺมา จ สาริปุตฺโต อายสฺมา จ มหาจุนฺโท เยนายสฺมา ฉนฺโน เตนุปสงฺกมิํสุ, อุปสงฺกมิตฺวา ปญฺญตฺเต อาสเน นิสีทิํสุ, นิสชฺช โข อายสฺมา สาริปุตฺโต อายสฺมนฺตํ ฉนฺนํ เอตทโวจ “กจฺจิ เต อาวุโส ฉนฺน ขมนียํ, กจฺจิ ยาปนียํ, กจฺจิ ทุกฺขา เวทนา ปฏิกฺกมนฺติ โน อภิกฺกมนฺติ, ปฏิกฺกโมสานํ ปญฺญายติ โน อภิกฺกโม”ติ.}

\prose{น เม อาวุโส สาริปุตฺต ขมนียํ, น ยาปนียํ, พาฬฺหา เม ทุกฺขา เวทนา อภิกฺกมนฺติ โน ปฏิกฺกมนฺติ, อภิกฺกโมสานํ ปญฺญายติ, โน ปฏิกฺกโม.\csromanspacer{} เสยฺยถาปิ อาวุโส พลวา ปุริโส ติเณฺหน สิขเรน\footnote{ขคฺเคน (ก)} มุทฺธนิ\footnote{มุทฺธานํ (สี, สฺยา, กํ, อิ)} อภิมตฺเถยฺย\footnote{อภิมนฺเถยฺย (สี)}.\csromanspacer{} เอวเมว โข อาวุโส อธิมตฺตา วาตา มุทฺธนิ2 อูหนนฺติ\footnote{อุปหนนฺติ (สี, สฺยา, กํ, อิ, ก), อุหนนฺติ (ก)}.\csromanspacer{} น เม อาวุโส ขมนียํ, น ยาปนียํ ฯเปฯ โน ปฏิกฺกโม.}

\gambhira{สฬายตนวคฺคสํยุตฺตปาฬิ}
\prose{\csromanfolio{282}เสยฺยถาปิ อาวุโส พลวา ปุริโส ทเฬฺหน วรตฺตกฺขณฺเฑน สีเส สีสเวฐํ ทเทยฺย, เอวเมว โข อาวุโส อธิมตฺตา สีเส สีสเวทนา.\csromanspacer{} น เม อาวุโส ขมนียํ, น ยาปนียํ ฯเปฯ โน ปฏิกฺกโม.\csromanspacer{} เสยฺยถาปิ อาวุโส ทกฺโข โคฆาตโก วา โคฆาตกนฺเตวาสี วา ติเณฺหน โควิกนฺตเนน กุจฺฉิํ ปริกนฺเตยฺย, เอวเมว โข อธิมตฺตา วาตา กุจฺฉิํ ปริกนฺตนฺติ.\csromanspacer{} น เม อาวุโส ขมนียํ, น ยาปนียํ ฯเปฯ โน ปฏิกฺกโม.\csromanspacer{} เสยฺยถาปิ อาวุโส เทฺว พลวนฺโต ปุริสา ทุพฺพลตรํ ปุริสํ นานาพาหาสุ คเหตฺวา องฺคารกาสุยา สนฺตาเปยฺยุํ สมฺปริตาเปยฺยุํ, เอวเมว โข อาวุโส อธิมตฺโต กายสฺมิํ ฑาโห, น เม อาวุโส ขมนียํ, น ยาปนียํ, พาฬฺหา เม ทุกฺขา เวทนา อภิกฺกมนฺติ โน ปฏิกฺกมนฺติ, อภิกฺกโมสานํ ปญฺญายติ โน ปฏิกฺกโม, สตฺถํ อาวุโส สาริปุตฺต อาหริสฺสามิ, นาวกงฺขามิ\footnote{นาปิ กงฺขามิ (ก)} ชีวิตนฺติ.}

\prose{มา อายสฺมา ฉนฺโน สตฺถํ อาหเรสิ, ยาเปตายสฺมา ฉนฺโน, ยาเปนฺตํ มยํ อายสฺมนฺตํ ฉนฺนํ อิจฺฉาม, สเจ อายสฺมโต ฉนฺนสฺส นตฺถิ สปฺปายานิ โภชนานิ, อหํ อายสฺมโต ฉนฺนสฺส สปฺปายานิ โภชนานิ ปริเยสิสฺสามิ.\csromanspacer{} สเจ อายสฺมโต ฉนฺนสฺส นตฺถิ สปฺปายานิ เภสชฺชานิ อหํ อายสฺมโต ฉนฺนสฺส สปฺปายานิ เภสชฺชานิ ปริเยสิสฺสามิ.\csromanspacer{} สเจ อายสฺมโต ฉนฺนสฺส นตฺถิ ปติรูปา อุปฏฺฐากา, อหํ อายสฺมนฺตํ ฉนฺนํ อุปฏฺฐหิสฺสามิ.\csromanspacer{} มา อายสฺมา ฉนฺโน สตฺถํ อาหเรสิ, ยาเปตายสฺมา ฉนฺโน, ยาเปนฺตํ มยํ อายสฺมนฺตํ ฉนฺนํ อิจฺฉามาติ.}

\prose{น เม อาวุโส สาริปุตฺต นตฺถิ สปฺปายานิ โภชนานิ, อตฺถิ เม สปฺปายานิ โภชนานิ.\csromanspacer{} นปิ เม นตฺถิ สปฺปายานิ เภสชฺชานิ, อตฺถิ เม สปฺปายานิ เภสชฺชานิ.\csromanspacer{} นปิ เม นตฺถิ ปติรูปา อุปฏฺฐากา, อตฺถิ เม ปติรูปา อุปฏฺฐากา.\csromanspacer{} อปิ จ เม อาวุโส สตฺถา ปริจิณฺโณ ทีฆรตฺตํ มนาเปเนว, โน อมนาเปน.\csromanspacer{} เอตํ หิ อาวุโส สาวกสฺส ปติรูปํ, ยํ สตฺถารํ ปริจเรยฺย มนาเปเนว โน อมนาเปน, อนุปวชฺชํ\footnote{ตํ อนุปวชฺชํ (พหูสุ)} ฉนฺโน ภิกฺขุ สตฺถํ อาหริสฺสตีติ เอวเมตํ อาวุโส สาริปุตฺต ธาเรหีติ.}

\csromanheader{2}{1. สฬายตนสํยุตฺต}
\prose{\csromanfolio{283}ปุจฺเฉยฺยาม มยํ อายสฺมนฺตํ ฉนฺนํ กญฺจิเทว\footnote{กิญฺจิเทว (สฺยา, กํ, อิ, ก)} เทสํ, สเจ อายสฺมา ฉนฺโน โอกาสํ กโรติ ปญฺหสฺส เวยฺยากรณายาติ.\csromanspacer{} ปุจฺฉาวุโส สาริปุตฺต สุตฺวา เวทิสฺสามาติ.}

\prose{จกฺขุํ อาวุโส ฉนฺน จกฺขุวิญฺญาณํ จกฺขุวิญฺญาณวิญฺญาตพฺเพ ธมฺเม “เอตํ มม, เอโสหมสฺมิ, เอโส เม อตฺตา”ติ สมนุปสฺสสิ ฯเปฯ.\csromanspacer{} ชิวฺหํ อาวุโส ฉนฺน ชิวฺหาวิญฺญาณํ ชิวฺหาวิญฺญาณวิญฺญาตพฺเพ ธมฺเม “เอตํ มม, เอโสหมสฺมิ, เอโส เม อตฺตา”ติ สมนุปสฺสติ ฯเปฯ.\csromanspacer{} มนํ อาวุโส ฉนฺน มโนวิญฺญาณํ มโนวิญฺญาณวิญฺญาตพฺเพ ธมฺเม “เอตํ มม, เอโสหมสฺมิ, เอโส เม อตฺตา”ติ สมนุปสฺสสีติ.}

\prose{จกฺขุํ อาวุโส สาริปุตฺต จกฺขุวิญฺญาณํ จกฺขุวิญฺญาณวิญฺญาตพฺเพ ธมฺเม “เนตํ มม, เนโสหมสฺมิ, น เมโส อตฺตา”ติ สมนุปสฺสามิ ฯเปฯ.\csromanspacer{} ชิวฺหํ อาวุโส สาริปุตฺต ชิวฺหาวิญฺญาณํ ชิวฺหาวิญฺญาณวิญฺญาตพฺเพ ธมฺเม “เนตํ มม, เนโสหมสฺมิ, น เมโส อตฺตา”ติ สมนุปสฺสามิ ฯเปฯ.\csromanspacer{} มนํ อาวุโส สาริปุตฺต มโนวิญฺญาณํ มโนวิญฺญาณวิญฺญาตพฺเพ ธมฺเม “เนตํ มม, เนโสหมสฺมิ, น เมโส อตฺตา”ติ สมนุปสฺสามีติ.}

\prose{จกฺขุสฺมิํ อาวุโส ฉนฺน จกฺขุวิญฺญาเณ จกฺขุวิญฺญาณวิญฺญาตพฺเพสุ ธมฺเมสุ กิํ ทิสฺวา กิํ อภิญฺญาย จกฺขุํ จกฺขุวิญฺญาณํ จกฺขุวิญฺญาณวิญฺญาตพฺเพ ธมฺเม “เนตํ มม, เนโสหมสฺมิ, น เมโส อตฺตา”ติ สมนุปสฺสสิ ฯเปฯ.\csromanspacer{} ชิวฺหาย อาวุโส ฉนฺน ชิวฺหาวิญฺญาเณ ชิวฺหาวิญฺญาณวิญฺญาตพฺเพสุ ธมฺเมสุ กิํ ทิสฺวา กิํ อภิญฺญาย ชิวฺหํ ชิวฺหาวิญฺญาณํ ชิวฺหาวิญฺญาณวิญฺญาตพฺเพ ธมฺเม “เนตํ มม, เนโสหมสฺมิ, น เมโส อตฺตา”ติ สมนุปสฺสสิ.\csromanspacer{} มนสฺมิํ อาวุโส ฉนฺน มโนวิญฺญาเณ มโนวิญฺญาณวิญฺญาตพฺเพสุ ธมฺเมสุ กิํ ทิสฺวา กิํ อภิญฺญาย มนํ มโนวิญฺญาณํ มโนวิญฺญาณวิญฺญาตพฺเพ ธมฺเม “เนตํ มม, เนโสหมสฺมิ, น เมโส อตฺตา”ติ สมนุปสฺสสีติ.}

\prose{จกฺขุสฺมิํ อาวุโส สาริปุตฺต จกฺขุวิญฺญาเณ จกฺขุวิญฺญาณวิญฺญาตพฺเพสุ ธมฺเมสุ นิโรธํ ทิสฺวา นิโรธํ อภิญฺญาย จกฺขุํ จกฺขุวิญฺญาณํ จกฺขุวิญฺญาณวิญฺญาตพฺเพ ธมฺเม “เนตํ มม, เนโสหมสฺมิ, น เมโส}

\gambhira{สฬายตนวคฺคสํยุตฺตปาฬิ}
\prose{\csromanfolio{284}อตฺตา”ติ สมนุปสฺสามิ ฯเปฯ.\csromanspacer{} ชิวฺหาย อาวุโส สาริปุตฺต ชิวฺหาวิญฺญาเณ ชิวฺหาวิญฺญาณวิญฺญาตพฺเพสุ ธมฺเมสุ นิโรธํ ทิสฺวา นิโรธํ อภิญฺญาย ชิวฺหํ ชิวฺหาวิญฺญาณํ ชิวฺหาวิญฺญาณวิญฺญาตพฺเพ ธมฺเม “เนตํ มม, เนโสหมสฺมิ, น เมโส อตฺตา”ติ สมนุปสฺสามิ ฯเปฯ.\csromanspacer{} มนสฺมิํ อาวุโส สาริปุตฺต มโนวิญฺญาเณ มโนวิญฺญาณวิญฺญาตพฺเพสุ ธมฺเมสุ นิโรธํ ทิสฺวา นิโรธํ อภิญฺญาย มนํ มโนวิญฺญาณํ มโนวิญฺญาณวิญฺญาตพฺเพ ธมฺเม “เนตํ มม, เนโสหมสฺมิ, น เมโส อตฺตา”ติสมนุปสฺสามีติ.}

\prose{เอวํ วุตฺเต อายสฺมา มหาจุนฺโท อายสฺมนฺตํ ฉนฺนํ เอตทโวจ “ตสฺมาติห อาวุโส ฉนฺน อิทมฺปิ ตสฺส ภควโต สาสนํ นิจฺจกปฺปํ สาธุกํ มนสิ กาตพฺพํ, นิสฺสิตสฺส จลิตํ อนิสฺสิตสฺส จลิตํ นตฺถิ, จลิเต อสติ ปสฺสทฺธิ โหติ, ปสฺสทฺธิยา สติ นติ น โหติ, นติยา อสติ อาคติคติ น โหติ, อาคติคติยา อสติ จุตูปปาโต น โหติ, จุตูปปาเต อสติ เนวิธ น หุรํ น อุภยมนฺตเรน เอเสวนฺโต ทุกฺขสฺสา”ติ.\csromanspacer{} อถ โข อายสฺมา จ สาริปุตฺโต อายสฺมา จ มหาจุนฺโท อายสฺมนฺตํ ฉนฺนํ อิมินา โอวาเทน โอวทิตฺวา อุฏฺฐายาสนา ปกฺกมิํสุ.\csromanspacer{} อถ โข อายสฺมา ฉนฺโน อจิรปกฺกนฺเตสุ เตสุ อายสฺมนฺเตสุ สตฺถํ อาหเรสิ.}

\prose{อถ โข อายสฺมา สาริปุตฺโต เยน ภควา เตนุปสงฺกมิ, อุปสงฺกมิตฺวา ภควนฺตํ อภิวาเทตฺวา เอกมนฺตํ นิสีทิ, เอกมนฺตํ นิสินฺโน โข อายสฺมา สาริปุตฺโต ภควนฺตํ เอตทโวจ “อายสฺมตา ภนฺเต ฉนฺเนน สตฺถํ อาหริตํ, ตสฺส กา คติ โก อภิสมฺปราโย”ติ.\csromanspacer{} นนุ เต สาริปุตฺต ฉนฺเนน ภิกฺขุนา สมฺมุขาเยว อนุปวชฺชตา พฺยากตาติ.\csromanspacer{} อตฺถิ ภนฺเต ปุพฺพวิชฺชนํ\footnote{ปุพฺพวิจิรํ (สี), ปุพฺพวิชฺฌนํ (อิ), ปุพฺพชิรํ (ม 3. 310 ปิฏฺเฐ)} นาม วชฺชิคาโม, ตตฺถายสฺมโต ฉนฺนสฺส มิตฺตกุลานิ สุหชฺชกุลานิ อุปวชฺชกุลานีติ.\csromanspacer{} โหนฺติ เหเต สาริปุตฺต ฉนฺนสฺส ภิกฺขุโน มิตฺตกุลานิ สุหชฺชกุลานิ อุปวชฺชกุลานิ, น โข ปนาหํ สาริปุตฺต เอตฺตาวตา “ส-อุปวชฺโช”ติ วทามิ, โย โข สาริปุตฺต ตญฺจ กายํ นิกฺขิปติ, อญฺญญฺจ กายํ อุปาทิยติ, ตมหํ “ส-อุปวชฺโช”ติ วทามิ, ตํ ฉนฺนสฺส ภิกฺขุโน นตฺถิ, อนุปวชฺชํ ฉนฺเนน ภิกฺขุนา สตฺถํ อาหริตนฺติ, เอวเมตํ สาริปุตฺต ธาเรหีติ.\csromanspacer{}\csromanspacer{}จตุตฺถํ.}

\csromanmatikaanchor{mtk.127}
\csromantocmarkref{subsection}{5. ปุณฺณสุตฺต}{mtk.127}
\bookmark[dest={mtk.127},level=2]{5. ปุณฺณสุตฺต}
\csromanheader{2}{1. สฬายตนสํยุตฺต \textbf{5. ปุณฺณสุตฺต}}
\proseitem{88}{\csromanfolio{285}อถ\footnote{สาวตฺถินิทานํ. อถ (?) ม 3. 311 ปสฺสิตพฺพํ.} โข อายสฺมา ปุณฺโณ เยน ภควา เตนุปสงฺกมิ, อุปสงฺกมิตฺวา ฯเปฯ เอกมนฺตํ นิสินฺโน โข อายสฺมา ปุณฺโณ ภควนฺตํ เอตทโวจ “สาธุ เม ภนฺเต ภควา สํขิตฺเตน ธมฺมํ เทเสตุ, ยมหํ ภควโต ธมฺมํ สุตฺวา เอโก วูปกฏฺโฐ อปฺปมตฺโต อาตาปี ปหิตตฺโต วิหเรยฺยนฺ”ติ.}

\prose{สนฺติ โข ปุณฺณ จกฺขุวิญฺเญยฺยา รูปา อิฏฺฐา กนฺตา มนาปา ปิยรูปา กามูปสํหิตา รชนียา, ตญฺเจ ภิกฺขุ อภินนฺทติ อภิวทติ อชฺโฌสาย ติฏฺฐติ, ตสฺส ตํ อภินนฺทโต อภิวทโต อชฺโฌสาย ติฏฺฐโต อุปฺปชฺชติ นนฺที, นนฺทิสมุทยา ทุกฺขสมุทโย ปุณฺณาติ วทามิ ฯเปฯ.\csromanspacer{} สนฺติ โข ปุณฺณ ชิวฺหาวิญฺเญยฺยา รสา ฯเปฯ.\csromanspacer{} สนฺติ โข ปุณฺณ มโนวิญฺเญยฺยา ธมฺมา อิฏฺฐา กนฺตา มนาปา ปิยรูปา กามูปสํหิตา รชนียา, ตญฺเจ ภิกฺขุ อภินนฺทติ อภิวทติ อชฺโฌสาย ติฏฺฐติ, ตสฺส ตํ อภินนฺทโต อภิวทโต อชฺโฌสาย ติฏฺฐโต อุปฺปชฺชติ นนฺที, นนฺทิสมุทยา ทุกฺขสมุทโย ปุณฺณาติ วทามิ.}

\prose{สนฺติ โข ปุณฺณ จกฺขุวิญฺเญยฺยา รูปา อิฏฺฐา กนฺตา มนาปา ปิยรูปา กามูปสํหิตา รชนียา, ตญฺเจ ภิกฺขุ นาภินนฺทติ นาภิวทติ นาชฺโฌสาย ติฏฺฐติ, ตสฺส ตํ อนภินนฺทโต อนภิวทโต อนชฺโฌสาย ติฏฺฐโต นิรุชฺฌติ นนฺที, นนฺทินิโรธา ทุกฺขนิโรโธ ปุณฺณาติ วทามิ ฯเปฯ.\csromanspacer{} สนฺติ โข ปุณฺณ มโนวิญฺเญยฺยา ธมฺมา อิฏฺฐา กนฺตา มนาปา ปิยรูปา กามูปสํหิตา รชนียา, ตญฺเจ ภิกฺขุ นาภินนฺทติ นาภิวทติ นาชฺโฌสาย ติฏฺฐติ, ตสฺส ตํ อนภินนฺทโต อนภิวทโต อนชฺโฌสาย ติฏฺฐโต นิรุชฺฌติ นนฺที, นนฺทินิโรธา ทุกฺขนิโรโธ ปุณฺณาติ วทามิ.}

\prose{อิมินา ตฺวํ\footnote{อิมินา จ ตฺวํ.} ปุณฺณ มยา สํขิตฺเตน โอวาเทน โอวทิโต กตมสฺมิํ\footnote{กตรสฺมิํ (ม 3. 312)} ชนปเท วิหริสฺสสีติ.\csromanspacer{} อตฺถิ ภนฺเต สุนาปรนฺโต นาม ชนปโท, ตตฺถาหํ วิหริสฺสามีติ.}

\prose{จณฺฑา โข ปุณฺณ สุนาปรนฺตกา มนุสฺสา, ผรุสา โข ปุณฺณ สุนาปรนฺตกา มนุสฺสา, สเจ ตํ ปุณฺณ สุนาปรนฺตกา มนุสฺสา อกฺโกสิสฺสนฺติ ปริภาสิสฺสนฺติ, ตตฺร เต ปุณฺณ กินฺติ ภวิสฺสตีติ.\csromanspacer{} สเจ มํ ภนฺเต}

\gambhira{สฬายตนวคฺคสํยุตฺตปาฬิ}
\prose{\csromanfolio{286}สุนาปรนฺตกา มนุสฺสา อกฺโกสิสฺสนฺติ ปริภาสิสฺสนฺติ, ตตฺร เม เอวํ ภวิสฺสติ “ภทฺทกา วติเม สุนาปรนฺตกา มนุสฺสา, สุภทฺทกา วติเม สุนาปรนฺตกา มนุสฺสา, ยํ เม\footnote{มํ (สพฺพตฺถ)} น ยิเม ปาณินา ปหารํ เทนฺตี”ติ, เอวเมตฺถ\footnote{เอวมฺเมตฺถ (?)} ภควา ภวิสฺสติ, เอวเมตฺถ สุคต ภวิสฺสตีติ.}

\prose{สเจ ปน เต ปุณฺณ สุนาปรนฺตกา มนุสฺสา ปาณินา ปหารํ ทสฺสนฺติ, ตตฺร ปน เต ปุณฺณ กินฺติ ภวิสฺสตีติ.\csromanspacer{} สเจ เม ภนฺเต สุนาปรนฺตกา มนุสฺสา ปาณินา ปหารํ ทสฺสนฺติ, ตตฺร เม เอวํ ภวิสฺสติ “ภทฺทกา วติเม สุนาปรนฺตกา มนุสฺสา, สุภทฺทกา วติเม สุนาปรนฺตกา มนุสฺสา, ยํ เม1 น ยิเม เลฑฺฑุนา ปหารํ เทนฺตี”ติ, เอวเมตฺถ ภควา ภวิสฺสติ, เอวเมตฺถ สุคต ภวิสฺสตีติ.}

\prose{สเจ ปน เต ปุณฺณ สุนาปรนฺตกา มนุสฺสา เลฑฺฑุนา ปหารํ ทสฺสนฺติ, ตตฺร ปน เต ปุณฺณ กินฺติ ภวิสฺสตีติ.\csromanspacer{} สเจ เม ภนฺเต สุนาปรนฺตกา มนุสฺสา เลฑฺฑุนา ปหารํ ทสฺสนฺติ, ตตฺร เม เอวํ ภวิสฺสติ “ภทฺทกา วติเม สุนาปรนฺตกา มนุสฺสา, สุภทฺทกา วติเม สุนาปรนฺตกา มนุสฺสา, ยํ เม น ยิเม ทณฺเฑน ปหารํ เทนฺตี”ติ, เอวเมตฺถ ภควา ภวิสฺสติ, เอวเมตฺถ สุคต ภวิสฺสตีติ.}

\prose{สเจ ปน ปุณฺณ สุนาปรนฺตกา มนุสฺสา ทณฺเฑน ปหารํ ทสฺสนฺติ, ตตฺร ปน เต ปุณฺณ กินฺติ ภวิสฺสตีติ.\csromanspacer{} สเจ เม ภนฺเต สุนาปรนฺตกา มนุสฺสา ทณฺเฑน ปหารํ ทสฺสนฺติ, ตตฺร เม เอวํ ภวิสฺสติ “ภทฺทกา วติเม สุนาปรนฺตกา มนุสฺสา, สุภทฺทกา วติเม สุนาปรนฺตกา มนุสฺสา, ยํ เม น ยิเม สตฺเถน ปหารํ เทนฺตี”ติ, เอวเมตฺถ ภควา ภวิสฺสติ, เอวเมตฺถ สุคต ภวิสฺสตีติ.}

\prose{สเจ ปน เต ปุณฺณ สุนาปรนฺตกา มนุสฺสา สตฺเถน ปหารํ ทสฺสนฺติ, ตตฺร ปน เต ปุณฺณ กินฺติ ภวิสฺสตีติ.\csromanspacer{} สเจ เม ภนฺเต สุนาปรนฺตกา มนุสฺสา สตฺเถน ปหารํ ทสฺสนฺติ, ตตฺร เม เอวํ ภวิสฺสติ “ภทฺทกา วติเม สุนาปรนฺตกา มนุสฺสา, สุภทฺทกา วติเม สุนาปรนฺตกา มนุสฺสา, ยํ มํ น ยิเม ติเณฺหน สตฺเถน ชีวิตา โวโรเปนฺตี”ติ, เอวเมตฺถ ภควา ภวิสฺสติ, เอวเมตฺถ สุคต ภวิสฺสตีติ.}

\csromanheader{2}{1. สฬายตนสํยุตฺต}
\prose{\csromanfolio{287}สเจ ปน ตํ ปุณฺณ สุนาปรนฺตกา มนุสฺสา ติเณฺหน สตฺเถน ชีวิตา โวโรเปสฺสนฺติ, ตตฺร ปน เต ปุณฺณ กินฺติ ภวิสฺสตีติ.\csromanspacer{} สเจ มํ ภนฺเต สุนาปรนฺตกา มนุสฺสา ติเณฺหน สตฺเถน ชีวิตา โวโรเปสฺสนฺติ, ตตฺร เม เอวํ ภวิสฺสติ “สนฺติ โข ตสฺส ภควโต สาวกา กาเยน จ ชีวิเตน จ อฏฺฏียมานา หรายมานา ชิคุจฺฉมานา สตฺถหารกํ ปริเยสนฺติ, ตํ เม อิทํ อปริยิฏฺฐญฺเญว สตฺถหารกํ ลทฺธนฺ”ติ, เอวเมตฺถ ภควา ภวิสฺสติ, เอวเมตฺถ สุคต ภวิสฺสตีติ.}

\prose{สาธุ สาธุ ปุณฺณ, สกฺขิสฺสสิ โข ตฺวํ ปุณฺณ อิมินา ทมูปสเมน สมนฺนาคโต สุนาปรนฺตสฺมิํ ชนปเท วตฺถุํ, ยสฺส ทานิ ตฺวํ ปุณฺณ กาลํ มญฺญสีติ.}

\prose{อถ โข อายสฺมา ปุณฺโณ ภควโต วจนํ อภินนฺทิตฺวา อนุโมทิตฺวา อุฏฺฐายาสนา ภควนฺตํ อภิวาเทตฺวา ปทกฺขิณํ กตฺวา เสนาสนํ สํสาเมตฺวา ปตฺตจีวรมาทาย เยน สุนาปรนฺโต ชนปโท เตน จาริกํ ปกฺกามิ, อนุปุพฺเพน จาริกํ จรมาโน เยน สุนาปรนฺโต ชนปโท ตทวสริ, ตตฺร สุทํ อายสฺมา ปุณฺโณ สุนาปรนฺตสฺมิํ ชนปเท วิหรติ.\csromanspacer{} อถ โข อายสฺมา ปุณฺโณ เตเนวนฺตรวสฺเสน ปญฺจมตฺตานิ อุปาสกสตานิ ปฏิเวเทสิ\footnote{ปฏิปาเทสิ (สี, อิ), ปฏิเทเสสิ (สฺยา, กํ)}, เตเนวนฺตรวสฺเสน ปญฺจมตฺตานิ อุปาสิกาสตานิ ปฏิเวเทสิ, เตเนวนฺตรวสฺเสน ติสฺโส วิชฺชา สจฺฉากาสิ, เตเนวนฺตรวสฺเสน ปรินิพฺพายิ.}

\prose{อถ โข สมฺพหุลา ภิกฺขู เยน ภควา เตนุปสงฺกมิํสุ ฯเปฯ เอกมนฺตํ นิสินฺนา โข เต ภิกฺขู ภควนฺตํ เอตทโวจุํ “โย โส ภนฺเต ปุณฺโณ นาม กุลปุตฺโต ภควตา สํขิตฺเตน โอวาเทน โอวทิโต, โส กาลงฺกโต, ตสฺส กา คติ โก อภิสมฺปราโย”ติ.\csromanspacer{} ปณฺฑิโต ภิกฺขเว ปุณฺโณ กุลปุตฺโต\footnote{กุลปุตฺโต อโหสิ (สพฺพตฺถ)}, ปจฺจปาทิ\footnote{สจฺจวาที (สฺยา, กํ, ก)} ธมฺมสฺสานุธมฺมํ, น จ มํ ธมฺมาธิกรณํ วิเหเสสิ\footnote{วิเหเฐสิ (สี, สฺยา, กํ)}, ปรินิพฺพุโต ภิกฺขเว ปุณฺโณ กุลปุตฺโตติ.\csromanspacer{}\csromanspacer{}ปญฺจมํ.}

\gambhira{สฬายตนวคฺคสํยุตฺตปาฬิ}
\csromanmatikaanchor{mtk.128}
\csromantocmarkref{subsection}{6. พาหิยสุตฺต}{mtk.128}
\bookmark[dest={mtk.128},level=2]{6. พาหิยสุตฺต}
\csromanheader{2}{6. พาหิยสุตฺต}
\proseitem{89}{\csromanfolio{288}อถ โข อายสฺมา พาหิโย เยน ภควา เตนุปสงฺกมิ ฯเปฯ เอกมนฺตํ นิสินฺโน โข อายสฺมา พาหิโย ภควนฺตํ เอตทโวจ “สาธุ เม ภนฺเต ภควา สํขิตฺเตน ธมฺมํ เทเสตุ, ยมหํ ภควโต ธมฺมํ สุตฺวา เอโก วูปกฏฺโฐ อปฺปมตฺโต อาตาปี ปหิตตฺโต วิหเรยฺยนฺ”ติ.}

\prose{ตํ กิํ มญฺญสิ พาหิย, จกฺขุ นิจฺจํ วา อนิจฺจํ วาติ.\csromanspacer{} อนิจฺจํ ภนฺเต.\csromanspacer{} ยํ ปนานิจฺจํ, ทุกฺขํ วา ตํ สุขํ วาติ.\csromanspacer{} ทุกฺขํ ภนฺเต.\csromanspacer{} ยํ ปนานิจฺจํ ทุกฺขํ วิปริณามธมฺมํ, กลฺลํ นุ ตํ สมนุปสฺสิตุํ “เอตํ มม, เอโสหมสฺมิ, เอโส เม อตฺตา”ติ.\csromanspacer{} โน เหตํ ภนฺเต.\csromanspacer{} รูปา นิจฺจา วา อนิจฺจา วาติ.\csromanspacer{} อนิจฺจา ภนฺเต ฯเปฯ.\csromanspacer{} จกฺขุวิญฺญาณํ.\csromanspacer{} จกฺขุสมฺผสฺโส ฯเปฯ.\csromanspacer{} ยมฺปิทํ มโนสมฺผสฺสปจฺจยา อุปฺปชฺชติ เวทยิตํ สุขํ วา ทุกฺขํ วา อทุกฺขมสุขํ วา, ตมฺปิ นิจฺจํ วา อนิจฺจํ วาติ.\csromanspacer{} อนิจฺจํ ภนฺเต.\csromanspacer{} ยํ ปนานิจฺจํ, ทุกฺขํ วา ตํ สุขํ วาติ.\csromanspacer{} ทุกฺขํ ภนฺเต.\csromanspacer{} ยํ ปนานิจฺจํ ทุกฺขํ วิปริณามธมฺมํ, กลฺลํ นุ ตํ สมนุปสฺสิตุํ “เอตํ มม, เอโสหมสฺมิ, เอโส เม อตฺตา”ติ.\csromanspacer{} โน เหตํ ภนฺเต.\csromanspacer{} เอวํ ปสฺสํ พาหิย สุตวา อริยสาวโก จกฺขุสฺมิมฺปิ นิพฺพินฺทติ, รูเปสุปิ นิพฺพินฺทติ, จกฺขุวิญฺญาเณปิ นิพฺพินฺทติ, จกฺขุสมฺผสฺเสปิ นิพฺพินฺทติ ฯเปฯ.\csromanspacer{} ยมฺปิทํ มโนสมฺผสฺสปจฺจยา อุปฺปชฺชติ เวทยิตํ สุขํ วา ทุกฺขํ วา อทุกฺขมสุขํ วา, ตสฺมิมฺปิ นิพฺพินฺทติ, นิพฺพินฺทํ วิรชฺชติ, วิราคา วิมุจฺจติ, วิมุตฺตสฺมิํ “วิมุตฺตมฺ”อิติ ญาณํ โหติ, “ขีณา ชาติ, วุสิตํ พฺรหฺมจริยํ, กตํ กรณียํ, นาปรํ อิตฺถตฺตายา”ติ ปชานาตีติ.}

\prose{อถ โข อายสฺมา พาหิโย ภควโต ภาสิตํ อภินนฺทิตฺวา อนุโมทิตฺวา อุฏฺฐายาสนา ภควนฺตํ อภิวาเทตฺวา ปทกฺขิณํ กตฺวา ปกฺกามิ.\csromanspacer{} อถ โข อายสฺมา พาหิโย เอโก วูปกฏฺโฐ อปฺปมตฺโต อาตาปี ปหิตตฺโต วิหรนฺโต นจิรสฺเสว ยสฺสตฺถาย กุลปุตฺตา สมฺมเทว อคารสฺมา อนคาริยํ ปพฺพชนฺติ, ตทนุตฺตรํ พฺรหฺมจริยปริโยสานํ ทิฏฺเฐว ธมฺเม สยํ อภิญฺญา สจฺฉิกตฺวา อุปสมฺปชฺช วิหาสิ, “ขีณา ชาติ, วุสิตํ พฺรหฺมจริยํ, กตํ กรณียํ, นาปรํ อิตฺถตฺตายา”ติ อพฺภญฺญาสิ, อญฺญตโร จ ปนายสฺมา พาหิโย อรหตํ อโหสีติ.\csromanspacer{}\csromanspacer{}ฉฏฺฐํ.}

\csromanmatikaanchor{mtk.129}
\csromantocmarkref{subsection}{7-8. เอชาสุตฺต}{mtk.129}
\bookmark[dest={mtk.129},level=2]{7-8. เอชาสุตฺต}
\csromanheader{2}{1. สฬายตนสํยุตฺต \textbf{7. ปฐม-เอชาสุตฺต}}
\proseitem{90}{\csromanfolio{289}เอชา ภิกฺขเว โรโค, เอชา คณฺโฑ, เอชา สลฺลํ, ตสฺมาติห ภิกฺขเว ตถาคโต อเนโช วิหรติ วีตสลฺโล, ตสฺมาติห ภิกฺขเว ภิกฺขุ เจปิ อากงฺเขยฺย “อเนโช วิหเรยฺยํ\footnote{วิหเรยฺย (สี, อิ, ก)} วีตสลฺโล”ติ.\csromanspacer{} จกฺขุํ น มญฺเญยฺย, จกฺขุสฺมิํ น มญฺเญยฺย, จกฺขุโต น มญฺเญยฺย, จกฺขุ เมติ น มญฺเญยฺย.\csromanspacer{} รูเป น มญฺเญยฺย, รูเปสุ น มญฺเญยฺย, รูปโต น มญฺเญยฺย, รูปา เมติ น มญฺเญยฺย.\csromanspacer{} จกฺขุวิญฺญาณํ น มญฺเญยฺย, จกฺขุวิญฺญาณสฺมิํ น มญฺเญยฺย, จกฺขุวิญฺญาณโต น มญฺเญยฺย, จกฺขุวิญฺญาณํ เมติ น มญฺเญยฺย.\csromanspacer{} จกฺขุสมฺผสฺสํ น มญฺเญยฺย, จกฺขุสมฺผสฺสสฺมิํ น มญฺเญยฺย, จกฺขุสมฺผสฺสโต น มญฺเญยฺย, จกฺขุสมฺผสฺโส เมติ น มญฺเญยฺย.\csromanspacer{} ยมฺปิทํ จกฺขุสมฺผสฺสปจฺจยา อุปฺปชฺชติ เวทยิตํ สุขํ วา ทุกฺขํ วา อทุกฺขมสุขํ วา, ตมฺปิ น มญฺเญยฺย, ตสฺมิมฺปิ น มญฺเญยฺย, ตโตปิ น มญฺเญยฺย, ตํ เมติ น มญฺเญยฺย.}

\prose{โสตํ น มญฺเญยฺย.\csromanspacer{} ฆานํ น มญฺเญยฺย.\csromanspacer{} ชิวฺหํ น มญฺเญยฺย, ชิวฺหาย น มญฺเญยฺย, ชิวฺหาโต น มญฺเญยฺย, ชิวฺหา เมติ น มญฺเญยฺย.\csromanspacer{} รเส น มญฺเญยฺย.\csromanspacer{} ชิวฺหาวิญฺญาณํ น มญฺเญยฺย.\csromanspacer{} ชิวฺหาสมฺผสฺสํ น มญฺเญยฺย.\csromanspacer{} ยมฺปิทํ ชิวฺหาสมฺผสฺสปจฺจยา อุปฺปชฺชติ เวทยิตํ สุขํ วา ทุกฺขํ วา อทุกฺขมสุขํ วา, ตมฺปิ น มญฺเญยฺย, ตสฺมิมฺปิ น มญฺเญยฺย, ตโตปิ น มญฺเญยฺย, ตํ เมติ น มญฺเญยฺย.}

\prose{กายํ น มญฺเญยฺย.\csromanspacer{} มนํ น มญฺเญยฺย, มนสฺมิํ น มญฺเญยฺย, มนโต น มญฺเญยฺย, มโน เมติ น มญฺเญยฺย.\csromanspacer{} ธมฺเม น มญฺเญยฺย.\csromanspacer{} มโน วิญฺญาณํ.\csromanspacer{} มโนสมฺผสฺสํ.\csromanspacer{} ยมฺปิทํ มโนสมฺผสฺสปจฺจยา อุปฺปชฺชติ เวทยิตํ สุขํ วา ทุกฺขํ วา อทุกฺขมสุขํ วา, ตมฺปิ น มญฺเญยฺย, ตสฺมิมฺปิ น มญฺเญยฺย, ตโตปิ น มญฺเญยฺย, ตํ เมติ น มญฺเญยฺย.\csromanspacer{} สพฺพํ น มญฺเญยฺย, สพฺพสฺมิํ น มญฺเญยฺย, สพฺพโต น มญฺเญยฺย, สพฺพํ เมติ น มญฺเญยฺย.}

\prose{โส เอวํ อมญฺญมาโน น กิญฺจิปิ โลเก อุปาทิยติ, อนุปาทิยํ น ปริตสฺสติ, อปริตสฺสํ ปจฺจตฺตญฺเญว ปรินิพฺพายติ, “ขีณา ชาติ, วุสิตํ พฺรหฺมจริยํ, กตํ กรณียํ, นาปรํ อิตฺถตฺตายา”ติ ปชานาตีติ.\csromanspacer{}\csromanspacer{}สตฺตมํ.}

\gambhira{สฬายตนวคฺคสํยุตฺตปาฬิ}
\csromanheader{2}{8. ทุติย-เอชาสุตฺต}
\proseitem{91}{\csromanfolio{290}เอชา ภิกฺขเว โรโค, เอชา คณฺโฑ, เอชา สลฺลํ, ตสฺมาติห ภิกฺขเว ตถาคโต อเนโช วิหรติ วีตสลฺโล, ตสฺมาติห ภิกฺขเว ภิกฺขุ เจปิ อากงฺเขยฺย “อเนโช วิหเรยฺยํ วีตสลฺโล”ติ.\csromanspacer{} จกฺขุํ น มญฺเญยฺย, จกฺขุสฺมิํ น มญฺเญยฺย, จกฺขุโต น มญฺเญยฺย, จกฺขุ เมติ น มญฺเญยฺย.\csromanspacer{} รูเป น มญฺเญยฺย.\csromanspacer{} จกฺขุวิญฺญาณํ.\csromanspacer{} จกฺขุสมฺผสฺสํ.\csromanspacer{} ยมฺปิทํ จกฺขุสมฺผสฺสปจฺจยา อุปฺปชฺชติ เวทยิตํ สุขํ วา ทุกฺขํ วา อทุกฺขมสุขํ วา, ตมฺปิ น มญฺเญยฺย, ตสฺมิมฺปิ น มญฺเญยฺย, ตโตปิ น มญฺเญยฺย, ตํ เมติ น มญฺเญยฺย.\csromanspacer{} ยํ หิ ภิกฺขเว มญฺญติ, ยสฺมิํ มญฺญติ, ยโต มญฺญติ, ยํ เมติ มญฺญติ, ตโต ตํ โหติ อญฺญถา, อญฺญถาภาวี ภวสตฺโต โลโก ภวเมว อภินนฺทติ ฯเปฯ.}

\prose{ชิวฺหํ น มญฺเญยฺย, ชิวฺหาย น มญฺเญยฺย, ชิวฺหาโต น มญฺเญยฺย, ชิวฺหา เมติ น มญฺเญยฺย.\csromanspacer{} รเส น มญฺเญยฺย.\csromanspacer{} ชิวฺหาวิญฺญาณํ.\csromanspacer{} ชิวฺหาสมฺผสฺสํ.\csromanspacer{} ยมฺปิทํ ชิวฺหาสมฺผสฺสปจฺจยา อุปฺปชฺชติ เวทยิตํ สุขํ วา ทุกฺขํ วา อทุกฺขมสุขํ วา, ตมฺปิ น มญฺเญยฺย, ตสฺมิมฺปิ น มญฺเญยฺย, ตโตปิ น มญฺเญยฺย, ตํ เมติ น มญฺเญยฺย.\csromanspacer{} ยํ หิ ภิกฺขเว มญฺญติ, ยสฺมิํ มญฺญติ, ยโต มญฺญติ, ยํ เมติ มญฺญติ, ตโต ตํ โหติ อญฺญถา, อญฺญถาภาวี ภวสตฺโต โลโก ภวเมว อภินนฺทติ ฯเปฯ.}

\prose{มนํ น มญฺเญยฺย, มนสฺมิํ น มญฺเญยฺย, มนโต น มญฺเญยฺย, มโน เมติ น มญฺเญยฺย.\csromanspacer{} มโนวิญฺญาณํ.\csromanspacer{} มโนสมฺผสฺสํ.\csromanspacer{} ยมฺปิทํ มโนสมฺผสฺสปจฺจยา อุปฺปชฺชติ เวทยิตํ สุขํ วา ทุกฺขํ วา อทุกฺขมสุขํ วา, ตมฺปิ น มญฺเญยฺย, ตสฺมิมฺปิ มญฺเญยฺย, ตโตปิ น มญฺเญยฺย, ตํ เมติ น มญฺเญยฺย.\csromanspacer{} ยํ หิ ภิกฺขเว มญฺญติ, ยสฺมิํ มญฺญติ, ยโต มญฺญติ, ยํ เมติ มญฺญติ, ตโต ตํ โหติ อญฺญถา, อญฺญถาภาวี ภวสตฺโต โลโก ภวเมว อภินนฺทติ.}

\prose{ยาวตา ภิกฺขเว ขนฺธธาตุ-อายตนา, ตมฺปิ น มญฺเญยฺย, ตสฺมิมฺปิ น มญฺเญยฺย, ตโตปิ น มญฺเญยฺย, ตํ เมติ น มญฺเญยฺย.\csromanspacer{} โส เอวํ อมญฺญมาโน น กิญฺจิ โลเก อุปาทิยติ, อนุปาทิยํ น ปริตสฺสติ, อปริตสฺสํ ปจฺจตฺตญฺเญว ปรินิพฺพายติ, “ขีณา ชาติ, วุสิตํ พฺรหฺมจริยํ, กตํ กรณียํ, นาปรํ อิตฺถตฺตายา”ติ ปชานาตีติ.\csromanspacer{}\csromanspacer{}อฏฺฐมํ.}

\csromanheader{2}{1. สฬายตนสํยุตฺต \textbf{9. ปฐมทฺวยสุตฺต}}
\csromanmatikaanchor{mtk.130}
\csromantocmarkref{subsection}{9-10. ทฺวยสุตฺต}{mtk.130}
\bookmark[dest={mtk.130},level=2]{9-10. ทฺวยสุตฺต}
\proseitem{92}{\csromanfolio{291}ทฺวยํ โว ภิกฺขเว เทเสสฺสามิ, ตํ สุณาถ.\csromanspacer{} กิญฺจ ภิกฺขเว ทฺวยํ, จกฺขุญฺเจว รูปา จ, โสตญฺเจว สทฺทา จ, ฆานญฺเจว คนฺธา จ, ชิวฺหา เจว รสา จ, กาโย เจว โผฏฺฐพฺพา จ, มโน เจว ธมฺมา จ.\csromanspacer{} อิทํ วุจฺจติ ภิกฺขเว ทฺวยํ.}

\prose{โย ภิกฺขเว เอวํ วเทยฺย “อหเมตํ ทฺวยํ ปจฺจกฺขาย อญฺญํ ทฺวยํ ปญฺญเปสฺสามี”ติ.\csromanspacer{} ตสฺส วาจาวตฺถุกเมวสฺส, ปุฏฺโฐ จ น สมฺปาเยยฺย, อุตฺตริํ จ วิฆาตํ อาปชฺเชยฺย.\csromanspacer{} ตํ กิสฺส เหตุ, ยถา ตํ ภิกฺขเว อวิสยสฺมินฺติ.\csromanspacer{}\csromanspacer{}นวมํ.}

\csromanheader{2}{10. ทุติยทฺวยสุตฺต}
\proseitem{93}{ทฺวยํ ภิกฺขเว ปฏิจฺจ วิญฺญาณํ สมฺโภติ.\csromanspacer{} กถญฺจ ภิกฺขเว ทฺวยํ ปฏิจฺจ วิญฺญาณํ สมฺโภติ, จกฺขุญฺจ ปฏิจฺจ รูเป จ อุปฺปชฺชติ จกฺขุวิญฺญาณํ.\csromanspacer{} จกฺขุ อนิจฺจํ วิปริณามิ อญฺญถาภาวิ.\csromanspacer{} รูปา อนิจฺจา วิปริณามิโน อญฺญถาภาวิโน.\csromanspacer{} อิตฺเถตํ ทฺวยํ จลญฺเจว พฺยถญฺจ อนิจฺจํ วิปริณามิ อญฺญถาภาวิ.\csromanspacer{} จกฺขุวิญฺญาณํ อนิจฺจํ วิปริณามิ อญฺญถาภาวิ.\csromanspacer{} โยปิ เหตุ โยปิ ปจฺจโย จกฺขุวิญฺญาณสฺส อุปฺปาทาย, โสปิ เหตุ โสปิ ปจฺจโย อนิจฺโจ วิปริณามี อญฺญถาภาวี.\csromanspacer{} อนิจฺจํ โข ปน ภิกฺขเว ปจฺจยํ ปฏิจฺจ อุปฺปนฺนํ จกฺขุวิญฺญาณํ, กุโต นิจฺจํ ภวิสฺสติ.\csromanspacer{} ยา โข ภิกฺขเว อิเมสํ ติณฺณํ ธมฺมานํ สงฺคติ สนฺนิปาโต สมวาโย, อยํ วุจฺจติ จกฺขุสมฺผสฺโส.\csromanspacer{} จกฺขุสมฺผสฺโสปิ อนิจฺโจ วิปริณามี อญฺญถาภาวี.\csromanspacer{} โยปิ เหตุ โยปิ ปจฺจโย จกฺขุสมฺผสฺสสฺส อุปฺปาทาย, โสปิ เหตุ โสปิ ปจฺจโย อนิจฺโจ วิปริณามี อญฺญถาภาวี.\csromanspacer{} อนิจฺจํ โข ปน ภิกฺขเว ปจฺจยํ ปฏิจฺจ อุปฺปนฺโน จกฺขุสมฺผสฺโส, กุโต นิจฺโจ ภวิสฺสติ.\csromanspacer{} ผุฏฺโฐ ภิกฺขเว เวเทติ, ผุฏฺโฐ เจเตติ, ผุฏฺโฐ สญฺชานาติ.\csromanspacer{} อิตฺเถเตปิ ธมฺมา จลา เจว พฺยถา จ อนิจฺจา วิปริณามิโน อญฺญถาภาวิโน ฯเปฯ.}

\prose{ชิวฺหญฺจ ปฏิจฺจ รเส จ อุปฺปชฺชติ ชิวฺหาวิญฺญาณํ, ชิวฺหา อนิจฺจา วิปริณามี อญฺญถาภาวี\footnote{วิปริณามินี อญฺญถาภาวินี (?)}.\csromanspacer{} รสา อนิจฺจา วิปริณามิโน อญฺญถาภาวิโน.\csromanspacer{} อิตฺเถตํ ทฺวยํ จลญฺเจว พฺยถญฺจ อนิจฺจํ วิปริณามิ อญฺญถาภาวิ.\csromanspacer{} ชิวฺหาวิญฺญาณํ อนิจฺจํ วิปริณามิ อญฺญถาภาวิ.\csromanspacer{} โยปิ เหตุ โยปิ ปจฺจโย ชิวฺหาวิญฺญาณสฺส}

\gambhira{สฬายตนวคฺคสํยุตฺตปาฬิ}
\prose{\csromanfolio{292}อุปฺปาทาย, โสปิ เหตุ โสปิ ปจฺจโย อนิจฺโจ วิปริณามี อญฺญถาภาวี.\csromanspacer{} อนิจฺจํ โข ปน ภิกฺขเว ปจฺจยํ ปฏิจฺจ อุปฺปนฺนํ ชิวฺหาวิญฺญาณํ, กุโต นิจฺจํ ภวิสฺสติ.\csromanspacer{} ยา โข ภิกฺขเว อิเมสํ ติณฺณํ ธมฺมานํ สงฺคติ สนฺนิปาโต สมวาโย, อยํ วุจฺจติ ชิวฺหาสมฺผสฺโส, ชิวฺหาสมฺผสฺโสปิ อนิจฺโจ วิปริณามี อญฺญถาภาวี.\csromanspacer{} โยปิ เหตุ โยปิ ปจฺจโย ชิวฺหาสมฺผสฺสสฺส อุปฺปาทาย, โสปิ เหตุ โสปิ ปจฺจโย อนิจฺโจ วิปริณามี อญฺญถาภาวี.\csromanspacer{} อนิจฺจํ โข ปน ภิกฺขเว ปจฺจยํ ปฏิจฺจ อุปฺปนฺโน ชิวฺหาสมฺผสฺโส, กุโต นิจฺโจ ภวิสฺสติ.\csromanspacer{} ผุฏฺโฐ ภิกฺขเว เวเทติ, ผุฏฺโฐ เจเตติ, ผุฏฺโฐ สญฺชานาติ, อิตฺเถเตปิ ธมฺมา จลา เจว พฺยถา จ อนิจฺจา วิปริณามิโน อญฺญถาภาวิโน ฯเปฯ.}

\prose{มนญฺจ ปฏิจฺจ ธมฺเม จ อุปฺปชฺชติ มโนวิญฺญาณํ, มโน อนิจฺโจ วิปริณามี อญฺญถาภาวี.\csromanspacer{} ธมฺมา อนิจฺจา วิปริณามิโน อญฺญถาภาวิโน.\csromanspacer{} อิตฺเถตํ ทฺวยํ จลญฺเจว พฺยถญฺจ อนิจฺจํ วิปริณามิ อญฺญถาภาวิ.\csromanspacer{} มโนวิญฺญาณํ อนิจฺจํ วิปริณามิ อญฺญถาภาวิ.\csromanspacer{} โยปิ เหตุ โยปิ ปจฺจโย มโนวิญฺญาณสฺส อุปฺปาทาย, โสปิ เหตุ โสปิ ปจฺจโย อนิจฺโจ วิปริณามี อญฺญถาภาวี.\csromanspacer{} อนิจฺจํ โข ปน ภิกฺขเว ปจฺจยํ ปฏิจฺจ อุปฺปนฺนํ มโนวิญฺญาณํ, กุโต นิจฺจํ ภวิสฺสติ.\csromanspacer{} ยา โข ภิกฺขเว อิเมสํ ติณฺณํ ธมฺมานํ สงฺคติ สนฺนิปาโต สมวาโย, อยํ วุจฺจติ มโนสมฺผสฺโส, มโนสมฺผสฺโสปิ อนิจฺโจ วิปริณามี อญฺญถาภาวี.\csromanspacer{} โยปิ เหตุ โยปิ ปจฺจโย มโนสมฺผสฺสสฺส อุปฺปาทาย, โสปิ เหตุ โสปิ ปจฺจโย อนิจฺโจ วิปริณามี อญฺญถาภาวี.\csromanspacer{} อนิจฺจํ โข ปน ภิกฺขเว ปจฺจยํ ปฏิจฺจ อุปฺปนฺโน มโนสมฺผสฺโส, กุโต นิจฺโจ ภวิสฺสติ.\csromanspacer{} ผุฏฺโฐ ภิกฺขเว เวเทติ, ผุฏฺโฐ เจเตติ, ผุฏฺโฐ สญฺชานาติ.\csromanspacer{} อิตฺเถเตปิ ธมฺมา จลา เจว พฺยถา จ อนิจฺจา วิปริณามิโน อญฺญถาภาวิโน.\csromanspacer{} เอวํ โข ภิกฺขเว ทฺวยํ ปฏิจฺจ วิญฺญาณํ สมฺโภตีติ.\csromanspacer{}\csromanspacer{}ทสมํ.\csromanspacer{} ฉนฺนวคฺโค จตุตฺโถ.}

\titlehead{\textbf{ตสฺสุทฺทานํ}}
\csromangathameasure{ปโลกสุญฺญา สํขิตฺตํ, ฉนฺโน ปุณฺโณ จ พาหิโย. \\ เอเชน จ ทุเว วุตฺตา, ทฺวเยหิ อปเร ทุเวติ.}
\csromangathagroup{\csromangathastanza{ปโลกสุญฺญา สํขิตฺตํ, ฉนฺโน ปุณฺโณ จ พาหิโย. \\ เอเชน จ ทุเว วุตฺตา, ทฺวเยหิ อปเร ทุเวติ.}}

\csromanmatikaanchor{mtk.131}
\csromantocmarknopage{section}{(10) 5. สฬวคฺค}{mtk.131}
\bookmark[dest={mtk.131},level=1]{(10) 5. สฬวคฺค}
\csromanheader{1}{1. สฬายตนสํยุตฺต \textbf{(10) 5. สฬวคฺค}}
\csromanmatikaanchor{mtk.132}
\csromantocmarkref{subsection}{1. อทนฺต-อคุตฺตสุตฺต}{mtk.132}
\bookmark[dest={mtk.132},level=2]{1. อทนฺต-อคุตฺตสุตฺต}
\csromanheader{2}{1. อทนฺต-อคุตฺตสุตฺต}
\proseitem{94}{\csromanfolio{293}สาวตฺถินิทานํ.\csromanspacer{} ฉยิเม ภิกฺขเว ผสฺสายตนา อทนฺตา อคุตฺตา อรกฺขิตา อสํวุตา ทุกฺขาธิวาหา โหนฺติ.\csromanspacer{} กตเม ฉ, จกฺขุ ภิกฺขเว ผสฺสายตนํ อทนฺตํ อคุตฺตํ อรกฺขิตํ อสํวุตํ ทุกฺขาธิวาหํ โหติ ฯเปฯ.\csromanspacer{} ชิวฺหา ภิกฺขเว ผสฺสายตนํ อทนฺตํ อคุตฺตํ อรกฺขิตํ อสํวุตํ ทุกฺขาธิวาหํ โหติ ฯเปฯ.\csromanspacer{} มโน ภิกฺขเว ผสฺสายตนํ อทนฺตํ อคุตฺตํ อรกฺขิตํ อสํวุตํ ทุกฺขาธิวาหํ โหติ.\csromanspacer{} อิเม โข ภิกฺขเว ฉ ผสฺสายตนา อทนฺตา อคุตฺตา อรกฺขิตา อสํวุตา ทุกฺขาธิวาหา โหนฺติ.}

\prose{ฉยิเม ภิกฺขเว ผสฺสายตนา สุทนฺตา สุคุตฺตา สุรกฺขิตา สุสํวุตา สุขาธิวาหา โหนฺติ.\csromanspacer{} กตเม ฉ, จกฺขุ ภิกฺขเว ผสฺสายตนํ สุทนฺตํ สุคุตฺตํ สุรกฺขิตํ สุสํวุตํ สุขาธิวาหํ โหติ วฺ.\csromanspacer{} ชิวฺหา ภิกฺขเว ผสฺสายตนํ สุทนฺตํ สุคุตฺตํ สุรกฺขิตํ สุสํวุตํ สุขาธิวาหํ โหติ ฯเปฯ.\csromanspacer{} มโน ภิกฺขเว ผสฺสายตนํ สุทนฺตํ สุคุตฺตํ สุรกฺขิตํ สุสํวุตํ สุขาธิวาหํ โหติ.\csromanspacer{} อิเม โข ภิกฺขเว ฉ ผสฺสายตนา สุทนฺตา สุคุตฺตา สุรกฺขิตา สุสํวุตา สุขาธิวาหา โหนฺตีติ.\csromanspacer{} อิทมโวจ ภควา ฯเปฯ เอตทโวจ สตฺถา\csromandash{}}

\csromangathameasure{“สเฬว ผสฺสายตนานิ ภิกฺขโว, \\ อสํวุโต ยตฺถ ทุกฺขํ นิคจฺฉติ. \\ เตสญฺจ เย สํวรณํ อเวทิสุํ, \\ สทฺธาทุติยา วิหรนฺตานวสฺสุตา. \\ ทิสฺวาน รูปานิ มโนรมานิ, \\ อโถปิ ทิสฺวาน อมโนรมานิ. \\ มโนรเม ราคปถํ วิโนทเย, \\ น จาปฺปิยํ เมติ มนํ ปโทสเย. \\ สทฺทญฺจ สุตฺวา ทุภยํ ปิยาปฺปิยํ, \\ ปิยมฺหิ สทฺเท น สมุจฺฉิโต สิยา. \\ อโถปฺปิเย โทสคตํ วิโนทเย, \\ น จาปฺปิยํ เมติ มนํ ปโทสเย.}
\csromangathagroup{\csromangathastanza{“สเฬว\footnote{ฉเฬว (ก)} ผสฺสายตนานิ ภิกฺขโว, \\ อสํวุโต ยตฺถ ทุกฺขํ นิคจฺฉติ. \\ เตสญฺจ เย สํวรณํ อเวทิสุํ, \\ สทฺธาทุติยา วิหรนฺตานวสฺสุตา.}\csromangathastanza{ทิสฺวาน รูปานิ มโนรมานิ, \\ อโถปิ ทิสฺวาน อมโนรมานิ. \\ มโนรเม ราคปถํ วิโนทเย, \\ น จาปฺปิยํ เมติ มนํ ปโทสเย.}\csromangathastanza{สทฺทญฺจ สุตฺวา ทุภยํ ปิยาปฺปิยํ, \\ ปิยมฺหิ สทฺเท น สมุจฺฉิโต สิยา. \\ อโถปฺปิเย โทสคตํ วิโนทเย, \\ น จาปฺปิยํ เมติ มนํ ปโทสเย.}}

\gambhira{สฬายตนวคฺคสํยุตฺตปาฬิ}
\csromangathameasure{คนฺธญฺจ ฆตฺวา สุรภิํ มโนรมํ, \\ อโถปิ ฆตฺวา อสุจิํ อกนฺติยํ. \\ อกนฺติยสฺมิํ ปฏิฆํ วิโนทเย, \\ ฉนฺทานุนีโต น จ กนฺติเย สิยา. \\ รสญฺจ โภตฺวาน อสาทิตญฺจ สาทุํ, \\ อโถปิ โภตฺวาน อสาทุเมกทา. \\ สาทุํ รสํ นาชฺโฌสาย ภุญฺเช, \\ วิโรธมาสาทุสุ โนปทํสเย. \\ ผสฺเสน ผุฏฺโฐ น สุเขน มชฺเช, \\ ทุกฺเขน ผุฏฺโฐปิ น สมฺปเวเธ. \\ ผสฺสทฺวยํ สุขทุกฺเข อุเปกฺเข. \\ อนานุรุทฺโธ อวิรุทฺธ เกนจิ. \\ ปปญฺจสญฺญา อิตรีตรา นรา, \\ ปปญฺจยนฺตา อุปยนฺติ สญฺญิโน. \\ มโนมยํ เคหสิตญฺจ สพฺพํ, \\ ปนุชฺช เนกฺขมฺมสิตํ อิรียติ. \\ เอวํ มโน ฉสฺสุ ยทา สุภาวิโต, \\ ผุฏฺฐสฺส จิตฺตํ น วิกมฺปเต กฺวจิ. \\ เต ราคโทเส อภิภุยฺย ภิกฺขโว. \\ ภวตฺถ ชาติมรณสฺส ปารคา”ติ. ปฐมํ.}
\csromangathagroup{\csromangathastanza{\csromanfolio{294}คนฺธญฺจ ฆตฺวา สุรภิํ มโนรมํ, \\ อโถปิ ฆตฺวา อสุจิํ อกนฺติยํ. \\ อกนฺติยสฺมิํ ปฏิฆํ วิโนทเย, \\ ฉนฺทานุนีโต น จ กนฺติเย สิยา.}\csromangathastanza{รสญฺจ โภตฺวาน อสาทิตญฺจ สาทุํ, \\ อโถปิ โภตฺวาน อสาทุเมกทา. \\ สาทุํ รสํ นาชฺโฌสาย ภุญฺเช, \\ วิโรธมาสาทุสุ โนปทํสเย.}\csromangathastanza{ผสฺเสน ผุฏฺโฐ น สุเขน มชฺเช\footnote{มชฺเฌ (สฺยา, กํ, อิ)}, \\ ทุกฺเขน ผุฏฺโฐปิ น สมฺปเวเธ. \\ ผสฺสทฺวยํ สุขทุกฺเข อุเปกฺเข. \\ อนานุรุทฺโธ อวิรุทฺธ เกนจิ.}\csromangathastanza{ปปญฺจสญฺญา อิตรีตรา นรา, \\ ปปญฺจยนฺตา อุปยนฺติ สญฺญิโน. \\ มโนมยํ เคหสิตญฺจ สพฺพํ, \\ ปนุชฺช เนกฺขมฺมสิตํ อิรียติ.}\csromangathastanza{เอวํ มโน ฉสฺสุ ยทา สุภาวิโต, \\ ผุฏฺฐสฺส จิตฺตํ น วิกมฺปเต กฺวจิ. \\ เต ราคโทเส อภิภุยฺย ภิกฺขโว. \\ ภวตฺถ\footnote{ภวถ (สี, สฺยา, กํ)} ชาติมรณสฺส ปารคา”ติ.\csromanspacer{} ปฐมํ.}}

\csromanmatikaanchor{mtk.133}
\csromantocmarkref{subsection}{2. มาลุกฺยปุตฺตสุตฺต}{mtk.133}
\bookmark[dest={mtk.133},level=2]{2. มาลุกฺยปุตฺตสุตฺต}
\csromanheader{2}{2. มาลุกฺยปุตฺตสุตฺต}
\proseitem{95}{อถ โข อายสฺมา มาลุกฺยปุตฺโต\footnote{มาลุงฺกฺยปุตฺโต (สี)} เยน ภควา เตนุปสงฺกมิ ฯเปฯ เอกมนฺตํ นิสินฺโน โข อายสฺมา มาลุกฺยปุตฺโต ภควนฺตํ เอตทโวจ “สาธุ เม ภนฺเต ภควา สํขิตฺเตน ธมฺมํ เทเสตุ, ยมหํ ภควโต ธมฺมํ สุตฺวา เอโก วูปกฏฺโฐ อปฺปมตฺโต อาตาปี ปหิตตฺโต วิหเรยฺยนฺ”ติ.}

\csromanheader{2}{1. สฬายตนสํยุตฺต}
\prose{\csromanfolio{295}เอตฺถ ทานิ มาลุกฺยปุตฺต กิํ ทหเร ภิกฺขู วกฺขาม “ยตฺร หิ นาม ตฺวํ ภิกฺขุ ชิณฺโณ วุทฺโธ มหลฺลโก อทฺธคโต วโย-อนุปฺปตฺโต สํขิตฺเตน โอวาทํ ยาจสี”ติ.\csromanspacer{} กิญฺจาปาหํ ภนฺเต ชิณฺโณ วุทฺโธ มหลฺลโก อทฺธคโต วโย-อนุปฺปตฺโต, เทเสตุ เม ภนฺเต ภควา สํขิตฺเตน ธมฺมํ เทเสตุ สุคโต สํขิตฺเตน ธมฺมํ, อปฺเปว นามาหํ ภควโต ภาสิตสฺส อตฺถํ อาชาเนยฺยํ, อปฺเปว นามาหํ ภควโต ภาสิตสฺส ทายาโท อสฺสนฺติ.}

\prose{ตํ กิํ มญฺญสิ มาลุกฺยปุตฺต, เย เต จกฺขุวิญฺเญยฺยา รูปา อทิฏฺฐา อทิฏฺฐปุพฺพา, น จ ปสฺสสิ, น จ เต โหติ ปสฺเสยฺยนฺติ, อตฺถิ เต ตตฺถ ฉนฺโท วา ราโค วา เปมํ วาติ.\csromanspacer{} โน เหตํ ภนฺเต.}

\prose{เย เต โสตวิญฺเญยฺยา สทฺทา อสฺสุตา อสฺสุตปุพฺพา, น จ สุณาสิ, น จ เต โหติ สุเณยฺยนฺติ, อตฺถิ เต ตตฺถ ฉนฺโท วา ราโค วา เปมํ วาติ.\csromanspacer{} โน เหตํ ภนฺเต.}

\prose{เย เต ฆานวิญฺเญยฺยา คนฺธา อฆายิตา อฆายิตปุพฺพา, น จ ฆายสิ, น จ เต โหติ ฆาเยยฺยนฺติ, อตฺถิ เต ตตฺถ ฉนฺโท วา ราโค วา เปมํ วาติ.\csromanspacer{} โน เหตํ ภนฺเต.}

\prose{เย เต ชิวฺหาวิญฺเญยฺยา รสา อสายิตา อสายิตปุพฺพา, น จ สายสิ, น จ เต โหติ สาเยยฺยนฺติ, อตฺถิ เต ตตฺถ ฉนฺโท วา ราโค วา เปมํ วาติ.\csromanspacer{} โน เหตํ ภนฺเต.}

\prose{เย เต กายวิญฺเญยฺยา โผฏฺฐพฺพา อสํผุฏฺฐา อสํผุฏฺฐปุพฺพา, น จ ผุสสิ, น จ เต โหติ ผุเสยฺยนฺติ, อตฺถิ เต ตตฺถ ฉนฺโท วา ราโค วา เปมํ วาติ.\csromanspacer{} โน เหตํ ภนฺเต.}

\prose{เย เต มโนวิญฺเญยฺยา ธมฺมา อวิญฺญาตา อวิญฺญาตปุพฺพา, น จ วิชานาสิ, น จ เต โหติ วิชาเนยฺยนฺติ, อตฺถิ เต ตตฺถ ฉนฺโท วา ราโค วา เปมํ วาติ.\csromanspacer{} โน เหตํ ภนฺเต.}

\prose{เอตฺถ จ เต มาลุกฺยปุตฺต ทิฏฺฐสุตมุตวิญฺญาตพฺเพสุ ธมฺเมสุ ทิฏฺเฐ ทิฏฺฐมตฺตํ ภวิสฺสติ, สุเต สุตมตฺตํ ภวิสฺสติ, มุเต มุตมตฺตํ ภวิสฺสติ, วิญฺญาเต วิญฺญาตมตฺตํ ภวิสฺสติ.\csromanspacer{} ยโต โข เต มาลุกฺยปุตฺต ทิฏฺฐสุตมุตวิญฺญาตพฺเพสุ ธมฺเมสุ ทิฏฺเฐ ทิฏฺฐมตฺตํ ภวิสฺสติ, สุเต สุตมตฺตํ}

\gambhira{สฬายตนวคฺคสํยุตฺตปาฬิ}
\prose{\csromanfolio{296}ภวิสฺสติ, มุเต มุตมตฺตํ ภวิสฺสติ, วิญฺญาเต วิญฺญาตมตฺตํ ภวิสฺสติ, ตโต ตฺวํ มาลุกฺยปุตฺต น เตน.\csromanspacer{} ยโต ตฺวํ มาลุกฺยปุตฺต น เตน, ตโต ตฺวํ มาลุกฺยปุตฺต น ตตฺถ.\csromanspacer{} ยโต ตฺวํ มาลุกฺยปุตฺต น ตตฺถ, ตโต ตฺวํ มาลุกฺยปุตฺต เนวิธ น หุรํ น อุภยมนฺตเรน, เอเสวนฺโต ทุกฺขสฺสาติ.}

\prose{อิมสฺส ขฺวาหํ ภนฺเต ภควตา สํขิตฺเตน ภาสิตสฺส วิตฺถาเรน อตฺถํ อาชานามิ\csromandash{}}

\csromangathameasure{รูปํ ทิสฺวา สติ มุฏฺฐา, ปิยํ นิมิตฺตํ มนสิ กโรโต. \\ สารตฺตจิตฺโต เวเทติ, ตญฺจ อชฺโฌส ติฏฺฐติ. \\ ตสฺส วฑฺฒนฺติ เวทนา, อเนกา รูปสมฺภวา. \\ อภิชฺฌา จ วิเหสา จ, จิตฺตมสฺสูปหญฺญติ. \\ เอวํ อาจินโต ทุกฺขํ, อารา นิพฺพาน วุจฺจติ. \\ สทฺทํ สุตฺวา สติ มุฏฺฐา, ปิยํ นิมิตฺตํ มนสิ กโรโต. \\ สารตฺตจิตฺโต เวเทติ, ตญฺจ อชฺโฌส ติฏฺฐติ. \\ ตสฺส วฑฺฒนฺติ เวทนา, อเนกา สทฺทสมฺภวา. \\ อภิชฺฌา จ วิเหสา จ, จิตฺตมสฺสูปหญฺญติ. \\ เอวํ อาจินโต ทุกฺขํ, อารา นิพฺพาน วุจฺจติ. \\ คนฺธํ ฆตฺวา สติ มุฏฺฐา, ปิยํ นิมิตฺตํ มนสิ กโรโต. \\ สารตฺตจิตฺโต เวเทติ, ตญฺจ อชฺโฌส ติฏฺฐติ. \\ ตสฺส วฑฺฒนฺติ เวทนา, อเนกา คนฺธสมฺภวา. \\ อภิชฺฌา จ วิเหสา จ, จิตฺตมสฺสูปหญฺญติ. \\ เอวํ อาจินโต ทุกฺขํ, อารา นิพฺพาน วุจฺจติ. \\ รสํ โภตฺวา สติ มุฏฺฐา, ปิยํ นิมิตฺตํ มนสิ กโรโต. \\ สารตฺตจิตฺโต เวเทติ, ตญฺจ อชฺโฌส ติฏฺฐติ. \\ ตสฺส วฑฺฒนฺติ เวทนา, อเนกา รสสมฺภวา. \\ อภิชฺฌา จ วิเหสา จ, จิตฺตมสฺสูปหญฺญติ. \\ เอวํ อาจินโต ทุกฺขํ, อารา นิพฺพาน วุจฺจติ.}
\csromangathagroup{\csromangathastanza{รูปํ ทิสฺวา สติ มุฏฺฐา, ปิยํ นิมิตฺตํ มนสิ กโรโต. \\ สารตฺตจิตฺโต เวเทติ, ตญฺจ อชฺโฌส\footnote{อชฺโฌสาย (สี)} ติฏฺฐติ.}\csromangathastanza{ตสฺส วฑฺฒนฺติ เวทนา, อเนกา รูปสมฺภวา. \\ อภิชฺฌา จ วิเหสา จ, จิตฺตมสฺสูปหญฺญติ.}\csromangathastanza{เอวํ อาจินโต ทุกฺขํ, อารา นิพฺพาน วุจฺจติ. \\ สทฺทํ สุตฺวา สติ มุฏฺฐา, ปิยํ นิมิตฺตํ มนสิ กโรโต.}\csromangathastanza{สารตฺตจิตฺโต เวเทติ, ตญฺจ อชฺโฌส ติฏฺฐติ. \\ ตสฺส วฑฺฒนฺติ เวทนา, อเนกา สทฺทสมฺภวา.}\csromangathastanza{อภิชฺฌา จ วิเหสา จ, จิตฺตมสฺสูปหญฺญติ. \\ เอวํ อาจินโต ทุกฺขํ, อารา นิพฺพาน วุจฺจติ.}\csromangathastanza{คนฺธํ ฆตฺวา สติ มุฏฺฐา, ปิยํ นิมิตฺตํ มนสิ กโรโต. \\ สารตฺตจิตฺโต เวเทติ, ตญฺจ อชฺโฌส ติฏฺฐติ.}\csromangathastanza{ตสฺส วฑฺฒนฺติ เวทนา, อเนกา คนฺธสมฺภวา. \\ อภิชฺฌา จ วิเหสา จ, จิตฺตมสฺสูปหญฺญติ.}\csromangathastanza{เอวํ อาจินโต ทุกฺขํ, อารา นิพฺพาน วุจฺจติ. \\ รสํ โภตฺวา สติ มุฏฺฐา, ปิยํ นิมิตฺตํ มนสิ กโรโต.}\csromangathastanza{สารตฺตจิตฺโต เวเทติ, ตญฺจ อชฺโฌส ติฏฺฐติ. \\ ตสฺส วฑฺฒนฺติ เวทนา, อเนกา รสสมฺภวา.}\csromangathastanza{อภิชฺฌา จ วิเหสา จ, จิตฺตมสฺสูปหญฺญติ. \\ เอวํ อาจินโต ทุกฺขํ, อารา นิพฺพาน วุจฺจติ.}}

\csromanheader{2}{1. สฬายตนสํยุตฺต}
\csromangathameasure{ผสฺสํ ผุสฺส สติ มุฏฺฐา, ปิยํ นิมิตฺตํ มนสิ กโรโต. \\ สารตฺตจิตฺโต เวเทติ, ตญฺจ อชฺโฌส ติฏฺฐติ. \\ ตสฺส วฑฺฒนฺติ เวทนา, อเนกา ผสฺสสมฺภวา. \\ อภิชฺฌา จ วิเหสา จ, จิตฺตมสฺสูปหญฺญติ. \\ เอวํ อาจินโต ทุกฺขํ, อารา นิพฺพาน วุจฺจติ. \\ ธมฺมํ ญตฺวา สติ มุฏฺฐา, ปิยํ นิมิตฺตํ มนสิ กโรโต. \\ สารตฺตจิตฺโต เวเทติ, ตญฺจ อชฺโฌส ติฏฺฐติ. \\ ตสฺส วฑฺฒนฺติ เวทนา, อเนกา ธมฺมสมฺภวา. \\ อภิชฺฌา จ วิเหสา จ, จิตฺตมสฺสูปหญฺญติ. \\ เอวํ อาจินโต ทุกฺขํ, อารา นิพฺพาน วุจฺจติ. \\ น โส รชฺชติ รูเปสุ, รูปํ ทิสฺวา ปฏิสฺสโต. \\ วิรตฺตจิตฺโต เวเทติ, ตญฺจ นาชฺโฌส ติฏฺฐติ. \\ ยถาสฺส ปสฺสโต รูปํ, เสวโต จาปิ เวทนํ. \\ ขียติ โนปจียติ, เอวํ โส จรตี สโต. \\ เอวํ อปจินโต ทุกฺขํ, สนฺติเก นิพฺพาน วุจฺจติ. \\ น โส รชฺชติ สทฺเทสุ, สทฺทํ สุตฺวา ปฏิสฺสโต. \\ วิรตฺตจิตฺโต เวเทติ, ตญฺจ นาชฺโฌส ติฏฺฐติ. \\ ยถาสฺส สุณโต สทฺทํ, เสวโต จาปิ เวทนํ. \\ ขียติ โนปจียติ, เอวํ โส จรตี สโต. \\ เอวํ อปจินโต ทุกฺขํ, สนฺติเก นิพฺพาน วุจฺจติ. \\ น โส รชฺชติ คนฺเธสุ, คนฺธํ ฆตฺวา ปฏิสฺสโต. \\ วิรตฺตจิตฺโต เวเทติ, ตญฺจ นาชฺโฌส ติฏฺฐติ. \\ ยถาสฺส ฆายโต คนฺธํ, เสวโต จาปิ เวทนํ. \\ ขียติ โนปจียติ, เอวํ โส จรตี สโต. \\ เอวํ อปจินโต ทุกฺขํ, สนฺติเก นิพฺพาน วุจฺจติ. \\ น โส รชฺชติ รเสสุ, รสํ โภตฺวา ปฏิสฺสโต. \\ วิรตฺตจิตฺโต เวเทติ, ตญฺจ นาชฺโฌส ติฏฺฐติ.}
\csromangathagroup{\csromangathastanza{\csromanfolio{297}ผสฺสํ ผุสฺส สติ มุฏฺฐา, ปิยํ นิมิตฺตํ มนสิ กโรโต. \\ สารตฺตจิตฺโต เวเทติ, ตญฺจ อชฺโฌส ติฏฺฐติ.}\csromangathastanza{ตสฺส วฑฺฒนฺติ เวทนา, อเนกา ผสฺสสมฺภวา. \\ อภิชฺฌา จ วิเหสา จ, จิตฺตมสฺสูปหญฺญติ.}\csromangathastanza{เอวํ อาจินโต ทุกฺขํ, อารา นิพฺพาน วุจฺจติ. \\ ธมฺมํ ญตฺวา สติ มุฏฺฐา, ปิยํ นิมิตฺตํ มนสิ กโรโต.}\csromangathastanza{สารตฺตจิตฺโต เวเทติ, ตญฺจ อชฺโฌส ติฏฺฐติ. \\ ตสฺส วฑฺฒนฺติ เวทนา, อเนกา ธมฺมสมฺภวา.}\csromangathastanza{อภิชฺฌา จ วิเหสา จ, จิตฺตมสฺสูปหญฺญติ. \\ เอวํ อาจินโต ทุกฺขํ, อารา นิพฺพาน วุจฺจติ.}\csromangathastanza{น โส รชฺชติ รูเปสุ, รูปํ ทิสฺวา ปฏิสฺสโต. \\ วิรตฺตจิตฺโต เวเทติ, ตญฺจ นาชฺโฌส ติฏฺฐติ.}\csromangathastanza{ยถาสฺส ปสฺสโต รูปํ, เสวโต จาปิ เวทนํ. \\ ขียติ โนปจียติ, เอวํ โส จรตี สโต.}\csromangathastanza{เอวํ อปจินโต ทุกฺขํ, สนฺติเก นิพฺพาน วุจฺจติ. \\ น โส รชฺชติ สทฺเทสุ, สทฺทํ สุตฺวา ปฏิสฺสโต.}\csromangathastanza{วิรตฺตจิตฺโต เวเทติ, ตญฺจ นาชฺโฌส ติฏฺฐติ. \\ ยถาสฺส สุณโต สทฺทํ, เสวโต จาปิ เวทนํ.}\csromangathastanza{ขียติ โนปจียติ, เอวํ โส จรตี สโต. \\ เอวํ อปจินโต ทุกฺขํ, สนฺติเก นิพฺพาน วุจฺจติ.}\csromangathastanza{น โส รชฺชติ คนฺเธสุ, คนฺธํ ฆตฺวา ปฏิสฺสโต. \\ วิรตฺตจิตฺโต เวเทติ, ตญฺจ นาชฺโฌส ติฏฺฐติ.}\csromangathastanza{ยถาสฺส ฆายโต คนฺธํ, เสวโต จาปิ เวทนํ. \\ ขียติ โนปจียติ, เอวํ โส จรตี สโต.}\csromangathastanza{เอวํ อปจินโต ทุกฺขํ, สนฺติเก นิพฺพาน วุจฺจติ. \\ น โส รชฺชติ รเสสุ, รสํ โภตฺวา ปฏิสฺสโต. \\ วิรตฺตจิตฺโต เวเทติ, ตญฺจ นาชฺโฌส ติฏฺฐติ.}}

\gambhira{สฬายตนวคฺคสํยุตฺตปาฬิ}
\csromangathameasure{ยถาสฺส สายโต รสํ, เสวโต จาปิ เวทนํ. \\ ขียติ โนปจียติ, เอวํ โส จรตี สโต. \\ เอวํ อปจินโต ทุกฺขํ, สนฺติเก นิพฺพาน วุจฺจติ. \\ น โส รชฺชติ ผสฺเสสุ, ผสฺสํ ผุสฺส ปฏิสฺสโต. \\ วิรตฺตจิตฺโต เวเทติ, ตญฺจ นาชฺโฌส ติฏฺฐติ. \\ ยถาสฺส ผุสโต ผสฺสํ, เสวโต จาปิ เวทนํ. \\ ขียติ โนปจียติ, เอวํ โส จรตี สโต. \\ เอวํ อปจินโต ทุกฺขํ, สนฺติเก นิพฺพาน วุจฺจติ. \\ น โส รชฺชติ ธมฺเมสุ, ธมฺมํ ญตฺวา ปฏิสฺสโต. \\ วิรตฺตจิตฺโต เวเทติ, ตญฺจ นาชฺโฌส ติฏฺฐติ. \\ ยถาสฺสชานโต ธมฺมํ, เสวโต จาปิ เวทนํ. \\ ขียติ โนปจียติ, เอวํ โส จรตี สโต. \\ เอวํ อปจินโต ทุกฺขํ, สนฺติเก นิพฺพาน วุจฺจตีติ.}
\csromangathagroup{\csromangathastanza{\csromanfolio{298}ยถาสฺส สายโต รสํ, เสวโต จาปิ เวทนํ. \\ ขียติ โนปจียติ, เอวํ โส จรตี สโต.}\csromangathastanza{เอวํ อปจินโต ทุกฺขํ, สนฺติเก นิพฺพาน วุจฺจติ. \\ น โส รชฺชติ ผสฺเสสุ, ผสฺสํ ผุสฺส ปฏิสฺสโต.}\csromangathastanza{วิรตฺตจิตฺโต เวเทติ, ตญฺจ นาชฺโฌส ติฏฺฐติ. \\ ยถาสฺส ผุสโต ผสฺสํ, เสวโต จาปิ เวทนํ.}\csromangathastanza{ขียติ โนปจียติ, เอวํ โส จรตี สโต. \\ เอวํ อปจินโต ทุกฺขํ, สนฺติเก นิพฺพาน วุจฺจติ.}\csromangathastanza{น โส รชฺชติ ธมฺเมสุ, ธมฺมํ ญตฺวา ปฏิสฺสโต. \\ วิรตฺตจิตฺโต เวเทติ, ตญฺจ นาชฺโฌส ติฏฺฐติ.}\csromangathastanza{ยถาสฺสชานโต ธมฺมํ, เสวโต จาปิ เวทนํ. \\ ขียติ โนปจียติ, เอวํ โส จรตี สโต. \\ เอวํ อปจินโต ทุกฺขํ, สนฺติเก นิพฺพาน วุจฺจตีติ.}}

\prose{อิมสฺส ขฺวาหํ ภนฺเต ภควตา สํขิตฺเตน ภาสิตสฺส เอวํ วิตฺถาเรน อตฺถํ อาชานามีติ.\csromanspacer{} สาธุ สาธุ มาลุกฺยปุตฺต, สาธุ โข ตฺวํ มาลุกฺยปุตฺต มยา สํขิตฺเตน ภาสิตสฺส วิตฺถาเรน อตฺถํ อาชานาสิ\csromandash{}}

\csromangathameasure{รูปํ ทิสฺวา สติ มุฏฺฐา, ปิยํ นิมิตฺตํ มนสิ กโรโต. \\ สารตฺตจิตฺโต เวเทติ, ตญฺจ อชฺโฌส ติฏฺฐติ. \\ ตสฺส วฑฺฒนฺติ เวทนา, อเนกา รูปสมฺภวา. \\ อภิชฺฌา จ วิเหสา จ, จิตฺตมสฺสูปหญฺญติ. \\ เอวํ อาจินโต ทุกฺขํ, อารา นิพฺพาน วุจฺจติ ฯเปฯ. \\ น โส รชฺชติ ธมฺเมสุ, ธมฺมํ ญตฺวา ปฏิสฺสโต. \\ วิรตฺตจิตฺโต เวเทติ, ตญฺจ นาชฺโฌส ติฏฺฐติ. \\ ยถาสฺส วิชานโต ธมฺมํ, เสวโต จาปิ เวทนํ. \\ ขียติ โนปจียติ, เอวํ โส จรตี สโต. \\ เอวํ อปจินโต ทุกฺขํ, สนฺติเก นิพฺพาน วุจฺจตีติ.}
\csromangathagroup{\csromangathastanza{รูปํ ทิสฺวา สติ มุฏฺฐา, ปิยํ นิมิตฺตํ มนสิ กโรโต. \\ สารตฺตจิตฺโต เวเทติ, ตญฺจ อชฺโฌส ติฏฺฐติ.}\csromangathastanza{ตสฺส วฑฺฒนฺติ เวทนา, อเนกา รูปสมฺภวา. \\ อภิชฺฌา จ วิเหสา จ, จิตฺตมสฺสูปหญฺญติ.}\csromangathastanza{เอวํ อาจินโต ทุกฺขํ, อารา นิพฺพาน วุจฺจติ ฯเปฯ. \\ น โส รชฺชติ ธมฺเมสุ, ธมฺมํ ญตฺวา ปฏิสฺสโต.}\csromangathastanza{วิรตฺตจิตฺโต เวเทติ, ตญฺจ นาชฺโฌส ติฏฺฐติ. \\ ยถาสฺส วิชานโต ธมฺมํ, เสวโต จาปิ เวทนํ.}\csromangathastanza{ขียติ โนปจียติ, เอวํ โส จรตี สโต. \\ เอวํ อปจินโต ทุกฺขํ, สนฺติเก นิพฺพาน วุจฺจตีติ.}}

\prose{อิมสฺส โข มาลุกฺยปุตฺต มยา สํขิตฺเตน ภาสิตสฺส เอวํ วิตฺถาเรน อตฺโถ ทฏฺฐพฺโพติ.}

\csromanheader{2}{1. สฬายตนสํยุตฺต}
\prose{\csromanfolio{299}อถ โข อายสฺมา มาลุกฺยปุตฺโต ภควโต ภาสิตํ อภินนฺทิตฺวา อนุโมทิตฺวา อุฏฺฐายาสนา ภควนฺตํ อภิวาเทตฺวา ปทกฺขิณํ กตฺวา ปกฺกามิ.\csromanspacer{} อถ โข อายสฺมา มาลุกฺยปุตฺโต เอโก วูปกฏฺโฐ อปฺปมตฺโต อาตาปี ปหิตตฺโต วิหรนฺโต นจิรสฺเสว ยสฺสตฺถาย กุลปุตฺตา สมฺมเทว อคารสฺมา อนคาริยํ ปพฺพชนฺติ, ตทนุตฺตรํ พฺรหฺมจริยปริโยสานํ ทิฏฺเฐว ธมฺเม สยํ อภิญฺญา สจฺฉิกตฺวา อุปสมฺปชฺช วิหาสิ, “ขีณา ชาติ, วุสิตํ พฺรหฺมจริยํ, กตํ กรณียํ, นาปรํ อิตฺถตฺตายา”ติ อพฺภญฺญาสิ.\csromanspacer{} อญฺญตโร จ ปนายสฺมา มาลุกฺยปุตฺโต อรหตํ อโหสีติ.\csromanspacer{}\csromanspacer{}ทุติยํ.}

\csromanmatikaanchor{mtk.134}
\csromantocmarkref{subsection}{3. ปริหานธมฺมสุตฺต}{mtk.134}
\bookmark[dest={mtk.134},level=2]{3. ปริหานธมฺมสุตฺต}
\csromanheader{2}{3. ปริหานธมฺมสุตฺต}
\proseitem{96}{ปริหานธมฺมญฺจ โว ภิกฺขเว เทเสสฺสามิ อปริหานธมฺมญฺจ ฉ จ อภิภายตนานิ, ตํ สุณาถ.\csromanspacer{} กถญฺจ ภิกฺขเว ปริหานธมฺโม โหติ.\csromanspacer{} อิธ ภิกฺขเว ภิกฺขุโน จกฺขุนา รูปํ ทิสฺวา อุปฺปชฺชนฺติ ปาปกา อกุสลา สรสงฺกปฺปา\footnote{อกุสลา ธมฺมา สรสงฺกปฺปา (สฺยา, กํ, อิ, ก) อุปริ อาสีวิสวคฺเค สตฺตมสุตฺต ปน “อกุสลา สรสงฺกปฺปา”เตฺว ว สพฺพตฺถ ทิสฺสติ.} สํโยชนิยา, ตญฺเจ ภิกฺขุ อธิวาเสติ, นปฺปชหติ, น วิโนเทติ, น พฺยนฺตีกโรติ\footnote{พฺยนฺติกโรติ (อิ), พฺยนฺติํ กโรติ (ก)}, น อนภาวํ คเมติ, เวทิตพฺพเมตํ ภิกฺขเว ภิกฺขุนา “ปริหายามิ กุสเลหิ ธมฺเมหิ, ปริหานํ เหตํ วุตฺตํ ภควตา”ติ ฯเปฯ.}

\prose{ปุน จปรํ ภิกฺขเว ภิกฺขุโน ชิวฺหาย รสํ สายิตฺวา อุปฺปชฺชนฺติ ฯเปฯ.\csromanspacer{} ปุน จปรํ ภิกฺขเว ภิกฺขุโน มนสา ธมฺมํ วิญฺญาย อุปฺปชฺชนฺติ ปาปกา อกุสลา สรสงฺกปฺปา สํโยชนิยา, ตญฺเจ ภิกฺขุ อธิวาเสติ, นปฺปชหติ, น วิโนเทติ, น พฺยนฺตีกโรติ, น อนภาวํ คเมติ, เวทิตพฺพเมตํ ภิกฺขเว ภิกฺขุนา “ปริหายามิ กุสเลหิ ธมฺเมหิ, ปริหานํ เหตํ วุตฺตํ ภควตา”ติ.\csromanspacer{} เอวํ โข ภิกฺขเว ปริหานธมฺโม โหติ.}

\prose{กถญฺจ ภิกฺขเว อปริหานธมฺโม โหติ.\csromanspacer{} อิธ ภิกฺขเว ภิกฺขุโน จกฺขุนา รูปํ ทิสฺวา อุปฺปชฺชนฺติ ปาปกา อกุสลา สรสงฺกปฺปา สํโยชนิยา, ตญฺเจ ภิกฺขุ นาธิวาเสติ, ปชหติ, วิโนเทติ, พฺยนฺตีกโรติ, อนภาวํ คเมติ, เวทิตพฺพเมตํ ภิกฺขเว ภิกฺขุนา “น ปริหายามิ กุสเลหิ ธมฺเมหิ, อปริหานํ เหตํ วุตฺตํ ภควตา”ติ ฯเปฯ.}

\gambhira{สฬายตนวคฺคสํยุตฺตปาฬิ}
\prose{\csromanfolio{300}ปุน จปรํ ภิกฺขเว ภิกฺขุโน ชิวฺหาย รสํ สายิตฺวา อุปฺปชฺชนฺติ ฯเปฯ.\csromanspacer{} ปุน จปรํ ภิกฺขเว ภิกฺขุโน มนสา ธมฺมํ วิญฺญาย อุปฺปชฺชนฺติ ปาปกา อกุสลา สรสงฺกปฺปา สํโยชนิยา, ตญฺเจ ภิกฺขุ นาธิวาเสติ, ปชหติ, วิโนเทติ, พฺยนฺตีกโรติ, อนภาวํ คเมติ, เวทิตพฺพเมตํ ภิกฺขเว ภิกฺขุนา “น ปริหายามิ กุสเลหิ ธมฺเมหิ, อปริหานํ เหตํ วุตฺตํ ภควตา”ติ.\csromanspacer{} เอวํ โข ภิกฺขเว อปริหานธมฺโม โหติ.}

\prose{กตมานิ จ ภิกฺขเว ฉ อภิภายตนานิ.\csromanspacer{} อิธ ภิกฺขเว ภิกฺขุโน จกฺขุนา รูปํ ทิสฺวา นุปฺปชฺชนฺติ ปาปกา อกุสลา สรสงฺกปฺปา สํโยชนิยา, เวทิตพฺพเมตํ ภิกฺขเว ภิกฺขุนา อภิภูตเมตํ อายตนํ, อภิภายตนํ เหตํ วุตฺตํ ภควตาติ ฯเปฯ.\csromanspacer{} ปุน จปรํ ภิกฺขเว ภิกฺขุโน มนสา ธมฺมํ วิญฺญาย นุปฺปชฺชนฺติ ปาปกา อกุสลา ธมฺมา สรสงฺกปฺปา สํโยชนิยา เวทิตพฺพเมตํ ภิกฺขเว ภิกฺขุนา “อภิภูตเมตํ อายตนํ, อภิภายตนํ เหตํ วุตฺตํ ภควตา”ติ.\csromanspacer{} อิมานิ วุจฺจนฺติ ภิกฺขเว ฉ อภิภายตนานีติ.\csromanspacer{}\csromanspacer{}ตติยํ.}

\csromanmatikaanchor{mtk.135}
\csromantocmarkref{subsection}{4. ปมาทวิหารีสุตฺต}{mtk.135}
\bookmark[dest={mtk.135},level=2]{4. ปมาทวิหารีสุตฺต}
\csromanheader{2}{4. ปมาทวิหารีสุตฺต}
\proseitem{97}{ปมาทวิหาริญฺจ โว ภิกฺขเว เทเสสฺสามิ อปฺปมาทวิหารญฺจ, ตํ สุณาถ.\csromanspacer{} กถญฺจ ภิกฺขเว ปมาทวิหารี โหติ, จกฺขุนฺทฺริยํ อสํวุตสฺส ภิกฺขเว วิหรโต จิตฺตํ พฺยาสิญฺจติ\footnote{พฺยาสิจฺจติ (สี, สฺยา, กํ)} จกฺขุวิญฺเญยฺเยสุ รูเปสุ, ตสฺส พฺยาสิตฺตจิตฺตสฺส ปาโมชฺชํ น โหติ, ปาโมชฺเช อสติ ปีติ น โหติ, ปีติยา อสติ ปสฺสทฺธิ น โหติ, ปสฺสทฺธิยา อสติ ทุกฺขํ โหติ, ทุกฺขิโน จิตฺตํ น สมาธิยติ, อสมาหิเต จิตฺเต ธมฺมา น ปาตุภวนฺติ, ธมฺมานํ อปาตุภาวา ปมาทวิหารี เตฺวว สงฺขํ คจฺฉติ ฯเปฯ ชิวฺหินฺทฺริยํ อสํวุตสฺส ภิกฺขเว วิหรโต จิตฺตํ พฺยาสิญฺจติ ชิวฺหาวิญฺเญยฺเยสุ รเสสุ, ตสฺส พฺยสิตฺตจิตฺตสฺส ฯเปฯ ปมาทวิหารี เตฺวว สงฺขํ คจฺฉติ ฯเปฯ มนินฺทฺริยํ อสํวุตสฺส ภิกฺขเว วิหรโต จิตฺตํ พฺยาสิญฺจติ มโนวิญฺเญยฺเยสุ ธมฺเมสุ, ตสฺส พฺยาสิตฺตจิตฺตสฺส ปาโมชฺชํ น โหติ, ปาโมชฺเช อสติ ปีติ น โหติ, ปีติยา อสติ ปสฺสทฺธิ น โหติ, ปสฺสทฺธิยา อสติ ทุกฺขํ โหติ, ทุกฺขิโน จิตฺตํ น สมาธิยติ, อสมาหิเต จิตฺเต ธมฺมา น}

\vspace{\tipitakaparskip}
\csromancenter{สฬายตนสํยุตฺต ปาตุภวนฺติ, ธมฺมานํ อปาตุภาวา ปมาทวิหารี เตฺวว สงฺขํ คจฺฉติ.\csromanspacer{} เอวํ โข ภิกฺขเว ปมาทวิหารี โหติ.}
\prose{\csromanfolio{301}กถญฺจ ภิกฺขเว อปฺปมาทวิหารี โหติ.\csromanspacer{} จกฺขุนฺทฺริยํ สํวุตสฺส ภิกฺขเว วิหรโต จิตฺตํ น พฺยาสิญฺจติ จกฺขุวิญฺเญยฺเยสุ รูเปสุ, ตสฺส อพฺยาสิตฺตจิตฺตสฺส ปาโมชฺชํ ชายติ, ปมุทิตสฺส ปีติ ชายติ, ปีติมนสฺส กาโย ปสฺสมฺภติ, ปสฺสทฺธกาโย สุขํ วิหรติ, สุขิโน จิตฺตํ สมาธิยติ, สมาหิเต จิตฺเต ธมฺมา ปาตุภวนฺติ, ธมฺมานํ ปาตุภาวา อปฺปมาทวิหารี เตฺวว สงฺขํ คจฺฉติ ฯเปฯ.\csromanspacer{} ชิวฺหินฺทฺริยํ สํวุตสฺส ภิกฺขเว วิหรโต จิตฺตํ น พฺยาสิญฺจติ ฯเปฯ อปฺปมาทวิหารี เตฺวว สงฺขํ คจฺฉติ.\csromanspacer{} มนินฺทฺริยํ สํวุตสฺส ภิกฺขเว วิหรโต จิตฺตํ น พฺยาสิญฺจติ มโนวิญฺเญยฺเยสุ ธมฺเมสุ, ตสฺส อพฺยาสิตฺตจิตฺตสฺส ปาโมชฺชํ ชายติ, ปมุทิตสฺส ปีติ ชายติ, ปีติมนสฺส กาโย ปสฺสมฺภติ, ปสฺสทฺธกาโย สุขํ วิหรติ, สุขิโน จิตฺตํ สมาธิยติ, สมาหิเต จิตฺเต ธมฺมา ปาตุภวนฺติ, ธมฺมานํ ปาตุภาวา อปฺปมาทวิหารี เตฺวว สงฺขํ คจฺฉติ.\csromanspacer{} เอวํ โข ภิกฺขเว อปฺปมาทวิหารี โหตีติ.\csromanspacer{}\csromanspacer{}จตุตฺถํ.}

\csromanmatikaanchor{mtk.136}
\csromantocmarkref{subsection}{5. สํวรสุตฺต}{mtk.136}
\bookmark[dest={mtk.136},level=2]{5. สํวรสุตฺต}
\csromanheader{2}{5. สํวรสุตฺต}
\proseitem{98}{สํวรญฺจ โว ภิกฺขเว เทเสสฺสามิ อสํวรญฺจ, ตํ สุณาถ.\csromanspacer{} กถญฺจ ภิกฺขเว อสํวโร โหติ.\csromanspacer{} สนฺติ ภิกฺขเว จกฺขุวิญฺเญยฺยา รูปา อิฏฺฐา กนฺตา มนาปา ปิยรูปา กามูปสํหิตา รชนียา, ตญฺเจ ภิกฺขุ อภินนฺทติ อภิวทติ อชฺโฌสาย ติฏฺฐติ, เวทิตพฺพเมตํ ภิกฺขเว ภิกฺขุนา “ปริหายามิ กุสเลหิ ธมฺเมหิ, ปริหานํ เหตํ วุตฺตํ ภควตา”ติ ฯเปฯ.\csromanspacer{} สนฺติ ภิกฺขเว ชิวฺหาวิญฺเญยฺยา รสา ฯเปฯ.\csromanspacer{} สนฺติ ภิกฺขเว มโนวิญฺเญยฺยา ธมฺมา อิฏฺฐา กนฺตา มนาปา ปิยรูปา กามูปสํหิตา รชนียา, ตญฺเจ ภิกฺขุ อภินนฺทติ อภิวทติ อชฺโฌสาย ติฏฺฐติ, เวทิตพฺพเมตํ ภิกฺขเว ภิกฺขุนา “ปริหายามิ กุสเลหิ ธมฺเมหิ, ปริหานํ เหตํ วุตฺตํ ภควตา”ติ.\csromanspacer{} เอวํ โข ภิกฺขเว อสํวโร โหติ.}

\prose{กตญฺจ ภิกฺขเว สํวโร โหติ.\csromanspacer{} สนฺติ ภิกฺขเว จกฺขุวิญฺเญยฺยา รูปา อิฏฺฐา กนฺตา มนาปา ปิยรูปา กามูปสํหิตา รชนียา, ตญฺเจ ภิกฺขุ นาภินนฺทติ นาภิวทติ นาชฺโฌสาย ติฏฺฐติ, เวทิตพฺพเมตํ ภิกฺขเว ภิกฺขุนา “น ปริหายามิ กุสเลหิ ธมฺเมหิ, อปริหานํ เหตํ วุตฺตํ}

\gambhira{สฬายตนวคฺคสํยุตฺตปาฬิ}
\csromanmatikaanchor{mtk.137}
\csromanmatikaanchor{mtk.139}
\csromantocmarkref{subsection}{6. สมาธิสุตฺต}{mtk.137}
\bookmark[dest={mtk.137},level=2]{6. สมาธิสุตฺต}
\prose{\csromanfolio{302}ภควตา”ติ ฯเปฯ.\csromanspacer{} สนฺติ ภิกฺขเว ชิวฺหาวิญฺเญยฺยา รสา ฯเปฯ.\csromanspacer{} สนฺติ ภิกฺขเว มโนวิญฺเญยฺยา ธมฺมา อิฏฺฐา กนฺตา มนาปา ปิยรูปา กามูปสํหิตา รชนียา, ตญฺเจ ภิกฺขุ นาภินนฺทติ นาภิวทติ นาชฺโฌสาย ติฏฺฐติ, เวทิตพฺพเมตํ ภิกฺขุนา “น ปริหายามิ กุสเลหิ ธมฺเมหิ, อปริหานํ เหตํ วุตฺตํ ภควตา”ติ.\csromanspacer{} เอวํ โข ภิกฺขเว สํวโร โหตีติ.\csromanspacer{}\csromanspacer{}ปญฺจมํ.}

\csromanheader{2}{6. สมาธิสุตฺต}
\proseitem{99}{สมาธิํ ภิกฺขเว ภาเวถ, สมาหิโต ภิกฺขเว ภิกฺขุ ยถาภูตํ ปชานาติ.\csromanspacer{} กิญฺจ ยถาภูตํ ปชานาติ, “จกฺขุ อนิจฺจนฺ”ติ ยถาภูตํ ปชานาติ, “รูปา อนิจฺจา”ติ ยถาภูตํ ปชานาติ, “จกฺขุวิญฺญาณํ อนิจฺจนฺ”ติ ยถาภูตํ ปชานาติ, “จกฺขุสมฺผสฺส อนิจฺโจ”ติ ยถาภูตํ ปชานาติ.\csromanspacer{} ยมฺปิทํ จกฺขุสมฺผสฺสปจฺจยา อุปฺปชฺชติ เวทยิตํ สุขํ วา ทุกฺขํ วา อทุกฺขมสุขํ วา, ตมฺปิ อนิจฺจนฺติ ยถาภูตํ ปชานาติ ฯเปฯ “มโน อนิจฺจนฺ”ติ ยถาภูตํ ปชานาติ.\csromanspacer{} ธมฺมา.\csromanspacer{} มโนวิญฺญาณํ.\csromanspacer{} มโนสมฺผสฺโส.\csromanspacer{} ยมฺปิทํ มโนสมฺผสฺสปจฺจยา อุปฺปชฺชติ เวทยิตํ สุขํ วา ทุกฺขํ วา อทุกฺขมสุขํ วา, ตมฺปิ อนิจฺจนฺติ ยถาภูตํ ปชานาติ.\csromanspacer{} สมาธิํ ภิกฺขเว ภาเวถ, สมาหิโต ภิกฺขเว ภิกฺขุ ยถาภูตํ ปชานาตีติ.\csromanspacer{}\csromanspacer{}ฉฏฺฐํ.}

\csromanmatikaanchor{mtk.138}
\csromantocmarkref{subsection}{7. ปฏิสลฺลานสุตฺต}{mtk.138}
\bookmark[dest={mtk.138},level=2]{7. ปฏิสลฺลานสุตฺต}
\csromantocmarkref{subsection}{8-9. นตุมฺหากสุตฺต}{mtk.139}
\bookmark[dest={mtk.139},level=2]{8-9. นตุมฺหากสุตฺต}
\csromanheader{2}{7. ปฏิสลฺลานสุตฺต}
\proseitem{100}{ปฏิสลฺลาเน\footnote{ปฏิสลฺลานํ (สี, อิ, ก), ปฏิสลฺลีนา (สฺยา, กํ)} ภิกฺขเว โยคมาปชฺชถ, ปฏิสลฺลีโน ภิกฺขเว ภิกฺขุ ยถาภูตํ ปชานาติ.\csromanspacer{} กิญฺจ ยถาภูตํ ปชานาติ, “จกฺขุ อนิจฺจนฺ”ติ ยถาภูตํ ปชานาติ, “รูปา อนิจฺจา”ติ ยถาภูตํ ปชานาติ, “จกฺขุวิญฺญาณํ อนิจฺจนฺ”ติ ยถาภูตํ ปชานาติ, “จกฺขุสมฺผสฺโส อนิจฺโจ”ติ ยถาภูตํ ปชานาติ ฯเปฯ.\csromanspacer{} ยมฺปิทํ มโนสมฺผสฺสปจฺจยา อุปฺปชฺชติ เวทยิตํ สุขํ วา ทุกฺขํ วา อทุกฺขมสุขํ วา, ตมฺปิ อนิจฺจนฺติ ยถาภูตํ ปชานาติ.\csromanspacer{} ปฏิสลฺลาเน ภิกฺขเว โยคมาปชฺชถ, ปฏิสลฺลีโน ภิกฺขเว ภิกฺขุ ยถาภูตํ ปชานาตีติ.\csromanspacer{}\csromanspacer{}สตฺตมํ.}

\csromanheader{2}{8. ปฐมนตุมฺหากสุตฺต}
\proseitem{101}{ยํ\footnote{ยมฺปิ (อิ, ก)} ภิกฺขเว น ตุมฺหากํ ตํ ปชหถ, กํ โว ปหีนํ หิตาย สุขาย ภวิสฺสติ.\csromanspacer{} กิญฺจ ภิกฺขเว น ตุมฺหากํ, จกฺขุ ภิกฺขเว น}

\vspace{\tipitakaparskip}
\csromancenter{สฬายตนสํยุตฺต ตุมฺหากํ ตํ ปชหถ, ตํ โว ปหีนํ หิตาย สุขาย ภวิสฺสติ.\csromanspacer{} รูปา น ตุมฺหากํ เต ปชหถ, เต โว ปหีนา หิตาย สุขาย ภวิสฺสนฺติ.\csromanspacer{} จกฺขุวิญฺญาณํ น ตุมฺหากํ ตํ ปชหถ, ตํ โว ปหีนํ หิตาย สุขาย ภวิสฺสติ.\csromanspacer{} จกฺขุสมฺผสฺโส น ตุมฺหากํ ตํ ปชหถ, โส โว ปหีโน หิตาย สุขาย ภวิสฺสติ.\csromanspacer{} ยมฺปิทํ จกฺขุสมฺผสฺสปจฺจยา อุปฺปชฺชติ เวทยิตํ สุขํ วา ทุกฺขํ วา อทุกฺขมสุขํ วา, ตมฺปิ น ตุมฺหากํ ตํ ปชหถ, ตํ โว ปหีนํ หิตาย สุขาย ภวิสฺสติ ฯเปฯ.}
\prose{\csromanfolio{303}ชิวฺหา น ตุมฺหากํ ตํ ปชหถ, สา โว ปหีนา หิตาย สุขาย ภวิสฺสติ.\csromanspacer{} รสา น ตุมฺหากํ เต ปชหถ, เต โว ปหีนา หิตาย สุขาย ภวิสฺสนฺติ.\csromanspacer{} ชิวฺหาวิญฺญาณํ น ตุมฺหากํ ตํ ปชหถ, ตํ โว ปหีนํ หิตาย สุขาย ภวิสฺสติ.\csromanspacer{} ชิวฺหาสมฺผสฺโส น ตุมฺหากํ ตํ ปชหถ, โส โว ปหีโน หิตาย สุขาย ภวิสฺสติ.\csromanspacer{} ยมฺปิทํ ชิวฺหาสมฺผสฺสปจฺจยา อุปฺปชฺชติ เวทยิตํ สุขํ วา ทุกฺขํ วา อทุกฺขมสุขํ วา, ตมฺปิ น ตุมฺหากํ ตํ ปชหถ, ตํ โว ปหีนํ หิตาย สุขาย ภวิสฺสติ ฯเปฯ.}

\prose{มโน น ตุมฺหากํ ตํ ปชหถ, โส โว ปหีโน หิตาย สุขาย ภวิสฺสติ.\csromanspacer{} ธมฺมา น ตุมฺหากํ เต ปชหถ, เต โว ปหีนา หิตาย สุขาย ภวิสฺสติ.\csromanspacer{} มโนวิญฺญาณํ น ตุมฺหากํ ตํ ปชหถ, ตํ โว ปหีนํ หิตาย สุขาย ภวิสฺสติ.\csromanspacer{} มโนสมฺผสฺโส น ตุมฺหากํ ตํ ปชหถ, โส โว ปหีโน หิตาย สุขาย ภวิสฺสติ.\csromanspacer{} ยมฺปิทํ มโนสมฺผสฺสปจฺจยา อุปฺปชฺชติ เวทยิตํ สุขํ วา ทุกฺขํ วา อทุกฺขมสุขํ วา, ตมฺปิ น ตุมฺหากํ ตํ ปชหถ, ตํ โว ปหีนํ หิตาย สุขาย ภวิสฺสติ.}

\prose{เสยฺยถาปิ ภิกฺขเว ยํ อิมสฺมิํ เชตวเน ติณกฏฺฐสาขาปลาสํ, ตํ ชโน หเรยฺย วา ฑเหยฺย วา ยถาปจฺจยํ วา กเรยฺย, อปิ นุ ตุมฺหากํ เอวมสฺส “อเมฺห ชโน หรติ วา ฑหติ วา ยถาปจฺจยํ วา กโรตี”ติ.\csromanspacer{} โน เหตํ ภนฺเต.\csromanspacer{} ตํ กิสฺส เหตุ, น หิ โน เอตํ ภนฺเต อตฺตา วา อตฺตนิยํ วาติ.\csromanspacer{} เอวเมว โข ภิกฺขเว จกฺขุ น ตุมฺหากํ ตํ ปชหถ, ตํ โว ปหีนํ หิตาย สุขาย ภวิสฺสติ.\csromanspacer{} รูปา น ตุมฺหากํ.\csromanspacer{} จกฺขุวิญฺญาณํ.\csromanspacer{} จกฺขุสมฺผสฺโส ฯเปฯ.\csromanspacer{} ยมฺปิทํ มโนสมฺผสฺสปจฺจยา อุปฺปชฺชติ เวทยิตํ สุขํ วา ทุกฺขํ วา อทุกฺขมสุขํ วา, ตมฺปิ น ตุมฺหากํ ตํ ปชหถ, ตํ โว ปหีนํ หิตาย สุขาย ภวิสฺสตีติ.\csromanspacer{}\csromanspacer{}อฏฺฐมํ.}

\gambhira{สฬายตนวคฺคสํยุตฺตปาฬิ}
\csromanheader{2}{9. ทุติยนตุมฺหากสุตฺต}
\proseitem{102}{\csromanfolio{304}ยํ ภิกฺขเว น ตุมฺหากํ ตํ ปชหถ, ตํ โว ปหีนํ หิตาย สุขาย ภวิสฺสติ.\csromanspacer{} กิญฺจ ภิกฺขเว น ตุมฺหากํ, จกฺขุ ภิกฺขเว น ตุมฺหากํ ตํ ปชหถ, ตํ โว ปหีนํ หิตาย สุขาย ภวิสฺสติ.\csromanspacer{} รูปา น ตุมฺหากํ เต ปชหถ, เต โว ปหีนา หิตาย สุขาย ภวิสฺสนฺติ.\csromanspacer{} จกฺขุวิญฺญาณํ น ตุมฺหากํ ตํ ปชหถ, ตํ โว ปหีนํ หิตาย สุขาย ภวิสฺสติ.\csromanspacer{} จกฺขุสมฺผสฺโส น ตุมฺหากํ ตํ ปชหถ, โส โว ปหีโน หิตาย สุขาย ภวิสฺสติ ฯเปฯ.\csromanspacer{} ยมฺปิทํ มโนสมฺผสฺสปจฺจยา อุปฺปชฺชติ เวทยิตํ สุขํ วา ทุกฺขํ วา อทุกฺขมสุขํ วา, ตมฺปิ น ตุมฺหากํ ตํ ปชหถ, ตํ โว ปหีนํ หิตาย สุขาย ภวิสฺสติ.\csromanspacer{} ยมฺปิ ภิกฺขเว น ตุมฺหากํ ตํ ปชหถ, ตํ โว ปหีนํ หิตาย สุขาย ภวิสฺสตีติ.\csromanspacer{}\csromanspacer{}นวมํ.}

\csromanmatikaanchor{mtk.140}
\csromantocmarkref{subsection}{10. อุทกสุตฺต}{mtk.140}
\bookmark[dest={mtk.140},level=2]{10. อุทกสุตฺต}
\csromanheader{2}{10. อุทกสุตฺต}
\proseitem{103}{อุทโก\footnote{อุทฺทโก (สี, อิ)} สุทํ ภิกฺขเว รามปุตฺโต เอวํ วาจํ ภาสติ “อิทํ ชาตุ เวทคู, อิทํ ชาตุ สพฺพชี\footnote{สพฺพชิ (อิ)}, อิทํ ชาตุ อปลิขตํ คณฺฑมูลํ ปลิขณินฺ”ติ.\csromanspacer{} ตํ โข ปเนตํ ภิกฺขเว อุทโก รามปุตฺโต อเวทคูเยว สมาโน “เวทคูสฺมี”ติ ภาสติ, อสพฺพชีเยว สมาโน “สพฺพชีสฺมี”ติ ภาสติ, อปลิขตํเยว คณฺฑมูลํ “ปลิขตํ เม คณฺฑมูลนฺ”ติ ภาสติ.\csromanspacer{} อิธ โข ตํ ภิกฺขเว ภิกฺขุ สมฺมา วทมาโน วเทยฺย “อิทํ ชาตุ เวทคู, อิทํ ชาตุ สพฺพชี, อิทํ ชาตุ อปลิขตํ คณฺฑมูลํ ปลิขณินฺ”ติ.}

\prose{กถญฺจ ภิกฺขเว เวทคู โหติ, ยโต โข ภิกฺขเว ภิกฺขุ ฉนฺนํ ผสฺสายตนานํ สมุทยญฺจ อตฺถงฺคมญฺจ อสฺสาทญฺจ อาทีนวญฺจ นิสฺสรณญฺจ ยถาภูตํ ปชานาติ.\csromanspacer{} เอวํ โข ภิกฺขเว ภิกฺขุ เวทคู โหติ.}

\prose{กถญฺจ ภิกฺขเว ภิกฺขุ สพฺพชี โหติ, ยโต โข ภิกฺขเว ภิกฺขุ ฉนฺนํ ผสฺสายตนานํ สมุทยญฺจ อตฺถงฺคมญฺจ อสฺสาทญฺจ อาทีนวญฺจ นิสฺสรณญฺจ ยถาภูตํ วิทิตฺวา อนุปาทาวิมุตฺโต โหติ.\csromanspacer{} เอวํ โข ภิกฺขเว ภิกฺขุ สพฺพชี โหติ.}

\prose{กถญฺจ ภิกฺขเว ภิกฺขุโน อปลิขตํ คณฺฑมูลํ ปลิขตํ โหติ. “คณฺโฑ”ติ โข ภิกฺขเว อิมสฺเสตํ จาตุมหาภูติกสฺส กายสฺส}

\vspace{\tipitakaparskip}
\csromancenter{สฬายตนสํยุตฺต อธิวจนํ มาตาเปตฺติกสมฺภวสฺส โอทนกุมฺมาสูปจยสฺส อนิจฺจุจฺฉาทนปริมทฺทนเภทนวิทฺธํสนธมฺมสฺส. “คณฺฑมูลนฺ”ติ โข ภิกฺขเว ตณฺหาเยตํ อธิวจนํ.\csromanspacer{} ยโต โข ภิกฺขเว ภิกฺขุโน ตณฺหา ปหีนา โหติ, อุจฺฉินฺนมูลา ตาลาวตฺถุกตา อนภาวํกตา อายติํ อนุปฺปาทธมฺมา.\csromanspacer{} เอวํ โข ภิกฺขเว ภิกฺขุโน อปลิขตํ คณฺฑมูลํ ปลิขตํ โหติ.}
\prose{\csromanfolio{305}อุทโก สุทํ ภิกฺขเว รามปุตฺโต เอวํ วาจํ ภาสติ “อิทํ ชาตุ เวทคู, อิทํ ชาตุ สพฺพชี, อิทํ ชาตุ อปลิขตํ คณฺฑมูลํ ปลิขณินฺ”ติ.\csromanspacer{} ตํ โข ปเนตํ ภิกฺขเว อุทโก รามปุตฺโต อเวทคูเยว สมาโน “เวทคูสฺมี”ติ ภาสติ, อสพฺพชีเยว สมาโน “สพฺพชีสฺมี”ติ ภาสติ, อปลิขตํเยว คณฺฑมูลํ “ปลิขตํ เม คณฺฑมูลนฺ”ติ ภาสติ.\csromanspacer{} อิธ โข ตํ ภิกฺขเว ภิกฺขุ สมฺมา วทมาโน วเทยฺย “อิทํ ชาตุ เวทคู, อิทํ ชาตุ สพฺพชี, อิทํ ชาตุ อปลิขตํ คณฺฑมูลํ ปลิขณินฺ”ติ.\csromanspacer{}\csromanspacer{}ทสมํ.\csromanspacer{} สฬวคฺโค ปญฺจโม.}

\titlehead{\textbf{ตสฺสุทฺทานํ}}
\csromangathameasure{เทฺว สํคยฺหา ปริหานํ, ปมาทวิหารี จ สํวโร. \\ สมาธิ ปฏิสลฺลานํ, เทฺว นตุมฺหาเกน อุทฺทโกติ. สฬายตนวคฺเค ทุติยปณฺณสโก สมตฺโต.}
\csromangathagroup{\csromangathastanza{เทฺว สํคยฺหา ปริหานํ, ปมาทวิหารี จ สํวโร. \\ สมาธิ ปฏิสลฺลานํ, เทฺว นตุมฺหาเกน อุทฺทโกติ.\csromanspacer{} สฬายตนวคฺเค ทุติยปณฺณสโก สมตฺโต.}}

\csromanheader{2}{\textbf{ตสฺส วคฺคุทฺทานํ}}
\csromangathameasure{อวิชฺชา มิคชาลญฺจ, คิลานํ ฉนฺนํ จตุตฺถกํ. \\ สฬวคฺเคน ปญฺญาสํ, ทุติโย ปณฺณาสโก อยนฺติ. ปฐมสตกํ.}
\csromangathagroup{\csromangathastanza{อวิชฺชา มิคชาลญฺจ, คิลานํ ฉนฺนํ จตุตฺถกํ. \\ สฬวคฺเคน ปญฺญาสํ, ทุติโย ปณฺณาสโก อยนฺติ.\csromanspacer{} ปฐมสตกํ.}}

\csromanmatikaanchor{mtk.141}
\csromantocmarknopage{section}{(11) 1. โยคกฺเขมิวคฺค}{mtk.141}
\bookmark[dest={mtk.141},level=1]{(11) 1. โยคกฺเขมิวคฺค}
\csromanheader{1}{\textbf{(11) 1. โยคกฺเขมิวคฺค}}
\csromanmatikaanchor{mtk.142}
\csromantocmarkref{subsection}{1. โยคกฺเขมิสุตฺต}{mtk.142}
\bookmark[dest={mtk.142},level=2]{1. โยคกฺเขมิสุตฺต}
\csromanheader{2}{1. โยคกฺเขมิสุตฺต}
\proseitem{104}{สาวตฺถินิทานํ.\csromanspacer{} โยคกฺเขมิปริยายํ โว ภิกฺขเว ธมฺมปริยายํ เทเสสฺสามิ, ตํ สุณาถ.\csromanspacer{} กตโม จ ภิกฺขเว โยคกฺเขมิปริยาโย ธมฺมปริยาโย.\csromanspacer{} สนฺติ ภิกฺขเว จกฺขุวิญฺเญยฺยา รูปา อิฏฺฐา กนฺตา}

\gambhira{สฬายตนวคฺคสํยุตฺตปาฬิ}
\prose{\csromanfolio{306}มนาปา ปิยรูปา กามูปสํหิตา รชนียา, เต ตถาคตสฺส ปหีนา อุจฺฉินฺนมูลา ตาลาวตฺถุกตา อนภาวํกตา อายติํ อนุปฺปาทธมฺมา, เตสญฺจ ปหานาย อกฺขาสิ โยคํ, ตสฺมา “ตถาคโต โยคกฺเขมี”ติ วุจฺจติ ฯเปฯ.\csromanspacer{} สนฺติ ภิกฺขเว มโนวิญฺเญยฺยา ธมฺมา อิฏฺฐา กนฺตา มนาปา ปิยรูปา กามูปสํหิตา รชนียา, เต ตถาคตสฺส ปหีนา อุจฺฉินฺนมูลา ตาลาวตฺถุกตา อนภาวํกตา อายติํ อนุปฺปาทธมฺมา, เตสญฺจ ปหานาย อกฺขาสิ โยคํ, ตสฺมา “ตถาคโต โยคกฺเขมี”ติ วุจฺจติ.\csromanspacer{} อยํ โข ภิกฺขเว โยคกฺเขมิปริยาโย ธมฺมปริยาโยติ.\csromanspacer{}\csromanspacer{}ปฐมํ.}

\csromanmatikaanchor{mtk.143}
\csromantocmarkref{subsection}{2. อุปาทายสุตฺต}{mtk.143}
\bookmark[dest={mtk.143},level=2]{2. อุปาทายสุตฺต}
\csromanheader{2}{2. อุปาทายสุตฺต}
\proseitem{105}{กิสฺมิํ นุ โข ภิกฺขเว สติ กิํ อุปาทาย อุปฺปชฺชติ อชฺฌตฺตํ สุขํ ทุกฺขนฺติ.\csromanspacer{} ภควํมูลกา โน ภนฺเต ธมฺมา.\csromanspacer{} จกฺขุสฺมิํ โข ภิกฺขเว สติ จกฺขุํ อุปาทาย อุปฺปชฺชติ อชฺฌตฺตํ สุขํ ทุกฺขํ ฯเปฯ.\csromanspacer{} มนสฺมิํ สติ มนํ อุปาทาย อุปฺปชฺชติ อชฺฌตฺตํ สุขํ ทุกฺขํ.\csromanspacer{} ตํ กิํ มญฺญถ ภิกฺขเว, จกฺขุ นิจฺจํ วา อนิจฺจํ วาติ.\csromanspacer{} อนิจฺจํ ภนฺเต.\csromanspacer{} ยํ ปนานิจฺจํ, ทุกฺขํ วา ตํ สุขํ วาติ.\csromanspacer{} ทุกฺขํ ภนฺเต.\csromanspacer{} ยํ ปนานิจฺจํ ทุกฺขํ วิปริณามธมฺมํ, อปิ นุ ตํ อนุปาทาย อุปฺปชฺเชยฺย อชฺฌตฺตํ สุขํ ทุกฺขนฺติ.\csromanspacer{} โน เหตํ ภนฺเต ฯเปฯ.\csromanspacer{} ชิวฺหา นิจฺจา วา อนิจฺจา วาติ.\csromanspacer{} อนิจฺจา ภนฺเต.\csromanspacer{} ยํ ปนานิจฺจํ, ทุกฺขํ วา ตํ สุขํ วาติ.\csromanspacer{} ทุกฺขํ ภนฺเต.\csromanspacer{} ยํ ปนานิจฺจํ ทุกฺขํ วิปริณามธมฺมํ, อปิ นุ ตํ อนุปาทาย อุปฺปชฺเชยฺย อชฺฌตฺตํ สุขํ ทุกฺขนฺติ.\csromanspacer{} โน เหตํ ภนฺเต ฯเปฯ.\csromanspacer{} มโน นิจฺโจ วา อนิจฺโจ วาติ.\csromanspacer{} อนิจฺโจ ภนฺเต.\csromanspacer{} ยํ ปนานิจฺจํ, ทุกฺขํ วา ตํ สุขํ วาติ.\csromanspacer{} ทุกฺขํ ภนฺเต.\csromanspacer{} ยํ ปนานิจฺจํ ทุกฺขํ วิปริณามธมฺมํ, อปิ นุ ตํ อนุปาทาย อุปฺปชฺเชยฺย อชฺฌตฺตํ สุขํ ทุกฺขนฺติ.\csromanspacer{} โน เหตํ ภนฺเต.\csromanspacer{} เอวํ ปสฺสํ ภิกฺขเว สุตวา อริยสาวโก จกฺขุสฺมิมฺปิ นิพฺพินฺทติ ฯเปฯ มนสฺมิมฺปิ นิพฺพินฺทติ, นิพฺพินฺทํ วิรชฺชติ, วิราคา วิมุจฺจติ, วิมุตฺตสฺมิํ “วิมุตฺตมฺ”อิติ ญาณํ โหติ, “ขีณา ชาติ, วุสิตํ พฺรหฺมจริยํ, กตํ กรณียํ, นาปรํ อิตฺถตฺตายา”ติ ปชานาตีติ.\csromanspacer{}\csromanspacer{}ทุติยํ.}

\csromanmatikaanchor{mtk.144}
\csromantocmarkref{subsection}{3. ทุกฺขสมุทยสุตฺต}{mtk.144}
\bookmark[dest={mtk.144},level=2]{3. ทุกฺขสมุทยสุตฺต}
\csromanheader{2}{3. ทุกฺขสมุทยสุตฺต}
\proseitem{106}{ทุกฺขสฺส ภิกฺขเว สมุทยญฺจ อตฺถงฺคมญฺจ เทเสสฺสามิ, ตํ สุณาถ.\csromanspacer{} กตโม จ ภิกฺขเว ทุกฺขสฺส สมุทโย, จกฺขุญฺจ ปฏิจฺจ รูเป จ อุปฺปชฺชติ}

\vspace{\tipitakaparskip}
\csromancenter{สฬายตนสํยุตฺต จกฺขุวิญฺญาณํ, ติณฺณํ สงฺคติ ผสฺโส, ผสฺสปจฺจยา เวทนา, เวทนาปจฺจยา ตณฺหา.\csromanspacer{} อยํ ทุกฺขสฺส สมุทโย ฯเปฯ.\csromanspacer{} ชิวฺหญฺจ ปฏิจฺจ รเส จ อุปฺปชฺชติ ชิวฺหาวิญฺญาณํ, ติณฺณํ สงฺคติ ผสฺโส, ผสฺสปจฺจยา เวทนา, เวทนาปจฺจยา ตณฺหา.\csromanspacer{} อยํ ทุกฺขสฺส สมุทโย ฯเปฯ.\csromanspacer{} มนญฺจ ปฏิจฺจ ธมฺเม จ อุปฺปชฺชติ มโนวิญฺญาณํ, ติณฺณํ สงฺคติ ผสฺโส, ผสฺสปจฺจยา เวทนา, เวทนาปจฺจยา ตณฺหา.\csromanspacer{} อยํ โข ภิกฺขเว ทุกฺขสฺส สมุทโย.}
\prose{\csromanfolio{307}กตโม จ ภิกฺขเว ทุกฺขสฺส อตฺถงฺคโม, จกฺขุญฺจ ปฏิจฺจ รูเป จ อุปฺปชฺชติ จกฺขุวิญฺญาณํ, ติณฺณํ สงฺคติ ผสฺโส, ผสฺสปจฺจยา เวทนา, เวทนาปจฺจยา ตณฺหา.\csromanspacer{} ตสฺสาเยว ตณฺหาย อเสสวิราคนิโรธา อุปาทานนิโรโธ, อุปาทานนิโรธา ภวนิโรโธ, ภวนิโรธา ชาตินิโรโธ, ชาตินิโรธา ชรามรณํ โสกปริเทวทุกฺขโทมนสฺสุปายาสา นิรุชฺฌนฺติ, เอวเมตสฺส เกวลสฺส ทุกฺขกฺขนฺธสฺส นิโรโธ โหติ.\csromanspacer{} อยํ ทุกฺขสฺส อตฺถงฺคโม ฯเปฯ.\csromanspacer{} ชิวฺหญฺจ ปฏิจฺจ รเส จ อุปฺปชฺชติ ชิวฺหาวิญฺญาณํ ฯเปฯ.\csromanspacer{} มนญฺจ ปฏิจฺจ ธมฺเม จ อุปฺปชฺชติ มโนวิญฺญาณํ, ติณฺณํ สงฺคติ ผสฺโส, ผสฺสปจฺจยา เวทนา, เวทนาปจฺจยา ตณฺหา.\csromanspacer{} ตสฺสาเยว ตณฺหาย อเสสวิราคนิโรธา อุปาทานนิโรโธ, อุปาทานนิโรธา ภวนิโรโธ, ภวนิโรธา ชาตินิโรโธ, ชาตินิโรธา ชรามรณํ โสกปริเทวทุกฺขโทมนสฺสุปายาสา นิรุชฺฌนฺติ, เอวเมตสฺส เกวลสฺส ทุกฺขกฺขนฺธสฺส นิโรโธ โหติ.\csromanspacer{} อยํ โข ภิกฺขเว ทุกฺขสฺส อตฺถงฺคโมติ.\csromanspacer{}\csromanspacer{}ตติยํ.}

\csromanmatikaanchor{mtk.145}
\csromantocmarkref{subsection}{4. โลกสมุทยสุตฺต}{mtk.145}
\bookmark[dest={mtk.145},level=2]{4. โลกสมุทยสุตฺต}
\csromanheader{2}{4. โลกสมุทยสุตฺต}
\proseitem{107}{โลกสฺส ภิกฺขเว สมุทยญฺจ อตฺถงฺคมญฺจ เทเสสฺสามิ, ตํ สุณาถ.\csromanspacer{} กตโม จ ภิกฺขเว โลกสฺส สมุทโย, จกฺขุญฺจ ปฏิจฺจ รูเป จ อุปฺปชฺชติ จกฺขุวิญฺญาณํ, ติณฺณํ สงฺคติ ผสฺโส, ผสฺสปจฺจยา เวทนา, เวทนาปจฺจยา ตณฺหา, ตณฺหาปจฺจยา อุปาทานํ, อุปาทานปจฺจยา ภโว, ภวปจฺจยา ชาติ, ชาติปจฺจยา ชรามรณํ โสกปริเทวทุกฺขมโทมนสฺสุปายาสา สมฺภวนฺติ.\csromanspacer{} อยํ โข ภิกฺขเว โลกสฺส สมุทโย ฯเปฯ.\csromanspacer{} ชิวฺหญฺจ ปฏิจฺจ รเส จ อุปฺปชฺชติ ชิวฺหาวิญฺญาณํ ฯเปฯ.\csromanspacer{} มนญฺจ ปฏิจฺจ ธมฺเม จ อุปฺปชฺชติ มโนวิญฺญาณํ, ติณฺณํ สงฺคติ ผสฺโส, ผสฺสปจฺจยา เวทนา, เวทนาปจฺจยา ตณฺหา, ตณฺหาปจฺจยา อุปาทานํ, อุปาทานปจฺจยา}

\gambhira{สฬายตนวคฺคสํยุตฺตปาฬิ}
\prose{\csromanfolio{308}ภโว, ภวปจฺจยา ชาติ, ชาติปจฺจยา ชรามรณํ โสก ปริเทวทุกฺขโทมนสฺสุปายาสา สมฺภวนฺติ.\csromanspacer{} อยํ โข ภิกฺขเว โลกสฺส สมุทโย.}

\prose{กตโม จ ภิกฺขเว โลกสฺส อตฺถงฺคโม, จกฺขุญฺจ ปฏิจฺจ รูเป จ อุปฺปชฺชติ จกฺขุวิญฺญาณํ, ติณฺณํ สงฺคติ ผสฺโส, ผสฺสปจฺจยา เวทนา, เวทนาปจฺจยา ตณฺหา.\csromanspacer{} ตสฺสาเยว ตณฺหาย อเสสวิราคนิโรธา อุปาทานนิโรโธ, อุปาทานนิโรธา ภวนิโรโธ, ภวนิโรธา ชาตินิโรโธ, ชาตินิโรธา ชรามรณํ โสก ปริเทว ทุกฺขโทมนสฺสุปายาสา นิรุชฺฌนฺติ, เอวเมตสฺส เกวลสฺส ทุกฺขกฺขนฺธสฺส นิโรโธ โหติ.\csromanspacer{} อยํ โข ภิกฺขเว โลกสฺส อตฺถงฺคโม ฯเปฯ.\csromanspacer{} ชิวฺหญฺจ ปฏิจฺจ รเส จ อุปฺปชฺชติ ฯเปฯ.\csromanspacer{} มนญฺจ ปฏิจฺจ ธมฺเม จ อุปฺปชฺชติ มโนวิญฺญาณํ, ติณฺณํ สงฺคติ ผสฺโส, ผสฺสปจฺจยา เวทนา, เวทนาปจฺจยา ตณฺหา.\csromanspacer{} ตสฺสาเยว ตณฺหาย อเสสวิราคนิโรธา อุปาทานนิโรโธ, อุปาทานนิโรธา ฯเปฯ เอวเมตสฺส เกวลสฺส ทุกฺขกฺขนฺธสฺส นิโรโธ โหติ.\csromanspacer{} อยํ โข ภิกฺขเว โลกสฺส อตฺถงฺคโมติ.\csromanspacer{}\csromanspacer{}จตุตฺถํ.}

\csromanmatikaanchor{mtk.146}
\csromantocmarkref{subsection}{5. เสยฺโยหมสฺมิสุตฺต}{mtk.146}
\bookmark[dest={mtk.146},level=2]{5. เสยฺโยหมสฺมิสุตฺต}
\csromanheader{2}{5. เสยฺโยหมสฺมิสุตฺต}
\proseitem{108}{กิสฺมิํ นุ โข ภิกฺขเว สติ กิํ อุปาทาย กิํ อภินิวิสฺส “เสยฺโยหมสฺมี”ติ วา โหติ, “สทิโสหมสฺมี”ติ วา โหติ, “หีโนหมสฺมี”ติ วา โหตีติ.\csromanspacer{} ภควํมูลกา โน ภนฺเต ธมฺมา.\csromanspacer{} จกฺขุสฺมิํ โข ภิกฺขเว สติ จกฺขุํ อุปาทาย จกฺขุํ อภินิวิสฺส “เสยฺโยหมสฺมี”ติ วา โหติ, “สทิโสหมสฺมี”ติ วา โหติ, “หีโนหมสฺมี”ติ วา โหติ ฯเปฯ.\csromanspacer{} ชิวฺหาย สติ ฯเปฯ.\csromanspacer{} มนสฺมิํ สติ มนํ อุปาทาย มนํ อภินิวิสฺส “เสยฺโยหมสฺมี”ติ วา โหติ, “สทิโสหมสฺมี”ติ วา โหติ, “หีโนหมสฺมี”ติ วา โหติ.\csromanspacer{} ตํ กิํ มญฺญถ ภิกฺขเว, จกฺขุ นิจฺจํ วา อนิจฺจํ วาติ.\csromanspacer{} อนิจฺจํ ภนฺเต.\csromanspacer{} ยํ ปนานิจฺจํ, ทุกฺขํ วา ตํ สุขํ วาติ.\csromanspacer{} ทุกฺขํ ภนฺเต.\csromanspacer{} ยํ ปนานิจฺจํ ทุกฺขํ วิปริณามธมฺมํ, อปิ นุ ตํ อนุปาทาย “เสยฺโยหมสฺมี”ติ วา อสฺส, “สทิโสหมสฺมี”ติ วา อสฺส, “หีโนหมสฺมี”ติ วา อสฺสาติ.\csromanspacer{} โน เหตํ ภนฺเต ฯเปฯ.\csromanspacer{} ชิวฺหา นิจฺจา วา อนิจฺจา วาติ.\csromanspacer{} อนิจฺจา ภนฺเต ฯเปฯ.\csromanspacer{} มโน นิจฺโจ วา อนิจฺโจ วาติ.\csromanspacer{} อนิจฺโจ ภนฺเต.\csromanspacer{} ยํ ปนานิจฺจํ, ทุกฺขํ}

\vspace{\tipitakaparskip}
\csromancenter{สฬายตนสํยุตฺต วา ตํ สุขํ วาติ.\csromanspacer{} ทุกฺขํ ภนฺเต.\csromanspacer{} ยํ ปนานิจฺจํ ทุกฺขํ วิปริณามธมฺมํ, อปิ นุ ตํ อนุปาทาย “เสยฺโยหมสฺมี”ติ วา อสฺส, “สทิโสหมสฺมี”ติ วา อสฺส, “หีโนหมสฺมี”ติ วา อสฺสาติ.\csromanspacer{} โน เหตํ ภนฺเต.}
\prose{\csromanfolio{309}เอวํ ปสฺสํ ภิกฺขเว สุตวา อริยสาวโก จกฺขุสฺมิมฺปิ นิพฺพินฺทติ ฯเปฯ มนสฺมิมฺปิ นิพฺพินฺทติ, นิพฺพินฺทํ วิรชฺชติ, วิราคา วิมุจฺจติ, วิมุตฺตสฺมิํ “วิมุตฺตมฺ”อิติ ญาณํ โหติ, “ขีณา ชาติ, วุสิตํ พฺรหฺมจริยํ, กตํ กรณียํ, นาปรํ อิตฺถตฺตายา”ติ ปชานาตีติ.\csromanspacer{}\csromanspacer{}ปญฺจมํ.}

\csromanmatikaanchor{mtk.147}
\csromantocmarkref{subsection}{6. สํโยชนิยสุตฺต}{mtk.147}
\bookmark[dest={mtk.147},level=2]{6. สํโยชนิยสุตฺต}
\csromanheader{2}{6. สํโยชนิยสุตฺต}
\proseitem{109}{สํโยชนิเย จ ภิกฺขเว ธมฺเม เทเสสฺสามิ สํโยชนญฺจ, ตํ สุณาถ.\csromanspacer{} กตเม จ ภิกฺขเว สํโยชนิยา ธมฺมา, กตมญฺจ สํโยชนํ.\csromanspacer{} จกฺขุํ ภิกฺขเว สํโยชนิโย ธมฺโม.\csromanspacer{} โย ตตฺถ ฉนฺทราโค ตํ ตตฺถ สํโยชนํ ฯเปฯ.\csromanspacer{} ชิวฺหา สํโยชนิโย ธมฺโม ฯเปฯ.\csromanspacer{} มโน สํโยชนิโย ธมฺโม.\csromanspacer{} โย ตตฺถ ฉนฺทราโค, ตํ ตตฺถ สํโยชนํ.\csromanspacer{} อิเม วุจฺจนฺติ ภิกฺขเว สํโยชนิยา ธมฺมา, อิทํ สํโยชนนฺติ.\csromanspacer{}\csromanspacer{}ฉฏฺฐํ.}

\csromanmatikaanchor{mtk.148}
\csromantocmarkref{subsection}{7. อุปาทานิยสุตฺต}{mtk.148}
\bookmark[dest={mtk.148},level=2]{7. อุปาทานิยสุตฺต}
\csromanheader{2}{7. อุปาทานิยสุตฺต}
\proseitem{110}{อุปาทานิเย จ ภิกฺขเว ธมฺเม เทเสสฺสามิ อุปาทานญฺจ, ตํ สุณาถ.\csromanspacer{} กตเม จ ภิกฺขเว อุปาทานิยา ธมฺมา, กตมญฺจ อุปาทานํ.\csromanspacer{} จกฺขุํ ภิกฺขเว อุปาทานิโย ธมฺโม.\csromanspacer{} โย ตตฺถ ฉนฺทราโค, ตํ ตตฺถ อุปาทานํ ฯเปฯ.\csromanspacer{} ชิวฺหา อุปาทานิโย ธมฺโม ฯเปฯ.\csromanspacer{} มโน อุปาทานิโย ธมฺโม.\csromanspacer{} โย ตตฺถ ฉนฺทราโค, ตํ ตตฺถ อุปาทานํ.\csromanspacer{} อิเม วุจฺจนฺติ ภิกฺขเว อุปาทานิยา ธมฺมา, อิทํ อุปาทานนฺติ.\csromanspacer{}\csromanspacer{}สตฺตมํ.}

\csromanmatikaanchor{mtk.149}
\csromantocmarkref{subsection}{8. อชฺฌตฺติกายตนปริชานนสุตฺต}{mtk.149}
\bookmark[dest={mtk.149},level=2]{8. อชฺฌตฺติกายตนปริชานนสุตฺต}
\csromanheader{2}{8. อชฺฌตฺติกายตนปริชานนสุตฺต}
\proseitem{111}{จกฺขุํ ภิกฺขเว อนภิชานํ อปริชานํ อวิราชยํ อปฺปชหํ อภพฺโพ ทุกฺขกฺขยาย.\csromanspacer{} โสตํ.\csromanspacer{} ฆานํ.\csromanspacer{} ชิวฺหํ.\csromanspacer{} กายํ.\csromanspacer{} มนํ อนภิชานํ อปริชานํ อวิราชยํ อปฺปชหํ อภพฺโพ ทุกฺขกฺขยาย.\csromanspacer{} จกฺขุญฺจ โข ภิกฺขเว อภิชานํ ปริชานํ วิราชยํ ปชหํ ภพฺโพ ทุกฺขกฺขยาย ฯเปฯ.\csromanspacer{} ชิวฺหํ.\csromanspacer{} กายํ.\csromanspacer{} มนํ อภิชานํ ปริชานํ วิราชยํ ปชหํ ภพฺโพ ทุกฺขกฺขยายาติ.\csromanspacer{}\csromanspacer{}อฏฺฐมํ.}

\gambhira{สฬายตนวคฺคสํยุตฺตปาฬิ}
\csromanmatikaanchor{mtk.150}
\csromantocmarkref{subsection}{9. พาหิรายตนปริชานนสุตฺต}{mtk.150}
\bookmark[dest={mtk.150},level=2]{9. พาหิรายตนปริชานนสุตฺต}
\csromanheader{2}{9. พาหิรายตนปริชานนสุตฺต}
\proseitem{112}{\csromanfolio{310}รูเป ภิกฺขเว อนภิชานํ อปริชานํ อวิราชยํ อปฺปชหํ อภพฺโพ ทุกฺขกฺขยาย.\csromanspacer{} สทฺเท.\csromanspacer{} คนฺเธ.\csromanspacer{} รเส.\csromanspacer{} โผฏฺฐพฺเพ.\csromanspacer{} ธมฺเม อนภิชานํ อปริชานํ อวิราชยํ อปฺปชหํ อภพฺโพ ทุกฺขกฺขยาย.\csromanspacer{} รูเป จ โข ภิกฺขเว อภิชานํ ปริชานํ วิราชยํ ปชหํ ภพฺโพ ทุกฺขกฺขยาย.\csromanspacer{} สทฺเท.\csromanspacer{} คนฺเธ.\csromanspacer{} รเส.\csromanspacer{} โผฏฺฐพฺเพ.\csromanspacer{} ธมฺเม อภิชานํ ปริชานํ วิราชยํ ปชหํ ภพฺโพ ทุกฺขกฺขยายาติ.\csromanspacer{}\csromanspacer{}นวมํ.}

\csromanmatikaanchor{mtk.151}
\csromantocmarkref{subsection}{10. อุปสฺสุติสุตฺต}{mtk.151}
\bookmark[dest={mtk.151},level=2]{10. อุปสฺสุติสุตฺต}
\csromanheader{2}{10. อุปสฺสุติสุตฺต}
\proseitem{113}{เอกํ สมยํ ภควา นาติเก\footnote{ญาติเก (สี, สฺยา, กํ)} วิหรติ คิญฺชกาวสเถ.\csromanspacer{} อถ โข ภควา รโหคโต ปฏิสลฺลีโน อิมํ ธมฺมปริยายํ อภาสิ “จกฺขุํ จ ปฏิจฺจ รูเป จ อุปฺปชฺชติ จกฺขุวิญฺญาณํ, ติณฺณํ สงฺคติ ผสฺโส, ผสฺสปจฺจยา เวทนา, เวทนาปจฺจยา ตณฺหา, ตณฺหาปจฺจยา อุปาทานํ, อุปาทานปจฺจยา ภโว, ภวปจฺจยา ชาติ, ชาติปจฺจยา ชรามรณํ โสกปริเทวทุกฺขโทมนสฺสุปายาสา สมฺภวนฺติ.\csromanspacer{} เอวเมตสฺส เกวลสฺส ทุกฺขกฺขนฺธสฺส สมุทโย โหติ.\csromanspacer{} ชิวฺหญฺจ ปฏิจฺจ รเส จ อุปฺปชฺชติ ฯเปฯ.\csromanspacer{} มนญฺจ ปฏิจฺจ ธมฺเม จ อุปฺปชฺชติ มโนวิญฺญาณํ, ติณฺณํ สงฺคติ ผสฺโส, ผสฺสปจฺจยา เวทนา, เวทนาปจฺจยา ตณฺหา, ตณฺหาปจฺจยา อุปาทานํ, อุปาทานปจฺจยา ภโว, ภวปจฺจยา ชาติ, ชาติปจฺจยา ชรามรณํ โสกปริเทวทุกฺขมนสฺสุปายาสา สมฺภวนฺติ.\csromanspacer{} เอวเมตสฺส เกวลสฺส ทุกฺขกฺขนฺธสฺส สมุทโย โหติ.}

\prose{จกฺขุญฺจ ปฏิจฺจ รูเป จ อุปฺปชฺชติ จกฺขุวิญฺญาณํ, ติณฺณํ สงฺคติ ผสฺโส, ผสฺสปจฺจยา เวทนา, เวทนาปจฺจยา ตณฺหา.\csromanspacer{} ตสฺสาเยว ตณฺหาย อเสสวิราคนิโรธา อุปาทานนิโรโธ, อุปาทานนิโรธา ภวนิโรโธ, ภวนิโรธา ชาตินิโรโธ, ชาตินิโรธา ชรามรณํ โสกปริเทวทุกฺขโทมนสฺสุปายาสา นิรุชฺฌนฺติ.\csromanspacer{} เอวเมตสฺส เกวลสฺส ทุกฺขกฺขนฺธสฺส นิโรโธ โหติ ฯเปฯ.\csromanspacer{} ชิวฺหญฺจ ปฏิจฺจ รเส จ อุปฺปชฺชติ ฯเปฯ.\csromanspacer{} มนญฺจ ปฏิจฺจ ธมฺเม จ อุปฺปชฺชติ มโนวิญฺญาณํ, ติณฺณํ สงฺคติ ผสฺโส, ผสฺสปจฺจยา เวทนา, เวทนาปจฺจยา ตณฺหา.\csromanspacer{} ตสฺสาเยว ตณฺหาย อเสสวิราคนิโรธา}

\vspace{\tipitakaparskip}
\csromancenter{สฬายตนสํยุตฺต อุปาทานนิโรโธ, อุปาทานนิโรธา ฯเปฯ.\csromanspacer{} เอวเมตสฺส เกวลสฺส ทุกฺขกฺขนฺธสฺส นิโรโธ โหตี”ติ.}
\csromanmatikaanchor{mtk.152}
\csromanmatikaanchor{mtk.153}
\csromantocmarknopage{section}{(12) 2. โลกกามคุณวคฺค}{mtk.152}
\bookmark[dest={mtk.152},level=1]{(12) 2. โลกกามคุณวคฺค}
\csromantocmarkref{subsection}{1-2. มารปาสสุตฺต}{mtk.153}
\bookmark[dest={mtk.153},level=2]{1-2. มารปาสสุตฺต}
\prose{\csromanfolio{311}เตน โข ปน สมเยน อญฺญตโร ภิกฺขุ ภควโต อุปสฺสุติ\footnote{อุปสฺสุติํ (สี, ก)} ฐิโต โหติ.\csromanspacer{} อทฺทสา โข ภควา ตํ ภิกฺขุํ อุปสฺสุติ ฐิตํ, ทิสฺวาน ตํ ภิกฺขุํ เอตทโวจ “อสฺโสสิ โน ตฺวํ ภิกฺขุ อิมํ ธมฺมปริยายนฺ”ติ.\csromanspacer{} เอวํ ภนฺเต.\csromanspacer{} อุคฺคณฺหาหิ ตฺวํ ภิกฺขุ อิมํ ธมฺมปริยายํ, ปริยาปุณาหิ ตฺวํ ภิกฺขุ อิมํ ธมฺมปริยายํ, ธาเรหิ ตฺวํ ภิกฺขุ อิมํ ธมฺมปริยายํ, อตฺถสํหิโตยํ ภิกฺขุ ธมฺมปริยาโย อาทิพฺรหฺมจริยโกติ.\csromanspacer{}\csromanspacer{}ทสมํ.\csromanspacer{} โยคกฺเขมิวคฺโค ปฐโม.}

\titlehead{\textbf{ตสฺสุทฺทานํ}}
\csromangathameasure{โยคกฺเขมิ อุปาทาย, ทุกฺขํ โลโก จ เสยฺโย จ. \\ สํโยชนํ อุปาทานํ, เทฺว ปริชานํ อุปสฺสุตีติ.}
\csromangathagroup{\csromangathastanza{โยคกฺเขมิ อุปาทาย, ทุกฺขํ โลโก จ เสยฺโย จ. \\ สํโยชนํ อุปาทานํ, เทฺว ปริชานํ อุปสฺสุตีติ.}}

\csromanheader{1}{\textbf{(12) 2. โลกกามคุณวคฺค}}
\csromanheader{2}{1. ปฐมมารปาสสุตฺต}
\proseitem{114}{สนฺติ ภิกฺขเว จกฺขุวิญฺเญยฺยา รูปา อิฏฺฐา กนฺตา มนาปา ปิยรูปา กามูปสํหิตา รชนียา.\csromanspacer{} ตญฺเจ ภิกฺขุ อภินนฺทติ อภิวทติ อชฺโฌสาย ติฏฺฐติ.\csromanspacer{} อยํ วุจฺจติ ภิกฺขเว ภิกฺขุ อาวาสคโต มารสฺส, มารสฺส วสํ คโต\footnote{วสคโต (สี-ฏฺฐ, สฺยา-ฏฺฐ)}, ปฏิมุกฺก’สฺส มารปาโส.\csromanspacer{} พทฺโธ โส มารพนฺธเนน, ยถากามกรณีโย ปาปิมโต ฯเปฯ.}

\prose{สนฺติ ภิกฺขเว ชิวฺหาวิญฺเญยฺยา รสา อิฏฺฐา กนฺตา มนาปา ปิยรูปา กามูปสํหิตา รชนียา, ตญฺเจ ภิกฺขุ อภินนฺทติ อภิวทติ อชฺโฌสาย ติฏฺฐติ.\csromanspacer{} อยํ วุจฺจติ ภิกฺขเว ภิกฺขุ อาวาสคโต มารสฺส, มารสฺส วสํ คโต, ปฏิมุกฺก’สฺส มารปาโส.\csromanspacer{} พทฺโธ โส มารพนฺธเนน ฯเปฯ.}

\gambhira{สฬายตนวคฺคสํยุตฺตปาฬิ}
\prose{\csromanfolio{312}สนฺติ ภิกฺขเว มโนวิญฺเญยฺยา ธมฺมา อิฏฺฐา กนฺตา มนาปา ปิยรูปา กามูปสํหิตา รชนียา.\csromanspacer{} ตญฺเจ ภิกฺขุ อภินนฺทติ อภิวทติ อชฺโฌสาย ติฏฺฐติ.\csromanspacer{} อยํ วุจฺจติ ภิกฺขเว ภิกฺขุ อาวาสคโต มารสฺส, มารสฺส วสํ คโต, ปฏิมุกฺก’สฺส มารปาโส.\csromanspacer{} พทฺโธ โส มารพนฺธเนน, ยถากามกรณีโย ปาปิมโต ฯเปฯ.}

\prose{สนฺติ จ โข ภิกฺขเว จกฺขุวิญฺเญยฺยา รูปา อิฏฺฐา กนฺตา มนาปา ปิยรูปา กามูปสํหิตา รชนียา.\csromanspacer{} ตญฺเจ ภิกฺขุ นาภินนฺทติ นาภิวทติ นาชฺโฌสาย ติฏฺฐติ.\csromanspacer{} อยํ วุจฺจติ ภิกฺขเว ภิกฺขุ นาวาสคโต มารสฺส, น มารสฺส วสํ คโต, อุมฺมุกฺก’สฺส มารปาโส.\csromanspacer{} มุตฺโต โส มารพนฺธเนน, น ยถากามกรณีโย ปาปิมโต ฯเปฯ.}

\prose{สนฺติ ภิกฺขเว ชิวฺหาวิญฺเญยฺยา รสา อิฏฺฐา กนฺตา มนาปา ปิยรูปา กามูปสํหิตา รชนียา.\csromanspacer{} ตญฺเจ ภิกฺขุ นาภินนฺทติ นาภิวทติ นาชฺโฌสาย ติฏฺฐติ.\csromanspacer{} อยํ วุจฺจติ ภิกฺขเว ภิกฺขุ นาวาสคโต มารสฺส, น มารสฺส วสํ คโต, อุมฺมุกฺก’สฺส มารปาโส.\csromanspacer{} มุตฺโต โส มารพนฺธเนน, น ยถากามกรณีโย ปาปิมโต ฯเปฯ.}

\prose{สนฺติ ภิกฺขเว มโนวิญฺเญยฺยา ธมฺมา อิฏฺฐา กนฺตา มนาปา ปิยรูปา กามูปสํหิตา รชนียา.\csromanspacer{} ตญฺเจ ภิกฺขุ นาภินนฺทติ นาภิวทติ นาชฺโฌสายา ติฏฺฐติ.\csromanspacer{} อยํ วุจฺจติ ภิกฺขเว ภิกฺขุ นาวาสคโต มารสฺส, น มารสฺส วสํ คโต, อุมฺมุกฺก’สฺส มารปาโส.\csromanspacer{} มุตฺโต โส มารพนฺธเนน, น ยถากามกรณีโย ปาปิมโตติ.\csromanspacer{}\csromanspacer{}ปฐมํ.}

\csromanheader{2}{2. ทุติยมารปาสสุตฺต}
\proseitem{115}{สนฺติ ภิกฺขเว จกฺขุวิญฺเญยฺยา รูปา อิฏฺฐา กนฺตา มนาปา ปิยรูปา กามูปสํหิตา รชนียา.\csromanspacer{} ตญฺเจ ภิกฺขุ อภินนฺทติ อภิวทติ อชฺโฌสาย ติฏฺฐติ.\csromanspacer{} อยํ วุจฺจติ ภิกฺขเว ภิกฺขุ พทฺโธ จกฺขุวิญฺเญยฺเยสุ รูเปสุ อาวาสคโต มารสฺส, มารสฺส วสํ คโต, ปฏิมุกฺก’สฺส มารปาโส.\csromanspacer{} พทฺโธ โส มารพนฺธเนน, ยถากามกรณีโย ปาปิมโต ฯเปฯ.}

\prose{สนฺติ ภิกฺขเว ชิวฺหาวิญฺเญยฺยา รสา ฯเปฯ.\csromanspacer{} สนฺติ ภิกฺขเว มโนวิญฺเญยฺยา ธมฺมา อิฏฺฐา กนฺตา มนาปา ปิยรูปา กามูปสํหิตา รชนียา.\csromanspacer{} ตญฺเจ ภิกฺขุ}

\vspace{\tipitakaparskip}
\csromancenter{สฬายตนสํยุตฺต อภินนฺทติ อภิวทติ อชฺโฌสาย ติฏฺฐติ.\csromanspacer{} อยํ วุจฺจติ ภิกฺขเว ภิกฺขุ พทฺโธ มโนวิญฺเญยฺเยสุ ธมฺเมสุ อาวาสคโต มารสฺส, มารสฺส วสํ คโต, ปฏิมุกฺก’สฺส มารปาโส.\csromanspacer{} พทฺโธ โส มารพนฺธเนน, ยถากามกรณีโย ปาปิมโต ฯเปฯ.}
\prose{\csromanfolio{313}สนฺติ จ โข ภิกฺขเว จกฺขุวิญฺเญยฺยา รูปา อิฏฺฐา กนฺตา มนาปา ปิยรูปา กามูปสํหิตา รชนียา.\csromanspacer{} ตญฺเจ ภิกฺขุ นาภินนฺทติ นาภิวทติ นาชฺโฌสาย ติฏฺฐติ.\csromanspacer{} อยํ วุจฺจติ ภิกฺขเว ภิกฺขุ มุตฺโต จกฺขุวิญฺเญยฺเยหิ รูเปหิ, นาวาสคโต มารสฺส, น มารสฺส วสํ คโต, อุมฺมุกฺก’สฺส มารปาโส.\csromanspacer{} มุตฺโต โส มารพนฺธเนน, น ยถากามกรณีโย ปาปิมโต ฯเปฯ.}

\prose{สนฺติ ภิกฺขเว ชิวฺหาวิญฺเญยฺยา รสา ฯเปฯ.\csromanspacer{} สนฺติ ภิกฺขเว มโนวิญฺเญยฺยา ธมฺมา อิฏฺฐา กนฺตา มนาปา ปิยรูปา กามูปสํหิตา รชนียา.\csromanspacer{} ตญฺเจ ภิกฺขุ นาภินนฺทติ นาภิวทติ นาชฺโฌสาย ติฏฺฐติ.\csromanspacer{} อยํ วุจฺจติ ภิกฺขเว ภิกฺขุ มุตฺโต มโนวิญฺเญยฺเยหิ ธมฺเมหิ, นาวาสคโต มารสฺส, น มารสฺส วสํ คโต, อุมฺมุกฺก’สฺส มารปาโส.\csromanspacer{} มุตฺโต โส มารพนฺธเนน, น ยถากามกรณีโย ปาปิมโตติ.\csromanspacer{}\csromanspacer{}ทุติยํ.}

\csromanmatikaanchor{mtk.154}
\csromantocmarkref{subsection}{3. โลกนฺตคมนสุตฺต}{mtk.154}
\bookmark[dest={mtk.154},level=2]{3. โลกนฺตคมนสุตฺต}
\csromanheader{2}{3. โลกนฺตคมนสุตฺต}
\proseitem{116}{นาหํ ภิกฺขเว คมเนน โลกสฺส อนฺตํ ญาเตยฺยํ ทฏฺเฐยฺยํ\footnote{ทิฏฺเฐยฺยํ (สฺยา, กํ, ก)} ปตฺเตยฺยนฺติ วทามิ, น จ ปนาหํ ภิกฺขเว อปฺปตฺวา โลกสฺส อนฺตํ ทุกฺขสฺส อนฺตกิริยํ วทามีติ.\csromanspacer{} อิทํ วตฺวา ภควา อุฏฺฐายาสนา วิหารํ ปาวิสิ.\csromanspacer{} อถ โข เตสํ ภิกฺขูนํ อจิรปกฺกนฺตสฺส ภควโต เอตทโหสิ “อิทํ โข โน อาวุโส ภควา สํขิตฺเตน อุทฺเทสํ อุทฺทิสิตฺวา วิตฺถาเรน อตฺถํ อวิภชิตฺวา อุฏฺฐายาสนา วิหารํ ปวิฏฺโฐ, ‘นาหํ ภิกฺขเว คมเนน โลกสฺส อนฺตํ ญาเตยฺยํ ทฏฺเฐยฺยํ ปตฺเตยฺยนฺติ วทามิ, น จ ปนาหํ ภิกฺขเว อปฺปตฺวา โลกสฺส อนฺตํ ทุกฺขสฺส อนฺตกิริยํ วทามี’ติ, โก นุ โข อิมสฺส ภควตา สํขิตฺเตน อุทฺเทสสฺส อุทฺทิฏฺฐสฺส วิตฺถาเรน อตฺถํ อวิภตฺตสฺส วิตฺถาเรน อตฺถํ วิภเชยฺยา”ติ.}

\gambhira{สฬายตนวคฺคสํยุตฺตปาฬิ}
\prose{\csromanfolio{314}อถ โข เตสํ ภิกฺขูนํ เอตทโหสิ “อยํ โข อายสฺมา อานนฺโท สตฺถุ เจว สํวณฺณิโต, สมฺภาวิโต จ วิญฺญูนํ สพฺรหฺมจารีนํ, ปโหติ จายสฺมา อานนฺโท อิมสฺส ภควตา สํขิตฺเตน อุทฺเทสสฺส อุทฺทิฏฺฐสฺส วิตฺถาเรน อตฺถํ อวิภตฺตสฺส วิตฺถาเรน อตฺถํ วิภชิตุํ.\csromanspacer{} ยํนูน มยํ เยนายสฺมา อานนฺโท เตนุปสงฺกเมยฺยาม, อุปสงฺกมิตฺวา อายสฺมนฺตํ อานนฺทํ เอตมตฺถํ ปฏิปุจฺเฉยฺยามา”ติ.}

\prose{อถ โข เต ภิกฺขู เยนายสฺมา อานนฺโท เตนุปสงฺกมิํสุ, อุปสงฺกมิตฺวา อายสฺมตา อานนฺเทน สทฺธิํ สมฺโมทิํสุ, สมฺโมทนียํ กถํ สารณียํ วีติสาเรตฺวา เอกมนฺตํ นิสีทิํสุ, เอกมนฺตํ นิสินฺโน โข เต ภิกฺขู อายสฺมนฺตํ อานนฺทํ เอตทโวจุํ\csromandash{}}

\prose{อิทํ โข โน อาวุโส อานนฺท ภควา สํขิตฺเตน อุทฺเทสํ อุทฺทิสิตฺวา วิตฺถาเรน อตฺถํ อวิภชิตฺวา อุฏฺฐายาสนา วิหารํ ปวิฏฺโฐ, “นาหํ ภิกฺขเว คมเนน โลกสฺส อนฺตํ ญาเตยฺยํ ทฏฺเฐยฺยํ ปตฺเตยฺยนฺติ วทามิ, น จ ปนาหํ ภิกฺขเว อปฺปตฺวา โลกสฺส อนฺตํ ทุกฺขสฺส อนฺตกิริยํ วทามี”ติ.\csromanspacer{} เตสํ โน อาวุโส อมฺหากํ อจิรปกฺกนฺตสฺส ภควโต เอตทโหสิ “อิทํ โข โน อาวุโส ภควา สํขิตฺเตน อุทฺเทสํ อุทฺทิสิตฺวา วิตฺถาเรน อตฺถํ อวิภชิตฺวา อุฏฺฐายาสนา วิหารํ ปวิฏฺโฐ, ‘นาหํ ภิกฺขเว คมเนน โลกสฺส อนฺตํ ญาเตยฺยํ ทฏฺเฐยฺยํ ปตฺเตยฺยนฺติ วทามิ, น จ ปนาหํ ภิกฺขเว อปฺปตฺวา โลกสฺส อนฺตํ ทุกฺขสฺส อนฺตกิริยํ วทามี’ติ, โก นุ โข อิมสฺส ภควตา สํขิตฺเตน อุทฺเทสสฺส อุทฺทิฏฺฐสฺส วิตฺถาเรน อตฺถํ อวิภตฺตสฺส วิตฺถาเรน อตฺถํ วิภเชยฺยา”ติ.\csromanspacer{} เตสํ โน อาวุโส อมฺหากํ เอตทโหสิ “อยํ โข อาวุโส อายสฺมา อานนฺโท สตฺถุ เจว สํวณฺณิโต, สมฺภาวิโต จ วิญฺญูนํ สพฺรหฺมจารีนํ, ปโหติ จายสฺมา อานนฺโท อิมสฺส ภควตา สํขิตฺเตน อุทฺเทสสฺส อุทฺทิฏฺฐสฺส วิตฺถาเรน อตฺถํ อวิภตฺตสฺส วิตฺถาเรน อตฺถํ วิภชิตุํ.\csromanspacer{} ยํนูน มยํ เยนายสฺมา อานนฺโท เตนุปสงฺกเมยฺยาม, อุปสงฺกมิตฺวา อายสฺมนฺตํ อานนฺทํ เอตมตฺถํ ปฏิปุจฺเฉยฺยามา”ติ.\csromanspacer{} วิภชตายสฺมา อานนฺโทติ.}

\prose{เสยฺยถาปิ อาวุโส ปุริโส สารตฺถิโก สารคเวสี สารปริเยสนํ จรมาโน มหโต รุกฺขสฺส ติฏฺฐโต สารวโต อติกฺกมฺเมว มูลํ อติกฺกมฺเมว ขนฺธํ สาขาปลาเส สารํ ปริเยสิตพฺพํ มญฺเญยฺย, เอวํสมฺปทมิทํ อายสฺมนฺตานํ สตฺถริ สมฺมุขีภูเต ตํ ภควนฺตํ}

\vspace{\tipitakaparskip}
\csromancenter{สฬายตนสํยุตฺต อติสิตฺวา อเมฺห เอตมตฺถํ ปฏิปุจฺฉิตพฺพํ มญฺญถ\footnote{มญฺเญถ (อิ, ก)}.\csromanspacer{} โส หาวุโส ภควา ชานํ ชานาติ ปสฺสํ ปสฺสติ จกฺขุภูโต ญาณภูโต ธมฺมภูโต พฺรหฺมภูโต วตฺตา ปวตฺตา อตฺถสฺส นินฺเนตา อมตสฺส ทาตา ธมฺมสฺสามี ตถาคโต, โส เจว ปเนตสฺส กาโล อโหสิ, ยํ ภควนฺตํเยว เอตมตฺถํ ปฏิปุจฺเฉยฺยาถ.\csromanspacer{} ยถา โว ภควา พฺยากเรยฺย, ตถา โว ธาเรยฺยาถาติ.}
\prose{\csromanfolio{315}อทฺธาวุโส อานนฺท ภควา ชานํ ชานาติ ปสฺสํ ปสฺสติ จกฺขุภูโต ญาณภูโต ธมฺมภูโต พฺรหฺมภูโต วตฺตา ปวตฺตา อตฺถสฺส นินฺเนตา อมตสฺส ทาตา ธมฺมสฺสามี ตถาคโต, โส เจว ปเนตสฺส กาโล อโหสิ, ยํ ภควนฺตํเยว เอตมตฺถํ ปฏิปุจฺเฉยฺยาม.\csromanspacer{} ยถา โน ภควา พฺยากเรยฺย, ตถา นํ ธาเรยฺยาม.\csromanspacer{} อปิ จายสฺมา อานนฺโท สตฺถุ เจว สํวณฺณิโต, สมฺภาวิโต จ วิญฺญูนํ สพฺรหฺมจารีนํ, ปโหติ จายสฺมา อานนฺโท อิมสฺส ภควตา สํขิตฺเตน อุทฺเทสสฺส อุทฺทิฏฺฐสฺส วิตฺถาเรน อตฺถํ อวิภตฺตสฺส วิตฺถาเรน อตฺถํ วิภชิตุํ, วิภชตายสฺมา อานนฺโท อครุํ กริตฺวาติ.}

\prose{เตนหาวุโส สุณาถ สาธุกํ มนสิ กโรถ, ภาสิสฺสามีติ. “เอวมาวุโสติ โข เต ภิกฺขู อายสฺมโต อานนฺทสฺส ปจฺจสฺโสสุํ.\csromanspacer{} อายสฺมา อานนฺโท เอตทโวจ\csromandash{}}

\prose{ยํ โข โว อาวุโส ภควา สํขิตฺเตน อุทฺเทสํ อุทฺทิสิตฺวา วิตฺถาเรน อตฺถํ อวิภชิตฺวา อุฏฺฐายาสนา วิหารํ ปวิฏฺโฐ, “นาหํ ภิกฺขเว คมเนน โลกสฺส อนฺตํ ญาเตยฺยํ ทฏฺเฐยฺยํ ปตฺเตยฺยนฺติ วทามิ, น จ ปนาหํ ภิกฺขเว อปฺปตฺวา โลกสฺส อนฺตํ ทุกฺขสฺส อนฺตกิริยํ วทามี”ติ.\csromanspacer{} อิมสฺส ขฺวาหํ อาวุโส ภควตา สํขิตฺเตน อุทฺเทสสฺส อุทฺทิฏฺฐสฺส วิตฺถาเรน อตฺถํ อวิภตฺตสฺส วิตฺถาเรน อตฺถํ อาชานามิ.\csromanspacer{} เยน โข อาวุโส โลกสฺมิํ โลกสญฺญี โหติ โลกมานี, อยํ วุจฺจติ อริยสฺส วินเย โลโก.\csromanspacer{} เกน จาวุโส โลกสฺมิํ โลกสญฺญี โหติ โลกมานี, จกฺขุนา โข อาวุโส โลกสฺมิํ โลกสญฺญี โหติ โลกมานี.\csromanspacer{} โสเตน โข อาวุโส.\csromanspacer{} ฆาเนน โข อาวุโส.\csromanspacer{} ชิวฺหาย โข อาวุโส โลกสฺมิํ โลกสญฺญี โหติ}

\gambhira{สฬายตนวคฺคสํยุตฺตปาฬิ}
\prose{\csromanfolio{316}โลกมานี.\csromanspacer{} กาเยน โข อาวุโส.\csromanspacer{} มเนน โข อาวุโส โลกสฺมิํ โลกสญฺญี โหติ โลกมานี.\csromanspacer{} เยน โข อาวุโส โลกสฺมิํ โลกสญฺญี โหติ โลกมานี, อยํ วุจฺจติ อริยสฺส วินเย โลโก.\csromanspacer{} ยํ โข โว อาวุโส ภควา สํขิตฺเตน อุทฺเทสํ อุทฺทิสิตฺวา วิตฺถาเรน อตฺถํ อวิภชิตฺวา อุฏฺฐายาสนา วิหารํ ปวิฏฺโฐ, “นาหํ ภิกฺขเว คมเนน โลกสฺส อนฺตํ ญาเตยฺยํ ทฏฺเฐยฺยํ ปตฺเตยฺยนฺติ วทามิ, น จ ปนาหํ ภิกฺขเว อปฺปตฺวา โลกสฺส อนฺตํ ทุกฺขสฺส อนฺตกิริยํ วทามี”ติ.\csromanspacer{} อิมสฺส ขฺวาหํ อาวุโส ภควตา สํขิตฺเตน อุทฺเทสสฺส อุทฺทิฏฺฐสฺส วิตฺถาเรน อตฺถํ อวิภตฺตสฺส เอวํ วิตฺถาเรน อตฺถํ อาชานามิ.\csromanspacer{} อากงฺขมานา จ ปน ตุเมฺห อายสฺมนฺโต ภควนฺตํเยว อุปสงฺกมิตฺวา เอตมตฺถํ ปฏิปุจฺเฉยฺยาถ.\csromanspacer{} ยถา โว ภควา พฺยากโรติ, ตถา นํ ธาเรยฺยาถาติ.}

\prose{“เอวมาวุโส”ติ โข เต ภิกฺขู อายสฺมโต อานนฺทสฺส ปฏิสฺสุตฺวา อุฏฺฐายาสนา เยน ภควา เตนุปสงฺกมิํสุ, อุปสงฺกมิตฺวา ภควนฺตํ อภิวาเทตฺวา เอกมนฺตํ นิสีทิํสุ, เอกมนฺตํ นิสินฺนา โข เต ภิกฺขู ภควนฺตํ เอตทโวจุํ\csromandash{}}

\prose{ยํ โข โน ภนฺเต ภควา สํขิตฺเตน อุทฺเทสํ อุทฺทิสิตฺวา วิตฺถาเรน อตฺถํ อวิภชิตฺวา อุฏฺฐายาสนา วิหารํ ปวิฏฺโฐ, “นาหํ ภิกฺขเว คมเนน โลกสฺส อนฺตํ ญาเตยฺยํ ทฏฺเฐยฺยํ ปตฺเตยฺยนฺติ วทามิ, น จ ปนาหํ ภิกฺขเว อปฺปตฺวา โลกสฺส อนฺตํ ทุกฺขสฺส อนฺตกิริยํ วทามี”ติ.\csromanspacer{} เตสํ โน ภนฺเต อมฺหากํ อจิรปกฺกนฺตสฺส ภควโต เอตทโหสิ “อิทํ โข โน อาวุโส ภควา สํขิตฺเตน อุทฺเทสํ อุทฺทิสิตฺวา วิตฺถาเรน อตฺถํ อวิภชิตฺวา อุฏฺฐายาสนา วิหารํ ปวิฏฺโฐ, ‘นาหํ ภิกฺขเว คมเนน โลกสฺส อนฺตํ ญาเตยฺยํ ทฏฺเฐยฺยํ ปตฺเตยฺยนฺติ วทามิ, น จ ปนาหํ ภิกฺขเว อปฺปตฺวา โลกสฺส อนฺตํ ทุกฺขสฺส อนฺตกิริยํ วทามี’ติ.\csromanspacer{} โก นุ โข อิมสฺส ภควตา สํขิตฺเตน อุทฺเทสสฺส อุทฺทิฏฺฐสฺส วิตฺถาเรน อตฺถํ อวิภตฺตสฺส วิตฺถาเรน อตฺถํ วิภเชยฺยา”ติ.\csromanspacer{} เตสํ โน ภนฺเต อมฺหากํ เอตทโหสิ “อยํ โข อายสฺมา อานนฺโท สตฺถุ เจว สํวณฺณิโต, สมฺภาวิโต จ วิญฺญูนํ สพฺรหฺมจารีนํ, ปโหติ จายสฺมา อานนฺโท อิมสฺส ภควตา สํขิตฺเตน อุทฺเทสสฺส อุทฺทิฏฺฐสฺส วิตฺถาเรน อตฺถํ อวิภตฺตสฺส วิตฺถาเรน อตฺถํ วิภชิตุํ, ยํนูน มยํ เยนายสฺมา อานนฺโท เตนุปสงฺกเมยฺยาม, อุปสงฺกมิตฺวา อายสฺมนฺตํ อานนฺทํ เอตมตฺถํ ปฏิปุจฺเฉยฺยามา”ติ.\csromanspacer{} อถ โข}

\vspace{\tipitakaparskip}
\csromancenter{สฬายตนสํยุตฺต มยํ ภนฺเต เยนายสฺมา อานนฺโท เตนุปสงฺกมิมฺห, อุปสงฺกมิตฺวา อายสฺมนฺตํ อานนฺทํ เอตมตฺถํ ปฏิปุจฺฉิมฺห, เตสํ โน ภนฺเต อายสฺมตา อานนฺเทน อิเมหิ อากาเรหิ อิเมหิ ปเทหิ อิเมหิ พฺยญฺชเนหิ อตฺโถ วิภตฺโตติ.}
\prose{\csromanfolio{317}ปณฺฑิโต ภิกฺขเว อานนฺโท, มหาปญฺโญ ภิกฺขเว อานนฺโท.\csromanspacer{} มํ เจปิ ตุเมฺห ภิกฺขเว เอตมตฺถํ ปฏิปุจฺเฉยฺยาถ, อหมฺปิ ตํ เอวเมวํ พฺยากเรยฺยํ.\csromanspacer{} ยถา ตํ อานนฺเทน พฺยากตํ, เอโส เจเวตสฺส อตฺโถ, เอวญฺจ นํ ธาเรยฺยาถาติ.\csromanspacer{}\csromanspacer{}ตติยํ.}

\csromanmatikaanchor{mtk.155}
\csromantocmarkref{subsection}{4. กามคุณสุตฺต}{mtk.155}
\bookmark[dest={mtk.155},level=2]{4. กามคุณสุตฺต}
\csromanheader{2}{4. กามคุณสุตฺต}
\proseitem{117}{ปุพฺเพว เม ภิกฺขเว สมฺโพธา อนภิสมฺพุทฺธสฺส โพธิสตฺตสฺเสว สโต เอตทโหสิ “เยเม ปญฺจ กามคุณา เจตโส สมฺผุฏฺฐปุพฺพา อตีตา นิรุทฺธา วิปริณตา, ตตฺร เม จิตฺตํ พหุลํ คจฺฉมานํ คจฺเฉยฺย ปจฺจุปฺปนฺเนสุ วา อปฺปํ วา อนาคเตสุ”.\csromanspacer{} ตสฺส มยฺหํ ภิกฺขเว เอตทโหสิ “เยเม ปญฺจ กามคุณา เจตโส สมฺผุฏฺฐปุพฺพา อตีตา นิรุทฺธา วิปริณตา, ตตฺร เม อตฺตรูเปน อปฺปมาโท สติ เจตโส อารกฺโข กรณีโย”.\csromanspacer{} ตสฺมาติห ภิกฺขเว ตุมฺหากมฺปิ เย เต ปญฺจ กามคุณา เจตโส สมฺผุฏฺฐปุพฺพา อตีตา นิรุทฺธา วิปริณตา, ตตฺร โว จิตฺตํ พหุลํ คจฺฉมานํ คจฺเฉยฺย ปจฺจุปฺปนฺเนสุ วา อปฺปํ วา อนาคเตสุ.\csromanspacer{} ตสฺมาติห ภิกฺขเว ตุมฺหากมฺปิ เย เต ปญฺจ กามคุณา เจตโส สมฺผุฏฺฐปุพฺพา อตีตา นิรุทฺธา วิปริณตา, ตตฺร โว อตฺตรูเปหิ อปฺปมาโท สติ เจตโส อารกฺโข กรณีโย.\csromanspacer{} ตสฺมาติห ภิกฺขเว เส อายตเน เวทิตพฺเพ, ยตฺถ จกฺขุ จ นิรุชฺฌติ รูปสญฺญา จ นิรุชฺฌติ เส อายตเน เวทิตพฺเพ ฯเปฯ ยตฺถ ชิวฺหา จ นิรุชฺฌติ รสสญฺญา จ นิรุชฺฌติ เส อายตเน เวทิตพฺเพ ฯเปฯ ยตฺถ มโน จ นิรุชฺฌติ ธมฺมสญฺญา จ นิรุชฺฌติ เส อายตเน เวทิตพฺเพติ.\csromanspacer{} อิทํ วตฺวา ภควา อุฏฺฐายาสนา วิหารํ ปาวิสิ.}

\prose{อถ โข เตสํ ภิกฺขูนํ อจิรปกฺกนฺตสฺส ภควโต เอตทโหสิ “อิทํ โข โน อาวุโส ภควา สํขิตฺเตน อุทฺเทสํ อุทฺทิสิตฺวา วิตฺถาเรน อตฺถํ อวิภชิตฺวา อุฏฺฐายาสนา วิหารํ ปวิฏฺโฐ, ‘ตสฺมาติห ภิกฺขเว เส อายตเน เวทิตพฺเพ, ยตฺถ จกฺขุ จ นิรุชฺฌติ รูปสญฺญา จ นิรุชฺฌติ เส อายตเน เวทิตพฺเพ ฯเปฯ ยตฺถ ชิวฺหา จ นิรุชฺฌติ รสสญฺญา จ นิรุชฺฌติ}

\gambhira{สฬายตนวคฺคสํยุตฺตปาฬิ}
\prose{\csromanfolio{318}เส อายตเน เวทิตพฺเพ ฯเปฯ ยตฺถ มโน จ นิรุชฺฌติ ธมฺมสญฺญา จ นิรุชฺฌติ เส อายตเน เวทิตพฺเพ’ติ.\csromanspacer{} โก นุ โข อิมสฺส ภควตา สํขิตฺเตน อุทฺเทสสฺส อุทฺทิฏฺฐสฺส วิตฺถาเรน อตฺถํ อวิภตฺตสฺส วิตฺถาเรน อตฺถํ วิภเชยฺยา”ติ.}

\prose{อถ โข เตสํ ภิกฺขูนํ เอตทโหสิ “อยํ โข อายสฺมา อานนฺโท สตฺถุ เจว สํวณฺณิโต, สมฺภาวิโต จ วิญฺญูนํ สพฺรหฺมจารีนํ, ปโหติ จายสฺมา อานนฺโท อิมสฺส ภควตา สํขิตฺเตน อุทฺเทสสฺส อุทฺทิฏฺฐสฺส วิตฺถาเรน อตฺถํ อวิภตฺตสฺส วิตฺถาเรน อตฺถํ วิภชิตุํ.\csromanspacer{} ยํนูน มยํ เยนายสฺมา อานนฺโท เตนุปสงฺกเมยฺยาม, อุปสงฺกมิตฺวา อายสฺมนฺตํ อานนฺทํ เอตมตฺถํ ปฏิปุจฺเฉยฺยามา”ติ.}

\prose{อถ โข เต ภิกฺขู เยนายสฺมา อานนฺโท เตนุปสงฺกมิํสุ, อุปสงฺกมิตฺวา อายสฺมตา อานนฺเทน สทฺธิํ สมฺโมทิํสุ, สมฺโมทนียํ กถํ สารณียํ วีติสาเรตฺวา เอกมนฺตํ นิสีทิํสุ, เอกมนฺตํ นิสินฺนา โข เต ภิกฺขู อายสฺมนฺตํ อานนฺทํ เอตทโวจุํ\csromandash{}}

\prose{อิทํ โข โน อาวุโส อานนฺท ภควา สํขิตฺเตน อุทฺเทสํ อุทฺทิสิตฺวา วิตฺถาเรน อตฺถํ อวิภชิตฺวา อุฏฺฐายาสนา วิหารํ ปวิฏฺโฐ, “ตสฺมาติห ภิกฺขเว เส อายตเน เวทิตพฺเพ.\csromanspacer{} ยตฺถ จกฺขุ จ นิรุชฺฌติ รูปสญฺญา จ นิรุชฺฌติ เส อายตเน เวทิตพฺเพ ฯเปฯ.\csromanspacer{} ยตฺถ ชิวฺหา จ นิรุชฺฌติ รสสญฺญา จ นิรุชฺฌติ เส อายตเนน เวทิตพฺเพ ฯเปฯ.\csromanspacer{} ยตฺถ มโน จ นิรุชฺฌติ ธมฺมสญฺญา จ นิรุชฺฌติ เส อายตเน เวทิตพฺเพ”ติ.\csromanspacer{} เตสํ โน อาวุโส อมฺหากํ อจิรปกฺกนฺตสฺส ภควโต เอตทโหสิ “อิทํ โข โน อาวุโส ภควา สํขิตฺเตน อุทฺเทสํ อุทฺทิสิตฺวา วิตฺถาเรน อตฺถํ อวิภชิตฺวา อุฏฺฐายาสนา วิหารํ ปวิฏฺโฐ, ‘ตสฺมาติห ภิกฺขเว เส อายตเน เวทิตพฺเพ.\csromanspacer{} ยตฺถ จกฺขุ จ นิรุชฺฌติ รูปสญฺญา จ นิรุชฺฌติ เส อายตเน เวทิตพฺเพ ฯเปฯ.\csromanspacer{} ยตฺถ ชิวฺหา จ นิรุชฺฌติ รสสญฺญา จ นิรุชฺฌติ เส อายตเน เวทิตพฺเพ ฯเปฯ.\csromanspacer{} ยตฺถ มโน จ นิรุชฺฌติ ธมฺมสญฺญา จ นิรุชฺฌติ เส อายตเน เวทิตพฺเพ’ติ.\csromanspacer{} โก นุ โข อิมสฺส ภควตา สํขิตฺเตน อุทฺเทสสฺส อุทฺทิฏฺฐสฺส วิตฺถาเรน อตฺถํ อวิภตฺตสฺส วิตฺถาเรน อตฺถํ วิภเชยฺยา”ติ.\csromanspacer{} เตสํ โน อาวุโส อมฺหากํ เอตทโหสิ “อยํ โข อายสฺมา อานนฺโท สตฺถุ เจว สํวณฺณิโต, สมฺภาวิโต จ วิญฺญูนํ สพฺรหฺมจารีนํ, ปโหติ จายสฺมา อานนฺโท อิมสฺส ภควตา}

\vspace{\tipitakaparskip}
\csromancenter{สฬายตนสํยุตฺต สํขิตฺเตน อุทฺเทสสฺส อุทฺทิฏฺฐสฺส วิตฺถาเรน อตฺถํ อวิภตฺตสฺส วิตฺถาเรน อตฺถํ วิภชิตุํ, ยํนูน มยํ เยนายสฺมา อานนฺโท เตนุปสงฺกเมยฺยาม, อุปสงฺกมิตฺวา อายสฺมนฺตํ อานนฺทํ เอตมตฺถํ ปฏิปุจฺเฉยฺยามา”ติ.\csromanspacer{} วิภชตายสฺมา อานนฺโทติ.}
\prose{\csromanfolio{319}เสยฺยถาปิ อาวุโส ปุริโส สารตฺถิโก สารคเวสี สารปริเยสนํ จรมาโน มหโต รุกฺขสฺส ฯเปฯ วิภชตายสฺมา อานนฺโท อครุํ กริตฺวาติ.}

\prose{เตนหาวุโส สุณาถ สาธุกํ มนสิ กโรถ, ภาสิสฺสามีติ. “เอวมาวุโส”ติ โข เต ภิกฺขู อายสฺมโต อานนฺทสฺส ปจฺจสฺโสสุํ.\csromanspacer{} อายสฺมา อานนฺโท เอตทโวจ\csromandash{}}

\prose{“ยํ โข โว อาวุโส ภควา สํขิตฺเตน อุทฺเทสํ อุทฺทิสิตฺวา วิตฺถาเรน อตฺถํ อวิภชิตฺวา อุฏฺฐายาสนา วิหารํ ปวิฏฺโฐ, ‘ตสฺมาติห ภิกฺขเว เส อายตเน เวทิตพฺเพ.\csromanspacer{} ยตฺถ จกฺขุ จ นิรุชฺฌติ รูปสญฺญา จ นิรุชฺฌติ เส อายตเน เวทิตพฺเพ ฯเปฯ.\csromanspacer{} ยตฺถ มโน จ นิรุชฺฌติ ธมฺมสญฺญา จ นิรุชฺฌติ เส อายตเน เวทิตพฺเพ’ติ.\csromanspacer{} อิมสฺส ขฺวาหํ อาวุโส ภควตา สํขิตฺเตน อุทฺเทสสฺส อุทฺทิฏฺฐสฺส วิตฺถาเรน อตฺถํ อวิภตฺตสฺส วิตฺถาเรน อตฺถํ อาชานามิ, สฬายตนนิโรธํ โน เอตํ อาวุโส ภควตา สนฺธาย ภาสิตํ ‘ตสฺมาติห ภิกฺขเว เส อายตเน เวทิตพฺเพ.\csromanspacer{} ยตฺถ จกฺขุ จ นิรุชฺฌติ รูปสญฺญา จ นิรุชฺฌติ เส อายตเน เวทิตพฺเพ ฯเปฯ.\csromanspacer{} ยตฺถ มโน จ นิรุชฺฌติ ธมฺมสญฺญา จ นิรุชฺฌติ เส อายตเน เวทิตพฺเพ’ติ.\csromanspacer{} อยํ โข อาวุโส ภควา สํขิตฺเตน อุทฺเทสํ อุทฺทิสิตฺวา วิตฺถาเรน อตฺถํ อวิภชิตฺวา อุฏฺฐายาสนา วิหารํ ปวิฏฺโฐ, ‘ตสฺมาติห ภิกฺขเว เส อายตเน เวทิตพฺเพ.\csromanspacer{} ยตฺถ จกฺขุ จ นิรุชฺฌติ รูปสญฺญา จ นิรุชฺฌติ เส อายตเน เวทิตพฺเพ ฯเปฯ.\csromanspacer{} ยตฺถ มโน จ นิรุชฺฌติ ธมฺมสญฺญา จ นิรุชฺฌติ เส อายตเน เวทิตพฺเพ’ติ.\csromanspacer{} อิมสฺส ขฺวาหํ อาวุโส ภควตา สํขิตฺเตน อุทฺเทสสฺส อุทฺทิฏฺฐสฺส วิตฺถาเรน อตฺถํ อวิภตฺตสฺส เอวํ วิตฺถาเรน อตฺถํ อาชานามิ.\csromanspacer{} อากงฺขมานา จ ปน ตุเมฺห อายสฺมนฺโต ภควนฺตํเยว อุปสงฺกมถ, อุปสงฺกมิตฺวา เอตมตฺถํ ปุจฺเฉยฺยาถ.\csromanspacer{} ยถา โว ภควา พฺยากโรติ, ตถา นํ ธาเรยฺยาถา”ติ.}

\gambhira{สฬายตนวคฺคสํยุตฺตปาฬิ}
\prose{\csromanfolio{320}“เอวมาวุโส”ติ โข เต ภิกฺขู อายสฺมโต อานนฺทสฺส ปฏิสฺสุตฺวา อุฏฺฐายาสนา เยน ภควา เตนุปสงฺกมิํสุ, อุปสงฺกมิตฺวา ภควนฺตํ อภิวาเทตฺวา เอกมนฺตํ นิสีทิํสุ, เอกมนฺตํ นิสินฺนา โข เต ภิกฺขู ภควนฺตํ เอตทโวจุํ\csromandash{}}

\prose{ยํ โข โน ภนฺเต ภควา สํขิตฺเตน อุทฺเทสํ อุทฺทิสิตฺวา วิตฺถาเรน อตฺถํ อวิภชิตฺวา อุฏฺฐายาสนา วิหารํ ปวิฏฺโฐ, “ตสฺมาติห ภิกฺขเว เส อายตเน เวทิตพฺเพ.\csromanspacer{} ยตฺถ จกฺขุ จ นิรุชฺฌติ รูปสญฺญา จ นิรุชฺฌติ เส อายตเน เวทิตพฺเพ ฯเปฯ.\csromanspacer{} ยตฺถ ชิวฺหา จ นิรุชฺฌติ รสสญฺญา จ นิรุชฺฌติ เส อายตเน เวทิตพฺเพ ฯเปฯ.\csromanspacer{} ยตฺถ มโน จ นิรุชฺฌติ ธมฺมสญฺญา จ นิรุชฺฌติ เส อายตเน เวทิตพฺเพ”ติ.\csromanspacer{} เตสํ โน ภนฺเต อมฺหากํ อจิรปกฺกนฺตสฺส ภควโต เอตทโหสิ “อิทํ โข โน อาวุโส ภควา สํขิตฺเตน อุทฺเทสํ อุทฺทิสิตฺวา วิตฺถาเรน อตฺถํ อวิภชิตฺวา อุฏฺฐายาสนา วิหารํ ปวิฏฺโฐ, ‘ตสฺมาติห ภิกฺขเว เส อายตเน เวทิตพฺเพ.\csromanspacer{} ยตฺถ จกฺขุ จ นิรุชฺฌติ รูปสญฺญา จ นิรุชฺฌติ เส อายตเน เวทิตพฺเพ ฯเปฯ.\csromanspacer{} ยตฺถ มโน จ นิรุชฺฌติ ธมฺมสญฺญา จ นิรุชฺฌติ เส อายตเน เวทิตพฺเพ’ติ.\csromanspacer{} โก นุ โข อิมสฺส ภควตา สํขิตฺเตน อุทฺเทสสฺส อุทฺทิฏฺฐสฺส วิตฺถาเรน อตฺถํ อวิภตฺตสฺส วิตฺถาเรน อตฺถํ วิภเชยฺยา”ติ.\csromanspacer{} เตสํ โน ภนฺเต อมฺหากํ เอตทโหสิ “อยํ โข อายสฺมา อานนฺโท สตฺถุ เจว สํวณฺณิโต, สมฺภาวิโต จ วิญฺญูนํ สพฺรหฺมจารีนํ, ปโหติ จายสฺมา อานนฺโท อิมสฺส ภควตา สํขิตฺเตน อุทฺเทสสฺส อุทฺทิฏฺฐสฺส วิตฺถาเรน อตฺถํ อวิภตฺตสฺส วิตฺถาเรน อตฺถํ วิภชิตุํ, ยํนูน มยํ เยนายสฺมา อานนฺโท เตนุปสงฺกเมยฺยาม, อุปสงฺกมิตฺวา อายสฺมนฺตํ อานนฺทํ เอตมตฺถํ ปฏิปุจฺเฉยฺยามา”ติ.\csromanspacer{} อถ โข มยํ ภนฺเต เยนายสฺมา อานนฺโท เตนุปสงฺกมิมฺห, อุปสงฺกมิตฺวา อายสฺมนฺตํ อานนฺทํ เอตมตฺถํ ปฏิปุจฺฉิมฺห, เตสํ โน ภนฺเต อายสฺมตา อานนฺเทน อิเมหิ อากาเรหิ อิเมหิ ปเทหิ อิเมหิ พฺยญฺชเนหิ อตฺโถ วิภตฺโตติ.}

\prose{ปณฺฑิโต ภิกฺขเว อานนฺโท, มหาปญฺโญ ภิกฺขเว อานนฺโท.\csromanspacer{} มํ เจปิ ตุเมฺห ภิกฺขเว เอตมตฺถํ ปฏิปุจฺเฉยฺยาถ, อหมฺปิ หํ เอวเมวํ พฺยากเรยฺยํ.\csromanspacer{} ยถา ตํ อานนฺเทน พฺยากตํ, เอโส เจเวตสฺส อตฺโถ, เอวญฺจ นํ ธาเรยฺยาถาติ.\csromanspacer{}\csromanspacer{}จตุตฺถํ.}

\csromanmatikaanchor{mtk.156}
\csromantocmarkref{subsection}{5. สกฺกปญฺหสุตฺต}{mtk.156}
\bookmark[dest={mtk.156},level=2]{5. สกฺกปญฺหสุตฺต}
\csromanheader{2}{1. สฬายตนสํยุตฺต \textbf{5. สกฺกปญฺหสุตฺต}}
\proseitem{118}{\csromanfolio{321}เอกํ สมยํ ภควา ราชคเห วิหรติ คิชฺฌกูเฏ ปพฺพเต.\csromanspacer{} อถ โข สกฺโก เทวานมินฺโท เยน ภควา เตนุปสงฺกมิ, อุปสงฺกมิตฺวา ภควนฺตํ อภิวาเทตฺวา เอกมนฺตํ อฏฺฐาสิ, เอกมนฺตํ ฐิโต โข สกฺโก เทวานมินฺโท ภควนฺตํ เอตทโวจ “โก นุ โข ภนฺเต เหตุ โก ปจฺจโย, เยน มิเธกจฺเจ สตฺตา ทิฏฺเฐว ธมฺเม โน ปรินิพฺพายนฺติ.\csromanspacer{} โก ปน ภนฺเต เหตุ โก ปจฺจโย, เยน มิเธกจฺเจ สตฺตา ทิฏฺเฐว ธมฺเม ปรินิพฺพายนฺตี”ติ.\csromanspacer{} สนฺติ โข เทวานมินฺท จกฺขุวิญฺเญยฺยา รูปา อิฏฺฐา กนฺตา มนาปา ปิยรูปา กามูปสํหิตา รชนียา, ตญฺเจ ภิกฺขุ อภินนฺทติ อภิวทติ อชฺโฌสาย ติฏฺฐติ, ตสฺส ตํ อภินนฺทโต อภิวทโต อชฺโฌสาย ติฏฺฐโต ตนฺนิสฺสิตํ วิญฺญาณํ โหติ ตทุปาทานํ, ส-อุปาทาโน เทวานมินฺท ภิกฺขุ โน ปรินิพฺพายติ ฯเปฯ.}

\prose{สนฺติ โข เทวานมินฺท ชิวฺหาวิญฺเญยฺยา รสา ฯเปฯ.\csromanspacer{} สนฺติ โข เทวานมินฺท มโนวิญฺเญยฺยา ธมฺมา อิฏฺฐา กนฺตา มนาปา ปิยรูปา กามูปสํหิตา รชนียา, ตญฺเจ ภิกฺขุ อภินนฺทติ อภิวทติ อชฺโฌสาย ติฏฺฐติ, ตสฺส ตํ อภินนฺทโต อภิวทโต อชฺโฌสาย ติฏฺฐโต ตนฺนิสฺสิตํ วิญฺญาณํ โหติ ตทุปาทานํ, ส-อุปาทาโน เทวานมินฺท ภิกฺขุ โน ปรินิพฺพายติ.\csromanspacer{} อยํ โข เทวานมินฺท เหตุ อยํ ปจฺจโย, เยน มิเธกจฺเจ สตฺตา ทิฏฺเฐว ธมฺเม โน ปรินิพฺพายนฺติ.}

\prose{สนฺติ จ โข เทวานมินฺท จกฺขุวิญฺเญยฺยา รูปา อิฏฺฐา กนฺตา มนาปา ปิยรูปา กามูปสํหิตา รชนียา, ตญฺเจ ภิกฺขุ นาภินนฺทติ นาภิวทติ นาชฺโฌสาย ติฏฺฐติ, ตสฺส ตํ อนภินนฺทโต อนภิวทโต อนชฺโฌสาย ติฏฺฐโต น ตนฺนิสฺสิตํ วิญฺญาณํ โหติ น ตทุปาทานํ, อนุปาทาโน เทวามินฺท ภิกฺขุ ปรินิพฺพายติ ฯเปฯ.}

\prose{สนฺติ โข เทวานมินฺท ชิวฺหาวิญฺเญยฺยา รสา ฯเปฯ.\csromanspacer{} สนฺติ โข เทวานมินฺท มโนวิญฺเญยฺยา ธมฺมา อิฏฺฐา กนฺตา มนาปา ปิยรูปา กามูปสํหิตา รชนียา, ตญฺเจ ภิกฺขุ นาภินนฺทติ นาภิวทติ นาชฺโฌสาย ติฏฺฐติ, ตสฺส ตํ อนภินนฺทโต อนภิวทโต อนชฺโฌสาย ติฏฺฐโต น ตนฺนิสฺสิตํ วิญฺญาณํ โหติ น ตทุปาทานํ, อนุปาทาโน เทวานมินฺท}

\gambhira{สฬายตนวคฺคสํยุตฺตปาฬิ}
\prose{\csromanfolio{322}ภิกฺขุ ปรินิพฺพายติ.\csromanspacer{} อยํ โข เทวานมินฺท เหตุ อยํ ปจฺจโย, เยน มิเธกจฺเจ สตฺตา ทิฏฺเฐว ธมฺเม ปรินิพฺพายนฺตีติ.\csromanspacer{}\csromanspacer{}ปญฺจมํ.}

\csromanmatikaanchor{mtk.157}
\csromantocmarkref{subsection}{6. ปญฺจสิขสุตฺต}{mtk.157}
\bookmark[dest={mtk.157},level=2]{6. ปญฺจสิขสุตฺต}
\csromanheader{2}{6. ปญฺจสิขสุตฺต}
\proseitem{119}{เอกํ สมยํ ภควา ราชคเห วิหรติ คิชฺฌกูเฏ ปพฺพเต.\csromanspacer{} อถ โข ปญฺจสิโข คนฺธพฺพเทวปุตฺโต เยน ภควา เตนุปสงฺกมิ, อุปสงฺกมิตฺวา ภควนฺตํ อภิวาเทตฺวา เอกมนฺตํ อฏฺฐาสิ, เอกมนฺตํ ฐิโต โข ปญฺจสิโข คนฺธพฺพเทวปุตฺโต ภควนฺตํ เอตทโวจ “โก นุ โข ภนฺเต เหตุ โก ปจฺจโย, เยน มิเธกจฺเจ สตฺตา ทิฏฺเฐว ธมฺเม โน ปรินิพฺพายนฺติ.\csromanspacer{} โก ปน ภนฺเต เหตุ โก ปจฺจโย, เยน มิเธกจฺเจ สตฺตา ทิฏฺเฐว ธมฺเม ปรินิพฺพายนฺตี”ติ.\csromanspacer{} สนฺติ โข ปญฺจสิข จกฺขุวิญฺเญยฺยา รูปา ฯเปฯ.\csromanspacer{} สนฺติ โข ปญฺจสิข มโนวิญฺเญยฺยา ธมฺมา อิฏฺฐา กนฺตา มนาปา ปิยรูปา กามูปสํหิตา รชนียา, ตญฺเจ ภิกฺขุ อภินนฺทติ อภิวทติ อชฺโฌสาย ติฏฺฐติ, ตสฺส ตํ อภินนฺทโต อภิวทโต อชฺโฌสาย ติฏฺฐโต ตนฺนิสฺสิตํ วิญฺญาณํ โหติ ตทุปาทานํ, ส-อุปาทาโน ปญฺจสิข ภิกฺขุ โน ปรินิพฺพายติ.\csromanspacer{} อยํ โข ปญฺจสิข เหตุ อยํ ปจฺจโย, เยน มิเธกจฺเจ สตฺตา ทิฏฺเฐว ธมฺเม โน ปรินิพฺพายนฺติ.}

\prose{สนฺติ จ โข ปญฺจสิข จกฺขุวิญฺเญยฺยา รูปา อิฏฺฐา กนฺตา มนาปา ฯเปฯ.\csromanspacer{} สนฺติ โข ปญฺจสิข มโนวิญฺเญยฺยา ธมฺมา อิฏฺฐา กนฺตา มนาปา ปิยรูปา กามูปสํหิตา รชนียา, ตญฺเจ ภิกฺขุ นาภินนฺทติ นาภิวทติ นาชฺโฌสาย ติฏฺฐติ, ตสฺส ตํ อนภินนฺทโต อนภิวทโต อนชฺโฌสาย ติฏฺฐโต น ตนฺนิสฺสิตํ วิญฺญาณํ โหติ น ตทุปาทานํ, อนุปาทาโน ปญฺจสิข ภิกฺขุ ปรินิพฺพายติ.\csromanspacer{} อยํ โข ปญฺจสิข เหตุ อยํ ปจฺจโย, เยน มิเธกจฺเจ สตฺตา ทิฏฺเฐว ธมฺเม ปรินิพฺพายนฺตีติ.\csromanspacer{}\csromanspacer{}ฉฏฺฐํ.}

\csromanmatikaanchor{mtk.158}
\csromantocmarkref{subsection}{7. สาริปุตฺตสทฺธิวิหาริกสุตฺต}{mtk.158}
\bookmark[dest={mtk.158},level=2]{7. สาริปุตฺตสทฺธิวิหาริกสุตฺต}
\csromanheader{2}{7. สาริปุตฺตสทฺธิวิหาริกสุตฺต}
\proseitem{120}{เอกํ สมยํ อายสฺมา สาริปุตฺโต สาวตฺถิยํ วิหรติ เชตวเน อนาถปิณฺฑิกสฺส อาราเม.\csromanspacer{} อถ โข อญฺญตโร ภิกฺขุ เยนายสฺมา สาริปุตฺโต เตนุปสงฺกมิ, อุปสงฺกมิตฺวา อายสฺมตา สาริปุตฺเตน สทฺธิํ สมฺโมทิ, สมฺโมทนียํ กถํ สารณียํ วีติสาเรตฺวา เอกมนฺตํ นิสีทิ, เอกมนฺตํ นิสินฺโน โข โส ภิกฺขุ อายสฺมนฺตํ สาริปุตฺตํ เอตทโวจ}

\vspace{\tipitakaparskip}
\csromancenter{สฬายตนสํยุตฺต “สทฺธิวิหาริโก อาวุโส สาริปุตฺต ภิกฺขุ สิกฺขํ ปจฺจกฺขาย หีนายาวตฺโต”ติ.\csromanspacer{} เอวเมตํ อาวุโส โหติ อินฺทฺริเยสุ อคุตฺตทฺวารสฺส โภชเน อมตฺตญฺญุโน ชาคริยํ อนนุยุตฺตสฺส, โส วตาวุโส ภิกฺขุ อินฺทฺริเยสุ อคุตฺตทฺวาโร โภชเน อมตฺตญฺญู ชาคริยํ อนนุยุตฺโต ยาวชีวํ ปริปุณฺณํ ปริสุทฺธํ พฺรหฺมจริยํ สนฺตาเนสฺสตีติ เนตํ ฐานํ วิชฺชติ.\csromanspacer{} โส วตาวุโส ภิกฺขุ อินฺทฺริเยสุ คุตฺตทฺวาโร โภชเน มตฺตญฺญู ชาคริยํ อนุยุตฺโต ยาวชีวํ ปริปุณฺณํ ปริสุทฺธํ พฺรหฺมจริยํ สนฺตาเนสฺสตีติ ฐานเมตํ วิชฺชติ.}
\prose{\csromanfolio{323}กถญฺจาวุโส อินฺทฺริเยสุ คุตฺตทฺวาโร โหติ, อิธาวุโส ภิกฺขุ จกฺขุนา รูปํ ทิสฺวา น นิมิตฺตคฺคาหี โหติ นานุพฺยญฺจนคฺคาหี, ยตฺวาธิกรณเมนํ จกฺขุนฺทฺริยํ อสํวุตํ วิหรนฺตํ อภิชฺฌาโทมนสฺสา ปาปกา อกุสลา ธมฺมา อนฺวาสฺสเวยฺยุํ, ตสฺส สํวราย ปฏิปชฺชติ, รกฺขติ จกฺขุนฺทฺริยํ, จกฺขุนฺทฺริเย สํวรํ อาปชฺชติ.\csromanspacer{} โสเตน สทฺทํ สุตฺวา.\csromanspacer{} ฆาเนน คนฺธํ ฆายิตฺวา.\csromanspacer{} ชิวฺหาย รสํ สายิตฺวา.\csromanspacer{} กาเยน โผฏฺฐพฺพํ ผุสิตฺวา.\csromanspacer{} มนสา ธมฺมํ วิญฺญาย นิมิตฺตคฺคาหี โหติ นานุพฺยญฺชนคฺคาหี, ยตฺวาธิกรณเมนํ มนินฺทฺริยํ อสํวุตํ วิหรนฺตํ อภิชฺฌาโทมนสฺสา ปาปกา อกุสลา ธมฺมา อนฺวาสฺสเวยฺยุํ, ตสฺส สํวราย ปฏิปชฺชติ, รกฺขติ มนินฺทฺริยํ, มนินฺทฺริเย สํวรํ อาปชฺชติ.\csromanspacer{} เอวํ โข อาวุโส อินฺทฺริเยสุ คุตฺตทฺวาโร โหติ.}

\prose{กถญฺจาวุโส โภชเน มตฺตญฺญู โหติ, อิธาวุโส ภิกฺขุ ปฏิสงฺขา โยนิโส อาหารํ อาหาเรติ เนว ทวาย น มทาย น มณฺฑนาย น วิภูสนาย, ยาวเทว อิมสฺส กายสฺส ฐิติยา ยาปนาย วิหิํสูปรติยา พฺรหฺมจริยานุคฺคหาย อิติ ปุราณญฺจ เวทนํ ปฏิหงฺขามิ นวญฺจ เวทนํ น อุปฺปาเทสฺสามิ, ยาตฺรา จ เม ภวิสฺสติ อนวชฺชตา จ ผาสุวิหาโร จาติ.\csromanspacer{} เอวํ โข อาวุโส โภชเน มตฺตญฺญู โหติ.}

\prose{กถญฺจาวุโส ชาคริยํ อนุยุตฺโต โหติ, อิธาวุโส ภิกฺขุ ทิวสํ จงฺกเมน นิสชฺชาย อาวรณีเยหิ ธมฺเมหิ จิตฺตํ ปริโสเธติ, รตฺติยา ปฐมํ ยามํ จงฺกเมน นิสชฺชาย อาวรณีเยหิ ธมฺเมหิ จิตฺตํ ปริโสเธติ, รตฺติยา มชฺฌิมํ ยามํ ทกฺขิเณน ปสฺเสน สีหเสยฺยํ กปฺเปติ ปาเท ปาทํ อจฺจาธาย สโต สมฺปชาโน อุฏฺฐานสญฺญํ มนสิ กริตฺวา รตฺติยา ปจฺฉิมํ ยามํ ปจฺจุฏฺฐาย จงฺกเมน นิสชฺชาย อาวรณีเยหิ ธมฺเมหิ จิตฺตํ ปริโสเธติ.\csromanspacer{} เอวํ โข อาวุโส}

\gambhira{สฬายตนวคฺคสํยุตฺตปาฬิ}
\prose{\csromanfolio{324}ชาคริยํ อนุยุตฺโต โหติ.\csromanspacer{} ตสฺมาติหาวุโส เอวํ สิกฺขิตพฺพํ “อินฺทฺริเยสุ คุตฺตทฺวารา ภวิสฺสาม, โภชเน มตฺตญฺญุโน ชาคริยํ อนุยุตฺตา”ติ.\csromanspacer{} เอวํ หิ โว อาวุโส สิกฺขิตพฺพนฺติ.\csromanspacer{}\csromanspacer{}สตฺตมํ.}

\csromanmatikaanchor{mtk.159}
\csromantocmarkref{subsection}{8. ราหุโลวาทสุตฺต}{mtk.159}
\bookmark[dest={mtk.159},level=2]{8. ราหุโลวาทสุตฺต}
\csromanheader{2}{8. ราหุโลวาทสุตฺต}
\proseitem{121}{เอกํ สมยํ ภควา สาวตฺถิยํ วิหรติ เชตวเน อนาถปิณฺฑิกสฺส อาราเม.\csromanspacer{} อถ โข ภควโต รโหคตสฺส ปฏิสลฺลีนสฺส เอวํ เจตโส ปริวิตกฺโก อุทปาทิ “ปริปกฺกา โข ราหุลสฺส วิมุตฺติปริปาจนิยา ธมฺมา, ยํนูนาหํ ราหุลํ อุตฺตริํ อาสวานํ ขเย วิเนยฺยนฺ”ติ.\csromanspacer{} อถ โข ภควา ปุพฺพณฺหสมยํ นิวาเสตฺวา ปตฺตจีวรมาทาย สาวตฺถิยํ ปิณฺฑาย จริตฺวา ปจฺฉาภตฺตํ ปิณฺฑปาตปฏิกฺกนฺโต อายสฺมนฺตํ ราหุลํ อามนฺเตสิ “คณฺหาหิ ราหุล นิสีทนํ, เยน อนฺธวนํ เตนุปสงฺกมิสฺสาม ทิวาวิหารายา”ติ. “เอวํ ภนฺเต”ติ โข อายสฺมา ราหุโล ภควโต ปฏิสฺสุตฺวา นิสีทนํ อาทาย ภควนฺตํ ปิฏฺฐิโต ปิฏฺฐิโต อนุพนฺธิ.}

\prose{เตน โข ปน สมเยน อเนกานิ เทวตาสหสฺสานิ ภควนฺตํ อนุพนฺธานิ โหนฺติ “อชฺช ภควา อายสฺมนฺตํ ราหุลํ อุตฺตริํ อาสวานํ ขเย วิเนสฺสตี”ติ.\csromanspacer{} อถ โข ภควา อนฺธวนํ อชฺโฌคาเหตฺวา อญฺญตรสฺมิํ รุกฺขมูเล ปญฺญตฺเต อาสเน นิสีทิ.\csromanspacer{} อายสฺมาปิ โข ราหุโล ภควนฺตํ อภิวาเทตฺวา เอกมนฺตํ นิสีทิ, เอกมนฺตํ นิสินฺนํ โข อายสฺมนฺตํ ราหุลํ ภควา เอตทโวจ\csromandash{}}

\prose{“ตํ กิํ มญฺญสิ ราหุล, จกฺขุ นิจฺจํ วา อนิจฺจํ วา”ติ.\csromanspacer{} อนิจฺจํ ภนฺเต.\csromanspacer{} ยํ ปนานิจฺจํ, ทุกฺขํ วา ตํ สุขํ วาติ.\csromanspacer{} ทุกฺขํ ภนฺเต.\csromanspacer{} ยํ ปนานิจฺจํ ทุกฺขํ วิปริณามธมฺมํ, กลฺลํ นุ ตํ สมนุปสฺสิตุํ “เอตํ มม, เอโสหมสฺมิ, เอโส เม อตฺตา”ติ.\csromanspacer{} โน เหตํ ภนฺเต. ( )\footnote{(ตํ กิํ มญฺญสิ) เอวมิตเรสุปิ. (ม 3. 325)} รูปา นิจฺจา วา อนิจฺจา วาติ.\csromanspacer{} อนิจฺจา ภนฺเต ฯเปฯ.\csromanspacer{} จกฺขุวิญฺญาณํ นิจฺจํ วา อนิจฺจํ วาติ.\csromanspacer{} อนิจฺจํ ภนฺเต.\csromanspacer{} จกฺขุสมฺผสฺโส นิจฺโจ วา อนิจฺโจ วาติ.\csromanspacer{} อนิจฺโจ ภนฺเต.\csromanspacer{} ยมฺปิทํ จกฺขุสมฺผสฺสปจฺจยา อุปฺปชฺชติ เวทนาคตํ สญฺญาคตํ สงฺขารคตํ วิญฺญาณคตํ, ตมฺปิ นิจฺจํ วา อนิจฺจํ วาติ.\csromanspacer{} อนิจฺจํ ภนฺเต.\csromanspacer{} ยํ ปนานิจฺจํ, ทุกฺขํ วา ตํ สุขํ วาติ.\csromanspacer{} ทุกฺขํ ภนฺเต.\csromanspacer{} ยํ ปนานิจฺจํ ทุกฺขํ วิปริณามธมฺมํ, กลฺลํ นุ ตํ สมนุปสฺสิตุํ “เอตํ มม,}

\vspace{\tipitakaparskip}
\csromancenter{สฬายตนสํยุตฺต เอโสหมสฺมิ, เอโส เม อตฺตา”ติ.\csromanspacer{} โน เหตํ ภนฺเต ฯเปฯ.\csromanspacer{} ชิวฺหา นิจฺจา วา อนิจฺจา วาติ.\csromanspacer{} อนิจฺจา ภนฺเต ฯเปฯ.\csromanspacer{} ชิวฺหาวิญฺญาณํ นิจฺจํ วา อนิจฺจํ วาติ.\csromanspacer{} อนิจฺจํ ภนฺเต.\csromanspacer{} ชิวฺหาสมฺผสฺโส นิจฺโจ วา อนิจฺโจ วาติ.\csromanspacer{} อนิจฺโจ ภนฺเต.\csromanspacer{} ยมฺปิทํ ชิวฺหาสมฺผสฺสปจฺจยา อุปฺปชฺชติ เวทนาคตํ สญฺญาคตํ สงฺขารคตํ วิญฺญาณคตํ, ตมฺปิ นิจฺจํ วา อนิจฺจํ วาติ.\csromanspacer{} อนิจฺจํ ภนฺเต.\csromanspacer{} ยํ ปนานิจฺจํ, ทุกฺขํ วา ตํ สุขํ วาติ.\csromanspacer{} ทุกฺขํ ภนฺเต.\csromanspacer{} ยํ ปนานิจฺจํ ทุกฺขํ วิปริณามธมฺมํ, กลฺลํ นุ ตํ สมนุปสฺสิตุํ “เอตํ มม, เอโสหมสฺมิ, เอโส เม อตฺตา”ติ.\csromanspacer{} โน เหตํ ภนฺเต ฯเปฯ.\csromanspacer{} มโน นิจฺโจ วา อนิจฺโจ วาติ.\csromanspacer{} อนิจฺโจ ภนฺเต.\csromanspacer{} ยํ ปนานิจฺจํ, ทุกฺขํ วา ตํ สุขํ วาติ.\csromanspacer{} ทุกฺขํ ภนฺเต.\csromanspacer{} ยํ ปนานิจฺจํ ทุกฺขํ วิปริณามธมฺมํ, กลฺลํ นุ ตํ สมนุปสฺสิตุํ “เอตํ มม, เอโสหมสฺมิ, เอโส เม อตฺตา”ติ.\csromanspacer{} โน เหตํ ภนฺเต.\csromanspacer{} ธมฺมา นิจฺจา วา อนิจฺจา วาติ.\csromanspacer{} อนิจฺจา ภนฺเต ฯเปฯ.\csromanspacer{} มโนวิญฺญาณํ นิจฺจํ วา อนิจฺจํ วาติ.\csromanspacer{} อนิจฺจํ ภนฺเต ฯเปฯ.\csromanspacer{} มโนสมฺผสฺโส นิจฺโจ วา อนิจฺโจ วาติ.\csromanspacer{} อนิจฺโจ ภนฺเต.\csromanspacer{} ยมฺปิทํ มโนสมฺผสฺสปจฺจยา อุปฺปชฺชติ เวทนาคตํ สญฺญาคตํ สงฺขารคตํ วิญฺญาณคตํ, ตมฺปิ นิจฺจํ วา อนิจฺจํ วาติ.\csromanspacer{} อนิจฺจํ ภนฺเต.\csromanspacer{} ยํ ปนานิจฺจํ, ทุกฺขํ วา ตํ สุขํ วาติ.\csromanspacer{} ทุกฺขํ ภนฺเต.\csromanspacer{} ยํ ปนานิจฺจํ ทุกฺขํ วิปริณามธมฺมํ, กลฺลํ นุ ตํ สมนุปสฺสิตุํ “เอตํ มม, เอโสหมสฺมิ, เอโส เม อตฺตา”ติ.\csromanspacer{} โน เหตํ ภนฺเต.}
\prose{\csromanfolio{325}เอวํ ปสฺสํ ราหุล สุตวา อริยสาวโก จกฺขุสฺมิมฺปิ นิพฺพินฺทติ, รูเปสุปิ นิพฺพินฺทติ, จกฺขุวิญฺญาเณปิ นิพฺพินฺทติ, จกฺขุสมฺผสฺเสปิ นิพฺพินฺทติ, ยมฺปิทํ จกฺขุสมฺผสฺสปจฺจยา อุปฺปชฺชติ เวทนาคตํ สญฺญาคตํ สงฺขารคตํ วิญฺญาณคตํ, ตสฺมิมฺปิ นิพฺพินฺทติ ฯเปฯ ชิวฺหายปิ นิพฺพินฺทติ, รเสสุปิ นิพฺพินฺทติ, ชิวฺหาวิญฺญาเณปิ นิพฺพินฺทติ, ชิวฺหาสมฺผสฺเสปิ นิพฺพินฺทติ, ยมฺปิทํ ชิวฺหาสมฺผสฺสปจฺจยา อุปฺปชฺชติ เวทนาคตํ สญฺญาคตํ สงฺขารคตํ วิญฺญาณคตํ, ตสฺมิมฺปิ นิพฺพินฺทติ ฯเปฯ มนสฺมิมฺปิ นิพฺพินฺทติ, ธมฺเมสุปิ นิพฺพินฺทติ, มโนวิญฺญาเณปิ นิพฺพินฺทติ, มโนสมฺผสฺเสปิ นิพฺพินฺทติ, ยมฺปิทํ มโนสมฺผสฺสปจฺจยา อุปฺปชฺชติ เวทนาคตํ สญฺญาคตํ สงฺขารคตํ วิญฺญาณคตํ, ตสฺมิมฺปิ นิพฺพินฺทติ, นิพฺพินฺทํ วิรชฺชติ, วิราคา วิมุจฺจติ, วิมุตฺตสฺมิํ “วิมุตฺตมฺ”อิติ ญาณํ โหติ, “ขีณา ชาติ, วุสิตํ พฺรหฺมจริยํ, กตํ กรณียํ, นาปรํ อิตฺถตฺตายา”ติ ปชานาตีติ.}

\prose{อิทมโวจ ภควา, อตฺตมโน อายสฺมา ราหุโล ภควโต ภาสิตํ อภินนฺทิ.\csromanspacer{} อิมสฺมิํ จ ปน เวยฺยากรณสฺมิํ ภญฺญมาเน อายสฺมโต ราหุลสฺส อนุปาทาย อาสเวหิ จิตฺตํ วิมุจฺจิ,}

\gambhira{สฬายตนวคฺคสํยุตฺตปาฬิ}
\csromanmatikaanchor{mtk.160}
\csromanmatikaanchor{mtk.161}
\csromantocmarkref{subsection}{9. สํโยชนิยธมฺมสุตฺต}{mtk.160}
\bookmark[dest={mtk.160},level=2]{9. สํโยชนิยธมฺมสุตฺต}
\csromantocmarkref{subsection}{10. อุปาทานิยสุตฺต}{mtk.161}
\bookmark[dest={mtk.161},level=2]{10. อุปาทานิยสุตฺต}
\prose{\csromanfolio{326}อเนกานญฺจ เทวตาสหสฺสานํ วิรชํ วีตมลํ ธมฺมจกฺขุํ อุทปาทิ “ยํ กิญฺจิ สมุทยธมฺมํ, สพฺพํ ตํ นิโรธธมฺมนฺ”ติ.\csromanspacer{}\csromanspacer{}อฏฺฐมํ.}

\csromanheader{2}{9. สํโยชนิยธมฺมสุตฺต}
\proseitem{122}{สํโยชนิเย จ ภิกฺขเว ธมฺเม เทเสสฺสามิ สํโยชนญฺจ, ตํ สุณาถ.\csromanspacer{} กตเม จ ภิกฺขเว สํโยชนิยา ธมฺมา, กตมํ จ สํโยชนํ.\csromanspacer{} สนฺติ ภิกฺขเว จกฺขุวิญฺเญยฺยา รูปา อิฏฺฐา กนฺตา มนาปา ปิยรูปา กามูปสํหิตา รชนียา, อิเม วุจฺจนฺติ ภิกฺขเว สํโยชนิยา ธมฺมา.\csromanspacer{} โย ตตฺถ ฉนฺทราโค, ตํ ตตฺถ สํโยชนํ ฯเปฯ.\csromanspacer{} สนฺติ ภิกฺขเว ชิวฺหาวิญฺเญยฺยา รสา ฯเปฯ.\csromanspacer{} สนฺติ ภิกฺขเว มโนวิญฺเญยฺยา ธมฺมา อิฏฺฐา กนฺตา มนาปา ปิยรูปา กามูปสํหิตา รชนียา, อิเม วุจฺจนฺติ ภิกฺขเว สํโยชนิยา ธมฺมา.\csromanspacer{} โย ตตฺถ ฉนฺทราโค, ตํ ตตฺถ สํโยชนนฺติ.\csromanspacer{}\csromanspacer{}นวมํ.}

\csromanheader{2}{10. อุปาทานิยธมฺมสุตฺต}
\proseitem{123}{อุปาทานิเย จ ภิกฺขเว ธมฺเม เทเสสฺสามิ อุปาทานญฺจ, ตํ สุณาถ.\csromanspacer{} กตเม จ ภิกฺขเว อุปาทานิยา ธมฺมา, กตมํ จ อุปาทานํ.\csromanspacer{} สนฺติ ภิกฺขเว จกฺขุวิญฺเญยฺยา รูปา อิฏฺฐา กนฺตา มนาปา ปิยรูปา กามูปสํหิตา รชนียา, อิเม วุจฺจนฺติ ภิกฺขเว อุปาทานิยา ธมฺมา.\csromanspacer{} โย ตตฺถ ฉนฺทราโค, ตํ ตตฺถ อุปาทานํ ฯเปฯ.\csromanspacer{} สนฺติ ภิกฺขเว ชิวฺหาวิญฺเญยฺยา รสา ฯเปฯ.\csromanspacer{} สนฺติ ภิกฺขเว มโนวิญฺเญยฺยา ธมฺมา อิฏฺฐา กนฺตา มนาปา ปิยรูปา กามูปสํหิตา รชนียา, อิเม วุจฺจนฺติ ภิกฺขเว อุปาทานิยา ธมฺมา.\csromanspacer{} โย ตตฺถ ฉนฺทราโค, ตํ ตตฺถ อุปาทานนฺติ.\csromanspacer{}\csromanspacer{}ทสมํ.\csromanspacer{} โลกกามคุณาวคฺโค ทุติโย.}

\titlehead{\textbf{ตสฺสุทฺทานํ}}
\csromangathameasure{มารปาเสน เทฺว วุตฺตา, โลกกามคุเณน จ. \\ สกฺโก ปญฺจสิโข เจว, สาริปุตฺโต จ ราหุโล. \\ สํโยชนํ อุปาทานํ, วคฺโค เตน ปวุจฺจตีติ.}
\csromangathagroup{\csromangathastanza{มารปาเสน เทฺว วุตฺตา, โลกกามคุเณน จ. \\ สกฺโก ปญฺจสิโข เจว, สาริปุตฺโต จ ราหุโล. \\ สํโยชนํ อุปาทานํ, วคฺโค เตน ปวุจฺจตีติ.}}

\csromanmatikaanchor{mtk.162}
\csromantocmarknopage{section}{(13) 3. คหปติวคฺค}{mtk.162}
\bookmark[dest={mtk.162},level=1]{(13) 3. คหปติวคฺค}
\csromanheader{1}{1. สฬายตนสํยุตฺต \textbf{(13) 3. คหปติวคฺค}}
\csromanheader{2}{1. เวสาลีสุตฺต}
\csromanmatikaanchor{mtk.163}
\csromantocmarkref{subsection}{1. เวสาลิสุตฺต}{mtk.163}
\bookmark[dest={mtk.163},level=2]{1. เวสาลิสุตฺต}
\proseitem{124}{\csromanfolio{327}เอกํ สมยํ ภควา เวสาลิยํ วิหรติ มหาวเน กูฏาคารสาลายํ.\csromanspacer{} อถ โข อุคฺโค คหปติ เวสาลิโก เยน ภควา เตนุปสงฺกมิ, อุปสงฺกมิตฺวา เอกมนฺตํ นิสีทิ, เอกมนฺตํ นิสินฺโน โข อุคฺโค คหปติ เวสาลิโก ภควนฺตํ เอตทโวจ “โก นุ โข ภนฺเต เหตุ โก ปจฺจโย, เยน มิเธกจฺเจ สตฺตา ทิฏฺเฐว ธมฺเม โน ปรินิพฺพายนฺติ.\csromanspacer{} โก ปน ภนฺเต เหตุ โก ปจฺจโย, เยน มิเธกจฺเจ สตฺตา ทิฏฺเฐว ธมฺเม ปรินิพฺพายนฺตี”ติ.}

\prose{สนฺติ โข คหปติ จกฺขุวิญฺเญยฺยา รูปา อิฏฺฐา กนฺตา มนาปา ปิยรูปา กามูปสํหิตา รชนียา, ตญฺเจ ภิกฺขุ อภินนฺทติ อภิวทติ อชฺโฌสาย ติฏฺฐติ, ตสฺส ตํ อภินนฺทโต อภิวทโต อชฺโฌสาย ติฏฺฐโต ตนฺนิสฺสิตํ วิญฺญาณํ โหติ ตทุปาทานํ, ส-อุปาทาโน คหปติ ภิกฺขุ โน ปรินิพฺพายติ ฯเปฯ.\csromanspacer{} สนฺติ โข คหปติ ชิวฺหาวิญฺเญยฺยา รสา ฯเปฯ.\csromanspacer{} สนฺติ โข คหปติ มโนวิญฺเญยฺยา ธมฺมา อิฏฺฐา กนฺตา มนาปา ปิยรูปา กามูปสํหิตา รชนียา, ตญฺเจ ภิกฺขุ อภินนฺทติ อภิวทติ อชฺโฌสาย ติฏฺฐติ, ตสฺส ตํ อภินนฺทโต อภิวทโต อชฺโฌสาย ติฏฺฐโต ตนฺนิสฺสิตํ วิญฺญาณํ โหติ ตทุปาทานํ, ส-อุปาทาโน คหปติ ภิกฺขุ โน ปรินิพฺพายติ.\csromanspacer{} อยํ โข คหปติ เหตุ อยํ ปจฺจโย, เยน มิเธกจฺเจ สตฺตา ทิฏฺเฐว ธมฺเม โน ปรินิพฺพายนฺติ.}

\prose{สนฺติ จ โข คหปติ จกฺขุวิญฺเญยฺยา รูปา อิฏฺฐา กนฺตา มนาปา ปิยรูปา กามูปสํหิตา รชนียา, ตญฺเจ ภิกฺขุ นาภินนฺทติ นาภิวทติ นาชฺโฌสาย ติฏฺฐติ, ตสฺส ตํ อนภินนฺทโต อนภิวทโต อนชฺโฌสาย ติฏฺฐโต น ตนฺนิสฺสิตํ วิญฺญาณํ โหติ น ตทุปาทานํ, อนุปาทาโน คหปติ ภิกฺขุ ปรินิพฺพายติ ฯเปฯ.\csromanspacer{} สนฺติ โข คหปติ ชิวฺหาวิญฺเญยฺยา รสา ฯเปฯ.\csromanspacer{} สนฺติ โข คหปติ มโนวิญฺเญยฺยา ธมฺมา อิฏฺฐา กนฺตา มนาปา ปิยรูปา กามูปสํหิตา รชนียา, ตญฺเจ ภิกฺขุ นาภินนฺทติ นาภิวทติ นาชฺโฌสาย ติฏฺฐติ, ตสฺส ตํ อนภินนฺทโต อนภิวทโต อนชฺโฌสาย ติฏฺฐโต น ตนฺนิสฺสิตํ วิญฺญาณํ โหติ, น ตทุปาทานํ,}

\gambhira{สฬายตนวคฺคสํยุตฺตปาฬิ}
\prose{\csromanfolio{328}อนุปาทาโน คหปติ ภิกฺขุ ปรินิพฺพายติ.\csromanspacer{} อยํ โข คหปติ เหตุ อยํ ปจฺจโย, เยน มิเธกจฺเจ สตฺตา ทิฏฺเฐว ธมฺเม ปรินิพฺพายนฺตีติ.\csromanspacer{}\csromanspacer{}ปฐมํ}

\csromanmatikaanchor{mtk.164}
\csromantocmarkref{subsection}{2. วชฺชีสุตฺต}{mtk.164}
\bookmark[dest={mtk.164},level=2]{2. วชฺชีสุตฺต}
\csromanheader{2}{2. วชฺชีสุตฺต}
\proseitem{125}{เอกํ สมยํ ภควา วชฺชีสุ วิหรติ หตฺถิคาเม.\csromanspacer{} อถ โข อุคฺโค คหปติ หตฺถิคามโก เยน ภควา เตนุปสงฺกมิ, อุปสงฺกมิตฺวา เอกมนฺตํ นิสีทิ, เอกมนฺตํ นิสินฺโน โข อุคฺโค คหปติ หตฺถิคามโก ภควนฺตํ เอตทโวจ “โก นุ โข ภนฺเต เหตุ โก ปจฺจโย, เยน มิเธกจฺเจ สตฺตา ทิฏฺเฐว ธมฺเม โน ปรินิพฺพายนฺติ.\csromanspacer{} โก ปน ภนฺเต เหตุ โก ปจฺจโย, เยน มิเธกจฺเจ สตฺตา ทิฏฺเฐว ธมฺเม ปรินิพฺพายนฺตี”ติ. (ยถา ปุริมสุตฺตนฺตํ, เอวํ วิตฺตาเรตพฺพํ.) อยํ โข คหปติ เหตุ อยํ ปจฺจโย, เยน มิเธกจฺเจ สตฺตา ทิฏฺเฐว ธมฺเม ปรินิพฺพายนฺตีติ.\csromanspacer{}\csromanspacer{}ทุติยํ.}

\csromanmatikaanchor{mtk.165}
\csromantocmarkref{subsection}{3. นาฬนฺทสุตฺต}{mtk.165}
\bookmark[dest={mtk.165},level=2]{3. นาฬนฺทสุตฺต}
\csromanheader{2}{3. นาฬนฺทสุตฺต}
\proseitem{126}{เอกํ สมยํ ภควา นาฬนฺทายํ วิหรติ ปาวาริกมฺพวเน.\csromanspacer{} อถ โข อุปาลิ คหปติ เยน ภควา เตนุปสงฺกมิ ฯเปฯ เอกมนฺตํ นิสินฺโน โข อุปาลิ คหปติ ภควนฺตํ เอตทโวจ “โก นุ โข ภนฺเต เหตุ โก ปจฺจโย, เยน มิเธกจฺเจ สตฺตา ทิฏฺเฐว ธมฺเม โน ปรินิพฺพายนฺติ.\csromanspacer{} โก ปน ภนฺเต เหตุ โก ปจฺจโย, เยน มิเธกจฺเจ สตฺตา ทิฏฺเฐว ธมฺเม ปรินิพฺพายนฺตี”ติ. (ยถา ปุริมสุตฺตนฺตํ, เอวํ วิตฺถาเรตพฺพํ.) อยํ โข คหปติ เหตุ อยํ ปจฺจโย, เยน มิเธกจฺเจ สตฺตา ทิฏฺเฐว ธมฺเม ปรินิพฺพายนฺตีติ.\csromanspacer{}\csromanspacer{}ตติยํ.}

\csromanmatikaanchor{mtk.166}
\csromantocmarkref{subsection}{4. ภารทฺวาชสุตฺต}{mtk.166}
\bookmark[dest={mtk.166},level=2]{4. ภารทฺวาชสุตฺต}
\csromanheader{2}{4. ภารทฺวาชสุตฺต}
\proseitem{127}{เอกํ สมยํ อายสฺมา ปิณฺโฑลภารทฺวาโช โกสมฺพิยํ วิหรติ โฆสิตาราเม.\csromanspacer{} อถ โข ราชา อุเทโน เยนายสฺมา ปิณฺโฑลภารทฺวาโช เตนุปสงฺกมิ, อุปสงฺกมิตฺวา อายสฺมตา ปิณฺโฑลภารทฺวาเชน สทฺธิํ สมฺโมทิ, สมฺโมทนียํ กถํ สารณียํ วีติสาเรตฺวา เอกมนฺตํ นิสีทิ, เอกมนฺตํ นิสินฺโน โข ราชา อุเทโน อายสฺมนฺตํ ปิณฺโฑลภารทฺวาชํ เอตทโวจ “โก นุ โข โภ ภารทฺวาช}

\vspace{\tipitakaparskip}
\csromancenter{สฬายตนสํยุตฺต เหตุ โก ปจฺจโย, เยนิเม ทหรา ภิกฺขู สุสู\footnote{สุสุ (สี, ก)} กาฬเกสา ภเทฺรน โยพฺพเนน สมนฺนาคตา ปฐเมน วยสา อนิกีฬิตาวิโน กาเมสุ ยาวชีวํ ปริปุณฺณํ ปริสุทฺธํ พฺรหฺมจริยํ จรนฺติ, อทฺธานญฺจ อาปาเทนฺตี”ติ.\csromanspacer{} วุตฺตํ โข เอตํ มหาราช เตน ภควตา ชานตา ปสฺสตา อรหตา สมฺมาสมฺพุทฺเธน “เอถ ตุเมฺห ภิกฺขเว มาตุมตฺตีสุ มาตุจิตฺตํ อุปฏฺฐเปถ, ภคินิมตฺตีสุ ภคินิจิตฺตํ อุปฏฺฐเปถ, ธีตุมตฺตีสุ ธีตุจิตฺตํ อุปฏฺฐเปถา”ติ.\csromanspacer{} อยํ โข มหาราช เหตุ อยํ ปจฺจโย, เยนิเม ทหรา ภิกฺขู สุสู กาฬเกสา ภเทฺรน โยพฺพเนน สมนฺนาคตา ปฐเมน วยสา อนิกีฬิตาวิโน กาเมสุ ยาวชีวํ ปริปุณฺณํ ปริสุทฺธํ พฺรหฺมจริยํ จรนฺติ, อทฺธานญฺจ อาปาเทนฺตีติ.}
\prose{\csromanfolio{329}โลลํ\footnote{โลฬํ (สฺยา, กํ)} โข โภ ภารทฺวาช จิตฺตํ อปฺเปกทา มาตุมตฺตีสุปิ โลภธมฺมา อุปฺปชฺชนฺติ, ภคินิมตฺตีสุปิ โลภธมฺมา อุปฺปชฺชนฺติ, ธีตุมตฺตีสุปิ โลภธมฺมา อุปฺปชฺชนฺติ, อตฺถิ นุ โข โภ ภารทฺวาช อญฺโญ จ เหตุ อญฺโญ จ ปจฺจโย, เยนิเม ทหรา ภิกฺขู สุสู กาฬเกสา ฯเปฯ อทฺธานญฺจ อาปาเทนฺตีติ. วุตฺตํ โข เอตํ มหาราช เตน ภควตา ชานตา ปสฺสตา อรหตา สมฺมาสมฺพุทฺเธน, เอถ ตุเมฺห ภิกฺขเว อิมเมว กายํ อุทฺธํ ปาทตลา อโธ เกสมตฺถกา ตจปริยนฺตํ ปูรํ นานปฺปการสฺส อสุจิโน ปจฺจเวกฺขถ.\csromanspacer{} อตฺถิ อิมสฺมิํ กาเย เกสา โลมา นขา ทนฺตา ตโจ, มํสํ นฺหารุ\footnote{นหารุ (สี, สฺยา, กํ, อิ)} อฏฺฐิ อฏฺฐิมิญฺชํ\footnote{อฏฺฐิมิญฺชา (สี)} วกฺกํ, หทยํ ยกนํ กิโลมกํ ปิหกํ ปปฺผาสํ, อนฺตํ อนฺตคุณํ อุทริยํ กรีสํ, ปิตฺตํ เสมฺหํ ปุพฺโพ โลหิตํ เสโท เมโท, อสฺสุ วสา เขโฬ สิงฺฆาณิกา ลสิกา มุตฺตนฺติ.\csromanspacer{} อยมฺปิ โข มหาราช เหตุ อยํ ปจฺจโย, เยนิเม ทหรา ภิกฺขู สุสู กาฬเกสา ฯเปฯ อทฺธานญฺจ อาปาเทนฺตีติ.\csromanspacer{} เย เต โภ ภารทฺวาช ภิกฺขู ภาวิตกายา ภาวิตสีลา ภาวิตจิตฺตา ภาวิตปญฺญา, เตสํ ตํ สุกรํ โหติ.\csromanspacer{} เย จ โข เต โภ ภารทฺวาช ภิกฺขู อภาวิตกายา อภาวิตสีลา อภาวิตจิตฺตา อภาวิตปญฺญา, เตสํ ตํ ทุกฺกรํ โหติ.\csromanspacer{} อปฺเปกทา โภ ภารทฺวาช อสุภโต}

\gambhira{สฬายตนวคฺคสํยุตฺตปาฬิ}
\prose{\csromanfolio{330}มนสิ กริสฺสามีติ\footnote{มนสิ กริสฺสามาติ (สี, สฺยา, กํ, อิ)} สุภโตว\footnote{สุภโต วา (สี), สุภโต จ (สฺยา, กํ)} อาคจฺฉติ.\csromanspacer{} อตฺถิ นุ โข โภ ภารทฺวาช อญฺโญ จ โข เหตุ อญฺโญ จ ปจฺจโย, เยนิเม ทหรา ภิกฺขู สุสู กาฬเกสา ฯเปฯ อทฺธานญฺจ อาปาเทนฺตีติ.}

\prose{วุตฺตํ โข เอตํ มหาราช เตน ภควตา ชานตา ปสฺสตา อรหตา สมฺมาสมฺพุทฺเธน, เอถ ตุเมฺห ภิกฺขเว อินฺทฺริเยสุ คุตฺตทฺวารา วิหรถ, จกฺขุนา รูปํ ทิสฺวา มา นิมิตฺตคฺคาหิโน อหุวตฺถ, มานุพฺยญฺชนคฺคาหิโน, ยตฺวาธิกรณเมนํ จกฺขุนฺทฺริยํ อสํวุตํ วิหรนฺตํ อภิชฺฌาโทมนสฺสา ปาปกา อกุสลา ธมฺมา อนฺวาสฺสเวยฺยุํ, ตสฺส สํวราย ปฏิปชฺชถ, รกฺขถ จกฺขุนฺทฺริยํ, จกฺขุนฺทฺริเย สํวรํ อาปชฺชถ.\csromanspacer{} โสเตน สทฺทํ สุตฺวา ฯเปฯ.\csromanspacer{} ฆาเนน คนฺธํ ฆายิตฺวา.\csromanspacer{} ชิวฺหาย รสํ สายิตฺวา.\csromanspacer{} กาเยน โผฏฺฐพฺพํ ผุสิตฺวา.\csromanspacer{} มนสา ธมฺมํ วิญฺญาย มา นิมิตฺตคฺคาหิโน อหุวตฺถ, มานุพฺยญฺชนคฺคาหิโน, ยตฺวาธิกรณเมนํ มนินฺทฺริยํ อสํวุตํ วิหรนฺตํ อภิชฺฌาโทมนสฺสา ปาปกา อกุสลา ธมฺมา อนฺวาสฺสเวยฺยุํ, ตสฺส สํวราย ปฏิปชฺชถ, รกฺขถ มนินฺทฺริยํ, มนินฺทฺริเย สํวรํ อาปชฺชถาติ.\csromanspacer{} อยมฺปิ โข มหาราช เหตุ อยํ ปจฺจโย, เยนิเม ทหรา ภิกฺขู สุสู กาฬเกสา ภเทฺรน โยพฺพเนน สมนฺนาคตา ปฐเมน วยสา อนิกีฬิตาวิโน กาเมสุ ยาวชีวํ ปริปุณฺณํ ปริสุทฺธํ พฺรหฺมจริยํ จรนฺติ, อทฺธานญฺจ อาปาเทนฺตีติ.}

\prose{อจฺฉริยํ โภ ภารทฺวาช, อพฺภุตํ โภ ภารทฺวาช, ยาว สุภาสิตํ จิทํ\footnote{ยาว สุภาสิตมิทํ (สี)} โภ ภารทฺวาช เตน ภควตา ชานตา ปสฺสตา อรหตา สมฺมาสมฺพุทฺเธน.\csromanspacer{} เอโสว โข โภ ภารทฺวาช เหตุ เอส ปจฺจโย, เยนิเม ทหรา ภิกฺขู สุสู กาฬเกสา ภเทฺรน โยพฺพเนน สมนฺนาคตา ปฐเมน วยสา อนิกีฬิตาวิโน กาเมสุ ยาวชีวํ ปริปุณฺณํ ปริสุทฺธํ พฺรหฺมจริยํ จรนฺติ, อทฺธานญฺจ อาปาเทนฺตีติ.\csromanspacer{} อหมฺปิ โข โภ\footnote{อหมฺปิ โภ (สี, อิ)} ภารทฺวาช ยสฺมิํ สมเย อรกฺขิเตเนว กาเยน อรกฺขิตาย วาจาย อรกฺขิเตน จิตฺเตน อนุปฏฺฐิตาย สติยา อสํวุเตหิ อินฺทฺริเยหิ อนฺเตปุรํ ปวิสามิ, อติวิย มํ ตสฺมิํ สมเย โลภธมฺมา ปริสหนฺติ.\csromanspacer{} ยสฺมิํ จ ขฺวาหํ โภ ภารทฺวาช สมเย รกฺขิเตเนว กาเยน รกฺขิตาย วาจาย รกฺขิเตน จิตฺเตน อุปฏฺฐิตาย สติยา สํวุเตหิ อินฺทฺริเยหิ อนฺเตปุรํ ปวิสามิ, น มํ}

\vspace{\tipitakaparskip}
\csromancenter{สฬายตนสํยุตฺต ตถา ตสฺมิํ สมเย โลภธมฺมา ปริสหนฺติ.\csromanspacer{} อภิกฺกนฺตํ โภ ภารทฺวาช, อภิกฺกนฺตํ โภ ภารทฺวาช, เสยฺยถาปิ โภ ภารทฺวาช นิกฺกุชฺชิตํ\footnote{นิกุชฺชิตํ (อิ)} วา อุกฺกุชฺเชยฺย, ปฏิจฺฉนฺนํ วา วิวเรยฺย, มูฬฺหสฺส วา มคฺคํ อาจิกฺเขยฺย, อนฺธกาเร วา เตลปชฺโชตํ ธาเรยฺย “จกฺขุมนฺโต รูปานิ ทกฺขนฺตี”ติ.\csromanspacer{} เอวเมวํ โภตา ภารทฺวาเชน อเนกปริยาเยน ธมฺโม ปกาสิโต, เอสาหํ โภ ภารทฺวาช ตํ ภควนฺตํ สรณํ คจฺฉามิ ธมฺมญฺจ ภิกฺขุสํฆญฺจ อุปาสกํ มํ ภวํ ภารทฺวาโช ธาเรตุ อชฺชตคฺเค ปาณุเปตํ สรณํ คตนฺติ.\csromanspacer{}\csromanspacer{}จตุตฺถํ.}
\csromanmatikaanchor{mtk.167}
\csromantocmarkref{subsection}{5. โสณสุตฺต}{mtk.167}
\bookmark[dest={mtk.167},level=2]{5. โสณสุตฺต}
\csromanheader{2}{5. โสณสุตฺต}
\proseitem{128}{\csromanfolio{331}เอกํ สมยํ ภควา ราชคเห วิหรติ เวฬุวเน กลนฺทกนิวาเป.\csromanspacer{} อถ โข โสโณ คหปติปุตฺโต เยน ภควา เตนุปสงฺกมิ, อุปสงฺกมิตฺวา ภควนฺตํ อภิวาเทตฺวา เอกมนฺตํ นิสีทิ, เอกมนฺตํ นิสินฺโน โข โสโณ คหปติปุตฺโต ภควนฺตํ เอตทโวจ “โก นุ โข ภนฺเต เหตุ, โก ปจฺจโย เยน มิเธกจฺเจ สตฺตา ทิฏฺเฐว ธมฺเม โน ปรินิพฺพายนฺติ, โก ปน ภนฺเต เหตุ โก ปจฺจโย, เยน มิเธกจฺเจ สตฺตา ทิฏฺเฐว ธมฺเม ปรินิพฺพายนฺตี”ติ. (ยถา ปุริมสุตฺตนฺตํ, เอวํ วิตฺถาเรตพฺพํ.) อยํ โข โสณ เหตุ อยํ ปจฺจโย, เยน มิเธกจฺเจ สตฺตา ทิฏฺเฐว ธมฺเม ปรินิพฺพายนฺตีติ.\csromanspacer{}\csromanspacer{}ปญฺจมํ.}

\csromanmatikaanchor{mtk.168}
\csromantocmarkref{subsection}{6. โฆสิตสุตฺต}{mtk.168}
\bookmark[dest={mtk.168},level=2]{6. โฆสิตสุตฺต}
\csromanheader{2}{6. โฆสิตสุตฺต}
\proseitem{129}{เอกํ สมยํ อายสฺมา อานนฺโท โกสมฺพิยํ วิหรติ โฆสิตาราเม.\csromanspacer{} อถ โข โฆสิโต คหปติ เยนายสฺมา อานนฺโท เตนุปสงฺกมิ ฯเปฯ เอกมนฺตํ นิสินฺโน โข โฆสิโต คหปติ อายสฺมนฺตํ อานนฺทํ เอตทโวจ “ธาตุนานตฺตํ ธาตุนานตฺตนฺติ ภนฺเต อานนฺท วุจฺจติ, กิตฺตาวตา นุ โข ภนฺเต ธาตุนานตฺตํ วุตฺตํ ภควตา”ติ.\csromanspacer{} สํวิชฺชติ โข คหปติ จกฺขุธาตุ รูปา จ มนาปา จกฺขุวิญฺญาณญฺจ, สุขเวทนิยํ ผสฺสํ ปฏิจฺจ อุปฺปชฺชติ สุขา เวทนา, สํวิชฺชติ โข คหปติ จกฺขุธาตุ รูปา จ อมนาปา จกฺขุวิญฺญาณญฺจ, ทุกฺขเวทนิยํ ผสฺสํ ปฏิจฺจ อุปฺปชฺชติ ทุกฺขา เวทนา, สํวิชฺชติ โข คหปติ จกฺขุธาตุ รูปา จ มนาปา อุเปกฺขาเวทนิยา จกฺขุวิญฺญาณญฺจ อทุกฺขมสุขเวทนิยํ ผสฺสํ ปฏิจฺจ อุปฺปชฺชติ}

\gambhira{สฬายตนวคฺคสํยุตฺตปาฬิ}
\prose{\csromanfolio{332}อทุกฺขมสุขา เวทนา ฯเปฯ.\csromanspacer{} สํวิชฺชติ โข คหปติ ชิวฺหาธาตุ รสา จ มนาปา ชิวฺหาวิญฺญาณญฺจ, สุขเวทนิยํ ผสฺสํ ปฏิจฺจ อุปฺปชฺชติ สุขา เวทนา, สํวิชฺชติ โข คหปติ ชิวฺหาธาตุ รสา จ อมนาปา ชิวฺหาวิญฺญาณญฺจ, ทุกฺขเวทนิยํ ผสฺสํ ปฏิจฺจ อุปฺปชฺชติ ทุกฺขา เวทนา, สํวิชฺชติ โข คหปติ ชิวฺหาธาตุ รสา จ อุเปกฺขาเวทนิยา ชิวฺหาวิญฺญาณญฺจ, อทุกฺขมสุขเวทนิยํ ผสฺสํ ปฏิจฺจ อุปฺปชฺชติ อทุกฺขมสุขา เวทนา ฯเปฯ.\csromanspacer{} สํวิชฺชติ โข คหปติ มโนธาตุ ธมฺมา จ มนาปา มโนวิญฺญาณญฺจ, สุขเวทนิยํ ผสฺสํ ปฏิจฺจ อุปฺปชฺชติ สุขา เวทนา, สํวิชฺชติ โข คหปติ มโนธาตุ ธมฺมา จ อมนาปา มโนวิญฺญาณญฺจ, ทุกฺขเวทนิยํ ผสฺสํ ปฏิจฺจ อุปฺปชฺชติ ทุกฺขา เวทนา, สํวิชฺชติ โข คหปติ มโนธาตุ ธมฺมา จ อุเปกฺขาเวทนิยา มโนวิญฺญาณญฺจ, อทุกฺขมสุขเวทนิยํ ผสฺสํ ปฏิจฺจ อุปฺปชฺชติ อทุกฺขมสุขา เวทนา.\csromanspacer{} เอตฺตาวตา โข คหปติ ธาตุนานตฺตํ วุตฺตํ ภควตาติ.\csromanspacer{}\csromanspacer{}ฉฏฺฐํ.}

\csromanmatikaanchor{mtk.169}
\csromantocmarkref{subsection}{7. หาลิทฺทิกานิสุตฺต}{mtk.169}
\bookmark[dest={mtk.169},level=2]{7. หาลิทฺทิกานิสุตฺต}
\csromanheader{2}{7. หาลิทฺทิกานิสุตฺต}
\proseitem{130}{เอกํ สมยํ อายสฺมา มหากจฺจาโน อวนฺตีสุ วิหรติ กุรรฆเร\footnote{กุลฆเร (สฺยา, ก)} ปปาเต\footnote{ปวตฺเต (สี,อิ), สมฺปวตฺเต (สฺยา, กํ, ก) เอตฺเถว อฏฺฐมปิฏฺเฐปิ.} ปพฺพเต.\csromanspacer{} อถ โข หาลิทฺทิกานิ\footnote{หาลิทฺทกานิ (สี, สฺยา, กํ)} คหปติ เยนายสฺมา มหากจฺจาโน เตนุปสงฺกมิ ฯเปฯ เอกมนฺตํ นิสินฺโน โข หาลิทฺทิกานิ คหปติ อายสฺมนฺตํ มหากจฺจานํ เอตทโวจ “วุตฺตมิทํ ภนฺเต ภควตา ธาตุนานตฺตํ ปฏิจฺจ อุปฺปชฺชติ ผสฺสนานตฺตํ, ผสฺสนานตฺตํ ปฏิจฺจ อุปฺปชฺชติ เวทนานานตฺตนฺ”ติ.\csromanspacer{} กถํ นุ โข ภนฺเต ธาตุนานตฺตํ ปฏิจฺจ อุปฺปชฺชติ ผสฺสนานตฺตํ, ผสฺสนานตฺตํ ปฏิจฺจ อุปฺปชฺชติ เวทนานานตฺตนฺติ.\csromanspacer{} อิธ คหปติ ภิกฺขุ จกฺขุนา รูปํ ทิสฺวา มนาปํ อิตฺเถตนฺติ ปชานาติ, จกฺขุวิญฺญาณํ สุขเวทนิยญฺจ\footnote{สุขเวทนิยํ, สุขเวทนิยํ (สี, อิ), สุขเวทนิยญฺจ, สุขเวทนิยํ (สฺยา, กํ, ก) เอวํ “ทุกฺขเวทนิยญฺจ อทุกฺขมสุขเวทนิยญฺจา”ติ ปเทสุปิ. อฏฺฐกถาฏีกา โอโลเกตพฺพา.} ผสฺสํ ปฏิจฺจ อุปฺปชฺชติ สุขา เวทนา.\csromanspacer{} จกฺขุนา โข ปเนว\footnote{ปเนวํ (สฺยา, กํ, ก)} รูปํ ทิสฺวา อมนาปํ อิตฺเถตนฺติ ปชานาติ, จกฺขุวิญฺญาณํ ทุกฺขเวทนิยญฺจ ผสฺสํ ปฏิจฺจ อุปฺปชฺชติ ทุกฺขา เวทนา.\csromanspacer{} จกฺขุนา โข ปเนว รูปํ ทิสฺวา}

\vspace{\tipitakaparskip}
\csromancenter{สฬายตนสํยุตฺต อุเปกฺขาฏฺฐานิยํ\footnote{อุเปกฺขาเวทนิยํ (ก)} อิตฺเถตนฺติ ปชานาติ, จกฺขุวิญฺญาณํ อทุกฺขมสุขเวทนิยญฺจ ผสฺสํ ปฏิจฺจ อุปฺปชฺชติ อทุกฺขมสุขา เวทนา.}
\prose{\csromanfolio{333}ปุน จปรํ คหปติ ภิกฺขุ โสเตน สทฺทํ สุตฺวา ฯเปฯ ฆาเนน คนฺธํ ฆายิตฺวา ฯเปฯ ชิวฺหาย รสํ สายิตฺวา ฯเปฯ กาเยน โผฏฺฐพฺพํ ผุสิตฺวา ฯเปฯ มนสา ธมฺมํ วิญฺญาย มนาปํ อิตฺเถตนฺติ ปชานาติ, มโนวิญฺญาณํ สุขเวทนิยญฺจ ผสฺสํ ปฏิจฺจ อุปฺปชฺชติ สุขา เวทนา.\csromanspacer{} มนสา โข ปเนว ธมฺมํ วิญฺญาย อมนาปํ อิตฺเถตนฺติ ปชานาติ, มโนวิญฺญาณํ ทุกฺขเวทนิยญฺจ ผสฺสํ ปฏิจฺจ อุปฺปชฺชติ ทุกฺขา เวทนา.\csromanspacer{} มนสา โข ปเนว ธมฺมํ วิญฺญาย อุเปกฺขาฏฺฐานิยํ อิตฺเถตนฺติ ปชานาติ, มโนวิญฺญาณํ อทุกฺขมสุขเวทนิยญฺจ ผสฺสํ ปฏิจฺจ อุปฺปชฺชติ อทุกฺขมสุขา เวทนา.\csromanspacer{} เอวํ โข คหปติ ธาตุนานตฺตํ ปฏิจฺจ อุปฺปชฺชติ ผสฺสนานตฺตํ, ผสฺสนานตฺตํ ปฏิจฺจ อุปฺปชฺชติ เวทนานานตฺตนฺติ.\csromanspacer{}\csromanspacer{}สตฺตมํ.}

\csromanmatikaanchor{mtk.170}
\csromantocmarkref{subsection}{8. นกุลปิตุสุตฺต}{mtk.170}
\bookmark[dest={mtk.170},level=2]{8. นกุลปิตุสุตฺต}
\csromanheader{2}{8. นกุลปิตุสุตฺต}
\proseitem{131}{เอกํ สมยํ ภควา ภคฺเคสุ วิหรติ สุสุมารคิเร เภสกฬาวเน มิคทาเย.\csromanspacer{} อถ โข นกุลปิตา คหปติ เยน ภควา เตนุปสงฺกมิ ฯเปฯ เอกมนฺตํ นิสินฺโน โข นกุลปิตา คหปติ ภควนฺตํ เอตทโวจ”โก นุ โข ภนฺเต เหตุ โก ปจฺจโย, เยน มิเธกจฺเจ สตฺตา ทิฏฺเฐว ธมฺเม โน ปรินิพฺพายนฺติ, โก ปน ภนฺเต เหตุ โก ปจฺจโย, เยน มิเธกจฺเจ สตฺตา ทิฏฺเฐว ธมฺเม ปรินิพฺพายนฺตี”ติ.\csromanspacer{} สนฺติ โข คหปติ จกฺขุวิญฺเญยฺยา รูปา อิฏฺฐา กนฺตา มนาปา ปิยรูปา กามูปสํหิตา รชนียา, ตญฺเจ ภิกฺขุ อภินนฺทติ อภิวทติ อชฺโฌสาย ติฏฺฐติ, ตสฺส ตํ อภินนฺทโต อภิวทโต อชฺโฌสาย ติฏฺฐโต ตนฺนิสฺสิตํ วิญฺญาณํ โหติ, ตทุปาทานํ, ส-อุปาทาโน คหปติ ภิกฺขุ โน ปรินิพฺพายติ ฯเปฯ.\csromanspacer{} สนฺติ โข คหปติ ชิวฺหาวิญฺเญยฺยา รสา ฯเปฯ.\csromanspacer{} สนฺติ โข คหปติ มโนวิญฺเญยฺยา ธมฺมา อิฏฺฐา กนฺตา มนาปา ปิยรูปา กามูปสํหิตา รชนียา, ตญฺเจ ภิกฺขุ อภินนฺทติ อภิวทติ อชฺโฌสาย ติฏฺฐติ, ตสฺส ตํ อภินนฺทโต อภิวทโต อชฺโฌสาย ติฏฺฐโต ตนฺนิสฺสิตํ วิญฺญาณํ โหติ, ตทุปาทานํ, ส-อุปาทาโน คหปติ ภิกฺขุ โน}

\gambhira{สฬายตนวคฺคสํยุตฺตปาฬิ}
\prose{\csromanfolio{334}ปรินิพฺพายติ.\csromanspacer{} อยํ โข คหปติ เหตุ อยํ ปจฺจโย, เยน มิเธกจฺเจ สตฺตา ทิฏฺเฐว ธมฺเม โน ปรินิพฺพายนฺติ.}

\prose{สนฺติ จ โข คหปติ จกฺขุวิญฺเญยฺยา รูปา อิฏฺฐา กนฺตา มนาปา ปิยรูปา กามูปสํหิตา รชนียา, ตญฺเจ ภิกฺขุ นาภินนฺทติ นาภิวทติ นาชฺโฌสาย ติฏฺฐติ, ตสฺส ตํ อนภินนฺทโต อนภิวทโต อนชฺโฌสาย ติฏฺฐโต น ตนฺนิสฺสิตํ วิญฺญาณํ โหติ, น ตทุปาทานํ, อนุปาทาโน คหปติ ภิกฺขุ ปรินิพฺพายติ ฯเปฯ.\csromanspacer{} สนฺติ โข คหปติ ชิวฺหาวิญฺเญยฺยา รสา ฯเปฯ.\csromanspacer{} สนฺติ โข คหปติ มโนวิญฺเญยฺยา ธมฺมา อิฏฺฐา กนฺตา มนาปา ปิยรูปา กามูปสํหิตา รชนียา, ตญฺเจ ภิกฺขุ นาภินนฺทติ นาภิวทติ นาชฺโฌสาย ติฏฺฐติ, ตสฺส ตํ นาภินนฺทโต นาภิวทโต อนชฺโฌสาย ติฏฺฐโต น ตนฺนิสฺสิตํ วิญฺญาณํ โหติ, น ตทุปาทานํ, อนุปาทาโน คหปติ ภิกฺขุ ปรินิพฺพายติ.\csromanspacer{} อยํ โข คหปติ เหตุ อยํ ปจฺจโย, เยน มิเธกจฺเจ สตฺตา ทิฏฺเฐว ธมฺเม ปรินิพฺพายนฺตีติ.\csromanspacer{}\csromanspacer{}อฏฺฐมํ.}

\csromanmatikaanchor{mtk.171}
\csromantocmarkref{subsection}{9. โลหิจฺจสุตฺต}{mtk.171}
\bookmark[dest={mtk.171},level=2]{9. โลหิจฺจสุตฺต}
\csromanheader{2}{9. โลหิจฺจสุตฺต}
\proseitem{132}{เอกํ สมยํ อายสฺมา มหากจฺจาโน อวนฺตีสุ วิหรติ มกฺกรกเต\footnote{มกฺกรกเฏ (สี, สฺยา, กํ, อิ)} อรญฺญกุฏิกายํ.\csromanspacer{} อถ โข โลหิจฺจสฺส พฺราหฺมณสฺส สมฺพหุลา อนฺเตวาสิกา กฏฺฐหารกา มาณวกา เยนายสฺมโต มหากจฺจานสฺส อรญฺญกุฏิกา เตนุปสงฺกมิํสุ, อุปสงฺกมิตฺวา ปริโต ปริโต กุฏิกาย อนุจงฺกมนฺติ อนุวิจรนฺติ อุจฺจาสทฺทา มหาสทฺทา กานิจิ กานิจิ เสเลยฺยกานิ กโรนฺติ\footnote{เสลิสฺสกานิ กโรนฺตา (สี)} “อิเม ปน มุณฺฑกา สมณกา อิพฺภา กณฺหา\footnote{กิณฺหา (สี, อิ)} พนฺธุปาทาปจฺจา อิเมสํ ภรตกานํ สกฺกตา ครุกตา มานิตา ปูชิตา อปจิตา”ติ.\csromanspacer{} อถ โข อายสฺมา มหากจฺจาโน วิหารา นิกฺขมิตฺวา เต มาณวเก เอตทโวจ “มา มาณวกา สทฺทมกตฺถ, ธมฺมํ โว ภาสิสฺสามี”ติ.\csromanspacer{} เอวํ วุตฺเต เต มาณวกา ตุณฺหี อเหสุํ.\csromanspacer{} อถ โข อายสฺมา มหากจฺจาโน เต มาณวเก คาถาหิ อชฺฌภาสิ\csromandash{}}

\csromangathameasure{“สีลุตฺตมา ปุพฺพตรา อเหสุํ, \\ เต พฺราหฺมณา เย ปุราณํ สรนฺติ. \\ คุตฺตานิ ทฺวารานิ สุรกฺขิตานิ, \\ อเหสุํ เตสํ อภิภุยฺย โกธํ.}
\csromangathagroup{\csromangathastanza{“สีลุตฺตมา ปุพฺพตรา อเหสุํ, \\ เต พฺราหฺมณา เย ปุราณํ สรนฺติ. \\ คุตฺตานิ ทฺวารานิ สุรกฺขิตานิ, \\ อเหสุํ เตสํ อภิภุยฺย โกธํ.}}

\csromanheader{2}{1. สฬายตนสํยุตฺต}
\csromangathameasure{ธมฺเม จ ฌาเน จ รตา อเหสุํ, \\ เต พฺราหฺมณา เย ปุราณํ สรนฺติ. \\ อิเม จ โวกฺกมฺม ชปามเสติ, \\ โคตฺเตน มตฺตา วิสมํ จรนฺติ. \\ โกธาภิภูตา ปุถุ-อตฺตทณฺฑา, \\ วิรชฺชมานา สตณฺหาตเณฺหสุ. \\ อคุตฺตทฺวารสฺส ภวนฺติ โมฆา, \\ สุปิเนว ลทฺธํ ปุริสสฺส วิตฺตํ. \\ อนาสกา ถณฺฑิลสายิกา จ, \\ ปาโต สินานญฺจ ตโย จ เวทา. \\ ขราชินํ ชฏาปงฺโก, \\ มนฺตา สีลพฺพตํ ตโป. \\ กุหนา วงฺกทณฺฑา จ, อุทกาจมนานิ จ. \\ วณฺณา เอเต พฺราหฺมณานํ, กตา กิญฺจิกฺขภาวนา. \\ จิตฺตญฺจ สุสมาหิตํ, วิปฺปสนฺนมนาวิลํ. \\ อขิลํ สพฺพภูเตสุ, โส มคฺโค พฺรหฺมปตฺติยา”ติ.}
\csromangathagroup{\csromangathastanza{\csromanfolio{335}ธมฺเม จ ฌาเน จ รตา อเหสุํ, \\ เต พฺราหฺมณา เย ปุราณํ สรนฺติ. \\ อิเม จ โวกฺกมฺม ชปามเสติ, \\ โคตฺเตน มตฺตา วิสมํ จรนฺติ.}\csromangathastanza{โกธาภิภูตา ปุถุ-อตฺตทณฺฑา\footnote{โกธาภิภูตาสุปุถุตฺตทณฺฑา (สฺยา, กํ, ก)}, \\ วิรชฺชมานา สตณฺหาตเณฺหสุ. \\ อคุตฺตทฺวารสฺส ภวนฺติ โมฆา, \\ สุปิเนว ลทฺธํ ปุริสสฺส วิตฺตํ.}\csromangathastanza{อนาสกา ถณฺฑิลสายิกา จ, \\ ปาโต สินานญฺจ ตโย จ เวทา. \\ ขราชินํ ชฏาปงฺโก, \\ มนฺตา สีลพฺพตํ ตโป. \\ กุหนา วงฺกทณฺฑา จ, อุทกาจมนานิ จ. \\ วณฺณา เอเต พฺราหฺมณานํ, กตา กิญฺจิกฺขภาวนา.}\csromangathastanza{จิตฺตญฺจ สุสมาหิตํ, วิปฺปสนฺนมนาวิลํ. \\ อขิลํ สพฺพภูเตสุ, โส มคฺโค พฺรหฺมปตฺติยา”ติ.}}

\prose{อถ โข เต มาณวกา กุปิตา อนตฺตมนา เยน โลหิจฺโจ พฺราหฺมโณ เตนุปสงฺกมิํสุ, อุปสงฺกมิตฺวา โลหิจฺจํ พฺราหฺมณํ เอตทโวจุํ “ยคฺเฆ ภวํ ชาเนยฺย สมโณ มหากจฺจาโน พฺราหฺมณานํ มนฺเต\footnote{มนฺตํ (ก)} เอกํเสน อปวทติ ปฏิกฺโกสตี”ติ.\csromanspacer{} เอวํ วุตฺเต โลหิจฺโจ พฺราหฺมโณ กุปิโต อโหสิ อนตฺตมโน.\csromanspacer{} อถ โข โลหิจฺจสฺส พฺราหฺมณสฺส เอตทโหสิ “น โข ปน เมตํ ปติรูปํ, โยหํ อญฺญทตฺถุ มาณวกานํเยว สุตฺวา สมณํ มหากจฺจานํ อกฺโกเสยฺยํ\footnote{อกฺโกเสยฺยํ วิรุชฺเฌยฺยํ (สฺยา, กํ, ก)} ปริภาเสยฺยํ, ยํนูนาหํ อุปสงฺกมิตฺวา ปุจฺเฉยฺยนฺ”ติ.}

\prose{อถ โข โลหิจฺโจ พฺราหฺมโณ เตหิ มาณวเกหิ สทฺธิํ เยนายสฺมา มหากจฺจาโน เตนุปสงฺกมิ, อุปสงฺกมิตฺวา อายสฺมตา มหากจฺจาเนน สทฺธิํ สมฺโมทิ, สมฺโมทนียํ กถํ สารณียํ วีติสาเรตฺวา เอกมนฺตํ นิสีทิ, เอกมนฺตํ นิสินฺโน โข โลหิจฺโจ}

\gambhira{สฬายตนวคฺคสํยุตฺตปาฬิ}
\prose{\csromanfolio{336}พฺราหฺมโณ อายสฺมนฺตํ มหากจฺจานํ เอตทโวจ “อาคมํสุ นุ ขฺวิธ โภ กจฺจาน อมฺหากํ สมฺพหุลา อนฺเตวาสิกา กฏฺฐหารกา มาณวกา”ติ.\csromanspacer{} อาคมํสุ ขฺวิธ เต พฺราหฺมณ สมฺพหุลา อนฺเตวาสิกา กฏฺฐหารกา มาณวกาติ.\csromanspacer{} อหุ ปน โภโต กจฺจานสฺส เตหิ มาณวเกหิ สทฺธิํ โกจิเทว กถาสลฺลาโปติ.\csromanspacer{} อหุ โข เม พฺราหฺมณ เตหิ มาณวเกหิ สทฺธิํ โกจิเทว กถาสลฺลาโปติ.\csromanspacer{} ยถา กถํ ปน โภโต กจฺจานสฺส เตหิ มาณวเกหิ สทฺธิํ อโหสิ กถาสลฺลาโปติ.\csromanspacer{} เอวํ โข เม พฺราหฺมณ เตหิ มาณวเกหิ สทฺธิํ อโหสิ กถาสลฺลาโป\csromandash{}}

\csromangathameasure{สีลุตฺตมา ปุพฺพตรา อเหสุํ, \\ เต พฺราหฺมณา เย ปุราณํ สรนฺติ ฯเปฯ. \\ อขิลํ สพฺพภูเตสุ, \\ โส มคฺโค พฺรหฺมปตฺติยาติ.}
\csromangathagroup{\csromangathastanza{สีลุตฺตมา ปุพฺพตรา อเหสุํ, \\ เต พฺราหฺมณา เย ปุราณํ สรนฺติ ฯเปฯ. \\ อขิลํ สพฺพภูเตสุ, \\ โส มคฺโค พฺรหฺมปตฺติยาติ.}}

\prose{เอวํ โข เม พฺราหฺมณ เตหิ มาณวเกหิ สทฺธิํ อโหสิ กถาสลฺลาโปติ.}

\prose{อคุตฺตทฺวาโรติ\footnote{อคุตฺตทฺวาโร อคุตฺตทฺวาโรติ (ก)} ภวํ กจฺจาโน อาห, กิตฺตาวตา นุ โข โภ กจฺจาน อคุตฺตทฺวาโร โหตีติ.\csromanspacer{} อิธ พฺราหฺมณ เอกจฺโจ จกฺขุนา รูปํ ทิสฺวา ปิยรูเป รูเป อธิมุจฺจติ, อปฺปิยรูเป รูเป พฺยาปชฺชติ, อนุปฏฺฐิตกายสฺสติ\footnote{อนุปฏฺฐิตาย สติยา (สฺยา, กํ, อิ, ก) อุปริ อาสีวิสวคฺเค อวสฺสุตสุตฺเต ปน “อนุปฏฺฐิตกายสฺสตี”เตฺวว สพฺพตฺถ ทิสฺสติ.} จ วิหรติ ปริตฺตเจตโส, ตญฺจ เจโตวิมุตฺติํ ปญฺญาวิมุตฺติํ ยถาภูตํ นปฺปชานาติ, ยตฺถสฺส เต อุปฺปนฺนา ปาปกา อกุสลา ธมฺมา อปริเสสา นิรุชฺฌนฺติ.\csromanspacer{} โสเตน สทฺทํ สุตฺวา.\csromanspacer{} ฆาเนน คนฺธํ ฆายิตฺวา.\csromanspacer{} ชิวฺหาย รสํ สายิตฺวา.\csromanspacer{} กาเยน โผฏฺฐพฺพํ ผุสิตฺวา.\csromanspacer{} มนสา ธมฺมํ วิญฺญาย ปิยรูเป ธมฺเม อธิมุจฺจติ.\csromanspacer{} อปฺปิยรูเป จ ธมฺเม พฺยาปชฺชติ.\csromanspacer{} อนุปฏฺฐิตกายสฺสติ จ วิหรติ ปริตฺตเจตโส, ตญฺจ เจโตวิมุตฺติํ ปญฺญาวิมุตฺติํ ยถาภูตํ นปฺปชานาติ, ยตฺถสฺส เต อุปฺปนฺนา ปาปกา อกุสลา ธมฺมา อปริเสสา นิรุชฺฌนฺติ, เอวํ โข พฺราหฺมณ อคุตฺตทฺวาโร โหตีติ.\csromanspacer{} อจฺฉริยํ โภ กจฺจาน, อพฺภุตํ โภ กจฺจาน, ยาวญฺจิทํ โภตา กจฺจาเนน อคุตฺตทฺวาโรว สมาโน อคุตฺตทฺวาโรติ อกฺขาโต.}

\csromanheader{2}{1. สฬายตนสํยุตฺต}
\prose{\csromanfolio{337}คุตฺตทฺวาโรติ ภวํ กจฺจาโน อาห, กิตฺตาวตา นุ โข โภ กจฺจาน คุตฺตทฺวาโร โหตีติ.\csromanspacer{} อิธ พฺราหฺมณ ภิกฺขุ จกฺขุนา รูปํ ทิสฺวา ปิยรูเป รูเป นาธิมุจฺจติ, อปฺปิยรูเป รูเป น พฺยาปชฺชติ, อุปฏฺฐิตกายสฺสติ จ วิหรติ อปฺปมาณเจตโส, ตญฺจ เจโตวิมุตฺติํ ปญฺญาวิมุตฺติํ ยถาภูตํ ปชานาติ, ยตฺถสฺส เต อุปฺปนฺนา ปาปกา อกุสลา ธมฺมา อปริเสสา นิรุชฺฌนฺติ.\csromanspacer{} โสเตน สทฺทํ สุตฺวา.\csromanspacer{} ฆาเนน คนฺธํ ฆายิตฺวา.\csromanspacer{} ชิวฺหาย รสํ สายิตฺวา.\csromanspacer{} กาเยน โผฏฺฐพฺพํ ผุสิตฺวา.\csromanspacer{} มนสา ธมฺมํ วิญฺญาย ปิยรูเป ธมฺเม นาธิมุจฺจติ, อปฺปิยรูเป ธมฺเม น พฺยาปชฺชติ, อุปฏฺฐิตกายสฺสติ จ วิหรติ อปฺปมาณเจตโส, ตญฺจ เจโตวิมุตฺติํ ปญฺญาวิมุตฺติํ ยถาภูตํ ปชานาติ, ยตฺถสฺส เต อุปฺปนฺนา ปาปกา อกุสลา ธมฺมา อปริเสสา นิรุชฺฌนฺติ, เอวํ โข พฺราหฺมณ คุตฺตทฺวาโร โหตีติ.}

\prose{อจฺฉริยํ โภ กจฺจาน, อพฺภุตํ โภ กจฺจาน, ยาวญฺจิทํ โภตา กจฺจาเนน คุตฺตทฺวาโรว สมาโน คุตฺตทฺวาโรติ อกฺขาโต.\csromanspacer{} อภิกฺกนฺตํ โภ กจฺจาน, อภิกฺกนฺตํ โภ กจฺจาน, เสยฺยถาปิ โภ กจฺจาน นิกฺกุชฺชิตํ วา อุกฺกุชฺเชยฺย, ปฏิจฺฉนฺนํ วา วิวเรยฺย มูฬฺหสฺส วา มคฺคํ อาจิกฺเขยฺย, อนฺธกาเร วา เตลปชฺโชตํ ธาเรยฺย “จกฺขุมนฺโต รูปานิ ทกฺขนฺตี”ติ.\csromanspacer{} เอวเมวํ โภตา กจฺจาเนน อเนกปริยาเยน ธมฺโม ปกาสิโต, เอสาหํ โภ กจฺจาน ตํ ภควนฺตํ สรณํ คจฺฉามิ, ธมฺมญฺจ ภิกฺขุสํฆญฺจ, อุปาสกํ มํ ภวํ กจฺจาโน ธาเรตุ อชฺชตคฺเค ปาณุเปตํ สรณํ คตํ.\csromanspacer{} ยถา จ ภวํ กจฺจาโน มกฺกรกเต อุปาสกกุลานิ อุปสงฺกมติ, เอวเมว โลหิจฺจกุลํ อุปสงฺกมตุ, ตตฺถ เย มาณวกา วา มาณวิกา วา ภวนฺตํ กจฺจานํ อภิวาเทสฺสนฺติ, ปจฺจุฏฺฐิสฺสนฺติ, อาสนํ วา อุทกํ วา ทสฺสนฺติ, เตสํ ตํ ภวิสฺสติ ทีฆรตฺตํ หิตาย สุขายาติ.\csromanspacer{}\csromanspacer{}นวมํ.}

\csromanmatikaanchor{mtk.172}
\csromantocmarkref{subsection}{10. เวรหจฺจานิสุตฺต}{mtk.172}
\bookmark[dest={mtk.172},level=2]{10. เวรหจฺจานิสุตฺต}
\csromanheader{2}{10. เวรหจฺจานิสุตฺต}
\proseitem{133}{เอกํ สมยํ อายสฺมา อุทายี กามณฺฑายํ วิหรติ โตเทยฺยสฺส พฺราหฺมณสฺส อมฺพวเน.\csromanspacer{} อถ โข เวรหจฺจานิโคตฺตาย พฺราหฺมณิยา อนฺเตวาสี มาณวโก เยนายสฺมา อุทายี เตนุปสงฺกมิ, อุปสงฺกมิตฺวา อายสฺมตา อุทายินา สทฺธิํ สมฺโมทิ, สมฺโมทนียํ กถํ สารณียํ วีติสาเรตฺวา เอกมนฺตํ นิสีทิ, เอกมนฺตํ นิสินฺนํ โข ตํ}

\gambhira{สฬายตนวคฺคสํยุตฺตปาฬิ}
\prose{\csromanfolio{338}มาณวกํ อายสฺมา อุทายี ธมฺมิยา กถาย สนฺทสฺเสสิ สมาทเปสิ สมุตฺเตเชสิ สมฺปหํเสสิ.\csromanspacer{} อถ โข โส มาณวโก อายสฺมตา อุทายินา ธมฺมิยา กถาย สนฺทสฺสิโต สมาทปิโต สมุตฺเตชิโต สมฺปหํสิโต อุฏฺฐายาสนา เยน เวรหจฺจานิโคตฺตา พฺราหฺมณี, เตนุปสงฺกมิ, อุปสงฺกมิตฺวา เวรหจฺจานิโคตฺตํ พฺราหฺมณิํ เอตทโวจ “ยคฺเฆ โภติ ชาเนยฺยาสิ\footnote{โภติ ชาเนยฺย (สี, อิ, ก), โภตี ชาเนยฺย (สฺยา, กํ)}, สมโณ อุทายี ธมฺมํ เทเสติ อาทิกลฺยาณํ มชฺเฌกลฺยาณํ ปริโยสานกลฺยาณํ สาตฺถํ สพฺยญฺชนํ, เกวลปริปุณฺณํ ปริสุทฺธํ พฺรหฺมจริยํ ปกาเสตี”ติ.\csromanspacer{} เตน หิ ตฺวํ มาณวก มม วจเนน สมณํ อุทายิํ นิมนฺเตหิ สฺวาตนาย ภตฺเตนาติ. “เอวํ โภตี”ติ โข โส มาณวโก เวรหจฺจานิโคตฺตาย พฺราหฺมณิยา ปฏิสฺสุตฺวา เยนายสฺมา อุทายี เตนุปสงฺกมิ, อุปสงฺกมิตฺวา อายสฺมนฺตํ อุทายิํ เอตทโวจ “อธิวาเสตุ กิร ภวํ อุทายี อมฺหากํ อาจริยภริยาย เวรหจฺจานิโคตฺตาย พฺราหฺมณิยา สฺวาตนาย ภตฺตนฺ”ติ, อธิวาเสสิ โข อายสฺมา อุทายี ตุณฺหีภาเวน.\csromanspacer{} อถ โข อายสฺมา อุทายี ตสฺสา รตฺติยา อจฺจเยน ปุพฺพณฺหสมยํ นิวาเสตฺวา ปตฺตจีวรมาทาย เยน เวรหจฺจานิโคตฺตาย พฺราหฺมณิยา นิเวสนํ เตนุปสงฺกมิ, อุปสงฺกมิตฺวา ปญฺญตฺเต อาสเน นิสีทิ.\csromanspacer{} อถ โข เวรหจฺจานิโคตฺตา พฺราหฺมณี อายสฺมนฺตํ อุทายิํ ปณีเตน ขาทนีเยน โภชนีเยน สหตฺถา สนฺตปฺเปสิ สมฺปวาเรสิ.\csromanspacer{} อถ โข เวรหจฺจานิโคตฺตา พฺราหฺมณี อายสฺมนฺตํ อุทายิํ ภุตฺตาวิํ โอนีตปตฺตปาณิํ ปาทุกา อาโรหิตฺวา อุจฺเจ อาสเน นิสีทิตฺวา สีสํ โอคุณฺฐิตฺวา อายสฺมนฺตํ อุทายิํ เอตทโวจ “ภณ สมณธมฺมนฺ”ติ. “ภวิสฺสติ ภคินิ สมโย”ติ วตฺวา อุฏฺฐายาสนา ปกฺกมิ\footnote{ปกฺกามิ (สฺยา, กํ, อิ)}.}

\prose{ทุติยมฺปิ โข โส มาณวโก เยนายสฺมา อุทายี เตนุปสงฺกมิ, อุปสงฺกมิตฺวา อายสฺมตา อุทายินา สทฺธิํ สมฺโมทิ, สมฺโมทนียํ กถํ สารณียํ วีติสาเรตฺวา เอกมนฺตํ นิสีทิ, เอกมนฺตํ นิสินฺนํ โข ตํ มาณวกํ อายสฺมา อุทายี ธมฺมิยา กถาย สนฺทสฺเสสิ สมาทเปสิ สมุตฺเตเชสิ สมฺปหํเสสิ.\csromanspacer{} ทุติยมฺปิ โข โส มาณวโก อายสฺมตา อุทายินา ธมฺมิยา กถาย สนฺทสฺสิโต สมาทปิโต สมุตฺเตชิโต สมฺปหํสิโต อุฏฺฐายาสนา เยน เวรหจฺจานิโคตฺตา พฺราหฺมณี เตนุปสงฺกมิ,}

\vspace{\tipitakaparskip}
\csromancenter{สฬายตนสํยุตฺต อุปสงฺกมิตฺวา เวรหจฺจานิโคตฺตํ พฺราหฺมณิํ เอตทโวจ “ยคฺเฆ โภติ ชาเนยฺยาสิ, สมโณ อุทายี ธมฺมํ เทเสติ อาทิกลฺยาณํ มชฺเฌกลฺยาณํ ปริโยสานกลฺยาณํ สาตฺถํ สพฺยญฺชนํ เกวลปริปุณฺณํ ปริสุทฺธํ พฺรหฺมจริยํ ปกาเสตี”ติ.\csromanspacer{} เอวเมวํ ปน ตฺวํ มาณวก สมณสฺส อุทายิสฺส วณฺณํ ภาสสิ, สมโณ ปนุทายี “ภณ สมณธมฺมนฺ”ติ วุตฺโต สมาโน “ภวิสฺสติ ภคินิ สมโย”ติ วตฺวา อุฏฺฐายาสนา ปกฺกนฺโตติ.\csromanspacer{} ตถา หิ ปน ตฺวํ โภติ ปาทุกา อาโรหิตฺวา อุจฺเจ อาสเน นิสีทิตฺวา สีสํ โอคุณฺฐิตฺวา เอตทโวจ “ภณ สมณธมฺมนฺ”ติ.\csromanspacer{} ธมฺมครุโน หิ เต ภวนฺโต ธมฺมคารวาติ.\csromanspacer{} เตน หิ ตฺวํ มาณวก มม วจเนน สมณํ อุทายิํ นิมนฺเตหิ สฺวาตนาย ภตฺเตนาติ. “เอวํ โภตี”ติ โข โส มาณวโก เวรหจฺจานิโคตฺตาย พฺราหฺมณิยา ปฏิสฺสุตฺวา เยนายสฺมา อุทายี เตนุปสงฺกมิ, อุปสงฺกมิตฺวา อายสฺมนฺตํ อุทายิํ เอตทโวจ “อธิวาเสตุ กิร ภวํ อุทายี อมฺหากํ อาจริยภริยาย เวรหจฺจานิโคตฺตาย พฺราหฺมณิยา สฺวาตนาย ภตฺตนฺ”ติ.\csromanspacer{} อธิวาเสสิ โข อายสฺมา อุทายี ตุณฺหีภาเวน.}
\prose{\csromanfolio{339}อถ โข อายสฺมา อุทายี ตสฺสา รตฺติยา อจฺจเยน ปุพฺพณฺหสมยํ นิวาเสตฺวา ปตฺตจีวรมาทาย เยน เวรหจฺจานิโคตฺตาย พฺราหฺมณิยา นิเวสนํ เตนุปสงฺกมิ, อุปสงฺกมิตฺวา ปญฺญตฺเต อาสเน นิสีทิ.\csromanspacer{} อถ โข เวรหจฺจานิโคตฺตา พฺราหฺมณี อายสฺมนฺตํ อุทายิํ ปณิเตน ขาทนีเยน โภชนีเยน สหตฺถา สนฺตปฺเปสิ สมฺปวาเรสิ.\csromanspacer{} อถ โข เวรหจฺจานิโคตฺตา พฺราหฺมณี อายสฺมนฺตํ อุทายิํ ภุตฺตาวิํ โอนีตปตฺตปาณิํ ปาทุกา โอโรหิตฺวา นีเจ อาสเน นิสีทิตฺวา สีสํ วิวริตฺวา อายสฺมนฺตํ อุทายิํ เอตทโวจ “กิสฺมิํ นุ โข ภนฺเต สติ อรหนฺโต สุขทุกฺขํ ปญฺญเปนฺติ, กิสฺมิํ อสติ อรหนฺโต สุขทุกฺขํ น ปญฺญเปนฺตี”ติ.\csromanspacer{} จกฺขุสฺมิํ โข ภคินิ สติ อรหนฺโต สุขทุกฺขํ ปญฺญเปนฺติ, จกฺขุสฺมิํ อสติ อรหนฺโต สุขทุกฺขํ น ปญฺญเปนฺติ ฯเปฯ.\csromanspacer{} ชิวฺหาย สติ อรหนฺโต สุขทุกฺขํ ปญฺญเปนฺติ, ชิวฺหาย อสติ อรหนฺโต สุขทุกฺขํ น ปญฺญเปนฺติ ฯเปฯ.\csromanspacer{} มนสฺมิํ สติ อรหนฺโต สุขทุกฺขํ ปญฺญเปนฺติ, มนสฺมิํ อสติ อรหนฺโต สุขทุกฺขํ น ปญฺญเปนฺตีติ.}

\prose{เอวํ วุตฺเต เวรหจฺจานิโคตฺตา พฺราหฺมณี อายสฺมนฺตํ อุทายิํ เอตทโวจ “อภิกฺกนฺตํ ภนฺเต, อภิกฺกนฺตํ ภนฺเต, เสยฺยถาปิ ภนฺเต นิกฺกุชฺชิตํ วา}

\gambhira{สฬายตนวคฺคสํยุตฺตปาฬิ}
\prose{\csromanfolio{340}อุกฺกุชฺเชยฺย, ปฏิจฺฉนฺนํ วา วิวเรยฺย, มูฬฺหสฺส วา มคฺคํ อาจิกฺเขยฺย, อนฺธกาเร วา เตลปชฺโชตํ ธาเรยฺย ‘จกฺขุมนฺโต รูปานิ ทกฺขนฺตี’ติ, เอวเมวํ อยฺเยน อุทายินา อเนกปริยาเยน ธมฺโม ปกาสิโต, เอสาหํ อยฺย อุทายิ ตํ ภควนฺตํ สรณํ คจฺฉามิ, ธมฺมญฺจ ภิกฺขุสํฆญฺจ, อุปาสิกํ มํ อยฺโย อุทายี ธาเรตุ อชฺชตคฺเค ปาณุเปตํ สรณํ คตนฺ”ติ.\csromanspacer{}\csromanspacer{}ทสมํ.\csromanspacer{} คหปติวคฺโค ตติโย.}

\titlehead{\textbf{ตสฺสุทฺทานํ}}
\csromangathameasure{เวสาลี วชฺชิ นาฬนฺทา, ภารทฺวาช โสโณ จ โฆสิโต. \\ หาลิทฺทิโก นกุลปิตา, โลหิจฺโจ เวรหจฺจานีติ.}
\csromangathagroup{\csromangathastanza{เวสาลี วชฺชิ นาฬนฺทา, ภารทฺวาช โสโณ จ โฆสิโต. \\ หาลิทฺทิโก นกุลปิตา, โลหิจฺโจ เวรหจฺจานีติ.}}

\csromanmatikaanchor{mtk.173}
\csromantocmarknopage{section}{(14) 4. เทวทหวคฺค}{mtk.173}
\bookmark[dest={mtk.173},level=1]{(14) 4. เทวทหวคฺค}
\csromanheader{1}{\textbf{(14) 4. เทวทหวคฺค}}
\csromanmatikaanchor{mtk.174}
\csromantocmarkref{subsection}{1. เทวทหสุตฺต}{mtk.174}
\bookmark[dest={mtk.174},level=2]{1. เทวทหสุตฺต}
\csromanheader{2}{1. เทวทหสุตฺต}
\proseitem{134}{เอกํ สมยํ ภควา สกฺเกสุ วิหรติ เทวทหํ นาม สกฺยานํ นิคโม, ตตฺร โข ภควา ภิกฺขู อามนฺเตสิ นาหํ ภิกฺขเว สพฺเพสํเยว ภิกฺขูนํ ฉสุ ผสฺสายตเนสุ อปฺปมาเทน กรณียนฺติ วทามิ, น จ ปนาหํ ภิกฺขเว สพฺเพสํเยว ภิกฺขูนํ ฉสุ ผสฺสายตเนสุ นปฺปมาเทน กรณียนฺติ วทามิ.\csromanspacer{} เย เต ภิกฺขเว ภิกฺขู อรหนฺโต ขีณาสวา วุสิตวนฺโต กตกรณียา โอหิตภารา อนุปฺปตฺตสทตฺถา ปริกฺขีณภวสํโยชนา สมฺมทญฺญา วิมุตฺตา, เตสาหํ ภิกฺขเว ภิกฺขูนํ ฉสุ ผสฺสายตเนสุ นาปฺปมาเทน กรณียนฺติ วทามิ.\csromanspacer{} ตํ กิสฺส เหตุ, กตํ เตสํ อปฺปมาเทน, อภพฺพา เต ปมชฺชิตุํ.\csromanspacer{} เย จ โข เต ภิกฺขเว ภิกฺขู เสกฺขา\footnote{เสขา (สี, สฺยา, กํ, อิ, ก)} อปฺปตฺตมานสา อนุตฺตรํ โยคกฺเขมํ ปตฺถยมานา วิหรนฺติ, เตสาหํ ภิกฺขเว ภิกฺขูนํ ฉสุ ผสฺสายตเนสุ อปฺปมาเทน กรณียนฺติ วทามิ.\csromanspacer{} ตํ กิสฺส เหตุ, สนฺติ ภิกฺขเว จกฺขุวิญฺเญยฺยา รูปา มโนรมาปิ อมโนรมาปิ, ตฺยาสฺส ผุสฺส ผุสฺส จิตฺตํ น ปริยาทาย ติฏฺฐนฺติ, เจตโส อปริยาทานา อารทฺธํ โหติ วีริยํ อสลฺลีนํ, อุปฏฺฐิตา สติ อสมฺมุฏฺฐา\footnote{อปมฺมุฏฺฐา (สี), อปฺปมุฏฺฐา (สฺยา, กํ)}, ปสฺสทฺโธ กาโย อสารทฺโธ, สมาหิตํ จิตฺตํ เอกคฺคํ.\csromanspacer{} อิมํ ขฺวาหํ ภิกฺขเว}

\vspace{\tipitakaparskip}
\csromancenter{สฬายตนสํยุตฺต อปฺปมาทผลํ สมฺปสฺสมาโน เตสํ ภิกฺขูนํ ฉสุ ผสฺสายตเนสุ อปฺปมาเทน กรณียนฺติ วทามิ ฯเปฯ.\csromanspacer{} สนฺติ ภิกฺขเว มโนวิญฺเญยฺยา ธมฺมา มโนรมาปิ อมโนรมาปิ, ตฺยาสฺส ผุสฺส ผุสฺส จิตฺตํ น ปริยาทาย ติฏฺฐนฺติ, เจตโส อปริยาทานา อารทฺธํ โหติ วีริยํ อสลฺลีนํ, อุปฏฺฐิตา สติ อสมฺมุฏฺฐา, ปสฺสทฺโธ กาโย อสารทฺโธ, สมาหิตํ จิตฺตํ เอกคฺคํ.\csromanspacer{} อิมํ ขฺวาหํ ภิกฺขเว อปฺปมาทผลํ สมฺปสฺสมาโน เตสํ ภิกฺขูนํ ฉสุ\footnote{ฉสฺสุ (สี)} ผสฺสายตเนสุ อปฺปมาเทน กรณียนฺติ วทามีติ.\csromanspacer{}\csromanspacer{}ปฐมํ.}
\csromanmatikaanchor{mtk.175}
\csromantocmarkref{subsection}{2. ขณสุตฺต}{mtk.175}
\bookmark[dest={mtk.175},level=2]{2. ขณสุตฺต}
\csromanheader{2}{2. ขณสุตฺต}
\csromanmatikaanchor{mtk.176}
\csromantocmarkref{subsection}{3-4. รูปารามสุตฺต}{mtk.176}
\bookmark[dest={mtk.176},level=2]{3-4. รูปารามสุตฺต}
\proseitem{135}{\csromanfolio{341}ลาภา โว ภิกฺขเว, สุลทฺธํ โว ภิกฺขเว, ขโณ โว ปฏิลทฺโธ พฺรหฺมจริยวาสาย.\csromanspacer{} ทิฏฺฐา มยา ภิกฺขเว ฉผสฺสายตนิกา นาม นิรยา, ตตฺถ ยํ กิญฺจิ จกฺขุนา รูปํ ปสฺสติ, อนิฏฺฐรูปํเยว ปสฺสติ, โน อิฏฺฐรูปํ, อกนฺตรูปํเยว ปสฺสติ, โน กนฺตรูปํ, อมนาปรูปํเยว ปสฺสติ, โน มนาปรูปํ.\csromanspacer{} ยํ กิญฺจิ โสเตน สทฺทํ สุณาติ ฯเปฯ.\csromanspacer{} ยํ กิญฺจิ ฆาเนน คนฺธํ ฆายติ.\csromanspacer{} ยํ กิญฺจิ ชิวฺหาย รสํ สายติ.\csromanspacer{} ยํ กิญฺจิ กาเยน โผฏฺฐพฺพํ ผุสติ.\csromanspacer{} ยํ กิญฺจิ มนสา ธมฺมํ วิชานาติ, อนิฏฺฐรูปํเยว วิชานาติ, โน อิฏฺฐรูปํ, อกนฺตรูปํเยว วิชานาติ, โน กนฺตรูปํ, อมนาปรูปํเยว วิชานาติ, โน มนาปรูปํ, ลาภา โว ภิกฺขเว, สุลทฺธํ โว ภิกฺขเว, ขโณ โว ปฏิลทฺโธ พฺรหฺมจริยวาสาย.\csromanspacer{} ทิฏฺฐา มยา ภิกฺขเว ฉผสฺสายตนิกา นาม สคฺคา, ตตฺถ ยํ กิญฺจิ จกฺขุนา รูปํ ปสฺสติ, อิฏฺฐรูปํเยว ปสฺสติ, โน อนิฏฺฐรูปํ, กนฺตรูปํเยว ปสฺสติ, โน อกนฺตรูปํ มนาปรูปํเยว ปสฺสติ, โน อมนาปรูปํ ฯเปฯ.\csromanspacer{} ยํ กิญฺจิ ชิวฺหาย รสํ สายติ ฯเปฯ.\csromanspacer{} ยํ กิญฺจิ มนสา ธมฺมํ วิชานาติ, อิฏฺฐรูปํเยว วิชานาติ, โน อนิฏฺฐรูปํ, กนฺตรูปํเยว วิชานาติ, โน อกนฺตรูปํ.\csromanspacer{} มนาปรูปํเยว วิชานาติ, โน อมนาปรูปํ.\csromanspacer{} ลาภา โว ภิกฺขเว, สุลทฺธํ โว ภิกฺขเว, ขโณ โว ปฏิลทฺโธ พฺรหฺมจริยวาสายาติ.\csromanspacer{}\csromanspacer{}ทุติยํ.}

\csromanheader{2}{3. ปฐมรูปารามสุตฺต}
\proseitem{136}{รูปารามา ภิกฺขเว เทวมนุสฺสา รูปรตา รูปสมฺมุทิตา, รูปวิปริณามวิราคนิโรธา ทุกฺขา ภิกฺขเว เทวมนุสฺสา วิหรนฺติ.\csromanspacer{} สทฺทารามา}

\gambhira{สฬายตนวคฺคสํยุตฺตปาฬิ}
\prose{\csromanfolio{342}ภิกฺขเว เทวมนุสฺสา สทฺทรตา สทฺทสมฺมุทิตา, สทฺทวิปริณามวิราคนิโรธา ทุกฺขา ภิกฺขเว เทวมนุสฺสา วิหรนฺติ.\csromanspacer{} คนฺธารามา.\csromanspacer{} รสารามา.\csromanspacer{} โผฏฺฐพฺพารามา.\csromanspacer{} ธมฺมารามา ภิกฺขเว เทวมนุสฺสา ธมฺมรตา ธมฺมสมฺมุทิตา, ธมฺมวิปริณามวิราคนิโรธา ทุกฺขา ภิกฺขเว เทวมนุสฺสา วิหรนฺติ.\csromanspacer{} ตถาคโต จ โข ภิกฺขเว อรหํ สมฺมาสมฺพุทฺโธ รูปานํ สมุทยญฺจ อตฺถงฺคมญฺจ อสฺสาทญฺจ อาทีนวญฺจ นิสฺสรณญฺจ ยถาภูตํ วิทิตฺวา น รูปาราโม น รูปรโต น รูปสมฺมุทิโต, รูปวิปริณามวิราคนิโรธา สุโข ภิกฺขเว ตถาคโต วิหรติ.\csromanspacer{} สทฺทานํ.\csromanspacer{} คนฺธานํ.\csromanspacer{} รสานํ.\csromanspacer{} โผฏฺฐพฺพานํ.\csromanspacer{} ธมฺมานํ สมุทยญฺจ อตฺถงฺคมญฺจ อสฺสาทญฺจ อาทีนวญฺจ นิสฺสรณญฺจ ยถาภูตํ วิทิตฺวา น ธมฺมาราโม น ธมฺมรโต น ธมฺมสมฺมุทิโต, ธมฺมวิปริณามวิราคนิโรธา สุโข ภิกฺขเว ตถาคโต วิหรติ.\csromanspacer{} อิทมโวจ ภควา, อิทํ วตฺวาน สุคโต อถาปรํ เอตทโวจ สตฺถา\csromandash{}}

\csromangathameasure{รูปา สทฺทา รสา คนฺธา, ผสฺสา ธมฺมา จ เกวลา. \\ อิฏฺฐา กนฺตา มนาปา จ, ยาวตตฺถีติ วุจฺจติ. \\ สเทวกสฺส โลกสฺส, เอเต โว สุขสมฺมตา. \\ ยตฺถ เจเต นิรุชฺฌนฺติ, ตํ เตสํ ทุกฺขสมฺมตํ. \\ สุขํ ทิฏฺฐมริเยภิ, สกฺกายสฺส นิโรธนํ. \\ ปจฺจนีกมิทํ โหติ, สพฺพโลเกน ปสฺสตํ. \\ ยํ ปเร สุขโต อาหุ, ตทริยา อาหุ ทุกฺขโต. \\ ยํ ปเร ทุกฺขโต อาหุ, ตทริยา สุขโต วิทู. \\ ปสฺส ธมฺมํ ทุราชานํ, สมฺมูเฬฺหตฺถ อวิทฺทสุ. \\ นิวุตานํ ตโม โหติ, อนฺธกาโร อปสฺสตํ. \\ สตญฺจ วิวฏํ โหติ, อาโลโก ปสฺสตามิว. \\ สนฺติเก น วิชานนฺติ, มคฺคา ธมฺมสฺส อโกวิทา. \\ ภวราคปเรเตภิ, ภวราคานุสารีภิ. \\ มารเธยฺยานุปนฺเนหิ, นายํ ธมฺโม สุสมฺพุโธ.}
\csromangathagroup{\csromangathastanza{รูปา สทฺทา รสา คนฺธา, ผสฺสา ธมฺมา จ เกวลา. \\ อิฏฺฐา กนฺตา มนาปา จ, ยาวตตฺถีติ วุจฺจติ.}\csromangathastanza{สเทวกสฺส โลกสฺส, เอเต โว สุขสมฺมตา. \\ ยตฺถ เจเต นิรุชฺฌนฺติ, ตํ เตสํ ทุกฺขสมฺมตํ.}\csromangathastanza{สุขํ\footnote{สุขนฺติ (สี)} ทิฏฺฐมริเยภิ, สกฺกายสฺส นิโรธนํ. \\ ปจฺจนีกมิทํ โหติ, สพฺพโลเกน ปสฺสตํ.}\csromangathastanza{ยํ ปเร สุขโต อาหุ, ตทริยา อาหุ ทุกฺขโต. \\ ยํ ปเร ทุกฺขโต อาหุ, ตทริยา สุขโต วิทู.}\csromangathastanza{ปสฺส ธมฺมํ ทุราชานํ, สมฺมูเฬฺหตฺถ อวิทฺทสุ. \\ นิวุตานํ ตโม โหติ, อนฺธกาโร อปสฺสตํ.}\csromangathastanza{สตญฺจ วิวฏํ โหติ, อาโลโก ปสฺสตามิว. \\ สนฺติเก น วิชานนฺติ, มคฺคา\footnote{มคา (สี)} ธมฺมสฺส อโกวิทา.}\csromangathastanza{ภวราคปเรเตภิ, ภวราคานุสารีภิ\footnote{ภวโสตานุสาริภิ (สฺยา, กํ, อิ), ภวโสตานุสาริหิ (สี)}. \\ มารเธยฺยานุปนฺเนหิ, นายํ ธมฺโม สุสมฺพุโธ.}}

\csromanheader{2}{1. สฬายตนสํยุตฺต}
\csromanmatikaanchor{mtk.177}
\csromantocmarkref{subsection}{5-6. นตุมฺหากสุตฺต}{mtk.177}
\bookmark[dest={mtk.177},level=2]{5-6. นตุมฺหากสุตฺต}
\csromangathameasure{โก นุ อญฺญตฺร มริเยภิ, ปทํ สมฺพุทฺธุมรหติ. \\ ยํ ปทํ สมฺมทญฺญาย, ปรินิพฺพนฺติ อนาสวาติ. ตติยํ.}
\csromangathagroup{\csromangathastanza{\csromanfolio{343}โก นุ อญฺญตฺร มริเยภิ, ปทํ สมฺพุทฺธุมรหติ. \\ ยํ ปทํ สมฺมทญฺญาย, ปรินิพฺพนฺติ อนาสวาติ.\csromanspacer{} ตติยํ.}}

\csromanheader{2}{4. ทุติยรูปารามสุตฺต}
\proseitem{137}{รูปารามา ภิกฺขเว เทวมนุสฺสา รูปรตา รูปสมฺมุทิตา, รูปวิปริณามวิราคนิโรธา ทุกฺขา ภิกฺขเว เทวมนุสฺสา วิหรนฺติ.\csromanspacer{} สทฺทารามา.\csromanspacer{} คนฺธารามา.\csromanspacer{} รสารามา.\csromanspacer{} โผฏฺฐพฺพารามา.\csromanspacer{} ธมฺมารามา ภิกฺขเว เทวมนุสฺสา ธมฺมรตา ธมฺมสมฺมุทิตา, ธมฺมวิปริณามวิราคนิโรธา ทุกฺขา ภิกฺขเว เทวมนุสฺสา วิหรนฺติ.\csromanspacer{} ตถาคโต จ ภิกฺขเว อรหํ สมฺมาสมฺพุทฺโธ รูปานํ สมุทยญฺจ อตฺถงฺคมญฺจ อสฺสาทญฺจ อาทีนวญฺจ นิสฺสรณญฺจ ยถาภูตํ วิทิตฺวา น รูปาราโม น รูปรโต น รูปสมฺมุทิโต, รูปวิปริณามวิราคนิโรธา สุโข ภิกฺขเว ตถาคโต วิหรติ.\csromanspacer{} สทฺทานํ.\csromanspacer{} คนฺธานํ.\csromanspacer{} รสานํ.\csromanspacer{} โผฏฺฐพฺพานํ.\csromanspacer{} ธมฺมานํ สมุทยญฺจ อตฺถงฺคมญฺจ อสฺสาทญฺจ อาทีนวญฺจ นิสฺสรณญฺจ ยถาภูตํ วิทิตฺวา น ธมฺมาราโม น ธมฺมรโต น ธมฺมสมฺมุทิโต, ธมฺมวิปริณามวิราคนิโรธา สุโข ภิกฺขเว ตถาคโต วิหรตีติ.\csromanspacer{}\csromanspacer{}จตุตฺถํ.}

\csromanheader{2}{5. ปฐมนตุมฺหากสุตฺต}
\proseitem{138}{ยํ ภิกฺขเว น ตุมฺหากํ, ตํ ปชหถ, ตํ โว ปหีนํ หิตาย สุขาย ภวิสฺสติ.\csromanspacer{} กิญฺจ ภิกฺขเว น ตุมฺหากํ.\csromanspacer{} จกฺขุ ภิกฺขเว น ตุมฺหากํ, ตํ ปชหถ, ตํ โว ปหีนํ หิตาย สุขาย ภวิสฺสติ ฯเปฯ.\csromanspacer{} ชิวฺหา น ตุมฺหากํ, ตํ ปชหถ, สา โว ปหีนา หิตาย สุขาย ภวิสฺสติ ฯเปฯ.\csromanspacer{} มโน น ตุมฺหากํ, ตํ ปชหถ, โส โว ปหีโน หิตาย สุขาย ภวิสฺสติ.\csromanspacer{} เสยฺยถาปิ ภิกฺขเว ยํ อิมสฺมิํ เชตวเน ติณกฏฺฐสาขาปลาสํ, ตํ ชโน หเรยฺย วา ฑเหยฺย วา ยถาปจฺจยํ วา กเรยฺย, อปิ นุ ตุมฺหากํ เอวมสฺส “อเมฺห ชโน หรติ วา ฑหติ วา ยถาปจฺจยํ วา กโรตี”ติ.\csromanspacer{} โน เหตํ ภนฺเต.\csromanspacer{} ตํ กิสฺส เหตุ.\csromanspacer{} น หิ โน เอตํ ภนฺเต อตฺตา วา อตฺตนิยํ วาติ.\csromanspacer{} เอวเมว โข ภิกฺขเว จกฺขุ น ตุมฺหากํ, ตํ ปชหถ, ตํ โว ปหีนํ หิตาย สุขาย ภวิสฺสติ ฯเปฯ.\csromanspacer{} ชิวฺหา น ตุมฺหากํ, ตํ ปชหถ, สา โว ปหีนา หิตาย สุขาย ภวิสฺสติ ฯเปฯ.\csromanspacer{} มโน น}

\gambhira{สฬายตนวคฺคสํยุตฺตปาฬิ}
\csromanmatikaanchor{mtk.178}
\csromantocmarkref{subsection}{7-9. อชฺชตฺต-อนิจฺจเหตุสุตฺตาทิ}{mtk.178}
\bookmark[dest={mtk.178},level=2]{7-9. อชฺชตฺต-อนิจฺจเหตุสุตฺตาทิ}
\prose{\csromanfolio{344}ตุมฺหากํ, ตํ ปชหถ, โส โว ปหีโน หิตาย สุขาย ภวิสฺสตีติ.\csromanspacer{}\csromanspacer{}ปญฺจมํ.}

\csromanheader{2}{6. ทุติยนตุมฺหากสุตฺต}
\proseitem{139}{ยํ ภิกฺขเว น ตุมฺหากํ, ตํ ปชหถ, ตํ โว ปหีนํ หิตาย สุขาย ภวิสฺสติ.\csromanspacer{} กิญฺจ ภิกฺขเว น ตุมฺหากํ.\csromanspacer{} รูปา ภิกฺขเว น ตุมฺหากํ, เต ปชหถ, เต โว ปหีนา หิตาย สุขาย ภวิสฺสนฺติ.\csromanspacer{} สทฺทา.\csromanspacer{} คนฺธา.\csromanspacer{} รสา.\csromanspacer{} โผฏฺฐพฺพา.\csromanspacer{} ธมฺมา น ตุมฺหากํ, เต ปชหถ, เต โว ปหีนา หิตาย สุขาย ภวิสฺสนฺติ.\csromanspacer{} เสยฺยถาปิ ภิกฺขเว ยํ อิมสฺมิํ เชตวเน ติณกฏฺฐ ฯเปฯ เอวเมว โข ภิกฺขเว รูปา น ตุมฺหากํ, เต ปชหถ, เต โว ปหีนา หิตาย สุขาย ภวิสฺสนฺตีติ.\csromanspacer{}\csromanspacer{}ฉฏฺฐํ.}

\csromanheader{2}{7. อชฺฌตฺต-อนิจฺจเหตุสุตฺต}
\proseitem{140}{จกฺขุํ ภิกฺขเว อนิจฺจํ, โยปิ เหตุ โยปิ ปจฺจโย จกฺขุสฺส อุปฺปาทาย, โสปิ อนิจฺโจ, อนิจฺจสมฺภูตํ ภิกฺขเว จกฺขุ กุโต นิจฺจํ ภวิสฺสติ ฯเปฯ.\csromanspacer{} ชิวฺหา อนิจฺจา, โยปิ เหตุ โยปิ ปจฺจโย ชิวฺหาย อุปฺปาทาย, โสปิ อนิจฺโจ, อนิจฺจสมฺภูตา ภิกฺขเว ชิวฺหา กุโต นิจฺจา ภวิสฺสติ ฯเปฯ.\csromanspacer{} มโน อนิจฺโจ, โยปิ ภิกฺขเว เหตุ โยปิ ปจฺจโย มนสฺส อุปฺปาทาย, โสปิ อนิจฺโจ, อนิจฺจสมฺภูโต ภิกฺขเว มโน กุโต นิจฺโจ ภวิสฺสติ.\csromanspacer{} เอวํ ปสฺสํ ภิกฺขเว สุตวา อริยสาวโก จกฺขุสฺมิมฺปิ นิพฺพินฺทติ ฯเปฯ ชิวฺหายปิ นิพฺพินฺทติ ฯเปฯ มนสฺมิมฺปิ นิพฺพินฺทติ, นิพฺพินฺทํ วิรชฺชติ, วิราคา วิมุจฺจติ, วิมุตฺตสฺมิํ “วิมุตฺตมฺ”อิติ ญาณํ โหติ, “ขีณา ชาติ, วุสิตํ พฺรหฺมจริยํ, กตํ กรณียํ, นาปรํ อิตฺถตฺตายา”ติ ปชานาตีติ.\csromanspacer{}\csromanspacer{}สตฺตมํ.}

\csromanheader{2}{8. อชฺฌตฺตทุกฺขเหตุสุตฺต}
\proseitem{141}{จกฺขุํ ภิกฺขเว ทุกฺขํ, โยปิ เหตุ โยปิ ปจฺจโย จกฺขุสฺส อุปฺปาทาย, โสปิ ทุกฺโข, ทุกฺขสมฺภูตํ ภิกฺขเว จกฺขุ กุโต สุขํ ภวิสฺสติ ฯเปฯ.\csromanspacer{} ชิวฺหา ทุกฺขา, โยปิ เหตุ โยปิ ปจฺจโย ชิวฺหาย อุปฺปาทาย, โสปิ ทุกฺโข, ทุกฺขสมฺภูตา ภิกฺขเว ชิวฺหา กุโต สุขา ภวิสฺสติ ฯเปฯ.\csromanspacer{} มโน ทุกฺโข, โยปิ เหตุ โยปิ ปจฺจโย มนสฺส อุปฺปาทาย, โสปิ ทุกฺโข, ทุกฺขสมฺภูโต ภิกฺขเว มโน กุโต สุโข ภวิสฺสติ.\csromanspacer{} เอวํ ปสฺสํ ฯเปฯ นาปรํ อิตฺถตฺตายาติ ปชานาตีติ.\csromanspacer{}\csromanspacer{}อฏฺฐมํ.}

\csromanheader{2}{1. สฬายตนสํยุตฺต \textbf{9. อชฺฌตฺตานตฺตเหตุสุตฺต}}
\csromanmatikaanchor{mtk.179}
\csromantocmarkref{subsection}{10-12. พาหิรานิจฺจเหตุสุตฺตาทิ}{mtk.179}
\bookmark[dest={mtk.179},level=2]{10-12. พาหิรานิจฺจเหตุสุตฺตาทิ}
\proseitem{142}{\csromanfolio{345}จกฺขุํ ภิกฺขเว อนตฺตา, โยปิ เหตุ โยปิ ปจฺจโย จกฺขุสฺส อุปฺปาทาย, โสปิ อนตฺตา, อนตฺตสมฺภูตํ ภิกฺขเว จกฺขุ กุโต อตฺตา ภวิสฺสติ ฯเปฯ.\csromanspacer{} ชิวฺหา อนตฺตา, โยปิ เหตุ โยปิ ปจฺจโย ชิวฺหาย อุปฺปาทาย, โสปิ อนตฺตา, อนตฺตสมฺภูตา ภิกฺขเว ชิวฺหา กุโต อตฺตา ภวิสฺสติ ฯเปฯ.\csromanspacer{} มโน อนตฺตา, โยปิ เหตุ โยปิ ปจฺจโย มนสฺส อุปฺปาทาย, โสปิ อนตฺตา, อนตฺตสมฺภูโต ภิกฺขเว มโน กุโต อตฺตา ภวิสฺสติ.\csromanspacer{} เอวํ ปสฺสํ ฯเปฯ นาปรํ อิตฺถตฺตายาติ ปชานาตีติ.\csromanspacer{}\csromanspacer{}นวมํ.}

\csromanheader{2}{10. พาหิรานิจฺจเหตุสุตฺต}
\proseitem{143}{รูปา ภิกฺขเว อนิจฺจา, โยปิ เหตุ โยปิ ปจฺจโย รูปานํ อุปฺปาทาย, โสปิ อนิจฺโจ, อนิจฺจสมฺภูตา ภิกฺขเว รูปา กุโตนิจฺจา ภวิสฺสนฺติ.\csromanspacer{} สทฺทา.\csromanspacer{} คนฺธา.\csromanspacer{} รสา.\csromanspacer{} โผฏฺฐพฺพา.\csromanspacer{} ธมฺมา อนิจฺจา, โยปิ เหตุ โยปิ ปจฺจโย ธมฺมานํ อุปฺปาทาย, โสปิ อนิจฺโจ, อนิจฺจสมฺภูตา ภิกฺขเว ธมฺมา กุโต นิจฺจา ภวิสฺสนฺติ.\csromanspacer{} เอวํ ปสฺสํ ฯเปฯ นาปรํ อิตฺถตฺตายาติ ปชานาตีติ.\csromanspacer{}\csromanspacer{}ทสมํ.}

\csromanheader{2}{11. พาหิรทุกฺขเหตุสุตฺต}
\proseitem{144}{รูปา ภิกฺขเว ทุกฺขา, โยปิ เหตุ โยปิ ปจฺจโย รูปานํ อุปฺปาทาย, โสปิ ทุกฺโข, ทุกฺขสมฺภูตา ภิกฺขเว รูปา กุโต สุขา ภวิสฺสนฺติ.\csromanspacer{} สทฺทา.\csromanspacer{} คนฺธา.\csromanspacer{} รสา.\csromanspacer{} โผฏฺฐพฺพา.\csromanspacer{} ธมฺมา ทุกฺขา, โยปิ เหตุ โยปิ ปจฺจโย ธมฺมานํ อุปฺปาทาย, โสปิ ทุกฺโข, ทุกฺขสมฺภูตา ภิกฺขเว ธมฺมา กุโต สุขา ภวิสฺสนฺติ.\csromanspacer{} เอวํ ปสฺสํ ฯเปฯ นาปรํ อิตฺถตฺตายาติ ปชานาตีติ.\csromanspacer{}\csromanspacer{}เอกาทสมํ.}

\csromanheader{2}{12. พาหิรานตฺตเหตุสุตฺต}
\proseitem{145}{รูปา ภิกฺขเว อนตฺตา, โยปิ เหตุ โยปิ ปจฺจโย รูปานํ อุปฺปาทาย, โสปิ อนตฺตา, อนตฺตสมฺภูตา ภิกฺขเว รูปา กุโต อตฺตา ภวิสฺสนฺติ.\csromanspacer{} สทฺทา.\csromanspacer{} คนฺธา.\csromanspacer{} รสา.\csromanspacer{} โผฏฺฐพฺพา.\csromanspacer{} ธมฺมา อนตฺตา, โยปิ เหตุ โยปิ ปจฺจโย ธมฺมานํ อุปฺปาทาย, โสปิ อนตฺตา, อนตฺตสมฺภูตา ภิกฺขเว}

\gambhira{สฬายตนวคฺคสํยุตฺตปาฬิ}
\prose{\csromanfolio{346}ธมฺมา กุโต อตฺตา ภวิสฺสนฺติ.\csromanspacer{} เอวํ ปสฺสํ ภิกฺขเว สุตวา อริยสาวโก รูเปสุปิ นิพฺพินฺทติ.\csromanspacer{} สทฺเทสุปิ.\csromanspacer{} คนฺเธสุปิ.\csromanspacer{} รเสสุปิ.\csromanspacer{} โผฏฺฐพฺเพสุปิ.\csromanspacer{} ธมฺเมสุปิ นิพฺพินฺทติ, นิพฺพินฺทํ วิรชฺชติ, วิราคา วิมุจฺจติ, วิมุตฺตสฺมิํ “วิมุตฺตมฺ”อิติ ญาณํ โหติ, “ขีณา ชาติ, วุสิตํ พฺรหฺมจริยํ, กตํ กรณียํ, นาปรํ อิตฺถตฺตายา”ติ ปชานาตีติ.\csromanspacer{}\csromanspacer{}ทฺวาทสมํ.\csromanspacer{} เทวทหวคฺโค จตุตฺโถ.}

\titlehead{\textbf{ตสฺสุทฺทานํ}}
\csromangathameasure{เทวทโห ขโณ รูปา, เทฺว นตุมฺหากเมว จ. \\ เหตุนาปิ ตโย วุตฺตา, ทุเว อชฺฌตฺตพาหิราติ.}
\csromangathagroup{\csromangathastanza{เทวทโห ขโณ รูปา, เทฺว นตุมฺหากเมว จ. \\ เหตุนาปิ ตโย วุตฺตา, ทุเว อชฺฌตฺตพาหิราติ.}}

\csromanmatikaanchor{mtk.180}
\csromantocmarknopage{section}{(15) 5. นวปุราณวคฺค}{mtk.180}
\bookmark[dest={mtk.180},level=1]{(15) 5. นวปุราณวคฺค}
\csromanheader{1}{\textbf{(15) 5. นวปุราณวคฺค}}
\csromanmatikaanchor{mtk.181}
\csromantocmarkref{subsection}{1. กมฺมนิโรธสุตฺต}{mtk.181}
\bookmark[dest={mtk.181},level=2]{1. กมฺมนิโรธสุตฺต}
\csromanheader{2}{1. กมฺมนิโรธสุตฺต}
\proseitem{146}{นวปุราณานิ ภิกฺขเว กมฺมานิ เทเสสฺสามิ กมฺมนิโรธํ กมฺมนิโรธคามินิํ จ ปฏิปทํ, ตํ สุณาถ สาธุกํ มนสิ กโรถ, ภาสิสฺสามีติ.\csromanspacer{} กตมญฺจ ภิกฺขเว ปุราณกมฺมํ.\csromanspacer{} จกฺขุ ภิกฺขเว ปุราณกมฺมํ อภิสงฺขตํ อภิสญฺเจตยิตํ เวทนิยํ ทฏฺฐพฺพํ ฯเปฯ.\csromanspacer{} ชิวฺหา ปุราณกมฺมา อภิสงฺขตา อภิสญฺเจตยิตา เวทนิยา ทฏฺฐพฺพา ฯเปฯ.\csromanspacer{} มโน ปุราณกมฺโม อภิสงฺขโต อภิสญฺเจตยิโต เวทนิโย ทฏฺฐพฺโพ.\csromanspacer{} อิทํ วุจฺจติ ภิกฺขเว ปุราณกมฺมํ.\csromanspacer{} กตมญฺจ ภิกฺขเว นวกมฺมํ.\csromanspacer{} ยํ โข ภิกฺขเว เอตรหิ กมฺมํ กโรติ กาเยน วาจาย มนสา, อิทํ วุจฺจติ ภิกฺขเว นวกมฺมํ.\csromanspacer{} กตโม จ ภิกฺขเว กมฺมนิโรโธ.\csromanspacer{} โย โข ภิกฺขเว กายกมฺมวจีกมฺมมโนกมฺมสฺส นิโรธา วิมุตฺติํ ผุสติ, อยํ วุจฺจติ ภิกฺขเว กมฺมนิโรโธ.\csromanspacer{} กตมา จ ภิกฺขเว กมฺมนิโรธคามินี ปฏิปทา.\csromanspacer{} อยเมว อริโย อฏฺฐงฺคิโก มคฺโค, เสยฺยถิทํ, สมฺมาทิฏฺฐิ สมฺมาสงฺกปฺโป สมฺมาวาจา สมฺมากมฺมนฺโต สมฺมา-อาชีโว สมฺมาวายาโม สมฺมาสติ สมฺมาสมาธิ, อยํ วุจฺจติ ภิกฺขเว กมฺมนิโรธคามินี ปฏิปทา.\csromanspacer{} อิติ โข ภิกฺขเว เทสิตํ มยา ปุราณกมฺมํ,}

\vspace{\tipitakaparskip}
\csromancenter{สฬายตนสํยุตฺต เทสิตํ นวกมฺมํ, เทสิโต กมฺมนิโรโธ, เทสิตา กมฺมนิโรธคามินี ปฏิปทา, ยํ โข ภิกฺขเว สตฺถารา กรณียํ สาวกานํ หิเตสินา อนุกมฺปเกน อนุกมฺปํ อุปาทาย, กตํ โว ตํ มยา.\csromanspacer{} เอตานิ ภิกฺขเว รุกฺขมูลานิ, เอตานิ สุญฺญาคารานิ, ฌายถ ภิกฺขเว, มา ปมาทตฺถ, มา ปจฺฉาวิปฺปฏิสาริโน อหุวตฺถ, อยํ โว อมฺหากํ อนุสาสนีติ.\csromanspacer{}\csromanspacer{}ปฐมํ.}
\csromanheader{2}{2. อนิจฺจนิพฺพานสปฺปายสุตฺต}
\csromanmatikaanchor{mtk.182}
\csromantocmarkref{subsection}{2-4. อนิจฺจนิพฺพานสปฺปายสุตฺตาทิ}{mtk.182}
\bookmark[dest={mtk.182},level=2]{2-4. อนิจฺจนิพฺพานสปฺปายสุตฺตาทิ}
\proseitem{147}{\csromanfolio{347}นิพฺพานสปฺปายํ โว ภิกฺขเว ปฏิปทํ เทเสสฺสามิ, ตํ สุณาถ ฯเปฯ.\csromanspacer{} กตมา จ สา ภิกฺขเว นิพฺพานสปฺปายา ปฏิปทา.\csromanspacer{} อิธ ภิกฺขเว ภิกฺขุ จกฺขุํ อนิจฺจนฺติ ปสฺสติ, รูปา อนิจฺจาติ ปสฺสติ, จกฺขุวิญฺญาณํ อนิจฺจนฺติ ปสฺสติ, จกฺขุสมฺผสฺโส อนิจฺโจติ ปสฺสติ, ยมฺปิทํ จกฺขุสมฺผสฺสปจฺจยา อุปฺปชฺชติ เวทยิตํ สุขํ วา ทุกฺขํ วา อทุกฺขมสุขํ วา, ตมฺปิ อนิจฺจนฺติ ปสฺสติ ฯเปฯ.\csromanspacer{} ชิวฺหา อนิจฺจาติ ปสฺสติ, รสา อนิจฺจาติ ปสฺสติ, ชิวฺหาวิญฺญาณํ อนิจฺจนฺติ ปสฺสติ, ชิวฺหาสมฺผสฺโส อนิจฺโจติ ปสฺสติ, ยมฺปิทํ ชิวฺหาสมฺผสฺสปจฺจยา อุปฺปชฺชติ เวทยิตํ สุขํ วา ทุกฺขํ วา อทุกฺขมสุขํ วา, ตมฺปิ อนิจฺจนฺติ ปสฺสติ ฯเปฯ.\csromanspacer{} มโน อนิจฺโจติ ปสฺสติ, ธมฺมา อนิจฺจาติ ปสฺสติ, มโนวิญฺญาณํ อนิจฺจนฺติ ปสฺสติ, มโนสมฺผสฺโส อนิจฺโจติ ปสฺสติ, ยมฺปิทํ มโนสมฺผสฺสปจฺจยา อุปฺปชฺชติ เวทยิตํ สุขํ วา ทุกฺขํ วา อทุกฺขมสุขํ วา, ตมฺปิ อนิจฺจนฺติ ปสฺสติ.\csromanspacer{} อยํ โข สา ภิกฺขเว นิพฺพานสปฺปายา ปฏิปทาติ.\csromanspacer{}\csromanspacer{}ทุติยํ.}

\csromanheader{2}{3. ทุกฺขนิพฺพานสปฺปายสุตฺต}
\proseitem{148}{นิพฺพานสปฺปายํ โว ภิกฺขเว ปฏิปทํ เทเสสฺสามิ, ตํ สุณาถ ฯเปฯ.\csromanspacer{} กตมา จ สา ภิกฺขเว นิพฺพานสปฺปายา ปฏิปทา.\csromanspacer{} อิธ ภิกฺขเว จกฺขุํ ทุกฺขนฺติ ปสฺสติ, รูปา ทุกฺขาติ ปสฺสติ, จกฺขุวิญฺญาณํ ทุกฺขนฺติ ปสฺสติ, จกฺขุสมฺผสฺโส ทุกฺโขติ ปสฺสติ, ยมฺปิทํ จกฺขุสมฺผสฺสปจฺจยา อุปฺปชฺชติ เวทยิตํ สุขํ วา ทุกฺขํ วา อทุกฺขมสุขํ วา, ตมฺปิ ทุกฺขนฺติ ปสฺสติ ฯเปฯ.\csromanspacer{} ชิวฺหา ทุกฺขาติ ปสฺสติ ฯเปฯ มโน ทุกฺโขติ ปสฺสติ, ธมฺมา ทุกฺขาติ ปสฺสติ, มโนวิญฺญาณํ ทุกฺขนฺติ ปสฺสติ, มโนสมฺผสฺโส ทุกฺโขติ ปสฺสติ, ยมฺปิทํ มโนสมฺผสฺสปจฺจยา อุปฺปชฺชติ เวทยิตํ สุขํ วา ทุกฺขํ วา อทุกฺขมสุขํ วา,}

\gambhira{สฬายตนวคฺคสํยุตฺตปาฬิ}
\prose{\csromanfolio{348}ตมฺปิ ทุกฺขนฺติ ปสฺสติ.\csromanspacer{} อยํ โข สา ภิกฺขเว นิพฺพานสปฺปายา ปฏิปทาติ.\csromanspacer{}\csromanspacer{}ตติยํ.}

\csromanheader{2}{4. อนตฺตนิพฺพานสปฺปายสุตฺต}
\proseitem{149}{นิพฺพานสปฺปายํ โว ภิกฺขเว ปฏิปทํ เทเสสฺสามิ, ตํ สุณาถ ฯเปฯ.\csromanspacer{} กตมา จ สา ภิกฺขเว นิพฺพานสปฺปายา ปฏิปทา.\csromanspacer{} อิธ ภิกฺขเว ภิกฺขุ จกฺขุํ อนตฺตาติ ปสฺสติ, รูปา อนตฺตาติ ปสฺสติ, จกฺขุวิญฺญาณํ อนตฺตาติ ปสฺสติ, จกฺขุสมฺผสฺโส อนตฺตาติ ปสฺสติ, ยมฺปิทํ จกฺขุสมฺผสฺสปจฺจยา อุปฺปชฺชติ เวทยิตํ สุขํ วา ทุกฺขํ วา อทุกฺขมสุขํ วา, ตมฺปิ อนตฺตาติ ปสฺสติ ฯเปฯ มโน อนตฺตาติ ปสฺสติ, ธมฺมา อนตฺตาติ ปสฺสติ, มโนวิญฺญาณํ อนตฺตาติ ปสฺสติ, มโนสมฺผสฺโส อนตฺตาติ ปสฺสติ, ยมฺปิทํ มโนสมฺผสฺสปจฺจยา อุปฺปชฺชติ เวทยิตํ สุขํ วา ทุกฺขํ วา อทุกฺขมสุขํ วา, ตมฺปิ อนตฺตาติ ปสฺสติ.\csromanspacer{} อยํ โข สา ภิกฺขเว นิพฺพานสปฺปายา ปฏิปทาติ.\csromanspacer{}\csromanspacer{}จตุตฺถํ.}

\csromanmatikaanchor{mtk.183}
\csromantocmarkref{subsection}{5. นิพฺพานสปฺปายปฏิปทาสุตฺต}{mtk.183}
\bookmark[dest={mtk.183},level=2]{5. นิพฺพานสปฺปายปฏิปทาสุตฺต}
\csromanheader{2}{5. นิพฺพานสปฺปายปฏิปทาสุตฺต}
\proseitem{150}{นิพฺพานสปฺปายํ โว ภิกฺขเว ปฏิปทํ เทเสสฺสามิ, ตํ สุณาถ ฯเปฯ.\csromanspacer{} กตมา จ สา ภิกฺขเว นิพฺพานสปฺปายา ปฏิปทา.\csromanspacer{} ตํ กิํ มญฺญถ ภิกฺขเว จกฺขุ นิจฺจํ วา อนิจฺจํ วาติ.\csromanspacer{} อนิจฺจํ ภนฺเต.\csromanspacer{} ยํ ปนานิจฺจํ, ทุกฺขํ วา ตํ สุขํ วาติ.\csromanspacer{} ทุกฺขํ ภนฺเต.\csromanspacer{} ยํ ปนานิจฺจํ ทุกฺขํ วิปริณามธมฺมํ, กลฺลํ นุ ตํ สมนุปสฺสิตุํ “เอตํ มม, เอโสหมสฺมิ, เอโส เม อตฺตา”ติ.\csromanspacer{} โน เหตํ ภนฺเต.\csromanspacer{} รูปา นิจฺจา วา อนิจฺจา วาติ.\csromanspacer{} อนิจฺจา ภนฺเต.\csromanspacer{} จกฺขุวิญฺญาณํ, จกฺขุสมฺผสฺโส ฯเปฯ.\csromanspacer{} ยมฺปิทํ มโนสมฺผสฺสปจฺจยา อุปฺปชฺชติ เวทยิตํ สุขํ วา ทุกฺขํ วา อทุกฺขมสุขํ วา, ตมฺปิ นิจฺจํ วา อนิจฺจํ วาติ.\csromanspacer{} อนิจฺจํ ภนฺเต.\csromanspacer{} ยํ ปนานิจฺจํ, ทุกฺขํ วา ตํ สุขํ วาติ.\csromanspacer{} ทุกฺขํ ภนฺเต.\csromanspacer{} ยํ ปนานิจฺจํ ทุกฺขํ วิปริณามธมฺมํ, กลฺลํ นุ ตํ สมนุปสฺสิตุํ “เอตํ มม, เอโสหมสฺมิ, เอโส เม อตฺตา”ติ.\csromanspacer{} โน เหตํ ภนฺเต.\csromanspacer{} เอวํ ปสฺสํ ภิกฺขเว สุตวา อริยสาวโก จกฺขุสฺมิมฺปิ นิพฺพินฺทติ, รูเปสุปิ นิพฺพินฺทติ, จกฺขุวิญฺญาเณปิ นิพฺพินฺทติ, จกฺขุสมฺผสฺเสปิ นิพฺพินฺทติ ฯเปฯ.\csromanspacer{} ยมฺปิทํ มโนสมฺผสฺสปจฺจยา อุปฺปชฺชติ เวทยิตํ สุขํ วา ทุกฺขํ วา อทุกฺขมสุขํ วา, ตสฺมิมฺปิ นิพฺพินฺทติ, นิพฺพินฺทํ วิรชฺชติ, วิราคา วิมุจฺจติ ฯเปฯ นาปรํ อิตฺถตฺตายาติ ปชานาติ.\csromanspacer{} อยํ โข สา ภิกฺขเว นิพฺพานสปฺปายา ปฏิปทาติ.\csromanspacer{}\csromanspacer{}ปญฺจมํ.}

\csromanmatikaanchor{mtk.184}
\csromantocmarkref{subsection}{6. อนฺเตวาสิกสุตฺต}{mtk.184}
\bookmark[dest={mtk.184},level=2]{6. อนฺเตวาสิกสุตฺต}
\csromanheader{2}{1. สฬายตนสํยุตฺต \textbf{6. อนฺเตวาสิกสุตฺต}}
\proseitem{151}{\csromanfolio{349}อนนฺเตวาสิกมิทํ ภิกฺขเว พฺรหฺมจริยํ วุสฺสติ อนาจริยกํ, สนฺเตวาสิโก ภิกฺขเว ภิกฺขุ สาจริยโก ทุกฺขํ น ผาสุ\footnote{ผาสุํ (สี, อิ)} วิหรติ, อนนฺเตวาสิโก ภิกฺขเว ภิกฺขุ อนาจริยโก สุขํ ผาสุ วิหรติ.\csromanspacer{} กถญฺจ ภิกฺขุ สนฺเตวาสิโก สาจริยโก ทุกฺขํ น ผาสุ วิหรติ.\csromanspacer{} อิธ ภิกฺขเว ภิกฺขุโน จกฺขุนา รูปํ ทิสฺวา อุปฺปชฺชนฺติ ปาปกา อกุสลา ธมฺมา สรสงฺกปฺปา สญฺโญชนิยา, ตฺยาสฺส อนฺโต วสนฺติ.\csromanspacer{} อนฺตสฺส วสนฺติ ปาปกา อกุสลา ธมฺมาติ ตสฺมา สนฺเตวาสิโกติ วุจฺจติ.\csromanspacer{} เต นํ สมุทาจรนฺติ.\csromanspacer{} สมุทาจรนฺติ นํ ปาปกา อกุสลา ธมฺมาติ ตสฺมา สาจริยโกติ วุจฺจติ ฯเปฯ.}

\prose{ปุน จปรํ ภิกฺขเว ภิกฺขุโน ชิวฺหาย รสํ สายิตฺวา อุปฺปชฺชนฺติ ปาปกา อกุสลา ธมฺมา สรสงฺกปฺปา สญฺโญชนิยา, ตฺยาสฺส อนฺโต วสนฺติ.\csromanspacer{} อนฺตสฺส วสนฺติ ปาปกา อกุสลา ธมฺมาติ ตสฺมา สนฺเตวาสิโกติ วุจฺจติ.\csromanspacer{} เต นํ สมุทาจรนฺติ.\csromanspacer{} สมุทาจรนฺติ นํ ปาปกา อกุสลา ธมฺมาติ ตสฺมา สาจริยโกติ วุจฺจติ ฯเปฯ.}

\prose{ปุน จปรํ ภิกฺขเว ภิกฺขุโน มนสา ธมฺมํ วิญฺญาย อุปฺปชฺชนฺติ ปาปกา อกุสลา ธมฺมา สรสงฺกปฺปา สญฺโญชนิยา, ตฺยาสฺส อนฺโต วสนฺติ.\csromanspacer{} อนฺตสฺส วสนฺติ ปาปกา อกุสลา ธมฺมาติ ตสฺมา สนฺเตวาสิโกติ วุจฺจติ.\csromanspacer{} เต นํ สมุทาจรนฺติ.\csromanspacer{} สมุทาจรนฺติ นํ ปาปกา อกุสลา ธมฺมาติ ตสฺมา สาจริยโกติ วุจฺจติ.\csromanspacer{} เอวํ โข ภิกฺขเว ภิกฺขุ สนฺเตวาสิโก สาจริยโก ทุกฺขํ น ผาสุ วิหรติ.}

\prose{กถญฺจ ภิกฺขเว ภิกฺขุ อนนฺเตวาสิโก อนาจริยโก สุขํ ผาสุ วิหรติ.\csromanspacer{} อิธ ภิกฺขเว ภิกฺขุโน จกฺขุนา รูปํ ทิสฺวา น อุปฺปชฺชนฺติ ปาปกา อกุสลา ธมฺมา สรสงฺกปฺปา สญฺโญชนิยา, ตฺยาสฺส น อนฺโต วสนฺติ.\csromanspacer{} นาสฺส อนฺโต วสนฺติ ปาปกา อกุสลา ธมฺมาติ ตสฺมา อนนฺเตวาสิโกติ วุจฺจติ.\csromanspacer{} เต นํ น สมุทาจรนฺติ.\csromanspacer{} น สมุทาจรนฺติ นํ ปาปกา อกุสลา ธมฺมาติ ตสฺมา อนาจริยโกติ วุจฺจติ ฯเปฯ.}

\prose{ปุน จปรํ ภิกฺขเว ภิกฺขุโน ชิวฺหาย รสํ สายิตฺวา น อุปฺปชฺชนฺติ ปาปกา อกุสลา ธมฺมา สรสงฺกปฺปา สญฺโญชนิยา, ตฺยาสฺส น อนฺโต วสนฺติ.}

\gambhira{สฬายตนวคฺคสํยุตฺตปาฬิ}
\prose{\csromanfolio{350}นาสฺส อนฺโต วสนฺติ ปาปกา อกุสลา ธมฺมาติ ตสฺมา อนนฺเตวาสิโกติ วุจฺจติ.\csromanspacer{} เต นํ น สมุทาจรนฺติ.\csromanspacer{} น สมุทาจรนฺติ นํ ปาปกา อกุสลา ธมฺมาติ ตสฺมา อนาจริยโกติ วุจฺจติ ฯเปฯ.}

\prose{ปุน จปรํ ภิกฺขเว ภิกฺขุโน มนสา ธมฺมํ วิญฺญาย น อุปฺปชฺชนฺติ ปาปกา อกุสลา ธมฺมา สรสงฺกปฺปา สญฺโญชนิยา, ตฺยาสฺส น อนฺโต วสนฺติ.\csromanspacer{} นาสฺส อนฺโต วสนฺติ ปาปกา อกุสลา ธมฺมาติ ตสฺมา อนนฺเตวาสิโกติ วุจฺจติ.\csromanspacer{} เต นํ น สมุทาจรนฺติ.\csromanspacer{} น สมุทาจรนฺติ นํ ปาปกา อกุสลา ธมฺมาติ ตสฺมา อนาจริยโกติ วุจฺจติ.\csromanspacer{} เอวํ โข ภิกฺขเว ภิกฺขุ อนนฺเตวาสิโก อนาจริยโก สุขํ ผาสุ วิหรติ.\csromanspacer{} อนนฺเตวาสิกมิทํ ภิกฺขเว พฺรหฺมจริยํ วุสฺสติ อนาจริยกํ, สนฺเตวาสิโก ภิกฺขเว ภิกฺขุ สาจริยโก ทุกฺขํ น ผาสุ วิหรติ, อนนฺเตวาสิโก ภิกฺขเว ภิกฺขุ อนาจริยโก สุขํ ผาสุ วิหรตีติ.\csromanspacer{}\csromanspacer{}ฉฏฺฐํ.}

\csromanmatikaanchor{mtk.185}
\csromantocmarkref{subsection}{7. กิมตฺถิยพฺรหฺมจริยสุตฺต}{mtk.185}
\bookmark[dest={mtk.185},level=2]{7. กิมตฺถิยพฺรหฺมจริยสุตฺต}
\csromanheader{2}{7. กิมตฺถิยพฺรหฺมจริยสุตฺต}
\proseitem{152}{สเจ โว ภิกฺขเว อญฺญติตฺถิยา ปริพฺพาชกา เอวํ ปุจฺเฉยฺยุํ “กิมตฺถิยํ อาวุโส สมเณ โคตเม พฺรหฺมจริยํ วุสฺสตี”ติ.\csromanspacer{} เอวํ ปุฏฺฐา ตุเมฺห ภิกฺขเว เตสํ อญฺญติตฺถิยานํ ปริพฺพาชกานํ เอวํ พฺยากเรยฺยาถ “ทุกฺขสฺส โข อาวุโส ปริญฺญาย ภควติ พฺรหฺมจริยํ วุสฺสตี”ติ.\csromanspacer{} สเจ ปน โว ภิกฺขเว อญฺญติตฺถิยา ปริพฺพาชกา เอวํ ปุจฺเฉยฺยุํ “กตมํ ปนาวุโส ทุกฺขํ, ยสฺส ปริญฺญาย สมเณ โคตเม พฺรหฺมจริยํ วุสฺสตี”ติ.\csromanspacer{} เอวํ ปุฏฺฐา ตุเมฺห ภิกฺขเว เตสํ อญฺญติตฺถิยานํ ปริพฺพาชกานํ เอวํ พฺยากเรยฺยาถ\csromandash{}}

\prose{จกฺขุ โข อาวุโส ทุกฺขํ, ตสฺส ปริญฺญาย ภควติ พฺรหฺมจริยํ วุสฺสติ.\csromanspacer{} รูปา ทุกฺขา, เตสํ ปริญฺญาย ภควติ พฺรหฺมจริยํ วุสฺสติ.\csromanspacer{} จกฺขุวิญฺญาณํ ทุกฺขํ, ตสฺส ปริญฺญาย ภควติ พฺรหฺมจริยํ วุสฺสติ.\csromanspacer{} จกฺขุสมฺผสฺโส ทุกฺโข, ตสฺส ปริญฺญาย ภควติ พฺรหฺมจริยํ วุสฺสติ.\csromanspacer{} ยมฺปิทํ จกฺขุสมฺผสฺสปจฺจยา อุปฺปชฺชติ เวทยิตํ สุขํ วา ทุกฺขํ วา อทุกฺขมสุขํ วา, ตมฺปิ ทุกฺขํ, ตสฺส ปริญฺญาย ภควติ พฺรหฺมจริยํ วุสฺสติ ฯเปฯ.\csromanspacer{} ชิวฺหา ทุกฺขา.\csromanspacer{} มโน ทุกฺโข, ตสฺส ปริญฺญาย ภควติ พฺรหฺมจริยํ วุสฺสติ.\csromanspacer{} ยมฺปิทํ มโนสมฺผสฺสปจฺจยา อุปฺปชฺชติ เวทยิตํ สุขํ วา ทุกฺขํ วา อทุกฺขมสุขํ วา, ตมฺปิ ทุกฺขํ, ตสฺส ปริญฺญาย ภควติ พฺรหฺมจริยํ วุสฺสติ.\csromanspacer{} อิทํ โข}

\vspace{\tipitakaparskip}
\csromancenter{สฬายตนสํยุตฺต อาวุโส ทุกฺขํ ยสฺส ปริญฺญาย ภควติ พฺรหฺมจริยํ วุสฺสตีติ.\csromanspacer{} เอวํ ปุฏฺฐา ตุเมฺห ภิกฺขเว เตสํ อญฺญติตฺถิยานํ ปริพฺพาชกานํ เอวํ พฺยากเรยฺยาถาติ.\csromanspacer{}\csromanspacer{}สตฺตมํ.}
\csromanmatikaanchor{mtk.186}
\csromantocmarkref{subsection}{8. อตฺถินุโขปริยายสุตฺต}{mtk.186}
\bookmark[dest={mtk.186},level=2]{8. อตฺถินุโขปริยายสุตฺต}
\csromanheader{2}{8. อตฺถินุโขปริยายสุตฺต}
\proseitem{153}{\csromanfolio{351}อตฺถิ นุ โข ภิกฺขเว ปริยาโย, ยํ ปริยายํ อาคมฺม ภิกฺขุ อญฺญเตฺรว สทฺธาย อญฺญตฺร รุจิยา อญฺญตฺร อนุสฺสวา อญฺญตฺร อาการปริวิตกฺกา อญฺญตฺร ทิฏฺฐินิชฺฌานกฺขนฺติยา อญฺญํ พฺยากเรยฺย “ขีณา ชาติ, วุสิตํ พฺรหฺมจริยํ, กตํ กรณียํ, นาปรํ อิตฺถตฺตายาติ ปชานามี”ติ\footnote{ปชานาตีติ (สฺยา, กํ, อิ, ก)}.\csromanspacer{} ภควํมูลกา โน ภนฺเต ธมฺมา ภควํเนตฺติกา ภควํปฏิสรณา, สาธุ วต ภนฺเต ภควนฺตํเยว ปฏิภาตุ เอตสฺส ภาสิตสฺส อตฺโถ, ภควโต สุตฺวา ภิกฺขู ธาเรสฺสนฺตีติ.\csromanspacer{} เตน หิ ภิกฺขเว สุณาถ สาธุกํ มนสิ กโรถ, ภาสิสฺสามีติ. “เอวํ ภนฺเต”ติ โข เต ภิกฺขู ภควโต ปจฺจสฺโสสุํ.\csromanspacer{} ภควา เอตทโวจ\csromandash{}อตฺถิ ภิกฺขเว ปริยาโย ยํ ปริยายํ อาคมฺม ภิกฺขุ อญฺญเตฺรว สทฺธาย อญฺญตฺร รุจิยา อญฺญตฺร อนุสฺสวา อญฺญตฺร อาการปริวิตกฺกา อญฺญตฺร ทิฏฺฐินิชฺฌานกฺขนฺติยา อญฺญํ พฺยากเรยฺย “ขีณา ชาติ, วุสิตํ พฺรหฺมจริยํ, กตํ กรณียํ, นาปรํ อิตฺถตฺตายาติ ปชานามี”ติ1.}

\prose{กตโม จ ภิกฺขเว ปริยาโย.\csromanspacer{} ยํ ปริยายํ อาคมฺม ภิกฺขุ อญฺญเตฺรว สทฺธาย ฯเปฯ อญฺญตฺร ทิฏฺฐินิชฺฌานกฺขนฺติยา อญฺญํ พฺยากโรติ “ขีณา ชาติ, วุสิตํ พฺรหฺมจริยํ, กตํ กรณียํ, นาปรํ อิตฺถตฺตายาติ ปชานามี”ติ1.\csromanspacer{} อิธ ภิกฺขเว ภิกฺขุ จกฺขุนา รูปํ ทิสฺวา สนฺตํ วา อชฺฌตฺตํ ราคโทสโมหํ อตฺถิ เม อชฺฌตฺตํ ราคโทสโมโหติ ปชานาติ, อสนฺตํ วา อชฺฌตฺตํ ราคโทสโมหํ นตฺถิ เม อชฺฌตฺตํ ราคโทสโมโหติ ปชานาติ, ยํ ตํ ภิกฺขเว ภิกฺขุ จกฺขุนา รูปํ ทิสฺวา สนฺตํ วา อชฺฌตฺตํ ราคโทสโมหํ อตฺถิ เม อชฺฌตฺตํ ราคโทสโมโหติ ปชานาติ, อสนฺตํ วา อชฺฌตฺตํ ราคโทสโมหํ นตฺถิ เม อชฺฌตฺตํ ราคโทสโมโหติ ปชานาติ.\csromanspacer{} อปิ นุเม ภิกฺขเว ธมฺมา สทฺธาย วา เวทิตพฺพา, รุจิยา วา เวทิตพฺพา, อนุสฺสเวน วา เวทิตพฺพา, อาการปริวิตกฺเกน วา เวทิตพฺพา, ทิฏฺฐินิชฺฌานกฺขนฺติยา วา เวทิตพฺพาติ.\csromanspacer{} โน}

\gambhira{สฬายตนวคฺคสํยุตฺตปาฬิ}
\prose{\csromanfolio{352}เหตํ ภนฺเต.\csromanspacer{} นนุเม ภิกฺขเว ธมฺมา ปญฺญาย ทิสฺวา เวทิตพฺพาติ.\csromanspacer{} เอวํ ภนฺเต.\csromanspacer{} อยํ โข ภิกฺขเว ปริยาโย, ยํ ปริยายํ อาคมฺม ภิกฺขุ อญฺญเตฺรว สทฺธาย อญฺญตฺร รุจิยา อญฺญตฺร อนุสฺสวา อญฺญตฺร อาการปริวิตกฺกา อญฺญตฺร ทิฏฺฐินิชฺฌานกฺขนฺติยา อญฺญํ พฺยากโรติ “ขีณา ชาติ, วุสิตํ พฺรหฺมจริยํ, กตํ กรณียํ, นาปรํ อิตฺถตฺตายาติ ปชานามี”ติ\footnote{ปชานาตีติ (สฺยา, กํ, อิ, ก)} ฯเปฯ.}

\prose{ปุน จปรํ ภิกฺขเว ภิกฺขุ ชิวฺหาย รสํ สายิตฺวา สนฺตํ วา อชฺฌตฺตํ ฯเปฯ.\csromanspacer{} ราคโทสโมโหติ ปชานาติ, อสนฺตํ วา อชฺฌตฺตํ ราคโทสโมหํ นตฺถิ เม อชฺฌตฺตํ ราคโทสโมโหติ ปชานาติ, ยํ ตํ ภิกฺขเว ชิวฺหาย รสํ สายิตฺวา สนฺตํ วา อชฺฌตฺตํ ราคสโมหํ อตฺถิ เม อชฺฌตฺตํ ราคโทสโมโหติ ปชานาติ, อสนฺตํ วา อชฺฌตฺตํ ราคโทสโมหํ นตฺถิ เม อชฺฌตฺตํ ราคโทสโมโหติ ปชานาติ.\csromanspacer{} อปิ นุเม ภิกฺขเว ธมฺมา สทฺธาย วา เวทิตพฺพา, รุจิยา วา เวทิตพฺพา, อนุสฺสเวน วา เวทิตพฺพา, อาการปริวิตกฺเกน วา เวทิตพฺพา, ทิฏฺฐินิชฺฌานกฺขนฺติยา วา เวทิตพฺพาติ.\csromanspacer{} โน เหตํ ภนฺเต.\csromanspacer{} นนุเม ภิกฺขเว ธมฺมา ปญฺญาย ทิสฺวา เวทิตพฺพาติ.\csromanspacer{} เอวํ ภนฺเต.\csromanspacer{} อยมฺปิ โข ภิกฺขเว ปริยาโย, ยํ ปริยายํ อาคมฺม ภิกฺขุ อญฺญเตฺรว สทฺธาย อญฺญตฺร รุจิยา อญฺญตฺร อนุสฺสวา อญฺญตฺร อาการปริวิตกฺกา อญฺญตฺร ทิฏฺฐิชฺฌานกฺขนฺติยา อญฺญํ พฺยากโรติ “ขีณา ชาติ, วุสิตํ พฺรหฺมจริยํ, กตํ กรณียํ, นาปรํ อิตฺถตฺตายาติ ปชานามี”ติ1 ฯเปฯ.}

\prose{ปุน จปรํ ภิกฺขเว ภิกฺขุ มนสา ธมฺมํ วิญฺญาย สนฺตํ วา อชฺฌตฺตํ ราคโทสโมหํ อตฺถิ เม อชฺฌตฺตํ ราคโทสโมโหติ ปชานาติ, อสนฺตํ วา อชฺฌตฺตํ ราคโทสโมหํ นตฺถิ เม อชฺฌตฺตํ ราคโทสโมโหติ ปชานาติ, ยํ ตํ ภิกฺขเว ภิกฺขุ มนสา ธมฺมํ วิญฺญาย สนฺตํ วา อชฺฌตฺตํ ราคโทสโมหํ อตฺถิ เม อชฺฌตฺตํ ราคโทสโมโหติ ปชานาติ, อสนฺตํ วา อชฺฌตฺตํ ราคโทสโมหํ นตฺถิ เม อชฺฌตฺตํ ราคโทสโมโหติ ปชานาติ.\csromanspacer{} อปิ นุเม ภิกฺขเว ธมฺมา สทฺธาย วา เวทิตพฺพา, รุจิยา วา เวทิตพฺพา, อนุสฺสเวน วา เวทิตพฺพา, อาการปริวิตกฺเกน วา เวทิตพฺพา, ทิฏฺฐินิชฺฌานกฺขนฺติยา วา เวทิตพฺพาติ.\csromanspacer{} โน เหตํ ภนฺเต.\csromanspacer{} นนุเม ภิกฺขเว ธมฺมา ปญฺญาย ทิสฺวา เวทิตพฺพาติ.\csromanspacer{} เอวํ ภนฺเต.\csromanspacer{} อยมฺปิ โข ภิกฺขเว ปริยาโย, ยํ ปริยายํ อาคมฺม}

\vspace{\tipitakaparskip}
\csromancenter{สฬายตนสํยุตฺต ภิกฺขุ อญฺญเตฺรว สทฺธาย อญฺญตฺร รุจิยา อญฺญตฺร อนุสฺสวา อญฺญตฺร อาการปริวิตกฺกา อญฺญตฺร ทิฏฺฐินิชฺฌานกฺขนฺติยา อญฺญํ พฺยากโรติ “ขีณา ชาติ, วุสิตํ พฺรหฺมจริยํ, กตํ กรณียํ, นาปรํ อิตฺถตฺตายาติ ปชานามี”ติ.\csromanspacer{}\csromanspacer{}อฏฺฐมํ.}
\csromanmatikaanchor{mtk.187}
\csromantocmarkref{subsection}{9. อินฺทฺริยสมฺปนฺนสุตฺต}{mtk.187}
\bookmark[dest={mtk.187},level=2]{9. อินฺทฺริยสมฺปนฺนสุตฺต}
\csromanheader{2}{9. อินฺทฺริยสมฺปนฺนสุตฺต}
\csromanmatikaanchor{mtk.188}
\csromantocmarkref{subsection}{10. ธมฺมกถิกปุจฺฉาสุตฺต}{mtk.188}
\bookmark[dest={mtk.188},level=2]{10. ธมฺมกถิกปุจฺฉาสุตฺต}
\proseitem{154}{\csromanfolio{353}อถ โข อญฺญตโร ภิกฺขุ เยน ภควา เตนุปสงฺกมิ ฯเปฯ เอกมนฺตํ นิสินฺโน โข โส ภิกฺขเว ภควนฺตํ เอตทโวจ “อินฺทฺริยสมฺปนฺโน อินฺทฺริยสมฺปนฺโน”ติ ภนฺเต วุจฺจติ, กิตฺตาวตา นุ โข ภนฺเต อินฺทฺริยสมฺปนฺโน โหตีติ.}

\prose{จกฺขุนฺทฺริเย เจ ภิกฺขุ อุทยพฺพยานุปสฺสี วิหรนฺโต จกฺขุนฺทฺริเย นิพฺพินฺทติ ฯเปฯ.\csromanspacer{} ชิวฺหินฺทฺริเย เจ ภิกฺขุ อุทยพฺพยานุปสฺสี วิหรนฺโต ชิวฺหินฺทฺริเย นิพฺพินฺทติ ฯเปฯ.\csromanspacer{} มนินฺทฺริเย เจ ภิกฺขุ อุทยพฺพยานุปสฺสี วิหรนฺโต มนินฺทฺริเย นิพฺพินฺทติ, นิพฺพินฺทํ วิรชฺชติ ฯเปฯ วิมุตฺตสฺมิํ “วิมุตฺตมฺ”อิติ ญาณํ โหติ, “ขีณา ชาติ, วุสิตํ พฺรหฺมจริยํ, กตํ กรณียํ, นาปรํ อิตฺถตฺตายา”ติ ปชานาติ.\csromanspacer{} เอตฺตาวตา โข ภิกฺขุ อินฺทฺริยสมฺปนฺโน โหตีติ.\csromanspacer{}\csromanspacer{}นวมํ.}

\csromanheader{2}{10. ธมฺมกถิกปุจฺฉสุตฺต}
\proseitem{155}{อถ โข อญฺญตโร ภิกฺขุ เยน ภควา เตนุปสงฺกมิ ฯเปฯ เอกมนฺตํ นิสินฺโน โข โส ภิกฺขุ ภควนฺตํ เอตทโวจ “ธมฺมกถิโก ธมฺมกถิโก”ติ ภนฺเต วุจฺจติ, กิตฺตาวตา นุ โข ภนฺเต ธมฺมกถิโก โหตีติ.}

\prose{จกฺขุสฺส เจ ภิกฺขุ นิพฺพิทาย วิราคาย นิโรธาย ธมฺมํ เทเสติ, ธมฺมกถิโก ภิกฺขูติ อลํวจนาย.\csromanspacer{} จกฺขุสฺส เจ ภิกฺขุ นิพฺพิทาย วิราคาย นิโรธาย ปฏิปนฺโน โหติ, ธมฺมานุธมฺมปฺปฏิปนฺโน ภิกฺขูติ อลํวจนาย.\csromanspacer{} จกฺขุสฺส เจ ภิกฺขุ นิพฺพิทา วิราคา นิโรธา อนุปาทาวิมุตฺโต โหติ, ทิฏฺฐธมฺมนิพฺพานปฺปตฺโต ภิกฺขูติ อลํวจนาย ฯเปฯ.\csromanspacer{} ชิวฺหาย เจ ภิกฺขุ นิพฺพิทาย วิราคาย นิโรธาย ธมฺมํ เทเสติ, ธมฺมกถิโก ภิกฺขูติ อลํวจนาย ฯเปฯ.\csromanspacer{} มนสฺส เจ ภิกฺขุ นิพฺพิทาย วิราคาย นิโรธาย ธมฺมํ เทเสติ, ธมฺมกถิโก ภิกฺขูติ อลํวจนาย.\csromanspacer{} มนสา เจ ภิกฺขุ นิพฺพิทาย วิราคาย นิโรธาย ปฏิปนฺโน โหติ, ธมฺมานุธมฺมปฺปฏิปนฺโน ภิกฺขูติ อลํวจนาย.}

\gambhira{สฬายตนวคฺคสํยุตฺตปาฬิ}
\csromanmatikaanchor{mtk.189}
\csromanmatikaanchor{mtk.190}
\csromantocmarknopage{section}{(16) 1. นนฺทิกฺขยวคฺค}{mtk.189}
\bookmark[dest={mtk.189},level=1]{(16) 1. นนฺทิกฺขยวคฺค}
\csromantocmarkref{subsection}{1-4. อชฺฌตฺตนนฺทิกฺขยสุตฺตาทิ}{mtk.190}
\bookmark[dest={mtk.190},level=2]{1-4. อชฺฌตฺตนนฺทิกฺขยสุตฺตาทิ}
\prose{\csromanfolio{354}มนสฺส เจ ภิกฺขุ นิพฺพิทา วิราคา นิโรธา อนุปาทาวิมุตฺโต โหติ, ทิฏฺฐธมฺมนิพฺพานปฺปตฺโต ภิกฺขูติ อลํวจนายาติ.\csromanspacer{}\csromanspacer{}ทสมํ.\csromanspacer{} นวปุราณวคฺโค ปญฺจโม.}

\titlehead{\textbf{ตสฺสุทฺทานํ}}
\csromangathameasure{กมฺมํ จตฺตาริ สปฺปายา, อนนฺเตวาสิ กิมตฺถิยา. \\ อตฺถิ นุ โข ปริยาโย, อินฺทฺริยกถิเกน จาติ. สฬายตนวคฺเค ตติยปณฺณาสโก สมตฺโต.}
\csromangathagroup{\csromangathastanza{กมฺมํ จตฺตาริ สปฺปายา, อนนฺเตวาสิ กิมตฺถิยา. \\ อตฺถิ นุ โข ปริยาโย, อินฺทฺริยกถิเกน จาติ.\csromanspacer{} สฬายตนวคฺเค ตติยปณฺณาสโก สมตฺโต.}}

\csromanheader{2}{\textbf{ตสฺส วคฺคุทฺทานํ}}
\csromangathameasure{โยคกฺเขมิ จ โลโก จ, คหปติ เทวทเหน จ. \\ นวปุราเณน ปณฺณาโส, ตติโย เตน วุจฺจตีติ.}
\csromangathagroup{\csromangathastanza{โยคกฺเขมิ จ โลโก จ, คหปติ เทวทเหน จ. \\ นวปุราเณน ปณฺณาโส, ตติโย เตน วุจฺจตีติ.}}

\csromanheader{1}{\textbf{(16) 1. นนฺทิกฺขยวคฺค}}
\csromanheader{2}{1. อชฺฌตฺตนนฺทิกฺขยสุตฺต}
\proseitem{156}{อนิจฺจํเยว ภิกฺขเว ภิกฺขุ จกฺขุํ อนิจฺจนฺติ ปสฺสติ, สาสฺส\footnote{สายํ (อิ, ก)} โหติ สมฺมาทิฏฺฐิ, สมฺมา ปสฺสํ นิพฺพินฺทติ, นนฺทิกฺขยา ราคกฺขโย, ราคกฺขยา นนฺทิกฺขโย, นนฺทิราคกฺขยา จิตฺตํ สุวิมุตฺตนฺติ วุจฺจติ ฯเปฯ.\csromanspacer{} อนิจฺจํเยว ภิกฺขเว ภิกฺขุ ชิวฺหํ อนิจฺจนฺติ ปสฺสติ, สาสฺส โหติ สมฺมาทิฏฺฐิ, สมฺมา ปสฺสํ นิพฺพินฺทติ, นนฺทิกฺขยา ราคกฺขโย, ราคกฺขยา ฯเปฯ จิตฺตํ สุวิมุตฺตนฺติ วุจฺจติ.\csromanspacer{} อนิจฺจํเยว ภิกฺขเว ภิกฺขุ มนํ อนิจฺจนฺติ ปสฺสติ, สาสฺส โหติ สมฺมาทิฏฺฐิ, สมฺมา ปสฺสํ นิพฺพินฺทติ, นนฺทิกฺขยา ราคกฺขโย, ราคกฺขยา นนฺทิกฺขโย, นนฺทิราคกฺขยา จิตฺตํ สุวิมุตฺตนฺติ วุจฺจตีติ.\csromanspacer{}\csromanspacer{}ปฐมํ.}

\csromanheader{2}{2. พาหิรนนฺทิกฺขยสุตฺต}
\proseitem{157}{อนิจฺเจเยว ภิกฺขเว ภิกฺขุ รูเป อนิจฺจาติ ปสฺสติ, สาสฺส โหติ สมฺมาทิฏฺฐิ, สมฺมา ปสฺสํ นิพฺพินฺทติ, นนฺทิกฺขยา ราคกฺขโย, ราคกฺขยา}

\vspace{\tipitakaparskip}
\csromancenter{สฬายตนสํยุตฺต นนฺทิกฺขโย, นนฺทิราคกฺขยา จิตฺตํ สุวิมุตฺตนฺติ วุจฺจติ.\csromanspacer{} อนิจฺเจเยว ภิกฺขเว ภิกฺขุ สทฺเท.\csromanspacer{} คนฺเธ.\csromanspacer{} รเส.\csromanspacer{} โผฏฺฐพฺเพ.\csromanspacer{} ธมฺเม อนิจฺจาติ ปสฺสติ, สาสฺส โหติ สมฺมาทิฏฺฐิ, สมฺมา ปสฺสํ นิพฺพินฺทติ, นนฺทิกฺขยา ราคกฺขโย, ราคกฺขยานนฺทิกฺขโย, นนฺทิราคกฺขยา จิตฺตํ สุวิมุตฺตนฺติ วุจฺจตีติ.\csromanspacer{}\csromanspacer{}ทุติยํ.}
\csromanheader{2}{3. อชฺฌตฺต-อนิจฺจนนฺทิกฺขยสุตฺต}
\proseitem{158}{\csromanfolio{355}จกฺขุํ ภิกฺขเว โยนิโส มนสิ กโรถ, จกฺขานิจฺจตญฺจ ยถาภูตํ สมนุปสฺสถ, จกฺขุํ ภิกฺขเว ภิกฺขุ โยนิโส มนสิกโรนฺโต จกฺขานิจฺจตญฺจ ยถาภูตํ สมนุปสฺสนฺโต จกฺขุสฺมิมฺปิ นิพฺพินฺทติ, นนฺทิกฺขยา ราคกฺขโย, ราคกฺขยา นนฺทิกฺขโย, นนฺทิราคกฺขยา จิตฺตํ สุวิมุตฺตนฺติ วุจฺจติ.\csromanspacer{} โสตํ ภิกฺขเว โยนิโส มนสิ กโรถ.\csromanspacer{} ฆานํ.\csromanspacer{} ชิวฺหํ ภิกฺขเว โยนิโส มนสิ กโรถ, ชิวฺหานิจฺจตญฺจ ยถาภูตํ สมนุปสฺสถ, ชิวฺหํ ภิกฺขเว ภิกฺขุ โยนีโส มนสิกโรนฺโต ชิวฺหานิจฺจตญฺจ ยถาภูตํ สมนุปสฺสนฺโต ชิวฺหายปิ นิพฺพินฺทติ นนฺทิกฺขยา ราคกฺขโย, ราคกฺขยา นนฺทิกฺขโย, นนฺทิราคกฺขยา จิตฺตํ สุวิมุตฺตนฺติ วุจฺจติ.\csromanspacer{} กายํ.\csromanspacer{} มนํ ภิกฺขเว โยนิโส มนสิ กโรถ, มนานิจฺจตญฺจ ยถาภูตํ สมนุปสฺสถ, มนํ ภิกฺขเว ภิกฺขุ โยนิโส มนสิกโรนฺโต มนานิจฺจตญฺจ ยถาภูตํ สมนุปสฺสนฺโต มนสฺมิมฺปิ นิพฺพินฺทติ, นนฺทิกฺขยา ราคกฺขโย, ราคกฺขยา นนฺทิกฺขโย, นนฺทิราคกฺขยา จิตฺตํ สุวิมุตฺตนฺติ วุจฺจตีติ.\csromanspacer{}\csromanspacer{}ตติยํ.}

\vspace{\tipitakaparskip}
\csromancenter{พาหิร-อนิจฺจนนฺทิกฺขยสุตฺต 159.\csromanspacer{} รูเป ภิกฺขเว โยนิโส มนสิ กโรถ, รูปานิจฺจตญฺจ ยถาภูตํ สมนุปสฺสถ, รูเป ภิกฺขเว ภิกฺขุ โยนิโส มนสิกโรนฺโต รูปานิจฺจตญฺจ ยถาภูตํ สมนุปสฺสนฺโต รูเปสุปิ นิพฺพินฺทติ, นนฺทิกฺขยา ราคกฺขโย, ราคกฺขยา นินฺทิกฺขโย, นินฺทิราคกฺขยา จิตฺตํ สุวิมุตฺตนฺติ วุจฺจติ.\csromanspacer{} สทฺเท.\csromanspacer{} คนฺเธ.\csromanspacer{} รเส.\csromanspacer{} โผฏฺฐพฺเพ.\csromanspacer{} ธมฺเม ภิกฺขเว โยนิโส มนสิ กโรถ, ธมฺมานิจฺจตญฺจ ยถาภูตํ สมนุปสฺสถ, ธมฺเม ภิกฺขเว ภิกฺขุ โยนิโส มนสิกโรนฺโต ธมฺมานิจฺจตญฺจ ยถาภูตํ สมนุปสฺสนฺโต ธมฺเมสุปิ นิพฺพินฺทติ, นนฺทิกฺขยา ราคกฺขโย, ราคกฺขยา นนฺทิกฺขโย, นนฺทิราคกฺขยา จิตฺตํ สุวิมุตฺตนฺติ วุจฺจตีติ.\csromanspacer{}\csromanspacer{}จตุตฺถํ.}
\gambhira{สฬายตนวคฺคสํยุตฺตปาฬิ}
\csromanmatikaanchor{mtk.191}
\csromantocmarkref{subsection}{5. ชีวกมฺพวนสมาธิสุตฺต}{mtk.191}
\bookmark[dest={mtk.191},level=2]{5. ชีวกมฺพวนสมาธิสุตฺต}
\csromanheader{2}{5. ชีวกมฺพวนสมาธิสุตฺต}
\proseitem{160}{\csromanfolio{356}เอกํ สมยํ ภควา ราชคเห วิหรติ ชีวกมฺพวเน.\csromanspacer{} ตตฺร โข ภควา ภิกฺขู อามนฺเตสิ “ภิกฺขโว”ติ ฯเปฯ.\csromanspacer{} สมาธิํ ภิกฺขเว ภาเวถ, สมาหิตสฺส ภิกฺขเว ภิกฺขุโน ยถาภูตํ โอกฺขายติ.\csromanspacer{} กิญฺจ ยถาภูตํ โอกฺขายติ.\csromanspacer{} จกฺขุํ อนิจฺจนฺติ ยถาภูตํ โอกฺขายติ, รูปา อนิจฺจาติ ยถาภูตํ โอกฺขายติ, จกฺขุวิญฺญาณํ อนิจฺจนฺติ ยถาภูตํ โอกฺขายติ, จกฺขุสมฺผสฺโส อนิจฺโจติ ยถาภูตํ โอกฺขายติ.\csromanspacer{} ยมฺปิทํ จกฺขุสมฺผสฺสปจฺจยา อุปฺปชฺชติ เวทยิตํ สุขํ วา ทุกฺขํ วา อทุกฺขมสุขํ วา, ตมฺปิ อนิจฺจนฺติ ยถาภูตํ โอกฺขายติ ฯเปฯ.\csromanspacer{} ชิวฺหา อนิจฺจาติ ยถาภูตํ โอกฺขายติ ฯเปฯ.\csromanspacer{} มโน อนิจฺโจติ ยถาภูตํ โอกฺขายติ, ธมฺมา อนิจฺจาติ ยถาภูตํ โอกฺขายติ ฯเปฯ.\csromanspacer{} ยมฺปิทํ มโนสมฺผสฺสปจฺจยา อุปฺปชฺชติ เวทยิตํ สุขํ วา ทุกฺขํ วา อทุกฺขมสุขํ วา, ตมฺปิ อนิจฺจนฺติ ยถาภูตํ โอกฺขายติ.\csromanspacer{} สมาธิํ ภิกฺขเว ภาเวถ, สมาหิตสฺส ภิกฺขเว ภิกฺขุโน ยถาภูตํ โอกฺขายตีติ.\csromanspacer{}\csromanspacer{}ปญฺจมํ.}

\csromanmatikaanchor{mtk.192}
\csromantocmarkref{subsection}{6. ชีวกมฺพวนปฏิสลฺลานสุตฺต}{mtk.192}
\bookmark[dest={mtk.192},level=2]{6. ชีวกมฺพวนปฏิสลฺลานสุตฺต}
\csromanheader{2}{6. ชีวกมฺพวนปฏิสลฺลานสุตฺต}
\proseitem{161}{เอกํ สมยํ ภควา ราชคเห วิหรติ ชีวกมฺพวเน.\csromanspacer{} ตตฺร โข ภควา ภิกฺขู อามนฺเตสิ ฯเปฯ.\csromanspacer{} ปฏิสลฺลาเน ภิกฺขเว โยคมาปชฺชถ, ปฏิสลฺลีนสฺส ภิกฺขเว ภิกฺขุโน ยถาภูตํ โอกฺขายติ.\csromanspacer{} กิญฺจ ยถาภูตํ โอกฺขายติ.\csromanspacer{} จกฺขุํ อนิจฺจนฺติ ยถาภูตํ โอกฺขายติ, รูปา อนิจฺจาติ ยถาภูตํ โอกฺขายติ, จกฺขุวิญฺญาณํ อนิจฺจนฺติ ยถาภูตํ โอกฺขายติ, จกฺขุสมฺผสฺโส อนิจฺโจติ ยถาภูตํ โอกฺขายติ.\csromanspacer{} ยมฺปิทํ จกฺขุสมฺผสฺสปจฺจยา อุปฺปชฺชติ เวทยิตํ สุขํ วา ทุกฺขํ วา อทุกฺขมสุขํ วา, ตมฺปิ อนิจฺจนฺติ ยถาภูตํ โอกฺขายติ ฯเปฯ.\csromanspacer{} มโน อนิจฺโจติ ยถาภูตํ โอกฺขายติ.\csromanspacer{} ธมฺมา.\csromanspacer{} มโนวิญฺญาณํ.\csromanspacer{} มโนสมฺผสฺโส.\csromanspacer{} ยมฺปิทํ มโนสมฺผสฺสปจฺจยา อุปฺปชฺชติ เวทยิตํ สุขํ วา ทุกฺขํ วา อทุกฺขมสุขํ วา, ตมฺปิ อนิจฺจนฺติ ยถาภูตํ โอกฺขายติ.\csromanspacer{} ปฏิสลฺลาเน ภิกฺขเว โยคมาปชฺชถ, ปฏิสลฺลีนสฺส ภิกฺขเว ภิกฺขุโน ยถาภูตํ โอกฺขายตีติ.\csromanspacer{}\csromanspacer{}ฉฏฺฐํ.}

\csromanmatikaanchor{mtk.193}
\csromantocmarkref{subsection}{7-9. โกฏฺฐิก-อนิจฺจสุตฺตาทิ}{mtk.193}
\bookmark[dest={mtk.193},level=2]{7-9. โกฏฺฐิก-อนิจฺจสุตฺตาทิ}
\csromanheader{2}{7. โกฏฺฐิก-อนิจฺจสุตฺต}
\proseitem{162}{อถ โข อายสฺมา มหาโกฏฺฐิโก เยน ภควา เตนุปสงฺกมิ ฯเปฯ เอกมนฺตํ นิสินฺโน โข อายสฺมา โกฏฺฐิโก ภควนฺตํ}

\vspace{\tipitakaparskip}
\csromancenter{สฬายตนสํยุตฺต เอตทโวจ “สาธุ เม ภนฺเต ภควา สํขิตฺเตน ธมฺมํ เทเสตุ, ยมหํ ภควโต ธมฺมํ สุตฺวา เอโก วูปกฏฺโฐ อปฺปมตฺโต อาตาปี ปหิตตฺโต วิหเรยฺยนฺ”ติ.}
\prose{\csromanfolio{357}ยํ โข โกฏฺฐิก อนิจฺจํ, ตตฺร เต ฉนฺโท ปหาตพฺโพ.\csromanspacer{} กิญฺจ โกฏฺฐิก อนิจฺจํ.\csromanspacer{} จกฺขุ โข โกฏฺฐิก อนิจฺจํ, ตตฺร เต ฉนฺโท ปหาตพฺโพ.\csromanspacer{} รูปา อนิจฺจา, ตตฺร เต ฉนฺโท ปหาตพฺโพ.\csromanspacer{} จกฺขุวิญฺญาณํ อนิจฺจํ, ตตฺร เต ฉนฺโท ปหาตพฺโพ.\csromanspacer{} จกฺขุสมฺผสฺโส อนิจฺโจ, ตตฺร เต ฉนฺโท ปหาตพฺโพ.\csromanspacer{} ยมฺปิทํ จกฺขุสมฺผสฺสปจฺจยา อุปฺปชฺชติ เวทยิตํ สุขํ วา ทุกฺขํ วา อทุกฺขมสุขํ วา, ตมฺปิ อนิจฺจํ, ตตฺร เต ฉนฺโท ปหาตพฺโพ ฯเปฯ.\csromanspacer{} ชิวฺหา อนิจฺจา, ตตฺร เต ฉนฺโท ปหาตพฺโพ.\csromanspacer{} รสา อนิจฺจา, ตตฺร เต ฉนฺโท ปหาตพฺโพ.\csromanspacer{} ชิวฺหาวิญฺญาณํ อนิจฺจํ, ตตฺร เต ฉนฺโท ปหาตพฺโพ.\csromanspacer{} ชิวฺหาสมฺผสฺโส อนิจฺโจ, ตตฺร เต ฉนฺโท ปหาตพฺโพ.\csromanspacer{} ยมฺปิทํ ชิวฺหาสมฺผสฺสปจฺจยา อุปฺปชฺชติ เวทยิตํ สุขํ วา ทุกฺขํ วา อทุกฺขมสุขํ วา, ตมฺปิ อนิจฺจํ, ตตฺร เต ฉนฺโท ปหาตพฺโพ ฯเปฯ.\csromanspacer{} มโน อนิจฺโจ, ตตฺร เต ฉนฺโท ปหาตพฺโพ.\csromanspacer{} ธมฺมา อนิจฺจา, ตตฺร เต ฉนฺโท ปหาตพฺโพ.\csromanspacer{} มโนวิญฺญาณํ อนิจฺจํ, ตตฺร เต ฉนฺโท ปหาตพฺโพ.\csromanspacer{} มโนสมฺผสฺโส อนิจฺโจ, ตตฺร เต ฉนฺโท ปหาตพฺโพ.\csromanspacer{} ยมฺปิทํ มโนสมฺผสฺสปจฺจยา อุปฺปชฺชติ เวทยิตํ สุขํ วา ทุกฺขํ วา อทุกฺขมสุขํ วา, ตมฺปิ อนิจฺจํ, ตตฺร เต ฉนฺโท ปหาตพฺโพ.\csromanspacer{} ยํ โข โกฏฺฐิก อนิจฺจํ, ตตฺร เต ฉนฺโท ปหาตพฺโพติ.\csromanspacer{}\csromanspacer{}สตฺตมํ.}

\csromanheader{2}{8. โกฏฺฐิกทุกฺขสุตฺต}
\proseitem{163}{อถ โข อายสฺมา มหาโกฏฺฐิโก ฯเปฯ ภควนฺตํ เอตทโวจ “สาธุ เม ภนฺเต ฯเปฯ วิหเรยฺยนฺ”ติ.\csromanspacer{} ยํ โข โกฏฺฐิก ทุกฺขํ, ตตฺร เต ฉนฺโท ปหาตพฺโพ.\csromanspacer{} กิญฺจ โกฏฺฐิก ทุกฺขํ.\csromanspacer{} จกฺขุ โข โกฏฺฐิก ทุกฺขํ, ตตฺร เต ฉนฺโท ปหาตพฺโพ.\csromanspacer{} รูปา ทุกฺขา, ตตฺร เต ฉนฺโท ปหาตพฺโพ.\csromanspacer{} จกฺขุวิญฺญาณํ ทุกฺขํ, ตตฺร เต ฉนฺโท ปหาตพฺโพ.\csromanspacer{} จกฺขุสมฺผสฺโส ทุกฺโข, ตตฺร เต ฉนฺโท ปหาตพฺโพ.\csromanspacer{} ยมฺปิทํ จกฺขุสมฺผสฺสปจฺจยา อุปฺปชฺชติ เวทยิตํ สุขํ วา ทุกฺขํ วา อทุกฺขมสุขํ วา, ตมฺปิ ทุกฺขํ, ตตฺร เต ฉนฺโท ปหาตพฺโพ ฯเปฯ.\csromanspacer{} ชิวฺหา ทุกฺขา, ตตฺร เต ฉนฺโท ปหาตพฺโพ ฯเปฯ.\csromanspacer{} มโน ทุกฺโข, ตตฺร เต ฉนฺโท ปหาตพฺโพ.}

\gambhira{สฬายตนวคฺคสํยุตฺตปาฬิ}
\csromanmatikaanchor{mtk.194}
\csromantocmarkref{subsection}{10-12. มิจฺฉาทิฏฺฐิปหานสุตฺตาทิ}{mtk.194}
\bookmark[dest={mtk.194},level=2]{10-12. มิจฺฉาทิฏฺฐิปหานสุตฺตาทิ}
\prose{\csromanfolio{358}ธมฺมา ทุกฺขา, ตตฺร เต ฉนฺโท ปหาตพฺโพ.\csromanspacer{} มโนวิญฺญาณํ ทุกฺขํ, ตตฺร เต ฉนฺโท ปหาตพฺโพ.\csromanspacer{} มโนสมฺผสฺโส ทุกฺโข, ตตฺร เต ฉนฺโท ปหตาตพฺโพ.\csromanspacer{} ยมฺปิทํ มโนสมฺผสฺสปจฺจยา อุปฺปชฺชติ เวทยิตํ สุขํ วา ทุกฺขํ วา อทุกฺขมสุขํ วา, ตมฺปิ ทุกฺขํ, ตตฺร เต ฉนฺโท ปหาตพฺโพ.\csromanspacer{} ยํ โข โกฏฺฐิก ทุกฺขํ, ตตฺร เต ฉนฺโท ปหาตพฺโพติ.\csromanspacer{}\csromanspacer{}อฏฺฐมํ.}

\csromanheader{2}{9. โกฏฺฐิก-อนตฺตสุตฺต}
\proseitem{164}{เอกมนฺตํ ฯเปฯ วิหเรยฺยนฺติ.\csromanspacer{} โย โข โกฏฺฐิก อนตฺตา, ตตฺร เต ฉนฺโท ปหาตพฺโพ.\csromanspacer{} โก จ โกฏฺฐิก อนตฺตา.\csromanspacer{} จกฺขุ โข โกฏฺฐิก อนตฺตา, ตตฺร เต ฉนฺโท ปหาตพฺโพ.\csromanspacer{} รูปา อนตฺตา, ตตฺร เต ฉนฺโท ปหาตพฺโพ.\csromanspacer{} จกฺขุวิญฺญาณํ อนตฺตา, ตตฺร เต ฉนฺโท ปหาตพฺโพ.\csromanspacer{} จกฺขุสมฺผสฺโส อนตฺตา, ตตฺร เต ฉนฺโท ปหาตพฺโพ.\csromanspacer{} ยมฺปิทํ จกฺขุสมฺผสฺสปจฺจยา อุปฺปชฺชติ เวทยิตํ สุขํ วา ทุกฺขํ วา อทุกฺขมสุขํ วา, ตมฺปิ อนตฺตา, ตตฺร เต ฉนฺโท ปหาตพฺโพ ฯเปฯ.\csromanspacer{} ชิวฺหา อนตฺตา, ตตฺร เต ฉนฺโท ปหาตพฺโพ ฯเปฯ.\csromanspacer{} มโน อนตฺตา, ตตฺร เต ฉนฺโท ปหาตพฺโพ.\csromanspacer{} ธมฺมา อนตฺตา, ตตฺร เต ฉนฺโท ปหาตพฺโพ.\csromanspacer{} มโนวิญฺญาณํ.\csromanspacer{} มโนสมฺผสฺโส.\csromanspacer{} ยมฺปิทํ มโนสมฺผสฺสปจฺจยา อุปฺปชฺชติ เวทยิตํ สุขํ วา ทุกฺขํ วา อทุกฺขมสุขํ วา, ตมฺปิ อนตฺตา, ตตฺร เต ฉนฺโท ปหาตพฺโพ.\csromanspacer{} โย โข โกฏฺฐิก อนตฺตา, ตตฺร เต ฉนฺโท ปหาตพฺโพติ.\csromanspacer{}\csromanspacer{}นวมํ.}

\csromanheader{2}{10. มิจฺฉาทิฏฺฐิปหานสุตฺต}
\proseitem{165}{อถ โข อญฺญตโร ภิกฺขุ เยน ภควา เตนุปสงฺกมิ ฯเปฯ เอกมนฺตํ นิสินฺโน โส ภิกฺขุ ภควนฺตํ เอตทโวจ “กถํ นุ โข ภนฺเต ชานโต กถํ ปสฺสโต มิจฺฉาทิฏฺฐิ ปหียตี”ติ.}

\prose{จกฺขุํ โข ภิกฺขุ อนิจฺจโต ชานโต ปสฺสโต มิจฺฉาทิฏฺฐิ ปหียติ.\csromanspacer{} รูเป อนิจฺจโต ชานโต ปสฺสโต มิจฺฉาทิฏฺฐิ ปหียติ.\csromanspacer{} จกฺขุวิญฺญาณํ อนิจฺจโต ชานโต ปสฺสโต มิจฺฉาทิฏฺฐิ ปหียติ.\csromanspacer{} จกฺขุสมฺผสฺสํ อนิจฺจโต ชานโต ปสฺสโต มิจฺฉาทิฏฺฐิ ปหียติ ฯเปฯ.\csromanspacer{} ยมฺปิทํ มโนสมฺผสฺสปจฺจยา อุปฺปชฺชติ เวทยิตํ สุขํ วา ทุกฺขํ วา อทุกฺขมสุขํ วา, ตมฺปิ อนิจฺจโต ชานโต ปสฺสโต มิจฺฉาทิฏฺฐิ ปหียติ.\csromanspacer{} เอวํ โข ภิกฺขุ ชานโต เอวํ ปสฺสโต มิจฺฉาทิฏฺฐิ ปหียตีติ.\csromanspacer{}\csromanspacer{}ทสมํ.}

\csromanheader{2}{1. สฬายตนสํยุตฺต \textbf{11. สกฺกายทิฏฺฐิปหานสุตฺต}}
\proseitem{166}{\csromanfolio{359}อถ โข อญฺญตโร ภิกฺขุ ฯเปฯ เอตทโวจ “กถํ นุ โข ภนฺเต ชานโต กถํ ปสฺสโต สกฺกายทิฏฺฐิ ปหียตี”ติ.\csromanspacer{} จกฺขุํ โข ภิกฺขุ ทุกฺขโต ชานโต ปสฺสโต สกฺกายทิฏฺฐิ ปหียติ.\csromanspacer{} รูเป ทุกฺขโต ชานโต ปสฺสโต สกฺกายทิฏฺฐิ ปหียติ.\csromanspacer{} จกฺขุวิญฺญาณํ ทุกฺขโต ชานโต ปสฺสโต สกฺกายทิฏฺฐิ ปหียติ.\csromanspacer{} จกฺขุสมฺผสฺสํ ทุกฺขโต ชานโต ปสฺสโต สกฺกายทิฏฺฐิ ปหียติ ฯเปฯ.\csromanspacer{} ยมฺปิทํ มโนสมฺผสฺสปจฺจยา อุปฺปชฺชติ เวทยิตํ สุขํ วา ทุกฺขํ วา อทุกฺขมสุขํ วา, ตมฺปิ ทุกฺขโต ชานโต ปสฺสโต สกฺกายทิฏฺฐิ ปหียติ.\csromanspacer{} เอวํ โข ภิกฺขุ ชานโต เอวํ ปสฺสโต สกฺกายทิฏฺฐิ ปหียตีติ.\csromanspacer{}\csromanspacer{}เอกาทสมํ.}

\csromanheader{2}{12. อตฺตานุทิฏฺฐิปหานสุตฺต}
\proseitem{167}{อถ โข อญฺญตโร ภิกฺขุ ฯเปฯ เอตทโวจ “กถํ นุ โข ภนฺเต ชานโต กถํ ปสฺสโต อตฺตานุทิฏฺฐิ ปหียตี”ติ.\csromanspacer{} จกฺขุํ โข ภิกฺขุ อนตฺตโต ชานโต ปสฺสโต อตฺตานุทิฏฺฐิ ปหียติ.\csromanspacer{} รูเป อนตฺตโต ชานโต ปสฺสโต อตฺตานุทิฏฺฐิ ปหียติ.\csromanspacer{} จกฺขุวิญฺญาณํ อนตฺตโต ชานโต ปสฺสโต อตฺตานุทิฏฺฐิ ปหียติ.\csromanspacer{} จกฺขุสมฺผสฺสํ อนตฺตโต ชานโต ปสฺสโต อตฺตานุทิฏฺฐิ ปหียติ.\csromanspacer{} ยมฺปิทํ จกฺขุสมฺผสฺสปจฺจยา อุปฺปชฺชติ เวทยิตํ สุขํ วา ทุกฺขํ วา อทุกฺขมสุขํ วา, ตมฺปิ อนตฺตโต ชานโต ปสฺสโต อตฺตานุทิฏฺฐิ ปหียติ ฯเปฯ.\csromanspacer{} ชิวฺหํ อนตฺตโต ชานโต ปสฺสโต อตฺตานุทิฏฺฐิ ปหียติ ฯเปฯ.\csromanspacer{} มนํ อนตฺตโต ชานโต ปสฺสโต อตฺตานุทิฏฺฐิ ปหียติ.\csromanspacer{} ธมฺเม.\csromanspacer{} มโนวิญฺญาณํ.\csromanspacer{} มโนสมฺผสฺสํ.\csromanspacer{} ยมฺปิทํ มโนสมฺผสฺสปจฺจยา อุปฺปชฺชติ เวทยิตํ สุขํ วา ทุกฺขํ วา อทุกฺขมสุขํ วา, ตมฺปิ อนตฺตโต ชานโต ปสฺสโต อตฺตานุทิฏฺฐิ ปหียตีติ.\csromanspacer{}\csromanspacer{}ทฺวาทสมํ.\csromanspacer{} นนฺทิกฺขยวคฺโค ปฐโม.}

\titlehead{\textbf{ตสฺสุทฺทานํ}}
\csromangathameasure{นนฺทิกฺขเยน จตฺตาโร, ชีวกมฺพวเน ทุเว. \\ โกฏฺฐิเกน ตโย วุตฺตา, มิจฺฉา สกฺกาย อตฺตโนติ.}
\csromangathagroup{\csromangathastanza{นนฺทิกฺขเยน จตฺตาโร, ชีวกมฺพวเน ทุเว. \\ โกฏฺฐิเกน ตโย วุตฺตา, มิจฺฉา สกฺกาย อตฺตโนติ.}}

\gambhira{สฬายตนวคฺคสํยุตฺตปาฬิ}
\csromanmatikaanchor{mtk.195}
\csromantocmarknopage{section}{(17) 2. สฏฺฐิเปยฺยาลวคฺค}{mtk.195}
\bookmark[dest={mtk.195},level=1]{(17) 2. สฏฺฐิเปยฺยาลวคฺค}
\csromanheader{1}{\textbf{(17) 2. สฏฺฐิเปยฺยาลวคฺค}}
\csromanmatikaanchor{mtk.196}
\csromantocmarkref{subsection}{1-60. อชฺฌตฺต-อนิจฺจฉนฺทสุตฺต}{mtk.196}
\bookmark[dest={mtk.196},level=2]{1-60. อชฺฌตฺต-อนิจฺจฉนฺทสุตฺต}
\csromanheader{2}{1. อชฺฌตฺต-อนิจฺจฉนฺทสุตฺต}
\proseitem{168}{\csromanfolio{360}ยํ ภิกฺขเว อนิจฺจํ, ตตฺร โว ฉนฺโท ปหาตพฺโพ.\csromanspacer{} กิญฺจ ภิกฺขเว อนิจฺจํ.\csromanspacer{} จกฺขุ ภิกฺขเว อนิจฺจํ, ตตฺร โว ฉนฺโท ปหาตพฺโพ ฯเปฯ.\csromanspacer{} ชิวฺหา อนิจฺจา, ตตฺร โว ฉนฺโท ปหาตพฺโพ ฯเปฯ.\csromanspacer{} มโน อนิจฺโจ, ตตฺร โว ฉนฺโท ปหาตพฺโพ.\csromanspacer{} ยํ ภิกฺขเว อนิจฺจํ, ตตฺร โว ฉนฺโท ปหาตพฺโพติ.}

\csromanheader{2}{2. อชฺฌตฺต-อนิจฺจราคสุตฺต}
\proseitem{169}{ยํ ภิกฺขเว อนิจฺจํ, ตตฺร โว ราโค ปหาตพฺโพ.\csromanspacer{} กิญฺจ ภิกฺขเว อนิจฺจํ.\csromanspacer{} จกฺขุ ภิกฺขเว อนิจฺจํ, ตตฺร โว ราโค ปหาตพฺโพ ฯเปฯ.\csromanspacer{} ชิวฺหา อนิจฺจา, ตตฺร โว ราโค ปหาตพฺโพ ฯเปฯ.\csromanspacer{} มโน อนิจฺโจ, ตตฺร โว ราโค ปหาตพฺโพ.\csromanspacer{} ยํ ภิกฺขเว อนิจฺจํ, ตตฺร โว ราโค ปหาตพฺโพติ.}

\csromanheader{2}{3. อชฺฌตฺต-อนิจฺจฉนฺทราคสุตฺต}
\proseitem{170}{ยํ ภิกฺขเว อนิจฺจํ, ตตฺร โว ฉนฺทราโค ปหาตพฺโพ.\csromanspacer{} กิญฺจ ภิกฺขเว อนิจฺจํ.\csromanspacer{} จกฺขุ ภิกฺขเว อนิจฺจํ, ตตฺร โว ฉนฺทราโค ปหาตพฺโพ ฯเปฯ.\csromanspacer{} ชิวฺหา อนิจฺจา, ตตฺร โว ฉนฺทราโค ปหาตพฺโพฯเปฯ.\csromanspacer{} มโน อนิจฺโจ, ตตฺร โว ฉนฺทราโค ปหาตพฺโพ.\csromanspacer{} ยํ ภิกฺขเว อนิจฺจํ, ตตฺร โว ฉนฺทราโค ปหาตพฺโพติ.}

\csromanheader{2}{4-6. ทุกฺขฉนฺทาทิสุตฺต}
\proseitem{171-173}{ยํ ภิกฺขเว ทุกฺขํ, ตตฺร โว ฉนฺโท ปหาตพฺโพ.\csromanspacer{} ราโค ปหาตพฺโพ.\csromanspacer{} ฉนฺทราโค ปหาตพฺโพ.\csromanspacer{} กิญฺจ ภิกฺขเว ทุกฺขํ.\csromanspacer{} จกฺขุ ภิกฺขเว ทุกฺขํ, ตตฺร โว ฉนฺโท ปหาตพฺโพ.\csromanspacer{} ราโค ปหาตพฺโพ.\csromanspacer{} ฉนฺทราโค ปหาตพฺโพ ฯเปฯ.\csromanspacer{} ชิวฺหา ทุกฺขา ฯเปฯ.\csromanspacer{} มโน ทุกฺโข, ตตฺร โว ฉนฺโท ปหาตพฺโพ.\csromanspacer{} ราโค ปหาตพฺโพ.\csromanspacer{} ฉนฺทราโค ปหาตพฺโพ.\csromanspacer{} ยํ ภิกฺขเว ทุกฺขํ, ตตฺร โว ฉนฺโท ปหาตพฺโพ.\csromanspacer{} ราโค ปหาตพฺโพ.\csromanspacer{} ฉนฺทราโค ปหาตพฺโพติ.}

\csromanheader{2}{1. สฬายตนสํยุตฺต \textbf{7-9. อนตฺตฉนฺทาทิสุตฺต}}
\proseitem{174-176}{\csromanfolio{361}โย ภิกฺขเว อนตฺตา, ตตฺร โว ฉนฺโท ปหาตพฺโพ.\csromanspacer{} ราโค ปหาตพฺโพ.\csromanspacer{} ฉนฺทราโค ปหาตพฺโพ.\csromanspacer{} โก จ ภิกฺขเว อนตฺตา.\csromanspacer{} จกฺขุ ภิกฺขเว อนตฺตา, ตตฺร โว ฉนฺโท ปหาตพฺโพ.\csromanspacer{} ราโค ปหาตพฺโพ.\csromanspacer{} ฉนฺทราโค ปหาตพฺโพ ฯเปฯ.\csromanspacer{} ชิวฺหา อนตฺตา, ตตฺร โว ฉนฺโท ปหาตพฺโพ.\csromanspacer{} ราโค ปหาตพฺโพ.\csromanspacer{} ฉนฺทราโค ปหาตพฺโพ ฯเปฯ.\csromanspacer{} มโน อนตฺตา, ตตฺร โว ฉนฺโท ปหาตพฺโพ.\csromanspacer{} ราโค ปหาตพฺโพ.\csromanspacer{} ฉนฺทราโค ปหาตพฺโพ.\csromanspacer{} โย ภิกฺขเว อนตฺตา, ตตฺร โว ฉนฺโท ปหาตพฺโพ.\csromanspacer{} ราโค ปหาตพฺโพ.\csromanspacer{} ฉนฺทราโค ปหาตพฺโพติ.}

\csromanheader{2}{10-12. พาหิรานิจฺจฉนฺทาทิสุตฺต}
\proseitem{177-179}{ยํ ภิกฺขเว อนิจฺจํ, ตตฺร โว ฉนฺโท ปหาตพฺโพ.\csromanspacer{} ราโค ปหาตพฺโพ.\csromanspacer{} ฉนฺทราโค ปหาตพฺโพ.\csromanspacer{} กิญฺจ ภิกฺขเว อนิจฺจํ.\csromanspacer{} รูปา ภิกฺขเว อนิจฺจา, ตตฺร โว ฉนฺโท ปหาตพฺโพ.\csromanspacer{} ราโค ปหาตพฺโพ.\csromanspacer{} ฉนฺทราโค ปหาตพฺโพ.\csromanspacer{} สทฺทา อนิจฺจา, ตตฺร โว ฉนฺโท ปหาตพฺโพ.\csromanspacer{} ราโค ปหาตพฺโพ.\csromanspacer{} ฉนฺทราโค ปหาตพฺโพ.\csromanspacer{} คนฺธา อนิจฺจา, ตตฺร โว ฉนฺโท ปหาตพฺโพ.\csromanspacer{} ราโค ปหาตพฺโพ.\csromanspacer{} ฉนฺทราโค ปหาตพฺโพ.\csromanspacer{} รสา อนิจฺจา, ตตฺร โว ฉนฺโท ปหาตพฺโพ.\csromanspacer{} ราโค ปหาตพฺโพ.\csromanspacer{} ฉนฺทราโค ปหาตพฺโพ.\csromanspacer{} โผฏฺฐพฺพา อนิจฺจา, ตตฺร โว ฉนฺโท ปหาตพฺโพ.\csromanspacer{} ราโค ปหาตพฺโพ.\csromanspacer{} ฉนฺทราโค ปหาตพฺโพ.\csromanspacer{} ธมฺมา อนิจฺจา, ตตฺร โว ฉนฺโท ปหาตพฺโพ.\csromanspacer{} ราโค ปหาตพฺโพ.\csromanspacer{} ฉนฺทราโค ปหาตพฺโพ.\csromanspacer{} ยํ ภิกฺขเว อนิจฺจํ, ตตฺร โว ฉนฺโท ปหาตพฺโพ.\csromanspacer{} ราโค ปหาตพฺโพ.\csromanspacer{} ฉนฺทราโค ปหาตพฺโพติ.}

\csromanheader{2}{13-15. พาหิรทุกฺขฉนฺทาทิสุตฺต}
\proseitem{180-182}{ยํ ภิกฺขเว ทุกฺขํ, ตตฺร โว ฉนฺโท ปหาตพฺโพ.\csromanspacer{} ราโค ปหาตพฺโพ.\csromanspacer{} ฉนฺทราโค ปหาตพฺโพ.\csromanspacer{} กิญฺจ ภิกฺขเว ทุกฺขํ.\csromanspacer{} รูปา ภิกฺขเว ทุกฺขา, ตตฺร โว ฉนฺโท ปหาตพฺโพ.\csromanspacer{} ราโค ปหาตพฺโพ.\csromanspacer{} ฉนฺทราโค ปหาตพฺโพ.\csromanspacer{} สทฺทา.\csromanspacer{} คนฺธา.\csromanspacer{} รสา.\csromanspacer{} โผฏฺฐพฺพา.\csromanspacer{} ธมฺมา ทุกฺขา, ตตฺร โว ฉนฺโท ปหาตพฺโพ.\csromanspacer{} ราโค ปหาตพฺโพ.\csromanspacer{} ฉนฺทราโค ปหาตพฺโพ.\csromanspacer{} ยํ ภิกฺขเว ทุกฺขํ, ตตฺร โว ฉนฺโท ปหาตพฺโพ.\csromanspacer{} ราโค ปหาตพฺโพ.\csromanspacer{} ฉนฺทราโค ปหาตพฺโพติ.}

\gambhira{สฬายตนวคฺคสํยุตฺตปาฬิ}
\csromanheader{2}{16-18. พาหิรานตฺตฉนฺทาทิสุตฺต}
\proseitem{183-185}{\csromanfolio{362}โย ภิกฺขเว อนตฺตา, ตตฺร โว ฉนฺโท ปหาตพฺโพ.\csromanspacer{} ราโค ปหาตพฺโพ.\csromanspacer{} ฉนฺทราโค ปหาตพฺโพ.\csromanspacer{} โก จ ภิกฺขเว อนตฺตา.\csromanspacer{} รูปา ภิกฺขเว อนตฺตา, ตตฺร โว ฉนฺโท ปหาตพฺโพ.\csromanspacer{} ราโค ปหาตพฺโพ.\csromanspacer{} ฉนฺทราโค ปหาตพฺโพ.\csromanspacer{} สทฺทา.\csromanspacer{} คนฺธา.\csromanspacer{} รสา.\csromanspacer{} โผฏฺฐพฺพา.\csromanspacer{} ธมฺมา อนตฺตา, ตตฺร โว ฉนฺโท ปหาตพฺโพ.\csromanspacer{} ราโค ปหาตพฺโพ.\csromanspacer{} ฉนฺทราโค ปหาตพฺโพ.\csromanspacer{} โย ภิกฺขเว อนตฺตา, ตตฺร โว ฉนฺโท ปหาตพฺโพ.\csromanspacer{} ราโค ปหาตพฺโพ.\csromanspacer{} ฉนฺทราโค ปหาตพฺโพติ.}

\csromanheader{2}{19. อชฺฌตฺตาตีตานิจฺจสุตฺต}
\proseitem{186}{จกฺขุ ภิกฺขเว อนิจฺจํ อตีตํ ฯเปฯ.\csromanspacer{} ชิวฺหา อนิจฺจา อตีตา ฯเปฯ.\csromanspacer{} มโน อนิจฺโจ อตีโต.\csromanspacer{} เอวํ ปสฺสํ ภิกฺขเว สุตวา อริยสาวโก จกฺขุสฺมิมฺปิ นิพฺพินฺทติ ฯเปฯ ชิวฺหายปิ นิพฺพินฺทติ ฯเปฯ มนสฺมิมฺปิ นิพฺพินฺทติ, นิพฺพินฺทํ วิรชฺชติ, วิราคา วิมุจฺจติ, วิมุตฺตสฺมิํ “วิมุตฺตมฺ”อิติ ญาณํ โหติ, “ขีณา ชาติ, วุสิตํ พฺรหฺมจริยํ, กตํ กรณียํ, นาปรํ อิตฺถตฺตายา”ติ ปชานาตีติ.}

\csromanheader{2}{20. อชฺฌตฺตานาคตานิจฺจสุตฺต}
\proseitem{187}{จกฺขุ ภิกฺขเว อนิจฺจํ อนาคตํ ฯเปฯ.\csromanspacer{} ชิวฺหา อนิจฺจา อนาคตา ฯเปฯ.\csromanspacer{} มโน อนิจฺโจ อนาคโต.\csromanspacer{} เอวํ ปสฺสํ ฯเปฯ นาปรํ อิตฺถตฺตายาติ ปชานาตีติ.}

\csromanheader{2}{21. อชฺฌตฺตปจฺจุปฺปนฺนานิจฺจสุตฺต}
\proseitem{188}{จกฺขุ ภิกฺขเว อนิจฺจํ ปจฺจุปฺปนฺนํ ฯเปฯ.\csromanspacer{} ชิวฺหา อนิจฺจา ปจฺจุปฺปนฺนา ฯเปฯ.\csromanspacer{} มโน อนิจฺโจ ปจฺจุปฺปนฺโน.\csromanspacer{} เอวํ ปสฺสํ ฯเปฯ นาปรํ อิตฺถตฺตายาติ ปชานาตีติ.}

\csromanheader{2}{22-24. อชฺฌตฺตาตีตาทิทุกฺขสุตฺต}
\proseitem{189-191}{จกฺขุ ภิกฺขเว ทุกฺขํ อตีตํ.\csromanspacer{} อนาคตํ.\csromanspacer{} ปจฺจุปฺปนฺนํ ฯเปฯ.\csromanspacer{} ชิวฺหา ทุกฺขา อตีตา.\csromanspacer{} อนาคตา.\csromanspacer{} ปจฺจุปฺปนฺนา ฯเปฯ.\csromanspacer{} มโน ทุกฺโข อตีโต.\csromanspacer{} อนาคโต.\csromanspacer{} ปจฺจุปฺปนฺโน.\csromanspacer{} เอวํ ปสฺสํ ภิกฺขเว ฯเปฯ นาปรํ อิตฺถตฺตายาติ ปชานาตีติ.}

\csromanheader{2}{1. สฬายตนสํยุตฺต \textbf{25-27. อชฺฌตฺตาตีตาทิ-อนตฺตสุตฺต}}
\proseitem{192-194}{\csromanfolio{363}จกฺขุ ภิกฺขเว อนตฺตา อตีตํ.\csromanspacer{} อนาคตํ.\csromanspacer{} ปจฺจุปฺปนฺนํ ฯเปฯ.\csromanspacer{} ชิวฺหา อนตฺตา ฯเปฯ.\csromanspacer{} มโน อนตฺตา อตีโต.\csromanspacer{} อนาคโต.\csromanspacer{} ปจฺจุปฺปนฺโน.\csromanspacer{} เอวํ ปสฺสํ ฯเปฯ นาปรํ อิตฺถตฺตายาติ ปชานาตีติ.}

\csromanheader{2}{28-30. พาหิราตีตาทิ-อนิจฺจสุตฺต}
\proseitem{195-197}{รูปา ภิกฺขเว อนิจฺจา อตีตา.\csromanspacer{} อนาคตา.\csromanspacer{} ปจฺจุปฺปนฺนา.\csromanspacer{} สทฺทา.\csromanspacer{} คนฺธา.\csromanspacer{} รสา.\csromanspacer{} โผฏฺฐพฺพา.\csromanspacer{} ธมฺมา อนิจฺจา อตีตา.\csromanspacer{} อนาคตา.\csromanspacer{} ปจฺจุปฺปนฺนา.\csromanspacer{} เอวํ ปสฺสํ ฯเปฯ นาปรํ อิตฺถตฺตายาติ ปชานาตีติ.}

\csromanheader{2}{31-33. พาหิราตีตาทิทุกฺขสุตฺต}
\proseitem{198-200}{รูปา ภิกฺขเว ทุกฺขา อตีตา.\csromanspacer{} อนาคตา.\csromanspacer{} ปจฺจุปฺปนฺนา.\csromanspacer{} สทฺทา.\csromanspacer{} คนฺธา.\csromanspacer{} รสา.\csromanspacer{} โผฏฺฐพฺพา.\csromanspacer{} ธมฺมา ทุกฺขา อตีตา.\csromanspacer{} อนาคตา.\csromanspacer{} ปจฺจุปฺปนฺนา.\csromanspacer{} เอวํ ปสฺสํ ฯเปฯ นาปรํ อิตฺถตฺตายาติ ปชานาตีติ.}

\csromanheader{2}{34-36. พาหิราตีตาทิ-อนตฺตสุตฺต}
\proseitem{201-203}{รูปา ภิกฺขเว อนตฺตา อตีตา.\csromanspacer{} อนาคตา.\csromanspacer{} ปจฺจุปฺปนฺนา.\csromanspacer{} สทฺทา.\csromanspacer{} คนฺธา.\csromanspacer{} รสา.\csromanspacer{} โผฏฺฐพฺพา.\csromanspacer{} ธมฺมา อนตฺตา อตีตา.\csromanspacer{} อนาคตา.\csromanspacer{} ปจฺจุปฺปนฺนา.\csromanspacer{} เอวํ ปสฺสํ ฯเปฯ นาปรํ อิตฺถตฺตายาติ ปชานาตีติ.}

\csromanheader{2}{37. อชฺฌตฺตาตีตยทนิจฺจสุตฺต}
\proseitem{204}{จกฺขุ ภิกฺขเว อนิจฺจํ อตีตํ, ยทนิจฺจํ ตํ ทุกฺขํ, ยํ ทุกฺขํ ตทนตฺตา, ยทนตฺตา ตํ “เนตํ มม, เนโสหมสฺมิ, น เมโส อตฺตา”ติ เอวเมตํ ยถาภูตํ สมฺมปฺปญฺญาย ทฏฺฐพฺพํ ฯเปฯ.\csromanspacer{} ชิวฺหา อนิจฺจา อตีตา, ยทนิจฺจํ ตํ ทุกฺขํ, ยํ ทุกฺขํ ตทนตฺตา, ยทนตฺตา ตํ “เนตํ มม, เนโสหมสฺมิ, น เมโส อตฺตา”ติ เอวเมตํ ยถาภูตํ สมฺมปฺปญฺญาย ทฏฺฐพฺพํ ฯเปฯ.\csromanspacer{} มโน อนิจฺโจ อตีโต, ยทนิจฺจํ ตํ ทุกฺขํ, ยํ ทุกฺขํ ตทนตฺตา, ยทนตฺตา ตํ “เนตํ มม, เนโสหมสฺมิ, น เมโส อตฺตา”ติ เอวเมตํ ยถาภูตํ สมฺมปฺปญฺญาย ทฏฺฐพฺพํ.\csromanspacer{} เอวํ ปสฺสํ ฯเปฯ นาปรํ อิตฺถตฺตายาติ ปชานาตีติ.}

\gambhira{สฬายตนวคฺคสํยุตฺตปาฬิ}
\csromanheader{2}{38. อชฺฌตฺตานาคตยทนิจฺจสุตฺต}
\proseitem{205}{\csromanfolio{364}จกฺขุ ภิกฺขเว อนิจฺจํ อนาคตํ, ยทนิจฺจํ ตํ ทุกฺขํ, ยํ ทุกฺขํ ตทนตฺตา, ยทนตฺตา ตํ “เนตํ มม, เนโสหมสฺมิ, น เมโส อตฺตา”ติ เอวเมตํ ยถาภูตํ สมฺมปฺปญฺญาย ทฏฺฐพฺพํ ฯเปฯ.\csromanspacer{} ชิวฺหา อนิจฺจา อนาคตา, ยทนิจฺจํ ตํ ทุกฺขํ, ยํ ทุกฺขํ ตทนตฺตา, ยทนตฺตา ตํ “เนตํ มม, เนโสหมสฺมิ, น เมโส อตฺตา”ติ เอวเมตํ ยถาภูตํ สมฺมปฺปญฺญาย ทฏฺฐพฺพํ ฯเปฯ.\csromanspacer{} มโน อนิจฺโจ อนาคโต, ยทนิจฺจํ ตํ ทุกฺขํ, ยํ ทุกฺขํ ตทนตฺตา, ยทนตฺตา ตํ “เนตํ มม, เนโสหมสฺมิ, น เมโส อตฺตา”ติ เอวเมตํ ยถาภูตํ สมฺมปฺปญฺญาย ทฏฺฐพฺพํ ฯเปฯ.\csromanspacer{} เอวํ ปสฺสํ ภิกฺขเว ฯเปฯ นาปรํ อิตฺถตฺตายาติ ปชานาตีติ.}

\csromanheader{2}{39. อชฺฌตฺตปจฺจุปฺปนฺนยทนิจฺจสุตฺต}
\proseitem{206}{จกฺขุ ภิกฺขเว อนิจฺจํ ปจฺจุปฺปนฺนํ, ยทนิจฺจํ ตํ ทุกฺขํ, ยํ ทุกฺขํ ตทนตฺตา, ยทนตฺตา ตํ “เนตํ มม, เนโสหมสฺมิ, น เมโส อตฺตา”ติ เอวเมตํ ยถาภูตํ สมฺมปฺปญฺญาย ทฏฺฐพฺพํ ฯเปฯ ชิวฺหา อนิจฺจา ปจฺจุปฺปนฺนา, ยทนิจฺจํ ตํ ทุกฺขํ, ยํ ทุกฺขํ ตทนตฺตา, ยทนตฺตา ตํ “เนตํ มม, เนโสหมสฺมิ, น เมโส อตฺตา”ติ เอวเมตํ ยถาภูตํ สมฺมปฺปญฺญาย ทฏฺฐพฺพํ ฯเปฯมโน อนิจฺโจ ปจฺจุปฺปนฺโน, ยทนิจฺจํ ตํ ทุกฺขํ, ยํ ทุกฺขํ ตทนตฺตา, ยทนตฺตา ตํ “เนตํ มม, เนโสหมสฺมิ, น เมโส อตฺตา”ติ เอวเมตํ ยถาภูตํ สมฺมปฺปญฺญาย ทฏฺฐพฺพํ.\csromanspacer{} เอวํ ปสฺสํ ฯเปฯ นาปรํ อิตฺถตฺตายาติ ปชานาตีติ.}

\csromanheader{2}{40-42. อชฺฌตฺตาตีตาทิยํทุกฺขสุตฺต}
\proseitem{207-209}{จกฺขุ ภิกฺขเว ทุกฺขํ อตีตํ.\csromanspacer{} อนาคตํ.\csromanspacer{} ปจฺจุปฺปนฺนํ, ยํ ทุกฺขํ ตทนตฺตา, ยทนตฺตา ตํ “เนตํ มม, เนโสหมสฺมิ, น เมโส อตฺตา”ติ เอวเมตํ ยถาภูตํ สมฺมปฺปญฺญาย ทฏฺฐพฺพํ ฯเปฯ.\csromanspacer{} ชิวฺหา ทุกฺขา ฯเปฯ.\csromanspacer{} มโน ทุกฺโข อตีโต.\csromanspacer{} อนาคโต.\csromanspacer{} ปจฺจุปฺปนฺโน, ยํ ทุกฺขํ ตทนตฺตา, ยทนตฺตา ตํ “เนตํ มม, เนโสหมสฺมิ, น เมโส อตฺตา”ติ เอวเมตํ ยถาภูตํ สมฺมปฺปญฺญาย ทฏฺฐพฺพํ.\csromanspacer{} เอวํ ปสฺสํ ฯเปฯ นาปรํ อิตฺถตฺตายาติ ปชานาตีติ.}

\csromanheader{2}{43-45. อชฺฌตฺตาตีตาทิยทนตฺตสุตฺต}
\proseitem{210-212}{จกฺขุ ภิกฺขเว อนตฺตา อตีตํ.\csromanspacer{} อนาคตํ.\csromanspacer{} ปจฺจุปฺปนฺนํ, ยทนตฺตา ตํ “เนตํ มม, เนโสหมสฺมิ, น เมโส อตฺตา”ติ}

\vspace{\tipitakaparskip}
\csromancenter{สฬายตนสํยุตฺต เอวเมตํ ยถาภูตํ สมฺมปฺปญฺญาย ทฏฺฐพฺพํ ฯเปฯ.\csromanspacer{} ชิวฺหา อนตฺตา ฯเปฯ.\csromanspacer{} มโน อนตฺตา อตีโต.\csromanspacer{} อนาคโต.\csromanspacer{} ปจฺจุปฺปนฺโน, ยทนตฺตา ตํ “เนตํ มม, เนโสหมสฺมิ, น เมโส อตฺตา”ติ เอวเมตํ ยถาภูตํ สมฺมปฺปญฺญาย ทฏฺฐพฺพํ.\csromanspacer{} เอวํ ปสฺสํ ฯเปฯ นาปรํ อิตฺถตฺตายาติ ปชานาตีติ.}
\csromanheader{2}{46-48. พาหิราตีตาทิยทนิจฺจสุตฺต}
\proseitem{213-215}{\csromanfolio{365}รูปา ภิกฺขเว อนิจฺจา อตีตา.\csromanspacer{} อนาคตา.\csromanspacer{} ปจฺจุปฺปนฺนา, ยทนิจฺจํ ตํ ทุกฺขํ, ยํ ทุกฺขํ ตทนตฺตา, ยทนตฺตา ตํ “เนตํ มม, เนโสหมสฺมิ, น เมโส อตฺตา”ติ เอวเมตํ ยถาภูตํ สมฺมปฺปญฺญาย ทฏฺฐพฺพํ.\csromanspacer{} สทฺทา.\csromanspacer{} คนฺธา.\csromanspacer{} รสา.\csromanspacer{} โผฏฺฐพฺพา.\csromanspacer{} ธมฺมา อนิจฺจา อตีตา.\csromanspacer{} อนาคตา.\csromanspacer{} ปจฺจุปฺปนฺนา, ยทนิจฺจํ ตํ ทุกฺขํ, ยํ ทุกฺขํ ตทนตฺตา, ยทนตฺตา ตํ “เนตํ มม, เนโสหมสฺมิ, น เมโส อตฺตา”ติ เอวเมตํ ยถาภูตํ สมฺมปฺปญฺญาย ทฏฺฐพฺพํ.\csromanspacer{} เอวํ ปสฺสํ ฯเปฯ นาปรํ อิตฺถตฺตายาติ ปชานาตีติ.}

\csromanheader{2}{49-51. พาหิราตีตาทิยํทุกฺขสุตฺต}
\proseitem{216-218}{รูปา ภิกฺขเว ทุกฺขา อตีตา.\csromanspacer{} อนาคตา.\csromanspacer{} ปจฺจุปฺปนฺนา, ยํ ทุกฺขํ ตทนตฺตา, ยทนตฺตา ตํ “เนตํ มม, เนโสหมสฺมิ, น เมโส อตฺตา”ติ เอวเมตํ ยถาภูตํ สมฺมปฺปญฺญาย ทฏฺฐพฺพํ.\csromanspacer{} สทฺทา.\csromanspacer{} คนฺธา.\csromanspacer{} รสา.\csromanspacer{} โผฏฺฐพฺพา.\csromanspacer{} ธมฺมา ทุกฺขา อตีตา.\csromanspacer{} อนาคตา.\csromanspacer{} ปจฺจุปฺปนฺนา, ยํ ทุกฺขํ ตทนตฺตา, ยทนตฺตา ตํ “เนตํ มม, เนโสหมสฺมิ, น เมโส อตฺตา”ติ เอวเมตํ ยถาภูตํ สมฺมปฺปญฺญาย ทฏฺฐพฺพํ.\csromanspacer{} เอวํ ปสฺสํ ฯเปฯ นาปรํ อิตฺถตฺตายาติ ปชานาตีติ.}

\csromanheader{2}{52-54. พาหิราตีตาทิยทนตฺตสุตฺต}
\proseitem{219-221}{รูปา ภิกฺขเว อนตฺตา อตีตา.\csromanspacer{} อนาคตา.\csromanspacer{} ปจฺจุปฺปนฺนา, ยทนตฺตา ตํ “เนตํ มม, เนโสหมสฺมิ, น เมโส อตฺตา”ติ เอวเมตํ ยถาภูตํ สมฺมปฺปญฺญาย ทฏฺฐพฺพํ.\csromanspacer{} สทฺทา.\csromanspacer{} คนฺธา.\csromanspacer{} รสา.\csromanspacer{} โผฏฺฐพฺพา.\csromanspacer{} ธมฺมา อนตฺตา อตีตา.\csromanspacer{} อนาคตา.\csromanspacer{} ปจฺจุปฺปนฺนา, ยทนตฺตา ตํ “เนตํ มม, เนโสหมสฺมิ, น เมโส อตฺตา”ติ เอวเมตํ ยถาภูตํ สมฺมปฺปญฺญาย ทฏฺฐพฺพํ.\csromanspacer{} เอวํ ปสฺสํ ฯเปฯ นาปรํ อิตฺถตฺตายาติ ปชานาตีติ.}

\gambhira{สฬายตนวคฺคสํยุตฺตปาฬิ}
\csromanheader{2}{55. อชฺฌตฺตายตน-อนิจฺจสุตฺต}
\proseitem{222}{\csromanfolio{366}จกฺขุ ภิกฺขเว อนิจฺจํ ฯเปฯ.\csromanspacer{} ชิวฺหา อนิจฺจา ฯเปฯ.\csromanspacer{} มโน อนิจฺโจ.\csromanspacer{} เอวํ ปสฺสํ ฯเปฯ นาปรํ อิตฺถตฺตายาติ ปชานาตีติ.}

\csromanheader{2}{56. อชฺฌตฺตายตนทุกฺขสุตฺต}
\proseitem{223}{จกฺขุ ภิกฺขเว ทุกฺขํ ฯเปฯ.\csromanspacer{} ชิวฺหา ทุกฺขา ฯเปฯ.\csromanspacer{} มโน ทุกฺโข.\csromanspacer{} เอวํ ปสฺสํ ฯเปฯ นาปรํ อิตฺถตฺตายาติ ปชานาตีติ.}

\csromanheader{2}{57. อชฺฌตฺตายตน-อนตฺตสุตฺต}
\proseitem{224}{จกฺขุ ภิกฺขเว อนตฺตา ฯเปฯ.\csromanspacer{} ชิวฺหา อนตฺตา ฯเปฯ.\csromanspacer{} มโน อนตฺตา.\csromanspacer{} เอวํ ปสฺสํ ฯเปฯ นาปรํ อิตฺถตฺตายาติ ปชานาตีติ.}

\csromanheader{2}{58. พาหิรายตน-อนิจฺจสุตฺต}
\proseitem{225}{รูปา ภิกฺขเว อนิจฺจา.\csromanspacer{} สทฺทา.\csromanspacer{} คนฺธา.\csromanspacer{} รสา.\csromanspacer{} โผฏฺฐพฺพา.\csromanspacer{} ธมฺมา อนิจฺจา.\csromanspacer{} เอวํ ปสฺสํ ฯเปฯ นาปรํ อิตฺถตฺตายาติ ปชานาตีติ.}

\csromanheader{2}{59. พาหิรายตนทุกฺขสุตฺต}
\proseitem{226}{รูปา ภิกฺขเว ทุกฺขา.\csromanspacer{} สทฺทา.\csromanspacer{} คนฺธา.\csromanspacer{} รสา.\csromanspacer{} โผฏฺฐพฺพา.\csromanspacer{} ธมฺมา ทุกฺขา.\csromanspacer{} เอวํ ปสฺสํ ฯเปฯ นาปรํ อิตฺถตฺตายาติ ปชานาตีติ.}

\csromanheader{2}{60. พาหิรายตน-อนตฺตสุตฺต}
\proseitem{227}{รูปา ภิกฺขเว อนตฺตา.\csromanspacer{} สทฺทา.\csromanspacer{} คนฺธา.\csromanspacer{} รสา.\csromanspacer{} โผฏฺฐพฺพา.\csromanspacer{} ธมฺมา อนตฺตา.\csromanspacer{} เอวํ ปสฺสํ ฯเปฯ นาปรํ อิตฺถตฺตายาติ ปชานาตีติ.\csromanspacer{} สฏฺฐิเปยฺยาโล สมตฺโต.}

\titlehead{\textbf{ตสฺสุทฺทานํ}}
\csromangathameasure{ฉนฺเทนฏฺฐารส โหนฺติ, อตีเตน จ เทฺว นว. \\ ยทนิจฺจาฏฺฐารส วุตฺตา, ตโย อชฺฌตฺตพาหิรา. \\ เปยฺยาโล สฏฺฐิโก วุตฺโต, พุทฺเธนาทิจฺจพนฺธุนาติ. สุตฺตนฺตานิ สฏฺฐิ.}
\csromangathagroup{\csromangathastanza{ฉนฺเทนฏฺฐารส โหนฺติ, อตีเตน จ เทฺว นว. \\ ยทนิจฺจาฏฺฐารส วุตฺตา, ตโย อชฺฌตฺตพาหิรา. \\ เปยฺยาโล สฏฺฐิโก วุตฺโต, พุทฺเธนาทิจฺจพนฺธุนาติ.\csromanspacer{} สุตฺตนฺตานิ สฏฺฐิ.}}

\csromanmatikaanchor{mtk.197}
\csromantocmarknopage{section}{(18) 3. สมุทฺทวคฺค}{mtk.197}
\bookmark[dest={mtk.197},level=1]{(18) 3. สมุทฺทวคฺค}
\csromanheader{1}{1. สฬายตนสํยุตฺต \textbf{(18) 3. สมุทฺทวคฺค}}
\csromanheader{2}{1. ปฐมสมุทฺทสุตฺต}
\csromanmatikaanchor{mtk.198}
\csromantocmarkref{subsection}{1-2. สมุทฺทสุตฺต}{mtk.198}
\bookmark[dest={mtk.198},level=2]{1-2. สมุทฺทสุตฺต}
\proseitem{228}{\csromanfolio{367}“สมุทฺโท สมุทฺโท”ติ ภิกฺขเว อสฺสุตวา ปุถุชฺชโน ภาสติ, เนโส ภิกฺขเว อริยสฺส วินเย สมุทฺโท, มหา เอโส ภิกฺขเว อุทกราสิ มหา-อุทกณฺณโว.\csromanspacer{} จกฺขุ ภิกฺขเว ปุริสสฺส สมุทฺโท, ตสฺส รูปมโย เวโค.\csromanspacer{} โย ตํ รูปมยํ เวคํ สหติ, อยํ วุจฺจติ ภิกฺขเว อตริ จกฺขุสมุทฺทํ ส-อูมิํ สาวฏฺฏํ สคาหํ สรกฺขสํ, ติณฺโณ ปารงฺคโต ถเล ติฏฺฐติ พฺราหฺมโณ ฯเปฯ.\csromanspacer{} ชิวฺหา ภิกฺขเว ปุริสสฺส สมุทฺโท, ตสฺส รสมโย เวโค.\csromanspacer{} โย ตํ รสมยํ เวคํ สหติ, อยํ วุจฺจติ ภิกฺขเว อตริ ชิวฺหาสมุทฺทํ ส-อูมิํ สาวฏฺฏํ สคาหํ สรกฺขสํ, ติณฺโณ ปารงฺคโต ถเล ติฏฺฐติ พฺราหฺมโณ ฯเปฯ.\csromanspacer{} มโน ภิกฺขเว ปุริสสฺส สมุทฺโท, ตสฺส ธมฺมมโย เวโค.\csromanspacer{} โย ตํ ธมฺมมยํ เวคํ สหติ, อยํ วุจฺจติ ภิกฺขเว อตริ มโนสมุทฺทํ ส-อูมิํ สาวฏฺฏํ สคาหํ สรกฺขสํ, ติณฺโณ ปารงฺคโต ถเล ติฏฺฐติ พฺราหฺมโณติ.\csromanspacer{} อิทมโวจ ฯเปฯ สตฺถา\csromandash{}}

\vspace{\tipitakaparskip}
\csromancenter{“โย อิมํ สมุทฺทํ สคาหํ สรกฺขสํ.}
\csromancenter{ส-อูมิํ สาวฏฺฏํ สภยํ ทุตฺตรํ อจฺจตริ.}
\prose{ส เวทคู วุสิตพฺรหฺมจริโย.}

\prose{โลกนฺตคู ปารคโตติ วุจฺจติ”ติ.\csromanspacer{} ปฐมํ.}

\csromanheader{2}{2. ทุติยสมุทฺทสุตฺต}
\proseitem{229}{“สมุทฺโท สมุทฺโท”ติ ภิกฺขเว อสฺสุตวา ปุถุชฺชโน ภาสติ, เนโส ภิกฺขเว อริยสฺส วินเย สมุทฺโท, มหา เอโส ภิกฺขเว อุทกราสิ มหา-อุทกณฺณโว.\csromanspacer{} สนฺติ ภิกฺขเว จกฺขุวิญฺเญยฺยา รูปา อิฏฺฐา กนฺตา มนาปา ปิยรูปา กามูปสํหิตา รชนียา, อยํ วุจฺจติ ภิกฺขเว อริยสฺส วินเย สมุทฺโท.\csromanspacer{} เอตฺถายํ สเทวโก โลโก สมารโก สพฺรหฺมโก สสฺสมณพฺราหฺมณี ปชา สเทวมนุสฺสา เยภุยฺเยน สมุนฺนา ตนฺตากุลกชาตา กุลคณฺฐิกชาตา\footnote{คุฬาคุณฺฐิกชาตา (สี), กุลคุณฺฑิกชาตา (สฺยา, กํ), คุณคุณิกชาตา (อิ), กุลาคุณฺฑิกชาตา (ก)} มุญฺชปพฺพชภูตา อปายํ ทุคฺคติํ วินิปาตํ สํสารํ นาติวตฺตติ ฯเปฯ.}

\gambhira{สฬายตนวคฺคสํยุตฺตปาฬิ}
\prose{\csromanfolio{368}สนฺติ ภิกฺขเว ชิวฺหาวิญฺเญยฺยา รสา ฯเปฯ.\csromanspacer{} สนฺติ ภิกฺขเว มโนวิญฺเญยฺยา ธมฺมา อิฏฺฐา กนฺตา มนาปา ปิยรูปา กามูปสํหิตา รชนียา, อยํ วุจฺจติ ภิกฺขเว อริยสฺส วินเย สมุทฺโท.\csromanspacer{} เอตฺถายํ สเทวโก โลโก สมารโก สพฺรหฺมโก สสฺสมณพฺราหฺมณี ปชา สเทวมนุสฺสา เยภุยฺเยน สมุนฺนา ตนฺตากุลกชาตา กุลคณฺฐิกชาตา มุญฺชปพฺพชภูตา อปายํ ทุคฺคติํ วินิปาตํ สํสารํ นาติวตฺตตีติ.}

\csromangathameasure{ยสฺส ราโค จ โทโส จ, อวิชฺชา จ วิราชิตา. \\ โส อิมํ สมุทฺทํ สคาหํ สรกฺขสํ, \\ ส-อูมิภยํ ทุตฺตรํ อจฺจตริ. \\ สงฺคาติโค มจฺจุชโห นิรุปธิ, ปหาสิ ทุกฺขํ อปุนพฺภวาย. \\ อตฺถงฺคโต โส น ปุเนติ, อโมหยี มจฺจุราชนฺติ พฺรูมีติ. ทุติยํ.}
\csromangathagroup{\csromangathastanza{ยสฺส ราโค จ โทโส จ, อวิชฺชา จ วิราชิตา.}\csromangathastanza{โส อิมํ สมุทฺทํ สคาหํ สรกฺขสํ, \\ ส-อูมิภยํ ทุตฺตรํ อจฺจตริ.}\csromangathastanza{สงฺคาติโค มจฺจุชโห นิรุปธิ, ปหาสิ ทุกฺขํ อปุนพฺภวาย. \\ อตฺถงฺคโต โส น ปุเนติ\footnote{น ปมาณเมติ (สี, สฺยา, กํ, อิ)}, อโมหยี มจฺจุราชนฺติ พฺรูมีติ.\csromanspacer{} ทุติยํ.}}

\csromanheader{2}{3. พาฬิสิโกปมสุตฺต}
\csromanmatikaanchor{mtk.199}
\csromantocmarkref{subsection}{3. พาฬิสิโกปมสุตฺต}{mtk.199}
\bookmark[dest={mtk.199},level=2]{3. พาฬิสิโกปมสุตฺต}
\proseitem{230}{เสยฺยถาปิ ภิกฺขเว พาฬิสิโก อามิสคตพฬิสํ คมฺภีเร อุทกรหเท ปกฺขิเปยฺย, ตเมนํ อญฺญตโร อามิสจกฺขุ มจฺโฉ คิเลยฺย.\csromanspacer{} เอวํ หิ โส ภิกฺขเว มจฺโฉ คิลิตพฬิโส พาฬิสิกสฺส อนยํ อาปนฺโน พฺยสนํ อาปนฺโน ยถากามกรณีโย พาฬิสิกสฺส.}

\prose{เอวเมว โข ภิกฺขเว ฉยิเม พฬิสา โลกสฺมิํ อนยาย สตฺตานํ วธาย\footnote{พฺยาพาธาย (สี, อิ)} ปาณินํ.\csromanspacer{} กตเม ฉ, สนฺติ ภิกฺขเว จกฺขุวิญฺเญยฺยา รูปา อิฏฺฐา กนฺตา มนาปา ปิยรูปา กามูปสํหิตา รชนียา, ตญฺเจ ภิกฺขุ อภินนฺทติ อภิวทติ อชฺโฌสาย ติฏฺฐติ.\csromanspacer{} อยํ วุจฺจติ ภิกฺขเว ภิกฺขุ คิลิตพฬิโส มารสฺส, อนยํ อาปนฺโน พฺยสนํ อาปนฺโน ยถากามกรณีโย ปาปิมโต ฯเปฯ.\csromanspacer{} สนฺติ ภิกฺขเว ชิวฺหาวิญฺเญยฺยา รสา ฯเปฯ.}

\prose{สนฺติ ภิกฺขเว มโนวิญฺเญยฺยา ธมฺมา อิฏฺฐา กนฺตา มนาปา ปิยรูปา กามูปสํหิตา รชนียา, ตญฺเจ ภิกฺขุ อภินนฺทติ อภิวทติ อชฺโฌสาย ติฏฺฐติ.\csromanspacer{} อยํ วุจฺจติ ภิกฺขเว ภิกฺขุ คิลิตพฬิโส มารสฺส, อนยํ อาปนฺโน พฺยสนํ อาปนฺโน ยถากามกรณีโย ปาปิมโต.}

\csromanheader{2}{1. สฬายตนสํยุตฺต}
\prose{\csromanfolio{369}สนฺติ จ ภิกฺขเว จกฺขุวิญฺเญยฺยา รูปา อิฏฺฐา กนฺตา มนาปา ปิยรูปา กามูปสํหิตา รชนียา, ตญฺเจ ภิกฺขุ นาภินนฺทติ นาภิวทติ นาชฺโฌสาย ติฏฺฐติ.\csromanspacer{} อยํ วุจฺจติ ภิกฺขเว ภิกฺขุ น คิลิตพฬิโส มารสฺส, อเภทิ พฬิสํ, ปริเภทิ พฬิสํ, น อนยํ อาปนฺโน น พฺยสนํ อาปนฺโน น ยถากามรณีโย ปาปิมโต ฯเปฯ.}

\prose{สนฺติ ภิกฺขเว ชิวฺหาวิญฺเญยฺยา รสา ฯเปฯ.\csromanspacer{} สนฺติ ภิกฺขเว มโนวิญฺเญยฺยา ธมฺมา อิฏฺฐา กนฺตา มนาปา ปิยรูปา กามูปสํหิตา รชนียา, ตญฺเจ ภิกฺขุ นาภินนฺทติ นาภิวทติ นาชฺโฌสาย ติฏฺฐติ.\csromanspacer{} อยํ วุจฺจติ ภิกฺขเว ภิกฺขุ น คิลิตพฬิโส มารสฺส, อเภทิ พฬิสํ, ปริเภทิ พฬิสํ, น อนยํ อาปนฺโน น พฺยสนํ อาปนฺโน น ยถากามกรณีโย ปาปิมโตติ.\csromanspacer{}\csromanspacer{}ตติยํ.}

\csromanheader{2}{4. ขีรรุกฺโขปมสุตฺต}
\csromanmatikaanchor{mtk.200}
\csromantocmarkref{subsection}{4. ขีรรุกฺโขปมสุตฺต}{mtk.200}
\bookmark[dest={mtk.200},level=2]{4. ขีรรุกฺโขปมสุตฺต}
\proseitem{231}{ยสฺส กสฺสจิ ภิกฺขเว ภิกฺขุสฺส วา ภิกฺขุนิยา วา จกฺขุวิญฺเญยฺเยสุ รูเปสุ โย ราโค, โส อตฺถิ, โย โทโส, โส อตฺถิ, โย โมโห, โส อตฺถิ.\csromanspacer{} โย ราโค, โส อปฺปหีโน, โย โทโส, โส อปฺปหีโน, โย โมโห, โส อปฺปหีโน.\csromanspacer{} ตสฺส ปริตฺตา เจปิ จกฺขุวิญฺเญยฺยา รูปา จกฺขุสฺส อาปาถํ อาคจฺฉนฺติ, ปริยาทิยนฺเตวสฺส จิตฺตํ, โก ปน วาโท อธิมตฺตานํ.\csromanspacer{} ตํ กิสฺส เหตุ, โย ภิกฺขเว ราโค, โส อตฺถิ, โย โทโส, โส อตฺถิ, โย โมโห, โส อตฺถิ.\csromanspacer{} โย ราโค, โส อปฺปหีโน, โย โทโส, โส อปฺปหีโน, โย โมโห, โส อปฺปหีโน ฯเปฯ.}

\prose{ยสฺส กสฺสจิ ภิกฺขเว ภิกฺขุสฺส วา ภิกฺขุนิยา วา ชิวฺหาวิญฺเญยฺเยสุ รเสสุ โย ราโค โส อตฺถิ ฯเปฯ.}

\prose{ยสฺส กสฺสจิ ภิกฺขเว ภิกฺขุสฺส วา ภิกฺขุนิยา วา มโนวิญฺเญยฺเยสุ ธมฺเมสุ โย ราโค, โส อตฺถิ, โย โทโส, โส อตฺถิ, โย โมโห, โส อตฺถิ.\csromanspacer{} โย ราโค, โส อปฺปหีโน, โย โทโส, โส อปฺปหีโน, โย โมโห, โส อปฺปหีโน.\csromanspacer{} ตสฺส ปริตฺตา เจปิ มโนวิญฺเญยฺยา ธมฺมา มนสฺส อาปาถํ อาคจฺฉนฺติ, ปริยาทิยนฺเตวสฺส จิตฺตํ, โก ปน วาโท อธิมตฺตานํ.\csromanspacer{} ตํ กิสฺส เหตุ, โย ภิกฺขเว ราโค, โส อตฺถิ, โย โทโส, โส อตฺถิ, โย โมโห,}

\gambhira{สฬายตนวคฺคสํยุตฺตปาฬิ}
\prose{\csromanfolio{370}โส อตฺถิ.\csromanspacer{} โย ราโค, โส อปฺปหีโน, โย โทโส, โส อปฺปหีโน, โย โมโห, โส อปฺปหีโน.}

\prose{เสยฺยถาปิ ภิกฺขเว ขีรรุกฺโข อสฺสตฺโถ วา นิโคฺรโธ วา ปิลกฺโข วา อุทุมฺพโร วา ทหโร ตรุโณ โกมารโก, ตเมนํ ปุริโส ติณฺหาย กุฐาริยา ยโต ยโต อาภินฺเทยฺย\footnote{ภินฺเทยฺย (สฺยา, กํ, สี-ฏฺฐ), อภินฺเทยฺย (กตฺถจิ)} อาคจฺเฉยฺย ขีรนฺติ.\csromanspacer{} เอวํ ภนฺเต.\csromanspacer{} ตํ กิสฺส เหตุ, ยํ หิ ภนฺเต ขีรํ, ตํ อตฺถีติ.}

\prose{เอวเมว โข ภิกฺขเว ยสฺส กสฺสจิ ภิกฺขุสฺส วา ภิกฺขุนิยา วา จกฺขุวิญฺเญยฺเยสุ รูเปสุ โย ราโค, โส อตฺถิ, โย โทโส, โส อตฺถิ, โย โมโห, โส อตฺถิ.\csromanspacer{} โย ราโค, โส อปฺปหีโน, โย โทโส, โส อปฺปหีโน, โย โมโห, โส อปฺปหีโน.\csromanspacer{} ตสฺส ปริตฺตา เจปิ จกฺขุวิญฺเญยฺยา รูปา จกฺขุสฺส อาปาถํ อาคจฺฉนฺติ, ปริยาทิยนฺเตวสฺส จิตฺตํ, โก ปน วาโท อธิมตฺตานํ.\csromanspacer{} ตํ กิสฺส เหตุ, โย ภิกฺขเว ราโค, โส อตฺถิ, โย โทโส, โส อตฺถิ, โย โมโห โส อตฺถิ.\csromanspacer{} โย ราโค, โส อปฺปหีโน, โย โทโส, โส อปฺปหีโน, โย โมโห, โส อปฺปหีโน ฯเปฯ.}

\prose{ยสฺส กสฺสจิ ภิกฺขเว ภิกฺขุสฺส วา ภิกฺขุนิยา วา ชิวฺหาวิญฺเญยฺเยสุ รเสสุ โย ราโค, โส อตฺถิ ฯเปฯ.}

\prose{ยสฺส กสฺสจิ ภิกฺขเว ภิกฺขุสฺส วา ภิกฺขุนิยา วา มโนวิญฺเญยฺเยสุ ธมฺเมสุ โย ราโค, โส อตฺถิ, โย โทโส, โส อตฺถิ, โย โมโห, โส อตฺถิ.\csromanspacer{} โย ราโค, โส อปฺปหีโน, โย โทโส, โส อปฺปหีโน, โย โมโห โส อปฺปหีโน.\csromanspacer{} ตสฺส ปริตฺตา เจปิ มโนวิญฺเญยฺยา ธมฺมา มนสฺส อาปาถํ อาคจฺฉนฺติ, ปริยาทิยนฺเตวสฺส จิตฺตํ, โก ปน วาโท อธิมตฺตานํ.\csromanspacer{} ตํ กิสฺส เหตุ, โย ภิกฺขเว ราโค, โส อตฺถิ, โย โทโส, โส อตฺถิ, โย โมโห, โส อตฺถิ.\csromanspacer{} โย ราโค, โส อปฺปหีโน, โย โทโส, โส อปฺปหีโน, โย โมโห, โส อปฺปหีโน.}

\prose{ยสฺส กสฺสจิ ภิกฺขเว ภิกฺขุสฺส วา ภิกฺขุนิยา วา จกฺขุวิญฺเญยฺเยสุ รูเปสุ โย ราโค, โส นตฺถิ, โย โทโส, โส นตฺถิ, โย โมโห, โส นตฺถิ.\csromanspacer{} โย ราโค, โส ปหีโน, โย โทโส,}

\vspace{\tipitakaparskip}
\csromancenter{สฬายตนสํยุตฺต โส ปหีโน, โย โมโห, โส ปหีโน.\csromanspacer{} ตสฺส อธิมตฺตา เจปิ จกฺขุวิญฺเญยฺยา รูปา จกฺขุสฺส อาปาถํ อาคจฺฉนฺติ, เนวสฺส จิตฺตํ ปริยาทิยนฺติ, โก ปน วาโท ปริตฺตานํ.\csromanspacer{} ตํ กิสฺส เหตุ, โย ภิกฺขเว ราโค, โส นตฺถิ, โย โทโส, โส นตฺถิ, โย โมโห, โส นตฺถิ.\csromanspacer{} โย ราโค, โส ปหีโน, โย โทโส, โส ปหีโน, โย โมโห, โส ปหีโน ฯเปฯ.}
\prose{\csromanfolio{371}ยสฺส กสฺสจิ ภิกฺขเว ภิกฺขุสฺส วา ภิกฺขุนิยา วา ชิวฺหาวิญฺเญยฺเยสุ รเสสุ ฯเปฯ มโนวิญฺเญยฺเยสุ ธมฺเมสุ โย ราโค, โส นตฺถิ, โย โทโส, โส นตฺถิ, โย โมโห, โส นตฺถิ.\csromanspacer{} โย ราโค, โส ปหีโน, โย โทโส, โส ปหีโน, โย โมโห, โส ปหีโน.\csromanspacer{} ตสฺส อธิมตฺตา เจปิ มโนวิญฺเญยฺยา ธมฺมา มนสฺส อาปาถํ อาคจฺฉนฺติ, เนวสฺส จิตฺตํ ปริยาทิยนฺติ, โก ปน วาโท ปริตฺตานํ.\csromanspacer{} ตํ กิสฺส เหตุ, โย ภิกฺขเว ราโค, โส นตฺถิ, โย โทโส, โส นตฺถิ, โย โมโห, โส นตฺถิ.\csromanspacer{} โย ราโค, โส ปหีโน, โย โทโส, โส ปหีโน, โย โมโห, โส ปหีโน.\csromanspacer{} เสยฺยถาปิ ภิกฺขเว ขีรรุกฺโข อสฺสตฺโถ วา นิโคฺรโธ วา ปิลกฺโข วา อุทุมฺพโร วา สุกฺโข โกลาโป เตโรวสฺสิโก, ตเมนํ ปุริโส ติณฺหาย กุฐาริยา ยโต ยโต อาภินฺเทยฺย, อาคจฺเฉยฺย ขีรนฺติ.\csromanspacer{} โน เหตํ ภนฺเต.\csromanspacer{} ตํ กิสฺส เหตุ, ยํ หิ ภนฺเต ขีรํ, ตํ นตฺถีติ.}

\prose{เอวเมว โข ภิกฺขเว ยสฺส กสฺสจิ ภิกฺขุสฺส วา ภิกฺขุนิยา วา จกฺขุวิญฺเญยฺเยสุ รูเปสุ โย ราโค, โส นตฺถิ, โย โทโส, โส นตฺถิ, โย โมโห, โส นตฺถิ.\csromanspacer{} โย ราโค, โส ปหีโน, โย โทโส, โส ปหีโน, โย โมโห, โส ปหีโน.\csromanspacer{} ตสฺส อธิมตฺตา เจปิ จกฺขุวิญฺเญยฺยา รูปา จกฺขุสฺส อาปาถํ อาคจฺฉนฺติ, เนวสฺส จิตฺตํ ปริยาทิยนฺติ, โก ปน วาโท ปริตฺตานํ.\csromanspacer{} ตํ กิสฺส เหตุ, โย ภิกฺขเว ราโค, โส นตฺถิ, โย โทโส, โส นตฺถิ, โย โมโห, โส นตฺถิ.\csromanspacer{} โย ราโค, โส ปหีโน, โย โทโส, โส ปหีโน, โย โมโห, โส ปหีโนฯเปฯ.}

\prose{ยสฺส กสฺสจิ ภิกฺขเว ภิกฺขุสฺส วา ภิกฺขุนิยา วา ชิวฺหาวิญฺเญยฺเยสุ รเสสุ ฯเปฯ.}

\gambhira{สฬายตนวคฺคสํยุตฺตปาฬิ}
\prose{\csromanfolio{372}ยสฺส กสฺสจิ ภิกฺขเว ภิกฺขุสฺส วา ภิกฺขุนิยา วา มโนวิญฺเญยฺเยสุ ธมฺเมสุ โย ราโค, โส นตฺถิ, โย โทโส, โส นตฺถิ, โย โมโห, โส นตฺถิ.\csromanspacer{} โย ราโค, โส ปหีโน, โย โทโส, โส ปหีโน, โย โมโห, โส ปหีโน.\csromanspacer{} ตสฺส อธิมตฺตา เจปิ มโนวิญฺเญยฺยา ธมฺมา มนสฺส อาปาถํ อาคจฺฉนฺติ, เนวสฺส จิตฺตํ ปริยาทิยนฺติ, โก ปน วาโท ปริตฺตานํ.\csromanspacer{} ตํ กิสฺส เหตุ, โย ภิกฺขเว ราโค, โส นตฺถิ, โย โทโส, โส นตฺถิ, โย โมโห, โส นตฺถิ.\csromanspacer{} โย ราโค, โส ปหีโน, โย โทโส, โส ปหีโน, โย โมโห, โส ปหีโนติ.\csromanspacer{}\csromanspacer{}จตุตฺถํ.}

\csromanheader{2}{5. โกฏฺฐิกสุตฺต}
\csromanmatikaanchor{mtk.201}
\csromantocmarkref{subsection}{5. โกฏฺฐิกสุตฺต}{mtk.201}
\bookmark[dest={mtk.201},level=2]{5. โกฏฺฐิกสุตฺต}
\proseitem{232}{เอกํ สมยํ อายสฺมา จ สาริปุตฺโต อายสฺมา จ มหาโกฏฺฐิโก พาราณสิยํ วิหรนฺติ อิสิปตเน มิคทาเย.\csromanspacer{} อถ โข อายสฺมา มหาโกฏฺฐิโก สายนฺหสมยํ ปฏิสลฺลาโน วุฏฺฐิโต เยนายสฺมา สาริปุตฺโต เตนุปสงฺกมิ, อุปสงฺกมิตฺวา อายสฺมตา สาริปุตฺเตน สทฺธิํ สมฺโมทิ, สมฺโมทนียํ กถํ สารณียํ วีติสาเรตฺวา เอกมนฺตํ นิสีทิ, เอกมนฺตํ นิสินฺโน โข อายสฺมา มหาโกฏฺฐิโก อายสฺมนฺตํ สาริปุตฺตํ เอตทโวจ\csromandash{}}

\prose{กิํ นุ โข อาวุโส สาริปุตฺต จกฺขุ รูปานํ สํโยชนํ, รูปา จกฺขุสฺส สํโยชนํ ฯเปฯ ชิวฺหา รสานํ สํโยชนํ, รสา ชิวฺหาย สํโยชนํ ฯเปฯ มโน ธมฺมานํ สํโยชนํ, ธมฺมา มนสฺส สํโยชนนฺติ.}

\prose{น โข อาวุโส โกฏฺฐิก จกฺขุ รูปานํ สํโยชนํ, น รูปา จกฺขุสฺส สํโยชนํ, ยญฺจ ตตฺถ ตทุภยํ ปฏิจฺจ อุปฺปชฺชติ ฉนฺทราโค, ตํ ตตฺถ สํโยชนํ ฯเปฯ.\csromanspacer{} น ชิวฺหา รสานํ สํโยชนํ, น รสา ชิวฺหาย สํโยชนํ, ยญฺจ ตตฺถ ตทุภยํ ปฏิจฺจ อุปฺปชฺชติ ฉนฺทราโค, ตํ ตตฺถ สํโยชนํ ฯเปฯ.\csromanspacer{} น มโน ธมฺมานํ สํโยชนํ, น ธมฺมา มนสฺส สํโยชนํ, ยญฺจ ตตฺถ ตทุภยํ ปฏิจฺจ อุปฺปชฺชติ ฉนฺทราโค, ตํ ตตฺถ สํโยชนํ.}

\prose{เสยฺยถาปิ อาวุโส กาโฬ จ พลีพทฺโท\footnote{พลิวทฺโท (สี, อิ), พลิพทฺโท (สฺยา, กํ, ก)} โอทาโต จ พลีพทฺโท เอเกน ทาเมน วา โยตฺเตน วา สํยุตฺตา อสฺสุ, โย นุ โข}

\vspace{\tipitakaparskip}
\csromancenter{สฬายตนสํยุตฺต เอวํ วเทยฺย “กาโฬ พลีพทฺโท โอทาตสฺส พลีพทฺทสฺส สํโยชนํ, โอทาโต พลีพทฺโท กาฬสฺส พลีพทฺทสฺส สํโยชนนฺ”ติ.\csromanspacer{} สมฺมา นุ โข โส วทมาโน วเทยฺยาติ.\csromanspacer{} โน เหตํ อาวุโส.\csromanspacer{} น โข อาวุโส กาโฬ พลีพทฺโท โอทาตสฺส พลีพทฺทสฺส สํโยชนํ, น โอทาโต พลีพทฺโท กาฬสฺส พลีพทฺทสฺส สํโยชนํ, เยน จ โข เต เอเกน ทาเมน วา โยตฺเตน วา สํยุตฺตา, ตํ ตตฺถ สํโยชนํ.}
\prose{\csromanfolio{373}เอวเมว โข อาวุโส น จกฺขุ รูปานํ สํโยชนํ, น รูปา จกฺขุสฺส สํโยชนํ, ยญฺจ ตตฺถ ตทุภยํ ปฏิจฺจ อุปฺปชฺชติ ฉนฺทราโค, ตํ ตตฺถ สํโยชนํ ฯเปฯ.\csromanspacer{} น ชิวฺหา รสานํ สํโยชนํ ฯเปฯ.\csromanspacer{} น มโน ธมฺมานํ สํโยชนํ, น ธมฺมา มนสฺส สํโยชนํ, ยญฺจ ตตฺถ ตทุภยํ ปฏิจฺจ อุปฺปชฺชติ ฉนฺทราโค, ตํ ตตฺถ สํโยชนํ.}

\prose{จกฺขุ วา อาวุโส รูปานํ สํโยชนํ อภวิสฺส รูปา วา จกฺขุสฺส สํโยชนํ, นยิทํ พฺรหฺมจริยวาโส ปญฺญาเยถ\footnote{ปญฺญายติ (ก)} สมฺมา ทุกฺขกฺขยาย.\csromanspacer{} ยสฺมา จ โข อาวุโส น จกฺขุ รูปานํ สํโยชนํ, น รูปา จกฺขุสฺส สํโยชนํ, ยญฺจ ตตฺถ ตทุภยํ ปฏิจฺจ อุปฺปชฺชติ ฉนฺทราโค, ตํ ตตฺถ สํโยชนํ, ตสฺมา พฺรหฺมจริยวาโส ปญฺญายติ สมฺมา ทุกฺขกฺขยาย ฯเปฯ.}

\prose{ชิวฺหา อาวุโส รสานํ สํโยชนํ อภวิสฺส รสา วา ชิวฺหาย สํโยชนํ, นยิทํ พฺรหฺมจริยวาโส ปญฺญาเยถ สมฺมา ทุกฺขกฺขยาย.\csromanspacer{} ยสฺมา จ โข อาวุโส น ชิวฺหา รสานํ สํโยชนํ, น รสา ชิวฺหาย สํโยชนํ, ยญฺจ ตตฺถ ตทุภยํ ปฏิจฺจ อุปฺปชฺชติ ฉนฺทราโค, ตํ ตตฺถ สํโยชนํ, ตสฺมา พฺรหฺมจริยวาโส ปญฺญายติ สมฺมา ทุกฺขกฺขยาย ฯเปฯ.}

\prose{มโน วา อาวุโส ธมฺมานํ สํโยชนํ อภวิสฺส ธมฺมา วา มนสฺส สํโยชนํ, นยิทํ พฺรหฺมจริยวาโส ปญฺญาเยถ สมฺมา ทุกฺขกฺขยาย.\csromanspacer{} ยสฺมา จ โข อาวุโส น มโน ธมฺมานํ สํโยชนํ, น ธมฺมา มนสฺส สํโยชนํ, ยญฺจ ตตฺถ ตทุภยํ ปฏิจฺจ อุปฺปชฺชติ ฉนฺทราโค, ตํ ตตฺถ สํโยชนํ, ตสฺมา พฺรหฺมจริยวาโส ปญฺญายติ สมฺมา ทุกฺขกฺขยาย.}

\gambhira{สฬายตนวคฺคสํยุตฺตปาฬิ}
\prose{\csromanfolio{374}อิมินาเปตํ อาวุโส ปริยาเยน เวทิตพฺพํ, ยถา น จกฺขุ รูปานํ สํโยชนํ, น รูปา จกฺขุสฺส สํโยชนํ, ยญฺจ ตตฺถ ตทุภยํ ปฏิจฺจ อุปฺปชฺชติ ฉนฺทราโค, ตํ ตตฺถ สํโยชนํ ฯเปฯ.\csromanspacer{} น ชิวฺหา รสานํ สํโยชนํ ฯเปฯ.\csromanspacer{} น มโน ธมฺมานํ สํโยชนํ, น ธมฺมา มนสฺส สํโยชนํ, ยญฺจ ตตฺถ ตทุภยํ ปฏิจฺจ อุปฺปชฺชติ ฉนฺทราโค, ตํ ตตฺถ สํโยชนํ.}

\prose{สํวิชฺชติ โข อาวุโส ภควโต จกฺขุ, ปสฺสติ ภควา จกฺขุนา รูปํ, ฉนฺทราโค ภควโต นตฺถิ, สุวิมุตฺตจิตฺโต ภควา.\csromanspacer{} สํวิชฺชติ โข อาวุโส ภควโต โสตํ, สุณาติ ภควา โสเตน สทฺทํ, ฉนฺทราโค ภควโต นตฺถิ, สุวิมุตฺตจิตฺโต ภควา.\csromanspacer{} สํวิชฺชติ โข อาวุโส ภควโต ฆานํ, ฆายติ ภควา ฆาเนน คนฺธํ, ฉนฺทราโค ภควโต นตฺถิ, สุวิมุตฺตจิตฺโต ภควา.\csromanspacer{} สํวิชฺชติ โข อาวุโส ภควโต ชิวฺหา, สายติ ภควา ชิวฺหาย รสํ, ฉนฺทราโค ภควโต นตฺถิ, สุวิมุตฺตจิตฺโต ภควา.\csromanspacer{} สํวิชฺชติ โข อาวุโส ภควโต กาโย, ผุสติ ภควา กาเยน โผฏฺฐพฺพํ, ฉนฺทราโค ภควโต นตฺถิ, สุวิมุตฺตจิตฺโต ภควา.\csromanspacer{} สํวิชฺชติ โข อาวุโส ภควโต มโน, วิชานาติ ภควา มนสา ธมฺมํ, ฉนฺทราโค ภควโต นตฺถิ, สุวิมุตฺตจิตฺโต ภควา.}

\prose{อิมินา โข เอตํ อาวุโส ปริยาเยน เวทิตพฺพํ, ยถา น จกฺขุ รูปานํ สํโยชนํ, น รูปา จกฺขุสฺส สํโยชนํ, ยญฺจ ตตฺถ ตทุภยํ ปฏิจฺจ อุปฺปชฺชติ ฉนฺทราโค, ตํ ตตฺถ สํโยชนํ.\csromanspacer{} น โสตํ.\csromanspacer{} น ฆานํ.\csromanspacer{} น ชิวฺหา รสานํ สํโยชนํ, น รสา ชิวฺหาย สํโยชนํ, ยญฺจ ตตฺถ ตทุภยํ ปฏิจฺจ อุปฺปชฺชติ ฉนฺทราโค, ตํ ตตฺถ สํโยชนํ.\csromanspacer{} น กาโย.\csromanspacer{} น มโน ธมฺมานํ สํโยชนํ, น ธมฺมา มนสฺส สํโยชนํ, ยญฺจ ตตฺถ ตทุภยํ ปฏิจฺจ อุปฺปชฺชติ ฉนฺทราโค, ตํ ตตฺถ สํโยชนนฺติ.\csromanspacer{}\csromanspacer{}ปญฺจมํ.}

\csromanheader{2}{6. กามภูสุตฺต}
\csromanmatikaanchor{mtk.202}
\csromantocmarkref{subsection}{6. กามภูสุตฺต}{mtk.202}
\bookmark[dest={mtk.202},level=2]{6. กามภูสุตฺต}
\proseitem{233}{เอกํ สมยํ อายสฺมา จ อานนฺโท อายสฺมา จ กามภู โกสมฺพิยํ วิหรนฺติ โฆสิตาราเม.\csromanspacer{} อถ โข อายสฺมา กามภู สายนฺหสมยํ ปฏิสลฺลานา วุฏฺฐิโต เยนายสฺมา อานนฺโท เตนุปสงฺกมิ, อุปสงฺกมิตฺวา อายสฺมตา อานนฺเทน สทฺธิํ สมฺโมทิ,}

\vspace{\tipitakaparskip}
\csromancenter{สฬายตนสํยุตฺต สมฺโมทนียํ กถํ สารณียํ วีติสาเรตฺวา เอกมนฺตํ นิสีทิ, เอกมนฺตํ นิสินฺโน โข อายสฺมา กามภู อายสฺมนฺตํ อานนฺทํ เอตทโวจ\csromandash{}}
\prose{\csromanfolio{375}กิํ นุ โข อาวุโส อานนฺท จกฺขุ รูปานํ สํโยชนํ, รูปา จกฺขุสฺส สํโยชนํ ฯเปฯ.\csromanspacer{} ชิวฺหา รสานํ สํโยชนํ, รสา ชิวฺหาย สํโยชนํ ฯเปฯ.\csromanspacer{} มโน ธมฺมานํ สํโยชนํ, ธมฺมา มนสฺส สํโยชนนฺติ.}

\prose{น โข อาวุโส กามภู\footnote{กามภุ (สี) โมคฺคลฺลาเน 65 เค วาติ สุตฺตํ ปสฺสิตพฺพํ.} จกฺขุ รูปานํ สํโยชนํ, น รูปา จกฺขุสฺส สํโยชนํ, ยญฺจ ตตฺถ ตทุภยํ ปฏิจฺจ อุปฺปชฺชติ ฉนฺทราโค, ตํ ตตฺถ สํโยชนํ ฯเปฯ.\csromanspacer{} น ชิวฺหา รสานํ สํโยชนํ, น รสา ชิวฺหาย สํโยชนํ ฯเปฯ.\csromanspacer{} น มโน ธมฺมานํ สํโยชนํ, น ธมฺมา มนสฺส สํโยชนํ, ยญฺจ ตตฺถ ตทุภยํ ปฏิจฺจ อุปฺปชฺชติ ฉนฺทราโค, ตํ ตตฺถ สํโยชนํ.}

\prose{เสยฺยถาปิ อาวุโส กาโฬ จ พลีพทฺโท โอทาโต จ พลีพทฺโท เอเกน ทาเมน วา โยตฺเตน วา สํยุตฺตา อสฺสุ, โย นุ โข เอวํ วเทยฺย “กาโฬ พลีพทฺโท โอทาตสฺส พลีพทฺทสฺส สํโยชนํ, โอทาโต พลีพทฺโท กาฬสฺส พลีพทฺทสฺส สํโยชนนฺ”ติ.\csromanspacer{} สมฺมา นุ โข โส วทมาโน วเทยฺยาติ.\csromanspacer{} โน เหตํ อาวุโส.\csromanspacer{} น โข อาวุโส กาโฬ พลีพทฺโท โอทาตสฺส พลีพทฺทสฺส สํโยชนํ, นปิ โอทาโต พลีพทฺโท กาฬสฺส พลีพทฺทสฺส สํโยชนํ, เยน จ โข เต เอเกน ทาเมน วา โยตฺเตน วา สํยุตฺตา, ตํ ตตฺถ สํโยชนํ.\csromanspacer{} เอวเมว โข อาวุโส น จกฺขุ รูปานํ สํโยชนํ, น รูปา จกฺขุสฺส สํโยชนํ ฯเปฯ.\csromanspacer{} น ชิวฺหา ฯเปฯ.\csromanspacer{} น มโน ฯเปฯ.\csromanspacer{} ยญฺจ ตตฺถ ตทุภยํ ปฏิจฺจ อุปฺปชฺชติ ฉนฺทราโค, ตํ ตตฺถ สํโยชนนฺติ.\csromanspacer{}\csromanspacer{}ฉฏฺฐํ.}

\csromanheader{2}{7. อุทายีสุตฺต}
\csromanmatikaanchor{mtk.203}
\csromantocmarkref{subsection}{7. อุทายีสุตฺต}{mtk.203}
\bookmark[dest={mtk.203},level=2]{7. อุทายีสุตฺต}
\proseitem{234}{เอกํ สมยํ อายสฺมา จ อานนฺโท อายสฺมา จ อุทายี โกสมฺพิยํ วิหรนฺติ โฆสิตาราเม.\csromanspacer{} อถ โข อายสฺมา อุทายี สายนฺหสมยํ ปฏิสลฺลานา วุฏฺฐิโต เยนายสฺมา อานนฺโท เตนุปสงฺกมิ, อุปสงฺกมิตฺวา อายสฺมตา อานนฺเทน สทฺธิํ สมฺโมทิ, สมฺโมทนียํ กถํ สารณียํ วีติสาเรตฺวา เอกมนฺตํ นิสีทิ, เอกมนฺตํ นิสินฺโน โข อายสฺมา อุทายี อายสฺมนฺตํ อานนฺทํ เอตทโวจ\csromandash{}}

\gambhira{สฬายตนวคฺคสํยุตฺตปาฬิ}
\prose{\csromanfolio{376}ยเถว นุ โข อาวุโส อานนฺท อยํ กาเย ภควตา อเนกปริยาเยน อกฺขาโต วิวโฏ ปกาสิโต “อิติปายํ กาโย อนตฺตา”ติ.\csromanspacer{} สกฺกา เอวเมว วิญฺญาณมฺปิทํ อาจิกฺขิตุํ เทเสตุํ ปญฺญเปตุํ ปฏฺฐเปตุํ วิวริตุํ วิภชิตุํ อุตฺตานีกาตุํ “อิติปิทํ วิญฺญาณํ อนตฺตา”ติ.}

\prose{ยเถว โข อาวุโส อุทายี อยํ กาโย ภควตา อเนกปริยาเยน อกฺขาโต วิวโฏ ปกาสิโต “อิติปายํ กาโย อนตฺตา”ติ.\csromanspacer{} สกฺกา เอวเมว วิญฺญาณมฺปิทํ อาจิกฺขิตุํ เทเสตุํ ปญฺญเปตุํ ปฏฺฐเปตุํ วิวริตุํ วิภชิตุํ อุตฺตานีกาตุํ “อิติปิทํ วิญฺญาณํ อนตฺตา”ติ.}

\prose{จกฺขุญฺจ อาวุโส ปฏิจฺจ รูเป จ อุปฺปชฺชติ จกฺขุวิญฺญาณนฺติ.\csromanspacer{} เอวมาวุโสติ.\csromanspacer{} โย จาวุโส เหตุ โย จ ปจฺจโย จกฺขุวิญฺญาณสฺส อุปฺปาทาย, โส จ เหตุ โส จ ปจฺจโย สพฺเพน สพฺพํ สพฺพถา สพฺพํ อปริเสสํ นิรุชฺเฌยฺย, อปิ นุ โข จกฺขุวิญฺญาณํ ปญฺญาเยถาติ.\csromanspacer{} โน เหตํ อาวุโส.\csromanspacer{} อิมินาปิ โข เอตํ อาวุโส ปริยาเยน ภควตา อกฺขาตํ วิวฏํ ปกาสิตํ “อิติปิทํ วิญฺญาณํ อนตฺตา”ติ ฯเปฯ.}

\prose{ชิวฺหญฺจาวุโส ปฏิจฺจ รเส จ อุปฺปชฺชติ ชิวฺหาวิญฺญาณนฺติ.\csromanspacer{} เอวมาวุโสติ.\csromanspacer{} โย จาวุโส เหตุ โย จ ปจฺจโย ชิวฺหาวิญฺญาณสฺส อุปฺปาทาย, โส จ เหตุ โส จ ปจฺจโย สพฺเพน สพฺพํ สพฺพถา สพฺพํ อปริเสสํ นิรุชฺเฌยฺย, อปิ นุ โข ชิวฺหาวิญฺญาณํ ปญฺญาเยถาติ.\csromanspacer{} โน เหตํ อาวุโส.\csromanspacer{} อิมินาปิ โข เอตํ อาวุโส ปริยาเยน ภควตา อกฺขาตํ วิวฏํ ปกาสิตํ “อิติปิทํ วิญฺญาณํ อนตฺตา”ติ ฯเปฯ.}

\prose{มนญฺจาวุโส ปฏิจฺจ ธมฺเม จ อุปฺปชฺชติ มโนวิญฺญาณนฺติ.\csromanspacer{} เอวมาวุโสติ.\csromanspacer{} โย จาวุโส เหตุ โย จ ปจฺจโย มโนวิญฺญาณสฺส อุปฺปาทาย, โส จ เหตุ โส จ ปจฺจโย สพฺเพน สพฺพํ สพฺพถา สพฺพํ อปริเสสํ นิรุชฺเฌยฺย, อปิ นุ โข มโนวิญฺญาณํ ปญฺญาเยถาติ.\csromanspacer{} โน เหตํ อาวุโส.\csromanspacer{} อิมินาปิ โข เอตํ อาวุโส ปริยาเยน ภควตา อกฺขาตํ วิวฏํ ปกาสิตํ “อิติปิทํ วิญฺญาณํ อนตฺตา”ติ.}

\prose{เสยฺยถาปิ อาวุโส ปุริโส สารตฺถิโก สารคเวสี สารปริเยสนํ จรมาโน ติณฺหํ กุฐาริํ อาทาย วนํ ปวิเสยฺย, โส}

\vspace{\tipitakaparskip}
\csromancenter{สฬายตนสํยุตฺต ตตฺถ ปสฺเสยฺย มหนฺตํ กทลิกฺขนฺธํ อุชุํ นวํ อกุกฺกุกชาตํ\footnote{อกุกฺกุฏกชาตํ (สฺยา, กํ), อกุกฺกชฏชาตํ (ก)}, ตเมนํ มูเล ฉินฺเทยฺย, มูเล เฉตฺวา อคฺเค ฉินฺเทยฺย, อคฺเค เฉตฺวา ปตฺตวฏฺฏิํ วินิพฺภุเชยฺย\footnote{วินิพฺภุชฺเชยฺย (อิ), วินิพฺภชฺเชยฺย (สฺยา, กํ)}, โส ตตฺถ เผคฺคุมฺปิ นาธิคจฺเฉยฺย, กุโต สารํ.\csromanspacer{} เอวเมว โข อาวุโส ภิกฺขุ ฉสุ ผสฺสายตเนสุ เนวตฺตานํ น อตฺตนิยํ สมนุปสฺสติ, โส เอวํ อสมนุปสฺสนฺโต\footnote{เอวํ สมนุปสฺสนฺโต (สฺยา, กํ, ก)} น กิญฺจิ โลเก อุปาทิยติ, อนุปาทิยํ น ปริตสฺสติ, อปริตสฺสํ ปจฺจตฺตญฺเญว ปรินิพฺพายติ, “ขีณา ชาติ, วุสิตํ พฺรหฺมจริยํ, กตํ กรณียํ, นาปรํ อิตฺถตฺตายา”ติ ปชานาตีติ.\csromanspacer{}\csromanspacer{}สตฺตมํ.}
\csromanheader{2}{8. อาทิตฺตปริยายสุตฺต}
\csromanmatikaanchor{mtk.204}
\csromantocmarkref{subsection}{8. อาทิตฺตปริยายสุตฺต}{mtk.204}
\bookmark[dest={mtk.204},level=2]{8. อาทิตฺตปริยายสุตฺต}
\proseitem{235}{\csromanfolio{377}อาทิตฺตปริยายํ โว ภิกฺขเว ธมฺมปริยายํ เทเสสฺสามิ, ตํ สุณาถ.\csromanspacer{} กตโม จ ภิกฺขเว อาทิตฺตปริยาโย ธมฺมปริยาโย.\csromanspacer{} วรํ ภิกฺขเว ตตฺตา ย อโยสลากาย อาทิตฺตาย สมฺปชฺชลิตาย สโชติภูตาย\footnote{สญฺโชติภูตาย (สฺยา, กํ)} จกฺขุนฺทฺริยํ สมฺปลิมฏฺฐํ, น เตฺวว จกฺขุวิญฺเญยฺเยสุ รูเปสุ อนุพฺยญฺชนโส นิมิตฺตคฺคาโห.\csromanspacer{} นิมิตฺตสฺสาทคถิตํ\footnote{นิมิตฺตสฺสาทคธิตํ (สฺยา, กํ, ก) ม 3. 268 ปิฏฺเฐปิ.} วา ภิกฺขเว วิญฺญาณํ ติฏฺฐมานํ ติฏฺเฐยฺย อนุพฺยญฺชนสฺสาทคถิตํ วา, ตสฺมิํ เจ สมเย กาลํ กเรยฺย, ฐานเมตํ วิชฺชติ, ยํ ทฺวินฺนํ คตีนํ อญฺญตรํ คติํ คจฺเฉยฺย นิรยํ วา ติรจฺฉานโยนิํ วา, อิมํ ขฺวาหํ ภิกฺขเว อาทีนวํ ทิสฺวา เอวํ วทามิ.}

\prose{วรํ ภิกฺขเว ติเณฺหน อโยสงฺกุนา อาทิตฺเตน สมฺปชฺชลิเตน สโชติภูเตน โสตินฺทฺริยํ สมฺปลิมฏฺฐํ, น เตฺวว โสตวิญฺเญยฺเยสุ สทฺเทสุ อนุพฺยญฺชนโส นิมิตฺตคฺคาโห.\csromanspacer{} นิมิตฺตสฺสาทคถิตํ วา ภิกฺขเว วิญฺญาณํ ติฏฺฐมานํ ติฏฺเฐยฺย อนุพฺยญฺชนสฺสาทคถิตํ วา, ตสฺมิํ เจ สมเย กาลํ กเรยฺย, ฐานเมตํ วิชฺชติ, ยํ ทฺวินฺนํ คตีนํ อญฺญตรํ คติํ คจฺเฉยฺย นิรยํ วา ติรจฺฉานโยนิํ วา, อิมํ ขฺวาหํ ภิกฺขเว อาทีนวํ ทิสฺวา เอวํ วทามิ.}

\gambhira{สฬายตนวคฺคสํยุตฺตปาฬิ}
\prose{\csromanfolio{378}วรํ ภิกฺขเว ติเณฺหน นขจฺเฉทเนน อาทิตฺเตน สมฺปชฺชลิเตน สโชติภูเตน ฆานินฺทฺริยํ สมฺปลิมฏฺฐํ, น เตฺวว ฆานวิญฺเญยฺเยสุ คนฺเธสุ อนุพฺยญฺชนโส นิมิตฺตคฺคาโห.\csromanspacer{} นิมิตฺตสฺสาทคถิตํ วา ภิกฺขเว วิญฺญาณํ ติฏฺฐมานํ ติฏฺเฐยฺย อนุพฺยญฺชนสฺสาทคถิตํ วา, ตสฺมิํ เจ สมเย กาลํ กเรยฺย, ฐานเมตํ วิชฺชติ, ยํ ทฺวินฺนํ คตีนํ อญฺญตรํ คติํ คจฺเฉยฺย นิรยํ วา ติรจฺฉานโยนิํ วา, อิมํ ขฺวาหํ ภิกฺขเว อาทีนวํ ทิสฺวา เอวํ วทามิ.}

\prose{วรํ ภิกฺขเว ติเณฺหน ขุเรน อาทิตฺเตน สมฺปชฺชลิเตน สโชติภูเตน ชิวฺหินฺทฺริยํ สมฺปลิมฏฺฐํ, น เตฺวว ชิวฺหาวิญฺเญยฺเยสุ รเสสุ อนุพฺยญฺชนโส นิมิตฺตคฺคาโห.\csromanspacer{} นิมิตฺตสฺสาทคถิตํ วา ภิกฺขเว วิญฺญาณํ ติฏฺฐมานํ ติฏฺเฐยฺย อนุพฺยญฺชนสฺสาทคถิตํ วา, ตสฺมิํ เจ สมเย กาลํ กเรยฺย, ฐานเมตํ วิชฺชติ, ยํ ทฺวินฺนํ คตีนํ อญฺญตรํ คติํ คจฺเฉยฺย นิรยํ วา ติรจฺฉานโยนิํ วา, อิมํ ขฺวาหํ ภิกฺขเว อาทีนวํ ทิสฺวา เอวํ วทามิ.}

\prose{วรํ ภิกฺขเว ติณฺหาย สตฺติยา อาทิตฺตาย สมฺปชฺชลิตาย สโชติภูตาย กายินฺทฺริยํ สมฺปลิมฏฺฐํ, น เตฺวว กายวิญฺเญยฺเยสุ โผฏฺฐพฺเพสุ อนุพฺยญฺชนโส นิมิตฺตคฺคาโห.\csromanspacer{} นิมิตฺตสฺสาทคถิตํ วา ภิกฺขเว วิญฺญาณํ วิญฺญาณํ ติฏฺฐมานํ ติฏฺเฐยฺย อนุพฺยญฺชนสฺสาทคถิตํ วา, ตสฺมิํ เจ สมเย กาลํ กเรยฺย, ฐานเมตํ วิชฺชติ, ยํ ทฺวินฺนํ คตีนํ อญฺญตรํ คติํ คจฺเฉยฺย นิรยํ วา ติรจฺจานโยนิํ วา, อิมํ ขฺวาหํ ภิกฺขเว อาทีนวํ ทิสฺวา เอวํ วทามิ.}

\prose{วรํ ภิกฺขเว โสตฺตํ.\csromanspacer{} โสตฺตํ โข ปนาหํ ภิกฺขเว วญฺฌํ ชีวิตานํ วทามิ, อผลํ ชีวิตานํ วทามิ, โมมูหํ ชีวิตานํ วทามิ, น เตฺวว ตถารูเป วิตกฺเก วิตกฺเกยฺย, ยถารูปานํ วิตกฺกานํ วสํ คโต สํฆํ ภินฺเทยฺย, อิมํ ขฺวาหํ ภิกฺขเว วญฺฌํ ชีวิตานํ อาทีนวํ ทิสฺวา เอวํ วทามิ.}

\prose{ตตฺถ ภิกฺขเว สุตวา อริยสาวโก อิติ ปฏิสญฺจิกฺขติ\textbf{\csromandash{}}ติฏฺฐตุ ตาว ตตฺตาย อโยสลากาย อาทิตฺตาย สมฺปชฺชลิตาย สโชติภูตาย จกฺขุนฺทฺริยํ สมฺปลิมฏฺฐํ, หนฺทาหํ อิทเมว มนสิ กโรมิ “อิติ จกฺขุ อนิจฺจํ, รูปา อนิจฺจา, จกฺขุวิญฺญาณํ อนิจฺจํ, จกฺขุสมฺผสฺโส อนิจฺโจ.\csromanspacer{} ยมฺปิทํ จกฺขุสมฺผสฺสปจฺจยา อุปฺปชฺชติ เวทยิตํ สุขํ วา ทุกฺขํ วา อทุกฺขมสุขํ วา, ตมฺปิ อนิจฺจํ”\footnote{อนิจฺจนฺ”ติ (?)}.}

\csromanheader{2}{1. สฬายตนสํยุตฺต}
\prose{\csromanfolio{379}ติฏฺฐตุ ตาว ติเณฺหน อโยสงฺกุนา อาทิตฺเตน สมฺปชฺชลิเตน สโชติภูเตน โสตินฺทฺริยํ สมฺปลิมฏฺฐํ, หนฺทาหํ อิทเมว มนสิ กโรมิ “อิติ โสตํ อนิจฺจํ, สทฺทา อนิจฺจา, โสตวิญฺญาณํ อนิจฺจํ, โสตสมฺผสฺโส อนิจฺโจ.\csromanspacer{} ยมฺปิทํ โสตสมฺผสฺสปจฺจยา อุปฺปชฺชติ เวทยิตํ สุขํ วา ทุกฺขํ วา อทุกฺขมสุขํ วา, ตมฺปิ อนิจฺจํ”.}

\prose{ติฏฺฐตุ ตาว ติเณฺหน นขจฺเฉทเนน อาทิตฺเตน สมฺปชฺชลิเตน สโชติภูเตน ฆานินฺทฺริยํ สมฺปิลิมฏฺฐํ, หนฺทาหํ อิทเมว มนสิ กโรมิ “อิติ ฆานํ อนิจฺจํ, คนฺธา อนิจฺจา, ฆานวิญฺญาณํ อนิจฺจํ, ฆานสมฺผสฺโส อนิจฺโจ.\csromanspacer{} ยมฺปิทํ ฆานสมฺผสฺสปจฺจยา อุปฺปชฺชติ เวทยิตํ ฯเปฯ.\csromanspacer{} ตมฺปิ อนิจฺจํ”.}

\prose{ติฏฺฐตุ ตาว ติเณฺหน ขุเรน อาทิตฺเตน สมฺปชฺชลิเตน สโชติภูเตน ชิวฺหินฺทฺริยํ สมฺปลิมฏฺฐํ, หนฺทาหํ อิทเมว มนสิ กโรมิ “อิติ ชิวฺหา อนิจฺจา, รสา อนิจฺจา, ชิวฺหาวิญฺญาณํ อนิจฺจํ, ชิวฺหาสมฺผสฺโส อนิจฺโจ.\csromanspacer{} ยมฺปิทํ ชิวฺหาสมฺผสฺสปจฺจยา อุปฺปชฺชติ ฯเปฯ.\csromanspacer{} ตมฺปิ อนิจฺจํ”.}

\prose{ติฏฺฐตุ ตาว ติณฺหาย สตฺติยา อาทิตฺตาย สมฺปชฺชลิตาย สโชติภูตาย กายินฺทฺริยํ สมฺปลิมฏฺฐํ, หนฺทาหํ อิทเมว มนสิ กโรมิ “อิติ กาโย อนิจฺโจ, โผฏฺฐาพฺพา อนิจฺจา, กายวิญฺญาณํ อนิจฺจํ, กายสมฺผสฺโส อนิจฺโจ.\csromanspacer{} ยมฺปิทํ กายสมฺผสฺสปจฺจยา อุปฺปชฺชติ เวทยิตํ ฯเปฯ.\csromanspacer{} ตมฺปิ อนิจฺจํ”.}

\prose{ติฏฺฐตุ ตาว โสตฺตํ.\csromanspacer{} หนฺทาหํ อิทเมว มนสิ กโรมิ “อิติ มโน อนิจฺโจ, ธมฺมา อนิจฺจา, มโนวิญฺญาณํ อนิจฺจํ, มโนสมฺผสฺโส อนิจฺโจ.\csromanspacer{} ยมฺปิทํ มโนสมฺผสฺสปจฺจยา อุปฺปชฺชติ เวทยิตํ สุขํ วา ทุกฺขํ วา อทุกฺขมสุขํ วา, ตมฺปิ อนิจฺจํ”.}

\prose{เอวํ ปสฺสํ ภิกฺขเว สุตวา อริยสาวโก จกฺขุสฺมิมฺปิ นิพฺพินฺทติ, รูเปสุปิ นิพฺพินฺทติ, จกฺขุวิญฺญาเณปิ นิพฺพินฺทติ, จกฺขุสมฺผสฺเสปิ นิพฺพินฺทติ ฯเปฯ.\csromanspacer{} ยมฺปิทํ มโนสมฺผสฺสปจฺจยา อุปฺปชฺชติ เวทยิตํ สุขํ วา ทุกฺขํ วา อทุกฺขมสุขํ วา, ตสฺมิมฺปิ นิพฺพินฺทติ, นิพฺพินฺทํ วิรชฺชติ, วิราคา วิมุจฺจติ, วิมุตฺตสฺมิํ “วิมุตฺตมฺ”อิติ ญาณํ โหติ, “ขีณา ชาติ, วุสิตํ พฺรหฺมจริยํ, กตํ กรณียํ, นาปรํ อิตฺถตฺตายา”ติ ปชานาติ.\csromanspacer{} อยํ โข ภิกฺขเว อาทิตฺตปริยาโย ธมฺมปริยาโยติ.\csromanspacer{}\csromanspacer{}อฏฺฐมํ.}

\gambhira{สฬายตนวคฺคสํยุตฺตปาฬิ}
\csromanheader{2}{9. ปฐมหตฺถปาโทปมสุตฺต}
\csromanmatikaanchor{mtk.205}
\csromantocmarkref{subsection}{9-10. หตฺถปาโทปมสุตฺต}{mtk.205}
\bookmark[dest={mtk.205},level=2]{9-10. หตฺถปาโทปมสุตฺต}
\proseitem{236}{\csromanfolio{380}หตฺเถสุ ภิกฺขเว สติ อาทานนิกฺเขปนํ ปญฺญายติ, ปาเทสุ สติ อภิกฺกมปฏิกฺกโม ปญฺญายติ, ปพฺเพสุ สติ สมิญฺชนปสารณํ ปญฺญายติ, กุจฺฉิสฺมิํ สติ ชิฆจฺฉา ปิปาสา ปญฺญายติ.\csromanspacer{} เอวเมว โข ภิกฺขเว จกฺขุสฺมิํ สติ จกฺขุสมฺผสฺสปจฺจยา อุปฺปชฺชติ อชฺฌตฺตํ สุขํ ทุกฺขํ ฯเปฯ.\csromanspacer{} ชิวฺหาย สติ ชิวฺหาสมฺผสฺสปจฺจยา อุปฺปชฺชติ อชฺฌตฺตํ สุขํ ทุกฺขํ ฯเปฯ.\csromanspacer{} มนสฺมิํ สติ มโนสมฺผสฺสปจฺจยา อุปฺปชฺชติ อชฺฌตฺตํ สุขํ ทุกฺขํ ฯเปฯ.}

\prose{หตฺเถสุ ภิกฺขเว อสติ อาทานนิกฺเขปนํ น ปญฺญายติ, ปาเทสุ อสติ อภิกฺขมปฏิกฺกโม น ปญฺญายติ, ปพฺเพสุ อสติ สมิญฺชนปสารณํ น ปญฺญายติ, กุจฺฉิสฺมิํ อสติ ชิฆจฺฉา ปิปาสา น ปญฺญายติ.\csromanspacer{} เอวเมว โข ภิกฺขเว จกฺขุสฺมิํ อสติ จกฺขุสมฺผสฺสปจฺจยา นุปฺปชฺชติ อชฺฌตฺตํ สุขํ ทุกฺขํ ฯเปฯ.\csromanspacer{} ชิวฺหาย อสติ ชิวฺหาสมฺผสฺสปจฺจยา นุปฺปชฺชติ ฯเปฯ.\csromanspacer{} มนสฺมิํ อสติ มโนสมฺผสฺสปจฺจยา นุปฺปชฺชติ อชฺฌตฺตํ สุขํ ทุกฺขนฺติ.\csromanspacer{}\csromanspacer{}นวมํ.}

\csromanheader{2}{10. ทุติยหตฺถปาโทปมสุตฺต}
\proseitem{237}{หตฺเถสุ ภิกฺขเว สติ อาทานนิกฺเขปนํ โหติ, ปาเทสุ สติ อภิกฺกมปฏิกฺกโม โหติ, ปพฺเพสุ สติ สมิญฺชนปสารณํ โหติ, กุจฺฉิสฺมิํ สติ ชิฆจฺฉา ปิปาสา โหติ.\csromanspacer{} เอวเมว โข ภิกฺขเว จกฺขุสฺมิํ สติ จกฺขุสมฺผสฺสปจฺจยา อุปฺปชฺชติ อชฺฌตฺตํ สุขํ ทุกฺขํ ฯเปฯ.\csromanspacer{} ชิวฺหาย สติ ฯเปฯ.\csromanspacer{} มนสฺมิํ สติ มโนสมฺผสฺสปจฺจยา อุปฺปชฺชติ อชฺฌตฺตํ สุขํ ทุกฺขํ ฯเปฯ.}

\prose{หตฺเถสุ ภิกฺขเว อสติ อาทานนิกฺเขปนํ น โหติ, ปาเทสุ อสติ อภิกฺกมปฏิกฺกโม น โหติ, ปพฺเพสุ อสติ สมิญฺชนปสารณํ น โหติ, กุจฺฉิสฺมิํ อสติ ชิฆจฺฉา ปิปาสา น โหติ.\csromanspacer{} เอวเมว โข ภิกฺขเว จกฺขุสฺมิํ อสติ จกฺขุ สมฺผสฺสปจฺจยา นุปฺปชฺชติ อชฺฌตฺตํ สุขํ ทุกฺขํ ฯเปฯ.\csromanspacer{} ชิวฺหาย อสติ ชิวฺหาสมฺผสฺสปจฺจยา นุปฺปชฺชติ ฯเปฯ.\csromanspacer{} มนสฺมิํ อสติ มโนสมฺผสฺสปจฺจยา นุปฺปชฺชติ อชฺฌตฺตํ สุขํ ทุกฺขนฺติ.\csromanspacer{}\csromanspacer{}ทสมํ.\csromanspacer{} สมุทฺทวคฺโค.}

\titlehead{1. สฬายตนสํยุตฺต \textbf{ตสฺสุทฺทานํ}}
\csromangathameasure{เทฺว สมุทฺทา พาฬิสิโก, ขีรรุกฺเขน โกฏฺฐิโก. \\ กามภู อุทายี เจว, อาทิตฺเตน จ อฏฺฐมํ. \\ หตฺถปาทูปมา เทฺวติ, วคฺโค เตน ปวุจฺจตีติ.}
\csromangathagroup{\csromangathastanza{\csromanfolio{381}เทฺว สมุทฺทา พาฬิสิโก, ขีรรุกฺเขน โกฏฺฐิโก. \\ กามภู อุทายี เจว, อาทิตฺเตน จ อฏฺฐมํ. \\ หตฺถปาทูปมา เทฺวติ, วคฺโค เตน ปวุจฺจตีติ.}}

\csromanheader{1}{\textbf{(19) 4. อาสีวิสวคฺค}}
\csromanmatikaanchor{mtk.206}
\csromantocmarknopage{section}{(19) 4. อาสีวิสวคฺค}{mtk.206}
\bookmark[dest={mtk.206},level=1]{(19) 4. อาสีวิสวคฺค}
\csromanheader{2}{1. อาสีวิโสปมสุตฺต}
\csromanmatikaanchor{mtk.207}
\csromantocmarkref{subsection}{1. อาสีวิโสปมสุตฺต}{mtk.207}
\bookmark[dest={mtk.207},level=2]{1. อาสีวิโสปมสุตฺต}
\proseitem{238}{เสยฺยถาปิ ภิกฺขเว จตฺตาโร อาสีวิสา อุคฺคเตชา โฆรวิสา, อถ ปุริโส อาคจฺเฉยฺย ชีวิตุกาโม อมริตุกาโม สุขกาโม ทุกฺขปฺปฏิกูโล.\csromanspacer{} ตเมนํ เอวํ วเทยฺยุํ “อิเม เต อมฺโภ ปุริส จตฺตาโร อาสีวิสา อุคฺคเตชา โฆรวิสา, กาเลน กาลํ วุฏฺฐาเปตพฺพา, กาเลน กาลํ นฺหาเปตพฺพา, กาเลน กาลํ โภเชตพฺพา, กาเลน กาลํ สํเวเสตพฺพา\footnote{ปเวเสตพฺพา (สฺยา, กํ, อิ, ก)}.\csromanspacer{} ยทา จ โข เต อมฺโภ ปุริส อิเมสํ จตุนฺนํ อาสีวิสานํ อุคฺคเตชานํ โฆรวิสานํ อญฺญตโร วา อญฺญตโร วา กุปฺปิสฺสติ, ตโต ตฺวํ อมฺโภ ปุริส มรณํ วา นิคจฺฉสิ, มรณมตฺตํ วา ทุกฺขํ, ยํ เต อมฺโภ ปุริส กรณียํ, ตํ กโรหี”ติ.}

\prose{อถ โข โส ภิกฺขเว ปุริโส ภีโต จตุนฺนํ อาสีวิสานํ อุคฺคเตชานํ โฆรวิสานํ เยน วา เตน วา ปลาเยถ.\csromanspacer{} ตเมนํ เอวํ วเทยฺยุํ “อิเม โข อมฺโภ ปุริส ปญฺจ วธกา ปจฺจตฺถิกา ปิฏฺฐิโต ปิฏฺฐิโต อนุพนฺธา ‘ยตฺเถว นํ ปสฺสิสฺสาม, ตตฺเถว ชีวิตา โวโรเปสฺสามา’ติ, ยํ เต อมฺโภ ปุริส กรณียํ, ตํ กโรหี”ติ.}

\prose{อถ โข โส ภิกฺขเว ปุริโส ภีโต จตุนฺนํ อาสีวิสานํ อุคฺคเตชานํ โฆรวิสานํ ภีโต ปญฺจนฺนํ วธกานํ ปจฺจตฺถิกานํ เยน วา เตน วา ปลาเยถ.\csromanspacer{} ตเมนํ เอวํ วเทยฺยุํ “อยํ เต อมฺโภ ปุริส ฉฏฺโฐ อนฺตรจโร วธโก อุกฺขิตฺตาสิโก ปิฏฺฐิโต ปิฏฺฐิโต อนุพนฺโธ ‘ยตฺเตว นํ ปสฺสิสฺสามิ, ตตฺเถว สิโร ปาเตสฺสามี’ติ, ยํ เต อมฺโภ ปุริส กรณียํ, ตํ กโรหี”ติ.}

\gambhira{สฬายตนวคฺคสํยุตฺตปาฬิ}
\prose{\csromanfolio{382}อถ โข โส ภิกฺขเว ปุริโส ภีโต จตุนฺนํ อาสีวิสานํ อุคฺคเตชานํ โฆรวิสานํ ภีโต ปญฺจนฺนํ วธกานํ ปจฺจตฺถิกานํ ภีโต ฉฏฺฐสฺส อนฺตรจรสฺส วธกสฺส อุกฺขิตฺตาสิกสฺส เยน วา เตน วา ปลาเยถ.\csromanspacer{} โส ปสฺเสยฺย สุญฺญํ คามํ, ยญฺญเทว ฆรํ ปวิเสยฺย, ริตฺตกญฺเญว ปวิเสยฺย, ตุจฺฉกญฺเญว ปวิเสยฺย, สุญฺญกญฺเญว ปวิเสยฺย.\csromanspacer{} ยญฺญเทว ภาชนํ ปริมเสยฺย, ริตฺตกญฺเญว ปริมเสยฺย, ตุจฺฉกญฺเญว ปริมเสยฺย, สุญฺญกญฺเญว ปริมเสยฺย.\csromanspacer{} ตเมนํ เอวํ วเทยฺยุํ “อิทานิ อมฺโภ ปุริส อิมํ สุญฺญํ คามํ โจรา คามฆาตกา ปวิสนฺติ\footnote{วธิสฺสนฺติ (สี, อิ)}, ยํ เต อมฺโภ ปุริส กรณียํ, ตํ กโรหี”ติ.}

\prose{อถ โข โส ภิกฺขเว ปุริโส ภีโต จตุนฺนํ อาสีวิสานํ อุคฺคเตชานํ โฆรวิสานํ, ภีโต ปญฺจนฺนํ วธกานํ ปจฺจตฺถิกานํ ภีโต ฉฏฺฐสฺส อนฺตรจรสฺส วธกสฺส อุกฺขิตฺตาสิกสฺส ภีโต โจรานํ คามฆาตกานํ เยน วา เตน วา ปลาเยถ, โส ปสฺเสยฺย มหนฺตํ อุทกณฺณวํ โอริมํ ตีรํ สาสงฺกํ สปฺปฏิภยํ, ปาริมํ ตีรํ เขมํ อปฺปฏิภยํ.\csromanspacer{} น จสฺส นาวา สนฺตารณี อุตฺตรเสตุ วา อปารา ปรํ คมนาย.\csromanspacer{} อถ โข ภิกฺขเว ตสฺส ปุริสสฺส เอวมสฺส “อยํ โข มหา-อุทกณฺณโว โอริมํ ตีรํ สาสงฺกํ สปฺปฏิภยํ, ปาริมํ ตีรํ เขมํ อปฺปฏิภยํ.\csromanspacer{} นตฺถิ จ\footnote{น จสฺส (สี, ก), นตฺถสฺส (สฺยา, กํ)} นาวา สนฺตารณี อุตฺตรเสตุ วา อปารา ปารํ คมนาย.\csromanspacer{} ยํนูนาหํ ติณกฏฺฐสาขาปลาสํ สํกฑฺฒิตฺวา กุลฺลํ พนฺธิตฺวา ตํ กุลฺลํ นิสฺสาย หตฺเถหิ จ ปาเทหิ จ วายมมาโน โสตฺถินา ปารํ คจฺเฉยฺยนฺ”ติ.}

\prose{อถ โข โส ภิกฺขเว ปุริโส ติณกฏฺฐสาขาปลาสํ สํกฑฺฒิตฺวา กุลฺลํ พนฺธิตฺวา ตํ กุลฺลํ นิสฺสาย หตฺเถหิ จ ปาเทหิ จ วายมมาโน โสตฺถินา ปารํ คจฺเฉยฺย, ติณฺโณ ปารงฺคโต\footnote{ปารคโต (สี, สฺยา, กํ)} ถเล ติฏฺฐติ พฺราหฺมโณ.}

\prose{อุปมา โข มฺยายํ ภิกฺขเว กตา อตฺถสฺส วิญฺญาปนาย, อยํ เจตฺถ\footnote{อยํ เจเวตฺถ (สี)} อตฺโถ. “จตฺตาโร อาสีวิสฺส อุคฺคเตชา โฆรวิสา”ติ โข ภิกฺขเว จตุนฺเนตํ มหาภูตานํ อธิวจนํ, ปถวีธาตุยา อาโปธาตุยา เตโชธาตุยา วาโยธาตุยา.}

\csromanheader{2}{1. สฬายตนสํยุตฺต}
\prose{\csromanfolio{383}“ปญฺจ วธกา ปจฺจตฺถิกา”ติ โข ภิกฺขเว ปญฺจนฺเนตํ อุปาทานกฺขนฺธานํ อธิวจนํ.\csromanspacer{} เสยฺยถิทํ, รูปุปาทานกฺขนฺทสฺส เวทนุปาทานกฺขนฺธสฺส สญฺญุปาทานกฺขนฺธสฺส สงฺขารุปาทานกฺขนฺธสฺส วิญฺญาณุปาทานกฺขนฺธสฺส.}

\prose{“ฉฏฺโฐ อนฺตรจโร วธโก อุกฺขิตฺตาสิโก”ติ โข ภิกฺขเว นนฺทีราคสฺเสตํ อธิวจนํ.}

\prose{“สุญฺโญ คาโม”ติ โข ภิกฺขเว ฉนฺเนตํ อชฺฌตฺติกานํ อายตนานํ อธิวจนํ.\csromanspacer{} จกฺขุโต เจปิ นํ ภิกฺขเว ปณฺฑิโต พฺยตฺโต เมธาวี อุปปริกฺขติ, ริตฺตกญฺเญว ขายติ, ตุจฺฉกญฺเญว ขายิติ, สุญฺญกญฺเญว ขายติ ฯเปฯ.\csromanspacer{} ชิวฺหาโต เจปิ นํ ภิกฺขเว ฯเปฯ.\csromanspacer{} มนโต เจปิ นํ ภิกฺขเว ปณฺฑิโต พฺยตฺโต เมธาวี อุปปริกฺขติ, ริตฺตกญฺเญว ขายติ, ตุจฺฉกญฺเญว ขายติ, สุญฺญกญฺเญว ขายติ.}

\prose{“โจรา คามฆาตกา”ติ โข ภิกฺขเว ฉนฺเนตํ พาหิรานํ อายตนานํ อธิวจนํ.\csromanspacer{} จกฺขุ ภิกฺขเว หญฺญติ มนาปามนาเปสุ รูเปสุ.\csromanspacer{} โสตํ ภิกฺขเว ฯเปฯ.\csromanspacer{} ฆานํ ภิกฺขเว ฯเปฯ.\csromanspacer{} ชิวฺหา ภิกฺขเว หญฺญติ มนาปามนาเปสุ รเสสุ.\csromanspacer{} กาโย ภิกฺขเว ฯเปฯ.\csromanspacer{} มโน ภิกฺขเว หญฺญติ มนาปามนาเปสุ ธมฺเมสุ.}

\prose{“มหา อุทกณฺณโว”ติ โข ภิกฺขเว จตุนฺเนตํ โอฆานํ อธิวจนํ.\csromanspacer{} กาโมฆสฺส ภโวฆสฺส ทิฏฺโฐฆสฺส อวิชฺโชฆสฺส.}

\prose{“โอริมํ ตีรํ สาสงฺกํ สปฺปฏิภยนฺ”ติ โข ภิกฺขเว สกฺกายสฺเสตํ อธิวจนํ.}

\prose{“ปาริมํ ตีรํ เขมํ อปฺปฏิภยนฺ”ติ โข ภิกฺขเว นิพฺพานสฺเสตํ อธิวจนํ.}

\prose{“กุลฺลนฺ”ติ โข ภิกฺขเว อริยสฺเสตํ อฏฺฐงฺคิกสฺส มคฺคสฺส อธิวจนํ.\csromanspacer{} เสยฺยถิทํ, สมฺมาทิฏฺฐิ ฯเปฯ สมฺมาสมาธิ.}

\prose{“ตสฺส หตฺเถหิ จ ปาเทหิ จ วายาโม”ติ โข ภิกฺขเว วีริยารมฺภสฺเสตํ อธิวจนํ.}

\prose{“ติณฺโณ ปารงฺคโต ถเล ติฏฺฐติ พฺราหฺมโณ”ติ โข ภิกฺขเว อรหโต เอตํ อธิวจนนฺติ.\csromanspacer{}\csromanspacer{}ปฐมํ.}

\gambhira{สฬายตนวคฺคสํยุตฺตปาฬิ}
\csromanheader{2}{2. รโถปมสุตฺต}
\csromanmatikaanchor{mtk.208}
\csromantocmarkref{subsection}{2. รโถปมสุตฺต}{mtk.208}
\bookmark[dest={mtk.208},level=2]{2. รโถปมสุตฺต}
\proseitem{239}{\csromanfolio{384}ตีหิ ภิกฺขเว ธมฺเมหิ สมนฺนาคโต ภิกฺขุ ทิฏฺเฐว ธมฺเม สุขโสมนสฺสพหุโล วิหรติ, โยนิ จสฺส อารทฺธา โหติ อาสวานํ ขยาย.\csromanspacer{} กตเมหิ ตีหิ, อินฺทฺริเยสุ คุตฺตทฺวาโร โหติ โภชเน มตฺตญฺญู ชาคริยํ อนุยุตฺโต.}

\prose{กถญฺจ ภิกฺขเว ภิกฺขุ อินฺทฺริเยสุ คุตฺตทฺวาโร โหติ.\csromanspacer{} อิธ ภิกฺขเว ภิกฺขุ จกฺขุนา รูปํ ทิสฺวา น นิมิตฺตคฺคาหี โหติ นานุพฺยญฺชนคฺคาหี, ยตฺวาธิกรณเมนํ จกฺขุนฺทฺริยํ อสํวุตํ วิหรนฺตํ อภิชฺฌาโทมนสฺสา ปาปกา อกุสลา ธมฺมา อนฺวาสฺสเวยฺยุํ, ตสฺส สํวราย ปฏิปชฺชติ, รกฺขติ จกฺขุนฺทฺริยํ, จกฺขุนฺทฺริเย สํวรํ อาปชฺชติ.\csromanspacer{} โสเตน สทฺทํ สุตฺวา.\csromanspacer{} ฆาเนน คนฺธํ ฆายิตฺวา.\csromanspacer{} ชิวฺหาย รสํ สายิตฺวา.\csromanspacer{} กาเยน โผฏฺฐพฺพํ ผุสิตฺวา.\csromanspacer{} มนสา ธมฺมํ วิญฺญาย น นิมิตฺตคฺคาหี โหติ นานุพฺยญฺชนคฺคาหี, ยตฺวาธิกรณเมนํ มนินฺทฺริยํ อสํวุตํ วิหรนฺตํ อภิชฺฌาโทมนสฺสา ปาปกา อกุสลา ธมฺมา อนฺวาสฺสเวยฺยุํ, ตสฺส สํวราย ปฏิปชฺชติ, รกฺขติ มนินฺทฺริยํ, มนินฺทฺริเย สํวรํ อาปชฺชติ.\csromanspacer{} เสยฺยถาปิ ภิกฺขเว สุภูมิยํ จาตุมหาปเถ อาชญฺญรโถ ยุตฺโต อสฺส ฐิโต โอธสฺตปโตโท\footnote{โอธตปโตโท (สฺยา, กํ), โอธสตปโตโท (อิ)}, ตเมนํ ทกฺโข โยคฺคาจริโย อสฺสทมฺมสารถิ อภิรุหิตฺวา วาเมน หตฺเถน รสฺมิโย คเหตฺวา ทกฺขิเณน หตฺเถน ปโตทํ คเหตฺวา เยนิจฺฉกํ ยทิจฺฉกํ สาเรยฺยปิ ปจฺจาสาเรยฺยปิ.\csromanspacer{} เอวเมว โข ภิกฺขเว ภิกฺขุ อิเมสํ ฉนฺนํ อินฺทฺริยานํ อารกฺขาย สิกฺขติ, สํยมาย สิกฺขติ, ทมาย สิกฺขติ, อุปสมาย สิกฺขติ.\csromanspacer{} เอวํ โข ภิกฺขเว ภิกฺขุ อินฺทฺริเยสุ คุตฺตทฺวาโร โหติ.}

\prose{กถญฺจ ภิกฺขเว ภิกฺขุ โภชเน มตฺตญฺญู โหติ.\csromanspacer{} อิธ ภิกฺขเว ภิกฺขุ ปฏิสงฺขา โยนิโส อาหารํ อาหาเรติ “เนวทวาย น มทาย น มณฺฑนาย น วิภูสนาย, ยาวเทว อิมสฺส กายสฺส ฐิติยา ยาปนาย วิหิํสูปรติยา พฺรหฺมจริยานุคฺคหาย, อิติ ปุราณญฺจ เวทนํ ปฏิหงฺขามิ, นวญฺจ เวทนํ น อุปฺปาเทสฺสามิ, ยาตฺรา จ เม ภวิสฺสติ อนวชฺชตา จ ผาสุวิหาโร จา”ติ.\csromanspacer{} เสยฺยถาปิ ภิกฺขเว ปุริโส วณํ}

\vspace{\tipitakaparskip}
\csromancenter{สฬายตนสํยุตฺต อาลิมฺเปยฺย ยาวเทว โรหนตฺถาย\footnote{โรปนตฺถาย (สี, อิ), เสวนตฺถาย (สฺยา, กํ), โคปนตฺถาย (ก)}, เสยฺยถา วา ปน อกฺขํ อพฺภญฺเชยฺย ยาวเทว ภารสฺส นิตฺถรณตฺถาย, เอวํ โข ภิกฺขเว ภิกฺขุ ปฏิสงฺขา โยนิโส อาหารํ อาหาเรติ “เนว ทวาย น มทาย น มณฺฑนาย น วิภูสนาย ยาวเทว อิมสฺส กายสฺส ฐิติยา ยาปนาย วิหิํสูปรติยา พฺรหฺมจริยานุคฺคหาย อิติ ปุราณญฺจ เวทนํ ปฏิหงฺขามิ, นวญฺจ เวทนํ น อุปฺปาเทสฺสามิ, ยาตฺรา จ เม ภวิสฺสติ อนวชฺชตา จ ผาสุวิหาโร จา”ติ.\csromanspacer{} เอวํ โข ภิกฺขเว ภิกฺขุ โภชเน มตฺตญฺญู โหติ.}
\prose{\csromanfolio{385}กถญฺจ ภิกฺขเว ภิกฺขุ ชาคริยํ อนุยุตฺโต โหติ.\csromanspacer{} อิธ ภิกฺขเว ภิกฺขุ ทิวสํ จงฺกเมน นิสชฺชาย อาวรณีเยหิ ธมฺเมหิ จิตฺตํ ปริโสเธติ, รตฺติยา ปฐมํ ยามํ จงฺกเมน นิสชฺชาย อาวรณีเยหิ ธมฺเมหิ จิตฺตํ ปริโสเธติ, รตฺติยา มชฺฌิมํ ยามํ ทกฺขิเณน ปสฺเสน สีหเสยฺยํ กปฺเปติ ปาเท ปาทํ อจฺจาธาย สโต สมฺปชาโน อุฏฺฐานสญฺญํ มนสิ กริตฺวา, รตฺติยา ปจฺฉิมํ ยามํ ปจฺจุฏฺฐาย จงฺกเมน นิสชฺชาย อาวรณีเยหิ ธมฺเมหิ จิตฺตํ ปริโสเธติ.\csromanspacer{} เอวํ โข ภิกฺขเว ภิกฺขุ ชาคริยํ อนุยุตฺโต โหติ.\csromanspacer{} อิเมหิ โข ภิกฺขเว ตีหิ ธมฺเมหิ สมนฺนาคโต ภิกฺขุ ทิฏฺเฐว ธมฺเม สุขโสมนสฺสพหุโล วิหรติ, โยนิ จสฺส อารทฺธา โหติ อาสวานํ ขยายาติ.\csromanspacer{}\csromanspacer{}ทุติยํ.}

\csromanheader{2}{3. กุมฺโมปมสุตฺต}
\csromanmatikaanchor{mtk.209}
\csromantocmarkref{subsection}{3. กุมฺโมปมสุตฺต}{mtk.209}
\bookmark[dest={mtk.209},level=2]{3. กุมฺโมปมสุตฺต}
\proseitem{240}{ภูตปุพฺพํ ภิกฺขเว กุมฺโม กจฺฉโป สายนฺหสมยํ อนุนทีตีเร โคจรปสุโต อโหสิ.\csromanspacer{} สิงฺคาโลปิ\footnote{สิคาโลปิ (สี, สฺยา, กํ, อิ)} โข ภิกฺขเว สายนฺหสมยํ อนุนทีตีเร โคจรปสุโต อโหสิ.\csromanspacer{} อทฺทสา โข ภิกฺขเว กุมฺโม กจฺฉโป สิงฺคาลํ ทูเรโตว โคจรปสุตํ, ทิสฺวาน โสณฺฑิปญฺจมานิ องฺคานิ สเก กปาเล สโมทหิตฺวา อปฺโปสฺสุกฺโก ตุณฺหีภูโต สงฺกสายติ.\csromanspacer{} สิงฺคาโลปิ โข ภิกฺขเว อทฺทส กุมฺมํ กจฺฉปํ ทูรโตว โคจรปสุตํ, ทิสฺวาน เยน กุมฺโม กจฺฉโป เตนุปสงฺกมิ, อุปสงฺกมิตฺวา กุมฺมํ กจฺฉปํ ปจฺจุปฏฺฐิโต อโหสิ “ยทายํ กุมฺโม กจฺฉโป โสณฺฑิปญฺจมานํ องฺคานํ อญฺญตรํ วา อญฺญตรํ วา องฺคํ อภินินฺนาเมสฺสติ, ตตฺเถว นํ คเหตฺวา อุทฺทาลิตฺวา ขาทิสฺสามี”ติ.\csromanspacer{} ยทา โข ภิกฺขเว กุมฺโม}

\gambhira{สฬายตนวคฺคสํยุตฺตปาฬิ}
\csromanmatikaanchor{mtk.210}
\csromantocmarkref{subsection}{4-5. ทารุกฺขนฺโธปมสุตฺต}{mtk.210}
\bookmark[dest={mtk.210},level=2]{4-5. ทารุกฺขนฺโธปมสุตฺต}
\prose{\csromanfolio{386}กจฺฉโป โสณฺฑิปญฺจมานํ องฺคานํ อญฺญตรํ วา อญฺญตรํ วา องฺคํ น อภินินฺนามิ.\csromanspacer{} อถ สิงฺคาโล กุมฺมมฺหา นิพฺพิชฺช ปกฺกามิ โอตารํ อลภมาโน.}

\prose{เอวเมว โข ภิกฺขเว ตุเมฺหปิ มาโร ปาปิมา สตตํ สมิตํ ปจฺจุปฏฺฐิโต “อปฺเปว นามาหํ อิเมสํ จกฺขุโต วา โอตารํ ลเภยฺยํ ฯเปฯ ชิวฺหาโต วา โอตารํ ลเภยฺยํ ฯเปฯ มนโต วา โอตารํ ลเภยฺยนฺ”ติ.\csromanspacer{} ตสฺมาติห ภิกฺขเว อินฺทฺริเยสุ คุตฺตทฺวารา วิหรถ, จกฺขุนา รูปํ ทิสฺวา มา นิมิตฺตคฺคาหิโน อหุวตฺถ, มา อนุพฺยญฺชนคฺคาหิโน, ยตฺวาธิกรณเมนํ จกฺขุนฺทฺริยํ อสํวุตํ วิหรนฺตํ อภิชฺฌาโทมนสฺสา ปาปกา อกุสลา ธมฺมา อนฺวาสฺสเวยฺยุํ.\csromanspacer{} ตสฺส สํวราย ปฏิปชฺชถ, รกฺขถ จกฺขุนฺทฺริยํ, จกฺขุนฺทฺริเย สํวรํ อาปชฺชถ.\csromanspacer{} โสตน สทฺทํ สุตฺวา.\csromanspacer{} ฆาเนน คนฺธํ ฆายิตฺวา.\csromanspacer{} ชิวฺหาย รสํ สายิตฺวา.\csromanspacer{} กาเยน โผฏฺฐพฺพํ ผุสิตฺวา.\csromanspacer{} มนสา ธมฺมํ วิญฺญาย มา นิมิตฺตคฺคาหิโน อหุวตฺถ, มา อนุพฺยญฺชนคฺคาหิโน, ยตฺวาธิกรณเมนํ มนินฺทฺริยํ อสํวุตํ วิหรนฺตํ อภิชฺฌาโทมนสฺสา ปาปกา อกุสลา ธมฺมา อนฺวาสฺสเวยฺยุํ.\csromanspacer{} ตสฺส สํวราย ปฏิปชฺชถ, รกฺขถ มนินฺทฺริยํ, มนินฺทฺริเย สํวรํ อาปชฺชถ.\csromanspacer{} ยโต ตุเมฺห ภิกฺขเว อินฺทฺริเยสุ คุตฺตทฺวารา วิหริสฺสถ, อถ ตุเมฺหหิปิ มาโร ปาปิมา นิพฺพิชฺช ปกฺกมิสฺสติ โอตารํ อลภมาโน กุมฺมมฺหาว สิงฺคาโลติ.}

\csromangathameasure{กุมฺโมว องฺคานิ สเก กปาเล, \\ สโมทหํ ภิกฺขุ มโนวิตกฺเก. \\ อนิสฺสิโต อญฺญมเหฐยาโน, \\ ปรินิพฺพุโตนูปวเทยฺย กิญฺจีติ. ตติยํ.}
\csromangathagroup{\csromangathastanza{กุมฺโมว องฺคานิ สเก กปาเล, \\ สโมทหํ ภิกฺขุ มโนวิตกฺเก. \\ อนิสฺสิโต อญฺญมเหฐยาโน, \\ ปรินิพฺพุโตนูปวเทยฺย กิญฺจีติ.\csromanspacer{} ตติยํ.}}

\csromanheader{2}{4. ปฐมทารุกฺขนฺโธปมสุตฺต}
\proseitem{241}{เอกํ สมยํ ภควา โกสมฺพิยํ วิหรติ คงฺคาย นทิยา ตีเร.\csromanspacer{} อทฺทสา โข ภควา มหนฺตํ ทารุกฺขนฺธํ คงฺคาย นทิยา โสเตน วุยฺหมานํ, ทิสฺวาน ภิกฺขู อามนฺเตสิ “ปสฺสถ โน ตุเมฺห ภิกฺขเว อมุํ มหนฺตํ ทารุกฺขนฺธํ คงฺคาย นทิยา โสเตน วุยฺหมานนฺ”ติ.\csromanspacer{} เอวํ ภนฺเต.\csromanspacer{} สเจ โส ภิกฺขเว ทารุกฺขนฺโธ น โอริมํ ตีรํ อุปคจฺฉติ, น ปาริมํ ตีรํ อุปคจฺฉติ, น มชฺเฌ สํสีทิสฺสติ, น ถเล อุสฺสีทิสฺสติ, น มนุสฺสคฺคาโห คเหสฺสติ,}

\vspace{\tipitakaparskip}
\csromancenter{สฬายตนสํยุตฺต น อมนุสฺสคฺคาโห คเหสฺสติ, น อาวฏฺฏคฺคาโห คเหสฺสติ, น อนฺโตปูติ ภวิสฺสติ.\csromanspacer{} เอวํ หิ โส ภิกฺขเว ทารุกฺขนฺโธ สมุทฺทนินฺโน ภวิสฺสติ สมุทฺทโปโณ สมุทฺทปพฺภาโร.\csromanspacer{} ตํ กิสฺส เหตุ, สมุทฺทนินฺโน ภิกฺขเว คงฺคาย นทิยา โสโต สมุทฺทโปโณ สมุทฺทปพฺภาโร.}
\prose{\csromanfolio{387}เอวเมว โข ภิกฺขเว สเจ ตุเมฺหปิ น โอริมํ ตีรํ อุปคจฺฉถ, น ปาริมํ ตีรํ อุปคจฺฉถ, น มชฺเฌ สํสีทิสฺสถ, น ถเล อุสฺสีทิสฺสถ, น มนุสฺสคฺคาโห คเหสฺสติ, น อมนุสฺสคฺคาโห คเหสฺสติ, น อาวฏฺฏคฺคาโห คเหสฺสติ, น อนฺโตปูตี ภวิสฺสถ.\csromanspacer{} เอวํ ตุเมฺห ภิกฺขเว นิพฺพานนินฺนา ภวิสฺสถ, นิพฺพานโปณา นิพฺพานปพฺภารา.\csromanspacer{} ตํ กิสฺส เหตุ, นิพฺพานนินฺนา ภิกฺขเว สมฺมาทิฏฺฐิ นิพฺพานโปณา นิพฺพานปพฺภาราติ.\csromanspacer{} เอวํ วุตฺเต อญฺญตโร ภิกฺขุ ภควนฺตํ เอตทโวจ “กิํ นุ โข ภนฺเต โอริมํ ตีรํ กิํ ปาริมํ ตีรํ โก มชฺเฌ สํสาโท\footnote{สํสีโท (ก), สํสีทิโต (สฺยา, กํ)} โก ถเล อุสฺสาโท โก มนุสฺสคฺคาโห โก อมนุสฺสคฺคาโห โก อาวฏฺฏคฺคาโห โก อนฺโตปูติภาโว”ติ.}

\prose{“โอริมํ ตีรนฺ”ติ โข ภิกฺขุ ฉนฺเนตํ อชฺฌตฺติกานํ อายตนานํ อธิวจนํ. “ปาริมํ ตีรนฺ”ติ โข ภิกฺขุ ฉนฺเนตํ พาหิรานํ อายตนานํ อธิวจนํ. “มชฺเฌ สํสาโท”ติ โข ภิกฺขุ นนฺทีราคสฺเสตํ อธิวจนํ. “ถเล อุสฺสาโท”ติ โข ภิกฺขุ อสฺมิมานสฺเสตํ อธิวจนํ.}

\prose{กตโม จ ภิกฺขุ มนุสฺสคฺคาโห.\csromanspacer{} อิธ ภิกฺขุ คิหีหิ สํสฏฺโฐ\footnote{คิหิสํสฏฺโฐ (ก)} วิหรติ, สหนนฺที สหโสกี สุขิเตสุ สุขิโต ทุกฺขิเตสุ ทุกฺขิโต อุปฺปนฺเนสุ กิจฺจกรณีเยสุ อตฺตนา เตสุ โยคํ อาปชฺชติ.\csromanspacer{} อยํ วุจฺจติ ภิกฺขุ มนุสฺสคฺคาโห.}

\prose{กตโม จ ภิกฺขุ อมนุสฺสคฺคาโห.\csromanspacer{} อิธ ภิกฺขุ เอกจฺโจ อญฺญตรํ เทวนิกายํ ปณิธาย พฺรหฺมจริยํ จรติ, อิมินาหํ สีเลน วา วเตน วา ตเปน วา พฺรหฺมจริเยน วา เทโว วา ภวิสฺสามิ เทวญฺญตโร วาติ.\csromanspacer{} อยํ วุจฺจติ ภิกฺขุ อมนุสฺสคฺคาโห. “อาวฏฺฏคฺคาโห”ติ โข ภิกฺขุ ปญฺจนฺเนตํ กามคุณานํ อธิวจนํ.}

\prose{กตโม จ ภิกฺขุ อนฺโตปูติภาโว.\csromanspacer{} อิธ ภิกฺขุ เอกจฺโจ ทุสฺสีโล โหติ, ปาปธมฺโม อสุจิสงฺกสฺสรสมาจาโร ปฏิจฺฉนฺนกมฺมนฺโต}

\gambhira{สฬายตนวคฺคสํยุตฺตปาฬิ}
\prose{\csromanfolio{388}อสฺสมโณ สมณปฏิญฺโญ อพฺรหฺมจารี พฺรหฺมจาริปฏิญฺโญ อนฺโตปูติ อวสฺสุโต กสมฺพุชาโต.\csromanspacer{} อยํ วุจฺจติ ภิกฺขุ อนฺโตปูติภาโวติ.}

\prose{เตน โข ปน สมเยน นนฺโท โคปาลโก ภควโต อวิทูเร ฐิโต โหติ.\csromanspacer{} อถ โข นนฺโท โคปาลโก ภควนฺตํ เอตทโวจ “อหํ โข ภนฺเต น โอริมํ ตีรํ อุปคจฺฉามิ, น ปาริมํ ตีรํ อุปคจฺฉามิ, น มชฺเฌ สํสีทิสฺสามิ, น ถเล อุสฺสีทิสฺสามิ, น มํ มนุสฺสคฺคาโห คเหสฺสติ, น อมนุสฺสคฺคาโห คเหสฺสติ, น อาวฏฺฏคฺคาโห คเหสฺสติ, น อนฺโตปูติ ภวิสฺสามิ.\csromanspacer{} ลเภยฺยาหํ ภนฺเต ภควโต สนฺติเก ปพฺพชฺชํ ลเภยฺยํ อุปสมฺปทนฺ”ติ.\csromanspacer{} เตน หิ ตฺวํ นนฺท สามิกานํ คาโว นิยฺยาเตหีติ\footnote{นียฺยาเทหีติ (สี), นิยฺยาเทหีติ (สฺยา, กํ, อิ)}.\csromanspacer{} คมิสฺสนฺติ ภนฺเต คาโว วจฺฉคิทฺธินิโยติ.\csromanspacer{} นิยฺยาเตเหว ตฺวํ นนฺท สามิกานํ คาโวติ.\csromanspacer{} อถ โข นนฺโท โคปาลโก สามิกานํ คาโว นิยฺยาเตตฺวา เยน ภควา เตนุปสงฺกมิ, อุปสงฺกมิตฺวา ภควนฺตํ เอตทโวจ “นิยฺยาติตา\footnote{นิยฺยาตา (สฺยา, กํ, ก, สี-ฐ)} ภนฺเต สามิกานํ คาโว, ลเภยฺยาหํ ภนฺเต ภควโต สนฺติเก ปพฺพชฺชํ, ลเภยฺยํ อุปสมฺปทนฺ”ติ.\csromanspacer{} อลตฺถ โข นนฺโท โคปาลโก ภควโต สนฺติเก ปพฺพชฺชํ, อลตฺถ อุปสมฺปทํ, อจิรูปสมฺปนฺโน จ ปนายสฺมา นนฺโท เอโก วูปกฏฺโฐ ฯเปฯ.\csromanspacer{} อญฺญตโร จ ปนายสฺมา นนฺโท อรหตํ อโหสีติ.\csromanspacer{}\csromanspacer{}จตุตฺถํ.}

\csromanheader{2}{5. ทุติยทารุกฺขนฺโธปมสุตฺต}
\proseitem{242}{เอกํ สมยํ ภควา กิมิลายํ\footnote{กิมฺพิลายํ (สี, อิ), กิมฺมิลายํ (สฺยา, กํ)} วิหรติ คงฺคาย นิทิยา ตีเร.\csromanspacer{} อทฺทสา โข ภควา มหนฺตํ ทารุกฺขนฺธํ คงฺคาย นทิยา โสเตน วุยฺหมานํ, ทิสฺวาน ภิกฺขู อามนฺเตสิ “ปสฺสถ โน ตุเมฺห ภิกฺขเว อมุํ มหนฺตํ ทารุกฺขนฺธํ คงฺคาย นทิยา โสเตน วุยฺหมานนฺ”ติ.\csromanspacer{} เอวํ ภนฺเต ฯเปฯ.\csromanspacer{} เอวํ วุตฺเต อายสฺมา กิมิโล ภควนฺตํ เอตทโวจ\csromandash{} กิํ นุ โข ภนฺเต โอริมํ ตีรํ ฯเปฯ.\csromanspacer{} กตโม จ กิมิล อนฺโตปูติภาโว.\csromanspacer{} อิธ กิมิล ภิกฺขุ อญฺญตรํ สํกิลิฏฺฐํ อาปตฺติํ อาปนฺโน โหติ, ยถารูปาย อาปตฺติยา น วุฏฺฐานํ ปญฺญายติ.\csromanspacer{} อยํ วุจฺจติ กิมิล อนฺโตปูติภาโวติ.\csromanspacer{}\csromanspacer{}ปญฺจมํ.}

\csromanheader{2}{1. สฬายตนสํยุตฺต \textbf{6. อวสฺสุตปริยายสุตฺต}}
\csromanmatikaanchor{mtk.211}
\csromantocmarkref{subsection}{6. อวสฺสุตปริยายสุตฺต}{mtk.211}
\bookmark[dest={mtk.211},level=2]{6. อวสฺสุตปริยายสุตฺต}
\proseitem{243}{\csromanfolio{389}เอกํ สมยํ ภควา สกฺเกสุ วิหรติ กปิลวตฺถุสฺมิํ นิโคฺรธาราเม.\csromanspacer{} เตน โข ปน สมเยน กาปิลวตฺถวานํ สกฺยานํ นวํ สนฺถาคารํ\footnote{สนฺธาคารํ (ก)} อจิรการิตํ โหติ อนชฺฌาวุฏฺฐํ\footnote{อนชฺฌาวุตฺถํ (สี, สฺยา, กํ, อิ)} สมเณน วา พฺราหฺมเณน วา เกนจิ วา มนุสฺสาภูเตน.\csromanspacer{} อถ โข กาปิลวตฺถวา สกฺยา เยน ภควา เตนุปสงฺกมิํสุ, อุปสงฺกมิตฺวา ภควนฺตํ อภิวาเทตฺวา เอกมนฺตํ นิสีทิํสุ, เอกมนฺตํ นิสินฺนา โข กาปิลวตฺถวา สกฺยา ภควนฺตํ เอตทโวจุํ “อิธ ภนฺเต กาปิลวตฺถวานํ สกฺยานํ นวํ สนฺถาคารํ อจิรการิตํ\footnote{อจิรการิตํ โหติ (ก)} อนชฺฌาวุฏฺฐํ สมเณน วา พฺราหฺมเณน วา เกนจิ วา มนุสฺสภูเตน, ตํ ภนฺเต ภควา ปฐมํ ปริภุญฺชตุ, ภควตา ปฐมํ ปริภุตฺตํ, ปจฺฉา กาปิลวตฺถวา สกฺยา ปริภุญฺชิสฺสนฺติ, ตทสฺส กาปิลวตฺถวานํ สกฺยานํ ทีฆรตฺตํ หิตาย สุขายา”ติ.\csromanspacer{} อธิวาเสสิ ภควา ตุณฺหีภาเวน.}

\prose{อถ โข กาปิลวตฺถวา สกฺยา ภควโต อธิวาสนํ วิทิตฺวา อุฏฺฐายาสนา ภควนฺตํ อภิวาเทตฺวา ปทกฺขิณํ กตฺวา เยน นวํ สนฺถาคารํ เตนุปสงฺกมิํสุ, อุปสงฺกมิตฺวา สพฺพสนฺถริํ\footnote{สพฺพสนฺถริํ สนฺถตํ (ก)} สนฺถคารํ สนฺถริตฺวา อาสนานิ ปญฺญาเปตฺวา อุทกมณิกํ ปติฏฺฐาเปตฺวา เตลปฺปทีปํ อาโรเปตฺวา เยน ภควา เตนุปสงฺกมิํสุ, อุปสงฺกมิตฺวา ภควนฺตํ เอตทโวจุํ “สพฺพสนฺถริสนฺถตํ\footnote{สพฺพสนฺถริํ สนฺถตํ (สี, อิ, ก)} ภนฺเต สนฺถาคารํ, อาสนานิ ปญฺญตฺตานิ, อุทกมณิโก ปติฏฺฐาปิโต, เตลปฺปทีโป อาโรปิโต, ยสฺส ทานิ ภนฺเต ภควา กาลํ มญฺญตี”ติ.\csromanspacer{} อถ โข ภควา นิวาเสตฺวา ปตฺตจีวรมาทาย สทฺธิํ ภิกฺขุสํเฆน เยน นวํ สนฺถาคารํ เตนุปสงฺกมิ, อุปสงฺกมิตฺวา ปาเท ปกฺขาเลตฺวา สนฺถาคารํ ปวิสิตฺวา มชฺฌิมํ ถมฺภํ นิสฺสาย ปุรตฺถาภิมุโข นิสีทิ.\csromanspacer{} ภิกฺขุสํโฆปิ โข ปาเท ปกฺขาเลตฺวา สนฺถาคารํ ปวิสิตฺวา ปจฺฉิมํ ภิตฺติํ นิสฺสาย ปุรตฺถาภิมุโข นิสีทิ ภควนฺตํเยว ปุรกฺขตฺวา.\csromanspacer{} กาปิลวตฺถวา สกฺยา ปาเท ปกฺขาเลตฺวา สนฺถาคารํ ปวิสิตฺวา ปุรตฺถิมํ ภิตฺติํ นิสฺสาย ปจฺฉิมาภิมุขา นิสีทิํสุ ภควนฺตํเยว ปุรกฺขตฺวา.\csromanspacer{} อถ โข ภควา กาปิลวตฺถเว สเกฺย พหุเทว รตฺติํ ธมฺมิยา กถาย สนฺทสฺเสตฺวา สมาทเปตฺวา สมุตฺเตเชตฺวา สมฺปหํเสตฺวา อุยฺโยเชสิ “อภิกฺกนฺตา โข โคตมา รตฺติ, ยสฺส ทานิ กาลํ}

\gambhira{สฬายตนวคฺคสํยุตฺตปาฬิ}
\prose{\csromanfolio{390}มญฺญถา”ติ. “เอวํ ภนฺเต”ติ โข กาปิลวตฺถวา สกฺยา ภควโต ปฏิสฺสุตฺวา อุฏฺฐายาสนา ภควนฺตํ อภิวาเทตฺวา ปทกฺขิณํ กตฺวา ปกฺกมิํสุ.}

\prose{อถ โข ภควา อจิรปกฺกนฺเตสุ กาปิลวตฺถเวสุ สเกฺยสุ อายสฺมนฺตํ มหาโมคฺคลฺลานํ อามนฺเตสิ “วิคตถินมิทฺโธ โข โมคฺคลฺลาน ภิกฺขุสํโฆ, ปฏิภาตุ ตํ โมคฺคลฺลาน ภิกฺขูนํ ธมฺมี กถา, ปิฏฺฐิ เม อาคิลายติ, ตมหํ อายมิสฺสามี”ติ. “เอวํ ภนฺเต”ติ โข อายสฺมา มหาโมคฺคลฺลาโน ภควโต ปจฺจสฺโสสิ.\csromanspacer{} อถ โข ภควา จตุคฺคุณํ สงฺฆาฏิํ ปญฺญเปตฺวา ทกฺขิเณน ปสฺเสน สีหเสยฺยํ กปฺเปสิ ปาเท ปาทํ อจฺจาธาย สโต สมฺปชาโน อุฏฺฐานสญฺญํ มนสิ กริตฺวา.\csromanspacer{} ตตฺร โข อายสฺมา มหาโมคฺคลฺลาโน ภิกฺขู อามนฺเตสิ “อาวุโส ภิกฺขเว”ติ. “อาวุโส”ติ โข เต ภิกฺขู อายสฺมโต มหาโมคฺคลฺลานสฺส ปจฺจสฺโสสุํ.\csromanspacer{} อายสฺมา มหาโมคฺคลฺลาโน เอตทโวจ “อวสฺสุตปริยายญฺจ โว อาวุโส เทเสสฺสามิ อนวสฺสุตปริยายญฺจ, ตํ สุณาถ สาธุกํ มนสิ กโรถ, ภาสิสฺสามี”ติ. “เอวมาวุโส”ติ โข เต ภิกฺขู อายสฺมโต มหาโมคฺคลฺลานสฺส ปจฺจสฺโสสุํ.\csromanspacer{} อายสฺมา มหาโมคฺคลฺลาโน เอตทโวจ\csromandash{}}

\prose{กถํ อาวุโส อวสฺสุโต โหติ.\csromanspacer{} อิธาวุโส ภิกฺขุ จกฺขุนา รูปํ ทิสฺวา ปิยรูเป รูเป อธิมุจฺจติ, อปฺปิยรูเป รูเป พฺยาปชฺชติ, อนุปฏฺฐิตกายสฺสตี วิหรติ ปริตฺตเจตโส, ตญฺจ เจโตวิมุตฺติํ ปญฺญาวิมุตฺติํ ยถาภูตํ นปฺปชานาติ, ยตฺถสฺส เต อุปฺปนฺนา ปาปกา อกุสลา ธมฺมา อปริเสสา นิรุชฺฌนฺติ ฯเปฯ.\csromanspacer{} ชิวฺหาย รสํ สายิตฺวา ฯเปฯ.\csromanspacer{} มนสา ธมฺมํ วิญฺญาย ปิยรูเป ธมฺเม อธิมุจฺจติ, อปฺปิยรูเป ธมฺเม พฺยาปชฺชติ, อนุปฏฺฐิตกายสฺสติ จ วิหรติ ปริตฺตเจตโส, ตญฺจ เจโตวิมุตฺติํ ปญฺญาวิมุตฺติํ ยถาภูตํ นปฺปชานาติ, ยตฺถสฺส เต อุปฺปนฺนา ปาปกา อกุสลา ธมฺมา อปริเสสา นิรุชฺฌนฺติ.\csromanspacer{} อยํ วุจฺจติ อาวุโส ภิกฺขุ อวสฺสุโต จกฺขุวิญฺเญยฺเยสุ รูเปสุ ฯเปฯ อวสฺสุโต ชิวฺหาวิญฺเญยฺเยสุ รเสสุ ฯเปฯ อวสฺสุโต มโนวิญฺเญยฺเยสุ ธมฺเมสุ.\csromanspacer{} เอวํวิหาริํ จาวุโส ภิกฺขุํ จกฺขุโต เจปิ นํ มาโร อุปสงฺกมติ, ลภเตว\footnote{ลเภถ (ก)} มาโร}

\vspace{\tipitakaparskip}
\csromancenter{สฬายตนสํยุตฺต โอตารํ, ลภติ\footnote{ลเภถ (ก)} มาโร อารมฺมณํ ฯเปฯ.\csromanspacer{} ชิวฺหาโต เจปิ นํ มาโร อุปสงฺกมติ, ลภเตว มาโร โอตารํ, ลภติ มาโร อารมฺมณํ ฯเปฯ.\csromanspacer{} มนโต เจปิ นํ มาโร อุปสงฺกมติ, ลภเตว มาโร โอตารํ, ลภติ มาโร อารมฺมณํ.}
\prose{\csromanfolio{391}เสยฺยถาปิ อาวุโส นฬาคารํ วา ติณาคารํ วา สุกฺขํ โกลาปํ เตโรวสฺสิกํ, ปุริตฺถิมาย เจปิ นํ ทิสาย ปุริโส อาทิตฺตาย ติณุกฺกาย อุปสงฺกเมยฺย, ลเภเถว\footnote{ลเภถ (ก)} อคฺคิ โอตารํ, ลเภถ อคฺคิ อารมฺมณํ.\csromanspacer{} ปจฺฉิมาย เจปิ นํ ทิสาย ปุริโส อาทิตฺตาย ติณุกฺกาย อุปสงฺกเมยฺย ฯเปฯ.\csromanspacer{} อุตฺตราย เจปิ นํ ทิสาย.\csromanspacer{} ทกฺขิณาย เจปิ นํ ทิสาย.\csromanspacer{} เหฏฺฐิมโต เจปิ นํ.\csromanspacer{} อุปริมโต เจปิ นํ.\csromanspacer{} ยโต กุโตจิ เจปิ นํ ปุริโส อาทิตฺตาย ติณุกฺกาย อุปสงฺกเมยฺย, ลเภเถว อคฺคิ โอตารํ, ลเภถ อคฺคิ อารมฺมณํ.\csromanspacer{} เอวเมว โข อาวุโส เอวํวิหาริํ ภิกฺขุํ จกฺขุโต เจปิ นํ มาโร อุปสงฺกมติ, ลภเตว มาโร โอตารํ, ลภติ มาโร อารมฺมณํ ฯเปฯ ชิวฺหาโต เจปิ นํ มาโร อุปสงฺกมติ ฯเปฯ มนโต เจปิ นํ มาโร อุปสงฺกมติ, ลภเตว มาโร โอตารํ, ลภติ มาโร อารมฺมณํ.\csromanspacer{} เอวํวิหาริํ จาวุโส ภิกฺขุํ รูปา อธิภํสุ, น ภิกฺขุ รูเป อธิโภสิ.\csromanspacer{} สทฺทา ภิกฺขุํ อธิภํสุ, น ภิกฺขุ สทฺเท อธิโภสิ.\csromanspacer{} คนฺธา ภิกฺขุํ อธิภํสุ, น ภิกฺขุ คนฺเธ อธิโภสิ.\csromanspacer{} รสา ภิกฺขุํ อธิภุํสุ, น ภิกฺขุ รเส อธิโภสิ.\csromanspacer{} โผฏฺฐพฺพา ภิกฺขุํ อธิภํสุ, น ภิกฺขุ โผฏฺฐพฺเพ อธิโภสิ.\csromanspacer{} ธมฺมา ภิกฺขุํ อธิภํสุ, น ภิกฺขุ ธมฺเม อธิโภสิ.\csromanspacer{} อยํ วุจฺจตาวุโส ภิกฺขุ รูปาธิภูโต สทฺทาธิภูโต คนฺธาธิภูโต รสาธิภูโต โผฏฺฐพฺพาธิภูโต ธมฺมาธิภูโต อธิภูโต อนธิภู\footnote{อนธิภูโต (สี, สฺยา, กํ, ก)}, อธิภํสุ นํ ปาปกา อกุสลา ธมฺมา สํกิเลสิกา โปโนภวิกา สทรา ทุกฺขวิปากา อายติํ ชาติชรามรณิยา.\csromanspacer{} เอวํ โข อาวุโส อวสฺสุโต โหติ.}

\prose{กถญฺจาวุโส อนวสฺสุโต โหติ, อิธาวุโส ภิกฺขุ จกฺขุนา รูปํ ทิสฺวา ปิยรูเป รูเป นาธิมุจฺจติ, อปฺปิยรูเป รูเป น พฺยาปชฺชติ, อุปฏฺฐิตกายสฺสติ จ วิหรติ อปฺปมาณเจตโส.\csromanspacer{} ตญฺจ เจโตวิมุตฺติํ ปญฺญาวิมุตฺติํ ยถาภูตํ ปชานาติ, ยตฺถสฺส เต อุปฺปนฺนา ปาปกา อกุสลา ธมฺมา อปริเสสา นิรุชฺฌนฺติ ฯเปฯ ชิวฺหาย รสํ สายิตฺวา ฯเปฯ}

\gambhira{สฬายตนวคฺคสํยุตฺตปาฬิ}
\prose{\csromanfolio{392}มนสา ธมฺมํ วิญฺญาย ปิยรูเป ธมฺเม นาธิมุจฺจติ, อปฺปิยรูเป ธมฺเม น พฺยาปชฺชติ, อุปฏฺฐิตกายสฺสติ จ วิหรติ อปฺปมาณเจตโส.\csromanspacer{} ตญฺจ เจโตวิมุตฺติํ ปญฺญาวิมุตฺติํ ยถาภูตํ ปชานาติ, ยตฺถสฺส เต อุปฺปนฺนา ปาปกา อกุสลา ธมฺมา อปริเสสา นิรุชฺฌนฺติ.\csromanspacer{} อยํ วุจฺจตาวุโส ภิกฺขุ อนวสฺสุโต จกฺขุวิญฺเญยฺเยสุ รูเปสุ ฯเปฯ อนวสฺสุโต มโนวิญฺเญยฺเยสุ ธมฺเมสุ.\csromanspacer{} เอวํวิหาริํ จาวุโส ภิกฺขุํ จกฺขุโต เจปิ นํ มาโร อุปสงฺกมติ, เนว ลภติ มาโร โอตารํ, น ลภติ มาโร อารมฺมณํ ฯเปฯ ชิวฺหาโต เจปิ นํ มาโร อุปสงฺกมติ ฯเปฯ มนโต เจปิ นํ มาโร อุปสงฺกมติ, เนว ลภติ มาโร โอตารํ, น ลภติ มาโร อารมฺมณํ.}

\prose{เสยฺยถาปิ อาวุโส กูฏาคารํ วา สาลา วา พหลมตฺติกา อทฺทาวเลปนา, ปุรตฺถิมาย เจปิ นํ ทิสาย ปุริโส อาทิตฺตาย ติณุกฺกาย อุปสงฺกเมยฺย, เนว ลเภถ อคฺคิ โอตารํ, น ลเภถ อคฺคิ อารมฺมณํ ฯเปฯ ปจฺฉิมาย เจปิ นํ.\csromanspacer{} อุตฺตราย เจปิ นํ.\csromanspacer{} ทกฺขิณาย เจปิ นํ.\csromanspacer{} เหฏฺฐิมโต เจปิ นํ.\csromanspacer{} อุปริมโต เจปิ นํ.\csromanspacer{} ยโต กุโตจิ เจปิ นํ ปุริโส อาทิตฺตาย ติณุกฺกาย อุปสงฺกเมยฺย, เนว ลเภถ อคฺคิ โอตารํ, น ลเภถ อคฺคิ อารมฺมณํ.\csromanspacer{} เอวเมว โข อาวุโส เอวํวิหาริํ ภิกฺขุํ จกฺขุโต เจปิ นํ มาโร อุปสงฺกมติ, เนว ลภติ มาโร โอตารํ, น ลภติ มาโร อารมฺมณํ ฯเปฯ มนโต เจปิ นํ มาโร อุปสงฺกมติ, เนว ลภติ มาโร โอตารํ, น ลภติ มาโร อารมฺมณํ.\csromanspacer{} เอวํวิหารี จาวุโส ภิกฺขุ รูเป อธิโภสิ, น รูปา ภิกฺขุํ อธิภํสุ, สทฺเท ภิกฺขุ อธิโภสิ, น สทฺทา ภิกฺขุํ อธิภํสุ.\csromanspacer{} คนฺเธ ภิกฺขุ อธิโภสิ, น คนฺธา ภิกฺขุํ อธิภํสุ.\csromanspacer{} รเส ภิกฺขุ อธิโภสิ, น รสา ภิกฺขุํ อธิภํสุ.\csromanspacer{} โผฏฺฐพฺเพ ภิกฺขุ อธิโภสิ, น โผฏฺฐพฺพา ภิกฺขุํ อธิภํสุ.\csromanspacer{} ธมฺเม ภิกฺขุ อธิโภสิ, น ธมฺมา ภิกฺขุํ อธิภํสุ.\csromanspacer{} อยํ วุจฺจตาวุโส ภิกฺขุ รูปาธิภู สทฺทาธิภู คนฺธาธิภู รสาธิภู โผฏฺฐพฺพาธิภู ธมฺมาธิภู อธิภู อนธิภูโต\footnote{อนธิภูโต เกหิจิ กิเลเสหิ (ก)}, อธิโภสิ เต ปาปเก อกุสเล ธมฺเม สํกิเลสิเก โปโนภวิเก สทเร ทุกฺขวิปาเก อายติํ ชาติชรามรณิเย.\csromanspacer{} เอวํ โข อาวุโส อนวสฺสุโต โหตีติ.}

\csromanheader{2}{1. สฬายตนสํยุตฺต}
\prose{\csromanfolio{393}อถ โข ภควา อุฏฺฐหิตฺวา อายสฺมนฺตํ มหาโมคฺคลฺลานํ อามนฺเตสิ “สาธุ สาธุ โมคฺคลฺลาน, สาธุ โข ตฺวํ โมคฺคลฺลาน ภิกฺขูนํ อวสฺสุตปริยายญฺจ อนวสฺสุตปริยายญฺจ อภาสี”ติ.}

\prose{อิทมโวจ อายสฺมา มหาโมคฺคลฺลาโน, สมนุญฺโญ สตฺถา อโหสิ.\csromanspacer{} อตฺตมนา เต ภิกฺขู อายสฺมโต มหาโมคฺคลฺลานสฺส ภาสิตํ อภินนฺทุนฺติ.\csromanspacer{}\csromanspacer{}ฉฏฺฐํ.}

\csromanheader{2}{7. ทุกฺขธมฺมสุตฺต}
\csromanmatikaanchor{mtk.212}
\csromantocmarkref{subsection}{7. ทุกฺขธมฺมสุตฺต}{mtk.212}
\bookmark[dest={mtk.212},level=2]{7. ทุกฺขธมฺมสุตฺต}
\proseitem{244}{ยโต โข ภิกฺขเว ภิกฺขุ สพฺเพสํเยว ทุกฺขธมฺมานํ สมุทยญฺจ อตฺถงฺคมญฺจ ยถาภูตํ ปชานาติ, ตถา โข ปนสฺส กามา ทิฏฺฐา โหนฺติ, ยถาสฺส กาเม ปสฺสโต โย กาเมสุ กามจฺฉนฺโท กามสฺเนโห กามมุจฺฉา กามปริฬาโห, โส นานุเสติ.\csromanspacer{} ตถา โข ปนสฺส จาโร จ วิหาโร จ อนุพุทฺโธ โหติ, ยถา จรนฺตํ วิหรนฺตํ อภิชฺฌาโทมนสฺสา ปาปกา อกุสลา ธมฺมา นานุเสนฺติ.}

\prose{กถญฺจ ภิกฺขเว สพฺเพสํเยว ทุกฺขธมฺมานํ สมุทยญฺจ อตฺถงฺคมญฺจ ยถาภูตํ ปชานาติ, “อิติ รูปํ, อิติ รูปสฺส สมุทโย, อิติ รูปสฺส อตฺถงฺคโม.\csromanspacer{} อิติ เวทนา ฯเปฯ อิติ สญฺญา.\csromanspacer{} อิติ สงฺขารา.\csromanspacer{} อิติ วิญฺญาณํ, อิติ วิญฺญาณสฺส สมุทโย, อิติ วิญฺญาณสฺส อตฺถงฺคโม”ติ.\csromanspacer{} เอวํ โข ภิกฺขเว ภิกฺขุ สพฺเพสํเยว ทุกฺขธมฺมานํ สมุทยญฺจ อตฺถงฺคมญฺจ ยถาภูตํ ปชานาติ.}

\prose{กถญฺจ ภิกฺขเว ภิกฺขุโน กามา ทิฏฺฐา โหนฺติ.\csromanspacer{} ยถาสฺส กาเม ปสฺสโต โย กาเมสุ กามจฺฉนฺโท กามสฺเนโห กามมุจฺฉา กามปริฬาโห, โส นานุเสติ.\csromanspacer{} เสยฺยถาปิ ภิกฺขเว องฺคารกาสุ สาธิกโปริสา ปุณฺณา องฺคารานํ วีตจฺจิกานํ วีตธูมานํ.\csromanspacer{} อถ ปุริโส อาคจฺเฉยฺย ชีวิตุกาโม อมริตุกาโม สุขกาโม ทุกฺขปฏิกูโล.\csromanspacer{} ตเมนํ เทฺว พลวนฺโต ปุริสา นานาพาหาสุ คเหตฺวา ตํ องฺคารกาสุํ อุปกฑฺเฒยฺยุํ, โส อิติจีติเจว กายํ สนฺนาเมยฺย.\csromanspacer{} ตํ กิสฺส เหตุ, ญาตํ\footnote{ญาณํ (ก)} หิ ภิกฺขเว ตสฺส ปุริสสฺส\footnote{ปุริสสฺส โหติ (สี, สฺยา, กํ, อิ), ปุริสสฺส เหตุ โหติ (ก) ม. 2 (28) ปิฏฺเฐปิ.} “อิมํ จาหํ องฺคารกาสุํ ปปติสฺสามิ, ตโตนิทานํ มรณํ วา นิคจฺฉิสฺสามิ มรณมตฺตํ วา ทุกฺขนฺ”ติ.\csromanspacer{} เอวเมว}

\gambhira{สฬายตนวคฺคสํยุตฺตปาฬิ}
\prose{\csromanfolio{394}โข ภิกฺขเว ภิกฺขุโน องฺคารกาสูปมา กามา ทิฏฺฐา โหนฺติ, ยถาสฺส กาเม ปสฺสโต โย กาเมสุ กามจฺฉนฺโท กามสฺเนโห กามมุจฺฉา กามปริฬาโห, โส นานุเสติ.}

\prose{กถญฺจ ภิกฺขเว ภิกฺขุโน จาโร จ วิหาโร จ อนุพุทฺโธ โหติ, ยถา จรนฺตํ วิหรนฺตํ อภิชฺฌาโทมนสฺสา ปาปกา อกุสลา ธมฺมา นานุสฺสวนฺติ\footnote{นานุเสนฺติ (ก)}.\csromanspacer{} เสยฺยถาปิ ภิกฺขเว ปุริโส พหุกณฺฏกํ ทายํ ปวิเสยฺย, ตสฺส ปุรโตปิ กณฺฏโก ปจฺฉโตปิ กณฺฏโก อุตฺตรโตปิ กณฺฏโก ทกฺขิณโตปิ กณฺฏโก หฏฺฐโตปิ กณฺฏโก อุปริโตปิ กณฺฏโก, โส สโตว อภิกฺกเมยฺย, สโตว ปฏิกฺกเมยฺย “มามํ กณฺฏโก”ติ.\csromanspacer{} เอวเมว โข ภิกฺขเว ยํ โลเก ปิยรูปํ สาตรูปํ, อยํ วุจฺจติ อริยสฺส วินเย กณฺฏโกติ.\csromanspacer{} อิติ วิทิตฺวา\footnote{กณฺฏโก. ตํ กณฺฏโกติ อิติ วิทิตฺวา (สี)} สํวโร จ อสํวโร จ เวทิตพฺโพ.}

\prose{กถญฺจ ภิกฺขเว อสํวโร โหติ, อิธ ภิกฺขเว ภิกฺขุ จกฺขุนา รูปํ ทิสฺวา ปิยรูเป รูเป อธิมุจฺจติ, อปฺปิยรูเป รูเป พฺยาปชฺชติ, อนุปฏฺฐิตกายสฺสติ จ วิหรติ ปริตฺตเจตโส.\csromanspacer{} ตญฺจ เจโตวิมุตฺติํ ปญฺญาวิมุตฺติํ ยถาภูตํ นปฺปชานาติ, ยตฺถสฺส เต อุปฺปนฺนา ปาปกา อกุสลา ธมฺมา อปริเสสา นิรุชฺฌนฺติ ฯเปฯ.\csromanspacer{} ชิวฺหาย รสํ สายิตฺวา ฯเปฯ.\csromanspacer{} มนสา ธมฺมํ วิญฺญาย ปิยรูเป ธมฺเม อธิมุจฺจติ, อปฺปิยรูเป ธมฺเม พฺยาปชฺชติ, อนุปฏฺฐิตกายสฺสติ จ วิหรติ ปริตฺตเจตโส.\csromanspacer{} ตญฺจ เจโตวิมุตฺติํ ปญฺญาวิมุตฺติํ ยถาภูตํ นปฺปชานาติ, ยตฺถสฺส เต อุปฺปนฺนา ปาปกา อกุสลา ธมฺมา อปริเสสา นิรุชฺฌนฺติ.\csromanspacer{} เอวํ โข ภิกฺขเว อสํวโร โหติ.}

\prose{กถญฺจ ภิกฺขเว สํวโร โหติ, อิธ ภิกฺขเว ภิกฺขุ จกฺขุนา รูปํ ทิสฺวา ปิยรูเป รูเป นาธิมุจฺจติ, อปฺปิยรูเป รูเป น พฺยาปชฺชติ, อุปฏฺฐิตกายสฺสติ จ วิหรติ อปฺปมาณเจตโส.\csromanspacer{} ตญฺจ เจโตวิมุตฺติํ ปญฺญาวิมุตฺติํ ยถาภูตํ ปชานาติ, ยตฺถสฺส เต อุปฺปนฺนา ปาปกา อกุสลา ธมฺมา อปริเสสา นิรุชฺฌนฺติ ฯเปฯ.\csromanspacer{} ชิวฺหาย รสํ สายิตฺวา ฯเปฯ.\csromanspacer{} มนสา ธมฺมํ วิญฺญาย ปิยรูเป ธมฺเม นาธิมุจฺจติ, อปฺปิยรูเป ธมฺเม น พฺยาปชฺชติ, อุปฏฺฐิตกายสฺสติ จ วิหรติ อปฺปมาณเจตโส.\csromanspacer{} ตญฺจ เจโตวิมุตฺติํ ปญฺญาวิมุตฺติํ ยถาภูตํ ปชานาติ, ยตฺถสฺส เต อุปฺปนฺนา}

\vspace{\tipitakaparskip}
\csromancenter{สฬายตนสํยุตฺต ปาปกา อกุสลา ธมฺมา อปริเสสา นิรุชฺฌนฺติ.\csromanspacer{} เอวํ โข ภิกฺขเว สํวโร โหติ.}
\prose{\csromanfolio{395}ตสฺส เจ ภิกฺขเว ภิกฺขุโน เอวํ จรโต เอวํ วิหรโต กทาจิ กรหจิ สติสมฺโมสา อุปฺปชฺชนฺติ ปาปกา อกุสลา สรสงฺกปฺปา สํโยชนิยา.\csromanspacer{} ทนฺโธ ภิกฺขเว สตุปฺปาโท, อถ โข นํ ขิปฺปเมว ปชหติ วิโนเทติ พฺยนฺตีกโรติ อนภาวํ คเมติ.}

\prose{เสยฺยถาปิ ภิกฺขเว ปุริโส ทิวสํสนฺตตฺเต\footnote{ทิวสสนฺตตฺเต (สี)} อโยกฏาเห เทฺว วา ตีณิ วา อุทกผุสิตานิ นิปาเตยฺย.\csromanspacer{} ทนฺโธ ภิกฺขเว อุทกผุสิตานํ นิปาโต, อถ โข นํ ขิปฺปเมว ปริกฺขยํ ปริยาทานํ คจฺเฉยฺย.\csromanspacer{} เอวเมว โข ภิกฺขเว ตสฺส เจ ภิกฺขุโน เอวํ จรโต เอวํ วิหรโต กทาจิ กรหจิ สติสมฺโมสา อุปฺปชฺชนฺติ ปาปกา อกุสลา สรสงฺกปฺปา สํโยชนิยา.\csromanspacer{} ทนฺโธ ภิกฺขเว สตุปฺปาโท, อถ โข นํ ขิปฺปเมว ปชหติ วิโนเทติ พฺยนฺตีกโรติ อนภาวํ คเมติ.\csromanspacer{} เอวํ โข ภิกฺขเว ภิกฺขุโน จาโร จ วิหาโร จ อนุพุทฺโธ โหติ, ยถา จรนฺตํ วิหรนฺตํ อภิชฺฌาโทมนสฺสา ปาปกา อกุสลา ธมฺมา นานุสฺสวนฺติ.\csromanspacer{} ตญฺเจ ภิกฺขเว ภิกฺขุํ เอวํ จรนฺตํ เอวํ วิหรนฺตํ ราชาโน วา ราชมหามตฺตา วา มิตฺตา วา อมจฺจา วา ญาตี วา สาโลหิตา วา โภเคหิ อภิหฏฺฐุํ ปวาเรยฺยุํ “เอหิ\footnote{เอวํ (สี)} โภ ปุริส กิํ เต อิเม กาสาวา อนุทหนฺติ, กิํ มุณฺโฑ กปาลมนุจรสิ, เอหิ หีนายาวตฺติตฺวา โภเค จ ภุญฺชสฺสุ, ปุญฺญานิ จ กโรหี”ติ.\csromanspacer{} โส วต ภิกฺขเว ภิกฺขุ เอวํ จรนฺโต เอวํ วิหรนฺโต สิกฺขํ ปจฺจกฺขาย หีนายาวตฺติสฺสตีติ เนตํ ฐานํ วิชฺชติ.}

\prose{เสยฺยถาปิ ภิกฺขเว คงฺคา นที ปาจีนนินฺนา ปาจีนโปณา ปาจีนปพฺภารา.\csromanspacer{} อถ มหาชนกาโย อาคจฺเฉยฺย กุทฺทาลปิฏกํ อาทาย “มยํ อิมํ คงฺคํ นทิํ ปจฺฉานินฺนํ กริสฺสาม ปจฺฉาโปณํ ปจฺฉาปพฺภารนฺ”ติ.\csromanspacer{} ตํ กิํ มญฺญถ ภิกฺขเว อปิ นุ โข โส มหาชนกาโย คงฺคํ นทิํ ปจฺฉานินฺนํ กเรยฺย ปจฺฉาโปณํ ปจฺฉาปพฺภารนฺติ.\csromanspacer{} โน เหตํ ภนฺเต.\csromanspacer{} ตํ กิสฺส เหตุ, คงฺคา ภนฺเต นที ปาจีนนินฺนา ปาจีนโปณา ปาจีนปพฺภารา, สา น สุกรา ปจฺฉานินฺนา กาตุํ ปจฺฉาโปณา ปจฺฉาปพฺภารา.\csromanspacer{} ยาวเทว จ ปน โส มหาชนกาโย กิลมถสฺส วิฆาตสฺส ภาคี อสฺสาติ.\csromanspacer{} เอวเมว โข ภิกฺขเว ตญฺเจ ภิกฺขุํ เอวํ จรนฺตํ เอวํ วิหรนฺตํ ราชาโน วา}

\gambhira{สฬายตนวคฺคสํยุตฺตปาฬิ}
\prose{\csromanfolio{396}ราชมหามตฺตา วา มิตฺตา วา อมจฺจา วา ญาตี วา สาโลหิตา วา โภเคหิ อภิหฏฺฐุํ ปวาเรยฺยุํ “เอหิ โภ ปุริส กิํ เต อิเม กาสาวา อนุทหนฺติ, กิํ มุณฺโฑ กปาลมนุจรสิ, เอหิ หีนายาวตฺติตฺวา โภเค จ ภุญฺชสฺสุ, ปุญฺญานิ จ กโรหี”ติ.\csromanspacer{} โส วต ภิกฺขเว ภิกฺขุ เอวํ จรนฺโต เอวํ วิหรนฺโต สิกฺขํ ปจฺจกฺขาย หีนายาวตฺติสฺสตีติ เนตํ ฐานํ วิชฺชติ.\csromanspacer{} ตํ กิสฺส เหตุ, ยํ หิ ตํ ภิกฺขเว จิตฺตํ ทีฆรตฺตํ วิเวกนินฺนํ วิเวกโปณํ วิเวกปพฺภารํ, ตถา\footnote{ตญฺจ (สฺยา, กํ, ก)} หีนายาวตฺติสฺสตีติ เนตํ ฐานํ วิชฺชตีติ.\csromanspacer{}\csromanspacer{}สตฺตมํ.}

\csromanheader{2}{8. กิํสุโกปมสุตฺต}
\csromanmatikaanchor{mtk.213}
\csromantocmarkref{subsection}{8. กิํสุโกปมสุตฺต}{mtk.213}
\bookmark[dest={mtk.213},level=2]{8. กิํสุโกปมสุตฺต}
\proseitem{245}{อถ โข อญฺญตโร ภิกฺขุ เยนญฺญตโร ภิกฺขุ เตนุปสงฺกมิ, อุปสงฺกมิตฺวา ตํ ภิกฺขุํ เอตทโวจ “กิตฺตาวตา นุ โข อาวุโส ภิกฺขุโน ทสฺสนํ สุวิสุทฺธํ โหตี”ติ.\csromanspacer{} ยโต โข อาวุโส ภิกฺขุ ฉนฺนํ ผสฺสายตนานํ สมุทยญฺจ อตฺถงฺคมญฺจ ยถาภูตํ ปชานาติ, เอตฺตาวตา โข อาวุโส ภิกฺขุโน ทสฺสนํ สุวิสุทฺธํ โหตีติ.}

\prose{อถ โข โส ภิกฺขุ อสนฺตุฏฺโฐ ตสฺส ภิกฺขุสฺส ปญฺหเวยฺยากรเณน\footnote{ปญฺหาเวยฺยากรเณน (สฺยา, กํ, ก)} เยนญฺญตโร ภิกฺขุ เตนุปสงฺกมิ, อุปสงฺกมิตฺวา ตํ ภิกฺขุํ เอตทโวจ “กิตฺตาวตา นุ โข อาวุโส ภิกฺขุโน ทสฺสนํ สุวิสุทฺธํ โหตี”ติ.\csromanspacer{} ยโต โข อาวุโส ภิกฺขุ ปญฺจนฺนํ อุปาทานกฺขนฺธานํ สมุทยญฺจ อตฺถงฺคมญฺจ ยถาภูตํ ปชานาติ, เอตฺตาวตา โข อาวุโส ภิกฺขุโน ทสฺสนํ สุวิสุทฺธํ โหตีติ.}

\prose{อถ โข โส ภิกฺขุ อสนฺตุฏฺโฐ ตสฺส ภิกฺขุสฺส ปญฺหเวยฺยากรเณน เยนญฺญตโร ภิกฺขุ เตนุปสงฺกมิ, อุปสงฺกมิตฺวา ตํ ภิกฺขุํ เอตทโวจ “กิตฺตาวตา นุ โข อาวุโส ภิกฺขุโน ทสฺสนํ สุวิสุทฺธํ โหตี”ติ.\csromanspacer{} ยโต โข อาวุโส ภิกฺขุ จตุนฺนํ มหาภูตานํ สมุทยญฺจ อตฺถงฺคมญฺจ ยถาภูตํ ปชานาติ.\csromanspacer{} เอตฺตาวตา โข อาวุโส ภิกฺขุโน ทสฺสนํ สุวิสุทฺธํ โหตีติ.}

\prose{อถ โข โส ภิกฺขุ อสนฺตุฏฺโฐ ตสฺส ภิกฺขุสฺส ปญฺหเวยฺยากรเณน เยนญฺญตโร ภิกฺขุ เตนุปสงฺกมิ, อุปสงฺกมิตฺวา ตํ ภิกฺขุํ เอตทโวจ}

\vspace{\tipitakaparskip}
\csromancenter{สฬายตนสํยุตฺต “กิตฺตาวตา นุ โข อาวุโส ภิกฺขุโน ทสฺสนํ สุวิสุทฺธํ โหตี”ติ.\csromanspacer{} ยโต โข อาวุโส ภิกฺขุ “ยํ กิญฺจิ สมุทยธมฺมํ, สพฺพํ ตํ นิโรธธมฺมนฺ”ติ ยถาภูตํ ปชานาติ.\csromanspacer{} เอตฺตาวตา โข อาวุโส ภิกฺขุโน ทสฺสนํ สุวิสุทฺธํ โหตีติ.}
\prose{\csromanfolio{397}อถ โข โส ภิกฺขุ อสนฺตุฏฺโฐ ตสฺส ภิกฺขุสฺส ปญฺหเวยฺยากรเณน เยน ภควา เตนุปสงฺกมิ, อุปสงฺกมิตฺวา ภควนฺตํ เอตทโวจ “อิธาหํ ภนฺเต เยนญฺญตโร ภิกฺขุ เตนุปสงฺกมิํ, อุปสงฺกมิตฺวา ตํ ภิกฺขุํ เอตทโวจํ ‘กิตฺตาวตา นุ โข อาวุโส ภิกฺขุโน ทสฺสนํ สุวิสุทฺธํ โหตี’ติ.\csromanspacer{} เอวํ วุตฺเต ภนฺเต โส ภิกฺขุ มํ เอตทโวจ ‘ยโต โข อาวุโส ภิกฺขุ ฉนฺนํ ผสฺสายตนานํ สมุทยญฺจ อตฺถงฺคมญฺจ ยถาภูตํ ปชานาติ, เอตฺตาวตา โข อาวุโส ภิกฺขุโน ทสฺสนํ สุวิสุทฺธํ โหตี’ติ.\csromanspacer{} อถ ขฺวาหํ ภนฺเต อสนฺตุฏฺโฐ ตสฺส ภิกฺขุสฺส ปญฺหเวยฺยากรเณน เยนญฺญตโร ภิกฺขุ เตนุปสงฺกมิํ, อุปสงฺกมิตฺวา ตํ ภิกฺขุํ เอตทโวจํ ‘กิตฺตาวตา นุ โข อาวุโส ภิกฺขุโน ทสฺสนํ สุวิสุทฺธํ โหตี’ติ.\csromanspacer{} เอวํ วุตฺเต ภนฺเต โส ภิกฺขุ มํ เอตทโวจ ‘ยโต โข อาวุโส ภิกฺขุ ปญฺจนฺนํ อุปาทานกฺขนฺธานํ ฯเปฯ จตุนฺนํ มหาภูตํ สมุทยญฺจ อตฺถงฺคมญฺจ ยถาภูตํ ปชานาติ ฯเปฯ ‘ยํ กิญฺจิ สมุทยธมฺมํ, สพฺพํ ตํ นิโรธธมฺมนฺ’ติ ยถาภูตํ ปชานาติ.\csromanspacer{} เอตฺตาวตา โข อาวุโส ภิกฺขุโน ทสฺสนํ สุวิสุทฺธํ โหตี’ติ.\csromanspacer{} อถ ขฺวาหํ ภนฺเต อสนฺตุฏฺโฐ ตสฺส ภิกฺขุสฺส ปญฺหเวยฺยากรเณน เยน ภควา เตนุปสงฺกมิํ ( )\footnote{(อุปสงฺกมิตฺวา ภควนฺตํ เอตทโวจํ) (ก)}, กิตฺตาวตา นุ โข ภนฺเต ภิกฺขุโน ทสฺสนํ สุวิสุทฺธํ โหตี”ติ.}

\prose{เสยฺยถาปิ ภิกฺขุ ปุริสสฺส กิํสุโก อทิฏฺฐปุพฺโพ อสฺส, โส เยนญฺญตโร ปุริโส กิํสุกสฺส ทสฺสาวี เตนุปสงฺกเมยฺย, อุปสงฺกมิตฺวา ตํ ปุริสํ เอวํ วเทยฺย “กีทิโส โภ ปุริส กิํสุโก”ติ, โส เอวํ วเทยฺย “กาฬโก โข อมฺโภ ปุริส กิํสุโก, เสยฺยถาปิ ฌามขาณู”ติ.\csromanspacer{} เตน โข ปน ภิกฺขุ สมเยน ตาทิโสวสฺส กิํสุโก, ยถาปิ\footnote{ยถา (สี, สฺยา, กํ) ทุติยวาราทีสุ ปน “ยถาปิ”เตฺวว ทิสฺสติ,} ตสฺส ปุริสสฺส ทสฺสนํ.\csromanspacer{} อถ โข โส ภิกฺขุ ปุริโส อสนฺตุฏฺฐา ตสฺส ปุริสสฺส ปญฺหเวยฺยากรเณน เยนญฺญตโร ปุริโส กิํสุกสฺส ทสฺสาวี เตนุปสงฺกเมยฺย, อุปสงฺกมิตฺวา ตํ ปุริสํ เอวํ วเทยฺย}

\gambhira{สฬายตนวคฺคสํยุตฺตปาฬิ}
\prose{\csromanfolio{398}“กีทิโส โภ ปุริส กิํสุโก”ติ.\csromanspacer{} โส เอวํ วเทยฺย “โลหิตโก โข อมฺโภ ปุริส กิํสุโก, เสยฺยถาปิ มํสเปสีติ.\csromanspacer{} เตน โข ปน ภิกฺขุ สมเยน ตาทิโสวสฺส กิํสุโก, ยถาปิ ตสฺส ปุริสสฺส ทสฺสนํ.\csromanspacer{} อถ โข โส ภิกฺขุ ปุริโส อสนฺตุฏฺโฐ ตสฺส ปุริสสฺส ปญฺหเวยฺยากรเณน เยนญฺญตโร ปุริโส กิํสุกสฺส ทสฺสาวี เตนุปสงฺกเมยฺย, อุปสงฺกมิตฺวา ตํ ปุริสํ เอวํ วเทยฺย “กีทิโส โภ ปุริส กิํสุโก”ติ.\csromanspacer{} โส เอวํ วเทยฺย “โอจีรกชาโต\footnote{โอชีรกชาโต (สี), โอทีรกชาโต (อิ)} โข อมฺโภ ปุริส กิํสุโก อาทินฺนสิปาฏิโก, เสยฺยถาปิ สิรีโส”ติ.\csromanspacer{} เตน โข ปน ภิกฺขุ สมเยน ตาทิโสวสฺส กิํสุโก, ยถาปิ ตสฺส ปุริสสฺส ทสฺสนํ.\csromanspacer{} อถ โข โส ภิกฺขุ ปุริโส อสนฺตุฏฺโฐ ตสฺส ปุริสสฺส ปญฺหเวยฺยากรเณน เยนญฺญตโร ปุริโส กิํสุกสฺส ทสฺสาวี เตนุปสงฺกเมยฺย, อุปสงฺกมิตฺวา ตํ ปุริสํ เอวํ วเทยฺย “กีทิโส โภ ปุริส กิํสุโกติ.\csromanspacer{} โส เอวํ วเทยฺย “พหลปตฺตปลาโส สนฺทจฺฉาโย\footnote{สณฺฑจฺฉาโย (สฺยา, กํ)} โข อมฺโภ ปุริส กิํสุโก, เสยฺยถาปิ นิโคฺรโธ”ติ.\csromanspacer{} เตน โข ปน ภิกฺขุ สมเยน ตาทิโสวสฺส กิํสุโก, ยถาปิ ตสฺส ปุริสสฺส ทสฺสนํ.\csromanspacer{} เอวเมว โข ภิกฺขุ ยถา ยถา อธิมุตฺตานํ เตสํ สปฺปุริสานํ ทสฺสนํ สุวิสุทฺธํ โหติ, ตถา ตถา โข เตหิ สปฺปุริเสหิ พฺยากตํ.}

\prose{เสยฺยถาปิ ภิกฺขุ รญฺโญ ปจฺจนฺติมํ นครํ ทฬฺหุทฺธาปํ\footnote{ทฬฺหุทฺทาปํ (สี, อิ)} ทฬฺหปาการโตรณํ ฉทฺวารํ, ตตฺรสฺส โทวาริโก ปณฺฑิโต พฺยตฺโต เมธาวิ อญฺญาตานํ นิวาเรตา, ญาตานํ ปเวเสตา.\csromanspacer{} ปุรตฺถิมาย ทิสาย อาคนฺตฺวา สีฆํ ทูตยุคํ ตํ โทวาริกํ เอวํ วเทยฺย “กหํ โภ ปุริส อิมสฺส นครสฺส นครสฺสามี”ติ.\csromanspacer{} โส เอวํ วเทยฺย “เอโส ภนฺเต มชฺเฌ สิงฺฆาฏเก นิสินฺโน”ติ.\csromanspacer{} อถ โข ตํ สีฆํ ทูตยุคํ นครสฺสามิกสฺส ยถาภูตํ วจนํ นิยฺยาเตตฺวา ยถาคตมคฺคํ ปฏิปชฺเชยฺย.\csromanspacer{} ปจฺฉิมาย ทิสาย อาคนฺตฺวา สีฆํ ทูตยุคํ ฯเปฯ.\csromanspacer{} อุตฺตราย ทิสาย.\csromanspacer{} ทกฺขิณาย ทิสาย อาคนฺตฺวา สีฆํ ทูตยุคํ ตํ โทวาริกํ เอวํ วเทยฺย “กหํ โภ ปุริส อิมสฺส นครสฺสามี”ติ.\csromanspacer{} โส เอวํ วเทยฺย “เอโส ภนฺเต มชฺเฌ สิงฺฆาฏเก นิสินฺโน”ติ.\csromanspacer{} อถ โข ตํ สีฆํ ทูตยุคํ นครสฺสามิกสฺส ยถาภูตํ วจนํ นิยฺยาเตตฺวา ยถาคตมคฺคํ ปฏิปชฺเชยฺย.}

\csromanheader{2}{1. สฬายตนสํยุตฺต}
\prose{\csromanfolio{399}อุปมา โข มฺยายํ ภิกฺขุ กตา อตฺถสฺส วิญฺญาปนาย.\csromanspacer{} อยญฺเจตฺถ อตฺโถ, “นครนฺ”ติ โข ภิกฺขุ อิมสฺเสตํ จาตุมหาภูติกสฺส กายสฺส อธิวจนํ มาตาเปตฺติกสมฺภวสฺส โอทนกุมฺมาสูปจยสฺส อนิจฺจุจฺฉาทนปริมทฺทนเภทน วิทฺธํสนธมฺมสฺส. “ฉ ทฺวารา”ติ โข ภิกฺขุ ฉนฺเนตํ อชฺฌตฺติกานํ อายตนานํ อธิวจนํ. “โทวาริโก”ติ โข ภิกฺขุ สติยา เอตํ อธิวจนํ. “สีฆํ ทูตยุคนฺ”ติ โข ภิกฺขุ สมถวิปสฺสนาเนตํ อธิวจนํ. “นครสฺสามี”ติ โข ภิกฺขุ วิญฺญาณสฺเสตํ อธิวจนํ. “มชฺเฌ สิงฺฆาฏโก”ติ โข ภิกฺขุ จตุนฺเนตํ มหาภูตานํ อธิวจนํ ปถวีธาตุยา อาโปธาตุยา เตโชธาตุยา วาโยธาตุยา. “ยถาภูตํ วจนนฺ”ติ โข ภิกฺขุ นิพฺพานสฺเสตํ อธิวจนํ. “ยถาคตมคฺโค”ติ โข ภิกฺขุ อริยสฺเสตํ อฏฺฐงฺคิกสฺส มคฺคสฺส อธิวจนํ.\csromanspacer{} เสยฺยถิทํ, สมฺมาทิฏฺฐิยา ฯเปฯ สมฺมาสมาธิสฺสาติ.\csromanspacer{}\csromanspacer{}อฏฺฐมํ.}

\csromanheader{2}{9. วีโณปมสุตฺต}
\csromanmatikaanchor{mtk.214}
\csromantocmarkref{subsection}{9. วีโณปมสุตฺต}{mtk.214}
\bookmark[dest={mtk.214},level=2]{9. วีโณปมสุตฺต}
\proseitem{246}{ยสฺส กสฺสจิ ภิกฺขเว ภิกฺขุสฺส วา ภิกฺขุนิยา วา จกฺขุวิญฺเญยฺเยสุ รูเปสุ อุปฺปชฺเชยฺย ฉนฺโท วา ราโค วา โทโส วา โมโห วา ปฏิฆํ วาปิ\footnote{ปฏิฆํ วา (สี)} เจตโส, ตโต จิตฺตํ นิวาเรยฺย “สภโย เจโส มคฺโค สปฺปฏิภโย จ สกณฺฏโก จ สคหโน จ อุมฺมคฺโค จ กุมฺมคฺโค จ ทุหิติโก จ, อสปฺปุริสเสวิโต เจโต มคฺโค, น เจโส มคฺโค สปฺปุริเสหิ เสวิโต, น ตฺวํ เอตํ อรหสี”ติ ตโต จิตฺตํ นิวารเย จกฺขุวิญฺเญยฺเยหิ รูเปหิ ฯเปฯ.\csromanspacer{} ยสฺส กสฺสจิ ภิกฺขเว ภิกฺขุสฺส วา ภิกฺขุนิยา วา ชิวฺหาวิญฺเญยฺเยสุ รเสสุ ฯเปฯ มโนวิญฺเญยฺเยสุ ธมฺเมสุ อุปฺปชฺเชยฺย ฉนฺโท วา ราโค วา โทโส วา โมโห วา ปฏิฆํ วาปิ เจตโส, ตโต จิตฺตํ นิวาเรยฺย “สภโย เจโส มคฺโค สปฺปฏิภโย จ สกณฺฏโก จ สคหโน จ อุมฺมคฺโค จ กุมฺมคฺโค จ ทุหิติโก จ, อสปฺปุริสเสวิโต เจโส มคฺโค, น เจโส มคฺโค สปฺปุริเสหิ เสวิโต, น ตฺวํ เอตํ อรหสี”ติ ตโต จิตฺตํ นิวารเย มโนวิญฺเญยฺเยหิ ธมฺเมหิ.}

\prose{เสยฺยถาปิ ภิกฺขเว กิฏฺฐํ สมฺปนฺนํ, กิฏฺฐารกฺโข\footnote{กิฏฺฐารกฺขโก (สี)} จ ปมตฺโต, โคโณ จ กิฏฺฐาโท, อทุํ กิฏฺฐํ โอตริตฺวา ยาวทตฺถํ มทํ อาปชฺเชยฺย, ปมาทํ}

\gambhira{สฬายตนวคฺคสํยุตฺตปาฬิ}
\prose{\csromanfolio{400}อาปชฺเชยฺย.\csromanspacer{} เอวเมว โข ภิกฺขเว อสฺสุตวา ปุถุชฺชโน ฉสุ ผสฺสายตเนสุ อสํวุตการี ปญฺจสุ กามคุเณสุ ยาวทตฺถํ มทํ อาปชฺชติ, ปมาทํ อาปชฺชติ.}

\prose{เสยฺยถาปิ ภิกฺขเว กิฏฺฐํ สมฺปนฺนํ, กิฏฺฐารกฺโข จ อปฺปมตฺโต, โคโณ จ กิฏฺฐาโท อทุํ กิฏฺฐํ โอตเรยฺย, ตเมนํ กิฏฺฐารกฺโข นาสายํ สุคฺคหิตํ คเณฺหยฺย, นาสายํ สุคฺคหิตํ คเหตฺวา อุปริฆฏายํ สุนิคฺคหิตํ นิคฺคเณฺหยฺย, อุปริฆฏายํ สุนิคฺคหิตํ นิคฺคเหตฺวา ทณฺเฑน สุตาฬิตํ ตาเฬยฺย, ทณฺเฑน สุตาฬิตํ ตาเฬตฺวา โอสชฺเชยฺย.\csromanspacer{} ทุติยมฺปิ โข ภิกฺขเว ฯเปฯ.\csromanspacer{} ตติยมฺปิ โข ภิกฺขเว โคโณ กิฏฺฐาโท อทุํ กิฏฺฐํ โอตเรยฺย, ตเมนํ กิฏฺฐารกฺโข นาสายํ สุคฺคหิตํ คเณฺหยฺย, นาสายํ สุคฺคหิตํ คเหตฺวา อุปริฆฏายํ สุนิคฺคหิตํ นิคฺคเณฺหยฺย, อุปริฆฏายํ สุนิคฺคหิตํ นิคฺคเหตฺวา ทณฺเฑน สุตาฬิตํ ตาเฬยฺย, ทณฺเฑน สุตาฬิตํ ตาเฬตฺวา โอสชฺเชยฺย.\csromanspacer{} เอวํ หิ โส ภิกฺขเว โคโณ กิฏฺฐาโท คามคโต วา อรญฺญคโต วา ฐานพหุโล วา อสฺส นิสชฺชพหุโล วา, น ตํ กิฏฺฐํ ปุน โอตเรยฺย ตเมว ปุริมํ ทณฺฑสมฺผสฺสํ สมนุสฺสรนฺโต.\csromanspacer{} เอวเมว โข ภิกฺขเว ยโต โข ภิกฺขุโน ฉสุ ผสฺสายตเนสุ จิตฺตํ อุทุชิตํ โหติ สุทุชิตํ, อชฺฌตฺตเมว สนฺติฏฺฐติ สนฺนิสีทติ, เอโกทิ โหติ สมาธิยติ.}

\prose{เสยฺยถาปิ ภิกฺขเว รญฺโญ วา ราชมหามตฺตสฺส วา วีณาย สทฺโท อสฺสุตปุพฺโพ อสฺส, โส วีณาสทฺทํ สุเณยฺย, โส เอวํ วเทยฺย “อมฺโภ กสฺส\footnote{กิสฺส (สี, อิ)} นุ โข เอโส สทฺโท เอวํรชนีโย เอวํกมนีโย เอวํมทนีโย เอวํมุจฺฉนีโย เอวํพนฺธนีโย”ติ.\csromanspacer{} ตเมนํ เอวํ วเทยฺยุํ “เอสา โข ภนฺเต วีณา นาม, ยสฺสา เอโส สทฺโท เอวํรชนีโย เอวํกมนีโย เอวํมทนีโย เอวํมุจฺฉนีโย เอวํพนฺธนีโย”ติ.\csromanspacer{} โส เอวํ วเทยฺย “คจฺฉถ เม โภ ตํ วีณํ อาหรถา”ติ, ตสฺส ตํ วีณํ อาหเรยฺยุํ.\csromanspacer{} ตเมนํ เอวํ วเทยฺยุํ “อยํ โข สา ภนฺเต วีณา, ยสฺสา เอโส สทฺโท เอวํรชนีโย เอวํกมนีโย เอวํมทนีโย เอวํมุจฺฉนีโย เอวํพนฺธนีโย”ติ.\csromanspacer{} โส เอวํ วเทยฺย “อลํ เม โภ ตาย วีณาย, ตเมว เม สทฺทํ อาหรถา”ติ.\csromanspacer{} ตเมนํ เอวํ วเทยฺยุํ “อยํ โข ภนฺเต วีณา นาม อเนกสมฺภารา มหาสมฺภารา, อเนเกหิ สมฺภาเรหิ สมารทฺธา}

\vspace{\tipitakaparskip}
\csromancenter{สฬายตนสํยุตฺต วทติ.\csromanspacer{} เสยฺยถิทํ, โทณิญฺจ ปฏิจฺจ จมฺมญฺจ ทณฺฑญฺจ ปฏิจฺจ อุปธารเณ จ ปฏิจฺจ ตนฺติโย จ ปฏิจฺจ โกณญฺจ ปฏิจฺจ ปุริสสฺส จ ตชฺชํ วายามํ ปฏิจฺจ, เอวายํ ภนฺเต วีณา นาม อเนกสมฺภารา มหาสมฺภารา อเนเกหิ สมฺภาเรหิ สมารทฺธา วทตี”ติ.\csromanspacer{} โส ตํ วีณํ ทสธา วา สตธา วา ผาเลยฺย, ทสธา วา สตธา วา ตํ ผาเลตฺวา สกลิกํ สกลิกํ กเรยฺย, สกลิกํ สกลิกํ กริตฺวา อคฺคินา ฑเหยฺย, อคฺคินา ฑหิตฺวา มสิํ กเรยฺย, มสิํ กริตฺวา มหาวาเต วา โอผุเนยฺย\footnote{โอปุเนยฺย (สี, อิ), โอผุเณยฺย (?)}, นทิยา วา สีฆโสตาย ปวาเหยฺย, โส เอวํ วเทยฺย “อสตี กิรายํ โภ วีณา นาม, ยเถวํ ยํ\footnote{ยเถวายํ (สี), ยเถวยํ (อิ)} กิญฺจิ วีณา นาม, เอตฺถ จ ปนายํชโน\footnote{เอตฺถ ปนายํ ชโน (สฺยา, กํ), เอตฺถ จ มหาชโน (อิ, ก)} อติเวลํ ปมตฺโต ปลฬิโต”ติ.\csromanspacer{} เอวเมว โข ภิกฺขเว ภิกฺขุ รูปํ สมเนฺวสติ\footnote{สมนฺเนสติ (สี, สฺยา, กํ), สมเนสติ (อิ)} ยาวตา รูปสฺส คติ, เวทนํ สมเนฺวสติ ยาวตา เวทนาย คติ, สญฺญํ สมเนฺวสติ ยาวตา สญฺญาย คติ, สงฺขาเร สมเนฺวสติ ยาวตา สงฺขารานํ คติ, วิญฺญาณํ สมเนฺวสติ ยาวตา วิญฺญาณสฺส คติ, ตสฺส รูปํ สมเนฺวสโต ยาวตา รูปสฺส คติ, เวทนํ สมเนฺวสโต ฯเปฯ.\csromanspacer{} สญฺญํ.\csromanspacer{} สงฺขาเร.\csromanspacer{} วิญฺญาณํ สมเนฺวสโต ยาวตา วิญฺญาณสฺส คติ.\csromanspacer{} ยมฺปิสฺส ตํ โหติ “อหนฺ”ติ วา “มมนฺ”ติ วา “อสฺมี”ติ วา, ตมฺปิ ตสฺส น โหตีติ.\csromanspacer{}\csromanspacer{}นวมํ.}
\csromanheader{2}{10. ฉปฺปาณโกปมสุตฺต}
\csromanmatikaanchor{mtk.215}
\csromantocmarkref{subsection}{10. ฉปฺปาณโกปมสุตฺต}{mtk.215}
\bookmark[dest={mtk.215},level=2]{10. ฉปฺปาณโกปมสุตฺต}
\proseitem{247}{\csromanfolio{401}เสยฺยถาปิ ภิกฺขเว ปุริโส อรุคตฺโต ปกฺกคตฺโต สรวนํ ปวิเสยฺย, ตสฺส กุสกณฺฏกา เจว ปาเท วิชฺเฌยฺยุํ, สรปตฺตานิ จ คตฺตานิ\footnote{สรปตฺตานิ ปกฺกคตฺตานิ (สฺยา, กํ), อรุปกฺกานิ คตฺตานิ (อิ, ก)} วิเลเขยฺยุํ.\csromanspacer{} เอวํ หิ โส ภิกฺขเว ปุริโส ภิยฺโยโสมตฺตาย ตโตนิทานํ ทุกฺขํ โทมนสฺสํ ปฏิสํเวทิเยถ.\csromanspacer{} เอวเมว โข ภิกฺขเว อิเธกจฺโจ ภิกฺขุ คามคโต วา อรญฺญคโต วา ลภติ วตฺตารํ “อยญฺจ โส\footnote{อยญฺจ โข (อิ, ก), อยํ โส (?)} อายสฺมา เอวํการี เอวํสมาจาโร อสุจิคามกณฺฏโก”ติ.\csromanspacer{} ตํ “กณฺฏโก”ติ\footnote{ตํ “อสุจิคามกณฺฏโก”ติ (ก)} อิติ วิทิตฺวา สํวโร จ อสํวโร จ เวทิตพฺโพ.}

\gambhira{สฬายตนวคฺคสํยุตฺตปาฬิ}
\prose{\csromanfolio{402}กถญฺจ ภิกฺขเว อสํวโร โหติ, อิธ ภิกฺขเว ภิกฺขุ จกฺขุนา รูปํ ทิสฺวา ปิยรูเป รูเป อธิมุจฺจติ, อปฺปิยรูเป รูเป พฺยาปชฺชติ, อนุปฏฺฐิตกายสฺสติ จ วิหรติ ปริตฺตเจตโส.\csromanspacer{} ตญฺจ เจโตวิมุตฺติํ ปญฺญาวิมุตฺติํ ยถาภูตํ นปฺปชานาติ, ยตฺถสฺส เต อุปฺปนฺนา ปาปกา อกุสลา ธมฺมา อปริเสสา นิรุชฺฌนฺติ.\csromanspacer{} โสเตน สทฺทํ สุตฺวา.\csromanspacer{} ฆาเนน คนฺธํ ฆายิตฺวา.\csromanspacer{} ชิวฺหาย รสํ สายิตฺวา.\csromanspacer{} กาเยน โผฏฺฐพฺพํ ผุสิตฺวา.\csromanspacer{} มนสา ธมฺมํ วิญฺญาย ปิยรูเป ธมฺเม อธิมุจฺจติ, อปฺปิยรูเป ธมฺเม พฺยาปชฺชติ, อนุปฏฺฐิตกายสฺสติ จ วิหรติ ปริตฺตเจตโส.\csromanspacer{} ตญฺจ เจโตวิมุตฺติํ ปญฺญาวิมุตฺติํ ยถาภูตํ นปฺปชานาติ, ยตฺถสฺส เต อุปฺปนฺนา ปาปกา อกุสลา ธมฺมา อปริเสสา นิรุชฺฌนฺติ.}

\prose{เสยฺยถาปิ ภิกฺขเว ปุริโส ฉปฺปาณเก คเหตฺวา นานาวิสเย นานาโคจเร ทฬฺหาย รชฺชุยา พนฺเธยฺย, อหิํ คเหตฺวา ทฬฺหาย รชฺชุยา พนฺเธยฺย.\csromanspacer{} สุสุมารํ\footnote{สุํสุมารํ (สี, สฺยา, กํ, อิ)} คเหตฺวา ทฬฺหาย รชฺชุยา พนฺเธยฺย.\csromanspacer{} ปกฺขิํ คเหตฺวา ทฬฺหาย รชฺชุยา พนฺเธยฺย.\csromanspacer{} กุกฺกุรํ คเหตฺวา ทฬฺหาย รชฺชุยา พนฺเธยฺย.\csromanspacer{} สิงฺคาลํ คเหตฺวา ทฬฺหาย รชฺชุยา พนฺเธยฺย.\csromanspacer{} มกฺกฏํ คเหตฺวา ทฬฺหาย รชฺชุยา พนฺเธยฺย, ทฬฺหาย รชฺชุยา พนฺธิตฺวา มชฺเฌ คณฺฐีํ กริตฺวา โอสฺสชฺเชยฺย.\csromanspacer{} อถ โข เต ภิกฺขเว ฉปฺปาณกา นานาวิสยา นานาโคจรา สกํ สกํ โคจรวิสยํ อาวิญฺเฉยฺยุํ\footnote{อาวิญฺเชยฺยุํ (สี)}, อหิ อาวิญฺเฉยฺย “วมฺมิกํ ปเวกฺขามี”ติ, สุสุมาโร อาวิญฺเฉยฺย “อุทกํ ปเวกฺขามี”ติ, ปกฺขี อาวิญฺเฉยฺย “อากาสํ เฑสฺสามี”ติ, กุกฺกุโร อาวิญฺเฉยฺย “คามํ ปเวกฺขามี”ติ, สิงฺคาโล อาวิญฺเฉยฺย “สีวถิกํ\footnote{สิวถิกํ (ก)} ปเวกฺขามี”ติ, มกฺกโฏ อาวิญฺเฉยฺย “วนํ ปเวกฺขามี”ติ.\csromanspacer{} ยทา โข เต ภิกฺขเว ฉปฺปาณกา ฌตฺตา อสฺสุ กิลนฺตา, อถ โข โย เนสํ ปาณกานํ พลวตโร อสฺส, ตสฺส เต อนุวตฺเตยฺยุํ อนุวิธาเยยฺยุํ วสํ คจฺเฉยฺยุํ.\csromanspacer{} เอวเมว โข ภิกฺขเว ยสฺส กสฺสจิ ภิกฺขุโน กายคตาสติ อภาวิตา อพหุลีกตา, ตํ จกฺขุ อาวิญฺฉติ มนาปิเยสุ รูเปสุ, อมนาปิยา รูปา ปฏิกูลา โหนฺติ ฯเปฯ มโน อาวิญฺฉติ มนาปิเยสุ ธมฺเมสุ, อมนาปิยา ธมฺมา ปฏิกูลา โหนฺติ.\csromanspacer{} เอวํ โข ภิกฺขเว อสํวโร โหติ.}

\prose{กถญฺจ ภิกฺขเว สํวโร โหติ, อิธ ภิกฺขเว ภิกฺขุ จกฺขุนา รูปํ ทิสฺวา ปิยรูเป รูเป นาธิมุจฺจติ, อปฺปิยรูเป รูเป น พฺยาปชฺชติ, อุปฏฺฐิตกายสฺสติ}

\vspace{\tipitakaparskip}
\csromancenter{สฬายตนสํยุตฺต จ วิหรติ อปฺปมาณเจตโส.\csromanspacer{} ตญฺจ เจโตวิมุตฺติํ ปญฺญาวิมุตฺติํ ยถาภูตํ ปชานาติ, ยตฺถสฺส เต อุปฺปนฺนา ปาปกา อกุสลา ธมฺมา อปริเสสา นิรุชฺฌนฺติ ฯเปฯ.\csromanspacer{} ชิวฺหาย รสํ สายิตฺวา ฯเปฯ.\csromanspacer{} มนสา ธมฺมํ วิญฺญาย ปิยรูเป ธมฺเม นาธิมุจฺจติ, อปฺปิยรูเป ธมฺเม น พฺยาปชฺชติ, อุปฏฺฐิตกายสฺสติ จ วิหรติ อปฺปมาณเจตโส.\csromanspacer{} ตญฺจ เจโตวิมุตฺติํ ปญฺญาวิมุตฺติํ ยถาภูตํ ปชานาติ, ยตฺถสฺส เต อุปฺปนฺนา ปาปกา อกุสลา ธมฺมา อปริเสสา นิรุชฺฌนฺติ.}
\prose{\csromanfolio{403}เสยฺยถาปิ ภิกฺขเว ปุริโส ฉปฺปาณเก คเหตฺวา นานาวิสเย นานาโคจเร ทฬฺหาย รชฺชุยา พนฺเธยฺย, อหิํ คเหตฺวา ทฬฺหาย รชฺชุยา พนฺเธยฺย.\csromanspacer{} สุสุมารํ คเหตฺวา ทฬฺหาย รชฺชุยา พนฺเธยฺย.\csromanspacer{} ปกฺขิํ คเหตฺวา ฯเปฯ.\csromanspacer{} กุกฺกุรํ คเหตฺวา.\csromanspacer{} สิงฺคาลํ คเหตฺวา.\csromanspacer{} มกฺกฏํ คเหตฺวา ทฬฺหาย รชฺชุยา พนฺเธยฺย, ทฬฺหาย รชฺชุยา พนฺธิตฺวา ทเฬฺห ขีเล วา ถมฺเภ วา อุปนิพนฺเธยฺย.\csromanspacer{} อถ โข เต ภิกฺขเว ฉปฺปาณกา นานาวิสยา นานาโคจรา สกํ สกํ โคจรวิสยํ อาวิญฺเฉยฺยุํ, อหิ อาวิญฺเฉยฺย “วมฺมิกํ ปเวกฺขามี”ติ, สุสุมาโร อาวิญฺเฉยฺย “อุทกํ ปเวกฺขามี”ติ, ปกฺขี อาวิญฺเฉยฺย “อากาสํ เฑสฺสามี”ติ, กุกฺกุโร อาวิญฺเฉยฺย “คามํ ปเวกฺขามี”ติ, สิงฺคาโล อาวิญฺเฉยฺย “สีวถิกํ ปเวกฺขามี”ติ, มกฺกโฏ อาวิญฺเฉยฺย “วนํ ปเวกฺขามี”ติ.\csromanspacer{} ยทา โข เต ภิกฺขเว ฉปฺปาณกา ฌตฺตา อสฺสุ กิลนฺตา, อถ ตเมว ขีลํ วา ถมฺภํ วา อุปติฏฺเฐยฺยุํ อุปนิสีเทยฺยุํ อุปนิปชฺเชยฺยุํ.\csromanspacer{} เอวเมว โข ภิกฺขเว ยสฺส กสฺสจิ ภิกฺขุโน กายคตาสติ ภาวิตา พหุลีกตา, ตํ จกฺขุ นาวิญฺฉติ มนาปิเยสุ รูเปสุ, อมนาปิยา รูปา นปฺปฏิกูลา โหนฺติ ฯเปฯ ชิวฺหา นาวิญฺฉติ มนาปิเยสุ รเสสุ ฯเปฯ มโน นาวิญฺฉติ มนาปิเยสุ ธมฺเมสุ, อมนาปิยา ธมฺมา นปฺปฏิกูลา โหนฺติ.\csromanspacer{} เอวํ โข ภิกฺขเว สํวโร โหติ.}

\prose{“ทเฬฺห ขีเล วา ถมฺเภ วา”ติ โข ภิกฺขเว กายคตาย สติยา เอตํ อธิวจนํ.\csromanspacer{} ตสฺมาติห โว ภิกฺขเว เอวํ สิกฺขิตพฺพํ “กายคตา โน สติ ภาวิตา ภวิสฺสติ พหุลีกตา ยานีกตา วตฺถุกตา อนุฏฺฐิตา ปริจิตา สุสมารทฺธา”ติ.\csromanspacer{} เอวํ หิ โข ภิกฺขเว สิกฺขิตพฺพนฺติ.\csromanspacer{}\csromanspacer{}ทสมํ.}

\gambhira{สฬายตนวคฺคสํยุตฺตปาฬิ}
\csromanheader{2}{11. ยวกลาปิสุตฺต}
\csromanmatikaanchor{mtk.216}
\csromantocmarkref{subsection}{11. ยวกลาปิสุตฺต}{mtk.216}
\bookmark[dest={mtk.216},level=2]{11. ยวกลาปิสุตฺต}
\proseitem{248}{\csromanfolio{404}เสยฺยถาปิ ภิกฺขเว ยวกลาปี จาตุมหาปเถ นิกฺขิตฺตา อสฺส, อถ ฉ ปุริสา อาคจฺเฉยฺยุํ พฺยาภงฺคิหตฺถา, เต ยวกลาปิํ ฉหิ พฺยาภงฺคีหิ หเนยฺยุํ, เอวํ หิ สา ภิกฺขเว ยวกลาปี สุหตา อสฺส ฉหิ พฺยาภงฺคีหิ หญฺญมานา.\csromanspacer{} อถ สตฺตโม ปุริโส อาคจฺเฉยฺย พฺยาภงฺคิหตฺโถ, โส ตํ ยวกลาปิํ สตฺตมาย พฺยาภงฺคิยา หเนยฺย, เอวํ หิ สา ภิกฺขเว ยวกลาปี สุหตตรา อสฺส สตฺตมาย พฺยาภงฺคิยา หญฺญมานา.\csromanspacer{} เอวเมว โข ภิกฺขเว อสฺสุตวา ปุถุชฺชโน จกฺขุสฺมิํ หญฺญติ มนาปามนาเปหิ รูเปหิ ฯเปฯ ชิวฺหาย หญฺญติ มนาปามนาเปหิ รเสหิ ฯเปฯ มนสฺมิํ หญฺญติ มนาปามนาเปหิ ธมฺเมหิ.\csromanspacer{} สเจ โส ภิกฺขเว อสฺสุตวา ปุถุชฺชโน อายติํ ปุนพฺภวาย เจเตติ.\csromanspacer{} เอวํ หิ โส ภิกฺขเว โมฆปุริโส สุหตตโร โหติ, เสยฺยถาปิ สา ยวกลาปี สตฺตมาย พฺยาภงฺคิยา หญฺญมานา.}

\prose{ภูตปุพฺพํ ภิกฺขเว เทวาสุรสงฺคาโม สมุปพฺยุโฬฺห\footnote{สมุปพฺพูโฬฺห (สี, อิ)} อโหสิ.\csromanspacer{} อถ โข ภิกฺขเว เวปจิตฺติ อสุรินฺโท อสุเร อามนฺเตสิ “สเจ มาริสา เทวาสุรสงฺคาเม สมุปพฺยูเฬฺห อสุรา ชิเนยฺยุํ, เทวา ปราชิเนยฺยุํ, เยน นํ สกฺกํ เทวานมินฺทํ กณฺฐปญฺจเมหิ พนฺธเนหิ พนฺธิตฺวา มม สนฺติเก อาเนยฺยาถ อสุรปุรนฺ”ติ.\csromanspacer{} สกฺโกปิ โข ภิกฺขเว เทวานมินฺโท เทเว ตาวติํเส อามนฺเตสิ “สเจ มาริโส เทวาสุรสงฺคาเม สมุปพฺยูเฬฺห เทวา ชิเนยฺยุํ, อสุรา ปราชิเนยฺยุํ, เยน นํ เวปจิตฺติํ อสุรินฺทํ กณฺฐปญฺจเมหิ พนฺเธเนหิ พนฺธิตฺวา มม สนฺติเก อาเนยฺยาถ สุธมฺมํ เทวสภนฺ”ติ.\csromanspacer{} ตสฺมิํ โข ปน ภิกฺขเว สงฺคาเม เทวา ชินิํสุ, อสุรา ปราชินิํสุ.\csromanspacer{} อถ โข ภิกฺขเว เทวา ตาวติํสา เวปจิตฺติํ อสุรินฺทํ กณฺฐปญฺจเมหิ พนฺธเนหิ พนฺธิตฺวา สกฺกสฺส เทวานมินฺทสฺส สนฺติเก อาเนสุํ สุธมฺมํ เทวสตํ.\csromanspacer{} ตตฺร สุทํ ภิกฺขเว เวปจิตฺติ อสุรินฺโท กณฺฐปญฺจเมหิ พนฺธเนหิ พทฺโธ\footnote{พนฺโธ (สี, สฺยา, กํ, ก)} โหติ.\csromanspacer{} ยทา โข ภิกฺขเว เวปจิตฺติสฺส อสุรินฺทสฺส เอวํ โหติ “ธมฺมิกา โข เทวา, อธมฺมิกา อสุรา, อิเธว ทานาหํ เทวปุรํ คจฺฉามี”ติ.\csromanspacer{} อถ กณฺฐปญฺจเมหิ พนฺธเนหิ มุตฺตํ อตฺตานํ สมนุปสฺสติ, ทิพฺเพหิ จ ปญฺจหิ กามคุเณหิ สมปฺปิโต สมงฺคีภูโต ปริจาเรติ.\csromanspacer{} ยทา จ โข ภิกฺขเว เวปจิตฺติสฺส อสุรินฺทสฺส เอวํ โหติ “ธมฺมิกา โข}

\vspace{\tipitakaparskip}
\csromancenter{สฬายตนสํยุตฺต อสุรา, อธมฺมิกา เทวา, ตตฺเถว ทานาหํ อสุรปุรํ คมิสฺสามี”ติ.\csromanspacer{} อถ กณฺฐปญฺจเมหิ พนฺธเนหิ พทฺธํ อตฺตานํ สมนุปสฺสติ, ทิพฺเพหิ จ ปญฺจหิ กามคุเณหิ ปริหายติ.\csromanspacer{} เอวํ สุขุมํ โข ภิกฺขเว เวปจิตฺติพนฺธนํ, ตโต สุขุมตรํ มารพนฺธนํ.\csromanspacer{} มญฺญมาโน โข ภิกฺขเว พทฺโธ มารสฺส, อมญฺญมาโน มุตฺโต ปาปิมโต.}
\prose{\csromanfolio{405}“อสฺมี”ติ ภิกฺขเว มญฺญิตเมตํ, “อยมหมสฺมี”ติ มญฺญิตเมตํ, “ภวิสฺสนฺ”ติ มญฺญิตเมตํ, “น ภวิสฺสนฺ”ติ มญฺญิตเมตํ, “รูปี ภวิสฺสนฺ”ติ มญฺญิตเมตํ, “อรูปี ภวิสฺสนฺ”ติ มญฺญิตเมตํ, “สญฺญี ภวิสฺสนฺ”ติ มญฺญิตเมตํ, “อสญฺญี ภวิสฺสนฺ”ติ มญฺญิตเมตํ, “เนวสญฺญีนาสญฺญี ภวิสฺสนฺ”ติ มญฺญิตเมตํ.\csromanspacer{} มญฺญิตํ ภิกฺขเว โรโค, มญฺญิตํ คณฺโฑ, มญฺญิตํ สลฺลํ.\csromanspacer{} ตสฺมาติห ภิกฺขเว “อมญฺญมาเนน\footnote{อมญฺญิตมาเนน (อิ, ก)} เจตสา วิหริสฺสามา”ติ เอวํ หิ โว ภิกฺขเว สิกฺขิตพฺพํ.}

\prose{“อสฺมี”ติ ภิกฺขเว อิญฺชิตเมตํ, “อยมหมสฺมี”ติ อิญฺชิตเมตํ. “ภวิสฺสนฺ”ติ อิญฺชิตเมตํ, “น ภวิสฺสนฺ”ติ อิญฺชิตเมตํ, “รูปี ภวิสฺสนฺ”ติ อิญฺชิตเมตํ, “อรูปี ภวิสฺสนฺ”ติ อิญฺชิตเมตํ, “สญฺญี ภวิสฺสนฺ”ติ อิญฺชิตเมตํ, “อสญฺญี ภวิสฺสนฺ”ติ อิญฺชิตเมตํ, “เนวสญฺญีนาสญฺญี ภวิสฺสนฺ”ติ อิญฺชิตเมตํ.\csromanspacer{} อิญฺชิตํ ภิกฺขเว โรโค, อิญฺชิตํ คณฺโฑ, อิญฺชิตํ สลฺลํ.\csromanspacer{} ตสฺมาติห ภิกฺขเว “อนิญฺชมาเนน\footnote{อนิญฺชิยมาเนน (สฺยา, กํ, ก)} เจตสา วิหริสฺสามา”ติ เอวํ หิ โว ภิกฺขเว สิกฺขิตพฺพํ.}

\prose{“อสฺมี”ติ ภิกฺขเว ผนฺทิตเมตํ, “อยมหมสฺมี”ติ ผนฺทิตเมตํ, “ภวิสฺสนฺ”ติ ฯเปฯ “น ภวิสฺสนฺ”ติ. “รูปี ภวิสฺสนฺ”ติ. “อรูปี ภวิสฺสนฺ”ติ. “สญฺญี ภวิสฺสนฺ”ติ. “อสญฺญี ภวิสฺสนฺ”ติ. “เนวสญฺญีนาสญฺญี ภวิสฺสนฺ”ติ ผนฺทิตเมตํ.\csromanspacer{} ผนฺทิตํ ภิกฺขเว โรโค, ผนฺทิตํ คณฺโฑ, ผนฺทิตํ สลฺลํ.\csromanspacer{} ตสฺมาติห ภิกฺขเว “อผนฺทมาเนน\footnote{อผนฺทิยมาเนน (สฺยา, กํ, ก)} เจตสา วิหรสฺสามา”ติ เอวํ หิ โว ภิกฺขเว สิกฺขิตพฺพํ.}

\prose{“อสฺมี”ติ ภิกฺขเว ปปญฺจิตเมตํ, “อยมหมสฺมี”ติ ปปญฺจิตเมตํ. “ภวิสฺสนฺ”ติ ฯเปฯ “น ภวิสฺสนฺ”ติ. “รูปี ภวิสฺสนฺ”ติ. “อรูปี ภวิสฺสนฺ”ติ. “สญฺญี ภวิสฺสนฺ”ติ. “อสญฺญี ภวิสฺสนฺ”ติ. “เนวสญฺญีนาสญฺญี ภวิสฺสนฺ”ติ}

\gambhira{สฬายตนวคฺคสํยุตฺตปาฬิ}
\prose{\csromanfolio{406}ปปญฺจิตเมตํ.\csromanspacer{} ปปญฺจิตํ ภิกฺขเว โรโค, ปปญฺจิตํ คณฺโฑ, ปปญฺจิตํ สลฺลํ.\csromanspacer{} ตสฺมาติห ภิกฺขเว “นิปฺปปญฺเจน เจตสา วิหริสฺสามา”ติ เอวํ หิ โว ภิกฺขเว สิกฺขิตพฺพํ.}

\prose{“อสฺมี”ติ ภิกฺขเว มานคตเมตํ, “อยมหมสฺมี”ติ มานคตเมตํ, “ภวิสฺสนฺ”ติ มานคตเมตํ, “น ภวิสฺสนฺ”ติ มานคตเมตํ, “รูปี ภวิสฺสนฺ”ติ มานคตเมตํ, “อรูปี ภวิสฺสนฺ”ติ มานคตเมตํ, “สญฺญี ภวิสฺสนฺ”ติ มานคตเมตํ, “อสญฺญี ภวิสฺสนฺ”ติ มานคตเมตํ, “เนวสญฺญีนาสญฺญี ภวิสฺสนฺ”ติ มานคตเมตํ.\csromanspacer{} มานคตํ ภิกฺขเว โรโค, มานคตํ คณฺโฑ, มานคตํ สลฺลํ.\csromanspacer{} ตสฺมาติห ภิกฺขเว “นิหตมาเนน เจตสา วิหริสฺสามา”ติ เอวํ หิ โว ภิกฺขเว สิกฺขิตพฺพนฺติ.\csromanspacer{}\csromanspacer{}เอกาทสมํ.\csromanspacer{} อาสีวิสวคฺโค จตุตฺโถ.}

\titlehead{\textbf{ตสฺสุทฺทานํ}}
\csromangathameasure{อาสีวิโส รโถ กุมฺโม, เทฺว ทารุกฺขนฺธา อวสฺสุโต. \\ ทุกฺขธมฺมา กิํสุกา วีณา, ฉปฺปาณา ยวกลาปีติ. สฬายตนวคฺเค จตุตฺถปณฺณาสโก สมตฺโต.}
\csromangathagroup{\csromangathastanza{อาสีวิโส รโถ กุมฺโม, เทฺว ทารุกฺขนฺธา อวสฺสุโต. \\ ทุกฺขธมฺมา กิํสุกา วีณา, ฉปฺปาณา ยวกลาปีติ.\csromanspacer{} สฬายตนวคฺเค จตุตฺถปณฺณาสโก สมตฺโต.}}

\csromanheader{2}{\textbf{ตสฺส วคฺคุทฺทานํ}}
\csromangathameasure{นินฺทิกฺขโย สฏฺฐินโย, สมุทฺโท อุรเคน จ. \\ จตุปณฺณาสกา เอเต, นิปาเตสุ ปกาสิตาติ. \\ \textbf{สฬายตนสํยุตฺตํ สมตฺตํ.}}
\csromangathagroup{\csromangathastanza{นินฺทิกฺขโย สฏฺฐินโย, สมุทฺโท อุรเคน จ. \\ จตุปณฺณาสกา เอเต, นิปาเตสุ ปกาสิตาติ. \\ \textbf{สฬายตนสํยุตฺตํ สมตฺตํ.}}}

\chapterheadpageread{2. เวทนาสํยุตฺต}
\csromanmatikaanchor{mtk.217}
\csromantocmarknopage{chapter}{2. เวทนาสํยุตฺต}{mtk.217}
\bookmark[dest={mtk.217},level=0]{2. เวทนาสํยุตฺต}
\csromanheader{1}{1. สคาถาวคฺค}
\csromanmatikaanchor{mtk.218}
\csromantocmarknopage{section}{1. สคาถาวคฺค}{mtk.218}
\bookmark[dest={mtk.218},level=1]{1. สคาถาวคฺค}
\csromanheader{2}{1. สมาธิสุตฺต}
\csromanmatikaanchor{mtk.219}
\csromantocmarkref{subsection}{1. สมาธิสุตฺต}{mtk.219}
\bookmark[dest={mtk.219},level=2]{1. สมาธิสุตฺต}
\proseitem{249}{\csromanfolio{407}ติสฺโส อิมา ภิกฺขเว เวทนา.\csromanspacer{} กตมา ติสฺโส, สุขา เวทนา ทุกฺขา เวทนา อทุกฺขมสุขา เวทนา.\csromanspacer{} อิมา โข ภิกฺขเว ติสฺโส เวทนาติ.}

\csromangathameasure{สมาหิโต สมฺปชาโน, สโต พุทฺธสฺส สาวโก. \\ เวทนา จ ปชานาติ, เวทนานญฺจ สมฺภวํ. \\ ยตฺถ เจตา นิรุชฺฌนฺติ, มคฺคญฺจ ขยคามินํ. \\ เวทนานํ ขยา ภิกฺขุ, นิจฺฉาโต ปรินิพฺพุโตติ. ปฐมํ.}
\csromangathagroup{\csromangathastanza{สมาหิโต สมฺปชาโน, สโต พุทฺธสฺส สาวโก. \\ เวทนา จ ปชานาติ, เวทนานญฺจ สมฺภวํ.}\csromangathastanza{ยตฺถ เจตา นิรุชฺฌนฺติ, มคฺคญฺจ ขยคามินํ. \\ เวทนานํ ขยา ภิกฺขุ, นิจฺฉาโต ปรินิพฺพุโตติ.\csromanspacer{} ปฐมํ.}}

\csromanheader{2}{2. สุขสุตฺต}
\csromanmatikaanchor{mtk.220}
\csromantocmarkref{subsection}{2. สุขสุตฺต}{mtk.220}
\bookmark[dest={mtk.220},level=2]{2. สุขสุตฺต}
\proseitem{250}{ติสฺโส อิมา ภิกฺขเว เวทนา.\csromanspacer{} กตมา ติสฺโส, สุขา เวทนา ทุกฺขา เวทนา อทุกฺขมสุขา เวทนา.\csromanspacer{} อิมา โข ภิกฺขเว ติสฺโส เวทนาติ.}

\csromangathameasure{สุขํ วา ยทิ วา ทุกฺขํ, อทุกฺขมสุขํ สห. \\ อชฺฌตฺตญฺจ พหิทฺธา จ, ยํ กิญฺจิ อตฺถิ เวทิตํ. \\ เอตํ ทุกฺขนฺติ ญตฺวาน, โมสธมฺมํ ปโลกินํ. \\ ผุสฺส ผุสฺส วยํ ปสฺสํ, เอวํ ตตฺถ วิรชฺชตีติ. ทุติยํ.}
\csromangathagroup{\csromangathastanza{สุขํ วา ยทิ วา ทุกฺขํ, อทุกฺขมสุขํ สห. \\ อชฺฌตฺตญฺจ พหิทฺธา จ, ยํ กิญฺจิ อตฺถิ เวทิตํ.}\csromangathastanza{เอตํ ทุกฺขนฺติ ญตฺวาน, โมสธมฺมํ ปโลกินํ. \\ ผุสฺส ผุสฺส วยํ ปสฺสํ, เอวํ ตตฺถ วิรชฺชตีติ.\csromanspacer{} ทุติยํ.}}

\csromanheader{2}{3. ปหานสุตฺต}
\csromanmatikaanchor{mtk.221}
\csromantocmarkref{subsection}{3. ปหานสุตฺต}{mtk.221}
\bookmark[dest={mtk.221},level=2]{3. ปหานสุตฺต}
\proseitem{251}{ติสฺโส อิมา ภิกฺขเว เวทนา.\csromanspacer{} กตมา ติสฺโส, สุขา เวทนา ทุกฺขา เวทนา อทุกฺขมสุขา เวทนา.\csromanspacer{} สุขาย ภิกฺขเว เวทนาย ราคานุสโย ปหาตพฺโพ, ทุกฺขาย เวทนาย ปฏิฆานุสโย ปหาตพฺโพ, อทุกฺขมสุขาย เวทนาย อวิชฺชานุสโย ปหาตพฺโพ.\csromanspacer{} ยโต โข ภิกฺขเว ภิกฺขุโน สุขาย เวทนาย ราคานุสโย ปหีโน โหติ, ทุกฺขาย เวทนาย ปฏิฆานุสโย ปหีโน โหติ,}

\gambhira{สฬายตนวคฺคสํยุตฺตปาฬิ}
\prose{\csromanfolio{408}อทุกฺขมสุขาย เวทนาย อวิชฺชานุสโย ปหีโน โหติ.\csromanspacer{} อยํ วุจฺจติ ภิกฺขเว ภิกฺขุ นิรนุสโย สมฺมทฺทโส, อจฺเฉจฺฉิ\footnote{อจฺเฉชฺชิ (พหูสุ)} ตณฺหํ, วิวตฺตยิ\footnote{วาวตฺตยิ (สี)} สํโยชนํ, สมฺมา มานาภิสมยา อนฺตมกาสิ ทุกฺขสฺสาติ.}

\csromangathameasure{สุขํ เวทยมานสฺส, เวทนํ อปฺปชานโต. \\ โส ราคานุสโย โหติ, อนิสฺสรณทสฺสิโน. \\ ทุกฺขํ เวทยมานสฺส, เวทนํ อปฺปชานโต. \\ ปฏิฆานุสโย โหติ, อนิสฺสรณทสฺสิโน. \\ อทุกฺขมสุขํ สนฺตํ, ภูริปญฺเญน เทสิตํ. \\ ตญฺจาปิ อภินนฺทติ, เนว ทุกฺขา ปมุจฺจติ. \\ ยโต จ ภิกฺขุ อาตาปี, สมฺปชญฺญํ น ริญฺจติ. \\ ตโต โส เวทนา สพฺพา, ปริชานาติ ปณฺฑิโต. \\ โส เวทนา ปริญฺญาย, ทิฏฺเฐ ธมฺเม อนาสโว. \\ กายสฺส เภทา ธมฺมฏฺโฐ, สงฺขฺยํ โนเปติ เวทคูติ. ตติยํ.}
\csromangathagroup{\csromangathastanza{สุขํ เวทยมานสฺส\footnote{เวทิยมานสฺส (สี, อิ)}, เวทนํ อปฺปชานโต. \\ โส ราคานุสโย โหติ, อนิสฺสรณทสฺสิโน.}\csromangathastanza{ทุกฺขํ เวทยมานสฺส, เวทนํ อปฺปชานโต. \\ ปฏิฆานุสโย โหติ, อนิสฺสรณทสฺสิโน.}\csromangathastanza{อทุกฺขมสุขํ สนฺตํ, ภูริปญฺเญน เทสิตํ. \\ ตญฺจาปิ อภินนฺทติ, เนว ทุกฺขา ปมุจฺจติ.}\csromangathastanza{ยโต จ ภิกฺขุ อาตาปี, สมฺปชญฺญํ น ริญฺจติ. \\ ตโต โส เวทนา สพฺพา, ปริชานาติ ปณฺฑิโต.}\csromangathastanza{โส เวทนา ปริญฺญาย, ทิฏฺเฐ ธมฺเม อนาสโว. \\ กายสฺส เภทา ธมฺมฏฺโฐ, สงฺขฺยํ โนเปติ เวทคูติ.\csromanspacer{} ตติยํ.}}

\csromanheader{2}{4. ปาตาลสุตฺต}
\csromanmatikaanchor{mtk.222}
\csromantocmarkref{subsection}{4. ปาตาลสุตฺต}{mtk.222}
\bookmark[dest={mtk.222},level=2]{4. ปาตาลสุตฺต}
\proseitem{252}{อสฺสุตวา ภิกฺขเว ปุถุชฺชโน ยํ วาจํ ภาสติ “อตฺถิ มหาสมุทฺเท ปาตาโล”ติ, ตํ โข ปเนตํ ภิกฺขเว อสฺสุตวา ปุถุชฺชโน อสนฺตํ อวิชฺชมานํ เอวํ วาจํ ภาสติ “อตฺถิ มหาสมุทฺเท ปาตาโล”ติ.\csromanspacer{} สารีริกานํ โข เอตํ ภิกฺขเว ทุกฺขานํ เวทนานํ อธิวจนํ, ยทิทํ “ปาตาโล”ติ.\csromanspacer{} อสฺสุตวา ภิกฺขเว ปุถุชฺชโน สารีริกาย ทุกฺขาย เวทนาย ผุฏฺโฐ สมาโน โสจติ กิลมติ ปริเทวติ อุรตฺตาฬิํ กนฺทติ สมฺโมหํ อาปชฺชติ, อยํ วุจฺจติ ภิกฺขเว อสฺสุตวา ปุถุชฺชโน ปาตาเล น ปจฺจุฏฺฐาสิ, คาธญฺจ นาชฺฌคา.\csromanspacer{} สุตวา จ โข ภิกฺขเว อริยสาวโก สารีริกาย ทุกฺขาย เวทนาย ผุฏฺโฐ สมาโน เนว โสจติ น กิลมติ น ปริเทวติ น อุรตฺตาฬิํ กนฺทติ น สมฺโมหํ อาปชฺชติ, อยํ วุจฺจติ ภิกฺขเว สุตวา อริยสาวโก ปาตาเล ปจฺจุฏฺฐาสิ, คาธญฺจ อชฺฌคาติ.}

\csromanheader{2}{2. เวทนาสํยุตฺต}
\csromangathameasure{โย เอตา นาธิวาเสติ, อุปฺปนฺนา เวทนา ทุขา. \\ สารีริกา ปาณหรา, ยาหิ ผุฏฺโฐ ปเวธติ. \\ อกฺกนฺทติ ปโรทติ, ทุพฺพโล อปฺปถามโก. \\ น โส ปาตาเล ปจฺจุฏฺฐาสิ, อโถ คาธมฺปิ นาชฺฌคา. \\ โย เจตา อธิวาเสติ, อุปฺปนฺนา เวทนา ทุขา. \\ สารีริกา ปาณหรา, ยาหิ ผุฏฺโฐ น เวธติ. \\ ส เว ปาตาเล ปจฺจุฏฺฐาสิ, อโถ คาธมฺปิ อชฺฌคาติ. จตุตฺถํ.}
\csromangathagroup{\csromangathastanza{\csromanfolio{409}โย เอตา นาธิวาเสติ, อุปฺปนฺนา เวทนา ทุขา. \\ สารีริกา ปาณหรา, ยาหิ ผุฏฺโฐ ปเวธติ.}\csromangathastanza{อกฺกนฺทติ ปโรทติ, ทุพฺพโล อปฺปถามโก. \\ น โส ปาตาเล ปจฺจุฏฺฐาสิ, อโถ คาธมฺปิ นาชฺฌคา.}\csromangathastanza{โย เจตา อธิวาเสติ, อุปฺปนฺนา เวทนา ทุขา. \\ สารีริกา ปาณหรา, ยาหิ ผุฏฺโฐ น เวธติ. \\ ส เว ปาตาเล ปจฺจุฏฺฐาสิ, อโถ คาธมฺปิ อชฺฌคาติ.\csromanspacer{} จตุตฺถํ.}}

\csromanheader{2}{5. ทฏฺฐพฺพสุตฺต}
\csromanmatikaanchor{mtk.223}
\csromantocmarkref{subsection}{5. ทฏฺฐพฺพสุตฺต}{mtk.223}
\bookmark[dest={mtk.223},level=2]{5. ทฏฺฐพฺพสุตฺต}
\proseitem{253}{ติสฺโส อิมา ภิกฺขเว เวทนา.\csromanspacer{} กตมา ติสฺโส, สุขา เวทนา ทุกฺขา เวทนา อทุกฺขมสุขา เวทนา.\csromanspacer{} สุขา ภิกฺขเว เวทนา ทุกฺขโต ทฏฺฐพฺพา, ทุกฺขา เวทนา สลฺลโต ทฏฺฐพฺพา, อทุกฺขมสุขา เวทนา อนิจฺจโต ทฏฺฐพฺพา.\csromanspacer{} ยโต โข ภิกฺขเว ภิกฺขุโน สุขา เวทนา ทุกฺขโต ทิฏฺฐา โหติ, ทุกฺขา เวทนา สลฺลโต ทิฏฺฐา โหติ, อทุกฺขมสุขา เวทนา อนิจฺจโต ทิฏฺฐา โหติ, อยํ วุจฺจติ ภิกฺขเว ภิกฺขุ สมฺมทฺทโส, อจฺเฉจฺฉิ ตณฺหํ, วิวตฺตยิ สํโยชนํ, สมฺมา มานาภิสมยา อนฺตมกาสิ ทุกฺขสฺสาติ.}

\csromangathameasure{โย สุขํ ทุกฺขโต อทฺท, ทุกฺขมทฺทกฺขิ สลฺลโต. \\ อทุกฺขมสุขํ สนฺตํ, อทฺทกฺขิ นํ อนิจฺจโต. \\ ส เว สมฺมทฺทโส ภิกฺขุ, ปริชานาติ เวทนา. \\ โส เวทนา ปริญฺญาย, ทิฏฺเฐ ธมฺเม อนาสโว. \\ กายสฺส เภทา ธมฺมฏฺโฐ, สงฺขฺยํ โนเปติ เวทคูติ. ปญฺจมํ.}
\csromangathagroup{\csromangathastanza{โย สุขํ ทุกฺขโต อทฺท, ทุกฺขมทฺทกฺขิ สลฺลโต. \\ อทุกฺขมสุขํ สนฺตํ, อทฺทกฺขิ นํ อนิจฺจโต.}\csromangathastanza{ส เว สมฺมทฺทโส ภิกฺขุ, ปริชานาติ เวทนา. \\ โส เวทนา ปริญฺญาย, ทิฏฺเฐ ธมฺเม อนาสโว. \\ กายสฺส เภทา ธมฺมฏฺโฐ, สงฺขฺยํ โนเปติ เวทคูติ.\csromanspacer{} ปญฺจมํ.}}

\csromanheader{2}{6. สลฺลสุตฺต}
\csromanmatikaanchor{mtk.224}
\csromantocmarkref{subsection}{6. สลฺลสุตฺต}{mtk.224}
\bookmark[dest={mtk.224},level=2]{6. สลฺลสุตฺต}
\proseitem{254}{อสฺสุตวา ภิกฺขเว ปุถุชฺชโน สุขมฺปิ เวทนํ เวทยติ\footnote{เวทิยติ (สี, อิ)}, ทุกฺขมฺปิ เวทนํ เวทยติ, อทุกฺขมสุขมฺปิ เวทนํ เวทยติ.\csromanspacer{} สุตวา ภิกฺขเว อริยสาวโก สุขมฺปิ เวทนํ เวทยติ, ทุกฺขมฺปิ เวทนํ เวทยติ,}

\gambhira{สฬายตนวคฺคสํยุตฺตปาฬิ}
\prose{\csromanfolio{410}อทุกฺขมสุขมฺปิ เวทนํ เวทยติ.\csromanspacer{} ตตฺร ภิกฺขเว โก วิเสโส, โก อธิปฺปยาโส\footnote{อธิปฺปาโย (สี, ก), อธิปฺปายโส (สฺยา, กํ), อธิปฺปาโยโส (อิ)}, กิํ นานากรณํ สุตวโต อริยสาวกสฺส อสฺสุตวตา ปุถุชฺชเนนาติ.\csromanspacer{} ภควํมูลกา โน ภนฺเต ธมฺมา ฯเปฯ.\csromanspacer{} อสฺสุตวา ภิกฺขเว ปุถุชฺชโน ทุกฺขาย เวทนาย ผุฏฺโฐ สมาโน โสจติ กิลมติ ปริเทวติ อุรตฺตาฬิํ กนฺทติ สมฺโมหํ อาปชฺชติ, โส เทฺว เวทนา เวทยติ กายิกญฺจ เจตสิกญฺจ.\csromanspacer{} เสยฺยถาปิ ภิกฺขเว ปุริสํ สลฺเลน วิชฺเฌยฺย\footnote{สลฺเลน อนุวิชฺเฌยฺยุํ (สี, สฺยา, กํ, อิ)}, ตเมนํ ทุติเยน สลฺเลน อนุเวธํ วิชฺเฌยฺย\footnote{สลฺเลน อนุวิชฺเฌยฺยุํ (สี), สลฺเลน อนุเวธํ วิชฺเฌยฺยุํ (สฺยา, กํ), สลฺเลน วิชฺเฌยฺยุํ (อิ)}.\csromanspacer{} เอวํ หิ โส ภิกฺขเว ปุริโส ทฺวิสลฺเลน เวทนํ เวทยติ.\csromanspacer{} เอวเมว โข ภิกฺขเว อสฺสุตวา ปุถุชฺชโน ทุกฺขาย เวทนาย ผุฏฺโฐ สมาโน โสจติ กิลมติ ปริเทวติ อุรตฺตาฬิํ กนฺทติ สมฺโมหํ อาปชฺชติ, โส เทฺว เวทนา เวทยติ กายิกญฺจ เจตสิกญฺจ.\csromanspacer{} ตสฺสาเยว โข ปน ทุกฺขาย เวทนาย ผุฏฺโฐ สมาโน ปฏิฆวา โหติ, ตเมนํ ทุกฺขาย เวทนาย ปฏิฆวนฺตํ โย ทุกฺขาย เวทนาย ปฏิฆานุสโย, โส อนุเสติ.\csromanspacer{} โส ทุกฺขาย เวทนาย ผุฏฺโฐ สมาโน กามสุขํ อภินนฺทติ.\csromanspacer{} ตํ กิสฺส เหตุ, น หิ โส ภิกฺขเว ปชานาติ อสฺสุตวา ปุถุชฺชโน อญฺญตฺร กามสุขา ทุกฺขาย เวทนาย นิสฺสรณํ, ตสฺส กามสุขญฺจ อภินนฺทโต โย สุขาย เวทนาย ราคานุสโย, โส อนุเสติ.\csromanspacer{} โส ตาสํ เวทนานํ สมุทยญฺจ อตฺถงฺคมญฺจ อสฺสาทญฺจ อาทีนวญฺจ นิสฺสรณญฺจ ยถาภูตํ นปฺปชานาติ, ตสฺส ตาสํ เวทนานํ สมุทยญฺจ อตฺถงฺคมญฺจ อสฺสาทญฺจ อาทีนวญฺจ นิสฺสรณญฺจ ยถาภูตํ อปฺปชานโต โย อทุกฺขมสุขาย เวทนาย อวิชฺชานุสโย, โส อนุเสติ.\csromanspacer{} โส สุขญฺเจ เวทนํ เวทยติ, สญฺญุตฺโต นํ เวทยติ.\csromanspacer{} ทุกฺขญฺเจ เวทนํ เวทยติ, สญฺญุตฺโต นํ เวทยติ.\csromanspacer{} อทุกฺขมสุขญฺเจ เวทนํ เวทยติ, สญฺญุตฺโต นํ เวทยติ.\csromanspacer{} อยํ วุจฺจติ ภิกฺขเว อสฺสุตวา ปุถุชฺชโน สญฺญุตฺโต ชาติยา ชราย มรเณน โสเกหิ ปริเทเวหิ ทุกฺเขหิ โทมนสฺเสหิ อุปายาเสหิ, สญฺญุตฺโต ทุกฺขสฺมาติ วทามิ.}

\prose{สุตวา จ โข ภิกฺขเว อริยสาวโก ทุกฺขาย เวทนาย ผุฏฺโฐ สมาโน น โสจติ น กิลมติ น ปริเทวติ น อุรตฺตาฬิํ กนฺทติ}

\vspace{\tipitakaparskip}
\csromancenter{เวทนาสํยุตฺต น สมฺโมหํ อาปชฺชติ, โส เอกํ เวทนํ เวทยติ กายิกํ, น เจตสิกํ.}
\prose{\csromanfolio{411}เสยฺยถาปิ ภิกฺขเว ปุริสํ สลฺเลน วิชฺเฌยฺย, ตเมนํ ทุติเยน สลฺเลน อนุเวธํ น วิชฺเฌยฺย, เอวํ หิ โส ภิกฺขเว ปุริโส เอกสลฺเลน เวทนํ เวทยติ.\csromanspacer{} เอวเมว โข ภิกฺขเว สุตวา อริยสาวโก ทุกฺขาย เวทนาย ผุฏฺโฐ สมาโน น โสจติ น กิลมติ น ปริเทวติ น อุรตฺตาฬิํ กนฺทติ น สมฺโมหํ อาปชฺชติ, โส เอกํ เวทนํ เวทยติ กายิกํ, น เจตสิกํ.\csromanspacer{} ตสฺสาเยว โข ปน ทุกฺขาย เวทนาย ผุฏฺโฐ สมาโน ปฏิฆวา น โหติ, ตเมนํ ทุกฺขาย เวทนาย อปฺปฏิฆวนฺตํ โย ทุกฺขาย เวทนาย ปฏิฆานุสโย, โส นานุเสติ.\csromanspacer{} โส ทุกฺขาย เวทนาย ผุฏฺโฐ สมาโน กามสุขํ นาภินนฺทติ.\csromanspacer{} ตํ กิสฺส เหตุ, ปชานาติ หิ โส ภิกฺขเว สุตวา อริยสาวโก อญฺญตฺร กามสุขา ทุกฺขาย เวทนาย นิสฺสรณํ, ตสฺส กามสุขํ นาภินนฺทโต โย สุขาย เวทนาย ราคานุสโย, โส นานุเสติ.\csromanspacer{} โส ตาสํ เวทนานํ สมุทยญฺจ อตฺถงฺคมญฺจ อสฺสาทญฺจ อาทีนวญฺจ นิสฺสรณญฺจ ยถาภูตํ ปชานาติ, ตสฺส ตาสํ เวทนานํ สมุทยญฺจ อตฺถงฺคมญฺจ อสฺสาทญฺจ อาทีนวญฺจ นิสฺสรณญฺจ ยถาภูตํ ปชานโต โย อทุกฺขมสุขาย เวทนาย อวิชฺชานุสโย, โส นานุเสติ.\csromanspacer{} โส สุขญฺเจ เวทนํ เวทยติ, วิสญฺญุตฺโต นํ เวทยติ.\csromanspacer{} ทุกฺขญฺเจ เวทนํ เวทยติ, วิสญฺญุตฺโต นํ เวทยติ.\csromanspacer{} อทุกฺขมสุขญฺจ เวทนํ เวทยติ, วิสญฺญุตฺโต นํ เวทยติ.\csromanspacer{} อยํ วุจฺจติ ภิกฺขเว สุตวา อริยสาวโก วิสญฺญุตฺโต ชาติยา ชราย มรเณน โสเกหิ ปริเทเวหิ ทุกฺเขหิ โทมนสฺเสหิ อุปายาเสหิ, วิสญฺญุตฺโต ทุกฺขสฺมาติ วทามิ.\csromanspacer{} อยํ โข ภิกฺขเว วิเสโส, อยํ อธิปฺปยาโส, อิทํ นานากรณํ สุตวโต อริยสาวกสฺส อสฺสุตวตา ปุถุชฺชเนนาติ.}

\csromangathameasure{น เวทนํ เวทยติ สปญฺโญ, \\ สุขมฺปิ ทุกฺขมฺปิ พหุสฺสุโตปิ. \\ อยญฺจ ธีรสฺส ปุถุชฺชเนน, \\ มหา วิเสโส กุสลสฺส โหติ.}
\csromangathagroup{\csromangathastanza{น เวทนํ เวทยติ สปญฺโญ, \\ สุขมฺปิ ทุกฺขมฺปิ พหุสฺสุโตปิ. \\ อยญฺจ ธีรสฺส ปุถุชฺชเนน, \\ มหา\footnote{อยํ (สฺยา, กํ, ก)} วิเสโส กุสลสฺส โหติ.}}

\gambhira{สฬายตนวคฺคสํยุตฺตปาฬิ}
\csromanmatikaanchor{mtk.225}
\csromantocmarkref{subsection}{7-8. เคลญฺญสุตฺต}{mtk.225}
\bookmark[dest={mtk.225},level=2]{7-8. เคลญฺญสุตฺต}
\csromangathameasure{สงฺขาตธมฺมสฺส พหุสฺสุตสฺส, \\ วิปสฺสโต โลกมิมํ ปรญฺจ. \\ อิฏฺฐสฺส ธมฺมา น มเถนฺติ จิตฺตํ, \\ อนิฏฺฐโต โน ปฏิฆาตเมติ. \\ ตสฺสานุโรธา อถวา วิโรธา, \\ วิธูปิตา อตฺถคตา น สนฺติ. \\ ปทญฺจ ญตฺวา วิรชํ อโสกํ, \\ สมฺมา ปชานาติ ภวสฺส ปารคูติ. ฉฏฺฐํ.}
\csromangathagroup{\csromangathastanza{\csromanfolio{412}สงฺขาตธมฺมสฺส พหุสฺสุตสฺส, \\ วิปสฺสโต\footnote{สมฺปสฺสโต (สี, อิ)} โลกมิมํ ปรญฺจ. \\ อิฏฺฐสฺส ธมฺมา น มเถนฺติ จิตฺตํ, \\ อนิฏฺฐโต โน ปฏิฆาตเมติ.}\csromangathastanza{ตสฺสานุโรธา อถวา วิโรธา, \\ วิธูปิตา อตฺถคตา น สนฺติ. \\ ปทญฺจ ญตฺวา วิรชํ อโสกํ, \\ สมฺมา ปชานาติ ภวสฺส ปารคูติ.\csromanspacer{} ฉฏฺฐํ.}}

\csromanheader{2}{7. ปฐมเคลญฺญสุตฺต}
\proseitem{255}{เอกํ สมยํ ภควา เวสาลิยํ วิหรติ มหาวเน กูฏาคารสาลายํ.\csromanspacer{} อถ โข ภควา สายนฺหสมยํ ปฏิสลฺลานา วุฏฺฐิโต เยน คิลานสาลา เตนุปสงฺกมิ, อุปสงฺกมิตฺวา ปญฺญตฺเต อาสเน นิสีทิ, นิสชฺช โข ภควา ภิกฺขู อามนฺเตสิ\csromandash{}}

\prose{สโต ภิกฺขเว ภิกฺขุ สมฺปชาโน กาลํ อาคเมยฺย, อยํ โว อมฺหากํ อนุสาสนี.}

\prose{กถญฺจ ภิกฺขเว ภิกฺขุ สโต โหติ.\csromanspacer{} อิธ ภิกฺขเว ภิกฺขุ กาเย กายานุปสฺสี วิหรติ อาตาปี สมฺปชาโน สติมา วิเนยฺย โลเก อภิชฺฌาโทมนสฺสํ.\csromanspacer{} เวทนาสุ เวทนานุปสฺสี วิหรติ ฯเปฯ.\csromanspacer{} จิตฺเต จิตฺตานุปสฺสี วิหรติ.\csromanspacer{} ธมฺเมสุ ธมฺมานุปสฺสี วิหรติ อาตาปี สมฺปชาโน สติมา วิเนยฺย โลเก อภิชฺฌาโทมนสฺสํ.\csromanspacer{} เอวํ โข ภิกฺขเว ภิกฺขุ สโต โหติ.}

\prose{กถญฺจ ภิกฺขเว ภิกฺขุ สมฺปชาโน โหติ.\csromanspacer{} อิธ ภิกฺขเว ภิกฺขุ อภิกฺกนฺเต ปฏิกฺกนฺเต สมฺปชานการี โหติ, อาโลกิเต วิโลกิเต สมฺปชานการี โหติ, สมิญฺชิเต ปสาริเต สมฺปชานการี โหติ, สงฺฆาฏิปตฺตจีวรธารเณ สมฺปชานการี โหติ, อสิเต ปีเต ขายิเต สายิเต สมฺปชานการี โหติ, อุจฺจารปสฺสาวกมฺเม สมฺปชานการี โหติ, คเต ฐิเต นิสินฺเน สุตฺเต ชาคริเต ภาสิเต}

\vspace{\tipitakaparskip}
\csromancenter{เวทนาสํยุตฺต ตุณฺหีภาเว สมฺปชานการี โหติ.\csromanspacer{} เอวํ โข ภิกฺขเว ภิกฺขุ สมฺปชานการี โหติ.\csromanspacer{} สโต ภิกฺขเว ภิกฺขุ สมฺปชาโน กาลํ อาคเมยฺย, อยํ โว อมฺหากํ อนุสาสนี.}
\prose{\csromanfolio{413}ตสฺส เจ ภิกฺขเว ภิกฺขุโน เอวํ สตสฺส สมฺปชานสฺส อปฺปมตฺตสฺส อาตาปิโน ปหิตตฺตสฺส วิหรโต อุปฺปชฺชติ สุขา เวทนา, โส เอวํ ปชานาติ “อุปฺปนฺนา โข มฺยายํ สุขา เวทนา, สา จ โข ปฏิจฺจ, โน อปฺปฏิจฺจ.\csromanspacer{} กิํ ปฏิจฺจ, อิมเมว กายํ ปฏิจฺจ.\csromanspacer{} อยํ โข ปน กาโย อนิจฺโจ สงฺขโต ปฏิจฺจสมุปฺปนฺโน, อนิจฺจํ โข ปน สงฺขตํ ปฏิจฺจสมุปฺปนฺนํ กายํ ปฏิจฺจ อุปฺปนฺนา สุขา เวทนา กุโต นิจฺจา ภวิสฺสตี”ติ.\csromanspacer{} โส กาเย จ สุขาย จ เวทนาย อนิจฺจานุปสฺสี วิหรติ, วยานุปสฺสี วิหรติ, วิราคานุปสฺสี วิหรติ, นิโรธานุปสฺสี วิหรติ, ปฏินิสฺสคฺคานุปสฺสี วิหรติ.\csromanspacer{} ตสฺส กาเย จ สุขาย จ เวทนาย อนิจฺจานุปสฺสิโน วิหรโต วยานุปสฺสิโน วิหรโต วิราคานุปสฺสิโน วิหรโต นิโรธานุปสฺสิโน วิหรโต ปฏินิสฺสคฺคานุปสฺสิโน วิหรโต โย กาเย จ สุขาย จ เวทนาย ราคานุสโย โส ปหียติ.}

\prose{ตสฺส เจ ภิกฺขเว ภิกฺขุโน เอวํ สตสฺส สมฺปชานสฺส อปฺปมตฺตสฺส อาตาปิโน ปหิตตฺตสฺส วิหรโต อุปฺปชฺชติ ทุกฺขา เวทนา, โส เอวํ ปชานาติ “อุปฺปนฺนา โข มฺยายํ ทุกฺขา เวทนา, สา จ โข ปฏิจฺจ, โน อปฺปฏิจฺจ.\csromanspacer{} กิํ ปฏิจฺจ, อิมเมว กายํ ปฏิจฺจ.\csromanspacer{} อยํ โข ปน กาโย อนิจฺโจ สงฺขโต ปฏิจฺจสมุปฺปนฺโน, อนิจฺจํ โข ปน สงฺขตํ ปฏิจฺจสมุปฺปนฺนํ กายํ ปฏิจฺจ อุปฺปนฺนา ทุกฺขา เวทนา กุโต นิจฺจา ภวิสฺสตี”ติ.\csromanspacer{} โส กาเย จ ทุกฺขาย จ เวทนาย อนิจฺจานุปสฺสี วิหรติ, วยานุปสฺสี วิหรติ, วิราคานุปสฺสี วิหรติ, นิโรธานุปสฺสี วิหรติ, ปฏินิสฺสคฺคานุปสฺสี วิหรติ.\csromanspacer{} ตสฺส กาเย จ ทุกฺขาย จ เวทนาย อนิจฺจานุปสฺสิโน วิหรโต ฯเปฯ ปฏินิสฺสคฺคานุปสฺสิโน วิหรโต โย กาเย จ ทุกฺขาย จ เวทนาย ปฏิฆานุสโย โส ปหียติ.}

\prose{ตสฺส เจ ภิกฺขเว ภิกฺขุโน เอวํ สตสฺส สมฺปชานสฺส อปฺปมตฺตสฺส อาตาปิโน ปหิตตฺตสฺส วิหรโต อุปฺปชฺชติ อทุกฺขมสุขา เวทนา, โส เอวํ ปชานาติ “อุปฺปนฺนา โข มฺยายํ อทุกฺขมสุขา เวทนา, สา จ โข ปฏิจฺจ, โน อปฺปฏิจฺจ.\csromanspacer{} กิํ ปฏิจฺจ, อิมเมว กายํ ปฏิจฺจ.\csromanspacer{} อยํ โข ปน กาโย อนิจฺโจ สงฺขโต ปฏิจฺจสมุปฺปนฺโน, อนิจฺจํ โข ปน สงฺขตํ ปฏิจฺจสมุปฺปนฺนํ}

\gambhira{สฬายตนวคฺคสํยุตฺตปาฬิ}
\prose{\csromanfolio{414}กายํ ปฏิจฺจ อุปฺปนฺนา อทุกฺขมสุขา เวทนา กุโต นิจฺจา ภวิสฺสตี”ติ.\csromanspacer{} โส กาเย จ อทุกฺขมสุขาย จ เวทนาย อนิจฺจานุปสฺสี วิหรติ, วยานุปสฺสี วิหรติ, วิราคานุปสฺสี วิหรติ, นิโรธานุปสฺสี วิหรติ, ปฏินิสฺสคฺคานุปสฺสี วิหรติ, ตสฺส กาเย จ อทุกฺขมสุขา จ เวทนาย อนิจฺจานุปสฺสิโน วิหรโต ฯเปฯ ปฏินิคฺคานุปสฺสิโน วิหรโต โย กาเย จ อทุกฺขมสุขาย จ เวทนาย อวิชฺชานุสโย โส ปหียติ.}

\prose{โส สุขญฺเจ เวทนํ เวทยติ, สา อนิจฺจาติ ปชานาติ, อนชฺโฌสิตาติ ปชานาติ, อนภินนฺทิตาติ ปชานาติ.\csromanspacer{} ทุกฺขญฺเจ เวทนํ เวทยติ ฯเปฯ.\csromanspacer{} อทุกฺขมสุขญฺเจ เวทนํ เวทยติ, สา อนิจฺจาติ ปชานาติ, อนชฺโฌสิตาติ ปชานาติ, อนภินนฺทิตาติ ปชานาติ.\csromanspacer{} โส สุขญฺเจ เวทนํ เวทยติ, วิสญฺญุตฺโต นํ เวทยติ.\csromanspacer{} ทุกฺขญฺเจ เวทนํ เวทยติ, วิสญฺญุตฺโต นํ เวทยติ.\csromanspacer{} อทุกฺขมสุขญฺเจ เวทนํ เวทยติ, วิสญฺญุตฺโต นํ เวทยติ.\csromanspacer{} โส กายปริยนฺติกํ เวทนํ เวทยมาโน กายปริยนฺติกํ เวทนํ เวทยามีติ ปชานาติ, ชีวิตปริยนฺติกํ เวทนํ เวทยมาโน ชีวิตปริยนฺติกํ เวทนํ เวทยามีติ ปชานาติ, กายสฺส เภทา อุทฺธํ ชีวิตปริยาทานา อิเธว สพฺพเวทยิตานิ อนภินนฺทิตานิ สีตีภวิสฺสนฺตีติ\footnote{สีติภวิสฺสนฺตีติ (สี, อิ, ก)} ปชานาติ.}

\prose{เสยฺยถาปิ ภิกฺขเว เตลญฺจ ปฏิจฺจ วฏฺฏิํ จ ปฏิจฺจ เตลปฺปทีโป ฌาเยยฺย, ตสฺเสว เตลสฺส จ วฏฺฏิยา จ ปริยาทานา อนาหาโร นิพฺพาเยยฺย.\csromanspacer{} เอวเมว โข ภิกฺขเว ภิกฺขุ กายปริยนฺติกํ เวทนํ เวทยมาโน กายปริยนฺติกํ เวทนํ เวทยามีติ ปชานาติ, ชีวิตปริยนฺติกํ เวทนํ เวทยมาโน ชีวิตปริยนฺติกํ เวทนํ เวทยามีติ ปชานาติ, กายสฺส เภทา อุทฺธํ ชีวิตปริยาทานา อิเธว สพฺพเวทยิตานิ อนภินนฺทิตานิ สีตีภวิสฺสนฺตีติ ปชานาตีติ.\csromanspacer{}\csromanspacer{}สตฺตมํ.}

\csromanheader{2}{8. ทุติยเคลญฺญสุตฺต}
\proseitem{256}{เอกํ สมยํ ภควา เวสาลิยํ วิหรติ มหาวเน กูฏาคารสาลายํ.\csromanspacer{} อถ โข ภควา สายนฺหสมยํ ปฏิสลฺลานา วุฏฺฐิโต เยน คิลานสาลา เตนุปสงฺกมิ, อุปสงฺกมิตฺวา ปญฺญตฺเต อาสเน นิสีทิ, นิสชฺช โข ภควา ภิกฺขู อามนฺเตสิ\csromandash{}}

\csromanheader{2}{2. เวทนาสํยุตฺต}
\prose{\csromanfolio{415}สโต ภิกฺขเว ภิกฺขุ สมฺปชาโน กาลํ อาคเมยฺย, อยํ โว อมฺหากํ อนุสาสนี.}

\prose{กถญฺจ ภิกฺขเว ภิกฺขุ สโต โหติ.\csromanspacer{} อิธ ภิกฺขเว ภิกฺขุ กาเย กายานุปสฺสี วิหรติ อาตาปี สมฺปชาโน สติมา วิเนยฺย โลเก อภิชฺฌาโทมนสฺสํ.\csromanspacer{} เวทนาสุ เวทนานุปสฺสี วิหรติ ฯเปฯ.\csromanspacer{} จิตฺเต จิตฺตานุปสฺสี วิหรติ.\csromanspacer{} ธมฺเมสุ ธมฺมานุปสฺสี วิหรติ อาตาปี สมฺปชาโน สติมา วิเนยฺย โลเก อภิชฺฌาโทมนสฺสํ.\csromanspacer{} เอวํ โข ภิกฺขเว ภิกฺขุ สโต โหติ.}

\prose{กถญฺจ ภิกฺขเว ภิกฺขุ สมฺปชาโน โหติ.\csromanspacer{} อิธ ภิกฺขเว ภิกฺขุ อภิกฺกนฺเต ปฏิกฺกนฺเต สมฺปชานการี โหติ ฯเปฯ ภาสิเต ตุณฺหีภาเว สมฺปชานการี โหติ.\csromanspacer{} เอวํ โข ภิกฺขเว ภิกฺขุ สมฺปชาโน โหติ.\csromanspacer{} สโต ภิกฺขเว ภิกฺขุ สมฺปชาโน กาลํ อาคเมยฺย, อยํ โว อมฺหากํ อนุสาสนี.}

\prose{ตสฺส เจ ภิกฺขเว ภิกฺขุโน เอวํ สตสฺส สมฺปชานสฺส อปฺปมตฺตสฺส อาตาปิโน ปหิตตฺตสฺส วิหรโต อุปฺปชฺชติ สุขา เวทนา, โส เอวํ ปชานาติ “อุปฺปนฺนา โข มฺยายํ สุขา เวทนา, สา จ โข ปฏิจฺจ, โน อปฺปฏิจฺจ.\csromanspacer{} กิํ ปฏิจฺจ, อิมเมว ผสฺสํ ปฏิจฺจ.\csromanspacer{} อยํ โข ปน ผสฺโส อนิจฺโจ สงฺขโต ปฏิจฺจสมุปฺปนฺโน, อนิจฺจํ โข ปน สงฺขตํ ปฏิจฺจสมุปฺปนฺนํ ผสฺสํ ปฏิจฺจ อุปฺปนฺนา สุขา เวทนา กุโต นิจฺจา ภวิสฺสตี”ติ.\csromanspacer{} โส ผสฺเส จ สุขาย จ เวทนาย อนิจฺจานุปสฺสี วิหรติ, วยานุปสฺสี วิหรติ, วิราคานุปสฺสี วิหรติ, นิโรธานุปสฺสี วิหรติ, ปฏินิสฺสคฺคานุปสฺสี วิหรติ.\csromanspacer{} ตสฺส ผสฺเส จ สุขาย จ เวทนาย อนิจฺจานุปสฺสิโน วิหรโต วยานุปสฺสิโน วิหรโต วิราคานุปสฺสิโน วิหรโต นิโรธานุปสฺสิโน วิหรโต ปฏินิสฺสคฺคานุปสฺสิโน วิหรโต โย ผสฺเส จ สุขาย จ เวทนาย ราคานุสโย โส ปหียติ.}

\prose{ตสฺส เจ ภิกฺขเว ภิกฺขุโน เอวํ สตสฺส ฯเปฯ วิหรโต อุปฺปชฺชติ ทุกฺขา เวทนา ฯเปฯ อุปฺปชฺชติ อทุกฺขมสุขา เวทนา, โส เอวํ ปชานาติ อุปฺปนฺนา โข มฺยายํ อทุกฺขมสุขา เวทนา, สา จ โข ปฏิจฺจ, โน อปฺปฏิจฺจ.\csromanspacer{} กิํ ปฏิจฺจ, อิมเมว ผสฺสํ ปฏิจฺจ. (ยถา ปุริมสุตฺเต, ตถา วิตฺถาเรตพฺโพ.)}

\gambhira{สฬายตนวคฺคสํยุตฺตปาฬิ}
\prose{\csromanfolio{416}กายสฺส เภทา อุทฺธํ ชีวิตปริยาทานา อิเธว สพฺพเวทยิตานิ อนภินนฺทิตานิ สีตีภวิสฺสนฺตีติ ปชานาติ.}

\prose{เสยฺยถาปิ ภิกฺขเว เตลญฺจ ปฏิจฺจ วฏฺฏิํ จ ปฏิจฺจ เตลปฺปทีโป ฌาเยยฺย, ตสฺเสว เตลสฺส จ วฏฺฏิยา จ ปริยาทานา อนาหาโร นิพฺพาเยยฺย.\csromanspacer{} เอวเมว โข ภิกฺขเว ภิกฺขุ กายปริยนฺติกํ เวทนํ เวทยมาโน กายปริยนฺติกํ เวทนํ เวทยามีติ ปชานาติ.\csromanspacer{} ชีวิตปริยนฺติกํ เวทนํ เวทยมาโน ชีวิตปริยนฺติกํ เวทนํ เวทยามีติ ปชานาติ, กายสฺส เภทา อุทฺธํ ชีวิตปริยาทานา อิเธว สพฺพเวทยิตานิ อนภินนฺทิตานิ สีตีภวิสฺสนฺตีติ ปชานาตีติ.\csromanspacer{}\csromanspacer{}อฏฺฐมํ.}

\csromanheader{2}{9. อนิจฺจสุตฺต}
\csromanmatikaanchor{mtk.226}
\csromantocmarkref{subsection}{9. อนิจฺจสุตฺต}{mtk.226}
\bookmark[dest={mtk.226},level=2]{9. อนิจฺจสุตฺต}
\proseitem{257}{ติสฺโส อิมา ภิกฺขเว เวทนา อนิจฺจา สงฺขตา ปฏิจฺจสมุปฺปนฺนา ขยธมฺมา วยธมฺมา วิราคธมฺมา นิโรธธมฺมา.\csromanspacer{} กตมา ติสฺโส.\csromanspacer{} สุขา เวทนา ทุกฺขา เวทนา อทุกฺขมสุขา เวทนา.\csromanspacer{} อิมา โข ภิกฺขเว ติสฺโส เวทนา อนิจฺจา สงฺขตา ปฏิจฺจสมุปฺปนฺนา ขยธมฺมา วยธมฺมา วิราคธมฺมา นิโรธธมฺมาติ.\csromanspacer{}\csromanspacer{}นวมํ.}

\csromanheader{2}{10. ผสฺสมูลกสุตฺต}
\csromanmatikaanchor{mtk.227}
\csromantocmarkref{subsection}{10. ผสฺสมูลกสุตฺต}{mtk.227}
\bookmark[dest={mtk.227},level=2]{10. ผสฺสมูลกสุตฺต}
\proseitem{258}{ติสฺโส อิมา ภิกฺขเว เวทนา ผสฺสชา ผสฺสมูลกา ผสฺสนิทานา ผสฺสปจฺจยา.\csromanspacer{} กตมา ติสฺโส, สุขา เวทนา ทุกฺขา เวทนา อทุกฺขมสุขา เวทนา.\csromanspacer{} สุขเวทนิยํ ภิกฺขเว ผสฺสํ ปฏิจฺจ อุปฺปชฺชติ สุขา เวทนา.\csromanspacer{} ตสฺเสว สุขเวทนิยสฺส ผสฺสสฺส นิโรธา ยํ ตชฺชํ เวทยิตํ สุขเวทนิยํ ผสฺสํ ปฏิจฺจ อุปฺปนฺนา สุขา เวทนา, สา นิรุชฺฌติ, สา วูปสมฺมติ.\csromanspacer{} ทุกฺขเวทนิยํ ภิกฺขเว ผสฺสํ ปฏิจฺจ อุปฺปชฺชติ ทุกฺขา เวทนา.\csromanspacer{} ตสฺเสว ทุกฺขเวทนิยสฺส ผสฺสสฺส นิโรธา ยํ ตชฺชํ เวทยิตํ ทุกฺขเวทนิยํ ผสฺสํ ปฏิจฺจ อุปฺปนฺนา ทุกฺขา เวทนา, สา นิรุชฺฌติ, สา วูปสมฺมติ.\csromanspacer{} อทุกฺขมสุขเวทนิยํ ภิกฺขเว ผสฺสํ ปฏิจฺจ อุปฺปชฺชติ อทุกฺขมสุขา เวทนา.\csromanspacer{} ตสฺเสว อทุกฺขมสุขเวทนิยสฺส ผสฺสสฺส นิโรธา ยํ ตชฺชํ เวทยิตํ อทุกฺขมสุขเวทนิยํ ผสฺสํ ปฏิจฺจ อุปฺปนฺนา อทุกฺขมสุขา เวทนา, สา นิรุชฺฌติ, สา วูปสมฺมติ.}

\vspace{\tipitakaparskip}
\csromancenter{เวทนาสํยุตฺต เสยฺยถาปิ ภิกฺขเว ทฺวินฺนํ กฏฺฐานํ สงฺฆฏฺฏนสโมธานา\footnote{สงฺขตฺตา ตสฺส สโมธานา (สฺยา, กํ) สงฺฆาตฺตา ตสฺส สโมธานา (ก) สํ 1. 322 ปิฏฺเฐปิ ปสฺสิตพฺพํ.} อุสฺมา ชายติ, เตโช อภินิพฺพตฺตติ.\csromanspacer{} เตสํเยว กฏฺฐานํ นานาภาวา วินิกฺเขปา ยา ตชฺชา อุสฺมา, สา นิรุชฺฌติ, สา วูปสมฺมติ.\csromanspacer{} เอวเมว โข ภิกฺขเว อิมา ติสฺโส เวทนา ผสฺสชา ผสฺสมูลกา ผสฺสนิทานา ผสฺสปจฺจยา.\csromanspacer{} ตชฺชํ ผสฺสํ ปฏิจฺจ ตชฺชา เวทนา อุปฺปชฺชนฺติ.\csromanspacer{} ตชฺชสฺส ผสฺสสฺส นิโรธา ตชฺชา เวทนา นิรุชฺฌนฺตีติ.\csromanspacer{}\csromanspacer{}ทสมํ.\csromanspacer{} เวทนาสํยุตฺตสฺส สคาถาวคฺโค ปฐโม.}
\titlehead{\textbf{ตสฺสุทฺทานํ}}
\csromangathameasure{สมาธิ สุขํ ปหาเนน, ปาตาลํ ทฏฺฐพฺเพน จ. \\ สลฺเลน เจว เคลญฺญา, อนิจฺจ ผสฺสมูลกาติ.}
\csromangathagroup{\csromangathastanza{\csromanfolio{417}สมาธิ สุขํ ปหาเนน, ปาตาลํ ทฏฺฐพฺเพน จ. \\ สลฺเลน เจว เคลญฺญา, อนิจฺจ ผสฺสมูลกาติ.}}

\csromanheader{1}{2. รโหคตวคฺค}
\csromanmatikaanchor{mtk.228}
\csromantocmarknopage{section}{2. รโหคตวคฺค}{mtk.228}
\bookmark[dest={mtk.228},level=1]{2. รโหคตวคฺค}
\csromanheader{2}{1. รโหคตสุตฺต}
\csromanmatikaanchor{mtk.229}
\csromantocmarkref{subsection}{1. รโหคตสุตฺต}{mtk.229}
\bookmark[dest={mtk.229},level=2]{1. รโหคตสุตฺต}
\proseitem{259}{อถ โข อญฺญตโร ภิกฺขุ เยน ภควา เตนุปสงฺกมิ, อุปสงฺกมิตฺวา ภควนฺตํ อภิวาเทตฺวา เอกมนฺตํ นิสีทิ, เอกมนฺตํ นิสินฺโน โข โส ภิกฺขุ ภควนฺตํ เอตทโวจ\csromandash{}อิธ มยฺหํ ภนฺเต รโหคตสฺส ปฏิสลฺลีนสฺส เอวํ เจตโส ปริวิตกฺโก อุทปาทิ “ติสฺโส เวทนา วุตฺตา ภควตา สุขา เวทนา ทุกฺขา เวทนา อทุกฺขมสุขา เวทนา, อิมา ติสฺโส เวทนา วุตฺตา ภควตา.\csromanspacer{} วุตฺตํ โข ปเนตํ ภควตา ‘ยํ กิญฺจิ เวทยิตํ, ตํ ทุกฺขสฺมินฺ’ติ, กิํ นุ โข เอตํ ภควตา สนฺธาย ภาสิตํ, ‘ยํ กิญฺจิ เวทยิตํ, ตํ ทุกฺขสฺมินฺ’ติ”, สาธุ สาธุ ภิกฺขุ, ติสฺโส อิมา ภิกฺขุ เวทนา วุตฺตา มยา สุขา เวทนา ทุกฺขา เวทนา อทุกฺขมสุขา เวทนา, อิมา ติสฺโส เวทนา วุตฺตา มยา.\csromanspacer{} วุตฺตํ โข ปเนตํ ภิกฺขุ มยา “ยํ กิญฺจิ เวทยิตํ, ตํ ทุกฺขสฺมินฺ”ติ.\csromanspacer{} ตํ โข ปเนตํ ภิกฺขุ มยา สงฺขารานํเยว อนิจฺจตํ สนฺธาย ภาสิตํ “ยํ กิญฺจิ เวทยิตํ, ตํ}

\gambhira{สฬายตนวคฺคสํยุตฺตปาฬิ}
\prose{\csromanfolio{418}ทุกฺขสฺมินฺ”ติ.\csromanspacer{} ตํ โข ปเนตํ ภิกฺขุ มยา สงฺขารานํเยว ขยธมฺมตํ ฯเปฯ วยธมฺมตํ ฯเปฯ วิราคธมฺมตํ ฯเปฯ นิโรธธมฺมตํ ฯเปฯ วิปริณามธมฺมตํ สนฺธาย ภาสิตํ “ยํ กิญฺจิ เวทยิตํ, ตํ ทุกฺขสฺมินฺ”ติ.\csromanspacer{} อถ โข ปน ภิกฺขุ มยา อนุปุพฺพสงฺขารานํ นิโรโธ อกฺขาโต, ปฐมํ ฌานํ สมาปนฺนสฺส วาจา นิรุทฺธา โหติ, ทุติยํ ฌานํ สมาปนฺนสฺส วิตกฺกวิจารา นิรุทฺธา โหนฺติ, ตติยํ ฌานํ สมาปนฺนสฺส ปีติ นิรุทฺธา โหติ, จตุตฺถํ ฌานํ สมาปนฺนสฺส อสฺสาสปสฺสาสา นิรุทฺธา โหนฺติ, อากาสานญฺจายตนํ สมาปนฺนสฺส รูปสญฺญา นิรุทฺธา โหติ, วิญฺญาณญฺจายตนํ สมาปนฺนสฺส อากาสานญฺจายตนสญฺญา นิรุทฺธา โหติ, อากิญฺจญฺญายตนํ สมาปนฺนสฺส วิญฺญาณญฺจายตนสญฺญา นิรุทฺธา โหติ, เนวสญฺญานาสญฺญายตนํ สมาปนฺนสฺส อากิญฺจญฺญายตนสญฺญา นิรุทฺธา โหติ, สญฺญาเวทยิตนิโรธํ สมาปนฺนสฺส สญฺญา จ เวทนา จ นิรุทฺธา โหนฺติ.\csromanspacer{} ขีณาสวสฺส ภิกฺขุโน ราโค นิรุทฺโธ โหติ, โทโส นิรุทฺโธ โหติ, โมโห นิรุทฺโธ โหติ.\csromanspacer{} อถ โข ภิกฺขุ มยา อนุปุพฺพสงฺขารานํ วูปสโม อกฺขาโต, ปฐมํ ฌานํ สมาปนฺนสฺส วาจา วูปสนฺตา โหติ, ทุติยํ ฌานํ สมาปนฺนสฺส วิตกฺกวิจารา วูปสนฺตา โหนฺติ ฯเปฯ สญฺญาเวทยิตนิโรธํ สมาปนฺนสฺส สญฺญา จ เวทนา จ วูปสนฺตา โหนฺติ.\csromanspacer{} ขีณาสวสฺส ภิกฺขุโน ราโค วูปสนฺโต โหติ, โทโส วูปสนฺโต โหติ, โมโห วูปสนฺโต โหติ.\csromanspacer{} ฉยิมา ภิกฺขุ ปสฺสทฺธิโย, ปฐมํ ฌานํ สมาปนฺนสฺส วาจา ปฏิปฺปสฺสทฺธา โหติ, ทุติยํ ฌานํ สมาปนฺนสฺส วิตกฺกวิจารา ปฏิปฺปสฺสทฺธา โหติ, ตติยํ ฌานํ สมาปนฺนสฺส ปีติ ปฏิปฺปสฺสทฺธา โหนฺติ, จตุตฺถํ ฌานํ สมาปนฺนสฺส อสฺสาสปสฺสาสา ปฏิปฺปสฺสทฺธา โหนฺติ, สญฺญาเวทยิตนิโรธํ สมาปนฺนสฺส สญฺญา จ เวทนา จ ปฏิปฺปสฺสทฺธา โหนฺติ.\csromanspacer{} ขีณาสวสฺส ภิกฺขุโน ราโค ปฏิปฺปสฺสทฺโธ โหติ โทโส ปฏิปฺปสฺสทฺโธ โหติ, โมโห ปฏิปฺปสฺสทฺโธ โหตีติ.\csromanspacer{}\csromanspacer{}ปฐมํ.}

\csromanheader{2}{2. ปฐม-อากาสสุตฺต}
\csromanmatikaanchor{mtk.230}
\csromantocmarkref{subsection}{2-3. อากาสสุตฺต}{mtk.230}
\bookmark[dest={mtk.230},level=2]{2-3. อากาสสุตฺต}
\proseitem{260}{เสยฺยถาปิ ภิกฺขเว อากาเส วิวิธา วาตา วายนฺติ, ปุรตฺถิมาปิ วาตา วายนฺติ, ปจฺฉิมาปิ วาตา วายนฺติ, อุตฺตราปิ วาตา วายนฺติ, ทกฺขิณาปิ วาตา วายนฺติ, สรชาปิ วาตา วายนฺติ, อรชาปิ วาตา วายนฺติ, สีตาปิ วาตา วายนฺติ, อุณฺหาปิ วาตา วายนฺติ, ปริตฺตาปิ}

\vspace{\tipitakaparskip}
\csromancenter{เวทนาสํยุตฺต วาตา วายนฺติ, อธิมตฺตาปิ วาตา วายนฺติ, เอวเมว โข ภิกฺขเว อิมสฺมิํ กายสฺมิํ วิวิธา เวทนา อุปฺปชฺชนฺติ, สุขาปิ เวทนา อุปฺปชฺชติ, ทุกฺขาปิ เวทนา อุปฺปชฺชติ, อทุกฺขมสุขาปิ เวทนา อุปฺปชฺชตีติ.}
\vspace{0.5\baselineskip}
\csromangathameasure{ยถาปิ วาตา อากาเส, วายนฺติ วิวิธา ปุถู. \\ ปุรตฺถิมา ปจฺฉิมา จาปิ, อุตฺตรา อถ ทกฺขิณา. \\ สรชา อรชา จปิ, สีตา อุณฺหา จ เอกทา. \\ อธิมตฺตา ปริตฺตา จ, ปุถู วายนฺติ มาลุตา. \\ ตเถวิมสฺมิํ กายสฺมิํ, สมุปฺปชฺชนฺติ เวทนา. \\ สุขทุกฺขสมุปฺปตฺติ, อทุกฺขมสุขา จ ยา. \\ ยโต จ ภิกฺขุ อาตาปี, สมฺปชญฺญํ น ริญฺจติ. \\ ตโต โส เวทนา สพฺพา, ปริชานาติ ปณฺฑิโต. \\ โส เวทนา ปริญฺญาย, ทิฏฺเฐ ธมฺเม อนาสโว. \\ กายสฺส เภทา ธมฺมฏฺโฐ, สงฺขฺยํ โนเปติ เวทคูติ. ทุติยํ.}
\csromangathagroup{\csromangathastanza{\csromanfolio{419}ยถาปิ วาตา อากาเส, วายนฺติ วิวิธา ปุถู. \\ ปุรตฺถิมา ปจฺฉิมา จาปิ, อุตฺตรา อถ ทกฺขิณา.}\csromangathastanza{สรชา อรชา จปิ, สีตา อุณฺหา จ เอกทา. \\ อธิมตฺตา ปริตฺตา จ, ปุถู วายนฺติ มาลุตา.}\csromangathastanza{ตเถวิมสฺมิํ กายสฺมิํ, สมุปฺปชฺชนฺติ เวทนา. \\ สุขทุกฺขสมุปฺปตฺติ, อทุกฺขมสุขา จ ยา.}\csromangathastanza{ยโต จ ภิกฺขุ อาตาปี, สมฺปชญฺญํ น ริญฺจติ\footnote{สมฺปชาโน นิรูปธิ (ก)}. \\ ตโต โส เวทนา สพฺพา, ปริชานาติ ปณฺฑิโต.}\csromangathastanza{โส เวทนา ปริญฺญาย, ทิฏฺเฐ ธมฺเม อนาสโว. \\ กายสฺส เภทา ธมฺมฏฺโฐ, สงฺขฺยํ โนเปติ เวทคูติ.\csromanspacer{} ทุติยํ.}}

\csromanheader{2}{3. ทุติย-อากาสสุตฺต}
\proseitem{261}{เสยฺยถาปิ ภิกฺขเว อากาเส วิวิธา วาตา วายนฺติ, ปุรตฺถิมาปิ วาตา วายนฺติ ฯเปฯ อธิมตฺตาปิ วาตา วายนฺติ, เอวเมว โข ภิกฺขเว อิมสฺมิํ กายสฺมิํ วิวิธา เวทนา อุปฺปชฺชนฺติ, สุขาปิ เวทนา อุปฺปชฺชติ, ทุกฺขาปิ เวทนา อุปฺปชฺชติ, อทุกฺขมสุขาปิ เวทนา อุปฺปชฺชตีติ.\csromanspacer{}\csromanspacer{}ตติยํ.}

\csromanheader{2}{4. อคารสุตฺต}
\csromanmatikaanchor{mtk.231}
\csromantocmarkref{subsection}{4. อคารสุตฺต}{mtk.231}
\bookmark[dest={mtk.231},level=2]{4. อคารสุตฺต}
\proseitem{262}{เสยฺยถาปิ ภิกฺขเว อาคนฺตุกาคารํ, ตตฺถ ปุรตฺถิมายปิ ทิสาย อาคนฺตฺวา วาสํ กปฺเปนฺติ, ปจฺฉิมายปิ ทิสาย อาคนฺตฺวา วาสํ กปฺเปนฺติ, อุตฺตรายปิ ทิสาย อาคนฺตฺวา วาสํ กปฺเปนฺติ, ทกฺขิณายปิ ทิสาย อาคนฺตฺวา วาสํ กปฺเปนฺติ, ขตฺติยาปิ อาคนฺตฺวา วาสํ กปฺเปนฺติ, พฺราหฺมณาปิ อาคนฺตฺวา วาสํ กปฺเปนฺติ, เวสฺสาปิ อาคนฺตฺวา วาสํ กปฺเปนฺติ, สุทฺทาปิ อาคนฺตฺวา วาสํ กปฺเปนฺติ, เอวเมว โข ภิกฺขเว อิมสฺมิํ กายสฺมิํ วิวิธา}

\gambhira{สฬายตนวคฺคสํยุตฺตปาฬิ}
\prose{\csromanfolio{420}เวทนา อุปฺปชฺชนฺติ, สุขาปิ เวทนา อุปฺปชฺชติ, ทุกฺขาปิ เวทนา อุปฺปชฺชติ, อทุกฺขมสุขาปิ เวทนา อุปฺปชฺชติ.\csromanspacer{} สามิสาปิ สุขา เวทนา อุปฺปชฺชติ, สามิสาปิ ทุกฺขา เวทนา อุปฺปชฺชติ, สามิสาปิ อทุกฺขมสุขา เวทนา อุปฺปชฺชติ.\csromanspacer{} นิรามิสาปิ สุขา เวทนา อุปฺปชฺชติ, นิรามิสาปิ ทุกฺขา เวทนา อุปฺปชฺชติ, นิรามิสาปิ อทุกฺขมสุขา เวทนา อุปฺปชฺชตีติ.\csromanspacer{}\csromanspacer{}จตุตฺถํ.}

\csromanheader{2}{5. ปฐม-อานนฺทสุตฺต}
\csromanmatikaanchor{mtk.232}
\csromantocmarkref{subsection}{5-6. อานนฺทสุตฺต}{mtk.232}
\bookmark[dest={mtk.232},level=2]{5-6. อานนฺทสุตฺต}
\proseitem{263}{อถ โข อายสฺมา อานนฺโท เยน ภควา เตนุปสงฺกมิ, อุปสงฺกมิตฺวา เอกมนฺตํ นิสีทิ, เอกมนฺตํ นิสินฺโน โข อายสฺมา อานนฺโท ภควนฺตํ เอตทโวจ “กตมา นุ โข ภนฺเต เวทนา, กตโม เวทนาสมุทโย, กตโม เวทนานิโรโธ, กตมา เวทนานิโรธคามินี ปฏิปทา, โก เวทนาย อสฺสาโท, โก อาทีนโว, กิํ นิสฺสรณนฺ”ติ.\csromanspacer{} ติสฺโส อิมา อานนฺท เวทนา สุขา เวทนา ทุกฺขา เวทนา อทุกฺขมสุขา เวทนา, อิมา วุจฺจนฺติ อานนฺท เวทนา.\csromanspacer{} ผสฺสสมุทยา เวทนาสมุทโย, ผสฺสนิโรธา เวทนานิโรโธ.\csromanspacer{} อยเมว อริโย อฏฺฐงฺคิโก มคฺโค เวทนานิโรธคามินี ปฏิปทา.\csromanspacer{} เสยฺยถิทํ, สมฺมาทิฏฺฐิ ฯเปฯ สมฺมาสมาธิ.\csromanspacer{} ยํ เวทนํ ปฏิจฺจ อุปฺปชฺชติ สุขํ โสมนสฺสํ, อยํ เวทนาย อสฺสาโท.\csromanspacer{} ยํ เวทนา อนิจฺจา ทุกฺขา วิปริณามธมฺมา, อยํ เวทนาย อาทีนโว.\csromanspacer{} โย เวทนาย ฉนฺทราควินโย ฉนฺทราคปฺปหานํ, อิทํ เวทนาย นิสฺสรณํ.\csromanspacer{} อถ โข ปนานนฺท มยา อนุปุพฺพสงฺขารานํ นิโรโธ อกฺขาโต.\csromanspacer{} ปฐมํ ฌานํ สมาปนฺนสฺส วาจา นิรุทฺธา โหติ ฯเปฯ สญฺญาเวทยิตนิโรธํ สมาปนฺนสฺส สญฺญา จ เวทนา จ นิรุทฺธา โหนฺติ.\csromanspacer{} ขีณาสวสฺส ภิกฺขุโน ราโค นิรุทฺโธ โหติ, โทโส นิรุทฺโธ โหติ, โมโห นิรุทฺโธ โหติ.\csromanspacer{} อถ โข ปนานนฺท มยา อนุปุพฺพสงฺขารานํ วูปสโม อกฺขาโต.\csromanspacer{} ปฐมํ ฌานํ สมาปนฺนสฺส วาจา วูปสนฺตา โหติ ฯเปฯ สญฺญาเวทยิตนิโรธํ สมาปนฺนสฺส สญฺญา จ เวทนา จ วูปสนฺตา โหนฺติ.\csromanspacer{} ขีณาสวสฺส ภิกฺขุโน ราโค วูปสนฺโต โหติ, โทโส วูปสนฺโต โหติ, โมโห วูปสนฺโต โหติ.\csromanspacer{} อถ โข ปนานนฺท มยา อนุปุพฺพสงฺขารานํ ปฏิปฺปสฺสทฺธิ อกฺขาตา.\csromanspacer{} ปฐมํ ฌานํ สมาปนฺนสฺส วาจา ปฏิปฺปสฺสทฺธา โหติ ฯเปฯ.\csromanspacer{} อากาสานญฺจายตนํ สมาปนฺนสฺส รูปสญฺญา ปฏิปฺปสฺสทฺธา โหติ, วิญฺญาณญฺจายตนํ สมาปนฺนสฺส อากาสานญฺจายตนสญฺญา ปฏิปฺปสฺสทฺธา โหติ, อากิญฺจญฺญายตนํ}

\vspace{\tipitakaparskip}
\csromancenter{เวทนาสํยุตฺต สมาปนฺนสฺส วิญฺญาณญฺจายตนสญฺญา ปฏิปฺปสฺสทฺธา โหติ, เนวสญฺญานาสญฺญายตนํ สมาปนฺนสฺส อากิญฺจญฺญายตนสญฺญา ปฏิปฺปสฺสทฺธา โหติ, สญฺญาเวทยิตนิโรธํ สมาปนฺนสฺส สญฺญา จ เวทนา จ ปฏิปฺปสฺสทฺธา โหนฺติ.\csromanspacer{} ขีณาสวสฺส ภิกฺขุโน ราโค ปฏิปฺปสฺสทฺโธ โหติ, โทโส ปฏิปฺปสฺสทฺโธ โหติ, โมโห ปฏิปฺปสฺสทฺโธ โหตีติ.\csromanspacer{}\csromanspacer{}ปญฺจมํ.}
\csromanheader{2}{6. ทุติย-อานนฺทสุตฺต}
\csromanmatikaanchor{mtk.233}
\csromantocmarkref{subsection}{7-8. สมฺพหุลสุตฺต}{mtk.233}
\bookmark[dest={mtk.233},level=2]{7-8. สมฺพหุลสุตฺต}
\proseitem{264}{\csromanfolio{421}อถ โข อายสฺมา อานนฺโท เยน ภควา เตนุปสงฺกมิ, อุปสงฺกมิตฺวา ภควนฺตํ อภิวาเทตฺวา เอกมนฺตํ นิสีทิ, เอกมนฺตํ นิสินฺนํ โข อายสฺมนฺตํ อานนฺทํ ภควา เอตทโวจ “กตมา นุ โข อานนฺท เวทนา, กตโม เวทนาสมุทโย, กตโม เวทนานิโรโธ, กตมา เวทนานิโรธคามินี ปฏิปทา, โก เวทนาย อสฺสาโท, โก อาทีนโว, กิํ นิสฺสรณนฺ”ติ.\csromanspacer{} ภควํมูลกา โน ภนฺเต ธมฺมา, ภควํเนตฺติกา, ภควํปฏิสรณา.\csromanspacer{} สาธุ ภนฺเต ภควนฺตญฺเญว ปฏิภาตุ เอตสฺส ภาสิตสฺส อตฺโถ, ภควโต สุตฺวา ภิกฺขู ธาเรสฺสนฺตีติ.\csromanspacer{} เตน หิ อานนฺท สุโณหิ สาธุกํ มนสิ กโรหิ, ภาสิสฺสามีติ.\csromanspacer{} เอวํ ภนฺเตติ โข อายสฺมา อานนฺโท ภควโต ปจฺจสฺโสสิ.\csromanspacer{} ภควา เอตทโวจ “ติสฺโส อิมา อานนฺท เวทนา, สุขา เวทนา ทุกฺขา เวทนา อทุกฺขมสุขา เวทนา, อิมา วุจฺจนฺติ อานนฺท เวทนา ฯเปฯ.\csromanspacer{} ผสฺสสมุทยา ฯเปฯ.\csromanspacer{} ขีณาสวสฺส ภิกฺขุโน ราโค ปฏิปฺปสฺสทฺโธ โหติ, โทโส ปฏิปฺปสฺสทฺโธ โหติ, โมโห ปฏิปฺปสฺสทฺโธ โหตี”ติ.\csromanspacer{}\csromanspacer{}ฉฏฺฐํ.}

\csromanheader{2}{7. ปฐมสมฺพหุลสุตฺต}
\proseitem{265}{อถ โข สมฺพหุลา ภิกฺขู เยน ภควา เตนุปสงฺกมิํสุ, อุปสงฺกมิตฺวา ภควนฺตํ อภิวาเทตฺวา เอกมนฺตํ นิสีทิํสุ, เอกมนฺตํ นิสินฺนา โข เต ภิกฺขู ภควนฺตํ เอตทโวจุํ “กตมา นุ โข ภนฺเต เวทนา, กตโม เวทนาสมุทโย, กตโม เวทนานิโรโธ, กตมา เวทนานิโรธคามินี ปฏิปทา, โก เวทนาย อสฺสาโท, โก อาทีนโว, กิํ นิสฺสรณนฺ”ติ.\csromanspacer{} ติสฺโส อิมา ภิกฺขเว เวทนา, สุขา เวทนา ทุกฺขา เวทนา อทุกฺขมสุขา เวทนา, อิมา วุจฺจนฺติ ภิกฺขเว เวทนา.}

\gambhira{สฬายตนวคฺคสํยุตฺตปาฬิ}
\prose{\csromanfolio{422}ผสฺสสมุทยา เวทนาสมุทโย, ผสฺสนิโรธา เวทนานิโรโธ, อยเมว อริโย อฏฺฐงฺคิโก มคฺโค เวทนานิโรธคามินี ปฏิปทา.\csromanspacer{} เสยฺยถิทํ, สมฺมาทิฏฺฐิ ฯเปฯ สมฺมาสมาธิ.\csromanspacer{} ยํ เวทนํ ปฏิจฺจ อุปฺปชฺชติ สุขํ โสมนสฺสํ, อยํ เวทนาย อสฺสาโท.\csromanspacer{} ยํ เวทนา อนิจฺจา ทุกฺขา วิปริณามธมฺมา, อยํ เวทนาย อาทีนโว.\csromanspacer{} โย เวทนาย ฉนฺทราควินโย ฉนฺทราคปฺปหานํ, อิทํ เวทนาย นิสฺสรณํ.}

\prose{อถ โข ปน ภิกฺขเว มยา อนุปุพฺพสงฺขารานํ นิโรโธ อกฺขาโต, ปฐมํ ฌานํ สมาปนฺนสฺส วาจา นิรุทฺธา โหติ ฯเปฯ.\csromanspacer{} ขีณาสวสฺส ภิกฺขุโน ราโค นิรุทฺโธ โหติ, โทโส นิรุทฺโธ โหติ, โมโห นิรุทฺโธ โหติ.\csromanspacer{} อถ โข ปน ภิกฺขเว มยา อนุปุพฺพสงฺขารานํ วูปสโม อกฺขาโต.\csromanspacer{} ปฐมํ ฌานํ สมาปนฺนสฺส วาจา วูปสนฺตา โหติ ฯเปฯ.\csromanspacer{} ขีณาสวสฺส ภิกฺขุโน ราโค วูปสนฺโต โหติ, โทโส วูปสนฺโต โหติ, โมโห วูปสนฺโต โหติ.\csromanspacer{} ฉยิมา ภิกฺขเว ปสฺสทฺธิโย.\csromanspacer{} ปฐมํ ฌานํ สมาปนฺนสฺส วาจา ปฏิปฺปสฺสทฺธา โหติ, ทุติยํ ฌานํ สมาปนฺนสฺส วิตกฺกวิจารา ปฏิปฺปสฺสทฺธา โหนฺติ, ตติยํ ฌานํ สมาปนฺนสฺส ปีติ ปฏิปฺปสฺสทฺธา โหติ, จตุตฺถํ ฌานํ สมาปนฺนสฺส อสฺสาสปสฺสาสา ปฏิปฺปสฺสทฺธา โหนฺติ, สญฺญาเวทยิตนิโรธํ สมาปนฺนสฺส อสฺสาสปสฺสาสา ปฏิปฺปสฺสทฺธา โหนฺติ, สญฺญาเวทยิตนิโรธํ สมาปนฺนสฺส สญฺญา จ เวทนา จ ปฏิปฺปสฺสทฺธา โหนฺติ.\csromanspacer{} ขีณาสวสฺส ภิกฺขุโน ราโค ปฏิปฺปสฺสทฺโธ โหติ, โทโส ปฏิปฺปสฺสทฺโธ โหติ, โมโห ปฏิปฺปสฺสทฺโธ โหติติ.\csromanspacer{}\csromanspacer{}สตฺตมํ.}

\csromanheader{2}{8. ทุติยสมฺพหุลสุตฺต}
\proseitem{266}{อถ โข สมฺพหุลา ภิกฺขู เยน ภควา เตนุปสงฺกมิํสุ ฯเปฯ เอกมนฺตํ นิสินฺนา โข เต ภิกฺขู ภควา เอตทโวจ “กตมา นุ โข ภิกฺขเว เวทนา, กตโม เวทนาสมุทโย, กตโม เวทนานิโรโธ, กตมา เวทนานิโรธคามินี ปฏิปทา, โก เวทนาย อสฺสาโท, โก อาทีนโว, กิํ นิสฺสรณนฺ”ติ.\csromanspacer{} ภควํมูลกา โน ภนฺเต ธมฺมา ฯเปฯ.\csromanspacer{} ติสฺโส อิมา ภิกฺขเว เวทนา สุขา เวทนา ทุกฺขา เวทนา อทุกฺขมสุขา เวทนา, อิมา วุจฺจนฺติ ภิกฺขเว เวทนา.\csromanspacer{} ผสฺสสมุทยา. (ยถา ปุริมสุตฺตนฺเต, ตถา วิตฺถาเรตพฺโพ.).\csromanspacer{} อฏฺฐมํ.}

\csromanheader{2}{2. เวทนาสํยุตฺต \textbf{9. ปญฺจกงฺคสุตฺต}}
\csromanmatikaanchor{mtk.234}
\csromantocmarkref{subsection}{9. ปญฺจกงฺคสุตฺต}{mtk.234}
\bookmark[dest={mtk.234},level=2]{9. ปญฺจกงฺคสุตฺต}
\proseitem{267}{\csromanfolio{423}อถ โข ปญฺจกงฺโค ถปติ เยนายสฺมา อุทายี เตนุปสงฺกมิ, อุปสงฺกมิตฺวา อายสฺมนฺตํ อุทายิํ อภิวาเทตฺวา เอกมนฺตํ นิสีทิ, เอกมนฺตํ นิสินฺโน โข ปญฺจกงฺโค ถปติ อายสฺมนฺตํ อุทายิํ เอตทโวจ “กติ นุ โข ภนฺเต อุทายิ เวทนา วุตฺตา ภควตา”ติ.\csromanspacer{} ติสฺโส โข ถปติ เวทนา วุตฺตา ภควตา, สุขา เวทนา ทุกฺขา เวทนา อทุกฺขมสุขา เวทนา, อิมา โข ถปติ ติสฺโส เวทนา วุตฺตา ภควตาติ.\csromanspacer{} เอวํ วุตฺเต ปญฺจกงฺโค ถปติ อายสฺมนฺตํ อุทายิํ เอตทโวจ “น โข ภนฺเต อุทายิ ติสฺโส เวทนา วุตฺตา ภควตา, เทฺว เวทนา วุตฺตา ภควตา สุขา เวทนา ทุกฺขา เวทนา, ยายํ ภนฺเต อทุกฺขมสุขา เวทนา, สนฺตสฺมิํ เอสา ปณีเต สุเข วุตฺตา ภควตา”ติ.\csromanspacer{} ทุติยมฺปิ โข อายสฺมา อุทายี ปญฺจกงฺคํ ถปติํ เอตทโวจ “น โข ถปติ เทฺว เวทนา วุตฺตา ภควตา, ติสฺโส เวทนา วุตฺตา ภควตา สุขา เวทนา ทุกฺขา เวทนา อทุกฺขมสุขา เวทนา, อิมา ติสฺโส เวทนา วุตฺตา ภควตา”ติ.\csromanspacer{} ทุติยมฺปิ โข ปญฺจกงฺโค ถปติ อายสฺมนฺตํ อุทายิํ เอตทโวจ “น โข ภนฺเต อุทายิ ติสฺโส เวทนา วุตฺตา ภควตา, เทฺว เวทนา วุตฺตา ภควตา สุขา เวทนา ทุกฺขา เวทนา, ยายํ ภนฺเต อทุกฺขมสุขา เวทนา, สนฺตสฺมิํ เอสา ปณีเต สุเข วุตฺตา ภควตา”ติ.\csromanspacer{} ตติยมฺปิ โข อายสฺมา อุทายี ปญฺจกงฺคํ ถปติํ เอตทโวจ “น โข ถปติ เทฺว เวทนา วุตฺตา ภควตา, ติสฺโส เวทนา วุตฺตา ภควตา สุขา เวทนา ทุกฺขา เวทนา อทุกฺขมสุขา เวทนา, อิมา ติสฺโส เวทนา วุตฺตา ภควตา”ติ.\csromanspacer{} ตติยมฺปิ โข ปญฺจกงฺโค ถปติ อายสฺมนฺตํ อุทายิํ เอตทโวจ “น โข ภนฺเต อุทายิ ติสฺโส เวทนา วุตฺตา ภควตา, เทฺว เวทนา วุตฺตา ภควตา สุขา เวทนา ทุกฺขา เวทนา, ยายํ ภนฺเต อทุกฺขมสุขา เวทนา, สนฺตสฺมิํ เอสา ปณีเต สุเข วุตฺตา ภควตา”ติ.\csromanspacer{} เนว สกฺขิ อายสฺมา อุทายี ปญฺจกงฺคํ ถปติํ สญฺญาเปตุํ, น ปนาสกฺขิ ปญฺจกงฺโค ถปติ อายสฺมนฺตํ อุทายิํ สญฺญาเปตุํ, อสฺโสสิ โข อายสฺมา อานนฺโท อายสฺมโต อุทายิสฺส ปญฺจกงฺเคน ถปตินา สทฺธิํ อิมํ กถาสลฺลาปํ.}

\prose{อถ โข อายสฺมา อานนฺโท เยน ภควา เตนุปสงฺกมิ, อุปสงฺกมิตฺวา เอกมนฺตํ นิสีทิ, เอกมนฺตํ นิสินฺโน โข อายสฺมา อานนฺโท}

\gambhira{สฬายตนวคฺคสํยุตฺตปาฬิ}
\prose{\csromanfolio{424}ยาวตโก อายสฺมโต อุทายิสฺส ปญฺจกงฺเคน ถปตินา สทฺธิํ อโหสิ กถาสลฺลาโป, ตํ สพฺพํ ภควโต อาโรเจสิ.\csromanspacer{} สนฺตเมว อานนฺท ปริยายํ ปญฺจกงฺโค ถปติ อุทายิสฺส ภิกฺขุโน นาพฺภนุโมทิ, สนฺตญฺจ ปนานนฺท ปริยายํ อุทายี ภิกฺขุ ปญฺจกงฺคสฺส ถปติโน นาพฺภนุโมทิ.\csromanspacer{} เทฺวปิ มยา อานนฺท เวทนา วุตฺตา ปริยาเยน, ติสฺโสปิ มยา เวทนา วุตฺตา ปริยาเยน, ปญฺจปิ มยา เวทนา วุตฺตา ปริยาเยน, ฉปิ มยา เวทนา วุตฺตา ปริยาเยน, อฏฺฐารสาปิ มยา เวทนา วุตฺตา ปริยาเยน ฉตฺติํสาปิ มยา เวทนา วุตฺตา ปริยาเยน, อฏฺฐสตมฺปิ มยา เวทนา วุตฺตา ปริยาเยน.\csromanspacer{} เอวํ ปริยายเทสิโต โข อานนฺท มยา ธมฺโม, เอวํ ปริยายเทสิเต โข อานนฺท มยา ธมฺเม เย อญฺญมญฺญสฺส สุภาสิตํ สุลปิตํ น สมนุมญฺญิสฺสนฺติ, น สมนุชานิสฺสนฺติ, น สมนุโมทิสฺสนฺติ, เตสํ เอตํ ปาฏิกงฺขํ “ภณฺฑนชาตา กลหชาตา วิวาทาปนฺโน อญฺญมญฺญํ มุขสตฺตีหิ วิตุทนฺตา วิหริสฺสนฺตี”ติ\footnote{วิหริสฺสนฺติ (สี, อิ, ก)}.\csromanspacer{} เอวํ ปริยายเทสิโต โข อานนฺท มยา ธมฺโม, เอวํ ปริยาย เทสิเต โข อานนฺท มยา ธมฺเม, เย อญฺญมญฺญสฺส สุภาสิตํ สุลปิตํ สมนุมญฺญิสฺสนฺติ, สมนุชานิสฺสนฺติ, สมนุโมทิสฺสนฺติ, เตสํ เอตํ ปาฏิกงฺขํ “สมคฺคา สมฺโมทมานา อวิวทมานา ขีโรทกีภูตา อญฺญมญฺญํ ปิยจกฺขูหิ สมฺปสฺสนฺตา วิหรสฺสนฺตี”ติ.}

\prose{ปญฺจิเม อานนฺท กามคุณา.\csromanspacer{} กตเม ปญฺจ, จกฺขุวิญฺเญยฺยา รูปา อิฏฺฐา กนฺตา มนาปา ปิยรูปา กามูปสํหิตา รชนียา ฯเปฯ กายวิญฺเญยฺยา โผฏฺฐพฺพา อิฏฺฐา กนฺตา มนาปา ปิยรูปา กามูปสํหิตา รชนียา, อิเม โข อานนฺท ปญฺจ กามคุณา.\csromanspacer{} ยํ โข อานนฺท อิเม ปญฺจ กามคุเณ ปฏิจฺจ อุปฺปชฺชติ สุขํ โสมนสฺสํ, อิทํ วุจฺจติ กามสุขํ.\csromanspacer{} เย โข อานนฺท เอวํ วเทยฺยุํ “เอตํ ปรมํ สนฺตํ สุขํ โสมนสฺสํ ปฏิสํเวเทนฺตี”ติ, อิทํ เนสาหํ นานุชานามิ.\csromanspacer{} ตํ กิสฺส เหตุ, อตฺถานนฺท เอตมฺหา สุขา อญฺญํ สุขํ อภิกฺกนฺตตรญฺจ ปณีตตรญฺจ.}

\prose{กตมญฺจานนฺท เอตมฺหา สุขา อญฺญํ สุขํ อภิกฺกนฺตตรญฺจ ปณีตตรญฺจ.\csromanspacer{} อิธานนฺท ภิกฺขุ วิวิจฺเจว กาเมหิ วิวิจฺจ อกุสเลหิ ธมฺเมหิ สวิตกฺกํ สวิจารํ วิเวกชํ ปีติสุขํ ปฐมํ ฌานํ อุปสมฺปชฺช วิหรติ, อิทํ โข อานนฺท เอตมฺหา สุขา อญฺญํ สุขํ อภิกฺกนฺตตรญฺจ ปณีตตรญฺจ.\csromanspacer{} เย โข อานนฺท}

\vspace{\tipitakaparskip}
\csromancenter{เวทนาสํยุตฺต เอวํ วเทยฺยุํ “เอตํปรมํ สนฺตํ สุขํ โสมนสฺสํ ปฏิสํเวเทนฺตี”ติ, อิทํ เนสาหํ นานุชานามิ.\csromanspacer{} ตํ กิสฺส เหตุ, อตฺถานนฺท เอตมฺหา สุขา อญฺญํ สุขํ อภิกฺกนฺตตรญฺจ ปณีตตรญฺจ.}
\prose{\csromanfolio{425}กตมญฺจานนฺท เอตมฺหา สุขา อญฺญํ สุขํ อภิกฺกนฺตตรญฺจ ปณีตตรญฺจ.\csromanspacer{} อิธานนฺท ภิกฺขุ วิตกฺกวิจารานํ วูปสมา อชฺฌตฺตํ สมฺปสาทนํ เจตโส เอโกทิภาวํ อวิตกฺกํ อวิจารํ สมาธิชํ ปีติสุขํ ทุติยํ ฌานํ อุปสมฺปชฺช วิหรติ, อิทํ โข อานนฺท เอตมฺหา สุขา อญฺญํ สุขํ อภิกฺกนฺตตรญฺจ ปณีตตรญฺจ.\csromanspacer{} เย โข อานนฺท เอวํ วเทยฺยุํ “เอตํปรมํ สนฺตํ สุขํ โสมนสฺสํ ปฏิสํเวเทนฺตี”ติ, อิทํ เนสาหํ นานุชานามิ.\csromanspacer{} ตํ กิสฺส เหตุ, อตฺถานนฺท เอตมฺหา สุขา อญฺญํ สุขํ อภิกฺกนฺตตรญฺจ ปณีตตรญฺจ.}

\prose{กตมญฺจานนฺท เอตมฺหา สุขา อญฺญํ สุขํ อภิกฺกนฺตรญฺจ ปญีตตรญฺจ.\csromanspacer{} อิธานนฺท ภิกฺขุ ปีติยา จ วิราคา อุเปกฺขโก จ วิหรติ สโต จ สมฺปชาโน, สุขญฺจ กาเยน ปฏิสํเวเทติ, ยํ ตํ อริยา อาจิกฺขนฺติ “อุเปกฺขโก สติมา สุขวิหารี”ติ, ตติยํ ฌานํ อุปสมฺปชฺช วิหรติ, อิทํ โข อานนฺท เอตมฺหา สุขา อญฺญํ สุขํ อภิกฺกนฺตตรญฺจ ปณีตตรญฺจ.\csromanspacer{} เย โข อานนฺท เอวํ วเทยฺยุํ “เอตํปรมํ สนฺตํ สุขํ โสมนสฺสํ ปฏิสํเวเทนฺตี”ติ, อิทํ เนสาหํ นานุชานามิ.\csromanspacer{} ตํ กิสฺส เหตุ, อตฺถานนฺท เอตมฺหา สุขา อญฺญํ สุขํ อภิกฺกนฺตตรญฺจ ปณีตตรญฺจ.}

\prose{กตมญฺจานนฺท เอตมฺหา สุขา อญฺญํ สุขํ อภิกฺกนฺตตรญฺจ ปณีตตรญฺจ.\csromanspacer{} อิธานนฺท ภิกฺขุ สุขสฺส จ ปหานา ทุกฺขสฺส จ ปหานา ปุพฺเพว โสมนสฺสโทมนสฺสานํ อตฺถงฺคมา อทุกฺขมสุขํ อุเปกฺขาสติปาริสุทฺธิํ จตุตฺถํ ฌานํ อุปสมฺปชฺช วิหรติ, อิทํ โข อานนฺท เอตมฺหา สุขา อญฺญํ สุขํ อภิกฺกนฺตตรญฺจ ปณีตตรญฺจ.\csromanspacer{} เย โข อานนฺท เอวํ วเทยฺยุํ “เอตํปรมํ สนฺตํ สุขํ โสมนสฺสํ ปฏิสํเวเทนฺตี”ติ, อิทํ เนสาหํ นานุชานามิ.\csromanspacer{} ตํ กิสฺส เหตุ, อตฺถานนฺท เอตมฺหา สุขา อญฺญํ สุขํ อภิกฺกนฺตตรญฺจ ปณีตตรญฺจ.}

\prose{กตมญฺจานนฺท เอตมฺหา สุขา อญฺญํ สุขํ อภิกฺกนฺตตรญฺจ ปณีตตรญฺจ.\csromanspacer{} อิธานนฺท ภิกฺขุ สพฺพโส รูปสญฺญานํ สมติกฺกมา ปฏิฆสญฺญานํ อตฺถงฺคมา นานตฺตสญฺญานํ อมนสิการา “อนนฺโต อากาโส”ติ อากาสานญฺจยตนํ อุปสมฺปชฺช วิหรติ, อิทํ โข อานนฺท เอตมฺหา สุขา อญฺญํ สุขํ อภิกฺกนฺตตรญฺจ ปณีตตรญฺจ.\csromanspacer{} เย โข อานนฺท เอวํ วเทยฺยุํ “เอตํปรมํ}

\gambhira{สฬายตนวคฺคสํยุตฺตปาฬิ}
\prose{\csromanfolio{426}สนฺตํ สุขํ โสมนสฺสํ ปฏิสํเวเทนฺตี”ติ, อิทํ เนสาหํ นานุชานามิ.\csromanspacer{} ตํ กิสฺส เหตุ, อตฺถานนฺท เอตมฺหา สุขา อญฺญํ สุขํ อภิกฺกนฺตตรญฺจ ปณีตตรญฺจ.}

\prose{กตมญฺจานนฺท เอตมฺหา สุขา อญฺญํ สุขํ อภิกฺกนฺตตรญฺจ ปณีตตรญฺจ.\csromanspacer{} อิธานนฺท ภิกฺขุ สพฺพโส อากาสานญฺจายตนํ สมติกฺกมฺม “อนนฺตํ วิญฺญาณนฺ”ติ วิญฺญาณญฺจายตนํ อุปสมฺปชฺช วิหรติ, อิทํ โข อานนฺท เอตมฺหา สุขา อญฺญํ สุขํ อภิกฺกนฺตตรญฺจ ปณีตตรญฺจ.\csromanspacer{} เย โข อานนฺท เอวํ วเทยฺยุํ “เอตํปรมํ สนฺตํ สุขํ โสมนสฺสํ ปฏิสํเวเทนฺตี”ติ, อิทํ เนสาหํ นานุชานามิ.\csromanspacer{} ตํ กิสฺส เหตุ, อตฺถานนฺท เอตมฺหา สุขา อญฺญํ สุขํ อภิกฺกนฺตตรญฺจ ปณีตตรญฺจ.}

\prose{กตมญฺจานนฺท เอตมฺหา สุขา อญฺญํ สุขํ อภิกฺกนฺตตรญฺจ ปณีตตรญฺจ.\csromanspacer{} อิธานนฺท ภิกฺขุ สพฺพโส วิญฺญาณญฺจายตนํ สมติกฺกมฺม “นตฺถิ กิญฺจี”ติ อากิญฺจญฺญายตนํ อุปสมฺปชฺช วิหรติ, อิทํ โข อานนฺท เอตมฺหา สุขา อญฺญํ สุขํ อภิกฺกนฺตตรญฺจ ปณีตตรญฺจ.\csromanspacer{} เย โข อานนฺท เอวํ วเทยฺยุํ “เอตํปรมํ สนฺตํ สุขํ โสมนสฺสํ ปฏิสํเวเทนฺตี”ติ, อิทํ เนสาหํ นานุชานามิ.\csromanspacer{} ตํ กิสฺส เหตุ, อตฺถานนฺท เอตมฺหา สุขา อญฺญํ สุขํ อภิกฺกนฺตตรญฺจ ปณีตตรญฺจ.}

\prose{กตมญฺจานนฺท เอตมฺหา สุขา อญฺญํ สุขํ อภิกฺกนฺตตรญฺจ ปณีตตรญฺจ.\csromanspacer{} อิธานนฺท ภิกฺขุ สพฺพโส อากิญฺจญฺญายตนํ สมติกฺกมฺม เนวสญฺญานาสญฺญายตนํ อุปสมฺปชฺช วิหรติ, อิทํ โข อานนฺท เอตมฺหา สุขา อญฺญํ สุขํ อภิกฺกนฺตตรญฺจ ปณีตตรญฺจ.\csromanspacer{} เย โข อานนฺท เอวํ วเทยฺยุํ “เอตํปรมํ สนฺตํ สุขํ โสมนสฺสํ ปฏิสํเวเทนฺตี”ติ, อิทํ เนสาหํ นานุชานามิ.\csromanspacer{} ตํ กิสฺส เหตุ, อตฺถานนฺท เอตมฺหา สุขา อญฺญํ สุขํ อภิกฺกนฺตตรญฺจ ปณีตตรญฺจ.}

\prose{กตมญฺจานนฺท เอตมฺหา สุขา อญฺญํ สุขํ อภิกฺกนฺตตรญฺจ ปณีตตรญฺจ.\csromanspacer{} อิธานนฺท ภิกฺขุ สพฺพโส เนวสญฺญานาสญฺญายตนํ สมติกฺกมฺม สญฺญาเวทยิตนิโรธํ อุปสมฺปชฺช วิหรติ, อิทํ โข อานนฺท เอตมฺหา สุขา อญฺญํ สุขํ อภิกฺกนฺตตรญฺจ ปณีตตรญฺจ.}

\prose{ฐานํ โข ปเนตํ อานนฺท วิชฺชติ, ยํ อญฺญติตฺถิยา ปริพฺพาชกา เอวํ วเทยฺยุํ “สญฺญาเวทยิตนิโรธํ สมโณ โคตโม อาห, ตญฺจ}

\vspace{\tipitakaparskip}
\csromancenter{เวทนาสํยุตฺต สุขสฺมิํ ปญฺญเปติ, ตยิทํ กิํสุ, ตยิทํ กถํสู”ติ.\csromanspacer{} เอวํวาทิโน อานนฺท อญฺญติตฺถิยา ปริพฺพาชกา เอวมสฺสุ วจนียา “น โข อาวุโส ภควา สุขญฺเญว เวทนํ สนฺธาย สุขสฺมิํ ปญฺญเปติ, ยตฺถ ยตฺถ อาวุโส สุขํ อุปลพฺภติ ยหิํ ยหิํ\footnote{ยํ หิ ยํ สุขํ (สี, อิ), ยหิํ ยหิํ สุขํ (สฺยา, กํ, ก) ม 2. 62 ปิฏฺเฐปิ.}, ตํ ตํ ตถาคโต สุขสฺมิํ ปญฺญเปตี”ติ.\csromanspacer{}\csromanspacer{}นวมํ.}
\csromanheader{2}{10. ภิกฺขุสุตฺต}
\csromanmatikaanchor{mtk.235}
\csromantocmarkref{subsection}{10. ภิกฺขุสุตฺต}{mtk.235}
\bookmark[dest={mtk.235},level=2]{10. ภิกฺขุสุตฺต}
\proseitem{268}{\csromanfolio{427}เทฺวปิ มยา ภิกฺขเว เวทนา วุตฺตา ปริยาเยน, ติสฺโสปิ มยา เวทนา วุตฺตา ปริยาเยน, ปญฺจปิ มยา เวทนา วุตฺตา ปริยาเยน, ฉปิ มยา เวทนา วุตฺตา ปริยาเยน, อฏฺฐารสาปิ มยา เวทนา วุตฺตา ปริยาเยน, ฉตฺติํสาปิ มยา เวทนา วุตฺตา ปริยาเยน, อฏฺฐสตมฺปิ มยา เวทนา วุตฺตา ปริยาเยน.\csromanspacer{} เอวํ ปริยายเทสิโต ภิกฺขเว มยา ธมฺโม, เอวํ ปริยายเทสิเต โข ภิกฺขเว มยา ธมฺเม เย อญฺญมญฺญสฺส สุภาสิตํ สุลปิตํ น สมนุมญฺญิสฺสนฺติ, น สมนุชานิสฺสนฺติ, น สมนุโมทิสฺสนฺติ, เตสํ เอตํ ปาฏิกงฺขํ “ภณฺฑนชาตา กลหชาตา วิวาทาปนฺนา อญฺญมญฺญํ มุขสตฺตีหิ วิตุทนฺตา วิหริสฺสนฺตี”ติ.\csromanspacer{} เอวํ ปริยายเทสิโต ภิกฺขเว มยา ธมฺโม, เอวํ ปริยายเทสิเต โข ภิกฺขเว มยา ธมฺเม เย อญฺญมญฺญสฺส สุภาสิตํ สุลปิตํ สมนุมญฺญิสฺสนฺติ, สมนุชานิสฺสนฺติ, สมนุโมทิสฺสนฺติ, เตสํ เอตํ ปาฏิกงฺขํ “สมคฺคา สมฺโมทมานา อวิวทมานา ขีโรทกีภูตา อญฺญมญฺญํ ปิยจกฺขูหิ สมฺปสฺสนฺตา วิหริสฺสนฺตี”ติ.}

\prose{ปญฺจิเม ภิกฺขเว กามคุณา ฯเปฯ.\csromanspacer{} ฐานํ โข ปเนตํ ภิกฺขเว วิชฺชติ, ยํ อญฺญติตฺถิยา ปริพฺพาชกา เอวํ วเทยฺยุํ “สญฺญาเวทยิตนิโรธํ สมโณ โคตโม อาห, ตญฺจ สุขสฺมิํ ปญฺญเปติ, ตยิทํ กิํสุ, ตยิทํ กถํสู”ติ.\csromanspacer{} เอวํวาทิโน ภิกฺขเว อญฺญติตฺติยา ปริพฺพาชกา เอวมสฺสุ วจนียา “น โข อาวุโส ภควา สุขญฺเญว เวทนํ สนฺธาย สุขสฺมิํ ปญฺญเปติ, ยตฺถ ยตฺถ อาวุโส สุขํ อุปลพฺภติ ยหิํ ยหิํ\footnote{ยํ หิ ยํ หิ (สี, อิ)}, ตํ ตํ ตถาคโต สุขสฺมิํ ปญฺญเปตี”ติ.\csromanspacer{}\csromanspacer{}ทสมํ.\csromanspacer{} รโหคตวคฺโค ทุติโย.}

\gambhira{สฬายตนวคฺคสํยุตฺตปาฬิ}
\titlehead{\textbf{ตสฺสุทฺทานํ}}
\csromangathameasure{รโหคตํ เทฺว อากาสํ, อคารํ เทฺว จ อานนฺทา. \\ สมฺพหุลา ทุเว วุตฺตา, ปญฺจกงฺโค จ ภิกฺขุนาติ.}
\csromangathagroup{\csromangathastanza{\csromanfolio{428}รโหคตํ เทฺว อากาสํ, อคารํ เทฺว จ อานนฺทา. \\ สมฺพหุลา ทุเว วุตฺตา, ปญฺจกงฺโค จ ภิกฺขุนาติ.}}

\csromanheader{1}{3. อฏฺฐสตปริยายวคฺค}
\csromanmatikaanchor{mtk.236}
\csromantocmarknopage{section}{3. อฏฺฐสตปริยายวคฺค}{mtk.236}
\bookmark[dest={mtk.236},level=1]{3. อฏฺฐสตปริยายวคฺค}
\csromanheader{2}{1. สีวกสุตฺต}
\csromanmatikaanchor{mtk.237}
\csromantocmarkref{subsection}{1. สีวกสุตฺต}{mtk.237}
\bookmark[dest={mtk.237},level=2]{1. สีวกสุตฺต}
\proseitem{269}{เอกํ สมยํ ภควา ราชคเห วิหรติ เวฬุวเน กลนฺทกนิวาเป.\csromanspacer{} อถ โข โมฬิยสีวโก ปริพฺพาชโก เยน ภควา เตนุปสงฺกมิ, อุปสงฺกมิตฺวา ภควตา สทฺธิํ สมฺโมทิ, สมฺโมทนียํ กถํ สารณียํ วีติสาเรตฺวา เอกมนฺตํ นิสีทิ, เอกมนฺตํ นิสินฺโน โข โมฬิยสีวโก ปริพฺพาชโก ภควนฺตํ เอตทโวจ “สนฺติ โภ โคตม เอเก สมณพฺราหฺมณา เอวํวาทิโน เอวํทิฏฺฐิโน ‘ยํ กิญฺจายํ ปุริสปุคฺคโล ปฏิสํเวเทติ สุขํ วา ทุกฺขํ วา อทุกฺขมสุขํ วา, สพฺพํ ตํ ปุพฺเพกตเหตู’ติ, อิธ\footnote{อิธ ปน (สฺยา, กํ, อิ, ก)} ภวํ โคตโม กิมาหา”ติ.}

\prose{ปิตฺตสมุฏฺฐานานิปิ โข สีวก อิเธกจฺจานิ เวทยิตานิ อุปฺปชฺชนฺติ, สามมฺปิ โข เอตํ สีวก เวทิตพฺพํ\footnote{เอวํ เวทิตพฺพํ (สฺยา, กํ, ก)}, ยถา ปิตฺตสมุฏฺฐานานิปิ อิเธกจฺจานิ เวทยิตานิ อุปฺปชฺชนฺติ.\csromanspacer{} โลกสฺสปิ โข เอตํ สีวก สจฺจสมฺมตํ, ยถา ปิตฺตสมุฏฺฐานานิปิ อิเธกจฺจานิ เวทยิตานิ อุปฺปชฺชนฺติ.\csromanspacer{} ตตฺร สีวก เย เต สมณพฺราหฺมณา เอวํวาทิโน เอวํทิฏฺฐิโน ‘ยํ กิญฺจายํ ปุริสปุคฺคโล ปฏิสํเวเทติ สุขํ วา ทุกฺขํ วา อทุกฺขมสุขํ วา, สพฺพํ ตํ ปุพฺเพกตเหตู’ติ.\csromanspacer{} ยญฺจ สามํ ญาตํ, ตญฺจ อติธาวนฺติ, ยญฺจ โลเก สจฺจสมฺมตํ, ตญฺจ อติธาวนฺติ, ตสฺมา เตสํ สมณพฺราหฺมณานํ มิจฺฉาติ วทามิ.}

\prose{เสมฺหสมุฏฺฐานานิปิ โข สีวก ฯเปฯ.\csromanspacer{} วาตสมุฏฺฐานานิปิ โข สีวก ฯเปฯ.\csromanspacer{} สนฺนิปาติกานิปิ โข สีวก ฯเปฯ.\csromanspacer{} อุตุปริณามชานิปิ โข สีวก ฯเปฯ.\csromanspacer{} วิสมปริหารชานิปิ โข สีวก ฯเปฯ.\csromanspacer{} โอปกฺกมิกานิปิ โข สีวก ฯเปฯ.\csromanspacer{} กมฺมวิปากชานิปิ โข สีวก อิเธกจฺจานิ เวทยิตานิ อุปฺปชฺชนฺติ, สามมฺปิ โข เอตํ สีวก เวทิตพฺพํ, ยถา กมฺมวิปากชานิปิ อิเธกจฺจานิ เวทยิตานิ อุปฺปชฺชนฺติ.\csromanspacer{} โลกสฺสปิ โข เอตํ สีวก สจฺจสมฺมตํ, ยถา กมฺมวิปากชานิปิ}

\vspace{\tipitakaparskip}
\csromancenter{เวทนาสํยุตฺต อิเธกจฺจานิ เวทยิตานิ อุปฺปชฺชนฺติ.\csromanspacer{} ตตฺร สีวก เย เต สมณพฺราหฺมณา เอวํวาทิโน เอวํทิฏฺฐิโน ‘ยํ กิญฺจายํ ปุริสปุคฺคโล ปฏิสํเวเทติ สุขํ วา ทุกฺขํ วา อทุกฺขมสุขํ วา, สพฺพํ ตํ ปุพฺเพกตเหตู’ติ.\csromanspacer{} ยญฺจ สามํ ญาตํ, ตญฺจ อติธาวนฺติ, ยญฺจ โลเก สจฺจสมฺมตํ, ตญฺจ อติธาวนฺติ, ตสฺมา เตสํ สมณพฺราหฺมณานํ มิจฺฉาติ วทามีติ.\csromanspacer{} เอวํ วุตฺเต โมฬิยสีวโก ปริพฺพาชโก ภควนฺตํ เอตทโวจ “อภิกฺกนฺตํ โภ โคตม, อภิกฺกนฺตํ โภ โคตม ฯเปฯ.\csromanspacer{} อุปาสกํ มํ ภวํ โคตโม ธาเรตุ อชฺชตคฺเค ปาณุเปตํ สรณํ คตนฺ”ติ.}
\vspace{0.5\baselineskip}
\csromanmatikaanchor{mtk.238}
\csromantocmarkref{subsection}{2-10. อฏฺฐสตสุตฺตาทิ}{mtk.238}
\bookmark[dest={mtk.238},level=2]{2-10. อฏฺฐสตสุตฺตาทิ}
\csromangathameasure{ปิตฺตํ เสมฺหญฺจ วาโต จ, สนฺนิปาตา อุตูนิ จ. \\ วิสมํ โอปกฺกมิกํ, กมฺมวิปาเกน อฏฺฐมีติ. ปฐมํ.}
\csromangathagroup{\csromangathastanza{\csromanfolio{429}ปิตฺตํ เสมฺหญฺจ วาโต จ, สนฺนิปาตา อุตูนิ จ. \\ วิสมํ โอปกฺกมิกํ, กมฺมวิปาเกน อฏฺฐมีติ.\csromanspacer{} ปฐมํ.}}

\csromanheader{2}{2. อฏฺฐสตสุตฺต}
\proseitem{270}{อฏฺฐสตปริยายํ โว ภิกฺขเว ธมฺมปริยายํ เทเสสฺสามิ, ตํ สุณาถ.\csromanspacer{} กตโม จ ภิกฺขเว อฏฺฐสตปริยาโย ธมฺมปริยาโย.\csromanspacer{} เทฺวปิ มยา ภิกฺขเว เวทนา วุตฺตา ปริยาเยน, ติสฺโสปิ มยา เวทนา วุตฺตา ปริยาเยน, ปญฺจปิ มยา เวทนา วุตฺตา ปริยาเยน, ฉปิ มยา เวทนา วุตฺตา ปริยาเยน, อฏฺฐารสาปิ มยา เวทนา วุตฺตา ปริยาเยน, ฉตฺติํสาปิ มยา เวทนา วุตฺตา ปริยาเยน, อฏฺฐสตมฺปิ มยา เวทนา วุตฺตา ปริยาเยน.\csromanspacer{} กตมา จ ภิกฺขเว เทฺว เวทนา, กายิกา จ เจตสิกา จ, อิมา วุจฺจนฺติ ภิกฺขเว เทฺว เวทนา.\csromanspacer{} กตมา จ ภิกฺขเว ติสฺโส เวทนา, สุขา เวทนา ทุกฺขา เวทนา อทุกฺขมสุขา เวทนา, อิมา วุจฺจนฺติ ภิกฺขเว ติสฺโส เวทนา.\csromanspacer{} กตมา จ ภิกฺขเว ปญฺจ เวทนา, สุขินฺทฺริยํ ทุกฺขินฺทฺริยํ โสมนสฺสินฺทฺริยํ โทมนสฺสินฺทฺริยํ อุเปกฺขินฺทฺริยํ, อิมา วุจฺจนฺติ ภิกฺขเว ปญฺจ เวทนา.\csromanspacer{} กตมา จ ภิกฺขเว ฉ เวทนา, จกฺขุสมฺผสฺสชา เวทนา ฯเปฯ มโนสมฺผสฺสชา เวทนา, อิมา วุจฺจนฺติ ภิกฺขเว ฉ เวทนา.\csromanspacer{} กตมา จ ภิกฺขเว อฏฺฐารส เวทนา, ฉ โสมนสฺสูปวิจารา ฉ โทมนสฺสูปวิจารา ฉ อุเปกฺขูปวิจารา, อิมา วุจฺจนฺติ ภิกฺขเว อฏฺฐารส เวทนา.\csromanspacer{} กตมา จ ภิกฺขเว ฉตฺติํส เวทนา, ฉ เคหสิตานิ\footnote{เคหสฺสิตานิ (ฐ)} โสมนสฺสานิ, ฉ เนกฺขมฺมสิตานิ\footnote{เนกฺขมฺมสฺสิตานิ (ฐ)} โสมนสฺสานิ, ฉ เคหสิตานิ โทมนสฺสานิ, ฉ}

\gambhira{สฬายตนวคฺคสํยุตฺตปาฬิ}
\prose{\csromanfolio{430}เนกฺขมฺมสิตานิ โทมนสฺสานิ, ฉ เคหสิตา อุเปกฺขา, ฉ เนกฺขมฺมสิตา อุเปกฺขา, อิมา วุจฺจนฺติ ภิกฺขเว ฉตฺติํส เวทนา.\csromanspacer{} กตมญฺจ ภิกฺขเว อฏฺฐสตํ เวทนา, อตีตา ฉตฺติํส เวทนา, อนาคตา ฉตฺติํส เวทนา, ปจฺจุปฺปนฺนา ฉตฺติํส เวทนา, อิมา วุจฺจนฺติ ภิกฺขเว อฏฺฐสตํ เวทนา.\csromanspacer{} อยํ ภิกฺขเว อฏฺฐสตปริยาโย ธมฺมปริยาโยติ.\csromanspacer{}\csromanspacer{}ทุติยํ.}

\csromanheader{2}{3. อญฺญตรภิกฺขุสุตฺต}
\proseitem{271}{อถ โข อญฺญตโร ภิกฺขุ เยน ภควา เตนุปสงฺกมิ, อุปสงฺกมิตฺวา ภควนฺตํ อภิวาเทตฺวา เอกมนฺตํ นิสีทิ, เอกมนฺตํ นิสินฺโน โข โส ภิกฺขุ ภควนฺตํ เอตทโวจ “กตมา นุ โข ภนฺเต เวทนา, กตโม เวทนาสมุทโย, กตมา เวทนาสมุทยคามินี ปฏิปทา, กตโม เวทนานิโรโธ, กตมา เวทนานิโรธคามินี ปฏิปทา, โก เวทนาย อสฺสาโท, โก อาทีนโว, กิํ นิสฺสรณนฺ”ติ.}

\prose{ติสฺโส อิมา ภิกฺขุ เวทนา, สุขา เวทนา ทุกฺขา เวทนา อทุกฺขมสุขา เวทนา, อิมา วุจฺจนฺติ ภิกฺขุ เวทนา.\csromanspacer{} ผสฺสสมุทยา เวทนาสมุทโย.\csromanspacer{} ตณฺหา เวทนาสมุทยคามินี ปฏิปทา.\csromanspacer{} ผสฺสนิโรธา เวทนานิโรโธ.\csromanspacer{} อยเมว อริโย อฏฺฐงฺคิโก มคฺโค เวทนานิโรธคามินี ปฏิปทา.\csromanspacer{} เสยฺยถิทํ, สมฺมาทิฏฺฐิ ฯเปฯ สมฺมาสมาธิ.\csromanspacer{} ยํ เวทนํ ปฏิจฺจ อุปฺปชฺชติ สุขํ โสมนสฺสํ, อยํ เวทนาย อสฺสาโท.\csromanspacer{} ยํ เวทนา อนิจฺจา ทุกฺขา วิปริณามธมฺมา, อยํ เวทนาย อาทีนโว.\csromanspacer{} โย เวทนาย ฉนฺทราควินโย ฉนฺทราคปฺปหานํ, อิทํ เวทนาย นิสฺสรณนฺติ.\csromanspacer{}\csromanspacer{}ตติยํ.}

\csromanheader{2}{4. ปุพฺพสุตฺต}
\proseitem{272}{ปุพฺเพว เม ภิกฺขเว สมฺโพธา อนภิสมฺพุทฺธสฺส โพธิสตฺตสฺเสว สโต เอตทโหสิ “กตมา นุ โข เวทนา, กตโม เวทนาสมุทโย, กตมา เวทนาสมุทยคามินี ปฏิปทา, กตโม เวทนานิโรโธ, กตมา เวทนานิโรธคามินี ปฏิปทา, โก เวทนาย อสฺสาโท, โก อาทีนโว, กิํ นิสฺสรณนฺ”ติ.\csromanspacer{} ตสฺส มยฺหํ ภิกฺขเว เอตทโหสิ “ติสฺโส อิมา เวทนา, สุขา เวทนา ทุกฺขา เวทนา อทุกฺขมสุขา เวทนา, อิมา วุจฺจนฺติ เวทนา.\csromanspacer{} ผสฺสสมุทยา เวทนาสมุทโย.\csromanspacer{} ตณฺหา เวทนาสมุทยคามินี ปฏิปทา ฯเปฯ.\csromanspacer{} โย เวทนาย ฉนฺทราควินโย ฉนฺทราคปฺปหานํ, อิทํ เวทนาย นิสฺสรณนฺ”ติ.\csromanspacer{}\csromanspacer{}จตุตฺถํ.}

\csromanheader{2}{2. เวทนาสํยุตฺต \textbf{5. ญาณสุตฺต}}
\proseitem{273}{\csromanfolio{431}อิมา เวทนาติ เม ภิกฺขเว ปุพฺเพ อนนุสฺสุเตสุ ธมฺเมสุ จกฺขุํ อุทปาทิ, ญาณํ อุทปาทิ, ปญฺญา อุทปาทิ, วิชฺชา อุทปาทิ, อาโลโก อุทปาทิ.\csromanspacer{} อยํ เวทนาสมุทโยติ เม ภิกฺขเว ปุพฺเพ อนนุสฺสุเตสุ ธมฺเมสุ จกฺขุํ อุทปาทิ ฯเปฯ อาโลโก อุทปาทิ.\csromanspacer{} อยํ เวทนาสมุทยคามินี ปฏิปทาติ เม ภิกฺขเว ปุพฺเพ อนนุสฺสุเตสุ ธมฺเมสุ จกฺขุํ อุทปาทิ ฯเปฯ.\csromanspacer{} อยํ เวทนานิโรโธติ เม ภิกฺขเว ปุพฺเพ อนนุสฺสุเตสุ ธมฺเมสุ จกฺขุํ อุทปาทิ ฯเปฯ.\csromanspacer{} อยํ เวทนานิโรธคามินี ปฏิปทาติ เม ภิกฺขเว ปุพฺเพ อนนุสฺสุเตสุ ธมฺเมสุ จกฺขุํ อุทปาทิ ฯเปฯ.\csromanspacer{} อยํ เวทนาย อสฺสาโทติ เม ภิกฺขเว ปุพฺเพ อนนุสฺสุเตสุ ธมฺเมสุ ฯเปฯ.\csromanspacer{} อยํ เวทนาย อาทีนโวติ เม ภิกฺขเว ปุพฺเพ อนนุสฺสุเตสุ ธมฺเมสุ ฯเปฯ.\csromanspacer{} อิทํ โข นิสฺสรณนฺติ เม ภิกฺขเว ปุพฺเพ อนนุสฺสุเตสุ ธมฺเมสุ จกฺขุํ อุทปาทิ, ญาณํ อุทปาทิ, ปญฺญา อุทปาทิ, วิชฺชา อุทปาทิ, อาโลโก อุทปาทีติ.\csromanspacer{}\csromanspacer{}ปญฺจมํ.}

\csromanheader{2}{6. สมฺพหุลภิกฺขุสุตฺต}
\proseitem{274}{อถ โข สมฺพหุลา ภิกฺขู เยน ภควา เตนุปสงฺกมิํสุ, อุปสงฺกมิตฺวา ฯเปฯ เอกมนฺตํ นิสินฺนา โข เต ภิกฺขู ภควนฺตํ เอตทโวจุํ “กตมา นุ โข ภนฺเต เวทนา, กตโม เวทนาสมุทโย, กตมา เวทนาสมุทคามินี ปฏิปทา, กตโม เวทนานิโรโธ, กตมา เวทนานิโรธคามินี ปฏิปทา, โก เวทนาย อสฺสาโท, โก อาทีนโว, กิํ นิสฺสรณนฺ”ติ.\csromanspacer{} ติสฺโส อิมา ภิกฺขเว เวทนา, สุขา เวทนา ทุกฺขา เวทนา อทุกฺขมสุขา เวทนา, อิมา วุจฺจนฺติ ภิกฺขเว เวทนา.\csromanspacer{} ผสฺสสมุทยา เวทนาสมุทโย.\csromanspacer{} ตณฺหา เวทนาสมุทยคามินี ปฏิปทา.\csromanspacer{} ผสฺสนิโรธา ฯเปฯ.\csromanspacer{} โย เวทนาย ฉนฺทราควินโย ฉนฺทราคปฺปหานํ, อิทํ เวทนาย นิสฺสรณนฺติ.\csromanspacer{}\csromanspacer{}ฉฏฺฐํ.}

\csromanheader{2}{7. ปฐมสมณพฺราหฺมณสุตฺต}
\proseitem{275}{ติสฺโส อิมา ภิกฺขเว เวทนา.\csromanspacer{} กตมา ติสฺโส, สุขา เวทนา ทุกฺขา เวทนา อทุกฺขมสุขา เวทนา.\csromanspacer{} เย หิ เกจิ ภิกฺขเว}

\gambhira{สฬายตนวคฺคสํยุตฺตปาฬิ}
\prose{\csromanfolio{432}สมณา วา พฺราหฺมณา วา อิมาสํ ติสฺสนฺนํ เวทนานํ สมุทยญฺจ อตฺถงฺคมญฺจ อสฺสาทญฺจ อาทีนวญฺจ นิสฺสรณญฺจ ยถาภูตํ นปฺปชานนฺติ, น เม เต ภิกฺขเว สมณา วา พฺราหฺมณา วา สมเณสุ วา สมณสมฺมตา พฺราหฺมเณสุ วา พฺราหฺมณสมฺมตา, น จ ปน เต อายสฺมนฺโต สามญฺญตฺถํ วา พฺรหฺมญฺญตฺถํ วา ทิฏฺเฐว ธมฺเม สยํ อภิญฺญา สจฺฉิกตฺวา อุปสมฺปชฺช วิหรนฺติ.\csromanspacer{} เย จ โข เกจิ ภิกฺขเว สมณา วา พฺราหฺมณา วา อิมาสํ ติสฺสนฺนํ เวทนานํ สมุทยญฺจ อตฺถงฺคมญฺจ อสฺสาทญฺจ อาทีนวญฺจ นิสฺสรณญฺจ ยถาภูตํ ปชานนฺติ, เต โข เม ภิกฺขเว สมณา วา พฺราหฺมณา วา สมเณสุ เจว สมณสมฺมตา พฺราหฺมเณสุ จ พฺราหฺมณสมฺมตา, เต จ ปนายสฺมนฺโต สามญฺญตฺถญฺจ พฺรหฺมญฺญตฺถญฺจ ทิฏฺเฐว ธมฺเม สยํ อภิญฺญา สจฺฉิกตฺวา อุปสมฺปชฺช วิหรนฺตีติ.\csromanspacer{}\csromanspacer{}สตฺตมํ.}

\csromanheader{2}{8. ทุติยสมณพฺราหฺมณสุตฺต}
\proseitem{276}{ติสฺโส อิมา ภิกฺขเว เวทนา, กตมา ติสฺโส, สุขา เวทนา ทุกฺขา เวทนา อทุกฺขมสุขา เวทนา.\csromanspacer{} เย หิ เกจิ ภิกฺขเว สมณา วา พฺราหฺมณา วา อิมาสํ ติสฺสนฺนํ เวทนานํ สมุทยญฺจ อตฺถงฺคมญฺจ อสฺสาทญฺจ อาทีนวญฺจ นิสฺสรณญฺจ ยถาภูตํ นปฺปชานนฺติ ฯเปฯ.\csromanspacer{} ปชานนฺติ ฯเปฯ.\csromanspacer{} สยํ อภิญฺญา สจฺฉิกตฺวา อุปสมฺปชฺช วิหรนฺตีติ.\csromanspacer{}\csromanspacer{}อฏฺฐมํ.}

\csromanheader{2}{9. ตติยสมณพฺราหฺมณสุตฺต}
\proseitem{277}{เย หิ เกจิ ภิกฺขเว สมณา วา พฺราหฺมณา วา เวทนํ นปฺปชานนฺติ, เวทนาสมุทยํ นปฺปชานนฺติ, เวทนานิโรธํ นปฺปชานนฺติ, เวทนานิโรธคามินิํ ปฏิปทํ นปฺปชานนฺติ ฯเปฯ.\csromanspacer{} ปชานนฺติ ฯเปฯ.\csromanspacer{} สยํ อภิญฺญา สจฺฉิกตฺวา อุปสมฺปชฺช วิหรนฺตีติ.\csromanspacer{}\csromanspacer{}นวมํ.}

\csromanheader{2}{10. สุทฺธิกสุตฺต}
\proseitem{278}{ติสฺโส อิมา ภิกฺขเว เวทนา.\csromanspacer{} กตมา ติสฺโส, สุขา เวทนา ทุกฺขา เวทนา อทุกฺขมสุขา เวทนา, อิมา โข ภิกฺขเว ติสฺโส เวทนาติ.\csromanspacer{}\csromanspacer{}ทสมํ.}

\csromanheader{2}{2. เวทนาสํยุตฺต \textbf{11. นิรามิสสุตฺต}}
\csromanmatikaanchor{mtk.239}
\csromantocmarkref{subsection}{11. นิรามิสสุตฺต}{mtk.239}
\bookmark[dest={mtk.239},level=2]{11. นิรามิสสุตฺต}
\proseitem{279}{\csromanfolio{433}อตฺถิ ภิกฺขเว สามิสา ปีติ, อตฺถิ นิรามิสา ปีติ, อตฺถิ นิรามิสา นิรามิสตรา ปีติ, อตฺถิ สามิสํ สุขํ, อตฺถิ นิรามิสํ สุขํ, อตฺถิ นิรามิสา นิรามิสตรํ สุขํ, อตฺถิ สามิสา อุเปกฺขา, อตฺถิ นิรามิสา อุเปกฺขา, อตฺถิ นิรามิสา นิรามิสตรา อุเปกฺขา, อตฺถิ สามิโส วิโมกฺโข, อตฺถิ นิรามิโส วิโมกฺโข, อตฺถิ นิรามิสา นิรามิสตโร วิโมกฺโข.\csromanspacer{} กตมา จ ภิกฺขเว สามิสา ปีติ.\csromanspacer{} ปญฺจิเม ภิกฺขเว กามคุณา.\csromanspacer{} กตเม ปญฺจ, จกฺขุวิญฺเญยฺยา รูปา อิฏฺฐา กนฺตา มนาปา ปิยรูปา กามูปสํหิตา รชนียา ฯเปฯ.\csromanspacer{} กายวิญฺเญยฺยา โผฏฺฐพฺพา อิฏฺฐา กนฺตา มนาปา ปิยรูปา กามูปสํหิตา รชนียา.\csromanspacer{} อิเม โข ภิกฺขเว ปญฺจ กามคุณา.\csromanspacer{} ยา โข ภิกฺขเว อิเม ปญฺจ กามคุเณ ปฏิจฺจ อุปฺปชฺชติ ปีติ, อยํ วุจฺจติ ภิกฺขเว สามิสา ปีติ.}

\prose{กตมา จ ภิกฺขเว นิรามิสา ปีติ.\csromanspacer{} อิธ ภิกฺขเว ภิกฺขุ วิวิจฺเจว กาเมหิ วิวิจฺจ อกุสเลหิ ธมฺเมหิ สวิตกฺกํ สวิจารํ วิเวกชํ ปีติสุขํ ปฐมํ ฌานํ อุปสมฺปชฺช วิหรติ, วิตกฺกวิจารานํ วูปสมา อชฺฌตฺตํ สมฺปสาทนํ เจตโส เอโกทิภาวํ อวิตกฺกํ อวิจารํ สมาธิชํ ปีติสุขํ ทุติยํ ฌานํ อุปสมฺปชฺช วิหรติ, อยํ วุจฺจติ ภิกฺขเว นิรามิสา ปีติ.}

\prose{กตมา จ ภิกฺขเว นิรามิสา นิรามิสตรา ปีติ.\csromanspacer{} ยา โข ภิกฺขเว ขีณาสวสฺส ภิกฺขุโน ราคา จิตฺตํ วิมุตฺตํ ปจฺจเวกฺขโต โทสา จิตฺตํ วิมุตฺตํ ปจฺจเวกฺขโต โมหา จิตฺตํ วิมุตฺตํ ปจฺจเวกฺขโต อุปฺปชฺชติ ปีติ, อยํ วุจฺจติ ภิกฺขเว นิรามิสา นิรามิสตรา ปีติ.}

\prose{กตมญฺจ ภิกฺขเว สามิสํ สุขํ.\csromanspacer{} ปญฺจิเม ภิกฺขเว กามคุณา.\csromanspacer{} กตเม ปญฺจ, จกฺขุวิญฺเญยฺยา รูปา อิฏฺฐา กนฺตา มนาปา ปิยรูปา กามูปสํหิตา รชนียา ฯเปฯ กายวิญฺเญยฺยา โผฏฺฐพฺพา อิฏฺฐา กนฺตา มนาปา ปิยรูปา กามูปสํหิตา รชนียา.\csromanspacer{} อิเม โข ภิกฺขเว ปญฺจ กามคุณา.\csromanspacer{} ยํ โข ภิกฺขเว อิเม ปญฺจ กามคุเณ ปฏิจฺจ อุปฺปชฺชติ สุขํ โสมนสฺสํ, อิทํ วุจฺจติ ภิกฺขเว สามิสํ สุขํ.}

\gambhira{สฬายตนวคฺคสํยุตฺตปาฬิ}
\prose{\csromanfolio{434}กตมญฺจ ภิกฺขเว นิรามิสํ สุขํ.\csromanspacer{} อิธ ภิกฺขเว ภิกฺขุ วิวิจฺเจว กาเมหิ วิวิจฺจ อกุสเลหิ ธมฺเมหิ สวิตกฺกํ สวิจารํ วิเวกชํ ปีติสุขํ ปฐมํ ฌานํ อุปสมฺปชฺช วิหรติ, วิตกฺกวิจารานํ วูปสมา อชฺฌตฺตํ สมฺปสาทนํ เจตโส เอโกทิภาวํ อวิตกฺกํ อวิจารํ สมาธิชํ ปีติสุขํ ทุติยํ ฌานํ อุปสมฺปชฺช วิหรติ, ปีติยา จ วิราคา อุเปกฺขโก จ วิหรติ สโต จ สมฺปชาโน, สุขญฺจ กาเยน ปฏิสํเวเทติ, ยํ ตํ อริยา อาจิกฺขนฺติ “อุเปกฺขโก สติมา สุขวิหารี”ติ, ตติยํ ฌานํ อุปสมฺปชฺช วิหรติ, อิทํ วุจฺจติ ภิกฺขเว นิรามิสํ สุขํ.}

\prose{กตมญฺจ ภิกฺขเว นิรามิสา นิรามิสตรํ สุขํ.\csromanspacer{} ยํ โข ภิกฺขเว ขีณาสวสฺส ภิกฺขุโน ราคา จิตฺตํ วิมุตฺตํ ปจฺจเวกฺขโต โทสา จิตฺตํ วิมุตฺตํ ปจฺจเวกฺขโต โมหา จิตฺตํ วิมุตฺตํ ปจฺจเวกฺขโต อุปฺปชฺชติ สุขํ โสมนสฺสํ, อิทํ วุจฺจติ ภิกฺขเว นิรามิสา นิรามิสตรํ สุขํ.}

\prose{กตมา จ ภิกฺขเว สามิสา อุเปกฺขา.\csromanspacer{} ปญฺจิเม ภิกฺขเว กามคุณา.\csromanspacer{} กตเม ปญฺจ, จกฺขุวิญฺเญยฺยา รูปา อิฏฺฐา กนฺตา มนาปา ปิยรูปา กมูปสํหิตา รชนียา ฯเปฯ กายวิญฺเญยฺยา โผฏฺฐพฺพา อิฏฺฐา กนฺตา มนาปา ปิยรูปา กามุปสํหิตา รชนียา.\csromanspacer{} อิเม โข ภิกฺขเว ปญฺจ กามคุณา.\csromanspacer{} ยา โข ภิกฺขเว อิเม ปญฺจ กามคุเณ ปฏิจฺจ อุปฺปชฺชติ อุเปกฺขา, อยํ วุจฺจติ ภิกฺขเว สามิสา อุเปกฺขา.}

\prose{กตมา จ ภิกฺขเว นิรามิสา อุเปกฺขา.\csromanspacer{} อิธ ภิกฺขเว ภิกฺขุ สุขสฺส จ ปหานา ทุกฺขสฺส จ ปหานา ปุพฺเพว โสมนสฺสโทมสฺสานํ อตฺถงฺคมา อทุกฺขมสุขํ อุเปกฺขาสติปาริสุทฺธิํ จตุตฺถํ ฌานํ อุปสมฺปชฺช วิหรติ, อยํ วุจฺจติ ภิกฺขเว นิรามิสา อุเปกฺขา.}

\prose{กตมา จ ภิกฺขเว นิรามิสา นิรามิสตรา อุเปกฺขา.\csromanspacer{} ยา โข ภิกฺขเว ขีณาสวสฺส ภิกฺขุโน ราคา จิตฺตํ วิมุตฺตํ ปจฺจเวกฺขโต โทสา จิตฺตํ วิมุตฺตํ ปจฺจเวกฺขโต โมหา จิตฺตํ วิมุตฺตํ ปจฺจเวกฺขโต อุปฺปชฺชติ อุเปกฺขา, อยํ วุจฺจติ ภิกฺขเว นิรามิสา นิรามิสตรา อุเปกฺขา.}

\prose{กตโม จ ภิกฺขเว สามิโส วิโมกฺโข.\csromanspacer{} รูปปฺปฏิสํยุตฺโต วิโมกฺโข สามิโส วิโมกฺโข.}

\csromanheader{2}{2. เวทนาสํยุตฺต}
\prose{\csromanfolio{435}กตโม จ ภิกฺขเว นิรามิโส วิโมกฺโข.\csromanspacer{} อรูปปฺปฏิสํยุตฺโต วิโมกฺโข นิรามิโส วิโมกฺโข.}

\prose{กตโม จ ภิกฺขเว นิรามิสา นิรามิสตโร วิโมกฺโข.\csromanspacer{} โย โข ภิกฺขเว ขีณาสวสฺส ภิกฺขุโน ราคา จิตฺตํ วิมุตฺตํ ปจฺจเวกฺขโต โทสา จิตฺตํ วิมุตฺตํ ปจฺจเวกฺขโต โมหา จิตฺตํ วิมุตฺตํ ปจฺจเวกฺขโต อุปฺปชฺชติ วิโมกฺโข, อยํ วุจฺจติ ภิกฺขเว นิรามิสา นิรามิสตโร วิโมกฺโขติ.\csromanspacer{}\csromanspacer{}เอกาทสมํ.\csromanspacer{} อฏฺฐสตปริยายวคฺโค ตติโย.}

\titlehead{\textbf{ตสฺสุทฺทานํ}}
\csromangathameasure{สีวก อฏฺฐสตํ ภิกฺขุ, ปุพฺเพ ญาณญฺจ ภิกฺขุนา. \\ สมณพฺราหฺมณา ตีณิ, สุทฺธิกญฺจ นิรามิสนฺติ. \\ \textbf{เวทนาสํยุตฺตํ สมตฺตํ.}}
\csromangathagroup{\csromangathastanza{สีวก อฏฺฐสตํ ภิกฺขุ, ปุพฺเพ ญาณญฺจ ภิกฺขุนา. \\ สมณพฺราหฺมณา ตีณิ, สุทฺธิกญฺจ นิรามิสนฺติ. \\ \textbf{เวทนาสํยุตฺตํ สมตฺตํ.}}}

\chapterheadpageread{3. มาตุคามสํยุตฺต}
\csromanmatikaanchor{mtk.240}
\csromantocmarknopage{chapter}{3. มาตุคามสํยุตฺต}{mtk.240}
\bookmark[dest={mtk.240},level=0]{3. มาตุคามสํยุตฺต}
\csromanheader{1}{1. ปฐมเปยฺยาลวคฺค}
\csromanmatikaanchor{mtk.241}
\csromantocmarkref{section}{1. ปฐมเปยฺยาลวคฺค}{mtk.241}
\bookmark[dest={mtk.241},level=1]{1. ปฐมเปยฺยาลวคฺค}
\csromanheader{2}{1. มาตุคามสุตฺต}
\proseitem{280}{\csromanfolio{436}ปญฺจหิ ภิกฺขเว องฺเคหิ สมนฺนาคโต มาตุคาโม เอกนฺต-อมนาโป โหติ ปุริสสฺส.\csromanspacer{} กตเมหิ ปญฺจหิ, น จ รูปวา โหติ, น จ โภควา โหติ, น จ สีลวา โหติ, อลโส จ โหติ, ปชญฺจสฺส น ลภติ.\csromanspacer{} อิเมหิ โข ภิกฺขเว ปญฺจหิ องฺเคหิ สมนฺนาคโต มาตุคาโม เอกนฺต-อมนาโป โหติ ปุริสสฺส.\csromanspacer{} ปญฺจหิ ภิกฺขเว องฺเคหิ สมนฺนาคโต มาตุคาโม เอกนฺตมนาโป โหติ ปุริสสฺส.\csromanspacer{} กตเมหิ ปญฺจหิ, รูปวา จ โหติ, โภควา จ โหติ, สีลวา จ โหติ, ทกฺโข จ โหติ อนลโส, ปชญฺจสฺส ลภติ.\csromanspacer{} อิเมหิ โข ภิกฺขเว ปญฺจหิ องฺเคหิ สมนฺนาคโต มาตุคาโม เอกนฺตมนาโป โหติ ปุริสสฺสาติ.\csromanspacer{}\csromanspacer{}ปฐมํ.}

\csromanheader{2}{2. ปุริสสุตฺต}
\proseitem{281}{ปญฺจหิ ภิกฺขเว องฺเคหิ สมนฺนาคโต ปุริโส เอกนฺต-อมนาโป โหติ มาตุคามสฺส.\csromanspacer{} กตเมหิ ปญฺจหิ, น จ รูปวา โหติ, น จ โภควา โหติ, น จ สีลวา โหติ, อลโส จ โหติ, ปชญฺจสฺส น ลภติ.\csromanspacer{} อิเมหิ โข ภิกฺขเว ปญฺจหิ องฺเคหิ สมนฺนาคโต ปุริโส เอกนฺต-อมนาโป โหติ มาตุคามสฺส.\csromanspacer{} ปญฺจหิ ภิกฺขเว องฺเคหิ สมนฺนาคโต ปุริโส เอกนฺตมนาโป โหติ มาตุคามสฺส.\csromanspacer{} กตเมหิ ปญฺจหิ, รูปวา จ โหติ, โภควา จ โหติ, สีลวา จ โหติ, ทกฺโข จ โหติ อนลโส, ปชญฺจสฺส ลภติ.\csromanspacer{} อิเมหิ โข ภิกฺขเว ปญฺจหิ องฺเคหิ สมนฺนาคโต ปุริโส เอกนฺตมนาโป โหติ มาตุคามสฺสาติ.\csromanspacer{}\csromanspacer{}ทุติยํ.}

\csromanheader{2}{3. อาเวณิกทุกฺขสุตฺต}
\proseitem{282}{ปญฺจิมานิ ภิกฺขเว มาตุคามสฺส อาเวณิกานิ ทุกฺขานิ, ยานิ มาตุคาโม ปจฺจนุโภติ อญฺญเตฺรว ปุริเสหิ.\csromanspacer{} กตมานิ ปญฺจ, อิธ}

\vspace{\tipitakaparskip}
\csromancenter{มตุคามสํยุตฺต ภิกฺขเว มาตุคาโม ทหโรว สมาโน ปติกุลํ คจฺฉติ, ญาตเกหิ วินา โหติ, อิทํ ภิกฺขเว มาตุคามสฺส ปฐมํ อาเวณิกํ ทุกฺขํ, ยํ มาตุคาโม ปจฺจนุโภติ อญฺญเตฺรว ปุริเสหิ.\csromanspacer{} ปุน จปรํ ภิกฺขเว มาตุคาโม อุตุนี โหติ, อิทํ ภิกฺขเว มาตุคามสฺส ทุติยํ อาเวณิกํ ทุกฺขํ, ยํ มาตุคาโม ปจฺจนุโภติ อญฺญเตฺรว ปุริเสหิ.\csromanspacer{} ปุน จปรํ ภิกฺขเว มาตุคาโม คพฺภินี โหติ, อิทํ ภิกฺขเว มาตุคามสฺส ตติยํ อาเวณิกํ ทุกฺขํ, ยํ มาตุคาโม ปจฺจนุโภติ อญฺญเตฺรว ปุริเสหิ.\csromanspacer{} ปุน จปรํ ภิกฺขเว มาตุคาโม วิชายติ, อิทํ ภิกฺขเว มาตุคามสฺส จตุตฺถํ อาเวณิกํ ทุกฺขํ, ยํ มาตุคาโม ปจฺจนุโภติ อญฺญเตฺรว ปุริเสหิ.\csromanspacer{} ปุน จปรํ ภิกฺขเว มาตุคาโม ปุริสสฺส ปาริจริยํ อุเปติ, อิทํ โข ภิกฺขเว มาตุคามสฺส ปญฺจมํ อาเวณิกํ ทุกฺขํ, ยํ มาตุคาโม ปจฺจนุโภติ อญฺญเตฺรว ปุริเสหิ.\csromanspacer{} อิมานิ โข ภิกฺขเว ปญฺจ มาตุคามสฺส อาเวณิกานิ ทุกฺขานิ, ยานิ มาตุคาโม ปจฺจนุโภติ อญฺญเตฺรว ปุริเสหีติ.\csromanspacer{}\csromanspacer{}ตติยํ.}
\csromanheader{2}{4. ตีหิธมฺเมหิสุตฺต}
\proseitem{283}{\csromanfolio{437}ตีหิ ภิกฺขเว ธมฺเมหิ สมนฺนาคโต มาตุคาโม เยภุยฺเยน กายสฺส เภทา ปรํ มรณา อปายํ ทุคฺคติํ วินิปาตํ นิรยํ อุปปชฺชติ.\csromanspacer{} กตเมหิ ตีหิ, อิธ ภิกฺขเว มาตุคาโม ปุพฺพณฺหสมยํ มจฺเฉรมลปริยุฏฺฐิเตน เจตสา อคารํ อชฺฌาวสติ, มชฺฌนฺหิกสมยํ อิสฺสาปริยุฏฺฐิเตน เจตสา อคารํ อชฺฌาวสติ, สายนฺหสมยํ กามราคปริยุฏฺฐิเตน เจตสา อคารํ อชฺฌาวสติ.\csromanspacer{} อิเมหิ โข ภิกฺขเว ตีหิ ธมฺเมหิ สมนฺนาคโต มาตุคาโม เยภุยฺเยน กายสฺส เภทา ปรํ มรณา อปายํ ทุคฺคติํ วินิปาตํ นิรยํ อุปปชฺชตีติ.\csromanspacer{}\csromanspacer{}จตุตฺถํ.}

\csromanheader{2}{5. โกธนสุตฺต}
\proseitem{284}{อถ โข อายสฺมา อนุรุทฺโธ เยน ภควา เตนุปสงฺกมิ, อุปสงฺกมิตฺวา เอกมนฺตํ นิสีทิ, เอกมนฺตํ นิสินฺโน โข อายสฺมา อนุรุทฺโธ ภควนฺตํ เอตทโวจ “อิธาหํ ภนฺเต มาตุคามํ ปสฺสามิ ทิพฺเพน จกฺขุนา วิสุทฺเธน อติกฺกนฺตมานุสเกน กายสฺส เภทา ปรํ มรณา อปายํ ทุคฺคติํ วินิปาตํ นิรยํ อุปปชฺชนฺตํ, กติหิ นุ โข ภนฺเต ธมฺเมหิ สมนฺนาคโต}

\gambhira{สฬายตนวคฺคสํยุตฺตปาฬิ}
\prose{\csromanfolio{438}มาตุคาโม กายสฺส เภทา ปรํ มรณา อปายํ ทุคฺคติํ วินิปาตํ นิรยํ อุปปชฺชตี”ติ.}

\prose{ปญฺจหิ โข อนุรุทฺธ ธมฺเมหิ สมนฺนาคโต มาตุคาโม กายสฺส เภทา ปรํ มรณา อปายํ ทุคฺคติํ วินิปาตํ นิรยํ อุปปชฺชติ.\csromanspacer{} กตเมหิ ปญฺจหิ, อสฺสทฺโธ จ โหติ, อหิริโก จ โหติ, อโนตฺตปฺปี จ โหติ, โกธโน จ โหติ, ทุปฺปญฺโญ จ โหติ.\csromanspacer{} อิเมหิ โข อนุรุทฺธ ปญฺจหิ ธมฺเมหิ สมนฺนาคโต มาตุคาโม กายสฺส เภทา ปรํ มรณา อปายํ ทุคฺคติํ วินิปาตํ นิรยํ อุปปชฺชตีติ.\csromanspacer{}\csromanspacer{}ปญฺจมํ.}

\csromanheader{2}{6. อุปนาหีสุตฺต}
\proseitem{285}{ปญฺจหิ อนุรุทฺธ ธมฺเมหิ สมนฺนาคโต มาตุคาโม กายสฺส เภทา ปรํ มรณา อปายํ ทุคฺคติํ วินิปาตํ นิรยํ อุปปชฺชติ.\csromanspacer{} กตเมหิ ปญฺจหิ, อสฺสทฺโธ จ โหติ, อหิริโก จ โหติ, อโนตฺตปฺปี จ โหติ, อุปนาหี จ โหติ, ทุปฺปญฺโญ จ โหติ.\csromanspacer{} อิเมหิ โข อนุรุทฺธ ปญฺจหิ ธมฺเมหิ สมนฺนาคโต มาตุคาโม กายสฺส เภทา ปรํ มรณา อปายํ ทุคฺคติํ วินิปาตํ นิรยํ อุปปชฺชตีติ.\csromanspacer{}\csromanspacer{}ฉฏฺฐํ.}

\csromanheader{2}{7. อิสฺสุกีสุตฺต}
\proseitem{286}{ปญฺจหิ อนุรุทฺธ ธมฺเมหิ สมนฺนาคโต มาตุคาโม กายสฺส เภทา ปรํ มรณา อปายํ ทุคฺคติํ วินิปาตํ นิรยํ อุปปชฺชติ.\csromanspacer{} กตเมหิ ปญฺจหิ, อสฺสทฺโธ จ โหติ, อหิริโก จ โหติ, อโนตฺตปฺปี จ โหติ, อิสฺสุกี จ โหติ, ทุปฺปญฺโญ จ โหติ.\csromanspacer{} อิเมหิ โข อนุรุทฺธ ปญฺจหิ ธมฺเมหิ สมนฺนาคโต มาตุคาโม กายสฺส เภทา ปรํ มรณา อปายํ ทุคฺคติํ วินิปาตํ นิรยํ อุปปชฺชตีติ.\csromanspacer{}\csromanspacer{}สตฺตมํ.}

\csromanheader{2}{8. มจฺฉรีสุตฺต}
\proseitem{287}{ปญฺจหิ อนุรุทฺธ ธมฺเมหิ สมนฺนาคโต มาตุคาโม กายสฺส เภทา ปรํ มรณา อปายํ ทุคฺคติํ วินิปาตํ นิรยํ อุปปชฺชติ.\csromanspacer{} กตเมหิ ปญฺจหิ, อสฺสทฺโธ จ โหติ, อหิริโก จ โหติ, อโนตฺตปฺปี จ โหติ, มจฺฉรี จ โหติ, ทุปฺปญฺโญ จ โหติ.\csromanspacer{} อิเมหิ โข อนุรุทฺธ}

\vspace{\tipitakaparskip}
\csromancenter{มตุคามสํยุตฺต ปญฺจหิ ธมฺเมหิ สมนฺนาคโต มาตุคาโม ฯเปฯ อปายํ ทุคฺคติํ วินิปาตํ นิรยํ อุปปชฺชตีติ.\csromanspacer{}\csromanspacer{}อฏฺฐมํ.}
\csromanheader{2}{9. อติจารีสุตฺต}
\proseitem{288}{\csromanfolio{439}ปญฺจหิ อนุรุทฺธ ธมฺเมหิ สมนฺนาคโต มาตุคาโม ฯเปฯ อปายํ ทุคฺคติํ วินิปาตํ นิรยํ อุปปชฺชติ.\csromanspacer{} กตเมหิ ปญฺจหิ, อสฺสทฺโธ จ โหติ, อหิริโก จ โหติ, อโนตฺตปฺปี จ โหติ, อติจารี จ โหติ, ทุปฺปญฺโญ จ โหติ.\csromanspacer{} อิเมหิ โข อนุรุทฺธ ปญฺจหิ ธมฺเมหิ สมนฺนาคโต มาตุคาโม ฯเปฯ อุปปชฺชตีติ.\csromanspacer{}\csromanspacer{}นวมํ.}

\csromanheader{2}{10. ทุสฺสีลสุตฺต}
\proseitem{289}{ปญฺจหิ อนุรุทฺธ ธมฺเมหิ สมนฺนาคโต มาตุคาโม ฯเปฯ นิรยํ อุปปชฺชติ.\csromanspacer{} กตเมหิ ปญฺจหิ, อสฺสทฺโธ จ โหติ, อหิริโก จ โหติ, อโนตฺตปฺปี จ โหติ, ทุสฺสีโล จ โหติ, ทุปฺปญฺโญ จ โหติ.\csromanspacer{} อิเมหิ โข อนุรุทฺธ ปญฺจหิ ธมฺเมหิ สมนฺนาคโต มาตุคาโม ฯเปฯ นิรยํ อุปปชฺชตีติ.\csromanspacer{}\csromanspacer{}ทสมํ.}

\csromanheader{2}{11. อปฺปสฺสุตสุตฺต}
\proseitem{290}{ปญฺจหิ อนุรุทฺธ ธมฺเมหิ สมนฺนาคโต มาตุคาโม ฯเปฯ นิรยํ อุปปชฺชติ.\csromanspacer{} กตเมหิ ปญฺจหิ, อสฺสทฺโธ จ โหติ, อหิริโก จ โหติ, อโนตฺตปฺปี จ โหติ, อปฺปสฺสุโต จ โหติ, ทุปฺปญฺโญ จ โหติ.\csromanspacer{} อิเมหิ โข อนุรุทฺธ ปญฺจหิ ธมฺเมหิ สมนฺนาคโต มาตุคาโม ฯเปฯ นิรยํ อุปปชฺชตีติ.\csromanspacer{}\csromanspacer{}เอกาทสมํ.}

\csromanheader{2}{12. กุสีตสุตฺต}
\proseitem{291}{ปญฺจหิ อนุรุทฺธ ธมฺเมหิ สมนฺนาคโต มาตุคาโม ฯเปฯ นิรยํ อุปปชฺชติ.\csromanspacer{} กตเมหิ ปญฺจหิ, อสฺสทฺโธ จ โหติ, อหิริโก จ โหติ, อโนตฺตปฺปี จ โหติ, กุสีโต จ โหติ, ทุปฺปญฺโญ จ โหติ.\csromanspacer{} อิเมหิ โข อนุรุทฺธ ปญฺจหิ ธมฺเมหิ สมนฺนาคโต มาตุคาโม อปายํ ทุคฺคติํ วินิปาตํ นิรยํ อุปปชฺชตีติ.\csromanspacer{}\csromanspacer{}ทฺวาทสมํ.}

\gambhira{สฬายตนวคฺคสํยุตฺตปาฬิ}
\csromanheader{2}{13. มุฏฺฐสฺสติสุตฺต}
\proseitem{292}{\csromanfolio{440}ปญฺจหิ อนุรุทฺธ ธมฺเมหิ สมนฺนาคโต มาตุคาโม ฯเปฯ นิรยํ อุปปชฺชติ.\csromanspacer{} กตเมหิ ปญฺจหิ, อสฺสทฺโธ จ โหติ, อหิริโก จ โหติ, อโนตฺตปฺปี จ โหติ, มุฏฺฐสฺสติ จ โหติ, ทุปฺปญฺโญ จ โหติ.\csromanspacer{} อิเมหิ โข อนุรุทฺธ ปญฺจหิ ธมฺเมหิ สมนฺนาคโต มาตุคาโม ฯเปฯ นิรยํ อุปปชฺชตีติ.\csromanspacer{}\csromanspacer{}เตรสมํ.}

\csromanheader{2}{14. ปญฺจเวรสุตฺต}
\proseitem{293}{ปญฺจหิ อนุรุทฺธ ธมฺเมหิ สมนฺนาคโต มาตุคาโม ฯเปฯ นิรยํ อุปปชฺชติ.\csromanspacer{} กตเมหิ ปญฺจหิ, ปาณาติปาตี จ โหติ, อทินฺนาทายี จ โหติ, กาเมสุมิจฺฉาจารี จ โหติ, มุสาวาที จ โหติ, สุราเมรยมชฺชปฺปมาทฏฺฐายี จ โหติ.\csromanspacer{} อิเมหิ โข อนุรุทฺธ ปญฺจหิ ธมฺเมหิ สมนฺนาคโต มาตุคาโม กายสฺส เภทา ปรํ มรณา อปายํ ทุคฺคติํ วินิปาตํ นิรยํ อุปปชฺชตีติ.\csromanspacer{}\csromanspacer{}จุทฺทสมํ.\csromanspacer{} ปฐมเปยฺยาลวคฺโค.}

\titlehead{\textbf{ตสฺสุทฺทานํ}}
\csromangathameasure{มาตุคาโม ปุริโส จ, อาเวณิกา ติธมฺโม จ. \\ โกธโน อุปนาหี จ, อิสฺสุกี มจฺฉเรน จ. \\ อติจารี จ ทุสฺสีโล, อปฺปสฺสุโต จ กุสีโต. \\ มุฏฺฐสฺสติ ปญฺจเวรํ, กณฺหปกฺเข ปกาสิโต.}
\csromangathagroup{\csromangathastanza{มาตุคาโม ปุริโส จ, อาเวณิกา ติธมฺโม จ\footnote{เทฺว มนาปามนาปาจ, อาเวณิกา ตีหิ อนุรุทฺโธ (สพฺพตฺถ)}. \\ โกธโน อุปนาหี จ, อิสฺสุกี มจฺฉเรน จ.}\csromangathastanza{อติจารี จ ทุสฺสีโล, อปฺปสฺสุโต จ กุสีโต. \\ มุฏฺฐสฺสติ ปญฺจเวรํ, กณฺหปกฺเข ปกาสิโต.}}

\csromanheader{1}{2. ทุติยเปยฺยาลวคฺค}
\csromanmatikaanchor{mtk.242}
\csromantocmarkref{section}{2. ทุติยเปยฺยาลวคฺค}{mtk.242}
\bookmark[dest={mtk.242},level=1]{2. ทุติยเปยฺยาลวคฺค}
\csromanheader{2}{1. อกฺโกธนสุตฺต}
\proseitem{294}{อถ โข อายสฺมา อนุรุทฺโธ เยน ภควา เตนุปสงฺกมิ, อุปสงฺกมิตฺวา ฯเปฯ เอกมนฺตํ นิสินฺโน โข อายสฺมา อนุรุทฺโธ ภควนฺตํ เอตทโวจ “อิธาหํ ภนฺเต มาตุคามํ ปสฺสามิ ทิพฺเพน จกฺขุนา วิสุทฺเธน}

\vspace{\tipitakaparskip}
\csromancenter{มตุคามสํยุตฺต อติกฺกนฺตมานุสเกน กายสฺส เภทา ปรํ มรณา สุคติํ สคฺคํ โลกํ อุปปชฺชนฺตํ, กติหิ นุ โข ภนฺเต ธมฺเมหิ สมนฺนาคโต มาตุคาโม กายสฺส เภทา ปรํ มรณา สุคติํ สคฺคํ โลกํ อุปปชฺชตี”ติ.}
\prose{\csromanfolio{441}ปญฺจหิ โข อนุรุทฺธ ธมฺเมหิ สมนฺนาคโต มาตุคาโม กายสฺส เภทา ปรํ มรณา สุคติํ สคฺคํ โลกํ อุปปชฺชติ.\csromanspacer{} กตเมหิ ปญฺจหิ, สทฺโธ จ โหติ, หิริมา จ โหติ, โอตฺตปฺปี จ โหติ, อกฺโกธโน จ โหติ, ปญฺญวา จ โหติ.\csromanspacer{} อิเมหิ โข อนุรุทฺธ ปญฺจหิ ธมฺเมหิ สมนฺนาคโต มาตุคาโม กายสฺส เภทา ปรํ มรณา สุคติํ สคฺคํ โลกํ อุปปชฺชตีติ.\csromanspacer{}\csromanspacer{}ปฐมํ.}

\csromanheader{2}{2. อนุปนาหีสุตฺต}
\proseitem{295}{ปญฺจหิ อนุรุทฺธ ธมฺเมหิ สมนฺนาคโต มาตุคาโม กายสฺส เภทา ปรํ มรณา สุคติํ สคฺคํ โลกํ อุปปชฺชติ.\csromanspacer{} กตเมหิ ปญฺจหิ, สทฺโธ จ โหติ, หิริมา จ โหติ, โอตฺตปฺปี จ โหติ, อนุปนาหี จ โหติ, ปญฺญวา จ โหติ.\csromanspacer{} อิเมหิ โข อนุรุทฺธ ปญฺจหิ ธมฺเมหิ สมนฺนาคโต มาตุคาโม กายสฺส เภทา ปรํ มรณา สุคติํ สคฺคํ โลกํ อุปปชฺชตีติ.\csromanspacer{}\csromanspacer{}ทุติยํ.}

\csromanheader{2}{3. อนิสฺสุกีสุตฺต}
\proseitem{296}{ปญฺจหิ อนุรุทฺธ ธมฺเมหิ สมนฺนาคโต มาตุคาโม กายสฺส เภทา ปรํ มรณา สุคติํ สคฺคํ โลกํ อุปปชฺชติ.\csromanspacer{} กตเมหิ ปญฺจหิ, สทฺโธ จ โหติ, หิริมา จ โหติ, โอตฺตปฺปี จ โหติ, อนิสฺสุกี จ โหติ, ปญฺญวา จ โหติ ฯเปฯ.\csromanspacer{}\csromanspacer{}ตติยํ.}

\csromanheader{2}{4. อมจฺฉรีสุตฺต}
\vspace{\tipitakaparskip}
\csromancenter{อมจฺฉรี จ โหติ, ปญฺญวา จ โหติ ฯเปฯ.\csromanspacer{}\csromanspacer{}จตุตฺถํ.}
\csromanheader{2}{5. อนติจารีสุตฺต}
\vspace{\tipitakaparskip}
\csromancenter{อนติจารี จ โหติ, ปญฺญวา จ โหติ ฯเปฯ.\csromanspacer{}\csromanspacer{}ปญฺจมํ.}
\csromanheader{2}{6. สุสีลสุตฺต}
\proseitem{299}{สีลวา จ โหติ, ปญฺญวา จ โหติ ฯเปฯ.\csromanspacer{}\csromanspacer{}ฉฏฺฐํ.}

\gambhira{สฬายตนวคฺคสํยุตฺตปาฬิ}
\csromanheader{2}{7. พหุสฺสุตสุตฺต}
\vspace{\tipitakaparskip}
\csromancenter{พหุสฺสุโต จ โหติ, ปญฺญวา จ โหติ ฯเปฯ.\csromanspacer{}\csromanspacer{}สตฺตมํ.}
\csromanheader{2}{8. อารทฺธวีริยสุตฺต}
\vspace{\tipitakaparskip}
\csromancenter{อารทฺธวีริโย จ โหติ, ปญฺญวา จ โหติ ฯเปฯ.\csromanspacer{}\csromanspacer{}อฏฺฐมํ.}
\csromanheader{2}{9. อุปฏฺฐิตสฺสติสุตฺต}
\proseitem{302}{\csromanfolio{442}อุปฏฺฐิตสฺสติ จ โหติ, ปญฺญวา จ โหติ ฯเปฯ.\csromanspacer{} อิเมหิ โข อนุรุทฺธ ปญฺจหิ ธมฺเมหิ สมนฺนาคโต มาตุคาโม กายสฺส เภทา ปรํ มรณา สุคติํ สคฺคํ โลกํ อุปปชฺชตีติ.\csromanspacer{}\csromanspacer{}นวมํ.\csromanspacer{}\csromanspacer{}อิเม อฏฺฐ สุตฺตนฺตสงฺเขปา.}

\csromanheader{2}{10. ปญฺจสีลสุตฺต}
\proseitem{303}{ปญฺจหิ อนุรุทฺธ ธมฺเมหิ สมนฺนาคโต มาตุคาโม กายสฺส เภทา ปรํ มรณา สุคติํ สคฺคํ โลกํ อุปปชฺชติ.\csromanspacer{} กตเมหิ ปญฺจหิ, ปาณาติปาตา ปฏิวิรโต จ โหติ, อทินฺนาทานา ปฏิวิรโต จ โหติ, กาเมสุมิจฺฉาจารา ปฏิวิรโต จ โหติ, มุสาวาทา ปฏิวิรโต จ โหติ, สุราเมรยมชฺชปฺปมาทฏฺฐานา ปฏิวิรโต จ โหติ.\csromanspacer{} อิเมหิ โข อนุรุทฺธ ปญฺจหิ ธมฺเมหิ สมนฺนาคโต มาตุคาโม กายสฺส เภทา ปรํ มรณา สุคติํ สคฺคํ โลกํ อุปปชฺชตีติ.\csromanspacer{}\csromanspacer{}ทสมํ.\csromanspacer{} ทุติยเปยฺยาลวคฺโค.}

\titlehead{\textbf{ตสฺสุทฺทานํ}}
\csromangathameasure{ทุติเย จ อกฺโกธโน, อนุปนาหี อนิสฺสุกี. \\ อมจฺฉรี อนติจารี, สีลวา จ พหุสฺสุโต. \\ วีริยํ สติ สีลํ จ, สุกฺกปกฺเข ปกาสิโต.}
\csromangathagroup{\csromangathastanza{ทุติเย จ\footnote{อนุรุทฺโธ (สพฺพตฺถ)} อกฺโกธโน, อนุปนาหี อนิสฺสุกี. \\ อมจฺฉรี อนติจารี, สีลวา จ พหุสฺสุโต. \\ วีริยํ สติ สีลํ จ, สุกฺกปกฺเข ปกาสิโต.}}

\csromanheader{1}{3. มตุคามสํยุตฺต \textbf{3. พลวคฺค}}
\csromanmatikaanchor{mtk.243}
\csromantocmarkref{section}{3. พลวคฺค}{mtk.243}
\bookmark[dest={mtk.243},level=1]{3. พลวคฺค}
\csromanheader{2}{1. วิสารทสุตฺต}
\proseitem{304}{\csromanfolio{443}ปญฺจิมานิ ภิกฺขเว มาตุคามสฺส พลานิ.\csromanspacer{} กตมานิ ปญฺจ, รูปพลํ โภคพลํ ญาติพลํ ปุตฺตพลํ สีลพลํ.\csromanspacer{} อิมานิ โข ภิกฺขเว ปญฺจ มาตุคามสฺส พลานิ.\csromanspacer{} อิเมหิ โข ภิกฺขเว ปญฺจหิ พเลหิ สมนฺนาคโต มาตุคาโม วิสารโท อคารํ อชฺฌาวสตีติ.\csromanspacer{}\csromanspacer{}ปฐมํ.}

\csromanheader{2}{2. ปสยฺหสุตฺต}
\proseitem{305}{ปญฺจิมานิ ภิกฺขเว มาตุคามสฺส พลานิ.\csromanspacer{} กตมานิ ปญฺจ, รูปพลํ โภคพลํ ญาติพลํ ปุตฺตพลํ สีลพลํ.\csromanspacer{} อิมานิ โข ภิกฺขเว ปญฺจ มาตุคามสฺส พลานิ.\csromanspacer{} อิเมหิ โข ภิกฺขเว ปญฺจหิ พเลหิ สมนฺนาคโต มาตุคาโม สามิกํ ปสยฺห อคารํ อชฺฌาวสตีติ.\csromanspacer{}\csromanspacer{}ทุติยํ.}

\csromanheader{2}{3. อภิภุยฺยสุตฺต}
\proseitem{306}{ปญฺจิมานิ ภิกฺขเว มาตุคามสฺส พลานิ.\csromanspacer{} กตมานิ ปญฺจ, รูปพลํ โภคพลํ ญาติพลํ ปุตฺตพลํ สีลพลํ.\csromanspacer{} อิมานิ โข ภิกฺขเว ปญฺจ มาตุคามสฺส พลานิ.\csromanspacer{} อิเมหิ โข ภิกฺขเว ปญฺจหิ พเลหิ สมนฺนาคโต มาตุคาโม สามิกํ อภิภุยฺย วตฺตตีติ.\csromanspacer{}\csromanspacer{}ตติยํ.}

\csromanheader{2}{4. เอกสุตฺต}
\proseitem{307}{เอเกน จ โข ภิกฺขเว พเลน สมนฺนาคโต ปุริโส มาตุคามํ อภิภุยฺย วตฺตติ.\csromanspacer{} กตเมน เอเกน พเลน.\csromanspacer{} อิสฺสริยพเลน อภิภูตํ มาตุคามํ เนว รูปพลํ ตายติ, น โภคพลํ ตายติ, น ญาติพลํ ตายติ, น ปุตฺตพลํ ตายติ, น สีลพลํ ตายตีติ.\csromanspacer{}\csromanspacer{}จตุตฺถํ.}

\csromanheader{2}{5. องฺคสุตฺต}
\proseitem{308}{ปญฺจิมานิ ภิกฺขเว มาตุคามสฺส พลานิ.\csromanspacer{} กตมานิ ปญฺจ, รูปพลํ โภคพลํ ญาติพลํ ปุตฺตพลํ สีลพลํ.\csromanspacer{} รูปพเลน จ ภิกฺขเว มาตุคาโม สมนฺนาคโต โหติ น จ โภคพเลน, เอวํ โส เตนงฺเคน อปริปูโร โหติ.\csromanspacer{} ยโต จ โข ภิกฺขเว มาตุคาโม รูปพเลน จ สมนฺนาคโต โหติ โภคพเลน จ, เอวํ โส เตนงฺเคน ปริปูโร โหติ.}

\gambhira{สฬายตนวคฺคสํยุตฺตปาฬิ}
\prose{\csromanfolio{444}รูปพเลน จ ภิกฺขเว มาตุคาโม สมนฺนาคโต โหติ โภคพเลน จ, น จ ญาติพเลน, เอวํ โส เตนงฺเคน อปริปูโร โหติ.\csromanspacer{} ยโต จ โข ภิกฺขเว มาตุคาโม รูปพเลน จ สมนฺนาคโต โหติ โภคพเลน จ ญาติพเลน จ, เอวํ โส เตนงฺเคน ปริปูโร โหติ.\csromanspacer{} รูปพเลน จ ภิกฺขเว มาตุคาโม สมนฺนาคโต โหติ โภคพเลน จ ญาติพเลน จ, น จ ปุตฺตพเลน, เอวํ โส เตนงฺเคน อปริปูโร โหติ.\csromanspacer{} ยโต จ โข ภิกฺขเว มาตุคาโม รูปพเลน จ สมนฺนาคโต โหติ โภคพเลน จ ญาติพเลน จ ปุตฺตพเลน จ, เอวํ โส เตนงฺเคน ปริปูโร โหติ.\csromanspacer{} รูปพเลน จ ภิกฺขเว มาตุคาโม สมนฺนาคโต โหติ โภคพเลน จ ญาติพเลน จ ปุตฺตพเลน จ, น จ สีลพเลน, เอวํ โส เตนงฺเคน อปริปูโร โหติ.\csromanspacer{} ยโต จ โข ภิกฺขเว มาตุคาโม รูปพเลน จ สมนฺนาคโต โหติ โภคพเลน จ ญาติพเลน จ ปุตฺตพเลน จ สีลพเลน จ, เอวํ โส เตนงฺเคน ปริปูโร โหติ.\csromanspacer{} อิมานิ โข ภิกฺขเว ปญฺจ มาตุคามสฺส พลานีติ.\csromanspacer{}\csromanspacer{}ปญฺจมํ.}

\csromanheader{2}{6. นาเสนฺติสุตฺต}
\proseitem{309}{ปญฺจิมานิ ภิกฺขเว มาตุคามสฺส พลานิ.\csromanspacer{} กตมานิ ปญฺจ, รูปพลํ โภคพลํ ญาติพลํ ปุตฺตพลํ สีลพลํ.\csromanspacer{} รูปพเลน จ ภิกฺขเว มาตุคาโม สมนฺนาคโต โหติ, น จ สีลพเลน, นาเสนฺเตว นํ, กุเล น วาเสนฺติ.\csromanspacer{} รูปพเลน จ ภิกฺขเว มาตุคาโม สมนฺนาคโต โหติ, โภคพเลน จ, น จ สีลพเลน, นาเสนฺเตว นํ, กุเล น วาเสนฺติ.\csromanspacer{} รูปพเลน จ ภิกฺขเว มาตุคาโม สมนฺนาคโต โหติ โภคพเลน จ, ญาติพเลน จ, น จ สีลพเลน, นาเสนฺเตว นํ, กุเล น วาเสนฺติ.\csromanspacer{} รูปพเลน จ ภิกฺขเว มาตุคาโม สมนฺนาคโต โหติ โภคพเลน จ ญาติพเลน จ, ปุตฺตพเลน จ, น จ สีลพเลน, นาเสนฺเตว นํ, กุเล น วาเสนฺติ.}

\prose{สีลพเลน จ ภิกฺขเว มาตุคาโม สมนฺนาคโต โหติ, น จ รูปพเลน, วาเสนฺเตว นํ กุเล, น นาเสนฺติ.\csromanspacer{} สีลพเลน จ ภิกฺขเว มาตุคาโม สมนฺนาคโต โหติ, น จ โภคพเลน, วาเสนฺเตว นํ กุเล, น นาเสนฺติ.\csromanspacer{} สีลพเลน จ ภิกฺขเว มาตุคาโม สมนฺนาคโต โหติ, น จ ญาติพเลน, วาเสนฺเตว นํ กุเล, น นาเสนฺติ.\csromanspacer{} สีลพเลน จ ภิกฺขเว มาตุคาโม สมนฺนาคโต โหติ, น จ ปุตฺตพเลน, วาเสนฺเตว นํ กุเล, น นาเสนฺติ.\csromanspacer{} อิมานิ โข ภิกฺขเว ปญฺจ มาตุคามสฺส พลานีติ.\csromanspacer{}\csromanspacer{}ฉฏฺฐํ.}

\csromanheader{2}{3. มตุคามสํยุตฺต \textbf{7. เหตุสุตฺต}}
\proseitem{310}{\csromanfolio{445}ปญฺจิมานิ ภิกฺขเว มาตุคามสฺส พลานิ.\csromanspacer{} กตมานิ ปญฺจ, รูปพลํ โภคพลํ ญาติพลํ ปุตฺตพลํ สีลพลํ.\csromanspacer{} น ภิกฺขเว มาตุคาโม รูปพลเหตุ วา โภคพลเหตุ วา ญาติพลเหตุ วา ปุตฺตพลเหตุ วา กายสฺส เภทา ปรํ มรณา สุคติํ สคฺคํ โลกํ อุปปชฺชติ, สีลพลเหตุ โข ภิกฺขเว มาตุคาโม กายสฺส เภทา ปรํ มรณา สุคติํ สคฺคํ โลกํ อุปปชฺชติ.\csromanspacer{} อิมานิ โข ภิกฺขเว ปญฺจ มาตุวามสฺส พลานีติ.\csromanspacer{}\csromanspacer{}สตฺตมํ.}

\csromanheader{2}{8. ฐานสุตฺต}
\proseitem{311}{ปญฺจิมานิ ภิกฺขเว ฐานานิ ทุลฺลภานิ อกตปุญฺเญน มาตุคาเมน.\csromanspacer{} กตมานิ ปญฺจ, ปติรูเป กุเล ชาเยยฺยนฺติ, อิทํ ภิกฺขเว ปฐมํ ฐานํ ทุลฺลภํ อกตปุญฺเญน มาตุคาเมน.\csromanspacer{} ปติรูเป กุเล ชายิตฺวา ปติรูปํ กุลํ คจฺเฉยฺยนฺติ, อิทํ ภิกฺขเว ทุติยํ ฐานํ ทุลฺลภํ อกตปุญฺเญน มาตุคาเมน.\csromanspacer{} ปติรูเป กุเล ชายิตฺวา ปติรูปํ กุลํ คนฺตฺวา อสปตฺติ อคารํ อชฺฌาวเสยฺยนฺติ, อิทํ ภิกฺขเว ตติยํ ฐานํ ทุลฺลภํ อกตปุญฺเญน มาตุคาเมน.\csromanspacer{} ปติรูเป กุเล ชายิตฺวา ปติรูปํ กุลํ คนฺตฺวา อสปตฺติ อคารํ อชฺฌาวสนฺตี ปุตฺตวตี อสฺสนฺติ, อิทํ ภิกฺขเว จตุตฺถํ ฐานํ ทุลฺลภํ อกตปุญฺเญน มาตุคาเมน.\csromanspacer{} ปติรูเป กุเล ชายิตฺวา ปติรูปํ กุลํ คนฺตฺวา อสปตฺติ อคารํ อชฺฌาวสนฺตี ปุตฺตวตี สมานา สามิกํ อภิภุยฺย วตฺเตยฺยนฺติ, อิทํ ภิกฺขเว ปญฺจมํ ฐานํ ทุลฺลภํ อกตปุญฺเญน มาตุคาเมน.\csromanspacer{} อิมานิ โข ภิกฺขเว ปญฺจ ฐานานิ ทุลฺลภานิ อกตปุญฺเญน มาตุคาเมนาติ.}

\prose{ปญฺจิมานิ ภิกฺขเว ฐานานิ สุลภานิ กตปุญฺเญน มาตุคาเมน.\csromanspacer{} กตมานิ ปญฺจ, ปติรูเป กุเล ชาเยยฺยนฺติ, อิทํ ภิกฺขเว ปฐมํ ฐานํ สุลภํ กตปุญฺเญน มาตุคาเมน.\csromanspacer{} ปติรูเป กุเล ชายิตฺวา ปติรูปํ กุลํ คจฺเฉยฺยนฺติ, อิทํ ภิกฺขเว ทุติยํ ฐานํ สุลภํ กตปุญฺเญน มาตุคาเมน.\csromanspacer{} ปติรูเป กุเล ชายิตฺวา ปติรูปํ กุลํ คนฺตฺวา อสปตฺติ อคารํ อชฺฌาวเสยฺยนฺติ, อิทํ ภิกฺขเว ตติยํ ฐานํ สุลภํ กตปุญฺเญน มาตุคาเมน.\csromanspacer{} ปติรูเป กุเล ชายิตฺวา ปติรูปํ กุลํ คนฺตฺวา อสปตฺติ อคารํ อชฺฌาวสนฺตี ปุตฺตวตี อสฺสนฺติ, อิทํ ภิกฺขเว จตุตฺถํ ฐานํ สุลภํ กตปุญฺเญน มาตุคาเมน.\csromanspacer{} ปติรูเป กุเล ชายิตฺวา ปติรูปํ กุลํ คนฺตฺวา อสปตฺติ อคารํ}

\gambhira{สฬายตนวคฺคสํยุตฺตปาฬิ}
\prose{\csromanfolio{446}อชฺฌาวสนฺตี ปุตฺตวตี สมานา สามิกํ อภิภุยฺย วตฺเตยฺยนฺติ, อิทํ ภิกฺขเว ปญฺจมํ ฐานํ สุลภํ กตปุญฺเญน มาตุคาเมน.\csromanspacer{} อิมานิ โข ภิกฺขเว ปญฺจ ฐานานิ สุลภานิ กตปุญฺเญน มาตุคาเมนาติ.\csromanspacer{}\csromanspacer{}อฏฺฐมํ.}

\csromanheader{2}{9. ปญฺจสีลวิสารทสุตฺต}
\proseitem{312}{ปญฺจหิ ภิกฺขเว ธมฺเมหิ สมนฺนาคโต มาตุคาโม วิสารโท อคารํ อชฺฌาวสติ.\csromanspacer{} กตเมหิ ปญฺจหิ, ปาณาติปาตา ปฏิวิรโต จ โหติ, อทินฺนาทานา ปฏิวิรโต จ โหติ, กาเมสุมิจฺฉาจารา ปฏิวิรโต จ โหติ, มุสาวาทา ปฏิวิรโต จ โหติ, สุราเมรยมชฺชปฺปมาทฏฺฐานา ปฏิวิรโต จ โหติ.\csromanspacer{} อิเมหิ โข ภิกฺขเว ปญฺจหิ ธมฺเมหิ สมนฺนาคโต มาตุคาโม วิสารโท อคารํ อชฺฌาวสตีติ.\csromanspacer{}\csromanspacer{}นวมํ.}

\csromanheader{2}{10. วฑฺฒีสุตฺต}
\proseitem{313}{ปญฺจหิ ภิกฺขเว วฑฺฒีหิ วฑฺฒมานา อริยสาวิกา อริยาย วฑฺฒิยา วฑฺฒติ, สาราทายินี จ โหติ วราทายินี จ กายสฺส.\csromanspacer{} กตเมหิ ปญฺจหิ, สทฺธาย วฑฺฒติ, สีเลน วฑฺฒติ, สุเตน วฑฺฒติ, จาเคน วฑฺฒติ, ปญฺญาย วฑฺฒติ.\csromanspacer{} อิเมหิ โข ภิกฺขเว ปญฺจหิ วฑฺฒีหิ วฑฺฒมานา อริยสาวิกา อริยาย วฑฺฒิยา วฑฺฒติ, สาราทายินี จ โหติ วราทายินี จ กายสฺสาติ.}

\csromangathameasure{สทฺธาย สีเลน จ ยาธ วฑฺฒติ, \\ ปญฺญาย จาเคน สุเตน จูภยํ. \\ สา ตาทิสี สีลวตี อุปาสิกา, \\ อาทียติ สารมิเธว อตฺตโนติ. ทสมํ. พลวคฺโค ตติโย.}
\csromangathagroup{\csromangathastanza{สทฺธาย สีเลน จ ยาธ วฑฺฒติ, \\ ปญฺญาย จาเคน สุเตน จูภยํ. \\ สา ตาทิสี สีลวตี อุปาสิกา, \\ อาทียติ สารมิเธว อตฺตโนติ.\csromanspacer{} ทสมํ.\csromanspacer{} พลวคฺโค ตติโย.}}

\titlehead{\textbf{ตสฺสุทฺทานํ}}
\csromangathameasure{วิสารทา ปสยฺห อภิภุยฺย, เอกํ องฺเคน ปญฺจมํ. \\ นาเสนฺติ เหตุ ฐานญฺจ, วิสารโท วฑฺฒินา ทสาติ. \\ \textbf{มาตุคามสํยุตฺตํ สมตฺตํ.}}
\csromangathagroup{\csromangathastanza{วิสารทา ปสยฺห อภิภุยฺย, เอกํ องฺเคน ปญฺจมํ. \\ นาเสนฺติ เหตุ ฐานญฺจ, วิสารโท วฑฺฒินา ทสาติ. \\ \textbf{มาตุคามสํยุตฺตํ สมตฺตํ.}}}

\chapterheadpageread{4. ชมฺพุขาทกสํยุตฺต}
\csromanmatikaanchor{mtk.244}
\csromantocmarknopage{chapter}{4. ชมฺพุขาทกสํยุตฺต}{mtk.244}
\bookmark[dest={mtk.244},level=0]{4. ชมฺพุขาทกสํยุตฺต}
\csromanheader{1}{1. นิพฺพานปญฺหาสุตฺต}
\csromanmatikaanchor{mtk.245}
\csromantocmarkref{section}{1. นิพฺพานปญฺหาสุตฺต}{mtk.245}
\bookmark[dest={mtk.245},level=1]{1. นิพฺพานปญฺหาสุตฺต}
\proseitem{314}{\csromanfolio{447}เอกํ สมยํ อายสฺมา สาริปุตฺโต มคเธสุ วิหรติ นาลกคามเก.\csromanspacer{} อถ โข ชมฺพุขาทโก ปริพฺพาชโก เยนายสฺมา สาริปุตฺโต เตนุปสงฺกมิ, อุปสงฺกมิตฺวา อายสฺมตา สาริปุตฺเตน สทฺธิํ สมฺโมทิ, สมฺโมทนียํ กถํ สารณียํ วีติสาเรตฺวา เอกมนฺตํ นิสีทิ, เอกมนฺตํ นิสินฺโน โข ชมฺพุขาทโก ปริพฺพาชโก อายสฺมนฺตํ สาริปุตฺตํ เอตทโวจ\csromandash{}}

\prose{“นิพฺพานํ นิพฺพานนฺ”ติ อาวุโส สาริปุตฺต วุจฺจติ, กตมํ นุ โข อาวุโส นิพฺพานนฺติ.\csromanspacer{} โย โข อาวุโส ราคกฺขโย โทสกฺขโย โมหกฺขโย, อิทํ วุจฺจติ “นิพฺพานนฺ”ติ.\csromanspacer{} อตฺถิ ปนาวุโส มคฺโค, อตฺถิ ปฏิปทา เอตสฺส นิพฺพานสฺส สจฺฉิกิริยายาติ.\csromanspacer{} อตฺถิ โข อาวุโส มคฺโค, อตฺถิ ปฏิปทา เอตสฺส นิพฺพานสฺส สจฺฉิกิริยายาติ.\csromanspacer{} กตโม ปนาวุโส มคฺโค, กตมา ปฏิปทา เอตสฺส นิพฺพานสฺส สจฺฉิกิริยายาติ.\csromanspacer{} อยเมว โข อาวุโส อริโย อฏฺฐงฺคิโก มคฺโค เอตสฺส นิพฺพานสฺส สจฺฉิกิริยาย.\csromanspacer{} เสยฺยถิทํ, สมฺมาทิฏฺฐิ สมฺมาสงฺกปฺโป สมฺมาวาจา สมฺมากมฺมนฺโต สมฺมา-อาชีโว สมฺมาวายาโม สมฺมาสติ สมฺมาสมาธิ.\csromanspacer{} อยํ โข อาวุโส มคฺโค, อยํ ปฏิปทา เอตสฺส นิพฺพานสฺส สจฺฉิกิริยายาติ.\csromanspacer{} ภทฺทโก อาวุโส มคฺโค ภทฺทิกา ปฏิปทา เอตสฺส นิพฺพานสฺส สจฺฉิกิริยาย, อลญฺจ ปนาวุโส สาริปุตฺต อปฺปมาทายาติ.\csromanspacer{}\csromanspacer{}ปฐมํ.}

\csromanheader{1}{2. อรหตฺตปญฺหาสุตฺต}
\csromanmatikaanchor{mtk.246}
\csromantocmarkref{section}{2. อรหตฺตปญฺหาสุตฺต}{mtk.246}
\bookmark[dest={mtk.246},level=1]{2. อรหตฺตปญฺหาสุตฺต}
\proseitem{315}{“อรหตฺตํ อรหตฺตนฺ”ติ อาวุโส สาริปุตฺต วุจฺจติ, กตมํ นุ โข อาวุโส อรหตฺตนฺติ.\csromanspacer{} โย โข อาวุโส ราคกฺขโย โทสกฺขโย โมหกฺขโย, อิทํ วุจฺจติ “อรหตฺตนฺ”ติ.\csromanspacer{} อตฺถิ ปนาวุโส มคฺโค, อตฺถิ ปฏิปทา เอตสฺส อรหตฺตสฺส สจฺฉิกิริยายาติ.\csromanspacer{} อตฺถิ โข อาวุโส มคฺโค, อตฺถิ ปฏิปทา เอตสฺส อรหตฺตสฺส สจฺฉิกิริยายาติ.}

\gambhira{สฬายตนวคฺคสํยุตฺตปาฬิ}
\prose{\csromanfolio{448}กตโม ปนาวุโส มคฺโค, กตมา ปฏิปทา เอตสฺส อรหตฺตสฺส สจฺฉิกิริยายาติ.\csromanspacer{} อยเมว โข อาวุโส อริโย อฏฺฐงฺคิโก มคฺโค เอตสฺส อรหตฺตสฺส สจฺฉิกิริยาย.\csromanspacer{} เสยฺยถิทํ, สมฺมาทิฏฺฐิ ฯเปฯ สมฺมาสมาธิ.\csromanspacer{} อยํ โข อาวุโส มคฺโค, อยํ ปฏิปทา เอตสฺส อรหตฺตสฺส สจฺฉิกิริยายาติ.\csromanspacer{} ภทฺทโก อาวุโส มคฺโค ภทฺทิกา ปฏิปทา เอตสฺส อรหตฺตสฺส สจฺฉิกิริยาย, อลญฺจ ปนาวุโส สาริปุตฺต อปฺปมาทายาติ.\csromanspacer{}\csromanspacer{}ทุติยํ.}

\csromanheader{1}{3. ธมฺมวาทีปญฺหาสุตฺต}
\csromanmatikaanchor{mtk.247}
\csromantocmarkref{section}{3. ธมฺมวาทีปญฺหาสุตฺต}{mtk.247}
\bookmark[dest={mtk.247},level=1]{3. ธมฺมวาทีปญฺหาสุตฺต}
\proseitem{316}{เก นุ โข อาวุโส สาริปุตฺต โลเก ธมฺมวาทิโน, เก โลเก สุปฺปฏิปนฺนา, เก โลเก สุคตาติ.\csromanspacer{} เย โข อาวุโส ราคปฺปหานาย ธมฺมํ เทเสนฺติ, โทสปฺปหานาย ธมฺมํ เทเสนฺติ, โมหปฺปหานาย ธมฺมํ เทเสนฺติ, เต โลเก ธมฺมวาทิโน.\csromanspacer{} เย โข อาวุโส ราคสฺส ปหานาย ปฏิปนฺนา, โทสสฺส ปหานาย ปฏิปนฺนา, โมหสฺส ปหานาย ปฏิปนฺนา, เต โลเก สุปฺปฏิปนฺนา.\csromanspacer{} เยสํ โข อาวุโส ราโค ปหีโน อุจฺฉินฺนมูโล ตาลาวตฺถุกโต อนภาวํกโต อายติํ อนุปฺปาทธมฺโม, โทโส ปหีโน อุจฺฉินฺนมูโล ตาลาวตฺถุกโต อนภาวํกโต อายติํ อนุปฺปาทธมฺโม, โมโห ปหีโน อุจฺฉินฺนมูโล ตาลาวตฺถุกโต อนภาวํกโต อายติํ อนุปฺปาทธมฺโม, เต โลเก สุคตาติ.}

\prose{อตฺถิ ปนาวุโส มคฺโค, อตฺถิ ปฏิปทา เอตสฺส ราคสฺส โทสสฺส โมหสฺส ปหานายาติ.\csromanspacer{} อตฺถิ โข อาวุโส มคฺโค, อตฺถิ ปฏิปทา เอตสฺส ราคสฺส โทสสฺส โมหสฺส ปหานายาติ.\csromanspacer{} กตโม ปนาวุโส มคฺโค, กตมา ปฏิปทา เอตสฺส ราคสฺส โทสสฺส โมหสฺส ปหานายาติ.\csromanspacer{} อยเมว โข อาวุโส อริโย อฏฺฐงฺคิโก มคฺโค เอตสฺส ราคสฺส โทสสฺส โมหสฺส ปหานาย.\csromanspacer{} เสยฺยถิทํ, สมฺมาทิฏฺฐิ ฯเปฯ สมฺมาสมาธิ.\csromanspacer{} อยํ โข อาวุโส มคฺโค, อยํ ปฏิปทา เอตสฺส ราคสฺส โทสสฺส โมหสฺส ปหานายาติ.\csromanspacer{} ภทฺทโก อาวุโส มคฺโค, ภทฺทิกา ปฏิปทา เอตสฺส ราคสฺส โทสสฺส โมหสฺส ปหานาย, อลญฺจ ปนาวุโส สาริปุตฺต อปฺปมาทายาติ.\csromanspacer{}\csromanspacer{}ตติยํ.}

\csromanheader{1}{4. ชมฺพุขาทกสํยุตฺต \textbf{4. กิมตฺถิยสุตฺต}}
\csromanmatikaanchor{mtk.248}
\csromantocmarkref{section}{4. กิมตฺถิยสุตฺต}{mtk.248}
\bookmark[dest={mtk.248},level=1]{4. กิมตฺถิยสุตฺต}
\proseitem{317}{\csromanfolio{449}กิมตฺถิยํ อาวุโส สาริปุตฺต สมเณ โคตเม พฺรหฺมจริยํ วุสฺสตีติ.\csromanspacer{} ทุกฺขสฺส โข อาวุโส ปริญฺญตฺถํ ภควติ พฺรหฺมจริยํ วุสฺสตีติ.\csromanspacer{} อตฺถิ ปนาวุโส มคฺโค, อตฺถิ ปฏิปทา เอตสฺส ทุกฺขสฺส ปริญฺญายาติ.\csromanspacer{} อตฺถิ โข อาวุโส มคฺโค, อตฺถิ ปฏิปทา เอตสฺส ทุกฺขสฺส ปริญฺญายาติ.\csromanspacer{} กตโม ปนาวุโส มคฺโค, กตมา ปฏิปทา เอตสฺส ทุกฺขสฺส ปริญฺญายาติ.\csromanspacer{} อยเมว โข อาวุโส อริโย อฏฺฐงฺคิโก มคฺโค เอตสฺส ทุกฺขสฺส ปริญฺญาย.\csromanspacer{} เสยฺยถิทํ, สมฺมาทิฏฺฐิ ฯเปฯ สมฺมาสมาธิ.\csromanspacer{} อยํ โข อาวุโส มคฺโค, อยํ ปฏิปทา เอตสฺส ทุกฺขสฺส ปริญฺญายาติ.\csromanspacer{} ภทฺทโก อาวุโส มคฺโค, ภทฺทิกา ปฏิปทา เอตสฺส ทุกฺขสฺส ปริญฺญาย, อลญฺจ ปนาวุโส สาริปุตฺต อปฺปมาทายาติ.\csromanspacer{}\csromanspacer{}จตุตฺถํ.}

\csromanheader{1}{5. อสฺสาสปฺปตฺตสุตฺต}
\csromanmatikaanchor{mtk.249}
\csromantocmarkref{section}{5. อสฺสาสปฺปตฺตสุตฺต}{mtk.249}
\bookmark[dest={mtk.249},level=1]{5. อสฺสาสปฺปตฺตสุตฺต}
\proseitem{318}{“อสฺสาสปฺปตฺโต อสฺสาสปฺปตฺโต”ติ อาวุโส สาริปุตฺต วุจฺจติ, กิตฺตาวตา นุ โข อาวุโส อสฺสาสปฺปตฺโต โหตีติ.\csromanspacer{} ยโต โข อาวุโส ภิกฺขุ ฉนฺนํ ผสฺสายตนานํ สมุทยญฺจ อตฺถงฺคมญฺจ อสฺสาทญฺจ อาทีนวญฺจ นิสฺสรณญฺจ ยถาภูตํ ปชานาติ.\csromanspacer{} เอตฺตาวตา โข อาวุโส อสฺสาสปฺปตฺโต โหตีติ.\csromanspacer{} อตฺถิ ปนาวุโส มคฺโค, อตฺถิ ปฏิปทา เอตสฺส อสฺสาสสฺส สจฺฉิกิริยายาติ.\csromanspacer{} อตฺถิ โข อาวุโส มคฺโค, อตฺถิ ปฏิปทา เอตสฺส อสฺสาสสฺส สจฺฉิกิริยายาติ.\csromanspacer{} กตโม ปนาวุโส มคฺโค, กตมา ปฏิปทา เอตสฺส อสฺสาสสฺส สจฺฉิกิริยายาติ.\csromanspacer{} อยเมว โข อาวุโส อริโย อฏฺฐงฺคิโก มคฺโค เอตสฺส อสฺสาสสฺส สจฺฉิกิริยาย.\csromanspacer{} เสยฺยถิทํ, สมฺมาทิฏฺฐิ ฯเปฯ สมฺมาสมาธิ.\csromanspacer{} อยํ โข อาวุโส มคฺโค, อยํ ปฏิปทา เอตสฺส อสฺสาสสฺส สจฺฉิกิริยายาติ.\csromanspacer{} ภทฺทโก อาวุโส มคฺโค, ภทฺทิกา ปฏิปทา เอตสฺส อสฺสาสสฺส สจฺฉิกิริยาย, อลญฺจ ปนาวุโส สาริปุตฺต อปฺปมาทายาติ.\csromanspacer{}\csromanspacer{}ปญฺจมํ.}

\csromanheader{1}{6. ปรมสฺสาสปฺปตฺตสุตฺต}
\csromanmatikaanchor{mtk.250}
\csromantocmarkref{section}{6. ปรมสฺสาสปฺปตฺตสุตฺต}{mtk.250}
\bookmark[dest={mtk.250},level=1]{6. ปรมสฺสาสปฺปตฺตสุตฺต}
\proseitem{319}{“ปรมสฺสาสปฺปตฺโต ปรมสฺสาสปฺปตฺโต”ติ อาวุโส สาริปุตฺต วุจฺจติ.\csromanspacer{} กิตฺตาวตา นุ โข อาวุโส ปรมสฺสาสปฺปตฺโต โหตีติ.}

\gambhira{สฬายตนวคฺคสํยุตฺตปาฬิ}
\csromanmatikaanchor{mtk.251}
\csromantocmarkref{section}{7-15. เวทนาปญฺหาสุตฺตาทิ}{mtk.251}
\bookmark[dest={mtk.251},level=1]{7-15. เวทนาปญฺหาสุตฺตาทิ}
\prose{\csromanfolio{450}ยโต โข อาวุโส ภิกฺขุ ฉนฺนํ ผสฺสายตนานํ สมุทยญฺจ อตฺถงฺคมญฺจ อสฺสาทญฺจ อาทีนวญฺจ นิสฺสรณญฺจ ยถาภูตํ วิทิตฺวา อนุปาทาวิมุตฺโต โหติ, เอตฺตาวตา โข อาวุโส ปรมสฺสาสปฺปตฺโต โหตีติ.\csromanspacer{} อตฺถิ ปนาวุโส มคฺโค, อตฺถิ ปฏิปทา เอตสฺส ปรมสฺสาสสฺส สจฺฉิกิริยายาติ.\csromanspacer{} อตฺถิ โข อาวุโส มคฺโค, อตฺถิ ปฏิปทา เอตสฺส ปรมสฺสาสสฺส สจฺฉิกิริยายาติ.\csromanspacer{} กตโม ปน อาวุโส มคฺโค, กตมา ปฏิปทา เอตสฺส ปรมสฺสาสสฺส สจฺฉิกิริยายาติ.\csromanspacer{} อยเมว โข อาวุโส อริโย อฏฺฐงฺคิโก มคฺโค เอตสฺส ปรมสฺสาสสฺส สจฺฉิกิริยาย.\csromanspacer{} เสยฺยถิทํ, สมฺมาทิฏฺฐิ ฯเปฯ สมฺมาสมาธิ.\csromanspacer{} อยํ โข อาวุโส มคฺโค.\csromanspacer{} อยํ ปฏิปทา เอตสฺส ปรมสฺสาสสฺส สจฺฉิกิริยายาติ.\csromanspacer{} ภทฺทโก อาวุโส มคฺโค, ภทฺทิกา ปฏิปทา เอตสฺส ปรมสฺสาสสฺส สจฺฉิกิริยาย, อลญฺจ ปนาวุโส สาริปุตฺต อปฺปมาทายาติ.\csromanspacer{}\csromanspacer{}ฉฏฺฐํ.}

\csromanheader{2}{7. เวทนาปญฺหาสุตฺต}
\proseitem{320}{“เวทนา เวทนา”ติ อาวุโส สาริปุตฺต วุจฺจติ, กตมา นุ โข อาวุโส เวทนาติ.\csromanspacer{} ติสฺโส อิมาวุโส เวทนา.\csromanspacer{} กตมา ติสฺโส, สุขา เวทนา ทุกฺขา เวทนา อทุกฺขมสุขา เวทนา.\csromanspacer{} อิมา โข อาวุโส ติสฺโส เวทนาติ.\csromanspacer{} อตฺถิ ปนาวุโส มคฺโค, อตฺถิ ปฏิปทา เอตาสํ ติสฺสนฺนํ เวทนานํ ปริญฺญายาติ.\csromanspacer{} อตฺถิ โข อาวุโส มคฺโค, อตฺถิ ปฏิปทา เอตาสํ ติสฺสนฺนํ เวทนานํ ปริญฺญายาติ.\csromanspacer{} กตโม ปนาวุโส มคฺโค, กตมา ปฏิปทา เอตาสํ ติสฺสนฺนํ เวทนานํ ปริญฺญายาติ.\csromanspacer{} อยเมว โข อาวุโส อริโย อฏฺฐงฺคิโก มคฺโค เอตาสํ ติสฺสนฺนํ เวทนานํ ปริญฺญายาติ.\csromanspacer{} เสยฺยถิทํ, สมฺมาทิฏฺฐิ ฯเปฯ สมฺมาสมาธิ.\csromanspacer{} อยํ โข อาวุโส มคฺโค, อยํ ปฏิปทา เอตาสํ ติสฺสนฺนํ เวทนานํ ปริญฺญายาติ.\csromanspacer{} ภทฺทโก อาวุโส มคฺโค, ภทฺทิกา ปฏิปทา เอตาสํ ติสฺสนฺนํ เวทนานํ ปริญฺญาย.\csromanspacer{} อลญฺจ ปนาวุโส สาริปุตฺต อปฺปมาทายาติ.\csromanspacer{}\csromanspacer{}สตฺตมํ.}

\csromanheader{2}{8. อาสวปญฺหาสุตฺต}
\proseitem{321}{“อาสโว อาสโว”ติ อาวุโส สาริปุตฺต วุจฺจติ, กตโม นุ โข อาวุโส อาสโวติ.\csromanspacer{} ตโย เม อาวุโส}

\vspace{\tipitakaparskip}
\csromancenter{ชมฺพุขาทกสํยุตฺต อาสวา, กามาสโว ภวาสโว อวิชฺชาสโว.\csromanspacer{} อิเม โข อาวุโส ตโย อาสวาติ.\csromanspacer{} อตฺถิ ปนาวุโส มคฺโค, อตฺถิ ปฏิปทา เอเตสํ อาสวานํ ปหานายาติ.\csromanspacer{} อตฺถิ โข อาวุโส มคฺโค, อตฺถิ ปฏิปทา เอเตสํ อาสวานํ ปหานายาติ.\csromanspacer{} กตโม อาวุโส มคฺโค, กตมา ปฏิปทา เอเตสํ อาสวานํ ปหานายาติ.\csromanspacer{} อยเมว โข อาวุโส อริโย อฏฺฐงฺคิโก มคฺโค เอเตสํ อาสวานํ ปหานาย.\csromanspacer{} เสยฺยถิทํ, สมฺมาทิฏฺฐิ ฯเปฯ สมฺมาสมาธิ.\csromanspacer{} อยํ โข อาวุโส มคฺโค, อยํ ปฏิปทา เอเตสํ อาสวานํ ปหานายาติ.\csromanspacer{} ภทฺทโก อาวุโส มคฺโค, ภทฺทิกา ปฏิปทา เอเตสํ อาสวานํ ปหานาย, อลญฺจ ปนาวุโส สาริปุตฺต อปฺปมาทายาติ.\csromanspacer{}\csromanspacer{}อฏฺฐมํ.}
\csromanheader{2}{9. อวิชฺชาปญฺหาสุตฺต}
\proseitem{322}{\csromanfolio{451}“อวิชฺชา อวิชฺชา”ติ อาวุโส สาริปุตฺต วุจฺจติ, กตมา นุ โข อาวุโส อวิชฺชาติ.\csromanspacer{} ยํ โข อาวุโส ทุกฺเข อญฺญาณํ, ทุกฺขสมุทเย อญฺญาณํ, ทุกฺขนิโรเธ อญฺญาณํ, ทุกฺขนิโรธคามินียา ปฏิปทาย อญฺญาณํ.\csromanspacer{} อยํ วุจฺจตาวุโส อวิชฺชาติ.\csromanspacer{} อตฺถิ ปนาวุโส มคฺโค, อตฺถิ ปฏิปทา เอติสฺสา อวิชฺชาย ปหานายาติ.\csromanspacer{} อตฺถิ โข อาวุโส มคฺโค, อตฺถิ ปฏิปทา เอติสฺสา อวิชฺชาย ปหานายาติ.\csromanspacer{} กตโม ปนาวุโส มคฺโค, กตมา ปฏิปทา เอติสฺสา อวิชฺชาย ปหานายาติ.\csromanspacer{} อยเมว โข อาวุโส อริโย อฏฺฐงฺคิโก มคฺโค เอติสฺสา อวิชฺชาย ปหานาย.\csromanspacer{} เสยฺยถิทํ, สมฺมาทิฏฺฐิ ฯเปฯ สมฺมาสมาธิ.\csromanspacer{} อยํ โข อาวุโส มคฺโค.\csromanspacer{} อยํ ปฏิปทา เอติสฺสา อวิชฺชาย ปหานายาติ.\csromanspacer{} ภทฺทโก อาวุโส มคฺโค, ภทฺทิกา ปฏิปทา เอติสฺสา อวิชฺชาย ปหานาย, อลญฺจ ปนาวุโส สาริปุตฺต อปฺปมาทายาติ.\csromanspacer{}\csromanspacer{}นวมํ.}

\csromanheader{2}{10. ตณฺหาปญฺหาสุตฺต}
\proseitem{323}{“ตณฺหา ตณฺหา”ติ อาวุโส สาริปุตฺต วุจฺจติ, กตมา นุ โข อาวุโส ตณฺหาติ.\csromanspacer{} ติสฺโส อิมา อาวุโส ตณฺหา, กามตณฺหา ภวตณฺหา วิภวตณฺหา.\csromanspacer{} อิมา โข อาวุโส ติสฺโส ตณฺหาติ.\csromanspacer{} อตฺถิ ปนาวุโส มคฺโค, อตฺถิ ปฏิปทา เอตาสํ ตณฺหานํ ปหานายาติ.\csromanspacer{} อตฺถิ โข อาวุโส มคฺโค, อตฺถิ ปฏิปทา เอตาสํ ตณฺหานํ ปหานายาติ.}

\gambhira{สฬายตนวคฺคสํยุตฺตปาฬิ}
\prose{\csromanfolio{452}กตโม ปนาวุโส มคฺโค, กตมา ปฏิปทา เอตาสํ ตณฺหานํ ปหานายาติ.\csromanspacer{} อยเมว โข อาวุโส อริโย อฏฺฐงฺคิโก มคฺโค เอตาสํ ตณฺหานํ ปหานาย.\csromanspacer{} เสยฺยถิทํ, สมฺมาทิฏฺฐิ ฯเปฯ สมฺมาสมาธิ.\csromanspacer{} อยํ โข อาวุโส มคฺโค, อยํ ปฏิปทา เอตาสํ ตณฺหานํ ปหานายาติ.\csromanspacer{} ภทฺทโก อาวุโส มคฺโค, ภทฺทิกา ปฏิปทา เอตาสํ ตณฺหานํ ปหานาย, อลญฺจ ปนาวุโส สาริปุตฺต อปฺปมาทายาติ.\csromanspacer{}\csromanspacer{}ทสมํ.}

\csromanheader{2}{11. โอฆปญฺหาสุตฺต}
\proseitem{324}{“โอโฆ โอโฆ”ติ อาวุโส สาริปุตฺต วุจฺจติ, กตโม นุ โข อาวุโส โอโฆติ.\csromanspacer{} จตฺตาโรเม อาวุโส โอฆา, กาโมโฆ ภโวโฆ ทิฏฺโฐโฆ อวิชฺโชโฆ.\csromanspacer{} อิเม โข อาวุโส จตฺตาโร โอฆาติ.\csromanspacer{} อตฺถิ ปนาวุโส มคฺโค, อตฺถิ ปฏิปทา เอเตสํ โอฆานํ ปหานายาติ.\csromanspacer{} อตฺถิ โข อาวุโส มคฺโค, อตฺถิ ปฏิปทา เอเตสํ โอฆานํ ปหานายาติ.\csromanspacer{} กตโม ปนาวุโส มคฺโค, กตมา ปฏิปทา เอเตสํ โอฆานํ ปหานายาติ.\csromanspacer{} อยเมว โข อาวุโส อริโย อฏฺฐงฺคิโก มคฺโค เอเตสํ โอฆานํ ปหานาย.\csromanspacer{} เสยฺยถิทํ, สมฺมาทิฏฺฐิ ฯเปฯ สมฺมาสมาธิ.\csromanspacer{} อยํ โข อาวุโส มคฺโค, อยํ ปฏิปทา เอเตสํ โอฆานํ ปหานายาติ.\csromanspacer{} ภทฺทโก อาวุโส มคฺโค, ภทฺทิกา ปฏิปทา เอเตสํ โอฆานํ ปหานาย, อลญฺจ ปนาวุโส สาริปุตฺต อปฺปมาทายาติ.\csromanspacer{}\csromanspacer{}เอกาทสมํ.}

\csromanheader{2}{12. อุปาทานปญฺหาสุตฺต}
\proseitem{325}{“อุปาทานํ อุปาทานนฺ”ติ อาวุโส สาริปุตฺต วุจฺจติ, กตมํ นุ โข อาวุโส อุปาทานนฺติ.\csromanspacer{} จตฺตาริมานิ อาวุโส อุปาทานานิ, กามุปาทานํ ทิฏฺฐุปาทานํ สีลพฺพตุปาทานํ อตฺตวาทุปาทานํ.\csromanspacer{} อิมานิ โข อาวุโส จตฺตาริ อุปาทานานีติ.\csromanspacer{} อตฺถิ ปนาวุโส มคฺโค, อตฺถิ ปฏิปทา เอเตสํ อุปาทานานํ ปหานายาติ.\csromanspacer{} อตฺถิ ปนาวุโส มคฺโค, อตฺถิ ปฏิปทา เอเตสํ อุปาทานานํ ปหานายาติ.\csromanspacer{} กตโม ปนาวุโส มคฺโค, กตมา ปฏิปทา เอเตสํ อุปาทานานํ ปหานายาติ.\csromanspacer{} อยเมว โข อาวุโส อริโย อฏฺฐงฺคิโก มคฺโค เอเตสํ อุปาทานานํ ปหานาย.}

\vspace{\tipitakaparskip}
\csromancenter{ชมฺพุขาทกสํยุตฺต เสยฺยถิทํ, สมฺมาทิฏฺฐิ ฯเปฯ สมฺมาสมาธิ.\csromanspacer{} อยํ โข อาวุโส มคฺโค, อยํ ปฏิปทา เอเตสํ อุปาทานานํ ปหานายาติ.\csromanspacer{} ภทฺทโก อาวุโส มคฺโค, ภทฺทิกา ปฏิปทา เอเตสํ อุปาทานานํ ปหานาย, อลญฺจ ปนาวุโส สาริปุตฺต อปฺปมาทายาติ.\csromanspacer{}\csromanspacer{}ทฺวาทสมํ.}
\csromanheader{2}{13. ภวปญฺหาสุตฺต}
\proseitem{326}{\csromanfolio{453}“ภโว ภโว”ติ อาวุโส สาริปุตฺต วุจฺจติ, กตโม นุ โข อาวุโส ภโวติ.\csromanspacer{} ตโย เม อาวุโส ภวา, กามภโว รูปภโว อรูปภโว.\csromanspacer{} อิเม โข อาวุโส ตโย ภวาติ.\csromanspacer{} อตฺถิ ปนาวุโส มคฺโค, อตฺถิ ปฏิปทา เอเตสํ ภวานํ ปริญฺญายาติ.\csromanspacer{} อตฺถิ โข อาวุโส มคฺโค, อตฺถิ ปฏิปทา เอเตสํ ภวานํ ปริญฺญายาติ.\csromanspacer{} กตโม ปนาวุโส มคฺโค, กตมา ปฏิปทา เอเตสํ ภวานํ ปริญฺญายาติ.\csromanspacer{} อยเมว โข อาวุโส อริโย อฏฺฐงฺคิโก มคฺโค เอเตสํ ภวานํ ปริญฺญาย.\csromanspacer{} เสยฺยถิทํ, สมฺมาทิฏฺฐิ ฯเปฯ สมฺมาสมาธิ.\csromanspacer{} อยํ โข อาวุโส มคฺโค, อยํ ปฏิปทา เอเตสํ ภวานํ ปริญฺญายาติ.\csromanspacer{} ภทฺทโก อาวุโส มคฺโค, ภทฺทิกา ปฏิปทา เอเตสํ ภวานํ ปริญฺญาย, อลญฺจ ปนาวุโส สาริปุตฺต อปฺปมาทายาติ.\csromanspacer{}\csromanspacer{}เตรสมํ.}

\csromanheader{2}{14. ทุกฺขปญฺหาสุตฺต}
\proseitem{327}{“ทุกฺขํ ทุกฺขนฺ”ติ อาวุโส สาริปุตฺต วุจฺจติ, กตมํ นุ โข อาวุโส ทุกฺขนฺติ.\csromanspacer{} ติสฺโส อิมา อาวุโส ทุกฺขตา, ทุกฺขทุกฺขตา สงฺขารทุกฺขตา วิปริณามทุกฺขตา.\csromanspacer{} อิมา โข อาวุโส ติสฺโส ทุกฺขตาติ.\csromanspacer{} อตฺถิ ปนาวุโส มคฺโค, อตฺถิ ปฏิปทา เอตาสํ ทุกฺขตานํ ปริญฺญายาติ.\csromanspacer{} อตฺถิ โข อาวุโส มคฺโค, อตฺถิ ปฏิปทา เอตาสํ ทุกฺขตานํ ปริญฺญายาติ.\csromanspacer{} กตโม ปนาวุโส มคฺโค, กตมา ปฏิปทา เอตาสํ ทุกฺขตานํ ปริญฺญายาติ.\csromanspacer{} อยเมว โข อาวุโส อริโย อฏฺฐงฺคิโก มคฺโค เอตาสํ ทุกฺขตานํ ปริญฺญาย.\csromanspacer{} เสยฺยถิทํ, สมฺมาทิฏฺฐิ ฯเปฯ สมฺมาสมาธิ.\csromanspacer{} อยํ โข อาวุโส มคฺโค, อยํ ปฏิปทา เอตาสํ ทุกฺขตานํ ปริญฺญายาติ.\csromanspacer{} ภทฺทโก อาวุโส มคฺโค, ภทฺทิกา ปฏิปทา เอตาสํ ทุกฺขตานํ ปริญฺญาย, อลญฺจ ปนาวุโส สาริปุตฺต อปฺปมาทายาติ.\csromanspacer{}\csromanspacer{}จุทฺทสมํ.}

\gambhira{สฬายตนวคฺคสํยุตฺตปาฬิ}
\csromanheader{2}{15. สกฺกายปญฺหาสุตฺต}
\proseitem{328}{\csromanfolio{454}“สกฺกาโย สกฺกาโย”ติ อาวุโส สาริปุตฺต วุจฺจติ, กตโม นุ โข อาวุโส สกฺกาโยติ.\csromanspacer{} ปญฺจิเม อาวุโส อุปาทานกฺขนฺธา สกฺกาโย วุตฺโต ภควตา.\csromanspacer{} เสยฺยถิทํ, รูปุปาทานกฺขนฺโธ เวทนุปาทานกฺขนฺโธ สญฺญุปาทานกฺขนฺโธ สงฺขารุปาทานกฺขนฺโธ วิญฺญาณุปาทานกฺขนฺโธ.\csromanspacer{} อิเม โข อาวุโส ปญฺจุปาทานกฺขนฺธา สกฺกาโย วุตฺโต ภควตาติ.\csromanspacer{} อตฺถิ ปนาวุโส มคฺโค, อตฺถิ ปฏิปทา เอตสฺส สกฺกายสฺส ปริญฺญายาติ.\csromanspacer{} อตฺถิ โข อาวุโส มคฺโค, อตฺถิ ปฏิปทา เอตสฺส สกฺกายสฺส ปริญฺญายาติ.\csromanspacer{} กตโม ปนาวุโส มคฺโค, กตมา ปฏิปทา เอตสฺส สกฺกายสฺส ปริญฺญายาติ.\csromanspacer{} อยเมว โข อาวุโส อริโย อฏฺฐงฺคิโก มคฺโค เอตสฺส สกฺกายสฺส ปริญฺญาย.\csromanspacer{} เสยฺยถิทํ, สมฺมาทิฏฺฐิ ฯเปฯ สมฺมาสมาธิ.\csromanspacer{} อยํ โข อาวุโส มคฺโค, อยํ ปฏิปทา เอตสฺส สกฺกายสฺส ปริญฺญายาติ.\csromanspacer{} ภทฺทโก อาวุโส มคฺโค, ภทฺทิกา ปฏิปทา เอตสฺส สกฺกายสฺส ปริญฺญาย, อลญฺจ ปนาวุโส สาริปุตฺต อปฺปมาทายาติ.\csromanspacer{}\csromanspacer{}ปนฺนรสมํ.}

\csromanheader{1}{16. ทุกฺกรปญฺหาสุตฺต}
\csromanmatikaanchor{mtk.252}
\csromantocmarkref{section}{16. ทุกฺกรปญฺหาสุตฺต}{mtk.252}
\bookmark[dest={mtk.252},level=1]{16. ทุกฺกรปญฺหาสุตฺต}
\proseitem{329}{กิํ นุ โข อาวุโส สาริปุตฺต อิมสฺมิํ ธมฺมวินเย ทุกฺกรนฺติ.\csromanspacer{} ปพฺพชฺชา โข อาวุโส อิมสฺมิํ ธมฺมวินเย ทุกฺกราติ.\csromanspacer{} ปพฺพชิเตน ปนาวุโส กิํ ทุกฺกรนฺติ.\csromanspacer{} ปพฺพชิเตน โข อาวุโส อภิรติ ทุกฺกราติ.\csromanspacer{} อภิรเตน ปนาวุโส กิํ ทุกฺกรนฺติ.\csromanspacer{} อภิรเตน โข อาวุโส ธมฺมานุธมฺมปฺปฏิปตฺติ ทุกฺกราติ.\csromanspacer{} กีวจิรํ ปนาวุโส ธมฺมานุธมฺมปฺปฏิปนฺโน ภิกฺขุ อรหํ อสฺสาติ.\csromanspacer{} นจิรํ อาวุโสติ.\csromanspacer{}\csromanspacer{}โสฬสมํ.}

\vspace{\tipitakaparskip}
\csromancenter{\textbf{ชมฺพุขาทกสํยุตฺตํ สมตฺตํ.}}
\titlehead{\textbf{ตสฺสุทฺทานํ}}
\csromangathameasure{นิพฺพานํ อรหตฺตญฺจ, ธมฺมวาที กิมตฺถิยํ. \\ อสฺสาโส ปรมสฺสาโส, เวทนา อาสวาวิชฺชา. \\ ตณฺหา โอฆา อุปาทานํ, ภโว ทุกฺขญฺจ สกฺกาโย. \\ อิมสฺมิํ ธมฺมวินเย ทุกฺกรนฺติ.}
\csromangathagroup{\csromangathastanza{นิพฺพานํ อรหตฺตญฺจ, ธมฺมวาที กิมตฺถิยํ. \\ อสฺสาโส ปรมสฺสาโส, เวทนา อาสวาวิชฺชา.}\csromangathastanza{ตณฺหา โอฆา อุปาทานํ, ภโว ทุกฺขญฺจ สกฺกาโย. \\ อิมสฺมิํ ธมฺมวินเย ทุกฺกรนฺติ.}}

\chapterheadpageread{5. สามณฺฑกสํยุตฺต}
\csromanmatikaanchor{mtk.253}
\csromantocmarknopage{chapter}{5. สามณฺฑกสํยุตฺต}{mtk.253}
\bookmark[dest={mtk.253},level=0]{5. สามณฺฑกสํยุตฺต}
\csromanheader{1}{1. สามณฺฑกสุตฺต}
\csromanmatikaanchor{mtk.254}
\csromantocmarkref{section}{1. สามณฺฑกสุตฺต}{mtk.254}
\bookmark[dest={mtk.254},level=1]{1. สามณฺฑกสุตฺต}
\proseitem{330}{\csromanfolio{455}เอกํ สมยํ อายสฺมา สาริปุตฺโต วชฺชีสุ วิหรติ อุกฺกเจลายํ คงฺคาย นทิยา ตีเร.\csromanspacer{} อถ โข สามณฺฑโก\footnote{สามณฺฑกานิ (สี)} ปริพฺพาชโก เยนายสฺมา สาริปุตฺโต เตนุปสงฺกมิ, อุปสงฺกมิตฺวา อายสฺมตา สาริปุตฺเตน สทฺธิํ สมฺโมทิ, สมฺโมทนียํ กถํ สารณียํ วีติสาเรตฺวา เอกมนฺตํ นิสีทิ, เอกมนฺตํ นิสินฺโน โข สามณฺฑโก ปริพฺพาชโก อายสฺมนฺตํ สาริปุตฺตํ เอตทโวจ\csromandash{}}

\prose{“นิพฺพานํ นิพฺพานนฺ”ติ อาวุโส สาริปุตฺต วุจฺจติ, กตมํ นุ โข อาวุโส นิพฺพานนฺติ.\csromanspacer{} โย โข อาวุโส ราคกฺขโย โทสกฺขโย โมหกฺขโย, อิทํ วุจฺจติ “นิพฺพานนฺ”ติ.\csromanspacer{} อตฺถิ ปนาวุโส มคฺโค, อตฺถิ ปฏิปทา เอตสฺส นิพฺพานสฺส สจฺฉิกิริยายาติ.\csromanspacer{} อตฺถิ โข อาวุโส มคฺโค, อตฺถิ ปฏิปทา เอตสฺส นิพฺพานสฺส สจฺฉิกิริยายาติ.}

\prose{กตโม ปนาวุโส มคฺโค, กตมา ปฏิปทา เอตสฺส นิพฺพานสฺส สจฺฉิกิริยายาติ.\csromanspacer{} อยเมว โข อาวุโส อริโย อฏฺฐงฺคิโก มคฺโค เอตสฺส นิพฺพานสฺส สจฺฉิกิริยาย.\csromanspacer{} เสยฺยถิทํ, สมฺมาทิฏฺฐิ ฯเปฯ สมฺมาสมาธิ.\csromanspacer{} อยํ โข อาวุโส มคฺโค, อยํ ปฏิปทา เอตสฺส นิพฺพานสฺส สจฺฉิกิริยายาติ.\csromanspacer{} ภทฺทโก อาวุโส มคฺโค, ภทฺทิกา ปฏิปทา เอตสฺส นิพฺพานสฺส สจฺฉิกิริยาย, อลญฺจ ปนาวุโส สาริปุตฺต อปฺปมาทายาติ.\csromanspacer{}\csromanspacer{}ปฐมํ. (ยถา ชมฺพุขาทกสํยุตฺตํ, ตถา วิตฺถาเรตพฺพํ.)}

\gambhira{สฬายตนวคฺคสํยุตฺตปาฬิ}
\csromanheader{1}{2. ทุกฺกรสุตฺต}
\csromanmatikaanchor{mtk.255}
\csromantocmarkref{section}{2. ทุกฺกรสุตฺต}{mtk.255}
\bookmark[dest={mtk.255},level=1]{2. ทุกฺกรสุตฺต}
\proseitem{331}{\csromanfolio{456}กิํ นุ โข อาวุโส สาริปุตฺต อิมสฺมิํ ธมฺมวินเย ทุกฺกรนฺติ.\csromanspacer{} ปพฺพชฺชา โข อาวุโส อิมสฺมิํ ธมฺมวินเย ทุกฺกราติ.\csromanspacer{} ปพฺพชิเตน ปนาวุโส กิํ ทุกฺกรนฺติ.\csromanspacer{} ปพฺพชิเตน โข อาวุโส อภิรติ ทุกฺกราติ.\csromanspacer{} อภิรเตน ปนาวุโส กิํ ทุกฺกรนฺติ.\csromanspacer{} อภิรเตน โข อาวุโส ธมฺมานุธมฺมปฺปฏิปตฺติ ทุกฺกราติ.\csromanspacer{} กีวจิรํ ปนาวุโส ธมฺมานุธมฺมปฺปฏิปนฺโน ภิกฺขุ อรหํ อสฺสาติ.\csromanspacer{} นจิรํ อาวุโสติ.\csromanspacer{}\csromanspacer{}โสฬสมํ.\csromanspacer{} ปุริมกสทิสํ อุทฺทานํ.}

\vspace{\tipitakaparskip}
\csromancenter{\textbf{สามณฺฑกสํยุตฺตํ สมตฺตํ.}}
\chapterheadpageread{6. โมคฺคลฺลานสํยุตฺต}
\csromanmatikaanchor{mtk.256}
\csromantocmarknopage{chapter}{6. โมคฺคลฺลานสํยุตฺต}{mtk.256}
\bookmark[dest={mtk.256},level=0]{6. โมคฺคลฺลานสํยุตฺต}
\csromanheader{2}{1. ปฐมฌานปญฺหาสุตฺต}
\csromanmatikaanchor{mtk.257}
\csromantocmarkref{section}{1-9. ปฐมฌานปญฺหาสุตฺตาทิ}{mtk.257}
\bookmark[dest={mtk.257},level=1]{1-9. ปฐมฌานปญฺหาสุตฺตาทิ}
\proseitem{332}{\csromanfolio{457}เอกํ สมยํ อายสฺมา มหาโมคฺคลฺลาโน สาวตฺถิยํ วิหรติ เชตวเน อนาถปิณฺฑิกสฺส อาราเม.\csromanspacer{} ตตฺร โข อายสฺมา มหาโมคฺคลฺลาโน ภิกฺขู อามนฺเตสิ “อาวุโส ภิกฺขเว”ติ. “อาวุโส”ติ โข เต ภิกฺขู อายสฺมโต มหาโมคฺคลฺลานสฺส ปจฺจสฺโสสุํ.\csromanspacer{} อายสฺมา มหาโมคฺคลฺลาโน เอตทโวจ\csromandash{}}

\prose{อิธ มยฺหํ อาวุโส รโหคตสฺส ปฏิสลฺลีนสฺส เอวํ เจตโส ปริวิตกฺโก อุทปาทิ “ปฐมํ ฌานํ ปฐมํ ฌานนฺ”ติ วุจฺจติ, กตมํ นุ โข ปฐมํ ฌานนฺติ.\csromanspacer{} ตสฺส มยฺหํ อาวุโส เอตทโหสิ “อิธ ภิกฺขุ วิวิจฺเจว กาเมหิ วิวิจฺจ อกุสเลหิ ธมฺเมหิ สวิตกฺกํ สวิจารํ วิเวกชํ ปีติสุขํ ปฐมํ ฌานํ อุปสมฺปชฺช วิหรติ, อิทํ วุจฺจติ ปฐมา ฌานนฺ”ติ.\csromanspacer{} โส ขฺวาหํ อาวุโส วิวิจฺเจว กาเมหิ วิวิจฺจ อกุสเลหิ ธมฺเมหิ สวิตกฺกํ สวิจารํ วิเวกชํ ปีติสุขํ ปฐมํ ฌานํ อุปสมฺปชฺช วิหรามิ, ตสฺส มยฺหํ อาวุโส อิมินา วิหาเรน วิหรโต กามสหคตา สญฺญามนสิการา สมุทาจรนฺติ.}

\prose{อถ โข มํ อาวุโส ภควา อิทฺธิยา อุปสงฺกมิตฺวา เอตทโวจ “โมคฺคลฺลาน โมคฺคลฺลาน มา พฺราหฺมณ ปฐมํ ฌานํ ปมาโท, ปฐเม ฌาเน จิตฺตํ สณฺฐเปหิ, ปฐเม ฌาเน จิตฺตํ เอโกทิํ กโรหิ\footnote{เอโกทิกโรหิ (อิ)}, ปฐเม ฌาเน จิตฺตํ สมาทหา”ติ.\csromanspacer{} โส ขฺวาหํ อาวุโส อปเรน สมเยน วิวิจฺเจว กาเมหิ วิวิจฺจ อกุสเลหิ ธมฺเมหิ สวิตกฺกํ สวิจารํ วิเวกชํ ปีติสุขํ ปฐมํ ฌานํ อุปสมฺปชฺช วิหาสิํ.\csromanspacer{} ยํ หิ ตํ อาวุโส สมฺมา วทมาโน วเทยฺย “สตฺถารานุคฺคหิโต สาวโก มหาภิญฺญตํ ปตฺโต”ติ, มมํ ตํ สมฺมา วทมาโน วเทยฺย “สตฺถารานุคฺคหิโต สาวโก มหาภิญฺญตํ ปตฺโต”ติ.\csromanspacer{}\csromanspacer{}ปฐมํ.}

\gambhira{สฬายตนวคฺคสํยุตฺตปาฬิ}
\csromanheader{2}{2. ทุติยฌานปญฺหาสุตฺต}
\proseitem{333}{\csromanfolio{458}“ทุติยํ ฌานํ ทุติยํ ฌานนฺ”ติ วุจฺจติ, กตมํ นุ โข ทุติยํ ฌานนฺติ.\csromanspacer{} ตสฺส มยฺหํ อาวุโส เอตทโหสิ “อิธ ภิกฺขุ วิตกฺกวิจารานํ วูปสมา อชฺฌตฺตํ สมฺปสาทนํ เจตโส เอโกทิภาวํ อวิตกฺกํ อวิจารํ สมาธิชํ ปีติสุขํ ทุติยํ ฌานํ อุปสมฺปชฺช วิหรติ, อิทํ วุจฺจติ ทุติยํ ฌานนฺ”ติ.\csromanspacer{} โส ขฺวาหํ อาวุโส วิตกฺกวิจารานํ วูปสมา อชฺฌตฺตํ สมฺปสาทนํ เจตโส เอโกทิภาวํ อวิตกฺกํ อวิจารํ สมาธิชํ ปีติสุขํ ทุติยํ ฌานํ อุปสมฺปชฺช วิหรามิ, ตสฺส มยฺหํ อาวุโส อิมินา วิหาเรน วิหรโต วิตกฺกสหคตา สญฺญามนสิการา สมุทาจรนฺติ.}

\prose{อถ โข มํ อาวุโส ภควา อิทฺธิยา อุปสงฺกมิตฺวา เอตทโวจ “โมคฺคลฺลาน โมคฺคลฺลาน มา พฺราหฺมณ ทุติยํ ฌานํ ปมาโท, ทุติเย ฌาเน จิตฺตํ สณฺฐเปหิ, ทุติเย ฌาเน จิตฺตํ เอโกทิํ กโรหิ, ทุติเย ฌาเน จิตฺตํ สมาทหา”ติ.\csromanspacer{} โส ขฺวาหํ อาวุโส อปเรน สมเยน วิตกฺกวิจารานํ วูปสมา อชฺฌตฺตํ สมฺปสาทนํ เจตโส เอโกทิภาวํ อวิตกฺกํ อวิจารํ สมาธิชํ ปีติสุขํ ทุติยํ ฌานํ อุปสมฺปชฺช วิหาสิํ.\csromanspacer{} ยํ หิ ตํ อาวุโส สมฺมา วทมาโน วเทยฺย “สตฺถารานุคฺคหิโต สาวโก มหาภิญฺญตํ ปตฺโต”ติ, มมํ ตํ สมฺมา วทมาโน วเทยฺย “สตฺถารานุคฺคหิโต สาวโก มหาภิญฺญตํ ปตฺโต”ติ.\csromanspacer{}\csromanspacer{}ทุติยํ.}

\csromanheader{2}{3. ตติยฌานปญฺหาสุตฺต}
\proseitem{334}{“ตติยํ ฌานํ ตติยํ ฌานนฺ”ติ วุจฺจติ, กตมํ นุ โข ตติยํ ฌานนฺติ.\csromanspacer{} ตสฺส มยฺหํ อาวุโส เอตทโหสิ “อิธ ภิกฺขุ ปีติยา จ วิราคา อุเปกฺขโก จ วิหรติ, สโต จ สมฺปชาโน สุขญฺจ กาเยน ปฏิสํเวเทติ, ยํ ตํ อริยา อาจิกฺขนฺติ ‘อุเปกฺขโก สติมา สุขวิหารี’ติ, ตติยํ ฌานํ อุปสมฺปชฺช วิหรติ, อิทํ วุจฺจติ ตติยํ ฌานนฺ”ติ.\csromanspacer{} โส ขฺวาหํ อาวุโส ปีติยา จ วิราคา อุเปกฺขโก จ วิหรามิ, สโต จ สมฺปชาโน สุขญฺจ กาเยน ปฏิสํเวเทมิ, ยํ ตํ อริยา อาจิกฺขนฺติ “อุเปกฺขโก สติมา สุขวิหารี”ติ, ตติยํ ฌานํ อุปสมฺปชฺช วิหรามิ, ตสฺส มยฺหํ อาวุโส อิมินา วิหาเรน วิหรโต ปีติสหคตา สญฺญามนสิการา สมุทาจรนฺติ.}

\csromanheader{2}{6. โมคฺคลฺลานสํยุตฺต}
\prose{\csromanfolio{459}อถ โข มํ อาวุโส ภควา อิทฺธิยา อุปสงฺกมิตฺวา เอตทโวจ “โมคฺคลฺลาน โมคฺคลฺลาน มา พฺราหฺมณ ตติยํ ฌานํ ปมาโท, ตติเย ฌาเน จิตฺตํ สณฺฐเปหิ, ตติเย ฌาเน จิตฺตํ เอโกทิํ กโรหิ, ตติเย ฌาเน จิตฺตํ สมาทหา”ติ.\csromanspacer{} โส ขฺวาหํ อาวุโส อปเรน สมเยน ปีติยา จ วิราคา อุเปกฺขโก จ วิหรามิ, สโต จ สมฺปชาโน สุขญฺจ กาเยน ปฏิสํเวเทมิ.\csromanspacer{} ยํ ตํ อริยา อาจิกฺขนฺติ “อุเปกฺขโก สติมา สุขวิหารี”ติ, ตติยํ ฌานํ อุปสมฺปชฺช วิหาสิํ.\csromanspacer{} ยํ หิ ตํ อาวุโส สมฺมา วทมาโน วเทยฺย ฯเปฯ มหาภิญฺญตํ ปตฺโตติ.\csromanspacer{}\csromanspacer{}ตติยํ.}

\csromanheader{2}{4. จตุตฺถฌานปญฺหาสุตฺต}
\proseitem{335}{“จตุตฺถํ ฌานํ จตุตฺถํ ฌานนฺ”ติ วุจฺจติ, กตมํ นุ โข จตุตฺถํ ฌานนฺติ.\csromanspacer{} ตสฺส มยฺหํ อาวุโส เอตทโหสิ “อิธ ภิกฺขุ สุขสฺส จ ปหานา ทุกฺขสฺส จ ปหานา ปุพฺเพว โสมนสฺสโทมนสฺสานํ อตฺถงฺคมา อทุกฺขมสุขํ อุเปกฺขาสติปาริสุทฺธิํ จตุตฺถํ ฌานํ อุปสมฺปชฺช วิหรติ, อิทํ วุจฺจติ จตุตฺถํ ฌานนฺ”ติ.\csromanspacer{} โส ขฺวาหํ อาวุโส สุขสฺส จ ปหานา ทุกฺขสฺส จ ปหานา ปุพฺเพว โสมนสฺสโทมนสฺสานํ อตฺถงฺคมา อทุกฺขมสุขํ อุเปกฺขาสติปาริสุทฺธิํ จตุตฺถํ ฌานํ อุปสมฺปชฺช วิหรามิ, ตสฺส มยฺหํ อาวุโส อิมินา วิหาเรน วิหรโต สุขสหคตา สญฺญามนสิการา สมุทาจรนฺติ.}

\prose{อถ โข มํ อาวุโส ภควา อิทฺธิยา อุปสงฺกมิตฺวา เอตทโวจ “โมคฺคลฺลาน โมคฺคลฺลาน มา พฺราหฺมณ จตุตฺถํ ฌานํ ปมาโท จตุตฺเถ ฌาเน จิตฺตํ สณฺฐเปหิ, จตุตฺเถ ฌาเน จิตฺตํ เอโกทิํ กโรหิ, จตุตฺเถ ฌาเน จิตฺตํ สมาทหา”ติ.\csromanspacer{} โส ขฺวาหํ อาวุโส อปเรน สมเยน สุขสฺส จ ปหานา ทุกฺขสฺส จ ปหานา ปุพฺเพว โสมนสฺสโทมนสฺสานํ อตฺถงฺคมา อทุกฺขมสุขํ อุเปกฺขาสติปาริสุทฺธิํ จตุตฺถํ ฌานํ อุปสมฺปชฺช วิหาสิํ.\csromanspacer{} ยํ หิ ตํ อาวุโส สมฺมา วทมาโน วเทยฺย ฯเปฯ มหาภิญฺญตํ ปตฺโตติ.\csromanspacer{}\csromanspacer{}จตุตฺถํ.}

\csromanheader{2}{5. อากาสานญฺจายตนปญฺหาสุตฺต}
\proseitem{336}{“อากาสานญฺจายตนํ อากาสานญฺจายตนนฺ”ติ วุจฺจติ, กตมํ นุ โข อากาสานญฺจายตนนฺติ.\csromanspacer{} ตสฺส มยฺหํ อาวุโส เอตทโหสิ}

\gambhira{สฬายตนวคฺคสํยุตฺตปาฬิ}
\prose{\csromanfolio{460}“อิธ ภิกฺขุ สพฺพโส รูปสญฺญานํ สมาติกฺกมา ปฏิฆสญฺญานํ อตฺถงฺคมา นานตฺตสญฺญานํ อมนสิการา ‘อนนฺโต อากาโส’อากาสานญฺจายตนํ อุปสมฺปชฺช วิหรติ, อิทํ วุจฺจติ อากาสานญฺจายตนนฺ”ติ.\csromanspacer{} โส ขฺวาหํ อาวุโส สพฺพโส รูปสญฺญานํ สมติกฺกมา ปฏิฆสญฺญานํ อตฺถิงฺคมา นานตฺตสญฺญานํ อมนสิการา “อนนฺโต อากาโส”ติ อากาสานญฺจายตนํ อุปสมฺปชฺช วิหรามิ, ตสฺส มยฺหํ อาวุโส อิมินา วิหาเรน วิหรโต รูปสหคตา สญฺญามนสิการา สมุทาจรนฺติ.}

\prose{อถ โข มํ อาวุโส ภควา อิทฺธิยา อุปสงฺกมิตฺวา เอตทโวจ “โมคฺคลฺลาน โมคฺคลฺลาน มา พฺราหฺมณ อากาสานญฺจายตนํ ปมาโท, อากาสานญฺจายตเน จิตฺตํ สณฺฐเปหิ, อากาสานญฺจายตเน จิตฺตํ เอโกทิํ กโรหิ, อากาสานญฺจายตเน จิตฺตํ สมาทหา”ติ.\csromanspacer{} โส ขฺวาหํ อาวุโส อปเรน สมเยน สพฺพโส รูปสญฺญานํ สมติกฺกมา ปฏิฆสญฺญานํ อตฺถงฺคมา นานตฺตสญฺญานํ อมนสิการา “อนนฺโต อากาโส”ติ อากาสานญฺจายตนํ อุปสมฺปชฺช วิหาสิํ.\csromanspacer{} ยํ หิ ตํ อาวุโส สมฺมา วทมาโน วเทยฺย ฯเปฯ มหาภิญฺญตํ ปตฺโตติ.\csromanspacer{}\csromanspacer{}ปญฺจมํ.}

\csromanheader{2}{6. วิญฺญาณญฺจายตนปญฺหาสุตฺต}
\proseitem{337}{“วิญฺญาณญฺจายตนํ วิญฺญาณญฺจายตนนฺ”ติ วุจฺจติ, กตมํ นุ โข วิญฺญาณญฺจายตนฺติ.\csromanspacer{} ตสฺส มยฺหํ อาวุโส เอตทโหสิ “อิธ ภิกฺขุ สพฺพโส อากาสานญฺจายตนํ สมติกฺกมฺม ‘อนนฺตํ วิญฺญาณนฺ’ติ วิญฺญาณญฺจายตนํ อุปสมฺปชฺช วิหรติ, อิทํ วุจฺจติ วิญฺญาณญฺจายตนนฺ”ติ.\csromanspacer{} โส ขฺวาหํ อาวุโส สพฺพโส อากาสานญฺจายตนํ สมติกฺกมฺม “อนนฺตํ วิญฺญาณนฺ”ติ วิญฺญาณญฺจายตนํ อุปสมฺปชฺช วิหรามิ, ตสฺส มยฺหํ อาวุโส อิมินา วิหาเรน วิหรโต อากาสานญฺจายตนสหคตา สญฺญามนสิการา สมุทาจรนฺติ.}

\prose{อถ โข มํ อาวุโส ภควา อิทฺธิยา อุปสงฺกมิตฺวา เอตทโวจ “โมคฺคลฺลาน โมคฺคลฺลาน มา พฺราหฺมณ วิญฺญาณญฺจายตนํ ปมาโท, วิญฺญาณญฺจายตเน จิตฺตํ สณฺฐเปหิ, วิญฺญาณญฺจายตเน จิตฺตํ เอโกทิํ กโรหิ, วิญฺญาณญฺจายตเน จิตฺตํ สมาทหา”ติ.\csromanspacer{} โส ขฺวาหํ อาวุโส อปเรน สมเยน สพฺพโส อากาสานญฺจายตนํ สมติกฺกมฺม “อนนฺตํ}

\csromanheader{2}{6. โมคฺคลฺลานสํยุตฺต}
\prose{\csromanfolio{461}วิญฺญาณนฺ”ติ วิญฺญาณญฺจายตนํ อุปสมฺปชฺช วิหาสิํ.\csromanspacer{} ยํ หิ ตํ อาวุโส สมฺมา วทมาโน วเทยฺย ฯเปฯ มหาภิญฺญตํ ปตฺโตติ.\csromanspacer{}\csromanspacer{}ฉฏฺฐํ.}

\csromanheader{2}{7. อากิญฺจญฺญายตนปญฺหาสุตฺต}
\proseitem{338}{“อากิญฺจญฺญายตนํ อากิญฺจญฺญายตนนฺ”ติ วุจฺจติ, กตมํ นุ โข อากิญฺจญฺญายตนนฺติ.\csromanspacer{} ตสฺส มยฺหํ อาวุโส เอตทโหสิ “อิธ ภิกฺขุ สพฺพโส วิญฺญาณญฺจายตนํ สมติกฺกมฺม “นตฺถิ กิญฺจี”ติ อากิญฺจญฺญายตนํ อุปสมฺปชฺช วิหรติ, อิทํ วุจฺจติ อากิญฺจญฺญายตนนฺ”ติ.\csromanspacer{} โส ขฺวาหํ อาวุโส สพฺพโส วิญฺญาณญฺจายตนํ สมติกฺกมฺม “นตฺถิ กิญฺจี”ติ อากิญฺจญฺญายตนํ อุปสมฺปชฺช วิหรามิ, ตสฺส มยฺหํ อาวุโส อิมินา วิหาเรน วิหรโต วิญฺญาณญฺจายตนสหคตา สญฺญามนสิการา สมุทาจรนฺติ.}

\prose{อถ โข มํ อาวุโส ภควา อิทฺธิยา อุปสงฺกมิตฺวา เอตทโวจ “โมคฺคลฺลาน โมคฺคลฺลาน มา พฺราหฺมณ อากิญฺจญฺญายตนํ ปมาโท, อากิญฺจญฺญายตเน จิตฺตํ สณฺฐเปหิ, อากิญฺจญฺญายตเน จิตฺตํ เอโกทิํ กโรหิ, อากิญฺจญฺญายตเน จิตฺตํ สมาทหา”ติ.\csromanspacer{} โส ขฺวาหํ อาวุโส อปเรน สมเยน สพฺพโส วิญฺญาณญฺจายตนํ สมติกฺกมฺม “นตฺถิ กิญฺจี”ติ อากิญฺจญฺญายตนํ อุปสมฺปชฺช วิหาสิํ.\csromanspacer{} ยํ หิ ตํ อาวุโส สมฺมา วทมาโน วเทยฺย ฯเปฯ มหาภิญฺญตํ ปตฺโตติ.\csromanspacer{}\csromanspacer{}สตฺตมํ.}

\csromanheader{2}{8. เนวสญฺญานาสญฺญายตนปญฺหาสุตฺต}
\proseitem{339}{“เนวสญฺญานาสญฺญายตนํ เนวสญฺญานาสญฺญายตนนฺ”ติ วุจฺจติ, กตมํ นุ โข เนวสญฺญานาสญฺญายตนนฺติ.\csromanspacer{} ตสฺส มยฺหํ อาวุโส เอตทโหสิ “อิธ ภิกฺขุ สพฺพโส อากิญฺจญฺญายตนํ สมติกฺกมฺม เนวสญฺญานาสญฺญายตนํ อุปสมฺปชฺช วิหรติ, อิทํ วุจฺจติ เนวสญฺญานาสญฺญายตนนฺ”ติ.\csromanspacer{} โส ขฺวาหํ อาวุโส สพฺพโส อากิญฺจญฺญายตนํ สมติกฺกมฺม เนวสญฺญานาสญฺญายตนํ อุปสมฺปชฺช วิหรามิ, ตสฺส มยฺหํ อาวุโส อิมินา วิหาเรน วิหรโต อากิญฺจญฺญายตนสหคตา สญฺญามนสิการา สมุทาจรนฺติ.}

\gambhira{สฬายตนวคฺคสํยุตฺตปาฬิ}
\prose{\csromanfolio{462}อถ โข มํ อาวุโส ภควา อิทฺธิยา อุปสงฺกมิตฺวา เอตทโวจ “โมคฺคลฺลาน โมคฺคลฺลาน มา พฺราหฺมณ เนวสญฺญานาสญฺญายตนํ ปมาโท, เนวสญฺญานาสญฺญายตเน จิตฺตํ สณฺฐเปหิ, เนวสญฺญานาสญฺญายตเน จิตฺตํ เอโกทิํ กโรหิ, เนวสญฺญานาสญฺญายตเน จิตฺตํ สมาทหา”ติ.\csromanspacer{} โส ขฺวาหํ อาวุโส อปเรน สมเยน สพฺพโส อากิญฺจญฺญายตนํ สมติกฺกมฺม เนวสญฺญานาสญฺญายตนํ อุปสมฺปชฺช วิหาสิํ.\csromanspacer{} ยํ หิ ตํ อาวุโส สมฺมา วทมาโน วเทยฺย ฯเปฯ มหาภิญฺญตํ ปตฺโตติ.\csromanspacer{}\csromanspacer{}อฏฺฐมํ.}

\csromanheader{2}{9. อนิมิตฺตปญฺหาสุตฺต}
\proseitem{340}{“อนิมิตฺโต เจโตสมาธิ อนิมิตฺโต เจโตสมาธี”ติ วุจฺจติ, กตโม นุ โข อนิมิตฺโต เจโตสมาธีติ.\csromanspacer{} ตสฺส มยฺหํ อาวุโส เอตทโหสิ “อิธ ภิกฺขุ สพฺพนิมิตฺตานํ อมนสิการา อนิมิตฺตํ เจโตสมาธิํ อุปสมฺปชฺช วิหรติ, อยํ วุจฺจติ อนิมิตฺโต เจโตสมาธี”ติ.\csromanspacer{} โส ขฺวาหํ อาวุโส สพฺพนิมิตฺตานํ อมนสิการา อนิมิตฺตํ เจโตสมาธิํ อุปสมฺปชฺช วิหรามิ, ตสฺส มยฺหํ อาวุโส อิมินา วิหาเรน วิหรโต นิมิตฺตานุสาริ วิญฺญาณํ โหติ.}

\prose{อถ โข มํ อาวุโส ภควา อิทฺธิยา อุปสงฺกมิตฺวา เอตทโวจ “โมคฺคลฺลาน โมคฺคลฺลาน มา พฺราหฺมณ อนิมิตฺตํ เจโตสมาธิํ ปมาโท, อนิมิตฺเต เจโตสมาธิสฺมิํ จิตฺตํ สณฺฐเปหิ, อนิมิตฺเต เจโตสมาธิสฺมิํ จิตฺตํ เอโกทิํ กโรหิ, อนิมิตฺเต เจเตสมาธิสฺมิํ จิตฺตํ สมาทหา”ติ.\csromanspacer{} โส ขฺวาหํ อาวุโส อปเรน สมเยน สพฺพนิมิตฺตานํ อมนสิการา อนิมิตฺตํ เจโตสมาธิํ อุปสมฺปชฺช วิหาสิํ.\csromanspacer{} ยํ หิ ตํ อาวุโส สมฺมา วทมาโน วเทยฺย “สตฺถารานุคฺคหิโต สาวโก มหาภิญฺญตํ ปตฺโต”ติ.\csromanspacer{} มมํ ตํ สมฺมา วทมาโน วเทยฺย “สตฺถารานุคฺคหิโต สาวโก มหาภิญฺญตํ ปตฺโต”ติ.\csromanspacer{}\csromanspacer{}นวมํ.}

\csromanheader{1}{10. สกฺกสุตฺต}
\csromanmatikaanchor{mtk.258}
\csromantocmarkref{section}{10. สกฺกสุตฺต}{mtk.258}
\bookmark[dest={mtk.258},level=1]{10. สกฺกสุตฺต}
\proseitem{341}{อถ โข อายสฺมา มหาโมคฺคลฺลาโน เสยฺยถาปิ นาม พลวา ปุริโส สมิญฺชิตํ วา พาหํ ปสาเรยฺย, ปสาริตํ วา พาหํ สมิญฺเชยฺย, เอวเมว เชตวเน อนฺตรหิโต เทเวสุ ตาวติํเสสุ}

\csromanheader{2}{6. โมคฺคลฺลานสํยุตฺต}
\prose{\csromanfolio{463}ปาตุรโหสิ.\csromanspacer{} อถ โข สกฺโก เทวานมินฺโท ปญฺจหิ เทวตาสเตหิ สทฺธิํ เยนยสฺมา มหาโมคฺคลฺลาโน เตนุปสงฺกมิ, อุปสงฺกมิตฺวา อายสฺมนฺตํ มหาโมคฺคลฺลานํ อภิวาเทตฺวา เอกมนฺตํ อฏฺฐาสิ, เอกมนฺตํ ฐิตํ โข สกฺกํ เทวานมินฺทํ อายสฺมา มหาโมคฺคลฺลาโน เอตทโวจ\csromandash{}}

\prose{สาธุ โข เทวานมินฺท พุทฺธสรณคมนํ\footnote{พุทฺธํ สรณคมนํ (สี)} โหติ, พุทฺธสรณคมนเหตุ\footnote{พุทฺธํ สรณคมนเหตุ (สี)} โข เทวานมินฺท เอวมิเธกจฺเจ สตฺตา กายสฺส เภทา ปรํ มรณา สุคติํ สคฺคํ โลกํ อุปปชฺชนฺติ.\csromanspacer{} สาธุ โข เทวานมินฺท ธมฺมสรณคมนํ โหติ, ธมฺมสรคมนเหตุ โข เทวานมินฺท เอวมิเธกจฺเจ สตฺตา กายสฺส เภทา ปรํ มรณา สุคติํ สคฺคํ โลกํ อุปปชฺชนฺติ.\csromanspacer{} สาธุ โข เทวานมินฺท สํฆสรณคมนํ โหติ, สํฆสรณคมนเหตุ โข เทวานมินฺท เอวมิเธกจฺเจ สตฺตา กายสฺส เภทา ปรํ มรณา สุคติํ สคฺคํ โลกํ อุปปชฺชนฺตีติ.}

\prose{สาธุ โข มาริส โมคฺคลฺลาน พุทฺธสรณคมนํ โหติ, พุทฺธสรณคมนเหตุ โข มาริส โมคฺคลฺลาน เอวมิเธกจฺเจ สตฺตา กายสฺส เภทา ปรํ มรณา สุคติํ สคฺคํ โลกํ อุปปชฺชนฺติ.\csromanspacer{} สาธุ โข มาริส โมคฺคลฺลาน ธมฺมสรณคมนํ โหติ, ธมฺมสรณคมนเหตุ โข มาริส โมคฺคลฺลาน เอวมิเธกจฺเจ สตฺตา กายสฺส เภทา ปรํ มรณา สุคติํ สคฺคํ โลกํ อุปปชฺชนฺติ.\csromanspacer{} สาธุ โข มาริส โมคฺคลฺลาน สํฆ ฯเปฯ สุคติํ สคฺคํ โลกํ อุปปชฺชนฺตีติ.}

\prose{อถ โข สกฺโก เทวานมินฺโท ฉหิ เทวตาสเตหิ สทฺธิํ ฯเปฯ.\csromanspacer{} อถ โข สกฺโก เทวานมินฺโท สตฺตหิ เทวตาสเตหิ สทฺธิํ ฯเปฯ.\csromanspacer{} อถ โข สกฺโก เทวานมินฺโท อฏฺฐหิ เทวตาสเตหิ สทฺธิํ ฯเปฯ.\csromanspacer{} อถ โข สกฺโก เทวานมินฺโท อสีติยา เทวตาสหสฺเสหิ สทฺธิํ เยนายสฺมา มหาโมคฺคลฺลาโน เตนุปสงฺกมิ, อุปสงฺกมิตฺวา อายสฺมนฺตํ มหาโมคฺคลฺลานํ อภิวาเทตฺวา เอกมนฺตํ อฏฺฐาสิ, เอกมนฺตํ ฐิตํ โข สกฺกํ เทวานมินฺทํ อายสฺมา มหาโมคฺคลฺลาโน เอตทโวจ\csromandash{}}

\prose{สาธุ โข เทวานมินฺท พุทฺธสรณคมนํ โหติ, พุทฺธสรณคมนเหตุ โข เทวานมินฺท เอวมิเธกจฺเจ สตฺตา กายสฺส เภทา ปรํ มรณา สุคติํ สคฺคํ โลกํ อุปปชฺชนฺติ.\csromanspacer{} สาธุ โข เทวานมินฺท ธมฺมสรณคมนํ โหติ,}

\gambhira{สฬายตนวคฺคสํยุตฺตปาฬิ}
\prose{\csromanfolio{464}ธมฺมสรณคมนเหตุ โข เทวานมินฺท เอวมิเธกจฺเจ สตฺตา กายสฺส เภทา ปรํ มรณา สุคติํ สคฺคํ โลกํ อุปปชฺชนฺติ.\csromanspacer{} สาธุ โข เทวานมินฺท สํฆสรณคมนํ โหติ, สํฆสรณคมนเหตุ โข เทวานมินฺท เอวมิเธกจฺเจ สตฺตา กายสฺส เภทา ปรํ มรณา สุคติํ สคฺคํ โลกํ อุปปชฺชนฺตีติ.}

\prose{สาธุ โข มาริส โมคฺคลฺลาน พุทฺธสรณคมนํ โหติ, พุทฺธสรณคมนเหตุ โข มาริส โมคฺคลฺลาน เอวมิเธกจฺเจ สตฺตา กายสฺส เภทา ปรํ มรณา สุคติํ สคฺคํ โลกํ อุปปชฺชนฺติ.\csromanspacer{} สาธุ โข มาริส โมคฺคลฺลาน ธมฺมสรณคมนํ โหติ ฯเปฯ.\csromanspacer{} สาธุ โข มาริส โมคฺคลฺลาน สํฆสรณคมนํ โหติ, สํฆสรณคมนเหตุ โข มาริส โมคฺคลฺลาน เอวมิเธกจฺเจ สตฺตา กายสฺส เภทา ปรํ มรณา สุคติํ สคฺคํ โลกํ อุปปชฺชนฺตีติ.}

\prose{อถ โข สกฺโก เทวานมินฺโท ปญฺจหิ เทวตาสเตหิ สทฺธิํ เยนายสฺมา มหาโมคฺคลฺลาโน เตนุปสงฺกมิ, อุปสงฺกมิตฺวา อายสฺมนฺตํ มหาโมคฺคลฺลานํ อภิวาเทตฺวา เอกมนฺตํ อฏฺฐาสิ, เอกมนฺตํ ฐิตํ โข สกฺกํ เทวานมินฺทํ อายสฺมา มหาโมคฺคลฺลาโน เอตทโวจ\csromandash{}}

\prose{สาธุ โข เทวานมินฺท พุทฺเธ อเวจฺจปฺปสาเทน สมนฺนาคมนํ โหติ “อิติปิ โส ภควา อรหํ สมฺมาสมฺพุทฺโธ วิชฺชาจรณสมฺปนฺโน สุคโต โลกวิทู อนุตฺตโร ปุริสทมฺมสารถิ สตฺถา เทวมนุสฺสานํ พุทฺโธ ภควา”ติ, พุทฺเธ อเวจฺจปฺปสาเทน สมนฺนาคมนเหตุ โข เทวานมินฺท เอวมิเธกจฺเจ สตฺตา กายสฺส เภทา ปรํ มรณา สุคติํ สคฺคํ โลกํ อุปปชฺชนฺติ.}

\prose{สาธุ โข เทวานมินฺท ธมฺเม อเวจฺจปฺปสาเทน สมนฺนาคมนํ โหติ “สฺวากฺขาโต ภควตา ธมฺโม สนฺทิฏฺฐิโก อกาลิโก เอหิปสฺสิโก โอปเนยฺยิโก ปจฺจตฺตํ เวทิตพฺโพ วิญฺญูหี”ติ, ธมฺเม อเวจฺจปฺปสาเทน สมนฺนาคมนเหตุ โข เทวานมินฺท เอวมิเธกจฺเจ สตฺตา กายสฺส เภทา ปรํ มรณา สุคติํ สคฺคํ โลกํ อุปปชฺชนฺติ.}

\prose{สาธุ โข เทวานมินฺท สํเฆ อเวจฺจปฺปสาเทน สมนฺนาคมนํ โหติ “สุปฺปฏิปนฺโน ภควโต สาวกสํโฆ อุชุปฺปฏิปนฺโน ภควโต}

\csromanheader{2}{6. โมคฺคลฺลานสํยุตฺต}
\prose{\csromanfolio{465}สาวกสํโฆ ญายปฺปฏิปนฺโน ภควโต สาวกสํโฆ สามีจิปฺปฏิปนฺโน ภควโต สาวกสํโฆ ยทิทํ จตฺตาริ ปุริสยุคานิ อฏฺฐ ปุริสปุคฺคลา เอส ภควโต สาวกสํโฆ อาหุเนยฺโย ปาหุเนยฺโย ทกฺขิเณยฺโย อญฺชลิกรณีโย อนุตฺตรํ ปุญฺญกฺเขตฺตํ โลกสฺสา”ติ, สํเฆ อเวจฺจปฺปสาเทน สมนฺนาคมนเหตุ โข เทวานมินฺท เอวมิเธกจฺเจ สตฺตา กายสฺส เภทา ปรํ มรณา สุคติํ สคฺคํ โลกํ อุปปชฺชนฺติ.}

\prose{สาธุ โข เทวานมินฺท อริยกนฺเตหิ สีเลหิ สมนฺนาคมนํ โหติ อขณฺเฑหิ อจฺฉิทฺเทหิ อสพเลหิ อกมฺมาเสหิ ภุชิสฺเสหิ วิญฺญุปฺปสตฺเถหิ อปรามฏฺเฐหิ สมาธิสํวตฺตนิเกหิ, อริยกนฺเตหิ สีเลหิ สมนฺนาคมนเหตุ โข เทวานมินฺท เอวมิเธกจฺเจ สตฺตา กายสฺส เภทา ปรํ มรณา สุคติํ สคฺคํ โลกํ อุปปชฺชนฺตีติ.}

\prose{สาธุ โข มาริส โมคฺคลฺลาน พุทฺเธ อเวจฺจปฺปสาเทน สมนฺนาคมนํ โหติ “อิติปิ โส ฯเปฯ สตฺถา เทวมนุสฺสานํ พุทฺโธ ภควา”ติ, พุทฺเธ อเวจฺจปฺปสาเทน สมนฺนาคมนเหตุ โข มาริส โมคฺคลฺลาน เอวมิเธกจฺเจ สตฺตา กายสฺส เภทา ปรํ มรณา สุคติํ สคฺคํ โลกํ อุปปชฺชนฺติ.}

\prose{สาธุ โข มาริส โมคฺคลฺลาน ธมฺเม อเวจฺจปฺปสาเทน สมนฺนาคมนํ โหติ “สฺวากฺขาโต ภควโต ธมฺโม ฯเปฯ ปจฺจตฺตํ เวทิตพฺโพ วิญฺญูหี”ติ ธมฺเม อเวจฺจปฺปสาเทน สมนฺนาคมนเหตุ โข มาริส โมคฺคลฺลาน เอวมิเธกจฺเจ สตฺตา กายสฺส เภทา ปรํ มรณา สุคติํ สคฺคํ โลกํ อุปปชฺชนฺติ.}

\prose{สาธุ โข มาริส โมคฺคลฺลานํ สํเฆ อเวจฺจปฺปสาเทน สมนฺนาคมนํ โหติ “สุปฺปฏิปนฺโน ภควโต สาวกสํโฆ ฯเปฯ อนุตฺตรํ ปุญฺญกฺเขตฺตํ โลกสฺสา”ติ, สํเฆ อเวจฺจปฺปสาเทน สมนฺนาคมนเหตุ โข มาริส โมคฺคลฺลาน เอวมิเธกจฺเจ สตฺตา กายสฺส เภทา ปรํ มรณา สุคติํ สคฺคํ โลกํ อุปปชฺชนฺติ.}

\prose{สาธุ โข มาริส โมคฺคลฺลาน อริยกนฺเตหิ สีเลหิ สมนฺนาคมนํ โหติ อขณฺเฑหิ ฯเปฯ สมาธิสํวตฺตนิเกหิ, อริยกนฺเตหิ สีเลหิ สมนฺนาคมนเหตุ โข มาริส โมคฺคลฺลาน เอวมิเธกจฺเจ สตฺตา กายสฺส เภทา ปรํ มรณา สุคติํ สคฺคํ โลกํ อุปปชฺชนฺตีติ.}

\gambhira{สฬายตนวคฺคสํยุตฺตปาฬิ}
\prose{\csromanfolio{466}อถ โข สกฺโก เทวานมินฺโท ฉหิ เทวตาสเตหิ สทฺธิํ ฯเปฯ.\csromanspacer{} อถ โข สกฺโก เทวานมินฺโท สตฺตหิ เทวตาสเตหิ สทฺธิํ ฯเปฯ.\csromanspacer{} อถ โข สกฺโก เทวานมินฺโท อฏฺฐหิ เทวตาสเตหิ สทฺธิํ ฯเปฯ.\csromanspacer{} อถ โข สกฺโก เทวานมินฺโท อสีติยา เทวตาสหสฺเสหิ สทฺธิํ เยนายสฺมา มหาโมคฺคลฺลาโน เตนุปสงฺกมิ, อุปสงฺกมิตฺวา อายสฺมนฺตํ มหาโมคฺคลฺลานํ อภิวาเทตฺวา เอกมนฺตํ อฏฺฐาสิ.\csromanspacer{} เอกมนฺตํ ฐิตํ โข สกฺกํ เทวานมินฺทํ อายสฺมา มหาโมคฺคลฺลาโน เอตทโวจ\csromandash{}}

\prose{สาธุ โข เทวานมินฺท พุทฺเธ อเวจฺจปฺปสาเทน สมนฺนาคมนํ โหติ “อิติปิ โส ภควา ฯเปฯ สตฺถา เทวมนุสฺสานํ พุทฺโธ ภควา”ติ, พุทฺเธ อเวจฺจปฺปสาเทน สมนฺนาคมนเหตุ โข เทวานมินฺท เอวมิเธกจฺเจ สตฺตา กายสฺส เภทา ปรํ มรณา สุคติํ สคฺคํ โลกํ อุปปชฺชนฺติ.}

\prose{สาธุ โข เทวานมินฺท ธมฺเม อเวจฺจปฺปสาเทน สมนฺนาคมนํ โหติ “สฺวากฺขาโต ภควตา ธมฺโม ฯเปฯ ปจฺจตฺตํ เวทิตพฺโพ วิญฺญูหี”ติ, ธมฺเม อเวจฺจปฺปสาเทน สมนฺนาคมนเหตุ โข เทวานมินฺท เอวมิเธกจฺเจ สตฺตา กายสฺส เภทา ปรํ มรณา สุคติํ สคฺคํ โลกํ อุปปชฺชนฺติ.}

\prose{สาธุ โข เทวานมินฺท สํเฆ อเวจฺจปฺปสาเทน สมนฺนาคมนํ โหติ “สุปฺปฏิปนฺโน ภควโต สาวกสํโฆ ฯเปฯ อนุตฺตรํ ปุญฺญกฺเขตฺตํ โลกสฺสา”ติ, สํเฆ อเวจฺจปฺปสาเทน สมนฺนาคมนเหตุ โข เทวานมินฺท เอวมิเธกจฺเจ สตฺตา กายสฺส เภทา ปรํ มรณา สุคติํ สคฺคํ โลกํ อุปปชฺชนฺติ.}

\prose{สาธุ โข เทวานมินฺท อริยกนฺเตหิ สีเลหิ สมนฺนาคมนํ โหติ อขณฺเฑหิ ฯเปฯ สมาธิสํวตฺตนิเกหิ, อริยกนฺเตหิ สีเลหิ สมนฺนาคมนเหตุ โข เทวานมินฺท เอวมิเธกจฺเจ สตฺตา กายสฺส เภทา ปรํ มรณา สุคติํ สคฺคํ โลกํ อุปปชฺชนฺตีติ.}

\prose{สาธุ โข มาริส โมคฺคลฺลาน พุทฺเธ อเวจฺจปฺปสาเทน สมนฺนาคมนํ โหติ “อิติปิ โส ภควา ฯเปฯ สตฺถา เทวมนุสฺสานํ พุทฺโธ ภควา”ติ, พุทฺเธ อเวจฺจปฺปสาเทน สมนฺนาคมนเหตุ โข มาริส โมคฺคลฺลาน เอวมิเธกจฺเจ สตฺตา กายสฺส เภทา ปรํ มรณา สุคติํ สคฺคํ โลกํ อุปปชฺชนฺติ.}

\csromanheader{2}{6. โมคฺคลฺลานสํยุตฺต}
\prose{\csromanfolio{467}สาธุ โข มาริส โมคฺคลฺลาน ธมฺเม อเวจฺจปฺปสาเทน สมนฺนาคมนํ โหติ “สฺวากฺขาโต ภควตา ธมฺโม ฯเปฯ ปจฺจตฺตํ เวทิตพฺโพ วิญฺญูหี”ติ, ธมฺเม อเวจฺจปฺปสาเทน สมนฺนาคมนเหตุ โข มาริส โมคฺคลฺลาน เอวเธกจฺเจ สตฺตา กายสฺส เภทา ปรํ มรณา มรณา สุคติํ สคฺคํ โลกํ อุปปชฺชนฺติ.}

\prose{สาธุ โข มาริส โมคฺคลฺลาน สํเฆ อเวจฺจปฺปสาเทน สมนฺนาคมนํ โหติ “สุปฺปฏิปนฺโน ภควโต สาวกสํโฆ ฯเปฯ อนุตฺตรํ ปุญฺญกฺเขตฺตํ โลกสฺสา”ติ, สํเฆ อเวจฺจปฺปสาเทน สมนฺนาคมนเหตุ โข มาริส โมคฺคลฺลาน เอวมิเธกจฺเจ สตฺตา กายสฺส เภทา ปรํ มรณา สุคติํ สคฺคํ โลกํ อุปปชฺชนฺติ.}

\prose{สาธุ โข มาริส โมคฺคลฺลาน อริยกนฺเตหิ สีเลหิ สมนฺนาคมนํ โหติ อขณฺเฑหิ ฯเปฯ สมาธิสํวตฺตนิเกหิ, อริยกนฺเตหิ สีเลหิ สมนฺนาคมนเหตุ โข มาริส โมคฺคลฺลาน เอวมิเธกจฺเจ สตฺตา กายสฺส เภทา ปรํ มรณา สุคติํ สคฺคํ โลกํ อุปปชฺชนฺตีติ.}

\prose{อถ โข สกฺโก เทวานมินฺโท ปญฺจหิ เทวตาสเตหิ สทฺธิํ เยนายสฺมา มหาโมคฺคลฺลาโน เตนุปสงฺกมิ ฯเปฯ เอกมนฺตํ ฐิตํ โข สกฺกํ เทวานมินฺทํ อายสฺมา มหาโมคฺคลฺลาโน เอตทโวจ\csromandash{}}

\prose{สาธุ โข เทวานมินฺท พุทฺธสรณคมนํ โหติ, พุทฺธสรณคมนเหตุ โข เทวานมินฺท เอวมิเธกจฺเจ สตฺตา กายสฺส เภทา ปรํ มรณา สุคติํ สคฺคํ โลกํ อุปปชฺชนฺติ, เต อญฺเญ เทเว ทสหิ ฐาเนหิ อธิคณฺหนฺติ ทิพฺเพน อายุนา ทิพฺเพน วณฺเณน ทิพฺเพน สุเขน ทิพฺเพน ยเสน ทิพฺเพน อาธิปเตยฺเยน ทิพฺเพหิ รูเปหิ ทิพฺเพหิ สทฺเทหิ ทิพฺเพหิ คนฺเธหิ ทิพฺเพหิ รเสหิ ทิพฺเพหิ โผฏฺฐพฺเพหิ.}

\prose{สาธุ โข เทวานมินฺท ธมฺมสรณคมนํ โหติ, ธมฺมสรณคมนเหตุ โข เทวานมินฺท เอวมิเธกจฺเจ สตฺตา กายสฺส เภทา ปรํ มรณา สุคติํ สคฺคํ โลกํ อุปปชฺชนฺติ, เต อญฺเญ เทเว ทสหิ ฐาเนหิ อธิคณฺหนฺติ ทิพฺเพน อายุนา ทิพฺเพน วณฺเณน ทิพฺเพน สุเขน ทิพฺเพน ยเสน ทิพฺเพน อาธิปเตยฺเยน ทิพฺเพหิ รูเปหิ ทิพฺเพหิ สทฺเทหิ ทิพฺเพหิ คนฺเธหิ ทิพฺเพหิ รเสหิ ทิพฺเพหิ โผฏฺฐพฺเพหิ.}

\gambhira{สฬายตนวคฺคสํยุตฺตปาฬิ}
\prose{\csromanfolio{468}สาธุ โข เทวานมินฺท สํฆสรณคมนํ โหติ, สํฆสรณคมนเหตุ โข เทวานมินฺท เอวมิเธกจฺเจ สตฺตา กายสฺส เภทา ปรํ มรณา สุคติํ สคฺคํ โลกํ อุปปชฺชนฺติ, เต อญฺเญ เทเว ทสหิ ฐาเนหิ อธิคณฺหนฺติ ทิพฺเพน อายุนา ทิพฺเพน วณฺเณน ทิพฺเพน สุเขน ทิพฺเพน ยเสน ทิพฺเพน อาธิปเตยฺเยน ทิพฺเพหิ รูเปหิ ทิพฺเพหิ สทฺเทหิ ทิพฺเพหิ คนฺเธหิ ทิพฺเพหิ รเสหิ ทิพฺเพหิ โผฏฺฐพฺเพหีติ.}

\prose{สาธุ โข มาริส โมคฺคลฺลาน พุทฺธสรณคมนํ โหติ, พุทฺธสรณคมนเหตุ โข มาริส โมคฺคลฺลาน เอวมิเธกจฺเจ สตฺตา กายสฺส เภทา ปรํ มรณา สุคติํ สคฺคํ โลกํ อุปปชฺชนฺติ, เต อญฺเญ เทเว ทสหิ ฐาเนหิ อธิคณฺหนฺติ ทิพฺเพน อายุนา ฯเปฯ ทิพฺเพหิ โผฏฺฐพฺเพหิ.}

\prose{สาธุ โข มาริส โมคฺคลฺลาน ธมฺมสรณคมนํ โหติ ฯเปฯ.}

\prose{สาธุ โข มาริส โมคฺคลฺลาน สํฆสรณคมนํ โหติ, สํฆสรณคมนเหตุ โข มาริส โมคฺคลฺลาน เอวมิเธกจฺเจ สตฺตา กายสฺส เภทา ปรํ มรณา สุคติํ สคฺคํ โลกํ อุปปชฺชนฺติ, เต อญฺเญ เทเว ทสหิ ฐาเนหิ อธิคณฺหนฺติ ทิพฺเพน อายุนา ทิพฺเพน วณฺเณน ทิพฺเพน สุเขน ทิพฺเพน ยเสน ทิพฺเพน อาธิปเตยฺเยน ทิพฺเพหิ รูเปหิ ทิพฺเพหิ สทฺเทหิ ทิพฺเพหิ คนฺเธหิ ทิพฺเพหิ รเสหิ ทิพฺเพหิ โผฏฺฐพฺเพหีติ.}

\prose{อถ โข สกฺโก เทวานมินฺโท ฉหิ เทวตาสเตหิ สทฺธิํ ฯเปฯ.\csromanspacer{} อถ โข สกฺโก เทวานมินฺโท สตฺตหิ เทวตาสเตหิ ฯเปฯ.\csromanspacer{} อถ โข สกฺโก เทวานมินฺโท อฏฺฐหิ เทวตาสเตหิ สทฺธิํ ฯเปฯ.\csromanspacer{} อถ โข สกฺโก เทวานมินฺโท อสีติยา เทวตาสหสฺเสหิ สทฺธิํ เยนายสฺมา มหาโมคฺคลฺลาโน เตนุปสงฺกมิ, อุปสงฺกมิตฺวา อายสฺมนฺตํ มหาโมคฺคลฺลานํ อภิวาเทตฺวา เอกมนฺตํ อฏฺฐาสิ.\csromanspacer{} เอกมนฺตํ ฐิตํ โข สกฺกํ เทวานมินฺทํ อายสฺมา มหาโมคฺคลฺลาโน เอตทโวจ\csromandash{}}

\prose{สาธุ โข เทวานมินฺท พุทฺธสรณคมนํ โหติ, พุทฺธสรณคมนเหตุ โข เทวานมินฺท เอวมิเธกจฺเจ สตฺตา กายสฺส เภทา ปรํ มรณา สุคติํ สคฺคํ โลกํ อุปปชฺชนฺติ, เต อญฺเญ เทเว ทสหิ ฐาเนหิ อธิคณฺหนฺติ ทิพฺเพน อายุนา ฯเปฯ ทิพฺเพหิ โผฏฺฐพฺเพหิ.}

\prose{สาธุ โข เทวานมินฺท ธมฺมสรณคมนํ โหติ ฯเปฯ.}

\csromanheader{2}{6. โมคฺคลฺลานสํยุตฺต}
\prose{\csromanfolio{469}สาธุ โข เทวานมินฺท สํฆสรณคมนํ โหติ, สํฆสรณคมนเหตุ โข เทวานมินฺท เอวมิเธกจฺเจ สตฺตา กายสฺส เภทา ปรํ มรณา สุคติํ สคฺคํ โลกํ อุปปชฺชนฺติ, เต อญฺเญ เทเว ทสหิ ฐาเนหิ อธิคณฺหนฺติ ทิพฺเพน อายุนา ทิพฺเพน วณฺเณนา ทิพฺเพน สุเขน ทิพฺเพน ยเสน ทิพฺเพน อาธิปเตยฺเยน ทิพฺเพหิ รูเปหิ ทิพฺเพหิ สทฺเทหิ ทิพฺเพหิ คนฺเธหิ ทิพฺเพหิ รเสหิ ทิพฺเพหิ โผฏฺฐพฺเพหีติ.}

\prose{สาธุ โข มาริส โมคฺคลฺลาน พุทฺธสรณคมนํ โหติ ฯเปฯ.\csromanspacer{} สาธุ โข มาริส โมคฺคลฺลาน ธมฺมสรณคมนํ โหติ ฯเปฯ.\csromanspacer{} สาธุ โข มาริส โมคฺคลฺลาน สํฆสรณคมนํ โหติ, สํฆสรณคมนเหตุ โข มาริส โมคฺคลฺลาน เอวมิเธกจฺเจ สตฺตา กายสฺส เภทา ปรํ มรณา สุคติํ สคฺคํ โลกํ อุปปชฺชนฺติ, เต อญฺเญ เทเว ทสหิ ฐาเนหิ อธิคณฺหนฺติ ทิพฺเพน อายุนา ทิพฺเพน วณฺเณน ทิพฺเพน สุเขน ทิพฺเพน ยเสน ทิพฺเพน อาธิปเตยฺเยน ทิพฺเพหิ รูเปหิ ทิพฺเพหิ สทฺเทหิ ทิพฺเพหิ คนฺเธหิ ทิพฺเพหิ รเสหิ ทิพฺเพหิ โผฏฺฐพฺเพหีติ.}

\prose{อถ โข สกฺโก เทวานมินฺโท ปญฺจหิ เทวตาสเตหิ สทฺธิํ เยนยสฺมา มหาโมคฺคลฺลาโน เตนุปสงฺกมิ, อุปสงฺกมิตฺวา อายสฺมนฺตํ มหาโมคฺคลฺลานํ อภิวาเทตฺวา เอกมนฺตํ อฏฺฐาสิ.\csromanspacer{} เอกมนฺตํ ฐิตํ โข สกฺกํ เทวานมินฺทํ อายสฺมา มหาโมคฺคลฺลาโน เอตทโวจ\csromandash{}}

\prose{สาธุ โข เทวานมินฺท พุทฺเธ อเวจฺจปฺปสาเทน สมนฺนาคมนํ โหติ “อิติปิ โส ภควา ฯเปฯ สตฺถา เทวมนุสฺสานํ พุทฺโธ ภควา”ติ, พุทฺเธ อเวจฺจปฺปสาเทน สมนฺนาคมนเหตุ โข เทวานมินฺท เอวมิเธกจฺเจ สตฺตา กายสฺส เภทา ปรํ มรณา สุคติํ สคฺคํ โลกํ อุปปชฺชนฺติ, เต อญฺเญ เทเว ทสหิ ฐาเนหิ อธิคณฺหนฺติ ทิพฺเพน อายุนา ฯเปฯ ทิพฺเพหิ โผฏฺฐพฺเพหิ.}

\prose{สาธุ โข เทวานมินฺท ธมฺเม อเวจฺจปฺปสาเทน สมนฺนาคมนํ โหติ “สฺวากฺขาโต ภควตา ธมฺโม ฯเปฯ ปจฺจตฺตํ เวทิตพฺโพ วิญฺญูหี”ติ, ธมฺเม อเวจฺจปฺปสาเทน สมนฺนาคมนเหตุ โข เทวานมินฺท เอวมิเธกจฺเจ สตฺตา กายสฺส เภทา ปรํ มรณา สุคติํ สคฺคํ โลกํ อุปปชฺชนฺติ ฯเปฯ.}

\prose{สาธุ โข เทวานมินฺท สํเฆ อเวจฺจปฺปสาเทน สมนฺนาคมนํ โหติ “สุปฺปฏิปนฺโน ภควโต สาวกสํโฆ ฯเปฯ โลกสฺสา”ติ, สํเฆ}

\gambhira{สฬายตนวคฺคสํยุตฺตปาฬิ}
\prose{\csromanfolio{470}อเวจฺจปฺปสาเทน สมนฺนาคมนเหตุ โข เทวานมินฺท เอวมิเธกจฺเจ สตฺตา กายสฺส เภทา ปรํ มรณา สุคติํ สคฺคํ โลกํ อุปปชฺชนฺติ ฯเปฯ.}

\prose{สาธุ โข เทวานมินฺท อริยกนฺเตหิ สีเลหิ สมนฺนาคมนํ โหติ อขณฺเฑหิ ฯเปฯ สมาธิสํวตฺตนิเกหิ, อริยกนฺเตหิ สีเลหิ สมนฺนาคมนเหตุ โข เทวานมินฺท เอวมิเธกจฺเจ สตฺตา กายสฺส เภทา ปรํ มรณา สุคติํ สคฺคํ โลกํ อุปปชฺชนฺติ, เต อญฺเญ เทเว ทสหิ ฐาเนหิ อธิคณฺหนฺติ ทิพฺเพน อายุนา ฯเปฯ ทิพฺเพหิ โผฏฺฐพฺเพหีติ.}

\prose{สาธุ โข มาริส โมคฺคลฺลาน พุทฺเธ อเวจฺจปฺปสาเทน สมนฺนาคมนํ โหติ “อิติปิ โส ภควา ฯเปฯ สตฺถา เทวมนุสฺสานํ พุทฺโธ ภควา”ติ, พุทฺเธ อเวจฺจปฺปสาเทน สมนฺนาคมนเหตุ โข มาริส โมคฺคลฺลาน เอวมิเธกจฺเจ สตฺตา กายสฺส เภทา ปรํ มรณา สุคติํ สคฺคํ โลกํ อุปปชฺชนฺติ, เต อญฺเญ เทเว ทสหิ ฐาเนหิ อธิคณฺหนฺติ ทิพฺเพน อายุนา ฯเปฯ ทิพฺเพหิ โผฏฺฐพฺเพหิ.}

\prose{สาธุ โข มาริส โมคฺคลฺลาน ธมฺเม อเวจฺจปฺปสาเทน สมนฺนาคมนํ โหติ “สฺวากฺขาโต ภควตา ธมฺโม ฯเปฯ ปจฺจตฺตํ เวทิตพฺโพ วิญฺญูหี”ติ, ธมฺเม อเวจฺจปฺปสาเทน สมนฺนาคมนเหตุ โข มาริส โมคฺคลฺลาน เอวมิเธกจฺเจ สตฺตา กายสฺส เภทา ปรํ มรณา สุคติํ สคฺคํ โลกํ อุปปชฺชนฺติ, เต อญฺเญ เทเว ทสหิ ฐาเนหิ อธิคณฺหนฺติ ทิพฺเพน อายุนา ฯเปฯ ทิพฺเพหิ โผฏฺฐพฺเพหิ.}

\prose{สาธุ โข มาริส โมคฺคลฺลานํ สํเฆ อเวจฺจปฺปสาเทน สมนฺนาคมนํ โหติ “สุปฺปฏิปนฺโน ภควโต สาวกสํโฆ ฯเปฯ โลกสฺสา”ติ, สํเฆ อเวจฺจปฺปสาเทน สมนฺนาคมนเหตุ โข มาริส โมคฺคลฺลาน เอวมิเธกจฺเจ สตฺตา กายสฺส เภทา ปรํ มรณา สุคติํ สคฺคํ โลกํ อุปปชฺชนฺติ ฯเปฯ.}

\prose{สาธุ โข มาริส โมคฺคลฺลาน อริยกนฺเตหิ สีเลหิ สมนฺนาคมนํ โหติ อขณฺเฑหิ ฯเปฯ สมาธิสํวตฺตนิเกหิ, อริยกนฺเตหิ สีเลหิ สมนฺนาคมนเหตุ โข มาริส โมคฺคลฺลาน เอวมิเธกจฺเจ สตฺตา กายสฺส เภทา ปรํ มรณา สุคติํ สคฺคํ โลกํ อุปปชฺชนฺติ, เต อญฺเญ เทเว ทสหิ ฐาเนหิ อธิคณฺหนฺติ ทิพฺเพน อายุนา ฯเปฯ ทิพฺเพหิ โผฏฺฐพฺเพหีติ.}

\csromanheader{2}{6. โมคฺคลฺลานสํยุตฺต}
\prose{\csromanfolio{471}อถ โข สกฺโก เทวานมินฺโท ฉหิ เทวตาสเตหิ สทฺธิํ ฯเปฯ.\csromanspacer{} อถ โข สกฺโก เทวานมินฺโท สตฺตหิ เทวตาสเตหิ สทฺธิํ ฯเปฯ.\csromanspacer{} อถ โข สกฺโก เทวานมินฺโท อฏฺฐหิ เทวตาสเตหิ สทฺธิํ ฯเปฯ.\csromanspacer{} อถ โข สกฺโก เทวานมินฺโท อสีติยา เทวตาสหสฺเสหิ สทฺธิํ เยนายสฺมา มหาโมคฺคลฺลาโน เตนุปสงฺกมิ, อุปสงฺกมิตฺวา อายสฺมนฺตํ มหาโมคฺคลฺลานํ อภิวาเทตฺวา เอกมนฺตํ อฏฺฐาสิ.\csromanspacer{} เอกมนฺตํ ฐิตํ โข สกฺกํ เทวานมินฺทํ อายสฺมา มหาโมคฺคลฺลาโน เอตทโวจ\csromandash{}}

\prose{สาธุ โข เทวานมินฺท พุทฺเธ อเวจฺจปฺปสาเทน สมนฺนาคมนํ โหติ “อิติปิ โส ภควา อรหํ สมฺมาสมฺพุทฺโธ วิชฺชาจรณสมฺปนฺโน สุคโต โลกวิทู อนุตฺตโร ปุริสทมฺมสารถิ สตฺถา เทวมนุสฺสานํ พุทฺโธ ภควา”ติ, พุทฺเธ อเวจฺจปฺปสาเทน สมนฺนาคมนเหตุ โข เทวานมินฺท เอวมิเธกจฺเจ สตฺตา กายสฺส เภทา ปรํ มรณา สุคติํ สคฺคํ โลกํ อุปปชฺชนฺติ, เต อญฺเญ เทเว ทสหิ ฐาเนหิ อธิคณฺหนฺติ ทิพฺเพน อายุนา ทิพฺเพน วณฺเณน ทิพฺเพน สุเขน ทิพฺเพน ยเสน ทิพฺเพน อาธิปเตยฺเยน ทิพฺเพหิ รูเปหิ ทิพฺเพหิ สทฺเทหิ ทิพฺเพหิ คนฺเธหิ ทิพฺเพหิ รเสหิ ทิพฺเพหิ โผฏฺฐพฺเพหิ.}

\prose{สาธุ โข เทวานมินฺท ธมฺเม อเวจฺจปฺปสาเทน สมนฺนาคมนํ โหติ “สฺวากฺขาโต ภควตา ธมฺโม สนฺทิฏฺฐิโก อกาลิโก เอหิปสฺสิโก โอปเนยฺยิโก ปจฺจตฺตํ เวทิตพฺโพ วิญฺญูหี”ติ, ธมฺเม อเวจฺจปฺปสาเทน สมนฺนาคมนเหตุ โข เทวานมินฺท เอวมิเธกจฺเจ สตฺตา กายสฺส เภทา ปรํ มรณา สุคติํ สคฺคํ โลกํ อุปปชฺชนฺติ, เต อญฺเญ เทเว ทสหิ ฐาเนหิ อธิคณฺหนฺติ ทิพฺเพน อายุนา ฯเปฯ ทิพฺเพหิ โผฏฺฐพฺเพหิ.}

\prose{สาธุ โข เทวานมินฺท สํเฆ อเวจฺจปฺปสาเทน สมนฺนาคมนํ โหติ “สุปฺปฏิปนฺโน ภควโต สาวกสํโฆ อุชุปฺปฏิปนฺโน ภควโต สาวกสํโฆ ญายปฺปฏิปนฺโน ภควโต สาวกสํโฆ สามีจิปฺปฏิปนฺโน ภควโต สาวกสํโฆ ยทิทํ จตฺตาริ ปุริสยุคานิ อฏฺฐ ปุริสปุคฺคลา เอส ภควโต สาวกสํโฆ อาหุเนยฺโย ปาหุเนยฺโย ทกฺขิเณยฺโย อญฺชลิกรณีโย อนุตฺตรํ ปุญฺญกฺเขตฺตํ โลกสฺสา”ติ, สํเฆ อเวจฺจปฺปสาเทน สมนฺนาคมนเหตุ โข เทวานมินฺท เอวมิเธกจฺเจ สตฺตา กายสฺส เภทา ปรํ มรณา สุคติํ สคฺคํ โลกํ อุปปชฺชนฺติ,}

\gambhira{สฬายตนวคฺคสํยุตฺตปาฬิ}
\prose{\csromanfolio{472}เต อญฺเญ เทเว ทสหิ ฐาเนหิ อธิคณฺหนฺติ ทิพฺเพน อายุนา ฯเปฯ ทิพฺเพหิ โผฏฺฐพฺเพหิ.}

\prose{สาธุ โข เทวานมินฺท อริยกนฺเตหิ สีเลหิ สมนฺนาคมนํ โหติ อขณฺเฑหิ อจฺฉิทฺเทหิ อสพเลหิ อกมฺมาเสหิ ภุชิสฺเสหิ วิญฺญุปฺปสตฺเถหิ อปรามฏฺเฐหิ สมาธิสํวตฺตนิเกหิ, อริยกนฺเตหิ สีเลหิ สมนฺนาคมนเหตุ โข เทวานมินฺท เอวมิเธกจฺเจ สตฺตา กายสฺส เภทา ปรํ มรณา สุคติํ สคฺคํ โลกํ อุปปชฺชนฺติ, เต อญฺเญ เทเว ทสหิ ฐาเนหิ อธิคณฺหนฺติ ทิพฺเพน อายุนา ทิพฺเพน วณฺเณน ทิพฺเพน สุเขน ทิพฺเพน ยเสน ทิพฺเพน อาธิปเตยฺเยน ทิพฺเพหิ รูเปหิ ทิพฺเพหิ สทฺเทหิ ทิพฺเพหิ คนฺเธหิ ทิพฺเพหิ รเสหิ ทิพฺเพหิ โผฏฺฐพฺเพหีติ.}

\prose{สาธุ โข มาริส โมคฺคลฺลาน พุทฺเธ อเวจฺจปฺปสาเทน สมนฺนาคมนํ โหติ “อิติปิ โส ภควา ฯเปฯ สตฺถา เทวมนุสฺสานํ พุทฺโธ ภควา”ติ, พุทฺเธ อเวจฺจปฺปสาเทน สมนฺนาคมนเหตุ โข มาริส โมคฺคลฺลาน เอวมิเธกจฺเจ สตฺตา กายสฺส เภทา ปรํ มรณา สุคติํ สคฺคํ โลกํ อุปปชฺชนฺติ, เต อญฺเญ เทเว ทสหิ ฐาเนหิ อธิคณฺหนฺติ ทิพฺเพน อายุนา ฯเปฯ ทิพฺเพหิ โผฏฺฐพฺเพหิ.}

\prose{สาธุ โข มาริส โมคฺคลฺลาน ธมฺเม อเวจฺจปฺปสาเทน สมนฺนาคมนํ โหติ “สฺวากฺขาโต ภควตา ธมฺโม ฯเปฯ ปจฺจตฺตํ เวทิตพฺโพ วิญฺญูหี”ติ, ธมฺเม อเวจฺจปฺปสาเทน สมนฺนาคมนเหตุ โข มาริส โมคฺคลฺลาน เอวมิเธกจฺเจ สตฺตา กายสฺส เภทา ปรํ มรณา สุคติํ สคฺคํ โลกํ อุปปชฺชนฺติ, เต อญฺเญ เทเว ทสหิ ฐาเนหิ อธิคณฺหนฺติ ทิพฺเพน อายุนา ฯเปฯ ทิพฺเพหิ โผฏฺฐพฺเพหิ.}

\prose{สาธุ โข มาริส โมคฺคลฺลาน สํเฆ อเวจฺจปฺปสาเทน สมนฺนาคมนํ โหติ “สุปฺปฏิปนฺโน ภควโต สาวกสํโฆ ฯเปฯ อนุตฺตรํ ปุญฺญกฺเขตฺตํ โลกสฺสา”ติ, สํเฆ อเวจฺจปฺปสาเทน สมนฺนาคมนเหตุ โข มาริส โมคฺคลฺลาน เอวมิเธกจฺเจ สตฺตา กายสฺส เภทา ปรํ มรณา สุคติํ สคฺคํ โลกํ อุปปชฺชนฺติ, เต อญฺเญ เทเว ทสหิ ฐาเนหิ อธิคณฺหนฺติ ทิพฺเพน อายุนา ฯเปฯ ทิพฺเพหิ โผฏฺฐพฺเพหิ.}

\prose{สาธุ โข มาริส โมคฺคลฺลาน อริยกนฺเตหิ สีเลหิ สมนฺนาคมนํ โหติ อขณฺเฑหิ ฯเปฯ สมาธิสํวตฺตนิเกหิ, อริยกนฺเตหิ สีเลหิ}

\csromanheader{2}{6. โมคฺคลฺลานสํยุตฺต}
\prose{\csromanfolio{473}สมนฺนาคมนเหตุ โข มาริส โมคฺคลฺลาน เอวมิเธกจฺเจ สตฺตา กายสฺส เภทา ปรํ มรณา สุคติํ สคฺคํ โลกํ อุปปชฺชนฺติ, เต อญฺเญ เทเว ทสหิ ฐาเนหิ อธิคณฺหนฺติ ทิพฺเพน อายุนา ทิพฺเพน วณฺเณน ทิพฺเพน สุเขน ทิพฺเพน ยเสน ทิพฺเพน อาธิปเตยฺเยน ทิพฺเพหิ รูเปหิ ทิพฺเพหิ สทฺเทหิ คนฺเธหิ ทพฺเพหิ รเสหิ ทิพฺเพหิ โผฏฺฐพฺเพหีติ.\csromanspacer{}\csromanspacer{}ทสมํ.}

\csromanheader{1}{11. จนฺทนสุตฺต}
\csromanmatikaanchor{mtk.259}
\csromantocmarkref{section}{11. จนฺทนสุตฺต}{mtk.259}
\bookmark[dest={mtk.259},level=1]{11. จนฺทนสุตฺต}
\proseitem{342}{อถ โข จนฺทโน\footnote{นนฺทโน (สี)} เทวปุตฺโต ฯเปฯ.}

\prose{อถ โข สุยาโม เทวปุตฺโต ฯเปฯ.}

\prose{อถ โข สนฺตุสิโต เทวปุตฺโต ฯเปฯ.}

\prose{อถ โข สุนิมฺมิโต เทวปุตฺโต ฯเปฯ.}

\prose{อถ โข วสวตฺติ เทวปุตฺโต ฯเปฯ. (ยถา สกฺกสุตฺตํ, ตถา อิเม ปญฺจ เปยฺยาลา วิตฺถาเรตพฺพา.).\csromanspacer{} เอกาทสมํ.}

\vspace{\tipitakaparskip}
\csromancenter{\textbf{โมคฺคลฺลานสํยุตฺตํ สมตฺตํ.}}
\titlehead{\textbf{ตสฺสุทฺทานํ}}
\csromangathameasure{สวิตกฺกาวิตกฺกญฺจ, สุเขน จ อุเปกฺขโก. \\ อากาสญฺเจว วิญฺญาณํ, อากิญฺจํ เนวสญฺญินา. \\ อนิมิตฺโต จ สกฺโก จ, จนฺทเนกาทเสน จาติ.}
\csromangathagroup{\csromangathastanza{สวิตกฺกาวิตกฺกญฺจ, สุเขน จ อุเปกฺขโก. \\ อากาสญฺเจว วิญฺญาณํ, อากิญฺจํ เนวสญฺญินา. \\ อนิมิตฺโต จ สกฺโก จ, จนฺทเนกาทเสน จาติ.}}

\chapterheadpageread{7. จิตฺตสํยุตฺต}
\csromanmatikaanchor{mtk.260}
\csromantocmarknopage{chapter}{7. จิตฺตสํยุตฺต}{mtk.260}
\bookmark[dest={mtk.260},level=0]{7. จิตฺตสํยุตฺต}
\csromanheader{1}{1. สํโยชนสุตฺต}
\csromanmatikaanchor{mtk.261}
\csromantocmarkref{section}{1. สํโยชนสุตฺต}{mtk.261}
\bookmark[dest={mtk.261},level=1]{1. สํโยชนสุตฺต}
\proseitem{343}{\csromanfolio{474}เอกํ สมยํ สมฺพหุลา เถรา ภิกฺขู มจฺฉิกาสณฺเฑ วิหรนฺติ อมฺพาฏกวเน.\csromanspacer{} เตน โข ปน สมเยน สมฺพหุลานํ เถรานํ ภิกฺขูนํ ปจฺฉาภตฺตํ ปิณฺฑปาตปฏิกฺกนฺตานํ มณฺฑลมาเฬ สนฺนิสินฺนานํ สนฺนิปติตานํ อยมนฺตรากถา อุทปาทิ “สํโยชนนฺติ วา อาวุโส สํโยชนิยา ธมฺมาติ วา อิเม ธมฺมา นานตฺถา นานาพฺยญฺชนา, อุทาหุ เอกตฺถา พฺยญฺชนเมว นานนฺ”ติ.\csromanspacer{} ตเตฺรกจฺเจหิ เถเรหิ ภิกฺขูหิ เอวํ พฺยากตํ โหติ “สํโยชนนฺติ วา อาวุโส สํโยชนิยา ธมฺมาติ วา อิเม ธมฺมา นานตฺถา เจว นานาพฺยญฺชนา จา”ติ.\csromanspacer{} เอกจฺเจหิ เถเรหิ ภิกฺขูหิ เอวํ พฺยากตํ โหติ “สํโยชนนฺติ วา อาวุโส สํโยชนิยา ธมฺมาติ วา อิเม ธมฺมา เอกตฺถา พฺยญฺชนเมว นานนฺ”ติ.}

\prose{เตน โข ปน สมเยน จิตฺโต คหปติ มิคปถกํ อนุปฺปตฺโต โหติ เกนจิเทว กรณีเยน.\csromanspacer{} อสฺโสสิ โข จิตฺโต คหปติ “สมฺพหุลานํ กิร เถรานํ ภิกฺขูนํ ปจฺฉาภตฺตํ ปิณฺฑปาตปฏิกฺกนฺตานํ มณฺฑลมาเฬ สนฺนิสินฺนานํ สนฺนิปติตานํ อยมนฺตารากถา อุทปาทิ ‘สํโยชนนฺติ วา อาวุโส สํโยชนิยา ธมฺมาติ วา อิเม ธมฺมา นานตฺถา นานาพฺยญฺชนา, อุทาหุ เอกตฺถา พฺยญฺชนเมว นานนฺ’ติ.\csromanspacer{} ตเตฺรกจฺเจหิ เถเรหิ ภิกฺขูหิ เอวํ พฺยากตํ ‘สํโยชนนฺติ วา อาวุโส สํโยชนิยา ธมฺมาติ วา อิเม ธมฺมา นานตฺถา เจว นานาพฺยญฺชนา จา’ติ.\csromanspacer{} เอกจฺเจหิ เถเรหิ ภิกฺขูหิ เอวํ พฺยากตํ ‘สํโยชนนฺติ วา อาวุโส สํโยชนิยา ธมฺมาติ วา อิเม ธมฺมา เอกตฺถา พฺยญฺชนเมว นานนฺ’ติ”.\csromanspacer{} อถ โข จิตฺโต คหปติ เยน เถรา ภิกฺขู เตนุปสงฺกมิ, อุปสงฺกมิตฺวา เถเร ภิกฺขู อภิวาเทตฺวา เอกมนฺตํ นิสีทิ, เอกมนฺตํ นิสินฺโน โข จิตฺโต คหปติ เถเร ภิกฺขู เอตทโวจ “สุตํ เมตํ ภนฺเต สมฺพหุลานํ กิร เถรานํ ภิกฺขูนํ ปจฺฉาภตฺตํ ปิณฺฑปาตปฏิกฺกนฺตานํ มณฺฑลมาเฬ สนฺนิสินฺนานํ สนฺนิปติตานํ อยมนฺตรากถา อุทปาทิ ‘สํโยชนนฺติ วา อาวุโส สํโยชนิยา ธมฺมาติ วา อิเม ธมฺมา นานตฺถา นานาพฺยญฺชนา, อุทาหุ เอกตฺถา พฺยญฺชนเมว นานนฺ’ติ.\csromanspacer{} เอกจฺเจหิ เถเรหิ}

\csromanheader{2}{7. จิตฺตสํยุตฺต}
\prose{\csromanfolio{475}ภิกฺขูหิ เอวํ พฺยากตํ ‘สํโยชนนฺติ วา อาวุโส สํโยชนิยา ธมฺมาติ วา อิเม ธมฺมา นานตฺถา เจว นานาพฺยญฺชนา จา’ติ, เอกจฺเจหิ เถเรหิ ภิกฺขูหิ เอวํ พฺยากตํ ‘สํโยชนนฺติ วา อาวุโส สํโยชนิยา ธมฺมาติ วา อิเม ธมฺมา เอกตฺถา พฺยญฺชนเมว นานนฺ’ติ”.\csromanspacer{} เอวํ คหปตีติ.}

\prose{สํโยชนนฺติ วา ภนฺเต สํโยชนิยา ธมฺมาติ วา อิเม ธมฺมา นานตฺถา เจว นานาพฺยญฺชนา จ, เตน หิ ภนฺเต อุปมํ โว กริสฺสามิ, อุปมายปิเธกจฺเจ วิญฺญู ปุริสา ภาสิตสฺส อตฺถํ อาชานนฺติ.\csromanspacer{} เสยฺยถาปิ ภนฺเต กาโฬ จ พลีพทฺโท โอทาโต จ พลีพทฺโท เอเกน ทาเมน วา โยตฺเตน วา สํยุตฺตา อสฺสุ, โย นุ โข เอวํ วเทยฺย “กาโฬ พลีพทฺโท โอทาตสฺส พลีพทฺทสฺส สํโยชนํ, โอทาโต พลีพทฺโท กาฬสฺส พลีพทฺทสฺส สํโยชนนฺ”ติ, สมฺมา นุ โข โส วทมาโน วเทยฺยาติ.\csromanspacer{} โน เหตํ คหปติ, น โข คหปติ กาโฬ พลีพทฺโท โอทาตสฺส พลีพทฺทสฺส สํโยชนํ, นปิ โอทาโต พลีพทฺโท กาฬสฺส พลีพทฺทสฺส สํโยชนํ, เยน โข เต เอเกน ทาเมน วา โยตฺเตน วา สํยุตฺตา, ตํ ตตฺถ สํโยชนนฺติ.\csromanspacer{} เอวเมว โข ภนฺเต น จกฺขุ รูปานํ สํโยชนํ, น รูปา จกฺขุสฺส สํโยชนํ, ยญฺจ ตตฺถ ตทุภยํ ปฏิจฺจ อุปฺปชฺชติ ฉนฺทราโค, ตํ ตตฺถ สํโยชนํ.\csromanspacer{} น โสตํ สทฺทานํ.\csromanspacer{} น ฆานํ คนฺธานํ.\csromanspacer{} น ชิวฺหา รสานํ.\csromanspacer{} น กาโย โผฏฺฐพฺพานํ สํโยชนํ, น โผฏฺฐพฺพา กายสฺส สํโยชนํ, ยญฺจ ตตฺถ ตทุภยํ ปฏิจฺจ อุปฺปชฺชติ ฉนฺทราโค, ตํ ตตฺถ สํโยชนํ.\csromanspacer{} น มโน ธมฺมานํ สํโยชนํ, น ธมฺมา มนสฺส สํโยชนํ.\csromanspacer{} ยญฺจ ตตฺถ ตทุภยํ ปฏิจฺจ อุปฺปชฺชติ ฉนฺทราโค, ตํ ตตฺถ สํโยชนนฺติ.\csromanspacer{} ลาภา เต คหปติ, สุลทฺธํ เต คหปติ, ยสฺส เต คมฺภีเร พุทฺธวจเน ปญฺญาจกฺขุ กมตีติ.\csromanspacer{}\csromanspacer{}ปฐมํ.}

\csromanheader{1}{2. ปฐม-อิสิทตฺตสุตฺต}
\csromanmatikaanchor{mtk.262}
\csromantocmarkref{section}{2. ปฐม-อิสิทตฺตสุตฺต}{mtk.262}
\bookmark[dest={mtk.262},level=1]{2. ปฐม-อิสิทตฺตสุตฺต}
\proseitem{344}{เอกํ สมยํ สมฺพหุลา เถรา ภิกฺขู มจฺฉิกาสณฺเฑ วิหรนฺติ อมฺพาฏกวเน.\csromanspacer{} อถ โข จิตฺโต คหปติ เยน เถรา ภิกฺขู เตนุปสงฺกมิ, อุปสงฺกมิตฺวา เถเร ภิกฺขู อภิวาเทตฺวา เอกมนฺตํ นิสีทิ, เอกมนฺตํ นิสินฺโน โข จิตฺโต คหปติ เถเร ภิกฺขู เอตทโวจ “อธิวาเสนฺตุ เม ภนฺเต เถรา สฺวาตนาย ภตฺตนฺ”ติ.\csromanspacer{} อธิวาเสสุํ โข เถรา ภิกฺขู}

\gambhira{สฬายตนวคฺคสํยุตฺตปาฬิ}
\prose{\csromanfolio{476}ตุณฺหีภาเวน.\csromanspacer{} อถ โข จิตฺโต คหปติ เถรานํ ภิกฺขูนํ อธิวาสนํ วิทิตฺวา อุฏฺฐายาสนา เถเร ภิกฺขู อภิวาเทตฺวา ปทกฺขิณํ กตฺวา ปกฺกามิ.\csromanspacer{} อถ โข เถรา ภิกฺขู ตสฺสา รตฺติยา อจฺจเยน ปุพฺพณฺหสมยํ นิวาเสตฺวา ปตฺตจีวรมาทาย เยน จิตฺตสฺส คหปติสฺส นิเวสนํ เตนุปสงฺกมิํสุ, อุปสงฺกมิตฺวา ปญฺญตฺเต อาสเน นิสีทิํสุ.}

\prose{อถ โข จิตฺโต คหปติ เยน เถรา ภิกฺขู เตนุปสงฺกมิ, อุปสงฺกมิตฺวา เถเร ภิกฺขู อภิวาเทตฺวา เอกมนฺตํ นิสีทิ, เอกมนฺตํ นิสินฺโน โข จิตฺโต คหปติ อายสฺมนฺตํ เถรํ เอตทโวจ “ธาตุนานตฺตํ ธาตุนานตฺตนฺติ ภนฺเต เถร วุจฺจติ, กิตฺตาวตา นุ โข ภนฺเต ธาตุนานตฺตํ วุตฺตํ ภควตา”ติ.\csromanspacer{} เอวํ วุตฺเต อายสฺมา เถโร ตุณฺหี อโหสิ.\csromanspacer{} ทุติยมฺปิ โข จิตฺโต คหปติ อายสฺมนฺตํ เถรํ เอตทโวจ “ธาตุนานตฺตํ ธาตุนานตฺตนฺติ ภนฺเต เถร วุจฺจติ, กิตฺตาวตา นุ โข ภนฺเต ธาตุนานตฺตํ วุตฺตํ ภควตา”ติ.\csromanspacer{} ทุติยมฺปิ โข อายสฺมา เถโร ตุณฺหี อโหสิ.\csromanspacer{} ตติยมฺปิ โข จิตฺโต คหปติ อายสฺมนฺตํ เถรํ เอตทโวจ “ธาตุนานตฺตํ ธาตุนานตฺตนฺติ ภนฺเต เถร วุจฺจติ, กิตฺตาวตา นุ โข ภนฺเต ธาตุนานตฺตํ วุตฺตํ ภควตา”ติ.\csromanspacer{} ตติยมฺปิ โข อายสฺมา เถโร ตุณฺหี อโหสิ.}

\prose{เตน โข ปน สมเยน อายสฺมา อิสิทตฺโต ตสฺมิํ ภิกฺขุสํเฆ สพฺพนวโก โหติ.\csromanspacer{} อถ โข อายสฺมา อิสิทตฺโต อายสฺมนฺตํ เถรํ เอตทโวจ “พฺยากโรมหํ ภนฺเต เถร จิตฺตสฺส คหปติโน เอตํ ปญฺหนฺ”ติ.\csromanspacer{} พฺยากโรหิ ตฺวํ อาวุโส อิสิทตฺต จิตฺตสฺส คหปติโน เอตํ ปญฺหนฺติ.\csromanspacer{} เอวํ หิ ตฺวํ คหปติ ปุจฺฉสิ “ธาตุนานตฺตํ ธาตุนานนตฺตนฺติ ภนฺเต เถร วุจฺจติ, กิตฺตาวตา นุ โข ภนฺเต ธาตุนานตฺตํ วุตฺตํ ภควตา”ติ.\csromanspacer{} เอวํ ภนฺเต.\csromanspacer{} อิทํ โข คหปติ ธาตุนานตฺตํ วุตฺตํ ภควตา จกฺขุธาตุ รูปธาตุ จกฺขุวิญฺญาณธาตุ ฯเปฯ มโนธาตุ ธมฺมธาตุ มโนวิญฺญาณธาตุ.\csromanspacer{} เอตฺตาวตา โข คหปติ ธาตุนานตฺตํ วุตฺตํ ภควตาติ.}

\prose{อถ โข จิตฺโต คหปติ อายสฺมโต อิสิทตฺตสฺส ภาสิตํ อภินนฺทิตฺวา อนุโมทิตฺวา เถเร ภิกฺขู ปณีเตน ขาทนีเยน โภชนีเยน สหตฺถา สนฺตปฺเปสิ สมฺปวาเรสิ.\csromanspacer{} อถ โข เถรา ภิกฺขู ภุตฺตาวิโน โอนีตปตฺตปานีโน อุฏฺฐายาสนา ปกฺกมิํสุ.\csromanspacer{} อถ โข อายสฺมา เถโร อายสฺมนฺตํ อิสิทตฺตํ เอตทโวจ “สาธุ โข ตํ อาวุโส}

\csromanheader{2}{7. จิตฺตสํยุตฺต}
\prose{\csromanfolio{477}อิสิทตฺต เอโส ปโญฺห ปฏิภาสิ, เนโส ปโญฺห มํ ปฏิภาสิ.\csromanspacer{} เตนหาวุโส อิสิทตฺต ยทา อญฺญถาปิ\footnote{ยทา อญฺญทาปิ (สี, อิ) อญฺญทาปิ (?)} เอวรูโป ปโญฺห อาคจฺเฉยฺย, ตญฺเญเวตฺถ ปฏิภาเสยฺยา”ติ.\csromanspacer{}\csromanspacer{}ทุติยํ.}

\csromanheader{1}{3. ทุติย-อิสิทตฺตสุตฺต}
\csromanmatikaanchor{mtk.263}
\csromantocmarkref{section}{3. ทุติย-อิสิทตฺตสุตฺต}{mtk.263}
\bookmark[dest={mtk.263},level=1]{3. ทุติย-อิสิทตฺตสุตฺต}
\proseitem{345}{เอกํ สมยํ สมฺพหุลา เถรา ภิกฺขู มจฺฉิกาสณฺเฑ วิหรนฺติ อมฺพาฏกวเน.\csromanspacer{} อถ โข จิตฺโต คหปติ เยน เถรา ภิกฺขู เตนุปสงฺกมิ, อุปสงฺกมิตฺวา เถเร ภิกฺขู อภิวาเทตฺวา เอกมนฺตํ นิสีทิ, เอกมนฺตํ นิสินฺโน โข จิตฺโต คหปติ เถเร ภิกฺขู เอตทโวจ “อธิวาเสนฺตุ เม ภนฺเต เถรา สฺวาตนาย ภตฺตนฺ”ติ.\csromanspacer{} อธิวาเสสุํ โข เถรา ภิกฺขู ตุณฺหีภาเวน.\csromanspacer{} อถ โข จิตฺโต คหปติ เถรานํ ภิกฺขูนํ อธิวาสนํ วิทิตฺวา อุฏฺฐายาสนา เถเร ภิกฺขู อภิวาเทตฺวา ปทกฺขิณํ กตฺวา ปกฺกามิ.\csromanspacer{} อถ โข เถรา ภิกฺขู ตสฺสา รตฺติยา อจฺจเยน ปุพฺพณฺหสมยํ นิวาเสตฺวา ปตฺตจีวรมาทาย เยน จิตฺตสฺส คหปติสฺส นิเวสนํ เตนุปสงฺกมิํสุ, อุปสงฺกมิตฺวา ปญฺญตฺเต อาสเน นิสีทิํสุ. อถ โข จิตฺโต คหปติ เยน เถรา ภิกฺขู เตนุปสงฺกมิ, อุปสงฺกมิตฺวา เถเร ภิกฺขู อภิวาเทตฺวา เอกมนฺตํ นิสีทิ, เอกมนฺตํ นิสินฺโน โข จิตฺโต คหปติ อายสฺมนฺตํ เถรํ เอตทโวจ “ยา อิมา ภนฺเต เถร อเนกวิหิตา ทิฏฺฐิโย โลเก อุปฺปชฺชนฺติ ‘สสฺสโต โลโก’ติ วา, ‘อสสฺสโต โลโก’ติ วา, ‘อนฺตวา โลโก’ติ วา, ‘อนนฺตวา โลโก’ติ วา, ‘ตํ ชีวํ ตํ สรีรนฺ’ติ วา, ‘อญฺญํ ชีวํ อญฺญํ สรีรนฺ’ติ วา, ‘โหติ ตถาคโต ปรํ มรณา’ติ วา, ‘น โหติ ตถาคโต ปรํ มรณา’ติ วา, โหติ จ น จ โหติ ตถาคโต ปรํ มรณา’ติ วา, ‘เนว โหติ น น โหติ ตถาคโต ปรํ มรณา’ติ วา.\csromanspacer{} ยานิ จิมานิ ทฺวาสฏฺฐิ ทิฏฺฐิคตานิ พฺรหฺมชาเล ภณิตานิ, อิมา นุ โข ภนฺเต ทิฏฺฐิโย กิสฺมิํ สติ โหนฺติ, กิสฺมิํ อสติ น โหนฺตี”ติ.}

\prose{เอวํ วุตฺเต อายสฺมา เถโร ตุณฺหี อโหสิ.\csromanspacer{} ทุติยมฺปิ โข จิตฺโต คหปติ ฯเปฯ.\csromanspacer{} ตติยมฺปิ โข จิตฺโต คหปติ อายสฺมนฺตํ เถรํ เอตทโวจ “ยา อิมา ภนฺเต เถร อเนกวิหิตา ทิฏฺฐิโย}

\gambhira{สฬายตนวคฺคสํยุตฺตปาฬิ}
\prose{\csromanfolio{478}โลเก อุปฺปชฺชนฺติ ‘สสฺสโต โลโก’ติ วา, ‘อสสฺสโต โลโก’ติ วา, ‘อนฺตวา โลโก’ติ วา, ‘อนนฺตวา โลโก’ติ วา, ‘ตํ ชีวํ ตํ สรีรนฺ’ติ วา, ‘อญฺญํ ชีวํ อญฺญํ สรีรนฺ’ติ วา, ‘โหติ ตถาคโต ปรํ มรณา’ติ วา, ‘น โหติ ตถาคโต ปรํ มรณา’ติ วา, ‘โหติ จ น จ โหติ ตถาคโต ปรํ มรณา’ติ วา, ‘เนว โหติ น น โหติ ตถาคโต ปรํ มรณา’ติ วา.\csromanspacer{} ยานิ จิมานิ ทฺวาสฏฺฐิ ทิฏฺฐิคตานิ พฺรหฺมชาเล ภณิตานิ, อิมา นุ โข ภนฺเต ทิฏฺฐิโย กิสฺมิํ สติ โหนฺติ, กิสฺมิํ อสติ น โหนฺตี”ติ.\csromanspacer{} ตติยมฺปิ โข อายสฺมา เถโร ตุณฺหี อโหสิ.}

\prose{เตน โข ปน สมเยน อายสฺมา อิสิทตฺโต ตสฺมิํ ภิกฺขุสํเฆ สพฺพนวโก โหติ.\csromanspacer{} อถ โข อายสฺมา อิสิทตฺโต อายสฺมนฺตํ เถรํ เอตทโวจ “พฺยากโรมหํ ภนฺเต เถร จิตฺตสฺส คหปติโน เอตํ ปญฺหนฺ”ติ.\csromanspacer{} พฺยากโรหิ ตฺวํ อาวุโส อิสิทตฺต จิตฺตสฺส คหปติโน เอตํ ปญฺหนฺติ.\csromanspacer{} เอวํ หิ ตฺวํ คหปติ ปุจฺฉสิ “ยา อิมา ภนฺเต เถร อเนกวิหิตา ทิฏฺฐิโย โลเก อุปปชฺชนฺติ ‘สสฺสโต โลโก’ติ วา ฯเปฯ อิมา นุ โข ภนฺเต ทิฏฺฐิโย กิสฺมิํ สติ โหนฺติ, กิสฺมิํ อสติ น โหนฺตี’ติ”.\csromanspacer{} เอวํ ภนฺเต.\csromanspacer{} ยา อิมา คหปติ อเนกวิหิตา ทิฏฺฐิโย โลเก อุปฺปชฺชนฺติ “สสฺสโต โลโก”ติ วา, “อสสฺสโต โลโก”ติ วา, “อนฺตวา โลโก”ติ วา, “อนนฺตวา โลโก”ติ วา, “ตํ ชีวํ ตํ สรีรนฺ”ติ วา, “อญฺญํ ชีวํ อญฺญํ สรีรนฺ”ติ วา, “โหติ ตถาคโต ปรํ มรณา”ติ วา, “น โหติ ตถาคโต ปรํ มรณา”ติ วา, “โหติ จ น จ โหติ ตถาคโต ปรํ มรณา”ติ วา, “เนว โหติ น น โหติ ตถาคโต ปรํ มรณา”ติ วา.\csromanspacer{} ยานิ จิมานิ ทฺวาสฏฺฐิ ทิฏฺฐิคตานิ พฺรหฺมชาเล ภณิตานิ, อิมา โข คหปติ ทิฏฺฐิโย สกฺกายทิฏฺฐิยา สติ โหนฺติ, สกฺกายทิฏฺฐิยา อสติ น โหนฺตีติ.}

\prose{กถํ ปน ภนฺเต สกฺกายทิฏฺฐิ โหตีติ.\csromanspacer{} อิธ คหปติ อสฺสุตวา ปุถุชฺชโน อริยานํ อทสฺสาวี อริยธมฺมสฺส อโกวิโท อริยธมฺเม อวินีโต, สปฺปุริสานํ อทสฺสาวี สปฺปุริสธมฺมสฺส อโกวิโท สปฺปุริสธมฺเม อวินีโต รูปํ อตฺตโต สมนุปสฺสติ, รูปวนฺตํ วา อตฺตานํ, อตฺตนิ วา รูปํ, รูปสฺมิํ วา อตฺตานํ.\csromanspacer{} เวทนํ อตฺตโต สมนุปสฺสติ ฯเปฯ สญฺญํ.\csromanspacer{} สงฺขาเร.\csromanspacer{} วิญฺญาณํ อตฺตโต สมนุปสฺสติ, วิญฺญาณวนฺตํ วา อตฺตานํ,}

\csromanheader{2}{7. จิตฺตสํยุตฺต}
\prose{\csromanfolio{479}อตฺตนิ วา วิญฺญาณํ, วิญฺญาณสฺมิํ วา อตฺตานํ.\csromanspacer{} เอวํ โข คหปติ สกฺกายทิฏฺฐิ โหตีติ.}

\prose{กถํ ปน ภนฺเต สกฺกายทิฏฺฐิ น โหตีติ.\csromanspacer{} อิธ คหปติ สุตวา อริยสาวโก อริยานํ ทสฺสาวี อริยธมฺมสฺส โกวิโท อริยธมฺเม สุวินีโต, สปฺปุริสานํ ทสฺสาวี สปฺปุริสธมฺมสฺส โกวิโท สปฺปุริสธมฺเม สุวินีโต น รูปํ อตฺตโต สมนุปสฺสติ, น รูปวนฺตํ วา อตฺตานํ, น อตฺตนิ วา รูปํ, น รูปสฺมิํ วา อตฺตานํ.\csromanspacer{} น เวทนํ.\csromanspacer{} น สญฺญํ.\csromanspacer{} น สงฺขาเร.\csromanspacer{} น วิญฺญาณํ อตฺตโต สมนุปสฺสติ, น วิญฺญาณวนฺตํ วา อตฺตานํ, น อตฺตนิ วา วิญฺญาณํ, น วิญฺญาณสฺมิํ วา อตฺตานํ.\csromanspacer{} เอวํ โข คหปติ สกฺกายทิฏฺฐิ น โหตีติ.}

\prose{กุโต ภนฺเต อยฺโย อิสิทตฺโต อาคจฺฉตีติ.\csromanspacer{} อวนฺติยา โข คหปติ อาคจฺฉามีติ.\csromanspacer{} อตฺถิ ภนฺเต อวนฺติยา อิสิทตฺโต นาม กุลปุตฺโต อมฺหากํ อทิฏฺฐสหาโย ปพฺพชิโต, ทิฏฺโฐ โส อายสฺมตาติ.\csromanspacer{} เอวํ คหปตีติ.\csromanspacer{} กหํ นุ โข โส ภนฺเต อายสฺมา เอตรหิ วิหรตีติ.\csromanspacer{} เอวํ วุตฺเต อายสฺมา อิสิทตฺโต ตุณฺหี อโหสิ.\csromanspacer{} อยฺโย โน ภนฺเต อิสิทตฺโตติ.\csromanspacer{} เอวํ คหปตีติ.\csromanspacer{} อภิรมตุ ภนฺเต อยฺโย อิสิทตฺโต มจฺฉิกาสณฺเฑ รมณียํ อมฺพาฏกวนํ, อหํ อยฺยสฺส อิสิทตฺตสฺส อุสฺสุกฺกํ กริสฺสามิ จีวรปิณฺฑปาตเสนาสนคิลานปฺปจฺจยเภสชฺชปริกฺขารานนฺติ.\csromanspacer{} กลฺยาณํ วุจฺจติ คหปตีติ.}

\prose{อถ โข จิตฺโต คหปติ อายสฺมโต อิสิทตฺตสฺส ภาสิตํ อภินนฺทิตฺวา อนุโมทิตฺวา เถเร ภิกฺขู ปณีเตน ขาทนีเยน โภชนีเยน สหตฺถา สนฺตปฺเปสิ สมฺปวาเรสิ.\csromanspacer{} อถ โข เถรา ภิกฺขู ภุตฺตาวิโน โอนีตปตฺตปาณีโน อุฏฺฐายาสนา ปกฺกมิํสุ.\csromanspacer{} อถ โข อายสฺมา เถโร อายสฺมนฺตํ อิสีทตฺตํ เอตทโวจ “สาธุ โข ตํ อาวุโส อิสิทตฺต เอโส ปโญฺห ปฏิภาสิ, เนโส ปโญฺห มํ ปฏิภาสิ.\csromanspacer{} เตนหาวุโส อิสิทตฺต ยทา อญฺญถาปิ เอวรูโป ปโญฺห อาคจฺเฉยฺย, ตญฺเญวตฺถ ปฏิภาเสยฺยา”ติ.\csromanspacer{} อถ โข อายสฺมา อิสิทตฺโต เสนาสนํ สํสาเมตฺวา ปตฺตจีวรมาทาย มจฺฉิกาสณฺฑมฺหา ปกฺกามิ.\csromanspacer{} ยํ มจฺฉิกาสณฺฑมฺหา ปกฺกามิ, ตถา ปกฺกนฺโตว อโหสิ.\csromanspacer{} น ปุน ปจฺฉาคจฺฉีติ.\csromanspacer{}\csromanspacer{}ตติยํ.}

\gambhira{สฬายตนวคฺคสํยุตฺตปาฬิ}
\csromanheader{1}{4. มหกปาฏิหาริยสุตฺต}
\csromanmatikaanchor{mtk.264}
\csromantocmarkref{section}{4. มหกปาฏิหาริยสุตฺต}{mtk.264}
\bookmark[dest={mtk.264},level=1]{4. มหกปาฏิหาริยสุตฺต}
\proseitem{346}{\csromanfolio{480}เอกํ สมยํ สมฺพุหุลา เถรา ภิกฺขู มจฺฉิกาสณฺเฑ วิหรนฺติ อมฺพาฏกวเน.\csromanspacer{} อถ โข จิตฺโต คหปติ เยน เถรา ภิกฺขู เตนุปสงฺกมิ, อุปสงฺกมิตฺวา เถเร ภิกฺขู อภิวาเทตฺวา เอกมนฺตํ นิสีทิ, เอกมนฺตํ นิสินฺโน โข จิตฺโต คหปติ เถเร ภิกฺขู เอตทโวจ “อธิวาเสนฺตุ เม ภนฺเต เถรา ภิกฺขู สฺวาตนาย โคกุเล ภตฺตนฺ”ติ.\csromanspacer{} อธิวาเสสุํ โข เถรา ภิกฺขู ตุณฺหีภาเวน.\csromanspacer{} อถ โข จิตฺโต คหปติ เถรานํ ภิกฺขูนํ อธิวาสนํ วิทิตฺวา อุฏฺฐายาสนา เถเร ภิกฺขู อภิวาเทตฺวา ปทกฺขิณํ กตฺวา ปกฺกามิ.\csromanspacer{} อถ โข เถรา ภิกฺขู ตสฺสา รตฺติยา อจฺจเยน ปุพฺพณฺหสมยํ นิวาเสตฺวา ปตฺตจีวรมาทาย เยน จิตฺตสฺส คหปติโน โคกุลํ เตนุปสงฺกมิํสุ, อุปสงฺกมิตฺวา ปญฺญตฺเต อาสเน นิสีทิํสุ. อถ โข จิตฺโต คหปติ เถเร ภิกฺขู ปณีเตน สปฺปิปายาเสน สหตฺถา สนฺตปฺเปสิ สมฺปวาเรสิ.\csromanspacer{} อถ โข เถรา ภิกฺขู ภุตฺตาวิโน โอนีตปตฺตปาณิโน อุฏฺฐายาสนา ปกฺกมิํสุ.\csromanspacer{} จิตฺโตปิ โข คหปติ “เสสกํ วิสฺสชฺเชถา”ติ วตฺวา เถเร ภิกฺขู ปิฏฺฐิโต ปิฏฺฐิโต อนุพนฺธิ.\csromanspacer{} เตน โข ปน สมเยน อุณฺหํ โหติ กุถิตํ\footnote{กุฏฺฐิตํ (สี, สฺยา, กํ, อิ)}, เต จ เถรา ภิกฺขู ปเวลิยมาเนน มญฺเญ กาเยน คจฺฉนฺติ, ยถา ตํ โภชนํ ภุตฺตาวิโน.}

\prose{เตน โข ปน สมเยน อายสฺมา มหโก ตสฺมิํ ภิกฺขุสํเฆ สพฺพนวโก โหติ.\csromanspacer{} อถ โข อายสฺมา มหโก อายสฺมนฺตํ เถรํ เอตทโวจ “สาธุ ขฺวสฺส ภนฺเต เถร สีตโก จ วาโต วาเยยฺย, อพฺภสมฺปิลาโป\footnote{อพฺภสํพิลาโป (สี), อพฺภสํวิลาโป (อิ)} จ อสฺส, เทโว จ เอกเมกํ ผุสาเยยฺยา”ติ.\csromanspacer{} สาธุ ขฺวสฺส อาวุโส มหก ยํ สีตโก จ วาโต วาเยยฺย, อพฺภสมฺปิลาโป จ อสฺส, เทโว จ เอกเมกํ ผุสาเยยฺยาติ.\csromanspacer{} อถ โข อายสฺมา มหโก ตถารูปํ อิทฺธาภิสงฺขารํ อภิสงฺขริ\footnote{อภิสงฺขาสิ (สี)}, ยถา สีตโก จ วาโต วายิ, อพฺภสมฺปิลาโป จ อสฺส\footnote{อาสิ (?)}, เทโว จ เอกเมกํ ผุสิ.\csromanspacer{} อถ โข จิตฺตสฺส คหปติโน เอตทโหสิ “โย โข อิมสฺมิํ ภิกฺขุสํเฆ สพฺพนวโก ภิกฺขุ, ตสฺสายํ เอวรูโป}

\csromanheader{2}{7. จิตฺตสํยุตฺต}
\prose{\csromanfolio{481}อิทฺธานุภาโว”ติ.\csromanspacer{} อถ โข อายสฺมา มหโก อารามํ สมฺปาปุณิตฺวา อายสฺมนฺตํ เถรํ เอตทโวจ “อลเมตฺตาวตา ภนฺเต เถรา”ติ.\csromanspacer{} อลเมตฺตาวตา อาวุโส มหก, กตเมตฺตาวตา อาวุโส มหก, ปูชิตเมตฺตาวตา อาวุโส มหกาติ.\csromanspacer{} อถ โข เถรา ภิกฺขู ยถาวิหารํ อคมํสุ, อายสฺมาปิ มหโก สกํ วิหารํ อคมาสิ.}

\prose{อถ โข จิตฺโต คหปติ เยนายสฺมา มหโก เตนุปสงฺกมิ, อุปสงฺกมิตฺวา อายสฺมนฺตํ มหกํ อภิวาเทตฺวา เอกมนฺตํ นิสีทิ, เอกมนฺตํ นิสินฺโน โข จิตฺโต คหปติ อายสฺมนฺตํ มหกํ เอตทโวจ “สาธุ เม ภนฺเต อยฺโย มหโก อุตฺตริ มนุสฺสธมฺมํ อิทฺธิปาฏิหาริยํ ทสฺเสตู”ติ.\csromanspacer{} เตน หิ ตฺวํ คหปติ อาลินฺเท อุตฺตราสงฺคํ ปญฺญเปตฺวา ติณกลาปํ โอกาเสหีติ. “เอวํ ภนฺเต”ติ โข จิตฺโต คหปติ อายสฺมโต มหกสฺส ปฏิสฺสุตฺวา อาลินฺเท อุตฺตราสงฺคํ ปญฺญเปตฺวา ติณกลาปํ โอกาเสสิ.\csromanspacer{} อถ โข อายสฺมา มหโก วิหารํ ปวิสิตฺวา สูจิฆฏิกํ ทตฺวา ตถารูปํ อิทฺธาภิสงฺขารํ อภิสงฺขริ, ยถา ตาลจฺฉิคฺคเฬน จ อคฺคฬนฺตริกาย จ อจฺจิ นิกฺขมิตฺวา ติณานิ ฌาเปสิ, อุตฺตราสงฺคํ น ฌาเปสิ.\csromanspacer{} อถ โข จิตฺโต คหปติ อุตฺตราสงฺคํ ปปฺโผเฏตฺวา สํวิคฺโค โลมหฏฺฐชาโต เอกมนฺตํ อฏฺฐาสิ.\csromanspacer{} อถ โข อายสฺมา มหโก วิหารา นิกฺขมิตฺวา จิตฺตํ คหปติํ เอตทโวจ “อลเมตฺตาวตา คหปตี”ติ.\csromanspacer{} อลเมตฺตาวตา ภนฺเต มหก, กตเมตฺตาวตา ภนฺเต มหก, ปูชิตเมตฺตาวตา ภนฺเต มหก, อภิรมตุ ภนฺเต อยฺโย มหโก มจฺฉิกาสณฺเฑ รมณียํ อมฺพาฏกวนํ, อหํ อยฺยสฺส มหกสฺส อุสฺสุกฺกํ กริสฺสามิ จีวรปิณฺฑปาตเสนาสนคิลานปฺปจฺจยเภสชฺชปริกฺขารานนฺติ.\csromanspacer{} กลฺยาณํ วุจฺจติ คหปตีติ.\csromanspacer{} อถ โข อายสฺมา มหโก เสนาสนํ สํสาเมตฺวา ปตฺตจีวรมาทาย มจฺฉิกาสณฺฑมฺหา ปกฺกามิ.\csromanspacer{} ยํ มจฺฉิกาสณฺฑมฺหา ปกฺกามิ, ตถา ปกฺกนฺโตว อโหสิ.\csromanspacer{} น ปุน ปจฺฉาคจฺฉีติ.\csromanspacer{}\csromanspacer{}จตุตฺถํ.}

\csromanheader{1}{5. ปฐมกามภูสุตฺต}
\csromanmatikaanchor{mtk.265}
\csromantocmarkref{section}{5. ปฐมกามภูสุตฺต}{mtk.265}
\bookmark[dest={mtk.265},level=1]{5. ปฐมกามภูสุตฺต}
\proseitem{347}{เอกํ สมยํ อายสฺมา กามภู มจฺฉิกาสณฺเฑ วิหรติ อมฺพาฏกวเน.\csromanspacer{} อถ โข จิตฺโต คหปติ เยนายสฺมา กามภู เตนุปสงฺกมิ, อุปสงฺกมิตฺวา อายสฺมนฺตํ กามภุํ อภิวาเทตฺวา เอกมนฺตํ}

\gambhira{สฬายตนวคฺคสํยุตฺตปาฬิ}
\prose{\csromanfolio{482}นิสีทิ, เอกมนฺตํ นิสินฺนํ โข จิตฺตํ คหปติํ อายสฺมา กามภู เอตทโวจ “วุตฺตมิทํ คหปติ\csromandash{}}

\csromangathameasure{‘เนลงฺโค เสตปจฺฉาโท, เอกาโร วตฺตตี รโถ. \\ อนีฆํ ปสฺส อายนฺตํ, ฉินฺนโสตํ อพนฺธนนฺ’ติ.}
\csromangathagroup{\csromangathastanza{‘เนลงฺโค เสตปจฺฉาโท, เอกาโร วตฺตตี รโถ. \\ อนีฆํ ปสฺส อายนฺตํ\footnote{อปฺปตฺตํ (สฺยา, กํ, ก)}, ฉินฺนโสตํ อพนฺธนนฺ’ติ.}}

\prose{อิมสฺส นุ โข คหปติ สํขิตฺเตน ภาสิตสฺส กถํ วิตฺถาเรน อตฺโถ ทฏฺฐพฺโพ”ติ.\csromanspacer{} กิํ นุ โข เอตํ ภนฺเต ภควตา ภาสิตนฺติ.\csromanspacer{} เอวํ คหปตีติ.\csromanspacer{} เตน หิ ภนฺเต มุหุตฺตํ อาคเมหิ ยาวสฺส อตฺถํ เปกฺขามีติ.\csromanspacer{} อถ โข จิตฺโต คหปติ มุหุตฺตํ ตุณฺหี หุตฺวา อายสฺมนฺตํ กามภุํ เอตทโวจ\csromandash{}}

\prose{“เนลงฺคนฺ”ติ โข ภนฺเต สีลานเมตํ อธิวจนํ. “เสตปจฺฉาโท”ติ โข ภนฺเต วิมุตฺติยา เอตํ อธิวจนํ. “เอกาโร”ติ โข ภนฺเต สติยา เอตํ อธิวจนํ. “วตฺตตี”ติ โข ภนฺเต อภิกฺกมปฏิกฺกมสฺเสตํ อธิวจนํ. “รโถ”ติ โข ภนฺเต อิมสฺเสตํ จาตุมหาภูติกสฺส กายสฺส อธิวจนํ มาตาเปตฺติกสมฺภวสฺส โอทนกุมฺมาสูปจยสฺส อนิจฺจุจฺฉาทนปริมทฺทนเภทนวิทฺธํสนธมฺมสฺส.\csromanspacer{} ราโค โข ภนฺเต นีโฆ, โทโส นีโฆ, โมโห นีโฆ, เต ขีณาสวสฺส ภิกฺขุโน ปหีนา อุจฺฉินฺนมูลา ตาลาวตฺถุกตา อนภาวํกตา อายติํ อนุปฺปาทธมฺมา, ตสฺมา ขีณาสโว ภิกฺขุ “อนีโฆ”ติ วุจฺจติ. “อายนฺตนฺ”ติ โข ภนฺเต อรหโต เอตํ อธิวจนํ. “โสโต”ติ โข ภนฺเต ตณฺหาเยตํ อธิวจนํ.\csromanspacer{} สา ขีณาสวสฺส ภิกฺขุโน ปหีนา อุจฺฉินฺนมูลา ตาลาวตฺถุกตา อนภาวํกตา อายติํ อนุปฺปาทธมฺมา, ตสฺมา ขีณาสโว ภิกฺขุ “ฉินฺนโสโต”ติ วุจฺจติ.\csromanspacer{} ราโค โข ภนฺเต พนฺธนํ, โทโส พนฺธนํ, โมโห พนฺธนํ, เต ขีณาสวสฺส ภิกฺขุโน ปหีนา อุจฺฉินฺนมูลา ตาลาวตฺถุกตา อนภาวํกตา อายติํ อนุปฺปาทธมฺมา, ตสฺมา ขีณาสโว ภิกฺขุ “อพนฺธโน”ติ วุจฺจติ.\csromanspacer{} อิติ โข ภนฺเต ยํ ตํ ภควตา วุตฺตํ\csromandash{}}

\csromangathameasure{“เนลงฺโค เสตปจฺฉาโท, เอกาโร วตฺตตี รโถ. \\ อนีฆํ ปสฺส อายนฺตํ, ฉินฺนโสตํ อพนฺธนนฺ”ติ.}
\csromangathagroup{\csromangathastanza{“เนลงฺโค เสตปจฺฉาโท, เอกาโร วตฺตตี รโถ. \\ อนีฆํ ปสฺส อายนฺตํ, ฉินฺนโสตํ อพนฺธนนฺ”ติ.}}

\csromanheader{2}{7. จิตฺตสํยุตฺต}
\prose{\csromanfolio{483}อิมสฺส โข ภนฺเต ภควตา สํขิตฺเตน ภาสิตสฺส เอวํ วิตฺตาเรน อตฺถํ อาชานามีติ.\csromanspacer{} ลาภา เต คหปติ, สุลทฺธํ เต คหปติ, ยสฺส เต คมฺภีเร พุทฺธวจเน ปญฺญาจกฺขุ กมตีติ.\csromanspacer{}\csromanspacer{}ปญฺจมํ.}

\csromanheader{1}{6. ทุติยกามภูสุตฺต}
\csromanmatikaanchor{mtk.266}
\csromantocmarkref{section}{6. ทุติยกามภูสุตฺต}{mtk.266}
\bookmark[dest={mtk.266},level=1]{6. ทุติยกามภูสุตฺต}
\proseitem{348}{เอกํ สมยํ อายสฺมา กามภู มจฺฉิกาสณฺเฑ วิหรติ อมฺพาฏกวเน.\csromanspacer{} อถ โข จิตฺโต คหปติ เยนายสฺมา กามภู เตนุปสงฺกมิ, อุปสงฺกมิตฺวา เอกมนฺตํ นิสีทิ, เอกมนฺตํ นิสินฺโน โข จิตฺโต คหปติ อายสฺมนฺตํ กามภุํ เอตทโวจ “กติ นุ โข ภนฺเต สงฺขารา”ติ.\csromanspacer{} ตโย โข คหปติ สงฺขารา, กายสงฺขโร วจีสงฺขาโร จิตฺตสงฺขาโรติ. “สาธุ ภนฺเต”ติ โข จิตฺโต คหปติ อายสฺมโต กามภุสฺส ภาสิตํ อภินนฺทิตฺวา อนุโมทิตฺวา อายสฺมนฺตํ กามภุํ อุตฺตริํ ปญฺหํ อปุจฺฉิ “กตโม ปน ภนฺเต กายสงฺขาโร, กตโม วจีสงฺขาโร, กตโม จิตฺตสงฺขาโร”ติ.\csromanspacer{} อสฺสาสปสฺสาสา โข คหปติ กายสงฺขาโร, วิตกฺกวิจารา วจีสงฺขาโร, สญฺญา จ เวทนา จ จิตฺตสงฺขาโรติ.}

\prose{สาธุ ภนฺเตติ โข จิตฺโต คหปติ ฯเปฯ อุตฺตริํ ปญฺหํ อปุจฺฉิ “กสฺมา ปน ภนฺเต อสฺสาสปสฺสาสา กายสงฺขาโร, กสฺมา วิตกฺกวิจารา วจีสงฺขาโร, กสฺมา สญฺญา จ เวทนา จ จิตฺตสงฺขาโร”ติ.\csromanspacer{} อสฺสาสปสฺสาสา โข คหปติ กายิกา เอเต ธมฺมา กายปฺปฏิพทฺธา, ตสฺมา อสฺสาสปสฺสาสา กายสงฺขาโร.\csromanspacer{} ปุพฺเพ โข คหปติ วิตกฺเกตฺวา วิจาเรตฺวา ปจฺฉา วาจํ ภินฺทติ, ตสฺมา วิตกฺกวิจารา วจีสงฺขาโร.\csromanspacer{} สญฺญา จ เวทนา จ เจตสิกา เอเต ธมฺมา จิตฺตปฺปฏิพทฺธา, ตสฺมา สญฺญา จ เวทนา จ จิตฺตสงฺขาโร.}

\prose{สาธุ ฯเปฯ อุตฺตริํ ปญฺหํ อปุจฺฉิ “กถํ ปน ภนฺเต สญฺญาเวทยิตนิโรธสมาปตฺติ โหตี”ติ.\csromanspacer{} น โข คหปติ สญฺญาเวทยิตนิโรธํ สมาปชฺชนฺตสฺส ภิกฺขุโน เอวํ โหติ “อหํ สญฺญาเวทยิตนิโรธํ สมาปชฺชิสฺสนฺ”ติ วา, “อหํ สญฺญาเวทยิตนิโรธํ สมาปชฺชามี”ติ วา, “อหํ สญฺญาเวทยิตนิโรธํ สมาปนฺโน”ติ วา, อถ ขฺวสฺส ปุพฺเพว ตถา จิตฺตํ ภาวิตํ โหติ, ยํ ตํ ตถตฺตาย อุปเนตีติ.}

\gambhira{สฬายตนวคฺคสํยุตฺตปาฬิ}
\prose{\csromanfolio{484}สาธุ ฯเปฯ อุตฺตริํ ปญฺหํ อปุจฺฉิ “สญฺญาเวทยิตนิโรธํ สมาปชฺชนฺตสฺส ปน ภนฺเต ภิกฺขุโน กตเม ธมฺมา ปฐมํ นิรุชฺฌนฺติ, ยทิ วา กายสงฺขาโร ยทิ วา วจีสงฺขาโร ยทิ วา จิตฺตสงฺขาโร”ติ.\csromanspacer{} สญฺญาเวทยิตนิโรธํ สมาปชฺชนฺตสฺส โข คหปติ ภิกฺขุโน วจีสงฺขาโร ปฐมํ นิรุชฺฌติ, ตโต กายสงฺขาโร, ตโต จิตฺตสงฺขาโรติ.}

\prose{สาธุ ฯเปฯ อุตฺตริํ ปญฺหํ อปุจฺฉิ “ยฺวายํ ภนฺเต มโต กาลงฺกโต, โย จายํ ภิกฺขุ สญฺญาเวทยิตนิโรธํ สมาปนฺโน, อิเมสํ กิํ นานากรณนฺ”ติ.\csromanspacer{} ยฺวายํ คหปติ มโต กาลงฺกโต, ตสฺส กายสงฺขาโร นิรุทฺโธ ปฏิปฺปสฺสทฺโธ, วจีสงฺขาโร นิรุทฺโธ ปฏิปฺปสฺสทฺโธ, จิตฺตสงฺขาโร นิรุทฺโธ ปฏิปฺปสฺสทฺโธ, อายุ ปริกฺขีโณ, อุสฺมา วูปสนฺตา, อินฺทฺริยานิ วิปริภินฺนานิ.\csromanspacer{} โย จ ขฺวายํ คหปติ ภิกฺขุ สญฺญาเวทยิตนิโรธํ สมาปนฺโน, ตสฺสปิ กายสงฺขาโร นิรุทฺโธ ปฏิปฺปสฺสทฺโธ, วจีสงฺขาโร นิรุทฺโธ ปฏิปฺปสฺสทฺโธ, จิตฺตสงฺขาโร นิรุทฺโธ ปฏิสฺสทฺโธ, อายุ อปริกฺขีโณ, อุสฺมา อวูปสนฺตา, อินฺทฺริยานิ วิปฺปสนฺนานิ.\csromanspacer{} ยฺวายํ คหปติ มโต กาลงฺกโต, โย จายํ ภิกฺขุ สญฺญาเวทยิตนิโรธํ สมาปนฺโน, อิทํ เนสํ นานากรณนฺติ.}

\prose{สาธุ ฯเปฯ อุตฺตริํ ปญฺหํ อปุจฺฉิ “กถํ ปน ภนฺเต สญฺญาเวทยิตนิโรธสมาปตฺติยา วุฏฺฐานํ โหตี”ติ.\csromanspacer{} น โข คหปติ สญฺญาเวทยิตนิโรธสมาปตฺติยา วุฏฺฐหนฺตสฺส ภิกฺขุโน เอวํ โหติ “อหํ สญฺญาเวทยิตนิโรธสมาปตฺติยา วุฏฺฐหิสฺสนฺ”ติ วา, “อหํ สญฺญาเวทยิตนิโรธสมาปตฺติยา วุฏฺฐหามี”ติ วา, “อหํ สญฺญาเวทยิตนิโรธสมาปตฺติยา วุฏฺฐิโต”ติ วา, อถ ขฺวสฺส ปุพฺเพว ตถา จิตฺตํ ภาวิตํ โหติ, ยํ ตํ ตถตฺตาย อุปเนตีติ.}

\prose{สาธุ ภนฺเต ฯเปฯ อุตฺตริํ ปญฺหํ อปุจฺฉิ “สญฺญาเวทยิตนิโรธสมาปตฺติยา วุฏฺฐหนฺตสฺส ปน ภนฺเต ภิกฺขุโน กตเม ธมฺมา ปฐมํ อุปฺปชฺชนฺติ, ยทิ วา กายสงฺขาโร ยทิ วา วจีสงฺขาโร ยทิ วา จิตฺตสงฺขาโร”ติ.\csromanspacer{} สญฺญาเวทยิตนิโรธสมาปตฺติยา วุฏฺฐหนฺตสฺส คหปติ ภิกฺขุโน จิตฺตสงฺขาโร ปฐมํ อุปฺปชฺชติ, ตโต กายสงฺขาโร, ตโต วจีสงฺขาโรติ.}

\prose{สาธุ ฯเปฯ อุตฺตริํ ปญฺหํ อปุจฺฉิ “สญฺญาเวทยิตนิโรธสมาปตฺติยา วุฏฺฐิตํ ปน ภนฺเต ภิกฺขุํ กติ ผสฺสา ผุสนฺติ”.\csromanspacer{} สญฺญาเวทยิตนิโรธสมาปตฺติยา}

\csromanheader{2}{7. จิตฺตสํยุตฺต}
\prose{\csromanfolio{485}วุฏฺฐิตํ โข คหปติ ภิกฺขุํ ตโย ผสฺสา ผุสนฺติ สุญฺญโต ผสฺโส อนิมิตฺโต ผสฺโส อปฺปณิหิโต ผสฺโสติ.}

\prose{สาธุ ฯเปฯ อุตฺตริํ ปญฺหํ อปุจฺฉิ “สญฺญาเวทยิตนิโรธสมาปตฺติยา วุฏฺฐิตสฺส ปน ภนฺเต ภิกฺขุโน กิํนินฺนํ จิตฺตํ โหติ กิํโปณํ กิํปพฺภารนฺ”ติ.\csromanspacer{} สญฺญาเวทยิตนิโรธสมาปตฺติยา วุฏฺฐิตสฺส โข คหปติ ภิกฺขุโน วิเวกนินฺนํ จิตฺตํ โหติ วิเวกโปณํ วิเวกปพฺภารนฺติ.}

\prose{สาธุ ภนฺเตติ โข จิตฺโต คหปติ อายสฺมโต กามภุสฺส ภาสิตํ อภินนฺทิตฺวา อนุโมทิตฺวา อายสฺมนฺตํ กามภุํ อุตฺตริํ ปญฺหํ อปุจฺฉิ “สญฺญาเวทยิตนิโรธสมาปตฺติยา ปน ภนฺเต กติ ธมฺมา พหูปการา”ติ.\csromanspacer{} อทฺธา โข ตฺวํ คหปติ ยํ ปฐมํ ปุจฺฉิตพฺพํ ตํ ปุจฺฉสิ, อปิ จ ตฺยาหํ พฺยากริสฺสามิ “สญฺญาเวทยิตนิโรธสมาปตฺติยา โข คหปติ เทฺว ธมฺมา พหูปการา สมโถ จ วิปสฺสนา จา”ติ.\csromanspacer{}\csromanspacer{}ฉฏฺฐํ.}

\csromanheader{1}{7. โคทตฺตสุตฺต}
\csromanmatikaanchor{mtk.267}
\csromantocmarkref{section}{7. โคทตฺตสุตฺต}{mtk.267}
\bookmark[dest={mtk.267},level=1]{7. โคทตฺตสุตฺต}
\proseitem{349}{เอกํ สมยํ อายสฺมา โคทตฺโต มจฺฉิกาสณฺเฑ วิหรติ อมฺพาฏกวเน.\csromanspacer{} อถ โข จิตฺโต คหปติ เยนายสฺมา โคทตฺโต เตนุปสงฺกมิ, อุปสงฺกมิตฺวา อายสฺมนฺตํ โคทตฺตํ อภิวาเทตฺวา เอกมนฺตํ นิสีทิ, เอกมนฺตํ นิสินฺนํ โข จิตฺตํ คหปติํ อายสฺมา โคทตฺโต เอตทโวจ “ยา จายํ คหปติ อปฺปมาณา เจโตวิมุตฺติ, ยา จ อากิญฺจญฺญา เจโตวิมุตฺติ, ยา จ สุญฺญตา เจโตวิมุตฺติ, ยา จ อนิมิตฺตา เจโตวิมุตฺติ.\csromanspacer{} อิเม ธมฺมา นานตฺถา นานาพฺยญฺชนา, อุทาหุ เอกตฺถา พฺยญฺชนเมว นานนฺ”ติ.\csromanspacer{} อตฺถิ ภนฺเต ปริยาโย, ยํ ปริยายํ อาคมฺม อิเม ธมฺมา นานตฺถา เจว นานาพฺยญฺชนา จ.\csromanspacer{} อตฺถิ ปน ภนฺเต ปริยาโย, ยํ ปริยายํ อาคมฺม อิเม ธมฺมา เอกตฺถา พฺยญฺชนเมว นานนฺติ.}

\prose{กตโม จ ภนฺเต ปริยาโย, ยํ ปริยายํ อาคมฺม อิเม ธมฺมา นานตฺถา เจว นานาพฺยญฺชนา จ.\csromanspacer{} อิธ ภนฺเต ภิกฺขุ เมตฺตาสหคเตน เจตสา เอกํ ทิสํ ผริตฺวา วิหรติ.\csromanspacer{} ตถา ทุติยํ.\csromanspacer{} ตถา ตติยํ.\csromanspacer{} ตถา จตุตฺถํ\footnote{จตุตฺถิํ (?)}.\csromanspacer{} อิติ อุทฺธมโธ ติริยํ สพฺพธิ สพฺพตฺตตายา สพฺพาวนฺตํ}

\gambhira{สฬายตนวคฺคสํยุตฺตปาฬิ}
\prose{\csromanfolio{486}โลกํ เมตฺตาสหคเตน เจตสา วิปุเลน มหคฺคเตน อปฺปมาเณน อเวเรน อพฺยาปชฺเชน\footnote{อพฺยาปชฺเฌน (สี, สฺยา, กํ, อิ), อพฺยาพชฺเฌน (?)} ผริตฺวา วิหรติ.\csromanspacer{} กรุณาสหคเตน เจตสา ฯเปฯ.\csromanspacer{} มุทิตาสหคเตน เจตสา ฯเปฯ.\csromanspacer{} อุเปกฺขาสหคเตน เจตสา เอกํ ทิสํ ผริตฺวา วิหรติ.\csromanspacer{} ตถา ทุติยํ.\csromanspacer{} ตถา ตติยํ.\csromanspacer{} ตถา จตุตฺถํ.\csromanspacer{} อิติ อุทฺธมโธ ติริยํ สพฺพธิ สพฺพตฺตตาย สพฺพาวนฺตํ โลกํ อุเปกฺขาสหคเตน เจตสา วิปุเลน มหคฺคเตน อปฺปมาเณน อเวเรน อพฺยาปชฺเชน ผริตฺวา วิหรติ.\csromanspacer{} อยํ วุจฺจติ ภนฺเต อปฺปมาณา เจโตวิมุตฺติ.}

\prose{กตมา จ ภนฺเต อากิญฺจญฺญา เจโตวิมุตฺติ.\csromanspacer{} อิธ ภนฺเต ภิกฺขุ สพฺพโส วิญฺญาณญฺจายตนํ สมติกฺกมฺม “นตฺถิ กิญฺจี”ติ อากิญฺจญฺญายตนํ อุปสมฺปชฺช วิหรติ.\csromanspacer{} อยํ วุจฺจติ ภนฺเต อากิญฺจญฺญา เจโตวิมุตฺติ.}

\prose{กตมา จ ภนฺเต สุญฺญตา เจโตวิมุตฺติ.\csromanspacer{} อิธ ภนฺเต ภิกฺขุ อรญฺญคโต วา รุกฺขมูลคโต วา สุญฺญาคารคโต วา อิติ ปฏิสญฺจิกฺขติ “สุญฺญมิทํ อตฺเตน วา อตฺตนิเยน วา”ติ.\csromanspacer{} อยํ วุจฺจติ ภนฺเต สุญฺญตา เจโตวิมุตฺติ.}

\prose{กตมา จ ภนฺเต อนิมิตฺตา เจโตวิมุตฺติ.\csromanspacer{} อิธ ภนฺเต ภิกฺขุ สพฺพนิมิตฺตานํ อมนสิการา อนิมิตฺตํ เจโตสมาธิํ อุปสมฺปชฺช วิหรติ.\csromanspacer{} อยํ วุจฺจติ ภนฺเต อนิมิตฺตา เจโตวิมุตฺติ.\csromanspacer{} อยํ โข ภนฺเต ปริยาโย, ยํ ปริยายํ อาคมฺม อิเม ธมฺมา นานตฺถา เจว นานาพฺยญฺชนา จ.}

\prose{กตโม จ ภนฺเต ปริยาโย, ยํ ปริยายํ อาคมฺม อิเม ธมฺมา เอกตฺถา พฺยญฺชนเมว นานํ.\csromanspacer{} ราโค ภนฺเต ปมาณกรโณ, โทโส ปมาณกรโณ, โมโห ปมาณกรโณ, เต ขีณาสวสฺส ภิกฺขุโน ปหีนา อุจฺฉินฺนมูลา ตาลาวตฺถุกตา อนภาวํกตา อายติํ อนุปฺปาทธมฺมา.\csromanspacer{} ยาวตา โข ภนฺเต อปฺปมาณา เจโตวิมุตฺติโย อกุปฺปา ตาสํ เจโตวิมุตฺติ อคฺคมกฺขายติ.\csromanspacer{} สา โข ปน อกุปฺปา เจโตวิมุตฺติ สุญฺญา ราเคน สุญฺญา โทเสน สุญฺญา โมเหน.\csromanspacer{} ราโค โข ภนฺเต กิญฺจนํ, โทโส กิญฺจนํ, โมโห กิญฺจนํ, เต ขีณาสวสฺส ภิกฺขุโน ปหีนา อุจฺฉินฺนมูลา ตาลาวตฺถุกตา อนภาวํกตา อายติํ อนุปฺปาทธมฺมา.\csromanspacer{} ยาวตา โข ภนฺเต อากิญฺจญฺญา เจโตวิมุตฺติโย อกุปฺปา ตาสํ เจโตวิมุตฺติ อคฺคมกฺขายติ.\csromanspacer{} สา}

\csromanheader{2}{7. จิตฺตสํยุตฺต}
\prose{\csromanfolio{487}โข ปน อกุปฺปา เจโตวิมุตฺติ สุญฺญา ราเคน สุญฺญา โทเสน สุญฺญา โมเหน.\csromanspacer{} ราโค โข ภนฺเต นิมิตฺตกรโณ, โทโส นิมิตฺตกรโณ, โมโห นิมิตฺตกรโณ, เต ขีณาสวสฺส ภิกฺขุโน ปหีนา อุจฺฉินฺนมูลา ตาลาวตฺถุกตา อนภาวํกตา อายติํ อนุปฺปาทธมฺมา.\csromanspacer{} ยาวตา โข ภนฺเต อนิมิตฺตา เจโตวิมุตฺติโย อกุปฺปา ตาสํ เจโตวิมุตฺติ อคฺคมกฺขายติ.\csromanspacer{} สา โข ปน อกุปฺปา เจโตวิมุตฺติ สุญฺญา ราเคน สุญฺญา โทเสน สุญฺญา โมเหน.\csromanspacer{} อยํ โข ภนฺเต ปริยาโย, ยํ ปริยายํ อาคมฺม อิเม ธมฺมา เอกตฺถา พฺยญฺชนเมว นานนฺติ.\csromanspacer{} ลาภา เต คหปติ, สุลทฺธํ เต คหปติ, ยสฺส เต คมฺภีเร พุทฺธวจเน ปญฺญาจกฺขุ กมตีติ.\csromanspacer{}\csromanspacer{}สตฺตมํ.}

\csromanheader{1}{8. นิคณฺฐนาฏปุตฺตสุตฺต}
\csromanmatikaanchor{mtk.268}
\csromantocmarkref{section}{8. นิคณฺฐนาฏปุตฺตสุตฺต}{mtk.268}
\bookmark[dest={mtk.268},level=1]{8. นิคณฺฐนาฏปุตฺตสุตฺต}
\proseitem{350}{เตน โข ปน สมเยน นิคณฺโฐ นาฏปุตฺโต\footnote{นาตปุตฺโต (สี)} มจฺฉิกาสณฺฑํ อนุปฺปตฺโต โหติ มหติยา นิคณฺฐปริสาย สทฺธิํ.\csromanspacer{} อสฺโสสิ โข จิตฺโต คหปติ “นิคณฺโฐ กิร นาฏปุตฺโต มจฺฉิกาสณฺฑํ อนุปฺปตฺโต มหติยา นิคณฺฐปริสาย สทฺธินฺ”ติ.\csromanspacer{} อถ โข จิตฺโต คหปติ สมฺพหุเลหิ อุปาสเกหิ สทฺธิํ เยน นิคณฺโฐ นาฏปุตฺโต เตนุปสงฺกมิ, อุปสงฺกมิตฺวา นิคณฺเฐน นาฏปุตฺเตน สทฺธิํ สมฺโมทิ, สมฺโมทนียํ กถํ สารณียํ วีติสาเรตฺวา เอกมนฺตํ นิสีทิ, เอกมนฺตํ นิสินฺนํ โข จิตฺตํ คหปติํ นิคณฺโฐ นาฏปุตฺโต เอตทโวจ “สทฺทหสิ ตฺวํ คหปติ ‘สมณสฺส โคตมสฺส อตฺถิ อวิตกฺโก อวิจาโร สมาธิ, อตฺถิ วิตกฺกวิจารานํ นิโรโธ’ติ”.\csromanspacer{} น ขฺวาหํ เอตฺถ ภนฺเต ภควโต สทฺธาย คจฺฉามิ “อตฺถิ อวิตกฺโก อวิจาโร สมาธิ, อตฺถิ วิตกฺกวิจารานํ นิโรโธ”ติ.}

\prose{เอวํ วุตฺเต นิคณฺโฐ นาฏปุตฺโต อุลฺโลเกตฺวา\footnote{สกํ ปริสํ อปโลเกตฺวา (สี, สฺยา, กํ), โอโลเกตฺวา (สี-ฐ, สฺยา-ฐ)} เอตทโวจ “อิทํ ภวนฺโต ปสฺสนฺตุ, ยาว อุชุโก จายํ จิตฺโต คหปติ, ยาว อสโฐ จายํ จิตฺโต คหปติ, ยาว อมายาวี จายํ จิตฺโต คหปติ.\csromanspacer{} วาตํ วา โส ชาเลน พาเธตพฺพํ มญฺเญยฺย, โย วิตกฺกวิจาเร นิโรเธตพฺพํ}

\gambhira{สฬายตนวคฺคสํยุตฺตปาฬิ}
\prose{\csromanfolio{488}มญฺเญยฺย.\csromanspacer{} สกมุฏฺฐินา วา โส คงฺคาย โสตํ อาวาเรตพฺพํ มญฺเญยฺย, โย วิตกฺกวิจาเร นิโรเธตพฺพํ มญฺเญยฺยา”ติ.}

\prose{ตํ กิํ มญฺญสิ ภนฺเต, กตมํ นุ โข ปณีตตรํ ญาณํ วา สทฺธา วาติ.\csromanspacer{} สทฺธาย โข คหปติ ญาณํเยว ปณีตตรนฺติ.\csromanspacer{} อหํ โข ภนฺเต ยาวเทว อากงฺขามิ, วิวิจฺเจว กาเมหิ วิวิจฺจ อกุสเลหิ ธมฺเมหิ สวิตกฺกํ สวิจารํ วิเวกชํ ปีติสุขํ ปฐมํ ฌานํ อุปสมฺปชฺช วิหรามิ.\csromanspacer{} อหํ โข ภนฺเต ยาวเทว อากงฺขามิ, วิตกฺกวิจารานํ วูปสมา ฯเปฯ ทุติยํ ฌานํ อุปสมฺปชฺช วิหรามิ.\csromanspacer{} อหํ โข ภนฺเต ยาวเทว อากงฺขามิ, ปีติยา จ วิราคา ฯเปฯ ตติยํ ฌานํ อุปสมฺปชฺช วิหรามิ.\csromanspacer{} อหํ โข ภนฺเต ยาวเทว อากงฺขามิ, สุขสฺส จ ปหานา ฯเปฯ จตุตฺถํ ฌานํ อุปสมฺปชฺช วิหรามิ.\csromanspacer{} น โส ขฺวาหํ ภนฺเต เอวํ ชานนฺโต เอวํ ปสฺสนฺโต กสฺส อญฺญสฺส สมณสฺส วา พฺราหฺมณสฺส วา สทฺธาย คมิสฺสามิ “อตฺถิ อวิตกฺโก อวิจาโร สมาธิ, อตฺถิ วิตกฺกวิจารานํ นิโรโธ”ติ.}

\prose{เอวํ วุตฺเต นิคณฺโฐ นาฏปุตฺโต สกํ ปริสํ อปโลเกตฺวา เอตทโวจ “อิทํ ภวนฺโต ปสฺสนฺตุ, ยาว อนุชุโก จายํ จิตฺโต คหปติ, ยาว สโฐ จายํ จิตฺโต คหปติ, ยาว มายาวี จายํ จิตฺโต คหปตี”ติ.}

\prose{อิทาเนว โข เต\footnote{อิทาเนว จ ปน (สฺยา, กํ, ก)} มยํ ภนฺเต ภาสิตํ เอวํ อาชานาม “อิทํ ภวนฺโต ปสฺสนฺตุ, ยาว อุชุโก จายํ จิตฺโต คหปติ, ยาว อสโฐ จายํ จิตฺโต คหปติ, ยาว อมายาวี จายํ จิตฺโต คหปตี”ติ.\csromanspacer{} อิทาเนว จ ปน มยํ ภนฺเต ภาสิตํ เอวํ อาชานาม “อิทํ ภวนฺโต ปสฺสนฺตุ, ยาว อนุชุโก จายํ จิตฺโต คหปติ, ยาว สโฐ จายํ จิตฺโต คหปติ, ยาว มายาวี จายํ จิตฺโต คหปตี”ติ.\csromanspacer{} สเจ เต ภนฺเต ปุริมํ สจฺจํ, ปจฺฉิมํ เต มิจฺฉา.\csromanspacer{} สเจ ปน เต ภนฺเต ปุริมํ มิจฺฉา, ปจฺฉิมํ เต สจฺจํ.\csromanspacer{} อิเม โข ปน ภนฺเต ทส สหธมฺมิกา ปญฺหา อาคจฺฉนฺติ, ยทา เนสํ อตฺถํ อาชาเนยฺยาสิ, อถ มํ ปฏิหเรยฺยาสิ สทฺธิํ นิคณฺฐปริสาย.\csromanspacer{} เอโก ปโญฺห เอโก อุทฺเทโส เอกํ เวยฺยากรณํ, เทฺว ปญฺหา เทฺว อุทฺเทสา เทฺว เวยฺยากรณานิ, ตโย ปญฺหา ตโย อุทฺเทสา ตีณิ เวยฺยากรณานิ, จตฺตาโร ปญฺหา จตฺตาโร อุทฺเทสา}

\csromanheader{2}{7. จิตฺตสํยุตฺต}
\prose{\csromanfolio{489}จตฺตาริ เวยฺยากรณานิ, ปญฺจ ปญฺหา ปญฺจ อุทฺเทสา ปญฺจ เวยฺยากรณานิ, ฉ ปญฺหา ฉ อุทฺเทสา ฉ เวยฺยากรณานิ, สตฺต ปญฺหา สตฺต อุทฺเทสา สตฺต เวยฺยากรณานิ, อฏฺฐ ปญฺหา อฏฺฐ อุทฺเทสา อฏฺฐ เวยฺยากรณานิ, นว ปญฺหา นว อุทฺเทสา นว เวยฺยากรณานิ, ทส ปญฺหา ทส อุทฺเทสา ทส เวยฺยากรณานีติ.\csromanspacer{} อถ โข จิตฺโต คหปติ นิคณฺฐํ นาฏปุตฺตํ อิเม ทส สหธมฺมิเก ปเญฺห อาปุจฺฉิตฺวา อุฏฺฐายาสนา ปกฺกามีติ.\csromanspacer{}\csromanspacer{}อฏฺฐมํ.}

\csromanheader{1}{9. อเจลกสฺสปสุตฺต}
\csromanmatikaanchor{mtk.269}
\csromantocmarkref{section}{9. อเจลกสฺสปสุตฺต}{mtk.269}
\bookmark[dest={mtk.269},level=1]{9. อเจลกสฺสปสุตฺต}
\proseitem{351}{เตน โข ปน สมเยน อเจโล กสฺสโป มจฺฉิกาสณฺฑํ อนุปฺปตฺโต โหติ จิตฺตสฺส คหปติโน ปุราณคิหิสหาโย.\csromanspacer{} อสฺโสสิ โข จิตฺโต คหปติ “อเจโล กิร กสฺสโป มจฺฉิกาสณฺฑํ อนุปฺปตฺโต อมฺหากํ ปุราณคิหิสหาโย”ติ.\csromanspacer{} อถ โข จิตฺโต คหปติ เยน อเจโล กสฺสโป เตนุปสงฺกมิ, อุปสงฺกมิตฺวา อเจเลน กสฺสเปน สทฺธิํ สมฺโมทิ, สมฺโมทนียํ กถํ สารณียํ วีติสาเรตฺวา เอกมนฺตํ นิสีทิ, เอกมนฺตํ นิสินฺโน โข จิตฺโต คหปติ อเจลํ กสฺสปํ เอตทโวจ “กีวจิรํ ปพฺพชิตสฺส ภนฺเต กสฺสปา”ติ.\csromanspacer{} ติํสมตฺตานิ โข เม คหปติ วสฺสานิ ปพฺพชิตสฺสาติ.\csromanspacer{} อิเมหิ ปน เต ภนฺเต ติํสมตฺเตหิ วสฺเสหิ อตฺถิ โกจิ อุตฺตริ มนุสฺสธมฺมา\footnote{อุตฺตริมนุสฺสธมฺโม (สฺยา, กํ)} อลมริยญาณทสฺสนวิเสโส อธิคโต ผาสุวิหาโร.\csromanspacer{} อิเมหิ โข เม คหปติ ติํสมตฺเตหิ วสฺเสหิ ปพฺพชิตสฺส นตฺถิ โกจิ อุตฺตริ มนุสฺสธมฺมา อลมริยญาณทสฺสนวิเสโส อธิคโต ผาสุวิหาโร อญฺญตฺร นคฺเคยฺยา จ มุณฺเฑยฺยา จ ปาวฬนิปฺโผฏนาย จาติ.\csromanspacer{} เอวํ วุตฺเต จิตฺโต คหปติ อเจลํ กสฺสปํ เอตทโวจ “อจฺฉริยํ วต โภ, อพฺภุตํ วต โภ, ธมฺมสฺส สฺวากฺขาตตา\footnote{สพฺพตฺถปิ เอวเมว ทิสฺสติ.}, ยตฺร หิ นาม ติํสมตฺเตหิ วสฺเสหิ น โกจิ อุตฺตริ มนุสฺสธมฺมา อลมริยญาณทสฺสนวิเสโส อธิคโต อภวิสฺส ผาสุวิหาโร อญฺญตฺร นคฺเคยฺยา จ มุณฺเฑยฺยา จ ปาวฬนิปฺโผฏนาย จาติ.}

\prose{ตุยฺหํ ปน คหปติ กีวจิรํ อุปาสกตฺตํ อุปคตสฺสาติ.\csromanspacer{} มยฺหมฺปิ โข ปน ภนฺเต ติํสมตฺตานิ วสฺสานิ อุปาสกตฺตํ อุปคตสฺสาติ.\csromanspacer{} อิเมหิ ปน เต คหปติ ติํสมตฺเตหิ วสฺเสหิ อตฺถิ โกจิ อุตฺตริ มนุสฺสธมฺมา}

\gambhira{สฬายตนวคฺคสํยุตฺตปาฬิ}
\prose{\csromanfolio{490}อลมริยญาณทสฺสนวิเสโส อธิคโต ผาสุวิหาโรติ.\csromanspacer{} คิหิโนปิ สิยา ภนฺเต, อหํ หิ ภนฺเต ยาวเทว อากงฺขามิ, วิวิจฺเจว กาเมหิ วิวิจฺจ อกุสเลหิ ธมฺเมหิ สวิตกฺกํ สวิจารํ วิเวกชํ ปีติสุกํ ปฐมํ ฌานํ อุปสมฺปชฺช วิหรามิ.\csromanspacer{} อหํ หิ ภนฺเต ยาวเทว อากงฺขามิ, วิตกฺกวิจารานํ วูปสมา ทุติยํ ฌานํ อุปสมฺปชฺช วิหรามิ.\csromanspacer{} อหํ หิ ภนฺเต ยาวเทว อากงฺขามิ, ปีติยา จ วิราคา ฯเปฯ ตติยํ ฌานํ อุปสมฺปชฺช วิหรามิ.\csromanspacer{} อหํ หิ ภนฺเต ยาวเทว อากงฺขามิ, สุขสฺส จ ปหานา ฯเปฯ จตุตฺถํ ฌานํ อุปสมฺปชฺช วิหรามิ.\csromanspacer{} สเจ โข ปนาหํ ภนฺเต ภควโต\footnote{ภควตา (สฺยา, กํ)} ปฐมตรํ กาลํ กเรยฺยํ, อนจฺฉริยํ โข ปเนตํ ยํ มํ ภควา เอวํ พฺยากเรยฺย “นตฺถิ ตํ สํโยชนํ, เยน สํโยชเนน สํยุตฺโต จิตฺโต คหปติ ปุน อิมํ โลกํ อาคจฺเฉยฺยา”ติ.\csromanspacer{} เอวํ วุตฺเต อเจโล กสฺสโป จิตฺตํ คหปติํ เอตทโวจ “อจฺฉริยํ วต โภ, อพฺภุตํ วต โภ, ธมฺมสฺส สฺวากฺขาตตา, ยตฺร หิ นาม คิหี โอทาตวสโน\footnote{คิหี โอทาตวสนา (สี, อิ)} เอวรูปํ อุตฺตริ มนุสฺสธมฺมา อลมริยญาณทสฺสนวิเสสํ อธิคมิสฺสติ\footnote{อธิคมิสฺสนฺติ (สี, อิ)} ผาสุวิหารํ.\csromanspacer{} ลเภยฺยาหํ คหปติ อิมสฺมิํ ธมฺมวินเย ปพฺพชฺชํ, ลเภยฺยํ อุปสมฺปทนฺ”ติ.}

\prose{อถ โข จิตฺโต คหปติ อเจลํ กสฺสปํ อาทาย เยน เถรา ภิกฺขู เตนุปสงฺกมิ, อุปสงฺกมิตฺวา เถเร ภิกฺขู เอตทโวจ “อยํ ภนฺเต อเจโล กสฺสโป อมฺหากํ ปุราณคิหิสหาโย, อิมํ เถรา ปพฺพาเชนฺตุ อุปสมฺปาเทนฺตุ, อหมสฺส อุสฺสุกฺกํ กริสฺสามิ จีวรปิณฺฑปาตเสนาสนคิลานปฺปจฺจยเภสชฺชปริกฺขารานนฺ”ติ.\csromanspacer{} อลตฺถ โข อเจโล กสฺสโป อิมสฺมิํ ธมฺมวินเย ปพฺพชฺชํ, อลตฺถ อุปสมฺปทํ.\csromanspacer{} อจิรูปสมฺปนฺโน จ ปนายสฺมา กสฺสโป เอโก วูปกฏฺโฐ อปฺปมตฺโต อาตาปี ปหิตตฺโต วิหรนฺโต นจิรสฺเสว ยสฺสตฺถาย กุลปุตฺตา สมฺมเทว อคารสฺมา อนคาริยํ ปพฺพชนฺติ, ตทนุตฺตรํ พฺรหฺมจริยปริโยสานํ ทิฏฺเฐว ธมฺเม สยํ อภิญฺญา สจฺฉิกตฺวา อุปสมฺปชฺช วิหาสิ, “ขีณา ชาติ, วุสิตํ พฺรหฺมจริยํ, กตํ กรณียํ, นาปรํ อิตฺถตฺตายา”ติ อพฺภญฺญาสิ.\csromanspacer{} อญฺญตโร จ ปนายสฺมา กสฺสโป อรหตํ อโหสีติ.\csromanspacer{}\csromanspacer{}นวมํ.}

\csromanheader{2}{7. จิตฺตสํยุตฺต}
\csromanheader{1}{10. คิลานทสฺสนสุตฺต}
\csromanmatikaanchor{mtk.270}
\csromantocmarkref{section}{10. คิลานทสฺสนสุตฺต}{mtk.270}
\bookmark[dest={mtk.270},level=1]{10. คิลานทสฺสนสุตฺต}
\proseitem{352}{\csromanfolio{491}เตน โข ปน สมเยน จิตฺโต คหปติ อาพาธิโก โหติ ทุกฺขิโต พาฬฺหคิลาโน.\csromanspacer{} อถ โข สมฺพหุลา อารามเทวตา วนเทวตา รุกฺขเทวตา โอสธิติณวนปฺปตีสุ อธิวตฺถา เทวตา สงฺคมฺม สมาคมฺม จิตฺตํ คหปติํ เอตทโวจุํ “ปณิเธหิ คหปติ ‘อนาคตมทฺธานํ ราชา อสฺสํ จกฺกวตฺตี’ติ”.}

\prose{เอวํ วุตฺเต จิตฺโต คหปติ ตา อารามเทวตา วนเทวตา รุกฺขเทวตา โอสธิติณนปฺปตีสุ อธิวตฺถา เทวตา เอตทโวจ “ตมฺปิ อนิจฺจํ, ตมฺปิ อทฺธุวํ, ตมฺปิ ปหาย คมนียนฺ”ติ.\csromanspacer{} เอวํ วุตฺเต จิตฺตสฺส คหปติโน มิตฺตามจฺจา ญาติสาโลหิตา จิตฺตํ คหปติํ เอตทโวจุํ “สติํ อยฺยปุตฺต อุปฏฺฐเปหิ, มา วิปฺปลปี”ติ.\csromanspacer{} กินฺตาหํ วทามิ, ยํ มํ ตุเมฺห เอวํ วเทถ “สติํ อยฺยปุตฺต อุปฏฺฐเปหิ, มา วิปฺปลปี”ติ.\csromanspacer{} เอวํ โข ตฺวํ อยฺยปุตฺต วเทสิ “ตมฺปิ อนิจฺจํ, ตมฺปิ อทฺธุวํ, ตมฺปิ ปหาย คมนียนฺ”ติ.\csromanspacer{} ตถา หิ ปน มํ อารามเทวตา วนเทวตา รุกฺขเทวตา โอสธิติณวนปฺปตีสุ อธิวตฺตา เทวตา เอวมาหํสุ “ปณิเธหิ คหปติ ‘อนาคตมทฺธานํ ราชา อสฺสํ จกฺกวตฺตี’ติ”, ตาหํ เอวํ วทามิ “ตมฺปิ อนิจฺจํ ฯเปฯ ตมฺปิ ปหาย คมนียนฺ”ติ.\csromanspacer{} กิํ ปน ตา อยฺยปุตฺต อารามเทวตา วนเทวตา รุกฺขเทวตา โอสธิติณวนปฺปตีสุ อธิวตฺถา เทวตา อตฺถวสํ สมฺปสฺสามานา เอวมาหํสุ “ปณิเธหิ คหปติ ‘อนาคตมทฺธานํ ราชา อสฺสํ จกฺกวตฺตี’ติ”.\csromanspacer{} ตาสํ โข อารามเทวตานํ วนเทวตานํ รุกฺขเทวตานํ โอสธิติณวนปฺปตีสุ อธิวตฺถานํ เทวตานํ เอวํ โหติ “อยํ โข จิตฺโต คหปติ สีลวา\footnote{สีลวนฺโต (ก)} กลฺยาณธมฺโม, สเจ ปณิทหิสฺสติ ‘อนาคตมทฺธานํ ราชา อสฺสํ จกฺกวตฺตี’ติ, ตสฺส โข อยํ อิชฺฌิสฺสติ สีลวโต เจโตปณิธิ วิสุทฺธตฺตา.\csromanspacer{} ธมฺมิโก ธมฺมิกํ ผลํ อนุปสฺสตี”ติ.\csromanspacer{} อิมํ โข ตา อารามเทวตา วนเทวตา รุกฺขเทวตา โอสธิติณวนปฺปตีสุ อธิวตฺถา เทวตา อตฺถวสํ สมฺปสฺสมานา เอวมาหํสุ “ปณิเธหิ คหปติ ‘อนาคตมทฺธานํ ราชา อสฺสํ จกฺกวตฺตี’ติ”.\csromanspacer{} ตาหํ เอวํ วทามิ “ตมฺปิ อนิจฺจํ, ตมฺปิ อทฺธุ วํ, ตมฺปิ ปหาย คมนียนฺ”ติ.}

\gambhira{สฬายตนวคฺคสํยุตฺตปาฬิ}
\prose{\csromanfolio{492}เตน หิ อยฺยปุตฺต อเมฺหปิ โอวทาหีติ.\csromanspacer{} ตสฺมา หิ โว เอวํ สิกฺขิตพฺพํ “พุทฺเธ อเวจฺจสาเทน สมนฺนาคตา ภวิสฺสาม ‘อิติปิ โส ภควา อรหํ สมฺมาสมฺพุทฺโธ วิชฺชาจรณสมฺปนฺโน สุคโต โลกวิทู อนุตฺตโร ปุริสทมฺมสารถิ สตฺถา เทวมนุสฺสานํ พุทฺโธ ภควา’ติ.\csromanspacer{} ธมฺเม อเวจฺจปฺปสาเทน สมนฺนาคตา ภวิสฺสาม ‘สฺวากฺขาโต ภควตา ธมฺโม สนฺทิฏฺฐิโก อกาลิโก เอหิปสฺสิโก โอปเนยฺยิโก ปจฺจตฺตํ เวทิตพฺโพ วิญฺญูหี’ติ.\csromanspacer{} สํเฆ อเวจฺจปฺปสาเทน สมนฺนาคตา ภวิสฺสาม ‘สุปฺปฏิปนฺโน ภควโต สาวกสํโฆ อุชุปฺปฏิปนฺโน ภควโต สาวกสํโฆ ญายปฺปฏิปนฺโน ภควโต สาวกสํโฆ สามีจิปฺปฏิปนฺโน ภควโต สาวกสํโฆ ยทิทํ จตฺตาริ ปุริสยุคานิ อฏฺฐ ปุริสปุคฺคลา เอส ภควโต สาวกสํโฆ อาหุเนยฺโย ปาหุเนยฺโย ทกฺขิเณยฺโย อญฺชลิกรณีโย อนุตฺตรํ ปุญฺญกฺเขตฺตํ โลกสฺสา’ติ.\csromanspacer{} ยํ โข ปน กิญฺจิ กุเล เทยฺยธมฺมํ, สพฺพํ ตํ อปฺปฏิวิภตฺตํ ภวิสฺสติ สีลวนฺเตหิ กลฺยาณธมฺเมหี”ติ.\csromanspacer{} เอวํ หิ โว สิกฺขิตพฺพนฺติ.\csromanspacer{} อถ โข จิตฺโต คหปติ มิตฺตามจฺเจ ญาติสาโลหิเต พุทฺเธ จ ธมฺเม จ สํเฆ จ จาเค จ สมาทเปตฺวา กาลมกาสีติ.\csromanspacer{}\csromanspacer{}ทสมํ.}

\vspace{\tipitakaparskip}
\csromancenter{\textbf{จิตฺตสํยุตฺตํ สมตฺตํ.}}
\titlehead{\textbf{ตสฺสุทฺทานํ}}
\csromangathameasure{สํโยชนํ เทฺว อิสิทตฺตา, มหโก กามภูปิ จ. \\ โคทตฺโต จ นิคณฺโฐ จ, อเจเลน คิลานทสฺสนนฺติ.}
\csromangathagroup{\csromangathastanza{สํโยชนํ เทฺว อิสิทตฺตา, มหโก กามภูปิ จ. \\ โคทตฺโต จ นิคณฺโฐ จ, อเจเลน คิลานทสฺสนนฺติ.}}

\chapterheadpageread{8. คามณิสํยุตฺต}
\csromanmatikaanchor{mtk.271}
\csromantocmarknopage{chapter}{8. คามณิสํยุตฺต}{mtk.271}
\bookmark[dest={mtk.271},level=0]{8. คามณิสํยุตฺต}
\csromanheader{1}{1. จณฺฑสุตฺต}
\csromanmatikaanchor{mtk.272}
\csromantocmarkref{section}{1. จณฺฑสุตฺต}{mtk.272}
\bookmark[dest={mtk.272},level=1]{1. จณฺฑสุตฺต}
\proseitem{353}{\csromanfolio{493}สาวตฺถินิทานํ.\csromanspacer{} อถ โข จณฺโฑ คามณิ เยน ภควา เตนุปสงฺกมิ, อุปสงฺกมิตฺวา ภควนฺตํ อภิวาเทตฺวา เอกมนฺตํ นิสีทิ, เอกมนฺตํ นิสินฺโน โข จณฺโฑ คามณิ ภควนฺตํ เอตทโวจ “โก นุ โข ภนฺเต เหตุ โก ปจฺจโย, เยน มิเธกจฺโจ จณฺโฑ จณฺโฑเตฺวว\footnote{เยน มิเธกจฺโจ จณฺโฑเตว (สี, อิ)} สงฺขํ คจฺฉติ.\csromanspacer{} โก ปน ภนฺเต เหตุ โก ปจฺจโย, เยน มิเธกจฺโจ โสรโต โสรโตเตฺวว\footnote{เยน มิเธกจฺโจ สุรโตเตว (สี, อิ)} สงฺข คจฺฉตี”ติ.\csromanspacer{} อิธ คามณิ เอกจฺจสฺส ราโค อปฺปหีโน โหติ, ราคสฺส อปฺปหีนตฺตา ปเร โกเปนฺติ, ปเรหิ โกปิยมาโน โกปํ ปาตุกโรติ, โส จณฺโฑเตฺวว สงฺขํ คจฺฉติ.\csromanspacer{} โทโส อปฺปหีโน โหติ, โทสสฺส อปฺปหีนตฺตา ปเร โกเปนฺติ, ปเรหิ โกปิยมาโน โกปํ ปาตุกโรติ, โส จณฺโฑเตฺวว สงฺขํ คจฺฉติ.\csromanspacer{} โมโห อปฺปหีโน โหติ, โมหสฺส อปฺปหีนตฺตา ปเร โกเปนฺติ, ปเรหิ โกปิยมาโน โกปํ ปาตุกโรติ, โส จณฺโฑเตฺวว สงฺขํ คจฺฉติ.\csromanspacer{} อยํ โข คามณิ เหตุ อยํ ปจฺจโย, เยน มิเธกจฺโจ จณฺโฑ จณฺโฑเตฺวว สงฺข คจฺฉติ.}

\prose{อิธ ปน คามณิ เอกจฺจสฺส ราโค ปหีโน โหติ, ราคสฺส ปหีนตฺตา ปเร น โกเปนฺติ, ปเรหิ โกปิยมาโน โกปํ น ปาตุกโรติ, โส โสรโตเตฺวว สงฺขํ คจฺฉติ.\csromanspacer{} โทโส ปหีโน โหติ, โทสสฺส ปหีนตฺตา ปเร น โกเปนฺติ, ปเรหิ โกปิยมาโน โกปํ น ปาตุกโรติ, โส โสรโตเตฺวว สงฺขํ คจฺฉติ.\csromanspacer{} โมโห ปหีโน โหติ, โมหสฺส ปหีนตฺตา ปเร น โกเปนฺติ, ปเรหิ โกปิยมาโน โกปํ น ปาตุกโรติ, โส โสรโตเตฺวว สงฺขํ คจฺฉติ.\csromanspacer{} อยํ โข คามณิ เหตุ อยํ ปจฺจโย, เยน มิเธกจฺโจ โสรโต โสรโตเตฺวว สงฺขํ คจฺฉตีติ.}

\prose{เอวํ วุตฺเต จณฺโฑ คามณิ ภควนฺตํ เอตทโว “อภิกฺกนฺตํ ภนฺเต, อภิกฺกนฺตํ ภนฺเต, เสยฺยถาปิ ภนฺเต นิกฺกุชฺชิตํ วา อุกฺกุชฺเชยฺย, ปฏิจฺฉนฺนํ}

\gambhira{สฬายตนวคฺคสํยุตฺตปาฬิ}
\prose{\csromanfolio{494}วา วิวเรยฺย, มูฬฺหสฺส วา มคฺคํ อาจิกฺเขยฺย, อนฺธกาเร วา เตลปชฺโชตํ ธาเรยฺย ‘จกฺขุมนฺโต รูปานิ ทกฺขนฺตี’ติ.\csromanspacer{} เอวเมวํ ภควตา อเนกปริยาเยน ธมฺโม ปกาสิโต, เอสาหํ ภนฺเต ภควนฺตํ สรณํ คจฺฉามิ ธมฺมญฺจ ภิกฺขุสํฆญฺจ, อุปาสกํ มํ ภควา ธาเรตุ อชฺชตคฺเค ปาณุเปตํ สรณํ คตนฺ”ติ.\csromanspacer{}\csromanspacer{}ปฐมํ.}

\csromanheader{1}{2. ตาลปุฏสุตฺต}
\csromanmatikaanchor{mtk.273}
\csromantocmarkref{section}{2. ตาลปุฏสุตฺต}{mtk.273}
\bookmark[dest={mtk.273},level=1]{2. ตาลปุฏสุตฺต}
\proseitem{354}{เอกํ สมยํ ภควา ราชคเห วิหรติ เวฬุวเน กลนฺทกนิวาเป.\csromanspacer{} อถ โข ตาลปุโฏ\footnote{ตาลปุตฺโต (สี, สฺยา, กํ)} นฏคามณิ เยน ภควา เตนุปสงฺกมิ, อุปสงฺกมิตฺวา ภควนฺตํ อภิวาเทตฺวา เอกมนฺตํ นิสีทิ, เอกมนฺตํ นิสินฺโน โข ตาลปุโฏ นฏคามณิ ภควนฺตํ เอตทโวจ “สุตํ เมตํ ภนฺเต ปุพฺพกานํ อาจริยปาจริยานํ นฏานํ ภาสมานานํ ‘โย โส นโฏ รงฺคมชฺเฌ สมชฺชมชฺเฌ สจฺจาลิเกน ชนํ หาเสติ รเมติ, โส กายสฺส เภทา ปรํ มรณา ปหาสานํ เทวานํ สหพฺยตํ อุปปชฺชตี’ติ, อิธ ภควา กิมาหา”ติ.\csromanspacer{} อลํ คามณิ ติฏฺฐเตตํ, มา มํ เอตํ ปุจฺฉีติ.\csromanspacer{} ทุติยมฺปิ โข ตาลปุโฏ นฏคามณิ ภควนฺตํ เอตทโวจ “สุตํ เมตํ ภนฺเต ปุพฺพกานํ อาจริยปาจริยานํ นฏานํ ภาสมานานํ ‘โย โส นโฏ รงฺคมชฺเฌ สมชฺชมชฺเฌ สจฺจาลิเกน ชนํ หาเสติ รเมติ, โส กายสฺส เภทา ปรํ มรณา ปหาสานํ เทวานํ สหพฺยตํ อุปปชฺชตี’ติ, อิธ ภควา กิมาหา”ติ.\csromanspacer{} อลํ คามณิ ติฏฺฐเตตํ, มา มํ เอตํ ปุจฺฉีติ.\csromanspacer{} ตติยมฺปิ โข ตาลปุโฏ นฏคามณิ ภควนฺตํ เอตทโวจ “สุตํ เมตํ ภนฺเต ปุพฺพกานํ อาจริยปาจริยานํ นฏานํ ภาสมานานํ ‘โย โส นโฏ รงฺคมชฺเฌ สมชฺชมชฺเฌ สจฺจาลิเกน ชนํ หาเสติ รเมติ, โส กายสฺส เภทา ปรํ มรณา ปหาสานํ เทวานํ สหพฺยตํ อุปปชฺชตี’ติ, อิธ ภควา กิมาหา”ติ.}

\prose{อทฺธา โข ตฺยาหํ คามณิ น ลภามิ\footnote{นาลตฺถํ (สฺยา, กํ, อิ, ก)} “อลํ คามณิ ติฏฺฐเตตํ, มา มํ เอตํ ปุจฺฉี”ติ, อปิ จ ตฺยาหํ พฺยากริสฺสามิ.\csromanspacer{} ปุพฺเพ โข คามณิ สตฺตา อวีตราคา ราคพนฺธนพทฺธา, เตสํ นโฏ รงฺคมชฺเฌ สมชฺชมชฺเฌ เย ธมฺมา รชนียา, เต อุปสํหรติ ภิยฺโยโสมตฺตาย.\csromanspacer{} ปุพฺเพ โข คามณิ สตฺตา อวีตโทสา โทสพนฺธนพทฺธา, เตสํ นโฏ รงฺคมชฺเฌ สมชฺชมชฺเฌ เย}

\vspace{\tipitakaparskip}
\csromancenter{คามณิสํยุตฺต ธมฺมา โทสนียา, เต อุปสํหรติ ภิยฺโยโสมตฺตาย.\csromanspacer{} ปุพฺเพ โข คามณิ สตฺตา อวีตโมหา โมหพนฺธนพทฺธา, เตสํ นโฏ รงฺคมชฺเฌ สมชฺชมชฺเฌ เย ธมฺมา โมหนียา, เต อุปสํหรติ ภิยฺโยโสมตฺตาย โส อตฺตนา มตฺโต ปมตฺโต ปเร มเทตฺวา ปมาเทตฺวา กายสฺส เภทา ปรํ มรณา ปหาโส นาม นิรโย, ตตฺถ อุปปชฺชติ.\csromanspacer{} สเจ โข ปนสฺส เอวํทิฏฺฐิ โหติ “โย โส นโฏ รงฺคมชฺเฌ สมชฺชมชฺเฌ สจฺจาลิเกน ชนํ หาเสติ รเมติ, โส กายสฺส เภทา ปรํ มรณา ปหาสานํ เทวานํ สหพฺยตํ อุปปชฺชตี”ติ, สาสฺส โหติ มิจฺฉาทิฏฺฐิ.\csromanspacer{} มิจฺฉาทิฏฺฐกสฺส โข ปนาหํ คามณิ ปุริสปุคฺคลสฺส ทฺวินฺนํ คตีนํ อญฺญตรํ คติํ วทามิ นิรยํ วา ติรจฺฉานโยนิํ วาติ.}
\prose{\csromanfolio{495}เอวํ วุตฺเต ตาลปุโฏ นฏคามณิ ปโรทิ, อสฺสูนิ ปวตฺเตสิ.\csromanspacer{} เอตํ โข ตฺยาหํ คามณิ นาลตฺถํ “อลํ คามณิ ติฏฺฐเตตํ, มา มํ เอตํ ปุจฺฉี”ติ.\csromanspacer{} นาหํ ภนฺเต เอตํ โรทามิ, ยํ มํ ภควา เอวมาห.\csromanspacer{} อปิ จาหํ ภนฺเต ปุพฺพเกหิ อาจริยปาจริเยหิ นเฏหิ ทีฆรตฺตํ นิกโต วญฺจิโต ปลุทฺโธ “โย โส นโฏ รงฺคมชฺเฌ สมชฺชมชฺเฌ สจฺจาลิเกน ชนํ หาเสติ รเมติ, โส กายสฺส เภทา ปรํ มรณา ปหาสานํ เทวานํ สหพฺยตํ อุปปชฺชตี”ติ.}

\prose{อภิกฺกนฺตํ ภนฺเต, อภิกฺกนฺตํ ภนฺเต, เสยฺยถาปิ ภนฺเต นิกฺกุชฺชิตํ วา อุกฺกุชฺเชยฺย, ปฏิจฺฉนฺนํ วา วิวเรยฺย, มูฬฺหสฺส วา มคฺคํ อาจิกฺเขยฺย, อนฺธกาเร วา เตลปชฺโชตํ ธาเรยฺย “จกฺขุมนฺโต รูปานิ ทกฺขนฺตี”ติ, เอวเมวํ ภควตา อเนกปริยาเยน ธมฺโม ปกาสิโต, เอสาหํ ภนฺเต ภควนฺตํ สรณํ คจฺฉามิ ธมฺมญฺจ ภิกฺขุสํฆญฺจ, ลเภยฺยาหํ ภนฺเต ภควโต สนฺติเก ปพฺพชฺชํ, ลเภยฺยํ อุปสมฺปทนฺติ.\csromanspacer{} อลตฺถ โข ตาลปุโฏ นฏคามณิ ภควโต สนฺติเก ปพฺพชฺชํ, อลตฺถ อุปสมฺปทํ.\csromanspacer{} อจิรูปสมฺปนฺโน จ ปนายสฺมา ตาลปุโฏ ฯเปฯ อรหตํ อโหสีติ.\csromanspacer{}\csromanspacer{}ทุติยํ.}

\csromanheader{1}{3. โยธาชีวสุตฺต}
\csromanmatikaanchor{mtk.274}
\csromantocmarkref{section}{3. โยธาชีวสุตฺต}{mtk.274}
\bookmark[dest={mtk.274},level=1]{3. โยธาชีวสุตฺต}
\proseitem{355}{อถ โข โยธาชีโว คามณิ เยน ภควา เตนุปสงฺกมิ, อุปสงฺกมิตฺวา ฯเปฯ เอกมนฺตํ นิสินฺโน โข โยธาชีโว คามณิ ภควนฺตํ เอตทโวจ “สุตํ เมตํ ภนฺเต ปุพฺพกานํ อาจริยปาจริยานํ โยธาชีวานํ ภาสมานานํ ‘โย โส โยธาชีโว สงฺคาเม อุสฺสหติ วายมติ,}

\gambhira{สฬายตนวคฺคสํยุตฺตปาฬิ}
\prose{\csromanfolio{496}ตเมนํ อุสฺสหนฺตํ วายมนฺตํ ปเร หนนฺติ ปริยาปาเทนฺติ, โส กายสฺส เภทา ปรํ มรณา ปรชิตานํ เทวานํ สหพฺยตํ อุปปชฺชตี’ติ, อิธ ภควา กิมาหา”ติ.\csromanspacer{} อลํ คามณิ ติฏฺฐเตตํ, มา มํ เอตํ ปุจฺฉีติ.\csromanspacer{} ทุติยมฺปิ โข ฯเปฯ.\csromanspacer{} ตติยมฺปิ โข โยธาชีโว คามณิ ภควนฺตํ เอตทโวจ “สุตํ เมตํ ภนฺเต ปุพฺพกานํ อาจริยปาจริยานํ โยธาชีวานํ ภาสมานานํ ‘โย โส โยธาชีโว สงฺคาเม อุสฺสหติ วายมติ, ตเมนํ อุสฺสหนฺตํ วายมนฺตํ ปเร หนนฺติ ปริยาปาเทนฺติ, โส กายสฺส เภทา ปรํ มรณา ปรชิตานํ เทวานํ สหพฺยตํ อุปปชฺชตี’ติ, อิธ ภควา กิมาหา”ติ.}

\prose{อทฺธา โข ตฺยาหํ คามณิ น ลภามิ “อลํ คามณิ ติฏฺฐเตตํ, มา มํ เอตํ ปุจฺฉี”ติ, อปิ จ ตฺยาหํ พฺยากริสฺสามิ.\csromanspacer{} โย โส คามณิ โยธาชีโว สงฺคาเม อุสฺสหติ วายมติ, ตสฺส ตํ จิตฺตํ ปุพฺเพ คหิตํ\footnote{หีนํ (สี, อิ)} ทุกฺกฏํ ทุปฺปณิหิตํ “อิเม สตฺตา หญฺญนฺตุ วา พชฺฌนฺตุ วา อุจฺฉิชฺชนฺตุ วา วินสฺสนฺตุ วา มา วา อเหสุํ อิติ วา”ติ, ตเมนํ อุสฺสหนฺตํ วายมนฺตํ ปเร หนนฺติ ปริยาปาเทนฺติ, โส กายสฺส เภทา ปรํ มรณา ปรชิโต นาม นิรโย, ตตฺถ อุปปชฺชตี”ติ.\csromanspacer{} สเจ โข ปนสฺส เอวํ ทิฏฺฐิ โหติ “โย โส โยธาชีโว สงฺคาเม อุสฺสหติ วายมติ, ตเมนํ อุสฺสหนฺตํ วายมนฺตํ ปเร หนนฺติ ปริยาปาเทนฺติ, โส กายสฺส เภทา ปรํ มรณา ปรชิตานํ เทวานํ สหพฺยตํ อุปปชฺชตี”ติ, สาสฺส โหติ มิจฺฉาทิฏฺฐิ.\csromanspacer{} มิจฺฉาทิฏฺฐิกสฺส โข ปนาหํ คามณิ ปุริสปุคฺคลสฺส ทฺวินฺนํ คตีนํ อญฺญตรํ คติํ วทามิ นิรยํ วา ติรจฺฉานโยนิํ วาติ.}

\prose{เอวํ วุตฺเต โยธาชีโว คามณิ ปโรทิ, อสฺสูนิ ปวตฺเตสิ.\csromanspacer{} เอตํ โข ตฺยาหํ คามณิ นาลตฺถํ “อลํ คามณิ ติฏฺฐเตตํ, มา มํ เอตํ ปุจฺฉี”ติ.\csromanspacer{} นาหํ ภนฺเต เอตํ โรทามิ, ยํ มํ ภควา เอวมาห.\csromanspacer{} อปิ จาหํ ภนฺเต ปุพฺพเกหิ อาจริยปาจริเยหิ โยธาชีเวหิ ทีฆรตฺตํ นิกโต วญฺจิโต ปลุทฺโธ “โย โส โยธาชีโว สงฺคาเม อุสฺสหติ วายมติ, ตเมนํ อุสฺสหนฺตํ วายมนฺตํ ปเร หนนฺติ ปริยาปาเทนฺติ, โส กายสฺส เภทา ปรํ มรณา ปรชิตานํ เทวานํ สหพฺยตํ อุปปชฺชตี”ติ.\csromanspacer{} อภิกฺกนฺตํ ภนฺเต ฯเปฯ อชฺชตคฺเค ปาณุเปตํ สรณํ คตนฺติ.\csromanspacer{}\csromanspacer{}ตติยํ.}

\csromanheader{1}{8. คามณิสํยุตฺต \textbf{4. หตฺถาโรหสุตฺต}}
\csromanmatikaanchor{mtk.275}
\csromantocmarkref{section}{4. หตฺถาโรหสุตฺต}{mtk.275}
\bookmark[dest={mtk.275},level=1]{4. หตฺถาโรหสุตฺต}
\proseitem{356}{\csromanfolio{497}อถ โข หตฺถาโรโห คามณิ เยน ภควา เตนุปสงฺกมิ, อุปสงฺกมิตฺวา ฯเปฯ อชฺชตคฺเค ปาณุเปตํ สรณํ คตนฺติ.}

\csromanheader{1}{5. อสฺสาโรหสุตฺต}
\csromanmatikaanchor{mtk.276}
\csromantocmarkref{section}{5. อสฺสาโรหสุตฺต}{mtk.276}
\bookmark[dest={mtk.276},level=1]{5. อสฺสาโรหสุตฺต}
\proseitem{357}{อถ โข อสฺสาโรโห คามณิ เยน ภควา เตนุปสงฺกมิ, อุปสงฺกมิตฺวา เอกมนฺตํ นิสีทิ, เอกมนฺตํ นิสินฺโน โข อสฺสาโรโห คามณิ ภควนฺตํ เอตทโวจ “สุตํ เมตํ ภนฺเต ปุพฺพกานํ อาจริยปาจริยานํ อสฺสาโรหานํ ภาสมานานํ ‘โย โส อสฺสาโรโห สงฺคาเม อุสฺสหติ วายมติ, ตเมนํ อุสฺสหนฺตํ วายมนฺตํ ปเร หนนฺติ ปริยาปาเทนฺติ, โส กายสฺส เภทา ปรํ มรณา ปรชิตานํ เทวานํ สหพฺยตํ อุปปชฺชตี’ติ, อิธ ภควา กิมาหา”ติ.\csromanspacer{} อลํ คามณิ, ติฏฺเฐเตตํ, มา มํ เอตํ ปุจฺฉีติ.\csromanspacer{} ทุติยมฺปิ โข ฯเปฯ.\csromanspacer{} ตติยมฺปิ โข อสฺสาโรโห คามณิ ภควนฺตํ เอตทโวจ “สุตํ เมตํ ภนฺเต ปุพฺพกานํ อาจริยปาจริยานํ อสฺสาโรหานํ ภาสมานานํ ‘โย โส อสฺสาโรโห สงฺคาเม อุสฺสหติ วายมติ, ตเมนํ อุสฺสหนฺตํ วายมนฺตํ ปเร หนนฺติ ปริยาปาเทนฺติ, โส กายสฺส เภทา ปรํ มรณา ปรชิตานํ เทวานํ สหพฺยตํ อุปปชฺชตี’ติ, อิธ ภควา กิมาหา”ติ.}

\prose{อทฺธา โข ตฺยาหํ คามณิ น ลภามิ “อลํ คามณิ, ติฏฺฐเตตํ, มา มํ เอตํ ปุจฺฉี”ติ.\csromanspacer{} อปิ จ โข ตฺยาหํ พฺยากริสฺสามิ.\csromanspacer{} โย โส คามณิ อสฺสาโรโห สงฺคาเม อุสฺสหติ วายมติ, ตสฺส ตํ จิตฺตํ ปุพฺเพ คหิตํ ทุกฺกฏํ ทุปฺปณิหิตํ “อิเม สตฺตา หญฺญนฺตุ วา พชฺฌนฺตุ วา อุจฺฉิชฺชนฺตุ วา วินสฺสนฺตุ วา มา อเหสุํ อิติ วา”ติ, ตเมนํ อุสฺสหนฺตํ วายมนฺตํ ปเร หนนฺติ ปริยาปาเทนฺติ, โส กายสฺส เภทา ปรํ มรณา ปรชิโต นาม นิรโย, ตตฺถ อุปปชฺชติ.\csromanspacer{} สเจ โข ปนสฺส เอวํ ทิฏฺฐิ โหติ “โย โส อสฺสาโรโห สงฺคาเม อุสฺสหติ วายมติ, ตเมนํ อุสฺสหนฺตํ วายมนฺตํ ปเร หนนฺติ ปริยาปาเทนฺติ, โส กายสฺส เภทา ปรํ มรณา ปรชิตานํ เทวานํ สหพฺยตํ อุปปชฺชตี”ติ, สาสฺส โหติ มิจฺฉาทิฏฺฐิ.\csromanspacer{} มิจฺฉาทิฏฺฐิกสฺส โข ปนาหํ คามณิ ปุริสปุคฺคลสฺส ทฺวินฺนํ คตีนํ อญฺญตรํ คติํ วทามิ นิรยํ วา ติรจฺฉานโยนิํ วาติ.}

\gambhira{สฬายตนวคฺคสํยุตฺตปาฬิ}
\prose{\csromanfolio{498}เอวํ วุตฺเต อสฺสาโรโห คามณิ ปโรทิ, อสฺสูนิ ปวตฺเตสิ.\csromanspacer{} เอตํ โข ตฺยาหํ คามณิ นาลตฺถํ “อลํ คามณิ ติฏฺฐเตตํ, มา มํ เอตํ ปุจฺฉี”ติ.\csromanspacer{} นาหํ ภนฺเต เอตํ โรทามิ, ยํ มํ ภควา เอวมาห.\csromanspacer{} อปิ จาหํ ภนฺเต ปุพฺพเกหิ อาจริยปาจริเยหิ อสฺสาโรเหหิ ทีฆรตฺตํ นิกโต วญฺจิโต ปลุทฺโธ “โย โส อสฺสาโรโห สงฺคาเม อุสฺสหติ วายมติ, ตเมนํ อุสฺสหนฺตํ วายมนฺตํ ปเร หนนฺติ ปริยาปาเทนฺติ, โส กายสฺส เภทา ปรํ มรณา ปรชิตานํ เทวานํ สหพฺยตํ อุปปชฺชตี”ติ.\csromanspacer{} อภิกฺกนฺตํ ภนฺเต ฯเปฯ อชฺชตคฺเค ปาณุเปตํ สรณํ คตนฺติ.}

\csromanheader{1}{6. อสิพนฺธกปุตฺตสุตฺต}
\csromanmatikaanchor{mtk.277}
\csromantocmarkref{section}{6. อสิพนฺธกปุตฺตสุตฺต}{mtk.277}
\bookmark[dest={mtk.277},level=1]{6. อสิพนฺธกปุตฺตสุตฺต}
\proseitem{358}{เอกํ สมยํ ภควา นาฬนฺทายํ วิหรติ ปาวาริกมฺพวเน.\csromanspacer{} อถ โข อสิพนฺธกปุตฺโต คามณิ เยน ภควา เตนุปสงฺกมิ, อุปสงฺกมิตฺวา ภควนฺตํ อภิวาเทตฺวา เอกมนฺตํ นิสีทิ, เอกมนฺตํ นิสินฺโน โข อสิพนฺธกปุตฺโต คามณิ ภควนฺตํ เอตทโวจ “พฺราหฺมณา ภนฺเต ปจฺฉา ภูมกา กามณฺฑลุกา เสวาลมาลิกา อุทโกโรหกา อคฺคิปริจารกา, เต มตํ กาลงฺกตํ อุยฺยาเปนฺติ นาม, สญฺญาเปนฺติ นาม, สคฺคํ นาม โอกฺกาเมนฺติ.\csromanspacer{} ภควา ปน ภนฺเต อรหํ สมฺมาสมฺพุทฺโธ ปโหติ ตถา กาตุํ, ยถา สพฺโพ โลโก กายสฺส เภทา ปรํ มรณา สุคติํ สคฺคํ โลกํ อุปปชฺเชยฺยา”ติ.\csromanspacer{} เตน หิ คามณิ ตญฺเญเวตฺถ ปฏิปุจฺฉิสฺสามิ ยถา เต ขเมยฺย, ตถา นํ พฺยากเรยฺยาสีติ.}

\prose{ตํ กิํ มญฺญสิ คามณิ, อิธสฺส ปุริโส ปาณาติปาตี อทินฺนาทายี กาเมสุมิจฺฉาจารี มุสาวาที ปิสุณวาโจ ผรุสวาโจ สมฺผปฺปลาปี อภิชฺฌาลุ พฺยาปนฺนจิตฺโต มิจฺฉาทิฏฺฐิโก, ตเมนํ มหา ชนกาโย สงฺคมฺม สมาคมฺม อายาเจยฺย โถเมยฺย ปญฺชลิโก อนุปริสกฺเกยฺย “อยํ ปุริโส กายสฺส เภทา ปรํ มรณา สุคติํ สคฺคํ โลกํ อุปปชฺชตู”ติ.\csromanspacer{} ตํ กิํ มญฺญสิ คามณิ, อปิ นุ โส ปุริโส มหโต ชนกายสฺส อายาจนเหตุ วา โถมนเหตุ วา ปญฺชลิกา อนุปริสกฺกนเหตุ วา กายสฺส เภทา ปรํ มรณา สุคติํ สคฺคํ โลกํ อุปปชฺเชยฺยาติ.\csromanspacer{} โน เหตํ ภนฺเต.}

\prose{เสยฺยถาปิ คามณิ ปุริโส มหติํ ปุถุสิลํ คมฺภีเร อุทกรหเท ปกฺขิเปยฺย, ตเมนํ มหา ชนกาโย สงฺคมฺม สมาคมฺม อายาเจยฺย}

\vspace{\tipitakaparskip}
\csromancenter{คามณิสํยุตฺต โถเมยฺย ปญฺชลิโก อนุปริสกฺเกยฺย “อุมฺมุชฺช โภ ปุถุสิเล, อุปฺลว โภ ปุถุสิเล, ถลมุปฺลว โภ ปุถุสิเล”ติ.\csromanspacer{} ตํ กิํ มญฺญสิ คามณิ, อปิ นุ สา ปุถุสิลา มหโต ชนกายสฺส อายาจนเหตุ วา โถมนเหตุ วา ปญฺชลิกา อนุปริสกฺกนเหตุ วา อุมฺมุชฺเชยฺย วา อุปลเวยฺย วา ถลํ วา อุปฺลเวยฺยาติ.\csromanspacer{} โนเหตํ ภนฺเต.\csromanspacer{} เอวเมว โข คามณิ โย โส ปุริโส ปาณาติปาตี อทินฺนาทายี กาเมสุมิจฺฉาจารี มุสาวาที ปิสุณวาโจ ผรุสวาโจ สมฺผปฺปลาปี อภิชฺชาลุ พฺยาปนฺนจิตฺโต มิจฺฉาทิฏฺฐิโก.\csromanspacer{} กิญฺจาปิ ตํ มหา ชนกาโย สงฺคมฺม สมาคมฺม อายาเจยฺย โถเมยฺย ปญฺชลิโก อนุปริสกฺเกยฺย “อยํ ปุริโส กายสฺส เภทา ปรํ มรณา สุคติํ สคฺคํ โลกํ อุปปชฺชตู”ติ.\csromanspacer{} อถ โข โส ปุริโส กายสฺส เภทา ปรํ มรณา อปายํ ทุคฺคติํ วินิปาตํ นิรยํ อุปปชฺเชยฺย.}
\prose{\csromanfolio{499}ตํ กิํ มญฺญสิ คามณิ, อิธสฺส\footnote{อิธ (ก), อิธ จสฺส (?)} ปุริโส ปาณาติปาตา ปฏิวิรโต อทินฺนาทานา ปฏิวิรโต กาเมสุมิจฺฉาจารา ปฏิวิรโต มุสาวาทา ปฏิวิรโต ปิสุณาย วาจาย ปฏิวิรโต ผรุสาย วาจาย ปฏิวิรโต สมฺผปฺปลาปา ปฏิวิรโต อนภิชฺฌาลุ อพฺยาปนฺนจิตฺโต สมฺมาทิฏฺฐิโก, ตเมนํ มหา ชนกาโย สงฺคมฺม สมาคมฺม อายาเจยฺย โถเมยฺย ปญฺชลิโก อนุปริสกฺเกยฺย “อยํ ปุริโส กายสฺส เภทา ปรํ มรณา อปายํ ทุคฺคติํ วินิปาตํ นิรยํ อุปปชฺชตู”ติ.\csromanspacer{} ตํ กิํ มญฺญสิ คามณิ, อปิ นุ โส ปุริโส มหโต ชนกายสฺส อายาจนเหตุ วา โถมนเหตุ วา ปญฺชลิกา อนุปริสกฺกนเหตุ วา กายสฺส เภทา ปรํ มรณา อปายํ ทุคฺคติํ วินิปาตํ นิรยํ อุปปชฺเชยฺยาติ.\csromanspacer{} โน เหตํ ภนฺเต.}

\prose{เสยฺยถาปิ คามณิ ปุริโส สปฺปิกุมฺภํ วา เตลกุมฺภํ วา คมฺภีเร อุทกรหเท โอคาเหตฺวา ภินฺเทยฺย.\csromanspacer{} ตตฺร ยาสฺส สกฺขรา วา กฐลา วา, สา อโธคามี\footnote{อโธคามินี (?)} อสฺส.\csromanspacer{} ยญฺจ ขฺวสฺส ตตฺร สปฺปิ วา เตลํ วา, ตํ อุทฺธํ คามิ อสฺส, ตเมนํ มหา ชนกาโย สงฺคมฺม สมาคมฺม อายาเจยฺย โถเมยฺย ปญฺชลิโก อนุปริสกฺเกยฺย “โอสีท โภ สปฺปิเตล, สํสีท โภ สปฺปิเตล.\csromanspacer{} อโธ คจฺฉ โภ สปฺปิเตลา”ติ.\csromanspacer{} ตํ กิํ มญฺญสิ คามณิ,}

\gambhira{สฬายตนวคฺคสํยุตฺตปาฬิ}
\prose{\csromanfolio{500}อปิ นุ ตํ สปฺปิเตลํ มหโต ชนกายสฺส อายาจนเหตุ วา โถมนเหตุ วา ปญฺชลิกา อนุปริสกฺกนเหตุ วา โอสีเทยฺย วา สํสีเทยฺย วา อโธ วา คจฺเฉยฺยาติ.\csromanspacer{} โน เหตํ ภนฺเต.\csromanspacer{} เอวเมว โข คามณิ โย โส ปุริโส ปาณาติปาตา ปฏิวิรโต อทินฺนาทานา ปฏิวิรโต กาเมสุมิจฺฉาจารา ปฏิวิรโต มุสาวาทา ปฏิวิรโต ปิสุณาย วาจาย ปฏิวิรโต ผรุสาย วาจาย ปฏิวิรโต สมฺผปฺปลาปา ปฏิวิรโต อนภิชฺฌาลุ อพฺยาปนฺนจิตฺโต สมฺมาทิฏฺฐิโก กิญฺจาปิ ตํ มหา ชนกาโย สงฺคมฺม สมาคมฺม อายาเจยฺย โถเมยฺย ปญฺชลิโก อนุปริสกฺเกยฺย “อยํ ปุริโส กายสฺส เภทา ปรํ มรณา อปายํ ทุคฺคติํ วินิปาตํ นิรยํ อุปปชฺชตู”ติ.\csromanspacer{} อถ โข โส ปุริโส กายสฺส เภทา ปรํ มรณา สุคติํ สคฺคํ โลกํ อุปปชฺเชยฺยาติ.\csromanspacer{} เอวํ วุตฺเต อสิพนฺธกปุตฺโต คามณิ ภควนฺตํ เอตทโวจ “อภิกฺกนฺตํ ภนฺเต ฯเปฯ อชฺชตคฺเค ปาณุเปตํ สรณํ คตนฺ”ติ.}

\csromanheader{1}{7. เขตฺตูปมสุตฺต}
\csromanmatikaanchor{mtk.278}
\csromantocmarkref{section}{7. เขตฺตูปมสุตฺต}{mtk.278}
\bookmark[dest={mtk.278},level=1]{7. เขตฺตูปมสุตฺต}
\proseitem{359}{เอกํ สมยํ ภควา นาฬนฺทายํ วิหรติ ปาวาริกมฺพวเน.\csromanspacer{} อถ โข อสิพนฺธกปุตฺโต คามณิ เยน ภควา เตนุปสงฺกมิ, อุปสงฺกมิตฺวา ภควนฺตํ อภิวาเทตฺวา เอกมนฺตํ นิสีทิ, เอกมนฺตํ นิสินฺโน โข อสิพนฺธกปุตฺโต คามณิ ภควนฺตํ เอตทโวจ “นนุ ภนฺเต ภควา สพฺพปาณภูตหิตานุกมฺปี วิหรตี”ติ.\csromanspacer{} เอวํ คามณิ ตถาคโต สพฺพปาณภูตหิตานุกมฺปี วิหรตีติ.\csromanspacer{} อถ กิญฺจรหิ ภนฺเต ภควา เอกจฺจานํ สกฺกจฺจํ ธมฺมํ เทเสติ, เอกจฺจานํ โน ตถา สกฺกจฺจํ ธมฺมํ เทเสตีติ.\csromanspacer{} เตน หิ คามณิ ตญฺเญเวตฺถ ปฏิปุจฺฉิสฺสามิ.\csromanspacer{} ยถา เต ขเมยฺย, ตถา นํ พฺยากเรยฺยาสิ.\csromanspacer{} ตํ กิํ มญฺญสิ คามณิ, อิธสฺสุ\footnote{อิธ (สี, สฺยา, กํ, อิ)} กสฺสกสฺส คหปติโน ตีณิ เขตฺตานิ, เอกํ เขตฺตํ อคฺคํ, เอกํ เขตฺตํ มชฺฌิมํ, เอกํ เขตฺตํ หีนํ ชงฺคลํ อูสรํ ปาปภูมิ.\csromanspacer{} ตํ กิํ มญฺญสิ คามณิ, อสุ กสฺสโก คหปติ พีชานิ ปติฏฺฐาเปตุกาโม กตฺถ ปฐมํ ปติฏฺฐาเปยฺย, ยํ วา อทุํ เขตฺตํ อคฺคํ, ยํ วา อทุํ เขตฺตํ มชฺฌิมํ, ยํ วา อทุํ เขตฺตํ หีนํ ชงฺคลํ อูสรํ ปาปภูมีติ.\csromanspacer{} อสุ ภนฺเต กสฺสโก คหปติ พีชานิ ปติฏฺฐาเปตุกาโม ยํ อทุํ เขตฺตํ อคฺคํ, ตตฺถ ปติฏฺฐาเปยฺย.\csromanspacer{} ตตฺถ ปติฏฺฐาเปตฺวา ยํ อทุํ เขตฺตํ มชฺฌิมํ, ตตฺถ}

\vspace{\tipitakaparskip}
\csromancenter{คามณิสํยุตฺต ปติฏฺฐาเปยฺย.\csromanspacer{} ตตฺถ ปติฏฺฐาเปตฺวา ยํ อทุํ เขตฺตํ หีนํ ชงฺคลํ อูสรํ ปาปภูมิ, ตตฺถ ปติฏฺฐาเปยฺยปิ โนปิ ปติฏฺฐาเปยฺย.\csromanspacer{} ตํ กิสฺส เหตุ, อนฺตมโส โคภตฺตมฺปิ ภวิสฺสตีติ.\csromanspacer{} เสยฺยถาปิ คามณิ ยํ อทุํ เขตฺตํ อคฺคํ, เอวเมว มยฺหํ ภิกฺขุภิกฺขุนิโย, เตสาหํ ธมฺมํ เทเสมิ อาทิกลฺยาณํ มชฺเฌกลฺยาณํ ปริโยสานกลฺยาณํ สาตฺถํ สพฺยญฺชนํ เกวลปริปุณฺณํ ปริสุทฺธํ พฺรหฺมจริยํ ปกาเสมิ.\csromanspacer{} ตํ กิสฺส เหตุ, เอเต หิ คามณิ มํทีปา มํเลณา มํตาณา มํสรณา วิหรนฺติ.\csromanspacer{} เสยฺยถาปิ คามณิ ยํ อทุํ เขตฺตํ มชฺฌิมํ, เอวเมว มยฺหํ อุปาสก-อุปาสิกาโย, เตสํ ปาหํ ธมฺมํ เทเสมิ อาทิกลฺยาณํ มชฺเฌกลฺยาณํ ปริโยสานกลฺยาณํ สาตฺถํ สพฺยญฺชนํ เกวลปริปุณฺณํ ปริสุทฺธํ พฺรหฺมจริยํ ปกาเสมิ.\csromanspacer{} ตํ กิสฺส เหตุ, เอเต หิ คามณิ มํทีปา มํเลณา มํตาณา มํสรณา วิหรนฺติ.\csromanspacer{} เสยฺยถาปิ คามณิ ยํ อทุํ เขตฺตํ หีนํ ชงฺคลํ อูสรํ ปาปภูมิ, เอวเมว มยฺหํ อญฺญติตฺถิยา สมณพฺราหฺมณปริพฺพาชกา, เตสํ ปาหํ ธมฺมํ เทเสมิ อาทิกลฺยาณํ มชฺเฌกลฺยาณํ ปริโยสานกลฺยาณํ สาตฺถํ สพฺยญฺชนํ เกวลปริปุณฺณํ ปริสุทฺธํ พฺรหฺมจริยํ ปกาเสมิ.\csromanspacer{} ตํ กิสฺส เหตุ, อปฺเปว นาม เอกํ ปทมฺปิ อาชาเนยฺยุํ, ตํ เนสํ อสฺส ทีฆรตฺตํ หิตาย สุขายาติ.}
\prose{\csromanfolio{501}เสยฺยถาปิ คามณิ ปุริสสฺส ตโย อุทกมณิกา.\csromanspacer{} เอโก อุทกมณิโก อจฺฉิทฺโท อหารี อปริหารี, เอโก อุทกมณิโก อจฺฉิทฺโท หารี ปริหารี, เอโก อุทกมณิโก ฉิทฺโท หารี ปริหารี.\csromanspacer{} ตํ กิํ มญฺญสิ คามณิ, อสุ ปุริโส อุทกํ นิกฺขิปิตุกาโม กตฺถ ปฐมํ นิกฺขิเปยฺย.\csromanspacer{} โย วา โส อุทกมณิโก อจฺฉิทฺโท อหารี อปริหารี, โย วา โส อุทกมณิโก อจฺฉิทฺโท หารี ปริหารี, โย วา โส อุทกมณิโก ฉิทฺโท หารี ปริหารีติ.\csromanspacer{} อสุ ภนฺเต ปุริโส อุทกํ นิกฺขิปิตุกาโม โย โส อุทกมณิโก อจฺฉิทฺโท อหารี อปริหารี, ตตฺถ นิกฺขิเปยฺย.\csromanspacer{} ตตฺถ นิกฺขิปิตฺวา โย โส อุทกมณิโก อจฺฉิทฺโท หารี ปริหารี, ตตฺถ นิกฺขิเปยฺย, ตตฺถ นิกฺขิปิตฺวา โย โส อุทกมณิโก ฉิทฺโท หารี ปริหารี, ตตฺถ นิกฺขิเปยฺยปิ โนปิ นิกฺขิเปยฺย.\csromanspacer{} ตํ กิสฺส เหตุ, อนฺตมโส ภณฺฑโธวนมฺปิ ภวิสฺสตีติ.\csromanspacer{} เสยฺยถาปิ คามณิ โย โส อุทกมณิโก อจฺฉิทฺโท อหารี อปริหารี, เอวเมว มยฺหํ ภิกฺขุภิกฺขุนิโย, เตสาหํ ธมฺมํ เทเสมิ อาทิกลฺยาณํ มชฺเฌกลฺยาณํ}

\gambhira{สฬายตนวคฺคสํยุตฺตปาฬิ}
\prose{\csromanfolio{502}ปริโยสานกลฺยาณํ สาตฺถํ สพฺยญฺชนํ เกวลปริปุณฺณํ ปริสุทฺธํ พฺรหฺมจริยํ ปกาเสมิ, ตํ กิสฺส เหตุ, เอเต หิ คามณิ มํทีปา มํเลณา มํตาณา มํสรณา วิหรนฺติ.\csromanspacer{} เสยฺยถาปิ คามณิ โย โส อุทกมณิ โก อจฺฉิทฺโท หารี ปริหารี เอวเมว มยฺหํ อุปาสก-อุปาสิกาโย, เตสาหํ ธมฺมํ เทเสมิ อาทิกลฺยาณํ มชฺเฌกลฺยาณํ ปริโยสานกลฺยาณํ สาตฺถํ สพฺยญฺชนํ เกวลปริปุณฺณํ ปริสุทฺธํ พฺรหฺมจริยํ ปกาเสมิ, ตํ กิสฺส เหตุ, เอเต หิ คามณิ มํทีปา มํเลณา มํตาณา มํสรณา วิหรนฺติ.\csromanspacer{} เสยฺยถาปิ คามณิ โย โส อุทกมณิโก ฉิทฺโท หารี ปริหารี, เอวเมว มยฺหํ อญฺญติตฺถิยา สมณพฺราหฺมณปริพฺพาชกา, เตสาหํ ธมฺมํ เทเสมิ อาทิกลฺยาณํ มชฺเฌกลฺยาณํ ปริโยสานกลฺยาณํ สาตฺถํ สพฺยญฺชนํ เกวลปริปุณฺณํ ปุริสุทฺธํ พฺรหฺมจริยํ ปกาเสมิ, ตํ กิสฺส เหตุ, อปฺเปว นาม เอกํ ปทมฺปิ อาชาเนยฺยุํ, ตํ เนสํ อสฺส ทีฆรตฺตํ หิตาย สุขายาติ.\csromanspacer{} เอวํ วุตฺเต อสิพนฺธกปุตฺโต คามณิ ภควนฺตํ เอตโวจ “อภิกฺกนฺตํ ภนฺเต, อภิกฺกนฺตํ ภนฺเต ฯเปฯ อุปาสกํ มํ ภควา ธาเรตุ อชฺชตคฺเค ปาณุเปตํ สรณํ คตนฺ”ติ.}

\csromanheader{1}{8. สงฺขธมสุตฺต}
\csromanmatikaanchor{mtk.279}
\csromantocmarkref{section}{8. สงฺขธมสุตฺต}{mtk.279}
\bookmark[dest={mtk.279},level=1]{8. สงฺขธมสุตฺต}
\proseitem{360}{เอกํ สมยํ ภควา นาฬนฺทายํ วิหรติ ปาวาริกมฺพวเน.\csromanspacer{} อถ โข อสิพนฺธกปุตฺโต คามณิ นิคณฺฐสาวโก เยน ภควา เตนุปสงฺกมิ, อุปสงฺกมิตฺวา เอกมนฺตํ นิสีทิ.\csromanspacer{} เอกมนฺตํ นิสินฺนํ โข อสิพนฺธกปุตฺตํ คามณิํ ภควา เอตทโวจ “กถํ นุ โข คามณิ นิคณฺโฐ นาฏปุตฺโต สาวกานํ ธมฺมํ เทเสตี”ติ.\csromanspacer{} เอวํ โข ภนฺเต นิคณฺโฐ นาฏปุตฺโต สาวกานํ ธมฺมํ เทเสติ “โย โกจิ ปาณํ อติปาเตติ, สพฺโพ โส อาปายิโก เนรยิโก.\csromanspacer{} โย โกจิ อทินฺนํ อาทิยติ, สพฺโพ โส อาปายิโก เนรยิโก.\csromanspacer{} โย โกจิ กาเมสุ มิจฺฉา จรติ, สพฺโพ โส อาปายิโก เนรยิโก.\csromanspacer{} โย โกจิ มุสา ภณติ, สพฺโพ โส อาปายิโก เนรยิโก.\csromanspacer{} ยํพหุลํ ยํพหุลํ วิหรติ, เตน เตน นียตี”ติ.\csromanspacer{} เอวํ โข ภนฺเต นิคณฺโฐ นาฏปุตฺโต สาวกานํ ธมฺมํ เทเสตีติ.\csromanspacer{} ยํพหุลํ ยํพหุลญฺจ คามณิ วิหรติ, เตน เตน นียติ, เอวํ สนฺเต น โกจิ อาปายิโก เนรยิโก ภวิสฺสติ, ยถา นิคณฺฐสฺส นาฏปุตฺตสฺส วจนํ.}

\csromanheader{2}{8. คามณิสํยุตฺต}
\prose{\csromanfolio{503}ตํ กิํ มญฺญสิ คามณิ, โย โส ปุริโส ปาณาติปาตี รตฺติยา วา ทิวสสฺส วา สมยาสมยํ อุปาทาย กตโม พหุตโร สมโย ยํ วา โส ปาณมติปาเตติ, ยํ วา โส ปาณํ นาติปาเตตีติ.\csromanspacer{} โย โส ภนฺเต ปุริโส ปาณาติปาตี รตฺติยา วา ทิวสสฺส วา สมยาสมยํ อุปาทาย อปฺปตโร โส สมโย ยํ โส ปาณมติปาเตติ, อถ โข เสฺวว พหุตโร สมโย ยํ โส ปาณํ นาติปาเตตีติ.\csromanspacer{} ยํ พหุลํ ยํ พหุลญฺจ คามณิ วิหรติ เตน เตน นียตีติ, เอวํ สนฺเต น โกจิ อาปายิโก เนรยิโก ภวิสฺสติ, ยถา นิคณฺฐสฺส นาฏปุตฺตสฺส วจนํ.}

\prose{ตํ กิํ มญฺญสิ คามณิ, โย โส ปุริโส อทินฺนาทายี รตฺติยา วา ทิวสสฺส วา สมยาสมยํ อุปาทาย กตโม พหุตโร สมโย ยํ วา โส อทินฺนํ อาทิยติ, ยํ วา โส อทินฺนํ นาทิยตีติ.\csromanspacer{} โย โส ภนฺเต ปุริโส อทินฺนาทายี รตฺติยา วา ทิวสสฺส วา สมยาสมยํ อุปาทาย อปฺปตโร โส สมโย ยํ โส อทินฺนํ อาทิยติ, อถ โข เสฺวว พหุตโร สมโย ยํ โส อทินฺนํ นาทิยตีติ.\csromanspacer{} ยํพหุลํ ยํพหุลญฺจ คามณิ วิหรติ เตน เตน นียตีติ, เอวํ สนฺเต น โกจิ อาปายิโก เนรยิโก ภวิสฺสติ, ยถา นิคณฺฐสฺส นาฏปุตฺตสฺส วจนํ.}

\prose{ตํ กิํ มญฺญสิ คามณิ, โย โส ปุริโส กาเมสุมิจฺฉาจารี รตฺติยา วา ทิวสสฺส วา สมยาสมยํ อุปาทาย กตโม พหุตโร สมโย ยํ วา โส กาเมสุ มิจฺฉา จรติ, ยํ วา โส กาเมสุ มิจฺฉา น จรตีติ.\csromanspacer{} โย โส ภนฺเต ปุริโส กาเมสุมิจฺฉาจารี รตฺติยา วา ทิวสสฺส วา สมยาสมยํ อุปาทาย อปฺปตโร โส สมโย ยํ โส กาเมสุ มิจฺฉา จรติ, อถ โข เสฺวว พหุตโร สมโย ยํ โส กาเมสุ มิจฺฉา น จรตีติ.\csromanspacer{} ยํพหุลํ ยํพหุลญฺจ คามณิ วิหรติ เตน เตน นียตีติ, เอวํ สนฺเต น โกจิ อาปายิโก เนรยิโก ภวิสฺสติ, ยถา นิคณฺฐสฺส นาฏปุตฺตสฺส วจนํ.}

\prose{ตํ กิํ มญฺญสิ คามณิ, โย โส ปุริโส มุสาวาที รตฺติยา วา ทิวสสฺส วา สมยาสมยํ อุปาทาย กตโม พหุตโร สมโย ยํ วา โส มุสา ภณติ, ยํ วา โส มุสา น ภณตีติ.\csromanspacer{} โย โส ภนฺเต ปุริโส มุสาวาที รตฺติยา วา ทิวสสฺส วา สมยาสมยํ}

\gambhira{สฬายตนวคฺคสํยุตฺตปาฬิ}
\prose{\csromanfolio{504}อุปาทาย อปฺปตโร โส สมโย ยํ โส มุสา ภณติ, อถ โข เสฺวว พหุตโร สมโย ยํ โส มุสา น ภณตีติ.\csromanspacer{} ยํพหุลํ ยํ พหุลญฺจ คามณิ วิหรติ เตน เตน นียตีติ, เอวํ สนฺเต น โกจิ อาปายิโก เนรยิโก ภวิสฺสติ, ยถา นิคณฺฐสฺส นาฏปุตฺตสฺส วจนํ.}

\prose{อิธ คามณิ เอกจฺโจ สตฺถา เอวํวาที โหติ เอวํทิฏฺฐิ\footnote{เอวํทิฏฺฐี (ก)} “โย โกจิ ปาณมติปาเตติ, สพฺโพ โส อาปายิโก เนรยิโก.\csromanspacer{} โย โกจิ อทินฺนํ อาทิยติ, สพฺโพ โส อาปายิโก เนรยิโก.\csromanspacer{} โย โกจิ กาเมสุ มิจฺฉา จรติ, สพฺโพ โส อาปายิโก เนรยิโก.\csromanspacer{} โย โกจิ มุสา ภณติ, สพฺโพ โส อาปายิโก เนรยิโก”ติ.\csromanspacer{} ตสฺมิํ โข ปน คามณิ สตฺถริ สาวโก อภิปฺปสนฺโน โหติ, ตสฺส เอวํ โหติ “มยฺหํ โข สตฺถา เอวํวาที เอวํทิฏฺฐิ ‘โย โกจิ ปาณมติปาเตติ, สพฺโพ โส อาปายิโก เนรยิโก’ติ.\csromanspacer{} อตฺถิ โข ปน มยา ปาโณ อติปาติโต, อหมฺปมฺหิ อาปายิโก เนรยิโก”ติ ทิฏฺฐิํ ปฏิลภติ.\csromanspacer{} ตํ คามณิ วาจํ อปฺปหาย ตํ จิตฺตํ อปฺปหาย ตํ ทิฏฺฐิํ อปฺปฏินิสฺสชฺชิตฺวา ยถาภตํ นิกฺขิตฺโต เอวํ นิรเย. “มยฺหํ โข สตฺถา เอวํวาที เอวํทิฏฺฐิ ‘โย โกจิ อทินฺนํ อาทิยติ, สพฺโพ โส อาปายิโก เนรยิโก’ติ.\csromanspacer{} อตฺถิ โข ปน มยา อทินฺนํ อาทินฺนํ, อหมฺปมฺหิ อาปายิโก เนรยิโก’ติ ทิฏฺฐิํ ปฏิลภติ.\csromanspacer{} ตํ คามณิ วาจํ อปฺปหาย ตํ จิตฺตํ อปฺปหาย ตํ ทิฏฺฐิํ อปฺปฏินิสฺสชฺชิตฺวา ยถาภตํ นิกฺขิตฺโต เอวํ นิรเย. “มยฺหํ โข สตฺถา เอวํ วาที เอวํทิฏฺฐิ ‘โย โกจิ กาเมสุ มิจฺฉา จรติ, สพฺโพ โส อาปายิโก เนรยิโก’ติ.\csromanspacer{} อตฺถิ โข ปน มยา กาเมสุ มิจฺฉา จิณฺณํ, อหมฺปมฺหิ อาปายิโก เนรยิโก”ติ ทิฏฺฐิํ ปฏิลภติ.\csromanspacer{} ตํ คามณิ วาจํ อปฺปหาย ตํ จิตฺตํ อปฺปหาย ตํ ทิฏฺฐิํ อปฺปฏินิสฺสชฺชิตฺวา ยถาภตํ นิกฺขิตฺโต เอวํ นิรเย. “มยฺหํ โข สตฺถา เอวํวาที เอวํทิฏฺฐิ ‘โย โกจิ มุสา ภณติ สพฺโพ โส อาปายิโก เนรยิโก’ติ.\csromanspacer{} อตฺถิ โข ปน มยา มุสา ภณิตํ, อหมฺปมฺหิ อาปายิโก เนรยิโก”ติ ทิฏฺฐิํ ปฏิลภติ.\csromanspacer{} ตํ คามณิ วาจํ อปฺปหาย ตํ จิตฺตํ อปฺปหาย ตํ ทิฏฺฐิํ อปฺปฏินิสฺสชฺชิตฺวา ยถาภูตํ นิกฺขิตฺโต เอวํ นิรเย.}

\csromanheader{2}{8. คามณิสํยุตฺต}
\prose{\csromanfolio{505}อิธ ปน คามณิ ตถาคโต โลเก อุปฺปชฺชติ อรหํ สมฺมาสมฺพุทฺโธ วิชฺชาจรณสมฺปนฺโน สุคโต โลกวิทู อนุตฺตโร ปุริสทมฺมสารถิ สตฺถา เทวมนุสฺสานํ พุทฺโธ ภควา, โส อเนกปริยาเยน ปาณาติปาตํ ครหติ วิครหติ, “ปาณาติปาตา วิรมถา”ติ จาห.\csromanspacer{} อทินฺนาทานํ ครหติ วิครหติ, “อทินฺนาทานา วิรมถา”ติ จาห.\csromanspacer{} กาเมสุมิจฺฉาจารํ ครหติ วิครหติ, “กาเมสุจฺฉาจารา วิรมถา”ติ จาห.\csromanspacer{} มุสาวาทํ ครหติ วิครหติ, “มุสาวาทา วิรมถา”ติ จาห.\csromanspacer{} ตสฺมิํ โข ปน คามณิ สตฺถริ สาวโก อภิปฺปสนฺโน โหติ, โส อิติ ปฏิสญฺจิกฺขติ “ภควา โข อเนกปริยาเยน ปาณาติปาตํ ครหติ วิครหติ, ‘ปาณาติปาตา วิรมถา’ติ จาห.\csromanspacer{} อตฺถิ โข ปน มยา ปาโณ อติปาติโต ยาวตโก วา ตาวตโก วา, โย โข ปน มยา ปาโณ อติปาติโต ยาวตโก วา ตาวตโก วา, ตํ น สุฏฺฐุ ตํ น สาธุ.\csromanspacer{} อหญฺเจว\footnote{อหญฺเจ (?)} โข ปน ตปฺปจฺจยา วิปฺปฏิสารี อสฺสํ, น เมตํ ปาปํ กมฺมํ\footnote{ปาปกมฺมํ (สฺยา, กํ, อิ, ก)} อกตํ ภวิสฺสตี”ติ.\csromanspacer{} โส อิติปิ ปฏิสงฺขาย ตญฺเจว ปาณาติปาตํ ปชหติ, อายติํ จ ปาณาติปาตา ปฏิวิรโต โหติ.\csromanspacer{} เอวเมตสฺส ปาปสฺส กมฺมสฺส ปหานํ โหติ, เอวเมตสฺส ปาปสฺส กมฺมสฺส สมติกฺกโม โหติ.}

\prose{“ภควา โข อเนกปริยาเยน อทินฺนาทานํ ครหติ วิครหติ, ‘อทินฺนาทานา วิรมถา’ติ จาห.\csromanspacer{} อตฺถิ โข ปน มยา อทินฺนํ อาทินฺนํ ยาวตกํ วา ตาวตกํ วา, ยํ โข ปน มยา อทินฺนํ อาทินฺนํ ยาวตกํ วา ตาวตกํ วา, ตํ น สุฏฺฐุ ตํ น สาธุ.\csromanspacer{} อหญฺเจว โข ปน ตปฺปจฺจยา วิปฺปฏิสารี อสฺสํ, น เมตํ ปาปํ กมฺมํ อกตํ ภวิสฺสตี”ติ.\csromanspacer{} โส อิติ ปฏิสงฺขาย ตญฺเจว อทินฺนาทานํ ปชหติ, อายติํ จ อทินฺนาทานา ปฏิวิรโต โหติ.\csromanspacer{} เอวเมตสฺส ปาปสฺส กมฺมสฺส ปหานํ โหติ, เอวเมตสฺส ปาปสฺส กมฺมสฺส สมติกฺกโม โหติ.}

\prose{“ภควา โข ปน อเนกปริยาเยน กาเมสุมิจฺฉาจารํ ครหติ วิครหติ, ‘กาเมสุมิจฺฉาจารา วิรมถา’ติ จาห.\csromanspacer{} อตฺถิ โข ปน มยา กาเมสุ มิจฺฉา จิณฺณํ ยาวตกํ วา ตาวตกํ วา, ยํ โข ปน มยา กาเมสุ มิจฺฉา จิณฺณํ ยาวตกํ วา ตาวตกํ วา, ตํ น สุฏฺฐุ ตํ น สาธุ.\csromanspacer{} อหญฺเจว โข ปน ตปฺปจฺจยา วิปฺปฏิสารี อสฺสํ,}

\gambhira{สฬายตนวคฺคสํยุตฺตปาฬิ}
\prose{\csromanfolio{506}น เมตํ ปาปํ กมฺมํ อกตํ ภวิสฺสตี”ติ.\csromanspacer{} โส อิติ ปฏิสงฺขาย ตญฺเจว กาเมสุมิจฺฉาจารํ ปชหติ, อายติํ จ กาเมสุมิจฺฉาจารา ปฏิวิรโต โหติ.\csromanspacer{} เอวเมตสฺส ปาปสฺส กมฺมสฺส ปหานํ โหติ, เอวเมตสฺส ปาปสฺส กมฺมสฺส สมติกฺกโม โหติ.}

\prose{“ภควา โข ปน อเนกปริยาเยน มุสาวาทํ ครหติ วิครหติ, ‘มุสาวาทา วิรมถา’ติ จาห.\csromanspacer{} อตฺถิ โข ปน มยา มุสา ภณิตํ ยาวตกํ วา ตาวตกํ วา, ยํ โข ปน มยา มุสา ภณิตํ ยาวตกํ วา ตาวตกํ วา, ตํ น สุฏฺฐุ ตํ น สาธุ.\csromanspacer{} อหญฺเจว โข ปน ตปฺปจฺจยา วิปฺปฏิสารี อสฺสํ, น เมตํ ปาปํ กมฺมํ อกตํ ภวิสฺสตี”ติ.\csromanspacer{} โส อิติ ปฏิสงฺขาย ตญฺเจว มุสาวาทํ ปชหติ, อายติํ จ มุสาวาทา ปฏิวิรโต โหติ.\csromanspacer{} เอวเมตสฺส ปาปสฺส กมฺมสฺส ปหานํ โหติ, เอวเมตสฺส ปาปสฺส กมฺมสฺส สมติกฺกโม โหติ.}

\prose{โส ปาณาติปาตํ ปหาย ปาณาติปาตา ปฏิวิรโต โหติ.\csromanspacer{} อทินฺนาทานํ ปหาย อทินฺนาทานา ปฏิวิรโต โหติ.\csromanspacer{} กาเมสุมิจฺฉาจารํ ปหาย กาเมสุมิจฺฉาจารา ปฏิวิรโต โหติ.\csromanspacer{} มุสาวาทํ ปหาย มุสาวาทา ปฏิวิรโต โหติ.\csromanspacer{} ปิสุณํ วาจํ ปหาย ปิสุณาย วาจาย ปฏิวิรโต โหติ.\csromanspacer{} ผรุสํ วาจํ ปหาย ผรุสาย วาจาย ปฏิวิรโต โหติ.\csromanspacer{} สมฺผปฺปลาปํ ปหาย สมฺผปฺปลาปา ปฏิวิรโต โหติ.\csromanspacer{} อภิชฺฌํ ปหาย อนภิชฺฌาลุ โหติ.\csromanspacer{} พฺยาปาทปฺปโทสํ ปหาย อพฺยาปนฺนจิตฺโต โหติ.\csromanspacer{} มิจฺฉาทิฏฺฐิํ ปหาย สมฺมาทิฏฺฐิโก โหติ.}

\prose{ส โข โส คามณิ อริยสาวโก เอวํ วิคตาภิชฺโฌ วิคตพฺยาปาโท อสมฺมูโฬฺห สมฺปชาโน ปฏิสฺสโต เมตฺตาสหคเตน เจตสา เอกํ ทิสํ ผริตฺวา วิหรติ.\csromanspacer{} ตถา ทุติยํ.\csromanspacer{} ตถา ตติยํ.\csromanspacer{} ตถา จตุตฺถํ.\csromanspacer{} อิติ อุทฺธมโธ ติริยํ สพฺพธิ สพฺพตฺตตาย สพฺพาวนฺตํ โลกํ เมตฺตาสหคเตน เจตสา วิปุเลน มหคฺคเตน อปฺปมาเณน อเวเรน อพฺยาปชฺเชน ผริตฺวา วิหรติ.\csromanspacer{} เสยฺยถาปิ คามณิ พลวา สงฺขธมฺโม อปฺปกสิเรเนว จตุทฺทิสา วิญฺญาเปยฺย.\csromanspacer{} เอวเมว โข คามณิ เอวํ ภาวิตาย เมตฺตาย เจโตวิมุตฺติยา เอวํ พหุลีกตาย ยํ ปมาณกตํ กมฺมํ น ตํ ตตฺราวสิสฺสติ, น ตํ ตตฺราวติฏฺฐติ.}

\csromanheader{2}{8. คามณิสํยุตฺต}
\prose{\csromanfolio{507}ส โข โส คามณิ อริยสาวโก เอวํ วิคตาภิชฺโฌ วิคตพฺยาปาโท อสมฺมูโฬฺห สมฺปชาโน ปฏิสฺสโต กรุณาสหคเตน เจตสา ฯเปฯ มุทิตาสหคเตน เจตสา ฯเปฯ อุเปกฺขาสหคเตน เจตสา เอกํ ทิสํ ผริตฺวา วิหรติ.\csromanspacer{} ตถา ทุติยํ.\csromanspacer{} ตถา ตติยํ.\csromanspacer{} ตถา จตุตฺถํ.\csromanspacer{} อิติ อุทฺธมโธ ติริยํ สพฺพธิ สพฺพตฺตตาย สพฺพาวนฺตํ โลกํ อุเปกฺขาสหคเตน เจตสา วิปุเลน มหคฺคเตน อปฺปมาเณน อเวเรน อพฺยาปชฺเชน ผริตฺวา วิหรติ.\csromanspacer{} เสยฺยถาปิ คามณิ พลวา สงฺขธโม อปฺปกสิเรเนว จตุทฺทิสา วิญฺญาเปยฺย.\csromanspacer{} เอวเมว โข คามณิ เอวํ ภาวิตาย อุเปกฺขาย เจโตวิมุตฺติยา เอวํ พหุลีกตาย ยํ ปมาณกตํ กมฺมํ, น ตํ ตตฺราวสิสฺสติ, น ตํ ตตฺราวติฏฺฐตีติ.\csromanspacer{} เอวํ วุตฺเต อสิพนฺธกปุตฺโต คามณิ ภควนฺตํ เอตทโวจ “อภิกฺกมนฺตํ ภนฺเต, อภิกฺกนฺตํ ภนฺเต ฯเปฯ อุปาสกํ มํ ภควา ธาเรตุ อชฺชตคฺเค ปาณุเปตํ สรณํ คตนฺ”ติ.}

\csromanheader{1}{9. กุลสุตฺต}
\csromanmatikaanchor{mtk.280}
\csromantocmarkref{section}{9. กุลสุตฺต}{mtk.280}
\bookmark[dest={mtk.280},level=1]{9. กุลสุตฺต}
\proseitem{361}{เอกํ สมยํ ภควา โกสเลสุ จาริกํ จรมาโน มหตา ภิกฺขุสํเฆน สทฺธิํ เยน นาฬนฺทา ตทวสริ, ตตฺร สุทํ ภควา นาฬนฺทายํ วิหรติ ปาวาริกมฺพวเน.}

\prose{เตน โข ปน สมเยน นาฬนฺทา ทุพฺภิกฺขา โหติ ทฺวีหิติกา เสตฏฺฐิกา สลากาวุตฺตา.\csromanspacer{} เตน โข ปน สมเยน นิคณฺโฐ นาฏปุตฺโต นาฬนฺทายํ ปฏิวสติ มหติยา นิคณฺฐปริสาย สทฺธิํ.\csromanspacer{} อถ โข อสิพนฺธกปุตฺโต คามณิ นิคณฺฐสาวโก เยน นิคณฺโฐ นาฏปุตฺโต เตนุปสงฺกมิ, อุปสงฺกมิตฺวา นิคณฺฐํ นาฏปุตฺตํ อภิวาเทตฺวา เอกมนฺตํ นิสีทิ, เอกมนฺตํ นิสินฺนํ โข อสิพนฺธกปุตฺตํ คามณิํ นิคณฺโฐ นาฏปุตฺโต เอตทโวจ “เอหิ ตฺวํ คามณิ สมณสฺส โคตมสฺส วาทํ อาโรเปหิ, เอวํ เต กลฺยาโณ กิตฺติสทฺโท อพฺภุคฺคจฺฉิสฺสติ, อสิพนฺธกปุตฺเตน คามณินา สมณสฺส โคตมสฺส เอวํมหิทฺธิกสฺส เอวํมหานุภาวสฺส วาโท อาโรปิโต”ติ.}

\prose{กถํ ปนาหํ ภนฺเต สมณสฺส โคตมสฺส เอวํมหิทฺธิกสฺส เอวํมหานุภาวสฺส วาทํ อาโรเปสฺสามีติ.\csromanspacer{} เอหิ ตฺวํ คามณิ เยน สมโณ โคตโม เตนุปสงฺกมิ, อุปสงฺกมิตฺวา สมณํ โคตมํ}

\gambhira{สฬายตนวคฺคสํยุตฺตปาฬิ}
\prose{\csromanfolio{508}เอวํ วเทหิ “นนุ ภนฺเต ภควา อเนกปริยาเยน กุลานํ อนุทฺทยํ\footnote{อนุทยํ (สฺยา, กํ, อิ, ก)} วณฺเณติ, อนุรกฺขํ วณฺเณติ, อนุกมฺปํ วณฺเณตี”ติ.\csromanspacer{} สเจ โข คามณิ สมโณ โคตโม เอวํ ปุฏฺโฐ เอวํ พฺยากโรติ “เอวํ คามณิ ตถาคโต อเนกปริยาเยน กุลานํ อนุทฺทยํ วณฺเณติ, อนุรกฺขํ วณฺเณติ, อนุกมฺปํ วณฺเณตี”ติ.\csromanspacer{} ตเมนํ ตฺวํ เอวํ วเทยฺยาสิ “อถ กิญฺจรหิ ภนฺเต ภควา ทุพฺภิกฺเข ทฺวีหิติเก เสตฏฺฐิเก สลากาวุตฺเต มหตา ภิกฺขุสํเฆน สทฺธิํ จาริกํ จรติ, อุจฺเฉทาย ภควา กุลานํ ปฏิปนฺโน, อนยาย ภควา กุลานํ ปฏิปนฺโน, อุปฆาตาย ภควา กุลานํ ปฏิปนฺโน”ติ.\csromanspacer{} อิมํ โข เต คามณิ สมโณ โคตโม อุภโตโกฏิกํ ปญฺหํ ปุฏฺโฐ เนว สกฺขติ\footnote{สกฺขิติ (สี) ม 2. 55 ปิฏฺเฐปิ.} อุคฺคิลิตุํ, เนว สกฺขติ โอคิลิตุนฺติ.\csromanspacer{} เอวํ ภนฺเตติ โข อสิพนฺธกปุตฺโต คามณิ นิคณฺฐสฺส นาฏปุตฺตสฺส ปฏิสฺสุตฺวา อุฏฺฐายาสนา นิคณฺฐํ นาฏปุตฺตํ อภิวาเทตฺวา ปทกฺขิณํ กตฺวา เยน ภควา เตนุปสงฺกมิ, อุปสงฺกมิตฺวา ภควนฺตํ อภิวาเทตฺวา เอกมนฺตํ นิสีทิ, เอกมนฺตํ นิสินฺโน โข อสิพนฺธกปุตฺโต คามณิ ภควนฺตํ เอตทโวจ\csromandash{}}

\prose{นนุ ภนฺเต ภควา อเนกปริยาเยน กุลานํ อนุทฺทยํ วณฺเณติ, อนุรกฺขํ วณฺเณติ, อนุกมฺปํ วณฺเณตีติ.\csromanspacer{} เอวํ คามณิ ตถาคโต อเนกปริยาเยน กุลานํ อนุทฺทยํ วณฺเณติ, อนุรกฺขํ วณฺเณติ, อนุกมฺปํ วณฺเณตีติ.\csromanspacer{} อถ กิญฺจรหิ ภนฺเต ภควา ทุพฺภิกฺเข ทฺวีหิติเก เสตฏฺฐิเก สลากาวุตฺเต มหตา ภิกฺขุสํเฆน สทฺธิํ จาริกํ จรติ, อุจฺเฉทาย ภควา กุลานํ ปฏิปนฺโน, อนยาย ภควา กุลานํ ปฏิปนฺโน, อุปฆาตาย ภควา กุลานํ ปฏิปนฺโนติ.\csromanspacer{} อิโต โส คามณิ เอกนวุติกปฺเป\footnote{เอกนวุโต กปฺโป (สฺยา, กํ), เอกนวุติกปฺโป (อิ)} ยมหํ อนุสฺสารามิ นาภิชานามิ กิญฺจิ กุลํ ปกฺกภิกฺขานุปฺปทานมตฺเตน อุปหตปุพฺพํ.\csromanspacer{} อถ โข ยานิ ตานิ กุลานิ อฑฺฒานิ มหทฺธนานิ มหาโภคานิ ปหูตชาตรูปรชตานิ ปหูตวิตฺตูปกรณานิ ปหูตธนธญฺญานิ, สพฺพานิ ตานิ ทานสมฺภูตานิ เจว สจฺจสมฺภูตานิ จ สามญฺญสมฺภูตานิ จ.\csromanspacer{} อฏฺฐ โข คามณิ เหตู อฏฺฐ ปจฺจยา กุลานํ อุปฆาตาย, ราชโต วา กุลานิ อุปฆาตํ คจฺฉนฺติ, โจรโต วา กุลานิ อุปฆาตํ คจฺฉนฺติ, อคฺคิโต วา กุลานิ อุปฆาตํ}

\vspace{\tipitakaparskip}
\csromancenter{คามณิสํยุตฺต คจฺฉนฺติ, อุทกโต วา กุลานิ อุปฆาตํ คจฺฉนฺติ, นิหิตํ วา ฐานา วิคจฺฉติ\footnote{นิหิตํ วา นาธิคจฺฉนฺติ (สี, อิ)}, ทุปฺปยุตฺตา วา กมฺมนฺตา วิปชฺชนฺติ, กุเล วา กุลงฺคาโรติ\footnote{กุลานํ วา กุลงฺคาโร (สี)} อุปฺปชฺชติ โย เต โภเค วิกิรติ วิธมติ วิทฺธํเสติ, อนิจฺจตาเยว อฏฺฐมีติ, อิเม โข คามณิ อฏฺฐ เหตู อฏฺฐ ปจฺจยา กุลานํ อุปฆาตาย.\csromanspacer{} อิเมสุ โข คามณิ อฏฺฐสุ เหตูสุ อฏฺฐสุ ปจฺจเยสุ สํวิชฺชมาเนสุ โย มํ เอวํ วเทยฺย “อุจฺเฉทาย ภควา กุลานํ ปฏิปนฺโน, อนยาย ภควา กุลานํ ปฏิปนฺโน, อุปฆาตาย ภควา กุลานํ ปฏิปนฺโน”ติ.\csromanspacer{} ตํ คามณิ วาจํ อปฺปหาย ตํ จิตฺตํ อปฺปหาย ตํ ทิฏฺฐิํ อปฺปฏินิสฺสชฺชิตฺวา ยถาภตํ นิกฺขิตฺโต, เอวํ นิรเยติ.\csromanspacer{} เอวํ วุตฺเต อสิพนฺธกปุตฺโต คามณิ ภควนฺตํ เอตทโวจ “อภิกฺกนฺตํ ภนฺเต, อภิกฺกนฺตํ ภนฺเต ฯเปฯ อุปาสกํ มํ ภควา ธาเรตุ อชฺชตคฺเค ปาณุเปตํ สรณํ คตนฺ”ติ.\csromanspacer{}\csromanspacer{}นวมํ.}
\csromanheader{1}{10. มณิจูฬกสุตฺต}
\csromanmatikaanchor{mtk.281}
\csromantocmarkref{section}{10. มณิจูฬกสุตฺต}{mtk.281}
\bookmark[dest={mtk.281},level=1]{10. มณิจูฬกสุตฺต}
\proseitem{362}{\csromanfolio{509}เอกํ สมยํ ภควา ราชคเห วิหรติ เวฬุวเน กลนฺทกนิวาเป.\csromanspacer{} เตน โข ปน สมเยน ราชนฺเตปุเร ราชปริสาย สนฺนิสินฺนานํ สนฺนิปติตานํ อยมนฺตรากถา อุทปาทิ “กปฺปติ สมณานํ สกฺยปุตฺติยานํ ชาตรูปรชตํ, สาทิยนฺติ สมณา สกฺยปุตฺติยา ชาตรูปรชตํ, ปฏิคฺคณฺหนฺติ สมณา สกฺยปุตฺติยา ชาตรูปรชตนฺ”ติ.}

\prose{เตน โข ปน สมเยน มณิจูฬโก คามณิ ตสฺสํ ปริสายํ นิสินฺโน โหติ.\csromanspacer{} อถ โข มณิจูฬโก คามณิ ตํ ปริสํ เอตทโวจ “มา อยฺโย\footnote{อยฺยา (สี, อิ)} เอวํ อวจุตฺถ, น กปฺปติ สมณานํ สกฺยปุตฺติยานํ ชาตรูปรชตํ, น สาทิยนฺติ สมณา สกฺยปุตฺติยา ชาตรูปรชตํ, นปฺปฏิคฺคณฺหนฺติ สมณา สกฺยปุตฺติยา ชาตรูปรชตํ, นิกฺขิตฺตมณิสุวณฺณา สมณา สกฺยปุตฺติยา, อเปตชาตรูปรชตา”ติ.\csromanspacer{} อสกฺขิ โข มณิจูฬโก คามณิ ตํ ปริสํ สญฺญาเปตุํ.\csromanspacer{} อถ โข มณิจูฬโก คามณิ เยน ภควา เตนุปสงฺกมิ, อุปสงฺกมิตฺวา ภควนฺตํ อภิวาเทตฺวา เอกมนฺตํ นิสีทิ, เอกมนฺตํ นิสินฺโน โข มณิจูฬโก คามณิ ภควนฺตํ เอตทโวจ “อิธ ภนฺเต ราชนฺเตปุเร ราชปริสาย สนฺนิสินฺนานํ สนฺนิปติตานํ อยมนฺตรากถา อุทปาทิ ‘กปฺปติ สมณานํ สกฺยปุตฺติยานํ ชาตรูปรชตํ, สาทิยนฺติ สมณา}

\gambhira{สฬายตนวคฺคสํยุตฺตปาฬิ}
\prose{\csromanfolio{510}สกฺยปุตฺติยา ชาตรูปรชตํ, ปฏิคฺคณฺหนฺติ สมณา สกฺยปุตฺติยา ชาตรูปรชตนฺ’ติ.\csromanspacer{} เอวํ วุตฺเต อหํ ภนฺเต ตํ ปริสํ เอตทโวจํ ‘มา อยฺโย เอวํ อวจุตฺถ, น กปฺปติ สมณานํ สกฺยปุตฺติยานํ ชาตรูปรชตํ, น สาทิยนฺติ สมณา สกฺยปุตฺติยา ชาตรูปรชตํ, นปฺปฏิคฺคณฺหนฺติ สมณา สกฺยปุตฺติยา ชาตรูปรชตํ, นิกฺขิตฺตมณิสุวณฺณา สมณา สกฺยปุตฺติยา, อเปตชาตรูปรชตา’ติ.\csromanspacer{} อสกฺขิํ ขฺวาหํ ภนฺเต ตํ ปริสํ สญฺญาเปตุํ.\csromanspacer{} กจฺจาหํ ภนฺเต เอวํ พฺยากรมาโน วุตฺตวาที เจว ภควโต โหมิ, น จ ภควนฺตํ อภูเตน อพฺภาจิกฺขามิ, ธมฺมสฺส จานุธมฺมํ พฺยากโรมิ, น จ โกจิ สหธมฺมิโก วาทานุวาโท คารยฺหํ ฐานํ อาคจฺฉตี”ติ.}

\prose{ตคฺฆ ตฺวํ คามณิ เอวํ พฺยากรมาโน วุตฺตวาที เจว เม โหสิ, น จ มํ อภูเตน อพฺภจิกฺขสิ, ธมฺมสฺส จานุธมฺมํ พฺยากโรสิ, น จ โกจิ สหธมฺมิโก วาทานุวาโท คารยฺหํ ฐานํ อาคจฺฉติ.\csromanspacer{} น หิ คามณิ กปฺปติ สมณานํ สกฺยปุตฺติยานํ ชาตรูปรชตํ, น สาทิยนฺติ สมณา สกฺยปุตฺติยา ชาตรูปรชตํ, นปฺปฏิคฺคณฺหนฺติ สมณา สกฺยปุตฺติยา ชาตรูปรชตํ, นิกฺขิตฺตมณิสุวณฺณา สมณา สกฺยปุตฺติยา, อเปตชาตรูปรชตา.\csromanspacer{} ยสฺส โข คามณิ ชาตรูปรชตํ กปฺปติ, ปญฺจปิ ตสฺส กามคุณา กปฺปนฺติ, ยสฺส ปญฺจ กามคุณา กปฺปนฺติ ( )\footnote{(ตสฺสปิ ชาตรูปรชตํ กปฺปติ,) (สฺยา, กํ)}, เอกํเสเนตํ คามณิ ธาเรยฺยาสิ “อสฺสมณธมฺโม อสกฺยาปุตฺติยธมฺโม”ติ, อปิ จาหํ คามณิ เอวํ วทามิ “ติณํ ติณตฺถิเกน ปริเยสิตพฺพํ, ทารุ ทารุตฺถิเกน ปริเยสิตพฺพํ, สกฏํ สกฏตฺถิเกน ปริเยสิตพฺพํ, ปุริโส ปุริสตฺถิเกน ปริเยสตพฺโพ”\footnote{ปริเยสิตพฺโพ”ติ (?)} นเตฺววาหํ คามณิ “เกนจิ ปริยาเยน ชาตรูปรชตํ สาทิตพฺพํ ปริเยสิตพฺพนฺ”ติ วทามีติ.\csromanspacer{}\csromanspacer{}ทสมํ.}

\csromanheader{1}{11. ภทฺรกสุตฺต}
\csromanmatikaanchor{mtk.282}
\csromantocmarkref{section}{11. ภทฺรกสุตฺต}{mtk.282}
\bookmark[dest={mtk.282},level=1]{11. ภทฺรกสุตฺต}
\proseitem{363}{เอกํ สมยํ ภควา มลฺเลสุ วิหรติ อุรุเวลกปฺปํ นาม มลฺลานํ นิคโม.\csromanspacer{} อถ โข ภทฺรโก คามณิ เยน ภควา เตนุปสงฺกมิ, อุปสงฺกมิตฺวา ภควนฺตํ อภิวาเทตฺวา เอกมนฺตํ นิสีทิ, เอกมนฺตํ นิสินฺโน โข ภทฺรโก คามณิ ภควนฺตํ เอตทโวจ “สาธุ เม ภนฺเต ภควา ทุกฺขสฺส สมุทยญฺจ อตฺถงฺคมญฺจ เทเสตู”ติ.\csromanspacer{} อหญฺเจ\footnote{อหญฺจ (สฺยา, กํ, ก)} เต คามณิ อตีตมทฺธานํ}

\vspace{\tipitakaparskip}
\csromancenter{คามณิสํยุตฺต อารพฺภ ทุกฺขสฺส สมุทยญฺจ อตฺถงฺคมญฺจ เทเสยฺยํ “เอวํ อโหสิ อตีตมทฺธานนฺ”ติ, ตตฺร เต สิยา กงฺขา สิยา วิมติ.\csromanspacer{} อหญฺเจ\footnote{อหญฺจ (สฺยา, กํ, ก)} เต คามณิ อนาคตมทฺธานํ อารพฺภ ทุกฺขสฺส สมุทยญฺจ อตฺถงฺคมญฺจ เทเสยฺยํ “เอวํ ภวิสฺสติ อนาคตมทฺธานนฺ”ติ, ตตฺราปิ เต สิยา กงฺขา สิยา วิมติ.\csromanspacer{} อปิ จาหํ คามณิ อิเธว นิสินฺโน เอตฺเถว เต นิสินฺนสฺส ทุกฺขสฺส สมุทยญฺจ อตฺถงฺคมญฺจ เทเสสฺสามิ, ตํ สุณาหิ สาธุกํ มนสิ กโรหิ, ภาสิสฺสามีติ.\csromanspacer{} เอวํ ภนฺเตติ โข ภทฺรโก คามณิ ภควโต ปจฺจสฺโสสิ.\csromanspacer{} ภควา เอตทโวจ\csromandash{}}
\prose{\csromanfolio{511}ตํ กิํ มญฺญสิ คามณิ, อตฺถิ เต อุรุเวลกปฺเป มนุสฺสา, เยสํ เต วเธน วา พนฺเธน วา ชานิยา วา ครหาย วา อุปฺปชฺเชยฺยุํ โสกปริเทว ทุกฺขโทมนสฺสุปายาสาติ.\csromanspacer{} อตฺถิ เม ภนฺเต อุรุเวลกปฺเป มนุสฺสา, เยสํ เม วเธน วา พนฺเธน วา ชานิยา วา ครหาย วา อุปฺปชฺเชยฺยุํ โสกปริเทวทุกฺขโทมนุสฺสุปายาสาติ.\csromanspacer{} อตฺถิ ปน เต คามณิ อุรุเวลกปฺเป มนุสฺสา, เยสํ เต วเธน วา พนฺเธน วา ชานิยา วา ครหาย วา นุปฺปชฺเชยฺยุํ โสกปริเทวทุกฺขมนสฺสุปายาสาติ.\csromanspacer{} อตฺถิ เม ภนฺเต อุรุเวลกปฺเป มนุสฺสา, เยสํ เม วเธน วา พนฺเธน วา ชานิยา วา ครหาย วา นุปชฺเชยฺยุํ โสกปริเทวทุกฺขโทมนุสฺสุปายาสาติ.\csromanspacer{} โก นุ โข คามณิ เหตุ โก ปจฺจโย, เยน เต เอกจฺจานํ อุรุเวลกปฺปิยานํ มนุสฺสานํ วเธน วา พนฺเธน วา ชานิยา วา ครหาย วา อุปฺปชฺเชยฺยุํ โสกปริเทวทุกฺขมนุสฺสุปายาสาติ.\csromanspacer{} เยสํ เม ภนฺเต อุรุเวลกปฺปิยานํ มนุสฺสานํ วเธน วา พนฺเธน วา ชานิยา วา ครหาย วา อุปฺปชฺเชยฺยุํ โสกปริเทวทุกฺขโทมนสฺสุปายาสา, อตฺถิ เม เตสุ ฉนฺทราโค.\csromanspacer{} เยสํ ปน ภนฺเต อุรุเวลกปฺปิยานํ มนุสฺสานํ วเธน วา พนฺเธน วา ชานิยา วา ครหาย วา นุปฺปชฺเชยฺยุํ โสกปริเทวทุกฺขโทมนสฺสุปายาสา, นตฺถิ เม เตสุ ฉนฺทราโคติ.\csromanspacer{} อิมินา ตฺวํ คามณิ ธมฺเมน ทิฏฺเฐน วิทิเตน อกาลิเกน ปตฺเตน ปริโยคาเฬฺหน อตีตานาคเต นยํ เนหิ, ยํ โข กิญฺจิ อตีตมทฺธานํ ทุกฺขํ อุปฺปชฺชมานํ อุปฺปชฺชิ\footnote{อุปฺปชฺชติ (สพฺพตฺถ)}, สพฺพํ ตํ ฉนฺทมูลกํ ฉนฺทนิทานํ, ฉนฺโท หิ มูลํ ทุกฺขสฺส.\csromanspacer{} ยมฺปิ หิ กิญฺจิ อนาคตมทฺธานํ ทุกฺขํ อุปฺปชฺชมานํ อุปฺปชฺชิสฺสติ, สพฺพํ ตํ ฉนฺทมูลกํ ฉนฺทนิทานํ, ฉนฺโท หิ มูลํ}

\gambhira{สฬายตนวคฺคสํยุตฺตปาฬิ}
\prose{\csromanfolio{512}ทุกฺขสฺสาติ.\csromanspacer{} อจฺฉริยํ ภนฺเต, อพฺภุตํ ภนฺเต, ยาว สุภาสิตํ จิทํ\footnote{สุภาสิตมิทํ (สี, อิ) 2-2. “ยงฺกิญฺจิ อตีตมทฺธานํ ทุกฺขํ อุปฺปชฺชมานํ อุปฺปชฺชติ, สพฺพนฺตํ ฉนฺทมูลกํ ฉนฺทนิทานํ, ฉนฺโท หิ มูลํ ทุกฺขสฺสาติ, ยงฺกิญฺจิ อนาคตมทฺธานํ ทุกฺขํ อุปฺปชฺชมานํ อุปฺปชฺชิสฺสติ, สพฺพนฺติ ฉนฺทมูลกํ ฉนฺทนิทานํ, ฉนฺโท หิ มูลํ ทุกฺขสฺสา”ติ. (สฺยา, กํ)} ภนฺเต ภควตา2 “ยํ กิญฺจิ ทุกฺขํ อุปฺปชฺชมานํ อุปฺปชฺชติ, สพฺพํ ตํ ฉนฺทมูลกํ ฉนฺทนิทานํ, ฉนฺโท หิ มูลํ ทุกฺขสฺสา”ติ2.\csromanspacer{} อตฺถิ เม ภนฺเต จิรวาสี นาม กุมาโร พหิ อาวสเถ\footnote{อาเวสเน (?)} ปฏิวสติ.\csromanspacer{} โส ขฺวาหํ ภนฺเต กาลสฺเสว วุฏฺฐาย ปุริสํ อุยฺโยเชมิ\footnote{อุยฺโยเชสิํ (ก)} “คจฺฉ ภเณ จิรวาสิํ กุมารํ ชานาหี”ติ.\csromanspacer{} ยาวกีวญฺจ ภนฺเต โส ปุริโส นาคจฺฉติ, ตสฺส เม โหเตว อญฺญถตฺตํ “มา เหว จิรวาสิสฺส กุมารสฺส กิญฺจิ อาพาธยิตฺถา”ติ\footnote{อาพาธเยถาติ (สฺยา, กํ, อิ, ก)}.}

\prose{ตํ กิํ มญฺญสิ คามณิ, จิรวาสิสฺส กุมารสฺส วเธน วา พนฺเธน วา ชานิยา วา ครหาย วา อุปฺปชฺเชยฺยุํ โสกปริเทวทุกฺขโทมนสฺสุปายาสาติ.\csromanspacer{} จิรวาสิสฺส เม ภนฺเต กุมารสฺส วเธน วา พนฺเธน วา ชานิยา วา ครหาย วา ชีวิตสฺสปิ สิยา อญฺญถตฺถํ, กิํ ปน เม นุปฺปชฺชิสฺสนฺติ โสกปริเทวทุกฺขโทมนสฺสุปายาสาติ.\csromanspacer{} อิมินาปิ โข เอตํ คามณิ ปริยาเยน เวทิตพฺพํ “ยํ กิญฺจิ ทุกฺขํ อุปฺปชฺชมานํ อุปฺปชฺชติ, สพฺพํ ตํ ฉนฺทมูลกํ ฉนฺทนิทานํ, ฉนฺโท หิ มูลํ ทุกฺขสฺสา”ติ.}

\prose{ตํ กิํ มญฺญสิ คามณิ, ยทา เต จิรวาสิมาตา\footnote{จิรวาสิสฺส มาตา (สี, อิ)} อทิฏฺฐา อโหสิ อสฺสุตา, อโหสิ เต จิรวาสิมาตุยา ฉนฺโท วา ราโค วา เปมํ วาติ.\csromanspacer{} โน เหตํ ภนฺเต.\csromanspacer{} ทสฺสนํ วา เต คามณิ อาคมฺม สวนํ วา, เอวํ เต อโหสิ จิรวาสิมาตุยา ฉนฺโท วา ราโค วา เปมํ วาติ.\csromanspacer{} เอวํ ภนฺเต.}

\prose{ตํ กิํ มญฺญสิ คามณิ, จิรวาสิมาตุยา เต วเธน วา พนฺเธน วา ชานิยา วา ครหาย วา อุปฺปชฺเชยฺยุํ โสกปริเทวทุกฺขโทมนสฺสุปายาสาติ.\csromanspacer{} จิรวาสิมาตุยา เม ภนฺเต วา พนฺเธน วา ชานิยา วา ครหาย วา ชีวิตสฺสปิ สิยา อญฺญถตฺตํ, กิํ ปน เม นุปฺปชฺชิสฺสนฺติ โสกปริเทวทุกฺขโทมนุสฺสุปายาสาติ.\csromanspacer{} อิมินาปิ โข เอตํ คามณิ ปริยาเยน เวทิตพฺพํ “ยํ กิญฺจิ ทุกฺขํ อุปฺปชฺชมานํ อุปฺปชฺชติ, สพฺพํ ตํ ฉนฺทมูลกํ ฉนฺทนิทานํ, ฉนฺโท หิ มูลํ ทุกฺขสฺสา”ติ.\csromanspacer{}\csromanspacer{}เอกาทสมํ.}

\csromanheader{1}{8. คามณิสํยุตฺต \textbf{12. ราสิยสุตฺต}}
\csromanmatikaanchor{mtk.283}
\csromantocmarkref{section}{12. ราสิยสุตฺต}{mtk.283}
\bookmark[dest={mtk.283},level=1]{12. ราสิยสุตฺต}
\proseitem{364}{\csromanfolio{513}อถ โข ราสิโย คามณิ เยน ภควา เตนุปสงฺกมิ, อุปสงฺกมิตฺวา ภควนฺตํ อภิวาเทตฺวา เอกมนฺตํ นิสีทิ, เอกมนฺตํ นิสินฺโน โข ราสิโย คามณิ ภควนฺตํ เอตทโวจ “สุตํ เมตํ ภนฺเต ‘สมโณ โคตโม สพฺพํ ตปํ ครหติ, สพฺพํ ตปสฺสิํ ลูขชีวิํ เอกํเสน อุปวทติ อุปกฺโกสตี’ติ\footnote{อุปกฺโกสติ อุปวทตีติ (ที 1. 153 ปิฏฺเฐ)}.\csromanspacer{} เย เต ภนฺเต เอวมาหํสุ ‘สมโณ โคตโม สพฺพํ ตปํ ครหติ, สพฺพํ ตปสฺสิํ ลูขชีวิํ เอกํเสน อุปวทติ อุปกฺโกสตี’ติ.\csromanspacer{} กจฺจิ เต ภนฺเต ภควโต วุตฺตวาทิโน, น จ ภควนฺตํ อภูเตน อพฺภาจิกฺขนฺติ, ธมฺมสฺส จานุธมฺมํ พฺยากโรนฺติ, น จ โกจิ สหธมฺมิโก วาทานุวาโท คารยฺหํ ฐานํ อาคจฺฉตี”ติ.\csromanspacer{} เย เต คามณิ เอวมาหํสุ “สมโณ โคตโม สพฺพํ ตปํ ครหติ, สพฺพํ ตปสฺสิํ ลูขชีวิํ เอกํเสน อุปวทติ อุปกฺโกสตี”ติ, น เม เต วุตฺตวาทิโน, อพฺภาจิกฺขนฺติ จ ปน มํ เต อสตา ตุจฺฉา อภูเตน.}

\prose{เทฺวเม คามณิ อนฺตา ปพฺพชิเตน น เสวิตพฺพา.\csromanspacer{} โย จายํ กาเมสุ กามสุขลฺลิกานุโยโค หีโน คมฺโม โปถุชฺชนิโก อนริโย อนตฺถสํหิโต, โย จายํ อตฺตกิลมถานุโยโค ทุกฺโข อนริโย อนตฺถสํหิโต.\csromanspacer{} เอเต เต คามณิ อุโภ อนฺเต อนุปคมฺม มชฺฌิมา ปฏิปทา ตถาคเตน อภิสมฺพุทฺธา จกฺขุกรณี ญาณกรณี อุปสมาย อภิญฺญาย สมฺโพธาย นิพฺพานาย สํวตฺตติ.\csromanspacer{} กตมา จ สา คามณิ มชฺฌิมา ปฏิปทา ตถาคเตน อภิสมฺพุทฺธา จกฺขุกรณี ญาณกรณี อุปสมาย อภิญฺญาย สมฺโพธาย นิพฺพานาย สํวตฺตติ, อยเมว อริโย อฏฺฐงฺคิโก มคฺโค.\csromanspacer{} เสยฺยถิทํ, สมฺมาทิฏฺฐิ ฯเปฯ สมฺมาสมาธิ, อยํ โข สา คามณิ มชฺฌิมา ปฏิปทา ตถาคเตน อภิสมฺพุทฺธา จกฺขุกรณี ญาณกรณี อุปสมาย อภิญฺญาย สมฺโพธาย นิพฺพานาย สํวตฺตติ.}

\prose{ตโย โข เม คามณิ กามโภคิโน สนฺโต สํวิชฺชมานา โลกสฺมิํ.\csromanspacer{} กตเม ตโย.\csromanspacer{} อิธ คามณิ เอกจฺโจ กามโภคี อธมฺเมน โภเค ปริเยสติ สาหเสน, อธมฺเมน โภเค ปริเยสิตฺวา สาหเสน น อตฺตานํ สุเขติ น ปีเณติ, น สํวิภชติ น ปุญฺญานิ กโรติ.\csromanspacer{} อิธ ปน คามณิ เอกจฺโจ กามโภคี อธมฺเมน โภเค}

\gambhira{สฬายตนวคฺคสํยุตฺตปาฬิ}
\prose{\csromanfolio{514}ปริเยสติ สาหเสน, อธมฺเมน โภเค ปริเยสิตฺวา สาหเสน อตฺตานํ สุเขติ ปีเณติ, น สํวิภชติ น ปุญฺญานิ กโรติ.\csromanspacer{} อิธ ปน คามณิ เอกจฺโจ กามโภคี อธมฺเมน โภเค ปริเยสติ สาหเสน, อธมฺเมน โภเค ปริเยสิตฺวา สาหเสน อตฺตานํ สุเขติ ปีเณติ, สํวิภชติ ปุญฺญานิ กโรติ. (\footnote{คถิโต (สี)}- 2-3)}

\prose{อิธ ปน คามณิ เอกจฺโจ กามโภคี ธมฺมาธมฺเมน โภเค ปริเยสติ สาหเสนปิ อสาหเสนปิ, ธมฺมาธมฺเมน โภเค ปริเยสิตฺวา สาหเสนปิ อสาหเสนปิ น อตฺตานํ สุเขติ น ปีเณติ, น สํวิภชติ น ปุญฺญานิ กโรติ.\csromanspacer{} อิธ ปน คามณิ เอกจฺโจ กามโภคี ธมฺมาธมฺเมน โภเค ปริเยสติ สาหเสนปิ อสาหเสนปิ, ธมฺมาธมฺเมน โภเค ปริเยสิตฺวา สาหเสนปิ อสาหเสนปิ อตฺตานํ สุเขติ ปีเณติ, น สํวิภชติ น ปุญฺญานิ กโรติ.\csromanspacer{} อิธ ปน คามณิ เอกจฺโจ กามโภคี ธมฺมาธมฺเมน โภเค ปริเยสติ สาหเสนปิ อสาหเสนปิ, ธมฺมาธมฺเมน โภเค ปริเยสิตฺวา สาหเสนปิ อสาหเสนปิ อตฺตานํ สุเขติ ปีเณติ, สํวิภชติ ปุญฺญานิ กโรติ. (4-5-6)}

\prose{อิธ ปน คามณิ เอกจฺโจ กามโภคี ธมฺเมน โภเค ปริเยสติ อสาหเสน, ธมฺเมน โภเค ปริเยสิตฺวา อสาหเสน น อตฺตานํ สุเขติ น ปีเณติ, น สํวิภชติ น ปุญฺญานิ กโรติ.\csromanspacer{} อิธ ปน คามณิ เอกจฺโจ กามโภคี ธมฺเมน โภเค ปริเยสติ อสาหเสน, ธมฺเมน โภเค ปริเยสิตฺวา อสาหเสน อตฺตานํ สุเขติ ปีเณติ, น สํวิภชติ น ปุญฺญานิ กโรติ.\csromanspacer{} อิธ ปน คามณิ เอกจฺโจ กามโภคี ธมฺเมน โภเค ปริเยสติ อสาหเสน, ธมฺเมน โภเค ปริเยสิตฺวา อสาหเสน อตฺตานํ สุเขติ ปีเณติ, สํวิภชติ ปุญฺญานิ กโรติ, เต จ โภเค คธิโต1 มุจฺฉิโต อชฺโฌปนฺโน อนาทีนวทสฺสาวี อนิสฺสรณปญฺโญ ปริภุญฺชติ.\csromanspacer{} อิธ ปน คามณิ เอกจฺโจ กามโภคี ธมฺเมน โภเค ปริเยสติ อสาหเสน, ธมฺเมน โภเค ปริเยสิตฺวา อสาหเสน อตฺตานํ สุเขติ ปีเณติ, สํวิภชติ ปุญฺญานิ กโรติ, เต จ โภเค อคธิโต อมุจฺฉิโต อนชฺโฌปนฺโน อาทีนวทสฺสาวี นิสฺสรณปญฺโญ ปริภุญฺชติ. (7-8-9)}

\csromanheader{2}{8. คามณิสํยุตฺต}
\prose{\csromanfolio{515}ตตฺร คามณิ ยฺวายํ กามโภคี อธมฺเมน โภเค ปริเยสติ สาหเสน, อธมฺเมน โภเค ปริเยสิตฺวา สาหเสน น อตฺตานํ สุเขติ น ปีเณติ, น สํวิภชติ น ปุญฺญานิ กโรติ.\csromanspacer{} อยํ คามณิ กามโภคี ตีหิ ฐาเนหิ คารโยฺห.\csromanspacer{} กตเมหิ ตีหิ ฐาเนหิ คารโยฺห.\csromanspacer{} อธมฺเมน โภเค ปริเยสติ สาหเสนาติ อิมินา ปฐเมน ฐาเนน คารโยฺห, น อตฺตานํ สุเขติ น ปีเณตีติ อิมินา ทุติเยน ฐาเนน คารโยฺห, น สํวิภชติ น ปุญฺญานิ กโรตีติ อิมินา ตติเยน ฐาเนน คารโยฺห.\csromanspacer{} อยํ คามณิ กามโภคี อิเมหิ ตีหิ ฐาเนหิ คารโยฺห. (1)}

\prose{ตตฺร คามณิ ยฺวายํ กามโภคี อธมฺเมน โภเค ปริเยสติ สาหเสน, อธมฺเมน โภเค ปริเยสิตฺวา สาหเสน อตฺตานํ สุเขติ ปีเณติ, น สํวิภชติ น ปุญฺญานิ กโรติ.\csromanspacer{} อยํ คามณิ กามโภคี ทฺวีหิ ฐาเนหิ คารโยฺห, เอเกน ฐาเนน ปาสํโส.\csromanspacer{} กตเมหิ ทฺวีหิ ฐาเนหิ คารโยฺห.\csromanspacer{} อธมฺเมน โภเค ปริเยสติ สาหเสนาติ อิมินา ปฐเมน ฐาเนน คารโยฺห, น สํวิภชติ น ปุญฺญานิ กโรตีติ อิมินา ทุติเยน ฐาเนน คารโยฺห.\csromanspacer{} กตเมน เอเกน ฐาเนน ปาสํโส.\csromanspacer{} อตฺตานํ สุเขติ ปีเณตีติ อิมินา เอเกน ฐาเนน ปาสํโส.\csromanspacer{} อยํ คามณิ กามโภคี อิเมหิ ทฺวีหิ ฐาเนหิ คารโยฺห, อิมินา เอเกน ฐาเนน ปาสํโส. (2)}

\prose{ตตฺร คามณิ ยฺวายํ กามโภคี อธมฺเมน โภคี ปริเยสติ สาหเสน, อธมฺเมน โภเค ปริเยสิตฺวา สาหเสน อตฺตานํ สุเขติ ปีเณติ, สํวิภชติ ปุญฺญานิ กโรติ.\csromanspacer{} อยํ คามณิ กามโภคี เอเกน ฐาเนน คารโยฺห, ทฺวีหิ ฐาเนหิ ปาสํโส.\csromanspacer{} กตเมน เอเกน ฐาเนน คารโยฺห.\csromanspacer{} อธมฺเมน โภเค ปริเยสติ สาหเสนาติ อิมินา เอเกน ฐาเนน คารโยฺห.\csromanspacer{} กตเมหิ ทฺวีหิ ฐาเนหิ ปาสํโส.\csromanspacer{} อตฺตานํ สุเขติ ปีเณตีติ อิมินา ปฐเมน ฐาเนน ปาสํโส.\csromanspacer{} สํวิภชติ ปุญฺญานิ กโรตีติ อิมินา ทุติเยน ฐาเนน ปาสํโส.\csromanspacer{} อยํ คามณิ กามโภคี อิมินา เอเกน ฐาเนน คารโยฺห, อิเมหิ ทฺวีหิ ฐาเนหิ ปาสํโส. (3)}

\prose{ตตฺร คามณิ ยฺวายํ กามโภคี ธมฺมาธมฺเมน โภเค ปริเยสติ สาหเสนปิ อสาหเสนปิ, ธมฺมาธมฺเมน โภเค ปริเยสิตฺวา สาหเสนปิ อสาหเสนปิ น อตฺตานํ สุเขติ น ปีเณติ, น สํวิภชติ}

\gambhira{สฬายตนวคฺคสํยุตฺตปาฬิ}
\prose{\csromanfolio{516}น ปุญฺญานิ กโรติ.\csromanspacer{} อยํ คามณิ กามโภคี เอเกน ฐาเนน ปาสํโส, ตีหิ ฐาเนหิ คารโยฺห.\csromanspacer{} กตเมน เอเกน ฐาเนน ปาสํโส.\csromanspacer{} ธมฺเมน โภเค ปริเยสติ อสาหเสนาติ อิมินา เอเกน ฐาเนน ปาสํโส.\csromanspacer{} กตเมหิ ตีหิ ฐาเนหิ คารโยฺห.\csromanspacer{} อธมฺเมน โภเค ปริเยสติ สาหเสนาติ อิมินา ปฐเมน ฐาเนน คารโยฺห, น อตฺตานํ สุเขติ น ปีเณตีติ อิมินา ทุติเยน ฐาเนน คารโยฺห, น สํวิภชติ น ปุญฺญานิ กโรตีติ อิมินา ตติเยน ฐาเนน คารโยฺห.\csromanspacer{} อยํ คามณิ กามโภคี อิมินา เอเกน ฐาเนน ปาสํโส, อิเมหิ ตีหิ ฐาเนหิ คารโยฺห. (4)}

\prose{ตตฺร คามณิ ยฺวาหํ กามโภคี ธมฺมาธมฺเมน โภเค ปริเยสติ สาหเสนปิ อสาหเสนปิ, ธมฺมาธมฺเมน โภเค ปริเยสิตฺวา สาหเสนปิ อสาหเสนปิ อตฺตานํ สุเขติ ปีเณติ, น สํวิภชติ น ปุญฺญานิ กโรติ.\csromanspacer{} อยํ คามณิ กามโภคี ทฺวีหิ ฐาเนหิ ปาสํโส, ทฺวีหิ ฐาเนหิ คารโยฺห.\csromanspacer{} กตเมหิ ทฺวีหิ ฐาเนหิ ปาสํโส.\csromanspacer{} ธมฺเมน โภเค ปริเยสติ อสาหเสนาติ อิมินา ปฐเมน ฐาเนน ปาสํโส, อตฺตานํ สุเขติ ปีเณตีติ อิมินา ทุติเยน ฐาเนน ปาสํโส.\csromanspacer{} กตเมหิ ทฺวีหิ ฐาเนหิ คารโยฺห.\csromanspacer{} อธมฺเมน โภเค ปริเยสติ สาหเสนาติ อิมินา ปฐเมน ฐาเนน คารโยฺห, น สํวิภชติ น ปุญฺญานิ กโรตีติ อิมินา ทุติเยน ฐาเนน คารโยฺห.\csromanspacer{} อยํ คามณิ กามโภคี อิเมหิ ทฺวีหิ ฐาเนหิ ปาสํโส, อิเมหิ ทฺวีหิ ฐาเนหิ คารโยฺห. (5)}

\prose{ตตฺร คามณิ ยฺวายํ กามโภคี ธมฺมาธมฺเมน โภเค ปริเยสติ สาหเสนปิ อสาหเสนปิ, ธมฺมาธมฺเมน โภเค ปริเยสิตฺวา สาหเสนปิ อสาหเสนปิ อตฺตานํ สุเขติ ปีเณติ, สํวิภชติ ปุญฺญานิ กโรติ.\csromanspacer{} อยํ คามณิ กามโภคี ตีหิ ฐาเนหิ ปาสํโส, เอเกน ฐาเนน คารโยฺห.\csromanspacer{} กตเมหิ ตีหิ ฐาเนหิ ปาสํโส.\csromanspacer{} ธมฺเมน โภเค ปริเยสติ อสาหเสนาติ อิมินา ปฐเมน ฐาเนน ปาสํโส, อตฺตานํ สุเขติ ปีเณตีติ อิมินา ทุติเยน ฐาเนน ปาสํโส, สํวิภชติ ปุญฺญานิ กโรตีติ อิมินา ตติเยน ฐาเนน ปาสํโส.\csromanspacer{} กตเมน เอเกน ฐาเนน คารโยฺห.\csromanspacer{} อธมฺเมน โภเค ปริเยสติ}

\vspace{\tipitakaparskip}
\csromancenter{คามณิสํยุตฺต สาหเสนาติ อิมินา เอเกน ฐาเนน คารโยฺห.\csromanspacer{} อยํ คามณิ กามโภคี อิเมหิ ตีหิ ฐาเนหิ ปาสํโส, อิมินา เอเกน ฐาเนน คารโยฺห. (6)}
\prose{\csromanfolio{517}ตตฺร คามณิ ยฺวายํ กามโภคี ธมฺเมน โภเค ปริเยสติ อสาหเสน, ธมฺเมน โภเค ปริเยสิตฺวา อสาหเสน น อตฺตานํ สุเขติ น ปีเณติ, น สํวิภชติ น ปุญฺญานิ กโรติ.\csromanspacer{} อยํ คามณิ กามโภคี เอเกน ฐาเนน ปาสํโส, ทฺวีหิ ฐาเนหิ คารโยฺห.\csromanspacer{} กตเมน เอเกน ฐาเนน ปาสํโส.\csromanspacer{} ธมฺเมน โภเค ปริเยสติ อสาหเสนาติ อิมินา เอเกน ฐาเนน ปาสํโส.\csromanspacer{} กตเมหิ ทฺวีหิ ฐาเนหิ คารโยฺห.\csromanspacer{} น อตฺตานํ สุเขติ น ปีเณตีติ อิมินา ปฐเมน ฐาเนน คารโยฺห, น สํวิภชติ น ปุญฺญานิ กโรตีติ อิมินา ทุติเยน ฐาเนน คารโยฺห.\csromanspacer{} อยํ คามณิ กามโภคี อิมินา เอเกน ฐาเนน ปาสํโส, อิเมหิ ทฺวีหิ ฐาเนหิ คารโยฺห. (7)}

\prose{ตตฺร คามณิ ยฺวายํ กามโภคี ธมฺเมน โภเค ปริเยสติ อสาหเสน, ธมฺเมน โภเค ปริเยสิตฺวา อสาหเสน อตฺตานํ สุเขติ ปีเณติ, น สํวิภชติ น ปุญฺญานิ กโรติ.\csromanspacer{} อยํ คามณิ กามโภคี ทฺวีหิ ฐาเนหิ ปาสํโส, เอเกน ฐาเนน คารโยฺห.\csromanspacer{} กตเมหิ ทฺวีหิ ฐาเนหิ ปาสํโส.\csromanspacer{} ธมฺเมน โภเค ปริเยสติ อสาหเสนาติ อิมินา ปฐเมน ฐาเนน ปาสํโส, อตฺตานํ สุเขติ ปีเณตีติ อิมินา ทุติเยน ฐาเนน ปาสํโส.\csromanspacer{} กตเมน เอเกน ฐาเนน คารโยฺห.\csromanspacer{} น สํวิภชติ น ปุญฺญานิ กโรตีติ อิมินา เอเกน ฐาเนน คารโยฺห.\csromanspacer{} อยํ คามณิ กามโภคี อิเมหิ ทฺวีหิ ฐาเนหิ ปาสํโส, อิมินา เอเกน ฐาเนน คารโยฺห. (8)}

\prose{ตตฺร คามณิ ยฺวายํ กามโภคี ธมฺเมน โภเค ปริเยสติ อสาหเสน, ธมฺเมน โภเค ปริเยสิตฺวา อสาหเสน อตฺตานํ สุเขติ ปีเณติ, สํวิภชติ ปุญฺญานิ กโรติ, เต จ โภเค คธิโต มุจฺฉิโต อชฺโฌปนฺโน อนาทีนวทสฺสาวี อนิสฺสรณปญฺโญ ปริภุญฺชติ.\csromanspacer{} อยํ คามณิ กามโภคี ตีหิ ฐาเนหิ ปาสํโส, เอเกน ฐาเนน คารโยฺห.\csromanspacer{} กตเมหิ ตีหิ ฐาเนหิ ปาสํโส.\csromanspacer{} ธมฺเมน โภเค ปริเยสติ อสาหเสนาติ อิมินา ปฐเมน ฐาเนน ปาสํโส, อตฺตานํ สุเขติ ปีเณตีติ อิมินา ทุติเยน ฐาเนน}

\gambhira{สฬายตนวคฺคสํยุตฺตปาฬิ}
\prose{\csromanfolio{518}ปาสํโส, สํวิภชติ ปุญฺญานิ กโรตีติ อิมินา ตติเยน ฐาเนน ปาสํโส.\csromanspacer{} กตเมน เอเกน ฐาเนน คารโยฺห.\csromanspacer{} เต จ โภเค คธิโต มุจฺฉิโต อชฺโฌปนฺโน อนาทีนวทสฺสาวี อนิสฺสรณปญฺโญ ปริภุญฺชตีติ อิมินา เอเกน ฐาเนน คารโยฺห.\csromanspacer{} อยํ คามณิ กามโภคี อิเมหิ ตีหิ ฐาเนหิ ปาสํโส, อิมินา เอเกน ฐาเนน คารโยฺห. (9)}

\prose{ตตฺร คามณิ ยฺวายํ กามโภคี ธมฺเมน โภเค ปริเยสติ อสาหเสน, ธมฺเมน โภเค ปริเยสิตฺวา อสาหเสน อตฺตานํ สุเขติ ปีเณติ, สํวิภชติ ปุญฺญานิ กโรติ, เต จ โภเค อคธิโต อมุจฺฉิโต อนชฺโฌปนฺโน อาทีนวทสฺสาวี นิสฺสรณปญฺโญ ปริภุญฺชติ.\csromanspacer{} อยํ คามณิ กามโภคี จตูหิ ฐาเนหิ ปาสํโส.\csromanspacer{} กตเมหิ จตูหิ ฐาเนหิ ปาสํโส.\csromanspacer{} ธมฺเมน โภเค ปริเยสติ อสาหเสนาติ อิมินา ปฐเมน ฐาเนน ปาสํโส, อตฺตานํ สุเขติ ปีเณตีติ อิมินา ทุติเยน ฐาเนน ปาสํโส, สํวิภชติ ปุญฺญานิ กโรตีติ อิมินา ตติเยน ฐาเนน ปาสํโส, เต จ โภเค อคธิโต อมุจฺฉิโต อนชฺโฌปนฺโน อาทีนวทสฺสาวี นิสฺสรณปญฺโญ ปริภุญฺชตีติ อิมินา จตุตฺเถน ฐาเนน ปาสํโส.\csromanspacer{} อยํ คามณิ กามโภคี อิเมหิ จตูหิ ฐาเนหิ ปาสํโส. (10)}

\prose{ตโยเม คามณิ ตปสฺสิโน ลูขชีวิโน สนฺโต สํวิชฺชมานา โลกสฺมิํ.\csromanspacer{} กตเม ตโย.\csromanspacer{} อิธ คามณิ เอกจฺโจ ตปฺปสฺสี ลูขชีวี สทฺธา อคารสฺมา อนคาริยํ ปพฺพชิโต โหติ “อปฺเปว นาม กุสลํ ธมฺมํ อธิคจฺเฉยฺยํ, อปฺเปว นาม อุตฺตริ มนุสฺสธมฺมา อลมริยญาณทสฺสนวิเสสํ สจฺฉิกเรยฺยนฺ”ติ, โส อตฺตานํ อาตาเปติ ปริตาเปติ, กุสลญฺจ ธมฺมํ นาธิคจฺฉติ, อุตฺตริ จ มนุสฺสธมฺมา อลมริยญาณทสฺสนวิเสสํ น สจฺฉิกโรติ. (1)}

\prose{อิธ ปน คามณิ เอกจฺโจ ตปสฺสี ลูขชีวี สทฺธา อคารสฺมา อนคาริยํ ปพฺพชิโต โหติ “อปฺเปว นาม กุสลํ ธมฺมํ อธิคจฺเฉยฺยํ, อปฺเปว นาม อุตฺตริ มนุสฺสธมฺมา อลมริยญาณทสฺสนวิเสสํ สจฺฉิกเรยฺยนฺ”ติ, โส อตฺตานํ อาตาเปติ ปริตาเปติ, กุสลํ หิ โข ธมฺมํ อธิคจฺฉติ, อุตฺตริ มนุสฺสธมฺมา อลมริยญาณทสฺสนวิเสสํ น สจฺฉิกโรติ. (2)}

\csromanheader{2}{8. คามณิสํยุตฺต}
\prose{\csromanfolio{519}อิธ ปน คามณิ เอกจฺโจ ตปสฺสี ลูขชีวี สทฺธา อคารสฺมา อนคาริยํ ปพฺพชิโต โหติ “อปฺเปว นาม กุสลํ ธมฺมํ อธิคจฺเฉยฺยํ, อปฺเปว นาม อุตฺตริ มนุสฺสธมฺมา อลมริยญาณทสฺสนวิเสสํ สจฺฉิกเรยฺยนฺ”ติ, โส อตฺตานํ อาตาเปติ ปริตาเปติ, กุสลญฺจ ธมฺมํ อธิคจฺฉติ, อุตฺตริ จ มนุสฺสธมฺมา อลมริยญาณทสฺสนวิเสสํ สจฺฉิกโรติ. (3)}

\prose{ตตฺร คามณิ ยฺวายํ ตปสฺสี ลูขชีวี อตฺตานํ อาตาเปติ ปริตาเปติ, กุสลญฺจ ธมฺมํ นาธิคจฺฉติ, อุตฺตริ จ มนุสฺสธมฺมา อลมริยญาณทสฺสนวิเสสํ น สจฺฉิกโรติ, อยํ คามณิ ตปสฺสี ลูขชีวี ตีหิ ฐาเนหิ คารโยฺห.\csromanspacer{} กตเมหิ ตีหิ ฐาเนหิ คารโยฺห.\csromanspacer{} อตฺตานํ อาตาเปติ ปริตาเปตีติ อิมินา ปฐเมน ฐาเนน คารโยฺห, กุสลญฺจ ธมฺมํ นาธิคจฺฉตีติ อิมินา ทุติเยน ฐาเนน คารโยฺห, อุตฺตริ จ มนุสฺสธมฺมา อลมริยญาณทสฺสนวิเสสํ น สจฺฉิกโรตีติ อิมินา ตติเยน ฐาเนน คารโยฺห.\csromanspacer{} อยํ คามณิ ตปสฺสี ลูขชีวี อิเมหิ ตีหิ ฐาเนหิ คารโยฺห. (1)}

\prose{ตตฺร คามณิ ยฺวายํ ตปสฺสี ลูขชีวี อตฺตานํ อาตาเปติ ปริตาเปติ, กุสลํ หิ โข ธมฺมํ อธิคจฺฉติ, อุตฺตริ จ มนุสฺสธมฺมา อลมริยญาณทสฺสนวิเสสํ น สจฺฉิกโรติ, อยํ คามณิ ตปสฺสี ลูขชีวี ทฺวีหิ ฐาเนหิ คารโยฺห, เอเกน ฐาเนน ปาสํโส.\csromanspacer{} กตเมหิ ทฺวีหิ ฐาเนหิ คารโยฺห.\csromanspacer{} อตฺตานํ อาตาเปติ ปริตาเปตีติ อิมินา ปฐเมน ฐาเนน คารโยฺห, อุตฺตริ จ มนุสฺสธมฺมา อลมริยญาณทสฺสนวิเสสํ น สจฺฉิกโรตีติ อิมินา ทุติเยน ฐาเนน คารโยฺห.\csromanspacer{} กตเมน เอเกน ฐาเนน ปาสํโส.\csromanspacer{} กุสลํ หิ โข ธมฺมํ อธิคจฺฉตีติ อิมินา เอเกน ฐาเนน ปาสํโส.\csromanspacer{} อยํ คามณิ ตปสฺสี ลูขชีวี อิเมหิ ทฺวีหิ ฐาเนหิ คารโยฺห, อิมินา เอเกน ฐาเนน ปาสํโส. (2)}

\prose{ตตฺร คามณิ ยฺวายํ ตปสฺสี ลูขชีวี อตฺตานํ อาตาเปติ ปริตาเปติ, กุสลญฺจ ธมฺมํ อธิคจฺฉติ, อุตฺตริ จ มนุสฺสธมฺมา อลมริยญาณทสฺสนวิเสสํ สจฺฉิกโรติ.\csromanspacer{} อยํ คามณิ ตปสฺสี ลูขชีวี เอเกน ฐาเนน คารโยฺห, ทฺวีหิ ฐาเนหิ ปาสํโส.\csromanspacer{} กตเมน เอเกน ฐาเนน คารโยฺห.\csromanspacer{} อตฺตานํ อาตาเปติ ปริตาเปตีติ อิมินา เอเกน ฐาเนน คารโยฺห.\csromanspacer{} กตเมหิ ทฺวีหิ ฐาเนหิ ปาสํโส.\csromanspacer{} กุสลญฺจ ธมฺมํ อธิคจฺฉตีติ อิมินา ปฐเมน ฐาเนน ปาสํโส, อุตฺตริ จ มนุสฺสธมฺมา อลมริยญาณทสฺสนวิเสสํ สจฺฉิกโรตีติ อิมินา ทุติเยน ฐาเนน}

\gambhira{สฬายตนวคฺคสํยุตฺตปาฬิ}
\prose{\csromanfolio{520}ปาสํโส.\csromanspacer{} อยํ คามณิ ตปสฺสี ลูขชีวี อิมินา เอเกน ฐาเนน คารโยฺห, อิเมหิ ทฺวีหิ ฐาเนหิ ปาสํโส. (3)}

\prose{ติสฺโส อิมา คามณิ สนฺทิฏฺฐิกา นิชฺชรา อกาลิกา เอหิปสฺสิกา โอปเนยฺยิกา ปจฺจตฺตํ เวทิตพฺพา วิญฺญูหิ.\csromanspacer{} กตมา ติสฺโส.\csromanspacer{} ยํ รตฺโต ราคาธิกรณํ อตฺตพฺยาพาธายปิ เจเตติ, ปรพฺยาธายปิ เจเตติ, อุภยพฺยาพาธายปิ เจเตติ.\csromanspacer{} ราเค ปหีเน เนวตฺตพฺยาพาธาย เจเตติ, น ปรพฺยาพาธาย เจเตติ, น อุภยพฺยาพาธาย เจเตติ, สนฺทิฏฺฐิกา นิชฺชรา อกาลิกา เอหิปสฺสิกา โอปเนยฺยิกา ปจฺจตฺตํ เวทิตพฺพา วิญฺญูหิ.\csromanspacer{} ยํ ทุฏฺโฐ โทสาธิกรณํ อตฺตพฺยาพาธายปิ เจเตติ, ปรพฺยาพาธายปิ เจเตติ, อุภยพฺยาพาธายปิ เจเตติ.\csromanspacer{} โทเส ปหีเน เนวตฺตพฺยาพาธาย เจเตติ, น ปรพฺยาพาธาย เจเตติ, น อุภยพฺยาพาธาย เจเตติ, สนฺทิฏฺฐิกา นิชฺชรา อกาลิกา เอหิปสฺสิกา โอปเนยฺยิกา ปจฺจตฺตํ เวทิตพฺพา วิญฺญูหิ.\csromanspacer{} ยํ มูโฬฺห โมหาธิกรณํ อตฺตพฺยาพาธายปิ เจเตติ, ปรพฺยาพาธายปิ เจเตติ, อุภยพฺยาพาธายปิ เจเตติ.\csromanspacer{} โมเห ปหีเน เนวตฺตพฺยาพาธาย เจเตติ, น ปรพฺยาพาธาย เจเตติ, น อุภยพฺยาพาธาย เจเตติ, สนฺทิฏฺฐิกา นิชฺชรา อกาลิกา เอหิปสฺสิกา โอปเนยฺยิกา ปจฺจตฺตํ เวทิตพฺพา วิญฺญูหิ.\csromanspacer{} อิมินา โข คามณิ ติสฺโส สนฺทิฏฺฐิกา นิชฺชรา อกาลิกา เอหิปสฺสิกา โอปเนยฺยิกา ปจฺจตฺตํ เวทิตพฺพา วิญฺญูหีติ.}

\prose{เอวํ วุตฺเต ราสิโย คามณิ ภควนฺตํ เอตทโวจ “อภิกฺกนฺตํ ภนฺเต ฯเปฯ อุปาสกํ มํ ภควา ธาเรตุ อชฺชตคฺเค ปาณุเปตํ สรณํ คตนฺ”ติ.\csromanspacer{}\csromanspacer{}ทฺวาทสมํ.}

\csromanheader{1}{13. ปาฏลิยสุตฺต}
\csromanmatikaanchor{mtk.284}
\csromantocmarkref{section}{13. ปาฏลิยสุตฺต}{mtk.284}
\bookmark[dest={mtk.284},level=1]{13. ปาฏลิยสุตฺต}
\csromantocmarknopage{section}{(9) อสงฺขตสํยุตฺต}{mtk.285}
\bookmark[dest={mtk.285},level=1]{(9) อสงฺขตสํยุตฺต}
\proseitem{365}{เอกํ สมยํ ภควา โกลิเยสุ วิหรติ อุตฺตรํ นาม\footnote{อุตฺตรกํ นาม (สี)} โกลิยานํ นิคโม.\csromanspacer{} อถ โข ปาฏลิโย คามณิ เยน ภควา เตนุปสงฺกมิ, อุปสงฺกมิตฺวา ภควนฺตํ อภิวาเทตฺวา เอกมนฺตํ นิสีทิ, เอกมนฺตํ นิสินฺโน โข ปาฏลิโย คามณิ ภควนฺตํ เอตทโวจ “สุตํ เมตํ ภนฺเต ‘สมโณ โคตโม มายํ ชานาตี’ติ, เย เต ภนฺเต เอวมาหํสุ ‘สมโณ โคตโม มายํ ชานาตี’ติ, กจฺจิ เต ภนฺเต}

\vspace{\tipitakaparskip}
\csromancenter{คามณิสํยุตฺต ภควโต วุตฺตวาทิโน, น จ ภควนฺตํ อภูเตน อพฺภาจิกฺขนฺติ, ธมฺมสฺส จานุธมฺมํ พฺยากโรนฺติ, น จ โกจิ สหธมฺมิโก วาทานุวาโท คารยฺหํ ฐานํ อาคจฺฉติ, อนพฺภาจิกฺขิตุกามา หิ มยํ ภนฺเต ภควนฺตนฺ”ติ.\csromanspacer{} เย เต คามณิ เอวมาหํสุ “สมโณ โคตโม มายํ ชานาตี”ติ, วุตฺตวาทิโน เจว เม, เต น จ มํ อภูเตน อพฺภาจิกฺขนฺติ, ธมฺมสฺส จานุธมฺมํ พฺยากโรนฺติ, น จ โกจิ สหธมฺมิโก วาทานุวาโท คารยฺหํ ฐานํ อาคจฺฉตีติ.\csromanspacer{} สจฺจํเยว กิร โภ มยํ เตสํ สมณพฺราหฺมณานํ น สทฺทหาม “สมโณ โคตโม มายํ ชานาตี”ติ, “สมโณ ขลุ โภ โคตโม มายาวี”ติ.\csromanspacer{} โย นุ โข คามณิ เอวํ วเทติ “อหํ มายํ ชานามี”ติ, โส เอวํ วเทติ “อหํ มายาวี”ติ.\csromanspacer{} ตเถว ตํ ภควา โหติ, ตเถว ตํ สุคต โหตีติ.\csromanspacer{} เตน หิ คามณิ ตญฺเญเวตฺถ ปฏิปุจฺฉิสฺสามิ.\csromanspacer{} ยถา เต ขเมยฺย, ตถา ตํ พฺยากเรยฺยาสิ.}
\prose{\csromanfolio{521}ตํ กิํ มญฺญสิ คามณิ, ชานาสิ ตฺวํ โกลิยานํ ลมฺพจูฬเก ภเฏติ.\csromanspacer{} ชานามหํ ภนฺเต โกลิยานํ ลมฺพจูฬเก ภเฏติ.\csromanspacer{} ตํ กิํ มญฺญสิ คามณิ, กิมตฺถิยา โกลิยานํ ลมฺพจูฬกา ภฏาติ.\csromanspacer{} เย จ ภนฺเต โกลิยานํ โจรา, เต จ ปฏิเสเธตุํ, ยานิ จ โกลิยานํ ทูเตยฺยานิ, ตานิ จ วหาตุํ\footnote{ตานิ จ ปหาตุํ (สฺยา, กํ), ตานิ จ ยาตุํ (กตฺถจิ), ตานิ จาวหาตุํ (?)}, เอตทตฺถิยา ภนฺเต โกลิยานํ ลมฺพจูฬกา ภฏาติ.\csromanspacer{} ตํ กิํ มญฺญสิ คามณิ, ชานาสิ ตฺวํ โกลิยานํ ลมฺพจูฬกา ภเฏ สีลวนฺเต วา เต ทุสฺสีเล วาติ.\csromanspacer{} ชานามหํ ภนฺเต โกลิยานํ ลมฺพจูฬเก ภเฏ ทุสฺสีเล ปาปธมฺเม, เย จ โลเก ทุสฺสีลา ปาปธมฺมา, โกลิยานํ ลมฺพจูฬกา ภฏา เตสํ อญฺญตราติ.\csromanspacer{} โย นุ โข คามณิ เอวํ วเทยฺย “ปาฏลิโย คามณิ ชานาติ โกลิยานํ ลมฺพจูฬเก ภเฏ ทุสฺสีเล ปาปธมฺเม, ปาฏลิโยปิ คามณิ ทุสฺสีโล ปาปธมฺโม”ติ, สมฺมา นุ โข โส วทมาโน วเทยฺยาติ.\csromanspacer{} โน เหตํ ภนฺเต.\csromanspacer{} อญฺเญ ภนฺเต โกลิยานํ ลมฺพจูฬกา ภฏา, อญฺโญหมสฺมิ, อญฺญถาธมฺมา โกลิยานํ ลมฺพจูฬกา ภฏา, อญฺญถาธมฺโมหมสฺมีติ.\csromanspacer{} ตฺวํ หิ นาม คามณิ ลจฺฉสิ “ปาฏลิโย คามณิ ชานาติ โกลิยานํ ลมฺพจูฬเก ภเฏ ทุสฺสีเล ปาปธมฺเม, น จ ปาฏลิโย คามณิ ทุสฺสีโล ปาปธมฺโม”ติ.\csromanspacer{} กสฺมา ตถาคโต น ลจฺฉติ “ตถาคโต มายํ ชานาติ, น จ ตถาคโต มายาวี”ติ.\csromanspacer{} มายญฺจาหํ คามณิ ปชานามิ,}

\gambhira{สฬายตนวคฺคสํยุตฺตปาฬิ}
\prose{\csromanfolio{522}มายาย จ วิปากํ, ยถาปฏิปนฺโน จ มายาวี กายสฺส เภทา ปรํ มรณา อปายํ ทุคฺคติํ วินิปาตํ นิรยํ อุปปชฺชติ, ตญฺจ ปชานามิ.}

\prose{ปาณาติปาตญฺจาหํ คามณิ ปชานามิ ปาณาติปาตสฺส จ วิปากํ, ยถาปฏิปนฺโน จ ปาณาติปาตี กายสฺส เภทา ปรํ มรณา อปายํ ทุคฺคติํ วินิปาตํ นิรยํ อุปปชฺชติ, ตญฺจ ปชานามิ.\csromanspacer{} อทินฺนาทานญฺจาหํ คามณิ ปชานามิ อทินฺนาทานสฺส จ วิปากํ, ยถาปฏิปนฺโน จ อทินฺนาทายี กายสฺส เภทา ปรํ มรณา อปายํ ทุคฺคติํ วินิปาตํ นิรยํ อุปปชฺชติ, ตญฺจ ปชานามิ.\csromanspacer{} กาเมสุมิจฺฉาจารญฺจาหํ คามณิ ปชานามิ กาเมสุมิจฺฉาจารสฺส จ วิปากํ, ยถาปฏิปนฺโน จ กาเมสุมิจฺฉาจารี กายสฺส เภทา ปรํ มรณา อปายํ ทุคฺคติํ วินิปาตํ นิรยํ อุปปชฺชติ, ตญฺจ ปชานามิ.\csromanspacer{} มุสาวาทญฺจาหํ คามณิ ปชานามิ มุสาวาทสฺส จ วิปากํ, ยถาปฏิปนฺโน จ มุสาวาที กายสฺส เภทา ปรํ มรณา อปายํ ทุคฺคติํ วินิปาตํ นิรยํ อุปปชฺชติ, ตญฺจ ปชานามิ.\csromanspacer{} ปิสุณวาจญฺจาหํ คามณิ ปชานามิ ปิสุณวาจาย จ วิปากํ, ยถาปฏิปนฺโน จ ปิสุณวาโจ กายสฺส เภทา ปรํ มรณา อปายํ ทุคฺคติํ วินิปาตํ นิรยํ อุปปชฺชติ, ตญฺจ ปชานามิ.\csromanspacer{} ผรุสวาจญฺจาหํ คามณิ ปชานามิ ผรุสวาจาย จ วิปากํ, ยถาปฏิปนฺโน จ ผรุสวาโจ กายสฺส เภทา ปรํ มรณา อปายํ ทุคฺคติํ วินิปาตํ นิรยํ อุปปชฺชติ, ตญฺจ ปชานามิ.\csromanspacer{} สมฺผปฺปลาปญฺจาหํ คามณิ ปชานามิ สมฺผปฺปลาปสฺส จ วิปากํ, ยถาปฏิปนฺโน จ สมฺผปฺปลาปี กายสฺส เภทา ปรํ มรณา อปายํ ทุคฺคติํ วินิปาตํ นิรยํ อุปปชฺชติ, ตญฺจ ปชานามิ.\csromanspacer{} อภิชฺฌญฺจาหํ คามณิ ปชานามิ อภิชฺฌาย จ วิปากํ, ยถาปฏิปนฺโน จ อภิชฺฌาลุ กายสฺส เภทา ปรํ มรณา อปายํ ทุคฺคติํ วินิปาตํ นิรยํ อุปปชฺชติ, ตญฺจ ปชานามิ.\csromanspacer{} พฺยาปาทปโทสญฺจาหํ คามณิ ปชานามิ พฺยาปาทปโทสสฺส จ วิปากํ, ยถาปฏิปนฺโน จ พฺยาปนฺนจิตฺโต กายสฺส เภทา ปรํ มรณา อปายํ ทุคฺคติํ วินิปาตํ นิรยํ อุปปชฺชติ, ตญฺจ ปชานามิ.\csromanspacer{} มิจฺฉาทิฏฺฐิญฺจาหํ คามณิ ปชานามิ มิจฺฉาทิฏฺฐิยา จ วิปากํ, ยถาปฏิปนฺโน จ มิจฺฉาทิฏฺฐิโก กายสฺส เภทา ปรํ มรณา อปายํ ทุคฺคติํ วินิปาตํ นิรยํ อุปปชฺชติ, ตญฺจ ปชานามิ.}

\prose{สนฺติ หิ คามณิ เอเก สมณพฺราหฺมณา เอวํวาทิโน เอวํทิฏฺฐิโน “โย โกจิ ปาณมติปาเตติ, สพฺโพ โส ทิฏฺเฐว ธมฺเม ทุกฺขํ โทมนสฺสํ ปฏิสํเวทยติ.\csromanspacer{} โย โกจิ อทินฺนํ อาทิยติ, สพฺโพ โส ทิฏฺเฐว}

\vspace{\tipitakaparskip}
\csromancenter{คามณิสํยุตฺต ธมฺเม ทุกฺขํ โทมนสฺสํ ปฏิสํเวทยติ.\csromanspacer{} โย โกจิ กาเมสุ มิจฺฉา จรติ, สพฺโพ โส ทิฏฺเฐว ธมฺเม ทุกฺขํ โทมนสฺสํ ปฏิสํเวทยติ.\csromanspacer{} โย โกจิ มุสา ภณติ, สพฺโพ โส ทิฏฺเฐว ธมฺเม ทุกฺขํ โทมนสฺสํ ปฏิสํเวทยตี”ติ.}
\prose{\csromanfolio{523}ทิสฺสติ โข ปน คามณิ อิเธกจฺโจ มาลี กุณฺฑลี สุนฺหาโต\footnote{สุนหาโต (สี, สฺยา, กํ, อิ)} สุวิลิตฺโต กปฺปิตเกสมสฺสุ อิตฺถิกาเมหิ ราชา มญฺเญ ปริจาเรนฺโต.\csromanspacer{} ตเมนํ เอวมาหํสุ “อมฺโภ อยํ ปุริโส กิํ อกาสิ, มาลี กุณฺฑลี สุนฺหาโต สุวิลิตฺโต กปฺปิตเกสมสฺสุ อิตฺถิกาเมหิ ราชา มญฺเญ ปริจาเรตี”ติ.\csromanspacer{} ตเมนํ เอวมาหํสุ “อมฺโภ อยํ ปุริโส รญฺโญ ปจฺจตฺถิกํ ปสยฺห ชีวิตา โวโรเปสิ, ตสฺส ราชา อตฺตมโน อภิหารมทาสิ, เตนายํ ปุริโส มาลี กุณฺฑลี สุนฺหาโต สุวิลิตฺโต กปฺปิตเกสมสฺสุ อิตฺถิกาเมหิ ราชา มญฺเญ ปริจาเรตี”ติ.}

\prose{ทิสฺสติ โข คามณิ อิเธกจฺโจ ทฬฺหาย รชฺชุยา ปจฺฉาพาหํ คาฬฺหพนฺธนํ พนฺธิตฺวา ขุรมุณฺฑํ กรตฺวา ขรสฺสเรน ปณเวน รถิยาย รถิยํ\footnote{รถิกาย รถิกํ (สี)} สิงฺฆาฏเกน สิงฺฆาฏกํ ปริเนตฺวา ทกฺขิเณน ทฺวาเรน นิกฺขาเมตฺวา ทกฺขิณโต นครสฺส สีสํ ฉิชฺชมาโน.\csromanspacer{} ตเมนํ เอวมาหํสุ “อมฺโภ อยํ ปุริโส กิํ อกาสิ, ทฬฺหาย รชฺชุยา ปจฺฉาพาหํ คาฬฺหพนฺธนํ พนฺธิตฺวา ขุรมุณฺฑํ กริตฺวา ขรสฺสเรน ปณเวน รถิยาย รถิยํ สิงฺฆาฏเกน สิงฺฆาฏกํ ปริเนตฺวา ทกฺขิเณน ทฺวาเรน นิกฺขาเมตฺวา ทกฺขิณโต นครสฺส สีสํ ฉินฺทตี”ติ\footnote{ฉิชฺชตีติ (กตฺถจิ)}.\csromanspacer{} ตเมนํ เอวมาหํสุ “อมฺโภ อยํ ปุริโส ราชเวรี อิตฺถิํ วา ปุริสํ วา ชีวิตา โวโรเปสิ, เตน นํ ราชาโน คเหตฺวา เอวรูปํ กมฺมการณํ กาเรนฺตี”ติ.}

\prose{ตํ กิํ มญฺญสิ คามณิ, อปิ นุ เต เอวรูปํ ทิฏฺฐํ วา สุตํ วาติ.\csromanspacer{} ทิฏฺฐญฺจ โน ภนฺเต สุตญฺจ สุยฺยิสฺสติ จาติ.\csromanspacer{} ตตฺร คามณิ เย เต สมณพฺราหฺมณา เอวํวาทิโน เอวํทิฏฺฐิโน “โย โกจิ ปาณมติปาเตติ, สพฺโพ โส ทิฏฺเฐว ธมฺเม ทุกฺขํ โทมนสฺสํ ปฏิสํเวทยตี”ติ, สจฺจํ วา เต อาหํสุ มุสา วาติ.\csromanspacer{} มุสา ภนฺเต.\csromanspacer{} เย ปน เต ตุจฺฉํ มุสา วิลปนฺติ, สีลวนฺโต วา เต ทุสฺสีลา วาติ.\csromanspacer{} ทุสฺสีลา ภนฺเต.\csromanspacer{} เย ปน เต}

\gambhira{สฬายตนวคฺคสํยุตฺตปาฬิ}
\prose{\csromanfolio{524}ทุสฺสีลา ปาปธมฺมา, มิจฺฉาปฏิปนฺนา วา เต สมฺมาปฏิปนฺนา วาติ.\csromanspacer{} มิจฺฉาปฏิปนฺนา ภนฺเต.\csromanspacer{} เย ปน เต มิจฺฉาปฏิปนฺนา, มิจฺฉาทิฏฺฐิกา วา เต สมฺมาทิฏฺฐิกา วาติ.\csromanspacer{} มิจฺฉาทิฏฺฐิกา ภนฺเต.\csromanspacer{} เย ปน เต มิจฺฉาทิฏฺฐิกา, กลฺลํ นุ เตสุ ปสีทิตุนฺติ.\csromanspacer{} โน เหตํ ภนฺเต. (1)}

\prose{ทิสฺสติ โข ปน คามณิ อิเธกจฺโจ มาลี กุณฺฑลี ฯเปฯ อิตฺถิกาเมหิ ราชา มญฺเญ ปริจาเรนฺโต.\csromanspacer{} ตเมนํ เอวมาหํสุ “อมฺโภ อยํ ปุริโส กิํ อกาสิ, มาลี กุณฺฑลี ฯเปฯ อิตฺถิกาเมหิ ราชา มญฺเญ ปริจาเรตี”ติ.\csromanspacer{} ตเมนํ เอวมาหํสุ “อมฺโภ อยํ ปุริโส รญฺโญ ปจฺจตฺถิกสฺส ปสยฺห รตนํ อหาสิ\footnote{ปสยฺห อทินฺนํ รตนํ อาทิยิ (ก)}, ตสฺส ราชา อตฺตมโน อภิหารมทาสิ, เตนายํ ปุริโส มาลี กุณฺฑลี ฯเปฯ อิตฺถิกาเมหิ ราชา มญฺเญ ปริจาเรตี”ติ.}

\prose{ทิสฺสติ โข คามณิ อิเธกจฺโจ ทฬฺหาย รชฺชุยา ฯเปฯ ทกฺขิณโต นครสฺส สีสํ ฉิชฺชมาโน.\csromanspacer{} ตเมนํ เอวมาหํสุ “อมฺโภ อยํ ปุริโส กิํ อกาสิ, ทฬฺหาย รชฺชุยา ฯเปฯ ทกฺขิณโต นครสฺส สีสํ ฉินฺทตี”ติ.\csromanspacer{} ตเมนํ เอวมาหํสุ “อมฺโภ อยํ ปุริโส คามา วา อรญฺญา วา อทินฺนํ เถยฺยสงฺขาตํ อาทิยิ, เตน นํ ราชาโน คเหตฺวา เอวรูปํ กมฺมการณํ กาเรนฺตี”ติ.\csromanspacer{} ตํ กิํ มญฺญสิ คามณิ, อปิ นุ เต เอวรูปํ ทิฏฺฐํ วา สุตํ วาติ.\csromanspacer{} ทิฏฺฐญฺจ โน ภนฺเต สุตญฺจ สุยฺยิสฺสติ จาติ.\csromanspacer{} ตตฺร คามณิ เย เต สมณพฺราหฺมณา เอวํวาทิโน เอวํทิฏฺฐิโน “โย โกจิ อทินฺนํ อาทิยติ, สพฺโพ โส ทิฏฺเฐว ธมฺเม ทุกฺขํ โทมนสฺสํ ปฏิสํเวทยตี”ติ.\csromanspacer{} สจฺจํ วา เต อาหํสุ มุสา วาติ ฯเปฯ กลฺลํ นุ เตสุ ปสีทิตุนฺติ.\csromanspacer{} โน เหตํ ภนฺเต. (2)}

\prose{ทิสฺสติ โข ปน คามณิ อิเธกจฺโจ มาลี กุณฺฑลี ฯเปฯ อิตฺถิกาเมหิ ราชา มญฺเญ ปริจาเรนฺโต.\csromanspacer{} ตเมนํ เอวมาหํสุ “อมฺโภ อยํ ปุริโส กิํ อกาสิ, มาลี กุณฺฑลี ฯเปฯ อิตฺถิกาเมหิ ราชา มญฺเญ ปริจาเรตี”ติ.\csromanspacer{} ตเมนํ เอวมาหํสุ “อมฺโภ อยํ ปุริโส รญฺโญ ปจฺจตฺถิกสฺส ทาเรสุ จาริตฺตํ อาปชฺชิ, ตสฺส ราชา อตฺตมโน อภิหารมทาสิ, เตนายํ ปุริโส มาลี กุณฺฑลี ฯเปฯ อิตฺถิกาเมหิ ราชา มญฺเญ ปริจาเรตี”ติ.}

\csromanheader{2}{8. คามณิสํยุตฺต}
\prose{\csromanfolio{525}ทิสฺสติ โข คามณิ อิเธกจฺโจ ทฬฺหาย รชฺชุยา ฯเปฯ ทกฺขิณโต นครสฺส สีสํ ฉิชฺชมาโน.\csromanspacer{} ตเมนํ เอวมาหํสุ “อมฺโภ อยํ ปุริโส กิํ อกาสิ, ทฬฺหาย รชฺชุยา ฯเปฯ ทกฺขิณโต นครสฺส สีสํ ฉินฺทตี”ติ.\csromanspacer{} ตเมนํ เอวมาหํสุ “อมฺโภ อยํ ปุริโส กุลิตฺถีสุ กุลกุมารีสุ จาริตฺตํ อาปชฺชิ, เตน นํ ราชาโน คเหตฺวา เอวรูปํ กมฺมการณํ กาเรนฺตี”ติ.\csromanspacer{} ตํ กิํ มญฺญสิ คามณิ, อปิ นุ เต เอวรูปํ ทิฏฺฐํ วา สุตํ วาติ.\csromanspacer{} ทิฏฺฐญฺจ โน ภนฺเต สุตญฺจ สุยฺยิสฺสติ จาติ.\csromanspacer{} ตตฺร คามณิ เย เต สมณพฺราหฺมณา เอวํวาทิโน เอวํทิฏฺฐิโน “โย โกจิ กาเมสุ มิจฺฉา จรติ, สพฺโพ โส ทิฏฺเฐว ธมฺเม ทุกฺขํ โทมนสฺสํ ปฏิสํเวทยตี”ติ, สจฺจํ วา เต อาหํสุ มุสา วาติ ฯเปฯ กลฺลํ นุ เตสุ ปสีทิตุนฺติ.\csromanspacer{} โน เหตํ ภนฺเต. (3)}

\prose{ทิสฺสติ โข ปน คามณิ อิเธกจฺโจ มาลี กุณฺฑลี สุนฺหาโต สุวิลิตฺโต กปฺปิตเกสมสฺสุ อิตฺถิกาเมหิ ราชา มญฺเญ ปริจาเรนฺโต.\csromanspacer{} ตเมนํ เอวมาหํสุ “อมฺโภ อยํ ปุริโส กิํ อกาสิ, มาลี กุณฺฑลี สุนฺหาโต สุวิลิตฺโต กปฺปิตเกสมสฺสุ อิตฺถิกาเมหิ ราชา มญฺเญ ปริจาเรตี”ติ.\csromanspacer{} ตเมนํ เอวมาหํสุ “อมฺโภ อยํ ปุริโส ราชานํ มุสาวาเทน หาเสสิ, ตสฺส ราชา อตฺตมโน อภิหารมทาสิ, เตนายํ ปุริโส มาลี กุณฺฑลี สุนฺหาโต สุวิลิตฺโต กปฺปิตเกสมสฺสุ อิตฺถิกาเมหิ ราชา มญฺเญ ปริจาเรตี”ติ.}

\prose{ทิสฺสติ โข คามณิ อิเธกจฺโจ ทฬฺหาย รชฺชุยา ปจฺฉาพาหํ คาฬฺหพนฺธนํ พนฺธิตฺวา ขุรมุณฺฑํ กริตฺวา ขรสฺสเรน ปณเวน รถิยาย รถิยํ สิงฺฆาฏเกน สิงฺฆาฏกํ ปริเนตฺวา ทกฺขิเณน ทฺวาเรน นิกฺขาเมตฺวา ทกฺขิณโต นครสฺส สีสํ ฉิชฺชมาโน.\csromanspacer{} ตเมนํ เอวมาหํสุ “อมฺโภ อยํ ปุริโส กิํ อกาสิ, ทฬฺหาย รชฺชุยา ปจฺฉาพาหํ คาฬฺหพนฺธนํ พนฺธิตฺวา ขุรมุณฺฑํ กริตฺวา ขรสฺสเรน ปณเวน รถิยาย รถิยํ สิงฺฆาฏเกน สิงฺฆาฏกํ ปริเนตฺวา ทกฺขิเณน ทฺวาเรน นิกฺขาเมตฺวา ทกฺขิณโต นครสฺส สีสํ ฉินฺทตี”ติ.\csromanspacer{} ตเมนํ เอวมาหํสุ “อมฺโภ อยํ ปุริโส คหปติสฺส วา คหปติปุตฺตสฺส วา มุสาวาเทน อตฺถํ ภญฺชิ, เตน นํ ราชาโน คเหตฺวา เอวรูปํ กมฺมการณํ กาเรนฺตี”ติ.\csromanspacer{} ตํ กิํ มญฺญสิ คามณิ, อปิ นุ เต เอวรูปํ ทิฏฺฐํ วา สุตํ วาติ.\csromanspacer{} ทิฏฺฐญฺจ}

\gambhira{สฬายตนวคฺคสํยุตฺตปาฬิ}
\prose{\csromanfolio{526}โน ภนฺเต สุตญฺจ สุยฺยิสฺสติ จาติ.\csromanspacer{} ตตฺร คามณิ เย เต สมณพฺราหฺมณา เอวํวาทิโน เอวํทิฏฺฐิโน “โย โกจิ มุสา ภณติ, สพฺโพ โส ทิฏฺเฐว ธมฺเม ทุกฺขํ โทมนสฺสํ ปฏิสํเวทยตี”ติ, สจฺจํ วา เต อาหํสุ มุสา วาติ.\csromanspacer{} มุสา ภนฺเต.\csromanspacer{} เย ปน เต ตุจฺฉํ มุสา วิลปนฺติ, สีลวนฺโต วา เต ทุสฺสีลา วาติ.\csromanspacer{} ทุสฺสีลา ภนฺเต.\csromanspacer{} เย ปน เต ทุสฺสีลา ปาปธมฺมา, มิจฺฉาปฏิปนฺนา วา เต สมฺมาปฏิปนฺนา วาติ.\csromanspacer{} มิจฺฉาปฏิปนฺนา ภนฺเต.\csromanspacer{} เย ปน เต มิจฺฉาปฏิปนฺนา, มิจฺฉาทิฏฺฐิกา วา เต สมฺมาทิฏฺฐิกา วาติ.\csromanspacer{} มิจฺฉาทิฏฺฐิกา ภนฺเต.\csromanspacer{} เย ปน เต มิจฺฉาทิฏฺฐิกา, กลฺลํ นุ เตสุ ปสีทิตุนฺติ.\csromanspacer{} โน เหตํ ภนฺเต. (4)}

\prose{อจฺฉิริยํ ภนฺเต, อพฺภุตํ ภนฺเต, อตฺถิ เม ภนฺเต อาวสถาคารํ ตตฺถ อตฺถิ มญฺจกานิ, อตฺถิ อาสนานิ, อตฺถิ อุทกมณิโก, อตฺถิ เตลปฺปทีโป, ตตฺถ โย สมโณ วา พฺราหฺมโณ วา วาสํ อุเปติ, เตนาหํ ยถาสตฺติ ยถาพลํ สํวิภชามิ, ภูตปุพฺพํ ภนฺเต จตฺตาโร สตฺถาโร นานาทิฏฺฐิกา นานาขนฺติกา นานารุจิกา, ตสฺมิํ อาวสถาคาเร วาสํ อุปคจฺฉุํ.}

\prose{เอโก สตฺถา เอวํวาที เอวํทิฏฺฐิ “นตฺถิ ทินฺนํ, นตฺถิ ยิฏฺฐํ, นตฺถิ หุตํ, นตฺถิ สุกตทุกฺกฏานํ กมฺมานํ ผลํ วิปาโก, นตฺถิ อยํ โลโก, นตฺถิ ปโร โลโก, นตฺถิ มาตา, นตฺถิ ปิตา, นตฺถิ สตฺตา โอปปาติกา, นตฺถิ โลเก สมณพฺราหฺมณา สมฺมคฺคตา สมฺมาปฏิปนฺนา, เย อิมญฺจ โลกํ ปรญฺจ โลกํ สยํ อภิญฺญา สจฺฉิกตฺวา ปเวเทนฺตี”ติ.}

\prose{เอโก สตฺถา เอวํวาที เอวํทิฏฺฐิ “อตฺถิ ทินฺนํ, อตฺถิ ยิฏฺฐํ, อตฺถิ หุตํ, อตฺถิ สุกตทุกฺกฏานํ กมฺมานํ ผลํ วิปาโก, อตฺถิ อยํ โลโก, อตฺถิ ปโร โลโก, อตฺถิ มาตา, อตฺถิ ปิตา, อตฺถิ สตฺตา โอปปาติกา, อตฺถิ โลเก สมณพฺราหฺมณา สมฺมคฺคตา สมฺมาปฏิปนฺนา, เย อิมญฺจ โลกํ ปรญฺจ โลกํ สยํ อภิญฺญา สจฺฉิกตฺวา ปเวเทนฺตี”ติ.}

\prose{เอโก สตฺถา เอวํวาที เอวํทิฏฺฐิ “กโรโต การยโต, ฉินฺทโต เฉทาปยโต, ปจโต ปาจาปยโต, โสจยโต โสจาปยโต, กิลมโต กิลมาปยโต, ผนฺทโต ผนฺทาปยโต, ปาณมติปาตยโต, อทินฺนํ อาทิยโต, สนฺธิํ ฉินฺทโต, นิลฺโลปํ หรโต,}

\vspace{\tipitakaparskip}
\csromancenter{คามณิสํยุตฺต เอกาคาริกํ กโรโต, ปริปนฺเถ ติฏฺฐโต, ปรทารํ คจฺฉโต, มุสา ภณโต, กโรโต น กรียติ ปาปํ, ขุรปริยนฺเตน เจปิ จกฺเกน โย อิมิสฺสา ปถวิยา ปาเณ เอกํ มํสขลํ เอกํ มํสปุญฺชํ กเรยฺย, นตฺถิ ตโตนิทานํ ปาปํ, นตฺถิ ปาปสฺส อาคโม.\csromanspacer{} ทกฺขิณญฺเจปิ คงฺคาย ตีรํ คจฺเฉยฺย หนนฺโต ฆาเตนฺโต ฉินฺทนฺโต เฉทาเปนฺโต ปจนฺโต ปาจาเปนฺโต, นตฺถิ ตโตนิทานํ ปาปํ, นตฺถิ ปาปสฺส อาคโม.\csromanspacer{} อุตฺตรญฺเจปิ คงฺคาย ตีรํ คจฺเฉยฺย ททนฺโต ทาเปนฺโต ยชนฺโต ยชาเปนฺโต, นตฺถิ ตโตนิทานํ ปุญฺญํ, นตฺถิ ปุญฺญสฺส อาคโม.\csromanspacer{} ทาเนน ทเมน สํยเมน สจฺจวชฺเชน นตฺถิ ปุญฺญํ, นตฺถิ ปุญฺญสฺส อาคโม”ติ.}
\prose{\csromanfolio{527}เอโก สตฺถา เอวํวาที เอวํทิฏฺฐิ “กโรโต การยโต, ฉินฺทโต เฉทาปยโต, ปจโต ปาจาปยโต, โสจยโต โสจาปยโต, กิลมโต กิลมาปยโต, ผนฺทโต ผนฺทาปยโต, ปาณมติปาตยโต, อทินฺนํ อาทิยโต, สนฺธิํ ฉินฺทโต, นิลฺโลปํ หรโต, เอกาคาริกํ กโรโต, ปริปนฺเถ ติฏฺฐโต, ปรทารํ คจฺฉโต, มุสา ภณโต, กโรโต กรียติ ปาปํ, ขุรปริยนฺเตน เจปิ จกฺเกน โย อิมิสฺสา ปถวิยา ปาเณ เอกํ มํสขลํ เอกํ มํสปุญฺชํ กเรยฺย, อตฺถิ ตโตนิทานํ ปาปํ, อตฺถิ ปาปสฺส อาคโม.\csromanspacer{} ทกฺขิณญฺเจปิ คงฺคาย ตีรํ คจฺเฉยฺย หนนฺโต ฆาเตนฺโต ฉินฺทนฺโต เฉทาเปนฺโต ปจนฺโต ปาจาเปนฺโต, อตฺถิ ตโตนิทานํ ปาปํ, อตฺถิ ปาปสฺส อาคโม.\csromanspacer{} อุตฺตรญฺเจปิ คงฺคาย ตีรํ คจฺเฉยฺย ททนฺโต ทาเปนฺโต ยชนฺโต ยชาเปนฺโต, อตฺถิ ตโตนิทานํ ปุญฺญํ, อตฺถิ ปุญฺญสฺส อาคโม.\csromanspacer{} ทาเนน ทเมน สํยเมน สจฺจวชฺเชน อตฺถิ ปุญฺญํ, อตฺถิ ปุญฺญสฺส อาคโม”ติ.}

\prose{ตสฺส มยฺหํ ภนฺเต อหุเทว กงฺขา, อหุ วิจิกิจฺฉา “โกสุ นาม อิเมสํ ภวตํ สมณพฺราหฺมณานํ สจฺจํ อาห โก มุสา”ติ.}

\prose{อลํ หิ เต คามณิ กงฺขิตุํ, อลํ วิจิกิจฺฉิตุํ, กงฺขนีเย จ ปน เต ฐาเน วิจิกิจฺฉา อุปฺปนฺนาติ.\csromanspacer{} เอวํ ปสนฺโนหํ ภนฺเต ภควติ “ปโหติ เม ภควา ตถา ธมฺมํ เทเสตุํ, ยถาหํ อิมํ กงฺขาธมฺมํ ปชเหยฺยนฺ”ติ.}

\prose{อตฺถิ คามณิ ธมฺมสมาธิ, ตตฺร เจ ตฺวํ จิตฺตสมาธิํ ปฏิลเภยฺยาสิ, เอวํ ตฺวํ อิมํ กงฺขาธมฺมํ ปชเหยฺยาสิ.\csromanspacer{} กตโม จ คามณิ ธมฺมสมาธิ, อิธ คามณิ อริยสาวโก ปาณาติปาตํ ปหาย ปาณาติปาตา}

\gambhira{สฬายตนวคฺคสํยุตฺตปาฬิ}
\prose{\csromanfolio{528}ปฏิวิรโต โหติ, อทินฺนาทานํ ปหาย อทินฺนาทานา ปฏิวิรโต โหติ, กาเมสุมิจฺฉาจารํ ปหาย กาเมสุมิจฺฉาจารา ปฏิวิรโต โหติ, มุสาวาทํ ปหาย มุสาวาทา ปฏิวิรโต โหติ, ปิสุณํ วาจํ ปหาย ปิสุณาย วาจาย ปฏิวิรโต โหติ, ผรุสํ วาจํ ปหาย ผรุสาย วาจาย ปฏิวิรโต โหติ, สมฺผปฺปลาปํ ปหาย สมฺผปฺปลาปา ปฏิวิรโต โหติ, อภิชฺฌํ ปหาย อนภิชฺฌาลุ โหติ, พฺยาปาทปโทสํ ปหาย อพฺยาปนฺนจิตฺโต โหติ, มิจฺฉาทิฏฺฐิํ ปหาย สมฺมาทิฏฺฐิโก โหติ.}

\prose{ส โข โส คามณิ อริยสาวโก เอวํ วิคตาภิชฺโฌ วิคตพฺยาปาโท อสมฺมูโฬฺห สมฺปชาโน ปฏิสฺสโต เมตฺตาสหคเตน เจตสา เอกํ ทิสํ ผริตฺวา วิหรติ.\csromanspacer{} ตถา ทุติยํ.\csromanspacer{} ตถา ตติยํ.\csromanspacer{} ตถา จตุตฺถํ.\csromanspacer{} อิติ อุทฺธมโธ ติริยํ สพฺพธิ สพฺพตฺตตาย สพฺพาวนฺตํ โลกํ เมตฺตาสหคเตน เจตสา วิปุเลน มหคฺคเตน อปฺปมาเณน อเวเรน อพฺยาปชฺเชน ผริตฺวา วิหรติ, โส อิติ ปฏิสญฺจิกฺขติ “ยฺวายํ สตฺถา เอวํวาที เอวํทิฏฺฐิ ‘นตฺถิ ทินฺนํ, นตฺถิ ยิฏฺฐํ, นตฺถิ หุตํ, นตฺถิ สุกตทุกฺกฏานํ กมฺมานํ ผลํ วิปาโก, นตฺถิ อยํ โลโก, นตฺถิ ปโร โลโก, นตฺถิ มาตา, นตฺถิ ปิตา, นตฺถิ สตฺตา โอปปาติกา, นตฺถิ โลเก สมณพฺราหฺมณา สมฺมคฺคตา สมฺมาปฏิปนฺนา, เย อิมญฺจ โลกํ ปรญฺจ โลกํ สยํ อภิญฺญา สจฺฉิกตฺวา ปเวเทนฺตี’ติ.\csromanspacer{} สเจ ตสฺส โภโต สตฺถุโน สจฺจํ วจนํ, อปณฺณกตาย มยฺหํ, ยฺวาหํ\footnote{โยหํ (สี, สฺยา, กํ, อิ)} น กิญฺจิ\footnote{กญฺจิ (?)} พฺยาพาเธมิ ตสํ วา ถาวรํ วา, อุภยเมตฺถ\footnote{อุภยตฺถ เม (?) ม 2. 65 ปิฏฺเฐ ปาฬิยา สํสนฺเทตพฺพํ.} กฏคฺคาโห, ยญฺจมฺหิ กาเยน สํวุโต วาจาย สํวุโต มนสา สํวุโต, ยญฺจ กายสฺส เภทา ปรํ มรณา สุคติํ สคฺคํ โลกํ อุปปชฺชิสฺสามี”ติ\footnote{ปรํ มรณา น อุปปชฺชิสฺสามีติ (?)}.\csromanspacer{} ตสฺส ปาโมชฺชํ ชายติ, ปมุทิตสฺส ปีติ ชายติ, ปีติมนสฺส กาโย ปสฺสมฺภติ, ปสฺสทฺธกาโย สุขํ เวทยติ, สุขิโน จิตฺตํ สมาธิยติ.\csromanspacer{} อยํ โข คามณิ ธมฺมสมาธิ, ตตฺร เจ ตฺวํ จิตฺตสมาธิํ ปฏิลเภยฺยาสิ, เอวํ ตฺวํ อิมํ กงฺขาธมฺมํ ปชเหยฺยาสิ.}

\prose{ส โข โส คามณิ อริยสาวโก เอวํ วิคตาภิชฺโฌ วิคตพฺยาปาโท อสมฺมูโฬฺห สมฺปชาโน ปฏิสฺสโต เมตฺตาสหคเตน เจตสา เอกํ ทิสํ ผริตฺวา วิหรติ.\csromanspacer{} ตถา ทุติยํ.\csromanspacer{} ตถา ตติยํ.}

\vspace{\tipitakaparskip}
\csromancenter{คามณิสํยุตฺต ตถา จตุตฺถํ.\csromanspacer{} อิติ อุทฺธมโธ ติริยํ สพฺพธิ สพฺพตฺตตาย สพฺพาวนฺตํ โลกํ เมตฺตาสหคเตน เจตสา วิปุเลน มหคฺคเตน อปฺปมาเณน อเวเรน อพฺยาปชฺเชน ผริตฺวา วิหรติ, โส อิติ ปฏิสญฺจิกฺขติ “ยฺวายํ สตฺถา เอวํวาที เอวํทิฏฺฐิ ‘อตฺถิ ทินฺนํ, อตฺถิ ยิฏฺฐํ, อตฺถิ หุตํ, อตฺถิ สุกตทุกฺกฏานํ กมฺมานํ ผลํ วิปาโก, อตฺถิ อยํ โลโก, อตฺถิ ปโร โลโก, อตฺถิ มาตา, อตฺถิ ปิตา, อตฺถิ สตฺตา โอปปาติกา, อตฺถิ โลเก สมณพฺราหฺมณา สมฺมคฺคตา สมฺมาปฏิปนฺนา, เย อิมญฺจ โลกํ ปรญฺจ โลกํ สยํ อภิญฺญา สจฺฉิกตฺวา ปเวเทนฺตี’ติ.\csromanspacer{} สเจ ตสฺส โภโต สตฺถุโน สจฺจํ วจนํ, อปณฺณกตาย มยฺหํ, ยฺวาหํ น กิญฺจิ พฺยาพาเธมิ ตสํ วา ถาวรํ วา, อุภยเมตฺถ กฏคฺคาโห, ยญฺจมฺหิ กาเยน สํวุโต วาจาย สํวุโต มนสา สํวุโต, ยญฺจ กายสฺส เภทา ปรํ มรณา สุคติํ สคฺคํ โลกํ อุปปชฺชิสฺสามี”ติ.\csromanspacer{} ตสฺส ปาโมชฺชํ ชายติ, ปมุทิตสฺส ปีติ ชายติ, ปีติมนสฺส กาโย ปสฺสมฺภติ, ปสฺสทฺธกาโย สุขํ เวทยติ, สุขิโน จิตฺตํ สมาธิยติ.\csromanspacer{} อยํ โข คามณิ ธมฺมสมาธิ, ตตฺร เจ ตฺวํ จิตฺตสมาธิํ ปฏิลเภยฺยาสิ, เอวํ ตฺวํ อิมํ กงฺขาธมฺมํ ปชเหยฺยาสิ.}
\prose{\csromanfolio{529}ส โข โส คามณิ อริยสาวโก เอวํ วิคตาภิชฺโฌ วิคตพฺยาปาโท อสมฺมูโฬฺห สมฺปชาโน ปฏิสฺสโต เมตฺตาสหคเตน เจตสา เอกํ ทิสํ ผริตฺวา วิหรติ.\csromanspacer{} ตถา ทุติยํ.\csromanspacer{} ตถา ตติยํ.\csromanspacer{} ตถา จตุตฺถํ.\csromanspacer{} อิติ อุทฺธมโธ ติริยํ สพฺพธิ สพฺพตฺตตาย สพฺพาวนฺตํ โลกํ เมตฺตาสหคเตน เจตสา วิปุเลน มหคฺคเตน อปฺปมาเณน อเวเรน อพฺยาปชฺเชน ผริตฺวา วิหรติ, โส อิติ ปฏิสญฺจิกฺขติ “ยฺวายํ สตฺถา เอวํวาที เอวํทิฏฺฐิ ‘กโรโต การยโต, ฉินฺทโต เฉทาปยโต, ปจโต ปาจาปยโต, โสจยโต โสจาปยโต, กิลมโต กิลมาปยโต, ผนฺทโต ผนฺทาปยโต, ปาณมติปาตยโต, อทินฺนํ อาทิยโต, สนฺธิํ ฉินฺทโต, นิลฺโลปํ หรโต, เอกาคาริกํ กโรโต, ปริปนฺเถ ติฏฺฐโต, ปรทารํ คจฺฉโต, มุสา ภณโต, กโรโต น กรียติ ปาปํ, ขุรปริยนฺเตน เจปิ จกฺเกน โย อิมิสฺสา ปถวิยา ปาเณ เอกํ มํสขลํ เอกํ มํสปุญฺชํ กเรยฺย, นตฺถิ ตโตนิทานํ ปาปํ, นตฺถิ ปาปสฺส อาคโม.\csromanspacer{} ทกฺขิณญฺเจปิ คงฺคาย ตีรํ คจฺเฉยฺย หนนฺโต ฆาเตนฺโต ฉินฺทนฺโต เฉทาเปนฺโต ปจนฺโต ปาจาเปนฺโต, นตฺถิ}

\gambhira{สฬายตนวคฺคสํยุตฺตปาฬิ}
\prose{\csromanfolio{530}ตโตนิทานํ ปาปํ, นตฺถิ ปาปสฺส อาคโม.\csromanspacer{} อุตฺตริญฺเจปิ คงฺคาย ตีรํ คจฺเฉยฺย ททนฺโต ทาเปนฺโต ยชนฺโต ยชาเปนฺโต, นตฺถิ ตโตนิทานํ ปุญฺญํ, นตฺถิ ปุญฺญสฺส อาคโม.\csromanspacer{} ทาเนน ทเมน สํยเมน สจฺจวชฺเชน นตฺถิ ปุญฺญํ, นตฺถิ ปุญฺญสฺส อาคโม’ติ.\csromanspacer{} สเจ ตสฺส โภโต สตฺถุโน สจฺจํ วจนํ, อปณฺณกตาย มยฺหํ, ยฺวาหํ น กิญฺจิ พฺยาพาเธมิ ตสํ วา ถาวรํวา, อุภยเมตฺถ กฏคฺคาโห, ยญฺจมฺหิ กาเยน สํวุโต วาจาย สํวุโต มนสา สํวุโต, ยญฺจ กายสฺส เภทา ปรํ มรณา สุคติํ สคฺคํ โลกํ อุปปชฺชิสฺสามี”ติ\footnote{ปรํ มรณา น อุปปชฺชิสฺสามีติ (?)}.\csromanspacer{} ตสฺส ปาโมชฺชํ ชายติ, ปมุทิตสฺส ปีติ ชายติ, ปีติมนสฺส กาโย ปสฺสมฺภติ, ปสฺสทฺธกาโย สุขํ เวทยติ, สุขิโน จิตฺตํ สมาธิยติ.\csromanspacer{} อยํ โข คามณิ ธมฺมสมาธิ, ตตฺร เจ ตฺวํ จิตฺตสมาธิํ ปฏิลเภยฺยาสิ, เอวํ ตฺวํ อิมํ กงฺขาธมฺมํ ปชเหยฺยาสิ.}

\prose{ส โข โส คามณิ อริยสาวโก เอวํ วิคตาภิชฺโฌ วิคตพฺยาปาโท อสมฺมูโฬฺห สมฺปชาโน ปฏิสฺสโต เมตฺตาสหคเตน เจตสา เอกํ ทิสํ ผริตฺวา วิหรติ.\csromanspacer{} ตถา ทุติยํ.\csromanspacer{} ตถา ตติยํ.\csromanspacer{} ตถา จตุตฺถํ.\csromanspacer{} อิติ อุทฺธมโธ ติริยํ สพฺพธิ สพฺพตฺตตาย สพฺพาวนฺตํ โลกํ เมตฺตาสหคเตน เจตสา วิปุเลน มหคฺคเตน อปฺปมาเณน อเวเรน อพฺยาปชฺเชน ผริตฺวา วิหรติ.\csromanspacer{} โส อิติ ปฏิสญฺจิกฺขติ “ยฺวายํ สตฺถา เอวํวาที เอวํทิฏฺฐิ ‘กโรโต การยโต, ฉินฺทโต เฉทาปยโต, ปจโต ปาจาปยโต, โสจยโต โสจาปยโต, กิลมโต กิลมาปยโต, ผนฺทโต ผนฺทาปยโต, ปาณมติปาตยโต, อทินฺนํ อาทิยโต, สนฺธิํ ฉินฺทโต, นิลฺโลปํ หรโต, เอกาคาริกํ กโรโต, ปริปนฺเถ ติฏฺฐโต, ปรทารํ คจฺฉโต, มุสา ภณโต, กโรโต กรียติ ปาปํ, ขุรปริยนฺเถน เจปิ จกฺเกน โย อิมิสฺสา ปถวิยา ปาเณ เอกํ มํสขลํ เอกํ มํสปุญฺชํ กเรยฺย, อตฺถิ ตโตนิทานํ ปาปํ, อตฺถิ ปาปสฺส อาคโม.\csromanspacer{} ทกฺขิณญฺเจปิ คงฺคาย ตีรํ คจฺเฉยฺย หนนฺโต ฆาเตนฺโต ฉินฺทนฺโต เฉทาเปนฺโต ปจนฺโต ปาจาเปนฺโต, อตฺถิ ตโตนิทานํ ปาปํ, อตฺถิ ปาปสฺส อาคโม.\csromanspacer{} อุตฺตรญฺเจปิ คงฺคาย ตีรํ คจฺเฉยฺย ททนฺโต ทาเปนฺโต ยชนฺโต ยชาเปนฺโต, อตฺถิ}

\vspace{\tipitakaparskip}
\csromancenter{คามณิสํยุตฺต ตโตนิทานํ ปุญฺญํ, อตฺถิ ปุญฺญสฺส อาคโม.\csromanspacer{} ทาเนน ทเมน สํยเมน สจฺจวชฺเชน อตฺถิ ปุญฺญํ, อตฺถิ ปุญฺญสฺส อาคโม’ติ.\csromanspacer{} สเจ ตสฺส โภโต สตฺถุโน สจฺจํ วจนํ, อปณฺณกตาย มยฺหํ, ยฺวาหํ น กิญฺจิ พฺยาพาเธมิ ตสํ วา ถาวรํ วา, อุภยเมตฺถ กฏคฺคาโห, ยญฺจมฺหิ กาเยน สํวุโส วาจาย สํวุโต มนสา สํวุโต, ยญฺจ กายสฺส เภทา ปรํ มรณา สุคติํ สคฺคํ โลกํ อุปปชฺชิสฺสามี”ติ.\csromanspacer{} ตสฺส ปาโมชฺชํ ชายติ, ปมุทิตสฺส ปีติ ชายติ, ปีติมนสฺส กาโย ปสฺสมฺภติ, ปสฺสทฺธกาโย สุขํ เวทยิติ, สุขิโน จิตฺตํ สมาธิยติ.\csromanspacer{} อยํ โข คามณิ ธมฺมสมาธิ, ตตฺร เจ ตฺวํ จิตฺตสมาธิํ ปฏิลเภยฺยาสิ, เอวํ ตฺวํ อิมํ กงฺขาธมฺมํ ปชเหยฺยาสิ.}
\prose{\csromanfolio{531}ส โข โส คามณิ อริยสาวโก เอวํ วิคตาภิชฺโฌ วิคตพฺยาปาโท อสมฺมูโฬฺห สมฺปชาโน ปฏิสฺสโต กรุณาสหคเตน เจตสา เอกํ ทิสํ ผริตฺวา วิหรติ ฯเปฯ มุทิตาสหคเตน เจตสา เอกํ ทิสํ ผริตฺวา วิหรติ ฯเปฯ.}

\prose{ส โข โส คามณิ อริยสาวโก เอวํ วิคตาภิชฺโฌ วิคตพฺยาปาโท อสมฺมูโฬฺห สมฺปชาโน ปฏิสฺสโต อุเปกฺขาสหคเตน เจตสา เอกํ ทิสํ ผริตฺวา วิหรติ.\csromanspacer{} ตถา ทุติยํ.\csromanspacer{} ตถา ตติยํ.\csromanspacer{} ตถา จตุตฺถํ.\csromanspacer{} อิติ อุทฺธมโธ ติริยํ สพฺพธิ สพฺพตฺตตาย สพฺพาวนฺตํ โลกํ อุเปกฺขาสหคเตน เจตสา วิปุเลน มหคฺคเตน อปฺปมาเณน อเวเรน อพฺยาปชฺเชน ผริตฺวา วิหรติ, โส อิติ ปฏิสญฺจิกฺขติ “ยฺวายํ สตฺถา เอวํวาที เอวํทิฏฺฐิ ‘นตฺถิ ทินฺนํ, นตฺถิ ยิฏฺฐํ, นตฺถิ หุตํ, นตฺถิ สุกตทุกฺกฏานํ กมฺมานํ ผลํ วิปาโก, นตฺถิ อยํ โลโก, นตฺถิ ปโร โลโก, นตฺถิ มาตา, นตฺถิ ปิตา, นตฺถิ สตฺตา โอปปาติกา, นตฺถิ โลเก สมณพฺราหฺมณา สมฺมคฺคตา สมฺมาปฏิปนฺนา, เย อิมญฺจ โลกํ ปรญฺจ โลกํ สยํ อภิญฺญา สจฺฉิกตฺวา ปเวเทนฺตี’ติ.\csromanspacer{} สเจ ตสฺส โภโต สตฺถุโน สจฺจํ วจนํ, อปณฺณกตาย มยฺหํ, ยฺวาหํ น กิญฺจิ พฺยาพาเธมิ ตสํ วา ถาวรํ วา, อุภยเมตฺถ กฏคฺคาโห, ยญฺจมฺหิ กาเยน สํวุโส วาจาย สํวุโต มนสา สํวุโต, ยญฺจ กายสฺส เภทา ปรํ มรณา สุคติํ สคฺคํ โลกํ อุปปชฺชิสฺสามี”ติ.\csromanspacer{} ตสฺส ปาโมชฺชํ ชายติ, ปมุทิตสฺส ปีติ ชายติ, ปีติมนสฺส กาโย ปสฺสมฺภติ, ปสฺสทฺธกาโย สุขํ เวทยติ, สุขิโน จิตฺตํ สมาธิยติ.}

\gambhira{สฬายตนวคฺคสํยุตฺตปาฬิ}
\prose{\csromanfolio{532}อยํ โข คามณิ ธมฺมสมาธิ, ตตฺร เจ ตฺวํ จิตฺตสมาธิํ ปฏิลเภยฺยาสิ, เอวํ ตฺวํ อิมํ กงฺขาธมฺมํ ปชเหยฺยาสิ.}

\prose{ส โข โส คามณิ อริยสาวโก เอวํ วิคตาภิชฺโฌ วิคตพฺยาปาโท อสมฺมูโฬฺห สมฺปชาโน ปฏิสฺสโต อุเปกฺขาสหคเตน เจตสา เอกํ ทิสํ ผริตฺวา วิหรติ.\csromanspacer{} ตถา ทุติยํ.\csromanspacer{} ตถา ตติยํ.\csromanspacer{} ตถา จตุตฺถํ.\csromanspacer{} อิติ อุทฺธมโธ ติริยํ สพฺพธิ สพฺพตฺตตาย สพฺพาวนฺตํ โลกํ อุเปกฺขาสหคเตน เจตสา วิปุเลน มหคฺคเตน อปฺปมาเณน อเวเรน อพฺยาปชฺเชน ผริตฺวา วิหรติ.\csromanspacer{} โส อิติ ปฏิสญฺจิกฺขติ “ยฺวายํ สตฺถา เอวํวาที เอวํทิฏฺฐิ ‘อตฺถิ ทินฺนํ, อตฺถิ ยิฏฺฐํ, อตฺถิ หุตํ, อตฺถิ สุกตทุกฺกฏานํ กมฺมานํ ผลํ วิปาโก, อตฺถิ อยํ โลโก, อตฺถิ ปโร โลโก, อตฺถิ มาตา, อตฺถิ ปิตา, อตฺถิ สตฺตา โอปปาติกา, อตฺถิ โลเก สมณาพฺราหฺมณา สมฺมคฺคตา สมฺมาปฏิปนฺนา, เย อิมญฺจ โลกํ ปรญฺจ โลกํ สยํ อภิญฺญา สจฺฉิกตฺวา ปเวเทนฺตี’ติ.\csromanspacer{} สเจ ตสฺส โภโต สตฺถุโน สจฺจํ วจนํ, อปณฺณกตาย มยฺหํ, ยฺวาหํ น กิญฺจิ พฺยาพาเธมิ ตสํ วา ถาวรํ วา, อุภยเมตฺถ กฏคฺคาโห, ยญฺจมฺหิ กาเยน สํวุโต วาจาย สํวุโต มนสา สํวุโต, ยญฺจ กายสฺส เภทา ปรํ มรณา สุคติํ สคฺคํ โลกํ อุปปชฺชิสฺสามี”ติ.\csromanspacer{} ตสฺส ปาโมชฺชํ ชายติ, ปมุทิตสฺส ปีติ ชายติ, ปีติมนสฺส กาโย ปสฺสมฺภติ, ปสฺสทฺธากาโย สุขํ เวทยติ, สุขิโน จิตฺตํ สมาธิยติ.\csromanspacer{} อยํ โข คามณิ ธมฺมสมาธิ, ตตฺร เจ ตฺวํ จิตฺตสมาธิํ ปฏิลเภยฺยาสิ, เอวํ ตฺวํ อิมํ กงฺขาธมฺมํ ปชเหยฺยาสิ.}

\prose{ส โข โส คามณิ อริยสาวโก เอวํ วิคตาภิชฺโฌ วิคตพฺยาปาโท อสมฺมู โฬฺห สมฺปชาโน ปฏิสฺสโต อุเปกฺขาสหคเตน เจตสา เอกํ ทิสํ ผริตฺวา วิหรติ.\csromanspacer{} ตถา ทุติยํ.\csromanspacer{} ตถา ตติยํ.\csromanspacer{} ตถา จตุตฺถํ.\csromanspacer{} อิติ อุทฺธมโธ ติริยํ สพฺพธิ สพฺพตฺตตาย สพฺพาวนฺตํ โลกํ อุเปกฺขาสหคเตน เจตสา วิปุเลน มหคฺคเตน อปฺปมาเณน อเวเรน อพฺยาปชฺเชน ผริตฺวา วิหรติ.\csromanspacer{} โส อิติ ปฏิสญฺจิกฺขติ “ยฺวายํ สตฺถา เอวํวาที เอวํทิฏฺฐิ ‘กโรโต การยโต, เฉทโต เฉทาปยโต, ปจโต ปาจาปยโต, โสจยโต โสจาปยโต, กิลมโต กิลมาปยโต, ผนฺทโต ผนฺทาปยโต, ปาณมติปาตยโต, อทินฺนํ อาทิยโต, สนฺธิํ ฉินฺทโต, นิลฺโลปํ หรโต,}

\vspace{\tipitakaparskip}
\csromancenter{คามณิสํยุตฺต เอกาคาริกํ กโรโต, ปริปนฺเถ ติฏฺฐโต, ปรทารํ คจฺฉโต, มุสา ภณโต, กโรโต น กรียติ ปาปํ, ขุรปริยนฺเตน เจปิ จกฺเกน โย อิมิสฺสา ปถวิยา ปาเณ เอกํ มํสขลํ เอกํ มํสปุญฺชํ กเรยฺย, นตฺถิ ตโตนิทานํ ปาปํ, นตฺถิ ปาปสฺส อาคโม.\csromanspacer{} ทกฺขิณญฺเจปิ คงฺคาย ตีรํ คจฺเฉยฺย หนนฺโต ฆาเตนฺโต ฉินฺทนฺโต เฉทาเปนฺโต ปจนฺโต ปาจาเปนฺโต, นตฺถิ ตโตนิทานํ ปารํ, นตฺถิ ปาปสฺส อาคโม.\csromanspacer{} อุตฺตรญฺเจปิ คงฺคาย ตีรํ คจฺเฉยฺย ททนฺโต ทาเปนฺโต ยชนฺโต ยชาเปนฺโต, นตฺถิ ตโตนิทานํ ปุญฺญํ, นตฺถิ ปุญฺญสฺส อาคโม.\csromanspacer{} ทาเนน ทเมน สํยเมน สจฺจวชฺเชน นตฺถิ ปุญฺญํ, นตฺถิ ปุญฺญสฺส อาคโม’ติ.\csromanspacer{} สเจ ตสฺส โภโต สตฺถุโน สจฺจํ วจนํ, อปณฺณกตาย มยฺหํ, ยฺวาหํ น กิญฺจิ พฺยาพาเธมิ ตสํ วา ถาวรํ วา, อุภยเมตฺถ กฏคฺคาโห, ยญฺจมฺหิ กาเยน สํวุโส วาจาย สํวุโต มนสา สํวุโต, ยญฺจ กายสฺส เภทา ปรํ มรณา สุคติํ สคฺคํ โลกํ อุปปชฺชิสฺสามี”ติ.\csromanspacer{} ตสฺส ปาโมชฺชํ ชายติ, ปมุทิตสฺส ปีติ ชายติ, ปีติมนสฺส กาโย ปสฺสมฺภติ, ปสฺสทฺธกาโย สุขํ เวทยติ, สุขิโน จิตฺตํ สมาธิยติ.\csromanspacer{} อยํ โข คามณิ ธมฺมสมาธิ, ตตฺร เจ ตฺวํ จิตฺตสมาธิํ ปฏิลเภยฺยาสิ, เอวํ ตฺวํ อิมํ กงฺขาธมฺมํ ปชเหยฺยาสิ.}
\prose{\csromanfolio{533}ส โข โส คามณิ อริยสาวโก เอวํ วิคตาภิชฺโฌ วิคตพฺยาปาโท อสมฺมูโฬฺห สมฺปชาโน ปฏิสฺสโต อุเปกฺขาสหคเตน เจตสา เอกํ ทิสํ ผริตฺวา วิหรติ.\csromanspacer{} ตถา ทุติยํ.\csromanspacer{} ตถา ตติยํ.\csromanspacer{} ตถา จตุตฺถํ.\csromanspacer{} อิติ อุทฺธมโธ ติริยํ สพฺพธิ สพฺพตฺตตาย สพฺพาวนฺตํ โลกํ อุเปกฺขาสหคเตน เจตสา วิปุเลน มหคฺคเตน อปฺปมาเณน อเวเรน อพฺยาปชฺเชน ผริตฺวา วิหรติ.\csromanspacer{} โส อิติ ปฏิสญฺจิกฺขติ “ยฺวายํ สตฺถา เอวํวาที เอวํทิฏฺฐิ ‘กโรโต การยโต, ฉินฺทโต เฉทาปยโต, ปจโต ปาจาปยโต, โสจยโต โสจาปยโต, กิลมโต กิลมาปยโต, ผนฺทโต ผนฺทาปยโต, ปาณมติปาตยโต, อทินฺนํ อาทิยโต, สนฺธิํ ฉินฺทโต, นิลฺโลปํ หรโต, เอกาคาริกํ กโรโต, ปริปนฺเถ ติฏฺฐโต, ปรทารํ คจฺฉโต, มุสา ภณโต, กโรโต กรียติ ปาปํ, ขุรปริยนฺเตน เจปิ จกฺเกน โย อิมิสฺสา ปถวิยา ปาเณ เอกํ มํสขลํ เอกํ มํสปุญฺชํ กเรยฺย, อตฺถิ ตโตนิทานํ ปาปํ, อตฺถิ ปาปสฺส อาคโม.\csromanspacer{} ทกฺขิณญฺเจปิ คงฺคาย ตีรํ คจฺเฉยฺย}

\gambhira{สฬายตนวคฺคสํยุตฺตปาฬิ}
\prose{\csromanfolio{534}หนนฺโต ฆาเตนฺโต ฉินฺทนฺโต เฉทาเปนฺโต ปจนฺโต ปาจาเปนฺโต, อตฺถิ ตโตนิทานํ ปาปํ, อตฺถิ ปาปสฺส อาคโม.\csromanspacer{} อุตฺตรญฺเจปิ คงฺคาย ตีรํ คจฺเฉยฺย ททนฺโต ทาเปนฺโต ยชนฺโต ยชาเปนฺโต, อตฺถิ ตโตนิทานํ ปุญฺญํ, อตฺถิ ปุญฺญสฺส อาคโม.\csromanspacer{} ทาเนน ทเมน สํยเมนํ สจฺจวชฺเชน อตฺถิ ปุญฺญํ, อตฺถิ ปุญฺญสฺส อาคโม’ติ.\csromanspacer{} สเจ ตสฺส โภโต สตฺถุโน สจฺจํ วจนํ, อปณฺณกตาย มยฺหํ, ยฺวาหํ น กิญฺจิ พฺยาพาเธมิ ตสํ วา ถาวรํ วา, อุภยเมตฺถ กฏคฺคาโห, ยญฺจมฺหิ กาเยน สํวุโต วาจาย สํวุโต มนสา สํวุโต, ยญฺจ กายสฺส เภทา ปรํ มรณา สุคติํ สคฺคํ โลกํ อุปปชฺชิสฺสามี”ติ.\csromanspacer{} ตสฺส ปาโมชฺชํ ชายติ, ปมุทิตสฺส ปีติ ชายติ, ปีติมนสฺส กาโย ปสฺสมฺภติ, ปสฺสทฺธกาโย สุขํ เวทยติ, สุขโน จิตฺตํ สมาธิยติ.\csromanspacer{} อยํ โข คามณิ ธมฺมสมาธิ, ตตฺร เจ ตฺวํ จิตฺตสมาธิํ ปฏิลเภยฺยาสิ, เอวํ ตฺวํ อิมํ กงฺขาธมฺมํ ปชเหยฺยาสีติ.}

\prose{เอวํ วุตฺเต ปาฏิลิโย คามณิ ภควนฺตํ เอตทโวจ “อภิกฺกนฺตํ ภนฺเต, อภิกฺกนฺตํ ภนฺเต ฯเปฯ อชฺชตคฺเค ปาณุเปตํ สรณํ คตนฺ”ติ.\csromanspacer{}\csromanspacer{}เตรสมํ.}

\vspace{\tipitakaparskip}
\csromancenter{\textbf{คามณิสํยุตฺตํ สมตฺตํ.}}
\titlehead{\textbf{ตสฺสุทฺทานํ}}
\csromangathameasure{จณฺโฑ ปุโฏ โยธาชีโว, หตฺถสฺโส อสิพนฺธโก. \\ เทสนา สงฺขกุลํ มณิจูฬํ, ภทฺรราสิยปาฏลีติ.}
\csromangathagroup{\csromangathastanza{จณฺโฑ ปุโฏ โยธาชีโว, หตฺถสฺโส อสิพนฺธโก. \\ เทสนา สงฺขกุลํ มณิจูฬํ, ภทฺรราสิยปาฏลีติ.}}

\csromanheader{2}{9. อสงฺขตสํยุตฺต}
\csromanheader{2}{1. ปฐมวคฺค}
\csromanmatikaanchor{mtk.286}
\csromantocmarkref{subsection}{1. ปฐมวคฺค}{mtk.286}
\bookmark[dest={mtk.286},level=2]{1. ปฐมวคฺค}
\csromanheader{2}{1. กายคตาสติสุตฺต}
\proseitem{366}{\csromanfolio{535}สาวตฺถินิทานํ.\csromanspacer{} อสงฺขตญฺจ โว ภิกฺขเว เทเสสฺสามิ อสงฺขตคามิญฺจ มคฺคํ, ตํ สุณาถ.\csromanspacer{} กตมญฺจ ภิกฺขเว อสงฺขตํ, โย ภิกฺขเว ราคกฺขโย โทสกฺขโย โมหกฺขโย.\csromanspacer{} อิทํ วุจฺจติ ภิกฺขเว อสงฺขตํ.\csromanspacer{} กตโม จ ภิกฺขเว อสงฺขตคามิมคฺโค, กายคตาสติ.\csromanspacer{} อยํ วุจฺจติ ภิกฺขเว อสงฺขตคามิมคฺโค.}

\prose{อิติ โข ภิกฺขเว เทสิตํ โว มยา อสงฺขตํ, เทสิโต อสงฺขตคามิมคฺโค.\csromanspacer{} ยํ ภิกฺขเว สตฺถารา กรณียํ สาวกานํ หิเตสินา อนุกมฺปเกน อนุกมฺปํ อุปาทาย, กตํ โว ตํ มยา.\csromanspacer{} เอตานิ ภิกฺขเว รุกฺขมูลานิ เอตานิ สุญฺญาคารานิ, ฌายถ\footnote{นิชฺฌายถ (ก)} ภิกฺขเว, มา ปมาทตฺถ, มา ปจฺฉา วิปฺปฏิสาริโน อหุวตฺถ, อยํ โว อมฺหากํ อนุสาสนีติ.\csromanspacer{}\csromanspacer{}ปฐมํ.}

\csromanheader{2}{2. สมถวิปสฺสนาสุตฺต}
\proseitem{367}{อสงฺขตญฺจ โว ภิกฺขเว เทเสสฺสามิ อสงฺขตคามิญฺจ มคฺคํ, ตํ สุณาถ.\csromanspacer{} กตมญฺจ ภิกฺขเว อสงฺขตํ, โย ภิกฺขเว ราคกฺขโย โทสกฺขโย โมหกฺขโย.\csromanspacer{} อิทํ วุจฺจติ ภิกฺขเว อสงฺขตํ.\csromanspacer{} กตโม จ ภิกฺขเว อสงฺขตคามิมคฺโค, สมโถ จ วิปสฺสนา จ.\csromanspacer{} อยํ วุจฺจติ ภิกฺขเว อสงฺขตคามิมคฺโค ฯเปฯ.\csromanspacer{} ทุติยํ.}

\csromanheader{2}{3. สวิตกฺกสวิจารสุตฺต}
\proseitem{368}{กตโม จ ภิกฺขเว อสงฺขตคามิมคฺโค, สวิตกฺกสวิจาโร สมาธิ อวิตกฺกวิจารมตฺโต สมาธิ อวิตกฺก-อวิจาโร สมาธิ.\csromanspacer{} อยํ วุจฺจติ ภิกฺขเว อสงฺขตคามิมคฺโค ฯเปฯ.\csromanspacer{}\csromanspacer{}ตติยํ.}

\gambhira{สฬายตนวคฺคสํยุตฺตปาฬิ}
\csromanheader{2}{4. สุญฺญตสมาธิสุตฺต}
\proseitem{369}{\csromanfolio{536}กตโม จ ภิกฺขเว อสงฺขตคามิมคฺโค, สุญฺญโต สมาธิ อนิมิตฺโต สมาธิ อปฺปณิหิโต สมาธิ.\csromanspacer{} อยํ วุจฺจติ ภิกฺขเว อสงฺขาตคามิมคฺโค ฯเปฯ.\csromanspacer{} จตุตฺถํ.}

\csromanheader{2}{5. สติปฏฺฐานสุตฺต}
\proseitem{370}{กตโม จ ภิกฺขเว อสงฺขตคามิมคฺโค, จตฺตโร สติปฏฺฐานา.\csromanspacer{} อยํ วุจฺจติ ภิกฺขเว อสงฺขตคามิมคฺโค ฯเปฯ.\csromanspacer{}\csromanspacer{}ปญฺจมํ.}

\csromanheader{2}{6. สมฺมปฺปธานสุตฺต}
\proseitem{371}{กตโม จ ภิกฺขเว อสงฺขตคามิมคฺโค, จตฺตาโร สมฺมปฺปธานา.\csromanspacer{} อยํ วุจฺจติ ภิกฺขเว อสงฺขตคามิมคฺโค ฯเปฯ.\csromanspacer{}\csromanspacer{}ฉฏฺฐํ.}

\csromanheader{2}{7. อิทฺธิปาทสุตฺต}
\proseitem{372}{กตโม จ ภิกฺขเว อสงฺขตคามิมคฺโค, จตฺตาโร อิทฺธิปาทา.\csromanspacer{} อยํ วุจฺจติ ภิกฺขเว ภิกฺขเว อสงฺขตคามิมคฺโค ฯเปฯ.\csromanspacer{} สตฺตมํ.}

\csromanheader{2}{8. อินฺทฺริยสุตฺต}
\proseitem{373}{กตโม จ ภิกฺขเว อสงฺขตคามิมคฺโค, ปญฺจินฺทฺริยานิ.\csromanspacer{} อยํ วุจฺจติ ภิกฺขเว อสงฺขตคามิมคฺโค ฯเปฯ.\csromanspacer{}\csromanspacer{}อฏฺฐมํ.}

\csromanheader{2}{9. พลสุตฺต}
\proseitem{374}{กตโม จ ภิกฺขเว อสงฺขตคามิมคฺโค, ปญฺจ พลานิ.\csromanspacer{} อยํ วุจฺจติ ภิกฺขเว อสงฺขตคามิมคฺโค ฯเปฯ.\csromanspacer{}\csromanspacer{}นวมํ.}

\csromanheader{2}{10. โพชฺฌงฺคสุตฺต}
\proseitem{375}{กตโม จ ภิกฺขเว อสงฺขตคามิมคฺโค, สตฺต โภชฺฌงฺคา.\csromanspacer{} อยํ วุจฺจติ ภิกฺขเว อสงฺขตคามิมคฺโค ฯเปฯ.\csromanspacer{}\csromanspacer{}ทสมํ.}

\csromanheader{2}{9. อสงฺขตสํยุตฺต}
\csromanheader{2}{11. มคฺคงฺคสุตฺต}
\proseitem{376}{\csromanfolio{537}กตโม จ ภิกฺขเว อสงฺขตคามิมคฺโค, อริโย อฏฺฐงฺคิโก มคฺโค.\csromanspacer{} อยํ วุจฺจติ ภิกฺขเว อสงฺขตคามิมคฺโค.\csromanspacer{} อิติ โข ภิกฺขเว เทสิตํ โว มยา อสงฺขตํ, เทสิโต อสงฺขตคามิมคฺโค.\csromanspacer{} ยํ ภิกฺขเว สตฺถารา กรณียํ สาวกานํ หิเตสินา อนุกมฺปเกน อนุกมฺปํ อุปาทาย, กตํ โว ตํ มยา.\csromanspacer{} เอตานิ ภิกฺขเว รุกฺขมูลานิ เอตานิ สุญฺญาคารานิ, ฌายถ ภิกฺขเว, มา ปมาทตฺถ, มา ปจฺฉา วิปฺปฏิสาริโน อหุวตฺถ, อยํ โว อมฺหากํ อนุสาสนีติ.\csromanspacer{}\csromanspacer{}เอกาทสมํ.\csromanspacer{} อสงฺขตสํยุตฺตสฺส ปฐโม วคฺโค.}

\titlehead{\textbf{ตสฺสุทฺทานํ}}
\csromangathameasure{กาโย สมโถ สวิตกฺโก, สุญฺญโต สติปฏฺฐานา. \\ สมฺมปฺปธานา อิทฺธิปาทา, อินฺทฺริย พล โพชฺฌงฺคา. \\ มคฺเคน เอกาทสมํ, ตสฺสุทฺทานํ ปวุจฺจติ.}
\csromangathagroup{\csromangathastanza{กาโย สมโถ สวิตกฺโก, สุญฺญโต สติปฏฺฐานา. \\ สมฺมปฺปธานา อิทฺธิปาทา, อินฺทฺริย พล โพชฺฌงฺคา. \\ มคฺเคน เอกาทสมํ, ตสฺสุทฺทานํ ปวุจฺจติ.}}

\csromanheader{2}{2. ทุติยวคฺค}
\csromanmatikaanchor{mtk.287}
\csromantocmarkref{subsection}{2. ทุติยวคฺค}{mtk.287}
\bookmark[dest={mtk.287},level=2]{2. ทุติยวคฺค}
\csromanheader{2}{1. อสงฺขตสุตฺต}
\proseitem{377}{อสงฺขตญฺจ โว ภิกฺขเว เทเสสฺสามิ อสงฺขตคามิญฺจ มคฺคํ, ตํ สุณาถ.\csromanspacer{} กตมญฺจ ภิกฺขเว อสงฺขตํ, โย ภิกฺขเว ราคกฺขโย โทสกฺขโย โมหกฺขโย.\csromanspacer{} อิทํ วุจฺจติ ภิกฺขเว อสงฺขตํ.\csromanspacer{} กตโม จ ภิกฺขเว อสงฺขตคามิมคฺโค, สมโถ.\csromanspacer{} อยํ วุจฺจติ ภิกฺขเว อสงฺขตคามิมคฺโค.\csromanspacer{} อิติ โข ภิกฺขเว เทสิตํ โว มยา อสงฺขตํ, เทสิโต อสงฺขตคามิมคฺโค.\csromanspacer{} ยํ ภิกฺขเว สตฺถารา กรณียํ สาวกานํ หิเตสินา อนุกมฺปเกน อนุกมฺปํ อุปาทาย, กตํ โว ตํ มยา.\csromanspacer{} เอตานิ ภิกฺขเว รุกฺขมูลานิ เอตานิ สุญฺญาคารานิ, ฌายถ ภิกฺขเว, มา ปมาทตฺถ, มา ปจฺฉา วิปฺปฏิสาริโน อหุวตฺถ, อยํ โว อมฺหากํ อนุสาสนีติ. (1)}

\gambhira{สฬายตนวคฺคสํยุตฺตปาฬิ}
\prose{\csromanfolio{538}อสงฺขตญฺจ โว ภิกฺขเว เทเสสฺสามิ อสงฺขตคามิญฺจ มคฺคํ, ตํ สุณาถ.\csromanspacer{} กตมญฺจ ภิกฺขเว อสงฺขตํ, โย ภิกฺขเว ราคกฺขโย โทสกฺขโย โมหกฺขโย.\csromanspacer{} อิทํ วุจฺจติ ภิกฺขเว อสงฺขตํ.\csromanspacer{} กตโม จ ภิกฺขเว อสงฺขตคามิมคฺโค, วิปสฺสนา.\csromanspacer{} อยํ วุจฺจติ ภิกฺขเว อสงฺขตคามิมคฺโค.\csromanspacer{} อิติ โข ภิกฺขเว เทสิตํ โว มยา อสงฺขตํ ฯเปฯ อยํ โว อมฺหากํ อนุสาสนีติ. (2)}

\prose{กตโม จ ภิกฺขเว อสงฺขตคามิมคฺโค, สวิตกฺโก สวิจาโร สมาธิ.\csromanspacer{} อยํ วุจฺจติ ภิกฺขเว อสงฺขตคามิมคฺโค ฯเปฯ.\csromanspacer{} กตโม จ ภิกฺขเว อสงฺขตคามิมคฺโค, อวิตกฺโก วิจารมตฺโต สมาธิ.\csromanspacer{} อยํ วุจฺจติ ภิกฺขเว อสงฺขตคามิมคฺโค ฯเปฯ.\csromanspacer{} กตโม จ ภิกฺขเว อสงฺขตคามิมคฺโค, อวิตกฺโก อวิจาโร สมาธิ.\csromanspacer{} อยํ วุจฺจติ ภิกฺขเว อสงฺขตคามิมคฺโค ฯเปฯ. (3-5)}

\prose{กตโม จ ภิกฺขเว อสงฺขตคามิมคฺโค, สุญฺญโต สมาธิ.\csromanspacer{} อยํ วุจฺจติ ภิกฺขเว อสงฺขตคามิมคฺโค ฯเปฯ.\csromanspacer{} กตโม จ ภิกฺขเว อสงฺขตคามิมคฺโค, อนิมิตฺโต สมาธิ.\csromanspacer{} อยํ วุจฺจติ ภิกฺขเว อสงฺขตคามิมคฺโค ฯเปฯ.\csromanspacer{} กตโม จ ภิกฺขเว อสงฺขตคามิมคฺโค, อปฺปณิหิโต สมาธิ.\csromanspacer{} อยํ วุจฺจติ ภิกฺขเว อสงฺขตคามิมคฺโค ฯเปฯ. (6- 8)}

\prose{กตโม จ ภิกฺขเว อสงฺขตคามิมคฺโค, อิธ ภิกฺขเว ภิกฺขุ กาเย กายานุปสฺสี วิหรติ อาตาปี สมฺปชาโน สติมา วิเนยฺย โลเก อภิชฺฌาโทมนสฺสํ.\csromanspacer{} อยํ วุจฺจติ ภิกฺขเว อสงฺขตคามิมคฺโค ฯเปฯ.\csromanspacer{} กตโม จ ภิกฺขเว อสงฺขตคามิมคฺโค, อิธ ภิกฺขเว ภิกฺขุ เวทนาสุ เวทนานุปสฺสี วิหรติ ฯเปฯ.\csromanspacer{} อยํ วุจฺจติ ภิกฺขเว อสงฺขตคามิมคฺโค ฯเปฯ.\csromanspacer{} กตโม จ ภิกฺขเว อสงฺขตคามิมคฺโค, อิธ ภิกฺขเว ภิกฺขุ จิตฺเต จิตฺตานุปสฺสี วิหรติ ฯเปฯ.\csromanspacer{} อยํ วุจฺจติ ภิกฺขเว อสงฺขตคามิมคฺโค ฯเปฯ.\csromanspacer{} กตโม จ ภิกฺขเว อสงฺขตคามิมคฺโค, อิธ ภิกฺขเว ภิกฺขุ ธมฺเมสุ ธมฺมานุปสฺสี วิหรติ ฯเปฯ อยํ วุจฺจติ ภิกฺขเว อสงฺขตคามิมคฺโค ฯเปฯ.(9-12)}

\prose{กตโม จ ภิกฺขเว อสงฺขตคามิมคฺโค, อิธ ภิกฺขเว ภิกฺขุ อนุปฺปนฺนานํ ปาปกานํ อกุสลานํ ธมฺมานํ อนุปฺปาทาย ฉนฺทํ ชเนติ วายมติ, วีริยํ อารภติ, จิตฺตํ ปคฺคณฺหาติ ปทหติ.\csromanspacer{} อยํ วุจฺจติ ภิกฺขเว อสงฺขตคามิมคฺโค ฯเปฯ.\csromanspacer{} กตโม จ ภิกฺขเว อสงฺขตคามิมคฺโค, อิธ ภิกฺขเว ภิกฺขุ}

\csromanheader{2}{9. อสงฺขตสํยุตฺต}
\prose{\csromanfolio{539}อุปฺปนฺนานํ ปาปกานํ อกุสลานํ ธมฺมานํ ปหานาย ฉนฺทํ ชเนติ วายมติ, วีริยํ อารภติ, จิตฺตํ ปคฺคณฺหาติ ปทหติ.\csromanspacer{} อยํ วุจฺจติ ภิกฺขเว อสงฺขตคามิมคฺโค ฯเปฯ.\csromanspacer{} กตโม จ ภิกฺขเว อสงฺขตคามิมคฺโค, อิธ ภิกฺขเว ภิกฺขุ อนุปฺปนฺนานํ กุสลานํ ธมฺมานํ อุปฺปาทาย ฉนฺทํ ชเนติ วายมติ, วีริยํ อารภติ, จิตฺตํ ปคฺคณฺหาติ ปทหติ.\csromanspacer{} อยํ วุจฺจติ ภิกฺขเว อสงฺขตคามิมคฺโค ฯเปฯ.\csromanspacer{} กตโม จ ภิกฺขเว อสงฺขตคามิมคฺโค, อิธ ภิกฺขเว ภิกฺขุ อุปฺปนฺนานํ กุสลานํ ธมฺมานํ ฐิติยา อสมฺโมสาย ภิยฺโยภาวาย เวปุลฺลาย ภาวนาย ปาริปูริยา ฉนฺทํ ชเนติ วายมติ, วีริยํ อารภติ, จิตฺตํ ปคฺคณฺหาติ ปทหติ.\csromanspacer{} อยํ วุจฺจติ ภิกฺขเว อสงฺขตคามิมคฺโค ฯเปฯ. (13-16)}

\prose{กตโม จ ภิกฺขเว อสงฺขตคามิมคฺโค, อิธ ภิกฺขเว ภิกฺขุ ฉนฺทสมาธิปธานสงฺขารสมนฺนาคตํ อิธิปาทํ ภาเวติ.\csromanspacer{} อยํ วุจฺจติ ภิกฺขเว อสงฺขตคามิมคฺโค ฯเปฯ.\csromanspacer{} กตโม จ ภิกฺขเว อสงฺขตคามิมคฺโค, อิธ ภิกฺขเว ภิกฺขุ วีริยสมาธิปธานสงฺขารสมนฺนาคตํ อิทฺธิปาทํ ภาเวติ.\csromanspacer{} อยํ วุจฺจติ ภิกฺขเว อสงฺขตคามิมคฺโค ฯเปฯ.\csromanspacer{} กตโม จ ภิกฺขเว อสงฺขตคามิมคฺโค, อิธ ภิกฺขเว ภิกฺขุ จิตฺตสมาธิปธานสงฺขารสมนฺนาคตํ อิทฺธิปาทํ ภาเวติ.\csromanspacer{} อยํ วุจฺจติ ภิกฺขเว อสงฺขตคามิมคฺโค ฯเปฯ.\csromanspacer{} กตโม จ ภิกฺขเว อสงฺขตคามิมคฺโค, อิธ ภิกฺขเว ภิกฺขุ วีมํสสมาธิปธานสงฺขารสมนฺนาคตํ อิทฺธิปาทํ ภาเวติ.\csromanspacer{} อยํ วุจฺจติ ภิกฺขเว อสงฺขตคามิมคฺโค ฯเปฯ. (17-20)}

\prose{กตโม จ ภิกฺขเว อสงฺขตคามิมคฺโค, อิธ ภิกฺขเว ภิกฺขุ สทฺธินฺทฺริยํ ภาเวติ วิเวกนิสฺสิตํ วิราคนิสฺสิตํ นิโรธนิสฺสิตํ โวสฺสคฺคปริณามิํ.\csromanspacer{} อยํ วุจฺจติ ภิกฺขเว อสงฺขตคามิมคฺโค ฯเปฯ.\csromanspacer{} กตโม จ ภิกฺขเว อสงฺขตคามิมคฺโค, อิธ ภิกฺขเว ภิกฺขุ วีริยินฺทฺริยํ ภาเวติ วิเวกนิสฺสิตํ ฯเปฯ.\csromanspacer{} อยํ วุจฺจติ ภิกฺขเว อสงฺขตคามิมคฺโค ฯเปฯ.\csromanspacer{} กตโม จ ภิกฺขเว อสงฺขตคามิมคฺโค, อิธ ภิกฺขเว ภิกฺขุ สตินฺทฺริยํ ภาเวติ ฯเปฯ.\csromanspacer{} อยํ วุจฺจติ ภิกฺขเว อสงฺขตคามิมคฺโค ฯเปฯ.\csromanspacer{} กตโม จ ภิกฺขเว อสงฺขตคามิมคฺโค, อิธ ภิกฺขเว ภิกฺขุ สมาธินฺทฺริยํ ภาเวติ ฯเปฯ.\csromanspacer{} อยํ วุจฺจติ ภิกฺขเว อสงฺขตคามิมคฺโค ฯเปฯ.\csromanspacer{} กตโม จ ภิกฺขเว อสงฺขตคามิมคฺโค, อิธ ภิกฺขเว ภิกฺขุ ปญฺญินฺทฺริยํ ภาเวติ วิเวกนิสฺสิตํ}

\gambhira{สฬายตนวคฺคสํยุตฺตปาฬิ}
\prose{\csromanfolio{540}วิราคนิสฺสิตํ นิโรธนิสฺสิตํ โวสฺสคฺคปริณามิํ.\csromanspacer{} อยํ วุจฺจติ ภิกฺขเว อสงฺขตคามิมคฺโค ฯเปฯ. (21-25)}

\prose{กตโม จ ภิกฺขเว อสงฺขตคามิมคฺโค, อิธ ภิกฺขเว ภิกฺขุ สทฺธาพลํ ภาเวติ วิเวกนิสฺสิตํ ฯเปฯ.\csromanspacer{} อยํ วุจฺจติ ภิกฺขเว อสงฺขตคามิมคฺโค ฯเปฯ กตโม จ ภิกฺขเว อสงฺขตคามิมคฺโค, อิธ ภิกฺขเว ภิกฺขุ วีริยพลํ ภาเวติ ฯเปฯ อยํ วุจฺจติ ภิกฺขเว อสงฺขตคามิมคฺโค ฯเปฯ.\csromanspacer{} กตโม จ ภิกฺขเว อสงฺขตคามิมคฺโค, อิธ ภิกฺขเว ภิกฺขุ สติพลํ ภาเวติ ฯเปฯ.\csromanspacer{} อยํ วุจฺจติ ภิกฺขเว อสงฺขตคามิมคฺโค ฯเปฯ.\csromanspacer{} กตโม จ ภิกฺขเว อสงฺขตคามิมคฺโค, อิธ ภิกฺขเว ภิกฺขุ สมาธิพลํ ภาเวติ ฯเปฯ.\csromanspacer{} อยํ วุจฺจติ ภิกฺขเว อสงฺขตคามิมคฺโค ฯเปฯ.\csromanspacer{} กตโม จ ภิกฺขเว อสงฺขตคามิมคฺโค, อิธ ภิกฺขเว ภิกฺขุ ปญฺญาพลํ ภาเวติ วิเวกนิสฺสิตํ วิราคนิสฺสิตํ นิโรธนิสฺสิตํ โวสฺสคฺคปริณามิํ.\csromanspacer{} อยํ วุจฺจติ ภิกฺขเว อสงฺขตคามิมคฺโค ฯเปฯ. (26-30)}

\prose{กตโม จ ภิกฺขเว อสงฺขตคามิมคฺโค, อิธ ภิกฺขเว ภิกฺขุ สติสมฺโพชฺฌงฺคํ ภาเวติ วิเวกนิสฺสิตํ ฯเปฯ.\csromanspacer{} อยํ วุจฺจติ ภิกฺขเว อสงฺขตคามิมคฺโค ฯเปฯ.\csromanspacer{} กตโม จ ภิกฺขเว อสงฺขตคามิมคฺโค, อิธ ภิกฺขเว ภิกฺขุ ธมฺมวิจยสมฺโพชฺฌงฺคํ ภาเวติ ฯเปฯ.\csromanspacer{} วีริยสมฺโพชฺฌงฺคํ ภาเวติ ฯเปฯ.\csromanspacer{} ปีติสมฺโพชฺฌงฺคํ ภาเวติ ฯเปฯ.\csromanspacer{} ปสฺสทฺธิสมฺโพชฺฌงฺคํ ภาเวติ ฯเปฯ.\csromanspacer{} สมาธิสมฺโพชฺฌงฺคํ ภาเวติ ฯเปฯ.\csromanspacer{} อุเปกฺขาสมฺโพชฺฌงฺคํ ภาเวติ วิเวกนิสฺสิตํ วิราคนิสฺสิตํ นิโรคนิสฺสิตํ โวสฺสคฺคปริณามิํ.\csromanspacer{} อยํ วุจฺจติ ภิกฺขเว อสงฺขตคามิมคฺโค ฯเปฯ. (31-37)}

\prose{กตโม จ ภิกฺขเว อสงฺขตคามิมคฺโค, อิธ ภิกฺขเว ภิกฺขุ สมฺมาทิฏฺฐิํ ภาเวติ วิเวกนิสฺสิตํ วิราคนิสฺสิตํ นิโรธนิสฺสิตํ โวสฺสคฺคปริณามิํ.\csromanspacer{} อยํ วุจฺจติ ภิกฺขเว อสงฺขตคามิมคฺโค ฯเปฯ.\csromanspacer{} กตโม จ ภิกฺขเว อสงฺขตคามิมคฺโค, อิธ ภิกฺขเว ภิกฺขุ สมฺมาสงฺกปฺปํ ภาเวติ ฯเปฯ.\csromanspacer{} สมฺมาวาจํ ภาเวติ ฯเปฯ.\csromanspacer{} สมฺมากมฺมนฺตํ ภาเวติ ฯเปฯ.\csromanspacer{} สมฺมา-อาชีวํ ภาเวติ ฯเปฯ.\csromanspacer{} สมฺมาวายามํ ภาเวติ ฯเปฯ.\csromanspacer{} สมฺมาสติํ ภาเวติ ฯเปฯ.\csromanspacer{} อสงฺขตญฺจ โว ภิกฺขเว เทเสสฺสามิ อสงฺขตคามิญฺจ มคฺคํ, ตํ สุณาถ.\csromanspacer{} กตมญฺจ ภิกฺขเว อสงฺขตํ ฯเปฯ.\csromanspacer{} กตโม จ ภิกฺขเว อสงฺขตคามิมคฺโค, อิธ ภิกฺขเว ภิกฺขุ สมฺมาสมาธิํ ภาเวติ วิเวกนิสฺสิตํ วิราคนิสฺสิตํ นิโรธนิสฺสิตํ}

\csromanheader{2}{9. อสงฺขตสํยุตฺต}
\prose{\csromanfolio{541}โวสฺสคฺคปริณามิํ.\csromanspacer{} อยํ วุจฺจติ ภิกฺขเว อสงฺขตคามิมคฺโค.\csromanspacer{} อิติ โข ภิกฺขเว เทสิตํ โว มยา อสงฺขตํ, เทสิโต อสงฺขตคามิมคฺโค.\csromanspacer{} ยํ ภิกฺขเว สตฺถารา กรณียํ สาวกานํ หิเตสินา อนุกมฺปเกน อนุกมฺปํ อุปาทาย, กตํ โว ตํ มยา.\csromanspacer{} เอตานิ ภิกฺขเว รุกฺขมูลานิ เอตานิ สุญฺญาคารานิ, ฌายถ ภิกฺขเว, มา ปมาทตฺถ, มา ปจฺฉา วิปฺปฏิสาริโน อหุวตฺถ, อยํ โว อมฺหากํ อนุสาสนีติ.\csromanspacer{}\csromanspacer{}ปฐมํ. (38- 45)}

\csromanheader{2}{2. อนตสุตฺต}
\proseitem{378}{อนตญฺจ โว ภิกฺขเว เทเสสฺสามิ อนตคามิญฺจ มคฺคํ, ตํ สุณาถ.\csromanspacer{} กตมญฺจ ภิกฺขเว อนตํ ฯเปฯ. (ยถา อสงฺขตํ, ตถา วิตฺถาเรตพฺพํ.).\csromanspacer{} ทุติยํ.}

\csromanheader{2}{3-32. อนาสวาทิสุตฺต}
\proseitem{379-408}{อนาสวญฺจ โว ภิกฺขเว เทเสสฺสามิ อนาสวคามิญฺจ มคฺคํ, ตํ สุณาถ.\csromanspacer{} กตมญฺจ ภิกฺขเว อนาสวํ ฯเปฯ.\csromanspacer{} สจฺจญฺจ โว ภิกฺขเว เทเสสฺสามิ สจฺจคามิญฺจ มคฺคํ, ตํ สุณาถ.\csromanspacer{} กตมญฺจ ภิกฺขเว สจฺจํ ฯเปฯ.\csromanspacer{} ปารญฺจ โว ภิกฺขเว เทเสสฺสามิ ปารคามิญฺจ มคฺคํ, ตํ สุณาถ.\csromanspacer{} กตมญฺจ ภิกฺขเว ปารํ ฯเปฯ.\csromanspacer{} นิปุณญฺจ โว ภิกฺขเว เทเสสฺสามิ นิปุณคามิญฺจ มคฺคํ, ตํ สุณาถ.\csromanspacer{} กตมญฺจ ภิกฺขเว นิปุณํ ฯเปฯ.\csromanspacer{} สุทุทฺทสญฺจ โว ภิกฺขเว เทเสสฺสามิ สุทุทฺทสคามิญฺจ มคฺคํ, ตํ สุณาถ.\csromanspacer{} กตมญฺจ ภิกฺขเว สุทุทฺทสํ ฯเปฯ.\csromanspacer{} อชชฺชรญฺจ โว ภิกฺขเว เทเสสฺสามิ อชชฺชรคามิญฺจ มคฺคํ, ตํ สุณาถ.\csromanspacer{} กตมญฺจ ภิกฺขเว อชชฺชรํ ฯเปฯ.\csromanspacer{} ธุวญฺจ โว ภิกฺขเว เทเสสฺสามิ ธุวคามิญฺจ มคฺคํ, ตํ สุณาถ.\csromanspacer{} กตมญฺจ ภิกฺขเว ธุวํ ฯเปฯ.\csromanspacer{} อปโลกิตญฺจ โว ภิกฺขเว เทเสสฺสามิ อปโลกิตคามิญฺจ มคฺคํ, ตํ สุณาถ.\csromanspacer{} กตมญฺจ ภิกฺขเว อปโลกิตํ ฯเปฯ.\csromanspacer{} อนิทสฺสนญฺจ โว ภิกฺขเว เทเสสฺสามิ อนิทสฺสนคามิญฺจ มคฺคํ, ตํ สุณาถ.\csromanspacer{} กตมญฺจ ภิกฺขเว อนิทสฺสนํ ฯเปฯ.\csromanspacer{} นิปฺปปญฺจญฺจ โว ภิกฺขเว เทเสสฺสามิ นิปฺปปญฺจคามิญฺจ มคฺคํ, ตํ สุณาถ.\csromanspacer{} กตมญฺจ ภิกฺขเว นิปฺปปญฺจํ ฯเปฯ.}

\prose{สนฺตญฺจ โว ภิกฺขเว เทเสสฺสามิ สนฺตคามิญฺจ มคฺคํ, ตํ สุณาถ.\csromanspacer{} กตมญฺจ ภิกฺขเว สนฺตํ ฯเปฯ.\csromanspacer{} อมตญฺจ โว ภิกฺขเว เทเสสฺสามิ อมตคามิญฺจ มคฺคํ, ตํ สุณาถ.\csromanspacer{} กตมญฺจ ภิกฺขเว อมตํ ฯเปฯ.\csromanspacer{} ปณีตญฺจ โว ภิกฺขเว เทเสสฺสามิ ปณีตคามิญฺจ มคฺคํ, ตํ สุณาถ.\csromanspacer{} กตมญฺจ ภิกฺขเว ปณีตํ ฯเปฯ.\csromanspacer{} สิวญฺจ โว}

\gambhira{สฬายตนวคฺคสํยุตฺตปาฬิ}
\prose{\csromanfolio{542}ภิกฺขเว เทเสสฺสามิ สิวคามิญฺจ มคฺคํ, ตํ สุณาถ.\csromanspacer{} กตมญฺจ ภิกฺขเว สิวํ ฯเปฯ.\csromanspacer{} เขมญฺจ โว ภิกฺขเว เทเสสฺสามิ เขมคามิญฺจ มคฺคํ, ตํ สุณาถ.\csromanspacer{} กตมญฺจ ภิกฺขเว เขมํ ฯเปฯ.\csromanspacer{} ตณฺหากฺขยญฺจ โว ภิกฺขเว เทเสสฺสามิ ตณฺหากฺขยคามิญฺจ มคฺคํ, ตํ สุณาถ.\csromanspacer{} กตมญฺจ ภิกฺขเว ตณฺหากฺขยํ ฯเปฯ.}

\prose{อจฺฉริยญฺจ โว ภิกฺขเว เทเสสฺสามิ อจฺฉริยคามิญฺจ มคฺคํ, ตํ สุณาถ.\csromanspacer{} กตมญฺจ ภิกฺขเว อจฺฉริยํ ฯเปฯ.\csromanspacer{} อพฺภุตญฺจ โว ภิกฺขเว เทเสสฺสามิ อพฺภุตคามิญฺจ มคฺคํ, ตํ สุณาถ.\csromanspacer{} กตมญฺจ ภิกฺขเว อพฺภุตํ ฯเปฯ.\csromanspacer{} อนีติกญฺจ โว ภิกฺขเว เทเสสฺสามิ อนีติกคามิญฺจ มคฺคํ, ตํ สุณาถ.\csromanspacer{} กตมญฺจ ภิกฺขเว อนีติกํ ฯเปฯ.\csromanspacer{} อนีติกธมฺมญฺจ โว ภิกฺขเว เทเสสฺสามิ อนีติกธมฺมคามิญฺจ มคฺคํ, ตํ สุณาถ.\csromanspacer{} กตมญฺจ ภิกฺขเว อนีติกธมฺมํ ฯเปฯ.\csromanspacer{} นิพฺพานญฺจ โว ภิกฺขเว เทเสสฺสามิ นิพฺพานคามิญฺจ มคฺคํ, ตํ สุณาถ.\csromanspacer{} กตมญฺจ ภิกฺขเว นิพฺพานํ ฯเปฯ.\csromanspacer{} อพฺยาพชฺฌญฺจ\footnote{อพฺยาปชฺฌญฺจ (สี, สฺยา, กํ, อิ)} โว ภิกฺขเว เทเสสฺสามิ อพฺยาพชฺฌคามิญฺจ มคฺคํ, ตํ สุณาถ.\csromanspacer{} กตมญฺจ ภิกฺขเว อพฺยาพชฺฌํ ฯเปฯ.\csromanspacer{} วิราคญฺจ โว ภิกฺขเว เทเสสฺสามิ วิราคคามิญฺจ มคฺคํ, ตํ สุณาถ.\csromanspacer{} กตโม จ ภิกฺขเว วิราโค ฯเปฯ.}

\prose{สุทฺธิญฺจ โว ภิกฺขเว เทเสสฺสามิ สุทฺธิคามิญฺจ มคฺคํ, ตํ สุณาถ.\csromanspacer{} กตมา จ ภิกฺขเว สุทฺธิ ฯเปฯ.\csromanspacer{} มุตฺติญฺจ โว ภิกฺขเว เทเสสฺสามิ มุตฺติคามิญฺจ มคฺคํ, ตํ สุณาถ.\csromanspacer{} กตมา จ ภิกฺขเว มุตฺติ ฯเปฯ.\csromanspacer{} อนาลยญฺจ โว ภิกฺขเว เทเสสฺสามิ อนาลยคามิญฺจ มคฺคํ, ตํ สุณาถ.\csromanspacer{} กตโม จ ภิกฺขเว อนาลโย ฯเปฯ.\csromanspacer{} ทีปญฺจ โว ภิกฺขเว เทเสสฺสามิ ทีปคามิญฺจ มคฺคํ, ตํ สุณาถ.\csromanspacer{} กตมญฺจ ภิกฺขเว ทีปํ ฯเปฯ.\csromanspacer{} เลณญฺจ โว ภิกฺขเว เทเสสฺสามิ เลณคามิญฺจ มคฺคํ, ตํ สุณาถ.\csromanspacer{} กตมญฺจ ภิกฺขเว เลณํ ฯเปฯ.\csromanspacer{} ตาณญฺจ โว ภิกฺขเว เทเสสฺสามิ ตาณคามิญฺจ มคฺคํ, ตํ สุณาถ.\csromanspacer{} กตมญฺจ ภิกฺขเว ตาณํ ฯเปฯ.\csromanspacer{} สรณญฺจ โว ภิกฺขเว สรณคามิญฺจ มคฺคํ, ตํ สุณาถ.\csromanspacer{} กตมญฺจ ภิกฺขเว สรณํ ฯเปฯ.\csromanspacer{}\csromanspacer{}พาตฺติํสติมํ.}

\csromanheader{2}{33. ปรายนสุตฺต}
\proseitem{409}{ปรายนญฺจ\footnote{ปรายณญฺจ (อิ, สี-ฏฺฐ)} โว ภิกฺขเว เทเสสฺสามิ ปรายนคามิญฺจ มคฺคํ, ตํ สุณาถ.\csromanspacer{} กตมญฺจ ภิกฺขเว ปรายนํ, โย ภิกฺขเว ราคกฺขโย}

\csromanheader{2}{9. อสงฺขตสํยุตฺต}
\prose{\csromanfolio{543}โทสกฺขโย โมหกฺขโย.\csromanspacer{} อิทํ วุจฺจติ ภิกฺขเว ปรายนํ.\csromanspacer{} กตโม จ ภิกฺขเว ปรายนคามี มคฺโค, กายคตาสติ.\csromanspacer{} อยํ วุจฺจติ ภิกฺขเว ปรายนคามิมคฺโค.\csromanspacer{} อิติ โข ภิกฺขเว เทสิตํ โว มยา ปรายนํ เทสิโต ปรายนคามิมคฺโค.\csromanspacer{} ยํ ภิกฺขเว สตฺถารา กรณียํ สาวกานํ หิเตสินา อนุกมฺปเกน อนุกมฺปํ อุปาทาย, กตํ โว ตํ มยา.\csromanspacer{} เอตานิ ภิกฺขเว รุกฺขมูลานิ, เอตานิ สุญฺญาคารานิ, ฌายถา ภิกฺขเว, มา ปมาทตฺถ, มา ปจฺฉา วิปฺปฏิสาริโน อหุวตฺถ, อยํ โว อมฺหากํ อนุสาสนีติ. (ยถา อสงฺขตํ, ตถา วิตฺถาเรตพฺพํ.) .\csromanspacer{} เตตฺติํสติมํ.\csromanspacer{} ทุติโย วคฺโค.}

\titlehead{\textbf{ตสฺสุทฺทานํ}}
\csromangathameasure{อสงฺขตํ อนตํ อนาสวํ, \\ สจฺจญฺจ ปารํ นิปุณํ สุทุทฺทสํ. \\ อชชฺชรํ ธุวํ อปโลกิตํ, \\ อนิทสฺสนํ นิปฺปปญฺจ สนฺตํ. \\ อมตํ ปณีตญฺจ สิวญฺจ เขมํ, \\ ตณฺหากฺขโย อจฺฉริยญฺจ อพฺภุตํ. \\ อนีติกํ อนีติกธมฺมํ, \\ นิพฺพานเมตํ สุคเตน เทสิตํ. \\ อพฺยาพชฺโฌ วิราโค จ, สุทฺธิ มุตฺติ อนาลโย. \\ ทีโป เลณญฺจ ตาณญฺจ, สรณญฺจ ปรายนนฺติ. \\ \textbf{อสงฺขตสํยุตฺตํ สมตฺตํ.}}
\csromangathagroup{\csromangathastanza{อสงฺขตํ อนตํ อนาสวํ, \\ สจฺจญฺจ ปารํ นิปุณํ สุทุทฺทสํ. \\ อชชฺชรํ ธุวํ อปโลกิตํ, \\ อนิทสฺสนํ นิปฺปปญฺจ สนฺตํ.}\csromangathastanza{อมตํ ปณีตญฺจ สิวญฺจ เขมํ, \\ ตณฺหากฺขโย อจฺฉริยญฺจ อพฺภุตํ. \\ อนีติกํ อนีติกธมฺมํ, \\ นิพฺพานเมตํ สุคเตน เทสิตํ. \\ อพฺยาพชฺโฌ วิราโค จ, สุทฺธิ มุตฺติ อนาลโย. \\ ทีโป เลณญฺจ ตาณญฺจ, สรณญฺจ ปรายนนฺติ.}\csromangathastanza{\textbf{อสงฺขตสํยุตฺตํ สมตฺตํ.}}}

\csromanheader{1}{10. อพฺยากตสํยุตฺต}
\csromanmatikaanchor{mtk.288}
\csromantocmarknopage{section}{10. อพฺยากตสํยุตฺต}{mtk.288}
\bookmark[dest={mtk.288},level=1]{10. อพฺยากตสํยุตฺต}
\csromanheader{2}{1. เขมาสุตฺต}
\csromanmatikaanchor{mtk.289}
\csromantocmarkref{subsection}{1. เขมาสุตฺต}{mtk.289}
\bookmark[dest={mtk.289},level=2]{1. เขมาสุตฺต}
\proseitem{410}{\csromanfolio{544}เอกํ สมยํ ภควา สาวตฺถิยํ วิหรติ เชตวเน อนาถปิณฺฑิกสฺส อาราเม.\csromanspacer{} เตน โข ปน สมเยน เขมา ภิกฺขุนี โกสเลสุ จาริกํ จรมานา อนฺตรา จ สาวตฺถิํ อนฺตรา จ สาเกตํ โตรณวตฺถุสฺมิํ วาสํ อุปคตา โหติ.\csromanspacer{} อถ โข ราชา ปเสนทิ โกสโล สาเกตา สาวตฺถิํ คจฺฉนฺโต อนฺตรา จ สาเกตํ อนฺตรา จ สาวตฺถิํ โตรณวตฺถุสฺมิํ เอกรตฺติวาสํ อุปคจฺฉิ.\csromanspacer{} อถ โข ราชา ปเสนทิ โกสโล อญฺญตรํ ปุริสํ อามนฺเตสิ “เอหิ ตฺวํ อมฺโภ ปุริส โตรณวตฺถุสฺมิํ ตถารูปํ สมณํ วา พฺราหฺมณํ วา ชาน, ยมหํ อชฺช ปยิรุปาเสยฺยนฺ”ติ.}

\prose{“เอวํ เทวา”ติ โข โส ปุริโส รญฺโญ ปเสนทิสฺส โกสลสฺส ปฏิสฺสุตฺวา เกวลกปฺปํ โตรณวตฺถุํ อาหิณฺฑนฺโต\footnote{อนฺวาหิณฺฑนฺโต (สี)} นาทฺทส ตถารูปํ สมณํ วา พฺราหฺมณํ วา, ยํ ราชา ปเสนทิ โกสโล ปยิรุปาเสยฺย.\csromanspacer{} อทฺทสา โข โส ปุริโส เขมํ ภิกฺขุนิํ โตรณวตฺถุสฺมิํ วาสํ อุปคตํ, ทิสฺวาน เยน ราชา ปเสนทิ โกสโล เตนุปสงฺกมิ, อุปสงฺกมิตฺวา ราชานํ ปเสนทิํ โกสลํ เอตทโวจ\csromandash{}}

\prose{“นตฺถิ โข เทว โตรณวตฺถุสฺมิํ ตถารูโป สมโณ วา พฺราหฺมโณ วา, ยํ เทโว ปยิรุปาเสยฺย, อตฺถิ จ โข เทว เขมา นาม ภิกฺขุนี ตสฺส ภควโต สาวิกา อรหโต สมฺมาสมฺพุทฺธสฺส, ตสฺสา โข ปน อยฺยาย เอวํ กลฺยาโณ กิตฺติสทฺโท อพฺภุคฺคโต ‘ปณฺฑิตา วิยตฺตา เมธาวินี พหุสฺสุตา จิตฺตกถา กลฺยาณปฏิภานา’ติ.\csromanspacer{} ตํ เทโว ปยิรุปาสตู”ติ.}

\prose{อถ โข ราชา ปเสนทิ โกสโล เยน เขมา ภิกฺขุนี เตนุปสงฺกมิ, อุปสงฺกมิตฺวา เขมํ ภิกฺขุนิํ อภิวาเทตฺวา เอกมนฺตํ นิสีทิ,}

\csromanheader{2}{10. อพฺยากตสํยุตฺต}
\prose{\csromanfolio{545}เอกมนฺตํ นิสินฺโน โข ราชา ปเสนทิ โกสโล เขมํ ภิกฺขุนิํ เอตทโวจ “กิํ นุ โข อยฺเย โหติ ตถาคโต ปรํ มรณา”ติ.\csromanspacer{} อพฺยากตํ โข เอตํ มหาราช ภควตา “โหติ ตถาคโต ปรํ มรณา”ติ.\csromanspacer{} กิํ ปนยฺเย น โหติ ตถาคโต ปรํ มรณาติ.\csromanspacer{} เอตมฺปิ โข มหาราช อพฺยากตํ ภควตา “น โหติ ตถาคโต ปรํ มรณา”ติ.\csromanspacer{} กิํ นุ โข อยฺเย โหติ จ น จ โหติ ตถาคโต ปรํ มรณาติ.\csromanspacer{} อพฺยากตํ โข เอตํ มหาราช ภควตา “โหติ จ น จ โหติ ตถาคโต ปรํ มรณา”ติ.\csromanspacer{} กิํ ปนยฺเย เนว โหติ น น โหติ ตถาคโต ปรํ มรณาติ.\csromanspacer{} เอตมฺปิ โข มหาราช อพฺยากตํ ภควตา “เนว โหติ น น โหติ ตถาคโต ปรํ มรณา”ติ.}

\prose{กิํ นุ โข อยฺเย โหติ ตถาคโต ปรํ มรณาติ อิติ ปุฏฺฐาสมานา อพฺยากตํ โข เอตํ มหาราช ภควตา “โหติ ตถาคโต ปรํ มรณา”ติ วเทสิ.\csromanspacer{} กิํ ปนยฺเย น โหติ ตถาคโต ปรํ มรณาติ อิติ ปุฏฺฐา สมานา เอตมฺปิ โข มหาราช อพฺยากตํ ภควตา “น โหติ ตถาคโต ปรํ มรณา”ติ วเทสิ.\csromanspacer{} กิํ นุ โข อยฺเย โหติ จ น จ โหติ ตถาคโต ปรํ มรณาติ อิติ ปุฏฺฐา สมานา อพฺยากตํ โข เอตํ มหาราช ภควตา “โหติ จ น จ โหติ ตถาคโต ปรํ มรณา”ติ วเทสิ.\csromanspacer{} กิํ ปนยฺเย เนว โหติ น น โหติ ตถาคโต ปรํ มรณาติ อิติ ปุฏฺฐา สมานา เอตมฺปิ โข มหาราช อพฺยากตํ ภควตา “เนว โหติ น น โหติ ตถาคโต ปรํ มรณา”ติ วเทสิ.\csromanspacer{} โก นุ โข อยฺเย เหตุ โก ปจฺจโย, เยเนตํ อพฺยากตํ ภควตาติ.}

\prose{เตน หิ มหาราช ตญฺเญเวตฺถ ปฏิปุจฺฉิสฺสามิ, ยถา เต ขเมยฺย, ตถา นํ พฺยากเรยฺยาสิ.\csromanspacer{} ตํ กิํ มญฺญสิ มหาราช, อตฺถิ เต โกจิ คณโก วา มุทฺทิโก วา สงฺขายโก วา, โย ปโหติ คงฺคาย วาลุกํ\footnote{วาลิกํ (สี, สฺยา, กํ)} คเณตุํ “เอตฺตกา\footnote{เอตฺติกา (สี)} วาลุกา” อิติ วา, “เอตฺตกานิ วาลุกสตานิ” อิติ วา, “เอตฺตกานิ วาลุกสหสฺสานิ” อิติ วา, “เอตฺตกานิ วาลุกสตสหสฺสานิ” อิติ วาติ.\csromanspacer{} โน เหตํ อยฺเย.\csromanspacer{} อตฺถิ ปน เต โกจิ คณโก วา มุทฺทิโก วา สงฺขายโก วา, โย ปโหติ}

\gambhira{สฬายตนวคฺคสํยุตฺตปาฬิ}
\prose{\csromanfolio{546}มหาสมุทฺเท อุทกํ คเณตุํ “เอตฺตกานิ อุทกาฬฺหกานิ” อิติ วา, “เอตฺตกานิ อุทกาฬฺหกสตานิ” อิติ วา, “เอตฺตกานิ อุทกาฬฺหกสหสฺสานิ” อิติ วา, “เอตฺตกานิ อุทกาฬฺหกสตสหสฺสานิ” อิติ วาติ.\csromanspacer{} โน เหตํ อยฺเย.\csromanspacer{} ตํ กิสฺส เหตุ, มหายฺเย สมุทฺโท คมฺภีโร อปฺปเมยฺโย ทุปฺปริโยคาโหติ.\csromanspacer{} เอวเมว โข มหาราช เยน รูเปน ตถาคตํ ปญฺญาปยมาโน ปญฺญาเปยฺย, ตํ รูปํ ตถาคตสฺส ปหีนํ อุจฺฉินฺนมูลํ ตาลาวตฺถุกตํ อนภาวํกตํ อายติํ อนุปฺปาทธมฺมํ, รูปสงฺขยวิมุตฺโต โข มหาราช ตถาคโต คมฺภีโร อปฺปเมยฺโย ทุปฺปริโยคาโห เสยฺยถาปิ มหาสมุทฺโท, “โหติ ตถาคโต ปรํ มรณา”ติปิ น อุเปติ, “น โหติ ตถาคโต ปรํ มรณา”ติปิ น อุเปติ, “โหติ จ น จ โหติ ตถาคโต ปรํ มรณา”ติปิ น อุเปติ, “เนว โหติ น น โหติ ตถาคโต ปรํ มรณา”ติปิ น อุเปติ.}

\prose{ยาย เวทนาย ตถาคตํ ปญฺญาปยมาโน ปญฺญาเปยฺย, สา เวทนา ตถาคตสฺส ปหีนา อุจฺฉินฺนมูลา ตาลาวตฺถุกตา อนภาวํกตา อายติํ อนุปฺปาทธมฺมา, เวทนาสงฺขายวิมุตฺโต มหาราช ตถาคโต คมฺภีโร อปฺปเมยฺโย ทุปฺปริโยคาโห เสยฺยถาปิ มหาสมุทฺโท, “โหติ ตถาคโต ปรํ มรณา”ติปิ น อุเปติ, “น โหติ ตถาคโต ปรํ มรณา”ติปิ น อุเปติ, “โหติ จ น จ โหติ ตถาคโต ปรํ มรณา”ติปิ น อุเปติ, “เนว โหติ น น โหติ ตถาคโต ปรํ มรณา”ติปิ น อุเปติ.}

\prose{ยาย สญฺญาย ตถาคตํ ฯเปฯ.\csromanspacer{} เยหิ สงฺขาเรหิ ตถาคตํ ปญฺญาปยมาโน ปญฺญาเปยฺย, เต สงฺขารา ตถาคตสฺส ปหีนา อุจฺฉินฺนมูลา ตาลาวตฺถุกตา อนภาวํกตา อายติํ อนุปฺปาทธมฺมา, สงฺขารสงฺขายวิมุตฺโต โข มหาราช ตถาคโต คมฺภีโร อปฺปเมยฺโย ทุปฺปริโยคาโห เสยฺยถาปิ มหาสมุทฺโท, “โหติ ตถาคโต ปรํ มรณา”ติปิ น อุเปติ, “น โหติ ตถาคโต ปรํ มรณา”ติปิ น อุเปติ, “เนว โหติ น น โหติ ตถาคโต ปรํ มรณา”ติปิ น อุเปติ.}

\prose{เยน วิญฺญาเณน ตถาคตํ ปญฺญาปยมาโน ปญฺญาเปยฺย, ตํ วิญฺญาณํ ตถาคตสฺส ปหีนํ อุจฺฉินฺนมูลํ ตาลาวตฺถุกตํ อนภาวํกตํ}

\csromanheader{2}{10. อพฺยากตสํยุตฺต}
\prose{\csromanfolio{547}อายติํ อนุปฺปาทธมฺมํ, วิญฺญาณสงฺขายวิมุตฺโต โข มหาราช ตถาคโต คมฺภีโร อปฺปเมยฺโย ทุปฺปริโยคาโห เสยฺยถาปิ มหาสมุทฺโท, “โหติ ตถาคโต ปรํ มรณา”ติปิ น อุเปติ, “น โหติ ตถาคโต ปรํ มรณา”ติปิ น อุเปติ, “โหติ จ น จ โหติ ตถาคโต ปรํ มรณา”ติปิ น อุเปติ, “เนว โหติ น น โหติ ตถาคโต ปรํ มรณา”ติปิ น อุเปตีติ.\csromanspacer{} อถ โข ราชา ปเสนทิ โกสโล เขมาย ภิกฺขุนิยา ภาสิตํ อภินนฺทิตฺวา อนุโมทิตฺวา อุฏฺฐายาสนา เขมํ ภิกฺขุนิํ อภิวาเทตฺวา ปทกฺขิณํ กตฺวา ปกฺกามิ.}

\prose{อถ โข ราชา ปเสนทิ โกสโล อปเรน สมเยน เยน ภควา เตนุปสงฺกมิ, อุปสงฺกมิตฺวา ภควนฺตํ อภิวาเทตฺวา เอกมนฺตํ นิสีทิ, เอกมนฺตํ นิสินฺโน โข ราชา ปเสนทิ โกสโล ภควนฺตํ เอตทโวจ “กิํ นุ โข ภนฺเต โหติ ตถาคโต ปรํ มรณา”ติ.\csromanspacer{} อพฺยากตํ โข เอตํ มหาราช มยา “โหติ ตถาคโต ปรํ มรณา”ติ.\csromanspacer{} กิํ ปน ภนฺเต น โหติ ตถาคโต ปรํ มรณาติ.\csromanspacer{} เอตมฺปิ โข มหาราช อพฺยากตํ มยา “น โหติ ตถาคโต ปรํ มรณา”ติ.\csromanspacer{} กิํ นุ โข ภนฺเต โหติ จ น จ โหติ ตถาคโต ปรํ มรณาติ.\csromanspacer{} อพฺยากตํ โข เอตํ มหาราช มยา “โหติ จ น จ โหติ ตถาคโต ปรํ มรณา”ติ.\csromanspacer{} กิํ ปน ภนฺเต เนว โหติ น น โหติ ตถาคโต ปรํ มรณาติ.\csromanspacer{} เอตมฺปิ โข มหาราช อพฺยากตํ มยา “เนว โหติ น น โหติ ตถาคโต ปรํ มรณา”ติ.\csromanspacer{} กิํ นุ โข ภนฺเต โหติ ตถาคโต ปรํ มรณาติ อิติ ปุฏฺโฐ สมาโน อพฺยากตํ โข เอตํ มหาราช มยา “โหติ ตถาคโต ปรํ มรณา”ติ วเทสิ ฯเปฯ.\csromanspacer{} กิํ ปน ภนฺเต เนว โหติ น น โหติ ตถาคโต ปรํ มรณาติ อิติ ปุฏฺโฐ สมาโน เอตมฺปิ โข มหาราช อพฺยากตํ มยา “เนว โหติ น น โหติ ตถาคโต ปรํ มรณา”ติ วเทสิ.\csromanspacer{} โก นุ โข ภนฺเต เหตุ โก ปจฺจโย, เยเนตํ อพฺยากตํ ภควตาติ.}

\prose{เตน หิ มหาราช ตญฺเญเวตฺถ ปฏิปุจฺฉิสฺสามิ, ยถา เต ขเมยฺย, ตถา นํ พฺยากเรยฺยาสิ.\csromanspacer{} ตํ กิํ มญฺญสิ มหาราช, อตฺถิ เต โกจิ คณโก วา มุทฺทิโก วา สงฺขายโก วา, โย ปโหติ คงฺคาย วาลุกํ}

\gambhira{สฬายตนวคฺคสํยุตฺตปาฬิ}
\prose{\csromanfolio{548}คเณตุํ “เอตฺตกา วาลุกา” อิติ วา ฯเปฯ “เอตฺตกานิ วาลุกสตสหสฺสานิ” อิติ วาติ.\csromanspacer{} โน เหตํ ภนฺเต.\csromanspacer{} อตฺถิ ปน เต โกจิ คณโก วา มุทฺทิโก วา สงฺขายโก วา, โย ปโหติ มหาสมุทฺเท อุทกํ คเณตุํ “เอตฺตกานิ อุทกาฬฺหกานิ” อิติ วา ฯเปฯ “เอตฺตกานิ อุทกาฬฺหกสตสหสฺสานิ” อิติ วาติ.\csromanspacer{} โน เหตํ ภนฺเต.\csromanspacer{} ตํ กิสฺส เหตุ, มหา ภนฺเต สมุทฺโท คมฺภีโร อปฺปเมยฺโย ทุปฺปริโยคาโห.\csromanspacer{} เอวเมว โข มหาราช เยน รูเปน ตถาคตํ ปญฺญปยมาโน ปญฺญาเปยฺย, ตํ รูปํ ตถาคตสฺส ปหีนํ อุจฺฉินฺนมูลํ ตาลาวตฺถุกตํ อนภาวํกตํ อายติํ อนุปฺปาทธมฺมํ, รูปสงฺขายวิมุตฺโต โข มหาราช ตถาคโต คมฺภีโร อปฺปเมยฺโย ทุปฺปริโยคาโห เสยฺยถาปิ มหาสมุทฺโท, “โหติ ตถาคโต ปรํ มรณา”ติปิ น อุเปติ ฯเปฯ “เนว โหติ น น โหติ ตถาคโต ปรํ มรณา”ติปิ น อุเปติ.\csromanspacer{} ยาย เวทนาย.\csromanspacer{} ยาย สญฺญาย.\csromanspacer{} เยหิ สงฺขาเรหิ.}

\prose{เยน วิญฺญาเณน ตถาคตํ ปญฺญาปยมาโน ปญฺญาเปยฺย, ตํ วิญฺญาณํ ตถาคตสฺส ปหีนํ อุจฺฉินฺนมูลํ ตาลาวตฺถุกตํ อนภาวํกตํ อายติํ อนุปฺปาทธมฺมํ, วิญฺญาณสงฺขายวิมุตฺโต โข มหาราช ตถาคโต คมฺภีโร อปฺปเมยฺโย ทุปฺปริโยคาโห เสยฺยถาปิ มหาสมุทฺโท, “โหติ ตถาคโต ปรํ มรณา”ติปิ น อุเปติ, “น โหติ ตถาคโต ปรํ มรณา”ติปิ น อุเปติ, “โหติ น จ น โหติ ตถาคโต ปรํ มรณา”ติปิ น อุเปติ, “เนว โหติ น น โหติ ตถาคโต ปรํ มรณา”ติปิ น อุเปตีติ.}

\prose{อจฺฉริยํ ภนฺเต, อพฺภุตํ ภนฺเต, ยตฺร หิ นาม สตฺถุ เจว\footnote{สตฺถุโน เจว (สี)} สาวิกาย จ อตฺเถน อตฺโถ พฺยญฺชเนน พฺยญฺชนํ สํสนฺทิสฺสติ สเมสฺสติ น วิโรธยิสฺสติ\footnote{วิหายิสฺสติ (สี, สฺยา, กํ), วิคายิสฺสติ (ก)}, ยทิทํ อคฺคปทสฺมิํ.\csromanspacer{} เอกมิทาหํ ภนฺเต สมยํ เขมํ ภิกฺขุนิํ อุปสงฺกมิตฺวา เอตมตฺถํ อปุจฺฉิํ, สาปิ เม อยฺยา เอเตหิ ปเทหิ เอเตหิ พฺยญฺชเนหิ เอตมตฺถํ พฺยากาสิ เสยฺยถาปิ ภควา.\csromanspacer{} อจฺฉริยํ ภนฺเต, อพฺภุตํ ภนฺเต, ยตฺร หิ นาม สตฺถุ เจว สาวิกาย จ อตฺเถน อตฺโถ พฺยญฺชเนน พฺยญฺชนํ สํสนฺทิสฺสติ สเมสฺสติ น วิโรธยิสฺสติ, ยทิทํ อคฺคปทสฺมิํ.\csromanspacer{} หนฺท ทานิ มยํ ภนฺเต คจฺฉาม, พหุกิจฺจา มยํ พหุกรณียาติ.\csromanspacer{} ยสฺส ทานิ ตฺวํ มหาราช กาลํ มญฺญสีติ.\csromanspacer{} อถ โข ราชา ปเสนทิ โกสโล}

\csromanheader{2}{10. อพฺยากตสํยุตฺต}
\prose{\csromanfolio{549}ภควโต ภาสิตํ อภินนฺทิตฺวา อนุโมทิตฺวา อุฏฺฐายาสนา ภควนฺตํ อภิวาเทตฺวา ปทกฺขิณํ กตฺวา ปกฺกามีติ.\csromanspacer{}\csromanspacer{}ปฐมํ.}

\csromanheader{2}{2. อนุราธสุตฺต}
\csromanmatikaanchor{mtk.290}
\csromantocmarkref{subsection}{2. อนุราธสุตฺต}{mtk.290}
\bookmark[dest={mtk.290},level=2]{2. อนุราธสุตฺต}
\proseitem{411}{เอกํ สมยํ ภควา เวสาลิยํ วิหรติ มหาวเน กูฏาคารสาลายํ.\csromanspacer{} เตน โข ปน สมเยน อายสฺมา อนุราโธ ภควโต อวิทูเร อรญฺญกุฏิกายํ วิหรติ.\csromanspacer{} อถ โข สมฺพหุลา อญฺญติตฺถิยา ปริพฺพาชกา เยนายสฺมา อนุราโธ เตนุปสงฺกมิํสุ, อุปสงฺกมิตฺวา อายสฺมตา อนุราเธน สทฺธิํ สมฺโมทิํสุ, สมฺโมทนียํ กถํ สารณียํ วีติสาเรตฺวา เอกมนฺตํ นิสีทิํสุ, เอกมนฺตํ นิสินฺนา โข เต อญฺญติตฺถิยา ปริพฺพาชกา อายสฺมนฺตํ อนุราธํ เอตทโวจุํ “โย โส อาวุโส อนุราธ ตถาคโต อุตฺตมปุริโส ปรมปุริโส ปรมปตฺติปตฺโต, ตํ ตถาคโต อิเมสุ จตูสุ ฐาเนสุ ปญฺญาปยมาโน ปญฺญาเปติ ‘โหติ ตถาคโต ปรํ มรณา’ติ วา, ‘น โหติ ตถาคโต ปรํ มรณา’ติ วา, ‘โหติ จ น จ โหติ ตถาคโต ปรํ มรณา’ติ วา, ‘เนว โหติ น น โหติ ตถาคโต ปรํ มรณา’ติ วา”ติ.\csromanspacer{} โย โส อาวุโส ตถาคโต อุตฺตมปุริโส ปรมปุริโส ปรมปตฺติปตฺโต, ตํ ตถาคโต อญฺญตฺร อิเมหิ จตูหิ ฐาเนหิ ปญฺญาปยมาโน ปญฺญาเปติ “โหติ ตถาคโต ปรํ มรณา”ติ วา, “น โหติ ตถาคโต ปรํ มรณา”ติ วา, “โหติ จ น จ โหติ ตถาคโต ปรํ มรณา”ติ วา, “เนว โหติ น น โหติ ตถาคโต ปรํ มรณา”ติ วาติ.\csromanspacer{} เอวํ วุตฺเต เต อญฺญติตฺถิยา ปริพฺพาชกา อายสฺมนฺตํ อนุราธํ เอตทโวจุํ “โส จายํ\footnote{โย จายํ (สี)} ภิกฺขุ นโว ภวิสฺสติ อจิรปพฺพชิโต, เถโร วา ปน พาโล อพฺยตฺโต”ติ.\csromanspacer{} อถ โข เต อญฺญติตฺถิยา ปริพฺพาชกา อายสฺมนฺตํ อนุราธํ นววาเทน จ พาลวาเทน จ อปสาเทตฺวา อุฏฺฐายาสนา ปกฺกมิํสุ.}

\prose{อถ โข อายสฺมโต อนุราธสฺส อจิรปกฺกนฺเตสุ อญฺญติตฺถิเยสุ ปริพฺพาชเกสุ เอตทโหสิ “สเจ โข มํ เต อญฺญติตฺถิยา ปริพฺพาชกา อุตฺตริํ ปุจฺเฉยฺยุํ, กถํ พฺยากรมาโน นุ ขฺวาหํ เตสํ อญฺญติตฺถิยานํ ปริพฺพาชกานํ, วุตฺตวาที เจว ภควโต อสฺสํ,}

\gambhira{สฬายตนวคฺคสํยุตฺตปาฬิ}
\prose{\csromanfolio{550}น จ ภควนฺตํ อภูเตน อพฺภาจิกฺเขยฺยํ, ธมฺมสฺส จานุธมฺมํ พฺยากเรยฺยํ, น จ โกจิ สหธมฺมิโก วาทานุวาโท คารยฺหํ ฐานํ อาคจฺเฉยฺยา”ติ.\csromanspacer{} อถ โข อายสฺมา อนุราโธ เยน ภควา เตนุปสงฺกมิ, อุปสงฺกมิตฺวา ภควนฺตํ อภิวาเทตฺวา เอกมนฺตํ นิสีทิ, เอกมนฺตํ นิสินฺโน โข อายสฺมา อนุราโธ ภควนฺตํ เอตทโวจ “อิธาหํ ภนฺเต ภควโต อวิทูเร อรญฺญกุฏิกายํ วิหรามิ, อถ โข ภนฺเต สมฺพหุลา อญฺญติตฺถิยา ปริพฺพาชกา เยนาหํ เตนุปสงฺกมิํสุ, อุปสงฺกมิตฺวา มยา สทฺธิํ สมฺโมทิํสุ, สมฺโมทนียํ กถํ สารณียํ วีติสาเรตฺวา เอกมนฺตํ นิสีทิํสุ, เอกมนฺตํ นิสินฺนา โข ภนฺเต เต อญฺญติตฺถิยา ปริพฺพาชกา มํ เอตทโวจุํ “โย โส อาวุโส อนุราธ ตถาคโต อุตฺตมปุริโส ปรมปุริโส ปรมปตฺติปตฺโต, ตํ ตถาคโต อิเมสุ จตูสุ ฐาเนสุ ปญฺญาปยมาโน ปญฺญาเปติ ‘โหติ ตถาคโต ปรํ มรณา’ติ วา ฯเปฯ ‘เนว โหติ น น โหติ ตถาคโต ปรํ มรณา’ติ วา”ติ.\csromanspacer{} เอวํ วุตฺตาหํ ภนฺเต เต อญฺญติตฺถิเย ปริพฺพาชเก เอตทโวจํ “โย โส อาวุโส ตถาคโต อุตฺตมปุริโส ปรมปุริโส ปรมปตฺติปตฺโต, ตํ ตถาคโต อญฺญตฺร อิเมหิ จตูหิ ฐาเนหิ ปญฺญปยมาโน ปญฺญาเปติ ‘โหติ ตถาคโต ปรํ มรณา’ติ วา ฯเปฯ ‘เนว โหติ น น โหติ ตถาคโต ปรํ มรณา’ติ วา”ติ.\csromanspacer{} เอวํ วุตฺเต ภนฺเต เต อญฺญติตฺถิยา ปริพฺพาชกา มํ เอตทโวจุํ “โส จายํ ภิกฺขุ นโว ภวิสฺสติ อจิรปพฺพชิโต, เถโร วา ปน พาโล อพฺยตฺโต”ติ.\csromanspacer{} อถ โข มํ ภนฺเต เต อญฺญติตฺถิยา ปริพฺพาชกา นววาเทน จ พาลวาเทน จ อปสาเทตฺวา อุฏฺฐายาสนา ปกฺกมิํสุ.\csromanspacer{} ตสฺส มยฺหํ ภนฺเต อจิรปกฺกนฺเตสุ เตสุ อญฺญติตฺถิเยสุ ปริพฺพาชเกสุ เอตทโหสิ “สเจ โข มํ เต อญฺญติตฺถิยา ปริพฺพาชเกสุ เอตทโหสิ “สเจ โข มํ เต อญฺญติตฺถิยา ปริพฺพาชกา อุตฺตริํ ปุจฺเฉยฺยุํ, กถํ พฺยากรมาโน นุ ขฺวาหํ เตสํ อญฺญติตฺถิยานํ ปริพฺพาชกานํ, วุตฺตวาที เจว ภควโต อสฺสํ, น จ ภควนฺตํ อภูเตน อพฺภาจิกฺเขยฺยํ, ธมฺมสฺส จานุธมฺมํ พฺยากเรยฺยํ, น จ โกจิ สหธมฺมิโก วาทานุวาโท คารยฺหํ ฐานํ อาคจฺเฉยฺยา”ติ.}

\prose{ตํ กิํ มญฺญสิ อนุราธ, รูปํ นิจฺจํ วา อนิจฺจํ วาติ.\csromanspacer{} อนิจฺจํ ภนฺเต.\csromanspacer{} ยํ ปนานิจฺจํ, ทุกฺขํ วา ตํ สุขํ วาติ.\csromanspacer{} ทุกฺขํ ภนฺเต.\csromanspacer{} ยํ ปนานิจฺจํ ทุกฺขํ วิปริณามธมฺมํ, กลฺลํ นุ ตํ สมนุปสฺสิตุํ “เอตํ มม, เอโสหมสฺมิ, เอโส เม อตฺตา”ติ.\csromanspacer{} โน เหตํ ภนฺเต.\csromanspacer{} เวทนา นิจฺจา วา อนิจฺจา วาติ ฯเปฯ.}

\csromanheader{2}{10. อพฺยากตสํยุตฺต}
\prose{\csromanfolio{551}สญฺญา.\csromanspacer{} สงฺขารา.\csromanspacer{} วิญฺญาณํ นิจฺจํ วา อนิจฺจํ วาติ.\csromanspacer{} อนิจฺจํ ภนฺเต.\csromanspacer{} ยํ ปนานิจฺจํ, ทุกฺขํ วา ตํ สุขํ วาติ.\csromanspacer{} ทุกฺขํ ภนฺเต.\csromanspacer{} ยํ ปนานิจฺจํ ทุกฺขํ วิปริณามมมฺมํ, กลฺลํ นุ ตํ สมนุปสฺสิตุํ “เอตํ มม, เอโสหมสฺมิ, เอโส เม อตฺตา”ติ.\csromanspacer{} โน เหตํ ภนฺเต.\csromanspacer{} ตสฺมาติห อนุราธ ยํ กิญฺจิ รูปํ อตีตานาคตปจฺจุปฺปนฺนํ อชฺฌตฺตํ วา พหิทฺธา วา โอฬาริกํ วา สุขุมํ วา หีนํ วา ปณีตํ วา ยํ ทูเร สนฺติเก วา สพฺพํ รูปํ “เนตํ มม, เนโสหมสฺมิ, น เมโส อตฺตา”ติ เอวเมตํ ยถาภูตํ สมฺมปฺปญฺญาย ทฏฺฐพฺพํ.\csromanspacer{} ยา กาจิ เวทนา อตีตานาคตปจฺจุปฺปนฺนา.\csromanspacer{} ยา กาจิ สญฺญา.\csromanspacer{} เย เกจิ สงฺขารา.\csromanspacer{} ยํ กิญฺจิ วิญฺญาณํ อตีตานาคตปจฺจุปฺปนฺนํ อชฺฌตฺตํ วา พหิทฺธา วา โอฬาริกํ วา สุขุมํ วา หีนํ วา ปณีตํ วา ยํ ทูเร สนฺติเก วา สพฺพํ วิญฺญาณํ “เนตํ มม, เนโสหมสฺมิ, น เมโส อตฺตา”ติ เอวเมตํ ยถาภูตํ สมฺมปฺปญฺญาย ทฏฺฐพฺพํ.\csromanspacer{} เอวํ ปสฺสํ อนุราธ สุตวา อริยสาวโก รูปสฺมิมฺปิ นิพฺพินฺทติ, เวทนายปิ นิพฺพินฺทติ, สญฺญายปิ นิพฺพินฺทติ, สงฺขาเรสุปิ นิพฺพินฺทติ, วิญฺญาณสฺมิมฺปิ นิพฺพินฺทติ, นิพฺพินฺทํ วิรชฺชติ, วิราคา วิมุจฺจติ, วิมุตฺตสฺมิํ “วิมุตฺตมฺ”อิติ ญาณํ โหติ, “ขีณา ชาติ, วุสิตํ พฺรหฺมจริยํ, กภํ กรณียํ, นาปรํ อิตฺถตฺตายา”ติ ปชานาติ.}

\prose{ตํ กิํ มญฺญสิ อนุราธ, รูปํ ตถาคโตติ สมนุปสฺสสีติ.\csromanspacer{} โน เหตํ ภนฺเต.\csromanspacer{} เวทนํ ตถาคโตติ สมนุปสฺสสีติ.\csromanspacer{} โน เหตํ ภนฺเต.\csromanspacer{} สญฺญํ ตถาคโตติ สมนุปสฺสสีติ.\csromanspacer{} โน เหตํ ภนฺเต.\csromanspacer{} สงฺขาเร ตถาคโตติ สมนุปสฺสสีติ.\csromanspacer{} โน เหตํ ภนฺเต.\csromanspacer{} วิญฺญาณํ ตถาคโตติ สมนุปสฺสสีติ.\csromanspacer{} โน เหตํ ภนฺเต.\csromanspacer{} ตํ กิํ มญฺญสิ อนุราธ, รูปสฺมิํ ตถาคโตติ สมนุปสฺสสีติ.\csromanspacer{} โน เหตํ ภนฺเต.\csromanspacer{} อญฺญตฺร รูปา ตถาคโตติ สมนุปสฺสสีติ.\csromanspacer{} โน เหตํ ภนฺเต.\csromanspacer{} เวทนาย ฯเปฯ.\csromanspacer{} อญฺญตฺร เวทนาย ฯเปฯ.\csromanspacer{} สญฺญาย ฯเปฯ.\csromanspacer{} อญฺญตฺร สญฺญาย ฯเปฯ.\csromanspacer{} สงฺขาเรสุ ฯเปฯ.\csromanspacer{} อญฺญตฺร สงฺขาเรหิ ฯเปฯ.\csromanspacer{} วิญฺญาณสฺมิํ ตถาคโตติ สมนุปสฺสสีติ.\csromanspacer{} โน เหตํ ภนฺเต.\csromanspacer{} อญฺญตฺร วิญฺญาณา ตถาคโตติ สมนุปสฺสสีติ.\csromanspacer{} โน เหตํ ภนฺเต.}

\prose{ตํ กิํ มญฺญสิ อนุราธ, รูปํ เวทนํ สญฺญํ สงฺขาเร วิญฺญาณํ ตถาคโตติ สมนุปสฺสสีติ.\csromanspacer{} โน เหตํ ภนฺเต.\csromanspacer{} ตํ กิํ มญฺญสิ อนุราธ, อยํ โส อรูปี อเวทโน อสญฺญี อสงฺขาโร อวิญฺญาโณ ตถาคโตติ สมนุปสฺสสีติ.\csromanspacer{} โน เหตํ ภนฺเต.\csromanspacer{} เอตฺถ จ เต อนุราธ}

\gambhira{สฬายตนวคฺคสํยุตฺตปาฬิ}
\csromanmatikaanchor{mtk.291}
\csromantocmarkref{subsection}{3-6. สาริปุตฺตโกฏฺฐิกสุตฺต}{mtk.291}
\bookmark[dest={mtk.291},level=2]{3-6. สาริปุตฺตโกฏฺฐิกสุตฺต}
\prose{\csromanfolio{552}ทิฏฺเฐว ธมฺเม สจฺจโต เถตโต ตถาคเต อนุปลพฺภิยมาเน\footnote{ตถาคโต อนุปลพฺภิยมาโน (สฺยา, ก), ตถาคเต อนุปลพฺภมาเน (?)}, กลฺลํ นุ เต ตํ เวยฺยากรณํ\footnote{เวยฺยากรณาย (สี)} “โย โส อาวุโส ตถาคโต อุตฺตมปุริโส ปรมปุริโส ปรมปตฺติปตฺโต, ตํ ตถาคโต อญฺญตฺร อิเมหิ จตูหิ ฐาเนหิ ปญฺญาปยมาโน ปญฺญาเปติ ‘โหติ ตถาคโต ปรํ มรณา’ติ วา ฯเปฯ ‘เนว โหติ น น โหติ ตถาคโต ปรํ มรณา’ติ วา”ติ.\csromanspacer{} โน เหตํ ภนฺเต.\csromanspacer{} สาธุ สาธุ อนุราธ ปุพฺเพ จาหํ อนุราธ เอตรหิ จ ทุกฺขญฺเจว ปญฺญาเปมิ, ทุกฺขสฺส จ นิโรธนฺติ.\csromanspacer{}\csromanspacer{}ทุติยํ.}

\csromanheader{2}{3. ปฐมสาริปุตฺตโกฏฺฐิกสุตฺต}
\proseitem{412}{เอกํ สมยํ อายสฺมา จ สาริปุตฺโต อายสฺมา จ มหาโกฏฺฐิโก พาราณสิยํ วิหรนฺติ อิสิปตเน มิคทาเย.\csromanspacer{} อถ โข อายสฺมา มหาโกฏฺฐิโก สายนฺหสมยํ ปฏิสลฺลานา วุฏฺฐิโต เยนายสฺมา สาริปุตฺโต เตนุปสงฺกมิ, อุปสงฺกมิตฺวา อายสฺมตา สาริปุตฺเตน สทฺธิํ สมฺโมทิ, สมฺโมทนียํ กถํ สารณียํ วีติสาเรตฺวา เอกมนฺตํ นิสีทิ, เอกมนฺตํ นิสินฺโน โข อายสฺมา มหาโกฏฺฐิโก อายสฺมนฺตํ สาริปุตฺตํ เอตทโวจ\csromandash{}}

\prose{กิํ นุ โข อาวุโส สาริปุตฺต โหติ ตถาคโต ปรํ มรณาติ.\csromanspacer{} อพฺยากตํ โข เอตํ อาวุโส ภควตา “โหติ ตถาคโต ปรํ มรณา”ติ.\csromanspacer{} กิํ ปนาวุโส น โหติ ตถาคโต ปรํ มรณาติ, เอตมฺปิ โข อาวุโส อพฺยากตํ ภควตา “น โหติ ตถาคโต ปรํ มรณา”ติ.\csromanspacer{} กิํ นุ โข อาวุโส โหติ จ น จ โหติ ตถาคโต ปรํ มรณาติ.\csromanspacer{} อพฺยากตํ โข เอตํ อาวุโส ภควตา “โหติ จ น จ โหติ ตถาคโต ปรํ มรณา”ติ.\csromanspacer{} กิํ ปนาวุโส เนว โหติ น น โหติ ตถาคโต ปรํ มรณา”ติ.\csromanspacer{} กิํ ปนาวุโส เนว โหติ น น โหติ ตถาคโต ปรํ มรณาติ, เอตมฺปิ โข อาวุโส อพฺยากตํ ภควตา “เนว โหติ น น โหติ ตถาคโต ปรํ มรณา”ติ.}

\prose{“กิํ นุ โข อาวุโส โหติ ตถาคโต ปรํ มรณา”ติ อิติ ปุฏฺโฐ สมาโน “อพฺยากตํ โข เอตํ อาวุโส ภควตา ‘โหติ}

\csromanheader{2}{10. อพฺยากตสํยุตฺต}
\prose{\csromanfolio{553}ตถาคโต ปรํ มรณา’ติ” วเทสิ ฯเปฯ “กิํ ปนาวุโส เนว โหติ น น โหติ ตถาคโต ปรํ มรณา”ติ อิติ ปุฏฺโฐ สมาโน “เอตมฺปิ โข อาวุโส อพฺยากตํ ภควตา ‘เนว โหติ น น โหติ ตถาคโต ปรํ มรณา’ติ” วเทสิ.\csromanspacer{} โก นุ โข อาวุโส เหตุ โก ปจฺจโย, เยเนตํ อพฺยากตํ ภควตาติ.}

\prose{โหติ ตถาคโต ปรํ มรณาติ โข อาวุโส รูปคตเมตํ, น โหติ ตถาคโต ปรํ มรณาติ รูปคตเมตํ, โหติ จ น จ โหติ ตถาคโต ปรํ มรณาติ รูปคตเมตํ, เนว โหติ น น โหติ ตถาคโต ปรํ มรณาติ รูปคตเมตํ.\csromanspacer{} โหติ ตถาคโต ปรํ รมณาติ โข อาวุโส เวทนาคตเมตํ, น โหติ ตถาคโต ปรํ มรณาติ เวทนาคตเมตํ, โหติ จ น จ โหติ ตถาคโต ปรํ มรณาติ เวทนาคตเมตํ, เนว โหติ น น โหติ ตถาคโต ปรํ มรณาติ เวทนาคตเมตํ.\csromanspacer{} โหติ ตถาคโต ปรํ มรณาติ โข อาวุโส สญฺญาคตเมตํ, น โหติ ตถาคโต ปรํ มรณาติ สญฺญาคตเมตํ, โหติ จ น จ โหติ ตถาคโต ปรํ มรณาติ สญฺญาคตเมตํ, เนว โหติ ตถาคโต ปรํ มรณาติ โข อาวุโส สงฺขารคตเมตํ.\csromanspacer{} โหติ ตถาคโต ปรํ มรณาติ โข อาวุโส สงฺขารคตเมตํ, น โหติ ตถาคโต ปรํ มรณาติ สงฺขารคตเมตํ, โหติ จ น จ โหติ ตถาคโต ปรํ มรณาติ สงฺขารคตเมตํ, เนว โหติ น น โหติ ตถาคโต ปรํ มรณาติ สงฺขารคตเมตํ.\csromanspacer{} โหติ ตถาคโต ปรํ มรณาติ โข อาวุโส วิญฺญาณคตเมตํ, น โหติ ตถาคโต ปรํ มรณาติ วิญฺญาณคตเมตํ, โหติ จ น จ โหติ ตถาคโต ปรํ มรณาติ วิญฺญาณคตเมตํ, เนว โหติ น น โหติ ตถาคโต ปรํ มรณาติ วิญฺญาณคตเมตํ.\csromanspacer{} อยํ โข อาวุโส เหตุ อยํ ปจฺจโย, เยเนตํ อพฺยากตํ ภควตาติ.\csromanspacer{}\csromanspacer{}ตติยํ.}

\csromanheader{2}{4. ทุติยสาริปุตฺตโกฏฺฐิกสุตฺต}
\proseitem{413}{เอกํ สมยํ อายสฺมา จ สาริปุตฺโต อายสฺมา จ มหาโกฏฺฐิโก พาราณสิยํ วิหรนฺติ อิสิปเตน มิคทาเย ฯเปฯ.}

\gambhira{สฬายตนวคฺคสํยุตฺตปาฬิ}
\prose{\csromanfolio{554}(สาเยว ปุจฺฉา.) โก นุ โข อาวุโส เหตุ โก ปจฺจโย, เยเนตํ อพฺยากตํ ภควตาติ.\csromanspacer{} รูปํ โข อาวุโส อชานโต อปสฺสโต ยถาภูตํ, รูปสมุทยํ อชานโต อปสฺสโต ยถาภูตํ, รูปนิโรธํ อชานโต อปสฺสโต ยถาภูตํ, รูปนิโรธคามินิํ ปฏิปทํ อชานโต อปสฺสโต ยถาภูตํ “โหติ ตถาคโต ปรํ มรณา”ติปิสฺส โหติ, “น โหติ ตถาคโต ปรํ มรณา”ติปิสฺส โหติ, “โหติ จ น จ โหติ ตถาคโต ปรํ มรณา”ติปิสฺส โหติ, “เนว โหติ น น โหติ ตถาคโต ปรํ มรณา”ติปิสฺส โหติ.\csromanspacer{} เวทนํ.\csromanspacer{} สญฺญํ.\csromanspacer{} สงฺขาเร.\csromanspacer{} วิญฺญาณํ อชานโต อปสฺสโต ยถาภูตํ, วิญฺญาณสมุทยํ อชานโต อปสฺสโต ยถาภูตํ, วิญฺญาณนิโรธํ อชานโต อปสฺสโต ยถาภูตํ, วิญฺญาณนิโรธคามินิํ ปฏิปทํ อชานโต อปสฺสโต ยถาภูตํ “โหติ ตถาคโต ปรํ มรณา”ติปิสฺส โหติ, “น โหติ ตถาคโต ปรํ มรณา”ติปิสฺส โหติ, “โหติ จ น จ โหติ ตถาคโต ปรํ มรณา”ติปิสฺส โหติ, “เนว โหติ น น โหติ ตถาคโต ปรํ มรณา”ติปิสฺส โหติ.}

\prose{รูปญฺจ โข อาวุโส ชานโต ปสฺสโต ยถาภูตํ, รูปสมุทยํ ชานโต ปสฺสโต ยถาภูตํ, รูปนิโรธํ ชานโต ปสฺสโต ยถาภูตํ, รูปนิโรธคามินิํ ปฏิปทํ ชานโต ปสฺสโต ยถาภูตํ “โหติ ตถาคโต ปรํ มรณา”ติปิสฺส น โหติ ฯเปฯ “เนว โหติ น น โหติ ตถาคโต ปรํ มรณา”ติปิสฺส น โหติ.\csromanspacer{} เวทนํ.\csromanspacer{} สญฺญํ.\csromanspacer{} สงฺขาเร.\csromanspacer{} วิญฺญาณํ ชานโต ปสฺสโต ยถาภูตํ, วิญฺญาณสมุทยํ ชานโต ปสฺสโต ยถาภูตํ, วิญฺญาณนิโรธํ ชานโต ปสฺสโต ยถาภูตํ, วิญฺญาณนิโรธคามินิํ ปฏิปทํ ชานโต ปสฺสโต ยถาภูตํ “โหติ ตถาคโต ปรํ มรณา”ติปิสฺส น โหติ, “น โหติ ตถาคโต ปรํ มรณา”ติปิสฺส น โหติ, “โหติ จ น จ โหติ ตถาคโต ปรํ มรณา”ติปิสฺส น โหติ, “เนว โหติ น น โหติ ตถาคโต ปรํ มรณา”ติปิสฺส น โหติ.\csromanspacer{} อยํ โข อาวุโส เหตุ อยํ ปจฺจโย, เยเนตํ อพฺยากตํ ภควตาติ.\csromanspacer{}\csromanspacer{}จตุตฺถํ.}

\csromanheader{2}{10. อพฺยากตสํยุตฺต}
\csromanheader{2}{5. ตติยสาริปุตฺตโกฏฺฐิกสุตฺต}
\proseitem{414}{\csromanfolio{555}เอกํ สมยํ อายสฺมา จ สาริปุตฺโต อายสฺมา จ มหาโกฏฺฐิโก พาราณสิยํ วิหรนฺติ อิสิปตเน มิคทาเย ฯเปฯ. (สาเยว ปุจฺฉา.) โก นุ โข อาวุโส เหตุ โก ปจฺจโย, เยเนตํ อพฺยากตํ ภควตาติ.\csromanspacer{} รูเป โข อาวุโส อวิคตราคสฺส อวิคตจฺฉนฺทสฺส อวิคตเปมสฺส อวิคตปิปาสสฺส อวิคตปริฬาหสฺส อวิคตตณฺหสฺส “โหติ ตถาคโต ปรํ มรณา”ติปิสฺส โหติ ฯเปฯ “เนว โหติ น น โหติ ตถาคโต ปรํ มรณา”ติปิสฺส โหติ.\csromanspacer{} เวทนาย.\csromanspacer{} สญฺญาย.\csromanspacer{} สงฺขาเรสุ.\csromanspacer{} วิญฺญาเณ อวิคตราคสฺส อวิคตจฺฉนฺทสฺส อวิคตเปมสฺส อวิคตปิปาสสฺส อวิคตปริฬาหสฺส อวิคตตณฺหสฺส “โหติ ตถาคโต ปรํ มรณา”ติปิสฺส โหติ ฯเปฯ “เนว โหติ น น โหติ ตถาคโต ปรํ รมณา”ติปิสฺส โหติ.\csromanspacer{} รูเป จ โข อาวุโส วิคตราคสฺส ฯเปฯ.\csromanspacer{} เวทนาย.\csromanspacer{} สญฺญาย.\csromanspacer{} สงฺขาเรสุ.\csromanspacer{} วิญฺญาเณ วิคตราคสฺส วิคตจฺฉนฺทสฺส วิคตเปมสฺส วิคตปิปาสสฺส วิคตปริฬาหสฺส วิคตตณฺหสฺส “โหติ ตถาคโต ปรํ มรณา”ติปิสฺส น โหติ ฯเปฯ “เนว โหติ น น โหติ ตถาคโต มรณา”ติปิสฺส น โหติ.\csromanspacer{} อยํ โข อาวุโส เหตุ อยํ ปจฺจโย, เยเนตํ อพฺยากตํ ภควตาติ.\csromanspacer{}\csromanspacer{}ปญฺจมํ.}

\csromanheader{2}{6. จตุตฺถสาริปุตฺตโกฏฺฐิกสุตฺต}
\proseitem{415}{เอกํ สมยํ อายสฺมา จ สาริปุตฺโต อายสฺมา จ มหาโกฏฺฐิโก พาราณสิยํ วิหรนฺติ อิสิปตเน มิคทาเย.\csromanspacer{} อถ โข อายสฺมา สาริปุตฺโต สายนฺหสมยํ ปฏิสลฺลานา วุฏฺฐิโต เยนายสฺมา มหาโกฏฺฐิโก เตนุปสงฺกมิ, อุปสงฺกมิตฺวา อายสฺมตา มหาโกฏฺฐิเกน สทฺธิํ สมฺโมทิ, สมฺโมทนียํ กถํ สารณียํ วีติสาเรตฺวา เอกมนฺตํ นิสีทิ, เอกมนฺตํ นิสินฺโน โข อายสฺมา สาริปุตฺโต อายสฺมนฺตํ มหาโกฏฺฐิกํ เอตทโวจ “กิํ นุ โข อาวุโส โกฏฺฐิก โหติ ตถาคโต ปรํ มรณา”ติ ฯเปฯ “กิํ ปนาวุโส เนว โหติ น น โหติ ตถาคโต ปรํ มรณา”ติ อิติ ปุฏฺโฐ สมาโน “เอตมฺปิ โข อาวุโส อพฺยากตํ ภควตา ‘เนว โหติ น น โหติ}

\gambhira{สฬายตนวคฺคสํยุตฺตปาฬิ}
\prose{\csromanfolio{556}ตถาคโต ปรํ มรณา’ติ” วเทสิ.\csromanspacer{} โก นุ โข อาวุโส เหตุ โก ปจฺจโย, เยเนตํ อพฺยากตํ ภควตาติ.}

\prose{รูปารามสฺส โข อาวุโส รูปรตสฺส รูปสมฺมุทิตสฺส รูปนิโรธํ อชานโต อปสฺสโต ยถาภูตํ “โหติ ตถาคโต ปรํ มรณา”ติปิสฺส โหติ, “น โหติ ตถาคโต ปรํ มรณา”ติปิสฺส โหติ, “โหติ จ น จ โหติ ตถาคโต ปรํ มรณา”ติปิสฺส โหติ, “เนว โหติ น น โหติ ตถาคโต ปรํ มรณา”ติปิสฺส โหติ.\csromanspacer{} เวทนารามสฺส โข อาวุโส เวทนารตสฺส เวทนาสมฺมุทิตสฺส เวทนานิโรธํ อชานโต อปสฺสโต ยถาภูตํ “โหติ ตถาคโต ปรํ มรณา”ติปิสฺส โหติ ฯเปฯ.\csromanspacer{} สญฺญารามสฺส โข อาวุโส.\csromanspacer{} สงฺขารารามสฺส โข อาวุโส.\csromanspacer{} วิญฺญาณารามสฺส โข อาวุโส วิญฺญาณรตสฺส วิญฺญาณสมฺมุทิตสฺส วิญฺญาณนิโรธํ อชานโต อปสฺสโต ยถาภูตํ “โหติ ตถาคโต ปรํ มรณา”ติปิสฺส โหติ ฯเปฯ “เนว โหติ น น โหติ ตถาคโต ปรํ มรณา”ติปิสฺส โหติ.}

\prose{น รูปารามสฺส โข อาวุโส น รูปรตสฺส น รูปสมฺมุทิตสฺส รูปนิโรธํ ชานโต ปสฺสโต ยถาภูตํ “โหติ ตถาคโต ปรํ มรณา”ติปิสฺส น โหติ ฯเปฯ “เนว โหติ น น โหติ ตถาคโต ปรํ มรณา”ติปิสฺส น โหติ.\csromanspacer{} น เวทนารามสฺส โข อาวุโส ฯเปฯ.\csromanspacer{} น สญฺญารามสฺส โข อาวุโส.\csromanspacer{} น สงฺขารารามสฺส โข อาวุโส.\csromanspacer{} น วิญฺญาณารามสฺส โข อาวุโส น วิญฺญาณรตสฺส น วิญฺญาณสมฺมุทิตสฺส วิญฺญาณนิโรธํ ชานโต ปสฺสโต ยถาภูตํ “โหติ ตถาคโต ปรํ มรณา”ติปิสฺส น โหติ ฯเปฯ “เนว โหติ น น โหติ ตถาคโต ปรํ มรณา”ติปิสฺส น โหติ.\csromanspacer{} อยํ โข อาวุโส เหตุ อยํ ปจฺจโย, เยเนตํ อพฺยากตํ ภควตาติ.}

\prose{สิยา ปนาวุโส อญฺโญปิ ปริยาโย, เยเนตํ อพฺยากตํ ภควตาติ.\csromanspacer{} สิยา อาวุโส, ภวารามสฺส โข อาวุโส ภวรตสฺส ภวสมฺมุทิตสฺส ภวนิโรธํ อชานโต อปสฺสโต ยถาภูตํ “โหติ ตถาคโต ปรํ มรณา”ติปิสฺส โหติ ฯเปฯ “เนว โหติ น น โหติ ตถาคโต ปรํ มรณา”ติปิสฺส โหติ.}

\csromanheader{2}{10. อพฺยากตสํยุตฺต}
\prose{\csromanfolio{557}น ภวารามสฺส โข อาวุโส น ภวรตสฺส น ภวสมฺมุทิตสฺส ภวนิโรธํ ชานโต ปสฺสโต ยถาภูตํ “โหติ ตถาคโต ปรํ มรณา”ติปิสฺส น โหติ ฯเปฯ “เนว โหติ น น โหติ ตถาคโต ปรํ มรณา”ติปิสฺส น โหติ.\csromanspacer{} อยมฺปิ โข อาวุโส ปริยาโย, เยเนตํ อพฺยากตํ ภควตาติ.}

\prose{สิยา ปนาวุโส อญฺโญปิ ปริยาโย, เยเนตํ อพฺยากตํ ภควตาติ.\csromanspacer{} สิยา อาวุโส, อุปาทานารามสฺส โข อาวุโส อุปาทานรตสฺส อุปาทานสมฺมุทิตสฺส อุปาทานนิโรธํ อชานโต อปสฺสโต ยถาภูตํ “โหติ ตถาคโต ปรํ มรณา”ติปิสฺส โหติ ฯเปฯ “เนว โหติ น น โหติ ตถาคโต ปรํ มรณา”ติปิสฺส โหติ.\csromanspacer{} น อุปาทานารามสฺส โข อาวุโส น อุปาทานรตสฺส น อุปาทานสมฺมุทิตสฺส อุปาทานนิโรธํ ชานโต ปสฺสโต ยถาภูตํ “โหติ ตถาคโต ปรํ มรณา”ติปิสฺส น โหติ ฯเปฯ “เนว โหติ น น โหติ ตถาคโต ปรํ มรณา”ติปิสฺส น โหติ.\csromanspacer{} อยมฺปิ โข อาวุโส ปริยาโย, เยเนตํ อพฺยากตํ ภควตาติ.}

\prose{สิยา ปนาวุโส อญฺโญปิ ปริยาโย, เยเนตํ อพฺยากตํ ภควตาติ.\csromanspacer{} สิยา อาวุโส, ตณฺหารามสฺส โข อาวุโส ตณฺหารตสฺส ตณฺหาสมฺมุทิตสฺส ตณฺหานิโรธํ อชานโต อปสฺสโต ยถาภูตํ “โหติ ตถาคโต ปรํ มรณา”ติปิสฺส โหติ ฯเปฯ “เนว โหติ น น โหติ ตถาคโต ปรํ มรณา”ติปิสฺส โหติ.\csromanspacer{} น ตณฺหารามสฺส โข อาวุโส น ตณฺหารตสฺส น ตณฺหาสมฺมุทิตสฺส ตณฺหานิโรธํ ชานโต ปสฺสโต ยถาภูตํ “โหติ ตถาคโต ปรํ มรณา”ติปิสฺส น โหติ ฯเปฯ “เนว โหติ น น โหติ ตถาคโต ปรํ มรณา”ติปิสฺส น โหติ.\csromanspacer{} อยมฺปิ โข อาวุโส ปริยาโย, เยเนตํ อพฺยากตํ ภควตาติ.}

\prose{สิยา ปนาวุโส อญฺโญปิ ปริยาโย, เยเนตํ อพฺยากตํ ภควตาติ.\csromanspacer{} เอตฺถ ทานิ อาวุโส สาริปุตฺต อิโต อุตฺตริ กิํ อิจฺฉสิ.\csromanspacer{} ตณฺหาสงฺขยวิมุตฺตสฺส อาวุโส สาริปุตฺต ภิกฺขุโน วฏฺฏํ\footnote{วตฺตํ (สฺยา, กํ, ก), วทฺธํ (อิ)} นตฺถิ ปญฺญาปนายาติ.\csromanspacer{}\csromanspacer{}ฉฏฺฐํ.}

\gambhira{สฬายตนวคฺคสํยุตฺตปาฬิ}
\csromanheader{2}{7. โมคฺคลฺลานสุตฺต}
\csromanmatikaanchor{mtk.292}
\csromantocmarkref{subsection}{7. โมคฺคลฺลานสุตฺต}{mtk.292}
\bookmark[dest={mtk.292},level=2]{7. โมคฺคลฺลานสุตฺต}
\proseitem{416}{\csromanfolio{558}อถ โข วจฺฉโคตฺโต ปริพฺพาชโก เยนายสฺมา มหาโมคฺคลฺลาโน เตนุปสงฺกมิ, อุปสงฺกมิตฺวา อายสฺมตา มหาโมคฺคลฺลาเนน สทฺธิํ สมฺโมทิ, สมฺโมทนียํ กถํ สารณียํ วีติสาเรตฺวา เอกมนฺตํ นิสีทิ, เอกมนฺตํ นิสินฺโน โข วจฺฉโคตฺโต ปริพฺพาชโก อายสฺมนฺตํ มหาโมคฺคลฺลานํ เอตทโวจ\csromandash{}}

\prose{กิํ นุ โข โภ โมคฺคลฺลาน สสฺสโต โลโกติ.\csromanspacer{} อพฺยากตํ โข เอตํ วจฺฉ ภควตา “สสฺสโต โลโก”ติ.\csromanspacer{} กิํ ปน โภ โมคฺคลฺลาน อสสฺสโต โลโกติ.\csromanspacer{} เอตมฺปิ โข วจฺฉ อพฺยากตํ ภควตา “อสสฺสโต โลโก”ติ.\csromanspacer{} กิํ นุ โข โภ โมคฺคลฺลาน อนฺตวา โลโกติ.\csromanspacer{} อพฺยากตํ โข เอตํ วจฺฉ ภควตา “อนฺตวา โลโก”ติ.\csromanspacer{} กิํ ปน โภ โมคฺคลฺลาน อนนฺตวา โลโกติ.\csromanspacer{} เอตมฺปิ โข วจฺฉ อพฺยากตํ ภควตา “อนนฺตวา โลโก”ติ.\csromanspacer{} กิํ นุ โข โภ โมคฺคลฺลาน ตํ ชีวํ ตํ สรีรนฺติ.\csromanspacer{} อพฺยากตํ โข เอตํ วจฺฉ ภควตา “ตํ ชีวํ ตํ สรีรนฺ”ติ.\csromanspacer{} กิํ ปน โภ โมคฺคลฺลาน อญฺญํ ชีวํ อญฺญํ สรีรนฺติ.\csromanspacer{} เอตมฺปิ โข วจฺฉ อพฺยากตํ ภควตา “อญฺญํ ชีวํ อญฺญํ สรีรนฺ”ติ.\csromanspacer{} กิํ นุ โข โภ โมคฺคลฺลาน โหติ ตถาคโต ปรํ มรณาติ.\csromanspacer{} อพฺยากตํ โข เอตํ วจฺฉ ภควตา “โหติ ตถาคโต ปรํ มรณา”ติ.\csromanspacer{} กิํ ปน โภ โมคฺคลฺลาน น โหติ ตถาคโต ปรํ มรณาติ.\csromanspacer{} เอตมฺปิ โข วจฺฉ อพฺยากตํ ภควตา “น โหติ ตถาคโต ปรํ มรณา”ติ.\csromanspacer{} กิํ นุ โข โภ โมคฺคลฺลาน โหติ จ น จ โหติ ตถาคโต ปรํ มรณาติ.\csromanspacer{} อพฺยากตํ โข เอตํ วจฺฉ ภควตา “โหติ จ น จ โหติ ตถาคโต ปรํ มรณา”ติ.\csromanspacer{} กิํ ปน โภ โมคฺคลฺลาน เนว โหติ น น โหติ ตถาคโต ปรํ มรณาติ.\csromanspacer{} เอตมฺปิ โข วจฺฉ อพฺยากตํ ภควตา “เนว โหติ น น โหติ ตถาคโต ปรํ มรณา”ติ.}

\prose{โก นุ โข โภ โมคฺคลฺลาน เหตุ โก ปจฺจโย, เยน อญฺญติตฺถิยานํ ปริพฺพาชกานํ เอวํ ปุฏฺฐานํ เอวํ เวยฺยากรณํ โหติ “สสฺสโต โลโก”ติ วา, “อสสฺสโต โลโก”ติ วา, “อนฺตวา โลโก”ติ วา, “อนนฺตวา โลโก”ติ วา, “ตํ ชีวํ ตํ สรีรนฺ”ติ วา, “อญฺญํ ชีวํ อญฺญํ สรีรนฺ”ติ วา, “โหติ ตถาคโต}

\csromanheader{2}{10. อพฺยากตสํยุตฺต}
\prose{\csromanfolio{559}ปรํ มรณา”ติ วา, “น โหติ ตถาคโต ปรํ มรณา”ติ วา, “โหติ จ น จ โหติ ตถาคโต ปรํ มรณา”ติ วา, เนว โหติ น น โหติ ตถาคโต ปรํ มรณา”ติ วา.\csromanspacer{} โก ปน โภ โมคฺคลฺลาน เหตุ โก ปจฺจโย, เยน สมณสฺส โคตมสฺส เอวํ ปุฏฺฐสฺส น เอวํ เวยฺยากรณํ โหติ “สสฺสโต โลโก”ติปิ, “อสสฺสโต โลโก”ติปิ, “อนฺตวา โลโก”ติปิ, “อนนฺตวา โลโก”ติปิ, “ตํ ชีวํ ตํ สรีรนฺ”ติปิ, “อญฺญํ ชีวํ อญฺญํ สรีรนฺ”ติปิ, “โหติ ตถาคโต ปรํ มรณา”ติปิ, “น โหติ ตถาคโต ปรํ มรณา”ติปิ, “โหติ จ น จ โหติ ตถาคโต ปรํ มรณา”ติปิ, “เนว โหติ น น โหติ ตถาคโต ปรํ มรณา”ติปีติ.}

\prose{อญฺญติตฺถิยา โข วจฺฉ ปริพฺพาชกา จกฺขุํ “เอตํ มม, เอโสหมสฺมิ, เอโส เม อตฺตา”ติ สมนุปสฺสนฺติ ฯเปฯ ชิวฺหํ “เอตํ มม, เอโสหมสฺมิ, เอโส เม อตฺตา”ติ สมนุปสฺสนฺติ ฯเปฯ มนํ “เอตํ มม, เอโสหมสฺมิ, เอโส เม อตฺตา”ติ สมนุปสฺสนฺติ.\csromanspacer{} ตสฺมา อญฺญติตฺถิยานํ ปริพฺพาชกานํ เอวํ ปุฏฺฐานํ เอวํ เวยฺยากรณํ โหติ “สสฺสโต โลโก”ติ วา ฯเปฯ “เนว โหติ น น โหติ ตถาคโต ปรํ มรณา”ติ วา.\csromanspacer{} ตถาคโต จ โข วจฺฉ อรหํ สมฺมาสมฺพุทฺโธ จกฺขุํ “เนตํ มม, เนโสหมสฺมิ, น เมโส อตฺตา”ติ สมนุปสฺสติ ฯเปฯ ชิวฺหํ “เนตํ มม, เนโสหมสฺมิ, น เมโส อตฺตา”ติ สมนุปสฺสติ ฯเปฯ มนํ “เนตํ มม, เนโสหมสฺมิ, น เมโส อตฺตา”ติ สมนุปสฺสติ.\csromanspacer{} ตสฺมา ตถาคตสฺส เอวํ ปุฏฺฐสฺส น เอวํ เวยฺยากรณํ โหติ “สสฺสโต โลโก”ติปิ ฯเปฯ “เนว โหติ น น โหติ ตถาคโต ปรํ มรณา”ติปีติ.}

\prose{อถ โข วจฺฉโคตฺโต ปริพฺพาชโก อุฏฺฐายาสนา เยน ภควา เตนุปสงฺกมิ, อุปสงฺกมิตฺวา ภควตา สทฺธิํ สมฺโมทิ, สมฺโมทนียํ กถํ สารณียํ วีติสาเรตฺวา เอกมนฺตํ นิสีทิ, เอกมนฺตํ นิสินฺโน โข วจฺฉโคตฺโต ปริพฺพาชโก ภควนฺตํ เอตทโวจ “กิํ นุ โข โภ โคตม สสฺสโต โลโก”ติ.\csromanspacer{} อพฺยากตํ โข เอตํ วจฺฉ มยา “สสฺสโต โลโก”ติ ฯเปฯ.\csromanspacer{} กิํ ปน โภ โคตม เนว โหติ น น โหติ ตถาคโต ปรํ มรณาติ.\csromanspacer{} เอตมฺปิ โข วจฺฉ อพฺยากตํ มยา “เนว โหติ น น โหติ ตถาคโต ปรํ มรณา”ติ.}

\gambhira{สฬายตนวคฺคสํยุตฺตปาฬิ}
\prose{\csromanfolio{560}โก นุ โข โภ โคตม เหตุ โก ปจฺจโย, เยน อญฺญติตฺถิยานํ ปริพฺพาชกานํ เอวํ ปุฏฺฐานํ เอวํ เวยฺยากรณํ โหติ “สสฺสโต โลโก”ติ วา ฯเปฯ “เนว โหติ น น โหติ ตถาคโต ปรํ มรณา”ติ วา.\csromanspacer{} โก ปน โภ โคตม เหตุ โก ปจฺจโย, เยน โภโต โคตมสฺส เอวํ ปุฏฺฐสฺส น เอวํ เวยฺยากรณํ โหติ “สสฺสโต โลโก”ติปิ ฯเปฯ “เนว โหติ น น โหติ ตถาคโต ปรํ มรณา”ติปีติ.}

\prose{อญฺญติตฺถิยา โข วจฺฉ ปริพฺพาชกา จกฺขุํ “เอตํ มม, เอโสหมสฺมิ, เอโส เม อตฺตา”ติ สมนุปสฺสนฺติ ฯเปฯ ชิวฺหํ “เอตํ มม, เอโสหมสฺมิ, เอโส เม อตฺตา”ติ สมนุสฺสนฺติ ฯเปฯ มนํ “เอตํ มม, เอโสหมสฺมิ, เอโส เม อตฺตา”ติ สมนุปสฺสนฺติ.\csromanspacer{} ตสฺมา อญฺญตฺถิยานํ ปริพฺพาชกานํ เอวํ ปุฏฺฐานํ เอวํ เวยฺยากรณํ โหติ “สสฺสโต โลโก”ติ วา ฯเปฯ “เนว โหติ น น โหติ ตถาคโต ปรํ มรณา”ติ วา.\csromanspacer{} ตถาคโต จ โข วจฺฉ อรหํ สมฺมาสมฺพุทฺโธ จกฺขุํ “เนตํ มม, เนโสหมสฺมิ, น เมโส อตฺตา”ติ สมนุปสฺสติ ฯเปฯ ชิวฺหํ “เนตํ มม, เนโสหมสฺมิ, น เมโส อตฺตา”ติ สมนุปสฺสติ ฯเปฯ มนํ “เนตํ มม, เนโสหมสฺมิ, น เมโส อตฺตา”ติ สมนุปสฺสติ.\csromanspacer{} ตสฺมา ตถาคตสฺส เอวํ ปุฏฺฐสฺส น เอวํ เวยฺยากรณํ โหติ “สสฺสโต โลโก”ติปิ, “อสสฺสโต โลโก”ติปิ, “อนฺตวาโลโก”ติปิ, “อนนฺตวา โลโก”ติปิ, “ตํ ชีวํ ตํ สรีรนฺ”ติปิ, “อญฺญํ ชีวํ อญฺญํ สรีรนฺ”ติปิ, “โหติ ตถาคโต ปรํ มรณา”ติปิ, “น โหติ ตถาคโต ปรํ มรณา”ติปิ, “โหติ จ น จ โหติ ตถาคโต ปรํ มรณา”ติปิ, “เนว โหติ น น โหติ ตถาคโต ปรํ มรณา”ติปีติ.}

\prose{อจฺฉริยํ โภ โคตม, อพฺภุตํ โภ โคตม, ยตฺร หิ นาม สตฺถุ จ\footnote{สตฺถุสฺส จ (สี, อิ), สตฺถุ เจว (เขมาสุตฺเต)} สาวกสฺส จ อตฺเถน อตฺโถ พฺยญฺชเนน พฺยญฺชนํ สํสนฺทิสฺสติ สเมสฺสติ น วิโรธยิสฺสติ, ยทิทํ อคฺคปทสฺมิํ.\csromanspacer{} อิทานาหํ โภ โคตม สมณํ มหาโมคฺคลฺลานํ อุปสงฺกมิตฺวา เอตมตฺถํ อปุจฺฉิํ, สมโณปิ เม โมคฺคลฺลาโน เอเตหิ ปเทหิ เอเตหิ พฺยญฺชเนหิ ตมตฺถํ พฺยากาสิ, เสยฺยถาปิ ภวํ โคตโม.\csromanspacer{} อจฺฉริยํ โภ โคตม, อพฺภุตํ โภ}

\csromanheader{2}{10. อพฺยากตสํยุตฺต}
\prose{\csromanfolio{561}โคตม, ยตฺร หิ นาม สตฺถุ จ สาวกสฺส จ อตฺเถน อตฺโถ พฺยญฺชเนน พฺยญฺชนํ สํสนฺทิสฺสติ สเมสฺสติ น วิโรธยิสฺสติ, ยทิทํ อคฺคปทสฺมิํติ.\csromanspacer{}\csromanspacer{}สตฺตมํ.}

\csromanheader{2}{8. วจฺฉโคตฺตสุตฺต}
\csromanmatikaanchor{mtk.293}
\csromantocmarkref{subsection}{8. วจฺฉโคตฺตสุตฺต}{mtk.293}
\bookmark[dest={mtk.293},level=2]{8. วจฺฉโคตฺตสุตฺต}
\proseitem{417}{อถ โข วจฺฉโคตฺโต ปริพฺพาชโก เยน ภควา เตนุปสงฺกมิ, อุปสงฺกมิตฺวา ภควตา สทฺธิํ สมฺโมทิ, สมฺโมทนียํ กถํ สารณียํ วีติสาเรตฺวา เอกมนฺตํ นิสีทิ, เอกมนฺตํ นิสินฺโน โข วจฺฉโคตฺโต ปริพฺพาชโก ภควนฺตํ เอตทโวจ “กิํ นุ โข โภ โคตม สสฺสโต โลโก”ติ.\csromanspacer{} อพฺยากตํ โข เอตํ วจฺฉ มยา “สสฺสโต โลโก”ติ ฯเปฯ.\csromanspacer{} กิํ ปน โภ โคตม เนว โหติ น น โหติ ตถาคโต ปรํ มรณาติ.\csromanspacer{} เอตมฺปิ โข วจฺฉ อพฺยากตํ มยา “เนว โหติ น น โหติ ตถาคโต ปรํ มรณา”ติ.}

\prose{โก นุ โข โภ โคตม เหตุ โก ปจฺจโย, เยน อญฺญติตฺถิยานํ ปริพฺพาชกานํ เอวํ ปุฏฺฐานํ เอวํ เวยฺยากรณํ โหติ “สสฺสโต โลโก”ติ วา ฯเปฯ “เนว โหติ น น โหติ ตถาคโต ปรํ มรณา”ติ วา.\csromanspacer{} โก ปน โภ โคตม เหตุ โก ปจฺจโย, เยน โภโต โคตมสฺส เอวํ ปุฏฺฐสฺส น เอวํ เวยฺยากรณํ โหติ “สสฺสโต โลโก”ติปิ ฯเปฯ “เนว โหติ น น โหติ ตถาคโต ปรํ มรณา”ติปีติ.}

\prose{อญฺญติตฺถิยา โข วจฺฉ ปริพฺพาชกา รูปํ อตฺตโต สมนุปสฺสนฺติ, รูปวนฺตํ วา อตฺตานํ, อตฺตนิ วา รูปํ, รูปสฺมิํ วา อตฺตานํ.\csromanspacer{} เวทนํ อตฺตโต สมนุปสฺสนฺติ ฯเปฯ.\csromanspacer{} สญฺญํ.\csromanspacer{} สงฺขาเร.\csromanspacer{} วิญฺญาณํ อตฺตโต สมนุปสฺสนฺติ, วิญฺญาณวนฺตํ วา อตฺตานํ, อตฺตนิ วา วิญฺญาณํ, วิญฺญาณสฺมิํ วา อตฺตานํ.\csromanspacer{} ตสฺมา อญฺญติตฺถิยานํ ปริพฺพาชกานํ เอวํ ปุฏฺฐานํ เอวํ เวยฺยากรณํ โหติ “สสฺสโต โลโก”ติ วา ฯเปฯ “เนว โหติ น น โหติ ตถาคโต ปรํ มรณา”ติ วา.\csromanspacer{} ตถาคโต จ โข วจฺฉ อรหํ สมฺมาสมฺพุทฺโธ น รูปํ อตฺตโต สมนุปสฺสติ, น รูปวนฺตํ วา อตฺตานํ, น อตฺตนิ วา รูปํ, น รูปสฺมิํ วา อตฺตานํ.\csromanspacer{} น เวทนํ อตฺตโต สมนุปสฺสติ ฯเปฯ.\csromanspacer{} น สญฺญํ.\csromanspacer{} น สงฺขาเร.\csromanspacer{} น วิญฺญาณํ อตฺตโต สมนุปสฺสติ,}

\gambhira{สฬายตนวคฺคสํยุตฺตปาฬิ}
\prose{\csromanfolio{562}น วิญฺญาณวนฺตํ วา อตฺตานํ, น อตฺตนิ วา วิญฺญาณํ, น วิญฺญาณสฺมิํ วา อตฺตานํ.\csromanspacer{} ตสฺมา ตถาคตสฺส เอวํ ปุฏฺฐสฺส น เอวํ เวยฺยากรณํ โหติ “สสฺสโต โลโก”ติปิ ฯเปฯ “เนว โหติ น น โหติ ตถาคโต ปรํ มรณา”ติปีติ.}

\prose{อถ โข วจฺฉโคตฺโต ปริพฺพาชโก อุฏฺฐายาสนา เยนยสฺมา มหาโมคฺคลฺลาโน เตนุปสงฺกมิ, อุปสงฺกมิตฺวา อายสฺมตา มหาโมคฺคลฺลาเนน สทฺธิํ สมฺโมทิ, สมฺโมทนียํ กถํ สารณียํ วีติสาเรตฺวา เอกมนฺตํ นิสีทิ, เอกมนฺตํ นิสินฺโน โข วจฺฉโคตฺโต ปริพฺพาชโก อายสฺมนฺตํ มหาโมคฺคลฺลานํ เอตทโวจ “กิํ นุ โข โภ โมคฺคลฺลาน สสฺสโต โลโก”ติ.\csromanspacer{} อพฺยากตํ โข เอตํ วจฺฉ ภควตา “สสฺสโต โลโก”ติ ฯเปฯ.\csromanspacer{} กิํ ปน โภ โมคฺคลฺลาน เนว โหติ น น โหติ ตถาคโต ปรํ มรณาติ.\csromanspacer{} เอตมฺปิ โข วจฺฉ อพฺยากตํ ภควตา “เนว โหติ น น โหติ ตถาคโต ปรํ มรณา”ติ.}

\prose{โก นุ โข โภ โมคฺคลฺลาน เหตุ โก ปจฺจโย, เยน อญฺญติตฺถิยานํ ปริพฺพาชกานํ เอวํ ปุฏฺฐานํ เอวํ เวยฺยากรณํ โหติ “สสฺสโต โลโก”ติ วา ฯเปฯ “เนว โหติ น น โหติ ตถาคโต ปรํ มรณา”ติ วา.\csromanspacer{} โก ปน โภ โมคฺคลฺลาน เหตุ โก ปจฺจโย, เยน สมณสฺส โคตมสฺส เอวํ ปุฏฺฐสฺส น เอวํ เวยฺยากรณํ โหติ “สสฺสโต โลโก”ติปิ ฯเปฯ “เนว โหติ น น โหติ ตถาคโต ปรํ มรณา”ติปีติ.}

\prose{อญฺญติตฺถิยา โข วจฺฉ ปริพฺพาชกา รูปํ อตฺตโต สมนุปสฺสนฺติ, รูปวนฺตํ วา อตฺตานํ, อตฺตนิ วา รูปํ, รูปสฺมิํ วา อตฺตานํ.\csromanspacer{} เวทนํ อตฺตโต สมนุปสฺสนฺติ ฯเปฯ.\csromanspacer{} สญฺญํ.\csromanspacer{} สงฺขาเร.\csromanspacer{} วิญฺญาณํ อตฺตโต สมนุปสฺสนฺติ, วิญฺญาณวนฺตํ วา อตฺตานํ, อตฺตนิ วา วิญฺญาณํ, วิญฺญาณสฺมิํ วา อตฺตานํ.\csromanspacer{} ตสฺมา อญฺญติตฺถิยานํ ปริพฺพาชกานํ เอวํ ปุฏฺฐานํ เอวํ เวยฺยากรณํ โหติ “สสฺสโต โลโก”ติ วา ฯเปฯ “เนว โหติ น น โหติ ตถาคโต ปรํ มรณา”ติ วา.\csromanspacer{} ตถาคโต จ โข วจฺฉ อรหํ สมฺมาสมฺพุทฺโธ น รูปํ อตฺตโต สมนุปสฺสติ, น รูปวนฺตํ วา อตฺตานํ, น อตฺตนิ วา รูปํ, น รูปสฺมิํ วา อตฺตานํ.\csromanspacer{} น เวทนํ อตฺตโต สมนุปสฺสติ ฯเปฯ.}

\csromanheader{2}{10. อพฺยากตสํยุตฺต}
\prose{\csromanfolio{563}น สญฺญํ.\csromanspacer{} น สงฺขาเร.\csromanspacer{} น วิญฺญาณํ อตฺตโต สมนุปสฺสติ, น วิญฺญาณวนฺตํ วา อตฺตานํ, น อตฺตนิ วา วิญฺญาณํ, น วิญฺญาณสฺมิํ วา อตฺตานํ.\csromanspacer{} ตสฺมา ตถาคตสฺส เอวํ ปุฏฺฐสฺส น เอวํ เวยฺยากรณํ โหติ “สสฺสโต โลโก”ติปิ, “อสสฺสโต โลโก”ติปิ, “อนฺตวา โลโก”ติปิ, “อนนฺตวา โลโก”ติปิ, “ตํ ชีวํ ตํ สรีรนฺ”ติปิ, อญฺญํ ชีวํ อญฺญํ สรีรนฺ”ติปิ, “โหติ ตถาคโต ปรํ มรณา”ติปิ, “น โหติ ตถาคโต ปรํ มรณา”ติปิ, “โหติ จ น จ โหติ ตถาคโต ปรํ มรณา”ติปิ, “เนว โหติ น น โหติ ตถาคโต ปรํ มรณา”ติปีติ.}

\prose{อจฺฉริยํ โภ โมคฺคลฺลาน, อพฺภุตํ โภ โมคฺคลฺลาน, ยตฺร หิ นาม สตฺถุ จ สาวกสฺส จ อตฺเถน อตฺโถ พฺยญฺชเนน พฺยญฺชนํ สํสนฺทิสฺสติ สเมสฺสติ น วิโรธยิสฺสติ, ยทิทํ อคฺคปทสฺมิํ.\csromanspacer{} อิทานาหํ โภ โมคฺคลฺลาน สมณํ โคตมํ อุปสงฺกมิตฺวา เอตมตฺถํ อปุจฺฉิํ, สมโณปิ เม โคตโม เอเตหิ ปเทหิ เอเตหิ พฺยญฺชเนหิ เอตมตฺถํ พฺยากาสิ, เสยฺยถาปิ ภวํ โมคฺคลฺลาโน.\csromanspacer{} อจฺฉริยํ โภ โมคฺคลฺลาน, อพฺภุตํ โภ โมคฺคลฺลาน, ยตฺร หิ นาม สตฺถุ จ สาวกสฺส จ อตฺเถน อตฺโถ พฺยญฺชเนน พฺยญฺชนํ สํสนฺทิสฺสติ สเมสฺสติ น วิโรธยิสฺสติ, ยทิทํ อคฺคปทสฺมิํติ.\csromanspacer{}\csromanspacer{}อฏฺฐมํ.}

\csromanheader{2}{9. กุตูหลสาลาสุตฺต}
\csromanmatikaanchor{mtk.294}
\csromantocmarkref{subsection}{9. กุตูหลสาลาสุตฺต}{mtk.294}
\bookmark[dest={mtk.294},level=2]{9. กุตูหลสาลาสุตฺต}
\proseitem{418}{อถ โข วจฺฉโคตฺโต ปริพฺพาชโก เยน ภควา เตนุปสงฺกมิ, อุปสงฺกมิตฺวา ภควตา สทฺธิํ สมฺโมทิ, สมฺโมทนียํ กถํ สารณียํ วีติสาเรตฺวา เอกมนฺตํ นิสีทิ, เอกมนฺตํ นิสินฺโน โข วจฺฉโคตฺโต ปริพฺพาชโก ภควนฺตํ เอตทโวจ\csromandash{}}

\prose{“ปุริมานิ โภ โคตม ทิวสานิ ปุริมตรานิ สมฺพหุลานํ นานาติตฺถิยานํ สมณพฺราหฺมณานํ ปริพฺพาชกานํ กุตูหลสาลายํ สนฺนิสินฺนานํ สนฺนิปติตานํ อยมนฺตรากถา อุทปาทิ ‘อยํ โข ปูรโณ กสฺสโป สงฺฆี เจว คณี จ คณาจริโย จ ญาโต ยสสฺสี ติตฺถกโร สาธุสมฺมโต พหุชนสฺส, โสปิ สาวกํ อพฺภตีตํ กาลงฺกตํ อุปปตฺตีสุ พฺยากโรติ ‘อสุ อมุตฺร อุปปนฺโน, อสุ อมุตฺร อุปปนฺโน’ติ.}

\gambhira{สฬายตนวคฺคสํยุตฺตปาฬิ}
\prose{\csromanfolio{564}โยปิสฺส สาวโก อุตฺตมปุริโส ปรมปุริโส ปรมปตฺติปตฺโต, ตมฺปิ สาวกํ อพฺภตีตํ กาลงฺกตํ อุปปตฺตีสุ พฺยากโรติ ‘อสุ อมุตฺร อุปปนฺโน, อสุ อมุตฺร อุปปนฺโน’ติ.}

\prose{อยมฺปิ โข มกฺขลิ โคสาโล ฯเปฯ.\csromanspacer{} อยมฺปิ โข นิคณฺโฐ นาฏปุตฺโต.\csromanspacer{} อยมฺปิ โข สญฺจโย\footnote{สญฺชโย (สี, สฺยา, กํ, อิ)} เพลฏฺฐปุตฺโต.\csromanspacer{} อยมฺปิ โข ปกุโธ\footnote{ปกุทฺโธ (อิ)} กจฺจาโน.\csromanspacer{} อยมฺปิ โข อชิโต เกสกมฺพโล สงฺฆี เจว คณี จ คณาจริโย จ ญาโต ยสสฺสี ติตฺถกโร สาธุสมฺมโต พหุชนสฺส, โสปิ สาวกํ อพฺภตีตํ กาลงฺกตํ อุปปตฺตีสุ พฺยากโรติ ‘อสุ อมุตฺร อุปปนฺโน, อสุ อมุตฺร อุปปนฺโน’ติ.\csromanspacer{} โยปิสฺส สาวโก อุตฺตมปุริโส ปรมปุริโส ปรมปตฺติปตฺโต, ตมฺปิ สาวกํ อพฺภตีตํ กาลงฺกตํ อุปปตฺตีสุ พฺยากโรติ ‘อสุ อมุตฺร อุปปนฺโน, อสุ อมุตฺร อุปปนฺโน’ติ.}

\prose{อยมฺปิ โข สมโณ โคตโม สงฺฆี เจว คณี จ คณาจริโย จ ญาโต ยสสฺสี ติตฺถกโร สาธุสมฺมโต พหุชนสฺส, โสปิ สาวกํ อพฺภตีตํ กาลงฺกตํ อุปปตฺตีสุ พฺยากโรติ ‘อสุ อมุตฺร อุปปนฺโน, อสุ อมุตฺร อุปปนฺโน’ติ.\csromanspacer{} โยปิสฺส\footnote{โย จ ขฺวสฺส (อิ)} สาวโก อุตฺตมปุริโส ปรมปุริโส ปรมปตฺติปตฺโต, ตญฺจ สาวกํ อพฺภตีตํ กาลงฺกตํ อุปปตฺตีสุ น พฺยากโรติ ‘อสุ อมุตฺร อุปปนฺโน, อสุ อมุตฺร อุปปนฺโน’ติ, อปิ จ โข นํ เอวํ พฺยากโรติ ‘อจฺเฉจฺฉิ ตณฺหํ, วิวตฺตยิ สํโยชนํ, สมฺมา มานาภิสมยา อนฺตมกาสิ ทุกฺขสฺสา’ติ.\csromanspacer{} ตสฺส มยฺหํ โภ โคตม อหุ เทว กงฺขา อหุ วิจิกิจฺฉา ‘กถํ นาม\footnote{กถญฺหิ นาม (สฺยา, กํ, อิ, ก) กถํ กถํ นาม (ฉกฺกงฺคุตฺตเร ปญฺจมวคฺเค ทุติยสุตฺเต)} สมณสฺส โคตมสฺส ธมฺโม อภิญฺเญยฺโย’ติ”\footnote{ธมฺมาภิญฺเญยฺยาติ (อิ, ก) ธมฺโม ... อญฺเญยฺโย (ฉกฺกงฺคุตฺตเร)}.}

\prose{อลํ หิ เต วจฺฉ กงฺขิตุํ, อลํ วิจิกิจฺฉิตุํ.\csromanspacer{} กงฺขนีเย จ ปน เต ฐาเน วิจิกิจฺฉา อุปฺปนฺนา.\csromanspacer{} ส-อุปาทานสฺส ขฺวาหํ วจฺฉ อุปปตฺติํ ปญฺญาเปมิ, โน อนุปาทานสฺส.\csromanspacer{} เสยฺยถาปิ วจฺฉ อคฺคิ ส-อุปาทาโน ชลติ, โน อนุปาทาโน.\csromanspacer{} เอวเมว ขฺวาหํ วจฺฉ ส-อุปาทานสฺส อุปปตฺติํ ปญฺญาเปมิ, โน อนุปาทานสฺสาติ.}

\csromanheader{2}{10. อพฺยากตสํยุตฺต}
\prose{\csromanfolio{565}ยสฺมิํ โภ โคตม สมเย อจฺจิ วาเตน ขิตฺตา ทูรมฺปิ คจฺฉติ, อิมสฺส ปน ภวํ โคตโม กิํ อุปาทานสฺมิํ ปญฺญาเปตีติ.\csromanspacer{} ยสฺมิํ โข วจฺฉ สมเย อจฺจิ วาเตน ขิตฺตา ทูรมฺปิ คจฺฉติ, ตมหํ วาตูปาทานํ ปญฺญาเปมิ.\csromanspacer{} วาโต หิสฺส วจฺฉ ตสฺมิํ สมเย อุปาทานํ โหตีติ.\csromanspacer{} ยสฺมิํ จ ปน โภ โคตม สมเย อิมญฺจ กายํ นิกฺขิปติ, สตฺโต จ อญฺญตรํ กายํ อนุปปนฺโน โหติ, อิมสฺส ปน ภวํ โคตโม กิํ อุปาทานสฺมิํ ปญฺญาเปตีติ.\csromanspacer{} ยสฺมิํ โข วจฺฉ สมเย อิมญฺจ กายํ นิกฺขิปติ, สตฺโต จ อญฺญตรํ กายํ อนุปปนฺโน โหติ, ตมหํ ตณฺหูปาทานํ วทามิ.\csromanspacer{} ตณฺหา หิสฺส วจฺฉ ตสฺมิํ สมเย อุปาทานํ โหตีติ\footnote{โหตีติ ฯเปฯ (ก)}.\csromanspacer{}\csromanspacer{}นวมํ.}

\csromanheader{2}{10. อานนฺทสุตฺต}
\csromanmatikaanchor{mtk.295}
\csromantocmarkref{subsection}{10. อานนฺทสุตฺต}{mtk.295}
\bookmark[dest={mtk.295},level=2]{10. อานนฺทสุตฺต}
\proseitem{419}{อถ โข วจฺฉโคตฺโต ปริพฺพาชโก เยน ภควา เตนุปสงฺกมิ, อุปสงฺกมิตฺวา ภควตา สทฺธิํ สมฺโมทิ, สมฺโมทนียํ กถํ สารณียํ วีติสาเรตฺวา เอกมนฺตํ นิสีทิ, เอกมนฺตํ นิสินฺโน โข วจฺฉโคตฺโต ปริพฺพาชโก ภควนฺตํ เอตทโวจ “กิํ นุ โข โภ โคตม อตฺถตฺตา”ติ.\csromanspacer{} เอวํ วุตฺเต ภควา ตุณฺหี อโหสิ.\csromanspacer{} กิํ ปน โภ โคตม นตฺถตฺตาติ.\csromanspacer{} ทุติยมฺปิ โข ภควา ตุณฺหี อโหสิ.\csromanspacer{} อถ โข วจฺฉโคตฺโต ปริพฺพาชโก อุฏฺฐายาสนา ปกฺกามิ.}

\prose{อถ โข อายสฺมา อานนฺโท อจิรปกฺกนฺเต วจฺฉโคตฺเต ปริพฺพาชเก ภควนฺตํ เอตทโวจ “กิํ นุ โข ภนฺเต ภควา วจฺฉโคตฺตสฺส ปริพฺพาชกสฺส ปญฺหํ ปุฏฺโฐ น พฺยากาสี”ติ.\csromanspacer{} อหญฺจานนฺท วจฺฉโคตฺตสฺส ปริพฺพาชกสฺส “อตฺถตฺตา”ติ ปุฏฺโฐ สมาโน “อตฺถตฺตา”ติ พฺยากเรยฺยํ, เย เต อานนฺท สมณพฺราหฺมณา สสฺสตวาทา, เตสเมตํ สทฺธิํ\footnote{เตสเมตํ ลทฺธิ (สี)} อภวิสฺส.\csromanspacer{} อหญฺจานนฺท วจฺฉโคตฺตสฺส ปริพฺพาชกสฺส “นตฺถตฺตา”ติ ปุฏฺโฐ สมาโน “นตฺถตฺตา”ติ พฺยากเรยฺยํ, เย เต อานนฺท สมณพฺราหฺมณา อุจฺเฉทวาทา, เตสเมตํ สทฺธิํ อภวิสฺส.\csromanspacer{} อหญฺจานนฺท วจฺฉโคตฺตสฺส ปริพฺพาชกสฺส “อตฺถตฺตา”ติ ปุฏฺโฐ สมาโน “อตฺถตฺตา”ติ พฺยากเรยฺยํ, อปิ นุ เม ตํ อานนฺท อนุโลมํ อภวิสฺส ญาณสฺส อุปฺปาทาย “สพฺเพ ธมฺมา อนตฺตา”ติ.\csromanspacer{} โน เหตํ ภนฺเต.\csromanspacer{} อหญฺจานนฺท วจฺฉโคตฺตสฺส}

\gambhira{สฬายตนวคฺคสํยุตฺตปาฬิ}
\prose{\csromanfolio{566}ปริพฺพาชกสฺส “นตฺถตฺตา”ติ ปุฏฺโฐ สมาโน “นตฺถตฺตา”ติ พฺยากเรยฺยํ, สมฺมูฬฺหสฺส อานนฺท วจฺฉโคตฺตสฺส ปริพฺพาชกสฺส ภิยฺโย สมฺโมหาย อภวิสฺส “อหุวา เม นูน ปุพฺเพ อตฺตา, โส เอตรหิ นตฺถี”ติ.\csromanspacer{}\csromanspacer{}ทสมํ.}

\csromanheader{2}{11. สภิยกจฺจานสุตฺต}
\csromanmatikaanchor{mtk.296}
\csromantocmarkref{subsection}{11. สภิยกจฺจานสุตฺต}{mtk.296}
\bookmark[dest={mtk.296},level=2]{11. สภิยกจฺจานสุตฺต}
\proseitem{420}{เอกํ สมยํ อายสฺมา สภิโย กจฺจาโน ญาติเก วิหรติ คิญฺชกาวสเถ.\csromanspacer{} อถ โข วจฺฉโคตฺโต ปริพฺพาชโก เยนายสฺมา สภิโย กจฺจาโน เตนุปสงฺกมิ, อุปสงฺกมิตฺวา อายสฺมตา สภิเยน กจฺจาเนน สทฺธิํ สมฺโมทิ, สมฺโมทนียํ กถํ สารณียํ วีติสาเรตฺวา เอกมนฺตํ นิสีทิ, เอกมนฺตํ นิสินฺโน โข วจฺฉโคตฺโต ปริพฺพาชโก อายสฺมนฺตํ สภิยํ กจฺจานํ เอตทโวจ “กิํ นุ โข โภ กจฺจาน โหติ ตถาคโต ปรํ มรณา”ติ.\csromanspacer{} อพฺยากตํ โข เอตํ วจฺฉ ภควตา “โหติ ตถาคโต ปรํ มรณา”ติ.\csromanspacer{} กิํ ปน โภ กจฺจาน น โหติ ตถาคโต ปรํ มรณาติ.\csromanspacer{} เอตมฺปิ โข วจฺฉ อพฺยากตํ ภควตา “น โหติ ตถาคโต ปรํ มรณา”ติ.\csromanspacer{} กิํ นุ โข โภ กจฺจาน โหติ จ น จ โหติ ตถาคโต ปรํ มรณาติ.\csromanspacer{} อพฺยากตํ โข เอตํ วจฺฉ ภวคตา “โหติ จ น จ โหติ ตถาคโต ปรํ มรณา”ติ.\csromanspacer{} กิํ ปน โภ กจฺจาน เนว โหติ น น โหติ ตถาคโต ปรํ มรณาติ.\csromanspacer{} เอตมฺปิ โข วจฺฉ อพฺยากตํ ภควตา “เนว โหติ น น โหติ ตถาคโต ปรํ มรณา”ติ.}

\prose{“กิํ นุ โข โภ กจฺฉาน โหติ ตถาคโต ปรํ มรณา”ติ อิติ ปุฏฺโฐ สมาโน “อพฺยากตํ โข เอตํ วจฺฉ ภควตา ‘โหติ ตถาคโต ปรํ มรณา’ติ” วเทสิ. “กิํ ปน โภ กจฺจาน น โหติ ตถาคโต ปรํ มรณา”ติ อิติ ปุฏฺโฐ สมาโน “อพฺยากตํ โข เอตํ วจฺฉ ภควตา ‘น โหติ ตถาคโต ปรํ มรณา’ติ” วเทสิ. “กิํ นุ โข โภ กจฺจาน โหติ จ น จ โหติ ตถาคโต ปรํ มรณา”ติ อิติ ปุฏฺโฐ สมาโน “อพฺยากตํ โข เอตํ วจฺฉ ภควตา ‘โหติ จ น จ โหติ ตฺ0อถาคโต ปรํ มรณา’ติ” วเทสิ. “กิํ ปน โภ กจฺจาน เนว โหติ น น โหติ ตถาคโต ปรํ มรณา”ติ อิติ ปุฏฺโฐ สมาโน “เอตมฺปิ โข วจฺฉ อพฺยากตํ ภควตา ‘เนว โหติ}

\csromanheader{2}{10. อพฺยากตสํยุตฺต}
\prose{\csromanfolio{567}น น โหติ ตถาคโต ปรํ มรณา’ติ” วเทสิ.\csromanspacer{} โก นุ โข โภ กจฺจาน เหตุ โก ปจฺจโย, เยเนตํ อพฺยากตํ สมเณน โคตเมนาติ.\csromanspacer{} โย จ วจฺฉ เหตุ โย จ ปจฺจโย ปญฺญาปนาย “รูปีติ วา อรูปีติ วา สญฺญีติ วา อสญฺญีติ วา เนวสญฺญีนาสญฺญีติ วา”, โส จ เหตุ โส จ ปจฺจโย สพฺเพน สพฺพํ สพฺพถา สพฺพํ อปริเสสํ นิรุชฺเฌยฺย, เกน นํ ปญฺญาปยมาโน ปญฺญาเปยฺย “รูปีติ วา อรูปีติ วา สญฺญีติ วา อสญฺญีติ วา เนวสญฺญีนาสญฺญีติ วา”ติ.\csromanspacer{} กีวจิรํ ปพฺพชิโตสิ โภ กจฺจานาติ.\csromanspacer{} นจิรํ อาวุโส ตีณิ วสฺสานีติ.\csromanspacer{} ยสฺสป’สฺส อาวุโส เอตเมตฺตเกน เอตฺตกเมว, ตํป’สฺส พหุ.\csromanspacer{} โก ปน วาโท เอวํ\footnote{โก ปน วาโท เอว (สี, อิ)} อภิกฺกนฺเตติ.\csromanspacer{}\csromanspacer{}เอกาทสมํ.}

\vspace{\tipitakaparskip}
\csromancenter{\textbf{อพฺยากตสํยุตฺตํ สมตฺตํ.}}
\titlehead{\textbf{ตสฺสุทฺทานํ}}
\csromangathameasure{เขมาเถรี อนุราโธ, สาริปุตฺโตติ โกฏฺฐิโก. \\ โมคฺคลฺลาโน จ วจฺโฉ จ, กุตูหลสาลานนฺโท. \\ สภิโย เอกาทสมนฺติ. สฬายตนวคฺโค จตุตฺโถ.}
\csromangathagroup{\csromangathastanza{เขมาเถรี อนุราโธ, สาริปุตฺโตติ โกฏฺฐิโก. \\ โมคฺคลฺลาโน จ วจฺโฉ จ, กุตูหลสาลานนฺโท. \\ สภิโย เอกาทสมนฺติ.\csromanspacer{} สฬายตนวคฺโค จตุตฺโถ.}}

\titlehead{\textbf{ตสฺสุทฺทานํ}}
\csromangathameasure{สฬายตนเวทนา, มาตุคาโม ชมฺพุขาทโก. \\ สามณฺฑโก โมคฺคลฺลาโน, จิตฺโต คามณิ’สงฺขตํ. \\ อพฺยากตนฺติ ทสธาติ.}
\csromangathagroup{\csromangathastanza{สฬายตนเวทนา, มาตุคาโม ชมฺพุขาทโก. \\ สามณฺฑโก โมคฺคลฺลาโน, จิตฺโต คามณิ’สงฺขตํ. \\ อพฺยากตนฺติ ทสธาติ.}}

\nitthitam{\textbf{สฬายตนวคฺคสํยุตฺตปาฬิ นิฏฺฐิตา.}}
\orphannote{สนฺตุสฺสิตํ (ก-สี, อิ, ก)}

